% !TeX program = pdflatex
% Mathematical working note.
% This note records the finite-cutoff RG prescription and the regular reference
% perimeter/corner extraction and outgoing closed insertions.
% Version 4 (22 October 2026): V3 with the section on two disjoint squares
% (sec:two-loops) inserted before sec:next-step and Theorem thm:two-loop-theorem
% added to the main-result subsection; V3 is unchanged otherwise.
% Version 3 (19 October 2026): V2 (17 October 2026) with the sourced effective
% interaction at the first step (stage 1 draft) inserted as Section 3, the
% background estimates (stage 2 draft) inserted as Appendix B, and the insertion
% at the retreated carrier (stage 5 draft) inserted as the last numbered section
% (sec:one-step-insertion), the exterior response with its history average
% (stage 3 draft) inserted as Appendix D (app:exterior-chessboard,
% lem:exterior-response-bound), the signed integrated Wilson defects
% (stage 4 draft) inserted as Appendix E (app:signed-defects, prop:c4-12-3),
% and the regional scalar of the covered branch (stage 4C draft, Sections 1,
% 2, 5 and 8; its Sections 3-7 were already the proof of
% prop:regional-square-extraction) inserted as the last appendix
% (app:regional-scalar, prop:c4c-5-1); V2 is unchanged otherwise.
\documentclass[11pt]{article}
\usepackage[T1]{fontenc}
\usepackage[utf8]{inputenc}
\usepackage{amsmath,amssymb,amsthm}
\usepackage{mathtools}
\usepackage{mathrsfs}
\usepackage[a4paper,margin=28mm]{geometry}
\usepackage{enumitem}
\usepackage{graphicx}
\usepackage[hidelinks]{hyperref}
\newcommand{\dd}{\mathrm{d}}
\newcommand{\Hol}{\operatorname{Hol}}
\newcommand{\tr}{\operatorname{tr}}
\newcommand{\id}{\mathbf 1}
\makeatletter % Stolen from the internet to make a fat \cdot which isn't as fat as a \bullet
\newcommand*{\bigcdot}{}% Check if undefined
\DeclareRobustCommand*{\bigcdot}{%
  \mathbin{\mathpalette\bigcdot@{}}%
}
\newcommand*{\bigcdot@scalefactor}{.7}
\newcommand*{\bigcdot@widthfactor}{1.15}
\newcommand*{\bigcdot@}[2]{%
  % #1: math style
  % #2: unused
  \sbox0{$#1\vcenter{}$}% math axis
  \sbox2{$#1\cdot\m@th$}%
  \hbox to \bigcdot@widthfactor\wd2{%
    \hfil
    \raise\ht0\hbox{%
      \scalebox{\bigcdot@scalefactor}{%
        \lower\ht0\hbox{$#1\bullet\m@th$}%
      }%
    }%
    \hfil
  }%
}
\makeatother
\setlength{\parindent}{0pt}
\setlength{\parskip}{0.6em}
\emergencystretch=1.5em
\title{A Wilson Loop after the First Renormalisation Step}
\author{}
\date{22 October 2026}

\newtheorem{theorem}{Theorem}
\newtheorem{proposition}{Proposition}
\newtheorem{lemma}{Lemma}
\newtheorem{definition}{Definition}
\newtheorem{corollary}{Corollary}
\newtheorem{conjecture}{Conjecture}
\theoremstyle{remark}
\newtheorem{remark}{Remark}
\theoremstyle{plain}
\renewcommand{\qedsymbol}{\textit{q.e.d.}}
\begin{document}
\maketitle

\section{The Insertion at a Fixed Cutoff}

This paper follows a Wilson loop through the first renormalisation-group step of Ba{\l}aban's construction of four-dimensional lattice Yang--Mills theory \cite{BalabanCRE,BalabanRG1}. That step is the block-spin transformation $T$ of \eqref{eq:bare-blocking-operator}: the bond variables of a periodic fine lattice of spacing $\varepsilon$ are averaged over blocks of side $L$, and the bare measure is mapped to the induced measure of the bond variables of the coarse lattice of spacing $L\varepsilon$. The loop is inserted in the bare measure and carried through $T$; the result (Section~\ref{sec:main-result}) writes its bare expectation as a perimeter factor times a corner factor times an average, over the coarse ensemble, of the loop at the background field reconstructed from the coarse field, plus two error terms that it bounds.

The gauge group $G$ is a compact connected simple Lie group, $g_0$ is the bare coupling, and $\chi \colon G \to \mathbb C$ is a normalised irreducible character, so that $\chi(\id)=1$. Let $\varepsilon>0$ be the fine lattice spacing and $L>1$ the blocking factor. The fine and coarse bond fields are $U$ and $V$, respectively. Choose an oriented square contour $C$ on the fine lattice and an oriented square $\bar C$ on the lattice of spacing $L\varepsilon$, with each side of $\bar C$ represented by the corresponding straight fine path in $C$. This assumption fixes the path correspondence; the two contours describe the same square in the underlying torus.

For a field $U$ and an ordered contour $C=(b_1,\ldots,b_n)$, holonomy and the Wilson loop are
\begin{equation}\label{eq:wilson-definition}
    \Hol_U C= U(b_1)\cdots U(b_n),
    \qquad W_{\chi}(C;U)=\chi(\Hol_U C)\;.
\end{equation}
The bond field satisfies $U(b^{-1})= U(b)^{-1}$ when a contour traverses a bond backwards. The character makes $W_{\chi}$ gauge invariant.

Let $E$ and $E_1$ be the positive fine and coarse bonds of the four-dimensional periodic lattices, and let $\mathcal P_T$ be the unoriented fine plaquettes. Fix the action's unitary realisation $G\subset U(N)$ and $\tr_N= N^{-1}\operatorname{Tr}$. With $\partial p$ denoting the oriented boundary of $p$, the bare action and density are
\begin{equation}\label{eq:bare-action}
    S(U)=\sum_{p\in \mathcal P_T}
        (1-\operatorname{Re}\tr_N U(\partial p)),
    \qquad \rho_0(U)=\exp(-g_0^{-2}S(U)-E_{\varepsilon})\;.
\end{equation}
Here $E_{\varepsilon}$ is a field-independent vacuum scalar, retained in both numerator and denominator. Product normalised Haar measures on $G^E$ and $G^{E_1}$ are denoted by $\dd U$ and $\dd V$. The full first-step operator, before the small-field/large-field decomposition, is CRE (0.1):
\begin{equation}\label{eq:bare-blocking-operator}
    (T\rho)(V)=\int
        \left(\prod_{c\in E_1}\delta_G(M(U,c)V(c)^{-1})\right)
        \rho(U)\,\dd U\;.
\end{equation}
Here $\delta_G$ is the group Dirac measure normalised against Haar measure, $\int\delta_G(WV^{-1})\,\dd V=1$ for every $W\in G$, and $M$ is the averaging operation, defined as follows. For a positive coarse bond $c$ let $P_c(U)$ be the transport along the straight fine path from $c_-$ to $c_+$, let $\beta_c(U)$ be the logarithmic average of RG1 (0.4), displayed in \eqref{eq:nonlinear-average} below, and let $\mathcal O_c$ be the set of fine configurations $U$ at which every logarithm in $\beta_c(U)$ has its argument in the chart domain (the chart is the inverse of $A\mapsto e^{iA}$ on the Hilbert ball $h(A)<\varrho_G$, fixed in the conventions paragraph of Section~\ref{sec:sourced-effective-interaction}). Put
\begin{equation}\label{eq:global-average}
    M(U,c)=
    \begin{cases}
        e^{i\beta_c(U)}P_c(U), & U\in\mathcal O_c,\\
        P_c(U), & U\notin\mathcal O_c.
    \end{cases}
\end{equation}
CRE's (0.1) integrates over every fine configuration, so it needs an averaging operation on all of $G^E$; CRE's own formula is RG1 (0.4), defined where its logarithms are, and CRE states (p.243) that any averaging operation satisfying general properties analogous to RG1 (0.5)--(0.9) may be taken. Equation~\eqref{eq:global-average} is the choice made here; the properties used in this paper are the following, (ii) being the analogue of the gauge covariance RG1 (0.6), proved directly from \eqref{eq:global-average}. The properties (0.5)--(0.9) of RG1's group mean, which CRE's licence presupposes for its averaging operation, are not used for $M$ off $\mathcal O_c$: on the branches the averaging operation is evaluated on $\mathcal O_c$ only, where it is RG1 (0.4) ((B1), Section~\ref{sec:branch-data}), and its value off $\mathcal O_c$ enters only through the constraint on the rows of $\Omega_1^c$.

\begin{lemma}[the averaging operation]\label{lem:global-average}
(i) $\mathcal O_c$ is open, and $M(\bigcdot ,c)\colon G^E\to G$ is Borel, analytic on $\mathcal O_c$, where it is RG1 (0.4), and a function of the bonds with both endpoints in the two blocks of $c$ alone.
(ii) For every fine gauge transformation $g\colon T\to G$, $U^g(b)= g(b_-)U(b)g(b_+)^{-1}$, one has $U^g\in\mathcal O_c$ if and only if $U\in\mathcal O_c$, and $M(U^g,c)= g(c_-)M(U,c)g(c_+)^{-1}$; in particular $M(U^g)= M(U)$ when $g$ is the identity at the block centres.
\end{lemma}

\begin{proof}
(i) The chart domain is open, being the image of the open ball under $A\mapsto e^{iA}$, which is injective with invertible differential there (the conventions paragraph of Section~\ref{sec:sourced-effective-interaction}). Each loop $U(\gamma)U([x,x'])U(\gamma')^{-1}P_c(U)^{-1}$ of \eqref{eq:nonlinear-average} is a product of bond variables and their inverses, continuous in $U$, so $\mathcal O_c$, the set where finitely many such loops lie in the chart domain, is open; the chart is analytic on the chart domain (inverse of an analytic map with invertible differential), so $\beta_c$ and $e^{i\beta_c}P_c$ are analytic on $\mathcal O_c$, and $P_c$ is analytic on $G^E$. A map that is continuous on an open set and continuous on its complement is Borel. The paths $\gamma$, $[x,x']$, $\gamma'$ and the straight path of $c$ have all their bonds with both endpoints in $\mathcal B_{c_-}\cup\mathcal B_{c_+}$: $\gamma\subset\mathcal B_{c_-}$, $\gamma'\subset\mathcal B_{c_+}$, and the straight path from $x\in\mathcal B_{c_-}$ to $x+Le_{\mu}\in\mathcal B_{c_+}$ leaves $\mathcal B_{c_-}$ at the bond at which it enters $\mathcal B_{c_+}$.
(ii) Every loop is based at $c_-$, so under $g$ it is conjugated by $g(c_-)$. The Hilbert ball is $\operatorname{Ad}$-invariant, $h$ being $\operatorname{Ad}$-invariant, and $e^{i\,gAg^{-1}}= ge^{iA}g^{-1}$; so the chart domain is invariant under conjugation, and for $W$ in it $\frac1i\log(gWg^{-1})= g(\frac1i\log W)g^{-1}$, by the injectivity of $A\mapsto e^{iA}$ on the ball. Hence $U^g\in\mathcal O_c$ if and only if $U\in\mathcal O_c$, $\beta_c(U^g)= g(c_-)\beta_c(U)g(c_-)^{-1}$ and $e^{i\beta_c(U^g)}= g(c_-)e^{i\beta_c(U)}g(c_-)^{-1}$ there, while $P_c(U^g)= g(c_-)P_c(U)g(c_+)^{-1}$ everywhere. The product gives the covariance on $\mathcal O_c$, and off $\mathcal O_c$ it is the covariance of $P_c$.
\end{proof}

The right side of \eqref{eq:bare-blocking-operator} is read as follows. For bounded Borel $h$ on $G^{E_1}$ and $W\in G^{E_1}$, $\int\dd V\,h(V)\prod_{c\in E_1}\delta_G(W(c)V(c)^{-1})= h(W)$; so for a nonnegative or integrable Borel function $\rho$ on $G^E$, $T\rho$ is the image measure of $\rho\,\dd U$ under the Borel map $M\colon G^E\to G^{E_1}$ of Lemma~\ref{lem:global-average}, $(T\rho)(A)=\int_{M^{-1}(A)}\rho\,\dd U$ for Borel $A\subseteq G^{E_1}$, that is $\int h\,\dd(T\rho)=\int h(M(U))\rho(U)\,\dd U$. This is the meaning of \eqref{eq:bare-blocking-operator} and of every group $\delta$-function of a coarse variable in this paper: it is integrated against that variable first. CRE writes (0.1) as a Haar density; for integrable $\rho$ the finite measure $T\rho$ has a Haar density exactly when it is absolutely continuous with respect to $\dd V$, and then $(T\rho)(V)$ denotes a version of that density. No absolute continuity is asserted or needed: every identity below in which $\rho_1(V)$ or $N_1(V)$ appears is an identity of measures, in which $\rho_1(V)\,\dd V$ and $N_1(V)\,\dd V$ stand for the measures $\rho_1$ and $N_1$ of \eqref{eq:first-step}, and CRE's transformations of (0.1) into (1.11)--(1.15) are identities for the integrals of bounded Borel functions of $V$ against these measures under the branch hypotheses (B1)--(B4) of Section~\ref{sec:branch-data}, which record the facts about Ba{\l}aban's construction that these identities use. CRE's additional operation $\mathbf R$ (p.243), applied after each transformation $T$ and changing the densities on the large-field regions, is not applied here: the first-step objects are those of $T$ alone. The bare expectation is $\mathbb E_{\varepsilon}[F]\coloneqq \int F\rho_0\,\dd U/\int\rho_0\,\dd U$ for bounded Borel $F$ on $G^E$. Ba{\l}aban's cutoffs, gauge insertions and changes of variables give alternative representations of this same integral; their full history sum is retained \cite{BalabanCRE}: the partitions of unity CRE (1.1), (1.4) sum to one pointwise, so the sum over histories of the branch integrals is the full integral $\int\rho_0\,\dd U$ whatever the value of $M$ off $\bigcup_c\mathcal O_c$, the branch integrals using $M$ on $\mathcal O_c$ only (Section~\ref{sec:branch-data}). For a finite or complex measure $\mu$ on $G^E$, $T\mu\coloneqq M_*\mu$ is the same pushforward, so that $T$ acts on the measures defined next. The outgoing measure and the measure with the insertion are
\begin{equation}\label{eq:first-step}
    \rho_1= T[\rho_0],
    \qquad N_1= T[\rho_0 W_{\chi}(C;\bigcdot )]\;.
\end{equation}
Here $\rho_1$ is a finite positive measure on $G^{E_1}$ with $\rho_1(G^{E_1})=\int\rho_0\,\dd U$, and $N_1$ is a complex measure with $| N_1(A)|\leqslant \int_{M^{-1}(A)}\rho_0| W_{\chi}|\,\dd U\leqslant \rho_1(A)$ for every Borel $A$, so that $| N_1|\leqslant \rho_1$ for the total variation and $N_1\ll\rho_1$. The outgoing Wilson function is the Radon--Nikodym derivative
\begin{equation}\label{eq:outgoing-loop}
    W_1(C;V)=\frac{\dd N_1}{\dd\rho_1}(V)\;,
\end{equation}
defined for $\rho_1$-almost every $V$, with $| W_1|\leqslant 1$ $\rho_1$-almost everywhere; a Borel version with $| W_1|\leqslant 1$ everywhere is fixed. Where $\rho_1$ and $N_1$ have Haar densities, $W_1(C;V)= N_1(V)/\rho_1(V)$ for $\rho_1$-almost every $V$, the density-ratio form in which CRE writes such quantities. This definition carries the bare insertion through the existing finite-cutoff transformation. It does not introduce a Wilson-loop counterterm.

The intermediate integration over regular fluctuations does not yet produce the measures \eqref{eq:first-step} on $G^{E_1}$. It also retains the exterior fine field and the unintegrated fluctuations. Those variables are integrated by the outer operation in CRE (1.29); the measure $\rho_1$ is obtained from the same operation without the insertion. The vacuum construction is given in CRE (1.11)--(1.30) \cite{BalabanCRE}.

\subsection{Normalisation}
The bound $| N_1|\leqslant \rho_1$ is the positivity of the blocking kernel together with $| W_{\chi}|\leqslant 1$.
The physical partition function and the expectation after the transformation are
\begin{equation}\label{eq:one-step-physical-partition}
    \mathcal Z_0=\int\rho_0(U)\,\dd U
        =\rho_1(G^{E_1})>0\;,
\end{equation}
\begin{equation}
    \mathbb E_{\varepsilon}[W_{\chi}(C)]
        =\frac{1}{\mathcal Z_0}N_1(G^{E_1})
        =\frac{1}{\mathcal Z_0}\int W_1(C;V)\,\dd\rho_1(V)\;.
\end{equation}
Both integrals retain the sum over small-field and large-field regions. Conditional fluctuation integrals used below have their own normalising factors; those factors remain inside the construction of $N_1$ and $\rho_1$.

\subsection{The main result}\label{sec:main-result}

The main statement of the paper concerns $\mathbb E_{\varepsilon}[W_{\chi}(C)]$ for a square contour: after the first step the expectation is a perimeter factor times a corner factor, times the expectation, in the coarse ensemble enriched by the large-field branch data, of $\mathbf 1_{A_{\mathrm{cov}}}e^{E^{\circ}_{\xi}}W_{\chi}(C_n;U_1)$: the fine loop at the background $U_1$ reconstructed from the coarse field and the exterior fine field, with a finite-geometry exponent and the indicator of the covered event, plus two error terms. The objects are defined in Sections~\ref{sec:sourced-effective-interaction}--\ref{sec:remaining-physical-comparison} and the appendices; they are listed here so that the statement can be read on its own, and nothing in the list is assumed. Theorem~\ref{thm:two-loop-theorem} gives the same form for the product of the Wilson functions of two vertex-disjoint squares in planes parallel to that of $e_1,e_2$, with the product of the two coefficients, the two finite-geometry exponents and a cross exponent bounded in terms of the distance between the squares; its objects specific to the pair are defined in Section~\ref{sec:two-loops} and referenced from its statement.

\emph{Notation.} $C_n$ is a coarse translate of the counterclockwise square with vertices $0$, $nLe_1$, $nL(e_1+e_2)$, $nLe_2$ in fine-lattice units (Section~\ref{sec:regional-square-scalar}), $n\geqslant 1$; it has $n_C=4nL$ contour bonds, $4n$ coarse segments indexed cyclically by $j\in\{0,\ldots,4n-1\}$, and $P(C_n)$ is its perimeter as it enters \eqref{eq:extracted-loop-scalar} (Section~\ref{sec:c4c-1}). $W_{\chi}(C;U)$ is the normalised character $\chi$ of the holonomy of the fine field $U$ around $C$, $\varepsilon$ the fine lattice spacing, and $\sigma$ names the cutoff norm of \eqref{eq:literal-cutoff}, the operator or the normalised Hilbert--Schmidt norm (conventions of Section~\ref{sec:sourced-effective-interaction}). $T$ is the operator \eqref{eq:bare-blocking-operator} with the averaging operation \eqref{eq:global-average}; $\rho_1$, $N_1$, $W_1$ and $\mathcal Z_0$ are \eqref{eq:first-step}, \eqref{eq:outgoing-loop} and \eqref{eq:one-step-physical-partition}, and $\mathbb E_{\varepsilon}$ is the bare expectation defined after \eqref{eq:bare-blocking-operator}. $\xi=(\Omega_1,R_1,V,V_0,A_{\mathrm e})$ is the branch data of Section~\ref{sec:branch-data} (the small-field region $\Omega_1$, the large-fluctuation region $R_1$, which with $\Omega_1$ determines the filled region $\Lambda_1$, the coarse field, the exterior fine field and the exterior free fluctuation; at the original carrier there are no labels), $U_1= U_1(V,V_0)$ the background of Section~\ref{sec:branch-data}, $\mathbb E_{\mathrm h}$ the history expectation \eqref{eq:c3-6-1} (weight $\mathsf w_{\xi}= e^{c_{\xi}}I_{\xi;1}$, normalisation $\mathcal Z_0$; positive, with $\mathbb E_{\mathrm h}[1]=1$), and $A_{\mathrm{cov}}$ the covered event of Lemma~\ref{lem:covered}. $\kappa_{C_n}= e^{-\delta m^{\sigma}(g_0,\varepsilon)P(C_n)}(z^{\sigma}_{\mathrm c}(g_0,\pi/2))^4$ is the scalar \eqref{eq:extracted-loop-scalar}, with the perimeter coefficient $\delta m^{\sigma}$ of \eqref{eq:marginal-square-mass} and the corner factor $z^{\sigma}_{\mathrm c}$ of \eqref{eq:corner-square-factor}, in the cutoff prescription $\sigma$ fixed throughout. $E_{\mathrm{def}}$ is the signed-defect term \eqref{eq:c4-1-1}, $E_{\mathrm{bdry}}$ the exterior-response and uncovered term \eqref{eq:c3-6-2}, and $E^{\circ}_{\xi}= E^{\circ}_{I,C_n,T}$ the finite-geometry exponent \eqref{eq:c4c-5-2} of Proposition~\ref{prop:c4c-5-1}, a function of the history through the carrier $I=\Lambda_1^*$ alone. $d_j=\max(0,d_T(K_j,I^c)-R_0)$ and $\iota_T=\max(0,N_{\min}/4-R_0)$ are the distances \eqref{eq:c4c-1-1}: $K_j$ the union of the free-input stencils of the fine traversals of segment $j$, $I^c$ the complement of the carrier in the torus free carrier, $d_T$ the periodic fine-bond metric, $N_{\min}$ the shortest fine period, and $R_0=4(1+R_*+R_H+L)$ the fixed range of Section~\ref{sec:regional-identification} ($R_*$ the perimeter range, not the walk range $n_*R$ of Section~\ref{app:thresholds}; this $R_0$ is not CRE's scale $R_0$). The scales of the statement are $\tau$, the marked-forest parameter of \eqref{eq:marginal-fixed-forest-bound}, $\tau= C[g_0(1+r)^3+e^{-cr^2}]$ with the cutoff radius $r= A_1(\log g_0^{-2})^{p_0}$ of \eqref{eq:literal-cutoff}; $\pi_*=\min(p_0(g_0),p_0(g_1))$ with $p_0(g)= A_0(\log g^{-2})^{p_0}$ and $g_1= g_1(g_0)$ the first-step coupling of CRE (1.27); $\varepsilon_1= g_1p_0(g_1)$; $\rho$ of \eqref{eq:c4-4-8}, $\rho= O(g_0^2(1+r)^4)$; and $b=\lfloor LMR_1/16\rfloor$ with CRE's scale $R_1= R_1(g_0)$ (Section~\ref{app:thresholds}). The constants $C$, $c$ of Proposition~\ref{prop:c4c-5-1}, $C'_L$, $c_*$ of \eqref{eq:c4-2-2} and Proposition~\ref{prop:c3-7-5}, $C_{\chi}$, $c$ of \eqref{eq:c4-6-8}, $C_{\mathrm{resp},\chi}$ of \eqref{eq:c3-7-14} (the constant after the replacement made there, bearing the name of the constant of Corollary~\ref{cor:c3-6-1} it is built from), and $c_q$ of \eqref{eq:c5-2-3} ($c_q= c_*$) are fixed before the coupling and depend on $L$, the group, the character, $M_1$, $M_2$, $M$, the fixed expansion and chart choices and CRE's constants $A_0$, $A_1$, $p_0$, $B_3$, $\beta_0$ only (the last through $c_{\theta}$ of \eqref{eq:c3-7-4}, hence $c_*= c_q$, which involve $\delta_0=(A_1/A_0)\varepsilon_0$, $B_3$ and CRE's bound (2.6) on $g_1/g_0$), not on $g_0$, $n$, the torus or the history; the constants $C$, $c$ of different statements are distinct and are named by their statement; $K=256$ is the fixed constant bounding, in units of $\ell= LMR_1$, the distance of a witness from its tested bond (Sections~\ref{sec:c3-7} and~\ref{sec:c5-2}, Lemma~\ref{lem:c5-7-joint}; distinct from the periods $K_{\mu}$ and the stencils $K_j$), and $C_K=4(4K\ell+9)^4$ depends on $g_0$ through CRE's $R_1= R_1(g_0)$, a power of $L$ with $R_1\geqslant (\log g_1^{-2})^{\mathfrak r}$; $\delta^I_q$ of \eqref{eq:c5-7-expbound-I} depends on $g_0$ (through $\tau$, $C_K$ and the one-point rate $q$ of Section~\ref{sec:c3-7}), not on $n$, the torus or the history, and is smaller than every power of $g_0$ (Remark~\ref{rem:main-regime}(1)).

\emph{Hypotheses.} (H1) $G$ is a compact connected simple Lie group in its fixed unitary realisation $G\subset U(N)$ and $\chi$ a normalised irreducible character (the opening of this section and the conventions of Section~\ref{sec:sourced-effective-interaction}). (H2) The parameters satisfy the standing conditions \eqref{eq:c-standing-conditions}, the constants are fixed in the order of Section~\ref{app:thresholds}, and $g_0$ lies below the threshold of Section~\ref{app:thresholds}: its items 1--6 and the further conditions listed there, among them the threshold $g_*$ of Proposition~\ref{prop:c3-7-5} (which contains that of Proposition~\ref{prop:c3-9-1}), the thresholds of Proposition~\ref{prop:c4-12-3} (those of Lemmas~\ref{lem:c4-2-1}, \ref{lem:c4-6-2} and~\ref{lem:c4-6-3}, $R_{\Upsilon}\geqslant 2$, $u\leqslant u_*$, $v\leqslant 1$ and the small-coupling conditions of Sections~\ref{sec:c4-10-1} and~\ref{sec:c4-11-5}), the common small-coupling thresholds of Proposition~\ref{prop:regional-square-extraction} and \eqref{eq:c5-7-T}(ii), the thresholds of Lemmas~\ref{lem:c3-7-2} and~\ref{lem:c3-7-4} (required already for (ii), through $g_*$) and the further threshold of \eqref{eq:c5-7-expbound-I}; Corollaries~\ref{cor:main-factorisation} and~\ref{cor:main-leading-size} use in addition the further threshold stated in Corollary~\ref{cor:main-factorisation}. (H3) The torus is one of RG1's (p.251, where the period is $2L^m$ in physical units with the lattice spacing $\varepsilon= L^{-K}$, so that RG1's exponent is $m-K$ in terms of the $m$ here; the family of tori is the same): the fine periods are $N_{\mu}=2L^m$ fine-lattice units in every direction, with the integer $m$ so large that $N_{\mu}$ is a multiple of $\ell= LMR_1$ (a power of $L$: $M$ by \eqref{eq:c-standing-conditions}, $R_1$ by CRE p.246; the periods are then compatible and even, Section~\ref{sec:c3-7}) exceeding $200\ell$ (Section~\ref{sec:c5-3-1}), and every coarse period $K_{\mu}= N_{\mu}/L=2L^{m-1}$ satisfies $K_{\mu}\geqslant 3(n+2)$ and exceeds $\max(2,2R_{\mathrm p}+8)$ (\eqref{eq:c-standing-conditions}); for (iii) in addition $N_{\min}\geqslant 4R_0+(4/c)\log(4n)$ with the $c$ of Proposition~\ref{prop:c4c-5-1} (\eqref{eq:c5-7-T}(iii)). These are lower bounds on $m$ in terms of the fixed loop and the coupling; nothing is let tend to infinity. The collar conditions \eqref{eq:regional-collar-sizes} concern $M$, $M_2$ and $R_1$, not $m$, and are among the thresholds of (H2) (item 1 of Section~\ref{app:thresholds} and \eqref{eq:c5-7-T}(i)). The results of Section~\ref{sec:c3-7} and of the later sections are stated for general compatible even periods and are used here for these. (H4) Ba{\l}aban's results are used as stated and cited, in the readings recorded where CRE fixes no convention: the branch hypotheses (B1)--(B4) of Section~\ref{sec:branch-data} with CRE's derivation of (1.11)--(1.25) as recorded there (the readings that are this paper's are marked as such where they are made, and they are the following: the independence from the coupling and the branch of the constants that CRE does not display, $B_3$, $\delta_{\mathrm c}$ and the $O(1)$ of (B3)(iii), the $O(L^2)$ of p.247 behind the domain statement of (B1) and item 6 of Section~\ref{app:thresholds}, and the $O(1)$, $B_3$ and $\delta$ of the bound on $\mathbf H_{1,Ax}$ of p.251 behind the condition ``$MR_1$ large'' of item 5 of Section~\ref{app:thresholds} (the $O(1)$ of \cite[p.~25]{BalabanAVG} is displayed there, $d(2d+1)$, and is no reading); the reading of CRE's ``domain of integration in (1.11)'' at p.248 as the support of every characteristic factor of (1.11), under which $U_1\in D^{(11)}(V,V_0)$, (B1)(iii) and Remark~\ref{rem:c3-8-1}; the reading of CRE's $\delta_0$ as $(A_1/A_0)\varepsilon_0$ following p.265, which fixes the radius $r=\delta_0/g_0$ at which $\kappa_{C_n}$ is identified (Section~\ref{app:thresholds}); in (B1), that $U_1(\Omega_1\cap\Lambda_1^c,V)$ is the minimal configuration for the constraints named at p.251, and the dependence of (1.21) on $(V,V_0)$; the use, in (B1), of CRE's statements about the configurations $U_{1,\square'}(V)$ and $U_1(\Omega_1\cap\Lambda_1^c,V)$ for the Borel-chosen axial representatives, where the stabiliser of the coarse data leaves several; the convention $\chi_1(\square')=0$, $\chi_1^c(\square')=1$, $\chi_0'(\square')=0$, $\chi_0'^c(\square')=1$ where the configuration $U_{1,\square'}(V)$ does not exist, which keeps (1.4) and (1.8) partitions of unity; the outer-row reading of (B2); three corrections of CRE's printed text: $v$ read from (1.14) where CRE p.249 points to ``(1.15)'' (before \eqref{eq:retained-prechange-density}), CRE's ``$z\neq y$'' in (1.5) read as $x\neq y$, as in (1.7) (\eqref{eq:retained-outer-measure-identities}), and CRE's undefined ``compatible'' of partitions (p.245, where the partitions introduced there ``have to be compatible'', also with $\pi_0$ and with the partitions introduced later, and p.257, ``compatible with all other partitions'') read as: the cubes of the coarser of two compatible partitions are unions of cubes of the finer (Section~\ref{app:regular-covariance}); in (B4)(i), the identification of the defining equation of $\widetilde D$ with the row system, and of CRE's $\widetilde D$ in (1.13)--(1.14) with the chart-class family whose fine field lies in $D^{(11)}(V,V_0)$ (existence, uniqueness, analyticity and the positivity of the Jacobian (B4)(iii) are proved in Lemma~\ref{lem:c2-fibre-smallness}, from (B1) (with (B3)(i), through (B1)(i)) and (B2) with the readings listed here, the fine periods being at least $8M$, and threshold items 5 and 6); CRE's $\widetilde Q$ read as $LQ_{U_1}$ and CRE's $\sigma$ as the normalised Haar density, in the step list of Section~\ref{sec:branch-data}; the conventions of Section~\ref{sec:sourced-effective-interaction} on the bonds and rows of $\Omega_1^c$, on $X^{\sim}$ and on $M_1$; and, in Section~\ref{sec:c3-7}, ``regular'' read as membership in the class VAR (6) (Lemma~\ref{lem:c3-7-seed}) and the set $B(P_1^1)^c$ of Lemma~\ref{lem:c3-7-2}; the domain of CRE's integral (1.11) is otherwise CRE's: its averages and means are the principal logarithms of \cite[p.~21, (22)--(25)]{BalabanAVG}; on the whole domain of integration, the designated bonds included, the fluctuation lies in the chart ball of the conventions paragraph of Section~\ref{sec:sourced-effective-interaction} by CRE pp.247--248 (CRE's statements, which CRE justifies by the reasoning of Lemma~1 of \cite{BalabanREG}), and the arguments of the averages and means lie there by \cite[pp.~24--25, (44)--(46)]{BalabanAVG} with (B3)(i), p.248, both derived in (B1)), RG1 (0.4) and (0.10)--(0.11) for the averaging operation and the group means, Ba{\l}aban's transfer statements CRE p.256 (all the results of VAR used for CRE's averaging and determining sets), RG1 after (0.10) and RG1 after (0.12), through which Section~\ref{sec:c3-9} takes over VAR's constraint (3) and \cite[Proposition~1, (51)]{BalabanAVG} for the averaging RG1 (0.4), VAR's spacing condition (1) for CRE's determining sets $\mathbf B_1(Q_w^{[4]})$ and the restoring bound \eqref{eq:c3-9-7}, CRE (1.5)--(1.7) for the gauge-fixing functional, CRE (1.27) for $g_1$, p.246 for $R_1$ and $\varepsilon_1$, and CRE (2.4)--(2.9) for the relations between the scales $R_j$, $g_j$ and $\varepsilon_j$ (Section~\ref{app:thresholds}), CRE's conditions on the parameters as taken over in item 5 of Section~\ref{app:thresholds}, and the validity of CRE's construction for the values $M_1< M_2< M$ fixed there following CRE's prescription (the conventions after \eqref{eq:c-standing-conditions}). No property of this paper's own objects is a hypothesis: every lemma and proposition cited below is proved in the paper under (H1)--(H4).

\begin{theorem}[a Wilson loop after the first renormalisation step]\label{thm:main-theorem}
Under (H1)--(H4):
\begin{enumerate}[label=(\roman*)]
    \item (Form.)
    \begin{equation}\label{eq:thm-main-identity}
        \begin{gathered}
            \mathbb E_{\varepsilon}[W_{\chi}(C_n)]=\kappa_{C_n}\,\mathbb E_{\mathrm h}\bigl[\mathbf 1_{A_{\mathrm{cov}}}e^{E^{\circ}_{\xi}}W_{\chi}(C_n;U_1)\bigr]+E_{\mathrm{def}}+E_{\mathrm{bdry}}\;,\\
            \kappa_{C_n}= e^{-\delta m^{\sigma}(g_0,\varepsilon)P(C_n)}\bigl(z^{\sigma}_{\mathrm c}(g_0,\pi/2)\bigr)^4\neq0\;,
        \end{gathered}
    \end{equation}
    where the perimeter coefficient $\delta m^{\sigma}$ and the corner factor $z^{\sigma}_{\mathrm c}$ depend on $g_0$, $\varepsilon$, $L$, the group, the character and $\sigma$ only, the fixed averaging and mean prescriptions of RG1 (0.4) and (0.10)--(0.11), the gauge-fixing functional of CRE (1.5)--(1.7), the designated-bond convention (Section~\ref{sec:c4c-1}) and the cutoff test (1.18) at the radius $r= A_1(\log g_0^{-2})^{p_0}$ (CRE's constants $A_1$, $p_0$, under the reading of $\delta_0$ recorded in Section~\ref{app:thresholds} and (H4)) being part of the setting, so that $\kappa_{C_n}$ depends on $n$ through the perimeter $P(C_n)$ alone and not on the torus or the history, and every term is integrable against the history weight.
    \item (Bounds.)
    \begin{equation}\label{eq:thm-main-bounds}
        \begin{gathered}
            | E^{\circ}_{\xi}|\leqslant C\tau^2\Bigl[\sum_{j=0}^{4n-1}e^{-cd_j}+4ne^{-c\iota_T}+(n+1)e^{-cn}\Bigr]\leqslant C\tau^2(9n+1)\quad\text{on }A_{\mathrm{cov}}\;,\\
            | E_{\mathrm{def}}|\leqslant C'_Ln_Ce^{-(c_*/2)\pi_*^2}+e^{C_{\chi} n_C\rho(\varepsilon_1+e^{-cb/(LM)})}-1\;,\\
            | E_{\mathrm{bdry}}|\leqslant C_{\mathrm{resp},\chi}\,n_C(1+g_0r)\,e^{-c_*\pi_*^2}\;.
        \end{gathered}
    \end{equation}
    \item (Size of the leading term.) $|\mathbb E_{\mathrm h}[\mathbf 1_{A_{\mathrm{cov}}}e^{E^{\circ}_{\xi}}W_{\chi}(C_n;U_1)]|\leqslant e^{C\tau^2(9n+1)}$, and under the additional hypotheses of (H2), (H3) for (iii),
    \begin{equation}\label{eq:thm-main-moment}
        \mathbb E_{\mathrm h}\bigl[\mathbf 1_{A_{\mathrm{cov}}}e^{| E^{\circ}_{\xi}|}\bigr]\leqslant \exp\bigl[C\tau^2(2+1/c)+n_C\delta^I_q\bigr]\;,\qquad\delta^I_q\leqslant \exp\bigl[-c_q\pi_*^2/(2C_K)\bigr]\;,
    \end{equation}
    so that the leading term $\kappa_{C_n}\mathbb E_{\mathrm h}[\mathbf 1_{A_{\mathrm{cov}}}e^{E^{\circ}_{\xi}}W_{\chi}(C_n;U_1)]$ of \eqref{eq:thm-main-identity} has modulus at most $|\kappa_{C_n}|\exp[C\tau^2(2+1/c)+n_C\delta^I_q]$, and $\mathbb E_{\mathrm h}[\mathbf 1_{A_{\mathrm{cov}}}| e^{E^{\circ}_{\xi}}-1|]\leqslant \exp[C\tau^2(2+1/c)+n_C\delta^I_q]-1$.
\end{enumerate}
\end{theorem}

\begin{proof}
(i) is Corollary~\ref{cor:c4c-5-2}, whose hypotheses are the conventions of Section~\ref{sec:c4c-1} and those of Proposition~\ref{prop:regional-square-extraction} (the square geometry and collar conditions of (H3), the thresholds of (H2)); $\kappa_{C_n}\neq0$ is an exponential by \eqref{eq:marginal-square-mass} and \eqref{eq:corner-square-factor}, as stated in Corollary~\ref{cor:c5-7-2}(a) and after \eqref{eq:retained-outgoing-observable}, $\delta m^{\sigma}$ is the absolutely convergent sum \eqref{eq:marginal-square-mass}, characterised there as a limit, and $z^{\sigma}_{\mathrm c}$ the exponential \eqref{eq:corner-square-factor} of a convergent corner sum, from both of which $n$ and the torus are absent. In (ii) the first bound is Proposition~\ref{prop:c4c-5-1} with the elementary bounds ($e^{-cd_j}\leqslant 1$, $e^{-c\iota_T}\leqslant 1$, $(n+1)e^{-cn}\leqslant n+1$); the second is \eqref{eq:c4-12-8} of Proposition~\ref{prop:c4-12-3}, whose hypotheses are those stated at the start of Section~\ref{sec:c4-12} (Section~\ref{sec:sourced-effective-interaction}, Appendix~\ref{app:background-estimates}, the conventions of Section~\ref{sec:c4-9-2}) and the thresholds named in (H2); the third is \eqref{eq:c3-7-14} of Proposition~\ref{prop:c3-7-5}, under compatible even periods and \eqref{eq:c3-1-2} (a further condition of Section~\ref{app:thresholds}). In (iii) the first bound is the positivity of $\mathbb E_{\mathrm h}$ \eqref{eq:c3-6-1}, $\mathbb E_{\mathrm h}[1]=1$, $| W_{\chi}|\leqslant 1$ and the crude bound of (ii); \eqref{eq:thm-main-moment} and the two consequences are \eqref{eq:c5-7-expbound-I} of Corollary~\ref{cor:c5-7-4}(c) and the sentence following it, from Corollary~\ref{cor:c5-7-2}(c) applied to $E^{\circ}_{\xi}$ (whose bound has the form \eqref{eq:c5-7-2}) under $4ne^{-c\iota_T}\leqslant 1$, and from \eqref{eq:c5-7-expmoment-I} of Lemma~\ref{lem:c5-7-expmoment} at $\lambda= C\tau^2$, whose joint witness bound \eqref{eq:c5-7-joint-I} (Lemma~\ref{lem:c5-7-joint}) rests on Section~\ref{sec:c3-7}. The preamble of Corollary~\ref{cor:c5-7-4} (\eqref{eq:c5-7-T} and the hypotheses of Proposition~\ref{prop:c5-5R-4}, that is those of Proposition~\ref{prop:c5-5-9} and \eqref{eq:c5-5R-T}, and of Lemma~\ref{lem:c5-7-3}, the setting of Proposition~\ref{prop:c5-6-4}) consists of the square geometry and period conditions of (H3) and of thresholds of Sections~\ref{sec:c5-2}--\ref{sec:c5-5} listed among the further conditions of Section~\ref{app:thresholds}, so it is implied by (H2) and (H3); its item (c) at $D= I$ states the conditions it uses.
\end{proof}

\begin{remark}[regime of the statement]\label{rem:main-regime}
(1) The scales: as $g_0\to0$, $\tau= C[g_0(1+r)^3+e^{-cr^2}]\to0$ (with the single radius $r=\delta_0/g_0= A_1(\log g_0^{-2})^{p_0}$ of the notation paragraph and Appendix~\ref{app:thresholds}, the radius of the cutoff test $\Theta^{\sigma}$ as well), while $e^{-c_*\pi_*^2}$, $e^{-(c_*/2)\pi_*^2}$ and $\delta^I_q$ are smaller than every power of $g_0$ ($\pi_*^2\geqslant A_0^22^{-2p_0}(\log g_0^{-2})^{2p_0}$ by the bounds \eqref{eq:c5-3-1} as used in Corollary~\ref{cor:c5-7-4}(c), with $p_0\geqslant 5\mathfrak r\geqslant 10$, and for $\delta^I_q$ the last member of \eqref{eq:thm-main-moment} with the arithmetic preceding \eqref{eq:c5-7-delta-small}, $C_K\leqslant C(L^2M)^42^{4\mathfrak r}(\log g_0^{-2})^{4\mathfrak r}$ and $2p_0-4\mathfrak r\geqslant 12$), and $\rho\varepsilon_1\to0$. The $n$-dependence of the bounds is explicit: pointwise on $A_{\mathrm{cov}}$, $E^{\circ}_{\xi}$ is at most linear in $n$; under the torus condition $4ne^{-c\iota_T}\leqslant 1$ the growth is carried by the $4n$ boundary terms $e^{-cd_j}$, each at most one and exponentially small when the carrier complement is far from segment $j$, while the remaining terms $4ne^{-c\iota_T}+(n+1)e^{-cn}$ of (ii) are at most $2+1/c$ (Corollary~\ref{cor:c5-7-2}(c), applied to $E^{\circ}_{\xi}$ as in the proof above); in history average the exponential of $| E^{\circ}_{\xi}|$ carries the perimeter-rate factor $e^{n_C\delta^I_q}$ of (iii), where $\delta^I_q$ is governed by the $C_K$-th root $q^{1/C_K}$ of the one-point rate $q$ of Section~\ref{sec:c3-7}, while Corollary~\ref{cor:main-factorisation} below separates the deterministic torus and corner part of $E^{\circ}_{\xi}$ and pays its carrier part at the rate $q^{1/16}$, which gives (iii) with $\delta^{\sharp}_q$ in place of $\delta^I_q$; $E_{\mathrm{bdry}}$ and the first term of the bound \eqref{eq:thm-main-bounds} on $E_{\mathrm{def}}$ are perimeter rates of the same kind; and the second term of that bound is small only when $n_C\rho\varepsilon_1\ll1$ ($e^x-1\geqslant x$). The bounds \eqref{eq:thm-main-bounds} on $E_{\mathrm{def}}$ and $E_{\mathrm{bdry}}$ are absolute, not relative to $\kappa_{C_n}= e^{-\delta m^{\sigma} P(C_n)}(z^{\sigma}_{\mathrm c})^4$; whether they lie below $|\kappa_{C_n}|= e^{-\delta m^{\sigma} P(C_n)}| z^{\sigma}_{\mathrm c}|^4$ ($\delta m^{\sigma}$ is real, by the argument following \eqref{eq:marginal-square-mass}) depends on the sizes of $\delta m^{\sigma}$ and $z^{\sigma}_{\mathrm c}$, for which Sections~\ref{sec:marginal-perimeter} and~\ref{sec:marginal-corner} give two-sided bounds: $|\kappa^{\sigma}_v|\leqslant C\tau^2$ (the sentence following \eqref{eq:marginal-full-perimeter}, by \eqref{eq:marginal-cluster-tail}, with \eqref{eq:marginal-square-mass}) and $|\zeta^{\sigma}_{v,w}|\leqslant C\tau^2$ (\eqref{eq:corner-localisation-bound}, Section~\ref{sec:marginal-corner}) give $|\delta m^{\sigma}| P(C_n)= n|\sum_v\kappa^{\sigma}_v|\leqslant 4Cn\tau^2$ and $e^{-4C\tau^2}\leqslant | z^{\sigma}_{\mathrm c}|^4\leqslant e^{4C\tau^2}$, hence $e^{-4C\tau^2(n+1)}\leqslant |\kappa_{C_n}|\leqslant e^{4C\tau^2(n+1)}$, with $C$ independent of $g_0$, $n$ and the torus (an upper bound is also at hand from the expectation: at any configuration of $A_{\mathrm{cov}}$, $|\kappa_{C_n}|=| S^{\circ}_{\xi,C_n}| e^{-\Re E^{\circ}_{\xi}}\leqslant e^{C\tau^2(9n+1)}$ by (ii) and $| S^{\circ}_{\xi,C_n}|\leqslant 1$, the proof of Proposition~\ref{prop:c4-12-3}): $E_{\mathrm{bdry}}$ and the first term of the bound on $E_{\mathrm{def}}$ are smaller than every power of $g_0$ at fixed $n$, and the ratios of their bounds to $|\kappa_{C_n}|$ are, respectively, $C_{\mathrm{resp},\chi}n_C(1+g_0r)e^{-c_*\pi_*^2}/|\kappa_{C_n}|$ and $C'_Ln_Ce^{-(c_*/2)\pi_*^2}/|\kappa_{C_n}|$, below one exactly when $\log(C_{\mathrm{resp},\chi}n_C(1+g_0r))+\delta m^{\sigma} P(C_n)+4\log| z^{\sigma}_{\mathrm c}|^{-1}< c_*\pi_*^2$, respectively $\log(C'_Ln_C)+\delta m^{\sigma} P(C_n)+4\log| z^{\sigma}_{\mathrm c}|^{-1}<(c_*/2)\pi_*^2$, conditions that involve the length of the loop through $n_C$ and $P(C_n)$; by the lower bound on $|\kappa_{C_n}|$ they hold whenever $\log(C_{\mathrm{resp},\chi}n_C(1+g_0r))+4C\tau^2(n+1)< c_*\pi_*^2$, respectively $\log(C'_Ln_C)+4C\tau^2(n+1)<(c_*/2)\pi_*^2$, that is for $n$ up to a fixed multiple of $\pi_*^2/\tau^2$ at small coupling, and for every fixed $n$ once $g_0$ is small ($\tau\to0$, $\pi_*\to\infty$); when $\delta m^{\sigma}>0$ the exact conditions fail at fixed $g_0$ for all large $n$, which the hypotheses allow on correspondingly large tori. The second term of the bound on $E_{\mathrm{def}}$, at most $2x$ with $x= C_{\chi} n_C\rho(\varepsilon_1+e^{-cb/(LM)})$ when $x\leqslant 1$ ($e^x-1\leqslant 2x$ for $0\leqslant x\leqslant 1$), is bounded only by a quantity polynomially small in $g_0$ at fixed $n$ ($\rho= O(g_0^2(1+r)^4)$); its comparison with $|\kappa_{C_n}|$ at the original carrier is made by the relative form of \eqref{eq:c4-12-7}: the far-history term there retains $| S^{\circ}_{\xi,C}|=|\kappa_{C_n}|| e^{E^{\circ}_{\xi}}|$ on $A_{\mathrm{cov}}$ \eqref{eq:c4c-5-1}, so by \eqref{eq:thm-main-moment} (the far-history term is the $\mathbb E_{\mathrm h}$-average, with the weight $\mathsf w_{\xi}$ and the normaliser $\mathcal Z_0$ of (iii), of $\mathbf 1_{A_{\mathrm{cov}}\cap\mathsf N_b^c}| S^{\circ}_{\xi,C}|(e^x-1)\leqslant \mathbf 1_{A_{\mathrm{cov}}}|\kappa_{C_n}| e^{| E^{\circ}_{\xi}|}(e^x-1)$) it is at most $|\kappa_{C_n}| e^{C\tau^2(2+1/c)+n_C\delta^I_q}(e^x-1)$, a relative bound, also with $\delta^{\sharp}_q$ in place of $\delta^I_q$ (Corollary~\ref{cor:main-factorisation}); only the near-history term $C'_Ln_Ce^{-(c_*/2)\pi_*^2}$ is absolute. (2) At the retreated carrier of the labelled history terms, Section~\ref{sec:one-step-insertion} gives the same identification $S^{\circ}_{\Lambda}=\kappa_Ce^{E^{\circ}_{\Lambda}}$ on $A^{\mathrm{fin}}_{\mathrm{cov}}$ (Corollary~\ref{cor:c5-7-2}(a)), the bound on the source-containing boundary terms relative to the leading coefficient (Corollary~\ref{cor:c5-7-4}(a); its factor $e^{C\tau^2(9n+1)}$ is exponential in $n$ at fixed $g_0$, so that bound is small at fixed $n_C$ only) and the bound \eqref{eq:c5-6-def-2} on the labelled defect (Proposition~\ref{prop:c5-6-def}; its near-history part $C'_Ln_Ce^{-(c_*/2)\pi_*^2}$ carries no factor $|\kappa_C|$ and is absolute, while its far-history part retains $| S^{\circ}_{\Lambda}|=|\kappa_C|| e^{E^{\circ}_{\Lambda}}|$ on $A^{\mathrm{fin}}_{\mathrm{cov}}$ and is relative by Corollary~\ref{cor:c5-7-2}(a) and the exponential moment \eqref{eq:c5-7-expbound} of Corollary~\ref{cor:c5-7-4}(c), as at the original carrier in (1)), and the first and exponential moments of $E^{\circ}_{\Lambda}$ (Corollary~\ref{cor:c5-7-4}(b),(c)). (3) Not proved here: a bound on the history average of $e^{E^{\circ}_{\xi}}$ literally independent of $n$ (the proved form is the perimeter rate of (iii)); a lower bound on $|\kappa_{C_n}|$ uniform in $n$ ($|\kappa_{C_n}|= e^{-\delta m^{\sigma} P(C_n)}| z^{\sigma}_{\mathrm c}|^4$ is bounded below uniformly in $n$ if and only if $\delta m^{\sigma}\leqslant 0$, the sign of $\delta m^{\sigma}$ is not determined here, and the lower bound $e^{-4C\tau^2(n+1)}$ of (1) decays in $n$); a lower bound on the modulus of the leading expectation $\mathbb E_{\mathrm h}[\mathbf 1_{A_{\mathrm{cov}}}e^{E^{\circ}_{\xi}}W_{\chi}(C_n;U_1)]$ uniform in $n$ (Corollary~\ref{cor:main-leading-size} below gives one at a fixed loop, with an error $\eta_n= a_n+e^{n_C\delta^{\sharp}_q}-1$ whose leading part for $n\leqslant \varepsilon_1^{-1}$ and nontrivial $\chi$ is the area term $2\mathrm{Lip}_{\chi} c_{\mathrm h}(nL)^2\varepsilon_1$, and when $\eta_n<1$ ``below $|\kappa_{C_n}|$'' in (1) compares with the leading term up to the factor $(1-\eta_n)^{-1}e^{C\tau^2(2+1/c)}$); and the second renormalisation step: neither the coupling nor the scale is iterated, and the output of the step is the pair of the complex measure $\kappa_{C_n}^{-1}N_1$ and the measure $\rho_1$ as described in Section~\ref{sec:remaining-physical-comparison}, not a statement about a coarse-lattice Wilson loop against a local effective density (Conjecture~\ref{conj:iteration} of Section~\ref{sec:next-step} states the form expected after $k+1$ steps for $n= L^k$).
\end{remark}

\begin{corollary}[factorisation of the leading term]\label{cor:main-factorisation}
Under (H1)--(H4), including the additional hypotheses of (H2) and (H3) for (iii) of Theorem~\ref{thm:main-theorem}, write the exponent \eqref{eq:c4c-5-2} on $A_{\mathrm{cov}}$ as $E^{\circ}_{\xi}= E^{\partial}_{\xi}+E^{\circ}_{T,n}$, with the carrier part and the torus and corner part
\begin{equation}\label{eq:main-factorisation-parts}
    E^{\partial}_{\xi}= F_{I,C,T}-F_{C,T}\;,\qquad E^{\circ}_{T,n}=(F_{C,T}-F_{C,\infty}(n))+E^{\sigma}_n
\end{equation}
in the notation of Proposition~\ref{prop:c4c-5-1}. Then $E^{\circ}_{T,n}$ does not depend on the history, $\kappa_{C_n}e^{E^{\circ}_{T,n}}=\mathcal W^{\circ,\sigma}_T(C_n)$ is the flat reference loop \eqref{eq:marginal-literal-ratio} of $C_n$ on the torus, and
\begin{equation}\label{eq:main-factorisation}
    \begin{gathered}
        \kappa_{C_n}\mathbb E_{\mathrm h}\bigl[\mathbf 1_{A_{\mathrm{cov}}}e^{E^{\circ}_{\xi}}W_{\chi}(C_n;U_1)\bigr]=\mathcal W^{\circ,\sigma}_T(C_n)\bigl(\mathbb E_{\mathrm h}\bigl[\mathbf 1_{A_{\mathrm{cov}}}W_{\chi}(C_n;U_1)\bigr]+R_n\bigr)\;,\\
        | E^{\circ}_{T,n}|\leqslant C\tau^2\bigl[4ne^{-c\iota_T}+(n+1)e^{-cn}\bigr]\leqslant C\tau^2(2+1/c)\;,\\
        | R_n|\leqslant \mathbb E_{\mathrm h}\bigl[\mathbf 1_{A_{\mathrm{cov}}}| e^{E^{\partial}_{\xi}}-1|\bigr]\leqslant e^{n_C\delta^{\sharp}_q}-1\;,\qquad
        \delta^{\sharp}_q\coloneqq C\tau^2N'_{\mathrm{lat}}(c)\,e^{C\tau^2L^4C_K\Phi_c}\,q^{1/16}\;,
    \end{gathered}
\end{equation}
with $C$, $c$ of Proposition~\ref{prop:c4c-5-1}, $N'_{\mathrm{lat}}(c)$ and $\Phi_c$ of Lemma~\ref{lem:c5-7-expmoment} (they depend on $L$ and $c$ only, \eqref{eq:c5-7-boundary-marks}), $C_K$ of Lemma~\ref{lem:c5-7-joint} and the one-point rate $q\leqslant e^{-c_q\pi_*^2}$ of Section~\ref{sec:c3-7} ($g_0< g_*$, as in \eqref{eq:c5-7-expbound-I}). For $g_0$ below a further threshold, depending on $L$, $M$, $M_1$, $M_2$, the group, the character, the source constants and $c$, $C\tau^2L^4C_K\Phi_c\leqslant 1$ and $\delta^{\sharp}_q\leqslant eC\tau^2N'_{\mathrm{lat}}(c)\,e^{-c_q\pi_*^2/16}$: the carrier part is paid at a fixed fraction of the exponent of the one-point rate, with no loss growing as $g_0\to0$. Moreover $\mathbb E_{\mathrm h}[\mathbf 1_{A_{\mathrm{cov}}}e^{| E^{\circ}_{\xi}|}]\leqslant \exp[C\tau^2(2+1/c)+n_C\delta^{\sharp}_q]$, the bound \eqref{eq:thm-main-moment} of (iii) with $\delta^{\sharp}_q$ in place of $\delta^I_q$. Consequently $\mathbb E_{\varepsilon}[W_{\chi}(C_n)]=\mathcal W^{\circ,\sigma}_T(C_n)(\mathbb E_{\mathrm h}[\mathbf 1_{A_{\mathrm{cov}}}W_{\chi}(C_n;U_1)]+R_n)+E_{\mathrm{def}}+E_{\mathrm{bdry}}$ by \eqref{eq:thm-main-identity}.
\end{corollary}

\begin{proof}
The decomposition is \eqref{eq:c4c-5-2}. In the proof of Proposition~\ref{prop:regional-square-extraction} at $D= I$, from which Proposition~\ref{prop:c4c-5-1} is obtained, the three parts are bounded separately: $| E^{\partial}_{\xi}|\leqslant C\tau^2\sum_{j=0}^{4n-1}e^{-cd_j}$ by \eqref{eq:regional-square-boundary-comparison}, $| F_{C,T}-F_{C,\infty}(n)|\leqslant C\tau^2(4n)e^{-c\iota_T}$ by \eqref{eq:regional-square-periodic-comparison}, and $| E^{\sigma}_n|\leqslant C\tau^2(n+1)e^{-cn}$ by \eqref{eq:corner-square-extraction-bound}; the bound of Proposition~\ref{prop:c4c-5-1} is their sum, with the largest of the three constants and the smallest of the three rates; these are the $C$, $c$ of that proposition fixed in the notation paragraph, the $c$ of (H3) and \eqref{eq:c5-7-T}(iii) is this one, and each of the three bounds holds with them. $F_{C,T}$ is the logarithm, on the connected source branch, of $\mathcal W^{\circ,\sigma}_T(C_n)= N^{\sigma}_{C_n,T}/D^{\sigma}_T$ \eqref{eq:marginal-literal-ratio}, $F_{C,\infty}(n)$ is the bulk logarithm \eqref{eq:corner-fixed-square-limit} and $E^{\sigma}_n$ the corner-extraction error; none of them involves the carrier $I$ or any other datum of the history, so $E^{\circ}_{T,n}$ is a number fixed by $g_0$, $n$, the torus and the setting. By the identity $F_{I,C,T}=\log\kappa_{C_n}+E^{\circ}_{\xi}$ of Proposition~\ref{prop:c4c-5-1}, $\log\kappa_{C_n}+E^{\circ}_{T,n}= F_{I,C,T}-E^{\partial}_{\xi}= F_{C,T}$, whence $\kappa_{C_n}e^{E^{\circ}_{T,n}}=\mathcal W^{\circ,\sigma}_T(C_n)$. The second form of the bound on $E^{\circ}_{T,n}$ uses $4ne^{-c\iota_T}\leqslant 1$, which is \eqref{eq:c5-7-T}(iii) in (H3), and $(n+1)e^{-cn}\leqslant 1+1/c$. Since $E^{\circ}_{T,n}$ is a constant and the integrands are bounded ($| E^{\partial}_{\xi}|\leqslant 4C\tau^2n$ on $A_{\mathrm{cov}}$, $| W_{\chi}|\leqslant 1$), $\mathbb E_{\mathrm h}[\mathbf 1_{A_{\mathrm{cov}}}e^{E^{\circ}_{\xi}}W_{\chi}(C_n;U_1)]= e^{E^{\circ}_{T,n}}\mathbb E_{\mathrm h}[\mathbf 1_{A_{\mathrm{cov}}}e^{E^{\partial}_{\xi}}W_{\chi}(C_n;U_1)]$, and the identity holds with $R_n=\mathbb E_{\mathrm h}[\mathbf 1_{A_{\mathrm{cov}}}(e^{E^{\partial}_{\xi}}-1)W_{\chi}(C_n;U_1)]$. By the positivity of $\mathbb E_{\mathrm h}$ \eqref{eq:c3-6-1}, $| W_{\chi}|\leqslant 1$, $e^{| x|}-1\geqslant | e^x-1|$ and the pointwise bound of \eqref{eq:c5-7-expmoment-I}, with $\lambda= C\tau^2$ and the marks $m_b=\mathbf 1\{ b\in E(B_{\xi})\}$, $| R_n|\leqslant \mathbb E_{\mathrm h}[\mathbf 1_{A_{\mathrm{cov}}}| e^{E^{\partial}_{\xi}}-1|]\leqslant \mathbb E_{\mathrm h}[\exp(\lambda\sum_b\varphi_c(b)m_b)]-1$, using $\mathbb E_{\mathrm h}[1]=1$. Let $L^p$ be the smallest power of $L$ exceeding $\sigma_K=2K\ell+4$ (Lemma~\ref{lem:c5-7-joint}), so $L^p\leqslant L\sigma_K$, and put $s_{\mu}=\min(L^p,N_{\mu}/2)$; $N_{\mu}/s_{\mu}$ is even, as $N_{\mu}=2L^m$ by (H3). Partition the sites of the torus into the boxes $Q=\prod_{\mu}[k_{\mu} s_{\mu},(k_{\mu}+1)s_{\mu})$, $k_{\mu}\in\mathbb Z/(N_{\mu}/s_{\mu})$, and colour $Q$ by $(k_{\mu}\bmod2)_{\mu}$, sixteen colours. Two distinct boxes of one colour differ in some direction $\mu$ by a nonzero even index difference modulo $N_{\mu}/s_{\mu}$, which forces $N_{\mu}/s_{\mu}\geqslant 4$ and $s_{\mu}= L^p$, so their sites are at periodic $\ell^{\infty}$ distance at least $L^p>\sigma_K$. For a finite set $\mathsf S$ of bonds let $N(\mathsf S)$ be the number of boxes containing the initial site of a bond of $\mathsf S$; one bond of $\mathsf S$ in each such box of the most frequent colour gives a subset $\mathsf S'$ at pairwise $\ell^{\infty}$ distance $>\sigma_K$ with $|\mathsf S'|\geqslant N(\mathsf S)/16$, so the first inequality of \eqref{eq:c5-7-joint-I} and $q\leqslant 1$ give $\mathbb E_{\mathrm h}[\prod_{b\in\mathsf S}m_b]\leqslant q^{N(\mathsf S)/16}$. Expand $\exp(\lambda\sum_b\varphi_c(b)m_b)=\prod_b(1+y_bm_b)$, $y_b= e^{\lambda\varphi_c(b)}-1\geqslant 0$, into the sum over $\mathsf S$ of $\prod_{b\in\mathsf S}y_bm_b$, as in the proof of Lemma~\ref{lem:c5-7-expmoment} (finitely many nonnegative terms); since $N(\mathsf S)=\sum_Q\mathbf 1\{\mathsf S\text{ has a bond with initial site in }Q\}$, the sum factorises over the boxes:
\begin{equation*}
    \begin{aligned}
        \mathbb E_{\mathrm h}\Bigl[\prod_b(1+y_bm_b)\Bigr]&\leqslant \sum_{\mathsf S}\prod_{b\in\mathsf S}y_b\,q^{N(\mathsf S)/16}=\prod_Q\Bigl(1+q^{1/16}\Bigl(\prod_{b\in Q}(1+y_b)-1\Bigr)\Bigr)\\
        &\leqslant \exp\Bigl(q^{1/16}\sum_Q\bigl(e^{\lambda\sum_{b\in Q}\varphi_c(b)}-1\bigr)\Bigr)\;,
    \end{aligned}
\end{equation*}
the products and sums over $b\in Q$ running over the bonds with initial site in $Q$. A box has at most $L^{4p}\leqslant L^4\sigma_K^4$ sites, hence at most $4L^4\sigma_K^4\leqslant L^4C_K$ bonds, so $\lambda\sum_{b\in Q}\varphi_c(b)\leqslant \lambda L^4C_K\Phi_c$ by \eqref{eq:c5-7-boundary-marks}; with $e^x-1\leqslant xe^x$ and $\sum_b\varphi_c(b)\leqslant n_CN'_{\mathrm{lat}}(c)$ (proof of Lemma~\ref{lem:c5-7-expmoment}) the exponent is at most $n_C\lambda N'_{\mathrm{lat}}(c)e^{\lambda L^4C_K\Phi_c}q^{1/16}= n_C\delta^{\sharp}_q$. The same pointwise bound gives $\mathbb E_{\mathrm h}[\mathbf 1_{A_{\mathrm{cov}}}e^{| E^{\partial}_{\xi}|}]\leqslant \mathbb E_{\mathrm h}[\exp(\lambda\sum_b\varphi_c(b)m_b)]\leqslant e^{n_C\delta^{\sharp}_q}$, and with $| E^{\circ}_{\xi}|\leqslant | E^{\circ}_{T,n}|+| E^{\partial}_{\xi}|$ the moment bound follows. For the further threshold: $C_K\leqslant 4(4K+9)^4(L^2M)^42^{4\mathfrak r}(\log g_0^{-2})^{4\mathfrak r}$ and $\tau\leqslant C[g_0(1+r)^3+e^{-cr^2}]$ with $r= A_1(\log g_0^{-2})^{p_0}$ (the comparison in Corollary~\ref{cor:c5-7-4}(c)), so $\tau^2C_K\to0$ as $g_0\to0$, and $q^{1/16}\leqslant e^{-c_q\pi_*^2/16}$. The last statement is \eqref{eq:thm-main-identity} with the identity of \eqref{eq:main-factorisation} substituted.
\end{proof}

\begin{corollary}[the leading term at a fixed loop]\label{cor:main-leading-size}
Under the hypotheses of Corollary~\ref{cor:main-factorisation}, let $\pi_{\chi}$ be the irreducible unitary representation of $G$ with $\chi=(\dim\pi_{\chi})^{-1}\operatorname{Tr}\pi_{\chi}$, $\mathrm{Lip}_{\chi}$ a constant with $\|\pi_{\chi}(g)-\pi_{\chi}(h)\|_{\mathrm{op}}\leqslant \mathrm{Lip}_{\chi}\| g-h\|_{\mathrm{op}}$ for $g,h\in G$ (it exists because $\pi_{\chi}$ is smooth and $G\subset U(N)$ is a compact submanifold, and it depends on the group and the character only), $C_L$ the contour constant with $|\Gamma_C|\leqslant C_Ln_C$ (Section~\ref{sec:c3-7}, the paragraph before Lemma~\ref{lem:c3-7-1}), depending on $L$ only, and $q$ the one-point rate of Lemma~\ref{lem:c3-7-4}. Then
\begin{equation}\label{eq:main-leading-size}
    \begin{gathered}
        \bigl|\mathbb E_{\mathrm h}\bigl[\mathbf 1_{A_{\mathrm{cov}}}W_{\chi}(C_n;U_1)\bigr]-1\bigr|\leqslant a_n\coloneqq 2\mathrm{Lip}_{\chi} c_{\mathrm h}(nL)^2\varepsilon_1+\bigl(2(nL+1)^2+C_Ln_C\bigr)q\;,\\
        \kappa_{C_n}\mathbb E_{\mathrm h}\bigl[\mathbf 1_{A_{\mathrm{cov}}}e^{E^{\circ}_{\xi}}W_{\chi}(C_n;U_1)\bigr]=\mathcal W^{\circ,\sigma}_T(C_n)\bigl(1+R^{\sharp}_n\bigr)\;,\qquad| R^{\sharp}_n|\leqslant \eta_n\coloneqq a_n+e^{n_C\delta^{\sharp}_q}-1\;,
    \end{gathered}
\end{equation}
with $c_{\mathrm h}=\sqrt N$, $\varepsilon_1= g_1p_0(g_1)$ and $\delta^{\sharp}_q$ of \eqref{eq:main-factorisation}. Hence, when $\eta_n<1$, the leading term of \eqref{eq:thm-main-identity} is nonzero, $(1-\eta_n)|\mathcal W^{\circ,\sigma}_T(C_n)|\leqslant |\kappa_{C_n}\mathbb E_{\mathrm h}[\mathbf 1_{A_{\mathrm{cov}}}e^{E^{\circ}_{\xi}}W_{\chi}(C_n;U_1)]|\leqslant (1+\eta_n)|\mathcal W^{\circ,\sigma}_T(C_n)|$, and $e^{-C\tau^2(2+1/c)}|\kappa_{C_n}|\leqslant |\mathcal W^{\circ,\sigma}_T(C_n)|\leqslant e^{C\tau^2(2+1/c)}|\kappa_{C_n}|$ with the $C$, $c$ of Corollary~\ref{cor:main-factorisation}. For every fixed $n$, $\eta_n\to0$ as $g_0\to0$; at small coupling $\eta_n$ is small for $n$ up to a small multiple of $L^{-1}\varepsilon_1^{-1/2}$, and the bound is not uniform in $n$.
\end{corollary}

\begin{proof}
Let $S_n$ be the set of the $(nL+1)^2$ fine sites of the closed square bounded by $C_n$ in its plane, and $A_{\square}=\{ S_n\subseteq\Lambda_1\}$, an event of the regions. If $A_{\square}$ fails, some $x\in S_n$ lies in $B_{\xi}=\Lambda_1^c$ and the positive bond $z_x=(x,x+e_1)$ lies in $E(B_{\xi})$ (notation of Section~\ref{sec:c3-7}); so $\mathbf 1_{A_{\square}^c}\leqslant \sum_{x\in S_n}\mathbf 1\{ z_x\in E(B_{\xi})\}$, and $\mathbb P_{\mathrm h}(A_{\square}^c)\leqslant (nL+1)^2q$ by $p_z\leqslant q$ (Lemma~\ref{lem:c3-7-4}) and the positivity of $\mathbb E_{\mathrm h}$ \eqref{eq:c3-6-1}. On a branch in $A_{\square}$ every fine plaquette $p$ of the square has its vertices in $\Lambda_1$, a union of blocks: the block of a vertex $x$ lies in $\Lambda_1$, a bond with both endpoints in that block is free (a designated bond crosses between two blocks) and belongs to $I$, and its initial site is within $\ell^{\infty}$ distance $L-1$ of $x$, so every bond of $p$ is within distance $L\leqslant 3R_{\mathrm c}$ of $I$, and Lemma~\ref{lem:c2-plaquette-deviation} gives $\| U_1(\partial p)-\id\|_{\mathrm{op}}\leqslant 2c_{\mathrm h}\varepsilon_1$. Lattice Stokes: order the $(nL)^2$ plaquettes of the square by rows, starting at the base corner $x_0$ of $C_n$, the image of the vertex $0$ of the square of the notation paragraph under the coarse translation that gives $C_n$ (bottom row first, left to right in each row), and let $Q_k$ be the union of the first $k$ of them, whose boundary, oriented counterclockwise and based at $x_0$, is a closed lattice path; $Q_1= p_1$ is the plaquette at $x_0$, for $k\geqslant 1$ the shared path $s$ below does not contain $x_0$, and $\partial Q_{(nL)^2}$ is $C_n$ (another base point or orientation replaces the holonomy $h$ below by a conjugate or by $h^{-1}$, which changes neither $\| h-\id\|_{\mathrm{op}}$ nor $|\chi(h)-1|$, since $\chi(h^{-1})=\overline{\chi(h)}$ and $\chi(\id)=1$); adding the next plaquette $p$ replaces the one or two bonds $s$ it shares with $\partial Q_k$ (its lower and left sides that lie in $Q_k$) by the complementary path $s'$ of $\partial p$, and if the based boundary path of $Q_k$ is $\alpha s\beta$, then $\alpha s'\beta$ is that of $Q_{k+1}$ and its holonomy is $vU(\partial p)^{\pm1}v^{-1}$ times the holonomy of $\alpha s\beta$, with $v$ the holonomy of $\alpha s$ and $\partial p$ based at the endpoint of $s$. For unitaries, $\| vwv^{-1}-\id\|=\| w-\id\|=\| w^{-1}-\id\|$ and $\| w_1w_2-\id\|\leqslant \| w_1-\id\|+\| w_2-\id\|$, so on $A_{\square}$ the holonomy $h$ of $U_1$ around $C_n$ satisfies $\| h-\id\|_{\mathrm{op}}\leqslant 2c_{\mathrm h}(nL)^2\varepsilon_1$, and, $\chi$ being a class function with $|\operatorname{Tr}X|\leqslant \dim\pi_{\chi}\,\| X\|_{\mathrm{op}}$, $| W_{\chi}(C_n;U_1)-1|=|\chi(h)-\chi(\id)|\leqslant \|\pi_{\chi}(h)-\id\|_{\mathrm{op}}\leqslant 2\mathrm{Lip}_{\chi} c_{\mathrm h}(nL)^2\varepsilon_1$. Since $\mathbb E_{\mathrm h}$ is positive with $\mathbb E_{\mathrm h}[1]=1$,
\begin{equation*}
    \mathbb E_{\mathrm h}\bigl[\mathbf 1_{A_{\mathrm{cov}}}W_{\chi}(C_n;U_1)\bigr]-1=\mathbb E_{\mathrm h}\bigl[\mathbf 1_{A_{\mathrm{cov}}\cap A_{\square}}(W_{\chi}-1)\bigr]+\mathbb E_{\mathrm h}\bigl[\mathbf 1_{A_{\mathrm{cov}}\cap A_{\square}^c}(W_{\chi}-1)\bigr]-\mathbb P_{\mathrm h}(A_{\mathrm{cov}}^c)\;,
\end{equation*}
and the three terms are at most $2\mathrm{Lip}_{\chi} c_{\mathrm h}(nL)^2\varepsilon_1$, $2\mathbb P_{\mathrm h}(A_{\square}^c)\leqslant 2(nL+1)^2q$ ($| W_{\chi}|\leqslant 1$) and $C_Ln_Cq$ \eqref{eq:c3-7-3} in modulus, which is the first bound. By Corollary~\ref{cor:main-factorisation} the leading term is $\mathcal W^{\circ,\sigma}_T(C_n)(\mathbb E_{\mathrm h}[\mathbf 1_{A_{\mathrm{cov}}}W_{\chi}(C_n;U_1)]+R_n)=\mathcal W^{\circ,\sigma}_T(C_n)(1+R^{\sharp}_n)$ with $R^{\sharp}_n=(\mathbb E_{\mathrm h}[\mathbf 1_{A_{\mathrm{cov}}}W_{\chi}(C_n;U_1)]-1)+R_n$ and $| R_n|\leqslant e^{n_C\delta^{\sharp}_q}-1$, which is the second; the modulus bounds follow from $1-\eta_n\leqslant |1+R^{\sharp}_n|\leqslant 1+\eta_n$ and $\mathcal W^{\circ,\sigma}_T(C_n)=\kappa_{C_n}e^{E^{\circ}_{T,n}}$ with $| E^{\circ}_{T,n}|\leqslant C\tau^2(2+1/c)$. As $g_0\to0$, $\varepsilon_1\to0$ ($g_1\leqslant (1+\beta_0)g_0$ and $g_0\leqslant (1+\beta_0)g_1$, CRE (2.6) as quoted in Section~\ref{app:thresholds}; the second inequality keeps $\varepsilon_1^{-1}$ polynomial in $g_0^{-1}$), $q\leqslant e^{-c_*\pi_*^2}$ (Proposition~\ref{prop:c3-7-5}) and, below the further threshold of Corollary~\ref{cor:main-factorisation}, $\delta^{\sharp}_q\leqslant eC\tau^2N'_{\mathrm{lat}}(c)e^{-c_q\pi_*^2/16}$, with $\pi_*\to\infty$; so $\eta_n\to0$ at fixed $n$, and the terms other than $2\mathrm{Lip}_{\chi} c_{\mathrm h}(nL)^2\varepsilon_1$ are smaller than every power of $g_0$ for $n\leqslant \varepsilon_1^{-1}$.
\end{proof}

\begin{theorem}[two disjoint squares after the first renormalisation step]\label{thm:two-loop-theorem}
Let $C= C_1\sqcup C_2$, $\chi_1$, $\chi_2$, $n_C= n_{C_1}+n_{C_2}$ and $W_{\chi_1,\chi_2}(C;U)=\chi_1(\Hol_UC_1)\chi_2(\Hol_UC_2)$ be as in Section~\ref{sec:two-loops-setting}: $C_i$ is a coarse translate of the counterclockwise square of side $n_iL$ in a plane parallel to that of $e_1,e_2$, $\chi_i$ a normalised irreducible character, and $V(C_1)\cap V(C_2)=\varnothing$. Assume (H1)--(H4) for $(C_1,\chi_1)$ and for $(C_2,\chi_2)$ in place of $(C_n,\chi)$, and let $g_0$ lie below the thresholds of Section~\ref{sec:two-loops}: those of Lemmas~\ref{lem:two-loop-source}, \ref{lem:two-loop-pins} and~\ref{lem:two-loop-cross}, and for (iii) that of \eqref{eq:two-loop-moment}. Let $A_{\mathrm{cov}}$ be the covered event of the pair, $E_{\mathrm{def}}$ and $E_{\mathrm{bdry}}$ the terms \eqref{eq:c4-1-1} and \eqref{eq:c3-6-2} formed with the observable and the groups of the pair (Section~\ref{sec:two-loops}), $\kappa_{C_i}$ and $E^{\circ}_{\xi,i}$ the coefficient and the exponent of Theorem~\ref{thm:main-theorem} for $(C_i,\chi_i)$, $E^{\times}_{\xi}$ the cross exponent of Lemma~\ref{lem:two-loop-cross}(b), and $\Delta_{\times}$, $d$ as in \eqref{eq:two-loop-cross-sum}. Then:
\begin{enumerate}[label=(\roman*)]
    \item (Form.)
    \begin{equation}\label{eq:thm-two-identity}
        \mathbb E_{\varepsilon}[W_{\chi_1,\chi_2}(C)]=\kappa_{C_1}\kappa_{C_2}\,\mathbb E_{\mathrm h}\bigl[\mathbf 1_{A_{\mathrm{cov}}}e^{E^{\circ}_{\xi,1}+E^{\circ}_{\xi,2}+E^{\times}_{\xi}}W_{\chi_1,\chi_2}(C;U_1)\bigr]+E_{\mathrm{def}}+E_{\mathrm{bdry}}\;,
    \end{equation}
    every term being integrable against the history weight; $\kappa_{C_i}\neq0$ depends on $n_i$ through $P(C_i)$ alone, as in Theorem~\ref{thm:main-theorem}(i).
    \item (Bounds.) On $A_{\mathrm{cov}}$, $| E^{\circ}_{\xi,i}|\leqslant C\tau^2\bigl[\sum_{j\in J_{C_i}}e^{-cd_j}+4n_ie^{-c\iota_T}+(n_i+1)e^{-cn_i}\bigr]\leqslant C\tau^2(9n_i+1)$ for $i=1,2$, with $J_{C_i}$ the $4n_i$ coarse segments of $C_i$, and
    \begin{equation}\label{eq:thm-two-bounds}
        \begin{gathered}
            | E^{\times}_{\xi}|\leqslant C_{\times}\rho\Delta_{\times}\leqslant C_{\times}\rho\,n_Ce^{-c_{\times} d/(LM)}\;,\\
            | E_{\mathrm{def}}|\leqslant e^{C_{\times}\rho\Delta_{\times}}\Bigl[C'_Ln_Ce^{-(c_*/2)\pi_*^2}+e^{C_{\chi} n_C\rho(\varepsilon_1+e^{-cb/(LM)})+C_{\times}\rho\Delta_{\times}}-1\Bigr]\;,\\
            | E_{\mathrm{bdry}}|\leqslant C_{\mathrm{resp},\chi_1\chi_2}\,n_C(1+g_0r)\,e^{-c_*\pi_*^2}\;.
        \end{gathered}
    \end{equation}
    \item (Size of the leading term.) $|\mathbb E_{\mathrm h}[\mathbf 1_{A_{\mathrm{cov}}}e^{E^{\circ}_{\xi,1}+E^{\circ}_{\xi,2}+E^{\times}_{\xi}}W_{\chi_1,\chi_2}(C;U_1)]|\leqslant e^{C\tau^2(9n_1+9n_2+2)+C_{\times}\rho\Delta_{\times}}$, and under the condition of (H3) for (iii) for each square,
    \begin{equation}\label{eq:thm-two-moment}
        \mathbb E_{\mathrm h}\bigl[\mathbf 1_{A_{\mathrm{cov}}}e^{| E^{\circ}_{\xi,1}|+| E^{\circ}_{\xi,2}|}\bigr]\leqslant \exp\bigl[2C\tau^2(2+1/c)+n_C\delta^{I,2}_q\bigr]\;,\qquad\delta^{I,2}_q\leqslant \exp\bigl[-c_q\pi_*^2/(2C_K)\bigr]\;,
    \end{equation}
    with $\delta^{I,2}_q$ of \eqref{eq:two-loop-moment}, so that the leading term of \eqref{eq:thm-two-identity} has modulus at most $|\kappa_{C_1}\kappa_{C_2}|\exp[2C\tau^2(2+1/c)+n_C\delta^{I,2}_q+C_{\times}\rho\Delta_{\times}]$, and $\mathbb E_{\mathrm h}[\mathbf 1_{A_{\mathrm{cov}}}| e^{E^{\circ}_{\xi,1}+E^{\circ}_{\xi,2}+E^{\times}_{\xi}}-1|]\leqslant \exp[2C\tau^2(2+1/c)+n_C\delta^{I,2}_q+C_{\times}\rho\Delta_{\times}]-1$.
\end{enumerate}
Here $C$ and the $c$ of the bounds on $E^{\circ}_{\xi,i}$ and of (iii) are the larger, respectively the smaller, of the constants of Proposition~\ref{prop:c4c-5-1} for the two characters, and $d_j$ is the distance \eqref{eq:c4c-1-1} of segment $j$ of $C_i$ at the carrier $I=\Lambda_1^*$ (not the distance $d$ of \eqref{eq:two-loop-cross-sum}); $C_{\chi}=\max_iC_{\chi_i}$ and the $c$ of the bound on $E_{\mathrm{def}}$ are those of Lemma~\ref{lem:two-loop-cross}, formed from the constants of \eqref{eq:c4-6-8}, and $C_{\times}$, $c_{\times}$ the constants of that lemma; $\rho$ is \eqref{eq:c4-4-8} with $u= C^{\mathrm{src}}_{\chi_1\chi_2}g_0(1+r)$, the enlarged source constant of Lemma~\ref{lem:two-loop-source}, which is distinct from this $C_{\chi}$; $C_{\mathrm{resp},\chi_1\chi_2}= C_L\max((C_{\mathrm{resp},\chi_1}+C_{\mathrm{resp},\chi_2})S_c,1)$ with the constants $C_{\mathrm{resp},\chi_i}$ of Corollary~\ref{cor:c3-6-1} for $\chi_i$ (those of \eqref{eq:c3-2-3} as enlarged in Lemma~\ref{lem:exterior-response-bound}, before the replacement made in \eqref{eq:c3-7-14}) and $C_L$, $S_c$ of Proposition~\ref{prop:c3-7-5}; $C'_L$, $c_*$, $c_q$, $C_K$, $\tau$, $\pi_*$, $\varepsilon_1$, $b$ are those of Theorem~\ref{thm:main-theorem}, a constant that depends on the character being taken at its larger value for $\chi_1,\chi_2$ when it is a prefactor and at its smaller value when it is a decay rate. The constants $C$, $c$, $C_{\chi}$, $C^{\mathrm{src}}_{\chi_1\chi_2}$, $C_{\times}$, $c_{\times}$, $C_{\mathrm{resp},\chi_1\chi_2}$, $C'_L$, $c_*$, $c_q$ are fixed before the coupling; $C_K$ and the scales $\tau$, $\pi_*$, $\varepsilon_1$, $\rho$, $b$ depend on $g_0$ as in Theorem~\ref{thm:main-theorem}; none of them depends on $n_1$, $n_2$, the positions of the squares, the torus or the history.
\end{theorem}

\begin{proof}
Section~\ref{sec:two-loops-proof}.
\end{proof}

\begin{remark}[the cross terms]\label{rem:two-loop-regime}
No condition on the distance of the squares is imposed beyond $V(C_1)\cap V(C_2)=\varnothing$. The cross factor $e^{E^{\times}_{\xi}}$ and the factors $e^{C_{\times}\rho\Delta_{\times}}$ in \eqref{eq:thm-two-bounds} and in (iii) are close to one when $C_{\times}\rho\Delta_{\times}$ is small. Since $\Delta_{\times}\leqslant n_C$ and $\rho= O(g_0^2(1+r)^4)$, this holds at fixed $n_1$, $n_2$ for all positions of the squares once $g_0$ is small; at fixed coupling, $C_{\times}\rho\Delta_{\times}\leqslant \epsilon$ whenever $d/(LM)\geqslant c_{\times}^{-1}(\log(C_{\times}\rho n_C)+\log\epsilon^{-1})$, which the hypotheses allow on tori with correspondingly large periods. As $C_{\times}\rho\Delta_{\times}\to0$ the bound on $E^{\times}_{\xi}$ tends to zero, and those on $E_{\mathrm{def}}$ and $E_{\mathrm{bdry}}$ take the form of \eqref{eq:thm-main-bounds} with $n_C= n_{C_1}+n_{C_2}$ and the constants $C_{\chi}$, $c$, $\rho$, $C_{\mathrm{resp},\chi_1\chi_2}$ of Theorem~\ref{thm:two-loop-theorem}.
\end{remark}

\section{Extraction from the Effective Interaction}
\label{sec:extraction-procedure}

Ba{\l}aban first constructs the effective interaction as a function of the remaining field. He then extracts its relevant coefficients at the identity and subtracts the corresponding local terms from that field-dependent interaction. We use the convention in \cite[(1.6)--(1.7)]{BalabanRG1} and following \cite[(1.30)]{BalabanCRE}, in which the Gaussian normalisation is included in the effective interaction. The intermediate definition in RG1 (2.13) keeps that contribution separate.

In RG1's small-field construction, with no large-field regions, let $B$ be a small coarse Lie-algebra field, and let $U_1(e^{iB})$ be the constrained critical configuration with coarse field $e^{iB}$. Write $\mathbf E^{(1)}$ for the first-step effective interaction, including the Gaussian normalisation, and put
\begin{equation}
    \mathcal E_1(B)=\mathbf E^{(1)}(g_0,U_1(e^{iB}))\;.
\end{equation}
Choose an orthonormal basis of the Lie algebra for the action's invariant inner product. The index $a$ denotes a basis element, $\mu$ a bond direction and $x$ its initial coarse site. The second derivative is
\begin{equation}
    \Pi_{1,\mu\nu}^{ab}(x,y)
        =\left.
            \frac{\partial^2\mathcal E_1(B)}
                {\partial B_{\mu}^a(x)\,\partial B_{\nu}^b(y)}
          \right|_{B=0}\;.
\end{equation}
After the colour reduction and the translation-invariant infinite-lattice limit, let $\Pi_{1,\mu\nu}(x)$ denote the resulting scalar kernel with its second site fixed at the origin. In Ba{\l}aban's conventions the coupling counterterm is
\begin{equation}
    \beta_1(g_0)=\sum_{x\in\mathbb Z^4}x_1x_2\Pi_{1,12}(g_0,x)\;.
\end{equation}
The definition and the symmetry argument are given in \cite[(1.20)--(1.22), (5.37)--(5.43)]{BalabanRG1}.

In a proper regular region $\Lambda_1$, the subtraction has the form
\begin{equation}\label{eq:balaban-regular-subtraction}
    \mathbf E^{(1)}(\Lambda_1,g_0,U_1)
        -\mathbf E^{(1)}(\Lambda_1,g_0,\id)
        -\beta_1(g_0)S_{\phi_1}(U_1)\;.
\end{equation}
Here $\phi_1$ is the localisation cutoff in \cite[(1.27)]{BalabanCRE}, supported in $\Lambda_1$ and equal to one on its specified interior. If $x(p)$ is the initial site of an oriented representative of the plaquette $p$, then
\begin{equation}
    S_{\phi_1}(U)=
        \sum_p\phi_1(x(p))
            \bigl(1-\operatorname{Re}\tr_N U(\partial p)\bigr)\;.
\end{equation}
The boundary interaction and the exterior integrations remain in \cite[(1.29)--(1.30)]{BalabanCRE}. Ba{\l}aban states after (1.30) that the regular interaction is the RG1 interaction restricted to $\Lambda_1$. Within RG1's cube-localised construction, the changes of reconstruction and summation domain are controlled in \cite[(4.35)--(4.37)]{BalabanRG1}.

For the Wilson loop, the quantities to extract are the scalar density along a straight side and the scalar contribution of a corner. We denote them by $\delta m$ and $z_{\mathrm c}$. Their construction below uses connected terms of the normalised Wilson integral at the identity background. Section~\ref{sec:regional-identification} identifies the straight-side coefficient with the flat projection of the conditional Wilson kernel in the first RG step, and estimates the error caused by its integration boundary. Section~\ref{sec:remaining-physical-comparison} gives the outgoing insertion with the reconstructed closed flux factors and the remaining conditional source terms retained. Section~\ref{sec:sourced-effective-interaction} writes the conditional first-step integral with the insertion, at the actual background and exterior fields, as an exact sum of local closed insertions paired with signed local measures; the estimates it states as lemmas are proved in Appendix~\ref{app:background-estimates}.

\section{The Sourced Effective Interaction at the First Step}
\label{sec:sourced-effective-interaction}

This section carries out the first step of the procedure in Section~\ref{sec:extraction-procedure}: the conditional first-step integral with the Wilson insertion is written as an exact sum of local closed insertions paired with signed local measures, at the actual background and exterior fields. Every displayed identity is exact at fixed volume. What the section does not prove, and states instead as lemmas with proofs in Appendix~\ref{app:background-estimates}, is the existence of the objects in those identities and one cluster identity: the two-sided bounds on the young block $H_{\xi,II}$ of the regional Hessian (Lemma~\ref{lem:regular-bounds}), the analyticity of the reconstruction on the cutoff box (Lemma~\ref{lem:charts}), the finiteness of the selected integrals (Lemma~\ref{lem:gaussian-payment}) and the absolute convergence of the cluster sums (Lemmas~\ref{lem:activity-bounds} and~\ref{lem:vacuum-clusters}, the latter including the cluster identity \eqref{eq:part1-joint-vacuum-exponential}, whose proof is Appendix~\ref{app:background-estimates}, Section~\ref{app:kp}). Each identity below holds under those lemmas, under the standing conditions \eqref{eq:c-standing-conditions}, and under the branch hypotheses (B1)--(B4), which record what is taken from Ba{\l}aban's construction of the branch density. The other inputs are CRE's derivation of (1.11)--(1.25), as recorded after \eqref{eq:native-source-physical-ratio}, RG1 (0.4) and (0.10)--(0.11) for the averaging operation and the axial transports, Appendix~\ref{app:flat-covariance} for the flat values, and the $\operatorname{Ad}$-invariance \eqref{eq:marginal-cutoff-invariance} of the bond norm. The bounds on the terms of the final identity are Appendix~\ref{app:exterior-chessboard} and Section~\ref{sec:one-step-insertion}.

Conventions. $G$ is the compact connected simple gauge group fixed at the start of the paper, in its fixed unitary realisation, so that group elements near the identity are $U= e^{iA}$ with $A$ in the real vector space $\mathfrak g$ of Hermitian matrices $i^{-1}\times$(Lie algebra); $\langle X,Y\rangle=\tr_N(XY)$ is the fixed $\operatorname{Ad}$-invariant real inner product on $\mathfrak g$ of the flat calculation (the opening paragraph of Section~\ref{sec:flat}; it is the inner product for which every constant of Appendix~\ref{app:flat-covariance}, in particular $\gamma_L$, is derived), $h(X)=\langle X,X\rangle^{1/2}$, $d_G=\dim\mathfrak g$. Every quadratic form of this section is defined without reference to this choice; the operators representing the forms, the Hilbert norms, the Lebesgue measures $\dd a$, the $\delta$-functions of CRE (1.14), the Haar density $\sigma$ and the determinants all refer to this one inner product, and a constant said to depend on $L$ and the group only is a constant for it. CRE's identity (1.15) is an identity between the fine integral, which involves no inner product, and an expression whose measures $\dd A$, $\delta$-functions and densities $\sigma$, constants $\sigma_0$, $z^{(0)}$ are taken for one and the same inner product on $\mathfrak g$, whichever it is; it therefore holds for the fixed inner product with the function $\sigma$ and the constants $\sigma_0$, $z^{(0)}$ belonging to it, and these are the values meant below. $\log$ is the group chart: the inverse of $A\mapsto e^{iA}$ on the $\operatorname{Ad}$-invariant Hilbert ball $h(A)<\varrho_G\coloneqq \pi/\sqrt N$, whose range in $G$ is the chart domain. On this ball $\| A\|_{\mathrm{op}}\leqslant \sqrt N\,h(A)<\pi$ (for $h$ the normalised Hilbert--Schmidt norm, $h(A)^2= N^{-1}\sum_j\lambda_j^2\geqslant N^{-1}\| A\|_{\mathrm{op}}^2$), so the eigenvalues of $A$ lie in $(-\pi,\pi)$: $A\mapsto e^{iA}$ is injective there (an eigenvalue $\lambda$ of $A$ is recovered from the eigenvalue $e^{i\lambda}$ of $e^{iA}$ on the same eigenspace), its differential is invertible (the differential of $\exp$ at $iA$ is singular only when two eigenvalues of $iA$ differ by a nonzero multiple of $2\pi i$), and its inverse is the restriction to the chart domain of the principal logarithm $\frac1i\log$ used in the flat calculation (``the chart domain'' of $\mathcal O_c$, the sentence preceding \eqref{eq:nonlinear-average} and \eqref{eq:global-average}; ``the regular chart containing the identity'', the sentence following \eqref{eq:action-primitive}); every logarithm below is taken in this chart. The chart is Ba{\l}aban's: \cite[p.~21, (22)--(25)]{BalabanAVG} defines $\log U= iA$ for a unitary $U=\sum_je^{i\lambda_j}P_j$ (spectral projections $P_j$, distinct $\lambda_j\in\,]-\pi,\pi]$) by $A=\sum_j\lambda_jP_j$, Hermitian with $| A|\leqslant \pi$, with the bounds $| U-1|\leqslant |\log U|\leqslant \frac\pi2| U-1|$; the chart logarithm is the restriction of this principal logarithm to the chart domain, on which it is analytic, and the bound (25) holds for the normalised Hilbert--Schmidt norm $h$ as for the operator norm, since $| e^{i\lambda}-1|=2|\sin(\lambda/2)|\geqslant \frac2\pi|\lambda|$ for $\lambda\in\,]-\pi,\pi]$: $h(\frac1i\log W)\leqslant \frac\pi2h(W-\id)$, and a unitary $W\in G$ with $\| W-\id\|_{\mathrm{op}}\leqslant r_G$, in particular one with $h(W-\id)\leqslant r_G/\sqrt N$, lies in the chart domain, $r_G$ being the group constant of Appendix~\ref{app:background-estimates} (the eigenvalue bound alone does not place $W$ in the chart domain for a general compact $G\subset U(N)$: for $G= SU(N)$ with $N\geqslant 10$, the element $W= e^{2\pi i/N}\id$ has $h(W-\id)=2\sin(\pi/N)<2/\sqrt N$, while its principal logarithm $(2\pi/N)\id$ does not lie in $\mathfrak{su}(N)$). In particular ``$\mathsf M(U)(c)$ is defined'' means throughout that every logarithm in the averaging formula RG1 (0.4) for the row $c$ has its argument in the chart domain, that is $U\in\mathcal O_c$, where the averaging operation \eqref{eq:global-average} is RG1 (0.4) (it is defined for every $U$; off $\mathcal O_c$ it is the straight transport, a value that enters only through the constraint $\bar V_0= V$ on the rows of $\Omega_1^c$, Section~\ref{sec:branch-data}), and ``the local group means $\mathcal U_U(y,x)$ are defined'' that the transports along the path family from $y$ to $x$ lie in the small-diameter domain of the fixed group mean selected after RG1 (0.9), with the relative logarithms of (0.10) in the chart domain. We use that specified analytic mean, with the left-right covariance asserted in RG1 (0.6) and (0.10), and define $\mathcal U_U(y,x)$ by (0.11). No uniqueness among all solutions having logarithms in the chart is asserted. The implicit function theorem identifies the specified mean with its near-diagonal solution on a sufficiently small neighbourhood. The same convention defines the domain $D^{(11)}$ of Section~\ref{sec:branch-data}, the set $\mathcal X$ of the measurability paragraph and the objects of Section~\ref{sec:regular-hessian-stencils}. The chart being Ba{\l}aban's, the domain statement that on CRE's domain of integration in (1.11) the arguments of these logarithms and the fluctuation lie in the chart domain is CRE's smallness of the fluctuation (pp.247--248, by Lemma~1 of \cite{BalabanREG}) and of the contour variables (\cite[pp.~24--25]{BalabanAVG}), derived in (B1) for $\varepsilon_0$ small; what is this paper's in it is only the reading of the constants $O(L^2)$ and $O(1)$, which CRE and \cite{BalabanAVG} do not display, as independent of the coupling, recorded in (B1) and in the hypothesis (H4) of the main theorem. Appendix~\ref{app:background-estimates} uses besides $\varrho_G$ an operator-norm radius $r_G$ and the comparison constants $c_{\mathrm h}=\sqrt N$, $c_{\mathrm n}=1/N$ between $h$, the matrix norm $|\bigcdot |$ of the cutoff and $\|\bigcdot \|_{\mathrm{op}}$ ($h\leqslant |\bigcdot |\leqslant \|\bigcdot \|_{\mathrm{op}}\leqslant \sqrt N h$, where $h(Y)=(N^{-1}\tr Y^*Y)^{1/2}$ for an arbitrary matrix $Y$, so that $\| Y\|_{\mathrm{op}}\leqslant c_{\mathrm h}| Y|$ holds for every matrix, not only for Hermitian ones, and for either cutoff norm); they are recalled where used. The right inverse $h_U$ and the force $J_U$ of Section~\ref{sec:regular-hessian-stencils} always carry their background subscript, which distinguishes them from the norm $h$ and the carrier $J$. $| X|$ is CRE's bond norm of \eqref{eq:literal-cutoff}, \eqref{eq:marginal-two-cutoffs} (the operator or the normalised Hilbert--Schmidt norm, fixed once; both are $\operatorname{Ad}$-invariant by \eqref{eq:marginal-cutoff-invariance}). For $B\in\mathfrak g^S$ the Hilbert norm is $\| B\|^2=\sum_{b\in S}h(B_b)^2$ and the phase norm is $\| B\|_{\infty}=\max_bh(B_b)$. $E_*$ is the set of free bonds (all positive fine bonds except the designated bonds $b_0(c)$, RG1 p.267: $b_0(c)$ is the fine bond contained in the coarse bond $c$ whose endpoints lie in the two blocks of $c$, the sentence following \eqref{eq:constraint-primitive}); $\mathcal B_y$ is the block with centre $y$ and $b_{\pm}$ the endpoints of a bond $b$; distances of bonds are periodic $\ell^{\infty}$ distances of initial vertices. Bonds and rows of a set $X$ of sites follow CRE (1.8)--(1.9): the bonds of $X$ are the positive fine bonds with at least one endpoint in $X$, $X^{(1)}$ is the set of block centres in $X$, the rows of $X$ are the coarse bonds with at least one endpoint in $X^{(1)}$, written $X^{[1]}$, and $X^*$ is the set of bonds of $X$ other than the designated bonds $b_0(c)$, $c\in X^{[1]}$ (so a bond of $X^*$ may have one endpoint outside $X$, and the variables of a row of $X$ include its designated bond). The complement is treated differently, following this paper's reading of CRE (1.11)--(1.12), which say only that $V_0= U|_{\Omega_1^c}$ are the integration variables localised in $\Omega_1^c$ (p.248), read with the bonds crossing from $\Omega_1$ to $\Omega_1^c$ as variables of the integral over $\Omega_1$: the bonds of $\Omega_1^c$, written $b\in\Omega_1^c$, are the positive fine bonds with both endpoints off $\Omega_1$, the restriction $V_0= U|_{\Omega_1^c}$ is the restriction to these bonds, and the rows of $\Omega_1^c$ are the coarse bonds with both block centres off $\Omega_1$; so every positive fine bond is a bond of $\Omega_1$ or a bond of $\Omega_1^c$ and not both, and every coarse bond is a row of $\Omega_1$ or a row of $\Omega_1^c$ and not both. $X^{\sim}$ is CRE's enlargement of $X$ by all cubes touching $X$ (p.245; read as including the cubes meeting $X$ only at an edge or a vertex, this paper's reading, used for the distance statement at the end of this sentence and in Lemma~\ref{lem:collar-coverage}(i)) of the partition in terms of which $X$ is defined (p.245, where the partition into cubes $\square$ of side $LM_2R_0$ is introduced, $\square^{\sim}$ is $\square$ enlarged by one layer of $M_2R_0$-cubes, and ``usually'' $X^{\sim}$ enlarges by a layer of the cubes of the partition defining $X$; here $LMR_1$-cubes for $\Omega_1$, $R_1'$ and $\Lambda_1$, $L^2M_2R_1$-cubes for $\square'$, $Q_1$ and $R_1$, p.246), $X^{\sim-1}=((X^c)^{\sim})^c$ is the shrinking and $X^{\sim k}$ the $k$-fold iterate; a site of $X^{\sim-1}$ is at $\ell^{\infty}$ distance at least the cube side from $X^c$. CRE's exponent in the definition of the scales, $R_j$ the smallest power of $L$ with $R_j\geqslant (\log g_j^{-2})^{\mathfrak r}$, $\mathfrak r\geqslant 2$ and $p_0\geqslant 5\mathfrak r$ (p.245, p.246, CRE (2.5)), is written $\mathfrak r$ here, because $r$ is the cutoff radius of \eqref{eq:literal-cutoff}. CRE's scales satisfy $R_0\leqslant (L+1)R_1$ by CRE (2.9), p.256, hence $R_0\leqslant LR_1$ as both are powers of $L$; and $M_1$ is CRE's $M_1$ of p.245, the size of cubes for which the theorems of the previous papers, in particular those on the variational problem, are valid (read as a lower bound on the cube side; CRE takes $M_2= L^2M_1$ as an example, which is consistent with the conditions below). The identity of a matrix space is written $\id$; the action is $S$ (CRE's $A(U)$) and CRE's scalar normaliser $E$ is written $\mathsf E$: it is the vacuum scalar $E_{\varepsilon}$ of \eqref{eq:bare-action}, a constant independent of the configuration and of the branch, which cancels in every ratio of this paper and enters only the absolute identities for $\mathcal Z_0$ (Section~\ref{sec:c3-7}). The standing conditions on the parameters are
\begin{equation}\label{eq:c-standing-conditions}
    \begin{gathered}
    L>11\text{ odd (RG1 p.251)};\quad M_1< M_2< M\text{ odd powers of }L,\ M\geqslant 8LM_2,\ \text{\eqref{eq:fixed-kernel-collar-size}},\ \delta_M\leqslant 1,\ h_+^2\delta_M^2/\gamma\leqslant \tfrac18,\ 4C_{\mathrm d}\delta_M^2(LM+1)^4\leqslant \tfrac12;\\
    \text{every fine period a multiple of }LMR_1\text{ and at least }8M;\quad
    \text{every coarse period}>\max(2,2R_{\mathrm p}+8);\\
    g_0\text{ below the threshold of Appendix~\ref{app:background-estimates}, Section~\ref{app:thresholds}, which includes }LMR_1>4R_{\mathrm c}+4L^2M_2R_1+8\\
    \text{and }\widehat\mu= C_{L,R_{\mathrm c}}\varepsilon_1<\min\bigl(\tfrac12\tau_L,\ \gamma_L/(4\max(1,c_L))\bigr)\text{ (the hypotheses of Lemma~\ref{lem:regular-bounds})},
    \end{gathered}
\end{equation}
where $R_{\mathrm p}$ and $R_{\mathrm c}$ are the constants of that appendix (\eqref{eq:regular-covariance-radius-replay} and Lemma~\ref{lem:collar-coverage}), $\gamma$, $h_+$ those of \eqref{eq:fixed-kernel-inputs}, $\delta_M$, $C_{\mathrm d}$ those of Lemma~\ref{lem:gaussian-payment}, $\tau_L$, $c_L$ those of Lemma~\ref{lem:hessian-gauge-locality} and $\gamma_L$, $C_{L,R_{\mathrm c}}$ those of Lemma~\ref{lem:regular-bounds}, all depending on $L$ and the group only ($\delta_M$ also on $M$), and $\varepsilon_1= g_1p_0(g_1)$ CRE's quantity of p.246, which depends on the coupling; the threshold depends on $L$, the group, the character $\chi$ (through the collar constants $C'_{\chi}$, $K_{\chi}$ of item 4 of Section~\ref{app:thresholds}), $M_1$, $M_2$, $M$, on CRE's constants $A_0$, $A_1$, the exponents $p_0$ and $\mathfrak r$, on the constants $B_3$, $\delta_{\mathrm c}$ and the $O(1)$ of CRE's comparison after (1.19) recorded in (B3)(iii), and on the map $g_0\mapsto g_1$ of CRE (1.27), which are fixed as in CRE with $A_0/A_1$ large enough for (B3)(i); the order of choice of every constant and the list of conditions are those of Section~\ref{app:thresholds}, whose closing paragraph also lists the further conditions imposed where the statements about the square and the appendices state them (Sections~\ref{sec:c3-1}, \ref{sec:c3-7}, \ref{sec:c5-2}--\ref{sec:c5-7}). The threshold also takes over CRE's own conditions on the parameters (item 5 of Section~\ref{app:thresholds}). Those on the coupling, $\varepsilon_0$ small enough for Theorem 1 of \cite{BalabanVAR} (CRE's reference [15]; CRE's bracketed numbers are those of RG1's bibliography) (p.248), $MR_1$ large (p.251) and $M_2R_1$ large enough for (B3)(ii) (p.250), are conditions that $g_0$ lies below a value depending on $L$, the group, $M_1$, $M_2$, $M$ and CRE's constants, met for small $g_0$ at fixed $M_1,M_2,M$ since $R_1\to\infty$. Those on $M_1$, $M_2$, $M$ are not listed by CRE: $M_1$ is the cube size for which the theorems of the previous papers hold (p.245), $M_2$ is ``not too much larger'' than $M_1$, e.g. $M_2= L^2M_1$, and $M$ is ``much larger than $M_2$'', its choice ``dictated by many conditions involving this constant'' (p.245). They are taken over as the hypothesis that CRE's construction is valid for the values $M_1< M_2< M$ used here, which follow CRE's prescription (Section~\ref{app:thresholds}); this paper adds to them only the conditions displayed above. The scale integer $R_1$ is CRE's power of $L$ of p.246 (Section~\ref{sec:branch-data}; the same letter denotes CRE's large-fluctuation region, and in every product $LMR_1$, $L^2M_2R_1$ it is the scale), so the first-step cubes of sides $LMR_1$ and $L^2M_2R_1$, the $LM$-, $M$- and $M_1$-cubes of Section~\ref{sec:fixed-first-kernel} and the blocks form nested partitions of the fine torus (all sides are powers of $L$, and the $LM$-cubes nest in the $L^2M_2R_1$-cubes once $LM_2R_1\geqslant M$, threshold item 1 of Section~\ref{app:thresholds}, which holds for small $g_0$ since $R_1\to\infty$); the collar inequality holds once the scale satisfies $R_1\geqslant R_{\mathrm c}+3$, since $LMR_1\geqslant 8L^2M_2R_1$. The condition ``every coarse period $>2$'' is the operative hypothesis of Lemma~\ref{lem:two-block} (it makes $b_0(c)$ the only designated bond in $\mathcal E_c$: with a coarse period equal to two, two coarse bonds would join the same pair of blocks). Every statement of this section is made under \eqref{eq:c-standing-conditions}. When $I=\varnothing$ below, $\mathfrak g^I$ is a point, $\dd a$ the unit mass, every sum reduces to its empty term and every identity holds trivially.

Some letters carry more than one meaning, each fixed by its context, as follows. $R_*$ is the walk range $n_*R$ in Section~\ref{sec:fixed-first-kernel} and in the order of choice of Section~\ref{app:thresholds}; from Section~\ref{sec:marginal-perimeter} on it is the fixed range introduced there, bounding the label-support diameters and the source-to-stencil distances (fixed independently of the coupling), and it is this $R_*$ that enters $R_0=4(1+R_*+R_H+L)$. $R_0$ is CRE's scale where CRE's scales and partitions are quoted (the sentence on $R_0\leqslant (L+1)R_1$ above, p.245, and Lemma~\ref{lem:c3-7-seed}); from Section~\ref{sec:regional-identification} on it is this paper's range constant $R_0=4(1+R_*+R_H+L)$, as stated there. $R_1$ is the scale integer in every product $LMR_1$, $L^2M_2R_1$ and wherever it is called the scale, and CRE's large-fluctuation region where it is called a region. $M$ is the cube side; the group mean is the sans-serif $\mathsf M(U)(c)$; $M(W)$ is in Section~\ref{sec:remaining-physical-comparison} CRE's level-wise collection $M^j(W)$ of (2.11) at the first step, and from Section~\ref{sec:c5-5-0} on the positive group majorant of a group $W$, with $M_v(\mathfrak c)$ the tree majorant of a tuple (Section~\ref{sec:c4-6-2}). $\ell$ is the macro side $LMR_1$ in Section~\ref{sec:c5-3-1} and in Appendix~\ref{app:exterior-chessboard}, where it is so declared, and a selection label under summation signs ($\sum_{\ell}$, Sections~\ref{sec:probe-kernels} and~\ref{sec:c5-5-0}); $\ell_i(Y)= d_i(Y)+5$ are the shifted tree lengths of a domain (defined with \eqref{eq:c2-tree-length}). $m_L$ is the Hessian constant of Lemma~\ref{lem:hessian-gauge-locality}; $m$ alone is the tuple length in $\mathfrak c=(D_1,\ldots,D_m)$ (Section~\ref{sec:selected-integrals}) and the segment length in coarse units in \eqref{eq:regional-mass-identification}; $m(D)=\varepsilon_*e^{-\lambda_*\ell_1(D)}$ is the majorant of Lemma~\ref{lem:activity-bounds}; $\delta m^{\sigma}$ is the mass coefficient \eqref{eq:marginal-square-mass}. $\theta\geqslant 0$ is the Mayer constant of Section~\ref{app:mayer-constants} (Section~\ref{sec:selected-integrals} on); in Appendix~\ref{app:exterior-chessboard}, from \eqref{eq:c3-7-4} on, $\theta$ is the witness threshold defined there, $\theta_{\vartheta}$ its $\vartheta$-dependent form (Section~\ref{sec:c5-2}), and $\theta_{\mathrm h}$, $\theta_{\mathrm x}$ are the constants of Section~\ref{sec:c4-8-2}; $\vartheta$ and $\vartheta_*$ are different letters throughout. $C_L$ is the coercivity constant \eqref{eq:covariance-coercivity-constant} of Appendix~\ref{app:flat-covariance} in Lemma~\ref{lem:regular-bounds} ($\gamma_L= C_L^{-1}$) and that appendix; it is the contour constant $|\Gamma_C|\leqslant C_Ln_C$ in Section~\ref{sec:c3-7}, at the opening of Appendix~\ref{app:signed-defects}, in Section~\ref{sec:c4-2} (with $C'_L=8C_L$), in Corollary~\ref{cor:main-leading-size} and in Remark~\ref{rem:next-step}; elsewhere it is a constant depending on $L$ (in Section~\ref{sec:c3-9-1} also on $N$) as stated where it appears. $K=256$ is the fixed integer of $\sigma_K$, $C_K$ (Lemma~\ref{lem:c5-7-joint}); in (H3) $K$ is RG1's exponent in $\varepsilon= L^{-K}$; the number of steps in Section~\ref{sec:next-step} is $k+1$. The italic $B_{\xi}=\Lambda_1^c$ is the filled bad region (with $B^{\mathrm{fin}}_{\xi}=(\Lambda_1^{\mathrm{fin}})^c$), the calligraphic $\mathcal B_{\xi}$, $\mathcal B^{\Lambda}_{\xi}$ the exterior response functional, and the fraktur $\mathfrak B$ a compatible family of the expansion; $B_3$ is CRE's constant; the fraktur $\mathfrak B_0=(\Lambda_1^{(0)})^c$ is a set of sites and $\mathfrak B_Q$ a determining set (Section~\ref{sec:c3-9-1}), the calligraphic $\mathcal B_y$ the block with centre $y$, and $B$ without the subscript $\xi$ a generic vector of $\mathfrak g^S$, the ball $B_{\delta}$, the bound $B_C$ or the host constants $B_h$, $B_{\mathrm h}$; in Section~\ref{sec:branch-data} the letter $F$ denotes integrands, and the logarithms $F_{I,C,T}$, $F_{C,T}$, $F_{C,\infty}(n)$ and the integrals $F_{\xi}(t)$ are named with their subscripts; elsewhere $F_{\gamma}$ is a set of edges of a polymer and $F(\Gamma)$ a weight in a cluster sum.

\subsection{Branch data, the outer operation and the conditional integral}
\label{sec:branch-data}

Fix one term of the history sum of \cite[(1.11)--(1.30)]{BalabanCRE}: the regions $\Omega_1,R_1$ (CRE uses $R_1$ both for the large-fluctuation region of (1.18) and for the scale integer of p.246, the smallest power of $L$ with $R_1\geqslant (\log g_1^{-2})^{\mathfrak r}$, which grows as $g_0\to0$; both uses are kept, the scale appearing in the products $LMR_1$, $L^2M_2R_1$, $MR_1$, $M_2R_1$ and in its defining inequality $R_1\geqslant (\log g_1^{-2})^{\mathfrak r}$; here $g_1= g_1(g_0)$ is CRE's first-step coupling, the constant value on $\Lambda_1^{\sim-1}$ of the solution of (1.27), $g_0^{-2}= g_1^{-2}+\beta_1(g_0)$ with CRE's $\beta_1$, a function of $g_0$ alone, and $\varepsilon_1= g_1p_0(g_1)$ as on p.246), the derived $\Lambda_1,S_1$ of (1.20) with $\Lambda_1$ the filled domain of the modification following (1.20), the new field $V$, the integration variable $V_0$ on $\Omega_1^c$, and the background $U_1= U_1(V,V_0)$ of (1.12) in its axial mean gauge, equal to $V_0$ on $\Omega_1^c$. The young carrier is $I=\Lambda_1^*$, the retained free carrier is $J=\Omega_1^*\setminus I$, and $A_{\mathrm e}\in\mathfrak g^J$ is the retained fluctuation. The young variable is $a\in\mathfrak g^I$; both are extended by zero to $\Omega_1^*$ where needed. The branch data are $\xi=(\Omega_1,R_1,V,V_0,A_{\mathrm e})$. The constraint $\delta(\bar UV^{-1})$ of CRE (1.11) is a product over the coarse bonds, and it is split by rows: on the rows of $\Omega_1$ (coarse bonds with a block centre in $\Omega_1$) the constraint $\bar U= V$ is the critical-point constraint of (1.12), imposed on the integral over the bonds of $\Omega_1$ (which include the bonds crossing from $\Omega_1$ to $\Omega_1^c$; the rows with one block centre off $\Omega_1$ average also over bonds of $\Omega_1^c$, which enter as the constant $V_0$); read, like every group $\delta$-function here, by integrating against the coarse variables it constrains, it is the fibre measure of the pushforward of the fine measure under the map $\Phi(U)=(\bar U|_{\Omega_1^{[1]}},U|_{\Omega_1^c})$, a submersion on the open set $\mathcal O^{\Phi}$ of fine configurations at which the averages and means on the rows of $\Omega_1$ are defined and $D\Phi$ is onto, which contains the set $\mathcal X$ of the measurability paragraph below and, on every branch, $D(V,V_0)$ (shown there) (the fibre measures of a $C^1$ submersion $\mathsf f$ of an open set $\mathcal Y$ of a manifold with a smooth positive density, here $G^E$ with Haar measure or $\mathfrak g^E$ with Lebesgue measure, into such a manifold of dimension $k$, are defined at \emph{every} point $y$ of $\mathsf f(\mathcal Y)$, not only almost everywhere: by the submersion theorem $\mathcal Y$ is covered by open sets $\mathcal Y_j$ with $C^1$ diffeomorphisms $\psi_j$ from open sets of $\mathbb R^k\times\mathbb R^{n-k}$ onto $\mathcal Y_j$ and $\mathsf f(\psi_j(y,z))= y$; with a continuous partition of unity $(\theta_j)$ subordinate to $(\mathcal Y_j)$ put $\mu^{\mathsf f}_y(g)\coloneqq \sum_j\int(\theta_jg)(\psi_j(y,z))\,|\det D\psi_j(y,z)|\,\dd z$ for Borel $g\geqslant 0$ on $\mathcal Y$, the determinants taken with respect to the densities, that of the manifold on $\mathcal Y_j$ and, on $\mathbb R^k\times\mathbb R^{n-k}$, the product of the target density in the variable $y$ and Lebesgue measure $\dd z$, so that $\dd y$ in (i) is the target density. Then (i) $\int_{\mathcal Y}\varphi(\mathsf f(x))g(x)\,\dd x=\int\varphi(y)\,\mu^{\mathsf f}_y(g)\,\dd y$ for Borel $\varphi\geqslant 0$ (the change of variables formula and Fubini's theorem); (ii) $y\mapsto\mu^{\mathsf f}_y(g)$ is continuous for every continuous $g$ with compact support in $\mathcal Y$ (dominated convergence); (iii) the family does not depend on the choices made, two families with (i) and (ii) agreeing for almost every $y$ by (i) and hence for every $y$ by (ii); (iv) for every $y\in\mathsf f(\mathcal Y)$ and every such $g$, $\mu^{\mathsf f}_y(g)=\lim_{\delta\to0}| E_{\delta}|^{-1}\int_{\mathsf f^{-1}(E_{\delta})}g\,\dd x$ for every family of Borel sets $E_{\delta}$ with $B_{c\delta}(y)\subseteq E_{\delta}\subseteq B_{C\delta}(y)$ for constants $0< c\leqslant C$, by (i) with $\varphi=\mathbf 1_{E_{\delta}}$ and (ii) (the mean of a continuous function over $E_{\delta}$ tends to its value at $y$); so $\mu^{\mathsf f}_y$ is carried by $\mathsf f^{-1}(y)$, and the fibre measures of the restriction of $\mathsf f$ to an open subset of $\mathcal Y$ are the restrictions of the $\mu^{\mathsf f}_y$. From (iv) follows the transformation rule, at every level: if $\Xi$ is a $C^1$ diffeomorphism of an open set $\mathcal N_0$ onto an open subset of $\mathcal Y$, $\psi$ a $C^1$ diffeomorphism of an open neighbourhood of $\mathsf f(\Xi(\mathcal N_0))$ onto its image, and $\mathsf f'\coloneqq \psi\circ\mathsf f\circ\Xi$ on $\mathcal N_0$ (a submersion), then $\mu^{\mathsf f'}_{\psi(y)}(g\circ\Xi\,|\det D\Xi|)=|\det D\psi(y)|^{-1}\mu^{\mathsf f}_y(g)$ for every $y\in\mathsf f(\Xi(\mathcal N_0))$ and every Borel $g\geqslant 0$ on $\Xi(\mathcal N_0)$: $\int_{\mathcal N_0}\mathbf 1_{\mathsf f'(x')\in B_{\delta}(\psi(y))}\,g(\Xi(x'))\,|\det D\Xi(x')|\,\dd x'=\int_{\Xi(\mathcal N_0)}\mathbf 1_{\mathsf f(x)\in\psi^{-1}(B_{\delta}(\psi(y)))}\,g(x)\,\dd x$ exactly, for every $\delta$ (the change of variables formula for $\Xi$); the sets $E_{\delta}=\psi^{-1}(B_{\delta}(\psi(y)))$ satisfy the condition of (iv) with $| E_{\delta}|/| B_{\delta}|\to|\det D\psi(y)|^{-1}$; dividing by $| B_{\delta}|$ and letting $\delta\to0$ gives the identity for continuous $g$ of compact support, and two Radon measures on the fibre, a $C^1$ submanifold, agreeing on these agree. For $\Phi$ on $\mathcal O^{\Phi}$: for Borel $F\geqslant 0$ supported in $\mathcal O^{\Phi}$ and Borel $\varphi\geqslant 0$ (hence for bounded Borel $\varphi$ when $F$ is integrable), $\int\varphi(\Phi(U))F(U)\,\dd U=\int\varphi(V,V_0)\int_{\Phi^{-1}(V,V_0)}F\,\dd\mu_{V,V_0}\,\dd V|_{\Omega_1^{[1]}}\dd V_0$ with $\mu_{V,V_0}\coloneqq \mu^{\Phi}_{(V,V_0)}$, defined for every $(V,V_0)\in\Phi(\mathcal O^{\Phi})$, in particular for every admissible pair, and since $\Phi$ preserves the coordinate $U|_{\Omega_1^c}$ its straightening maps may be taken to preserve it, so that for every $(V,V_0)$ the measure $\mu^{\Phi}_{(V,V_0)}$ is the fibre measure over $V$ of the submersion $U|_{\Omega_1}\mapsto\bar U|_{\Omega_1^{[1]}}$ on the slice $U|_{\Omega_1^c}= V_0$, the Jacobians agreeing because the differentials of the straightening maps are block triangular with the identity in the $V_0$ block), and CRE's chart, linearisation and elimination steps listed after \eqref{eq:native-source-physical-ratio} compute $\mu_{V,V_0}$ in the variables of (1.14)--(1.15), in which the constraint is linear, by the transformation rule applied at the level of the pair itself and two facts. First, if $\Theta$ is a $C^1$ diffeomorphism of an open set $\mathcal N_{\Theta}$ onto its image and $\mathcal Q\circ\Theta=\Lambda$ on $\mathcal N_{\Theta}$ with $\Lambda$ linear and onto, then $\mathcal Q$ is a submersion on $\Theta(\mathcal N_{\Theta})$ and its fibre measures there are $\mu^{\mathcal Q}_y=\Theta_*(|\det D\Theta|\,\mu^{\Lambda}_y)$ for every $y\in\Lambda(\mathcal N_{\Theta})$, the transformation rule with $\Xi=\Theta$ and $\psi=\id$ (the change of variables formula $\int_{\Theta(\mathcal N_{\Theta})}\varphi(\mathcal Q(x))F(x)\,\dd x=\int_{\mathcal N_{\Theta}}\varphi(\Lambda(x'))F(\Theta(x'))\,|\det D\Theta(x')|\,\dd x'$ alone gives the identity for almost every $y$ with an exceptional set depending on $\mathcal Q$, here on $U_1(V,V_0)$, which would not cover the level $0$ needed at every pair; the rule at every level does). Second, the fibre measure over every $y$, in particular over $0$, of a linear surjection $\Lambda$ whose restriction $\Lambda_{\mathrm d}$ to the designated variables is invertible is Lebesgue measure on the free variables $x_{\mathrm f}$ times $|\det\Lambda_{\mathrm d}|^{-1}$, since $(y,x_{\mathrm f})\mapsto(x_{\mathrm f},\Lambda_{\mathrm d}^{-1}(y-\Lambda_{\mathrm f}x_{\mathrm f}))$ is a single straightening map of $\Lambda$ with Jacobian $|\det\Lambda_{\mathrm d}|^{-1}$ (the sifting of $\delta(\Lambda x)$ over the designated variables); here that restriction is $\bigoplus_c\alpha_c(U_1)$ up to the normalisation of the $\delta$-functions, which gives $z^{(0)}(U_1)$. The transformation rule is applied first on the slice $U|_{\Omega_1^c}= V_0$ with $\Xi$ the composition of the chart $A'\mapsto e^{iA'}$ and the right translation by $U_1$ on the bonds of $\Omega_1$ ($|\det D\Xi|=\prod_b\sigma(A'(b))$, the Haar density, with respect to Lebesgue and Haar measure) and $\psi$ the row chart $W\mapsto\frac1i\log(WV^{-1})$ on a neighbourhood of $V$ in $G^{\Omega_1^{[1]}}$ ($|\det D\psi(V)|=\sigma_0^{-|\Omega_1^{[1]}|}$ with respect to Haar and Lebesgue measure), under which the slice submersion becomes $\mathcal Q_{U_1}$ of \eqref{eq:native-source-constraint-normalisation} and the level $V$ becomes the level $0$: $\mu_{V,V_0}=\sigma_0^{-|\Omega_1^{[1]}|}\Xi_*(\prod_b\sigma(A'(b))\,\mu^{\mathcal Q_{U_1}}_0)$ on $\Xi$ of the chart domain (the rule is applied on $\mathcal N_0\coloneqq \Xi^{-1}(\mathsf f^{-1}(\operatorname{dom}\psi))$, $\mathsf f$ the slice submersion, the part of the chart domain whose row averages lie in the domain of the row chart; every point of the fibre over $V$ in $\Xi$ of the chart domain lies in $\Xi(\mathcal N_0)$, since there $\mathsf f= V\in\operatorname{dom}\psi$, and the fibre measure of a restriction is the restriction of the fibre measure, so nothing is lost), which is the conversion of the row $\delta$-functions in the step list, at the level of the pair itself. The two facts are then applied to $\mathcal Q=\mathcal Q_{U_1}$ at fixed $(V,V_0)$, at the level $0$, in the chart variable $A'$ of (1.13), with $\Theta(A)= A-h_{U_1}\widetilde D(A)$ on the open set $\mathcal N_{\Theta}$ of Lemma~\ref{lem:c2-fibre-smallness}(iv), on which $\widetilde D=\varphi$ is analytic, solves the row equations and has $\det(\id-h_{U_1}\frac{\delta}{\delta A}\widetilde D)>0$, on which $\Theta$ takes chart-valued phases (so that $\Theta(\mathcal N_{\Theta})\subseteq\mathcal N_0$, the set on which the transformation rule was applied above: the phases of $\Theta(A)$ are chart-valued, and at every row $c$ the row phase $\Theta(A)|_{\mathcal E_c}$ has norm below $\vartheta_*<2\varrho_*$, where $\mathcal N_{U_1,c}$, and hence $\mathcal Q_{U_1,c}$ of \eqref{eq:native-source-constraint-normalisation}, is defined and analytic (Lemma~\ref{lem:c2-uniform-chart}(i), which applies at every row by Lemma~\ref{lem:c2-fibre-smallness}(i)), so that the row averages lie in the domain of the row chart), and which contains the argument of every $U\in D^{(11)}(V,V_0)$ (Lemma~\ref{lem:c2-fibre-smallness}(ii)) and does not depend on the region $R_1$ (this step precedes the insertion of (1.18), which introduces $R_1$): $\mathcal Q_{U_1}\circ\Theta= LQ_{U_1}=\Lambda$ on $\mathcal N_{\Theta}$ (the row equations and $LQ_{U_1}h_{U_1}=\id$), and $\Theta$ is injective on $\mathcal N_{\Theta}$, since $\Theta(A)=\Theta(A'')= q$ gives $\widetilde D(A)=\mathcal N_{U_1}(q)=\widetilde D(A'')$ by the row equations $d_c=\mathcal N_{U_1,c}(A-h_{U_1}d)$, hence $A= q+h_{U_1}\widetilde D(A)= A''$; so $\Theta$ is a diffeomorphism of $\mathcal N_{\Theta}$ onto the open set $\Theta(\mathcal N_{\Theta})$ by the inverse function theorem; on the rows of $\Omega_1^c$ (both block centres off $\Omega_1$) it is the constraint $\bar V_0= V$ on $V_0= U|_{\Omega_1^c}$, written $\delta(\bar V_0V^{-1})$ below, and the constrained integral $\int\dd V_0|_{\Omega_1^c}\,\delta(\bar V_0V^{-1})\cdots$ in the form in which CRE writes it in (1.14)--(1.15) (in (1.11) it is part of $\int\dd U\,\delta(\bar UV^{-1})$) is taken as it stands there and read, like every group $\delta$-function of a coarse variable in this paper (the paragraph following Lemma~\ref{lem:global-average}), by integrating it against the coarse variables it constrains: for bounded Borel $h$ on $G^{E_1}$ and nonnegative Borel $F$, $\int\dd V\,h(V)\int\dd V_0\,\delta(\bar V_0V^{-1})F(V,V_0)=\int\dd V|_{\Omega_1^{[1]}}\int\dd V_0\,h(V)F(V,V_0)$ with $V(c)= M(V_0,c)$ of \eqref{eq:global-average} on the rows $c$ of $\Omega_1^c$ (a function of the bonds of $\Omega_1^c$ alone, Lemma~\ref{lem:global-average}(i)); so $\int\mathcal O[F](V)\,\dd V$ below is the integral of $F$ against the measure in $V$ that the outer operation defines, and $\mathcal O[F](V)$ is a Haar density only where that measure has one. Let $D^{(11)}= D^{(11)}(\Omega_1,V)$ be CRE's domain of integration in (1.11): the fine configurations $U$ with $\chi_1(\Omega_1)\zeta(\Omega_1^c)\chi_0'(\Omega_1)\ne0$ on which the block averages $\mathsf M(U)$ are defined on every row of $\Omega_1$ and the local group means $\mathcal U_U(y,x)$, $x\in\mathcal B_y$, $y\in\Omega_1^{(1)}$, of the gauge-fixing functional $\mathbf G(\Omega_1^{(1)},U)$ of (1.11) are defined (in the chart of the preamble, Ba{\l}aban's principal logarithm; by CRE p.247 and p.248 with Lemma~1 of \cite{BalabanREG} and \cite[pp.~24--25]{BalabanAVG}, derived in (B1), these two conditions are implied by the nonvanishing of the cutoffs and the constraint for $\varepsilon_0$ small, so that $D^{(11)}$ is CRE's domain of integration in (1.11), taken as it stands; they are kept explicit so that $D^{(11)}\subseteq\mathcal O_c$ on every row of $\Omega_1$ is visible), subject to the constraint $\bar U= V$ on the rows of $\Omega_1$ and $\bar V_0= V$ on the rows of $\Omega_1^c$. Its fibre over the restriction is $D^{(11)}(V,V_0)=\{ U\in D^{(11)}:U|_{\Omega_1^c}= V_0\}$. A pair $(V,V_0)$, written $(\Omega_1,V,V_0)$ where the region is to be displayed and called a pair throughout ($\Omega_1$ is fixed on the branch), is \emph{admissible} when $D^{(11)}(V,V_0)\neq\varnothing$; (B1) provides the background $U_1\in D^{(11)}(V,V_0)$ for every admissible pair. The domain of integration of (1.15) with the cutoff (1.25) is in general smaller than $D^{(11)}(V,V_0)$ (the small- and large-fluctuation tests of (1.18) in (1.25) are further conditions; on $R_1$ the large-fluctuation tests exclude the configurations near $U_1$); it is described in the variables of (1.15) as follows. For an admissible pair, a free vector $B\in\mathfrak g^{\Omega_1^*}$ and a family $d=(d_c)_{c\in\Omega_1^{[1]}}$ of solutions of the row equations \eqref{eq:native-source-local-implicit} at the argument $A= g_0\mathsf C_{U_1}B$ (that is, $d_c=\mathcal N_{U_1,c}(A-h_{U_1}d)$ for every row $c$), put $q_d(B)\coloneqq A-h_{U_1}d\in\mathfrak g^{\Omega_1}$, the phase \eqref{eq:native-source-local-phase} formed with this family (CRE's (1.13)--(1.14) at the background $U_1$; on the free bonds $q_d(B)(b)= g_0B(b)$), and let $U_d(B)$ be the fine field equal to $e^{iq_d(B)}U_1$ on the bonds of $\Omega_1$ and to $V_0$ on the bonds of $\Omega_1^c$. The family $d$ is \emph{of chart class} when $h(q_d(B)(b))<\varrho_G$ for every bond $b$ of $\Omega_1$, that is, when the whole phase $q_d(B)$ is chart-valued, the class of CRE's variable $A'$ of (1.13). Let
\begin{equation}\label{eq:c-free-variable-domain}
    \mathcal A_{\xi}\coloneqq \{ B\in\mathfrak g^{\Omega_1^*}:\ \text{there is a family }d\text{ of chart class at }g_0\mathsf C_{U_1}B\text{ with }U_d(B)\in D^{(11)}(V,V_0)\text{ and the cutoff (1.25) nonzero at }U_d(B)\}\;,
\end{equation}
the cutoff (1.25) being a function of the fine field and the branch data (its tests run over the free bonds $(\square'^{\sim2})^*$, at which the scaled variable is $g_0^{-1}\frac1i\log(UU_1^{-1})$, by (1.18)--(1.19), p.250, and (1.25), p.252, in the variables of (1.15), and through $\chi_0'(S_1)$ over the fine field; no test involves a designated value of the linearised variable, so the cutoff does not involve $\widetilde D$). For $B\in\mathcal A_{\xi}$ the family $d$ of chart class with $U_d(B)\in D^{(11)}(V,V_0)$ is unique and equals $\varphi(g_0\mathsf C_{U_1}B)$ (Lemma~\ref{lem:c2-fibre-smallness}(iii)), and $\widetilde D_{U_1}(g_0\mathsf C_{U_1}B)\coloneqq d$, $q_{U_1}(B)\coloneqq q_d(B)$ and $U(B)\coloneqq U_d(B)$ are defined on $\mathcal A_{\xi}$ by this uniqueness; the cutoff (1.25) is then a function of $B$, $U_1$ and the branch data. Since $q_{U_1}(B)$ is chart-valued, it is the chart logarithm of the fluctuation $U(B)U_1^{-1}$, so that $B\mapsto U(B)$ is injective on $\mathcal A_{\xi}$ ($U(B)$ determines $U(B)U_1^{-1}$, hence its chart logarithm $q_{U_1}(B)$, hence $B$ on the free bonds, where $q_{U_1}(B)(b)= g_0B(b)$). Conversely, every $U\in D^{(11)}(V,V_0)$ at which the cutoff (1.25) is nonzero is of the form $U(B)$ with $B\in\mathcal A_{\xi}$: by (B1) the fluctuation $UU_1^{-1}$ lies in the chart domain bond by bond on $\Omega_1$, so $A'\coloneqq \frac1i\log(UU_1^{-1})$ of (1.13) is defined and chart-valued; the constraint $\mathsf M(U)(c)= V(c)=\mathsf M(U_1)(c)$ on the rows of $\Omega_1$ gives $\mathcal Q_{U_1,c}(A')=0$ in \eqref{eq:native-source-constraint-normalisation}, that is $LQ_{U_1}A'=-\mathcal N_{U_1}(A')$; with $d\coloneqq \mathcal N_{U_1}(A')$ the vector $A'+h_{U_1}d$ lies in the kernel of $Q_{U_1}$ ($LQ_{U_1}h_{U_1}=\id$) and agrees with $A'$ on the free bonds ($h_{U_1}$ has its outputs at the designated bonds), so it equals $g_0\mathsf C_{U_1}B$ with $B\coloneqq g_0^{-1}A'|_{\Omega_1^*}$, the kernel of $Q_{U_1}$ being the graph of $\mathsf C_{U_1}$ over the free bonds by \eqref{eq:regular-stencil-elimination} ($\alpha_c(U_1)$ invertible, (B2) and Lemma~\ref{lem:two-block}); hence $A'= g_0\mathsf C_{U_1}B-h_{U_1}d$, $d_c=\mathcal N_{U_1,c}(A')=\mathcal N_{U_1,c}(A-h_{U_1}d)$ solves every row equation, $d$ is of chart class because $A'$ is chart-valued, and $U_d(B)= e^{iA'}U_1= U$ on $\Omega_1$ while $U|_{\Omega_1^c}= V_0$; so $B\in\mathcal A_{\xi}$ and $U(B)= U$. Thus $\mathcal A_{\xi}$ is the domain of integration of (1.15) with the cutoff (1.25) in the free variable $B$ of (1.15), and $B\mapsto U(B)$ is a bijection of $\mathcal A_{\xi}$ onto $D(V,V_0)\coloneqq \{ U\in D^{(11)}(V,V_0):\text{the cutoff (1.25) is nonzero at }U\}= U(\mathcal A_{\xi})$, the same domain in the fine variable; put $D(\Omega_1,R_1,V)\coloneqq \bigcup_{V_0}D(V,V_0)$. By (B4)(ii) the differential of $B\mapsto U(B)$ is injective at every point of $\mathcal A_{\xi}$, with image the tangent space of the fibre $\Phi^{-1}(V,V_0)$ at $U(B)$. For an admissible pair and $A_{\mathrm e}\in\mathfrak g^J$, $a\in\mathfrak g^I$ let $U_{\xi}(a)= U(\Psi_{\xi}(a+A_{\mathrm e}))$ be the reconstructed fine field \eqref{eq:c-reconstructed-field} (the inverse of CRE's changes of variables (1.13), (1.14), (1.22) at the background $U_1$, defined where those are: where $(a+A_{\mathrm e})(b)\in\mathcal W_b$ for every $b\in S_1^*$, Section~\ref{sec:boundary-law-collar}, and every row solution is defined at $g_0\mathsf C_{U_1}\Psi_{\xi}(a+A_{\mathrm e})$), and let
\begin{equation}\label{eq:c-new-variable-domain}
    \mathcal D_{\xi}\coloneqq \{(A_{\mathrm e},a)\in\mathfrak g^J\times\mathfrak g^I:\ \Psi_{\xi}(a+A_{\mathrm e})\text{ is defined and }\Psi_{\xi}(a+A_{\mathrm e})\in\mathcal A_{\xi}\}
\end{equation}
be the domain of integration of (1.15) with the cutoff (1.25) in the variables of (1.22); by (B4)(ii) (Lemma~\ref{lem:boundary-bijection}) $\Psi_{\xi}$ maps $\mathcal D_{\xi}$ bijectively onto $\mathcal A_{\xi}$, so the reconstruction $(A_{\mathrm e},a)\mapsto U_{\xi}(a)$ is a bijection of $\mathcal D_{\xi}$ onto $D(V,V_0)$. In particular $\mathcal D_{\xi}\neq\varnothing$ exactly when $\mathcal A_{\xi}\neq\varnothing$, that is when $D(V,V_0)\neq\varnothing$: no admissible pair with nonempty domain of integration is lost in passing to the variables of (1.22). A \emph{branch} is a value of the data $\xi$ such that $(\Omega_1,V,V_0)$ is admissible and $(A_{\mathrm e},a)\in\mathcal D_{\xi}$ for some $a$. On a branch the set $\{ a:(A_{\mathrm e},a)\in\mathcal D_{\xi}\}$ is the open box $\{\Theta_{\xi}=1\}$ of \eqref{eq:native-source-young-cutoff} (Lemma~\ref{lem:collar-coverage}(iv)), and the large-field factor $\zeta_{\xi}(A_{\mathrm e})$ of \eqref{eq:c-large-field-factor} and the retained cutoff $\chi_{\xi,\mathrm e}(A_{\mathrm e})$ of \eqref{eq:retained-cutoff-composition} are defined and nonzero. The branch indicator $\iota_{\xi}(A_{\mathrm e})$ is one when $\xi$ is a branch and zero otherwise; as the fibre over a branch value is the open box, which contains $a=0$, $\iota_{\xi}(A_{\mathrm e})=1$ exactly when $(A_{\mathrm e},0)\in\mathcal D_{\xi}$, a Borel condition (a section of the Borel set $\mathcal D_{\xi}$, measurability paragraph). It is a factor of the integrand of \eqref{eq:native-source-explicit-outer-operation}, which therefore vanishes off the branches, and every object below is defined on branches only. The sets $\mathcal A_{\xi}$, $D(V,V_0)$ and $\mathcal D_{\xi}$ depend on $\xi$ only through $(\Omega_1,R_1,V,V_0)$: the cutoff (1.25) in the variable $B$ of (1.15) and the map $\Psi_{\xi}$ do not involve $A_{\mathrm e}$, which is a variable of these sets.

For a branch function $\mathfrak M_{\xi}$, the outer operation is
\begin{equation}\label{eq:native-source-explicit-outer-operation}
    \mathcal O[\mathfrak M](V)
        =\sum_{\Omega_1,R_1}\chi_1(\Omega_1)
            \int\delta(\bar V_0V^{-1})
            \left[\int\iota_{\xi}(A_{\mathrm e})\zeta_{\xi}(A_{\mathrm e})\chi_{\xi,\mathrm e}(A_{\mathrm e})\kappa_{\xi}\mathfrak M_{\xi}\,\dd A_{\mathrm e}|_J\right]
            \dd V_0|_{\Omega_1^c}\;,
\end{equation}
with the prefactor
\begin{equation}\label{eq:native-source-outer-prefactor}
    \kappa_{\xi}
        = z^{(0)}(U_1)\,g_0^{d_G|\Omega_1^*|}\sigma_0^{|\Omega_1^*|}
            z^{-(L^4-1)|\Omega_1^{(1)}|}
        \exp\left(-g_0^{-2}S(U_1)-\mathsf E\right)\;.
\end{equation}
Here $\chi_1(\Omega_1)$ is the background test of CRE (2.16)--(2.17) at $k=1$, a function of $V$; $\delta(\bar V_0V^{-1})$ is the part of the constraint $\delta(\bar UV^{-1})$ of (1.11) on the rows of $\Omega_1^c$ (both blocks off $\Omega_1$), as in CRE (1.14)--(1.15), where it stands in front of the $V_0$ integral. The large-field factor $\zeta(\Omega_1^c)$ of (1.11) is defined by the characteristic functions and gauge-fixing exponentials of (1.6) and (1.8) whose cubes or blocks lie in $\Omega_1^c$, resummed over the admissible $P_0$, $P_1$, $Q_1$ (p.248). It is a sum over the admissible $(P_0,P_1,Q_1)$ compatible with $\Omega_1$ of products of nonnegative functions of the fine field $U$ on the bonds near their cubes or blocks and of constants; the factors of (1.4), (1.6) and (1.8) not displayed in (1.11), and the bonds involved, are the following. The large-field factors $\chi_0^c(\square)$ of (1.1) for the $LM_2R_0$-cubes $\square\subseteq P_0$, and the factors $\chi_0(\square)$, $\square\subseteq P_0^c$, that survive the removal of (B3)(i), those of the cubes $\square\not\subseteq\Omega_1^{\sim}$: every such cube is at $\ell^{\infty}$ distance at least $LMR_1$ from $\Omega_1$ ($\Omega_1^{\sim}$ adds a layer of $LMR_1$-cubes, the partitions are nested, and the large-field regions are at distance at least $2LMR_1$ from $\Omega_1$, p.248), and the factor tests the plaquettes of $\square^{\sim}$, whose bonds are within $M_2R_0\leqslant LM_2R_1\leqslant MR_1/8< LMR_1$ of $\square$ (conventions: $R_0\leqslant LR_1$ by CRE (2.9); $M\geqslant 8LM_2$), so it is a function of $V_0$. The factors $\chi_1^c(\square')$, $\square'\subseteq P_1$, and $\chi_1(\square')$ of (1.4) for the cubes $\square'\subseteq P_1^c$ not in $\Omega_1$ (those in $\Omega_1$ form $\chi_1(\Omega_1)$) are functions of the block field $V$ (through the local minima $U_{1,\square'}(V)$ of (1.2)). The gauge-fixing factors $\chi_{Ax}$ of the blocks $\mathcal B_y$, $y\in P_1^1$, not in $\Omega_1^{\sim}$ (those in $\Omega_1^{\sim}$ are removed by (B3)(i)), the exponentials $e^{-g_0^{-2}\mathbf G(P_1^1\setminus\Omega_1^{(1)},U)}$ of (1.6)--(1.7) of the blocks of $P_1^1$ off $\Omega_1$, and the constants $z^{-(L^4-1)| P_1^1\setminus\Omega_1^{(1)}|}$ of the normaliser of (1.5)--(1.6) for these blocks (the exponent of (1.11) counts the blocks of $\Omega_1^{(1)}\subseteq P_1^1$ only): a block off $\Omega_1$ lies in $\Omega_1^c$, since $\Omega_1$ is a union of blocks, and the axial transports $\mathcal U_U(y,x)$ use the bonds inside the block, which are bonds of $\Omega_1^c$, so these are functions of $V_0$. The factors $\chi_0'^c(\square')$, $\square'\subseteq Q_1$, and $\chi_0'(\square')$ of (1.8) for the $L^2M_2R_1$-cubes $\square'\subseteq\Omega_1^c$ not in $Q_1$ (those in $\Omega_1$ form $\chi_0'(\Omega_1)$) involve the fine field on $(\square'^{\sim2})^*$, the bonds within $2L^2M_2R_1+1$ of $\square'$, together with the local minima $U_{1,\square'}(V)$ of (1.2): for $\square'\subseteq Q_1$ these are bonds of $\Omega_1^c$ ($Q_1$ is at distance at least $2LMR_1$ from $\Omega_1$ by (1.10)), and for the other cubes they are bonds of $\Omega_1^c$ and bonds of $\Omega_1$ within $2L^2M_2R_1+1$ of $\Omega_1^c$, which lie in the layer $\Omega_1\setminus\Omega_1^{\sim-1}\subseteq S_1$ (CRE (1.20)), since $LMR_1\geqslant 8L^2M_2R_1$. So $\zeta(\Omega_1^c)$ is a nonnegative function of $V$, $V_0$ and of the fine field on the bonds of $\Omega_1$ within $2L^2M_2R_1+1$ of $\Omega_1^c$, the last dependence coming from the factors $\chi_0'(\square')$, $\square'\subseteq\Omega_1^c\setminus Q_1$, only. CRE's phrase ``localized in $\Omega_1^c$'' refers to the cubes and blocks, and $\zeta(\Omega_1^c)$ is not a function of $V,V_0$ alone. In the variables of (1.22) it is
\begin{equation}\label{eq:c-large-field-factor}
    \zeta_{\xi}(A_{\mathrm e})=\zeta(\Omega_1^c)\bigl(U_{\xi}(a)\bigr)\;,
\end{equation}
with $U_{\xi}(a)$ the reconstructed fine field \eqref{eq:c-reconstructed-field}; the right side does not depend on $a$, because $V$, $V_0$ are branch data and the bonds of $\Omega_1$ on which $\zeta(\Omega_1^c)$ depends are at distance at least $LMR_1-2L^2M_2R_1-2>2R_{\mathrm c}+2L$ from $\Lambda_1$, so their links are the links of far rows, functions of $A_{\mathrm e}$ and the branch data only (Lemma~\ref{lem:row-locality}, (B4)). It is therefore a factor of the $A_{\mathrm e}$ integrand, as written in \eqref{eq:native-source-explicit-outer-operation}; where exactly it stands is immaterial for the identities of this section, which concern the $a$-integral at fixed branch data. The constants: $\sigma_0$ is the Haar constant of (1.14), $z$ the gauge-fixing normaliser of (1.5) and $\mathsf E$ CRE's $E$, all independent of the branch; $z^{(0)}= z^{(0)}(U_1)$ is the factor of the $\delta$-function integration over the designated variables (p.249); it is $\prod_{c\in\Omega_1^{[1]}}|\det_{\mathfrak g}\alpha_c(U_1)|^{-1}$ times a constant fixed by CRE's normalisation of the $\delta$-functions and by $|\Omega_1^{[1]}|$, with $\alpha_c$ of \eqref{eq:regular-stencil-designated-diagonal}, so it depends on the branch through $U_1$ (the linearised constraint $(\widetilde QA')(c)=0$ of p.249 is the row $c$ of $Q_{U_1}A'=0$, whose coefficient at $A'(b_0(c))$ is $\alpha_c(U_1)$, Lemma~\ref{lem:two-block} and RG1 p.267). None of these is used except as a factor of $\kappa_{\xi}$, which is CRE's prefactor verbatim. The retained cutoff $\chi_{\xi,\mathrm e}$ is defined in \eqref{eq:retained-cutoff-composition}: it is the product of the retained factors of CRE (1.25), namely the good free-bond tests of (1.18) at the bonds of $J$, the large-fluctuation tests on $R_1$ and $\chi_0'(S_1)$, all composed with the change of variables (1.22), and it depends on $A_{\mathrm e}$ and the branch data only (Lemma~\ref{lem:collar-coverage}(ii)). The conditional integrand $\mathcal K_{\xi}(a)$, $a\in\mathfrak g^I$, is the integrand of CRE (1.15) in the variables of (1.22) at the actual branch data, divided by the factors already in $\kappa_{\xi}$, $\zeta_{\xi}$ and $\chi_{\xi,\mathrm e}$ and by the factors $\chi_1(\Omega_1)$, $\delta(\bar V_0V^{-1})$ standing outside the $A_{\mathrm e}$ integral; it is written out in \eqref{eq:native-source-actual-functional}, and $\mathcal J^{\mathrm a}_{\xi}[f]=\int f\mathcal K_{\xi}\,\dd a$ is the conditional functional.

\emph{Measurability and integrability.} The integrand of CRE (1.15) with the cutoff (1.25), in the variables $(V,V_0,A_{\mathrm e},a)$ of (1.22), is a finite composition of indicators of sets defined by inequalities between continuous functions (the tests), of the background $U_1(V,V_0)$, of the local minima $U_{1,\square'}(V)$ of (1.2) entering $\chi_0'$ and of the field $\mathbf H_{1,Ax}$ of (1.21) entering the bondwise maps (1.22), of CRE's analytic objects of (B4) and of the bondwise maps (1.22). The background is a Borel function of the admissible pairs, by the following argument, which uses (B1), (B2) and the smoothness of the constraint. Let $\mathcal X$ be the open set of fine configurations $U$ on which $\mathsf M(U)$ is defined on every row of $\Omega_1$, the local group means $\mathcal U_U(y,x)$, $x\in\mathcal B_y$, $y\in\Omega_1^{(1)}$, are defined (chart convention of the preamble; they are analytic there) and every $\alpha_c(U)$, $c\in\Omega_1^{[1]}$, is invertible (Section~\ref{sec:regular-hessian-stencils}), and let $\Phi(U)=(\bar U|_{\Omega_1^{[1]}},U|_{\Omega_1^c})$ on $\mathcal X$; $D\Phi(U)$ is onto, since $LQ_Uh_U=\id$ with $h_U$ supported on the designated bonds of $\Omega_1$, so each fibre $\Phi^{-1}(V,V_0)\cap\mathcal X$ is a $C^{\infty}$ submanifold and the tangent space of the fibre is the kernel of $Q_U$, the range of $\mathsf C_U$. Let $\mathrm{Crit}=\{ U\in\mathcal X:D(S+\mathbf G(\Omega_1^{(1)},\bigcdot ))(U)\circ\mathsf C_U=0\}$ be the set of critical points of CRE's constrained functional of (1.12) on the fibres; it is relatively closed in $\mathcal X$, being defined by the vanishing of a continuous function; the critical points of CRE's functional, the points of the fibre at which the differential of the functional vanishes on the tangent space of the fibre, are the points of $\mathrm{Crit}$ (on $\mathcal X$ the fibres are manifolds with tangent space the range of $\mathsf C_U$, and $U_1\in\mathcal X$ by (B2) and its axial mean gauge). The set $\{(V,U):U\in D^{(11)}(\Omega_1,V)\}$ is Borel in $(V,U)$: it is a finite union of sets defined by non-strict and strict inequalities between functions of $(V,U)$ that are continuous in $U$ and Borel in $V$ (such functions are jointly Borel: $G^E$ is separable metric, and a function continuous in $U$ and Borel in $V$ is a pointwise limit of jointly Borel functions, the standard argument for Carath\'eodory functions), and by the equalities $\bar U(c)= V(c)$ on the rows $c$ of $\Omega_1$ (continuous in $U$ on the open set where the averages are defined) and $\mathsf M(U|_{\Omega_1^c},c)= V(c)$ on the rows of $\Omega_1^c$ (the global average \eqref{eq:global-average}, defined for every $U$ and Borel in $U$: analytic on the open set $\mathcal O_c$ and, off it, the straight transport of the preamble, a continuous function) (the tests involve $V$, $U$ and the Borel functions $U_{1,\square'}(V)$ of (B1); each factor $\chi_0'(\square')$ is zero for $V$ outside the Borel set $\mathcal V_{\mathrm{ax}}(\square')$ of Lemma~\ref{lem:borel-minimiser} for the determining set of $\square'$, where the configuration it compares with does not exist, by the convention of (B1), so that $\{(V,U):\chi_0'(\square')\neq0\}=\{(V,U):V\in\mathcal V_{\mathrm{ax}}(\square'),\ \sup_{b\in(\square'^{\sim2})^*}| U(b)U_{1,\square'}(V)(b)^{-1}-\id|<2\delta_0\}$ is Borel, and $U_{1,\square'}(V)$ is never evaluated off $\mathcal V_{\mathrm{ax}}(\square')$; likewise the factors of (1.4), functions of $V$ alone, are Borel: by the convention of (B1), $\chi_1(\square')=\mathbf 1\{ V\in\mathcal V_{\mathrm{ax}}(\square')\}\,\mathbf 1\{\sup_{p\subset\square'^{\sim}}| U_{1,\square'}(V)(\partial p)-\id|<\varepsilon_1L^{-2}\}$ and $\chi_1^c(\square')=1-\chi_1(\square')$), intersected with the open set where the averages and the means are defined; and by (B1) for every admissible pair the fibre $D^{(11)}(V,V_0)$ contains exactly one point of $\mathrm{Crit}$, namely $U_1(V,V_0)$. So the graph $\{(V,V_0,U):U\in D^{(11)}(\Omega_1,V)\cap\mathrm{Crit},\ \Phi(U)=(V|_{\Omega_1^{[1]}},V_0)\}$ of $(V,V_0)\mapsto U_1(V,V_0)$ is a Borel set (the intersection of a Borel set with the graph of the continuous map $\Phi$ on the open set $\mathcal X$, a Borel set, relatively closed over $\mathcal X$), and a function whose graph is a Borel subset of a product of Polish spaces is Borel on a Borel domain: the domain is the projection of a Borel set with singleton sections, Borel by the Lusin--Novikov theorem \cite[Theorem 18.10]{Kechris}, and a function on a Borel domain whose graph is Borel is Borel, \cite[Theorem 14.12]{Kechris}. Hence the set of admissible pairs is Borel and $U_1$ is Borel on it. The set $\mathcal A_{\xi}$ of \eqref{eq:c-free-variable-domain} is Borel, and $q_{U_1}$ is Borel on it, without circularity between the existence and the uniqueness of the chart-class family: let $\mathsf P$ be the set of pairs $(B,d)\in\mathfrak g^{\Omega_1^*}\times\mathfrak g^{\Omega_1^{[1]}}$ such that $d$ solves every row equation at $g_0\mathsf C_{U_1}B$, is of chart class, and $U_d(B)\in D^{(11)}(V,V_0)$ with the cutoff (1.25) nonzero at $U_d(B)$. It is Borel, jointly with $(V,V_0)$ through the Borel background $U_1$: the row equations are equalities between functions continuous on the open set where the relative averages are defined, the chart-class condition is open, $U_d(B)$ is continuous in $(B,d,U_1)$, $D^{(11)}(V,V_0)$ is a section of the Borel set above, and the cutoff is a finite product of indicators of inequalities between Borel functions. By definition $\mathcal A_{\xi}$ is the projection of $\mathsf P$ to the first factor, and every section $\mathsf P(B)$ has at most one point, by the uniqueness of the chart-class family (Lemma~\ref{lem:c2-fibre-smallness}(iii)); so $\mathcal A_{\xi}$ is Borel by the Lusin--Novikov theorem \cite[Theorem 18.10]{Kechris} (the projection of a Borel set with countable sections), and $B\mapsto d$ is Borel on $\mathcal A_{\xi}$ by \cite[Theorem 14.12]{Kechris} (its graph is $\mathsf P$); hence $q_{U_1}$, $U(B)$ and the row solutions $\widetilde D_c$ of (B4)(i) are Borel on $\mathcal A_{\xi}$, jointly with the branch data; so is $\frac{\delta}{\delta A}\widetilde D$ at the arguments $A= g_0\mathsf C_{U_1}B$, $B\in\mathcal A_{\xi}$, by the following formula. With $A'= q_{U_1}(B)$, the identity $\mathcal Q_{U_1}(A-h_{U_1}\widetilde D(A))= LQ_{U_1}A$ holds on the open set $\mathcal N_{\Theta}\ni A$ of Lemma~\ref{lem:c2-fibre-smallness}(iv), on which $\widetilde D=\varphi$ solves the row equations and is analytic (by \eqref{eq:native-source-constraint-normalisation} and $LQ_{U_1}h_{U_1}=\id$), so by the chain rule $D\mathcal Q_{U_1}(A')\circ(\id-h_{U_1}\frac{\delta}{\delta A}\widetilde D(A))= LQ_{U_1}$. Applied to $h_{U_1}w$, $w\in\mathfrak g^{\Omega_1^{[1]}}$, this gives $D\mathcal Q_{U_1}(A')h_{U_1}\circ(\id-\frac{\delta}{\delta A}\widetilde D(A)h_{U_1})= LQ_{U_1}h_{U_1}=\id$ on $\mathfrak g^{\Omega_1^{[1]}}$, so the square map $D\mathcal Q_{U_1}(A')h_{U_1}$ is invertible, and then $\frac{\delta}{\delta A}\widetilde D(A)=(D\mathcal Q_{U_1}(A')h_{U_1})^{-1}\bigl(D\mathcal Q_{U_1}(A')-LQ_{U_1}\bigr)$, a continuous function of $(U_1,A')$ on the open set where the averages are defined, composed with the Borel maps $U_1$ and $q_{U_1}$; hence it is Borel. The same identity shows that $\id-\frac{\delta}{\delta A}\widetilde D(A)h_{U_1}$ is invertible, hence $\det(\id-h_{U_1}\frac{\delta}{\delta A}\widetilde D(A))=\det(\id-\frac{\delta}{\delta A}\widetilde D(A)h_{U_1})\neq0$ (Sylvester's determinant identity), which is the nonvanishing part of (B4)(iii) at every argument of the domain of integration. Moreover $D(V,V_0)\subseteq\mathcal O^{\Phi}$, the set on which $\Phi$ is a submersion (the paragraph defining $D^{(11)}$): for $B\in\mathcal A_{\xi}$, $A= g_0\mathsf C_{U_1}B$ and $A'= q_{U_1}(B)$, the identity $D\mathcal Q_{U_1}(A')\circ(\id-h_{U_1}\frac{\delta}{\delta A}\widetilde D(A))= LQ_{U_1}$ just derived has right side onto by (B2) and second factor invertible; hence $D\mathcal Q_{U_1}(A')$ is onto. Since $\mathcal Q_{U_1,c}(A')=\frac1i\log(\mathsf M(e^{iA'}U_1)(c)V(c)^{-1})$ with $\mathsf M(e^{iA'}U_1)(c)= V(c)$, where the differential of the chart is the identity, and $A'\mapsto e^{iA'}U_1$ has invertible differential ($A'$ is chart-valued), the differential of $U\mapsto\bar U|_{\Omega_1^{[1]}}$ at $U(B)= e^{iA'}U_1$ is onto, that is, $D\Phi(U(B))$ is onto. Borel as well is $\mathcal D_{\xi}$, the preimage of $\mathcal A_{\xi}$ under $\Psi_{\xi}$, which is continuous in its argument on its open domain and depends on the pair through the Borel field $\mathbf H_{1,Ax}$ below, hence is jointly Borel by the argument just given, and $D(V,V_0)= U(\mathcal A_{\xi})$ is Borel as the image of a Borel set under the injective Borel map $B\mapsto U(B)$ (injective, as shown after \eqref{eq:c-free-variable-domain}), by the Lusin--Souslin theorem for Borel injections \cite[Corollary 15.2]{Kechris}; the Jacobians and the bondwise maps are continuous in their arguments; the local minima $U_{1,\square'}(V)$ of (1.2) are Borel functions of $V$, and the configuration $U_1(\Omega_1\cap\Lambda_1^c,V)$ of (1.21) is a Borel function of $(V,V_0)$ (CRE states that it depends on $V|_{\Omega_1\cap\Lambda_1^c}$; its defining conditions as recorded in (B1) involve $V_0$ where the rows of $(\Lambda_1^{\sim})^c$ cross to $\Omega_1^c$, and the Borel choice is made in the pair): they are the Borel choices of Lemma~\ref{lem:borel-minimiser} fixed in (B1); hence $\mathbf H_{1,Ax}= L\,\frac1i\log\bigl(U_1U_1(\Omega_1\cap\Lambda_1^c,V)^{-1}\bigr)$ of (1.21) is Borel in $(V,V_0)$ on the bonds of $S_1^*$, where alone it is used (the chart logarithm is continuous on the chart domain, which contains the argument there by CRE's bound p.251 with the constants of item 5 of Section~\ref{app:thresholds}, (H4), \eqref{eq:boundary-bijection-inputs}). So the integrand is jointly Borel, and nonnegative. By Tonelli's theorem the section integral $\mathcal J^{\mathrm a}_{\xi}[1]=\int\mathcal K_{\xi}\,\dd a$ is a $[0,\infty]$-valued Borel function of the branch data, and $\int\mathcal O[\mathcal J^{\mathrm a}_{\xi}[1]](V)\,\dd V$ is the integral of the nonnegative integrand over all variables, equal to the fine integral $\mathcal Z_0$ of \eqref{eq:native-source-physical-ratio} below, which is finite (a bounded function on a compact group); hence the outer integrals converge, $\mathcal J^{\mathrm a}_{\xi}[1]<\infty$ for almost every value of the branch data, and on every branch $\mathcal J^{\mathrm a}_{\xi}[1]<\infty$, since the integrand \eqref{eq:native-source-actual-functional} is bounded on the bounded box $\{\Theta_{\xi}=1\}$ (its exponent is the sum of an analytic function bounded on the closed box, Lemma~\ref{lem:charts}, and of a constant of the branch, Lemma~\ref{lem:collar}). The Wilson function $f_{\xi}(a)= W_{\chi}(C;U_{\xi}(a))$ is jointly Borel for the same reason and $| f_{\xi}|\leqslant 1$, so $\mathcal J^{\mathrm a}_{\xi}[f_{\xi}]$ is a Borel function of the branch data with $|\mathcal J^{\mathrm a}_{\xi}[f_{\xi}]|\leqslant \mathcal J^{\mathrm a}_{\xi}[1]$, and Fubini's theorem applies to it. No bound on $\mathcal J^{\mathrm a}_{\xi}[1]$ uniform in the branch is claimed or needed here. Then
\begin{equation}\label{eq:native-source-physical-ratio}
    \mathcal Z_0=\int\mathcal O[\mathcal J^{\mathrm a}_{\xi}[1]](V)\,\dd V,
    \qquad
    \mathbb E_{\varepsilon}[W_{\chi}(C)]
        =\frac{1}{\mathcal Z_0}\int\mathcal O[\mathcal J^{\mathrm a}_{\xi}[f_{\xi}]](V)\,\dd V\;,
\end{equation}
with $f_{\xi}(a)= W_{\chi}(C;U_{\xi}(a))$ the Wilson function of the reconstructed fine field $U_{\xi}(a)$ of \eqref{eq:c-reconstructed-field}. The first identity is obtained by CRE's steps from the fine integral to (1.15) with the cutoff (1.25), performed with $\zeta(\Omega_1^c)$ kept inside the $V_0$ integral and the cutoff in its exact composed form (CRE's estimate-based removals are recorded in (B3)); the second is the same derivation with the bounded function $W_{\chi}(C;\bigcdot )$ of the fine field, invariant under all fine gauge transformations, inserted under the fine integral. The steps, in the following order, which is CRE's except at two places (p.248 drops $\chi_0$ and $\chi_{Ax}$ before the resummation (1.11), and p.249--250 integrates the $\delta$-functions over the designated variables before inserting (1.18); the operations concerned are multiplications of the integrand by functions of the fine field equal to one, regroupings of a finite sum and the evaluation of $\delta$-functions, and the resulting integrand, a function of the fine field, is the same in either order): the insertion of the block-average $\delta$-functions and of the partitions of unity (1.1), (1.4) (multiplications by functions equal to one, invariant under the fine gauge transformations of (1.5): the tests of (1.1) are plaquette tests and those of (1.4) are functions of $V$), with the group Dirac measure normalised against Haar measure, $\int\delta(\bar UV^{-1})\,\dd V=1$ for every $U$ (the sentence following \eqref{eq:bare-blocking-operator}; CRE (0.1) is the blocking operator whose $\delta$-functions these are; this is why $\mathcal Z_0=\int\varrho_1\,\dd V$ in CRE's notation); the Faddeev--Popov gauge fixing (1.5)--(1.6), an identity for the integral of any function invariant under the gauge transformations $u$ of (1.5), which introduces the normaliser $z^{-(L^4-1)| P_1^1|}$ of (1.5) for the blocks of the gauge-fixed region; the insertion of the partition of unity (1.8) after the gauge is fixed (its tests compare $U$ with the axial-gauge minima $U_{1,\square'}(V)$ and are not gauge invariant); the resummation (1.11) over $P_0,P_1,Q_1$ at fixed $\Omega_1$ (p.248, a regrouping of a finite sum; $R_1$ is introduced afterwards, in (1.18)), after which the gauge-fixed blocks displayed are those of $\Omega_1^{(1)}$ and the displayed normaliser is $z^{-(L^4-1)|\Omega_1^{(1)}|}$, the gauge-fixing factors and normalisers of the other blocks of $P_1^1$ being part of $\zeta(\Omega_1^c)$; the removal of (B3)(i), which replaces factors by one where the retained factors are nonzero; the change of variables $U= U'U_1$ on the bonds of $\Omega_1$ at fixed $(V,V_0)$ (a translation of the Haar measure) and $A'=\frac1i\log U'$ of (1.13), the chart being a bijection of the domain onto its image with the Haar density as Jacobian, $\dd U'=\sigma(A')\,\dd A'$ bond by bond with $\sigma(0)=\sigma_0$ (the domain lies in the chart domain by (B1)); the conversion of the group $\delta$-functions on the rows of $\Omega_1$: by \eqref{eq:native-source-constraint-normalisation} and $\mathsf M(U_1)(c)= V(c)$, $\mathsf M(e^{iA'}U_1)(c)V(c)^{-1}= e^{i\mathcal Q_{U_1,c}(A')}$, and $\delta_G(e^{iX})=\sigma(X)^{-1}\delta(X)=\sigma_0^{-1}\delta(X)$ on the support, so the row $\delta$-functions become $\sigma_0^{-|\Omega_1^{[1]}|}\prod_c\delta(\mathcal Q_{U_1,c}(A'))$ (the transformation rule of the paragraph preceding \eqref{eq:c-free-variable-domain}, applied at the level of the pair $(V,V_0)$ itself, as is every later change of variables of this list: the fibre measures are transported at that level, not for almost every level); the linearising substitution $A'= B-h_{U_1}\widetilde D(B)$ of (1.14) (CRE writes both variables as $A$; in this step list $B$ is the unscaled variable of (1.14) and $A$ the scaled one of p.249, whereas Section~\ref{sec:branch-data} writes $B$ for the scaled free vector and $A= g_0\mathsf C_{U_1}B$ for the unscaled kernel vector, and the lemmas follow Section~\ref{sec:branch-data}), the change of variables by the diffeomorphism $\Theta$ of the paragraph preceding \eqref{eq:c-free-variable-domain}, defined there on the set $\mathcal N_{\Theta}$ of Lemma~\ref{lem:c2-fibre-smallness}(iv), which does not depend on $R_1$, since this step precedes the insertion of (1.18), which introduces $R_1$ ((B4)(i), (iii) and Lemma~\ref{lem:c2-fibre-smallness}(iv); $\Theta$ is injective with $\mathcal Q_{U_1}\circ\Theta= LQ_{U_1}$), with the Jacobian $|\det(\id-h_{U_1}\frac{\delta}{\delta A}\widetilde D)|$ (the first fact of that paragraph, at the level $0$), after which $\mathcal Q_{U_1}(A')= LQ_{U_1}B$ (the equation of RG1 p.267, \eqref{eq:native-source-local-implicit}) and the constraint is the linear constraint $\delta(\widetilde QB)$ of (1.14), CRE's $\widetilde Q$ being read as $LQ_{U_1}$ (RG1's $\widetilde Q$ in the equation quoted at \eqref{eq:native-source-local-implicit} is $Q_U$, so that RG1's $L\widetilde Q$ is this operator; $\delta(LQ_{U_1}B)= L^{-d_G|\Omega_1^{[1]}|}\delta(Q_{U_1}B)$, and the constant $L^{-d_G|\Omega_1^{[1]}|}$ of this rescaling of the $\delta$-functions is part of $z^{(0)}$, p.249), while the Haar factor $\prod_{b\in\Omega_1}\sigma(A'(b))$ is written as $\sigma_0^{\#\{\text{bonds of }\Omega_1\}}$ times the normalised density $\prod_b\sigma(A'(b))/\sigma_0$, CRE's $\sigma(A-h\widetilde D(A))$ in (1.14)--(1.15) (reading CRE's $\sigma$ there as the normalised density, the reading consistent with the exponent $|\Omega_1^*|\log\sigma_0$ of (1.14)): the $\sigma_0^{|\Omega_1^{[1]}|}$ of this product cancels the $\sigma_0^{-|\Omega_1^{[1]}|}$ of the conversion, and $\sigma_0^{\#\{\text{bonds of }\Omega_1\}-|\Omega_1^{[1]}|}=\sigma_0^{|\Omega_1^*|}$ is CRE's factor in (1.14) and (1.15); the scaling $B= g_0A$ of p.249 (CRE: $A= g_0A'$), with Jacobian $g_0^{d_G\times(\text{number of bonds of }\Omega_1)}$ and $\delta(\widetilde Qg_0A)= g_0^{-d_G|\Omega_1^{[1]}|}\delta(\widetilde QA)$, so that the power of $g_0$ is $g_0^{d_G|\Omega_1^*|}$; the insertion of the partition of unity (1.18) in the scaled variable (p.250) and the removal of (B3)(ii), which together produce the cutoff (1.25); the integration of the $\delta$-functions $\delta(\widetilde QA)$ over the designated variables $A(b_0(c))$, $c\in\Omega_1^{[1]}$, which is the elimination of (B2) and produces the factor $z^{(0)}(U_1)$, after which the remaining Lebesgue measure is $\dd A|_{\Omega_1^*}$ and the normalised Haar density remains at every bond of $\Omega_1$, at the designated bonds evaluated at the eliminated values (it is the Haar primitive \eqref{eq:native-source-haar-primitive}, part of CRE's $v$); and the bondwise change of variables (1.22) on $S_1^*$, a change of variables by (B4)(ii) with the Jacobian $|\det D\Psi_{\xi}|$. Each step is an identity for the integral of any bounded measurable function of the fine field, and the fine field is expressed through the integration variables by the reconstruction \eqref{eq:c-reconstructed-field}. No factor is moved out of the history sum. The outer operation is CRE's expansion at the stage of (1.15) with the cutoff (1.25); CRE's conditional integration (1.26), its cluster expansion and the renormalisation (1.27)--(1.28) are replaced by the construction of this section and are not used, except that the first-step coupling $g_1(g_0)$ defined by (1.27) on $\Lambda_1^{\sim-1}$ enters, as a parameter depending on $g_0$ alone, the scale $R_1$ and the quantity $\varepsilon_1$ of Lemma~\ref{lem:regular-bounds} (Section~\ref{sec:branch-data}). The following facts about the branch density are taken from CRE as stated there, except where an item records a reading, a derivation or a proof (among them the Borel choices, the convention and the derivations in (B1), the Remark under (B2), the parts of (B4)(i) and (iii) proved in Lemma~\ref{lem:c2-fibre-smallness}, and the proofs in (B4)(ii)). Together with the lemmas proved in Appendix~\ref{app:background-estimates}, the standing conditions, and the further CRE statements named in this section (the definitions (1.1)--(1.10) and (1.18)--(1.22), the local minima $U_{1,\square'}$ of (1.2) and the field $\mathbf H_{1,Ax}$ of (1.21), the axial mean gauge of $U_1$ of p.248, the scale relations (2.6) and (2.9) of p.255--256, the coupling map (1.27) with its $\beta_1$, (2.16)--(2.17) for $\chi_1(\Omega_1)$, (2.4)--(2.5) and p.246 for $\varepsilon_j$ and $R_j$, (1.17) for the gauge action on phases, and, through (B1), Theorem 1 of \cite{BalabanVAR}, CRE's [15]; (26) of \cite{BalabanVAR} is cited by CRE at p.248 for the expansion of the action), they are the inputs of this section:
\begin{itemize}
\item[(B1)] CRE's (1.11) is taken as it stands, with its domain of integration $D^{(11)}$ (Section~\ref{sec:branch-data}); the local minima $U_{1,\square'}(V)$ of (1.2) are minimal configurations of CRE's variational problem (2.12) for the determining set of (1.2), elements of the unique minimal orbit (p.256) taken in the axial gauge (p.247); the configuration $U_1(\Omega_1\cap\Lambda_1^c,V)$ of (1.21) is determined, p.251, by the conditions of (1.12) on the domain $(\Lambda_1^{\sim})^c$ (the constraint $\bar U= V$, the rows crossing to $\Omega_1^c$ involving $V_0$ as in (B2), and the minimality of the constrained functional, whose critical point is its minimum, p.248) and on $\Lambda_1^{\sim}\setminus\Lambda_1$ by the determining set $\mathbf B_1(\Lambda_1^c)$ and the configuration $Q_1^{s*}V$, in the axial gauge; CRE names the determining data only, and that the configuration is the minimal configuration in the sense of (2.12) (p.256: the functional $A^{\eta}$ minimised over the regular configurations satisfying the constraints) for the constraints named, is this paper's reading of p.251, recorded as such; CRE's ``depends on $V$ restricted to $\Omega_1\cap\Lambda_1^c$'' is read as its statement about the dependence on $V$, the crossing rows involving $V_0$ as for $U_1$, so that the Borel choice below is made in the pair $(V,V_0)$. CRE does not specify the representatives further, and its integrals over $V$ presuppose a measurable choice. By Lemma~\ref{lem:borel-minimiser}, which covers conditions that are equalities between continuous functions and the minimality of a continuous functional over a set so defined, a Borel choice exists for each of the two problems, in $V$ for (1.2) and in $(V,V_0)$ for (1.21); one such choice is fixed for each, and CRE's statements about these configurations that are used here (the tests (1.4) and (1.8), (B3)(iii), \eqref{eq:boundary-bijection-inputs}) are used for the chosen configurations, CRE stating them for minimal configurations in the axial gauge without specifying the representative (where the stabiliser of the coarse data under the gauge transformations at the block centres is nontrivial the axial representative is not unique, and the tests of (1.4) and (1.8) may depend on it; that CRE's statements hold for the chosen representative is this paper's reading of CRE's definitions through an unspecified representative, named in (H4)). The domain of (1.11) is described here. The averages and means of (1.11) are the principal logarithms of \cite[p.~21, (22)--(25)]{BalabanAVG}, Ba{\l}aban's own convention, restricted in the preamble to the chart ball. That on the domain of integration every argument of a logarithm of RG1 (0.4) on a row of $\Omega_1$ lies in the chart domain, every family of transports entering a mean (0.10)--(0.11) at a block centre of $\Omega_1^{(1)}$ lies in the small-diameter domain of the mean, and the fluctuation $U'= UU_1^{-1}$ lies in the chart domain on every bond of $\Omega_1$, the designated bonds $b_0(c)$ included, so that $A'=\frac1i\log U'$ of (1.13) is defined, is CRE's smallness statement (pp.247--248), which CRE justifies by the reasoning of the proof of Lemma~1 of \cite{BalabanREG}; it is derived here from CRE's statements with that proof as the model of CRE's reasoning, not from REG's lemma applied as stated (its hypotheses, the exact axial condition and REG's (1.7) on the blocks concerned, are not literally those of CRE's domain), as follows. (i) On the domain of integration the cutoffs $\chi_0$ and $\chi_{Ax}$ of (1.5)--(1.6) localised in $\Omega_1^{\sim}$ equal one ((B3)(i), p.248): $| U(\partial p)-1|<\varepsilon_0$ for every plaquette $p\subset\Omega_1^{\sim}$ and $| U(y,x)-1|<\varepsilon_0$ for every mean (0.11) at $y\in\Omega_1^{\sim(1)}$, $x\in B(y)$. (ii) For such $U$ CRE states $| UU_{1,\square'}^{-1}-1|< O(L^2)\varepsilon_0$ on every bond of $\square'^{\sim2}$ (p.247, ``on $\square'^{\sim2}$'', the designated bonds not excluded, whereas the test (1.8) runs over $(\square'^{\sim2})^*$), ``by the same reasoning as in the proof of Lemma~1 [14]'', that is Lemma~1 of \cite{BalabanREG} (p.~79, (1.24)--(1.25)): if two configurations $V_0$ and $V'V_0$ have plaquette variables within $\alpha_0L^{-2}$ of one on $\Omega_1$ ((1.7) there), the same variables on the axial contours $\Gamma_{y,x}$, $x\in B(y)$, and block averages within $\alpha_1$ of each other on $\Omega_1^{(1)}$, then $| V'-1|<4d^2\alpha_0+\alpha_1$ on every bond of $\Omega_1$; its proof (pp.~79--80) bounds the bonds inside a block from the plaquettes and the axial conditions, the bonds joining two blocks from these, and the designated bond $b_0$ of each coarse bond $c$ from the block-average condition (pp.79--80), and the bound is local (p.80); CRE applies the reasoning with $V_0= U_{1,\square'}$ in the axial gauge, the block averages of both configurations equal to $V$ on the blocks of $\square'^{\sim4}$ ((1.2)--(1.3), p.246, and the constraint of (1.11)) and the axial condition on $U$ holding within $\varepsilon_0$ by (i). (iii) $U_1$ lies in the domain of integration (p.248; CRE's ``in the domain of integration in (1.11)'' is read as the support of every characteristic factor of (1.11), the reading named in (H4) and Remark~\ref{rem:c3-8-1}), so (ii) applies to $U_1$ too, and $U'=(UU_{1,\square'}^{-1})(U_{1,\square'}U_1^{-1})$ is within $2\,O(L^2)\varepsilon_0$ of one on every bond of $\square'^{\sim2}$, for every cube $\square'\subset\Omega_1$ (p.248: ``both fields in the parenthesis \dots\ are small''). (iv) The arguments of RG1 (0.4) on a row of $\Omega_1$ are closed contours in two adjacent blocks of $\Omega_1^{\sim}$, and two transports of a family $\mathbf G(y,x)$ differ by such a contour; by the reasoning of \cite[pp.~24--25]{BalabanAVG} ((44)--(46): in a local axial gauge every bond of the four blocks $\Delta(p')$ is within $dL\alpha_0$ of one and every contour $\Gamma$ in them within $|\Gamma| dL\alpha_0= O(1)L^2\alpha_0$, ``for $L^2\alpha_0$ sufficiently small''), with $\alpha_0=\sqrt N\varepsilon_0$ from (i) (AVG's $|\bigcdot |$ is the operator norm, which the norm of (i) bounds up to the factor $\sqrt N$), every such closed contour variable, a conjugation-invariant quantity, is within $O(1)L^2\varepsilon_0$ of one, and so is the diameter of each family of transports. The constant $O(1)$ of (iv) is displayed by \cite{BalabanAVG}, $|\Gamma|\leqslant (2d+1)L$ and the factor $d$, up to the elementary count of the contours of RG1 (0.4) and (0.10)--(0.11) in the four blocks and the factor $\sqrt N$; the constant $O(L^2)$ of (ii) is not displayed by CRE and, like the constants of (B3)(iii), is taken to be independent of the coupling and of the branch (this paper's reading, named in (H4)), so that for $\varepsilon_0$ small at fixed $L$ the quantities of (iii) and (iv) lie below $\varrho_G$, $r_G/\sqrt N$ and the small-diameter radius of the mean (CRE's ``$\varepsilon_0$ sufficiently small'', p.248; item 5 of Section~\ref{app:thresholds}). This is the domain $D^{(11)}$ of Section~\ref{sec:branch-data}. A convention of this paper: every factor of (1.4) and (1.8) defined through the configuration $U_{1,\square'}(V)$ of (1.2) is assigned a value for the $V$ for which that configuration does not exist ($V\notin\mathcal V_{\mathrm{ax}}(\square')$, Lemma~\ref{lem:borel-minimiser}), namely $\chi_1(\square')=0$, $\chi_1^c(\square')=1$, $\chi_0'(\square')=0$ and $\chi_0'^c(\square')=1$, so that (1.4) and (1.8) remain partitions of unity (CRE's (1.2)--(1.4) and (1.8), p.246--247, define these factors through $U_{1,\square'}(V)$ without an existence clause; the convention is this paper's), and $U_{1,\square'}(V)$ enters only where it exists. Under it, where $\chi_1(\Omega_1)\neq0$ the configuration $U_{1,\square'}(V)$ exists for every cube $\square'\subset\Omega_1$, and where a factor $\chi_0'(\square')$ is nonzero it exists for that cube. For every admissible $(\Omega_1,V,V_0)$ the background $U_1= U_1(V,V_0)$ of (1.12) exists, is the unique critical point of CRE's constrained functional in the domain of integration in (1.11), that is in the fibre $D^{(11)}(V,V_0)$, and is a minimum with strictly positive second differential (CRE p.248: ``there exists exactly one critical point'' of (1.12) in the domain of integration in (1.11), by Theorem 1 of \cite{BalabanVAR} (CRE's reference [15]: Ba{\l}aban, Commun. Math. Phys. \textbf{102}, 277--309 (1985), the variational problem of the averaging operation) for $\varepsilon_0$ small; ``it is a minimum of (1.12), with a strictly positive second order differential'', stated there without further citation); the fluctuation $U'= UU_1^{-1}$ lies in the chart domain for every $U\in D^{(11)}(V,V_0)$ (p.248, derived above), so that $A'=\frac1i\log U'$ of (1.13) is defined.
\item[(B2)] For every admissible pair $(\Omega_1,V,V_0)$ the constraint $\delta(\widetilde QA')$ of (1.14) is solved on every row $c\in\Omega_1^{[1]}$ for the designated variable $A'(b_0(c))$ as a linear function of the free variables of the two blocks of $c$ (CRE p.249, stated at the critical point $U_1$, which exists for every admissible pair by (B1); CRE states this for the rows of $\Omega_1$ without distinguishing the rows with one block off $\Omega_1$; this paper's reading, used below, is that on such a row the average involves $V_0$ on the outer block as a constant and $b_0(c)$, a bond of $\Omega_1$, is the eliminated variable, following CRE's choice of $A'(b_0(c))$ as the variable solved for on every row, p.249). In the notation of Section~\ref{sec:regular-hessian-stencils} this is the invertibility of $\alpha_c(U_1)$ for every $c\in\Omega_1^{[1]}$, so that $U_1$ lies in the open set $\mathcal X$ of the measurability paragraph for every admissible pair.
\emph{Remark.} The invertibility is also proved here, at every row $c\in\Omega_1^{[1]}$, near or far: $\alpha_c(U)$ depends on $U|_{\mathcal E_c}$ only, Lemma~\ref{lem:two-block} gives its invertibility for backgrounds $e^{iA}$ with $\| A\|_{\infty}\leqslant \tau_L$ on $\mathcal E_c$ and shows that a site gauge transformation multiplies it on both sides by invertible maps, and the first assertion of Lemma~\ref{lem:c2-fibre-smallness}(i) of Appendix~\ref{app:background-estimates}, proved from (B1)(i), (iii), (iv), the first condition of threshold item 6 (with \eqref{eq:c2-tau-final}), the fine periods being at least $8M$, the inclusion $\mathcal E_c\subseteq\mathfrak P_c$ given by \eqref{eq:regular-stencil-bonds} and Lemma~\ref{lem:c2-local-gauge}, without (B2), gives a site gauge on the two blocks of $c$ in which $U_1= e^{i\phi}$ on $\mathcal E_c$ with $\|\phi\|_{\infty}\leqslant \tau_L/2$. So of (B2) only its reading of the rows with one block off $\Omega_1$ (the outer block entering as the constant $V_0$, and $A'(b_0(c))$ the eliminated variable) is an input.
\item[(B3)] The estimate-based simplifications in CRE's derivation of (1.25) hold under CRE's parameter hierarchy, in the following form: the dropped factors are equal to one wherever the retained factors are nonzero. (i) The removal of the factors $\chi_0(\square)$ and $\chi_{Ax}$ of the cubes and blocks localised in $\Omega_1^{\sim}$ after (1.10) ($A_0$ sufficiently larger than $A_1$, p.248). (ii) The removal of the second factor of (1.19) on the good cubes ($M_2R_1$ large, p.250), which leaves $\chi_0'$ on the cubes of $S_1$ only (p.250); the condition ``$M_2R_1$ large enough'' is a threshold on the coupling (it makes $O(1)B_3e^{-\delta_{\mathrm c}LM_2R_1}\varepsilon_1$ much smaller than $\delta_0$, both sides depending on $g_0$ through $\varepsilon_1$, $\delta_0$ and $R_1$), taken over in item 5 of Section~\ref{app:thresholds}. (iii) The comparison on which (ii) rests, p.250: for every $L^2M_2R_1$-cube $\square'\subseteq\Omega_1^{\sim-1}\setminus R_1$ (CRE's $\square'\subset R_1^c$; then $\square'^{\sim4}\subseteq\Omega_1$, p.250), (by the convention of (B1), $U_{1,\square'}(V)$ exists for every cube $\square'\subset\Omega_1$ where $\chi_1(\Omega_1)\neq0$, so the comparison is between two existing configurations), $| U_1(b)U_{1,\square'}(b)^{-1}-\id|< O(1)B_3e^{-\delta_{\mathrm c}LM_2R_1}\varepsilon_1$ for the bonds $b$ of $\square'^{\sim2}$, with the constant $B_3$, the decay rate $\delta_{\mathrm c}$ (CRE's $\delta$) and the $O(1)$ independent of the coupling and of the branch (CRE states no dependence for them: it uses the same symbols $B_3$, $\delta$ at every step and derives the comparison (3.19) from these facts and the bounds (106)--(108), (159)--(163) of [12], p.268; only their independence from the coupling is used, in Appendix~\ref{app:background-estimates}, Lemma~\ref{lem:c2-plaquette-deviation}). Item (iii) is used in Appendix~\ref{app:background-estimates}, Lemma~\ref{lem:c2-plaquette-deviation}; item (i) is used only through its conclusion, the form (1.11) of the integrand and hence the form (1.25) of the cutoff, and item (ii) through (1.25) and directly in the proof of Lemma~\ref{lem:collar-coverage}(iv). CRE's rewriting of the retained tests in the new variables (p.252) is not used: the retained cutoff is defined in \eqref{eq:retained-cutoff-composition} as the exact composition.
\item[(B4)] For every admissible pair $(\Omega_1,V,V_0)$ and every region $R_1$ of the history sum, CRE's expressions in (1.14)--(1.15) and (1.22)--(1.24) are defined on the domain of integration: the linearising correction $\widetilde D$ of (1.14), p.249 (at the foot of p.248 CRE refers to the transformations analysed between (16)--(18) of [16], or between (I.1.8)--(I.1.9) of [I], where [I] is RG1 in CRE's reference list, p.285, Commun. Math. Phys. 109, 249--301; RG1's (1.8)--(1.9), p.261, concern the functions $J_j$, and the linearisation of RG1's unit-lattice step is the one following RG1 (2.10), p.267, where RG1 states for its correction that its defining equation has exactly one solution, analytic in its argument, with an expansion beginning with quadratic terms; here $\widetilde D$ is the small solution $\varphi$ of (i)); the phase $A-h_{U_1}\widetilde D(A)$; the Haar density $\sigma(A-h_{U_1}\widetilde D(A))$; the logarithm $\operatorname{Tr}\log(\id-h_{U_1}\frac{\delta}{\delta A}\widetilde D)$ of the Jacobian of (1.14); and the bondwise maps of (1.22) with their Jacobian. Three further properties of these objects are recorded as part of (B4); they are CRE's and RG1's statements, this paper's readings, marked as such, and facts proved in Lemma~\ref{lem:c2-fibre-smallness}. (i) This paper's reading of RG1 p.267 and CRE pp.248--249: the defining equation of $\widetilde D$ is the system \eqref{eq:native-source-local-implicit} over the rows $c\in\Omega_1^{[1]}$, which decouples row by row (Lemma~\ref{lem:row-locality}); the row component $\widetilde D_c(A)$ depends on $A|_{\mathcal E_c}$ only. RG1 p.267 states for the equation $L\widetilde QB'+\widetilde C(B')= L\widetilde QB$, $B'= B-h\widetilde D(B)$, that ``there exists exactly one solution of this equation, and that it is an analytic function of $B$'', in RG1's unit-lattice step, without naming the class in which the solution is unique; this statement is quoted as Ba{\l}aban's description of $\widetilde D$ and is not used beyond the identification just made. The row equation \eqref{eq:native-source-local-implicit} involves the eliminated phase only through $e^{iq}$, so its solutions are determined only up to the periods of the exponential; the class used here is the chart class of Section~\ref{sec:branch-data} (the whole phase in the chart ball $h(q)<\varrho_G$, the class of CRE's chart-valued variable $A'$ of (1.13)) with the fine field in $D^{(11)}(V,V_0)$. In this class the row solution is determined by Lemma~\ref{lem:c2-fibre-smallness} of Appendix~\ref{app:background-estimates}, proved from (B1) (with (B3)(i), through (B1)(i)), (B2), the reading in (H4) of CRE's compatible partitions as nested, the fine periods being at least $8M$, Lemmas~\ref{lem:c2-local-gauge}, \ref{lem:c2-uniform-chart} and~\ref{lem:two-block}, and threshold items 5 and 6: at every row $c\in\Omega_1^{[1]}$, near or far, Lemma~\ref{lem:c2-uniform-chart} applies to $U_1$ and gives the small solution $\varphi_c$, analytic on the ball $\| A|_{\mathcal E_c}\|<\vartheta_*$, real on real phases, with $\varphi_c(0)=0$ and $D\varphi_c(0)=0$, unique among the solutions $d_c$ with $h(d_c)\leqslant \rho_{\mathrm d}$ (Lemma~\ref{lem:c2-fibre-smallness}(i)); the argument $A$ of every $U\in D^{(11)}(V,V_0)$ lies in the ball $\| A|_{\mathcal E_c}\|<\vartheta_*/2$ at every row (Lemma~\ref{lem:c2-fibre-smallness}(ii)); and every chart-class family $d$ at $g_0\mathsf C_{U_1}B$ with $U_d(B)\in D^{(11)}(V,V_0)$ is $\varphi(g_0\mathsf C_{U_1}B)$, $\varphi=(\varphi_c)_c$ (Lemma~\ref{lem:c2-fibre-smallness}(iii)). The row solution of this paper is $\widetilde D_c=\varphi_c$ at every row, and $\widetilde D=\varphi$ is analytic and solves the row equations on the open set $\mathcal N_{\Theta}$ of Lemma~\ref{lem:c2-fibre-smallness}(iv), which contains all arguments of the domain of integration. CRE's $\widetilde D$ is $\varphi$ wherever CRE defines it as a solution of the row system with chart-valued phase and fine field in $D^{(11)}(V,V_0)$, by Lemma~\ref{lem:c2-fibre-smallness}(iii); that this is CRE's use of $\widetilde D$ in (1.13)--(1.14), pp.248--249, is this paper's reading, recorded as such. (ii) The substitutions (1.14) and (1.22) are changes of variables: the map $B\mapsto U(B)$ of Section~\ref{sec:branch-data}, the composition of (1.13) and (1.14), is a bijection of $\mathcal A_{\xi}$ onto $D(V,V_0)$ (shown there from (B1), (B2) and Lemma~\ref{lem:c2-fibre-smallness}), and its differential, that of $B\mapsto g_0\mathsf C_{U_1}B$ followed by the substitution $A'= A-h_{U_1}\widetilde D(A)$ and the exponential, is injective at every point of $\mathcal A_{\xi}$ with image the tangent space of the fibre $\Phi^{-1}(V,V_0)$ at $U(B)$ (the Jacobian $\det(\id-h_{U_1}\frac{\delta}{\delta A}\widetilde D)$ of (1.14) is nonzero there, (iii), $\mathsf C_{U_1}$ is injective, and the image of $(\id-h_{U_1}\frac{\delta}{\delta A}\widetilde D)\mathsf C_{U_1}$ is the kernel of $D\mathcal Q_{U_1}(A')$ by the identity $D\mathcal Q_{U_1}(A')\circ(\id-h_{U_1}\frac{\delta}{\delta A}\widetilde D)= LQ_{U_1}$ of the measurability paragraph, $LQ_{U_1}\mathsf C_{U_1}=0$, and a count of dimensions, $D\mathcal Q_{U_1}(A')$ being onto; the map $B\mapsto U(B)$ goes from $\mathfrak g^{\Omega_1^*}$ into the fibre, so its differential is not square); and $\Psi_{\xi}$ of \eqref{eq:native-source-exact-boundary-map}, acting bond by bond on $S_1^*$, is a $C^{\infty}$ map which maps $\mathcal D_{\xi}$ bijectively onto $\mathcal A_{\xi}$ with invertible differential at every point of $\mathcal D_{\xi}$: this is Lemma~\ref{lem:boundary-bijection} of Section~\ref{sec:boundary-law-collar}, proved from CRE's bound on $\mathbf H_{1,Ax}$ (p.251), cited there, from the membership $U_1\in D^{(11)}(V,V_0)$ of (B1), that is the cutoff $\chi_0'(\Omega_1)$ at $U_1$ (CRE p.248), and from the choice of the cutoff $g$ of (1.22) made there (pp.249 and 252 are where CRE performs the substitutions with the Jacobians written there; the invertibility of $D\Psi_{\xi}$ is used in Lemma~\ref{lem:collar}(ii)). The Jacobian of (1.22) is used here only as $|\det D\Psi_{\xi}|$, the exact Jacobian of \eqref{eq:native-source-exact-boundary-map}, not through CRE's expression (1.24). (iii) The determinant $\det(\id-h_{U_1}\frac{\delta}{\delta A}\widetilde D)$ is nonzero on $\mathcal A_{\xi}$, by (i) and Sylvester's identity (the measurability paragraph of Section~\ref{sec:branch-data}); CRE performs (1.14) as a change of variables and writes its Jacobian as $\exp\operatorname{Tr}\log(I-h(\frac{\delta}{\delta A}\widetilde D)(A))$ (p.249). The determinant is positive there, so that CRE's logarithm is the real logarithm and equals $\sum_c\log\det_{\mathfrak g}(\id-h_{U_1}(c)\partial_{A(b_0(c))}\widetilde D_c)$ (Lemma~\ref{lem:primitive-sum}): this is Lemma~\ref{lem:c2-fibre-smallness}(iv), \eqref{eq:c2-fibre-jacobian}, at every row; the identities of this paper use the absolute Jacobian $|\det(\id-h_{U_1}\frac{\delta}{\delta A}\widetilde D)|$, for which the nonvanishing suffices. For the rows within $2R_{\mathrm c}+2L$ of $I$ and the inputs of Lemma~\ref{lem:charts}, the expressions of (B4) coincide with the analytic objects of that lemma, by the uniqueness of the small solution stated there; elsewhere they enter only through $a$-independent factors (Lemma~\ref{lem:collar}, \eqref{eq:c-large-field-factor}).
\end{itemize}

The Borel choice fixed in (B1) exists:

\begin{lemma}[Borel choice of minimal configurations]\label{lem:borel-minimiser}
Let $E$ be a finite set of fine bonds, $\mathbf B$ a set of blocks, and $\mathcal U\subseteq G^{E}$ a finite intersection of open and closed subsets of the compact metric space $G^E$, contained in the open set on which the group means $U(y,x)$ of RG1 (0.11) at the block centres $y$ of $\mathbf B$ are defined (chart convention of the preamble). For CRE's determining sets $\mathbf B$, that of (1.2) and that of (1.21), $E= E(\mathbf B)$ is the set of fine bonds of the region determined and $\mathcal U=\mathcal U_{\mathbf B}$ the set of regular configurations, on which the averaging operation $M_{\mathbf B}$ of CRE (2.11) is defined: the configurations of the class (2) of \cite{BalabanVAR} (CRE's [15]), p.278: finitely many strict inequalities between continuous functions of $U$ (the plaquette variables $| U(\partial p)-1|$ and the covariant divergences $|(D_U^{\eta*}\partial U)(b)|$, polynomials in the matrix entries of the bond variables and their inverses), with the averages (the logarithmic averages of RG1 (0.4)) and the means $U(y,x)$ defined in the chart convention, an open set; VAR's space (6) is its intersection with the constraint set (3), which is the constraint $\mathfrak c=0$ below; $M_{\mathbf B}$ is continuous on $\mathcal U$. Let $P$ be a Polish space of parameters, $\mathcal W$ a relatively open subset of $\mathcal U\times P$ (the set on which the averages entering the conditions are defined; $\mathcal W=\mathcal U\times P$ is allowed), $\mathfrak c$ and $\mathfrak d$ continuous maps on $\mathcal W$ into metric spaces, and $S$ a continuous real function on $\mathcal W$ (the action, which may involve the parameter through boundary plaquettes). For $p\in P$ let $\mathcal W(p)=\{ U:(U,p)\in\mathcal W\}$, an open subset of $\mathcal U$, let $\mathcal K(p)=\{ U\in\mathcal W(p):\mathfrak c(U,p)=0\}$, let $\mathcal V$ be the set of $p$ with $\mathcal K(p)\neq\varnothing$ on which $S(\bigcdot ,p)$ attains its infimum $m(p)$ over $\mathcal K(p)$, and let
\begin{equation}
\mathcal M_{\mathrm{ax}}=\{(p,U): p\in\mathcal V,\ U\in\mathcal K(p),\ S(U,p)= m(p),\ \mathfrak d(U,p)=0,\ U(y,x)=\id\text{ for every centre }y\text{ of a block of }\mathbf B\text{ whose bonds lie in }E\text{ and }x\in\mathcal B_y\setminus\{ y\}\}
\end{equation}
be the set of minimal configurations satisfying the further equalities $\mathfrak d=0$, in the axial gauge (p.248), with projection $\mathcal V_{\mathrm{ax}}$ to the first factor. Then $\mathcal V$, $m$, $\mathcal M_{\mathrm{ax}}$ and $\mathcal V_{\mathrm{ax}}$ are Borel, and there is a Borel map $\mathcal V_{\mathrm{ax}}\ni p\mapsto U_{\mathbf B}(p)$ with $(p,U_{\mathbf B}(p))\in\mathcal M_{\mathrm{ax}}$ for every $p\in\mathcal V_{\mathrm{ax}}$. For the problem (1.2), $U_{1,\square'}(V)= U(\mathbf B_1(\square'^{\sim4}),M^{\bigcdot }(Q_1^{s*}V))$, the parameter is $p= V|_{\square'^{\sim4}}$ (p.246), $\mathcal W=\mathcal U\times P$, $\mathfrak c(U,V)= M_{\mathbf B}(U)\,(M^{\bigcdot }(Q_1^{s*}V))^{-1}-\id$, so that $\mathcal K(V)=\{ U\in\mathcal U_{\mathbf B}:M_{\mathbf B}(U)= M^{\bigcdot }(Q_1^{s*}V)\}$ is the constraint set of the variational problem CRE (2.12) for the data of (1.2), and $\mathfrak d\equiv0$; for a regular $V$, CRE p.256 gives the unique minimal orbit of (2.12), so $V\in\mathcal V$, and this orbit has a regular representative in the axial gauge, so that $V\in\mathcal V_{\mathrm{ax}}$ (the proof below, from the gauge invariance stated at VAR p.278), which is the representative CRE chooses, ``taken in the axial gauge'' (p.247). For the configuration of (1.21), $p=(V,V_0)$, $E$ the bonds of $\Omega_1\cap\Lambda_1^c$ (CRE's ``function $U_1(\Omega_1\cap\Lambda_1^c,V)$'', p.251), $\mathcal W$ the relatively open set of the $(U,p)$ at which the averages of the rows crossing to $\Omega_1^c$, which involve $V_0$ on the outer block, are defined, $\mathfrak c$ collects the constraints $\bar U= V$ on the rows of $(\Lambda_1^{\sim})^c$ (with $V_0$ on the blocks off $\Omega_1$) and $M_{\mathbf B_1(\Lambda_1^c)}(U)= M^{\bigcdot }(Q_1^{s*}V)$ on the coarse bonds of the determining set (the form of the constraint of (1.2); CRE names the determining data, ``the set $\mathbf B_1(\Lambda_1^c)$ and the configuration $Q_1^{s*}V$'', p.251), $\mathfrak d\equiv0$, and $S(U,p)$ is the functional $A^{\eta}$ of (2.12) for the determining set $\mathbf B_1(\Lambda_1^c)$, in which the plaquettes meeting $\Omega_1^c$ involve $V_0$, in the reading of p.251 recorded in (B1): the minimality $S(U)= m(p)$ over $\mathcal K(p)$ is the condition of (1.12) on $(\Lambda_1^{\sim})^c$, the critical point of (1.12) being its minimum (p.248), and no critical-point equation is imposed separately (the constrained critical-point equation written with $\mathsf C_U$ presupposes the invertibility of the $\alpha_c(U)$ and is a consequence of the minimality where it is defined, not a condition of the problem); only the form of these conditions is used, each being an equality between continuous functions of $(U,p)$ or the minimality of a continuous functional over a set defined by such equalities, and the existence of the configuration CRE describes, for every admissible pair $(\Omega_1,V,V_0)$ (Section~\ref{sec:branch-data}), is CRE's application of Theorem~1 of \cite{BalabanVAR} to the determining set $\mathbf B_1(\Lambda_1^c)$: CRE introduces the configuration, p.251, for the terms of its history sum, whose domain of integration is $D^{(11)}(V,V_0)$, and uses it wherever (1.15) is taken, and p.251 states its determining conditions and no existence theorem for each pair; its existence is the existence of the minimal orbit of VAR Theorem~1 for these determining data, the application CRE makes of VAR to all its determining sets (p.256, as recorded in Section~\ref{sec:c3-9-1}), taken over as CRE's. Theorem 1 of \cite{BalabanVAR} (CRE's [15]) is cited by CRE at p.248, for the critical point of (1.12).
\end{lemma}

\begin{proof}
The set $\mathcal K=\{(p,U):(U,p)\in\mathcal W,\ \mathfrak c(U,p)=0\}$ is Borel in $P\times G^{E}$, as the zero set of a continuous function on the Borel set $\mathcal W$ (relatively open in the Borel set $\mathcal U\times P$), and so is $\mathcal K_t=\{(p,U)\in\mathcal K: S(U,p)< t\}$ for every rational $t$. Every section $\mathcal K_t(p)=\{ U\in\mathcal W(p): \mathfrak c(U,p)=0,\ S(U,p)< t\}$ is the intersection of $\mathcal U$, a finite intersection of open and closed subsets of the compact metric space $G^{E}$, with the open set $\mathcal W(p)$, a closed set (the zero set of $\mathfrak c(\bigcdot ,p)$ is relatively closed in $\mathcal W(p)$, the trace of a closed set) and an open set (the trace of $\{ S(\bigcdot ,p)< t\}$, relatively open in $\mathcal W(p)$), hence $K_{\sigma}$: an open subset of a compact metric space is a countable union of compact sets, and the intersection of a $K_{\sigma}$ set with a closed set, or with an open set, is $K_{\sigma}$. By the Arsenin--Kunugui theorem (a Borel subset of a product of Polish spaces whose sections are $K_{\sigma}$ has a Borel projection and a Borel uniformisation, \cite[Theorem 35.46(ii)]{Kechris}), $\operatorname{proj}\mathcal K_t=\{ p: \inf_{\mathcal K(p)}S< t\}$ is Borel for every rational $t$, and so is $\operatorname{proj}\mathcal K=\{ p:\mathcal K(p)\neq\varnothing\}$; hence $\widetilde m(p)=\inf_{\mathcal K(p)}S$ is a Borel function on $\operatorname{proj}\mathcal K$. The set $\mathcal M=\{(p,U)\in\mathcal K: S(U,p)=\widetilde m(p)\}$ is Borel, its sections $\mathcal K(p)\cap\{ S=\widetilde m(p)\}$ are $K_{\sigma}$, so $\mathcal V=\operatorname{proj}\mathcal M$ is Borel by the same theorem, and $m=\widetilde m$ on $\mathcal V$. The axial gauge conditions $U(y,x)=\id$ are relatively closed conditions on the open set on which the means $U(y,x)$ of RG1 (0.11) are defined (RG1 p.253 defines the mean for sets of sufficiently small diameter, and it is continuous there in the chart convention), an open set containing $\mathcal U$ by hypothesis; and $\mathfrak d(U,p)=0$ is a relatively closed condition on $\mathcal W$. So $\mathcal M_{\mathrm{ax}}=\mathcal M\cap\{\mathfrak d=0\}\cap\{ U\text{ axial}\}$ is Borel, and its sections are $K_{\sigma}$, being intersections of $K_{\sigma}$ subsets of $\mathcal U$ with closed sets (the trace on $\mathcal U$ of a relatively closed subset of an open set containing $\mathcal U$ is the trace of a closed set); the theorem gives that $\mathcal V_{\mathrm{ax}}$ is Borel and that $\mathcal M_{\mathrm{ax}}$ has a Borel uniformisation, a Borel map $U_{\mathbf B}$ on $\mathcal V_{\mathrm{ax}}$ with graph in $\mathcal M_{\mathrm{ax}}$. The statements about CRE's problems in the lemma are CRE's: a regular $V$ determines the unique minimal orbit of (2.12), whose elements are the minimal configurations (p.256), and the configurations $U_{1,\square'}$ are taken in the axial gauge (p.247); a gauge transformation equal to $\id$ at the block centres and to the mean $U(y,x)$ elsewhere brings a minimal configuration to the axial gauge, leaving the block averages of RG1 (0.4) at the block centres unchanged (RG1 (0.4) transforms at the block centres, proof of Lemma~\ref{lem:two-block}) and $S(U)$ unchanged when $S$ is gauge invariant, as the action is; the configuration so obtained is regular, satisfies the constraint and is minimal: VAR states that its class (2) is gauge invariant and its constraint set (3) invariant under the gauge transformations equal to one at the coarse sites (p.278), which the transformation used is, being the identity at the block centres, and that the space (6) and the action (5) are invariant under these (p.278); so $\mathcal M_{\mathrm{ax}}(V)\neq\varnothing$; the description of the configuration of (1.21) is p.251, taken as stated.
\end{proof}

\subsection{The constrained Hessian at a regular background and its stencils}
\label{sec:regular-hessian-stencils}

Let $\mathsf M$ be Ba{\l}aban's averaging operation \eqref{eq:nonlinear-average} (written $M$ there; here $M$ is the cube side), and let $\Omega$ be a region, a union of blocks, and $U$ a configuration of the fine torus for which $\mathsf M(U)(c)$ is defined on every row $c\in\Omega^{[1]}$ of that region and the local group means $\mathcal U_U(y,x)$, $x\in\mathcal B_y$, are defined for every centre $y\in\Omega^{(1)}$ (chart convention of the preamble; for $U= U_1$ on $\Omega_1$ both hold, because $U_1\in D^{(11)}$ and $U_1$ is in its axial mean gauge, $\mathcal U_{U_1}(y,x)=\id$); the objects below are then defined for phases $q$ in a neighbourhood of zero depending on $U$, which Lemma~\ref{lem:two-block} and Appendix~\ref{app:background-estimates} quantify for the backgrounds that occur. The objects of this subsection are defined relative to the region $\Omega$: the rows are the coarse bonds $c\in\Omega^{[1]}$ (at least one block in $\Omega$), the block sums run over the centres $y\in\Omega^{(1)}$, the free carrier is $\Omega^*$, and the links of $U$ on bonds with both endpoints off $\Omega$ enter as fixed constants, with the phase variable equal to zero there. The torus itself ($\Omega= T$) is written without a superscript; the branch uses $\Omega=\Omega_1$ and $U= U_1$. Define, for $c\in\Omega^{[1]}$,
\begin{equation}\label{eq:native-source-constraint-normalisation}
    \mathcal Q_{U,c}(q)=\frac{1}{i}\log\left(\mathsf M(e^{iq}U)(c)\,\mathsf M(U)(c)^{-1}\right),\qquad
    Q_U=\frac{1}{L}D\mathcal Q_U(0),\qquad
    \mathcal N_U(q)=\mathcal Q_U(q)-LQ_Uq\;,
\end{equation}
so that $\mathcal N_U$ is the remainder denoted $\widetilde C$ in the linearisation following \cite[(2.10)]{BalabanRG1}; the logarithm is the group chart of the conventions above, defined when the relative average lies in the chart domain. For a coarse bond $c$ with designated fine bond $b_0(c)$, and a coarse centre $y$, put
\begin{equation}\label{eq:regular-stencil-bonds}
    \mathcal E_c=\{ r\, | \,\{ r_-,r_+\}\subset\mathcal B_{c_-}\cup\mathcal B_{c_+}\},
    \qquad
    \mathcal E_y=\{ r\, | \,\{ r_-,r_+\}\subset\mathcal B_y\}\;,
\end{equation}
$r$ ranging over the positive fine bonds.

\begin{lemma}[two-block structure]\label{lem:two-block}
Let $L\geqslant 3$ be odd and let every coarse period exceed two. The row components $\mathcal Q_{U,c}$ and $\mathcal N_{U,c}$ depend only on the phases and background bonds in $\mathcal E_c$, and $b_0(c)$ is the only designated bond in $\mathcal E_c$. Hence
\begin{equation}\label{eq:regular-stencil-designated-diagonal}
    Q_{0,U}(c,c')=\delta_{cc'}\alpha_c(U),\qquad \alpha_c(U)= Q_U(c,b_0(c))\;,
\end{equation}
where $Q_{0,U}$ is the restriction of $Q_U$ to the designated bonds, with $\alpha_c(\id)= L^{-4}\id$. There is a phase radius $\tau_L>0$, depending on $L$ and the group only, such that for $U= e^{iA}$ with $\| A\|_{\infty}\leqslant \tau_L$ every $\alpha_c(U)$ is analytic in $A$ with smallest singular value at least $L^{-4}/2$. Under a site gauge transformation $u$, $\alpha_c(U^u)=\operatorname{Ad}_{u(c_-)}\alpha_c(U)\operatorname{Ad}_{u(b_0(c)_-)}^{-1}$.
\end{lemma}

\begin{proof}
Every path in the RG1 (0.4) row for $c$ lies in the two endpoint blocks, since the row averages parallel transports from the centre of $\mathcal B_{c_-}$ to the centre of $\mathcal B_{c_+}$ along paths inside the two blocks; a designated bond crosses its own coarse interface, and two adjacent blocks share exactly the interface labelled by $c$, so $b_0(c)$ is the only designated bond with both endpoints in $\mathcal B_{c_-}\cup\mathcal B_{c_+}$ (for $L\geqslant 3$ odd the other designated bonds cross other interfaces). At the flat field, differentiating RG1 (0.4) gives \eqref{eq:covariance-two-averages} of Appendix~\ref{app:flat-covariance}: the straight-path term $(Q_1a)(c)= L^{-5}\sum_{x\in\mathcal B_{c_-}}a([x,x'])$ contains $b_0(c)$ in exactly $L$ of its $L^4$ translated straight paths of $L$ bonds each (the $L$ starting points on the line through $b_0(c)$), each with coefficient $L^{-5}$, and the mean-path term involves the paths $\Phi$ inside the blocks, which contain no designated bond; this gives $\alpha_c(\id)= L^{-4}\id$, the designated coefficient $Q_{0,\infty}= L^{-4}I$ computed in Appendix~\ref{app:flat-covariance} (paragraph ``Construction on the infinite lattice'', after \eqref{eq:bulk-operator-definition}; the computation is local and holds on every torus). For $U= e^{iA}$, $\alpha_c(U)$ is an analytic function of $A|_{\mathcal E_c}$ near zero (a composition of exponentials, the averaging and the chart logarithm) with $\alpha_c(\id)= L^{-4}\id$; the smallest singular value is continuous, so for each of the finitely many row types modulo $L\mathbb Z^4$ (translations commute with $\mathsf M$) there is a radius on which $\alpha_c$ is analytic with smallest singular value at least $L^{-4}/2$, and $\tau_L$ is the smallest of these radii; it depends on $L$ and the group only. The gauge covariance is $\mathsf M(U^u)(c)= u(c_-)\mathsf M(U)(c)u(c_+)^{-1}$ (RG1 (0.4) transforms at the block centres) together with $(R_uq)(b)=\operatorname{Ad}_{u(b_-)}q(b)$ for the phase (CRE (1.17)): $\mathcal Q_{U^u,c}(R_uq)=\operatorname{Ad}_{u(c_-)}\mathcal Q_{U,c}(q)$, hence $Q_{U^u}R_u=\operatorname{Ad}_{u(c_-)}Q_U$ row by row, and the $b_0(c)$ column gives the formula.
\end{proof}

By (B2) and Lemma~\ref{lem:two-block}, $\alpha_c(U)$ is invertible on every row of the regions used below. With $LQ_Uh_U=\id$ for the supported right inverse $h_U$, whose only nonzero outputs are at the designated bonds, the elimination maps are
\begin{equation}\label{eq:regular-stencil-elimination}
    (\mathsf C_UB)(b_0(c))
        =-\alpha_c(U)^{-1}\sum_{b\in\mathcal E_c\cap E_*}Q_U(c,b)B(b),
    \qquad
    (h_Uv)(b_0(c))= h_U(c)v(c),\qquad h_U(c)=\frac{1}{L}\alpha_c(U)^{-1}\;,
\end{equation}
and $(\mathsf C_UB)(r)= B(r)$ on free bonds $r\in E_*$; thus $Q_U\mathsf C_U=0$. The free-input and background stencils of a fine bond are
\begin{equation}\label{eq:regular-stencil-extension-support}
    (\mathcal F_r,\mathcal A_r)=
    \begin{cases}
        (\{ r\},\varnothing),&r\in E_*\;,\\
        (\mathcal E_c\cap E_*,\mathcal E_c),&r= b_0(c)\;.
    \end{cases}
\end{equation}

The constrained quadratic form is CRE (1.14)--(1.16) on the region $\Omega$. With $J_U$ and $\Delta_U$ the first derivative and Hessian at $q=0$ of $q\mapsto S(e^{iq}U)$, $q\in\mathfrak g^{\Omega}$ (the plaquettes with a bond in $\Omega$ contribute; links off $\Omega$ are the fixed constants), and the relative gauge functional
\begin{equation}\label{eq:regular-covariance-gauge}
    \mathscr G^{\Omega}_U(q)=\sum_{y\in\Omega}\sum_{x\in\mathcal B_y\setminus\{ y\}}
        \left(1-\operatorname{Re}\tr_N(\mathcal U_{e^{iq}U}(y,x)\mathcal U_U(y,x)^{-1})\right),
    \qquad G^{(2),\Omega}_U(a)= D^2\mathscr G^{\Omega}_U(0)[a,a]\;,
\end{equation}
where $\mathcal U_U(y,x)$ is the axial transport from the block centre $y$ to $x$ of CRE's gauge-fixing functional $\mathbf G$ in (1.11) and the sum is over the block centres of $\Omega$, the regional Hessian $\mathsf H^{\Omega}_U$ on $\mathfrak g^{\Omega^*}$ is the symmetric operator of the quadratic form ($B$ is extended by zero to the free bonds outside $\Omega^*$, so that $a=\mathsf C_UB$ is defined on every row of $\Omega^{[1]}$)
\begin{equation}\label{eq:regular-covariance-hessian}
    \langle B,\mathsf H^{\Omega}_UB\rangle
        =\langle a,\Delta_Ua\rangle+G^{(2),\Omega}_U(a)
            -\langle J_U,h_UD^2\mathcal N_U(0)[a,a]\rangle,
    \qquad a=\mathsf C_UB\;,
\end{equation}
with the rows of $\mathcal N_U$ and $h_U$ those of $\Omega^{[1]}$. The right side equals $D^2(S+\mathscr G^{\Omega}_U)(0)[a,a]-DS(0)[h_UD^2\mathcal N_U(0)[a,a]]$, which does not involve the inner product; $\mathsf H^{\Omega}_U$ is its operator for the fixed one, and its entries $\mathsf H^{\Omega}_U(b,b')\in\operatorname{End}(\mathfrak g)$ are those of the polarised bilinear form. This is CRE (1.16), $\langle A,\Delta_1A\rangle=\langle A,\Delta A\rangle-2\langle hC^{(2)}(A),J_1\rangle+\mathbf G^{(2)}(A)$, with CRE's $C^{(2)}(A)=\frac12D^2\mathcal N_U(0)[A,A]$ (the quadratic term of the linearising transformation, RG1 p.267, $\widetilde D^{(2)}=\widetilde C^{(2)}$) and $\mathbf G^{(2)}(A)= D^2\mathscr G(0)[A,A]$; at $\Omega=\Omega_1$, $U= U_1$ the relative gauge functional is CRE's absolute one, $\mathscr G^{\Omega_1}_{U_1}(q)=\mathbf G(\Omega_1^{(1)},e^{iq}U_1)$ with $\mathbf G$ of (1.7), (1.11), because $U_1$ is in its axial mean gauge, $\mathcal U_{U_1}(y,x)=\id$ for $x\in\mathcal B_y$, $y\in\Omega_1^{(1)}$ (Section~\ref{sec:branch-data}; CRE p.248), so that $\mathbf G(\Omega_1^{(1)},U_1)=0$; and $\mathsf H^{\Omega_1}_{U_1}$ is CRE's $C^*\Delta_1C$ of (1.15), the same quadratic form written as an operator for the fixed inner product. The precision of the branch is
\begin{equation}\label{eq:c-regional-hessian}
    H_{\xi}=\mathsf H^{\Omega_1}_{U_1}\quad\text{on }\mathfrak g^{\Omega_1^*},\qquad H_{\xi,II}=\iota_I^*H_{\xi}\iota_I,\qquad H_{\xi,IJ}=\iota_I^*H_{\xi}\iota_J\;,
\end{equation}
with $\iota_I,\iota_J$ the extensions by zero. With $s_{p,U}(q)=1-\operatorname{Re}\tr_N((e^{iq}U)(\partial p))$, $\mathscr G_{y,U}$ the $y$-summand of \eqref{eq:regular-covariance-gauge}, and $a=\mathsf C_UB$, the form decomposes into local terms,
\begin{equation}\label{eq:regular-stencil-quadratic-decomposition}
    \langle B,\mathsf H^{\Omega}_UB\rangle=\sum_{p\cap\Omega\ne\varnothing}\mathfrak a_p(B;U)+\sum_{y\in\Omega}\mathfrak g_y(B;U)-\sum_{c\in\Omega^{[1]}}\mathfrak t_c(B;U),\qquad
    \begin{aligned}
        \mathfrak a_p&= D^2s_{p,U}(0)[a,a],\\
        \mathfrak g_y&= D^2\mathscr G_{y,U}(0)[a,a],\\
        \mathfrak t_c&=\tfrac1L\langle J_U(b_0(c)),\alpha_c(U)^{-1}D^2\mathcal N_{U,c}(0)[a,a]\rangle\;.
    \end{aligned}
\end{equation}
Their free-input and background sets are read off from Lemma~\ref{lem:two-block} and \eqref{eq:regular-stencil-elimination}: $a(r)$ depends on $B|_{\mathcal F_r}$ and on the background on $\mathcal A_r$, so $\mathfrak a_p$ has free-input set $\bigcup_{r\subset\partial p}\mathcal F_r$ and background set $\bigcup_{r\subset\partial p}(\{ r\}\cup\mathcal A_r)$; $\mathfrak g_y$ involves the phases on the axial paths inside $\mathcal B_y$, which contain no designated bond (designated bonds cross interfaces), so its free-input and background sets are $\mathcal E_y\cap E_*$ and $\mathcal E_y$; $\mathfrak t_c$ involves $\mathcal N_{U,c}$, $\alpha_c$ and $J_U(b_0(c))$, the latter depending on the plaquettes through $b_0(c)$, which lie in $\mathcal B_{c_-}\cup\mathcal B_{c_+}$ for $L\geqslant 3$, so its sets are $\mathcal E_c\cap E_*$ and $\mathcal E_c$. The entries of $\mathsf H^{\Omega}_U$ are the sums of the polarised entries of these terms. For a pair of positive bonds $(b,b')$ let $\mathcal T_{bb'}$ be the set of terms of \eqref{eq:regular-stencil-quadratic-decomposition} (plaquettes, centres, rows of the torus) whose free-input set contains both $b$ and $b'$, and let the \emph{entry stencil} $\mathcal S_{bb'}$ be the union of the background sets of these terms (one of finitely many types modulo $L\mathbb Z^4$; it is empty unless $d(b,b')\leqslant 4L$, Lemma~\ref{lem:hessian-gauge-locality}). The entry function is
\begin{equation}\label{eq:c-entry-function}
    f_{bb'}(U)=\sum_{\mathfrak T\in\mathcal T_{bb'}}(\text{polarised entry of }\mathfrak T\text{ at }(b,b'))\in\operatorname{End}(\mathfrak g),\qquad f_{bb'}(A)\coloneqq f_{bb'}(e^{iA})\;,
\end{equation}
defined for every configuration $U$ on $\mathcal S_{bb'}$ for which $\mathsf M(U)$ is defined and $\alpha_c(U)$ is invertible on the rows of the terms in $\mathcal T_{bb'}$ and the local group means $\mathcal U_U(y,x)$ are defined for the centres $y$ of the terms in $\mathcal T_{bb'}$ (the terms involve $\mathsf C_U$ on those rows and the means of those centres), whatever the configuration elsewhere; the torus form itself is not needed for its definition. The regional form is the compression of the torus form: for a background $U$ on which $\mathsf H^{\Omega}_U$ is defined (Section~\ref{sec:regular-hessian-stencils}: $\mathsf M(U)$ defined and $\alpha_c(U)$ invertible on the rows of $\Omega^{[1]}$, the means $\mathcal U_U(y,x)$ defined for $y\in\Omega^{(1)}$) and every pair $(b,b')\in\Omega^*\times\Omega^*$,
\begin{equation}\label{eq:c-regional-compression}
    \mathsf H^{\Omega}_U(b,b')= f_{bb'}(U|_{\mathcal S_{bb'}})\;,
\end{equation}
because a term whose free-input set contains the bond $b$ of $\Omega^*$ has its index in $\Omega$ (a plaquette containing $b$, or containing a bond $r$ with $b\in\mathcal F_r$, meets $\Omega$; a centre $y$ with $b\in\mathcal E_y$ lies in $\Omega$; a row $c$ with $b\in\mathcal E_c\cap E_*$ has an endpoint block containing an endpoint of $b$, so $c\in\Omega^{[1]}$), hence every term of $\mathcal T_{bb'}$ is a term of the regional sum, and the regional form is defined under the weaker requirement on $U$ stated above (Lemma~\ref{lem:c2-regional-entry} records this in the appendix).

\begin{lemma}[gauge law and locality of the Hessian]\label{lem:hessian-gauge-locality}
Let $u$ be a site gauge transformation of the torus, acting on the background by $U\mapsto U^u$ and on phases by $(R_uq)(b)=\operatorname{Ad}_{u(b_-)}q(b)$ (CRE (1.17)), and let $R_{u,*}$ be the restriction of $R_u$ to the free carrier $\Omega^*$; it is orthogonal for the invariant inner product. Then, for every region $\Omega$,
\begin{equation}\label{eq:regular-gauge-constrained-hessian}
    \mathsf H^{\Omega}_{U^u}= R_{u,*}\mathsf H^{\Omega}_UR_{u,*}^{-1}\;.
\end{equation}
Every entry $\mathsf H^{\Omega}_U(b,b')$ is a finite composition of local analytic operations in the background bonds of the entry stencil $\mathcal S_{bb'}$, which lies within the range $\rho_L=4L$ of $b$, including the fibre inverses $\alpha_c^{-1}$, and it vanishes for $d(b,b')>\rho_L$ (the stencils are listed in Appendix~\ref{app:background-estimates}, Lemma~\ref{lem:c2-entry-range}). The entry function $f_{bb'}$ of \eqref{eq:c-entry-function} is entrywise gauge covariant: for a site gauge transformation $u$,
\begin{equation}\label{eq:c-entry-gauge-covariance}
    f_{bb'}(U^u)=\operatorname{Ad}_{u(b_-)}f_{bb'}(U)\operatorname{Ad}_{u(b'_-)}^{-1}\;.
\end{equation}
The norms are the operator norm for the Hilbert norms induced by $h$ and, for $Df_{bb'}(A)$, the norm $\sup_{\| X\|_{\infty}\leqslant 1}\| Df_{bb'}(A)[X]\|$. Let $\tau_L$ denote the radius of Lemma~\ref{lem:two-block}, reduced so that these operations are analytic on $\| A\|_{\infty}\leqslant \tau_L$ and reduced once more in \eqref{eq:c2-tau-final} of Appendix~\ref{app:background-estimates}; every statement is made for this final value. Define
\begin{equation}\label{eq:regular-covariance-entry-constants}
    m_L=\max_b\sum_{b'}\sup_{\| A\|_{\infty}\leqslant \tau_L}\| f_{bb'}(A)\|,
    \qquad
    c_L=\max_b\sum_{b'}\sup_{\| A\|_{\infty}\leqslant \tau_L}\| Df_{bb'}(A)\|\;.
\end{equation}
These depend on $L$, the group and the invariant inner product only.
\end{lemma}

\begin{proof}
The intertwinings are $\mathsf M(U^u)(c)= u(c_-)\mathsf M(U)(c)u(c_+)^{-1}$, hence $\mathcal Q_{U^u}\circ R_u= R^{(1)}_u\mathcal Q_U$ with $(R^{(1)}_uv)(c)=\operatorname{Ad}_{u(c_-)}v(c)$ (proof of Lemma~\ref{lem:two-block}), and therefore $Q_{U^u}R_u= R^{(1)}_uQ_U$, $\mathcal N_{U^u}\circ R_u= R^{(1)}_u\mathcal N_U$, $h_{U^u}= R_uh_U(R^{(1)}_u)^{-1}$ (from the covariance of $\alpha_c$ in Lemma~\ref{lem:two-block}) and $\mathsf C_{U^u}= R_u\mathsf C_UR_{u,*}^{-1}$. The three terms of \eqref{eq:regular-covariance-hessian} are then each preserved: $S((e^{iR_uq})U^u)= S((e^{iq}U)^u)= S(e^{iq}U)$ gives $\Delta_{U^u}= R_u\Delta_UR_u^{-1}$ and $J_{U^u}= R_uJ_U$; each relative mean $\mathcal U_{e^{iq}U}(y,x)\mathcal U_U(y,x)^{-1}$ is conjugated by $u(y)$, which preserves the trace; and $\langle J_{U^u},h_{U^u}D^2\mathcal N_{U^u}(0)[R_ua,R_ua]\rangle=\langle R_uJ_U,R_uh_UD^2\mathcal N_U(0)[a,a]\rangle$ equals the untransformed force term because $R_u$ is orthogonal for the invariant inner product. The same three identities hold term by term in \eqref{eq:regular-stencil-quadratic-decomposition} (each term is a function of $U$ and $a=\mathsf C_UB$ on its own stencil, and $\mathsf C_{U^u}R_{u,*}B= R_u\mathsf C_UB$), so the polarised entry of each term transforms as $\operatorname{Ad}_{u(b_-)}(\bigcdot )\operatorname{Ad}_{u(b'_-)}^{-1}$, which is \eqref{eq:c-entry-gauge-covariance}. Locality: the index of a term is the lower vertex $x(p)$ of its plaquette, the centre $y$ or the initial vertex of $b_0(c)$, and the distance of a bond to the index is that of its initial vertex. By the two-block structure of Lemma~\ref{lem:two-block} substituted into \eqref{eq:regular-stencil-quadratic-decomposition}, the free-input and background bonds of a block or row term have initial vertices in the one or two blocks of the index, within $\ell^{\infty}$ distance $2L-1$ of it (the sites of two adjacent blocks are within $2L-1$ of each other); a bond $r$ of $\partial p$ has initial vertex within $1$ of $x(p)$, and the bonds of $\mathcal F_r\cup\mathcal A_r$ are within $2L-1$ of $r$, so the bonds of a plaquette term are within $2L$ of its index. Hence a term whose free-input set contains $b$ and $b'$ has $d(b,b')\leqslant 4L=\rho_L$, and its background bonds are within $4L$ of $b$. The stencils are listed in Lemma~\ref{lem:c2-entry-range}.
\end{proof}

\begin{lemma}[regular-background bounds]\label{lem:regular-bounds}
Let $\gamma_L= C_L^{-1}$ be the flat lower bound \eqref{eq:covariance-uniform-gap} of Appendix~\ref{app:flat-covariance}, $\tau_L$ the final radius of Lemma~\ref{lem:hessian-gauge-locality}, and let $R_{\mathrm c}$ (the collar range, Lemma~\ref{lem:collar-coverage}) and $C_{L,R_{\mathrm c}}$ be the constants fixed in Appendix~\ref{app:background-estimates}, Section~\ref{app:regular-covariance}; they depend on $L$ and the group only. Whenever \eqref{eq:c-standing-conditions} holds and $\widehat\mu= C_{L,R_{\mathrm c}}\varepsilon_1$, $\varepsilon_1= g_1p_0(g_1)$ (CRE p.246, with $g_1$ the first-step coupling), satisfies $\widehat\mu<\tau_L/2$ and $\max(1,c_L)\widehat\mu<\gamma_L/4$, on every branch (where $H_{\xi}$ is defined, by (B1) and (B2), as the operator of a quadratic form at fixed volume)
\begin{equation}\label{eq:c-regular-two-sided}
    \frac{\gamma_L}{2}\,\id\leqslant H_{\xi,II}\leqslant m_L\,\id\;,
\end{equation}
and $H_{\xi,II}(b,b')=0$ for $d(b,b')>\rho_L$ (the entry stencil $\mathcal S_{bb'}$ is empty beyond that distance, Lemma~\ref{lem:c2-entry-range}). The same bounds hold for every principal restriction of $H_{\xi,II}$. Moreover, for every row $b\in I$ the row sum over all free bonds, those outside $I$ included, obeys $\sum_{b'\in\Omega_1^*}\|(H_{\xi})_{bb'}\|\leqslant m_L$, with $\|\bigcdot \|$ the operator norm for $h$ and $m_L$ of \eqref{eq:regular-covariance-entry-constants}.
\end{lemma}

The proof is Appendix~\ref{app:background-estimates}, Section~\ref{app:regular-covariance}: for $b,b'\in I$ the entries $(H_{\xi})_{bb'}$ are $f_{bb'}(U_1|_{\mathcal S_{bb'}})$ by \eqref{eq:c-regional-compression} (Lemma~\ref{lem:c2-regional-entry}); a partition of unity on the coarse grid, a local gauge on each patch bringing $U_1$ to $e^{iA}$ with $\| A\|_{\infty}\leqslant \widehat\mu$ (Lemmas~\ref{lem:c2-plaquette-deviation} and~\ref{lem:c2-local-gauge}), and the entrywise fundamental theorem of calculus compare the form with the flat one.

\subsection{The source phase and the local primitives}
\label{sec:source-phase-primitives}

On the branch, put $U= U_1$ and $\Omega=\Omega_1$. The definitions and Lemma~\ref{lem:row-locality} of this subsection use only (B2), (B4), Lemma~\ref{lem:charts} (for the row solutions on the near rows), Lemmas~\ref{lem:c2-uniform-chart} and~\ref{lem:c2-fibre-smallness} of Appendix~\ref{app:background-estimates} (for the row solutions on the far rows and the agreement of the definitions), Section~\ref{sec:regular-hessian-stencils}, the map $\Psi_{\xi}$ of \eqref{eq:native-source-exact-boundary-map} with \eqref{eq:c-bondwise-shift} and \eqref{eq:c-collar-distance}, and Lemma~\ref{lem:collar-coverage}(i),(iii). They are stated for the background $U= U_1$ of the branch. For the reference background $U=\id$ of Section~\ref{sec:joint-extraction} only the near rows and the germs occur (the reference inputs vanish off $I$), the row solutions are those of Lemma~\ref{lem:charts} in its reference form (Lemma~\ref{lem:c2-reference}), (B4) is not invoked, and $v$, $V$ below are defined by the same formulas. Lemma~\ref{lem:primitive-sum} uses in addition that $U$ is a critical point of $S$ under the block-average constraint and is in the axial mean gauge, which holds for $U_1$ by (B1) and Section~\ref{sec:branch-data} and for $U=\id$ because the flat configuration minimises $S$ and $\mathscr G^{\Omega_1}_{\id}$ outright and has $\mathcal U_{\id}(y,x)=\id$. For a free vector $B\in\mathfrak g^{\Omega_1^*}$ put $A= g_0\mathsf C_UB\in\mathfrak g^{\Omega_1}$. The first-step correction $\widetilde D_U$ (for $U= U_1$, CRE's correction of (1.13)--(1.14) under the reading of (B4)(i)) is the solution, selected below, of the system
\begin{equation}\label{eq:native-source-local-implicit}
    \widetilde D_{U,c}(A)=\mathcal N_{U,c}\bigl(A-h_U\widetilde D_U(A)\bigr),\qquad c\in\Omega_1^{[1]}\;,
\end{equation}
the form of the equation of RG1 p.267 for RG1's own step ($L\widetilde QB'+\widetilde C(B')= L\widetilde QB$ with $B'= B-h\widetilde D(B)$, $\widetilde C=\mathcal N_U$, $\widetilde Q= Q_U$) written row by row. Let $\rho_{\mathrm d}>0$ be the radius $2K_*\vartheta_*^2$ of Appendix~\ref{app:background-estimates}, Lemma~\ref{lem:c2-uniform-chart}(ii); it depends on $L$ and the group only. A solution $d_c$ of the row equation of $c$ with $h(d_c)<\rho_{\mathrm d}$ is called \emph{small}. The row solution $\widetilde D_{U,c}(A)$ is defined as follows: for $U= U_1$ on the rows not within $2R_{\mathrm c}+2L$ of $I$ (the far rows; a row is within $\rho$ of a bond set when every bond of its two-block set is, see below) it is the row solution $\widetilde D_c(A)=\varphi_c(A|_{\mathcal E_c})$ of (B4)(i), defined on the ball $\| A|_{\mathcal E_c}\|<\vartheta_*$ (Lemma~\ref{lem:c2-fibre-smallness}(i)); the arguments $A= g_0\mathsf C_{U_1}B$, $B\in\mathcal A_{\xi}$, of the domain of integration \eqref{eq:c-free-variable-domain} lie in the ball of radius $\vartheta_*/2$, $\| A|_{\mathcal E_c}\|<\vartheta_*/2$ (Lemma~\ref{lem:c2-fibre-smallness}(ii)), and there its value is the component at $c$ of the chart-class family of Section~\ref{sec:branch-data} (Lemma~\ref{lem:c2-fibre-smallness}(iii)); on the rows within $2R_{\mathrm c}+2L$ of $I$ (the near rows), for $U= U_1$ and for the reference background, it is the unique small solution of the row equation of $c$ at the argument $A|_{\mathcal E_c}$ given by Lemma~\ref{lem:charts} on the row polydisc $\mathcal U_c$; and on a fixed neighbourhood $\mathcal N_c$ of $A|_{\mathcal E_c}=0$, for every row, it is the germ of Lemma~\ref{lem:row-locality}, where $\mathcal N_c$ is the neighbourhood given by that lemma, reduced so that it lies in the ball $\| A|_{\mathcal E_c}\|<\vartheta_*$ and the germ takes values with $h<\rho_{\mathrm d}$ (both hold near zero, where the germ vanishes, and both conditions are open; the phase $A-h_U(c)d_c$ then has norm below $\vartheta_*+H_*\rho_{\mathrm d}<2\vartheta_*\leqslant \varrho_G$ and is chart-valued). Where two of these definitions apply they agree: each produces a solution $d_c$ of the row equation with $h(d_c)\leqslant \rho_{\mathrm d}$ at an argument in the ball $\| A|_{\mathcal E_c}\|<\vartheta_*$ (the far-row values by Lemma~\ref{lem:c2-fibre-smallness}(i), the near-row values by Lemma~\ref{lem:charts}, the germ by the reduction of $\mathcal N_c$), and Lemma~\ref{lem:c2-uniform-chart}(ii), which applies at every row for $U= U_1$ (Lemma~\ref{lem:c2-fibre-smallness}(i)) and for the reference background $U=\id$ (zero phases), gives exactly one such solution. The reconstructed phase is the function of the free vector
\begin{equation}\label{eq:native-source-local-phase}
    q_U(B)= A-h_U\widetilde D_U(A)\in\mathfrak g^{\Omega_1},\qquad A= g_0\mathsf C_UB\;,
\end{equation}
defined on free and designated bonds of $\Omega_1$ wherever every row solution is defined. The integration variable of the branch is $a+A_{\mathrm e}\in\mathfrak g^{\Omega_1^*}$ in the variables of CRE (1.22); with $\Psi_{\xi}$ the bondwise map \eqref{eq:native-source-exact-boundary-map} of Section~\ref{sec:boundary-law-collar} from these variables to those of (1.15), the reconstructed fine field is
\begin{equation}\label{eq:c-reconstructed-field}
    U_{\xi}(a)(b)= e^{iq_U(\Psi_{\xi}(a+A_{\mathrm e}))(b)}U_1(b)\quad(b\in\Omega_1),\qquad U_{\xi}(a)(b)= V_0(b)\quad(b\in\Omega_1^c)\;,
\end{equation}
the fine field $U= U'U_1$ of CRE (1.13) at the actual branch data, $U_{\xi}(a)= U(\Psi_{\xi}(a+A_{\mathrm e}))$ in the notation of Section~\ref{sec:branch-data}; it is defined for $(A_{\mathrm e},a)\in\mathcal D_{\xi}$ by the definition of $\mathcal D_{\xi}$ and $\mathcal A_{\xi}$, and by Lemma~\ref{lem:charts} and Lemma~\ref{lem:row-locality} its links on the rows within $2R_{\mathrm c}+2L$ of $I$ are defined for every $a$ in the enlarged box defined there, the other links not depending on $a$. Write $U^0_{\xi}= U_{\xi}(0)$. Since $I\cap S_1^*=\varnothing$ (\eqref{eq:c-collar-distance} below) and $\Psi_{\xi}$ acts bond by bond and is the identity off $S_1^*$,
\begin{equation}\label{eq:c-bondwise-shift}
    \Psi_{\xi}(a+A_{\mathrm e})= a+\Psi_{\xi}(A_{\mathrm e})\qquad(a\in\mathfrak g^I,\ A_{\mathrm e}\in\mathfrak g^J)\;.
\end{equation}
A row $c$ is said to lie within $\rho$ of a bond set $S$ when every bond of $\mathcal E_c$ is within distance $\rho$ of $S$; as the bonds of $\mathcal E_c$ have their initial vertices in two adjacent blocks, this holds as soon as one bond of $\mathcal E_c$ is within $\rho-(2L-1)$ of $S$.

\begin{lemma}[rowwise locality of the reconstruction]\label{lem:row-locality}
The equation of row $c$ in \eqref{eq:native-source-local-implicit} involves $\widetilde D_{U,c}$ alone and depends only on $A|_{\mathcal E_c}$ and $U|_{\mathcal E_c}$; so the system decouples row by row, and a solution $d_c$ of the row equation at the argument $A|_{\mathcal E_c}$ is a function of $A|_{\mathcal E_c}$ and $U|_{\mathcal E_c}$. The row equation has a unique analytic germ of solutions at $(A|_{\mathcal E_c},d_c)=(0,0)$, with $d_c(0)=0$ and $Dd_c(0)=0$; it is the germ at zero of every solution continuous at zero with value zero, in particular, for $U= U_1$, of the row solution $\widetilde D_c=\varphi_c$ of (B4)(i). Consequently, wherever the row solutions are defined (Section~\ref{sec:source-phase-primitives}): on a free bond $r$, $q_U(B)(r)= g_0B(r)$; on $b_0(c)$, $q_U(B)(b_0(c))$ depends only on $B|_{\mathcal E_c\cap E_*}$ and on $U|_{\mathcal E_c}$; and $q_U(B)$ vanishes on all of $\mathcal E_c$ whenever $B=0$ on $\mathcal E_c\cap E_*$ and the row solution is the germ value there. In particular, by \eqref{eq:c-bondwise-shift}, the link $U_{\xi}(a)(r)$ depends on $a$ only through $a|_{\mathcal F_r\cap I}$ for $r\in\Omega_1$, and not at all for $r\in\Omega_1^c$.
\end{lemma}

\begin{proof}
$\mathcal E_c$ contains no designated bond other than $b_0(c)$ (Lemma~\ref{lem:two-block}), so $h_U\widetilde D_U(A)$ restricted to $\mathcal E_c$ involves only $\widetilde D_{U,c}$, and $\mathcal N_{U,c}$ depends on its argument restricted to $\mathcal E_c$ (Lemma~\ref{lem:two-block}); by \eqref{eq:regular-stencil-elimination}, $A|_{\mathcal E_c}$ depends on $B|_{\mathcal E_c\cap E_*}$ only. At $(0,0)$ the derivative in $d_c$ of $d_c-\mathcal N_{U,c}(A-h_U(c)d_c)$ is the identity, since $D\mathcal N_U(0)= D\mathcal Q_U(0)-LQ_U=0$, and $\mathcal N_{U,c}(0)=0$ (the zero phase reproduces the averages of $U$). The implicit function theorem gives the germ, unique among continuous solutions near $(0,0)$, with $d_c(0)=0$ and $Dd_c(0)= D\mathcal N_{U,c}(0)(\ldots)=0$. For $U= U_1$ the row solution $\varphi_c$ of (B4)(i) is analytic with $\varphi_c(0)=0$ and solves the row equation (Lemma~\ref{lem:c2-fibre-smallness}(i)), so its germ at zero is this one. If $B=0$ on $\mathcal E_c\cap E_*$ then $A=0$ on $\mathcal E_c$, so the germ value is $\widetilde D_{U,c}(0)=0$ and $q_U(B)=0$ on $\mathcal E_c$. The last claim follows from \eqref{eq:c-reconstructed-field}, \eqref{eq:c-bondwise-shift} and $\mathcal F_{b_0(c)}=\mathcal E_c\cap E_*$: the links of $U_{\xi}(a)$ on the bonds of $\mathcal E_c$ involve the row solution of $c$ at the argument $g_0\mathsf C_U\Psi_{\xi}(a+A_{\mathrm e})|_{\mathcal E_c}$, which involves $a$ only through $a|_{\mathcal E_c\cap E_*\cap I}$.
\end{proof}

The analyticity of the reconstruction on the cutoff box is the following lemma, proved in Appendix~\ref{app:background-estimates} (Lemmas~\ref{lem:c2-uniform-chart} and~\ref{lem:c2-primitive-bounds}); the enlarged box is $\{ a\in\mathfrak g^I\, | \,| a(b)|\leqslant 2r\ \forall b\in I\}$ with $r$ the cutoff radius of \eqref{eq:native-source-young-cutoff}, $R_{\mathrm c}$ is the collar range of Lemma~\ref{lem:collar-coverage}, the row polydisc of a row $c$ is $\mathcal U_c=\{ B\in\mathfrak g^{\Omega_1^*}\, | \,| B(b)|<2+4r\ \forall b\in\mathcal E_c\cap E_*\}$, and the block polydisc of a centre $y$ is $\mathcal U_y=\{ B\, | \,| B(b)|<2+4r\ \forall b\in\mathcal E_y\}$ (every bond of $\mathcal E_y$ is free). A centre $y$ lies within $\rho$ of a bond set $S$ when every bond of $\mathcal E_y$ does.

\begin{lemma}[charts]\label{lem:charts}
Let $\vartheta_*>0$ be the phase radius of Appendix~\ref{app:background-estimates}, Lemma~\ref{lem:c2-uniform-chart} ($\vartheta_*\leqslant \tau_L/2$; it depends on $L$ and the group only). On every branch, and for the reference background $U=\id$ of Section~\ref{sec:joint-extraction}, for every row $c\in\Omega_1^{[1]}$ within $2R_{\mathrm c}+2L$ of $I$ and every real phase $A$ with $\| A|_{\mathcal E_c}\|\leqslant \vartheta_*/2$ (Hilbert norm of $\mathfrak g^{\mathcal E_c}$), the row equation of $c$ at $A$ has exactly one small solution, $\widetilde D_{U,c}(A)=\varphi_c(A|_{\mathcal E_c})$, where $\varphi_c$ is analytic on the complex ball $\{\| A|_{\mathcal E_c}\|<\vartheta_*\}$, real on real phases, with $\varphi_c(0)=0$, $D\varphi_c(0)=0$; under the threshold of Section~\ref{app:thresholds} (item 3) $g_0\mathsf C_U\mathcal U_c$ lies in the ball $\| A|_{\mathcal E_c}\|\leqslant \vartheta_*/2$, so the row solution of Section~\ref{sec:source-phase-primitives} is defined at $A= g_0\mathsf C_UB$ for every $B\in\mathcal U_c$; the determinant $\det(\id-h_U(c)\partial_{A(b_0(c))}\widetilde D_{U,c}(A))$ is positive there, and the following are analytic functions of $B|_{\mathcal E_c\cap E_*}$ on $\mathcal U_c$: $\widetilde D_{U,c}$, the phase $q_U$ on $\mathcal E_c$, the normalised Haar logarithm $\log(\sigma(q_U(\bigcdot )(b))/\sigma(0))$ for $b\in\mathcal E_c$, the real logarithm of that determinant, and, when $c$ is within $2R_{\mathrm c}$ of $I$, the plaquette terms $s_{p,U}(q_U(\bigcdot ))$ of the plaquettes $p$ with a bond in $\mathcal E_c$ (as functions of $B$ on the intersection of the polydiscs of the rows of $p$, which are rows within $2R_{\mathrm c}+2L$ of $I$); the first four vanish at $B=0$, $\widetilde D_{U,c}$ and $q_U(\bigcdot )(b_0(c))-g_0(\mathsf C_UB)(b_0(c))$ together with their first derivatives. For every centre $y\in\Omega_1^{(1)}$ within $2R_{\mathrm c}+2L$ of $I$ and every $B\in\mathcal U_y$, the relative means $\mathcal U_{e^{iq}U}(y,x)\mathcal U_U(y,x)^{-1}$, $x\in\mathcal B_y$, at $q= q_U(B)= g_0B$ on $\mathcal E_y$ exist and lie within $r_G$ of $\id$ in operator norm, in the domain of the chart logarithm $\log_G$ of Appendix~\ref{app:background-estimates}, and the gauge term $\mathscr G_{y,U}(q_U(\bigcdot ))$ is an analytic function of $B|_{\mathcal E_y}$ on $\mathcal U_y$ vanishing at $B=0$. All these functions are the compositions of analytic functions $\varphi$ of the real phase $A|_{\mathcal E_c}$ (or $q|_{\mathcal E_y}$) with the linear map $B\mapsto g_0\mathsf C_UB$, where $\varphi$ is analytic on a complex neighbourhood of the closed phase ball $\| A|_{\mathcal E_c}\|\leqslant \vartheta_*/2$ (which contains $g_0\mathsf C_U\mathcal U_c$; for the row solution and the phase, on the complex ball of radius $\vartheta_*$) and has derivatives up to order five bounded on that closed ball by constants depending on $L$ and the group only; hence their derivatives in $B$ up to order five are bounded on the polydiscs by such constants, and the primitives of this subsection formed from them satisfy, on the enlarged box at the inputs $a+t\Psi_{\xi}(A_{\mathrm e})$, $t\in[0,1]$, $|\mathfrak p(a+t\Psi_{\xi}(A_{\mathrm e}))|\leqslant C_{\mathrm p}g_0(1+2r)^3$ with $C_{\mathrm p}$ depending on $L$ and the group only (the prefactor $g_0^{-2}$ of $E_{\xi,\alpha}$ is absorbed by the third-order Taylor remainder, Lemma~\ref{lem:c2-primitive-bounds}, \eqref{eq:c2-primitive-sizes}). Where $(A_{\mathrm e},a)\in\mathcal D_{\xi}$, these objects coincide with the objects of CRE (1.14)--(1.15) under the reading of (B4)(i), by (B4)(i) and the uniqueness of the small solution (the phase $g_0\mathsf C_U\Psi_{\xi}(a+A_{\mathrm e})$ lies in the ball of radius $\vartheta_*/2$ on every such row, by the inclusion stated next with $t=1$, $a$ being in the open box by Lemma~\ref{lem:collar-coverage}(iv)). Every $a+t\Psi_{\xi}(A_{\mathrm e})$ with $a$ in the enlarged box and $t\in[0,1]$ lies in $\mathcal U_c$ and $\mathcal U_y$ for every such row and centre (Lemma~\ref{lem:collar-coverage}(i),(iii): $\Psi_{\xi}(A_{\mathrm e})= A_{\mathrm e}$ and $| A_{\mathrm e}(b)|< r$ on the bonds of $\mathcal E_c$, $\mathcal E_y$, which are within $2R_{\mathrm c}+4L-1\leqslant 3R_{\mathrm c}$ of $I$, the range of Lemma~\ref{lem:collar-coverage}(i),(iii)), so every near primitive of this subsection is an analytic function of $(a,t)$ there, with its value and its derivatives in $a$ up to order four bounded as in \eqref{eq:c2-primitive-sizes}--\eqref{eq:convex-primitive-derivatives} of Lemma~\ref{lem:c2-primitive-bounds} (under threshold item 4 of Section~\ref{app:thresholds}, $g_0s_{\xi}^2\leqslant 4g_0(1+r)^3\leqslant 4\omega_*/C_0$, so these are bounds by constants depending on $L$ and the group only).
\end{lemma}

The centred exponent and its primitives are those of CRE (1.14)--(1.16), written as functions of the free vector $B\in\mathfrak g^{\Omega_1^*}$. With $s_{p,U}(q)=1-\operatorname{Re}\tr_N((e^{iq}U)(\partial p))$ for a plaquette $p$ meeting $\Omega_1$ (links off $\Omega_1$ fixed at $V_0$),
\begin{equation}\label{eq:native-source-energy-primitives}
    E_{\xi,p}(B)=-\frac{1}{g_0^2}\bigl[s_{p,U}(q_U(B))-s_{p,U}(0)\bigr]\;,
\end{equation}
and $E_{\xi,y}(B)=-g_0^{-2}\mathscr G_{y,U}(q_U(B))$ from the $y$-summand of \eqref{eq:regular-covariance-gauge}, $y\in\Omega_1$. For either label $\alpha$,
\begin{equation}\label{eq:native-source-energy-residual}
    p_{\xi,\alpha}(B)= E_{\xi,\alpha}(B)-DE_{\xi,\alpha}(0)[B]-\frac{1}{2}D^2E_{\xi,\alpha}(0)[B,B]\;.
\end{equation}
With $\sigma(q)\,\dd q$ the pullback under $q\mapsto e^{iq}$ of the Haar measure $\dd U$ of CRE (1.11) (the density $\sigma$ of (1.14), for the fixed inner product as explained in the conventions; its normalisation enters only the constants $\kappa_{\xi}$ and $c_{\xi}$), for every bond $b$ of $\Omega_1$,
\begin{equation}\label{eq:native-source-haar-primitive}
    p^{\mathrm H}_{\xi,b}(B)=\log\sigma(q_U(B)(b))-\log\sigma(0)\;,
\end{equation}
and for every row $c\in\Omega_1^{[1]}$
\begin{equation}\label{eq:native-source-jacobian-primitive}
    p^{\mathrm J}_{\xi,c}(B)
        =\log\Bigl|\det_{\mathfrak g}\left(\id-h_U(c)\partial_{A(b_0(c))}\widetilde D_{U,c}(g_0\mathsf C_UB)\right)\Bigr|\;,
\end{equation}
where $\partial_{A(b_0(c))}\widetilde D_{U,c}$ is the partial derivative of the row solution, regarded as a function of the $|\mathcal E_c|$ independent variables $A(b)$, $b\in\mathcal E_c$, with respect to the variable at $b_0(c)$, evaluated at $A= g_0\mathsf C_UB$ (this is the designated block of CRE's $\frac{\delta}{\delta A}\widetilde D$, proof of Lemma~\ref{lem:primitive-sum}). A primitive is defined at $B$ when the row solutions of its rows are defined at $g_0\mathsf C_UB$ (Section~\ref{sec:source-phase-primitives}: on the domain of integration, on the polydiscs for the rows near $I$, and near $B|_{\text{inputs}}=0$); for a primitive whose row solutions are not defined at $B=0$ on the domain of integration (which happens when $R_1\ne\varnothing$, since the large-fluctuation tests exclude $A=0$ on $R_1$), the Taylor polynomial at $B=0$ in \eqref{eq:native-source-energy-residual} is formed from the germs of Lemma~\ref{lem:row-locality}; CRE's $V$ is the sum of the terms of order higher than two of the expansion at $A=0$ (p.249, RG1 p.267), and Lemma~\ref{lem:primitive-sum} identifies the two. The letter $\mathfrak p$ denotes a primitive of any of the three kinds. By Lemma~\ref{lem:row-locality} the free-input set of $E_{\xi,p}$ and $p_{\xi,p}$ is $\bigcup_{r\subset\partial p}\mathcal F_r$, that of $E_{\xi,y}$ and $p_{\xi,y}$ is $\mathcal E_y\cap E_*=\mathcal E_y$, that of $p^{\mathrm H}_{\xi,b}$ is $\mathcal F_b$ and that of $p^{\mathrm J}_{\xi,c}$ is $\mathcal E_c\cap E_*$; the background set of each is the set of bonds whose links enter (bonds of $\Omega_1$ and, for plaquettes and rows at the boundary of $\Omega_1$, bonds of $\Omega_1^c$ carrying $V_0$), as in \eqref{eq:regular-stencil-quadratic-decomposition}; the rows of a primitive are the rows $c$ with $b_0(c)$ in its background set (a gauge primitive $E_{\xi,y}$ has no rows: $\mathcal E_y$ contains no designated bond, and $q_U(B)= g_0B$ on it). A primitive is \emph{near} if its free-input set meets $I$ and \emph{far} otherwise; a far primitive is a function of $A_{\mathrm e}$ and the branch data alone (Lemma~\ref{lem:row-locality}, on the domain of (B4)). A near primitive has a bond $b_1\in I$ in its free-input set, and every bond of $\mathcal E_c$ for a row $c$ of the primitive, and every bond of $\mathcal E_y$ for a gauge primitive, is within $4L+3\leqslant 2R_{\mathrm c}$ of $b_1$ (two bonds of a two-block set are within $2L-1$ of each other, the bonds of a plaquette within $1$); so the rows of a near primitive, and the centre of a near gauge primitive, lie within $2R_{\mathrm c}$ of $I$, and by Lemma~\ref{lem:charts} every near primitive is defined and analytic in $(a,t)$ at the inputs $a+t\Psi_{\xi}(A_{\mathrm e})$, $a$ in the enlarged box, $t\in[0,1]$. An \emph{arc} $\gamma$ of the contour $C$ is a set of consecutive bonds of $C$ (with their orientations), and $\gamma_+$ is the set of their positive representatives; the free-input and background stencils of an arc are
\begin{equation}\label{eq:native-source-arc-stencils}
    F_{\gamma}=\bigcup_{r\in\gamma_+}\mathcal F_r,
    \qquad
    D_{\gamma}=\gamma_+\cup\bigcup_{c\, | \,b_0(c)\in\gamma_+}\mathcal E_c\;.
\end{equation}

\begin{lemma}[exact primitive sum]\label{lem:primitive-sum}
Let $v$ and $V$ be the functions of CRE (1.15), p.249: $v(A)$ is the sum of the logarithms of the Haar densities $\sigma(A-h\widetilde D(A))$ over the bonds of $\Omega_1$ and of $\log|\det(\id-h\frac{\delta}{\delta A}\widetilde D)|$, the logarithm of the absolute Jacobian of (1.14), which is CRE's $\operatorname{Tr}\log(\id-h\frac{\delta}{\delta A}\widetilde D)$ where the determinant is positive (on the near rows, factor by factor, by Lemma~\ref{lem:charts}; on all of $\mathcal A_{\xi}$ by (B4)(iii), Lemma~\ref{lem:c2-fibre-smallness}(iv)), the identities below using the absolute value throughout, and $V(A)$ is the remainder after the Taylor polynomial of degree two at $A=0$ of the action and the gauge-fixing functional at the phase $A-h\widetilde D(A)$ (CRE's ``terms of order higher than two''; no series expansion is used, and the Taylor coefficients at $A=0$ are those of the germs of Lemma~\ref{lem:row-locality}). For $B\in\mathcal A_{\xi}$, and more generally for every $B\in\mathfrak g^{\Omega_1^*}$ at which every row solution is defined with the row solutions of the near rows those of Lemma~\ref{lem:charts} and those of the far rows the germs at zero argument ($B=0$ on their two-block sets), with $A= g_0\mathsf C_UB$,
\begin{equation}\label{eq:native-source-primitive-sum}
    v(A)-v(0)-g_0^{-2}V(A)
        =\sum_{\alpha}p_{\xi,\alpha}(B)+\sum_{b}p^{\mathrm H}_{\xi,b}(B)+\sum_{c}p^{\mathrm J}_{\xi,c}(B)\;,
\end{equation}
the first sum over the plaquettes $p$ meeting $\Omega_1$ and the centres $y\in\Omega_1^{(1)}$, the second over the bonds $b$ of $\Omega_1$, the third over the rows $c\in\Omega_1^{[1]}$ (the same ranges in every sum $\sum_{\alpha}$, $\sum_b$, $\sum_c$ below), so that the exponent of the second factor of CRE (1.15) at the free vector $B$, with the quadratic form written through the regional Hessian \eqref{eq:c-regional-hessian},
\begin{equation}\label{eq:native-collar-full-exponent}
    \mathscr E_{\xi}(B)=-\frac{1}{2}\langle B,H_{\xi} B\rangle+v(g_0\mathsf C_UB)-g_0^{-2}V(g_0\mathsf C_UB)\;,
\end{equation}
satisfies $\mathscr E_{\xi}(0)= v(0)$ with $v(0)=\sum_b\log(\sigma(0)/\sigma_0)=0$ (sum over the bonds of $\Omega_1$; here $\sigma$ is the density appearing in CRE's $v$, read as normalised relative to the outer constant $\sigma_0$ of (1.14)--(1.15), the reading fixed in the step list after \eqref{eq:native-source-physical-ratio}; in the unnormalised reading $v(0)=\sum_b\log\sigma(0)$ would still be a constant of the branch depending on $\Omega_1$ only, and in both readings $\sigma_0^{|\Omega_1^*|}e^{v(0)}\prod_be^{p^{\mathrm H}_{\xi,b}}$ is CRE's full Haar factor exactly once, the primitives being differences) and
\begin{equation}\label{eq:c-centred-exponent}
    \mathscr E_{\xi}(B)-\mathscr E_{\xi}(0)=-\frac{1}{2}\langle B,H_{\xi} B\rangle+\sum_{\alpha} p_{\xi,\alpha}(B)+\sum_bp^{\mathrm H}_{\xi,b}(B)+\sum_cp^{\mathrm J}_{\xi,c}(B)
    =\sum_{\alpha} E_{\xi,\alpha}(B)+\sum_bp^{\mathrm H}_{\xi,b}(B)+\sum_cp^{\mathrm J}_{\xi,c}(B)\;.
\end{equation}
\end{lemma}

\begin{proof}
The sum $\sum_{\alpha} E_{\xi,\alpha}(B)$ is $-g_0^{-2}$ times the difference of the action plus gauge-fixing functional at the phase $q_U(B)$ and at $q=0$: here the gauge-fixing functional of CRE (1.11), (1.14) is $\mathbf G(\Omega_1^{(1)},e^{iq}U_1)=\mathscr G^{\Omega_1}_{U_1}(q)$ because $U_1$ is in the axial mean gauge (Section~\ref{sec:regular-hessian-stencils}, after \eqref{eq:regular-covariance-hessian}). Since $D\widetilde D_U(0)=0$, $Dq_U(0)[B]= g_0\mathsf C_UB$, and the linear term of the Taylor expansion in $B$ at zero is $\sum_{\alpha} DE_{\xi,\alpha}(0)[B]=-g_0^{-1}\bigl(DS(U_1)[\mathsf C_UB]+D\mathscr G^{\Omega_1}_U(0)[\mathsf C_UB]\bigr)=0$: the relative gauge functional is a sum of nonnegative terms vanishing at $q=0$, so its first variation vanishes; by CRE (1.12) $U_1$ is a critical point of $S+\mathbf G(\Omega_1^{(1)},\bigcdot )$ on the configurations with $\mathsf M(U)= V$ on $\Omega_1^{(1)}$ and $U= V_0$ on $\Omega_1^c$, and $\mathbf G(\Omega_1^{(1)},\bigcdot )\geqslant 0$ attains its minimum $0$ at $U_1$, so $D\mathbf G(\Omega_1^{(1)},U_1)=0$ and $DS(U_1)[\delta]=0$ for every phase $\delta$ supported on the bonds of $\Omega_1$ with $D\mathcal Q_U(0)\delta= LQ_U\delta=0$ (the tangent space of the constraint $\mathsf M(e^{i\delta}U_1)= V$); and $\mathsf C_UB$ is such a phase by $Q_U\mathsf C_U=0$ \eqref{eq:regular-stencil-elimination}. The quadratic term is $-\frac12\langle B,H_{\xi} B\rangle$ by \eqref{eq:regular-covariance-hessian}, with $A= g_0\mathsf C_UB$: the second variation of the action is $\langle A,\Delta_UA\rangle$ plus the force term $-2\langle J_U,h_U\frac12D^2\mathcal N_U(0)[A,A]\rangle$ from the second-order part $-h_U\widetilde D^{(2)}_U(A)$ of the phase, and that of the gauge functional is $G^{(2),\Omega_1}_U(A)$. Subtracting both from $\sum_{\alpha} E_{\xi,\alpha}$ gives $\sum_{\alpha} p_{\xi,\alpha}$, which is therefore $-g_0^{-2}V(A)$. The Haar part of $v$ is $\sum_b\log\sigma(q_U(B)(b))$ bond by bond, so $v$ minus its Haar part at $B=0$ gives $\sum_bp^{\mathrm H}_{\xi,b}$. For the Jacobian, $\frac{\delta}{\delta A}\widetilde D$ has rows indexed by $\Omega_1^{[1]}$ and columns by the bonds of $\Omega_1$, and $h_U$ maps row $c$ to the bond $b_0(c)$; in the ordering (free bonds, designated bonds) the matrix $\id-h_U\frac{\delta}{\delta A}\widetilde D_U$ is block triangular with identity block on the free bonds and designated block $\operatorname{diag}_c(\id-h_U(c)\partial_{A(b_0(c))}\widetilde D_{U,c})$ by Lemma~\ref{lem:two-block} (row $c$ depends on $A|_{\mathcal E_c}$, which contains the single designated bond $b_0(c)$), so its determinant is $\prod_c\det_{\mathfrak g}(\id-h_U(c)\partial_{A(b_0(c))}\widetilde D_{U,c})$, positive on $\mathcal A_{\xi}$ by (B4)(iii), and at the other points admitted in the statement because each near-row factor is positive by Lemma~\ref{lem:charts} and each far-row factor equals one at zero argument ($D\widetilde D_{U,c}(0)=0$); so $\operatorname{Tr}\log$ of the matrix, CRE's term in $v$, is the real logarithm of the positive determinant, and since each factor is a nonzero real, the logarithm of the absolute value of the product is the sum of the logarithms of the absolute values of the factors, $\sum_c\log|\det_{\mathfrak g}(\cdots)|=\sum_cp^{\mathrm J}_{\xi,c}(B)$; this vanishes at $B=0$ since $D\widetilde D_U(0)=0$; hence $v(0)$ is the Haar part at $B=0$. The quadratic form of (1.15) in the free variable, $\frac12\langle B,C^*\Delta_1CB\rangle$, is $\frac12\langle B,H_{\xi} B\rangle$ by the identification $\mathsf H^{\Omega_1}_{U_1}= C^*\Delta_1C$ of Section~\ref{sec:regular-hessian-stencils} ($\Delta_1$ of (1.16) already contains $\mathbf G^{(2)}$ and the force term; the $g_0^{-2}$ of the action cancels the $g_0^2$ of $A= g_0\mathsf C_UB$). Subtracting $\mathscr E_{\xi}(0)= v(0)$ from \eqref{eq:native-collar-full-exponent} and inserting \eqref{eq:native-source-primitive-sum} gives the first form of \eqref{eq:c-centred-exponent}; the second form restores the quadratic term to $\sum_{\alpha} p_{\xi,\alpha}$.
\end{proof}

Each primitive is assigned a host $Y_{\alpha}$ in the native localisation family $\mathbf D_0$: lift its free-input and background stencil $\widetilde{\mathcal S}_{\alpha}$ to $\mathbb Z^4$ about its base (plaquette lower vertex, block centre, initial vertex of its bond, or initial vertex of $b_0(c)$ for $p^{\mathrm J}_{\xi,c}$; the stencil has $\ell^{\infty}$ diameter at most $4L< M$, below every period, so the lift about the base is well defined), let $Z_{\alpha}$ be the $M$-cubes $Q_z= Mz+[-M/2,M/2]^4$ meeting it, and put
\begin{equation}\label{eq:part1-potential-host}
    Y_{\alpha}=\pi_T\Bigl(\bigcup_{z\in\prod_{\mu}([k^-_{\alpha,\mu}-1,k^+_{\alpha,\mu}+1]\cap\mathbb Z)}Q_z\Bigr),
    \qquad k^-_{\alpha,\mu}=\min_{z\in Z_{\alpha}}z_{\mu},\quad k^+_{\alpha,\mu}=\max_{z\in Z_{\alpha}}z_{\mu}\;,
\end{equation}
with $\pi_T$ the reduction modulo the periods; since the stencil has diameter at most $4L< M$, it meets at most two consecutive $M$-cubes in each direction, so the rectangle has side at most $4$ cubes, and every period being at least $8M$, $\pi_T$ is injective on it. $Y_{\alpha}$ is connected and contains the stencil in its interior. The potential is, for $Y\in\mathbf D_0$ and $a\in\mathfrak g^I$,
\begin{equation}\label{eq:c-potential-definition}
    \mathbf V_{\xi}(Y,a)=\sum_{\mathfrak p\ \mathrm{near},\ Y_{\mathfrak p}= Y}\mathfrak p(a)\quad\text{if }| a(b)|< r\ \forall b\in E(Y),\qquad\mathbf V_{\xi}(Y,a)=0\quad\text{otherwise},
\end{equation}
the sum over the near primitives $\mathfrak p$ of all three kinds assigned to $Y$, with $E(Y)=\{ b\in I\, | \,b\cap Y\ne\varnothing\}$ and the primitives at zero retained input ($t=0$). It depends on $a|_{E(Y)}$ only, since the free-input set of a primitive assigned to $Y$ lies in $Y$, hence in $E(Y)$ (Lemma~\ref{lem:row-locality}), and it is bounded on the local box by Lemma~\ref{lem:charts} (each summand by $C_{\mathrm p}g_0(1+2r)^3$). On the open box $\{| a(b)|< r\ \forall b\in I\}$, $\sum_Y\mathbf V_{\xi}(Y,a)=\sum_{\mathfrak p\ \mathrm{near}}\mathfrak p(a)$.

\subsection{The boundary law, the exterior interpolation and the collar}
\label{sec:boundary-law-collar}

\emph{The change of variables.} Let $S_1=(\Omega_1\setminus\Omega_1^{\sim-1})\cup R_1$ (CRE (1.20)) and let $S_1^*$ be its free bonds. CRE (1.22) replaces the variable $A$ of (1.15) on each bond $b\in S_1^*$ by a new variable $A'(b)$ and leaves the other bonds unchanged: $A=\Psi_{\xi} A'$ with
\begin{equation}\label{eq:native-source-exact-boundary-map}
    (\Psi_{\xi} A')(b)= g_0^{-1}\log\left(e^{ig_0A'(b)}e^{-i\omega_b(A'(b))\mathbf H_{1,Ax}(b)}\right)\quad(b\in S_1^*),\qquad
    (\Psi_{\xi} A')(b)= A'(b)\quad(b\in\Omega_1^*\setminus S_1^*)\;,
\end{equation}
where $\omega_b$ is the $C^{\infty}$ cutoff $g$ of CRE p.252 ($0\leqslant g\leqslant 1$, $g=0$ on $| A'|\leqslant \frac43g_0^{-1}\delta_0$, $g=1$ where $| e^{ig_0A'}-1|\geqslant \frac53\delta_0$) for $b$ in the set $(R_1^{\sim2}\cap S_1)^*$ named in the definition preceding (1.22) (p.251); CRE specifies $g$ by these properties only, and one such $g$ with $| Dg(A')|\leqslant 6g_0/\delta_0$ for every $A'$ (operator norm for the cutoff norm $|\bigcdot |$) is fixed once: the function $\ell(A')=\min(1,\max(0,(| g_0A'|-\frac{17}{12}\delta_0)/(\frac16\delta_0)))$ is Lipschitz with constant $6g_0/\delta_0$, vanishes for $| g_0A'|\leqslant \frac{17}{12}\delta_0$ and equals one for $| g_0A'|\geqslant \frac{19}{12}\delta_0$, and its convolution with a $C^{\infty}$ probability density supported in $\{| A'|\leqslant \delta_0/(12g_0)\}$ is $C^{\infty}$ with values in $[0,1]$, has the same Lipschitz constant, vanishes for $| g_0A'|\leqslant \frac43\delta_0$ and equals one for $| g_0A'|\geqslant \frac53\delta_0$, hence wherever $| e^{ig_0A'}-\id|\geqslant \frac53\delta_0$, as $| e^{iX}-\id|\leqslant | X|$; and $\omega_b=1$ for the other bonds of $S_1^*$ (p.252 describes the second set as $S_1^*\setminus(R_1^{\sim2})^*$, which differs from the complement of the first set in $S_1^*$ on the bonds joining $R_1^{\sim2}\setminus S_1$ to $S_1\setminus R_1^{\sim2}$; the definition of p.251 is used, and (B4)(ii) refers to the map so defined; on the bonds of the difference the map is CRE's second form, $g$ replaced by one, and CRE's remark that these changes of variables ``do not change the functions $\chi^{(0)}$'' (p.252), its rewriting (1.23) of the cutoff and its rewriting of the retained tests in the new variables, p.252, with the localisation property stated there ((B3)), are not used, the cutoff being composed exactly in \eqref{eq:retained-cutoff-composition}), $\mathbf H_{1,Ax}(b)\in\mathfrak g$ is the field $\mathbf H_{1,Ax}(\Omega_1\cap\Lambda_1^c,\bigcdot )$ of CRE (1.21) at the actual branch data, as it appears in (1.22), without the prefactor $L^{-1}$ of (1.21) (a function of the pair $(V,V_0)$ through $U_1= U_1(V,V_0)$ and $U_1(\Omega_1\cap\Lambda_1^c,V)$, the latter stated by CRE p.251 to depend on $V|_{\Omega_1\cap\Lambda_1^c}$; Borel in $(V,V_0)$ by the measurability paragraph of Section~\ref{sec:branch-data}; with (1.21) and $\omega_b=1$, $e^{ig_0(\Psi_{\xi} A')(b)}U_1(b)= e^{ig_0A'(b)}e^{-i(1-L^{-1})\mathbf H_{1,Ax}(b)}U_1(\Omega_1\cap\Lambda_1^c,b)$, which is not the second factor of CRE's (1.23): as printed, (1.22) carries $\mathbf H_{1,Ax}$ where (1.21) carries $L^{-1}\mathbf H_{1,Ax}$, and (1.24) has $+F$ against the $-F$ of (1.22); the identities of this paper use (1.22) as printed, with the exact Jacobian $|\det D\Psi_{\xi}|$ and the cutoff composed in the variable of (1.15), never (1.23)--(1.24), so they do not depend on the resolution of this discrepancy), and $\log$ is the group chart, the inverse of $A\mapsto e^{iA}$ on the chart ball (so that $\log$ is the chart logarithm of the fluctuation, written $\frac1i\log$ in CRE's notation for the matrix logarithm, and $g_0^{-1}\log(\cdots)$ is the variable of (1.15)). $\Psi_{\xi}$ acts bond by bond, and $\psi_b$ denotes its component at $b\in S_1^*$, a $C^{\infty}$ map ($g$ is $C^{\infty}$, not analytic) on the open set $\mathcal W_b\subseteq\mathfrak g$ of those $A'(b)$ for which $g_0A'(b)$ lies in the chart ball and the argument of the chart logarithm lies in the chart domain (on the bonds where (1.22) is the identity, the variable $A'$ of (1.22) is the variable $A$ of (1.15), which on the free bonds is $g_0^{-1}$ times the fluctuation variable of (1.13); on $S_1^*$ the new variable $A'$ is restricted to the values with $g_0A'$ in the chart ball as well, which loses no preimage, every preimage under $\psi_b$ of a point of $\mathcal A_{\xi}$ lying in $g_0^{-1}K$ with $K$ the ball of the proof of Lemma~\ref{lem:boundary-bijection}(b),(c); the change of variables is made on this domain). By Lemma~\ref{lem:boundary-bijection} below, which is the (1.22) part of (B4)(ii), $\Psi_{\xi}$ maps $\mathcal D_{\xi}$ bijectively onto $\mathcal A_{\xi}$, the domain of integration of (1.15) with the cutoff (1.25) (Section~\ref{sec:branch-data}), and $D\Psi_{\xi}$ is invertible there, so that the change of variables formula applies with the Jacobian $|\det D\Psi_{\xi}|=\prod_{b\in S_1^*}|\det D\psi_b|$; $\chi_{\xi,\mathrm e}(A_{\mathrm e})$ below is set to zero when $A_{\mathrm e}(b)\notin\mathcal W_b$ for some $b\in S_1^*$. Thus $\log|\det D\Psi_{\xi}(A')|$ depends on $A'|_{S_1^*}$ only, and it is $C^{\infty}$; what Lemma~\ref{lem:collar} uses of it is that it is finite at $(A_{\mathrm e},0)\in\mathcal D_{\xi}$ and independent of $a$ (a constant on the branch, by $I\cap S_1^*=\varnothing$ and \eqref{eq:c-bondwise-shift}). The integration variable of the branch is $A'= a+A_{\mathrm e}$ with $a\in\mathfrak g^I$ and $A_{\mathrm e}\in\mathfrak g^J$. A bond of $I$ has an endpoint in $\Lambda_1\subseteq\Omega_1^{\sim-2}\cap(R_1'^{\sim})^c$, a bond of $S_1^*$ has an endpoint in $S_1$, and $\Omega_1^{\sim-2}$ is separated from $\Omega_1\setminus\Omega_1^{\sim-1}$, as $(R_1'^{\sim})^c$ is from $R_1\subseteq R_1'$, by a full layer of $LMR_1$-cubes (the enlargement adjoins every touching cube, Section~\ref{sec:sourced-effective-interaction}); hence, with distances between initial vertices,
\begin{equation}\label{eq:c-collar-distance}
    I\cap S_1^*=\varnothing,\qquad d(I,S_1^*)\geqslant LMR_1-2\;.
\end{equation}
The filling modification following CRE (1.20) only removes cubes from $\Lambda_1$, so \eqref{eq:c-collar-distance} holds for the final $\Lambda_1$. This gives \eqref{eq:c-bondwise-shift}.

\begin{lemma}[the boundary change is a bijection]\label{lem:boundary-bijection}
Let $r_G\in(0,1]$ be the group constant of Appendix~\ref{app:background-estimates} (defined there with the chart logarithm $\log_G$): every $U\in G$ with $\| U-\id\|_{\mathrm{op}}\leqslant r_G$ lies in the chart domain $N_G$, the image of the chart ball $\{ h(A)<\varrho_G\}$ under $A\mapsto e^{iA}$. Let $\delta_0\leqslant \delta_G\coloneqq \min\bigl(\log(3/2)/(12\sqrt N),\,r_G/(10\sqrt N)\bigr)$, a condition on the coupling ($\delta_0=(A_1/A_0)\varepsilon_0\to0$). Let $(\Omega_1,V,V_0)$ be an admissible pair and $R_1$ a large-fluctuation region of the history sum (Section~\ref{sec:branch-data}), so that $U_1$ and the configuration $U_1(\Omega_1\cap\Lambda_1^c,V)$ of (1.21) exist ((B1) and Lemma~\ref{lem:borel-minimiser}) and $\mathbf H_{1,Ax}(b)\coloneqq L\log\bigl(U_1(b)\,U_1(\Omega_1\cap\Lambda_1^c,V)(b)^{-1}\bigr)$, formed with the representatives chosen in Lemma~\ref{lem:borel-minimiser}, is defined wherever this chart logarithm is; under \eqref{eq:boundary-bijection-inputs} it is CRE's analytic $\mathbf H_{1,Ax}$ of (1.21), since $L^{-1}\mathbf H_{1,Ax}(b)$ then lies in the chart ball and (1.21) determines it there as this logarithm. For $b\in S_1^*$ let $\square'(b)\subseteq S_1$ be an $L^2M_2R_1$-cube with $b\in(\square'(b)^{\sim2})^*$ (one exists: $S_1$ is a union of such cubes, $\Omega_1\setminus\Omega_1^{\sim-1}$ being a union of $LMR_1$-cubes, each a union of $L^2M_2R_1$-cubes (CRE's cube partitions, pp.245--246, are nested: each cube of a coarser partition is a union of cubes of each finer one, the $L$-blocks included, the side ratios $M/(LM_2)$, a power of $L$ by the standing choice, and $LM_2R_1$ being integers), and $R_1$ a union of them, p.250; $b$ has an endpoint in $S_1$, hence in such a cube, and is not designated for a row outside $S_1^{[1]}$: if $b= b_0(c)$, the block containing its endpoint in $S_1$ lies in $S_1$ by the nesting, so $c\in S_1^{[1]}$, which $b\in S_1^*$ excludes), put $Z_b\coloneqq U_1(b)U_{1,\square'(b)}(b)^{-1}$, and assume CRE's bound of p.251 at the tolerance
\begin{equation}\label{eq:boundary-bijection-inputs}
|\mathbf H_{1,Ax}(b)|\leqslant \frac{\delta_0}{40}\qquad(b\in S_1^*)\;,
\end{equation}
the input discussed after the proof; the inequality $| Z_b-\id|<2\delta_0$ ($b\in S_1^*$) used with it is (B1)'s, derived in the proof. Let $\mathcal T_b\coloneqq \{ A\in\mathfrak g: h(g_0A)<\varrho_G,\ | e^{ig_0A}Z_b-\id|<2\delta_0\}$. Then for every $b\in S_1^*$: (a) $B(b)\in\mathcal T_b$ for every $B\in\mathcal A_{\xi}$; (b) $\mathcal T_b\subseteq\psi_b(\mathcal W_b)$; (c) $\psi_b$ is injective on $\psi_b^{-1}(\mathcal T_b)$, and $D\psi_b$ is invertible at every point of $\psi_b^{-1}(\mathcal T_b)$. Consequently $\Psi_{\xi}$ maps $\mathcal D_{\xi}$ bijectively onto $\mathcal A_{\xi}$, and $D\Psi_{\xi}$ is invertible at every point of $\mathcal D_{\xi}$.
\end{lemma}

\begin{proof}
Norm facts, for the cutoff norm $|\bigcdot |$ (operator or normalised Hilbert--Schmidt) and Hermitian $X$ with $\| X\|_{\mathrm{op}}\leqslant \pi$: $| e^{iX}-\id|\leqslant | X|\leqslant \frac\pi2| e^{iX}-\id|$ (eigenvalue by eigenvalue, $2|\sin(\theta/2)|\leqslant |\theta|\leqslant \pi|\sin(\theta/2)|$ for $|\theta|\leqslant \pi$, in both norms); $| UY|=| YU|=| Y|$ for unitary $U$, so $| UU'-\id|\leqslant | U-\id|+| U'-\id|$ for unitary $U,U'$; a unitary $U\in G$ with $| U-\id|< r_G/\sqrt N$ lies in the chart domain, since $\| U-\id\|_{\mathrm{op}}\leqslant \sqrt N| U-\id|$ for either cutoff norm and by the definition of $r_G$ (the eigenvalue argument alone does not give this for a general compact $G\subset U(N)$: the Hermitian logarithm $Y$ of $U$ with eigenvalues in $(-\pi,\pi]$ need not lie in the Lie algebra $\mathfrak g$, so $r_G$ is a constant of the group); and on $\{| X|\leqslant 12\delta_0\}$, writing $e^{iX}=\id+iX+R(X)$, one has $\| DR(X)\|\leqslant e^{\| X\|_{\mathrm{op}}}-1\leqslant e^{12\sqrt N\delta_0}-1\leqslant \frac12$ (the derivative of $X^k$ is a sum of $k$ terms $X^jVX^{k-1-j}$, each of norm at most $\| X\|_{\mathrm{op}}^{k-1}| V|$), so $| e^{iX}-e^{iX'}|\geqslant \frac12| X-X'|$ for $| X|,| X'|\leqslant 12\delta_0$.

Fix $b\in S_1^*$, write $H=\mathbf H_{1,Ax}(b)$ and $\omega=\omega_b$, and pass to the group-scale variable $x= g_0A'$: $\psi_b(A')= g_0^{-1}F(g_0A')$ with $F(x)=\log(e^{ix}W(x))$, $W(x)= e^{-i\widetilde\omega(x)H}$, $\widetilde\omega(x)=\omega(x/g_0)$, on $\mathfrak B_b= g_0\mathcal W_b=\{ x: h(x)<\varrho_G,\ e^{ix}W(x)\text{ in the chart domain}\}$; $| D\widetilde\omega|\leqslant 6/\delta_0$, and $| W(x)-\id|\leqslant \widetilde\omega(x)| H|\leqslant | H|\leqslant \delta_0/40$. Moreover $| Z_b-\id|<2\delta_0$: by (B1) the background $U_1$ lies in $D^{(11)}(V,V_0)$, the domain of integration in (1.11) (CRE p.248), on which the factor $\chi_0'(\square'(b))$ of $\chi_0'(\Omega_1)$ is nonzero, that is $| U_1(b')U_{1,\square'(b)}(b')^{-1}-\id|<2\delta_0$ for every $b'\in(\square'(b)^{\sim2})^*$ (CRE (1.8)), and $b$ is one of these bonds.

(a) For $B\in\mathcal A_{\xi}$, $U(B)(b)= e^{ig_0B(b)}U_1(b)$ ($b$ is free, so $q_{U_1}(B)(b)= g_0B(b)$, Lemma~\ref{lem:row-locality}), and $g_0B(b)$ is the component at $b$ of the chart logarithm $\frac1i\log U'$ of (1.13) (the variable written $A'$ there, not the new variable $A'$ of (1.22)), so $h(g_0B(b))<\varrho_G$ by the chart-class condition in the definition \eqref{eq:c-free-variable-domain} of $\mathcal A_{\xi}$; the cutoff (1.25) is nonzero at $U(B)$, in particular its factor $\chi_0'(\square'(b))$, the test $| U(B)(b')U_{1,\square'(b)}(b')^{-1}-\id|<2\delta_0$ for $b'\in(\square'(b)^{\sim2})^*$ (CRE (1.8) and p.250; a function of the fine field, composed with the changes of variables (1.14) and $A= g_0A'$ of p.249, Section~\ref{sec:branch-data}), and at $b'= b$ this reads $| e^{ig_0B(b)}Z_b-\id|<2\delta_0$.

(b) Let $y= g_0A$ with $A\in\mathcal T_b$. Then $| e^{iy}-\id|\leqslant | e^{iy}Z_b-\id|+| Z_b-\id|<4\delta_0$, and $\| y\|_{\mathrm{op}}<\pi$ since $h(y)<\varrho_G$, so $| y|\leqslant \frac\pi2\cdot4\delta_0<6.3\,\delta_0$. Let $K=\{ x\in\mathfrak g: | x|\leqslant 7\delta_0\}$, compact and convex, with $K\subseteq\mathfrak B_b$: $h(x)\leqslant | x|\leqslant 7\delta_0<\varrho_G$ (as $7\delta_0\leqslant 7r_G/(10\sqrt N)<\pi/\sqrt N=\varrho_G$, $r_G\leqslant 1$) and $| e^{ix}W(x)-\id|\leqslant | x|+| H|\leqslant 7\delta_0+\delta_0/40<8\delta_0\leqslant r_G/\sqrt N$, so $e^{ix}W(x)$, an element of $G$, lies in the chart domain. For $x\in K$, $e^{iF(x)}= e^{ix}W(x)$ gives $| e^{iF(x)}-e^{ix}|=| W(x)-\id|\leqslant | H|$ and $| e^{iF(x)}-\id|\leqslant 7\delta_0+| H|$, so $| F(x)|\leqslant \frac\pi2(7\delta_0+\delta_0/40)<11.1\,\delta_0<12\delta_0$ ($\| F(x)\|_{\mathrm{op}}<\pi$, $F(x)$ lying in the chart ball) and $| F(x)-x|\leqslant 2| e^{iF(x)}-e^{ix}|\leqslant 2| H|\leqslant \delta_0/20$, both $F(x)$ and $x$ lying in $\{| X|\leqslant 12\delta_0\}$. The continuous map $T(x)= x-F(x)+y$ maps $K$ into itself, $| T(x)|\leqslant \delta_0/20+6.3\,\delta_0<7\delta_0$, so by Brouwer's fixed point theorem it has a fixed point $x\in K$, that is $F(x)= y$; then $A'= x/g_0\in\mathcal W_b$ and $\psi_b(A')= A$.

(c) Let $x,x'\in\mathfrak B_b$ with $F(x)= F(x')= y$ and $y/g_0\in\mathcal T_b$. Then $e^{ix}= e^{iy}W(x)^{-1}$ gives $| e^{ix}-\id|\leqslant | e^{iy}-\id|+| H|<4\delta_0+\frac1{40}\delta_0$, so $| x|\leqslant \frac\pi2\cdot\frac{161}{40}\delta_0<6.4\,\delta_0$, and likewise $| x'|<6.4\,\delta_0$; both lie in $K$. From $e^{ix}W(x)= e^{ix'}W(x')$, $e^{ix}-e^{ix'}= e^{ix'}\bigl(W(x')W(x)^{-1}-\id\bigr)$ with $W(x')W(x)^{-1}= e^{i(\widetilde\omega(x)-\widetilde\omega(x'))H}$, so $| e^{ix}-e^{ix'}|\leqslant |\widetilde\omega(x)-\widetilde\omega(x')|\,| H|\leqslant \frac6{\delta_0}| x-x'|\cdot\frac{\delta_0}{40}=\frac3{20}| x-x'|$, while $| e^{ix}-e^{ix'}|\geqslant \frac12| x-x'|$; hence $x= x'$. For the differential at such an $x\in K$: with $P(x)= e^{ix}W(x)$, $F=\log\circ P$, and $D\log$ is invertible on the chart domain, so it suffices that $DP(x)$ is injective; $DP(x)[v]= D(e^{i\bigcdot })(x)[v]\,W(x)+e^{ix}DW(x)[v]$ with $| e^{ix}DW(x)[v]|=| D\widetilde\omega(x)[v]|\,| H|\leqslant \frac3{20}| v|$ and $| D(e^{i\bigcdot })(x)[v]\,W(x)|=| iv+DR(x)[v]|\geqslant \frac12| v|$, so $DP(x)[v]=0$ forces $v=0$.

The consequences: $\Psi_{\xi}$ is bondwise and the identity off $S_1^*$. For $B\in\mathcal A_{\xi}$, (a) and (b) give $A'$ with $A'(b)\in\mathcal W_b$ for $b\in S_1^*$ and $\Psi_{\xi}(A')= B$, and $(A_{\mathrm e},a)=(A'|_J,A'|_I)\in\mathcal D_{\xi}$ by \eqref{eq:c-new-variable-domain}; so $\Psi_{\xi}(\mathcal D_{\xi})=\mathcal A_{\xi}$, the inclusion $\subseteq$ being the definition. If $\Psi_{\xi}(A')=\Psi_{\xi}(A'')= B\in\mathcal A_{\xi}$, then bond by bond $\psi_b(A'(b))=\psi_b(A''(b))= B(b)\in\mathcal T_b$ by (a), so $A'(b)= A''(b)$ by (c); and $D\Psi_{\xi}=\bigoplus_{b\in S_1^*}D\psi_b\oplus\id$ is invertible at every point of $\mathcal D_{\xi}$ by (c). Moreover, by (b), (c) and the inverse function theorem, $\Psi_{\xi}$ restricted to the open set $\prod_{b\in S_1^*}\psi_b^{-1}(\mathcal T_b)\times\mathfrak g^{\Omega_1^*\setminus S_1^*}$ ($\mathcal T_b$ is open, being defined by strict inequalities between continuous functions) is a $C^{\infty}$ diffeomorphism onto the open set $\prod_{b\in S_1^*}\mathcal T_b\times\mathfrak g^{\Omega_1^*\setminus S_1^*}$, which contains $\mathcal A_{\xi}$ by (a); so $\mathcal D_{\xi}$ lies in the former open set, on which $\log|\det D\Psi_{\xi}|$ is $C^{\infty}$.
\end{proof}

\emph{The inputs.} The bound \eqref{eq:boundary-bijection-inputs} is CRE's bound of p.251 (an inline formula), $|\mathbf H_{1,Ax}|< O(1)B_3\exp(-\delta MR_1)\varepsilon_1$, ``much smaller than $\delta_0$, for $MR_1$ large enough'', stated for $\mathbf H_{1,Ax}$ ``on $\square'$'' for the cubes $\square'\subset S_1$ (p.251), on the strength of the distance from $\square'$ to the support of the argument of $\mathbf H_{1,Ax}$ (one layer of $M_1$-cubes touching $\partial\Lambda_1$, p.251) being greater than $MR_1$. Every bond of $S_1^*$ is a bond of such a cube: $S_1$ is a union of $L^2M_2R_1$-cubes $\square'$ (the sentence defining $\square'(b)$ above), and a bond of $S_1^*$ has an endpoint in $S_1$, hence in one of them; so CRE's bound, taken as stated with its constants $B_3$, $\delta_{\mathrm c}$ (CRE's $\delta$) and $O(1)$ independent of the coupling and of the branch as in (B3)(iii), holds on every bond of $S_1^*$. The ratio of the bound to $\delta_0$ is $O(1)B_3e^{-\delta_{\mathrm c}MR_1}\varepsilon_1/\delta_0$ with $\varepsilon_1/\delta_0\leqslant (1+\beta_0)A_0/A_1$ (CRE (2.8) and $\delta_0=(A_1/A_0)\varepsilon_0$, as in threshold item 5 of Appendix~\ref{app:thresholds}), so the tolerance $\delta_0/40$ is a condition ``$MR_1$ large'', met for small $g_0$ at fixed $M$ since $R_1\to\infty$. The second inequality of the proof, $| Z_b-\id|<2\delta_0$ on every bond of $S_1^*$, is CRE's statement that the critical point $U_1$ of (1.12) lies in the domain of integration in (1.11) (p.248, by Theorem 1 of \cite{BalabanVAR}; recorded in (B1) as $U_1\in D^{(11)}(V,V_0)$): that domain carries the factor $\chi_0'(\Omega_1)=\prod_{\square'\subset\Omega_1}\chi_0'(\square')$ of (1.8) and (1.11), and $\chi_0'(\square'(b))\neq0$ at $U_1$ is the inequality $| U_1(b')U_{1,\square'(b)}(b')^{-1}-\id|<2\delta_0$ for every $b'\in(\square'(b)^{\sim2})^*$, $b$ among them. The tolerance is the cutoff's $2\delta_0$. CRE states the comparison of p.250 between $U_1$ and $U_{1,\square'}$ for the cubes $\square'\subset R_1^c$ with $\square'^{\sim4}\subseteq\Omega_1$ (p.250; recorded in (B3)(iii) for the cubes of $\Omega_1^{\sim-1}\setminus R_1$, for all of which this holds), and its derivation of (1.25) keeps $\chi_0'$ on the cubes of $S_1$ (p.250); so for $Z_b$ formed with a cube $\square'(b)$ of the layer $\Omega_1\setminus\Omega_1^{\sim-1}$ no tolerance smaller than the cutoff's is stated (for the layer bonds of $(R_1^{\sim2}\cap S_1)^*$, CRE's ``very small'' of p.252, discussed below, concerns the comparison with the cubes of $R_1$ through (1.21), not $Z_b$), and none is derived here: for a cube $\square'$ of the layer within four $L^2M_2R_1$-layers of $\Omega_1^c$, $\square'^{\sim4}\not\subseteq\Omega_1$, and $U_1$ is tied by (1.12) to the exterior field $V_0$, which $\zeta(\Omega_1^c)$ constrains near $\Omega_1$ through small-field tests only: on $\Omega_1^{\sim}$ only small-field characteristic functions remain (p.248), the $\chi_0$ and $\chi_{Ax}$ localised there being equal to one and dropped (p.248, (B3)(i)), and the factors $\chi_0'(\square')$ of (1.8) of the cubes $\square'\subseteq\Omega_1^c$ outside $Q_1$ that carry one test the bonds of $(\square'^{\sim2})^*$, among them bonds of $\Omega_1$ for the cubes adjacent to it, with the tolerance $2\delta_0$. CRE's ``very small'' of p.252 concerns $U_1(\Omega_1\cap\Lambda_1^c)U_{1,\square'}^{-1}$ on $\square'^{\sim2}$, in CRE's discussion of the bonds of $(R_1^{\sim2}\cap S_1)^*$, not $Z_b$, and serves CRE to place $g(A')=1$ where the second factor of (1.23) vanishes; neither it nor (1.23) is used here, the cutoff being composed exactly in \eqref{eq:retained-cutoff-composition}, and the proof above uses the inequality $| Z_b-\id|<2\delta_0$ in the same way on the bonds of both forms. The proof depends on that inequality only through the numerical constants $4\delta_0$, $7\delta_0$ and $12\delta_0$ of (b) and (c) (and $\delta_G$), and goes through with $2\delta_0$ replaced by any tolerance $\eta\leqslant C\delta_0$ after the corresponding change of these constants; so a regularity bound for $U_1U_{1,\square'}^{-1}$ of the kind of p.247 with a coupling-independent constant would serve in place of the cutoff reading of (B1). The condition $\delta_0\leqslant \delta_G$ and the choice of $g$ above are this paper's; the condition on $MR_1$ at the tolerance $\delta_0/40$ (with $\varepsilon_1/\delta_0$ bounded by CRE (2.8) and the constants independent of the coupling by (B3)(iii)) and CRE's condition on $M_2R_1$ are listed under item 5 of Section~\ref{app:thresholds}, and $\delta_0\leqslant \delta_G$ among the further conditions listed after its enumeration.

\emph{The cutoff.} The small-field test of CRE (1.18) for a cube $\square'$ of side $L^2M_2R_1$ is the indicator of $\{\sup_{b\in(\square'^{\sim2})^*}| A(b)|< r\}$, $r=\delta_0/g_0$, in CRE's bond norm with the normalisation \eqref{eq:cre-cutoff-normalisation}, and it is the product of one indicator per bond of $(\square'^{\sim2})^*$. Let
\begin{equation}\label{eq:c-good-region}
    \mathcal G_{\xi}=\bigcup_{\square'\subset R_1^c}(\square'^{\sim2})^*,\qquad R_1^c=\Omega_1^{\sim-1}\setminus R_1\;,
\end{equation}
the union over the good cubes of (1.18). Since $\Lambda_1\subseteq R_1^c$ and both are unions of such cubes, the first two factors of CRE (1.25) are the product of the tests over all good cubes, which by idempotence of indicators is $\prod_{b\in\mathcal G_{\xi}}\mathbf 1_{\{| A(b)|< r\}}$. The young cutoff is
\begin{equation}\label{eq:native-source-young-cutoff}
    \Theta_{\xi}(a)=\prod_{b\in I}\mathbf 1_{\{| a(b)|< r\}}\;,
\end{equation}
and the retained cutoff is defined as the remaining factors of (1.25), composed with the change of variables:
\begin{equation}\label{eq:retained-cutoff-composition}
    \chi_{\xi,\mathrm e}(A_{\mathrm e})=\prod_{b\in\mathcal G_{\xi}\cap J}\mathbf 1_{\{|(\Psi_{\xi} A_{\mathrm e})(b)|< r\}}\cdot\bigl[\chi^{(0)c}(R_1)\chi_0'(S_1)\bigr](\Psi_{\xi} A_{\mathrm e})\;,
\end{equation}
where $\chi^{(0)c}(R_1)$ is the product of the large-fluctuation tests of (1.18) over the cubes of $R_1$ and $\chi_0'(S_1)$ the product of the tests of (1.8) over the cubes of $S_1$, both as functions of the variable $A$ of (1.15) on $\Omega_1^*$ (through the free links $e^{ig_0A(b)}U_1(b)$ in the case of (1.8)) and of the branch data ($V_0$ on $\Omega_1^c$, $U_1$ and the local minima $U_{1,\square'}$). CRE's rewriting of these factors in the new variables (p.252, (1.23)) is not used; the exact compositions are carried as they stand. Lemma~\ref{lem:collar-coverage}(ii) shows that \eqref{eq:retained-cutoff-composition} is indeed the complement of $\Theta_{\xi}$ in the exactly transformed cutoff.

\emph{The conditional integrand.} With $\mathscr E_{\xi}$ the exponent of the second exponential of CRE (1.15), written out in \eqref{eq:native-collar-full-exponent}, the conditional integrand of Section~\ref{sec:sourced-effective-interaction} is
\begin{equation}\label{eq:native-source-actual-functional}
    \mathcal K_{\xi}(a)=\begin{cases}\exp\Bigl(\mathscr E_{\xi}\bigl(\Psi_{\xi}(a+A_{\mathrm e})\bigr)+\log|\det D\Psi_{\xi}(a+A_{\mathrm e})|\Bigr),&\Theta_{\xi}(a)=1,\\[2pt] 0,&\Theta_{\xi}(a)=0,\end{cases}\qquad
    \mathcal J^{\mathrm a}_{\xi}[f]=\int f(a)\mathcal K_{\xi}(a)\,\dd a\;,
\end{equation}
the exponent being defined on the box by (B4) and Lemma~\ref{lem:charts}: by Lemma~\ref{lem:collar-coverage}(ii) below, the integrand of CRE (1.15) in the variables of (1.22), with the cutoff (1.25) in its exact composed form, is $\kappa_{\xi}\,\chi_{\xi,\mathrm e}(A_{\mathrm e})\,\mathcal K_{\xi}(a)$, as used in \eqref{eq:native-source-physical-ratio} ($\mathscr E_{\xi}$ and $\kappa_{\xi}$ are CRE's exponent and prefactor as printed, with $\sigma$, $\sigma_0$ and $z^{(0)}$ for the fixed inner product as explained in the conventions; only the $a$-dependence of $\mathcal K_{\xi}$ is used below). The change-of-variables term, which CRE expands as $V^{(0)}(S_1)$ ((1.24) and p.252), is kept exact inside $\mathcal K_{\xi}$; Lemma~\ref{lem:collar} shows that it is a constant on the branch.

\emph{The exterior interpolation.} Write $H= H_{\xi,II}$. For $t\in[0,1]$ put
\begin{equation}\label{eq:native-source-boundary-interpolation}
    \begin{aligned}
        P_{\xi,t}(a)&=\sum_{\mathfrak p\ \mathrm{near}}\mathfrak p(a+tA_{\mathrm e})\;,\qquad
        j_{\xi}(a)=\langle a,H_{\xi,IJ}A_{\mathrm e}\rangle\;,\\
        K_{\xi;t}(a)&=\mathbf 1_{\mathcal D_I}(a)\exp\left(-\frac{1}{2}\langle a,Ha\rangle-tj_{\xi}(a)+P_{\xi,t}(a)\right)\quad(\mathcal D_I=\{ a\, | \,| a(b)|< r\ \forall b\in I\}\text{ the open box, }\mathbf 1_{\mathcal D_I}=\Theta_{\xi}),\qquad
        I_{\xi;t}=\int K_{\xi;t}(a)\,\dd a,\qquad
        \dd\nu_{\xi;t}=\frac{K_{\xi;t}}{I_{\xi;t}}\,\dd a\;,
    \end{aligned}
\end{equation}
the sum running over the near primitives of Section~\ref{sec:source-phase-primitives} (all three kinds). By Lemma~\ref{lem:charts} and Lemma~\ref{lem:collar}\,(i), each near primitive is defined and analytic in $(a,t)$ on the enlarged box at the inputs $a+tA_{\mathrm e}$, $t\in[0,1]$; $K_{\xi;t}$ (zero off the open box) is therefore bounded and Borel on $\mathfrak g^I$, positive on the open box, and $0< I_{\xi;t}<\infty$. On the open box $P_{\xi,0}(a)=\sum_Y\mathbf V_{\xi}(Y,a)$ is the centred potential \eqref{eq:c-potential-definition}, and $K_{\xi;0}$ is the centred kernel of Section~\ref{sec:fixed-first-kernel}. This is the interpolation of CRE (3.27)--(3.29) in the exterior field, restricted to the terms that depend on $a$. The Wilson function and the centred coefficient are
\begin{equation}\label{eq:convex-response-operands}
    f_{\xi}(a)= W_{\chi}(C;U_{\xi}(a)),\qquad F_{\xi}(t)=\int f_{\xi}\,\dd\nu_{\xi;t}\;,
\end{equation}
with $U_{\xi}(a)$ the reconstructed fine field \eqref{eq:c-reconstructed-field} at the actual retained data, including the change of variables; it is not interpolated. By Lemmas~\ref{lem:charts} and~\ref{lem:row-locality}, $f_{\xi}$ is defined and analytic on the enlarged box: each contour link $U_{\xi}(a)(r)$ depends on $a$ only through $a|_{\mathcal F_r\cap I}$, through a row within $2R_{\mathrm c}$ of $I$ when $\mathcal F_r$ meets $I$, and is constant in $a$ otherwise. Off the enlarged box $f_{\xi}$ is extended by zero, which affects no integral against $\nu_{\xi;t}$.

\begin{lemma}[the collar]\label{lem:collar}
On every branch:
\begin{enumerate}
\item $\Psi_{\xi}(A_{\mathrm e})(b)= A_{\mathrm e}(b)$ for every $b\in J$ with $d(b,I)\leqslant 2R_{\mathrm c}$; every near primitive and every entry of $H_{\xi,IJ}$ involves $A_{\mathrm e}$ only on such bonds, so $\mathfrak p(a+t\Psi_{\xi}(A_{\mathrm e}))=\mathfrak p(a+tA_{\mathrm e})$ for near $\mathfrak p$ and $\langle a,H_{\xi,IJ}\Psi_{\xi}(A_{\mathrm e})\rangle= j_{\xi}(a)$. For every row $c$ within $2R_{\mathrm c}$ of $I$ the source phase on the bonds of $\mathcal E_c$ is, componentwise,
\begin{equation}\label{eq:native-collar-source-phase}
    q_U\bigl(\Psi_{\xi}(a+A_{\mathrm e})\bigr)(b)= q_U(a+A_{\mathrm e})(b)=\bigl[g_0\mathsf C_U(a+A_{\mathrm e})-h_U\widetilde D_U\bigl(g_0\mathsf C_U(a+A_{\mathrm e})\bigr)\bigr](b)\qquad(b\in\mathcal E_c)\;.
\end{equation}
\item With the branch constant
\begin{equation}\label{eq:native-collar-boundary-constant}
    c_{\xi}=\mathscr E_{\xi}(0)-\frac12\langle\Psi_{\xi}(A_{\mathrm e}),H_{\xi}\Psi_{\xi}(A_{\mathrm e})\rangle+\sum_{\mathfrak p\ \mathrm{far}}\mathfrak p\bigl(\Psi_{\xi}(A_{\mathrm e})\bigr)+\log|\det D\Psi_{\xi}(A_{\mathrm e})|\;,
\end{equation}
where $\mathscr E_{\xi}(0)= v(0)$ and the far primitives are evaluated at the retained field alone,
\begin{equation}\label{eq:c-collar-identity}
    \mathcal K_{\xi}(a)= e^{c_{\xi}}K_{\xi;1}(a)\quad(a\in\mathfrak g^I),\qquad
    \mathcal J^{\mathrm a}_{\xi}[f]= e^{c_{\xi}}I_{\xi;1}\int f\,\dd\nu_{\xi;1},\qquad
    \frac{\mathcal J^{\mathrm a}_{\xi}[f_{\xi}]}{\mathcal J^{\mathrm a}_{\xi}[1]}= F_{\xi}(1)\;.
\end{equation}
\end{enumerate}
\end{lemma}

\begin{proof}
(i) By Lemma~\ref{lem:collar-coverage}(i), a bond within $2R_{\mathrm c}$ of $I$ is not in $S_1^*$, where $\Psi_{\xi}$ is the identity. A near primitive has all its inputs within $4L+3\leqslant 2R_{\mathrm c}$ of $I$ (Section~\ref{sec:source-phase-primitives}); an entry $(H_{\xi,IJ})_{bb'}$ with $b\in I$ vanishes unless $d(b,b')\leqslant \rho_L=4L<2R_{\mathrm c}$ (Lemma~\ref{lem:hessian-gauge-locality}), since $R_{\mathrm c}\geqslant 2R_{\mathrm{loc}}+1=4L+5$. The phase on a row $c$ within $2R_{\mathrm c}$ of $I$ depends on the free vector on $\mathcal E_c\cap E_*$ (Lemma~\ref{lem:row-locality}), and every bond of $\mathcal E_c$ is within $2R_{\mathrm c}$ of $I$ by the definition of the row distance (Section~\ref{sec:source-phase-primitives}), so $\Psi_{\xi}$ is the identity on these bonds; on $\mathcal E_c$ only the row $c$ contributes to $h_U\widetilde D_U$ (Lemma~\ref{lem:row-locality}), which gives \eqref{eq:native-collar-source-phase} as an identity between the $\mathcal E_c$-components.

(ii) Put $A^{\Psi}_{\mathrm e}=\Psi_{\xi}(A_{\mathrm e})\in\mathfrak g^J$, so that $\Psi_{\xi}(a+A_{\mathrm e})= a+A^{\Psi}_{\mathrm e}$ by \eqref{eq:c-bondwise-shift}. By Lemma~\ref{lem:primitive-sum},
\begin{equation*}
    \mathscr E_{\xi}(a+A^{\Psi}_{\mathrm e})=\mathscr E_{\xi}(0)-\frac12\langle a,Ha\rangle-\langle a,H_{\xi,IJ}A^{\Psi}_{\mathrm e}\rangle-\frac12\langle A^{\Psi}_{\mathrm e},H_{\xi} A^{\Psi}_{\mathrm e}\rangle+\sum_{\mathfrak p\ \mathrm{near}}\mathfrak p(a+A^{\Psi}_{\mathrm e})+\sum_{\mathfrak p\ \mathrm{far}}\mathfrak p(A^{\Psi}_{\mathrm e})\;,
\end{equation*}
where the far primitives do not depend on $a$ because their free-input sets miss $I$. By (i) the third term is $-j_{\xi}(a)$ and the near sum is $P_{\xi,1}(a)$. The Jacobian factor depends on $(a+A_{\mathrm e})|_{S_1^*}= A_{\mathrm e}|_{S_1^*}$ only. The constant $c_{\xi}$ is finite: $\det D\Psi_{\xi}(A_{\mathrm e})\neq0$ by (B4)(ii) (the differential of (1.22) is invertible on $\mathcal D_{\xi}$, and $D\Psi_{\xi}$ at $a+A_{\mathrm e}$ depends on $A_{\mathrm e}|_{S_1^*}$ only), and the far primitives are finite: their free-input sets miss $I$, so their values at $A^{\Psi}_{\mathrm e}$ are their values at $a_0+A^{\Psi}_{\mathrm e}=\Psi_{\xi}(a_0+A_{\mathrm e})\in\mathcal A_{\xi}$ for any $a_0$ with $(A_{\mathrm e},a_0)\in\mathcal D_{\xi}$ (a branch has one), where every row solution is defined and the primitives are finite by (B4). Collecting the $a$-independent terms gives $\mathcal K_{\xi}= e^{c_{\xi}}K_{\xi;1}$, and the two remaining identities follow by integration, dividing by $\mathcal J^{\mathrm a}_{\xi}[1]= e^{c_{\xi}}I_{\xi;1}>0$.
\end{proof}

The constant $c_{\xi}$ depends on the branch data (through $A_{\mathrm e}$, $U_1$, $V_0$ and $\Omega_1$) and is not a universal constant: it enters the outer weight $\mathcal J^{\mathrm a}_{\xi}[1]= e^{c_{\xi}}I_{\xi;1}$ of the branch in \eqref{eq:native-source-physical-ratio}. It is never estimated separately: Section~\ref{sec:one-step-insertion} compares $\mathcal O[\mathcal J^{\mathrm a}_{\xi}[1]\,F_{\xi}(1)]$ with $\mathcal O[\mathcal J^{\mathrm a}_{\xi}[1]]$, the same nonnegative weight on both sides.

\begin{lemma}[boundary response]\label{lem:boundary-response}
For bounded Borel $f$ on $\mathfrak g^I$, with $\operatorname{Cov}_{\nu}(f,l)=\int fl\,\dd\nu-\int f\,\dd\nu\int l\,\dd\nu$,
\begin{equation}\label{eq:fixed-probe-boundary-functional}
    \int f\,\dd\nu_{\xi;1}=\int f\,\dd\nu_{\xi;0}+\mathcal B_{\xi}[f],\qquad
    \mathcal B_{\xi}[f]=\int_0^1\operatorname{Cov}_{\nu_{\xi;t}}(f,-j_{\xi}+\partial_tP_{\xi,t})\,\dd t\;,
\end{equation}
where $\partial_tP_{\xi,t}(a)=\sum_{\mathfrak p\ \mathrm{near}}D\mathfrak p(a+tA_{\mathrm e})[A_{\mathrm e}]$.
\end{lemma}

\begin{proof}
For $a$ in the box the exponent $E_t(a)=-\frac12\langle a,Ha\rangle-tj_{\xi}(a)+P_{\xi,t}(a)$ is jointly analytic in $(a,t)$ on a neighbourhood of $\overline{\mathcal D_I}\times[0,1]$, since each near primitive is analytic in the free vector on its row polydiscs (Lemma~\ref{lem:charts}) and the input $a+tA_{\mathrm e}$ is affine in $(a,t)$ with values in those polydiscs; its $t$-derivative $-j_{\xi}(a)+\partial_tP_{\xi,t}(a)$, with $\partial_tP_{\xi,t}(a)=\sum_{\mathfrak p}D\mathfrak p(a+tA_{\mathrm e})[A_{\mathrm e}]$ the derivative in the direction $A_{\mathrm e}$, is continuous and hence bounded on $\overline{\mathcal D_I}\times[0,1]$. Hence $t\mapsto\int fK_{\xi;t}\,\dd a$ and $t\mapsto I_{\xi;t}$ are $C^1$ with derivatives $\int f\,\partial_tE_t\,K_{\xi;t}\,\dd a$ and $\int\partial_tE_t\,K_{\xi;t}\,\dd a$ (differentiation under the integral over the bounded box, dominated by a constant times $K_{\xi;t}\leqslant e^{\sup| E_t|}$), so $F(t)=\int f\,\dd\nu_{\xi;t}$ has $F'(t)=\operatorname{Cov}_{\nu_{\xi;t}}(f,\partial_tE_t)$. Integrate over $[0,1]$.
\end{proof}

By Lemmas~\ref{lem:collar} and~\ref{lem:boundary-response}, on every branch
\begin{equation}\label{eq:c-exterior-response}
    \frac{\mathcal J^{\mathrm a}_{\xi}[f_{\xi}]}{\mathcal J^{\mathrm a}_{\xi}[1]}= F_{\xi}(1)= F_{\xi}(0)+\mathcal B_{\xi}[f_{\xi}]\;,
\end{equation}
the first term being the centred expectation of Section~\ref{sec:probe-kernels} and the second the exterior response, estimated in Appendix~\ref{app:exterior-chessboard}.

\begin{lemma}[collar coverage]\label{lem:collar-coverage}
Let $R_{\mathrm{loc}}=2L+2$, which exceeds the $\ell^{\infty}$ diameter $2L-1$ of every two-block set $\mathcal E_c$, hence of the sets $\mathcal E_y$, $\mathcal F_r$, $\mathcal A_r$ and of the free-input and background sets of the block, row, Haar and Jacobian primitives (those of a plaquette primitive have diameter at most $4L<2R_{\mathrm{loc}}$, proof of Lemma~\ref{lem:hessian-gauge-locality}), and let $R_{\mathrm c}\geqslant 2R_{\mathrm{loc}}+1$ be the collar range fixed in Appendix~\ref{app:background-estimates} after \eqref{eq:regular-covariance-radius-replay}; it depends on $L$ and the group only. Under \eqref{eq:c-standing-conditions}, which include $M\geqslant 8LM_2$ and $LMR_1>4R_{\mathrm c}+4L^2M_2R_1+8$, on every branch:
\begin{enumerate}
\item Every site within $\ell^{\infty}$ distance $LMR_1-1$ of $\Lambda_1$ lies in $\Omega_1^{\sim-1}\setminus R_1'\subseteq R_1^c$; $I\subseteq\mathcal G_{\xi}\subseteq\Omega_1^*$; and every free bond $b$ of $\Omega_1$ with $d(b,I)\leqslant 3R_{\mathrm c}$ lies in $\mathcal G_{\xi}\setminus S_1^*$.
\item The cutoff (1.25) composed with the change of variables, as a function of $A'= a+A_{\mathrm e}$, equals $\Theta_{\xi}(a)\chi_{\xi,\mathrm e}(A_{\mathrm e})$ with $\chi_{\xi,\mathrm e}$ of \eqref{eq:retained-cutoff-composition}, and $\chi_{\xi,\mathrm e}$ depends on $A_{\mathrm e}$ and the branch data only.
\item $| A_{\mathrm e}(b)|< r$ for every $b\in J$ with $d(b,I)\leqslant 3R_{\mathrm c}$; that is,
\begin{equation}\label{eq:convex-native-radius}
    r_{\mathrm e,\xi}=\max\bigl(\{| A_{\mathrm e}(b)|\, | \,b\in J,\ d(b,I)\leqslant 3R_{\mathrm c}\}\cup\{0\}\bigr)< r\;.
\end{equation}
\item For an admissible pair and a branch value of $A_{\mathrm e}$, $\{ a\in\mathfrak g^I:(A_{\mathrm e},a)\in\mathcal D_{\xi}\}=\{ a:\Theta_{\xi}(a)=1\}$, the open box of \eqref{eq:native-source-young-cutoff}.
\end{enumerate}
\end{lemma}

\begin{proof}
(i) $\Omega_1$ and $R_1'$ are unions of $LMR_1$-cubes (CRE (1.10), p.251) and $\Lambda_1\subseteq\Omega_1^{\sim-2}\cap(R_1'^{\sim})^c$, so every site within $\ell^{\infty}$ distance $LMR_1-1$ of $\Lambda_1$ lies in $\Omega_1^{\sim-1}\setminus R_1'\subseteq R_1^c$. Let $b$ be a free bond of $\Omega_1$ with $d(b,I)\leqslant 3R_{\mathrm c}$: its initial vertex is within $3R_{\mathrm c}+1$ of a site of $\Lambda_1$, so the cube $\square'$ of side $L^2M_2R_1$ containing it lies within $3R_{\mathrm c}+1+L^2M_2R_1< LMR_1-1$ of $\Lambda_1$ (\eqref{eq:c-standing-conditions}), hence $\square'\subseteq R_1^c$ is a good cube; $b$ has an endpoint in $\square'\subseteq\square'^{\sim2}$ and is not designated, so $b\in(\square'^{\sim2})^*\subseteq\mathcal G_{\xi}$. By \eqref{eq:c-collar-distance}, $d(I,S_1^*)\geqslant LMR_1-2>3R_{\mathrm c}$, so $b\notin S_1^*$. Taking $d(b,I)=0$ gives $I\subseteq\mathcal G_{\xi}$ (a bond of $\Lambda_1$ is a bond of $\Omega_1$, and a designated bond with an endpoint in $\Lambda_1$ belongs to a row of $\Lambda_1$ because $\Lambda_1$ is a union of blocks, so $I\subseteq\Omega_1^*$). For $\mathcal G_{\xi}\subseteq\Omega_1^*$: a good cube $\square'\subseteq\Omega_1^{\sim-1}$ has $\square'^{\sim2}\subseteq\Omega_1$ since $2L^2M_2R_1< LMR_1$ ($M\geqslant 8LM_2$); a bond of $(\square'^{\sim2})^*$ has an endpoint in $\Omega_1$, and a designated bond with an endpoint in $\square'^{\sim2}$, a union of blocks, belongs to a row with an endpoint in $(\square'^{\sim2})^{(1)}$ and is excluded by the definition of $^*$ in (1.8).

(ii) The first two factors of (1.25) are $\prod_{b\in\mathcal G_{\xi}}\mathbf 1_{\{| A(b)|< r\}}$ at $A=\Psi_{\xi}(a+A_{\mathrm e})$. For $b\in I\subseteq\mathcal G_{\xi}$, $A(b)= a(b)$ by \eqref{eq:c-bondwise-shift}, giving $\Theta_{\xi}(a)$; for $b\in\mathcal G_{\xi}\setminus I=\mathcal G_{\xi}\cap J$ (by (i)), $A(b)=(\Psi_{\xi} A_{\mathrm e})(b)$, giving the first factor of \eqref{eq:retained-cutoff-composition}. The factor $\chi^{(0)c}(R_1)$ involves $A$ on $(\square'^{\sim2})^*$ for cubes $\square'\subseteq R_1$, and $\chi_0'(S_1)$ involves the free links $e^{ig_0A(b)}U_1(b)$ on bonds $b\in(\square'^{\sim2})^*\cap\Omega_1^*$ for cubes $\square'\subseteq S_1$, together with $V_0$ on the bonds of $(\square'^{\sim2})^*$ with both endpoints off $\Omega_1$, $U_1$ and $U_{1,\square'}$. All these bonds have both endpoints within $2L^2M_2R_1+1$ of $R_1\cup S_1$, while every site of $\Lambda_1$ is at distance at least $LMR_1>2L^2M_2R_1+1$ from $R_1\cup S_1$; so none of them has an endpoint in $\Lambda_1$, none lies in $I$, and on them $A=\Psi_{\xi}(A_{\mathrm e})$. Hence the last two factors of (1.25), composed with $\Psi_{\xi}$, are functions of $\Psi_{\xi}(A_{\mathrm e})$ and the branch data.

(iii) For $b\in J$ with $d(b,I)\leqslant 3R_{\mathrm c}$, (i) gives $b\in\mathcal G_{\xi}\cap J$ and $b\notin S_1^*$, so $\mathbf 1_{\{|(\Psi_{\xi} A_{\mathrm e})(b)|< r\}}=\mathbf 1_{\{| A_{\mathrm e}(b)|< r\}}$ is a factor of $\chi_{\xi,\mathrm e}(A_{\mathrm e})$, which is nonzero on a branch (the maximum over the empty set is read as $0$).

(iv) Let $a_0$ be such that $(A_{\mathrm e},a_0)\in\mathcal D_{\xi}$ (a branch), so $B_0=\Psi_{\xi}(a_0+A_{\mathrm e})\in\mathcal A_{\xi}$ and, by (ii), $\Theta_{\xi}(a_0)\chi_{\xi,\mathrm e}(A_{\mathrm e})\neq0$. Let $a$ be any point of the open box and $B=\Psi_{\xi}(a+A_{\mathrm e})= a+\Psi_{\xi}(A_{\mathrm e})$ (\eqref{eq:c-bondwise-shift}; $\Psi_{\xi}(A_{\mathrm e})$ is defined, since $B_0$ is). It is to be shown that $B\in\mathcal A_{\xi}$, by checking the three conditions of \eqref{eq:c-free-variable-domain}. The cutoff (1.25) at $B$ is $\Theta_{\xi}(a)\chi_{\xi,\mathrm e}(A_{\mathrm e})=\chi_{\xi,\mathrm e}(A_{\mathrm e})\neq0$, by (ii). Every row equation has a solution at $g_0\mathsf C_UB$: on the near rows by Lemma~\ref{lem:charts}, since $B= a+A_{\mathrm e}$ on their two-block sets by (i) and lies in the row polydiscs by (iii); on the far rows the argument $(g_0\mathsf C_UB)|_{\mathcal E_c}$ equals $(g_0\mathsf C_UB_0)|_{\mathcal E_c}$ (Lemma~\ref{lem:row-locality}: it involves $B$ only through $B|_{\mathcal E_c\cap E_*}$, and $\mathcal E_c$ misses $I$ for a far row, every bond of which is at distance more than $2R_{\mathrm c}+2L-(2L-1)>0$ from $I$), where the solution exists because $B_0\in\mathcal A_{\xi}$. So $U(B)$ is defined, and on the far rows and the bonds of $\Omega_1^c$ its links are those of $U(B_0)$. It remains to show $U(B)\in D^{(11)}(V,V_0)$. The restriction to $\Omega_1^c$ is $V_0$. The block averages are defined on every row of $\Omega_1$ and satisfy the constraint $\bar U= V$ there: $U(B)= e^{iq_U(B)}U_1$ with $q_U(B)= A-h_U\widetilde D_U(A)$, $A= g_0\mathsf C_UB$, and on every row $c$ the row equation gives $\mathcal Q_{U,c}(q_U(B))= LQ_Uq_U(B)+\mathcal N_{U,c}(q_U(B))= LQ_U(A-h_U\widetilde D_U(A))+\widetilde D_{U,c}(A)=-\widetilde D_{U,c}(A)+\widetilde D_{U,c}(A)=0$, using $Q_U\mathsf C_U=0$ and $LQ_Uh_U=\id$; so $\mathsf M(U(B))(c)=\mathsf M(U_1)(c)= V(c)$ by \eqref{eq:native-source-constraint-normalisation} and (B1). The constraint $\bar V_0= V$ on the rows of $\Omega_1^c$ involves $V_0$ and $V$ only and holds because the pair is admissible. The tests of (1.11): $\chi_1(\Omega_1)$ is a function of $V$, $\zeta(\Omega_1^c)$ does not depend on $a$ (\eqref{eq:c-large-field-factor}), and both are nonzero at $U(B_0)\in D^{(11)}(V,V_0)$. The factor $\chi_0'(\Omega_1)$ is the product of $\chi_0'(\square')$ over the $L^2M_2R_1$-cubes $\square'\subseteq\Omega_1$ (CRE (1.8)); the cubes of $S_1=(\Omega_1\setminus\Omega_1^{\sim-1})\cup R_1$ (CRE (1.20)) give the factor $\chi_0'(S_1)$ of (1.25), nonzero at $B$ by the first condition, and every other cube $\square'\subseteq\Omega_1$ lies in $\Omega_1^{\sim-1}\setminus R_1$, so it is a good cube with $\square'^{\sim4}\subseteq\Omega_1$ (as $4L^2M_2R_1< LMR_1$), and its factor $\chi_0'(\square')$ is the second factor of (1.19), equal to one at $B$ by (B3)(ii), the retained factors being nonzero there. Hence $U(B)\in D^{(11)}(V,V_0)$, so $B\in\mathcal A_{\xi}$ and $(A_{\mathrm e},a)\in\mathcal D_{\xi}$. Conversely $(A_{\mathrm e},a)\in\mathcal D_{\xi}$ gives $\Psi_{\xi}(a+A_{\mathrm e})\in\mathcal A_{\xi}$, so the cutoff (1.25) is nonzero there, and $\Theta_{\xi}(a)=1$ by (ii).
\end{proof}

The exterior response has Ba{\l}aban's integration-by-parts form (CRE (3.28)). Put
\begin{equation}\label{eq:native-collar-mean-direction}
    \ell_{\xi}= H_{\xi,IJ}A_{\mathrm e},\qquad m_{\xi}= H_{\xi,II}^{-1}\ell_{\xi},\qquad D_{m_{\xi}}f(a)=\frac{\dd}{\dd s}f(a+sm_{\xi})\Big|_{s=0}\;,
\end{equation}
let $\mathcal D_I$ be the open box of \eqref{eq:native-source-boundary-interpolation}, with outward unit normal $n(a)$ and surface measure $\dd S$ for the Euclidean structure of $\langle\bigcdot ,\bigcdot \rangle$ on $\mathfrak g^I$ (for which $\dd a$ is Lebesgue measure), defined almost everywhere on its Lipschitz boundary (a product of balls of the norm $|\bigcdot |$), and
\begin{equation}\label{eq:native-collar-flux}
    \mathsf Q_{\xi,t}(a)=-\frac{1}{2}\langle a,H_{\xi,II}a\rangle-t\langle\ell_{\xi},a\rangle+P_{\xi,t}(a),\qquad
    \mathcal T_{\xi,t}[f]=\frac{1}{I_{\xi;t}}\int_{\partial\mathcal D_I}f(a)e^{\mathsf Q_{\xi,t}(a)}\langle m_{\xi},n(a)\rangle\,\dd S(a)\;,
\end{equation}
so that $K_{\xi;t}=\mathbf 1_{\mathcal D_I}e^{\mathsf Q_{\xi,t}}$ and $j_{\xi}(a)=\langle\ell_{\xi},a\rangle$.

\begin{lemma}[integration by parts with the hard cutoff]\label{lem:ibp}
Assume Lemma~\ref{lem:regular-bounds}, so that $H_{\xi,II}$ is invertible and $m_{\xi}$ is defined. If $f$ is $C^1$ on a neighbourhood of $\overline{\mathcal D_I}$, then for every $t\in[0,1]$, with expectations under $\nu_{\xi;t}$,
\begin{equation}\label{eq:native-collar-source-ibp}
    \operatorname{Cov}(f,-j_{\xi}+\partial_tP_{\xi,t})
        =-\mathbb E[D_{m_{\xi}}f]+\operatorname{Cov}(f,\partial_tP_{\xi,t}-D_{m_{\xi}}P_{\xi,t})+\mathcal T_{\xi,t}[f]-\mathbb E[f]\mathcal T_{\xi,t}[1]\;.
\end{equation}
\end{lemma}

\begin{proof}
$e^{\mathsf Q_{\xi,t}}$ is $C^1$ on the enlarged box, a neighbourhood of $\overline{\mathcal D_I}$, by Lemma~\ref{lem:charts}. The divergence theorem on the bounded Lipschitz domain $\mathcal D_I$ gives $\int_{\mathcal D_I}D_{m_{\xi}}(fe^{\mathsf Q_{\xi,t}})\,\dd a= I_{\xi;t}\mathcal T_{\xi,t}[f]$, and $D_{m_{\xi}}\mathsf Q_{\xi,t}=-j_{\xi}-t\langle m_{\xi},\ell_{\xi}\rangle+D_{m_{\xi}}P_{\xi,t}$ because $\langle m_{\xi},H_{\xi,II}a\rangle=\langle\ell_{\xi},a\rangle= j_{\xi}(a)$. Dividing by $I_{\xi;t}$, $\mathbb E[D_{m_{\xi}}f]+\mathbb E[f(-j_{\xi}-t\langle m_{\xi},\ell_{\xi}\rangle+D_{m_{\xi}}P_{\xi,t})]=\mathcal T_{\xi,t}[f]$. Subtract the same identity for $f=1$ multiplied by $\mathbb E[f]$; the constant $-t\langle m_{\xi},\ell_{\xi}\rangle$ cancels, and the two covariances combine as stated.
\end{proof}

Since $f_{\xi}$ is analytic on the enlarged box (after \eqref{eq:convex-response-operands}), Lemma~\ref{lem:ibp} applies to $f_{\xi}$. The estimate of the right-hand side is Lemma~\ref{lem:exterior-response-bound} in Appendix~\ref{app:exterior-chessboard}.

\subsection{A fixed walk expansion of the centred kernel}
\label{sec:fixed-first-kernel}

Write $H= H_{\xi,II}$ and $\mathsf G= H^{-1}$. Notation of Sections~\ref{sec:fixed-first-kernel}--\ref{sec:joint-extraction}: $E_0= E(Z_0)$ denotes the seed bonds (the designated bonds $E_0$ of \eqref{eq:constraint-extension} do not occur here); $N_i(Y)$ is a cube count (the measure $N_1$ of \eqref{eq:first-step} does not occur); $\sigma_{\Delta}$ are weakening variables (the Haar density $\sigma$ does not occur); $B$ is the conditioned variable on $\mathfrak g^{E_0}$ (the free vector $B$ of Sections~\ref{sec:regular-hessian-stencils}--\ref{sec:boundary-law-collar} does not occur); $D$ without argument is a component; $J_D$, $J_Z$ are sets of young bonds (the retained carrier $J$ does not occur); $\delta_M$ is a decay constant (the Mayer constant $\delta$ of Appendix~\ref{app:background-estimates} does not occur); $\gamma=\gamma_L/2$ is a spectral bound (arcs $\gamma$ of the contour do not occur); a connected final domain is any nonempty wall-connected union of $LM$-cubes; and $a_j$, $a_*$, $a_W$ are coefficients and weights, distinct from the young variable $a$. By Lemma~\ref{lem:regular-bounds}, on every branch
\begin{equation}\label{eq:fixed-kernel-inputs}
    \gamma\id\leqslant H\leqslant h_+\id,\qquad\gamma=\frac{\gamma_L}{2},\quad h_+= m_L\;,
\end{equation}
and by Lemma~\ref{lem:hessian-gauge-locality}, whose entries $\mathsf H^{\Omega_1}_{U_1}(b,c)$ at $b,c\in I$ are those of $H$, $H(b,c)=0$ when $d_{\infty,T}(b_-,c_-)> R$, where $R\geqslant \max(\rho_L,2)$ is fixed so that $H(b,c)$ is determined by the background, by the region $\Omega_1$ and by membership in $I$ within distance $R$ of $b_-$. Norms on $\mathfrak g^I$ are those of the invariant real inner product, with Hermitian extension on the complexification. The constants $\gamma,h_+,R$ depend on $L$ and the group only (with the invariant inner product fixed in the conventions), and every choice below is made before the branch fields and the coupling, in the order $\gamma,h_+,R$; then $\kappa_1$; then $n_*,R_*,C_{\mathrm w}$; then $M_1$, $\rho_{\mathrm w}= R_*+R+1+2M_1$ and $M$; then $g_0$. None of $\gamma,h_+,R,\kappa_1,C_{\mathrm d}$ depends on $M_1$, $M$ or the coupling ($C_{\mathrm d}$ of Lemma~\ref{lem:gaussian-payment} is a function of $\gamma,h_+,d_G$).

Fix an aligned auxiliary cubical partition of side $M_1$ in fine units, with cubes $M_1z+[-M_1/2,M_1/2]^4$; the $M$-cube partition $\pi_0$ and the $LM$-cube partition $\pi_1$ are nested over it, $M/M_1$ and $L$ are odd integers, every fine period is a multiple of $LMR_1$ and at least $8M$, and $M_1$ and $M$ are odd (\eqref{eq:c-standing-conditions}: $M_1$, $M_2$, $M$ are powers of the odd number $L$), so that cube faces lie at half-integer coordinates and no lattice site lies on a face. Cubes and neighbourhoods are periodic images; all radii used below are smaller than $M$, so the periods $\geqslant 8M$ make every neighbourhood a lifted copy of its $\mathbb Z^4$ model. Set $s_H=2h_+$, $\mathsf A_H=\id-H/s_H$, $q_H=1-\gamma/s_H<1$ and $V_n=4(2nR+3)^4$. A block of $\mathsf A_H^n$ vanishes beyond distance $nR$, has operator norm at most $q_H^n$ and depends only on coefficients within distance $(n+1)R$ of its first bond, so its block-norm row and column sums are at most $V_nq_H^n$. Let $\kappa_1\geqslant 1$ be the constant fixed in Appendix~\ref{app:background-estimates}, Section~\ref{app:mayer-constants} (\eqref{eq:c2-mayer-constants}; it depends on $L$ and the group only), the spatial Cauchy radius being $e^{\kappa_1}$; put $a_*= e^{81\kappa_1}$ and choose $n_*$ with
\begin{equation}\label{eq:fixed-kernel-grouping}
    \beta= V_{n_*}q_H^{n_*},\qquad a_*\beta\leqslant \frac{1}{4}\;.
\end{equation}
Put $R_*= n_*R$, $\mathsf P_H=\mathsf A_H^{n_*}$, $V_*= V_{n_*}$ and $a_j=4^{-j}\binom{2j}{j}\in(0,1]$. The binomial series gives
\begin{equation}\label{eq:fixed-kernel-square-root-series}
    \mathsf G^{1/2}= s_H^{-1/2}\sum_{r=0}^{n_*-1}\sum_{m=0}^{\infty} a_{n_*m+r}\mathsf A_H^r\mathsf P_H^m\;.
\end{equation}
For a free bond $b$ let $D_b$ be the union of the $M_1$-cubes meeting the closed ball of radius $R_*+R+1$ about $b_-$, with one further $M_1$-layer; it is connected, lies in the ball of radius $\rho_{\mathrm w}= R_*+R+1+2M_1$ about $b_-$, and contains the bonds and coefficient stencils of every block starting at $b$ in \eqref{eq:fixed-kernel-square-root-series} with an interior margin. Choose
\begin{equation}\label{eq:fixed-kernel-collar-size}
    M\geqslant 4(\rho_{\mathrm w}+R+R_*+2)\;.
\end{equation}
Then $2\rho_{\mathrm w}< M\leqslant LM$, so $D_b$ meets at most $2^4=16\leqslant 81$ cubes of $\pi_1$. A walk is a tuple $\omega=(r;b_0,\ldots,b_{m+1})$ with nonzero block product
\begin{equation}\label{eq:fixed-kernel-walk-block}
    (\mathsf A_H^r)_{b_0b_1}\prod_{j=1}^m(\mathsf P_H)_{b_jb_{j+1}}\;,
\end{equation}
and localisation domain $D(\omega)=\bigcup_{j\leqslant m}D_{b_j}$, which is connected and contains both endpoint bonds in its interior.

For an $M$-cube $Q$ let $Q^{\sim}$ be the union of the $3^4$ cubes of $\pi_0$ meeting $Q$, so that $\operatorname{int}Q^{\sim}$ is the open $\ell^{\infty}$ neighbourhood of radius $M$ of $Q$. For a union $Z_0$ of $M$-cubes (the seed, selected in Section~\ref{sec:selected-integrals}) let $\widetilde Z_0=\bigcup_{Q\subseteq Z_0}Q^{\sim}$ be its one-$M$-layer enlargement and $Z_0'$ the union of the $LM$-cubes meeting $\operatorname{int}Q^{\sim}$ for some cube $Q$ of $Z_0$. Both are unions over the cubes of $Z_0$, so both commute with unions of seeds, and $\operatorname{int}Q^{\sim}\subseteq\operatorname{int}Z_0'$ for every cube $Q$ of $Z_0$, since every $LM$-cube incident to a point of the open set $\operatorname{int}Q^{\sim}$ meets it. Assign $\sigma_{\Delta}=1$ to the cubes $\Delta\in\pi_1$ inside $Z_0'$ and complex variables $\sigma_{\Delta}$ to the others, and weight each walk by
\begin{equation}\label{eq:fixed-kernel-walk-weakening}
    w_{\omega}(\sigma)=\prod_{\Delta\in\pi_1,\ \Delta\not\subseteq Z_0',\ \Delta\cap D(\omega)\ne\varnothing}\sigma_{\Delta}\;.
\end{equation}
The weighted series is
\begin{equation}\label{eq:c-weighted-walk-series}
    S_{Z_0}(\sigma)_{b_0b_{m+1}}= s_H^{-1/2}\sum_{r=0}^{n_*-1}\sum_{m=0}^{\infty} a_{n_*m+r}\sum_{b_1,\ldots,b_m}w_{\omega}(\sigma)\,(\mathsf A_H^r)_{b_0b_1}\prod_{j=1}^m(\mathsf P_H)_{b_jb_{j+1}},\qquad\omega=(r;b_0,\ldots,b_{m+1})\;,
\end{equation}
the inner sum running over the intermediate bonds; $S_{Z_0}(1)=\mathsf G^{1/2}$ by \eqref{eq:fixed-kernel-square-root-series}.

\begin{lemma}[the walk series]\label{lem:walk-series}
On $|\sigma_{\Delta}|\leqslant e^{\kappa_1}$, $| w_{\omega}(\sigma)|\leqslant a_*^{m+1}$, and the series $S_{Z_0}(\sigma)$ converges absolutely in both Schur norms, uniformly on the polydisc, with operator norm and absolute block row and column sums at most
\begin{equation}\label{eq:fixed-kernel-schur-bound}
    C_{\mathrm w}= s_H^{-1/2}\frac{n_*V_*a_*}{1-a_*\beta}\;.
\end{equation}
It is holomorphic on the polydisc of radii $e^{\kappa_1+(\log2)/81}$. The bound is independent of $Z_0$, the volume and the number of its components. Under a site gauge transformation $u$ of the background, $S_{Z_0}(\sigma)$ formed from $\iota_I^*\mathsf H^{\Omega_1}_{U_1^u}\iota_I$ in place of $H$ equals $R_{u,*}S_{Z_0}(\sigma)R_{u,*}^{-1}$, with $R_{u,*}$ restricted to $I$ and the same domains and weights.
\end{lemma}

\begin{proof}
$D(\omega)$ meets at most $81(m+1)$ cubes of $\pi_1$, which gives the weight bound. The unweakened grouped products have absolute block row and column sums at most $V_*\beta^m$ after summing the intermediate bonds, so the weighted series is dominated by $s_H^{-1/2}n_*V_*\sum_m(a_*\beta)^ma_*$. Increasing each radius by the factor $e^{(\log2)/81}$ multiplies $a_*$ by $2$ and keeps $a_*\beta\leqslant 1/2$. The gauge law is \eqref{eq:regular-gauge-constrained-hessian} restricted to $I$ (so $\mathsf A_H$ is conjugated by $R_{u,*}$) applied to each block; the weights and domains do not involve the background.
\end{proof}

Put $E_0= E(Z_0)=\{ b\in I\, | \,b\cap Z_0\ne\varnothing\}$, let $\iota_{E_0}$ be extension by zero and $\Pi_{I\setminus E_0}$ the complementary orthogonal projection, and set
\begin{equation}\label{eq:fixed-kernel-C-Gamma}
    C_{E_0}=(\iota_{E_0}^*H\iota_{E_0})^{-1},\qquad\Gamma_{E_0}(\sigma)=\iota_{E_0}^*H\Pi_{I\setminus E_0}S_{Z_0}(\sigma)\;.
\end{equation}
Pairings $\langle\bigcdot ,\bigcdot \rangle$ involving complex vectors are the complex-bilinear extensions of the real pairing throughout.

\begin{lemma}[Gaussian conditioning]\label{lem:conditioning}
Let $\mu_0$ be the standard Gaussian measure on $\mathfrak g^I$ and $\mu_{C_{E_0}}$ the centred Gaussian measure on $\mathfrak g^{E_0}$ of covariance $C_{E_0}$. For bounded Borel $F$ on $\mathfrak g^{E_0}$,
\begin{equation}\label{eq:c-conditioning}
    \int F(a|_{E_0})\,\dd\mu_{\mathsf G}(a)=\int\left[\int F(B)\,e^{-\frac12\langle\Gamma_{E_0}(1)X,C_{E_0}\Gamma_{E_0}(1)X\rangle}e^{-\langle B,\Gamma_{E_0}(1)X\rangle}\,\dd\mu_{C_{E_0}}(B)\right]\dd\mu_0(X)\;.
\end{equation}
\end{lemma}

\begin{proof}
Write $a=\mathsf G^{1/2}X$, so that $a$ has law $\mu_{\mathsf G}$ when $X$ has law $\mu_0$. Given $a|_{I\setminus E_0}$, the conditional law of $a|_{E_0}$ under $\mu_{\mathsf G}$ is Gaussian with covariance $C_{E_0}$ and mean $-C_{E_0}\iota_{E_0}^*H\Pi_{I\setminus E_0}a=-C_{E_0}\Gamma_{E_0}(1)X$ (complete the square in $\frac12\langle a,Ha\rangle$), and its density with respect to $\mu_{C_{E_0}}$ is the displayed factor. This is RG2 (2.5)--(2.6).
\end{proof}

The weakened density is
\begin{equation}\label{eq:fixed-kernel-density}
    e^{-\frac{1}{2}\langle\Gamma_{E_0}(\sigma)X,C_{E_0}\Gamma_{E_0}(\sigma)X\rangle}e^{-\langle B,\Gamma_{E_0}(\sigma)X\rangle}\,\dd\mu_{C_{E_0}}(B)\,\dd\mu_0(X)\;.
\end{equation}
$C_{E_0}$ is real, positive and independent of $\sigma$; its determinant is one fixed positive number on the whole polydisc.

\begin{lemma}[component blocks]\label{lem:component-blocks}
Let $Z\supseteq Z_0'$ be a union of $LM$-cubes and set $\sigma_{\Delta}=0$ for $\Delta\not\subseteq Z$. For each wall-component $D$ of $Z$ put $J_D=\{ b\in I\, | \,b\subset\operatorname{int}D\}$ and $E_D= E_0\cap J_D$. Then $E_0\subseteq\bigsqcup_DJ_D$, $E_D= E(Z_0\cap D)$, initial vertices in different $E_D$ have distance at least $2M-2> R$, and
\begin{equation}\label{eq:fixed-kernel-exact-blocks}
    C_{E_0}=\bigoplus_D(\iota_{E_D}^*H\iota_{E_D})^{-1},\qquad\Gamma_{E_0}(\sigma)=\bigoplus_D\Gamma_{E_D}(\sigma|_D)\;,
\end{equation}
where $\Gamma_{E_D}(\sigma|_D)$ is constructed from the seed $Z_0\cap D$ alone, with final domain $D$ (so $\sigma_{\Delta}=0$ for every cube $\Delta$ of $\pi_1$ not in $D$) and the same $I$, $H$ and background. Columns of $\Gamma_{E_0}(\sigma)$ outside $\bigsqcup_DJ_D$ vanish, and $\Gamma_{E_0}(\sigma)$ does not depend on $\sigma_{\Delta}$ for cubes $\Delta$ in a component containing no seed cube.
\end{lemma}

\begin{proof}
A point is interior to a union of closed grid cubes exactly when all its incident cubes belong to it, and the incident family is wall connected (any two cubes containing a point are joined by a chain of cubes containing that point, consecutive ones sharing a wall); hence $\operatorname{int}Z=\bigsqcup_D\operatorname{int}D$, a disjoint union of open sets. For an $M$-cube $Q$ of $Z_0$ the open neighbourhood $\operatorname{int}Q^{\sim}$ is connected and lies in $\operatorname{int}Z_0'\subseteq\operatorname{int}Z$, so it lies in one component interior; neighbourhoods in different components are disjoint, so their seed cubes have distance at least $2M$, and a bond meeting a seed cube lies in the same component interior. This gives $E_0\subseteq\bigsqcup J_D$, $E_D= E(Z_0\cap D)$ and the separation, which by finite range makes $\iota_{E_0}^*H\iota_{E_0}$ block diagonal. Consider an entry $(\Gamma_{E_0}(\sigma))_{eb}=\sum_{b'}H_{eb'}(\Pi_{I\setminus E_0})_{b'b'}S_{Z_0}(\sigma)_{b'b}$ with $e\in E_D$. A nonzero term has $d(e_-,b'_-)\leqslant R$; $e$ meets a seed cube $Q$, so $e_-$ is within $1$ of $Q$ and the open ball of radius $R+2$ about $e_-$ lies in the open $M$-neighbourhood $\operatorname{int}Q^{\sim}$ of $Q$, because $R+3\leqslant M$ by \eqref{eq:fixed-kernel-collar-size}, hence in $\operatorname{int}D$; so $b'\subset\operatorname{int}D$. A walk from $b'$ to $b$ with $w_{\omega}\ne0$ has every cube of $\pi_1$ meeting $D(\omega)$ inside $Z$, so $D(\omega)\subset\operatorname{int}Z$; $D(\omega)$ is connected and contains $b'$, so $D(\omega)\subset\operatorname{int}D$ and $b\in J_D$. Its coefficients, domains and weakening variables are independent of the seed; the cubes with $\sigma_{\Delta}=1$ inside $D$ are those of $Z_0'\cap D=(Z_0\cap D)'$ ($Z_0'$ is the union over the seed cubes $Q$ of the $LM$-cubes meeting $\operatorname{int}Q^{\sim}$; such a cube meets $\operatorname{int}Z$ at a point whose incident $LM$-cubes all lie in $Z$ and form a wall-connected family, so it lies in the component containing $Q$); and $\Pi_{I\setminus E_0}$ agrees with $\Pi_{I\setminus E_D}$ on $J_D$. This gives the block decomposition and the vanishing of the columns outside $\bigsqcup J_D$. Finally, a walk contributing to a row $e\in E_D$ has $D(\omega)\subset\operatorname{int}D$, so its weight involves only $\sigma_{\Delta}$ with $\Delta\subseteq D$; a component without seed cube has no row at all.
\end{proof}

The estimates on this kernel are collected in one statement.

\begin{lemma}[Gaussian payment]\label{lem:gaussian-payment}
Let $N_M=\lfloor M/(2R_*)\rfloor$, $\delta_M=2C_{\mathrm w}4^{-N_M}$ and $A_*= h_+(\gamma^{-1/2}+1)$; by \eqref{eq:c-standing-conditions}, $\delta_M\leqslant 1$ and $h_+^2\delta_M^2/\gamma\leqslant 1/8$. Write $\pi_i(X)$ for the set of $\pi_i$-cubes of a union $X$ of such cubes. On every branch, for every seed $Z_0$, every union $Z\supseteq Z_0'$ of $LM$-cubes with $\sigma_{\Delta}=0$ for $\Delta\not\subseteq Z$, and every $\sigma$ in the closed polydisc $\{|\sigma_{\Delta}|\leqslant e^{\kappa_1}\ \forall\Delta\in\pi_1(Z)\setminus\pi_1(Z_0')\}$:
\begin{enumerate}
\item $\|\Gamma_{E_0}(\sigma)-\Gamma_{E_0}(1)\|\leqslant h_+\delta_M$ and $\| C_{E_0}^{1/2}(\Gamma_{E_0}(\sigma)-\Gamma_{E_0}(1))\|\leqslant h_+\gamma^{-1/2}\delta_M$, in the operator norms for the Hilbert norms on $\mathfrak g^I_{\mathbb C}$ and $\mathfrak g^{E_0}_{\mathbb C}$.
\item For $0\leqslant \alpha_G\leqslant \min(\gamma/2,\gamma^2/(16A_*^2))$ and $J_Z=\bigsqcup_DJ_D$, the integral of the absolute value of \eqref{eq:fixed-kernel-density} against $e^{\alpha_G\| B\|^2/2}$ is at most $\exp(C_{\mathrm d}(\alpha_G+\delta_M^2)| J_Z|)$, $C_{\mathrm d}= C_{\mathrm d}(\gamma,h_+,d_G)$; with $4C_{\mathrm d}\delta_M^2(LM+1)^4\leqslant 1/2$ (\eqref{eq:c-standing-conditions}) and $\alpha_G\leqslant (8C_{\mathrm d}(LM+1)^4)^{-1}$ this is at most $e^{N_1(Z)}$, and the integral of the absolute value factors over the components $D$ of $Z$ (the signed density factors by Lemma~\ref{lem:component-blocks}).
\item For $0<\gamma_2$ and the constant $c_{\mathrm n}>0$ of the group with $h(X)^2\geqslant c_{\mathrm n}| X|^2$ on $\mathfrak g$, $\prod_{b\in\mathsf P}\mathbf 1_{\{| B_b|\geqslant r\}}\leqslant e^{-a_{\mathrm{bad}}|\mathsf P|}\exp(\frac{\gamma_2}{2}\sum_{b\in\mathsf P}h(B_b)^2)$ with $a_{\mathrm{bad}}= c_{\mathrm n}\gamma_2r^2/2$, for every $\mathsf P\subseteq E_0$ (a pointwise inequality, used in the appendix with $\gamma_2\leqslant \alpha_G$).
\end{enumerate}
\end{lemma}

The proof is Appendix~\ref{app:background-estimates}, Section~\ref{app:fixed-kernel-payment}.

\subsection{The selected integrals and the polymer representation}
\label{sec:selected-integrals}

Connectivity of cubes is wall connectivity (RG1 before (0.24)); two domains are compatible when they share neither a cube nor a wall (RG2 (2.11)); $N_i(Y)$ is the number of $\pi_i$-cubes of $Y$, $d_i(Y)$ its tree length (RG1 p.257, written out in \eqref{eq:c2-tree-length}) and $\ell_i(Y)= d_i(Y)+5$. The localisation family $\mathbf D_0$ consists of the potential hosts \eqref{eq:part1-potential-host} and the probe hosts \eqref{eq:c-probe-host} below; all are connected unions of $M$-cubes. A potential family indexed by $\mathbf D_0$ is a collection $\{\mathbf W(Y,\bigcdot )\}_{Y\in\mathbf D_0}$ of Borel functions $\mathfrak g^I\to\mathbb C$, each depending on $a$ only through $a|_{E(Y)}$, $E(Y)=\{ b\in I\, | \,b\cap Y\ne\varnothing\}$, and bounded on the local box $\{| a_b|< r\ \forall b\in E(Y)\}$; only the values on the local boxes enter below, because every selected host $Y$ comes with the factors $\chi_b$, $b\in E(Y)\subseteq E(Y_0)$. The potential $\mathbf V_{\xi}$ of \eqref{eq:c-potential-definition} is such a family; $\mathbf V_{\xi}(Y,\bigcdot )=0$ on hosts carrying no near primitive.

\begin{definition}[selections and selected integrals]\label{def:selection}
A selection is a triple $\mathfrak s=(\mathcal D,\mathsf P,Z)$: a finite $\mathcal D\subseteq\mathbf D_0$, a finite $\mathsf P\subseteq I\setminus E(Y_0)$ where $Y_0=\bigcup_{Y\in\mathcal D}Y$, and a final domain $Z$. Its seed is $Z_0= Y_0\cup\bigcup\{ Q\in\pi_0\, | \,Q\cap b\ne\varnothing\text{ for some }b\in\mathsf P\}$, with $E_0= E(Z_0)$ and $Z_0'$ as in Section~\ref{sec:fixed-first-kernel}, and a final domain is a union $Z\supseteq Z_0'$ of $LM$-cubes. The selected integral of a bounded Borel $F$ on $\mathfrak g^{E_0}$ is
\begin{equation}\label{eq:c-selected-integral}
    \begin{aligned}
    \mathcal A^{\mathrm f}_{\mathfrak s}[F]
        &=\prod_{\Delta\in\pi_1(Z)\setminus\pi_1(Z_0')}\int_0^1\dd s_{\Delta}\oint_{|\sigma_{\Delta}|= e^{\kappa_1}}\frac{\dd\sigma_{\Delta}}{2\pi i(\sigma_{\Delta}-s_{\Delta})^2}
          \prod_{Y\in\mathcal D}\int_0^1\dd t_Y\int\dd\mu_0(X)\int\dd\mu_{C_{E_0}}(B)\\
        &\quad\cdot e^{-\frac12\langle\Gamma_{E_0}(\sigma)X,C_{E_0}\Gamma_{E_0}(\sigma)X\rangle}e^{-\langle B,\Gamma_{E_0}(\sigma)X\rangle}
          (-1)^{|\mathsf P|}\prod_{b\in E(Y_0)}\chi_b(B)\prod_{b\in\mathsf P}(1-\chi_b(B))
          \prod_{Y\in\mathcal D}\mathbf W(Y,B)e^{t_Y\mathbf W(Y,B)}\,F(B)\;,
    \end{aligned}
\end{equation}
with $\chi_b(B)=\mathbf 1_{\{| B_b|< r\}}$ in CRE's bond norm (the indicator of \eqref{eq:native-source-young-cutoff}) and $\sigma_{\Delta}=0$ for $\Delta\not\subseteq Z$. The vacuum activity of a connected final domain $Z$ is
\begin{equation}\label{eq:c-vacuum-activity}
    k^{\mathrm f}_{\xi}(Z)=\sum_{(\mathcal D,\mathsf P)\, | \,Z_0\ne\varnothing,\ Z_0'\subseteq Z}\mathcal A^{\mathrm f}_{(\mathcal D,\mathsf P,Z)}[1]\;,
\end{equation}
and $\Xi^{\mathrm f}_{\xi}=\sum_{\mathcal P\text{ compatible}}\prod_{Z\in\mathcal P}k^{\mathrm f}_{\xi}(Z)$ (empty family included), $(\Xi^{\mathrm f}_{\xi})^{[Y]}$ the same sum over families compatible with $Y$.
\end{definition}

The selected integral is finite: the integrand is holomorphic in $\sigma$ on the polydisc of Lemma~\ref{lem:walk-series} for fixed $(X,B,t)$, its absolute value is integrable in $(X,B)$ uniformly on the closed polydisc by Lemma~\ref{lem:gaussian-payment}(ii) together with the boundedness of $F$ and of each $\mathbf W(Y,\bigcdot )$ on its local box (enforced by the factors $\chi_b$, $b\in E(Y_0)$), and the contours have positive distance from $[0,1]$. Every selected datum depends only on $B$, since $E(Y)\subseteq E_0$ for $Y\in\mathcal D$ and $\mathsf P\subseteq E_0$. A fine bond meets at most two $M$-cubes (no site lies on a face), so at most $2|\mathsf P|$ cubes of the seed are not cubes of $Y_0$. The seed of a union of selections (disjoint host families and bad-bond sets) is the union of the seeds, and $\widetilde Z_0$, $Z_0'$ are unions over the seed cubes (Section~\ref{sec:fixed-first-kernel}), so all three commute with unions.

\begin{lemma}[polymer representation]\label{lem:polymer}
Let $\mathbf W$ be a potential family indexed by $\mathbf D_0$ and $\mathfrak n_{\xi}=(2\pi)^{d_G| I|/2}\det(H)^{-1/2}$. Then
\begin{equation}\label{eq:c-polymer}
    \int\Theta_{\xi}(a)\exp\Bigl(-\frac12\langle a,Ha\rangle+\sum_{Y\in\mathbf D_0}\mathbf W(Y,a)\Bigr)\dd a=\mathfrak n_{\xi}\,\Xi^{\mathrm f}_{\xi}[\mathbf W]\;,
\end{equation}
with the activities of Definition~\ref{def:selection} formed from $\mathbf W$, and the sum $\Xi^{\mathrm f}_{\xi}[\mathbf W]$ is finite (finitely many selections at fixed volume). In particular, for $\mathbf W=\mathbf V_{\xi}$,
\begin{equation}\label{eq:fixed-kernel-source-numerator}
    I_{\xi;0}=\mathfrak n_{\xi}\,\Xi^{\mathrm f}_{\xi}\;,\qquad\Xi^{\mathrm f}_{\xi}=\frac{I_{\xi;0}}{\mathfrak n_{\xi}}>0\;.
\end{equation}
\end{lemma}

\begin{proof}
The left side is $\mathfrak n_{\xi}\int\Theta_{\xi} e^{\sum_Y\mathbf W(Y,\bigcdot )}\dd\mu_{\mathsf G}$. Expand each $e^{\mathbf W(Y)}=1+\int_0^1\mathbf W(Y)e^{t_Y\mathbf W(Y)}\dd t_Y$; this gives the sum over $\mathcal D$. Write $\Theta_{\xi}=\prod_{E(Y_0)}\chi_b\prod_{I\setminus E(Y_0)}(1-(1-\chi_b))$ and expand the second product; this gives the sum over $\mathsf P$ with the sign $(-1)^{|\mathsf P|}$. The remaining integrand depends on $a|_{E_0}$ only (each $\mathbf W(Y,\bigcdot )$, $Y\in\mathcal D$, through $a|_{E(Y)}$ with $E(Y)\subseteq E_0$, and it is evaluated on the local box of $Y$ because of the factors $\chi_b$, $b\in E(Y_0)$), so Lemma~\ref{lem:conditioning} applies and produces \eqref{eq:fixed-kernel-density} at $\sigma=1$. For fixed $(X,B,t)$ the density \eqref{eq:fixed-kernel-density} is holomorphic in $\sigma$ on the polydisc (Lemma~\ref{lem:walk-series}); for each $\Delta\in\pi_1$ not in $Z_0'$ write its value at $\sigma_{\Delta}=1$ as the value at $\sigma_{\Delta}=0$ plus $\int_0^1\partial_{s_{\Delta}}\dd s_{\Delta}$, expand the product over these cubes, label each term by $Z= Z_0'\cup\{\text{differentiated cubes}\}$, and represent each derivative by its Cauchy integral along $|\sigma_{\Delta}|= e^{\kappa_1}$; these are pointwise identities in $(X,B,t)$, and the resulting parameter integrals may be taken outside the Gaussian integrals by Fubini, since the absolute value of the density is integrable uniformly on the polydisc (Lemma~\ref{lem:gaussian-payment}(ii)) and the other factors are bounded. This is \eqref{eq:c-selected-integral}; every step is an exact identity at fixed volume. The infinite objects are the walk series $S_{Z_0}(\sigma)$, which converges absolutely by Lemma~\ref{lem:walk-series}, and the Gaussian integrals over $(X,B)$ at complex $\sigma$, whose absolute convergence is Lemma~\ref{lem:gaussian-payment}(ii) (the walk-series bound alone does not control the real part of the quadratic exponent). Every cube of $Z_0'$ is wall-connected inside $Z_0'$ to a cube containing a seed cube: a cube $\Delta$ of $Z_0'$ meets $\operatorname{int}Q^{\sim}$ for a seed cube $Q$ at a point $x$; since $\operatorname{int}Q^{\sim}$ is open, $x$ may be taken in the interior of $\Delta$, so $x$ lies in an $M$-cube $Q'\subseteq Q^{\sim}\cap\Delta$, and $Q'$ is joined to $Q$ by a chain of face-adjacent $M$-cubes of $Q^{\sim}$, whose $LM$-cubes start with $\Delta$, are consecutively equal or wall adjacent, and all belong to $Z_0'$ (each contains an $M$-cube of $Q^{\sim}$, whose interior lies in $\operatorname{int}Q^{\sim}$). Hence a wall-component of $Z$ either contains a seed cube or consists of differentiated cubes only; in the second case $\Gamma_{E_0}(\sigma)$ does not depend on its $\sigma$-variables (Lemma~\ref{lem:component-blocks}), the Cauchy integral of a function constant in $\sigma_{\Delta}$ vanishes, and the term is zero. So only final domains all of whose components contain a seed cube contribute. For these, Lemma~\ref{lem:component-blocks} factors the density and the Gaussian measures over the components $D$ of $Z$ ($\mu_{C_{E_0}}=\bigotimes_D\mu_{C_{E_D}}$ and the density is a product by the block decomposition; the $X$-integral factors because each block of $\Gamma_{E_0}$ reads $X$ on $J_D$ only, and the remaining Gaussian factor of $\mu_0$ integrates to one), and the parameter integrals are indexed by cubes and hosts, each inside one component. The cutoff, potential and sign factors factor because each selected host is a connected union of $M$-cubes inside $Z_0$, hence inside one component, and each bad bond lies in the component containing its (at most two, face-adjacent) incident cubes. Given $Z$ with components $D_1,\ldots,D_m$, the map $(\mathcal D,\mathsf P)\mapsto((\mathcal D_j,\mathsf P_j))_j$, with $\mathcal D_j$ the hosts inside $D_j$ and $\mathsf P_j$ the bad bonds inside $D_j$, is a bijection from the selections with final domain $Z$ (all $Z_0'\subseteq Z$, every component seeded) onto the tuples of selections $(\mathcal D_j,\mathsf P_j,D_j)$ with nonempty seeds and final domains $D_j$: the inverse is the union, because the seed rule, the enlargement and the coarse cover commute with unions and $(Z_0\cap D_j)'\subseteq D_j$ (it is wall-connected to its seed inside $Z_0'\subseteq Z$). Summing over $Z$ gives the compatible-family sum, in which the empty selection contributes $\int\dd\mu_{\mathsf G}=1$. Equation \eqref{eq:fixed-kernel-source-numerator} is the case $\mathbf W=\mathbf V_{\xi}$: $K_{\xi;0}=\Theta_{\xi}\exp(-\frac12\langle a,Ha\rangle+\sum_Y\mathbf V_{\xi}(Y,a))$ by \eqref{eq:native-source-boundary-interpolation} and \eqref{eq:c-potential-definition} (on the open box $\sum_Y\mathbf V_{\xi}(Y,a)= P_{\xi,0}(a)$), and positivity holds because $K_{\xi;0}>0$ on the open box.
\end{proof}

\begin{lemma}[activity bounds]\label{lem:activity-bounds}
With the constants $\lambda_*,\varepsilon_*>0$ and $\theta\geqslant 0$ of Appendix~\ref{app:background-estimates}, Section~\ref{app:mayer-constants} (fixed there before the coupling together with its auxiliary constants $\kappa,\delta,\alpha_6,\omega_*$, independently of $M_1$ and $M$), and $g_0$ below its threshold, for $\mathbf W=\mathbf V_{\xi}$ on every branch and every connected final domain $Z$:
\begin{equation}\label{eq:fixed-kernel-raw-activity-payment}
    | k^{\mathrm f}_{\xi}(Z)|\leqslant \varepsilon_*e^{-\lambda_*\ell_1(Z)},\qquad
    \sum_{Z'\not\sim Z}| k^{\mathrm f}_{\xi}(Z')| e^{(1+6\theta)N_1(Z')}\leqslant \frac{1}{8}N_1(Z)\;,
\end{equation}
where $Z'$ ranges over the connected final domains and $Z'\not\sim Z$ means incompatible (a domain is incompatible with itself). The same holds for the reference data of Section~\ref{sec:joint-extraction} (Lemma~\ref{lem:c2-reference}).
\end{lemma}

\begin{lemma}[vacuum clusters]\label{lem:vacuum-clusters}
Under Lemma~\ref{lem:activity-bounds}, for an ordered tuple $\mathfrak c=(D_1,\ldots,D_m)$, $m\geqslant 1$, of connected final domains with Ursell coefficient $\varphi^T(\mathfrak c)=\sum_{\mathcal G}(-1)^{|\mathcal G|}$, the sum over the connected spanning subgraphs $\mathcal G$ of the incompatibility graph on $\{1,\ldots,m\}$ ($j\ne j'$ joined when $D_j\not\sim D_{j'}$; a domain is incompatible with itself, so repeated domains are joined; $\varphi^T=1$ for $m=1$), put
\begin{equation}\label{eq:part1-vacuum-cluster}
    \Phi^{\mathrm v}_{\xi}(\mathfrak c)=\frac{\varphi^T(\mathfrak c)}{m!}\prod_{j=1}^mk^{\mathrm f}_{\xi}(D_j)\;.
\end{equation}
For every finite set $\mathcal A$ of connected final domains, $\sum_{\mathfrak c\text{ hits }\mathcal A}|\Phi^{\mathrm v}_{\xi}(\mathfrak c)| e^{2\theta\sum_jN_1(D_j)}\leqslant \frac18\sum_{D\in\mathcal A}N_1(D)$, where $\mathfrak c$ hits $\mathcal A$ when some $D_j$ is incompatible with a member of $\mathcal A$, and
\begin{equation}\label{eq:part1-joint-vacuum-exponential}
    \frac{(\Xi^{\mathrm f}_{\xi})^{[\cup\mathcal A]}}{\Xi^{\mathrm f}_{\xi}}=\exp\Bigl(-\sum_{\mathfrak c\text{ hits }\mathcal A}\Phi^{\mathrm v}_{\xi}(\mathfrak c)\Bigr)\;.
\end{equation}
\end{lemma}

The proof, including the identity \eqref{eq:part1-joint-vacuum-exponential}, is Appendix~\ref{app:background-estimates}, Section~\ref{app:kp}: the Kotecký--Preiss criterion with the pinned-tree bound.

\subsection{Probe kernels and the centred expectation}
\label{sec:probe-kernels}

For each $b\in I$ let $\mathcal R_b=\{ r\text{ a bond of }\Omega_1\, | \,b\in\mathcal F_r\}$ be its forward influence set (the bonds whose reconstructed link depends on $a(b)$, Lemma~\ref{lem:row-locality}); it lies within $\ell^{\infty}$ distance $2L+1$ of $b$. Lift $\{ b\}\cup\mathcal R_b$ to $\mathbb Z^4$ about $b_-$ (its diameter is below every period), and let
\begin{equation}\label{eq:c-probe-host}
    Y_b=\pi_T\bigl(\text{the smallest rectangle of $M$-cubes containing the one-fine-layer neighbourhood of the lift, with one $M$-layer adjoined}\bigr)\;;
\end{equation}
the lift together with its one-fine-layer neighbourhood has $\ell^{\infty}$ extent at most $2(2L+1)+3=4L+5< M$ in every direction ($M\geqslant 8LM_2\geqslant 8L^2$), so it meets at most two $M$-cubes per direction and the rectangle has side at most $4$ cubes, of length $4M<8M$, so $\pi_T$ is injective on it. $Y_b$ is a connected member of $\mathbf D_0$, contains $b$ and $\mathcal R_b$ in its interior, and $N_0(Y_b)\leqslant 4^4\leqslant s_{\mathrm p}=5^4$. For a union $R$ of $LM$-cubes its eligible labels are
\begin{equation}\label{eq:probe-eligible-carrier}
    \mathsf L_R=\{ b\in I\, | \,Y_b\subseteq R\}\;,
\end{equation}
and $\overline{\mathcal D}_{\mathsf L}$ is the closed box $| a_b|\leqslant r$, $b\in\mathsf L$, in $\mathfrak g^{\mathsf L}$. Each $Y_b$ is a nonempty wall-connected union of $M$-cubes, so $Y_b\subseteq R$ places it inside exactly one wall-component of $R$; hence $\mathsf L_R$ is the disjoint union of the $\mathsf L_D$ over the components $D$ of $R$. For $\lambda\in\mathfrak g^I_{\mathbb C}$, with the complex bilinear extension of the real pairing, the probe family is
\begin{equation}\label{eq:probe-combined-potential}
    L_Y(\lambda,a)=\sum_{b\, | \,Y_b= Y}\langle\lambda_b,a_b\rangle,\qquad\widetilde{\mathbf V}_{\xi,\lambda}(Y,a)=\mathbf V_{\xi}(Y,a)+L_Y(\lambda,a)\;,
\end{equation}
a potential family indexed by $\mathbf D_0$ for each $\lambda$ (both terms depend on $a|_{E(Y)}$ only, since $Y_b= Y$ gives $b\in E(Y)$, and both are bounded on the local box of $Y$), entire in $\lambda$. Its partition function is
\begin{equation}\label{eq:probe-partition}
    \mathscr Z^{\mathrm c}_{\xi}(\lambda)=\int e^{\sum_{b\in I}\langle\lambda_b,a_b\rangle}K_{\xi;0}(a)\,\dd a=\mathfrak n_{\xi}\sum_{\mathcal P\text{ compatible}}\prod_{D\in\mathcal P}k^{\mathrm f}_{\xi,D}(\lambda)\;,
\end{equation}
by Lemma~\ref{lem:polymer} applied to $\widetilde{\mathbf V}_{\xi,\lambda}$, where $k^{\mathrm f}_{\xi,D}(\lambda)$ denotes the activity formed from that family; it is entire in $\lambda$ (a finite sum of integrals of entire functions of $\lambda$ with integrable majorants on compact sets of $\lambda$) and $k^{\mathrm f}_{\xi,D}(0)= k^{\mathrm f}_{\xi}(D)$.

\begin{lemma}[probe measures]\label{lem:probe-measures}
On every branch, for every connected final domain $D$ there is a finite complex Borel measure $\nu^{\mathrm f}_{\xi,D}$ on $\overline{\mathcal D}_{\mathsf L_D}$ with
\begin{equation}\label{eq:probe-Laplace-identity}
    k^{\mathrm f}_{\xi,D}(\lambda)=\int e^{\sum_{b\in\mathsf L_D}\langle\lambda_b,a_b\rangle}\,\dd\nu^{\mathrm f}_{\xi,D}(a)\quad(\lambda\in\mathfrak g^{I}_{\mathbb C}),\qquad
    \nu^{\mathrm f}_{\xi,D}(\overline{\mathcal D}_{\mathsf L_D})= k^{\mathrm f}_{\xi}(D)\;;
\end{equation}
in particular $k^{\mathrm f}_{\xi,D}(\lambda)$ depends on $\lambda$ only through $\lambda|_{\mathsf L_D}$. Explicitly, splitting each probe label from its host's ordinary label, a selection becomes $\ell=(\mathcal D,\mathsf J,\mathsf P,Z)$ with $\mathsf J\subseteq\mathsf L_Z$ the selected probes, $Y_0=\bigcup_{Y\in\mathcal D}Y\cup\bigcup_{b\in\mathsf J}Y_b$, and seed, $E_0$ and $Z_0'$ formed from this $Y_0$ and $\mathsf P$; the Mayer variable of a probe integrates exactly,
\begin{equation}\label{eq:probe-selected-factor}
    \int_0^1\langle\lambda_b,B_b\rangle e^{t_b\langle\lambda_b,B_b\rangle}\,\dd t_b= e^{\langle\lambda_b,B_b\rangle}-1\;,
\end{equation}
and with $\Lambda^{\mathrm f}_{\xi;\ell}$ the complex measure on $\mathfrak g^{E_0}$ defined by the right side of \eqref{eq:c-selected-integral} for $\mathbf W=\mathbf V_{\xi}$, with $Y_0$, $Z_0$, $E_0$, $Z_0'$ the enlarged data of $\ell$ throughout (the factors $\chi_b$ over $E(Y_0)$, the conditional law and the weakening over this seed), with the Mayer factors $\int_0^1\dd t_Y\,\mathbf W(Y,B)e^{t_Y\mathbf W(Y,B)}$ for $Y\in\mathcal D$ only and no probe factor,
\begin{equation}\label{eq:probe-selection-measure}
    \nu^{\mathrm f}_{\xi,D}=\sum_{\ell\, | \,Z= D,\ Z_0\ne\varnothing}\int\Bigl[\bigotimes_{b\in\mathsf J}(\delta_{B_b}-\delta_0)\otimes\bigotimes_{b\in\mathsf L_D\setminus\mathsf J}\delta_0\Bigr]\,\dd\Lambda^{\mathrm f}_{\xi;\ell}(B)\;.
\end{equation}
\end{lemma}

\begin{proof}
At a host $Y$ carrying the ordinary potential $\mathbf V_{\xi}(Y)$ and the probe terms $F_b=\langle\lambda_b,\bigcdot \rangle$, $b$ with $Y_b= Y$, write $e^{\widetilde{\mathbf V}_{\xi,\lambda}(Y)}= e^{\mathbf V_{\xi}(Y)}\prod_be^{F_b}=\prod_{\ell'}(1+(e^{F_{\ell'}}-1))$ over the labels $\ell'$ of $Y$ (the ordinary label with $F_{\ell'}=\mathbf V_{\xi}(Y)$ and one label per probe); the sum over nonempty label subsets $A$ of $\prod_{\ell'\in A}(e^{F_{\ell'}}-1)$ equals $e^{\widetilde{\mathbf V}_{\xi,\lambda}(Y)}-1$. Every term with $A\ne\varnothing$ has the same active host $Y$, so the split-label expansion has the same $Y_0$, $\mathsf P$, seed, cover and kernel as the one-potential-per-host expansion of Lemma~\ref{lem:polymer}, and the two activities agree term by term after the Mayer variables $t_b$ of the probes are integrated by \eqref{eq:probe-selected-factor}. The Laplace transform of $\delta_{B_b}-\delta_0$ in the variable $\lambda_b$ is $e^{\langle\lambda_b,B_b\rangle}-1$, and an unselected label contributes the factor one, the transform of $\delta_0$; this proves \eqref{eq:probe-Laplace-identity} for the measure \eqref{eq:probe-selection-measure}. Each $\Lambda^{\mathrm f}_{\xi;\ell}$ is a finite complex measure: it is the image under $(X,B,s,\sigma,t)\mapsto B$ of a complex density whose absolute value is integrable by Lemma~\ref{lem:gaussian-payment}(ii) and the boundedness of each $\mathbf V_{\xi}(Y,\bigcdot )$ on its local box, the Cauchy kernels and the cutoff factors; the Dirac tensor has total variation at most $2^{|\mathsf J|}$, and there are finitely many $\ell$ at fixed volume. The support lies in the closed box because $\chi_b(B)$, $b\in E(Y_0)\supseteq\mathsf J$, forces $| B_b|< r$ on the selected probes. At $\lambda=0$ every term with $\mathsf J\ne\varnothing$ has zero mass, so $\nu^{\mathrm f}_{\xi,D}(\overline{\mathcal D}_{\mathsf L_D})= k^{\mathrm f}_{\xi,D}(0)= k^{\mathrm f}_{\xi}(D)$.
\end{proof}

For a compatible family $\mathcal P$ let $\mathsf e_{\mathcal P}$ embed $\prod_{D\in\mathcal P}\mathfrak g^{\mathsf L_D}$ into $\mathfrak g^I$ by zero on the remaining labels (the $\mathsf L_D$ are pairwise disjoint).

\begin{lemma}[measure identity]\label{lem:measure-identity}
On every branch,
\begin{equation}\label{eq:fixed-probe-measure}
    \frac{K_{\xi;0}(a)\,\dd a}{\mathfrak n_{\xi}}=\sum_{\mathcal P\text{ compatible}}(\mathsf e_{\mathcal P})_*\bigotimes_{D\in\mathcal P}\nu^{\mathrm f}_{\xi,D}\;.
\end{equation}
\end{lemma}

\begin{proof}
Both sides are finite complex measures with compact support (finitely many families at fixed volume), and by \eqref{eq:probe-partition} and \eqref{eq:probe-Laplace-identity} they have the same Laplace transform on $\mathfrak g^I_{\mathbb C}$; restricted to imaginary $\lambda$ this is the Fourier transform, which determines a finite measure.
\end{proof}

Let $\delta_{0,\mathsf L_D}$ be the Dirac mass at $0\in\mathfrak g^{\mathsf L_D}$ and
\begin{equation}\label{eq:fixed-probe-zero-mass}
    \widehat\nu^{\mathrm f}_{\xi,D}=\nu^{\mathrm f}_{\xi,D}-k^{\mathrm f}_{\xi}(D)\delta_{0,\mathsf L_D}\;,
\end{equation}
a measure of zero mass (the subtracted term is the sum of the selections with $\mathsf J=\varnothing$). For a nonempty compatible family $\mathcal A$ and bounded Borel $f$ on $\mathfrak g^I$ put
\begin{equation}\label{eq:fixed-probe-active-pairing}
    \mathfrak Q_{\xi}(\mathcal A;f)=\int f(\mathsf e_{\mathcal A}a)\,\dd\Bigl(\bigotimes_{D\in\mathcal A}\widehat\nu^{\mathrm f}_{\xi,D}\Bigr)(a),\qquad
    q_{\xi}(\mathcal A)=\frac{(\Xi^{\mathrm f}_{\xi})^{[\cup\mathcal A]}}{\Xi^{\mathrm f}_{\xi}}\;.
\end{equation}

\begin{proposition}[centred expectation]\label{prop:centred-expectation}
On every branch, for bounded Borel $f$ on $\mathfrak g^I$,
\begin{equation}\label{eq:fixed-probe-centred-source}
    \int f\,\dd\nu_{\xi;0}=\frac{\int f(a)K_{\xi;0}(a)\,\dd a}{I_{\xi;0}}= f(0)+\sum_{\varnothing\ne\mathcal A\text{ compatible}}q_{\xi}(\mathcal A)\mathfrak Q_{\xi}(\mathcal A;f)\;,
\end{equation}
and, with Lemmas~\ref{lem:collar} and~\ref{lem:boundary-response},
\begin{equation}\label{eq:c-physical-numerator}
    \mathcal J^{\mathrm a}_{\xi}[f]= e^{c_{\xi}}I_{\xi;1}\Bigl[f(0)+\sum_{\varnothing\ne\mathcal A\text{ compatible}}q_{\xi}(\mathcal A)\mathfrak Q_{\xi}(\mathcal A;f)+\mathcal B_{\xi}[f]\Bigr]\;.
\end{equation}
\end{proposition}

\begin{proof}
Integrate $f$ against \eqref{eq:fixed-probe-measure} and expand each $\nu^{\mathrm f}_{\xi,D}$ as $\widehat\nu^{\mathrm f}_{\xi,D}+k^{\mathrm f}_{\xi}(D)\delta_{0,\mathsf L_D}$. A term is indexed by $\mathcal P$ and the subfamily $\mathcal A\subseteq\mathcal P$ carrying the zero-mass factors; $f(\mathsf e_{\mathcal P}a)$ with $a=0$ on the labels of $\mathcal P\setminus\mathcal A$ is $f(\mathsf e_{\mathcal A}a)$. Summing over $\mathcal P\supseteq\mathcal A$ gives $(\Xi^{\mathrm f}_{\xi})^{[\cup\mathcal A]}\mathfrak Q_{\xi}(\mathcal A;f)$, and $\mathcal A=\varnothing$ gives $\Xi^{\mathrm f}_{\xi} f(0)$. Divide by $I_{\xi;0}=\mathfrak n_{\xi}\Xi^{\mathrm f}_{\xi}$. The second display is $\mathcal J^{\mathrm a}_{\xi}[f]= e^{c_{\xi}}I_{\xi;1}\int f\,\dd\nu_{\xi;1}$ \eqref{eq:c-collar-identity} with \eqref{eq:fixed-probe-boundary-functional}.
\end{proof}

With $f= f_{\xi}$ (extended by zero off the enlarged box; the measures live on the closed box, inside the enlarged box, where $f_{\xi}$ is the Wilson function \eqref{eq:convex-response-operands}), the first display is $F_{\xi}(0)$, and \eqref{eq:c-physical-numerator} is the exact decomposition of the conditional Wilson numerator on every branch: the loop at zero young field, the sum of closed insertions paired with the signed probe measures, and the exterior response.

\subsection{Joint extraction of the scalar reference}
\label{sec:joint-extraction}

Let $r_n$ be the positive representative of traversal $n\in\{1,\ldots,n_C\}$ of $C$, the traversals numbered cyclically from a fixed base point. For a contour bond $r$ that is not a bond of $\Omega_1$ put $\mathcal F_r=\varnothing$: its link is $V_0(r)$ by \eqref{eq:c-reconstructed-field}, independent of $a$. By Lemma~\ref{lem:row-locality} and this convention, the link $U_{\xi}(a)(r_n)$ depends on the young variable only through $a|_{\mathcal F_{r_n}\cap I}$. For a connected final domain $D$ the contact set is
\begin{equation}\label{eq:part1-contact-set}
    \mathsf T_D=\{ n\, | \,\mathcal F_{r_n}\cap\mathsf L_D\ne\varnothing\}\;,
\end{equation}
and $D$ is marked when $\mathsf T_D\ne\varnothing$. If a family $\mathcal A$ contains an unmarked domain $D$, then for $a$ in the closed box $f_{\xi}(\mathsf e_{\mathcal A}a)= W_{\chi}(C;U_{\xi}(\mathsf e_{\mathcal A}a))$ does not depend on $a|_{\mathsf L_D}$ (every contour link depends on $a|_{\mathcal F_{r_n}\cap I}$, and $\mathcal F_{r_n}\cap\mathsf L_D=\varnothing$ for every $n$), and $\mathfrak Q_{\xi}(\mathcal A;f_{\xi})=0$ because $\widehat\nu^{\mathrm f}_{\xi,D}$ has zero mass and the measures are supported in the closed box; so only families of marked domains contribute to \eqref{eq:fixed-probe-centred-source} for $f= f_{\xi}$. Expand $q_{\xi}(\mathcal A)$ by Lemma~\ref{lem:vacuum-clusters},
\begin{equation}\label{eq:c-cluster-expansion}
    q_{\xi}(\mathcal A)=\sum_{k\geqslant 0}\frac{(-1)^k}{k!}\sum_{\mathfrak c_1,\ldots,\mathfrak c_k\text{ each hits }\mathcal A}\prod_{j=1}^k\Phi^{\mathrm v}_{\xi}(\mathfrak c_j)\;,
\end{equation}
the sum being absolutely convergent: $\sum_{\mathfrak c\text{ hits }\mathcal A}|\Phi^{\mathrm v}_{\xi}(\mathfrak c)|\leqslant \frac18\sum_{D\in\mathcal A}N_1(D)$ by Lemma~\ref{lem:vacuum-clusters}, so the $k$-fold sums are dominated by the exponential series. A tuple whose incompatibility graph is disconnected has $\varphi^T(\mathfrak c)=0$ (it has no connected spanning subgraph), so the terms containing such a tuple vanish and are dropped: every cluster occurrence $\mathfrak c_j$ below has a connected incompatibility graph. For each term form the graph on the marked domains of $\mathcal A$ and the $k$ cluster occurrences: a marked vertex is joined to each cluster that hits it, two clusters are joined when some constituent domains are incompatible. Every component has a marked root; its contacts are the union of the contact sets of its roots. The hull of a nonempty set of traversals is its shortest cyclic interval when that interval has fewer than $n_C/2$ positions, and the whole contour otherwise; a shortest interval of fewer than $n_C/2$ positions is unique (two such intervals containing the set have together fewer than $n_C$ positions, so they intersect in a cyclic interval, which contains the set and is shorter than one of them unless they coincide), and if $S\subseteq S'$ and the hull of $S'$ is short then the hull of $S$ is contained in it (the hull of $S'$ is a short interval containing $S$, so the hull of $S$ is short and, by the uniqueness just proved, is the shortest interval containing $S$; the intersection of the two short hulls is again an interval, their union having fewer than $n_C$ positions, and it contains $S$, so it is the hull of $S$). Let $\approx$ be the transitive closure of the relation ``the hulls of the contacts intersect'' on the components, and for a class $K$ let $\mathcal H(K)$ be the hull of the union of the contacts of its members. If the union of the members' hulls has fewer than $n_C/2$ positions, $\mathcal H(K)$ equals that union: the union is a cyclic interval (the members' hulls are chained by intersections, and two intersecting intervals with together fewer than $n_C$ positions have an interval union), it contains the union of the contacts, so it contains the hull $\mathcal H(K)$ by minimality; conversely each member's contacts lie in the union of the contacts, whose hull is short, so by monotonicity each member's hull lies in $\mathcal H(K)$. If the union of the members' hulls has at least $n_C/2$ positions, $\mathcal H(K)$ is the whole contour: a short $\mathcal H(K)$ would contain every member's hull by monotonicity. The groups are the classes if every $\mathcal H(K)$ is short, and the single class of all components otherwise; the intervals of distinct groups are disjoint, since intersecting intervals would relate their classes.

\begin{definition}[groups]\label{def:groups}
A group $W$ consists of its nonempty compatible marked-root set $\mathcal A(W)$, its ordered list of cluster occurrences $(\mathfrak c_1,\ldots,\mathfrak c_{k(W)})$, $k(W)\geqslant 0$, each hitting a root of $W$, and its interval $\mathcal H(W)$, either a cyclic interval of fewer than $n_C/2$ traversals or the whole contour read from the base point; the data arise from a term of \eqref{eq:c-cluster-expansion} as above; conversely, the term of \eqref{eq:c-cluster-expansion} with root family $\mathcal A(W)$ and cluster list $(\mathfrak c_1,\ldots,\mathfrak c_{k(W)})$ has $W$ as its only group (its graph has the components of $W$, which form one class with the interval $\mathcal H(W)$). Its weight is
\begin{equation}\label{eq:part1-group-weight}
    w_{\xi}(W)=\frac{(-1)^{k(W)}}{k(W)!}\prod_{j=1}^{k(W)}\Phi^{\mathrm v}_{\xi}(\mathfrak c_j)\;.
\end{equation}
Two groups are compatible when all their constituent domains (roots and cluster members) are mutually compatible and their intervals are disjoint; a family containing a whole-contour group is a singleton. For a compatible group family $\mathfrak B$ and $\mathcal S\subseteq\mathfrak B$,
\begin{equation}\label{eq:part1-closed-family-pairing}
    \mathcal A(\mathcal S)=\bigcup_{W\in\mathcal S}\mathcal A(W),\qquad
    A_{\xi}(\mathcal S)=\Bigl(\prod_{W\in\mathcal S}w_{\xi}(W)\Bigr)\mathfrak Q_{\xi}(\mathcal A(\mathcal S);f_{\xi}),\qquad A_{\xi}(\varnothing)= f_{\xi}(0)\;.
\end{equation}
\end{definition}

\begin{lemma}[grouped source]\label{lem:grouped-source}
On every branch,
\begin{equation}\label{eq:part1-grouped-source}
    F_{\xi}(0)=\sum_{\mathfrak B\text{ compatible}}A_{\xi}(\mathfrak B)\;,
\end{equation}
an absolutely convergent sum over compatible group families (including the empty family; there are finitely many root families at fixed volume but infinitely many cluster lists).
\end{lemma}

\begin{proof}
Substitute \eqref{eq:c-cluster-expansion} into \eqref{eq:fixed-probe-centred-source} for $f= f_{\xi}$. Each term, indexed by a family $\mathcal A$ of marked domains and an ordered list of $k$ cluster occurrences, is cut into its groups; the groups are pairwise compatible (roots in different groups are compatible as members of $\mathcal A$, clusters in different groups are compatible with each other and with the other groups' roots because the graph has no edge between components, and the intervals are disjoint by construction). The regrouping is reversible: a compatible group family determines $\mathcal A=\mathcal A(\mathfrak B)$ and the multiset of cluster occurrences, and the term's graph has as components exactly the components of the groups' own graphs (no edges between groups, by compatibility), with the same hulls, so the classes and intervals are recovered; identical occurrences cannot split between groups because their domains are self-incompatible, so a group's occurrences are those of its term. The ordered list of the term is an interleaving of the groups' lists, and the interleavings number $k!/\prod_Wk(W)!$, which turns the factor $(-1)^k/k!$ into $\prod_Ww_{\xi}(W)/\prod_W\Phi$-products. The rearrangement is of an absolutely convergent sum: there are finitely many root families $\mathcal A$ at fixed volume, $|\mathfrak Q_{\xi}(\mathcal A;f_{\xi})|\leqslant \| f_{\xi}\|_{\infty}\prod_{D\in\mathcal A}\|\widehat\nu^{\mathrm f}_{\xi,D}\|$ with the total variations finite, and $\sum_{k}\frac{1}{k!}\bigl(\sum_{\mathfrak c\text{ hits }\mathcal A}|\Phi^{\mathrm v}_{\xi}(\mathfrak c)|\bigr)^k\leqslant \exp(\frac18\sum_{D\in\mathcal A}N_1(D))$; hence the term sum converges absolutely, and so does the group-family sum, being a rearrangement of it.
\end{proof}

A superscript $\circ$ denotes the reference data: the same $\Omega_1$, $\Lambda_1$, $I$, $J$, cubes, hosts, constants and cutoffs, with the background $U^{\circ}_1=\id$ on every bond, $V^{\circ}_0=\id$, zero retained field $A^{\circ}_{\mathrm e}=0$ and $\Psi^{\circ}_{\xi}=\id$; thus $H^{\circ}_{\xi}=\mathsf H^{\Omega_1}_{\id}$, the reference phase is $q_{\id}$, the reference primitives are those of Section~\ref{sec:source-phase-primitives} at $U=\id$, $f^{\circ}_{\xi}(a)= W_{\chi}(C;e^{iq_{\id}(a)})$ with the links off $\Omega_1$ equal to $\id$, and $f^{\circ}_{\xi}(0)=\chi(\id)=1$ since $q_{\id}(0)=0$. These data need not be a branch of the history sum. Lemmas~\ref{lem:two-block}, \ref{lem:hessian-gauge-locality}, \ref{lem:row-locality} and~\ref{lem:primitive-sum} apply to $U=\id$ directly (Section~\ref{sec:source-phase-primitives}, opening paragraph; Lemma~\ref{lem:two-block} at $A=0$; the row solutions are those of Lemma~\ref{lem:charts} in its reference form, on the row polydiscs $\mathcal U_c$, and (B4) is not invoked), and the estimate lemmas used below (Lemmas~\ref{lem:regular-bounds}, \ref{lem:charts}, \ref{lem:gaussian-payment}, \ref{lem:activity-bounds}, \ref{lem:vacuum-clusters}) hold for the reference data by Appendix~\ref{app:background-estimates}, Lemma~\ref{lem:c2-reference}; in particular the reference near primitives are real on the open box (Lemma~\ref{lem:charts}: the analytic branches of the logarithms are real at real arguments, the row solution being real), so $K^{\circ}_{\xi;0}>0$ there and $I^{\circ}_{\xi;0}>0$. Every lemma of Sections~\ref{sec:fixed-first-kernel}--\ref{sec:joint-extraction} stated on every branch (Lemmas~\ref{lem:walk-series}, \ref{lem:conditioning}, \ref{lem:component-blocks}, \ref{lem:polymer}, \ref{lem:probe-measures}, \ref{lem:measure-identity}, \ref{lem:grouped-source} and the first display of Proposition~\ref{prop:centred-expectation}) holds verbatim for the reference data: their proofs use only the field-independent geometry and the lemmas just listed. Every object of Sections~\ref{sec:fixed-first-kernel}--\ref{sec:joint-extraction} formed from the reference data carries the superscript $\circ$; in particular $K^{\circ}_{\xi;0}$, $I^{\circ}_{\xi;0}$, $k^{\mathrm f,\circ}_{\xi}$, $\widehat\nu^{\mathrm f,\circ}_{\xi,D}$, $w^{\circ}_{\xi}(W)$. The index set of groups is the same for the branch and for the reference data: roots, cluster lists (the clusters are tuples of domains), contact sets and intervals are field-independent data (the sets $\mathcal F_{r_n}$ and $\mathsf L_D$ do not involve the background), and only the weights $w_{\xi}$, $w^{\circ}_{\xi}$ and the measures $\widehat\nu^{\mathrm f}_{\xi,D}$, $\widehat\nu^{\mathrm f,\circ}_{\xi,D}$ differ. For a group $W$ let $P^{\circ}_W(a)$ be the ordered product of the reference links along the consecutive traversals of $\mathcal H(W)$, in the cyclic order of the contour starting from the first traversal of the interval (for a whole-contour group, from the base point), the link of a traversal being $e^{iq_{\id}(\mathsf e_{\mathcal A(W)}a)(r)}$ when its bond $r$ is a bond of $\Omega_1$ and $\id$ otherwise, as in $f^{\circ}_{\xi}$ (oriented representatives inverted), with the young variable $\mathsf e_{\mathcal A(W)}a$, zero off the labels of $\mathcal A(W)$; then $P^{\circ}_W(0)=\id$. Put
\begin{equation}\label{eq:part1-reference-group}
    a_W= w^{\circ}_{\xi}(W)\int\chi\bigl(P^{\circ}_W(a)\bigr)\,\dd\Bigl(\bigotimes_{D\in\mathcal A(W)}\widehat\nu^{\mathrm f,\circ}_{\xi,D}\Bigr)(a)\;.
\end{equation}

\begin{lemma}[Schur identity]\label{lem:schur}
For $h\in G$,
\begin{equation}\label{eq:part1-closed-schur}
    w^{\circ}_{\xi}(W)\int\chi\bigl(P^{\circ}_W(a)h\bigr)\,\dd\Bigl(\bigotimes_{D\in\mathcal A(W)}\widehat\nu^{\mathrm f,\circ}_{\xi,D}\Bigr)(a)= a_W\chi(h)\;.
\end{equation}
Consequently
\begin{equation}\label{eq:part1-reference-gas}
    S^{\circ}_{\xi,C}=\sum_{\mathfrak B\text{ compatible}}\prod_{W\in\mathfrak B}a_W=\frac{\int f^{\circ}_{\xi}(a)K^{\circ}_{\xi;0}(a)\,\dd a}{I^{\circ}_{\xi;0}}\;.
\end{equation}
\end{lemma}

\begin{proof}
For constant $u\in G$ let $(R_ua)_b=\operatorname{Ad}_ua_b$ on $\mathfrak g^I$, an orthogonal map for the invariant inner product. At the identity background the constant gauge transformation $u$ fixes the background, so Lemma~\ref{lem:hessian-gauge-locality} gives $H^{\circ}_{\xi}= R_uH^{\circ}_{\xi} R_u^{-1}$; hence all walk coefficients, $C^{\circ}_{E_0}$, $\Gamma^{\circ}_{E_0}(\sigma)$ (the domains and weights do not involve the fields) and the Gaussian density \eqref{eq:fixed-kernel-density} are invariant under $(X,B)\mapsto(R_uX,R_uB)$, and $\mu_0$, $\mu_{C^{\circ}_{E_0}}$ are invariant. The cutoff factors are invariant because the bond norm is $\operatorname{Ad}$-invariant (\eqref{eq:marginal-cutoff-invariance}). The reconstruction is covariant: by the intertwinings in the proof of Lemma~\ref{lem:hessian-gauge-locality} at $U=\id$, $\mathcal N_{\id}\circ R_u= R^{(1)}_u\mathcal N_{\id}$, $h_{\id}R^{(1)}_u= R_uh_{\id}$ and $\mathsf C_{\id}R_{u,*}= R_u\mathsf C_{\id}$, so the unique small solution on the row polydiscs of Lemma~\ref{lem:charts} (reference version; the polydiscs are $R_u$-invariant because the bond norm is $\operatorname{Ad}$-invariant) satisfies $\widetilde D_{\id}(R_uA)= R^{(1)}_u\widetilde D_{\id}(A)$ and $q_{\id}(R_uB)= R_uq_{\id}(B)$, that is, every reconstructed link is conjugated by $u$; the primitives are invariant, being traces of products of links, Haar densities of conjugated phases ($\sigma$ is conjugation invariant) and determinants of conjugated operators. The probe term transforms as $L_Y(R_u\lambda,R_ua)= L_Y(\lambda,a)$. Hence $k^{\mathrm f,\circ}_{\xi,D}(R_u\lambda)= k^{\mathrm f,\circ}_{\xi,D}(\lambda)$, and by Laplace uniqueness $(R_u)_*\widehat\nu^{\mathrm f,\circ}_{\xi,D}=\widehat\nu^{\mathrm f,\circ}_{\xi,D}$. Since the links are conjugated, $P^{\circ}_W(R_ua)= uP^{\circ}_W(a)u^{-1}$. Let $\pi$ be the irreducible unitary representation with $\chi=\dim(\pi)^{-1}\operatorname{tr}\pi$ (the normalised character, $\chi(\id)=1$). The left side of \eqref{eq:part1-closed-schur}, as a function of $h$, is $\dim(\pi)^{-1}\operatorname{tr}(\Upsilon\pi(h))$ for the matrix $\Upsilon= w^{\circ}_{\xi}(W)\int\pi(P^{\circ}_W(a))\,\dd(\bigotimes\widehat\nu^{\circ})(a)$; by the invariance of the measure and the covariance of $P^{\circ}_W$, $\Upsilon=\pi(u)\Upsilon\pi(u)^{-1}$ for all $u\in G$, so $\Upsilon$ is a multiple of the identity by Schur's lemma, and $h=\id$ identifies the multiple as $a_W$. For \eqref{eq:part1-reference-gas} apply Lemma~\ref{lem:grouped-source} to the reference data (its sum converges absolutely, as shown there). Let $\mathfrak B=\{ W_1,\ldots,W_m\}$ be a compatible family. If $m=1$ and $W_1$ is a whole-contour group, the reference Wilson function at $\mathsf e_{\mathcal A(W_1)}a$ is $\chi(P^{\circ}_{W_1}(a))$. Otherwise all intervals are short and pairwise disjoint; number the groups in cyclic order starting from $W_1$, and use the cyclicity of the trace to start the ordered link product at the first traversal of $\mathcal H(W_1)$. The reference Wilson function at $\mathsf e_{\mathcal A(\mathfrak B)}a$ is then $\chi(P^{\circ}_{W_1}(a)\cdots P^{\circ}_{W_m}(a))$: a traversal outside every interval is outside every contact set, so its link is $e^{iq_{\id}(0)}=\id$ (or $\id$ by convention when its bond is not a bond of $\Omega_1$); a traversal in $\mathcal H(W_j)$ lies in no contact set of a root of another group, so its link depends only on $a|_{\mathcal A(W_j)}$, and the links of the consecutive traversals of $\mathcal H(W_j)$ multiply to $P^{\circ}_{W_j}(a)$ in the order of the product. The integral against $\bigotimes_j\bigotimes_{D\in\mathcal A(W_j)}\widehat\nu^{\circ}_D$ of this character times $\prod_jw^{\circ}(W_j)$ is evaluated by Fubini one group at a time: integrating over the variables of $W_1$ with $h= P^{\circ}_{W_2}(a)\cdots P^{\circ}_{W_m}(a)$ fixed gives $a_{W_1}\chi(P^{\circ}_{W_2}(a)\cdots P^{\circ}_{W_m}(a))$ by \eqref{eq:part1-closed-schur}, and so on, ending with $\prod_ja_{W_j}\chi(\id)=\prod_ja_{W_j}$. Since $|\chi|\leqslant 1$, $| a_W|\leqslant | w^{\circ}_{\xi}(W)|\prod_{D\in\mathcal A(W)}\|\widehat\nu^{\mathrm f,\circ}_{\xi,D}\|$, and the sum in \eqref{eq:part1-reference-gas} converges absolutely by the bound of Lemma~\ref{lem:grouped-source}.
\end{proof}

\begin{proposition}[joint extraction]\label{prop:joint-extraction}
For a compatible group family $\mathfrak B$ define the joint defect
\begin{equation}\label{eq:part1-joint-defect}
    R_{\xi}(\mathfrak B)=\sum_{\mathcal S\subseteq\mathfrak B}(-1)^{|\mathfrak B\setminus\mathcal S|}\Bigl(\prod_{W\in\mathfrak B\setminus\mathcal S}a_W\Bigr)A_{\xi}(\mathcal S)\;,\qquad R_{\xi}(\varnothing)= f_{\xi}(0)\;.
\end{equation}
Then, with $S^{\circ,[\mathfrak B]}_{\xi,C}$ the reference sum over group families compatible with every member of $\mathfrak B$,
\begin{equation}\label{eq:part1-scalar-extraction}
    F_{\xi}(0)= S^{\circ}_{\xi,C}f_{\xi}(0)+\sum_{\varnothing\ne\mathfrak B\text{ compatible}}S^{\circ,[\mathfrak B]}_{\xi,C}R_{\xi}(\mathfrak B)\;,
\end{equation}
and on every branch
\begin{equation}\label{eq:c-one-step-branch-identity}
    \mathcal J^{\mathrm a}_{\xi}[f_{\xi}]= e^{c_{\xi}}I_{\xi;1}\Bigl[S^{\circ}_{\xi,C}f_{\xi}(0)+\sum_{\varnothing\ne\mathfrak B\text{ compatible}}S^{\circ,[\mathfrak B]}_{\xi,C}R_{\xi}(\mathfrak B)+\mathcal B_{\xi}[f_{\xi}]\Bigr]\;.
\end{equation}
\end{proposition}

\begin{proof}
Boolean inversion of \eqref{eq:part1-joint-defect} gives $A_{\xi}(\mathfrak B)=\sum_{\mathcal S\subseteq\mathfrak B}(\prod_{\mathfrak B\setminus\mathcal S}a_W)R_{\xi}(\mathcal S)$; substitute into \eqref{eq:part1-grouped-source} and sum over $\mathfrak B\supseteq\mathcal S$: the families $\mathfrak B\setminus\mathcal S$ range over the group families compatible with every member of $\mathcal S$, giving $S^{\circ,[\mathcal S]}_{\xi,C}$. The interchange of the sums is justified by absolute convergence. Put $\beta_W=(| w_{\xi}(W)|+| w^{\circ}_{\xi}(W)|)\prod_{D\in\mathcal A(W)}\max(\|\widehat\nu^{\mathrm f}_{\xi,D}\|,\|\widehat\nu^{\mathrm f,\circ}_{\xi,D}\|)$; then $| a_W|\leqslant \beta_W$, $| A_{\xi}(\mathcal S)|\leqslant \| f_{\xi}\|_{\infty}\prod_{W\in\mathcal S}\beta_W$ and $| R_{\xi}(\mathcal S)|\leqslant \| f_{\xi}\|_{\infty}\prod_{W\in\mathcal S}2\beta_W$. The sum $\sum_W\beta_W$ over all groups is finite: a group is determined by its root family, its cluster list and its interval, there are finitely many root families at fixed volume, and $\sum_{W\, | \,\mathcal A(W)=\mathcal A}| w_{\xi}(W)|\leqslant \exp(\frac18\sum_{D\in\mathcal A}N_1(D))$ by Lemma~\ref{lem:vacuum-clusters}, likewise for $w^{\circ}_{\xi}$ by Lemma~\ref{lem:c2-reference}. A compatible group family has at most $N_{\mathrm m}$ members, the number of marked domains (its roots are pairwise distinct marked domains), so $\sum_{\mathfrak B}\sum_{\mathcal S\subseteq\mathfrak B}\prod_{W\in\mathfrak B\setminus\mathcal S}| a_W|\,| R_{\xi}(\mathcal S)|\leqslant \| f_{\xi}\|_{\infty}\sum_{\mathfrak B}\prod_{W\in\mathfrak B}3\beta_W\leqslant \| f_{\xi}\|_{\infty}\sum_{n\leqslant N_{\mathrm m}}\bigl(3\sum_W\beta_W\bigr)^n<\infty$; in particular $S^{\circ}_{\xi,C}$ and every $S^{\circ,[\mathcal S]}_{\xi,C}$ converge absolutely. Then \eqref{eq:c-physical-numerator} with $f= f_{\xi}$.
\end{proof}

At the reference data every nonempty defect vanishes: there $A^{\circ}_{\xi}(\mathcal S)=\prod_{W\in\mathcal S}a_W$ for every compatible $\mathcal S$ (proof of Lemma~\ref{lem:schur}), so $R^{\circ}_{\xi}(\mathfrak B)=\prod_{W\in\mathfrak B}a_W\sum_{\mathcal S\subseteq\mathfrak B}(-1)^{|\mathfrak B\setminus\mathcal S|}=0$ for $\mathfrak B\ne\varnothing$. The identity uses no division by a character value and no norm of an open transport.

\begin{lemma}[covered branches]\label{lem:covered}
Call $\xi$ covered when every bond of the contour is a bond of $\Omega_1$ and every free-input stencil $F_{\gamma}$ of the contour \eqref{eq:native-source-arc-stencils} lies in $I$, that is, $\mathcal F_r\subseteq I$ for every contour bond $r$, with $\mathcal F_r$ of \eqref{eq:regular-stencil-extension-support} (the convention $\mathcal F_r=\varnothing$ off $\Omega_1$ of Section~\ref{sec:joint-extraction} is not needed here, the first condition placing every contour bond in $\Omega_1$). On a covered branch $U_{\xi}(0)(r)= U_1(r)$ for every contour bond $r$, hence $f_{\xi}(0)= W_{\chi}(C;U_1)$.
\end{lemma}

\begin{proof}
At $a=0$ the free vector is $B=\Psi_{\xi}(A_{\mathrm e})$, which vanishes on $I$. For a contour bond $r$, a bond of $\Omega_1$ on a covered branch, $U_{\xi}(0)(r)= e^{iq_U(B)(r)}U_1(r)$ by \eqref{eq:c-reconstructed-field}, and by Lemma~\ref{lem:row-locality} $q_U(B)(r)$ depends on $B|_{\mathcal F_r}$ only and vanishes when $B|_{\mathcal F_r}=0$ (for a free bond directly, for $r= b_0(c)$ because $\mathcal F_r=\mathcal E_c\cap E_*$). Since $\mathcal F_r\subseteq I$ on a covered branch, $U_{\xi}(0)= U_1$ on the contour.
\end{proof}

This ends the construction. Appendix~\ref{app:background-estimates} proves Lemmas~\ref{lem:regular-bounds}, \ref{lem:charts}, \ref{lem:gaussian-payment}, \ref{lem:activity-bounds} and \ref{lem:vacuum-clusters}, and Lemma~\ref{lem:c2-reference} there transfers them to the reference data; Appendix~\ref{app:exterior-chessboard} bounds $\mathcal B_{\xi}[f_{\xi}]$ and the uncovered branches; Section~\ref{sec:one-step-insertion} identifies $S^{\circ}_{\xi,C}$ with the flat coefficient $\kappa_Ce^{E^{\circ}_{\xi}}$ defined there and bounds the defects. That section does not work at the carrier $I$ directly: it adds a second cutoff test on the collar of $\Lambda_1$, retreats $\Lambda_1$ and $I$ to a labelled carrier $D(\Lambda)\subseteq I$ on the branches where the test fails, writes eq:c-one-step-branch-identity for each labelled term, and identifies the labelled coefficient $S^{\circ}_{\Lambda}$ with $\kappa_Ce^{E^{\circ}_{\Lambda}}$ at the retreated carrier, with the defects of every labelled term bounded relative to $\kappa_C$ up to a label-average error paid by one-point witness probabilities.

\section{The Flat Regular-Region Calculation}
\label{sec:flat}

The flat evaluation sets the reconstructed background to the identity. On an all-regular periodic lattice it gives the integral specified in this section. Section~\ref{sec:regional-identification} treats the flat straight-side ratio in a proper conditional integration region. The covariance, constraint substitution, Haar density and cutoff are fixed together; the coefficient depends on this prescription.
Use the fixed faithful unitary realisation $G\subset U(N)$. Identify its Lie algebra $\mathfrak g$ with Hermitian matrices, so its group exponential is $e^{iX}$. Write $\tr_N= N^{-1}\operatorname{Tr}$ and use $\langle X,Y\rangle=\tr_N(XY)$ as the real inner product. The observable is exactly \eqref{eq:wilson-definition}, with the normalised character $\chi(\id)=1$ specified in Section~1. Put $d=\dim\mathfrak g$. The first derivative of $\chi$ vanishes: the trace of its Lie algebra representation vanishes on commutators, and a semisimple Lie algebra is its own commutator algebra.

\subsection{The Covariance and the Normalised Insertion}

Use fine-lattice units in the operators in this section. Let $E$ be the positively oriented bonds of $T$, let $T^{(1)}$ be its coarse centres, and let $E_1$ be the positively oriented coarse bonds. Take the odd blocking factor $L>11$ used by Ba{\l}aban (RG1 p.251; the lemmas of this paper that need only $L\geqslant 3$ are stated so). For the covariance bounds below, $T=\prod_{\mu=1}^4(\mathbb Z/N_{\mu}\mathbb Z)$ with $N_{\mu}= LK_{\mu}$ and $K_{\mu}\geqslant 3$. The block $\mathcal B_y$ consists of the fine sites at displacement $n$ from $y\in T^{(1)}$, with $n\in\mathbb Z^4$ and $| n_{\mu}|\leqslant (L-1)/2$ in each direction. For $x\in\mathcal B_y$, the path family $\mathcal G(y,x)$ traverses each nonzero coordinate displacement in one uninterrupted run, in every ordering of those runs; the path is empty for $x= y$. These are the centred blocks and path family of RG1, Section~0.

The averaging operation is \eqref{eq:global-average}; the formula below is RG1 (0.4), its value on the set $\mathcal O_c$ where every logarithm has its argument in the chart domain (Lemma~\ref{lem:global-average}). CRE permits this choice in the paragraph following (0.1) (p.243); no identification with the different average RG1 (0.12) is made. For a positive coarse bond $c=(y,y+Le_{\mu})$ and $x\in\mathcal B_y$, put $x'= x+Le_{\mu}$. The vectors $e_{\mu}$ are the fine-lattice unit vectors, $[x,x']$ is the straight fine path, and $P_c(U)$ is the transport along the straight fine path from $c_-$ to $c_+$. On the regular logarithm domain, RG1 (0.4) is
\begin{equation}\label{eq:nonlinear-average}
    \begin{aligned}
        \beta_c(U)&= L^{-4}\sum_{x\in \mathcal B_{c_-}}
            \frac{1}{|\mathcal G(c_-,x)|\,|\mathcal G(c_+,x')|}
            \sum_{\substack{\gamma\in \mathcal G(c_-,x)\\
                \gamma'\in \mathcal G(c_+,x')}}
            \frac{1}{i}\log\left(
                U(\gamma)U([x,x'])U(\gamma')^{-1}P_c(U)^{-1}
                \right)\;,\\
        M(U,c)&= e^{i\beta_c(U)}P_c(U)\qquad(U\in\mathcal O_c)\;.
    \end{aligned}
\end{equation}
Every logarithm is applied to a closed relative transport based at $c_-$. Define its logarithmic constraint map and derivative by
\begin{equation}\label{eq:constraint-primitive}
    \mathcal Q(q)=\frac{1}{i}\log M(e^{iq}),
    \qquad Q=\frac{1}{L}D\mathcal Q(0),
    \qquad \mathcal N(q)=\mathcal Q(q)-LQq\;.
\end{equation}
The maps $\mathcal Q$ and $\mathcal N$ take a neighbourhood of zero in $\mathfrak g^E$ to $\mathfrak g^{E_1}$. The designated bond $b_0(c)$ is the fine bond on the straight coarse bond $c$ crossing between its two blocks. Put $E_0=\{ b_0(c)\, | \,c\in E_1\}$ and $E_*= E\setminus E_0$. Let $\iota_0$ and $\iota_*$ extend fields on these two label sets by zero to $E$, and put $Q_0= Q\iota_0$. The designated coefficient matrix $Q_0$ is invertible in this regular construction. The maps are
\begin{equation}\label{eq:constraint-extension}
    h=\frac{1}{L}\iota_0Q_0^{-1},
    \qquad
    \mathsf C=(I-\iota_0Q_0^{-1}Q)\iota_*\;.
\end{equation}
Thus $h \colon \mathfrak g^{E_1}\to\mathfrak g^E$, $\mathsf C \colon \mathfrak g^{E_*}\to\mathfrak g^E$, $LQh= I$ and $Q\mathsf C=0$. The free variable is $B\in\mathfrak g^{E_*}$, and $\mathsf C B$ agrees with $B$ on $E_*$. These are the designated-bond construction in RG1 following (2.10) and the elimination in CRE (1.14)--(1.15) \cite{BalabanCRE,BalabanRG1}.

The gauge-fixing functional is the exponential one in CRE (1.7), not a delta-function axial gauge. For the transports along $\mathcal G(y,x)$, let $\mathcal U_U(y,x)$ be their local group mean: it is the unique solution near their diagonal value of
\begin{equation}
    \sum_{\gamma\in\mathcal G(y,x)}
        \log\bigl(U(\gamma)\mathcal U_U(y,x)^{-1}\bigr)=0\;.
\end{equation}
This is RG1 (0.10)--(0.11). The gauge-fixing functional and the action used in the constrained reference integral are
\begin{equation}\label{eq:action-primitive}
    \begin{aligned}
        \mathscr G(U)&=\sum_{y\in T^{(1)}}
            \sum_{x\in\mathcal B_y\setminus\{ y\}}
                (1-\operatorname{Re}\tr_N\mathcal U_U(y,x))\;,\\
        \mathscr A(U)&= S(U)+\mathscr G(U)\;.
    \end{aligned}
\end{equation}
Here $\partial p$ is the oriented boundary of the plaquette $p$. All group means and logarithms are used in the regular chart containing the identity.

For $a,b\in\mathfrak g^E$, the flat Hessian is defined by
\begin{equation}
    \langle a,\Delta_1b\rangle
        =\left.\frac{\partial^2}{\partial s\,\partial t}
            \mathscr A(e^{i(sa+tb)})\right|_{s= t=0}\;.
\end{equation}
The first derivative is zero at the identity. This is the $J_1=0$ specialisation of CRE (1.16). CRE (1.12) states strict positivity of the second differential at the regular constrained minimum. With the exponential gauge fixing and the constrained extension above, it makes the quadratic form on the finite free carrier strictly positive; no zero-mode subtraction is introduced. The quadratic form and its inverse on that carrier are
\begin{equation}\label{eq:flat-covariance}
    \mathsf H=\mathsf C^*\Delta_1\mathsf C,
    \qquad \mathsf K_B=\mathsf H^{-1},
    \qquad \mathsf K=\mathsf C\mathsf K_B\mathsf C^*\;.
\end{equation}
The adjoint and Lebesgue measure use the stated real inner products, summed over the respective bond labels. Thus $\mathsf K$ is the covariance of the linearised fine field $\mathsf C B$, rather than an unconditioned lattice propagator. The normalised Gaussian measure on the free carrier is denoted by $\mu_{\mathsf K_B}$. In an orthonormal basis of real dimension $D= d| E_*|$, it has density
\begin{equation}
    \dd\mu_{\mathsf K_B}(B)
        =(2\pi)^{-D/2}(\det\mathsf K_B)^{-1/2}
            e^{-\frac{1}{2}\langle B,\mathsf H B\rangle}\,\dd B\;.
\end{equation}

Ba{\l}aban's nonlinear constraint solution gives the fine phase
\begin{equation}\label{eq:flat-phase}
    q_{g_0}(B)= g_0\mathsf C B-h\widetilde D(g_0\mathsf C B),
    \qquad U_{g_0}(b)= e^{iq_{g_0}(B)(b)}\;.
\end{equation}
The correction $\widetilde D \colon \mathfrak g^E\to\mathfrak g^{E_1}$ is defined near zero by the unique small analytic solution
\begin{equation}
    \widetilde D(A)=\mathcal N(A-h\widetilde D(A)),
    \qquad \widetilde D(0)=0\;.
\end{equation}
Its construction is RG1 following (2.10). Since $\mathcal N(0)= D\mathcal N(0)=0$, its expansion starts at degree two. Equation~\eqref{eq:flat-phase} therefore solves the nonlinear averaging constraint and has $q_{g_0}(B)= g_0\mathsf C B+O(g_0^2)$ at fixed $B$.

Fix $\sigma\in \{\mathrm{op},\mathrm{HS}\}$ and use the corresponding matrix norm for both group-matrix and Lie-algebra cutoff tests:
\begin{equation}\label{eq:marginal-two-norms}
    | X|_{\mathrm{op}}=\sup_{| v|_2=1}| Xv|_2,
    \qquad | X|_{\mathrm{HS}}=
        \left(\tr_N(X^*X)\right)^{\frac{1}{2}}\;.
\end{equation}
Here $| v|_2$ is the Euclidean vector norm. The dependence on the fixed $\sigma$ is suppressed in $\vartheta_{g_0}$ and $\mathcal W_{g_0}$ below.

In the all-regular branch, $U_1= U_{1,\square'}=\id$. CRE (1.19) removes the approximate relative-field tests on the good fluctuation cubes. The other retained chart tests in CRE (1.25) have support in $S_1=\varnothing$. The guard statement following \cite[(1.10)]{BalabanCRE} removes the earlier plaquette and gauge cutoffs under its prescribed parameter hierarchy. With that hierarchy retained, \cite[(1.18)--(1.19), (1.25)--(1.26)]{BalabanCRE} gives the following bond product when its bond norm is read as $|\bigcdot |_{\sigma}$:
\begin{equation}\label{eq:literal-cutoff}
    \vartheta_{g_0}(B)=\prod_{b\in E_*}
        \mathbf 1_{\{\,| B(b)|_{\sigma}< r(g_0)\,\}},
    \qquad r(g_0)= A_1(\log g_0^{-2})^{p_0}\;.
\end{equation}
Here $A_1>0$ and $p_0\geqslant 5\mathfrak r$ ($\mathfrak r\geqslant 2$ the logarithmic exponent of the macro-cube sizes, the convention of Section~\ref{sec:sourced-effective-interaction}; Section~\ref{app:thresholds}) are the first-step cutoff parameters, and $\delta_0= g_0r(g_0)$. We use the inductive normalisation following CRE (3.4), $\delta_k= g_kA_1/A_0\,p_0(g_k)$ as printed at p.265, together with CRE (2.4), also at the first step:
\begin{equation}\label{eq:cre-cutoff-normalisation}
    \begin{aligned}
        p_0(g_k)&= A_0(\log g_k^{-2})^{p_0}\;,\\
        \varepsilon_k&= g_kp_0(g_k)\;,\\
        \delta_k&= g_k\frac{A_1}{A_0}p_0(g_k)
            =\frac{A_1}{A_0}\varepsilon_k\;.
    \end{aligned}
\end{equation}
For the first step this gives $\delta_0/g_0= A_1(\log g_0^{-2})^{p_0}$.

For the application to CRE, assume its compatible cube partitions and admissible first-step parameters: sufficiently large cutoff constants, sufficiently small $A_1/A_0$ and coupling, and $p_0\geqslant 5\mathfrak r$, where $\mathfrak r\geqslant 2$ is the logarithmic exponent defining the macro-cube sizes (Section~\ref{sec:sourced-effective-interaction}). Under these hypotheses the guard statement following \cite[(1.10)]{BalabanCRE} removes $\chi_0$ and $\chi_{Ax}$ in the regular region. The reference integral below is specified separately for each of the two matrix-norm prescriptions; its identification with a conditional kernel uses that same prescription.

Let $\sigma_H(q)$ be the density of normalised Haar measure in exponential coordinates relative to the stated Lie-algebra Lebesgue measure, and put $j_H(q)=\sigma_H(q)/\sigma_H(0)$. Define the Jacobian and residual exponent directly by the measure and action pullback:
\begin{equation}\label{eq:literal-residual}
    \begin{aligned}
        J_{g_0}(B)&=
            \left|\det_{\mathfrak g^E}
                \bigl(I-hD\widetilde D(g_0\mathsf C B)\bigr)\right|
            \prod_{b\in E}j_H(q_{g_0}(B)(b))\;,\\
        \mathcal P_{g_0}(B)&=\log J_{g_0}(B)
            -\frac{1}{g_0^2}\mathscr A(U_{g_0})
            +\frac{1}{2}\langle B,\mathsf H B\rangle\;.
    \end{aligned}
\end{equation}
The determinant is on the real fine carrier. The nonlinear change of variables is made before the linear constraints are eliminated, as in CRE (1.14). The subsequent linear-constraint factor, Haar density at zero, coupling rescaling, gauge normaliser and vacuum scalar are common to the numerator and denominator in this branch. After the cutoff reductions above, \eqref{eq:literal-residual} retains the action, Haar and Jacobian factors of CRE (1.14)--(1.16), with the same exponential gauge fixing. Its constant coefficient at $g_0=0$ is zero.

The regular sourced ratio is
\begin{equation}\label{eq:flat-ratio}
    \mathcal W_{g_0}(C)=
        \frac{\displaystyle\int
            \vartheta_{g_0}(B)e^{\mathcal P_{g_0}(B)}
            W_{\chi}(C;U_{g_0})\,\dd\mu_{\mathsf K_B}(B)}
        {\displaystyle\int
            \vartheta_{g_0}(B)e^{\mathcal P_{g_0}(B)}
            \,\dd\mu_{\mathsf K_B}(B)}\;.
\end{equation}
Only common scalar prefactors of this same regular integration cancel here. The denominator is not replaced by an independently chosen Gaussian normaliser.

\section{The Straight-Edge Coefficient from Gaussian Marginals}
\label{sec:marginal-perimeter}

The normalised flat integral has a connected expansion. Its extensive contribution along a long straight side defines $\delta m$; the contributions remaining near the corners define $z_{\mathrm c}$. We construct these terms using Gaussian principal marginals and the local interaction \eqref{eq:perimeter-literal-local-sum}. Section~\ref{sec:regional-identification} then compares the straight-side coefficient with the conditional kernel in a proper regular region.

For a square $C_n$ on the compatible finite torus $T$, the flat reference numerator and denominator are
\begin{equation}\label{eq:perimeter-native-integrals}
    \begin{aligned}
        N_{n,T}&=\int\vartheta_{g_0}(B)e^{\mathcal P_{g_0}(B)}
            W_{\chi}(C_n;U_{g_0})\,\dd\mu_{\mathsf K_B}(B)\;,\\
        D_T&=\int\vartheta_{g_0}(B)e^{\mathcal P_{g_0}(B)}
            \,\dd\mu_{\mathsf K_B}(B)\;.
    \end{aligned}
\end{equation}
These are the identity-background regular integrals. In the two-norm construction below, read the cutoff as $\vartheta_{g_0}^{\sigma}$ and write $N_{n,T}= N_{C_n,T}^{\sigma}$, $D_T= D_T^{\sigma}$ for the chosen prescription $\sigma$. They are distinct from the full physical partition function $\mathcal Z_0$.
\subsection{Connected Terms of the Normalised Integral}

Let $E_*$ be the free-bond carrier and let $C_T=\mathsf K_B$ on a compatible finite torus $T$. The flat gap and exponential decay proved in Appendix~\ref{app:flat-covariance} give positive constants $\lambda,\Lambda,A,\mu$, independent of $T$, such that
\begin{equation}\label{eq:marginal-covariance-input}
    \lambda I\leqslant C_T\leqslant \Lambda I,
    \qquad \| C_T(b,c)\|\leqslant A e^{-\mu d_T(b,c)}\;.
\end{equation}
For positive bonds $b=(x,i)$ and $c=(y,j)$, fix $d_T(b,c)= d_{\infty,T}(x,y)+\mathbf1_{\{ i\ne j\}}$, where the first term is periodic maximum-norm site distance in fine-lattice units. The same bounds hold for every principal marginal. One may take $\lambda= h_L^{-1}$ and $\Lambda=\gamma_L^{-1}$ from \eqref{eq:covariance-uniform-gap} and \eqref{eq:bulk-uniform-gap}. The free-inverse Neumann series \eqref{eq:bulk-neumann-inverse} applies on the compatible tori as well, with the same gap and row-sum bound. Let $d_{\mathrm{block}}(b,c)$ be the periodic maximum-norm distance between the coarse indices of the blocks containing the initial sites of $b,c$. The block-distance decay converts to the fine distance using $d_T(b,c)\leqslant Ld_{\mathrm{block}}(b,c)+2L$; decreasing the decay exponent by the factor $L$ and enlarging $A$ suffices. The metric on each Lie algebra factor is the fixed invariant real metric. All constants below may depend on this metric, $G$, $L$, the fixed local reconstruction and $\chi$, but not on the square size or torus volume.

Throughout this construction, $r_0(g_0)= r(g_0)= A_1(\log g_0^{-2})^{p_0}$ is exactly the threshold in \eqref{eq:literal-cutoff}. Restoring the fixed norm index gives the two cutoff prescriptions
\begin{equation}\label{eq:marginal-two-cutoffs}
    \vartheta_{g_0}^{\sigma}(B)
        =\prod_{b\in E_*}
            \mathbf1_{\{\,| B_b|_{\sigma}< r_0(g_0)\,\}}\;.
\end{equation}
The result holds separately for each prescription, with that cutoff retained in both numerator and denominator. With CRE's bond tests interpreted in the same norm, the regular-branch reduction above identifies its reference kernel. For the following estimates only the stated norm and its adjoint invariance are used. For unitary $u$, the required invariance follows from
\begin{equation}\label{eq:marginal-cutoff-invariance}
    \begin{aligned}
        | uXu^{-1}|_{\mathrm{op}}
            &=\sup_{| v|_2=1}| uXu^{-1}v|_2
                =\sup_{| w|_2=1}| Xw|_2=| X|_{\mathrm{op}}\;,\\
        | uXu^{-1}|_{\mathrm{HS}}^2
            &=\tr_N\bigl(uX^*Xu^{-1}\bigr)
                =\tr_N(X^*X)=| X|_{\mathrm{HS}}^2\;.
    \end{aligned}
\end{equation}
Thus each good- or bad-bond test is invariant for the same fixed norm and threshold; a supremum over the bonds of a cube and products of such tests preserve this invariance. In the regular branch the starred carrier excludes the designated bonds, $A=\mathsf C B$ agrees with $B$ on that carrier, and $g_0^{-1}\delta_0= r_0$ by \eqref{eq:cre-cutoff-normalisation}. Repeated cube tests reduce to the bond product \eqref{eq:literal-cutoff} by idempotence, as in the stated regular-branch reduction. This dictionary does not rescale a cutoff ball. The norm-comparison constants used for Gaussian estimates likewise do not change the cutoff set. The Schur argument requires the invariant cutoff itself.

The numerator and denominator under consideration are
\begin{equation}\label{eq:marginal-literal-ratio}
    \begin{aligned}
        N_{C,T}^{\sigma}&=\int\vartheta_{g_0}^{\sigma}(B)e^{\mathcal P_{g_0}(B)}
            W_{\chi}(C;U_{g_0})\,\dd\mu_{C_T}(B)\;,\\
        D_T^{\sigma}&=\int\vartheta_{g_0}^{\sigma}(B)e^{\mathcal P_{g_0}(B)}
            \,\dd\mu_{C_T}(B)\;,
        \qquad \mathcal W_T^{\circ,\sigma}(C)= N_{C,T}^{\sigma}/D_T^{\sigma}\;.
    \end{aligned}
\end{equation}

For the proof, fix $\sigma$ and suppress its index on the cutoffs, integrals and cluster weights. Write $r= r_0(g_0)$ and $\vartheta_b=\mathbf1_{\{| B_b|_{\sigma}< r\}}$. Let $\mathcal V_T$ index the local terms $v$ in \eqref{eq:perimeter-literal-local-sum}, with supports $\Lambda_v$. Define bounded measurable functions on the full real carrier by
\begin{equation}\label{eq:marginal-local-vertices}
    f_v(B)=\left(\prod_{b\in \Lambda_v}\vartheta_b(B)\right)
        (e^{v(B)}-1),
    \qquad h_b(B)=-(1-\vartheta_b(B))\;.
\end{equation}
The first function is zero where its cutoff vanishes, so the local analytic function is never evaluated outside its chart. On the original cutoff the exact identity is
\begin{equation}\label{eq:marginal-original-weight}
    \vartheta_{g_0}(B)e^{\mathcal P_{g_0}(B)}
        =\prod_{b\in E_*}(1+h_b(B))
            \prod_{v\in \mathcal V_T}(1+f_v(B))\;.
\end{equation}
For a square $C=(b_a)_{a\in\mathcal I_C}$, let $\mathcal I_C$ be its ordered finite traversal set and let $U_0=\id$. For $R\subseteq\mathcal I_C$, let $H_R(B)$ be the ordered transport around $C$ using $U_{g_0}(b_a)$ on $R$ and $U_0(b_a)$ elsewhere. The closed source differences are
\begin{equation}\label{eq:finite-source-closed-differences}
    G_S(B)=\sum_{R\subseteq S}(-1)^{| S|-| R|}
        \chi\left(H_R(B)(\Hol_{U_0}C)^{-1}\right),
    \qquad \sum_{S\subseteq\mathcal I_C}G_S(B)
        = W_{\chi}(C;U_{g_0})\;.
\end{equation}
The last equality is Boolean inversion and uses $\Hol_{U_0}C=\id$. Every summand is a closed contraction, invariant under simultaneous gauge transformation of field and reference. The identity-gauge representative is used to evaluate its flat reference integral.
A selected traversal $a$ has its fixed local free-input support $\Lambda_a$. Multiplying $G_S$ by $\prod_{a\in S}\prod_{b\in \Lambda_a}\vartheta_b$ leaves its contribution in the original cutoff integral unchanged. Define this product to be zero off that protected box, evaluating only the reconstruction rows selected by the mixed word.

A label selection is $J=(D,P,S)$, with $D\subseteq \mathcal V_T$, $P\subseteq E_*$ and $S$ a subset of contour traversals. Its active carrier and scalar local factor are
\begin{equation}\label{eq:marginal-selection}
    \begin{aligned}
        A_J&=\bigcup_{v\in D}\Lambda_v\cup P
            \cup\bigcup_{a\in S}\Lambda_a\;,\\
        F_J(B)&=\prod_{v\in D}f_v(B)\prod_{b\in P}h_b(B)
            \prod_{a\in S}\prod_{b\in \Lambda_a}\vartheta_b(B)\;.
    \end{aligned}
\end{equation}
The corresponding numerator term is exactly $\int F_JG_S\,\dd\mu_{(C_T)_{A_J}}$. This is a marginal of the original Gaussian: its coordinates in $E_*\setminus A_J$ have been integrated out, while every interaction and cutoff is still present in the sum over $J$. No retained exterior field is given a new law. There are fixed integers $s_*,q_*$ and a fixed range $R_*$ such that every label support has at most $s_*$ bonds and diameter at most $R_*$, and at most $q_*$ possible labels meet any bond. Enlarge $R_*$ so that every bond in $\Lambda_a$ is within $R_*$ of the source traversal $a$; for a reversed traversal use the free-input stencil of its positive representative. These are consequences of the bounded overlap of the primitive supports and the fact that a square has no repeated traversal. The source cutoff is inserted before expanding the full product of cutoff factors in \eqref{eq:marginal-original-weight}; an individual expanded term need not satisfy the full cutoff.

For a finite carrier $A$, put $C_A=(C_T)_A$. Introduce one scalar variable $x_{bc}= x_{cb}$ for each pair of distinct bonds of $A$, and define $C_A(x)$ by keeping its diagonal blocks and multiplying its $bc$ block by $x_{bc}$. For a forest $F$ on $A$ and parameters $w\in [0,1]^F$, let $w^F_{bc}$ be the minimum of the parameters on the unique forest path from $b$ to $c$, or zero if there is no path; its diagonal is one. Put $\partial_F=\prod_{\{ b,c\}\in F}\partial_{x_{bc}}$. The Taylor forest formula gives
\begin{equation}\label{eq:marginal-forest-identity}
    \int F_JG_S\,\dd\mu_{C_{A_J}}
        =\sum_{F\ \mathrm{forest\ on}\ A_J}
            \int_{[0,1]^F}\int F_J(B)G_S(B)
                \left.\partial_F p_{C_{A_J}(x)}(B)\right|_{x= w^F}
                    \,\dd B\,\dd w\;.
\end{equation}
For a covariance on $\mathfrak g^A$, the density used here is
\begin{equation}\label{eq:marginal-density-definition}
    p_M(B)=(2\pi)^{-d| A|/2}(\det M)^{-1/2}
        \exp\left(-\frac{1}{2}B^{\mathsf T}M^{-1}B\right)\;.
\end{equation}
The determinant square root is positive for positive real $M$ and is continued analytically on the accretive covariance domain. Let $P_{\pi}$ be the matrix whose $bc$ entry is one when $b,c$ belong to the same block of a partition $\pi$ of $A$, and zero otherwise. Let $\mathcal K_A$ be the convex hull of these partition matrices. If $\pi_t$ is the partition into components of the forest edges with $w_e\geqslant t$, then $w^F=\int_0^1P_{\pi_t}\,\dd t\in \mathcal K_A$. Block restriction and convexity give $\lambda I\leqslant C_A(x)\leqslant \Lambda I$ throughout $\mathcal K_A$.

For the fixed integrand $F_JG_S$, its Gaussian integral is smooth on the open set $\{ x\, | \,C_A(x)>(\lambda/2)I\}$: the protected local factors are bounded, and derivatives of the density are polynomials times a uniformly integrable Gaussian on compact subsets. This set contains the compact set $\mathcal K_A$. A smooth multiplier with compact support inside this open set and equal to one near $\mathcal K_A$ gives a smooth extension by zero to the full interpolation space. The Taylor forest formula applies to this extension. At the all-one matrix and at every $w^F$, its value and all the required derivatives agree with those of the original Gaussian integral. The extension therefore gives exactly \eqref{eq:marginal-forest-identity}; its cutoff and derivatives enter none of the estimates. No hard cutoff or Wilson function is differentiated. This is the Taylor forest construction of Abdesselam--Rivasseau, Theorems III.1 and IV.5 and Section V.A \cite{AbdesselamRivasseau}.

Join two bonds also when they belong to the same selected label support, and denote this finite-range graph by $Q_J$. A connected component of $F\cup Q_J$, together with its selected labels and forest edges, is a component label $\gamma$. Let $A_{\gamma},D_{\gamma},P_{\gamma},S_{\gamma},F_{\gamma}$ be these restrictions. Define its scalar integration functional by
\begin{equation}\label{eq:marginal-component-functional}
    \mathcal I_{\gamma}[H]
        =\int_{[0,1]^{F_{\gamma}}}\int
            F_{J_{\gamma}}(B)H(B)
            \left.\partial_{F_{\gamma}}p_{C_{A_{\gamma}}(x)}(B)
            \right|_{x= w^{F_{\gamma}}}
                \,\dd B\,\dd w\;.
\end{equation}
No norm of an open Wilson object is assigned to this functional. The covariance and every differentiated edge in \eqref{eq:marginal-forest-identity} are block diagonal across distinct components of $F\cup Q_J$. Their normalised densities, including determinants, consequently factor exactly into \eqref{eq:marginal-component-functional}. Every component uses the principal marginal of the same $C_T$; it is independent of the labels outside its carrier.

At the flat reference $C_T$ commutes with simultaneous constant conjugation. Principal restriction, scalar covariance interpolation, inversion and differentiation preserve this law. Every local scalar factor in $F_J$ is invariant as well. Thus each component integration is equivariant without any assertion about seed-adapted auxiliary inverses. For $S_{\gamma}\ne \varnothing$, let $h_{\gamma}$ be the shortest cyclic interval of traversal positions containing $S_{\gamma}$, resolving ties by the first initial position encountered from the fixed contour base in its orientation. A singleton has an interval of length one. On a line, use the ordinary interval from the first to the last selected traversal. Merge components whose intervals intersect into webs; the web interval is their union, possibly the whole circle. This union is an interval because its interval-intersection graph is connected. Distinct web intervals are disjoint. A web is integrated as the single closed source $G_{\cup_{\gamma\in W}S_{\gamma}}$ against the product of its component functionals. Its scalar activity is
\begin{equation}\label{eq:marginal-web-activity}
    z_W=\left(\bigotimes_{\gamma\in W}\mathcal I_{\gamma}\right)
        \left[G_{\cup_{\gamma\in W}S_{\gamma}}\right]\;.
\end{equation}
The vacuum activity is $z_{\gamma}=\mathcal I_{\gamma}[1]$ when $S_{\gamma}=\varnothing$. Components inside a web have disjoint carriers. Distinct webs must also have disjoint contour intervals. Vacuum objects are incompatible only through intersection of their carriers with another object's carrier. Independent constant conjugation in the component integrals and irreducibility give the integrated Schur factorisation for distinct webs. It is applied inside the closed Wilson contraction. To specify the logarithm domain, attach a complex variable $t_a$ to each contour traversal and replace $G_S$ by $(\prod_{a\in S}t_a)G_S$. The source polydisc is $| t_a|\leqslant 2$. Its factor $2^{| S|}$ is included in the constants of the source bounds below. At zero source the numerator is $D_T$, and at $t_a=1$ it is $N_{C,T}$. The finite forest expansion and the unique component and web partitions consequently give
\begin{equation}\label{eq:marginal-exact-gases}
    N_{C,T}=\Xi_{C,T},\qquad D_T=\Xi_{0,T},\qquad
    \mathcal W_T^{\circ}(C)=\frac{\Xi_{C,T}}{\Xi_{0,T}}\;.
\end{equation}
Both gases are sums over unordered compatible families of the activities just defined, with the empty family contributing one. There is no additional vacuum ratio to identify.

For the converse in this identification, take a compatible family of vacuum components and webs. Their constituent carriers are disjoint, so their label selections and forests have unambiguous unions. Every support edge and forest edge remains within its specified component. The component decomposition of this union recovers the original components, and the interval-intersection decomposition recovers precisely the webs, since distinct web intervals are disjoint. Thus the reassembly is inverse to the decomposition. If a web interval crosses the contour base, cyclicity of the trace moves that base between two web intervals before the Schur factorisation; the resulting scalar is unchanged.

\subsection{The Integrated Density Estimate}

Fix a forest on $A$, let $n=| A|$, and let $d_b$ be its vertex degrees. For an edge $e=\{ b,c\}$ with nonzero covariance block, use the complex Cauchy radius
\begin{equation}\label{eq:marginal-edge-radius}
    r_e=\frac{\lambda}{8d_bd_c\| C_T(b,c)\|}\;.
\end{equation}
A zero covariance block gives a zero differentiated term. Perturb only the forest variables about $w^F$. The sum of the norms of the perturbation blocks in row $b$ is at most $\lambda/8$, since $\sum_{c:\{ b,c\}\in F}(d_bd_c)^{-1}\leqslant 1$. Thus the covariance perturbation has operator norm at most $\lambda/8$. If $M$ is any covariance on these Cauchy circles, then
\begin{equation}\label{eq:marginal-complex-density}
    \operatorname{Re}M\geqslant \frac{7\lambda}{8}I,
    \qquad \| M\|\leqslant \Lambda+\frac{\lambda}{8},
    \qquad
    \operatorname{Re}M^{-1}\geqslant
        \frac{7\lambda/8}{(\Lambda+\lambda/8)^2}I=\kappa_G I\;.
\end{equation}
The inverse bound follows from $\operatorname{Re}M^{-1}= M^{-*}(\operatorname{Re}M)M^{-1}$. The determinant branch is continued from the positive real covariance. Its modulus satisfies $|\det M|\geqslant (7\lambda/8)^{d n}$. Consequently
\begin{equation}\label{eq:marginal-density-majorant}
    | p_M(B)|\leqslant
        (2\pi)^{-dn/2}(7\lambda/8)^{-dn/2}
            e^{-\kappa_G\| B\|^2/2}\;.
\end{equation}
This keeps the normalising determinant and supplies a product Gaussian majorant. Norm equivalence and a Gaussian exponential moment give a factor $C e^{-c r^2}$ for each distinct selected bad bond. Each regular vertex is bounded on its own cutoff by $C g_0(1+r)^3$, using \eqref{eq:perimeter-literal-local-bound}. The closed source with $q$ selected traversals has bound $(C_{\chi} g_0r)^q$; for $q=1$ it has the improved bound $C_{\chi} g_0^2r^2$, since $\dd\chi_{\id}=0$. These bounds are applied to the closed source inside the density integral.

Cauchy's formula and \eqref{eq:marginal-density-majorant} therefore bound the integrated term for a fixed component by
\begin{equation}\label{eq:marginal-fixed-forest-bound}
    K^{| A_{\gamma}|}
        \prod_{\{ b,c\}\in F_{\gamma}}
            \left(\frac{8d_bd_c}{\lambda}\| C_T(b,c)\|\right)
        \tau^{| D_{\gamma}|+| P_{\gamma}|+| S_{\gamma}|
            +\mathbf1_{\{| S_{\gamma}|=1\}}},
    \qquad
    \tau= C\left[g_0(1+r)^3+e^{-cr^2}\right]\;.
\end{equation}
Here and below $g_0$ is small enough that $\tau<1$ and the original chart condition holds. The same majorants multiply for a web. To justify the extra power for a component with exactly one source mark inside a larger web, integrate that component first inside the closed contraction. Its matrix factor is a scalar by constant conjugation and irreducibility, and its scalar is the closed one-mark character difference. Components with two or more marks already supply two source powers. This uses an exact integrated identity; it assigns no activity norm to an open matrix.

The degree factors in \eqref{eq:marginal-fixed-forest-bound} are paid by spatial decay. The free-bond lattice has uniformly bounded polynomial volume growth of degree four, including on its finite tori. For $d$ distinct neighbours of a fixed bond, their distances have sum at least $c_1d^{5/4}-c_2d$. Hence, for each $\zeta>0$, a constant $K_{\zeta}$ satisfies
\begin{equation}\label{eq:marginal-degree-payment}
    d^d\leqslant K_{\zeta}^d
        \exp\left(\zeta\sum_{j=1}^d d_T(b,c_j)\right)\;.
\end{equation}
Indeed $d\log d$ is bounded by any fixed positive multiple of $d^{5/4}$ plus a linear function of $d$. Set $0^0=1$ here. Taking $\zeta$ smaller than $\mu/4$ absorbs $\prod_b d_b^{d_b}$ into part of the covariance decay. The remaining forest edge weights are bounded by $K e^{-\mu d_T(b,c)/2}$, with a further positive part of the decay available for support-length weights.

Here is the connected summation. On a carrier of $n$ bonds there are at most $2^{q_*n}$ choices of elementary labels. A selection covering that carrier contains at least $\lceil n/s_*\rceil$ labels. The edges of $Q_J$ have bounded range and bounded degree. For fixed $J$ and a fixed graph $G= F\cup Q_J$, the edges of $F$ outside $Q_J$ are determined; its edges inside $Q_J$ have at most $2^{| Q_J|}$ choices. The support counts give $| Q_J|\leqslant q_*(s_*-1)n/2$. Retain this exponential multiplicity in $K_1$ below. Each forest edge weight is bounded by the following pair weight, and this weight is at least one on every edge of $Q_J$. Thus the forest sum is bounded, with that multiplicity, by a sum over connected graphs with pair weights
\begin{equation}\label{eq:marginal-counting-weight}
    v(b,c)=\mathbf1_{\{ d_T(b,c)\leqslant R_*\}}
        +K e^{-\mu d_T(b,c)/4},
    \qquad \sup_b\sum_{c\ne b}v(b,c)= V_*<\infty\;.
\end{equation}
The unused graph edges cost at most $\prod_{b< c}(1+v(b,c))\leqslant e^{V_*n/2}$. Keeping one spanning tree gives the remaining tree sum. For $n\geqslant 2$ and a fixed root bond $b_0$, summing its possible distinct vertices gives
\begin{equation}\label{eq:marginal-rooted-tree-count}
    \sum_{\substack{A\ni b_0\\| A|= n}}
        \sum_{T\ \mathrm{tree\ on}\ A}\prod_{\{ b,c\}\in T}v(b,c)
        \leqslant \frac{n^{n-2}}{(n-1)!}V_*^{n-1}\;.
\end{equation}
The factorial accounts for unordered nonroot vertices; the bound follows by summing their positions along a rooted labelled tree. The right-hand side is bounded exponentially in $n$. Combining these payments yields convergent series bounded by constants times
\begin{equation}\label{eq:marginal-load-series}
    \sum_{n\geqslant 1}K_1^n\tau^{\lceil n/s_*\rceil}= O(\tau),
    \qquad
    \sum_{n\geqslant 1}K_1^n\tau^{\max(2,\lceil n/s_*\rceil)}= O(\tau^2)\;.
\end{equation}
The first applies to vacuum components and the second to source components, rooted either at a spatial bond or at a specified source traversal. Choose the coupling after the constants so that $K_1^{s_*}\tau$ is sufficiently small. The argument leaves exponential weights in the carrier size and in the sum of covariance-edge lengths.

\subsection{Extent Weights and the Square Geometry}\label{sec:extent-weights}

Define the extent cost of a component, a web and a cluster respectively by
\begin{equation}\label{eq:marginal-extent}
    \ell(\gamma)=| A_{\gamma}|+
        \sum_{\{ b,c\}\in F_{\gamma}}d_T(b,c),
    \qquad \ell(W)=\sum_{\gamma\in W}\ell(\gamma),
    \qquad \ell(\Gamma)=\sum_{p\in \Gamma}\ell(p)\;.
\end{equation}
A vacuum object uses its component cost. Write $A_p$ for an object's carrier and $H_p$ for its contour interval length, with $H_p=0$ for vacuum objects. On a straight line, connectivity of $F_{\gamma}\cup Q_{J_{\gamma}}$ gives $H_{\gamma}\leqslant c_h\ell(\gamma)$, since the support edges have length at most $R_*$ and the source positions are within $R_*$ of their supports; $c_h$ is fixed at \eqref{eq:marginal-hull-extent} below, where the square is treated, and the same constant covers the line, the interval length there being at most the spatial span of the marks plus one. ($H_p$, the contour interval length, is also called the contact-interval length below.)

Here the square $C_n$ has vertices $0,nLe_1,nL(e_1+e_2),nLe_2$ in a fixed lift, with all other coordinates zero. Its side length is $nL$ fine bonds, its physical perimeter is $4nL\varepsilon$, and its corners are coarse sites. The finite torus has fine periods $LK_i$, with $K_i\geqslant 3(n+2)$ and the native grid divisibilities. The limits taken below (the joint limit in $n$ and the periods of \eqref{eq:marginal-square-mass}, the torus limit at fixed $n$ and the iterated limit of the corner section) are over such compatible integers $n$ and periods. Between points of this square the periodic coordinate distances equal their distances in the stated lift. This fixes one coordinate plane, one position and one orientation. An oriented square contour $C$ as chosen at the start of the paper, with coarse corners and sides along two coordinate directions, is a coarse translate of the square with vertices $0,nLe_{\mu},nL(e_{\mu}+e_{\nu}),nLe_{\nu}$, traversed in this cyclic order or in the reverse one (the base point does not affect $\chi(\operatorname{Hol})$). The permutation of the coordinate axes carrying $(e_{\mu},e_{\nu})$ to $(e_1,e_2)$, with the periods permuted accordingly, is a symmetry of the fine lattice, of the Wilson measure $\rho_0(U)\,\dd U$ and of every hypothesis on the periods (all are imposed on every period), so $\mathbb E_{\varepsilon}[W_{\chi}(C)]$ equals the expectation of a coarse translate of $C_n$, which is the contour admitted in Section~\ref{sec:c5-5-0} and in \eqref{eq:c5-7-T}(i). Reversing the orientation inverts the holonomy, and $\chi(U^{-1})=\overline{\chi(U)}$ for the character of a unitary representation, so the expectation of the reversed square is the complex conjugate of that of the original. Every statement below about $\mathbb E_{\varepsilon}[W_{\chi}(C_n)]$ therefore holds for every such square, up to complex conjugation for the reversed orientation.

Let $D$ be the maximum-norm diameter of a nonempty set of square contacts and put $s= nL$. If $D< s/2$, its contacts occupy one side or two adjacent sides near their common corner, and an arc of length at most $4D+1$ contains them. If $D\geqslant s/2$, the whole contour has length $4s\leqslant 8D$. A connected support network for the component has length at most $R_*| A_{\gamma}|+\sum_{F_{\gamma}}d_T(b,c)$, and each contact is within $R_*$ of that network. With the stated bond-distance convention, $c_h=32(1+R_*)$ (the interval constant, italic subscript; distinct from the group constant $c_{\mathrm h}=\sqrt N$ of Appendix~\ref{app:background-estimates}) therefore gives
\begin{equation}\label{eq:marginal-hull-extent}
    H_{\gamma}\leqslant c_h\ell(\gamma),
    \qquad H_W\leqslant \sum_{\gamma\in W}H_{\gamma}
        \leqslant c_h\ell(W)\;.
\end{equation}
Let $d_C$ count the shortest number of traversal steps between two positions, cyclically for a square. For $\theta>0$, a position $a$ in $h_{\gamma}$ satisfies
\begin{equation}\label{eq:marginal-hull-root}
    \mathbf1_{\{ a\in h_{\gamma}\}}
        \leqslant \sum_{b\in S_{\gamma}}
            e^{-\theta d_C(a,b)}e^{\theta H_{\gamma}}\;.
\end{equation}
The sum of $e^{-\theta d_C(a,b)}$ over all traversal positions is bounded by $1+2/(e^{\theta}-1)$, uniformly in the contour length. The already proved traversal-root load, with the extra extent weight, therefore gives an interval-root load of order $\tau^2$.

For a source component let $u_{\gamma}$ be the positive majorant on the right-hand side of \eqref{eq:marginal-fixed-forest-bound}, with the constants fixed there, and let $H_{\gamma}=| h_{\gamma}|$ be its contact-interval length. To sum webs, majorise the indicator that their interval graph is connected by its spanning-tree sum. The rooted-tree recursion is bounded by $u_{\gamma} e^{aH_{\gamma}}$ whenever the interval-root load, with this extra weight, is at most $a$. This follows by iterating $T_{\gamma}= u_{\gamma}\exp(\sum_{\gamma':h_{\gamma'}\cap h_{\gamma}\ne \varnothing}T_{\gamma'})$ from $T_{\gamma}= u_{\gamma}$, since the sum in the exponent is bounded by $aH_{\gamma}$. Spatial incompatibilities only reduce this majorant. The resulting whole-web loads are $O(\tau^2)$. Together with the vacuum load $O(\tau)$, they satisfy the hard-core polymer convergence criterion \cite{KoteckyPreiss}, with size equal to carrier size plus a small positive multiple of contour interval length. This proves absolute convergence of the connected logarithm of \eqref{eq:marginal-exact-gases} for sufficiently small fixed $g_0$. In particular the numerator has no zero on the source-scaling polydisc used by that expansion. Its logarithm is fixed by continuation from zero source, where numerator and denominator agree.

More explicitly, choose $\beta=\mu/(1024(1+c_h))$ and $\eta=\mu/1024$, and put $a(p)=| A_p|+\beta H_p$. The density and graph payments above remain valid after multiplying the component majorants by $e^{4| A_{\gamma}|+8\beta H_{\gamma}+8\eta\ell(\gamma)}$: the added edge exponent is smaller than the unused covariance decay, and the added carrier exponent only changes $K_1$. Use $a=\beta$ in the web recursion, reserving the remaining weights for its interval-root conversion. It then gives a constant $C_*$ such that
\begin{equation}\label{eq:marginal-weighted-loads}
    \begin{aligned}
        \sup_b\sum_{p\, | \,b\in A_p}
            | z_p| e^{2a(p)+2\eta\ell(p)}&\leqslant C_*\tau\;,\\
        \sup_a\sum_{W\, | \,a\in h_W}
            | z_W| e^{2a(W)+2\eta\ell(W)}&\leqslant C_*\tau^2\;.
    \end{aligned}
\end{equation}
The source-only version of the first sum is also bounded by $C_*\tau^2$. Choose $g_0$ after these constants so that $C_*\tau\leqslant 1/4$ and $C_*\tau^2\leqslant \beta/4$. Carrier intersection and interval intersection then give the criterion
\begin{equation}\label{eq:marginal-KP-payment}
    \sum_{q\not \sim p}| z_q| e^{a(q)+\eta\ell(q)}
        \leqslant C_*\tau| A_p|+C_*\tau^2H_p
        \leqslant \frac{1}{4}a(p)\;.
\end{equation}
A cluster $\Gamma=(p_1,\ldots,p_j)$, with $j\geqslant 1$, is an ordered tuple of vacuum components or source webs with connected incompatibility graph; repetitions are allowed and an object is incompatible with itself. Define its scalar weight by
\begin{equation}\label{eq:marginal-cluster-weight}
    \Phi_C(\Gamma)=\frac{1}{j!}
        \left(\sum_{H\ \mathrm{connected\ on}\ \{1,\ldots,j\}}
            \prod_{\{ r,s\}\in H}(-\mathbf1_{\{ p_r\not \sim p_s\}})\right)
        \prod_{r=1}^j z_{p_r}\;.
\end{equation}
For precision, the pinned version used here retains an activity factor. Put $\widetilde z_p= z_pe^{\eta\ell(p)}$, and let $m_p(\Gamma)$ count the occurrences of $p$ in the ordered cluster tuple. Then
\begin{equation}\label{eq:marginal-pinned-bound}
    \sum_{\Gamma} m_p(\Gamma)|\Phi_C(\Gamma)| e^{\eta\ell(\Gamma)}
        \leqslant | z_p| e^{\eta\ell(p)+a(p)}\;.
\end{equation}
Here is a direct justification. For a fixed incompatibility graph on labelled vertices, partition its connected subgraphs by their breadth-first spanning tree rooted at vertex one, choosing the least-labelled parent at the previous level. A tree's interval of subgraphs allows all edges within a level and the edges to the previous level whose other endpoint is no earlier than the chosen parent. These intervals partition the connected graphs. The alternating sum on each interval is zero or has modulus one. Therefore the absolute Ursell coefficient is bounded by the sum over its spanning trees. The positive generating sum of rooted trees with root species $p$ and vertex weights $|\widetilde z_q|$ is the monotone limit of
\begin{equation}\label{eq:marginal-pinned-recursion}
    T_p^{(0)}=|\widetilde z_p|,
    \qquad
    T_p^{(k+1)}=|\widetilde z_p|
        \exp\left(\sum_{q\not \sim p}T_q^{(k)}\right)\;.
\end{equation}
Permutation symmetry of the ordered tuple gives $\sum_{\Gamma} m_p(\Gamma)F(\Gamma)= j\sum_{\Gamma\, | \,p_1= p}F(\Gamma)$ at each fixed length $j$ for its symmetric weight $F$. Hence the factor $1/j!$ becomes $1/(j-1)!$, including repeated species. The exponential in \eqref{eq:marginal-pinned-recursion} assembles an unordered family of child trees, with its factorial for repeated choices; $T_p^{(k)}$ enumerates rooted trees of height at most $k$. Equation~\eqref{eq:marginal-KP-payment} implies inductively $T_p^{(k)}\leqslant |\widetilde z_p| e^{a(p)}$, proving \eqref{eq:marginal-pinned-bound}. For a countable polymer family, first restrict the species to a finite subset. The same upper bound holds for every such restriction, and monotone convergence of the positive majorants proves the pinned bound for the countable family. No total infinite-volume vacuum partition function is asserted.

Apply this bound to source objects meeting a fixed coarse contour segment $x$. Their weighted activity sum is bounded by the source-only loads in \eqref{eq:marginal-weighted-loads}. Thus
\begin{equation}\label{eq:marginal-cluster-tail}
    \sum_{\Gamma\, | \,x\in J_{\Gamma}}
        |\Phi_C(\Gamma)| e^{\eta\ell(\Gamma)}
        \leqslant C\tau^2\;.
\end{equation}
Here $J_{\Gamma}$ consists of the coarse segments containing its source marks. A coarse segment contains $L$ fine traversal positions, which changes only the constant. The same proof applies on the infinite lattice, using its finite principal covariance marginals.

A spanning tree of $F_{\gamma}\cup Q_{J_{\gamma}}$ connects each component carrier. Its metric length is at most $R_*| A_{\gamma}|+\sum_{F_{\gamma}}d_T(b,c)$. Paths of length at most $R_*$ attach its source contacts; there are at most $q_*| A_{\gamma}|$ contacts. A web becomes connected after adjoining its contour intervals, whose total length is at most $c_h\ell(W)$ by \eqref{eq:marginal-hull-extent}. In a cluster, incompatible objects meet through carriers or contour intervals, so these networks form a connected union. The carrier bonds themselves add at most $| A_{\gamma}|$ per component, and conversion from maximum-norm distances to lattice path lengths costs at most four. Thus its length and diameter are bounded by $C_{\mathrm{net}}\ell(\Gamma)$ with $C_{\mathrm{net}}=8(1+R_*)(1+q_*+c_h)$. This remains an upper bound with repeated objects. Equation~\eqref{eq:marginal-cluster-tail} therefore controls spatial extent as well as contour span.

\subsection{The Perimeter Coefficient}

Use the infinite-lattice covariance constructed from the same periodic block grid. Its finite principal marginals satisfy the same spectral and decay bounds, so the density, counting and weighted KP estimates above apply with the same constants. For an orientation $v\in \mathscr V_{\square}=\{ e_1,e_2,-e_1,-e_2\}$, use the straight line through the coarse origin in that orientation and divide it into half-open segments of $L$ fine traversals, indexed by $\mathbb Z$. Every side of $C_n$ is a coarse translation of the corresponding oriented segment of this line. Each square traversal belongs to exactly one side and one such segment. The lattice spacing $\varepsilon$ is the fixed physical length of a fine bond.

Write $\Phi_v$ for the connected weight \eqref{eq:marginal-cluster-weight} on the straight line in orientation $v$.
For a finite set $S$ of traversal positions on the line, take any finite oriented arc $I$ containing it. For $R\subseteq S$, let $u_a^R= U_{g_0}(b_a)$ on $R$ and $u_a^R= U_0(b_a)$ otherwise, with $U_0=\id$. The closed relative source on the line is explicitly
\begin{equation}\label{eq:marginal-line-source}
    G_S^v(B)=\sum_{R\subseteq S}(-1)^{| S|-| R|}
        \chi\left(\left(\prod_{a\in I}^{\longrightarrow}u_a^R(B)\right)
            P_{U_0}(I)^{-1}\right)\;.
\end{equation}
The arrow means order along the arc. The reference reverse transport closes the word, and extending the arc beyond its selected marks adds only identity factors in the reference gauge. This is the line version of the finite closed source definition used above. Let $J_{\Gamma}\subset \mathbb Z$ be the set of coarse line segments containing at least one source mark of the cluster. The vacuum subtraction in $\log\Xi_C-\log\Xi_0$ retains exactly the clusters with $J_{\Gamma}\ne \varnothing$. Translation by one coarse segment preserves the reference construction. Define the absolutely convergent line density and its coefficient by
\begin{equation}\label{eq:marginal-full-perimeter}
    \boxed{\displaystyle
        \kappa_v(g_0)=\sum_{\Gamma\, | \,0\in J_{\Gamma}}
                    \frac{\Phi_v(\Gamma)}{| J_{\Gamma}|},
        \qquad \delta m_v(g_0,\varepsilon)
            =-\frac{\kappa_v(g_0)}{L\varepsilon}}
        \;.
\end{equation}
The factor $| J_{\Gamma}|^{-1}$ allocates each translated connected contribution once. The weights contain the full local exponentials and exact hard-cutoff factors. The weighted rooted estimate \eqref{eq:marginal-cluster-tail} proves absolute convergence and $|\kappa_v|\leqslant C\tau^2$.

\paragraph{Length allocation and corners.}
Put $I_m=\{0,\ldots,m-1\}$. The infinite-lattice logarithm for the finite straight relative word of length $mL$ is
\begin{equation}\label{eq:marginal-finite-line-log}
    F_v(m)=\sum_{\Gamma\, | \,\varnothing\ne J_{\Gamma}\subseteq I_m}
        \Phi_v(\Gamma)\;.
\end{equation}
It is the bulk limit of the logarithm of the same finite-word numerator divided by $D_T$. For fixed $m$, this follows by truncating \eqref{eq:marginal-cluster-tail}, taking the limit of finitely many covariance-marginal integrals, and then removing the truncation. No infinite-volume probability measure is used. Translation invariance gives the exact difference
\begin{equation}\label{eq:marginal-line-endpoints}
    F_v(m)-m\kappa_v
        =-\sum_{x\in I_m}
            \sum_{\substack{\Gamma\, | \,x\in J_{\Gamma}\\
                                      J_{\Gamma}\not \subseteq I_m}}
                \frac{\Phi_v(\Gamma)}{| J_{\Gamma}|}\;.
\end{equation}
A cluster in the inner sum joins a source in segment $x$ to a source outside $I_m$. With $R_{\mathrm{end}}=2L+4R_*$, the network bound gives $C_{\mathrm{net}}\ell(\Gamma)\geqslant (L\min(x+1,m-x)-R_{\mathrm{end}})_+$, where $t_+=\max(t,0)$. Applying \eqref{eq:marginal-cluster-tail} and summing that geometric tail yields
\begin{equation}\label{eq:marginal-line-density-limit}
    | F_v(m)-m\kappa_v|\leqslant C\tau^2\;.
\end{equation}

For the square in the infinite-lattice reference, clusters whose source marks all lie on one side have exactly the corresponding line weights: their contact intervals stay on that side, their closed words have the same ordered nonidentity factors, and every local primitive is translated by a multiple of $L$. The Gaussian marginal is unchanged by this translation. All other source clusters touch at least two sides. Allocate each such cluster over its source segments using $| J_{\Gamma}|^{-1}$. A root segment at distance $j$ from its nearest corner can be joined to another side only by a cluster with $C_{\mathrm{net}}\ell(\Gamma)\geqslant (Lj-R_{\mathrm{end}})_+$. Thus the sum over mixed-side clusters is bounded by $C\tau^2\sum_{j\geqslant 0}e^{-c'j}$, with $c'=\eta L/C_{\mathrm{net}}$ (the extent weight $e^{\eta\ell(\Gamma)}$ of \eqref{eq:marginal-cluster-tail}; the factor $e^{\eta R_{\mathrm{end}}/C_{\mathrm{net}}}$ is absorbed in $C$). With $\widehat F_{\square}(n)$ denoting its absolutely convergent infinite-lattice connected sum, this proves
\begin{equation}\label{eq:marginal-square-corners}
    \left|\widehat F_{\square}(n)
        -\sum_{v\in \mathscr V_{\square}}F_v(n)\right|
        \leqslant C\tau^2,
    \qquad
    \left|\widehat F_{\square}(n)
        -n\sum_{v\in \mathscr V_{\square}}\kappa_v\right|
        \leqslant C\tau^2\;.
\end{equation}
This retains the bounded corner contribution, whose extraction is carried out in Section~\ref{sec:marginal-corner}.

\paragraph{The buffered finite-torus limit.}
The perimeter coefficient also follows along the stated buffered tori, without a separate uniform norm on differences of all activities. Fix an integer extent cutoff $R$. A cluster rooted at a specified source segment with $\ell(\Gamma)\leqslant R$ has all its sites and marks in a neighbourhood of radius $C_{\mathrm{net}}R+L+2R_*$ in fine-lattice units around that segment. There are finitely many such cluster shapes, including multiplicities: each object has extent at least one, so a tuple has at most $R$ objects. Away from the square corners by more than $\left\lceil\frac{C_{\mathrm{net}}R+2L+2R_*}{L}\right\rceil$ coarse segments, these are the same shapes as on the corresponding straight line.

Translate each such root to segment zero. Let $\epsilon_R(T)$ be the sum, over the four orientations and these finitely many shapes, of $|\Phi^T_v(\Gamma)-\Phi_v(\Gamma)|/| J_{\Gamma}|$. The finite-torus weight here is the same finite Gaussian-marginal integral on the translated sites. For large enough periods the shape embeds without wraparound. Periodisation gives entrywise convergence of all its covariance blocks, uniformly under the allowed coarse translations. The accretive density bound supplies a common integrable majorant for its Cauchy derivatives, so each integral and each finite product defining the cluster weight converges. Consequently $\epsilon_R(T)\to0$ as the shortest period tends to infinity, for every fixed $R$.

Allocate the finite-square logarithm over its source segments. The clusters omitted by the extent truncation cost at most $C\tau^2e^{-\eta R}$ per root. There are at most $cR$ corner-near roots per side, and each has load at most $C\tau^2$. Thus
\begin{equation}\label{eq:marginal-buffered-density}
    \left|\frac{1}{4n}\log\frac{N_{C_n,T}}{D_T}
        -\frac{1}{4}\sum_{v\in \mathscr V_{\square}}\kappa_v\right|
        \leqslant C\tau^2\left(\frac{R}{n}+e^{-\eta R}\right)
            +\epsilon_R(T)\;.
\end{equation}
First let $n$ and the compatible buffered periods tend to infinity with $R$ fixed, then let $R\to\infty$. Every logarithm in this limit is the analytic branch continued along $t_a= t$, $0\leqslant t\leqslant 1$, from the value zero of the logarithm of the ratio at $t=0$. Restore the cutoff index: $\kappa_v^{\sigma}$ is the absolutely convergent line sum \eqref{eq:marginal-full-perimeter} with the prescription \eqref{eq:marginal-two-cutoffs}. For each $\sigma\in \{\mathrm{op},\mathrm{HS}\}$, the flat square perimeter coefficient exists and is
\begin{equation}\label{eq:marginal-square-mass}
    \boxed{\displaystyle
        \delta m^{\sigma}(g_0,\varepsilon)
            =-\frac{1}{4L\varepsilon}
                \sum_{v\in \mathscr V_{\square}}\kappa_v^{\sigma}(g_0)
            =-\lim_{n\to\infty}
                \frac{1}{4nL\varepsilon}
                    \log\frac{N_{C_n,T_n}^{\sigma}}{D_{T_n}^{\sigma}}}
        \;.
\end{equation}
No equality of the full coefficients in different orientations is presumed. Reversing a finite relative path conjugates its character, and translation by its coarse length identifies it with the oppositely oriented reference path. All other factors of the weights (the Gaussian density and its forest derivatives, $f_v$, $h_b$ and $\theta$) are real, the reindexing $a\mapsto-a-1$ carries the $-v$ line onto the reverse of the $v$ line with the stencils, protected boxes, intervals and incompatibilities in bijection, and $\kappa_v$ is an absolutely convergent sum, so $\kappa_{-v}=\overline{\kappa_v}$. Hence the square coefficient is real.

\section{The Corner Coefficient in the Same First Step}
\label{sec:marginal-corner}

Fix the lattice spacing, character, reconstruction and cutoff prescription $\sigma$ as in Section~\ref{sec:marginal-perimeter}. Choose the coupling below the common minimum of the straight and bent-path smallness thresholds of this section and the regional thresholds of Proposition~\ref{prop:regional-square-extraction}, after fixing all chart radii and majorant constants. All limits in this section vary reference contours and enclosing tori at this one RG step; neither the coupling nor the RG scale is iterated. The coefficients below use exactly $N_{C,T}^{\sigma}$ and $D_T^{\sigma}$ in \eqref{eq:marginal-literal-ratio}. Suppress $\sigma$ until the conclusion. Write $C_{\infty}=\mathsf H_{\infty}^{-1}$ for the free-field covariance on $\ell^2(E_{*,\infty};\mathfrak g)$, constructed in \eqref{eq:bulk-neumann-inverse}; it is distinct from the reconstructed fine-field covariance $\mathsf K_{\infty}=\mathsf C_{\infty} C_{\infty}\mathsf C_{\infty}^*$. The countable primitive index family $\mathcal V_{\infty}$ consists of labels $(0,\alpha)$, $(1,b)$ and $(2,c)$ for the functions $v_{\alpha},j_b,d_c$ in \eqref{eq:perimeter-literal-local-terms}. Here $b$ ranges over positive fine bonds, $c$ over positive coarse bonds, and $\alpha$ over positive plaquettes and the mean-transport pairs $(y,x)$ with $y\in L\mathbb Z^4$ and $x\in\mathcal B_y\setminus\{ y\}$. These are exactly the infinite-lattice translates of the finitely many local templates modulo the coarse grid.

\subsection{A Local Corner Subtraction}

Let $v,w\in \mathscr V_{\square}$ be perpendicular incoming and outgoing orientations. For $a\in \mathbb Z$, the bent reference path has fine vertices and oriented bonds
\begin{equation}\label{eq:corner-bent-path}
    x_a^{v,w}=
    \begin{cases}
        av,&a<0,\\
        aw,&a\geqslant 0,
    \end{cases}
    \qquad b_a^{v,w}=(x_a^{v,w},x_{a+1}^{v,w})\;.
\end{equation}
Its coarse segment $j\in \mathbb Z$ contains the fine traversal indices $jL,\ldots,(j+1)L-1$. Thus the bend is at a coarse site, between segments $-1$ and $0$.

For every finite set of selected traversals, use the closed relative source \eqref{eq:marginal-line-source} on the corresponding finite arc of this bent path. Keep the same local vertices, protected cutoffs and Gaussian principal marginals. Define its components and webs by \eqref{eq:marginal-component-functional}--\eqref{eq:marginal-web-activity}, with ordinary intervals in the ordered traversal indices, and its connected weight $\Phi_{v,w}(\Gamma)$ by \eqref{eq:marginal-cluster-weight}. The set $J_{\Gamma}\subset \mathbb Z$ again records the coarse segments containing source marks. Vacuum-only clusters are excluded. These definitions involve finite closed contractions and finite-dimensional integrals only, even on the infinite reference path.

The density, label-counting and KP estimates of Section~\ref{sec:marginal-perimeter} apply to this path with enlarged fixed constants. Indeed, it has no repeated traversal, and the path distance between two of its points is at most twice their maximum-norm spatial distance. Each ray is part of one of the two straight lines, so at most $2q_*$ primitive or source labels meet any bond; the support sizes $s_*$ and ranges $R_*$ are unchanged. Thus the original counting proof applies with $q_*'=2q_*$, and the interval/extent bound holds with $c_h'=2c_h$. The resulting rooted estimate is
\begin{equation}\label{eq:corner-bent-tail}
    \sum_{\Gamma\, | \,j\in J_{\Gamma}}
        |\Phi_{v,w}(\Gamma)| e^{\eta'\ell(\Gamma)}
        \leqslant C\tau^2\;.
\end{equation}
Take $\beta'=\mu/(2048(1+c_h'))$ and $\eta'=\eta/2$. The index-distance conversion \eqref{eq:marginal-hull-root} is unchanged, since the traversal indices still form $\mathbb Z$ and the sum of its geometric kernel is $1+2/(e^{\theta}-1)$. The carrier-root and interval-root loads therefore obey \eqref{eq:marginal-weighted-loads} with the bent path, size $a'(p)=| A_p|+\beta'H_p$, exponent $\eta'$ and a fixed enlarged constant $C_*'$ obtained from the same label and positive-tree sums with $q_*'$. Choose the coupling so that $C_*'\tau\leqslant 1/4$ and $C_*'\tau^2\leqslant \beta'/4$. The incompatibility sum is then at most $a'(p)/4$, and the pinned proof \eqref{eq:marginal-pinned-bound}--\eqref{eq:marginal-pinned-recursion} gives \eqref{eq:corner-bent-tail}. Both cutoff choices satisfy this argument. No open transport receives an activity norm: the relative word transforms by conjugation at its base when field and reference are transformed together, and its character, component-integrated closed contractions and scalar cluster weights are gauge invariant.

Define the rooted density and its straight reference by
\begin{equation}\label{eq:corner-rooted-density}
    a_{v,w}(j)=\sum_{\Gamma\, | \,j\in J_{\Gamma}}
        \frac{\Phi_{v,w}(\Gamma)}{| J_{\Gamma}|},
    \qquad
    k_{v,w}(j)=
    \begin{cases}
        \kappa_v,&j<0,\\
        \kappa_w,&j\geqslant 0.
    \end{cases}
    \;.
\end{equation}
The straight densities are exactly \eqref{eq:marginal-full-perimeter}, not newly fitted constants. The proposed local corner coefficient is
\begin{equation}\label{eq:corner-local-extraction}
    \boxed{\displaystyle
        \zeta_{v,w}=\sum_{j\in \mathbb Z}
            \left(a_{v,w}(j)-k_{v,w}(j)\right),
        \qquad z_{v,w}=\exp(\zeta_{v,w})}\;.
\end{equation}

Here is the cancellation and the quantitative payment for this subtraction. A cluster with all source marks at negative indices has exactly the weight of the same cluster on the incoming straight line: its ordered nonidentity factors, source stencils, ordinary intervals, local vacuum factors and Gaussian marginals agree. The same holds on the positive half-line with orientation $w$. Vacuum carriers may cross the bend; they are not cut or discarded in this comparison. Consequently only clusters with source marks on both sides of index zero survive the difference of the rooted sums. On the straight side of the subtraction these are the straight clusters straddling the same cut. A surviving cluster rooted in segment $j$ has extent at least $c| j|-C_0$, with $c= L/(2C_{\mathrm{net}})$ and $C_0= R_*/C_{\mathrm{net}}$, since its connected support/interval network reaches the other half-line, so has length at least $L| j|-2R_*$, while the bent-path network constant ($C_{\mathrm{net}}$ with $q_*'$, $c_h'$) is at most $2C_{\mathrm{net}}$. The rooted bound, with part of its extent weight retained, therefore gives
\begin{equation}\label{eq:corner-localisation-bound}
    | a_{v,w}(j)-k_{v,w}(j)|
        \leqslant C\tau^2e^{-c| j|},
    \qquad
    |\zeta_{v,w}|\leqslant C\tau^2,
    \qquad | z_{v,w}-1|\leqslant C\tau^2\;.
\end{equation}
This proves absolute convergence and localisation at the corner, rather than merely bounding a fixed square. As a direct check, for the trivial character every nonempty closed source difference is zero. The source cluster sums, straight densities and corner logarithms then vanish, giving $N_{C,T}= D_T$ and $z_{v,w}=1$, as required.

An explicit cluster formula makes the endpoint allocations visible. Let $J_{\Gamma}^-= J_{\Gamma}\cap\mathbb Z_{<0}$ and $J_{\Gamma}^+= J_{\Gamma}\cap\mathbb Z_{\geqslant 0}$. In each of the following three sums, the cluster has both $J_{\Gamma}^-\ne \varnothing$ and $J_{\Gamma}^+\ne \varnothing$, and its weight is that of the indicated bent or straight path. Then
\begin{equation}\label{eq:corner-mixed-cluster-formula}
    \begin{aligned}
        \zeta_{v,w}
            &=\sum_{\Gamma}^{v,w}\Phi_{v,w}(\Gamma)
                -\sum_{\Gamma}^{v}
                    \frac{| J_{\Gamma}^-|}{| J_{\Gamma}|}\Phi_v(\Gamma)\\
            &\quad-\sum_{\Gamma}^{w}
                    \frac{| J_{\Gamma}^+|}{| J_{\Gamma}|}\Phi_w(\Gamma)\;.
    \end{aligned}
\end{equation}
Indeed, summing the rooted allocations gives one for a mixed bent cluster and the displayed fractions for the two missing straight half-line contributions. Each sum is absolutely convergent by \eqref{eq:corner-bent-tail} and the distance-to-cut payment. In particular the corner factor includes the straight-edge endpoint deficits as well as interactions joining the two perpendicular rays.

\subsection{Identification with the Square Constant}

First make the fixed-square bulk limit explicit. For a fixed positive integer $n$, let $\widehat F_{\square}(n)$ be the infinite-lattice connected sum over clusters with source marks on $C_n$, using its fixed cyclic traversal order and base. Write $F_{\square,\infty}(n)$ for the bulk-limit logarithm whose existence will be proved below. Allocate the connected sum over the $4n$ coarse source segments by $| J_{\Gamma}|^{-1}$. Its absolute convergence follows from \eqref{eq:marginal-cluster-tail}. The same bound holds uniformly on all compatible tori, which satisfy $K_i\geqslant 3(n+2)$, so the square geometry and weighted KP payment used to prove that bound have their original constants, including at the corners.

Truncate both connected sums to $\ell(\Gamma)\leqslant R$. Once the shortest fine period exceeds a fixed multiple of $nL+R+R_*$, each such rooted cluster has a unique lift through its source root to the fixed neighbourhood of the lifted square. Its connected support/interval network has length at most a fixed multiple of $R$, so a network winding around a torus cannot occur in this truncation. The primitive functions and supports agree after lifting. The cyclic interval rule and its tie-breaking base depend only on the fixed ordered list of square traversals, so they agree as well, including intervals crossing the base. There are finitely many lifted shapes, labels and multiplicities: the root belongs to a finite contour, the extent bounds all spatial displacements and the number of component and cluster occurrences, and each site has finitely many primitive labels. Covariance entries converge by the periodisation of $\mathsf H_{\infty}^{-1}$ (Appendix~\ref{app:flat-covariance}: the periodisation of the free inverse solves the torus inverse equation, the paragraph preceding \eqref{eq:bulk-periodisation}, and the image bound \eqref{eq:bulk-image-error}, displayed there for $\mathsf K$, holds for $\mathsf H_{\infty}^{-1}$ by the same argument from the decay of its entries established before \eqref{eq:bulk-kernel-decay}). For each shape, the differentiated Gaussian density majorant \eqref{eq:marginal-density-majorant} is integrable and uniform in the torus. Dominated convergence therefore gives convergence of every component integral, web activity and finite product in its cluster weight. Let $\mathscr G_{n,R}$ be this finite set of lifted ordered clusters with at least one source mark. Define
\begin{equation}\label{eq:corner-finite-shape-difference}
    \epsilon_{n,R}(T)=\sum_{\Gamma\in\mathscr G_{n,R}}
        |\Phi_{\square}^T(\Gamma)-\Phi_{\square}^{\infty}(\Gamma)|\;.
\end{equation}
For the domination just used, every fixed shape has $r_e^{-1}\leqslant 8d_bd_cA/\lambda$ in \eqref{eq:marginal-edge-radius}, uniformly in $T$. Thus the Cauchy factors multiplying \eqref{eq:marginal-density-majorant} are uniformly bounded for that shape. The differentiated density is analytic, and hence continuous, in its covariance entries on a common accretive neighbourhood. This gives pointwise convergence as well as the stated integrable majorant, including when a covariance edge tends to zero. It follows that $\epsilon_{n,R}(T)\to 0$ at fixed $n,R$.

The omitted clusters, including any winding clusters, have total allocated weight at most $Cn\tau^2e^{-\eta R}$ in each sum. Consequently
\begin{equation}\label{eq:corner-fixed-square-limit-bound}
    \left|\log\frac{N_{C_n,T}}{D_T}
        -\widehat F_{\square}(n)\right|
        \leqslant Cn\tau^2e^{-\eta R}+\epsilon_{n,R}(T)\;.
\end{equation}
Here the finite-torus logarithm is the connected expansion of the ratio, on its branch continued from zero source; the vacuum subtraction and original denominator are unchanged. First taking $T\to\infty$ at fixed $n,R$ and then $R\to\infty$ proves
\begin{equation}\label{eq:corner-fixed-square-limit}
    F_{\square,\infty}(n)
        =\widehat F_{\square}(n)
        =\lim_{T\to\infty}\log\frac{N_{C_n,T}}{D_T}\;.
\end{equation}
Thus the infinite connected sum and the bulk-limit logarithm used below are the same object. This argument needs no convergence of an infinite-volume vacuum partition function.

List the four incoming/outgoing pairs of the counterclockwise square as
\begin{equation}\label{eq:corner-square-pairs}
    \mathscr C_{\square}=
        \{(e_1,e_2),(e_2,-e_1),(-e_1,-e_2),(-e_2,e_1)\}\;.
\end{equation}
Use the branch of $F_{\square,\infty}(n)$ continued from zero source. The same closed-cluster allocation now gives positive constants $C,c$, independent of $n$, such that
\begin{equation}\label{eq:corner-square-extraction-bound}
    \left| F_{\square,\infty}(n)
        -n\sum_{v\in \mathscr V_{\square}}\kappa_v
        -\sum_{(v,w)\in \mathscr C_{\square}}\zeta_{v,w}\right|
        \leqslant C\tau^2(n+1)e^{-cn}\;.
\end{equation}
To verify this estimate, allocate every square cluster over its source segments by $| J_{\Gamma}|^{-1}$. Divide each side into the segments nearest its initial corner and those nearest its terminal corner, breaking a midpoint tie consistently. A corner then receives the last $\lfloor n/2\rfloor$ incoming segments and the first $n-\lfloor n/2\rfloor$ outgoing segments. Translate that corner to the coarse origin. At the base corner $(-e_2,e_1)$, first use trace cyclicity to put the incoming factors before the outgoing factors, exactly as in the cyclic reassembly after \eqref{eq:marginal-exact-gases}. For a root in its assigned half-sides, every cluster of extent less than $c_1n-C_1$ is identical to the corresponding bent-path cluster: all marked traversals and their stencils stay away from the other three corners, cyclic intervals become the ordinary local intervals, and the same global Gaussian principal marginals and vacuum labels occur. Any failure of this identification costs extent at least $c_1n-C_1$ and hence at most $C\tau^2e^{-c_2n}$ per root by \eqref{eq:marginal-cluster-tail} and \eqref{eq:corner-bent-tail}. There are $4n$ roots. Extending the four truncated rooted differences to all $j\in\mathbb Z$ costs another $C\tau^2e^{-c_3n}$ by \eqref{eq:corner-localisation-bound}. This proves \eqref{eq:corner-square-extraction-bound}, after enlarging constants for the finitely many short squares. Here $c_1= L/(4C_{\mathrm{net}})$ and $C_1=2R_*/C_{\mathrm{net}}$ (a cluster rooted in its assigned half-sides and reaching within $R_*$ of another corner has network length at least $nL/2-2R_*$, and the bent-path network constant is at most $2C_{\mathrm{net}}$), $c_2=\eta'c_1$ and $c_3= c/2$ with the $c$ of \eqref{eq:corner-localisation-bound} (the extension runs over $| j|\geqslant n/2$), the factors $e^{\eta'C_1}$ and the like being absorbed in $C$.

Restore the cutoff index and define the square corner matching factor by
\begin{equation}\label{eq:corner-square-factor}
    \boxed{\displaystyle
        z_{\mathrm c}^{\sigma}\!\left(g_0,\frac{\pi}{2}\right)
            =\exp\left(\frac{1}{4}
                \sum_{(v,w)\in \mathscr C_{\square}}
                    \zeta_{v,w}^{\sigma}(g_0)\right)}\;.
\end{equation}
This definition fixes the fourth-root branch by the exponential of the convergent corner logarithm. It is the geometric mean of the four oriented corner factors. Equality of those four factors is not needed or assumed; if the reference construction preserves the corresponding rotations, they agree and \eqref{eq:corner-square-factor} is their common value. For a real-valued character all the real Gaussian integrals, cluster weights and corner logarithms above are real, so these factors are positive. No reality symmetry for a general complex character is imposed.

For each fixed $n$, the inner limit in the following expression is the limit over compatible tori with their shortest period tending to infinity. Then
\begin{equation}\label{eq:corner-ratio-limit}
    \boxed{\displaystyle
        \log z_{\mathrm c}^{\sigma}\!\left(g_0,\frac{\pi}{2}\right)
            =\frac{1}{4}\lim_{n\to\infty}
                \left(\lim_{T\to\infty}
                    \log\frac{N_{C_n,T}^{\sigma}}{D_T^{\sigma}}
                    +4nL\varepsilon\delta m^{\sigma}\right)}\;.
\end{equation}
The denominator is the original vacuum integral throughout the inner limit, whose existence is \eqref{eq:corner-fixed-square-limit}. The outer limit follows from \eqref{eq:corner-square-extraction-bound} and \eqref{eq:marginal-square-mass}. The order of limits is part of this reference prescription; a constant-term limit on every diagonally growing finite torus is not asserted here.

Equivalently, write $\mathcal W_{\infty}^{\circ,\sigma}(C_n)=\lim_{T\to\infty}N_{C_n,T}^{\sigma}/D_T^{\sigma}$. There is an error $E_n^{\sigma}$ such that
\begin{equation}\label{eq:corner-reference-matching}
    \begin{aligned}
        \mathcal W_{\infty}^{\circ,\sigma}(C_n)
            &= e^{-4nL\varepsilon\delta m^{\sigma}}
                \left(z_{\mathrm c}^{\sigma}
                    \!\left(g_0,\frac{\pi}{2}\right)\right)^4
                e^{E_n^{\sigma}}\;,\\
        | E_n^{\sigma}|&\leqslant C\tau^2(n+1)e^{-cn}\;.
    \end{aligned}
\end{equation}
The corner coefficient uses the same flat integral and cutoff prescription as the perimeter coefficient. Its identification in a conditional regular region, with boundary and periodic errors, is Proposition~\ref{prop:regional-square-extraction}.


\section{Identification in a Regular Integration Region}
\label{sec:regional-identification}

The coefficient in \eqref{eq:marginal-full-perimeter} was computed with the translation-invariant flat covariance. A conditional integration in the first RG step instead uses the inverse of a restricted Hessian. We compare the two connected expansions, retaining the action terms that cross the integration boundary. This provides a quantitative identification of their straight-side coefficients.

\subsection{The Wilson source before the flat evaluation}

Let $\mathcal I_C$ be the ordered set of traversals of $C$. For $R\subseteq\mathcal I_C$, the mixed field $U^R$ uses the integration link $U(b)= U'(b)U_1(b)$ on a traversal in $R$ and the background link $U_1(b)$ on every other traversal. Here $U_1= U_1(\Omega_1,V)$ is the reconstruction before the boundary split (with the exterior determining field retained). Reversed traversals use inverse links. Define
\begin{equation}\label{eq:regional-wilson-polynomial}
    \begin{aligned}
        \mathcal G_S(U,U_1)
            &=\sum_{R\subseteq S}(-1)^{| S|-| R|}
                \chi(\Hol_{U^R}C)\;,\\
        \mathcal W(t;U,U_1)
            &=\sum_{S\subseteq\mathcal I_C}
                \left(\prod_{a\in S}t_a\right)\mathcal G_S(U,U_1)\;.
    \end{aligned}
\end{equation}
Here $t=(t_a)_{a\in\mathcal I_C}$ is a family of scalar parameters. Boolean inversion gives
\begin{equation}
    \mathcal W(1;U,U_1)= W_{\chi}(C;U),
    \qquad \mathcal W(0;U,U_1)= W_{\chi}(C;U_1)\;.
\end{equation}
Every term is a closed contraction, invariant under simultaneous gauge transformations of $U$ and $U_1$.

The Wilson insertion starts in \eqref{eq:bare-blocking-operator}. Its gauge invariance preserves it under the Faddeev--Popov insertion in \cite[(1.5)--(1.6)]{BalabanCRE}. It is then pulled back under $U= U'U_1$ and the constraint linearisation leading to (1.14). At that stage we introduce \eqref{eq:regional-wilson-polynomial}, whose value at $t=1$ is the pulled-back Wilson insertion. All subsequent changes of variables are applied to every fluctuating link in this polynomial, including the boundary change (1.22).

Fix the domain labels, the coarse field, the exterior fine field and the exterior free fluctuation field in this construction. Denote this collection by $\xi$. Define $K_{\xi}(t)$ to be the conditional fluctuation integral in \cite[(1.26)]{BalabanCRE}, with the pulled-back polynomial inserted and the change-of-variables density from (1.22)--(1.24) retained once. The same integral with insertion $1$ is $K_{\xi,\mathrm{vac}}$. These definitions precede the separation of the effective interaction into its regular and boundary terms. The exterior integrations and their prefactors are those of (1.15), (1.25) and (1.29).

The flat projection of this conditional kernel sets its coarse and exterior determining fields to the identity and its exterior free fluctuation field to zero. The domain labels are held fixed. This evaluation is applied to $K_{\xi}$, rather than to the exterior indicators or history weights. In particular, a large-field indicator is not evaluated at zero. At the projection point $\mathcal W(0)=1$, so the source logarithm is continued from $K_{\xi}(0)= K_{\xi,\mathrm{vac}}$. No logarithm of $W_{\chi}(C;U_1)$ at general background is required.

\subsection{Regional integrals and their connected densities}

Use fine-lattice units. Write $H=\mathsf H_{\infty}$ for \eqref{eq:bulk-operator-definition}, and $C_{\infty}= H^{-1}$. For a finite set $D\subset E_{*,\infty}$, let $i_D$ extend a free field by zero and put
\begin{equation}\label{eq:regional-precision}
    H_D= i_D^*Hi_D,
    \qquad C_D= H_D^{-1}\;.
\end{equation}
The inverse exists by \eqref{eq:bulk-uniform-gap}. Let $\mathcal V_D$ consist of the primitive labels from \eqref{eq:perimeter-literal-local-terms} whose supports meet $D$. Retain their original labels and set $v_D(B)= v(i_DB)$. Every primitive missing $D$ vanishes at this field. The cutoff and interaction are
\begin{equation}\label{eq:regional-weight}
    \Theta_D^{\sigma}(B)=\prod_{b\in D}\mathbf1_{\{| B_b|_{\sigma}< r\}},
    \qquad P_D(B)=\sum_{v\in\mathcal V_D}v(i_DB)\;.
\end{equation}
The threshold is the same $r= r_0(g_0)$ as in \eqref{eq:marginal-two-cutoffs}.

Let $I$ be an interval of $m$ coarse segments on a coarse translate of the oriented line $v\in\mathscr V_{\square}$. Let $P_U(I)$ denote the ordered transport along its fine traversals. Its closed relative source uses the inverse reference transport, as in \eqref{eq:marginal-line-source}. The regional numerator and denominator are
\begin{equation}\label{eq:regional-integrals}
    \begin{aligned}
        N_{D,I}^{\sigma}
            &=\int\Theta_D^{\sigma}(B)e^{P_D(B)}
                \chi\bigl(P_{U_{g_0}(i_DB)}(I)P_{U_0}(I)^{-1}\bigr)
                    \,\dd\mu_{C_D}(B)\;,\\
        D_D^{\sigma}
            &=\int\Theta_D^{\sigma}(B)e^{P_D(B)}
                    \,\dd\mu_{C_D}(B)\;.
    \end{aligned}
\end{equation}
Here $U_0=\id$ in the chosen frame; retaining its reverse transport specifies the gauge-invariant closed contraction.

Apply the protected-label forest expansion and the web construction of Section~\ref{sec:marginal-perimeter}, replacing every free-input support $\Lambda$ by $\Lambda\cap D$. A source mark with empty restricted stencil is zero. The other labels keep their original identities. Denote the connected weights by $\Phi_{D,I}(\Gamma)$ and their marked coarse-segment set by $J_{\Gamma}$. The density assigned to segment $j$ is
\begin{equation}\label{eq:regional-density}
    \kappa_{D,I}(j)=
        \sum_{\Gamma\, | \,j\in J_{\Gamma}}
            \frac{\Phi_{D,I}(\Gamma)}{| J_{\Gamma}|}\;.
\end{equation}
We suppress $\sigma$ in these weights and densities. It is fixed throughout each comparison.

For positive bonds $b=(x,i)$ and $c=(y,k)$, use $d(b,c)=\| x-y\|_{\infty}+\mathbf1_{\{ i\ne k\}}$. Let $R_*$ bound the label diameters and source-to-stencil distances, and let $R_H$ be the range of $H$ in this metric. With the interval constant $c_h$ of \eqref{eq:marginal-hull-extent}, set
\begin{equation}
    R_0=4(1+R_*+R_H+L),
    \qquad c_{\mathrm{reg}}=16(1+R_*+c_h+L)\;.
\end{equation}
Number the coarse segments of $I$ by $0,\ldots,m-1$. Let $K_j$ be the union of their fine-traversal input stencils on segment $j$, and define
\begin{equation}\label{eq:regional-distance}
    d_j=\max\bigl(0,
        \min\bigl(d(K_j,D^c),L\min(j,m-1-j)\bigr)-R_0\bigr)\;.
\end{equation}

\begin{proposition}\label{prop:regional-coefficient}
For either fixed cutoff prescription $\sigma$, and for sufficiently small $g_0$ after the constants in the comparison have been chosen, there are constants $C,c>0$, independent of $D$, $I$ and $j$, such that
\begin{equation}\label{eq:regional-coefficient-bound}
    |\kappa_{D,I}(j)-\kappa_v|
        \leqslant C\tau^2e^{-cd_j}\;.
\end{equation}
Here $\tau$ is defined in \eqref{eq:marginal-fixed-forest-bound} and $\kappa_v$ in \eqref{eq:marginal-full-perimeter}. If every $K_j$ is at distance at least $b$ from $D^c$, then
\begin{equation}\label{eq:regional-summed-bound}
    \left|\log\frac{N_{D,I}^{\sigma}}{D_D^{\sigma}}-m\kappa_v^{\sigma}\right|
        \leqslant C\tau^2(1+me^{-cb})\;.
\end{equation}
The logarithm is continued from zero source.
\end{proposition}

\begin{proof}
Principal restriction preserves $\gamma_L I\leqslant H_D\leqslant h_L I$. The finite-range Neumann series in \eqref{eq:bulk-neumann-inverse} therefore gives exponential decay for $C_D$, uniformly in the shape of $D$. Projecting $HC_{\infty}= I$ onto the $D$ rows gives
\begin{equation}\label{eq:regional-resolvent}
    C_D-(C_{\infty})_{D,D}
        = C_DH_{D,D^c}(C_{\infty})_{D^c,D}\;.
\end{equation}
Let $\mu>0$ be a common inverse-decay exponent in the fine-bond metric. A term on the right joins $u\in D$ and $w\in D^c$ with $d(u,w)\leqslant R_H$. For $s= d(b,u)+d(w,c)$,
\begin{equation}
    d(b,D^c)+d(c,D^c)+d(b,c)\leqslant 2s+2R_H\;.
\end{equation}
Half of $e^{-\mu s}$ pays this distance with exponent $a=\mu/4$; the remaining half is summable over $u$ and its bounded number of neighbours $w$. Polynomial volume growth gives
\begin{equation}\label{eq:regional-covariance-difference}
    \| C_D(b,c)-C_{\infty}(b,c)\|
        \leqslant C e^{-a(d(b,D^c)+d(c,D^c)+d(b,c))}\;.
\end{equation}
There is no factor proportional to the area of the boundary.

On $D$, interpolate between $(C_{\infty})_{D,D}$ and $C_D$, and denote their difference by $\mathcal D$. Every interpolated covariance and each of its principal marginals has spectrum in $[h_L^{-1},\gamma_L^{-1}]$, uniform exponential decay and scalar colour blocks. Thus the protected-source bounds and the Schur cancellation used in \eqref{eq:marginal-fixed-forest-bound} hold throughout the interpolation.

Consider a forest term on carrier $A$. If $M_t= C_t\circ w^F$ is its interpolated covariance and $E_e(C_t)$ the off-diagonal direction associated with edge $e$, its density derivative has the product-rule decomposition
\begin{equation}\label{eq:regional-density-derivative}
    \begin{aligned}
        \partial_t D^{| F|}p_{M_t}[(E_e(C_t))_{e\in F}]
            &= D^{| F|+1}p_{M_t}
                [\mathcal D_A\circ w^F,(E_e(C_t))_{e\in F}]\\
            &\quad+\sum_{e\in F}D^{| F|}p_{M_t}
                [E_e(\mathcal D),(E_f(C_t))_{f\in F\setminus\{ e\}}]\;.
    \end{aligned}
\end{equation}
Each tuple lists the directions of the symmetric multilinear derivative of the density with respect to its covariance.

The two kinds of terms use separate Cauchy estimates. In the first, add the direction $s(\mathcal D_A\circ w^F)$ to the covariance centre while leaving every forest direction fixed. The positive Schur mask contracts operator norm. With $\lambda= h_L^{-1}$, the radius $\lambda/(32\|\mathcal D_A\|)$ uses at most $\lambda/32$ of the accretivity margin; a zero direction gives zero derivative. Reduce the forest radii by a fixed factor to keep the total below $\lambda/8$. In the second kind of term, replace one forest direction by $E_e(\mathcal D)$ and use \eqref{eq:regional-covariance-difference}. Choose its radius from the exponential entry majorant, without dividing by an entry of $C_t$. The normalised density, including its determinant, is bounded by \eqref{eq:marginal-density-majorant}.

Since $\|\mathcal D_A\|\leqslant C| A| e^{-a d(A,D^c)}$, the integrated differentiated term is bounded by the positive forest majorant, with enlarged constants and reduced edge decay, times
\begin{equation}
    C| A|^2e^{-a d(A,D^c)}\;.
\end{equation}
The exponent of $\tau$ is unchanged. In particular the differentiated functional is still equivariant under constant conjugation, so a component with one source mark retains its quadratic character gain.

For each component $\gamma$, let $u_{\gamma}$ be its positive fixed-forest bound with the uniform exponential entry majorants. This definition applies to vacuum and source components. For a web $W$, put $u_W=\prod_{\gamma\in W}u_{\gamma}$; summing these numbers over webs gives the web load. The positive majorants themselves satisfy the polymer loads. Restricting a support to $D$ preserves its size bound $s_*$, overlap bound $q_*$ and range $R_*$. On $n$ bonds, label choices cost at most $2^{q_*n}$, the forest-to-graph multiplicity costs at most $2^{q_*(s_*-1)n/2}$, unused graph edges at most $e^{V_1n/2}$, and the rooted spatial tree sum at most $n^{n-2}V_1^{n-1}/(n-1)!$, where $V_1$ is the finite row sum of the reduced-decay graph weights. Thus a constant $K_2\geqslant 1$ bounds the vacuum and source root loads by
\begin{equation}
    \sum_{n\geqslant 1}K_2^n\tau^{\lceil n/s_*\rceil},
    \qquad
    \sum_{n\geqslant 1}K_2^n\tau^{\max(2,\lceil n/s_*\rceil)}\;.
\end{equation}
Grouping $n$ into intervals of $s_*$ gives bounds $C\tau/(1-K_2^{s_*}\tau)$ and $C\tau^2/(1-K_2^{s_*}\tau)$, respectively. These sums bound the positive forest majorants before any signed activity is used. The interval-tree recursion then bounds the web loads.

Let $\mu_1>0$ be an exponent remaining after the degree payment, no greater than $\min(\mu,a)/8$. Choose
\begin{equation}
    \alpha_*=\frac{\mu_1}{1024(1+c_h)},
    \qquad \eta_1=\frac{\mu_1}{1024}\;.
\end{equation}
The same sums allow the reserve $e^{4| A|+8\alpha_*H+8\eta_1\ell}$. Let $C_{\mathrm{load}}$ bound their weighted spatial and interval root loads after division by $\tau$ and $\tau^2$, respectively. Choose $g_0$ so that $C_{\mathrm{load}}\tau\leqslant 1/4$ and $C_{\mathrm{load}}\tau^2\leqslant \alpha_*/4$. The incompatibility sum is then bounded by one quarter of $| A_p|+\alpha_*H_p$, so the pinned tree bound applies to the positive majorants $u_p$.

For a cluster $\Gamma$, let $U(\Gamma)$ be its absolute Ursell coefficient, including its factorial, multiplied by the product of these majorants. The product rule, counting repeated activity occurrences, yields
\begin{equation}
    |\partial_t\Phi_t(\Gamma)|
        \leqslant C\ell(\Gamma)^2
            e^{-a d(A_{\Gamma},D^c)}U(\Gamma),
    \qquad
    \sum_{\Gamma\, | \,j\in J_{\Gamma}}
        U(\Gamma)e^{\eta_1\ell(\Gamma)}\leqslant C\tau^2\;.
\end{equation}
The union of carriers is $A_{\Gamma}$. The inequality follows from $\sum_p\ell(p)^2\leqslant \ell(\Gamma)^2$ and the positive-majorant KP bound. Half the extent weight absorbs the polynomial factor.

For completeness, connect a component's carrier using its forest and at most $| A| -1$ stencil edges, each of length at most $R_*$. Attach its contour interval at one selected mark by a path of length at most $R_*$. The interval costs at most $c_h\ell$. Carrier or interval incompatibility joins these networks for a connected cluster. Their total length is at most $(1+2R_*+c_h)\ell(\Gamma)$; adding the root segment costs at most $L$. The chosen $c_{\mathrm{reg}}$ bounds the result.

For $d_j\leqslant 2(R_0+R_*+c_{\mathrm{reg}})$, the rooted bounds prove the comparison after enlarging $C$. For larger $d_j$, clusters with $\ell(\Gamma)\geqslant d_j/(2c_{\mathrm{reg}})$ are bounded by the extent tail in each gas. Every shorter cluster rooted at $j$ lies away from both the integration boundary and the ends of $I$. Its original stencils lie inside $D$, so its labels, closed sources, local functions and incompatibilities coincide in the two expansions. Only its covariance changes, and every carrier is at distance at least $d_j/2-R_0$ from $D^c$. The fixed loss is absorbed in the comparison constant. Integrating the differentiated bound gives \eqref{eq:regional-coefficient-bound}. Division by $| J_{\Gamma}|$ cannot enlarge it. Finally, allocation and the endpoint geometric series give
\begin{equation}
    \log\frac{N_{D,I}^{\sigma}}{D_D^{\sigma}}
        =\sum_{j=0}^{m-1}\kappa_{D,I}(j),
    \qquad
    \sum_{j=0}^{m-1}e^{-cd_j}\leqslant C(1+me^{-cb})\;,
\end{equation}
which proves \eqref{eq:regional-summed-bound}.
\end{proof}

\subsection{Periodic regions and the CRE kernel}

Let $E_{*,T}$ be the periodic free carrier, $H_T$ the Hessian in \eqref{eq:flat-covariance}, and $i_{D,T}$ zero extension from $D\subseteq E_{*,T}$. The periodic primitives, written $v_T$ collectively (the terms $v_{\alpha}$, $j_b$ and $d_c$ of \eqref{eq:perimeter-literal-local-terms} on $T$, all three kinds), and the reconstructed links $U_{g_0,T}$ of \eqref{eq:flat-phase} on $T$ are those of Appendix~\ref{sec:perimeter-attempt}. Define
\begin{equation}\label{eq:regional-periodic-weight}
    H_{D,T}= i_{D,T}^*H_Ti_{D,T},
    \qquad C_{D,T}= H_{D,T}^{-1},
    \qquad P_{D,T}(B)=
        \sum_{v_T\, | \,\operatorname{supp}v_T\cap D\ne\varnothing}
            v_T(i_{D,T}B)\;.
\end{equation}
The primitives of Section~\ref{sec:source-phase-primitives} at the reference background $U=\id$ (written $p_{\xi,\alpha}$, $p^{\mathrm H}_{\xi,b}$, $p^{\mathrm J}_{\xi,c}$ as there, with the reference background in place of the branch data) and the local terms of \eqref{eq:perimeter-literal-local-terms} are the same functions, label by label.

\begin{lemma}[primitives at the identity background]\label{lem:identity-primitives}
Let the background be $U=\id$ on the torus $T$ with region $\Omega= T$ (Section~\ref{sec:regular-hessian-stencils}), so that every row $c\in T^{[1]}$ and every centre $y\in T^{(1)}$ occurs, and let $B\in\mathfrak g^{E_*}$ be real with $g_0C_{\mathsf C}\max_b| B(b)|\leqslant \vartheta_*/2$ and $g_0\max_b| B(b)|< a_L$, with $C_{\mathsf C}$, $\vartheta_*$ of Lemma~\ref{lem:c2-uniform-chart} and $a_L$ the chart radius of Appendix~\ref{sec:perimeter-attempt}. Here $|\bigcdot |$ is the cutoff norm (so $h\leqslant |\bigcdot |$, the conventions of Section~\ref{sec:sourced-effective-interaction}), $\|\bigcdot \|$ the Hilbert norm of Lemma~\ref{lem:c2-uniform-chart} and $\| A\|_{\infty}=\max_bh(A_b)$ the phase norm, in which $a_{\mathrm{ph}}$ and $a_L$ are measured; the row solutions, the phase $q_{\id}$ and the primitives $p_{\xi,\alpha}$, $p^{\mathrm H}_{\xi,b}$, $p^{\mathrm J}_{\xi,c}$ at the background $\id$ are those given by the formulas of Sections~\ref{sec:regular-hessian-stencils} and~\ref{sec:source-phase-primitives} with $\Omega= T$ and $U=\id$, the row solution of every row $c$ being the unique small solution $\varphi_c$ of Lemma~\ref{lem:c2-uniform-chart}(ii), which applies to every row of the identity background (proof of (i)). Then:
\begin{enumerate}
\item $\mathcal Q_{\id}=\mathcal Q$, $Q_{\id}= Q$, $\mathcal N_{\id}=\mathcal N$ \eqref{eq:constraint-primitive}, $\mathsf C_{\id}=\mathsf C$ and $h_{\id}(c)= L^3\id$ for every row \eqref{eq:constraint-extension}; every row solution $\widetilde D_{\id,c}(g_0\mathsf CB)$ is defined, is the unique small solution of Lemma~\ref{lem:c2-uniform-chart}(ii), and equals the row $c$ of $\widetilde D(g_0\mathsf CB)$ of \eqref{eq:flat-phase}, the two functions agreeing on the convex open set $\{ A\in\mathfrak g^{\mathcal E_c}\text{ real}:\| A\|<\vartheta_*,\ \| A\|_{\infty}< a_{\mathrm{ph}}\}$ of row arguments; hence $q_{\id}(B)= q_{g_0}(B)$.
\item $p_{\xi,p}(B)= v_p(B)$ for every plaquette $p$, and $p_{\xi,y}(B)=\sum_{x\in\mathcal B_y\setminus\{ y\}}v_{(y,x)}(B)$ for every centre $y$.
\item $p^{\mathrm H}_{\xi,b}(B)= j_b(B)$ for every bond $b$.
\item $p^{\mathrm J}_{\xi,c}(B)= d_c(B)$ for every row $c$.
\end{enumerate}
Consequently $\sum_{\alpha} p_{\xi,\alpha}(B)+\sum_bp^{\mathrm H}_{\xi,b}(B)+\sum_cp^{\mathrm J}_{\xi,c}(B)=\mathcal P_{g_0}(B)$ \eqref{eq:perimeter-literal-local-sum}; and for every bond set $D\subseteq E_*$ and every real $B\in\mathfrak g^D$ with $g_0C_{\mathsf C}\max_b| B(b)|\leqslant \vartheta_*/2$ and $g_0\max_b| B(b)|< a_L$, the sum of the primitives whose free-input sets meet $D$, evaluated at $i_{D,T}B$, equals $P_{D,T}(B)$.
\end{lemma}

\begin{proof}
(i) At $U=\id$ every transport is $\id$, so $\beta_c(\id)=0$ and $\mathsf M(\id)(c)= P_c(\id)=\id$ in \eqref{eq:nonlinear-average}; hence $\mathcal Q_{\id,c}(q)=\frac1i\log\mathsf M(e^{iq})(c)$ is the row $c$ of $\mathcal Q(q)$ in \eqref{eq:constraint-primitive}, and $Q_{\id}= Q$, $\mathcal N_{\id}=\mathcal N$ follow from the definitions \eqref{eq:native-source-constraint-normalisation}. By Lemma~\ref{lem:two-block}, $Q_0= Q\iota_0$ is diagonal with entries $\alpha_c(\id)= L^{-4}\id$, so $h_{\id}(c)= L^{-1}\alpha_c(\id)^{-1}= L^3\id$ is the row $c$ of $h= L^{-1}\iota_0Q_0^{-1}$; and since row $c$ of $Q$ is supported on $\mathcal E_c$ (Lemma~\ref{lem:two-block}), $(\mathsf C_{\id}B)(b_0(c))=-\alpha_c(\id)^{-1}\sum_{b\in\mathcal E_c\cap E_*}Q(c,b)B(b)=-(Q_0^{-1}Q\iota_*B)(c)=(\mathsf CB)(b_0(c))$, while both maps are the identity on free bonds; so $\mathsf C_{\id}=\mathsf C$. The row equation \eqref{eq:native-source-local-implicit} at $U=\id$ is the row $c$ of the equation defining $\widetilde D$ after \eqref{eq:flat-phase}. Lemma~\ref{lem:c2-uniform-chart} applies to the identity background on every row (its gauge hypothesis holds with $u=\id$, $\phi=0$), and by its part (i) the hypothesis on $B$ gives $\| g_0\mathsf CB|_{\mathcal E_c}\|\leqslant g_0C_{\mathsf C}\max_b| B(b)|\leqslant \vartheta_*/2$ for every row, so its part (ii) provides the unique small solution $\varphi_c(g_0\mathsf CB|_{\mathcal E_c})$, which is the row solution of Section~\ref{sec:source-phase-primitives} for the reference background, analytic on the real ball $\| A|_{\mathcal E_c}\|<\vartheta_*$ and agreeing with the germ of Lemma~\ref{lem:row-locality} at $A=0$. The row $c$ of $\widetilde D$ of Appendix~\ref{sec:perimeter-attempt} is real analytic on $\{\| A\|_{\infty}< a_{\mathrm{ph}}\}$, a function of $A|_{\mathcal S_c}$, hence of $A|_{\mathcal E_c}$ ($\mathcal S_c$, the positive fine bonds with both endpoints in $\mathcal B_{c_-}\cup\mathcal B_{c_+}$, is contained in $\mathcal E_c$, Section~\ref{sec:regular-hessian-stencils}), and agrees with the same germ at $A=0$ (the paragraph ``Locality of the first-step interaction'' there). The set of real $A|_{\mathcal E_c}$ with $\| A|_{\mathcal E_c}\|<\vartheta_*$ and $\| A\|_{\infty}< a_{\mathrm{ph}}$ is convex, hence connected, open, and contains $g_0\mathsf CB|_{\mathcal E_c}$ by the hypothesis on $B$ ($\| g_0\mathsf CB\|_{\infty}\leqslant \|\mathsf C\|_{\infty\to\infty}g_0\max_b| B(b)|< a_{\mathrm{ph}}$ by the definition of $a_L$ and $h\leqslant |\bigcdot |$); the identity theorem identifies the two analytic functions on it. With $h_{\id}= h$ and $\mathsf C_{\id}=\mathsf C$, \eqref{eq:native-source-local-phase} becomes \eqref{eq:flat-phase}: $q_{\id}(B)= g_0\mathsf CB-h\widetilde D(g_0\mathsf CB)= q_{g_0}(B)$.

(ii) Write $\psi(A)= A-h\widetilde D(A)$ as in Appendix~\ref{sec:perimeter-attempt}, so that $q_{g_0}(B)=\psi(g_0\mathsf CB)$, $\psi(0)=0$ and $D\psi(0)=\id$ because $D\widetilde D(0)=0$ (Lemma~\ref{lem:c2-uniform-chart}(ii)). For a plaquette $p$, $s_{p,\id}(q)=1-\operatorname{Re}\tr_N(e^{iq}(\partial p))= f_p(q)$ is the $p$-summand of $S(e^{iq})$, the function $f_p$ of Appendix~\ref{sec:perimeter-attempt}, and $s_{p,\id}(0)=0$ since $\tr_N(\id)=1$; so $E_{\xi,p}(B)=-g_0^{-2}f_p(\psi(g_0\mathsf CB))$ by (i). Each $f_p\geqslant 0$ vanishes at $q=0$, so $Df_p(0)=0$. The chain rule then gives $DE_{\xi,p}(0)[B]=-g_0^{-2}Df_p(0)[D\psi(0)g_0\mathsf CB]=0$ and $D^2E_{\xi,p}(0)[B,B]=-g_0^{-2}\bigl(D^2f_p(0)[D\psi(0)g_0\mathsf CB,D\psi(0)g_0\mathsf CB]+Df_p(0)[D^2\psi(0)[g_0\mathsf CB,g_0\mathsf CB]]\bigr)=-D^2f_p(0)[\mathsf CB,\mathsf CB]$, so by \eqref{eq:native-source-energy-residual}
\begin{equation}
p_{\xi,p}(B)=-\frac1{g_0^2}f_p(\psi(g_0\mathsf CB))+\frac12D^2f_p(0)[\mathsf CB,\mathsf CB]= v_p(B)\;.
\end{equation}
For a centre $y$: at $U=\id$ every path transport is $\id$, so the local group mean $\mathcal U_{\id}(y,x)$, the unique solution near the diagonal value of the equation after \eqref{eq:constraint-extension}, is $\id$; at $q= q_{g_0}(B)$ the mean $\mathcal U_{e^{iq}}(y,x)$ exists by Lemma~\ref{lem:c2-uniform-chart}(iv) at the identity background, since for a row $c$ with endpoint block $\mathcal B_y$ one has $\mathcal E_y\subseteq\mathcal E_c$ and $\| q_{g_0}(B)|_{\mathcal E_c}\|\leqslant \| A\|+2H_*K_*\| A\|^2\leqslant \vartheta_*$ for $A= g_0\mathsf CB|_{\mathcal E_c}$, $\| A\|\leqslant \vartheta_*/2$, by (ii) of that lemma, $\| h_{\id}(c)\|\leqslant H_*$ and $4K_*H_*\vartheta_*<1$; and the $y$-summand of \eqref{eq:regular-covariance-gauge} is $\mathscr G_{y,\id}(q)=\sum_{x\in\mathcal B_y\setminus\{ y\}}(1-\operatorname{Re}\tr_N\mathcal U_{e^{iq}}(y,x))=\sum_{x\in\mathcal B_y\setminus\{ y\}}f_{(y,x)}(q)$, the sum of the $(y,x)$-summands of $\mathscr G(e^{iq})$ in \eqref{eq:action-primitive}. So $E_{\xi,y}(B)=-g_0^{-2}\mathscr G_{y,\id}(q_{g_0}(B))=\sum_x\bigl(-g_0^{-2}f_{(y,x)}(\psi(g_0\mathsf CB))\bigr)$. Each $f_{(y,x)}\geqslant 0$ vanishes at $q=0$, so $Df_{(y,x)}(0)=0$, and the computation just made for $f_p$ applies to each $f_{(y,x)}$; since the passage from a function to its degree-two Taylor residual at $B=0$ in \eqref{eq:native-source-energy-residual} is linear, $p_{\xi,y}(B)=\sum_{x\in\mathcal B_y\setminus\{ y\}}v_{(y,x)}(B)$.

(iii) In \eqref{eq:native-source-haar-primitive}, $\sigma(q)\,\dd q$ is the pullback under $q\mapsto e^{iq}$ of the Haar measure of CRE (1.11), and in \eqref{eq:literal-residual} $\sigma_H$ is the density of normalised Haar measure in the same coordinates relative to the same Lebesgue measure; the Haar measure of the compact group being unique up to a positive constant, $\sigma=\kappa\sigma_H$ with $\kappa>0$, both are positive on the chart ball (Lemma~\ref{lem:c2-uniform-chart}(iii) takes the real logarithm of $\sigma(q)/\sigma(0)$ there), and $\sigma(q)/\sigma(0)=\sigma_H(q)/\sigma_H(0)= j_H(q)$. With $q_{\id}(B)= q_{g_0}(B)$ from (i), $p^{\mathrm H}_{\xi,b}(B)=\log j_H(q_{g_0}(B)(b))= j_b(B)$, both sides being real logarithms of positive numbers: the branch of $j_b$ in \eqref{eq:perimeter-literal-local-terms} is the real logarithm on the convex domain of (i), on which the argument stays positive.

(iv) By (i) (the row solution being $\varphi_c$, which agrees with the row of $\widetilde D$ on an open set containing $g_0\mathsf CB|_{\mathcal E_c}$, so that their derivatives agree there, and $h_{\id}(c)= L^3\id$), $p^{\mathrm J}_{\xi,c}(B)=\log|\det_{\mathfrak g}(\id-L^3\partial_{A(b_0(c))}\widetilde D_c(g_0\mathsf CB))|$. The determinant is positive on $\| A|_{\mathcal E_c}\|\leqslant \vartheta_*/2$ (Lemma~\ref{lem:c2-uniform-chart}(iii)), so the absolute value is redundant and the expression is $d_c(B)$, whose branch in \eqref{eq:perimeter-literal-local-terms} is the real logarithm of the positive determinant (the sentence after that display).

The first consequence is (ii)--(iv) summed over all labels, with \eqref{eq:perimeter-literal-local-sum}. For the second, every primitive and every term vanishes at zero input: $p_{\xi,\alpha}(0)= E_{\xi,\alpha}(0)=0$ ($s_{p,\id}(0)=0$, $\mathscr G_{y,\id}(0)=0$), $p^{\mathrm H}_{\xi,b}(0)=0$, $p^{\mathrm J}_{\xi,c}(0)=0$ ($\widetilde D(0)=0$, $D\widetilde D(0)=0$), and $v_{\alpha}(0)= j_b(0)= d_c(0)=0$ likewise. A primitive depends on $B$ only through its free-input set (Lemma~\ref{lem:row-locality}, the list after \eqref{eq:native-source-jacobian-primitive}), a term of \eqref{eq:perimeter-literal-local-terms} only through its support. For $B\in\mathfrak g^D$ as in the statement, $i_{D,T}B$ satisfies the hypothesis of the lemma (the maximum norm does not increase), a primitive whose free-input set misses $D$ has zero input at $i_{D,T}B$ and vanishes, and a term whose support misses $D$ vanishes there as well; so the sum of the primitives whose free-input sets meet $D$ and $P_{D,T}(B)$ are both the full sums at $i_{D,T}B$, which agree by the first consequence.
\end{proof}

Define $N_{D,I,T}^{\sigma}$ and $D_{D,T}^{\sigma}$ by \eqref{eq:regional-integrals} with $C_D,P_D,i_D,U_{g_0}$ replaced by $C_{D,T},P_{D,T},i_{D,T},U_{g_0,T}$. Their cutoff is still the product over $D$. Apply the same forest, web and connected-weight definitions to these periodic integrals, obtaining weights $\Phi_{D,I,T}(\Gamma)$, and set
\begin{equation}
    \kappa_{D,I,T}(j)=
        \sum_{\Gamma\, | \,j\in J_{\Gamma}}
            \frac{\Phi_{D,I,T}(\Gamma)}{| J_{\Gamma}|}\;.
\end{equation}
Assume that $I$ has a lift with $mL< N_{\min}/3$, where $N_{\min}$ is the smallest fine period. For both periodic open-segment gases the positive cluster majorants obey a uniform rooted tail. Let $U_{D,I,T}(\Gamma)$ be the absolute Ursell coefficient, including its factorial, times the positive object majorants, and use $U_{I,T}$ for $D= E_{*,T}$. There is $\eta_0>0$, independent of the region and periods, such that for every marked segment $j$,
\begin{equation}\label{eq:regional-periodic-line-tail}
    \sum_{\Gamma\, | \,j\in J_{\Gamma}}
        U_{D,I,T}(\Gamma)e^{\eta_0\ell(\Gamma)}
        \leqslant C\tau^2\;.
\end{equation}
The majorants dominate the absolute connected weights. The proof is the positive-majorant argument of Proposition~\ref{prop:regional-coefficient}: restricting supports preserves the label overlap bounds, and the periodic shell sum is bounded by the infinite-lattice shell sum. Under this lift hypothesis, index separation along the segment is its periodic separation. Thus \eqref{eq:marginal-hull-extent} holds on the lifted segment. Its ordinary intervals have the straight-line interval tree sum. Hence the same load inequalities and extent reserve prove \eqref{eq:regional-periodic-line-tail} for the full and restricted periodic gases, including the covariance interpolations used below.

On a compatible torus with fine periods $N_{\mu}$, put $N_{\min}=\min_{\mu} N_{\mu}$ and assume that $I$ has a lift with $mL< N_{\min}/3$. The set $D$ may wrap around the torus. Define $\kappa_{I,T}(j)$ by the same allocation with $D= E_{*,T}$. In this paragraph $d_j$ is \eqref{eq:regional-distance} with $d_T$ in place of $d$, including its endpoint term $L\min(j,m-1-j)$. For bounded $d_j\leqslant 2(R_0+R_*+c_{\mathrm{reg}})$, \eqref{eq:regional-periodic-line-tail} and $|\kappa_v|\leqslant C\tau^2$ give \eqref{eq:regional-torus-bound} directly after enlarging $C$. For larger $d_j$, the periodic version of \eqref{eq:regional-resolvent} and its differentiated cluster estimate compare $\kappa_{D,I,T}(j)$ with $\kappa_{I,T}(j)$. Clusters of extent at least $d_j/(2c_{\mathrm{reg}})$ have total rooted weight at most $C\tau^2e^{-cd_j}$. For smaller clusters the labels agree, and their carriers are at distance at least $d_j/2-R_0$ from $D^c$. The covariance derivative therefore has the same bound, after absorbing the fixed factor $e^{aR_0}$. This proves
\begin{equation}
    |\kappa_{D,I,T}(j)-\kappa_{I,T}(j)|
        \leqslant C\tau^2e^{-cd_j}\;.
\end{equation}

To compare the second density with the infinite straight-line density, lift the root segment to the chosen copy of $I$. A cluster network of length $c_{\mathrm{reg}}\ell+R_0< N_{\min}/4$ lifts uniquely and preserves carrier intersections. The distinct translates of $I$ are separated by more than $2N_{\min}/3$ because $mL< N_{\min}/3$, so all its source marks lie on that copy. If also $c_{\mathrm{reg}}\ell+R_0< L\min(j,m-1-j)$, the network avoids the two ends and has the same labels and ordered source factors as the infinite straight-line expansion. Clusters failing either restriction are bounded by the rooted extent tail. For the remaining clusters only the covariance marginals change. Periodisation and lifted separation $d(b,c)\leqslant N_{\min}/4$ give
\begin{equation}
    \| C_T(b,c)-C_{\infty}(b,c)\|
        \leqslant C e^{-\mu(N_{\min}-d(b,c))}
        \leqslant C e^{-\mu N_{\min}/2}e^{-\mu d(b,c)}\;,
\end{equation}
after expressing the decay exponent in fine-bond units. The free-inverse periodisation follows by applying $H_T$ to the absolutely convergent periodised kernel of $H^{-1}$, as in the proof of \eqref{eq:bulk-periodisation}. The covariance interpolation and the density-direction estimates in the proof of Proposition~\ref{prop:regional-square-extraction} apply with $h_A= e^{-\mu N_{\min}/2}$. They retain edge decay and the quadratic source gain; their polynomial carrier factors are absorbed by the reserved extent weight. With $\iota_T=\max(0,N_{\min}/4-R_0)$, summing the short-cluster derivatives and the two tails gives
\begin{equation}
    |\kappa_{I,T}(j)-\kappa_v|
        \leqslant C\tau^2\left[
            e^{-c(L\min(j,m-1-j)-R_0)_+}+e^{-c\iota_T}\right]\;.
\end{equation}
Here $x_+=\max(x,0)$. Absolute convergence also gives $\log(N_{D,I,T}^{\sigma}/D_{D,T}^{\sigma})=\sum_{j=0}^{m-1}\kappa_{D,I,T}(j)$, since summing the allocation weights $| J_{\Gamma}|^{-1}$ over the marked segments gives one for each source cluster. Adding the two comparisons proves
\begin{equation}\label{eq:regional-torus-bound}
    |\kappa_{D,I,T}(j)-\kappa_v|
        \leqslant C\tau^2(e^{-cd_j}+e^{-c\iota_T})\;.
\end{equation}

For the identification with $K_{\xi}$, take the averaging prescription RG1 (0.4), the bond cutoff $|\bigcdot |_{\sigma}$, and the threshold $\delta_0= g_0r_0(g_0)$ on both sides. The comparison applies separately to either specified $\sigma$; it does not determine an otherwise unspecified norm convention. Let $I$ consist of consecutive traversals of one straight side of $C$. Let $E_{\Omega,*}$ be the free carrier after the designated-bond elimination in \cite[(1.14)--(1.15)]{BalabanCRE}. Let $D_-$ consist of these bonds with both endpoints in $\Lambda_1$, and $D_+$ of those with at least one endpoint there. The comparison holds for every $D_-\subseteq D\subseteq D_+$, assigning the complementary free variables to the exterior integral. It therefore covers either endpoint convention for the integration boundary.

We make the geometric parameter choices in this identification explicit. In CRE's notation let $M_2,M,R_1$ be the cube-scale integers, so the test cubes have fine side $L^2M_2R_1$ and the macro cubes have fine side $LMR_1$; here $R_1$ denotes the scale integer, not the large-fluctuation region also denoted by $R_1$ in CRE. Take aligned nested partitions, with all periods divisible by $LMR_1$, and choose the powers of $L$ so that
\begin{equation}\label{eq:regional-collar-sizes}
    M\geqslant 8LM_2,
    \qquad LMR_1>4R_0+4L^2M_2R_1+8\;.
\end{equation}
Here $R_0=4(1+R_*+R_H+L)$ is our fixed range constant, distinct from the scale integer in CRE. These are explicit choices within the compatible-partition and large-$M$ prescription of \cite[Section~1]{BalabanCRE}. The collar inequality is included among the common parameter conditions. One may enforce it by first choosing $M$ so that $LM-4L^2M_2>4R_0+8$; it then holds for every scale integer $R_1\geqslant 1$. For macro regions we use the collar-distance statement following \cite[(1.10)]{BalabanCRE}; it supplies the metric separation, including near cube corners. For test-cube enlargements we need only a conservative reach bound: each added layer has width at most $L^2M_2R_1$. This bound also covers a layer of the $L$-fold finer cubes in the enlargement convention of CRE, Section~1; no equality between test and macro enlargements is used. The parallelepiped modification after (1.20) removes macro cubes from $\Lambda_1$; the remaining region is still a union of macro cubes and hence of test cubes.

Using (1.20) and that collar-distance statement, a site within distance $LMR_1-1$ of $\Lambda_1$ lies in $\Omega_1^{\sim-1}\setminus R_1'$. Every stencil meeting $D\subseteq D_+$ therefore remains in the regular region. A bond with an endpoint in $\Lambda_1$ is tested by the test cube containing that endpoint, which is contained in $\Lambda_1$ and is good. A bad test or a test in $S_1$ uses only bonds within $2L^2M_2R_1+1$ of its cube. The collar and \eqref{eq:regional-collar-sizes} keep those bonds outside $D_+$, including the fixed stencil enlargement. At identity determining data the critical field is the identity: it satisfies the constraints, has zero action and gauge-fixing functional, and is unique in the regular chart by \cite[(1.12)]{BalabanCRE}. Its force is zero. The Hessian on $D$ is consequently $H_{D,T}$, including the diagonal contributions of plaquettes crossing its boundary.

The nonlinear constraint in \cite[(2.4)]{BalabanRG1} becomes $M(e^{iB'})$ at identity background. Its principal logarithm is $\mathcal Q(B')$ from \eqref{eq:constraint-primitive}. The linearisation following (2.10) has linear part $L\widetilde Q B'$, hence $\widetilde Q= Q$ there. Its designated-row right inverse satisfies $LQh= I$. Uniqueness of the small solution for $\widetilde D$ gives exactly \eqref{eq:flat-phase}. Thus the action, Haar and Jacobian terms become $v_T(i_{D,T}B)$ with all primitive labels meeting $D$ retained.

The split reconstruction in \cite[(1.21)]{BalabanCRE} also equals the identity at identity data, and $Q_1^{s*}\id=\id$ by (1.3). Its logarithmic argument is zero. Equation (1.21) then gives $e^{iL^{-1}\mathbf H_{1,Ax}(0)}=1$; the small principal chart implies $\mathbf H_{1,Ax}(0)=0$. Since the BCH remainder satisfies $\mathbf F(X,0)=0$, the boundary change (1.22) is the identity, its Jacobian is one and its additional exponent is zero at this projection.

For the cutoff, every free bond in $D$ has an endpoint in a compatible good cube contained in $\Lambda_1$. The twice-enlarged starred test in \cite[(1.18)]{BalabanCRE} covers it. The designated-bond exclusion is the verbal convention following (1.8). A good cube outside $\Lambda_1$ can also test bonds in $D$: split its product of bond indicators into its $D$ and exterior factors. The $D$ factors are redundant with the covered interior cutoff by idempotence; all exterior factors stay in the outer integral. Bad tests and changed background tests are separated from $D$ by the collar. At zero exterior free field the conditional cutoff is therefore $\Theta_D^{\sigma}$.

The same norm is used on both sides. The threshold is $r=\delta_0/g_0$ by the normalisation convention \eqref{eq:cre-cutoff-normalisation}. The exterior quadratic cross term and its completion-of-square compensator vanish at the projection. Gaussian and constraint normalisations cancel between the two conditional kernels. They are not removed from the exterior integration.

Finally set $t_a=1$ on $I$ and $t_a=0$ elsewhere. At the flat projection, \eqref{eq:regional-wilson-polynomial} is precisely the closed relative source in \eqref{eq:regional-integrals}. Denoting the projected conditional kernels by $K_{\xi}^{\mathrm{flat}}$, we have
\begin{equation}\label{eq:regional-kernel-identification}
    \frac{K_{\xi}^{\mathrm{flat}}(1_I)}{K_{\xi,\mathrm{vac}}^{\mathrm{flat}}}
        =\frac{N_{D,I,T}^{\sigma}}{D_{D,T}^{\sigma}}\;.
\end{equation}
If $d_T(K_j,D^c)\geqslant b$ for every source segment $j$, then \eqref{eq:regional-torus-bound} gives
\begin{equation}\label{eq:regional-mass-identification}
    \left|-\frac{1}{mL\varepsilon}
        \log\frac{K_{\xi}^{\mathrm{flat}}(1_I)}
                    {K_{\xi,\mathrm{vac}}^{\mathrm{flat}}}
            -\delta m_v^{\sigma}\right|
        \leqslant \frac{C\tau^2}{L\varepsilon}
            \left(\frac{1}{m}+e^{-cb}+e^{-c\iota_T}\right)\;.
\end{equation}
For every domain label $\xi$ for which the buffer fits inside the torus and $mL< N_{\min}/3$, \eqref{eq:regional-mass-identification} bounds the distance between the per-length scalar and $\delta m_v^{\sigma}$ by $C\tau^2(1/m+e^{-cb}+e^{-c\iota_T})/(L\varepsilon)$ on the fixed torus; along any family of such labels with $m$, $b$ and $\iota_T$ tending to infinity this bound tends to zero, so the per-length scalar converges to $\delta m_v^{\sigma}$, which identifies it with \eqref{eq:marginal-full-perimeter}. The limit is the one defining the coefficient, as in \eqref{eq:marginal-square-mass}; the estimate \eqref{eq:regional-mass-identification} itself is a statement on each fixed torus. The four oriented straight-side coefficients give the square coefficient by the average in \eqref{eq:marginal-square-mass}. This comparison is made inside the conditional kernels; the outer history weights and the physical denominator \eqref{eq:one-step-physical-partition} retain their definitions.

\subsection{The Regional Square Scalar}\label{sec:regional-square-scalar}

Fix either cutoff prescription $\sigma$. Let $C_n$ be a coarse translate of the counterclockwise square with vertices $0,nLe_1,nL(e_1+e_2),nLe_2$ in fine-lattice units, and assume $K_{\mu}\geqslant 3(n+2)$ for every coarse period. Let $D$ be the free bonds integrated in a covered conditional kernel: every free-input stencil of the contour lies in $D$. Retain the original primitive labels, including those crossing the boundary of $D$, and define
\begin{equation}\label{eq:regional-square-scalar}
    S^{\circ}_{D,C_n,T}=
    \frac{\displaystyle\int\Theta_D^{\sigma}(B)e^{P_{D,T}(B)}
        \chi\bigl(\Hol_{U_{g_0,T}(i_{D,T}B)}C_n\bigr)\,\dd\mu_{C_{D,T}}(B)}
    {\displaystyle\int\Theta_D^{\sigma}(B)e^{P_{D,T}(B)}\,\dd\mu_{C_{D,T}}(B)}\;,
    \qquad C_{D,T}=(i_{D,T}^*H_Ti_{D,T})^{-1}\;.
\end{equation}
Here $H_T$ is the full flat free-field precision on the torus, and $P_{D,T}$ is \eqref{eq:regional-periodic-weight}; $\Theta_D^{\sigma}$ is the bond cutoff of \eqref{eq:regional-weight}. For a conditional history $\xi$, write $S^{\circ}_{\xi,C_n}= S^{\circ}_{D,C_n,T}$ with its young integration carrier $D$. This is a ratio of conditional reference integrals; the exterior indicators and the physical history weight are not evaluated at the identity.

Index the square's $4n$ coarse segments cyclically by $j$. Let $K_j$ be the union of the free-input stencils of its fine traversals. With complements taken in the torus free carrier and with the periodic fine-bond metric of \eqref{eq:regional-distance}, put
\begin{equation}\label{eq:regional-square-distances}
    d_j=\max\bigl(0,d_T(K_j,D^c)-R_0\bigr),
    \qquad \iota_T=\max\bigl(0,N_{\min}/4-R_0\bigr)\;.
\end{equation}
The fixed $R_0$ includes the label, reconstruction and precision ranges. Distance to the empty set is infinite, with $e^{-\infty}=0$. There is no distance to a side endpoint in $d_j$.

\begin{proposition}\label{prop:regional-square-extraction}
Under these hypotheses, the collar conditions of the conditional kernel, and the common small-coupling thresholds of the flat and regional expansions, the scalar \eqref{eq:regional-square-scalar} is nonzero. Its logarithm is continued from zero in the traversal-source polynomial \eqref{eq:regional-wilson-polynomial}. There are $C,c>0$, independent of $D,n$ and the compatible periods, and an error $E^{\circ}_{D,C_n,T}$ such that
\begin{equation}\label{eq:regional-square-extraction}
    S^{\circ}_{D,C_n,T}
        = e^{-\delta m^{\sigma} P(C_n)}
          \bigl(z_{\mathrm c}^{\sigma}(g_0,\pi/2)\bigr)^4
          e^{E^{\circ}_{D,C_n,T}}\;,
\end{equation}
\begin{equation}\label{eq:regional-square-error}
    | E^{\circ}_{D,C_n,T}|
        \leqslant C\tau^2\left[
            \sum_{j=0}^{4n-1}e^{-cd_j}
            +4ne^{-c\iota_T}+(n+1)e^{-cn}\right]\;.
\end{equation}
The coefficients are exactly \eqref{eq:marginal-square-mass} and \eqref{eq:corner-square-factor}, with the same cutoff prescription. The parameter $\tau$ is \eqref{eq:marginal-fixed-forest-bound}; fixed enlargements of its majorant constants are absorbed in $C$ and the coupling threshold.
The collar hypotheses identify this integral with the conditional CRE scalar. The estimates for the defined flat integral use its covered stencils, spectral bounds and small-coupling conditions; the collar enters at the identification, not in the subsequent cluster comparison.
\end{proposition}

\begin{proof}
Write $F_{D,C,T}=\log S^{\circ}_{D,C,T}$, with its branch specified below, $F_{C,T}=\log(N_{C_n,T}^{\sigma}/D_T^{\sigma})$, and $F_{C,\infty}(n)= F_{\square,\infty}(n)$. We compare the first two quantities, then the latter two, before applying the bulk corner extraction. This retains errors from changing the kernel and summation domain, as in \cite[(4.35)--(4.37)]{BalabanRG1}; the exterior integration of \cite[(1.29)--(1.30)]{BalabanCRE} is unchanged.

At the conditional reference data the background and exterior determining fields are the identity and the retained free field is zero. The row precision entries are the flat torus entries on $D$, including the diagonal terms from crossing plaquettes. The collar contains every stencil touching $D$, so the precision is $i_{D,T}^*H_Ti_{D,T}$. The primitives agree with the terms of \eqref{eq:perimeter-literal-local-terms} label by label, with the designated inverse $L^3$ at the identity and the gauge terms of a centre summed over its pairs $(y,x)$, by Lemma~\ref{lem:identity-primitives}, whose hypothesis on $B$ is the pair of small-coupling conditions fixed in the next paragraph; by its second consequence the sum of the primitives whose free-input sets meet $D$, at $i_{D,T}B$, is $P_{D,T}(B)$, terms whose support misses $D$ vanishing at $i_{D,T}B$. The cutoff is the same product of bond tests. Thus the numerator and denominator in \eqref{eq:regional-square-scalar} are those of the conditional reference scalar, with their normalisers retained in the ratio.

The local implicit solutions used in the two representations have the same analytic germ. To fix their common domain, let $C_{\mathsf C}$ and $\vartheta_*$ be the constants of Lemma~\ref{lem:c2-uniform-chart} ($C_{\mathsf C}$ bounds the fixed row map from the maximum free-bond norm to the row Hilbert norm, $\vartheta_*$ is its implicit-solution phase radius), and let $a_L$ be the flat chart radius in Appendix~\ref{sec:perimeter-attempt}. The row map is finite dimensional with finitely many translated types, so $C_{\mathsf C}<\infty$; the analytic implicit-function theorem gives a uniform positive phase radius. Reduce $a_L$ if necessary to lie within those row charts and their nonvanishing Jacobian domains. Choose the coupling so that
\begin{equation}
    g_0C_{\mathsf C}(2+4r)<\frac{\vartheta_*}{2},
    \qquad g_0(2+4r)< a_L\;.
\end{equation}
The factor $2+4r$ specifies a larger free-variable polydisc containing the integration box $\| B\|_{\infty}\leqslant r$ with a positive margin. The convex intersection of this polydisc with the phase balls lies in the common analytic charts. The identity theorem therefore identifies the solutions throughout the integration box. The common chart radius is fixed before the coupling; reducing it preserves all previous estimates by restriction to a smaller domain. These small-coupling requirements and all polymer-load inequalities are imposed together, which is possible because $g_0(1+r)^3\to0$ and $\tau\to0$.

Apply the protected-label expansion with supports intersected with $D$, keeping every original label. Zero extension commutes with constant conjugation, so the restricted primitives, cutoffs and closed source differences are invariant. The restricted support graph and forest partition the labels into components. Merge source components using the manuscript's shortest cyclic intervals, with its fixed tie-breaking base; a web interval is their union, possibly the whole circle. Compatible components and webs have disjoint carriers and disjoint web intervals, so their union reverses this decomposition without losing or duplicating labels. If a web crosses the base, cyclicity of the trace moves the base between web intervals before Schur factorisation. This proves the regional counterpart of \eqref{eq:marginal-exact-gases}. The alternative grouping of the conditional kernel need not agree term by term: both constructions give the same integral ratio.

The required quantitative bounds hold uniformly for these restricted labels. Principal restriction and inversion preserve the spectral bounds, and the finite-range Neumann series gives uniform exponential covariance decay in the periodic metric. Let $\mu_1>0$ be a remaining edge-decay exponent after the degree payment, no larger than one eighth of the covariance and comparison exponents. Set
\begin{equation}
    \alpha'=\frac{\mu_1}{1024(1+c_h)},
    \qquad \eta'=\frac{\mu_1}{1024},
    \qquad a'(p)=| A_p|+\alpha'H_p\;.
\end{equation}
Restriction preserves the support-size, overlap and range bounds $s_*,q_*,R_*$. The same labelled-tree count as in \eqref{eq:marginal-load-series}, with a fixed enlarged $K_2$, bounds vacuum and source loads by
\begin{equation}
    \sum_{k\geqslant 1}K_2^k\tau^{\lceil k/s_*\rceil},
    \qquad
    \sum_{k\geqslant 1}K_2^k\tau^{\max(2,\lceil k/s_*\rceil)}\;.
\end{equation}
The bounds remain valid with the reserve $e^{4| A_{\gamma}|+8\alpha'H_{\gamma}+8\eta'\ell(\gamma)}$. Contacts remain within $R_*$ of the restricted support network, so \eqref{eq:marginal-hull-extent} applies to the square, including at its corners. The cyclic conversion \eqref{eq:marginal-hull-root} has the same uniform geometric sum $1+2/(e^{\alpha'}-1)$. The web-tree recursion therefore gives positive majorants $u_p$ with
\begin{equation}
    \begin{aligned}
        \sup_b\sum_{p\, | \,b\in A_p}u_pe^{2a'(p)+2\eta'\ell(p)}&\leqslant C_*'\tau\;,\\
        \sup_x\sum_{W\, | \,x\in h_W}u_We^{2a'(W)+2\eta'\ell(W)}&\leqslant C_*'\tau^2\;.
    \end{aligned}
\end{equation}
The source-only spatial load is also $O(\tau^2)$. Choose $g_0$ last so that $K_2^{s_*}\tau\leqslant 1/2$, $C_*'\tau\leqslant 1/4$ and $C_*'\tau^2\leqslant \alpha'/4$. The incompatibility sum is then at most $a'(p)/4$. The pinned-tree proof of \eqref{eq:marginal-pinned-bound} yields
\begin{equation}\label{eq:regional-square-rooted-majorant}
    \sum_{\Gamma\, | \,j\in J_{\Gamma}}U(\Gamma)e^{\eta'\ell(\Gamma)}\leqslant C\tau^2\;,
\end{equation}
where $U$ is the absolute Ursell coefficient times the positive object majorants. The factor $2^{| S|}$ for source variables $| t_a|\leqslant 2$ is included in $K_2$. Thus the regional source logarithm exists on that polydisc and vanishes at zero source. Its exponential is the ratio \eqref{eq:regional-square-scalar}, proving nonvanishing.

For the boundary comparison, the periodic resolvent identity is
\begin{equation}
    C_{D,T}-(C_T)_{D,D}
        = C_{D,T}(H_T)_{D,D^c}(C_T)_{D^c,D}\;.
\end{equation}
If $u\in D$, $w\in D^c$ and $d_T(u,w)\leqslant R_H$, put $s= d_T(b,u)+d_T(w,c)$. The periodic triangle inequalities give
\begin{equation}
    d_T(b,D^c)+d_T(c,D^c)+d_T(b,c)\leqslant 2s+2R_H\;.
\end{equation}
Half the inverse-decay exponent pays these distances; the remaining convolution is bounded by a fixed neighbour count times $\sup_b\sum_u e^{-\mu d_T(b,u)/2}$. Periodic shells have at most the infinite-lattice shell counts, so this sum is uniform even for wrapping or disconnected $D$. Therefore
\begin{equation}
    \| C_{D,T}(b,c)-C_T(b,c)\|
        \leqslant Ce^{-a(d_T(b,D^c)+d_T(c,D^c)+d_T(b,c))}\;.
\end{equation}

Interpolate the principal covariances as in \eqref{eq:regional-density-derivative}. Here is explicit control of the two derivative terms. On a carrier $A$ of size $k$, use a common positive edge majorant $B_{bc}= A_0e^{-\mu_0d_T(b,c)}$ and write $h_A= e^{-a d_T(A,D^c)}$. The covariance direction satisfies $\|\mathcal D_A\|\leqslant Ckh_A$ and $\|\mathcal D(b,c)\|\leqslant Ch_AB_{bc}$. For forest degrees $d_b$, use radii
\begin{equation}
    r_{\{ b,c\}}=\frac{\lambda}{32d_bd_cB_{bc}},
    \qquad r_0=\frac{\lambda}{32Ckh_A}\;.
\end{equation}
The forest perturbation has norm at most $\lambda/32$, and the extra centre direction $\mathcal D_A\circ w^F$ has the same bound on its circle, since a positive Schur mask with unit diagonal is contractive. The total is below $\lambda/8$, so the normalised density has the Gaussian majorant \eqref{eq:marginal-density-majorant}. Replacing one edge direction by $E_e(\mathcal D)$ instead uses radius $\lambda/(32d_bd_cCh_AB_{bc})$ and costs $Ch_A$ against the same forest majorant. Summing the centre and edge derivatives costs at most $C(k+| F|+1)^2h_A$. Zero directions give zero derivatives; no covariance entry is divided by. Constant-conjugation invariance survives differentiation. A single marked component therefore retains its quadratic character gain, also when it crosses a corner. The product rule over repeated object occurrences and \eqref{eq:regional-square-rooted-majorant} give the differentiated cluster bound used in the regional proposition, with the same $\tau^2$ gain.

Allocate each source cluster to its marked coarse segments with weight $| J_{\Gamma}|^{-1}$. A network joins its carrier, stencil edges and contour intervals with length at most $C_{\rm net}\ell(\Gamma)+R_0$. The interval may turn a corner, but its fine-path length is still bounded by $c_h\ell$; anchor it at the source mark in segment $j$, hence within $R_0$ of $K_j$. For $d_j>2R_*+2R_0$, split at $\ell= d_j/(2C_{\rm net})$. Short clusters have their original stencils inside $D$, identical labels and cyclic intervals in both expansions, and carriers at distance at least $d_j/2-R_0$ from $D^c$, whose fixed loss is absorbed in the constant. Their covariance derivative is bounded by the preceding estimate. Long clusters are paid by \eqref{eq:regional-square-rooted-majorant}. Bounded $d_j$ are covered by the same rooted bound after enlarging $C$. This proves
\begin{equation}\label{eq:regional-square-boundary-comparison}
    | F_{D,C,T}-F_{C,T}|\leqslant C\tau^2\sum_{j=0}^{4n-1}e^{-cd_j}\;.
\end{equation}
No side-endpoint terms occur because both contours are the same closed square.

For the periodic comparison, the Neumann estimate for the free inverse gives $\| C_{\infty}(e,f)\|\leqslant \gamma_L^{-1}e^{-\mu_L\|\ell(e)-\ell(f)\|_{\infty}}$. Since $d_T(e,f)\leqslant L\|\ell(e)-\ell(f)\|_{\infty}+2L$ for fixed lifts, it gives exponential decay in fine units with exponent $\nu=\mu_L/L$. Periodising this free inverse solves the torus inverse equation. For lifted separation $d\leqslant N_{\min}/4$, summing the image shells gives
\begin{equation}
    \| C_T(e,f)-C_{\infty}(e,f)\|
        \leqslant Ce^{-\nu(N_{\min}-d)}
        \leqslant Ce^{-\nu N_{\min}/2}e^{-\nu d}\;.
\end{equation}
The same density-direction radii apply, with $h_A= e^{-\nu N_{\min}/2}$ and a reduced common decay exponent. A rooted cluster with $C_{\rm net}\ell+R_0< N_{\min}/4$ lifts uniquely. Projection is injective on its network, preserves carrier intersections, and preserves the fixed ordered square traversal list and its cyclic intervals. Distinct square images are farther than $2N_{\min}/3$ under the stated period hypothesis, so all its source marks lift to the chosen square. Thus the short torus and bulk clusters correspond with identical combinatorics; only their covariance marginals differ. The other clusters, including winding ones, are paid by their extent tails. Summing over $4n$ roots gives
\begin{equation}\label{eq:regional-square-periodic-comparison}
    | F_{C,T}-F_{C,\infty}(n)|\leqslant C\tau^2(4n)e^{-c\iota_T}\;.
\end{equation}

Finally use \eqref{eq:corner-square-extraction-bound}. The straight sum equals $-\delta m^{\sigma} P(C_n)$ and the four oriented corner logarithms sum to $4\log z_{\mathrm c}^{\sigma}$, by their definitions. Define
\begin{equation}
    E^{\circ}_{D,C_n,T}=(F_{D,C,T}-F_{C,T})
        +(F_{C,T}-F_{C,\infty}(n))+E_n^{\sigma}\;.
\end{equation}
Equations~\eqref{eq:regional-square-boundary-comparison}, \eqref{eq:regional-square-periodic-comparison} and \eqref{eq:corner-reference-matching} give the result. Every logarithm is the connected source branch, so the exponential identity has no branch ambiguity. The adjacent-side clusters and straight endpoint deficits remain in the existing corner coefficient.
\end{proof}

For fixed geometry, $E^{\circ}_{\xi,C_n}$ depends only on the regional labels and is therefore measurable under the original history sum. Its bound is finite at fixed cutoff. The scalar factorisation consequently leaves the exterior operation, its conditional normalisers and $\mathcal Z_0$ in place. The boundary sum in \eqref{eq:regional-square-error} can be of order $n$; Proposition~\ref{prop:regional-square-extraction} does not assert that $e^{E^{\circ}}$ is close to one for a loop running near the integration boundary.

\section{The Insertion Retained after One Step}
\label{sec:remaining-physical-comparison}

The first step produces a function of the coarse field through \eqref{eq:outgoing-loop}. Its representation below retains the closed flux insertions generated by the block average, together with the other source terms. The scalar factor extracted from the reference integral is
\begin{equation}\label{eq:extracted-loop-scalar}
    \kappa_C= e^{-\delta m^{\sigma} P(C)}
        \bigl(z_{\mathrm c}^{\sigma}(g_0,\pi/2)\bigr)^4\;.
\end{equation}
We use the same cutoff prescription and the counterclockwise $(e_1,e_2)$ square and its coarse translates, with the period, collar and coupling hypotheses of Proposition~\ref{prop:regional-square-extraction}, throughout this section. The physical conditional cutoff and its flat projection use the same fixed $\sigma$ (the cutoff norm $|\bigcdot |_{\sigma}$ of \eqref{eq:literal-cutoff}, not the Haar density $\sigma$) and threshold $r=\delta_0/g_0$. The scale parameters satisfy \eqref{eq:regional-collar-sizes} and the standing conditions \eqref{eq:c-standing-conditions}. Closeness to the integration boundary remains in $E^{\circ}_{\xi}$, the error of Proposition~\ref{prop:regional-square-extraction} at the carrier $I$, defined at the end of this subsection.

The first step is the operator \eqref{eq:bare-blocking-operator} with the averaging operation \eqref{eq:global-average}, Borel and gauge covariant by Lemma~\ref{lem:global-average}, and $\rho_1$, $N_1$ are the measures \eqref{eq:first-step}. Besides CRE's construction as recorded in (B1)--(B4), the representation below uses the Borel choices of minimal configurations of Lemma~\ref{lem:borel-minimiser}, the Borel background $U_1$ and row solutions of the measurability paragraph of Section~\ref{sec:branch-data}, the bijection $\Psi_{\xi}\colon\mathcal D_{\xi}\to\mathcal A_{\xi}$ with invertible differential of Lemma~\ref{lem:boundary-bijection}, which is the (1.22) part of (B4)(ii) (its hypotheses: $\delta_0\leqslant \delta_G$ is listed among the further conditions of Appendix~\ref{app:thresholds} and holds for small $g_0$, as $\delta_0= g_0r(g_0)\to0$; CRE's bound \eqref{eq:boundary-bijection-inputs} on $\mathbf H_{1,Ax}$, p.251, at the tolerance $\delta_0/40$ is the condition ``$MR_1$ large'' of item 5 there, met for small $g_0$ at fixed $M$ since $R_1\to\infty$, with $\varepsilon_1/\delta_0\leqslant (1+\beta_0)A_0/A_1$ by CRE (2.8) and the constants independent of the coupling by (B3)(iii)), the invertible differential of $B\mapsto U(B)$ (the (1.14) part of (B4)(ii)) and the nonvanishing Jacobian of (1.14) ((B4)(iii); its positivity on $\mathcal A_{\xi}$, Lemma~\ref{lem:c2-fibre-smallness}(iv), is used below only to match CRE's $\operatorname{Tr}\log$ with $\log|\det|$), Lemma~\ref{lem:collar-coverage}(ii),(iv) and Lemma~\ref{lem:collar} for the $a$-independence of the exterior factors and the box form of the young fibre, Lemmas~\ref{lem:charts} and~\ref{lem:collar} for the boundedness of the integrand on the box, the row factorisation of the Jacobian of (1.14) (proof of Lemma~\ref{lem:primitive-sum}, from Lemmas~\ref{lem:two-block} and~\ref{lem:row-locality}), the list of operations after \eqref{eq:native-source-physical-ratio} with \eqref{eq:one-step-physical-partition}, the disintegration theorem for standard Borel spaces \cite[(17.35), Theorem 15.6]{Kechris}, and, for the coverage event, the stencils of Appendix~\ref{sec:perimeter-attempt}, Lemma~\ref{lem:covered} and Proposition~\ref{prop:regional-square-extraction}; these are the facts about the unrestricted first step that the formulas below need, and each is proved or cited where it is named.

\subsection{The Conditional Integral}

With these inputs, this section uses the unrestricted transformation of \cite{BalabanCRE} (0.1) and the first-step identities (1.11)--(1.15) with the cutoff (1.25), as recorded in Section~\ref{sec:branch-data} under (B1)--(B4), with the averaging choice permitted after (0.1); CRE's (1.26)--(1.29) are not used, except that (1.27) defines the first-step coupling $g_1(g_0)$ entering $R_1$ and $\varepsilon_1$ (Section~\ref{sec:branch-data}). The outer data $\xi$ consist of the region labels, the coarse field $V$, the exterior fine field $V_0$ and the exterior free fluctuation $A_{\mathrm e}$ in those identities. For the integrations, partition positive fine bonds into $E_{\Omega}=\{ b\in E\, | \,b_-\in\Omega_1\text{ or }b_+\in\Omega_1\}$ and $E\setminus E_{\Omega}$. This disjoint integration convention assigns crossing bonds to $E_{\Omega}$ and specifies $V_0$ only on bonds with both endpoints outside $\Omega_1$. Let $E_{\Omega,*}$ be the free bonds remaining after the designated-bond elimination on $\Omega_1$, and let $D= I=\Lambda_1^*$ be the bonds integrated in $\Lambda_1$, the carrier $D_+$ of Section~\ref{sec:regional-identification} (the comparison there holds for every $D_-\subseteq D\subseteq D_+$; this section is at $D= D_+= I$, the carrier of Section~\ref{sec:c4c-1}, which the definition of the young fibre \eqref{eq:retained-change-domain} as a section of $\mathcal D_{\xi}\subseteq\mathfrak g^J\times\mathfrak g^I$ requires). Put $J= E_{\Omega,*}\setminus D$. The variable $a\in\mathfrak g^D$ is integrated and $A_{\mathrm e}\in\mathfrak g^J$ is held fixed. With $i_D$ and $i_J$ denoting extension by zero into $\mathfrak g^{E_{\Omega,*}}$, set $z_{\xi}(a)= i_Da+i_JA_{\mathrm e}$. Lebesgue measure $\dd a$ is the product measure from the fixed Lie-algebra inner product.

Write $U_1(\xi)$ for the critical configuration of CRE (1.12), in its axial mean gauge ((B1)). Let $\mathsf C_{\xi}$ be the linear solution map from free fields to fields satisfying the linearised constraints in CRE (1.14)--(1.15), and set $\mathsf H_{\xi}=\mathsf C_{\xi}^*\Delta_{1,\xi}\mathsf C_{\xi}$, with $\Delta_{1,\xi}$ given by CRE (1.16). Let $h_{\xi}$ be the designated-bond right inverse of the linearised constraint and $\widetilde D_{\xi}$ the analytic correction of CRE (1.14), under the reading of (B4)(i) (the corresponding construction for RG1's own step is the one following RG1 (2.10), p.267): they are $h_{U_1}$ and $\widetilde D_{U_1}$ of Section~\ref{sec:source-phase-primitives}, with the row solutions and their domains defined there. The inverse boundary change $\Psi_{\xi}$ is CRE (1.22), acting on the free fields in $\mathfrak g^{E_{\Omega,*}}$ and equal to the identity off $S_1^*$. Define
\begin{equation}\label{eq:retained-reconstruction}
    \begin{aligned}
        q_{\xi}(s)&= s-h_{\xi}\widetilde D_{\xi}(s)\;,\\
        U_{\xi}(a,b)&=
            \exp\bigl(iq_{\xi}(g_0\mathsf C_{\xi}\Psi_{\xi}(z_{\xi}(a)))(b)\bigr)
                U_1(\xi,b)\;.
    \end{aligned}
\end{equation}
The second formula applies on $E_{\Omega}$; on $E\setminus E_{\Omega}$, $U_{\xi}= V_0$. The map $\mathsf C_{\xi}$ takes $\mathfrak g^{E_{\Omega,*}}$ into $\mathfrak g^{E_{\Omega}}$. The variable $a$ on $I$, and $A_{\mathrm e}$ on $J\setminus S_1^*$, are CRE's principal fluctuation variables of (1.13), the chart logarithms of $UU_1(\xi)^{-1}$ divided by $g_0$ on those free bonds ($\Psi_{\xi}$ is the identity there; $I\cap S_1^*=\varnothing$ by \eqref{eq:c-collar-distance}); on $S_1^*\subseteq J$, $A_{\mathrm e}$ is the variable of (1.22), the (1.13) variable there being $\psi_b(A_{\mathrm e}(b))$; and the domain of the $a$-integration is specified in these variables below. With $\chi_0'(\Omega_1)$ denoting the product of the group tests (1.8) retained in (1.11), define
\begin{equation}\label{eq:retained-source-domain}
    \begin{aligned}
        \mathcal U_{\xi}&=\bigl\{ U\in G^E\, | \,
            U|_{E\setminus E_{\Omega}}= V_0,
            M(U)|_{\Omega_1^{[1]}}= V|_{\Omega_1^{[1]}},
            \mathcal I_{\xi}(U)>0\bigr\}\;,\\
        t_{\xi}(U)&=\frac{1}{i}\log(UU_1(\xi)^{-1})|_{E_{\Omega}}\;,\\
        \widehat{\mathcal A}_{\xi}&=\bigl\{ g_0^{-1}t_{\xi}(U)|_{E_{\Omega,*}}
            \, | \,U\in\mathcal U_{\xi},\ UU_1(\xi)^{-1}\text{ in the chart domain on }E_{\Omega}\bigr\}\;.
    \end{aligned}
\end{equation}
Here $\mathcal I_{\xi}(U)$ is the nonnegative factor multiplying the group Dirac factors in CRE (1.11), including $\chi_1$, $\zeta$, $\chi_0'$ and the gauge-fixing and action exponential, and $t_{\xi}$ is defined where $UU_1(\xi)^{-1}$ lies in the chart domain bond by bond on $E_{\Omega}$, as the chart logarithm of CRE (1.13). Thus $\mathcal U_{\xi}\supseteq D^{(11)}(V,V_0)$, the fibre of CRE's full integration domain (Section~\ref{sec:branch-data}), the fibre defined by all factors of (1.11) and not by the starred group tests $\chi_0'$ alone, with equality at every outer datum satisfying the exterior constraint $V(c)= M(V_0,c)$ on the rows $c$ of $\Omega_1^c$ ($\bar V_0= V$, Section~\ref{sec:branch-data}), the one condition of $D^{(11)}(V,V_0)$ absent from $\mathcal U_{\xi}$ (the outer measure $\omega$ defined below is carried by that set; only the inclusion is used), when $\mathcal I_{\xi}(U)>0$ is read as including that the averages and means in (1.11) are defined (the domain convention recorded in (B1)); the fluctuations $UU_1^{-1}$ of $D^{(11)}(V,V_0)$ lie in the chart domain by (B1). The linearising change (1.14) is the map $A\mapsto A-h_{\xi}\widetilde D_{\xi}(A)$ on the open set $\mathcal N_{\Theta}$ of Lemma~\ref{lem:c2-fibre-smallness}(iv), which contains $g_0\mathsf C_{\xi}\mathcal A_{\xi}$, whose Jacobian is the determinant of (B4)(iii); restricted to the constraint surface, the range of $g_0\mathsf C_{\xi}$, it is $z\mapsto q_{\xi}(g_0\mathsf C_{\xi} z)$, and on $\mathcal A_{\xi}$ it is $q_{U_1}$ of \eqref{eq:native-source-local-phase}: $q_{\xi}(g_0\mathsf C_{\xi} B)= q_{U_1}(B)$ for $B\in\mathcal A_{\xi}$, with the row solutions of (B4)(i) and $\mathsf C_{\xi}=\mathsf C_{U_1}$ (Section~\ref{sec:branch-data}). The values of $q_{\xi}$ are needed on $g_0\mathsf C_{\xi}\mathcal A_{\xi}$ only, and its Jacobian needs $\widetilde D_{\xi}$ on a neighbourhood of that set in $\mathfrak g^{E_{\Omega}}$, the set $\mathcal N_{\Theta}$, on which the row solutions are defined and analytic ((B4)(i), Lemma~\ref{lem:c2-fibre-smallness}(iv)); data that are not branches have zero contribution.

The boundary change is $\Psi_{\xi}$ of \eqref{eq:native-source-exact-boundary-map}, CRE (1.22) with the cutoff $g$ on the bonds of $(R_1^{\sim2}\cap S_1)^*$, the set defined at p.251, and its second form ($g$ replaced by $1$, the form of p.252) on the other bonds of $S_1^*$, a set differing from the one printed at p.252 as recorded in Section~\ref{sec:boundary-law-collar}, from new free variables to old free variables. The change of variables on these branches is Lemma~\ref{lem:boundary-bijection} (the (1.22) part of (B4)(ii)): $\Psi_{\xi}$ maps $\mathcal D_{\xi}$ of \eqref{eq:c-new-variable-domain} bijectively onto $\mathcal A_{\xi}$ of \eqref{eq:c-free-variable-domain}, and $\mathcal A_{\xi}\subseteq\widehat{\mathcal A}_{\xi}$, since for $B\in\mathcal A_{\xi}$ the field $U(B)$ lies in $D^{(11)}(V,V_0)\subseteq\mathcal U_{\xi}$ with $t_{\xi}(U(B))= q_{U_1}(B)$ chart-valued and $g_0^{-1}q_{U_1}(B)|_{E_{\Omega,*}}= B$. The young fibre over $A_{\mathrm e}$ is
\begin{equation}\label{eq:retained-change-domain}
    \mathcal D_{\xi}(A_{\mathrm e})\coloneqq
        \bigl\{ a\in\mathfrak g^D\, | \,(A_{\mathrm e},a)\in\mathcal D_{\xi}\bigr\}\;,
\end{equation}
the open box $\{\Theta_{\xi}=1\}$ on a branch (Lemma~\ref{lem:collar-coverage}(iv)). The real coordinate changes and their inverses are restricted to these domains; no injectivity on the whole Lie algebra is required. Their transformed indicators are retained in $\Theta_D^{\sigma}$ and the exterior factors (1.25). Values outside the admitted domains contribute zero.

Let $v_{\xi},V_{\xi}$ be the functions in CRE (1.15)--(1.16): $v_{\xi}$ contains the Haar logarithms and the logarithm of the Jacobian of the linearising substitution (1.14) (CRE p.249 says that $v$ is the sum of the first two terms under the second exponential ``in (1.15)''; as printed this points to (1.15) itself, whose second exponential begins with the quadratic form and $v$; the display meant is (1.14), whose second exponential begins with $\log\sigma(A-h\widetilde D(A))+\operatorname{Tr}\log(I-h(\frac{\delta}{\delta A}\widetilde D)(A))$, and $v$ is read as the normalised Haar logarithm plus $\log|\det(\id-h_{\xi}\frac{\delta}{\delta A}\widetilde D_{\xi})|$, Section~\ref{sec:branch-data} and the paragraph following \eqref{eq:retained-positive-kernel}), and $V_{\xi}$ contains the remaining Taylor residual of the action and gauge fixing. These are the full functions, before the boundary change. For $z$ in their real domain put
\begin{equation}\label{eq:retained-prechange-density}
    p_{\xi}(z)=\exp\left[
        -\frac{1}{2}\langle z,\mathsf H_{\xi} z\rangle
        +v_{\xi}(g_0\mathsf C_{\xi} z)
        -g_0^{-2}V_{\xi}(g_0\mathsf C_{\xi} z)\right]\;.
\end{equation}
With $\Theta_D^{\sigma}(a)=\prod_{b\in D}\mathbf1_{\{| a(b)|_{\sigma}< r\}}$, in CRE's bond norm of (1.18), written $|\bigcdot |_{\sigma}$ as in \eqref{eq:literal-cutoff} with the fixed $\sigma$ of this section, which is the norm of the young cutoff \eqref{eq:native-source-young-cutoff} (so at $D= I$, $\Theta_D^{\sigma}=\Theta_{\xi}$ on $\mathfrak g^I$, the young fibre \eqref{eq:retained-change-domain} is the box $\{\Theta_D^{\sigma}=1\}$ on a branch, and $\mathbf 1_{\mathcal D_{\xi}(A_{\mathrm e})}\Theta_D^{\sigma}=\mathbf 1_{\mathcal D_{\xi}(A_{\mathrm e})}$ there), the conditional density and integral are
\begin{equation}\label{eq:retained-positive-kernel}
    \begin{aligned}
        \mathcal K_{\xi}(a)&=
            \mathbf1_{\mathcal D_{\xi}(A_{\mathrm e})}(a)\Theta_D^{\sigma}(a)
            p_{\xi}(\Psi_{\xi}(z_{\xi}(a)))
            |\det D\Psi_{\xi}(z_{\xi}(a))|\;,\\
        \mathcal J_{\xi}[f]&=\int
            f(U_{\xi}(a))\mathcal K_{\xi}(a)\,\dd a\;.
    \end{aligned}
\end{equation}
Here $\mathcal J_{\xi}$ and $\mathcal K_{\xi}$ are the $\mathcal J^{\mathrm a}_{\xi}$ and $\mathcal K_{\xi}$ of the measurability paragraph of Section~\ref{sec:branch-data}, written in the variable $a$ of the young fibre. The determinant inside $v_{\xi}$, the Jacobian of CRE's (1.14) (not $|\det D\Psi_{\xi}|$, the Jacobian of the boundary change in \eqref{eq:retained-positive-kernel}), is that of the map $A\mapsto A-h_{\xi}\widetilde D_{\xi}(A)$ on $\mathfrak g^{E_{\Omega}}$, the phases on all bonds of $\Omega_1$; it equals the determinant of its designated block, the product over the rows of the determinants in \eqref{eq:native-source-jacobian-primitive} (proof of Lemma~\ref{lem:primitive-sum}: $h_{\xi}$ has outputs on the designated bonds only, so the free rows of $\id-h_{\xi}\frac{\delta}{\delta A}\widetilde D_{\xi}$ are identity rows, $b_0(c)$ is the only designated bond of $\mathcal E_c$, Lemma~\ref{lem:two-block}, and $\widetilde D_{\xi,c}$ depends on $A|_{\mathcal E_c}$ only, Lemma~\ref{lem:row-locality}). Every change of variables uses its absolute real Jacobian: $v_{\xi}$ contains $\log|\det(\id-h_{\xi}\frac{\delta}{\delta A}\widetilde D_{\xi})|$, the sum over the rows of \eqref{eq:native-source-jacobian-primitive}, and the logarithms of the positive normalised Haar density, so $p_{\xi}>0$ where it is defined and $\mathcal K_{\xi}\geqslant 0$; the determinant is positive on $\mathcal A_{\xi}$ ((B4)(iii), Lemma~\ref{lem:c2-fibre-smallness}(iv)) and the real logarithm CRE writes agrees with this; the identities of this section use the absolute Jacobian and need only the nonvanishing. The definition retains $p_{\xi}\circ\Psi_{\xi}$ without a truncation of the boundary terms described after CRE (1.24).

CRE's exterior expression is the sum over region labels and integral over $(V_0,A_{\mathrm e})$, with the constraint $\delta(\bar V_0V^{-1})$, the retained region indicators, $\zeta(\Omega_1^c)$, Haar factors and the prefactors of (1.15), including $z^{(0)}$. Its exterior factors are evaluated at their transformed arguments. After splitting the good tests into their $D$ and exterior bond factors, the factors remaining from (1.25) are independent of $a$: Lemma~\ref{lem:collar-coverage}(ii), under the standing conditions \eqref{eq:c-standing-conditions} (its collar hypothesis is their condition $LMR_1>4R_{\mathrm c}+4L^2M_2R_1+8$; the collar \eqref{eq:regional-collar-sizes} is that of Proposition~\ref{prop:regional-square-extraction}), separates the bad and changed tests from $D$, and factors of good tests on $D$ are redundant by idempotence (CRE's own separation in (1.26), (1.29) is not used). Those exterior factors, including any dependence on $A_{\mathrm e}$, remain in $\mathcal O_V$. With \eqref{eq:retained-positive-kernel}, no boundary interaction or normaliser is discarded.

The maps in these integrals are jointly Borel in their parameters and integration variables, by the measurability paragraph of Section~\ref{sec:branch-data}: the background $U_1$ is Borel on the Borel set of admissible pairs (CRE's uniqueness statement for the critical point in the domain of integration in (1.11), (B1), and the theorems for Borel graphs with singleton sections); pairs that are not admissible have $D^{(11)}(V,V_0)=\varnothing$, hence zero source weight. The configuration of CRE (1.21) is the configuration described at p.251, on its prescribed determining data, the notation $M^{\bigcdot }(W)$ being the level-wise collection $M^j(W)$ of CRE (2.11), at the first step $W$ and $M(W)$; the measurable choices of $U_{1,\square'}(V)$ and of $U_1(\Omega_1\cap\Lambda_1^c,V)$ are the Borel choices of Lemma~\ref{lem:borel-minimiser}, fixed in (B1). The value of $\mathbf H_{1,Ax}$ used in (1.22) is $L\frac1i\log(U_1U_1(\Omega_1\cap\Lambda_1^c,V)^{-1})$ by (1.21), the chart logarithm of a relative configuration lying in the chart domain on $S_1^*$ (CRE's bound p.251), and is Borel in $(V,V_0)$. The maps $Q_1^{s*}$ in CRE (1.3) and the averages on these regular data are Borel as well.

The matrix coefficients of $\mathsf C_{\xi},h_{\xi}$ are obtained by inversion of the invertible designated coefficient operators $\alpha_c(U_1)$, so they are Borel in the background. The row solutions $\widetilde D_{\xi}$ of Section~\ref{sec:source-phase-primitives} are Borel on $\mathcal A_{\xi}$ jointly with the branch data, by the Lusin--Novikov argument of the measurability paragraph of Section~\ref{sec:branch-data} (the set $\mathsf P$ of pairs $(B,d)$ with $d$ a chart-class solution family at $g_0\mathsf C_{U_1}B$ for which $U_d(B)\in D^{(11)}(V,V_0)$ with the cutoff (1.25) nonzero is Borel with sections of at most one point, Lemma~\ref{lem:c2-fibre-smallness}(iii)), and so is $\frac{\delta}{\delta A}\widetilde D_{\xi}$ (measurability paragraph). In CRE (1.22), the cutoff $g$ is the fixed $C^{\infty}$ function of the paragraph defining \eqref{eq:native-source-exact-boundary-map}, each $\psi_b$ is $C^{\infty}$ on its open domain $\mathcal W_b$ (the chart logarithm and the exponential are analytic there), and the boundary data have just been shown Borel. Hence $\Psi_{\xi}(z)$ and its Jacobian are jointly Borel, and $C^{\infty}$ in $z$. The sets $\mathcal A_{\xi}$ and $\mathcal D_{\xi}$, and the young fibres \eqref{eq:retained-change-domain}, sections of the Borel set $\mathcal D_{\xi}$, are Borel (measurability paragraph). The Hessian coefficients, $v_{\xi}$, $V_{\xi}$ and the constraint factor $z^{(0)}$ are Borel: each is a finite sum or product of the Haar logarithm, the plaquette and gauge terms of the action and gauge fixing, the row determinants and the coefficients of $\alpha_c(U_1)^{-1}$, evaluated at the background, the phase and the row solutions, that is a continuous function of $(U_1,A,\widetilde D_{\xi}(A),\frac{\delta}{\delta A}\widetilde D_{\xi}(A))$ and of the Taylor coefficients at $A=0$ of the germs of Lemma~\ref{lem:row-locality}, themselves continuous in $U_1$ by the implicit function theorem, composed with maps that are Borel jointly in the branch data and $a$ (the background, the row solutions and their derivative, measurability paragraph of Section~\ref{sec:branch-data}). The vacuum exponential is a constant, and $\zeta(\Omega_1^c)$ is a finite product of small-field tests, indicators of sets defined by inequalities between continuous functions of $(V,V_0)$ on the open set where the local means entering them are defined, and of gauge-fixing exponentials of continuous functions of $(V,V_0)$ on that set, every factor being set to zero where a local mean it contains is undefined (the convention of this paper for $\zeta(\Omega_1^c)$, matching the chart convention of Section~\ref{sec:sourced-effective-interaction} for the tests); so $\zeta(\Omega_1^c)$ is Borel as well. Together with the indicators, these facts prove joint Borel measurability of $U_{\xi}(a)$, $\mathcal K_{\xi}(a)$ and the outer weights. For an empty carrier, its Lie-algebra space is a point with unit measure; empty sums equal zero, and empty products and zero-dimensional determinants equal one. If $\Omega_1$ is empty, set $U_1= U_{\xi}= V_0$ on every bond.

The exterior group Dirac factors have a kernel interpretation as well. Let $\omega$ be the joint outer measure in CRE (1.15), (1.25), retaining $V$, the region labels, $V_0$ and $A_{\mathrm e}$: in the notation of \eqref{eq:native-source-explicit-outer-operation}, $\int F\,\dd\omega=\int\mathcal O[F](V)\,\dd V$ for nonnegative Borel $F$ of the outer data, in the sense fixed in Section~\ref{sec:branch-data} (the integral of $F$ against the measure in $V$ that the outer operation defines; $\mathcal O[F](V)$ is not a density). It contains exactly the exterior factors specified above. Read its group Dirac factors by first integrating them against the corresponding coarse bond variables, with Haar measure on the remaining coarse variables. This defines a positive joint measure, rather than a pointwise density on a fibre. Its weighted measure $\bar\eta=\mathcal J_{\xi}[1]\omega$ is finite and its $V$ marginal is the measure $\rho_1$ of \eqref{eq:first-step}: the operations in CRE (1.1)--(1.25), performed as listed after \eqref{eq:native-source-physical-ratio} with a bounded Borel test $h(V)$ retained (signed; $h(V)= h(\mathsf M(U))$ is invariant under the gauge transformations $u$ of (1.5), which are the identity at the block centres (CRE's printed ``$z\neq y$'' in (1.5) read as $x\neq y$, as in (1.7)), Lemma~\ref{lem:global-average}(ii), so the Faddeev--Popov step applies with it inserted; every integral below is finite: $\int\mathcal J_{\xi}[1]\,\dd\omega=\mathcal Z_0<\infty$ is the case $h=1$ of the first identity below, both sides being the integral of one nonnegative Borel integrand over all variables by Tonelli, and $|\mathcal J_{\xi}[W_{\chi}]|\leqslant \mathcal J_{\xi}[1]$) and with insertion $1$ or $W_{\chi}$, give
\begin{equation}\label{eq:retained-outer-measure-identities}
    \begin{aligned}
        \int h(V)\,\dd\rho_1(V)
            &=\int h(V)\mathcal J_{\xi}[1]\,\dd\omega(V,\xi)\;,\\
        \int h(V)\,\dd N_1(V)
            &=\int h(V)\mathcal J_{\xi}[W_{\chi}(C;\bigcdot )]\,\dd\omega(V,\xi)\;.
    \end{aligned}
\end{equation}
Taking $h=1$ in the first equality gives $\bar\eta$ the total mass $\mathcal Z_0<\infty$. Let $Q_V$ be the regular conditional probability of $\bar\eta$ on the outer data given $V$, defined for $\rho_1$-almost every $V$ (disintegration of the finite measure $\bar\eta$ with respect to its marginal $\rho_1$ along the coordinate projection $\xi\mapsto V$: the measure disintegration theorem for standard Borel spaces, \cite[(17.35)]{Kechris}, stated there as an exercise with the standard proof, through a Borel isomorphism $\phi$ \cite[Theorem 15.6]{Kechris} of the outer space with $[0,1]$ and the Radon--Nikodym derivatives with respect to $\rho_1$ of the countably many measures $B\mapsto\bar\eta(\{\xi:V(\xi)\in B,\ \phi(\xi)\leqslant t\})$, $t$ rational, regularised to a Borel map $V\mapsto Q_V$ into the probability measures; the outer space, a finite union over the region labels of products of the compact groups $G^{E_1}$, $G^{E\setminus E_{\Omega}}$ and the Euclidean space $\mathfrak g^J$, is Polish, hence standard Borel, and on the branches $\bar\eta$ is carried by a bounded set, $g_0A_{\mathrm e}(b)$ lying in the chart ball for every $b\in J$, so the compact case would suffice. The theorem gives: $Q_V$ is a probability measure for every $V$ off a Borel $\rho_1$-null set, $V\mapsto Q_V(E)$ is Borel for every Borel $E$, $\bar\eta(E)=\int Q_V(E)\,\dd\rho_1(V)$, $Q_V$ is concentrated on the fibre $\{\xi:V(\xi)= V\}$ for $\rho_1$-almost every $V$, and two such maps agree for $\rho_1$-almost every $V$; the concentration is what replaces $h(V(\xi))$ by $h(V)$ in the identity displayed below, and the uniqueness makes $\mathcal O_V$ below independent of the version off a $\rho_1$-null set). For a nonnegative measurable outer function $F$, the exterior operator is
\begin{equation}\label{eq:retained-outer-kernel}
    \mathcal O_V[F]=
        \int_{\{0<\mathcal J_{\xi}[1]<\infty\}}
            \frac{F(\xi)}{\mathcal J_{\xi}[1]}\,\dd Q_V(\xi)\;.
\end{equation}
This defines a positive kernel for $\rho_1$-almost every $V$. It gives $\omega$ back on finite-positive-mass outer data: $\int h(V)\mathcal O_V[F]\,\dd\rho_1(V)=\int h(V)F(\xi)\mathbf 1_{\{0<\mathcal J_{\xi}[1]<\infty\}}\,\dd\omega(V,\xi)$ for nonnegative Borel $F$ and bounded Borel $h\geqslant 0$, since $\dd\bar\eta=\mathcal J_{\xi}[1]\,\dd\omega$, and for signed bounded Borel $h$ whenever $\int F\mathbf 1_{\{0<\mathcal J_{\xi}[1]<\infty\}}\,\dd\omega<\infty$; so every exterior weight is retained. Equation~\eqref{eq:retained-outer-kernel} is the definition of $\mathcal O_V$ throughout this section; the identity just displayed is the sense in which it represents CRE's exterior expression: integrated against bounded Borel $h(V)$, it reproduces the $\omega$-integral of every outer function supported on the branches, and no pointwise identity of kernels is asserted. Set $\mathcal O_V=0$ on the Borel $\rho_1$-null set where $Q_V$ is not defined, so that $V\mapsto\mathcal O_V[F]$ is Borel. Extend it linearly to complex $F$ whenever $\mathcal O_V[| F|]<\infty$. In particular $\mathcal O_V[\mathcal J_{\xi}[1]]= Q_V(\{0<\mathcal J_{\xi}[1]<\infty\})=1$ for $\rho_1$-almost every $V$, by the next paragraph. The changes of variables are those of the unrestricted CRE construction, recorded after \eqref{eq:native-source-physical-ratio}.

Only outer data with $0<\mathcal J_{\xi}[1]<\infty$ need a conditional quotient; these are exactly the branches in the sense of Section~\ref{sec:branch-data}: $\mathcal J_{\xi}[1]=0$ when the young fibre $\mathcal D_{\xi}(A_{\mathrm e})$ is empty, that is off the branches of that section, while on a branch $\mathcal J_{\xi}[1]<\infty$ (measurability paragraph) and $\mathcal J_{\xi}[1]>0$, since $\mathcal K_{\xi}>0$ on the nonempty open box $\mathcal D_{\xi}(A_{\mathrm e})$ ($p_{\xi}>0$, $|\det D\Psi_{\xi}|>0$ by Lemma~\ref{lem:boundary-bijection}, and $\Theta_D^{\sigma}=1$ there). Zero-mass data are $\bar\eta$-null, hence $Q_V$-null for $\rho_1$-almost every $V$. Data with infinite conditional mass are $\omega$-null, since $\int\mathcal J_{\xi}[1]\,\dd\omega=\mathcal Z_0<\infty$ (the case $h=1$ of \eqref{eq:retained-outer-measure-identities}, with $\mathcal Z_0=\rho_1(G^{E_1})$ by \eqref{eq:one-step-physical-partition}), hence $\bar\eta$-null and $Q_V$-null for $\rho_1$-almost every $V$. So $Q_V$ is carried by the branches for $\rho_1$-almost every $V$. All branch contractions below are set to zero off the branches. This convention loses no physical mass and asserts no pointwise finiteness at exceptional outer data.

Coverage uses the fixed geometric stencils $\Lambda_b$ of Appendix~\ref{sec:perimeter-attempt}: $\Lambda_b=\{ b\}$ for a free bond, and $\Lambda_{b_0(c)}=\mathcal S_c\cap E_*$ for a designated bond, where $\mathcal S_c$ contains the bonds with both endpoints in the two endpoint blocks of $c$. Let $A_{\mathrm{cov}}$ consist of the branches on which every fine contour bond is a bond of $\Omega_1$ and $\Lambda_b\subseteq D$ for every such bond. The geometric condition depends only on region labels, and is precisely the coverage used in Proposition~\ref{prop:regional-square-extraction}; at the carrier $D= I$ of Section~\ref{sec:c4c-1} it is the covered event of Lemma~\ref{lem:covered} and of Theorem~\ref{thm:main-theorem}, since $\Lambda_b=\mathcal F_b$ for every fine bond ($\mathcal S_c\cap E_*=\mathcal E_c\cap E_*$, both being the free positive bonds with both endpoints in $\mathcal B_{c_-}\cup\mathcal B_{c_+}$, \eqref{eq:regular-stencil-bonds} and \eqref{eq:regular-stencil-extension-support}). The period and coupling conditions are fixed throughout this section, and \eqref{eq:regional-collar-sizes} supplies the collar condition for every compatible set of region labels. Under the hypotheses of Proposition~\ref{prop:regional-square-extraction} fixed at the start of this section (its period, collar and coupling hypotheses, with \eqref{eq:regional-collar-sizes} supplying the collar condition at the carrier $D= I(\xi)$ of every branch), that proposition consequently defines $E^{\circ}_{\xi}$ on $A_{\mathrm{cov}}$.

\subsection{The Closed Flux Insertions}

Let $\bar C=(c_1,\ldots,c_m)$ be the oriented coarse square, with $m=4n$, and let $C$ be the concatenation of its straight fine paths. Fix a covered conditional history $\xi$ and write its constrained background as $U_1= U_1(\xi)$. For each positive coarse bond traversed by $\bar C$, the first and last fine bonds of its straight path lie inside its initial and final blocks, respectively, since $L\geqslant 3$. Coverage puts these bonds in $\Omega_1$. As $\Omega_1$ is a union of blocks, both endpoint blocks lie there. In particular $c$ is a bond of the coarse region $\Omega_1^{(1)}$, so \cite[(1.12)]{BalabanCRE} imposes $M(U_1,c)= V(c)$. Here a bond of a region has at least one endpoint in that region, as in \cite[Introduction]{BalabanVAR}; both endpoints belong in this case. Reversed traversals use the constraint on the underlying positive bond. All relative logarithms are in the regular chart used in \eqref{eq:nonlinear-average}.

For a positive coarse bond $c$, define $Y_c$ by the first formula below. The averaging constraint and \eqref{eq:nonlinear-average} then give the second:
\begin{equation}\label{eq:retained-flux}
    Y_c(\xi)= e^{-i\beta_c(U_1)},
    \qquad P_c(U_1)= Y_c(\xi)V(c)\;.
\end{equation}
Thus $Y_c$ depends on the closed relative contours inside the two endpoint blocks of $c$. Set $V(c^{-1})= V(c)^{-1}$. For a reversed traversal use
\begin{equation}\label{eq:retained-flux-reversal}
    Y_{c^{-1}}(\xi)= V(c)^{-1}Y_c(\xi)^{-1}V(c),
    \qquad P_{c^{-1}}(U_1)= Y_{c^{-1}}(\xi)V(c)^{-1}\;.
\end{equation}
The second identity follows by inverting $P_c= Y_cV(c)$. The factor is therefore based at the initial vertex of the reversed traversal as well.

Define the closed contraction
\begin{equation}\label{eq:retained-loop-insertion}
    \mathcal L_{\xi}(C;V)
        =\chi\left(Y_{c_1}(\xi)V(c_1)
            \cdots Y_{c_m}(\xi)V(c_m)\right)
        = W_{\chi}(C;U_1(\xi))\;.
\end{equation}
This equality follows by substituting \eqref{eq:retained-flux} and \eqref{eq:retained-flux-reversal} in the ordered fine holonomy. Under a site gauge $u$, the closed logarithms in \eqref{eq:nonlinear-average} and their exponentials are conjugated at $c_-$. Consequently
\begin{equation}
    Y_{c_i}\longmapsto u(c_{i,-})Y_{c_i}u(c_{i,-})^{-1},
    \qquad V(c_i)\longmapsto
        u(c_{i,-})V(c_i)u(c_{i,+})^{-1}\;.
\end{equation}
The intermediate gauge factors cancel in the ordered product and its character is invariant. All subsequent insertions in this section are closed contractions. No separate activity norm is assigned to a factor $Y_c$ or to an open transport. The dependence on $\xi$ is retained: different histories at the same $V$ can have different constrained backgrounds and different insertions.

\subsection{The Numerator Carried to the Next Transformation}

The exterior operator $\mathcal O_V$ of \eqref{eq:retained-outer-kernel} and the conditional integral $\mathcal J_{\xi}$ defined before it give, for $\rho_1$-almost every $V$,
\begin{equation}\label{eq:retained-conditional-law}
    \mathcal O_V[\mathcal J_{\xi}[1]]=1,
    \qquad W_1(C;V)=\mathcal O_V[\mathcal J_{\xi}[W_{\chi}(C;\bigcdot )]]\;,
\end{equation}
that is, $\dd N_1=\mathcal O_V[\mathcal J_{\xi}[W_{\chi}(C;\bigcdot )]]\,\dd\rho_1$.

The identities in \eqref{eq:retained-conditional-law} are identities of measures on $G^{E_1}$, tested against bounded Borel functions. Let $\mathfrak M$ be the image of $\rho_0(U)\,\dd U$ under $U\mapsto(M(U),U)$, a finite positive measure on $G^{E_1}\times G^E$ (the blocking kernel (0.1) with bare weight $\rho_0\,\dd U$, in the reading of \eqref{eq:bare-blocking-operator}). Its $V$ marginal is $\rho_1$. Disintegration on these standard Borel spaces supplies conditional probabilities $\nu_V$ on fine fields, for $\rho_1$-almost every $V$. For bounded Borel functions $f$ on $G^E$ and $h$ on $G^{E_1}$,
\begin{equation}\label{eq:retained-disintegration}
    \int h(V)f(U)\,\dd\mathfrak M(V,U)
        =\int h(M(U))f(U)\rho_0(U)\,\dd U
        =\int h(V)
            \left(\int f(U)\,\dd\nu_V(U)\right)\,\dd\rho_1(V)\;.
\end{equation}
In particular, with $f= W_{\chi}(C;\bigcdot )$ and \eqref{eq:first-step}, $\int h\,\dd N_1=\int h(V)\int W_{\chi}(C;U)\,\dd\nu_V(U)\,\dd\rho_1(V)$, so $W_1(C;V)=\int W_{\chi}(C;U)\,\dd\nu_V(U)$ for $\rho_1$-almost every $V$, the conditional form of \eqref{eq:outgoing-loop}; the operations below act on the middle member of \eqref{eq:retained-disintegration}, and $\nu_V$ is not used further. The averaging operation is \eqref{eq:global-average}, Borel and gauge covariant by Lemma~\ref{lem:global-average}; \eqref{eq:nonlinear-average} is its formula on $\mathcal O_c$, which contains every configuration of $D^{(11)}$ on the rows of $\Omega_1$ (the domain convention recorded in (B1)).

For $f=1$ or $W_{\chi}(C;\bigcdot )$, apply the operations of CRE (1.1)--(1.25) listed after \eqref{eq:native-source-physical-ratio} with $h(V)$ retained. The partitions are pointwise decompositions of unity. The Faddeev--Popov step (1.5)--(1.6) is applied to the operator (0.1) with the averaging choice \eqref{eq:global-average}, whose covariance on all of $G^E$ is Lemma~\ref{lem:global-average}(ii); the closed Wilson character ($| W_{\chi}|\leqslant 1$, $\chi$ being a normalised irreducible character) is invariant under its fine gauges. The changes of variables pull that character back to $W_{\chi}(C;U_{\xi}(a))$, and $h(V)$ stays fixed: the gauge transformations of (1.5) are the identity at the block centres, so $M(U)$, hence $h(M(U))$, is unchanged by them (Lemma~\ref{lem:global-average}(ii)), and the later changes of variables are at fixed $V$. Cutoff factors are removed only where CRE states they equal one, and the boundary change is retained as \eqref{eq:retained-positive-kernel}. Thus the operations supply \eqref{eq:retained-outer-measure-identities} against every bounded Borel $h$; since $|\mathcal J_{\xi}[W_{\chi}]|\leqslant \mathcal J_{\xi}[1]$ (nonnegativity of \eqref{eq:retained-positive-kernel}), $\mathcal J_{\xi}[W_{\chi}]$ vanishes where $\mathcal J_{\xi}[1]=0$, and $\{\mathcal J_{\xi}[1]=\infty\}$ is $\omega$-null, so both $\omega$-integrals are integrals over the branches, and with \eqref{eq:retained-outer-kernel} they give $\int h\,\dd\rho_1=\int h(V)\mathcal O_V[\mathcal J_{\xi}[1]]\,\dd\rho_1(V)$ and $\int h\,\dd N_1=\int h(V)\mathcal O_V[\mathcal J_{\xi}[W_{\chi}(C;\bigcdot )]]\,\dd\rho_1(V)$. Equality of these finite measures gives the $\rho_1$-almost-everywhere identities \eqref{eq:retained-conditional-law}, the second by the uniqueness of the Radon--Nikodym derivative \eqref{eq:outgoing-loop}.

Nonnegativity of \eqref{eq:retained-positive-kernel} gives $|\mathcal J_{\xi}[W_{\chi}]|\leqslant \mathcal J_{\xi}[1]$, hence $|\mathcal O_V[\mathcal J_{\xi}[W_{\chi}]]|\leqslant \mathcal O_V[\mathcal J_{\xi}[1]]=1$, the bound $| W_1|\leqslant 1$ of Section~1.

For the covered branches $A_{\mathrm{cov}}$ defined above, set
\begin{equation}\label{eq:retained-leading-source}
    \mathcal L_{\xi}^{(0)}(C;V)=
    \begin{cases}
        \kappa_C e^{E^{\circ}_{\xi}}\mathcal L_{\xi}(C;V),
            &\xi\in A_{\mathrm{cov}},\\
        0,&\xi\notin A_{\mathrm{cov}}.
    \end{cases}
\end{equation}
Proposition~\ref{prop:c4c-5-1} (Appendix~\ref{app:regional-scalar}) identifies its coefficient with $S^{\circ}_{\xi,C}$ on covered histories, with $E^{\circ}_{\xi}= E^{\circ}_{I,C_n,T}$ at the carrier $D= I= D_+$ (all free bonds with an endpoint in $\Lambda_1$, the convention of Section~\ref{sec:c4c-1}). Its period and coupling hypotheses are fixed for the square; its boundary-distance terms allow arbitrary covered regional labels. The error $E^{\circ}_{\xi}$ depends only on those discrete labels; the background and branch indicator are Borel by the preceding construction, so this leading term is measurable. The other source terms are retained as the following conditional integral, with the same denominator:
\begin{equation}\label{eq:retained-source-correction}
    \mathcal H_{\xi}(C;V)=
        \frac{1}{\mathcal J_{\xi}[1]}
        \int\left[W_{\chi}(C;U_{\xi}(a))-\mathcal L_{\xi}^{(0)}(C;V)\right]
            \mathcal K_{\xi}(a)\,\dd a\;.
\end{equation}
This definition retains the difference between the conditional Wilson integral and $\mathcal L_{\xi}^{(0)}$, with the boundary dependence of \eqref{eq:retained-positive-kernel} unchanged. On an uncovered branch, $\mathcal L_{\xi}^{(0)}=0$ and it is the normalised conditional Wilson integral. No bound on $\mathcal H_{\xi}$ beyond the trivial $|\mathcal H_{\xi}|\leqslant 1+B_C$ of the next paragraph is asserted here; its history average enters \eqref{eq:thm-main-identity} through the defect and boundary terms $E_{\mathrm{def}}$, $E_{\mathrm{bdry}}$ of Corollary~\ref{cor:c4c-5-2}, which Theorem~\ref{thm:main-theorem}(ii) bounds.

For this fixed square, \eqref{eq:regional-square-error} and $|\chi|\leqslant 1$ give $|\mathcal L_{\xi}^{(0)}|\leqslant B_C$, where $B_C=|\kappa_C|\exp(C\tau^2(9n+1))<\infty$ (the constant $C$ of \eqref{eq:regional-square-error}, not the contour). Consequently $\mathcal J_{\xi}[1]|\mathcal H_{\xi}|\leqslant (1+B_C)\mathcal J_{\xi}[1]$. The numerator terms below are therefore integrable for the complex extension of \eqref{eq:retained-outer-kernel} at $\rho_1$-almost every $V$.

Substitution in \eqref{eq:retained-conditional-law} gives the following first-step identity for $\rho_1$-almost every $V$:
\begin{equation}\label{eq:retained-first-step-numerator}
    W_1(C;V)=\mathcal O_V\left[
        \mathcal J_{\xi}[1]
        \bigl(\mathcal L_{\xi}^{(0)}(C;V)+\mathcal H_{\xi}(C;V)\bigr)\right]\;.
\end{equation}
More explicitly, let $\nu_{\xi}=(U_{\xi})_*(\mathcal K_{\xi}(a)\,\dd a/\mathcal J_{\xi}[1])$ be the probability measure on fine fields carried by the branch. Under a fine gauge transformation $u$ (acting on $V$ at the block centres, Lemma~\ref{lem:global-average}(ii)), use its pushforward $(U\mapsto U^u)_*\nu_{\xi}$ together with $U_1^u$ and $V^u$. Both closed characters in \eqref{eq:retained-source-correction} are invariant, $E^{\circ}_{\xi}$ depends only on the unchanged region labels, and change of variables preserves the normalised conditional integral. This is the gauge-invariant meaning of each branch contraction; an axial-gauge representative and its conditional measure are transported together. The outer operation retains the corresponding history weights. The physical denominator is still $\mathcal Z_0=\rho_1(G^{E_1})$.

For $\rho_1$-almost every $V$, the scalar-extracted outgoing observable is
\begin{equation}\label{eq:retained-outgoing-observable}
    \widehat W_1(C;V)=\frac{W_1(C;V)}{\kappa_C}
        =\frac{\mathcal O_V\left[\mathcal J_{\xi}[1]
            \bigl(\kappa_C^{-1}\mathcal L_{\xi}^{(0)}(C;V)
                +\kappa_C^{-1}\mathcal H_{\xi}(C;V)\bigr)\right]}
            {\mathcal O_V[\mathcal J_{\xi}[1]]}\;.
\end{equation}
The factor $\kappa_C$ is nonzero: \eqref{eq:marginal-square-mass} defines its finite perimeter exponent and \eqref{eq:corner-square-factor} writes $z_{\mathrm c}$ as the exponential of the convergent corner sum. The pair of the complex measure $\kappa_C^{-1}N_1$ and the measure $\rho_1$ is the input to the next blocking transformation ($T$ acting on measures by the pushforward $T\mu= M_*\mu$ of Section~1). Equation~\eqref{eq:retained-loop-insertion} displays its leading closed insertion; \eqref{eq:retained-source-correction} specifies the other source terms carried with it. Iterating this representation requires local estimates for their transformed activities, including correlations between attachments and their boundary dependence. The first step does not require replacing $\mathcal L_{\xi}(C;V)$ by $W_{\chi}(\bar C;V)$.

\section{The Insertion at the Retreated Carrier}
\label{sec:one-step-insertion}

This section bounds the terms of the branch identity \eqref{eq:c-one-step-branch-identity} against the flat coefficient $\kappa_C$ of Section~\ref{sec:remaining-physical-comparison}: the source-containing boundary terms relative to $|\kappa_C|$ (Corollary~\ref{cor:c5-7-4}(a)), and the labelled defect by \eqref{eq:c5-6-def-2}, whose far-history part is relative to $|\kappa_C|$ through $| S^{\circ}_{\Lambda}|=|\kappa_C|| e^{E^{\circ}_{\Lambda}}|$ and the exponential moment of Corollary~\ref{cor:c5-7-4}(c), while its near-history part is an absolute perimeter rate (Remark~\ref{rem:main-regime}(2)). It does not work at the carrier $I$ of Section~\ref{sec:sourced-effective-interaction} directly. A second cutoff test is placed on the collar of $\Lambda_1$; on the branches where it fails, $\Lambda_1$ and $I$ are retreated to a labelled carrier $D(\Lambda)\subseteq I$, the branch identity is written for each labelled term, and the labelled coefficient $S^{\circ}_{\Lambda}$ is identified with $\kappa_Ce^{E^{\circ}_{\Lambda}}$ at the retreated carrier by Proposition~\ref{prop:regional-square-extraction}. The bonds removed by the retreat are paid by the physical measure, through one-point witness probabilities of the kind of Appendix~\ref{app:exterior-chessboard}; no lower bound on the exterior Schur complement is used. The inputs are Sections~\ref{sec:sourced-effective-interaction}, \ref{sec:regional-identification} and~\ref{sec:remaining-physical-comparison}, and Appendices~\ref{app:background-estimates}, \ref{app:exterior-chessboard}, \ref{app:signed-defects} and~\ref{app:regional-scalar}. CRE is \cite{BalabanCRE}, with its equation and page numbers.

\subsection{Setting and inputs}
\label{sec:c5-0}

Interior (young) bonds $D= I=\Lambda_1^*$, the free bonds with an endpoint in $\Lambda_1$ (the young carrier of Section~\ref{sec:sourced-effective-interaction}; $\Lambda_1^*$ is the starred bond set of the conventions preceding \eqref{eq:c-standing-conditions}), exterior free bonds $J=\Omega_1^*\setminus D$, where $\Omega_1^*$ is the free carrier (the $E_{\Omega,*}$ of Section~\ref{sec:remaining-physical-comparison}), and $z= z_{\xi}(a)= i_Da+i_JA_{\mathrm e}$. The kernel of Section~\ref{sec:remaining-physical-comparison} is $K_{\xi}(a)=\mathbf 1_{\mathcal D_{\xi}(A_{\mathrm e})}(a)\Theta^{\sigma}_D(a)p_{\xi}(\Psi_{\xi}(z_{\xi}(a)))|\det D\Psi_{\xi}|$ with $p_{\xi}(z)=\exp[-\tfrac12\langle z,H_{\xi} z\rangle+v_{\xi}(g_0C_{\xi} z)-g_0^{-2}V_{\xi}(g_0C_{\xi} z)]$ of \eqref{eq:retained-prechange-density} (written there with $\mathsf H_{\xi}= H_{\xi}$, $\mathsf C_{\xi}= C_{\xi}$; $H_{\xi}$ is the self-adjoint operator of the quadratic form) and $\Theta^{\sigma}_D(a)=\prod_{b\in D}\mathbf 1\{| a_b|_{\sigma}< r\}$, $r=\delta_0/g_0$. $\Psi_{\xi}$ is the bondwise boundary change CRE (1.22), the identity off $S_1^*\subset\Omega_1\cap\Lambda_1^c$, and $D\cap S_1^*=\varnothing$ by \eqref{eq:c-collar-distance}, so it changes $A_{\mathrm e}$ only: write $A'_{\mathrm e}=\Psi_{\xi}(A_{\mathrm e})|_J$, so that $\Psi_{\xi}(z)= i_Da+i_JA'_{\mathrm e}$. Then, with $H_{JD}= H_{DJ}^{\mathsf T}$,
\begin{equation}\label{eq:c5-0-1}
    -\tfrac12\langle\Psi_{\xi}(z),H_{\xi}\Psi_{\xi}(z)\rangle=-\tfrac12\langle a,H_{DD}a\rangle-\langle a,H_{DJ}A'_{\mathrm e}\rangle-\tfrac12\langle A'_{\mathrm e},H_{JJ}A'_{\mathrm e}\rangle\;,
\end{equation}
with $H_{DD}$, $H_{DJ}$, $H_{JJ}$ the blocks of $H_{\xi}$. All three terms are inside $K_{\xi}$ (not in the exterior operator $\mathcal O_V$ of \eqref{eq:retained-outer-kernel}), by \eqref{eq:retained-prechange-density}.

Inputs (Lemma~\ref{lem:regular-bounds}, stated in Section~\ref{sec:sourced-effective-interaction} and proved in Appendix~\ref{app:background-estimates}; Lemmas~\ref{lem:c2-entry-range}, \ref{lem:c2-regional-entry} and~\ref{lem:hessian-gauge-locality}); $\|\bigcdot \|$ on an entry is the operator norm for $h$:
\begin{enumerate}
\item[(I1)] entry range: $(H_{\xi})_{bb'}=0$ unless $| z(b)-z(b')|_1\leqslant 2$, hence unless the initial sites are within $\ell^{\infty}$ distance $\rho_L=4L$ in fine units; each row has at most $164L^4$ nonzero entries;
\item[(I2)] row sums: $\sum_{b'\in\Omega_1^*}\|(H_{\xi})_{bb'}\|\leqslant m_L$ for every row $b\in D$, the row-sum statement of Lemma~\ref{lem:regular-bounds} ($D= I$). Column sums are not used;
\item[(I3)] two-sided bound on every branch in the support of the outer weight: $(\gamma_L/2)\id\leqslant H_{\xi,DD}\leqslant m_L\id$, and the same for every principal restriction (Appendix~\ref{app:background-estimates} writes $I$ for the carrier $D$).
\end{enumerate}
Put $\gamma=\gamma_L/2$ and $h_+= m_L$. The constants $\gamma_L$, $m_L$, $\rho_L$ depend on $L$ and the group only (Section~\ref{app:thresholds}); the bounds (I2)--(I3) hold under the hypotheses of Lemma~\ref{lem:regular-bounds}, that is under \eqref{eq:c-standing-conditions} and the coupling thresholds of Appendix~\ref{app:background-estimates} (\eqref{eq:convex-background-thresholds}, Section~\ref{app:thresholds}). Norms: $|\bigcdot |_{\sigma}$, $\sigma\in\{\mathrm{op},\mathrm{HS}\}$, is the cutoff norm of \eqref{eq:literal-cutoff}, the operator norm or the normalised Hilbert--Schmidt norm $h$, fixed once (the conventions of Section~\ref{sec:sourced-effective-interaction}); $h\leqslant |\bigcdot |_{\sigma}\leqslant c_{\mathrm h}h$ with $c_{\mathrm h}=\sqrt N$. Section~\ref{sec:sourced-effective-interaction} and Appendix~\ref{app:background-estimates} write $I$ for the carrier that Section~\ref{sec:regional-identification} calls $D$; in this subsection $D= I$, while the retreated carriers $D(\Lambda)\subseteq I$ of Sections~\ref{sec:c5-2}--\ref{sec:c5-7} are its principal restrictions: $D'\subset D$ denotes the carrier of a retreated region in Section~\ref{sec:c5-3-1}, and (I3) passes to it by principal restriction while (I2), a bound on the rows of the one operator $H_{\xi}$, is unchanged. Bond convention (Section~\ref{sec:sourced-effective-interaction}): the bonds of a set $X$ are the positive fine bonds with at least one endpoint in $X$, and the bonds of $\Omega_1^c$ are those with both endpoints off $\Omega_1$. $\ell\coloneqq LMR_1$ (the macroscopic scale of Section~\ref{sec:c3-7}; the even-period restriction of that section is not imported by this notation).

\subsection{The mean shift}
\label{sec:c5-1}

Put $m=-H_{DD}^{-1}H_{DJ}A'_{\mathrm e}\in\mathfrak g^D$, so that $-\tfrac12\langle a,H_{DD}a\rangle-\langle a,H_{DJ}A'_{\mathrm e}\rangle=-\tfrac12\langle a-m,H_{DD}(a-m)\rangle+\tfrac12\langle m,H_{DD}m\rangle$.

\begin{lemma}[entrywise decay of the interior inverse]\label{lem:c5-1-1}
On a branch in the support of the outer weight, under the hypotheses of Lemma~\ref{lem:regular-bounds} (so that (I1)--(I3) of Section~\ref{sec:c5-0} hold), for $b,b'\in D$, with $d(b,b')$ the $\ell^{\infty}$ distance of the initial sites in fine units,
\begin{equation*}
    \| H_{DD}^{-1}(b,b')\|\leqslant \gamma^{-1}(1-\gamma/h_+)^{\lceil d(b,b')/\rho_L\rceil}\leqslant \gamma^{-1}e^{-\mu_Hd(b,b')}\;,\qquad \mu_H\coloneqq \frac{\gamma}{h_+\rho_L}=\frac{\gamma_L}{8Lm_L}\;.
\end{equation*}
\end{lemma}

\begin{proof}
By (I3), $0\leqslant K\coloneqq \id-h_+^{-1}H_{DD}\leqslant (1-\gamma/h_+)\id$, so $H_{DD}^{-1}= h_+^{-1}\sum_{n\geqslant 0}K^n$ (norm-convergent, $\| K\|\leqslant 1-\gamma/h_+<1$). $K$ has the entry range $\rho_L$ of (I1) (the identity is diagonal), so $K^n(b,b')=0$ when $d(b,b')> n\rho_L$. Hence $H_{DD}^{-1}(b,b')= h_+^{-1}\sum_{n\geqslant \lceil d/\rho_L\rceil}K^n(b,b')$ and $\| K^n(b,b')\|\leqslant \| K^n\|\leqslant (1-\gamma/h_+)^n$, since the inner product on $\mathfrak g^D$ is the bondwise sum of the inner product $\langle X,Y\rangle=\tr_N(XY)$ of $\mathfrak g$ (Appendix~\ref{app:flat-covariance}), whose norm is $h$, so that an entry of an operator on $\mathfrak g^D$ has $h$-operator norm at most the norm of the operator; the geometric sum gives $h_+^{-1}(1-\gamma/h_+)^{\lceil d/\rho_L\rceil}/(\gamma/h_+)$. Finally $(1-x)^s\leqslant e^{-xs}$ with $x=\gamma/h_+$, $s\geqslant d/\rho_L$.
\end{proof}

\begin{lemma}[sup-norm shift bound]\label{lem:c5-1-2}
Under the hypotheses of Lemma~\ref{lem:c5-1-1}, let $J_1\coloneqq \{ e\in J\, | \,H(b',e)\neq0\text{ for some }b'\in D\}$; by (I1), $J_1\subset\{ e\in J\, | \,\operatorname{dist}(e,D)\leqslant \rho_L\}$, where $\operatorname{dist}(e,D)=\min_{b\in D}d(e,b)$ is the $\ell^{\infty}$ distance of initial sites of (I1). Then $m$ depends on $A'_{\mathrm e}|_{J_1}$ only (exterior bonds with $\operatorname{dist}(e,D)>\rho_L$ have a zero column in $H_{DJ}$; $m=0$ if $J_1=\varnothing$), and for every $b\in D$
\begin{equation}\label{eq:c5-1-1}
    h(m_b)\leqslant C_H^{\mathrm{near}}\max_{e\in J_1}h(A'_{\mathrm e}(e))\;,\qquad C_H^{\mathrm{near}}\coloneqq \gamma^{-1}m_LN_{\mathrm{lat}}(\mu_H)\;,\qquad N_{\mathrm{lat}}(\mu)\coloneqq \sum_{s\geqslant 0}4(2s+1)^4e^{-\mu s}<\infty\;.
\end{equation}
\end{lemma}

\begin{proof}
$m_b=-\sum_{b'\in D}H_{DD}^{-1}(b,b')\sum_{e\in J_1}H(b',e)A'_{\mathrm e}(e)$. Operator norms of entries are those induced by $h$ (the convention of Appendix~\ref{app:background-estimates}), so $h(H(b',e)A'_{\mathrm e}(e))\leqslant \| H(b',e)\| h(A'_{\mathrm e}(e))$. By Lemma~\ref{lem:c5-1-1} and the row-sum bound (I2) for the rows $b'\in D$,
\begin{equation*}
    h(m_b)\leqslant \sum_{b'\in D}\gamma^{-1}e^{-\mu_Hd(b,b')}\sum_{e\in J_1}\| H(b',e)\| h(A'_{\mathrm e}(e))\leqslant \gamma^{-1}m_L\max_{e\in J_1}h(A'_{\mathrm e}(e))\sum_{b'\in D}e^{-\mu_Hd(b,b')}\;,
\end{equation*}
and $\sum_{b'}e^{-\mu_Hd(b,b')}\leqslant \sum_{s\geqslant 0}4(2s+1)^4e^{-\mu_Hs}$: at most $(2s+1)^4$ sites lie at $\ell^{\infty}$ distance $s$ from the initial site of $b$, each carries four positive bonds, and projection to the torus only identifies sites. Only row sums over $b'\in D$ are used; no column-sum bound for columns $e\in J$ is needed.
\end{proof}

\begin{remark}[norms and the choice of $\vartheta$]\label{rem:c5-1-3}
The cutoff norm $|\bigcdot |_{\sigma}$ is the operator norm or $h$ (Section~\ref{sec:c5-0}; the conventions of Section~\ref{sec:sourced-effective-interaction} and Appendix~\ref{app:background-estimates}), and $h\leqslant |\bigcdot |_{\sigma}\leqslant c_{\mathrm h}h$ with $c_{\mathrm h}=\sqrt N$. Hence $| m_b|_{\sigma}\leqslant c_{\mathrm h}h(m_b)\leqslant c_{\mathrm h}C_H^{\mathrm{near}}\max_{e\in J_1}h(A'_{\mathrm e}(e))\leqslant c_{\mathrm h}C_H^{\mathrm{near}}\max_{e\in J_1}| A'_{\mathrm e}(e)|_{\sigma}$. Put
\begin{equation*}
    \vartheta\coloneqq \frac1{4c_{\mathrm h}C_H^{\mathrm{near}}}\;,
\end{equation*}
a constant depending on $L$ and the group only (through $\gamma_L$, $m_L$, $\rho_L=4L$ and $N$); $\vartheta\leqslant 1/16$, since $c_{\mathrm h}\geqslant 1$, $N_{\mathrm{lat}}\geqslant 4$ (its $s=0$ term) and $m_L\geqslant \gamma$ by (I3). Consequently, if $| A'_{\mathrm e}(e)|_{\sigma}<\vartheta r$ for every $e\in J$ with $\operatorname{dist}(e,D)\leqslant \rho_L$ (a set containing $J_1$), then
\begin{equation}\label{eq:c5-1-2}
    \sup_{b\in D}| m_b|_{\sigma}\leqslant r/4\;.
\end{equation}
No coupling threshold beyond those under which (I2)--(I3) hold (Section~\ref{sec:c5-0}), and no condition on the exterior bonds farther than $\rho_L$ from $D$, is involved: those bonds do not enter $m$ (Lemma~\ref{lem:c5-1-2}). The decay of $H_{DD}^{-1}$ (Lemma~\ref{lem:c5-1-1}) is used only to sum the $b'$-series; it plays no role in decoupling far exterior bonds, which are decoupled by the entry range (I1) alone.
\end{remark}

\subsection{The exterior Gaussian in the history weight, and the payment of a bad collar bond}
\label{sec:c5-2}

Two things are needed from the exterior Gaussian: the identification of its place inside the history weight, which is the exact identity \eqref{eq:c5-2-1}, and a payment for a collar bond on which the fluctuation is large. A lower bound $S\geqslant \gamma$ for the Schur complement $S$ of \eqref{eq:c5-2-1}, which would pay such a bond by $e^{-c\gamma\vartheta^2r^2}$ ($c=1/2$ for the Hilbert--Schmidt cutoff, $1/(2N)$ for the operator-norm cutoff), is not supplied by the sources (Remark~\ref{rem:c5-2-2}), so the payment is taken from the physical measure instead (Section~\ref{sec:c5-2-3}).

\subsubsection{Where the exterior Gaussian sits}
\label{sec:c5-2-1}
By \eqref{eq:c5-0-1} and completing the square (which needs only $H_{DD}$ invertible, (I3), and $|\det D\Psi_{\xi}|$ independent of $a$: $\log|\det D\Psi_{\xi}(A')|$ depends on $A'|_{S_1^*}$ only, Section~\ref{sec:boundary-law-collar}, \eqref{eq:native-source-exact-boundary-map}),
\begin{equation}\label{eq:c5-2-1}
    K_{\xi}(a)= e^{-\frac12\langle A'_{\mathrm e},SA'_{\mathrm e}\rangle}\cdot\mathbf 1_{\mathcal D_{\xi}(A_{\mathrm e})}(a)\Theta^{\sigma}_D(a)e^{-\frac12\langle a-m,H_{DD}(a-m)\rangle+v_{\xi}-g_0^{-2}V_{\xi}}|\det D\Psi_{\xi}|\;,\qquad S\coloneqq H_{JJ}-H_{JD}H_{DD}^{-1}H_{DJ}\;,
\end{equation}
using $\tfrac12\langle m,H_{DD}m\rangle-\tfrac12\langle A'_{\mathrm e},H_{JJ}A'_{\mathrm e}\rangle=-\tfrac12\langle A'_{\mathrm e},SA'_{\mathrm e}\rangle$ and $H_{JD}= H_{DJ}^{\mathsf T}$. The factor is written in $A'_{\mathrm e}=\Psi_{\xi}(A_{\mathrm e})|_J$, which differs from $A_{\mathrm e}$ on $S_1^*$ only. Off $S_1^*$ it is the quadratic form CRE collects in $T_1$, CRE (1.29); on $S_1^*$ the conventions differ (CRE moves first-order boundary corrections into $V^{(0)}(S_1)$, p.252, while \eqref{eq:retained-prechange-density} retains $p_{\xi}\circ\Psi_{\xi}$ without truncation), and nothing below uses this correspondence. In the convention of Section~\ref{sec:sourced-effective-interaction} ($\mathsf w_{\xi}= e^{c_{\xi}}I_{\xi;1}$, with $c_{\xi}$ of Lemma~\ref{lem:collar} carrying $\mathscr E_{\xi}(0)$, the exterior quadratic form, the far primitives and $\log|\det D\Psi_{\xi}|$, and $K_{\xi;1}$ carrying the $II$ and $IJ$ blocks and the near primitives) the same algebra gives $e^{c_{\xi}}I_{\xi;1}= e^{-\frac12\langle A'_{\mathrm e},SA'_{\mathrm e}\rangle}\times e^{c_{\xi}+\frac12\langle A'_{\mathrm e},H_{JJ}A'_{\mathrm e}\rangle}\times\int\mathbf 1_{\mathcal D}\Theta\,e^{-\frac12\langle a-m,H_{DD}(a-m)\rangle}\times(\text{near primitives and cutoffs})(a)\,\dd a$, the far primitives and $\mathscr E_{\xi}(0)$ staying in $c_{\xi}$. So in either convention the Schur factor is an exact factor of $\mathsf w_{\xi}$. It is independent of $a$, hence cancels in the conditional ratio $\mathcal J_{\xi}[f]/\mathcal J_{\xi}[1]$ inside $\mathcal O_V$ (\eqref{eq:retained-outer-kernel}) and acts only through the history weight.

\begin{remark}[no lower bound for $S$ on $\mathfrak g^J$ is used]\label{rem:c5-2-2}
$\langle x,Sx\rangle=\min_{y\in\mathfrak g^D}\langle(y,x),H_{\xi}(y,x)\rangle$ ($H_{DD}>0$ by (I3)), so $S\geqslant \gamma$ on $\mathfrak g^J$ is equivalent to $\langle(y,x),H_{\xi}(y,x)\rangle\geqslant \gamma| x|^2$ for all $(y,x)\in\mathfrak g^{D\cup J}$, a lower bound of $H_{\xi}$ in the $J$-directions on all of $\mathfrak g^{E_{\Omega,*}}$. The sources supply no such bound: (i) the lower bound of Lemma~\ref{lem:regular-bounds} (its proof, local comparison) covers $\operatorname{supp}B$ by patches within $2R_{\mathrm c}-1$ of $I$, where the plaquette bound of Lemma~\ref{lem:c2-plaquette-deviation} (stated for cubes $\square'\subset\Omega_1^{\sim-1}\setminus R_1$ with $\square'^{\sim4}\subset\Omega_1$) yields a gauge with phases $\leqslant \widehat\mu$ and the flat gap $\gamma_L$ applies after the perturbation $c_L\widehat\mu<\gamma_L/4$; (ii) $J$ contains the free bonds of $S_1=(\Omega_1\setminus\Omega_1^{\sim-1})\cup R_1$ (CRE p.251): on $R_1$ the CRE (1.18) fluctuation test fails, which says nothing about the background $U_1$ (the comparison $U_1\approx U_{1,\square'}$ holds there too, Proposition~\ref{prop:c3-9-1}), but the plaquette bound of Lemma~\ref{lem:c2-plaquette-deviation}, from which that gauge is built, is proved only for cubes of $\Omega_1^{\sim-1}\setminus R_1$; on $\Omega_1\setminus\Omega_1^{\sim-1}$ the tests $\chi_1(\Omega_1)$, $\chi_0'(\Omega_1)$ are in force (CRE (1.11), p.250) but the comparison $U_1\approx U_{1,\square'}$ is not available there ($\square'^{\sim4}\not\subset\Omega_1$), and no near-flat gauge is supplied; (iii) a bound $H_{\xi}\geqslant \gamma$ on the regular part $D\cup J_1$ alone gives only the compression $S_{J_1J_1}\geqslant \gamma$, because $S_{J_1,J\setminus J_1}= H_{J_1,J\setminus J_1}\neq0$ (entry range $\rho_L$). Hence the Gaussian factor \eqref{eq:c5-2-1} is not used for any bound; the exterior field on the large-field part is treated as CRE treats it, paid by the cutoff factors that put it there. No counterexample is claimed: what is established is that no proof of $S\geqslant \gamma$ on $\mathfrak g^J$ exists in the sources, and none is attempted. Nothing in Sections~\ref{sec:c5-2}--\ref{sec:c5-7} uses $S$.
\end{remark}

\subsubsection{The payment of a bad collar bond from the physical measure}
\label{sec:c5-2-3}
For a region $\Lambda$ (a union of $\ell$-cubes, $\Lambda\subset\Lambda_1$) the carrier is $D(\Lambda)\coloneqq \{ b\in I\, | \,b\text{ has an endpoint in }\Lambda\}\subset I$ and the exterior is $J(\Lambda)\coloneqq E_{\Omega,*}\setminus D(\Lambda)\supset J$; the collar is $J_{\mathrm{col}}(\Lambda)\coloneqq \{ e\in J(\Lambda)\, | \,\operatorname{dist}(e,D(\Lambda))\leqslant \rho_L\}$ ($\ell^{\infty}$ distance of initial sites); it contains $J_1(\Lambda)$ of Lemma~\ref{lem:c5-1-2} ($H_{\xi}$ does not depend on $\Lambda$). Fix $\vartheta\in(0,1]$. A collar bond $e$ is bad if $| A_{\mathrm e}(e)|_{\sigma}\geqslant \vartheta r$, good otherwise; $A_{\mathrm e}(e)$ denotes the fluctuation variable of CRE (1.13)--(1.16) at $e$ whether $e$ is exterior or interior in the original split (for $e\in I\setminus D(\Lambda)$ it is the young variable $a_e$ of the original split, exterior on the labelled terms where it is tested, Section~\ref{sec:c5-3-1}). In the construction of Section~\ref{sec:c5-3-1} a labelled term carries the regions $\Lambda_1\supset\Lambda_1^{(0)}=\operatorname{Fill}(\Lambda_1)\supset\Lambda_1^{(1)}\supset\dots\supset\Lambda_1^{\mathrm{fin}}$, unions of $\ell$-cubes, $\Lambda_1^{(0)}$ label-independent (Lemma~\ref{lem:c5-3-0}) and the others determined by the labels ($B^{(0)},B^{(1)},\dots$) by the construction of Section~\ref{sec:c5-3-1}, and the set of bonds tested on it is $T(\mathrm{labels})\coloneqq \bigcup_kJ_{\mathrm{col}}(\Lambda_1^{(k)})\subset J(\Lambda_1^{\mathrm{fin}})$. The mark of a bond $e$ on a labelled term is
\begin{equation*}
    M_e\coloneqq \mathbf 1\{ e\in T(\mathrm{labels})\}\cdot\mathbf 1\{| A_{\mathrm e}(e)|_{\sigma}\geqslant \vartheta r\}\;,
\end{equation*}
a function of the labels and of $U$ ($A_{\mathrm e}$ is a function of $U$, of $V$, the block average of $U$, and of $\Omega_1$ through $U_1(\Omega_1,V)$); $M_e=0$ on every labelled term on which $e$ is not tested, so $M_e$ is defined on every history and every labelled term; on the support of a term's indicators a tested bond is bad exactly when it belongs to some $B^{(k)}$, so there $M_e=\mathbf 1\{ e\in\bigcup_kB^{(k)}\}$ is a function of the labels alone. Let $\mathsf w^{\mathrm{lab}}_{\xi}$ be the weight of a labelled term: $\mathsf w_{\xi}=\int_{\mathfrak g^I}\rho_{\xi}(a_{\mathrm{old}})\,\dd a_{\mathrm{old}}$ with $\rho_{\xi}\coloneqq e^{c_{\xi}}K_{\xi;1}$ the integrand of Lemma~\ref{lem:collar} (the factors already placed in $\kappa_{\xi}$, $\zeta_{\xi}$, $\chi_{\xi,\mathrm e}$ of $\mathcal O$ are not duplicated), and $\mathsf w^{\mathrm{lab}}_{\xi}\coloneqq \int_{\mathfrak g^I}\rho_{\xi}(a_{\mathrm{old}})\mathbf 1_{\mathrm{lab}}(a_{\mathrm{old}},A_{\mathrm e})\,\dd a_{\mathrm{old}}$, where $\mathbf 1_{\mathrm{lab}}$ is the product of the term's collar indicators (functions of $A_{\mathrm e}$ and of $a_{\mathrm{old}}|_{I\setminus D(\Lambda_1^{\mathrm{fin}})}$, Section~\ref{sec:c5-3-1}); inside this integral the variables $a_{\mathrm{old}}|_{I\setminus D(\Lambda_1^{\mathrm{fin}})}$ are exterior for the re-split of Lemma~\ref{lem:c5-2-5} and $a_{\mathrm{old}}|_{D(\Lambda_1^{\mathrm{fin}})}$ is the young variable, so $\mathsf w^{\mathrm{lab}}_{\xi}=\int_{\mathfrak g^{I\setminus D^{\mathrm{fin}}}}\mathbf 1_{\mathrm{lab}}e^{c^{\Lambda^{\mathrm{fin}}}_{\xi}}\bigl(\int_{\mathfrak g^{D^{\mathrm{fin}}}}K^{\Lambda^{\mathrm{fin}}}_{\xi;1}\,\dd a\bigr)\,\dd a_{\mathrm{old}}|_{I\setminus D^{\mathrm{fin}}}$. Since the indicators are nonnegative measurable functions summing to one (Section~\ref{sec:c5-3-4}(c)), $\sum_{\mathrm{lab}}\mathsf w^{\mathrm{lab}}_{\xi}=\mathsf w_{\xi}$ (Tonelli). The marked physical probability is
\begin{equation*}
    p_e^{\vartheta}\coloneqq \mathcal Z_0^{-1}\int\mathcal O\Bigl[\sum_{\mathrm{lab}}\mathsf w^{\mathrm{lab}}_{\xi} M_e\Bigr](V)\,\dd V\;,
\end{equation*}
with the history sum inside $\mathcal O$ as in \eqref{eq:c3-7-1}; since $M_e$ is a function of the labels alone on the support of $\mathbf 1_{\mathrm{lab}}$, $\mathsf w^{\mathrm{lab}}_{\xi} M_e=\int_{\mathfrak g^I}\rho_{\xi}\mathbf 1_{\mathrm{lab}}M_e\,\dd a_{\mathrm{old}}$, so $p_e^{\vartheta}$ is well defined. For $e\in J$ this equals $\mathcal Z_0^{-1}\int\mathcal O[\mathsf w_{\xi}\sum_{\mathrm{lab}}\mathbf 1_{\mathrm{lab}}M_e](V)\,\dd V$ with $\sum_{\mathrm{lab}}\mathbf 1_{\mathrm{lab}}M_e$ independent of $a_{\mathrm{old}}$, because $T(\mathrm{labels})\cap J= J_{\mathrm{col}}(\Lambda_1^{(0)})\cap J$ is label-independent ($D^{(k)}\subset D^{(0)}$); for $e\in I$ the mark involves a variable that is young in the original split, and the labelled-term form is the definition.

\begin{lemma}[collar bonds are regular, eligible and in the chart]\label{lem:c5-2-3}
Let $e\in J_{\mathrm{col}}(\Lambda)$ for some $\Lambda\subset\Lambda_1$ as above. Then: (a) $d(e,S_1^*)\geqslant \ell-2-\rho_L>0$, so $\Psi_{\xi}$ is the identity at $e$ and $A'_{\mathrm e}(e)= A_{\mathrm e}(e)$ ($d(I,S_1^*)\geqslant LMR_1-2$, \eqref{eq:c-collar-distance}; $\Psi_{\xi}$ bondwise with support $S_1^*$); also $e\in\mathcal G_{\xi}\setminus S_1^*$ by Lemma~\ref{lem:collar-coverage}(i), since $\rho_L=4L\leqslant 3R_{\mathrm c}$. (b) On the branch $| A_{\mathrm e}(e)|_{\sigma}< r$ (for $e\in J$ by Lemma~\ref{lem:collar-coverage}(iii); for $e\in I\setminus D(\Lambda)$ by the box factor $\mathbf 1\{| a_e|_{\sigma}< r\}$ of the retained cutoff, Lemma~\ref{lem:collar-coverage}(ii) with $I\subset\mathcal G_{\xi}$; both from the CRE (1.25) good test $\prod_{b\in\mathcal G_{\xi}}\mathbf 1\{| A(b)|< r\}$), so a bad collar bond has $\vartheta r\leqslant | A_{\mathrm e}(e)|_{\sigma}< r$ and $g_0A_{\mathrm e}(e)$ lies in the principal chart (chart condition $h(g_0A_{\mathrm e}(e))<\pi/\sqrt N$: $h\leqslant |\bigcdot |_{\sigma}$ and $| g_0A_{\mathrm e}(e)|_{\sigma}<\delta_0$, with $\delta_0<\pi/\sqrt N$ for small coupling). (c) The $L^2M_2R_1$-cube $\square'$ of CRE (1.18) containing the initial vertex of $e$ lies within $3R_{\mathrm c}+1+L^2M_2R_1<\ell-1$ of $\Lambda_1$, hence in $\Omega_1^{\sim-1}\setminus R_1'$ (Lemma~\ref{lem:collar-coverage}(i)), and $e\in(\square'^{\sim2})^*$ (the carrier written $(Q_w^{[2]})^*$ in Proposition~\ref{prop:c3-9-1} and Lemma~\ref{lem:c3-7-seed}(c), with $Q=\square'$ and $w= L^2M_2R_1$, Section~\ref{sec:c3-7}): it is a good, eligible cube in the sense of Proposition~\ref{prop:c3-9-1}, and the comparison $\sup_{b\in(\square'^{\sim2})^*}| U_1(b)U_{1,\square'}(b)^{-1}-\id|_{\sigma}\leqslant C\varepsilon_1e^{-cLM_2R_1}$ holds for it (CRE p.250 after (1.19) with Lemma~\ref{lem:c2-plaquette-deviation}).
\end{lemma}

\begin{proof}
(a) and (b) are the cited statements of Section~\ref{sec:sourced-effective-interaction}, (c) those of Lemma~\ref{lem:c2-plaquette-deviation} and Proposition~\ref{prop:c3-9-1}, applied to $e$ using $D(\Lambda)\subset I$ and $\operatorname{dist}(e,D(\Lambda))\leqslant \rho_L$.
\end{proof}

\begin{lemma}[a bad collar bond forces a witness at threshold $\theta_{\vartheta}$]\label{lem:c5-2-4}
Put
\begin{equation*}
    \theta_{\vartheta}\coloneqq \min\bigl(\varepsilon_0,\varepsilon_1/(4B_3N^{3/2}L^2),\vartheta\delta_0/(2C_Q)\bigr)\;,\qquad C_Q=(9K_L+2)(1+2B_3N^{3/2}L^2)\;,
\end{equation*}
with $C_Q$ the constant of \eqref{eq:c3-7-4} ($k_N=\sqrt N$ as in Lemma~\ref{lem:c3-7-seed}). Since $\vartheta\leqslant 1$, $\theta_{\vartheta}\geqslant \vartheta\theta$ with $\theta=\min(\varepsilon_0,\varepsilon_1/(4B_3N^{3/2}L^2),\delta_0/(2C_Q))$ the threshold \eqref{eq:c3-7-4}. For $g_0$ below a threshold depending on $L$, $M_1$, $M_2$, the group, the source constants and $\vartheta$ (those of Lemma~\ref{lem:c3-7-2} and of Lemma~\ref{lem:c5-2-3}(b)): on every labelled term and for every tested bond $e$ with $M_e=1$, some plaquette $p$ or some active mean-transport variable $w$ within the neighbourhood $N(\square')$ of Lemma~\ref{lem:c3-7-seed} (the $6\ell$-neighbourhood of $\square'$, whose side $L^2M_2R_1\leqslant \ell/8$ by $M\geqslant 8LM_2$; its sites are within $7\ell$ of the initial vertex of $e$, the gauge-label centres within $O(L)$ more, hence within $K\ell$, $K=256$) satisfies $| U(\partial p)-\id|_{\sigma}\geqslant \theta_{\vartheta}$ or $| w-\id|_{\sigma}\geqslant \theta_{\vartheta}$.
\end{lemma}

\begin{proof}
This is a re-run of the $R_1$-step of Lemma~\ref{lem:c3-7-2} with $\delta_0$ replaced by $\vartheta\delta_0\leqslant \delta_0$; it uses the comparison \eqref{eq:c3-7-seed} of Lemma~\ref{lem:c3-7-seed}(c), which is linear in the threshold, with the constant $C_Q$ and its smallness hypotheses. By Lemma~\ref{lem:c5-2-3}(b), $| g_0A_{\mathrm e}(e)|_{\sigma}\geqslant \vartheta\delta_0$ at the free bond $e$, and the unitary logarithm inequality gives $|\widetilde U(e)U_1(e)^{-1}-\id|\geqslant (2/\pi)\vartheta\delta_0$ for the gauge-fixed field $\widetilde U$. The comparison of Lemma~\ref{lem:c5-2-3}(c) is $<\vartheta\delta_0/10$ for small coupling: $\varepsilon_1/\delta_0$ is bounded by fixed source constants (CRE (2.6), (2.8)), and $e^{-cLM_2R_1}\to0$ as $R_1\to\infty$. Hence $|\widetilde U(e)U_{1,\square'}(e)^{-1}-\id|\geqslant (2/\pi-1/10)\vartheta\delta_0>\vartheta\delta_0/2$. If every plaquette and every active gauge-variable witness in $N(\square')$ were below $\theta_{\vartheta}$, Lemma~\ref{lem:c3-7-seed}(c) at the threshold $\theta_{\vartheta}\leqslant \theta$ (its final sentence; $\square'$ is $R_1$-eligible by Lemma~\ref{lem:c5-2-3}(c)) would give $|\widetilde U(e)U_{1,\square'}(e)^{-1}-\id|\leqslant C_Q\theta_{\vartheta}\leqslant \vartheta\delta_0/2$, a contradiction. The smallness hypotheses of Lemma~\ref{lem:c3-7-seed} ($\theta\leqslant \eta_0$, $2B_3k_N^3L^2\theta\leqslant \eta_0$, $k_N\theta\leqslant c_2'$, $C_0L^2k_N\theta\leqslant 1$, $\alpha=2k_N^2L^2\theta\leqslant a_1$) are all upper bounds on the threshold and are preserved by $\theta_{\vartheta}\leqslant \theta$; the $P_0$, $P_1$ cutoffs at $\varepsilon_0$, $\varepsilon_1$ in the neighbourhood and the local determining-data problem are those of the branch, unchanged by the collar labels. The conclusion holds on the original averaging-constraint and earlier-cutoff support, which is the support on which it is used below.
\end{proof}

\begin{lemma}[the history weight of a retreated carrier]\label{lem:c5-2-5}
Let $\Lambda\subset\Lambda_1$ be a union of $\ell$-cubes, $D(\Lambda)\subset I$, $J(\Lambda)\supset J$; write $a\in\mathfrak g^{D(\Lambda)}$ for the young variable of the retreated split and $A^{\Lambda}_{\mathrm e}\in\mathfrak g^{J(\Lambda)}$ for its exterior variable, so $A^{\Lambda}_{\mathrm e}=(A_{\mathrm e},a_{\mathrm{old}}|_{I\setminus D(\Lambda)})$ in the original variables, and $A^{\Lambda,\Psi}\coloneqq \Psi_{\xi}(A^{\Lambda}_{\mathrm e})=(\Psi_{\xi}(A_{\mathrm e}),a_{\mathrm{old}}|_{I\setminus D(\Lambda)})$, $\Psi_{\xi}$ being the identity on $I\setminus D(\Lambda)\subset I$. Keep $S_1$, $R_1$, $\Psi_{\xi}$, $U_1(\Omega_1\cap\Lambda_1^c,V)$ and $\chi(\Omega_1\cap\Lambda_1^c,S_1)$ as built from the original $\Lambda_1$ (determined by $(\Omega_1,R_1)$, hence recoverable from the labels); call a primitive near if its free-input set meets $D(\Lambda)$, far otherwise (the near primitives of Section~\ref{sec:source-phase-primitives} are those whose free-input set meets $I$, so near relative to $D(\Lambda)$ implies near relative to $I$). With $\mathscr E_{\xi}(B)=\mathscr E_{\xi}(0)-\tfrac12\langle B,H_{\xi} B\rangle+\sum_{\alpha}\mathfrak p_{\alpha}(B)$ the pre-boundary-change exponent of CRE (1.15) in the primitives of Section~\ref{sec:source-phase-primitives} (\eqref{eq:native-collar-full-exponent}, written through the primitive sum \eqref{eq:native-source-primitive-sum} of Lemma~\ref{lem:primitive-sum}; $\alpha$ runs over the primitives of all three kinds), put
\begin{align*}
    c^{\Lambda}_{\xi}&\coloneqq \mathscr E_{\xi}(0)-\tfrac12\langle A^{\Lambda,\Psi},H_{J(\Lambda)J(\Lambda)}A^{\Lambda,\Psi}\rangle+\sum_{\alpha\ \mathrm{far}}\mathfrak p_{\alpha}(A^{\Lambda,\Psi})+\log|\det D\Psi_{\xi}(A_{\mathrm e})|\;,\\
    K^{\Lambda}_{\xi;1}(a)&\coloneqq \mathbf 1_{\mathcal D^{\Lambda}}(a)\exp\Bigl[-\tfrac12\langle a,H_{D(\Lambda)D(\Lambda)}a\rangle-\langle a,H_{D(\Lambda)J(\Lambda)}A^{\Lambda,\Psi}\rangle+\sum_{\alpha\ \mathrm{near}}\mathfrak p_{\alpha}(a+A^{\Lambda,\Psi})\Bigr]\;,
\end{align*}
where $\mathbf 1_{\mathcal D^{\Lambda}}(a)\coloneqq \mathbf 1\{\Psi_{\xi}(a+A^{\Lambda}_{\mathrm e})\in\mathcal A_{\xi}\}$ is the indicator of the young fibre of the retreated split (for $\Lambda=\Lambda_1$ it is $\mathbf 1_{\mathcal D_{\xi}(A_{\mathrm e})}(a)$ of \eqref{eq:retained-change-domain}, by the definition \eqref{eq:c-new-variable-domain} of $\mathcal D_{\xi}$; the clause of that definition that $\Psi_{\xi}$ be defined depends on $A_{\mathrm e}$ only and is carried by $\chi_{\xi,\mathrm e}$). Then, pointwise,
\begin{equation*}
    e^{c^{\Lambda}_{\xi}}K^{\Lambda}_{\xi;1}(a)=\mathbf 1_{\mathcal D^{\Lambda}}(a)|\det D\Psi_{\xi}(A_{\mathrm e})|\exp\mathscr E_{\xi}(\Psi_{\xi}(a+A^{\Lambda}_{\mathrm e}))\;,
\end{equation*}
and the composed retained cutoffs and $\zeta$ multiply both sides. The box factors $\mathbf 1\{| a_b|_{\sigma}< r\}$ of the bonds $b\in I\setminus D(\Lambda)$ are factors of the retained cutoff (Lemma~\ref{lem:collar-coverage}(ii), $I\subset\mathcal G_{\xi}$), now functions of exterior variables. For $\Lambda=\Lambda_1$ these are the $c_{\xi}$, $K_{\xi;1}$ of Lemma~\ref{lem:collar} and the identity \eqref{eq:c3-7-8} of Lemma~\ref{lem:c3-7-3}. ($|\det D\Psi_{\xi}(A_{\mathrm e})|=|\det D\Psi_{\xi}(A^{\Lambda}_{\mathrm e})|$, the determinant depending on the $S_1^*$-components only.)
\end{lemma}

\begin{proof}
$\Psi_{\xi}(a+A^{\Lambda}_{\mathrm e})= a+A^{\Lambda,\Psi}$, because $\Psi_{\xi}$ is bondwise with support $S_1^*$ and $d(I,S_1^*)\geqslant \ell-2>0$ (\eqref{eq:c-collar-distance}), $D(\Lambda)\subset I$. Write $z= a+A^{\Lambda,\Psi}$. Expanding $-\tfrac12\langle z,H_{\xi} z\rangle$ into the $D(\Lambda)D(\Lambda)$, $D(\Lambda)J(\Lambda)$, $J(\Lambda)J(\Lambda)$ blocks and splitting $\sum_{\alpha}\mathfrak p_{\alpha}(z)$ into near and far primitives, where a far primitive depends on $z$ only through $z|_{J(\Lambda)}= A^{\Lambda,\Psi}$ (its free-input set misses $D(\Lambda)$; Lemma~\ref{lem:row-locality}), gives $\mathscr E_{\xi}(z)=\mathscr E_{\xi}(0)-\tfrac12\langle a,H_{D(\Lambda)D(\Lambda)}a\rangle-\langle a,H_{D(\Lambda)J(\Lambda)}A^{\Lambda,\Psi}\rangle-\tfrac12\langle A^{\Lambda,\Psi},H_{J(\Lambda)J(\Lambda)}A^{\Lambda,\Psi}\rangle+\sum_{\mathrm{near}}\mathfrak p_{\alpha}(z)+\sum_{\mathrm{far}}\mathfrak p_{\alpha}(A^{\Lambda,\Psi})$; multiplying by $|\det D\Psi_{\xi}(A_{\mathrm e})|\mathbf 1_{\mathcal D^{\Lambda}}(a)$ gives the identity. The density $\mathbf 1_{\mathcal D^{\Lambda}}|\det D\Psi_{\xi}|\exp\mathscr E_{\xi}(\Psi_{\xi}(a+A^{\Lambda}_{\mathrm e}))$ is the CRE (1.15) density in the transformed variables whatever free bonds are called interior; the split only decides which factors are collected in $c^{\Lambda}_{\xi}$ and which in $K^{\Lambda}_{\xi;1}$.
\end{proof}

\begin{proposition}[normalised payment]\label{prop:c5-2-6}
On compatible even fine periods (the standing restriction of Section~\ref{sec:c3-7}), for $0< g_0< g_*(\vartheta)$, a threshold depending on $L$, $M_1$, $M_2$, $M$, the group, the source constants and $\vartheta$ (below the thresholds of Lemmas~\ref{lem:c5-2-3}(b), \ref{lem:c5-2-4}, \ref{lem:c3-7-2} and~\ref{lem:c3-7-4}, and small enough for \eqref{eq:c5-2-3}), and for every bond $e$,
\begin{equation}\label{eq:c5-2-2}
    p_e^{\vartheta}\leqslant (\mu_0\otimes q_g^{\otimes\mathcal I}\otimes\mathrm{Haar}^{\otimes T})(W_e^{\vartheta})\leqslant q_{\vartheta}\coloneqq n_W\max(z_P(\theta_{\vartheta}),z_G(\theta_{\vartheta}))\;,
\end{equation}
with $W_e^{\vartheta}\coloneqq \{\text{a plaquette or auxiliary gauge-variable witness }\geqslant \theta_{\vartheta}\text{ within }K\ell\text{ of }e\}$, $z_P(\bigcdot )$, $z_G(\bigcdot )$ the tail bounds \eqref{eq:c3-7-12}--\eqref{eq:c3-7-13} of Lemma~\ref{lem:c3-7-4}, stated there for every threshold in $(0,\varepsilon_0]$ and here applied at $\theta_{\vartheta}\leqslant \varepsilon_0$: $z_P(\theta_{\vartheta})=\exp[-b_G\theta_{\vartheta}^2/(4g_0^2)+4d_G\log g_0^{-1}+4\log c_H^{-1}+48]$, $z_G(\theta_{\vartheta})= e^{1/2}(c_Hg_0^{d_G})^{-1}e^{-b_G\theta_{\vartheta}^2/g_0^2}$, $b_G=1/(2N)$, $n_W=7(2K\ell+7)^4$. Moreover, with $\pi_*=\min(p_0(g_0),p_0(g_1))$ and $\theta\geqslant c_{\theta} g_0\pi_*$ (Proposition~\ref{prop:c3-7-5}, from CRE (2.6)--(2.9) and $\delta_0=(A_1/A_0)\varepsilon_0$),
\begin{equation}\label{eq:c5-2-3}
    q_{\vartheta}\leqslant \exp[-c_q\vartheta^2\pi_*^2]\;,\qquad c_q\coloneqq b_Gc_{\theta}^2/8\;,\qquad\text{for }0< g_0< g_*(\vartheta)\;.
\end{equation}
\end{proposition}

\begin{proof}
Fix a history and a labelled term. By Lemma~\ref{lem:c5-2-5} (with $\Lambda=\Lambda_1^{\mathrm{fin}}$ of the term) its weight $\mathsf w^{\mathrm{lab}}_{\xi}$ is the integral of the CRE (1.15) density in transformed variables, times the term's collar indicators, and the mark $M_e$ is a nonnegative function of the labels and of $U$. The reversal in the proof of Lemma~\ref{lem:c3-7-3} is a change of variables for integrals of (density $\times$ nonnegative function of $U$ and the labels), valid for every such integrand, so it applies to $\mathsf w^{\mathrm{lab}}_{\xi} M_e$ exactly as it applies to $\mathsf w_{\xi}\mathbf 1\{ z\in E(B_{\xi})\}$ there: undo $\Psi_{\xi}$, the constraint elimination (it changes designated bonds only; $e$ is free), the scaling, the linearisation, the chart and the translation by $U_1$. After the reversal $M_e$ equals $\mathbf 1\{ e\in T(\mathrm{labels})\}\cdot\mathbf 1\{|\tfrac1i\log(\widetilde U(e)U_1(\Omega_1,V)(e)^{-1})|_{\sigma}\geqslant \vartheta\delta_0\}$ ($\Psi_{\xi}$ is the identity at $e$ by Lemma~\ref{lem:c5-2-3}(a); the relevant background is $U_1(\Omega_1,V)$, as $e\notin S_1^*$). By Lemma~\ref{lem:c5-2-4}, $M_e\leqslant \mathbf 1_{W_e^{\vartheta}}$ on the support of the term's cutoffs, and $W_e^{\vartheta}$ is label-free (independent $q_g$ variables are added for inactive indices, Haar variables for centres and unused sites, as in Lemma~\ref{lem:c3-7-3}). All labelled terms are nonnegative and their partitions ($P_0$, $P_1$, $Q_1$, $R_1$ and the collar labels, Section~\ref{sec:c5-3-4}(c)) sum to one, so $\sum_{\mathrm{hist}}\sum_{\mathrm{lab}}(\text{term}\cdot M_e)\leqslant \sum_{\mathrm{hist}}\sum_{\mathrm{lab}}(\text{term}\cdot\mathbf 1_{W_e^{\vartheta}})\leqslant (\mu_0\otimes q_g^{\otimes\mathcal I}\otimes\mathrm{Haar}^{\otimes T})(W_e^{\vartheta})\cdot\mathcal Z_0$ (Lemmas~\ref{lem:c3-7-3} and~\ref{lem:c3-7-4}; only ``$\leqslant $'' is needed), which gives the first inequality of \eqref{eq:c5-2-2}; the second is the union bound of Lemma~\ref{lem:c3-7-4} over the $n_W$ witness labels with the single-witness tails $z_P$, $z_G$ at threshold $\theta_{\vartheta}$. For \eqref{eq:c5-2-3}: $\theta_{\vartheta}/g_0\geqslant \vartheta\theta/g_0\geqslant \vartheta c_{\theta}\pi_*$, so $-\log z_P(\theta_{\vartheta})\geqslant b_Gc_{\theta}^2\vartheta^2\pi_*^2/4-4d_G\log g_0^{-1}-C$ and $-\log z_G(\theta_{\vartheta})\geqslant b_Gc_{\theta}^2\vartheta^2\pi_*^2-d_G\log g_0^{-1}-C$; $\log n_W\leqslant C\log\log g_0^{-2}$ since $\ell\leqslant L^2M(2\log g_0^{-2})^{\mathfrak r}$ (from CRE (2.5)--(2.6), Section~\ref{sec:c5-3-4}); and $\pi_*^2\geqslant A_0^2(\tfrac12\log g_0^{-2})^{2p_0}$ with $p_0\geqslant 5\mathfrak r\geqslant 10$ dominates every logarithm for small $g_0$, leaving the exponent $b_Gc_{\theta}^2\vartheta^2\pi_*^2/8$.
\end{proof}

The Wilson insertion does not enter $p_e^{\vartheta}$: as in Appendix~\ref{app:exterior-chessboard}, the insertion is bounded inside the conditional estimate ($|\mathcal B_{\xi}[f_{\xi}]|\leqslant C_{\mathrm{resp},\chi}g_0rD_{\xi}$ on covered branches under \eqref{eq:c3-1-2}, Lemma~\ref{lem:exterior-response-bound}), and the history average of $D_{\xi}$ is insertion-free.

\subsection{The retreat, its termination and its costs}
\label{sec:c5-3}

Cubes are the aligned $\ell$-cubes of CRE (1.10), (1.20); $\Lambda_1$ is a union of such cubes, $D= D(\Lambda_1)$, $J= J(\Lambda_1)$, $J_{\mathrm{col}}$ as in Section~\ref{sec:c5-2-3}; for the regions $\Lambda_1^{(k)}$ and $\Lambda_1^{\mathrm{fin}}$ defined below, $D^{(k)}\coloneqq D(\Lambda_1^{(k)})$, $J^{(k)}\coloneqq J(\Lambda_1^{(k)})$, $D^{\mathrm{fin}}\coloneqq D(\Lambda_1^{\mathrm{fin}})$ and $J^{\mathrm{fin}}\coloneqq J(\Lambda_1^{\mathrm{fin}})$ (the carrier and exterior of Section~\ref{sec:c5-2-3}); $\vartheta$ as in Remark~\ref{rem:c5-1-3} ($\vartheta\leqslant 1/16$). All tests use $|\bigcdot |_{\sigma}$.

\subsubsection{The collar partition, the retreat and the filling}
\label{sec:c5-3-1}
Connected components of unions of $\ell$-cubes are taken with respect to a fixed adjacency of $\ell$-cubes (CRE's; the arguments below use only that a component of $X\cup Y$ is either a component of $X$ not adjacent to $Y$ or a set meeting $Y$). $\operatorname{Fill}_1(\Lambda)$ denotes one pass of CRE's modification after (1.20) (p.251): every connected component of $\Lambda^c$ contained in a cube of side $100\ell$ is replaced by the smallest rectangular parallelepiped of $\ell$-cubes containing it; $\operatorname{Fill}(\Lambda)$ is the fixed point of iterating $\operatorname{Fill}_1$ (each pass removes at least one cube from $\Lambda$ or stops, so the iteration ends after finitely many passes). A region $\Lambda$ is filled if $\operatorname{Fill}_1(\Lambda)=\Lambda$, i.e. every component of $\Lambda^c$ contained in a cube of side $100\ell$ is a rectangular parallelepiped; this is CRE's condition (2.3) (p.255, $j=1$: $100MR_1$ in the $L^{-1}$-lattice is $100\ell$ fine sites; the side-$100\ell$ version is used so that filled regions satisfy (2.3) literally, while p.251 fills components in cubes of side $<100\ell$). Adjacency: two $\ell$-cubes are adjacent when they share a face, so that a rectangular parallelepiped of $\ell$-cubes is connected (used in the equivalence ``filled $\Leftrightarrow\operatorname{Fill}_1(\Lambda)=\Lambda$''); CRE p.251 does not specify its adjacency, nor does the filling induction in the proof of Lemma~\ref{lem:c3-7-2}, and the proofs below hold for any symmetric adjacency under which parallelepipeds are connected. Standing hypothesis of this section and of Sections~\ref{sec:c5-4}--\ref{sec:c5-7}, which import it, additional to \eqref{eq:c-standing-conditions} (every fine period a multiple of $LMR_1$ and at least $8M$) and to the compatible even fine periods of Appendix~\ref{app:exterior-chessboard}: every fine period exceeds $200\ell$. It is coupling-dependent through $\ell= LMR_1(g_0)$ (on a fixed torus it fails as $g_0\to0$), so every statement using it holds, at a fixed small $g_0$, on every compatible torus whose periods exceed $200\ell$: a lower bound on the periods in terms of the coupling, not a limit (Section~\ref{app:thresholds}). Under it a set contained in a cube of side $100\ell$ lifts isometrically to $\mathbb Z^4$, and its smallest containing parallelepiped of $\ell$-cubes is well defined. Whether CRE's $\Lambda_1$ is filled is not used: round $0$ fills it.

\begin{lemma}[round $0$]\label{lem:c5-3-0}
Put $\Lambda_1^{(0)}\coloneqq \operatorname{Fill}(\Lambda_1)$ and $\mathfrak B_0\coloneqq (\Lambda_1^{(0)})^c\supset\Lambda_1^c$ (a set of sites; the bad bond sets of the labels are $B^{(k)}$). (a) $\Lambda_1^{(0)}\subset\Lambda_1$ is a filled union of $\ell$-cubes, determined by $(\Omega_1,R_1)$ like $\Lambda_1$, so label-independent; if $\Lambda_1$ already satisfies (2.3) then $\Lambda_1^{(0)}=\Lambda_1$. (b) Every point of $\mathfrak B_0$ is within $100\ell$ of a point of the unfilled bad region $Q_1'^{\sim4}\cup(B(P_1^1)^c)^{\sim5}\cup R_1'^{\sim1}$ of CRE (1.10), (1.20), so the witness argument of Lemma~\ref{lem:c3-7-2} applies to $\mathfrak B_0$ without change: every $z\in E(\mathfrak B_0)$ forces a witness $\geqslant \theta$ within $K\ell$ of $z$, $K=256$, with the same $n_W$, and $p_z^{(0)}\coloneqq \mathcal Z_0^{-1}\int\mathcal O[\mathsf w_{\xi}\mathbf 1\{ z\in E(\mathfrak B_0)\}](V)\,\dd V\leqslant q$ on compatible even fine periods, for $g_0$ below the thresholds of Lemmas~\ref{lem:c3-7-2} and~\ref{lem:c3-7-4} ($g_0\leqslant g_H$ there).
\end{lemma}

\begin{proof}
(a) is the definition ($\operatorname{Fill}(\Lambda)\subset\Lambda$, and the output is filled). (b) The induction of Lemma~\ref{lem:c3-7-2} (every current component contains an original bad point, and each fill is performed inside a cube of side less than $100\ell$ containing that component; this also holds if small merged components are filled again) is run with CRE's p.251 fill as the base and the passes of $\operatorname{Fill}_1$ as further steps: a $\operatorname{Fill}_1$ pass changes only components contained in a cube of side $100\ell$, each containing an original bad point, and two points of such a cube are within $100\ell-1<100\ell$ of each other, so every added point is within $100\ell$ of an original bad point of its component; mergers preserve the presence of an original point. A point of the unfilled region is within $10\ell$ of a bad test (proof of Lemma~\ref{lem:c3-7-2}), whose witness lies in the witness neighbourhood of the test cube (the $6\ell$-neighbourhood $N(\square')$ of Lemma~\ref{lem:c3-7-seed}, as in Lemma~\ref{lem:c5-2-4}, or the defining enlargement of $P_0$), within $7\ell$ of the cube, whose side is less than $\ell$; so the witness is within $118\ell+O(L)+1< K\ell$ of $z$. The mark $\mathbf 1\{ z\in E(\mathfrak B_0)\}$ depends on the regions only, so Lemmas~\ref{lem:c3-7-3} and~\ref{lem:c3-7-4} apply verbatim with the same $n_W$, giving $p_z^{(0)}\leqslant q$.
\end{proof}

On a labelled term with current region $\Lambda_1^{(k)}$ ($k=0$: $\Lambda_1^{(0)}=\operatorname{Fill}(\Lambda_1)$, Lemma~\ref{lem:c5-3-0}; the bonds of $D\setminus D^{(0)}$ are exterior on every term, with their box factors as exterior good tests, Lemma~\ref{lem:c5-2-5}) insert, inside the $a$-independent part of the integrand, $1=\prod_e[\mathbf 1\{| A_{\mathrm e}(e)|_{\sigma}<\vartheta r\}+\mathbf 1\{| A_{\mathrm e}(e)|_{\sigma}\geqslant \vartheta r\}]$ over the untested bonds $e$ of $J_{\mathrm{col}}(\Lambda_1^{(k)})$ (collar bonds have $A'_{\mathrm e}(e)= A_{\mathrm e}(e)$, Lemma~\ref{lem:c5-2-3}(a), so the factors are $a$-independent, and they remain exterior factors on every later term, since $J(\Lambda)$ only grows as $\Lambda$ shrinks), and expand. A term is labelled by its bad set $B^{(k)}$. If $B^{(k)}=\varnothing$ the iteration stops with $\Lambda_1^{\mathrm{fin}}=\Lambda_1^{(k)}$. Otherwise put
\begin{gather*}
    \Lambda_1^{(k+1)}\coloneqq \operatorname{Fill}\Bigl(\Lambda_1^{(k)}\setminus\bigcup_{e\in B^{(k)}}\mathcal Q(e)\Bigr)\;,\\
    \mathcal Q(e)\coloneqq \text{the union of the $\ell$-cubes meeting the closed $\ell^{\infty}$-ball of radius $\rho_L+1$}\\
    \text{around the initial site $x_e$ of $e$.}
\end{gather*}
Since $2(\rho_L+1)+1=8L+3<\ell$ (the standing condition $LMR_1>4R_{\mathrm c}+4L^2M_2R_1+8$ of \eqref{eq:c-standing-conditions}), $\mathcal Q(e)$ consists of at most $2^4=16$ cubes and every point of $\mathcal Q(e)$ is within $\ell+\rho_L+1\leqslant 2\ell$ of $x_e$. The bonds of $D^{(k)}\setminus D^{(k+1)}$ become exterior variables of the term; they are the same bond variables as before (the chart CRE (1.13)--(1.16) and $U_1$ depend on $\Omega_1$ and $V$, not on $\Lambda_1$: CRE (1.12)--(1.16) are defined on $\Omega_1$, and $\Lambda_1$ first enters in (1.20)), and their box factors $\mathbf 1\{| a_b|_{\sigma}< r\}$ become exterior good tests (Lemma~\ref{lem:c5-2-5}). Convention (Lemma~\ref{lem:c5-2-5}): $S_1$, $R_1$, $\Psi_{\xi}$, $U_1(\Omega_1\cap\Lambda_1^c,V)$ and $\chi(\Omega_1\cap\Lambda_1^c,S_1)$ are those of the original $\Lambda_1$ throughout (consistent, since $\Omega_1\cap\Lambda_1^c\subset\Omega_1\cap(\Lambda_1^{(k)})^c$); the retreated term is the same integral with a different split of the free bonds into integrated and exterior ones, and $\Psi_{\xi}$ remains an invertible change of variables with its Jacobian retained (\eqref{eq:retained-change-domain}). The weight $\mathsf w^{\mathrm{lab}}_{\xi}$ of Section~\ref{sec:c5-2-3} is the integral of $\rho_{\xi}\times$(the term's indicators) over all of $\mathfrak g^I$, the variables of $I\setminus D^{\mathrm{fin}}$ being integrated as exterior variables of the re-split.

\begin{lemma}[termination, final state, counts]\label{lem:c5-3-1}
(a) The iteration stops after at most $|\Lambda_1|/\ell^4$ rounds with $B^{(k)}\neq\varnothing$, followed by one terminal round. (b) On every final term, every $e\in J_{\mathrm{col}}(\Lambda_1^{\mathrm{fin}})$ is good: $| A_{\mathrm e}(e)|_{\sigma}<\vartheta r$. (c) Every region $\Lambda_1^{(k)}$ is filled, and every point of $\Lambda_1^{(0)}\setminus\Lambda_1^{\mathrm{fin}}$ is within $102\ell$ of $x_e$ for some bond $e$ found bad on the term. (d) The bonds with an endpoint within $102\ell$ of a given site number at most $N_{\mathrm{ret}}\coloneqq 4(207\ell)^4$; hence at most $N_{\mathrm{ret}}$ bonds are removed on account of one bad bond $e$, and for a given bond $z$ at most $N_{\mathrm{ret}}$ bonds $e$ have $x_e$ within $102\ell+1$ of an endpoint of $z$ (that count is $4(204\ell+4)(204\ell+3)^3\leqslant N_{\mathrm{ret}}$ for $\ell\geqslant 2$).
\end{lemma}

\begin{proof}
(a) A bad $e\in J_{\mathrm{col}}(\Lambda_1^{(k)})$ has $\operatorname{dist}(e,D^{(k)})\leqslant \rho_L$, so some bond of $D^{(k)}$ has an endpoint in $\Lambda_1^{(k)}$ within $\rho_L+1$ of $x_e$; the cube containing that endpoint lies in $\mathcal Q(e)\cap\Lambda_1^{(k)}$ and is removed. So each non-terminal round removes at least one cube. (b) $\Lambda_1^{\mathrm{fin}}=\Lambda_1^{(k)}$ for the terminal round $k$, whose untested collar bonds were all good; the collar bonds tested earlier were either good (the same variable, still good) or bad, and a bad bond $e$ cannot lie in $J_{\mathrm{col}}(\Lambda_1^{\mathrm{fin}})$: $\Lambda_1^{\mathrm{fin}}\subset\Lambda_1^{(j)}\setminus\mathcal Q(e)$ for the round $j$ in which $e$ was found bad ($\operatorname{Fill}$ only removes cubes), so every $b\in D^{\mathrm{fin}}$ has an endpoint in $\Lambda_1^{\mathrm{fin}}$, disjoint from the closed ball of radius $\rho_L+1$ around $x_e$; if that endpoint is $x_b$ then $| x_b-x_e|_{\infty}\geqslant \rho_L+2$, and if it is $x_b+e_{\mu}$ then $| x_b-x_e|_{\infty}\geqslant \rho_L+1$; in both cases $\operatorname{dist}(b,e)>\rho_L$. (c) Induction over the rounds and over the passes of $\operatorname{Fill}$, with the invariant: the current region is filled, and every point removed so far lies within $100\ell$ of a point of $\bigcup\mathcal Q(e)$ over the bonds $e$ found bad so far. The invariant holds for $\Lambda_1^{(0)}$ (filled by Lemma~\ref{lem:c5-3-0}(a); nothing removed yet). In a round, removing the cubes of $\mathcal Q(B^{(k)})$ gives a complement whose components are either components of $(\Lambda_1^{(k)})^c$ not adjacent to any removed cube, or sets meeting $\bigcup_{e\in B^{(k)}}\mathcal Q(e)$. A component of the first kind contained in a cube of side $100\ell$ is a parallelepiped ($\Lambda_1^{(k)}$ is filled), so a $\operatorname{Fill}_1$ pass does not change it; a pass changes only components of the second kind, each inside a cube of side $100\ell$ that contains a point of some $\mathcal Q(e)$, so every point added lies within $100\ell$ of that point. After the pass the components are again either untouched components of the first kind or sets meeting $\bigcup\mathcal Q(e)$ (a merger involving a filled hull contains a $\mathcal Q$-point). Iterating to the fixed point keeps the invariant, and the output is filled by definition. Since every point of $\mathcal Q(e)$ is within $2\ell$ of $x_e$, every removed point is within $102\ell$ of $x_e$ for some bad $e$. (d) Sites within $102\ell$ of a given site number at most $(204\ell+1)^4$; the bonds with an endpoint among them have their initial site within $102\ell+1$, at most $4(204\ell+3)^4\leqslant 4(207\ell)^4$ bonds. Sites within $102\ell+1$ of one of the two endpoints of $z$ lie in a box of $(204\ell+4)(204\ell+3)^3$ sites, each the initial site of four bonds, and $204\ell+4\leqslant 207\ell$ for $\ell\geqslant 2$.
\end{proof}

\begin{remark}[what the final state gives]\label{rem:c5-3-2}
On a final term every $e\in J^{\mathrm{fin}}$ within $\rho_L$ of $D^{\mathrm{fin}}$ is good (Lemma~\ref{lem:c5-3-1}(b)) and $A'_{\mathrm e}(e)= A_{\mathrm e}(e)$ there (for $e\in J$ by Lemma~\ref{lem:c5-2-3}(a); for $e\in I\setminus D^{\mathrm{fin}}$ because $I\cap S_1^*=\varnothing$), so \eqref{eq:c5-1-2} gives $\sup_{b\in D^{\mathrm{fin}}}| m_b|_{\sigma}\leqslant r/4$ with the same $C_H^{\mathrm{near}}$ and $\vartheta$: (I1) is unchanged, (I2) is a statement about the rows $b\in I\supset D^{\mathrm{fin}}$ with the sum over all $b'$, and (I3) passes to $D^{\mathrm{fin}}$ by principal restriction. This is the input for Sections~\ref{sec:c5-4} and~\ref{sec:c5-5}, with margin $3r/4$.
\end{remark}

\begin{lemma}[the regions and the retreated fibre]\label{lem:c5-3-3}
(i) $\Lambda_1^{\mathrm{fin}}\subset\Lambda_1^{(0)}\subset\Lambda_1\subset\Omega_1^{\sim-2}\cap(R_1'^{\sim})^c$ is a filled union of $\ell$-cubes (Lemmas~\ref{lem:c5-3-0}(a), \ref{lem:c5-3-1}(c)), so it satisfies CRE (2.3) with the adjacency of Section~\ref{sec:c5-3-1}. (ii) Lemma~\ref{lem:collar-coverage}(i) (every site within $\ell-1$ of $\Lambda_1^{\mathrm{fin}}$ lies in $\Omega_1^{\sim-1}\setminus R_1'$; every free bond of $\Omega_1$ within $3R_{\mathrm c}$ of $D^{\mathrm{fin}}$ lies in $\mathcal G_{\xi}\setminus S_1^*$) and \eqref{eq:c-collar-distance} ($d(D^{\mathrm{fin}},S_1^*)\geqslant \ell-2$) hold by the inclusions $\Lambda_1^{\mathrm{fin}}\subset\Lambda_1$, $D^{\mathrm{fin}}\subset I$. (iii) $| A^{\mathrm{fin}}_{\mathrm e}(b)|_{\sigma}< r$ for every $b\in J^{\mathrm{fin}}$ with $d(b,D^{\mathrm{fin}})\leqslant 3R_{\mathrm c}$: for $b\in J$ by Lemma~\ref{lem:collar-coverage}(iii) ($d(b,I)\leqslant d(b,D^{\mathrm{fin}})$), for $b\in I\setminus D^{\mathrm{fin}}$ by the retained box factor $\mathbf 1\{| a_b|_{\sigma}< r\}$ (Lemma~\ref{lem:c5-2-5}), nonzero on the term. (iv) On a final term with exterior data $A^{\mathrm{fin}}_{\mathrm e}$ for which the young fibre $\{ a\in\mathfrak g^{D^{\mathrm{fin}}}\, | \,\Psi_{\xi}(a+A^{\mathrm{fin}}_{\mathrm e})\in\mathcal A_{\xi}\}$ is nonempty, the fibre is the open box $\{ a\, | \,| a_b|_{\sigma}< r\text{ for all }b\in D^{\mathrm{fin}}\}$ (otherwise it is empty and the term vanishes). $\Omega_1$, $R_1$, $S_1$, $\Psi_{\xi}$ and $U_1$ are those of the original $\Lambda_1$ by the convention of Section~\ref{sec:c5-3-1}. The history label acquires the finite sequence $(B^{(0)},B^{(1)},\dots)$ of bad sets; its bookkeeping in the history sum of Section~\ref{sec:sourced-effective-interaction} is Section~\ref{sec:c5-6}.
\end{lemma}

\begin{proof}[Proof of (iv)]
The proof of Lemma~\ref{lem:collar-coverage}(iv) is re-run with $I$ replaced by $D^{\mathrm{fin}}$; its inputs are checked one by one. Let $a_0$ be a point of the fibre (nonempty on a term that survives) and $a$ any point of the box; $B_0=\Psi_{\xi}(a_0+A^{\mathrm{fin}}_{\mathrm e})\in\mathcal A_{\xi}$ and $B= a+\Psi_{\xi}(A^{\mathrm{fin}}_{\mathrm e})$ (Lemma~\ref{lem:c5-2-5}). The cutoff CRE (1.25) at $B$ factorises as $\Theta_{D^{\mathrm{fin}}}(a)\times$(a function of $\Psi_{\xi}(A^{\mathrm{fin}}_{\mathrm e})$ and the branch data): the proof of Lemma~\ref{lem:collar-coverage}(ii) uses only $D^{\mathrm{fin}}\subset\mathcal G_{\xi}$, the bondwise identity $\Psi_{\xi}(a+A^{\mathrm{fin}}_{\mathrm e})= a+\Psi_{\xi}(A^{\mathrm{fin}}_{\mathrm e})$, and that the bonds of the factors $\chi^{(0)c}(R_1)$, $\chi_0'(S_1)$ have no endpoint in $\Lambda_1\supset\Lambda_1^{\mathrm{fin}}$; so (1.25) at $B$ equals (1.25) at $B_0$ with $\Theta_{D^{\mathrm{fin}}}(a_0)=1$ replaced by $\Theta_{D^{\mathrm{fin}}}(a)=1$, nonzero. Every row equation has a solution at $g_0\mathsf C_UB$: on the rows within $2R_{\mathrm c}+2L$ of $D^{\mathrm{fin}}$ (near rows for $I$ as well, $D^{\mathrm{fin}}\subset I$) by Lemma~\ref{lem:charts}, whose polydisc condition $| B(b)|<2+4r$ on $\mathcal E_c$ holds by (iii), $| a_b|< r$ and $B= a+\Psi_{\xi}(A^{\mathrm{fin}}_{\mathrm e})$ with $\Psi_{\xi}$ the identity on $\mathcal E_c$ (the bonds of $\mathcal E_c$ are within $2R_{\mathrm c}+4L-1\leqslant 3R_{\mathrm c}$ of $D^{\mathrm{fin}}$, hence in $\mathcal G_{\xi}\setminus S_1^*$ by (ii)); on the other rows by Lemma~\ref{lem:row-locality}, since $\mathcal E_c$ misses $D^{\mathrm{fin}}$ and $B= B_0$ off $D^{\mathrm{fin}}$, where the solution exists because $B_0\in\mathcal A_{\xi}$. The constraint $\mathsf M(U(B))= V$ on every row and the restriction $V_0$ off $\Omega_1$ follow as in the proof of Lemma~\ref{lem:collar-coverage}(iv) (they use the row equations only). The tests of CRE (1.11) at $U(B)$: $\chi_1(\Omega_1)$ (a function of $V$) and $\zeta(\Omega_1^c)$ (\eqref{eq:c-large-field-factor}) do not depend on $a$, so they take at $U(B)$ their values at $U(B_0)$, which are nonzero; $\chi_0'(S_1)$ is nonzero by the first condition; every other cube $\square'\subset\Omega_1$ lies in $\Omega_1^{\sim-1}\setminus R_1$ and its factor is the second factor of (1.19), equal to one by (B3)(ii). Hence $B\in\mathcal A_{\xi}$. Conversely a point of the fibre has nonzero (1.25), hence $\Theta_{D^{\mathrm{fin}}}(a)=1$.
\end{proof}

\subsubsection{Costs of the retreat and the comparison}
\label{sec:c5-3-4}
(a) Coverage and attachment. Appendix~\ref{app:exterior-chessboard} defines $D_{\xi}=\sum_{b\in I_C}e^{-\mu d(b,J)/2}$ (\eqref{eq:c3-5-1}), with $I_C\subset I$ the young bonds in the free-input stencils of the contour traversals, and the covered event $A_{\mathrm{cov}}$ (contour bonds in $\Omega_1$, free inputs in $I$; Section~\ref{sec:c3-6}). For a final term with carrier $D^{\mathrm{fin}}$ the corresponding quantities are $D^{\mathrm{fin}}_{\xi}\coloneqq \sum_{b\in I_C\cap D^{\mathrm{fin}}}e^{-\mu d(b,J^{\mathrm{fin}})/2}$ and $A^{\mathrm{fin}}_{\mathrm{cov}}\coloneqq \{\text{all contour bonds lie in }\Omega_1\text{ and every free input lies in }D^{\mathrm{fin}}\}$; both depend on the labels only through the regions. Put $B^{\mathrm{fin}}_{\xi}\coloneqq (\Lambda_1^{\mathrm{fin}})^c$ (the bad region $B_{\xi}$ of Section~\ref{sec:c3-7} is that of the original $\Lambda_1$). Pointwise on the term, with $c=\mu/2$,
\begin{equation*}
    D^{\mathrm{fin}}_{\xi}\leqslant \sum_{b\in\Gamma_C}\sum_ze^{-cd(b,z)}\mathbf 1\{ z\in E(B^{\mathrm{fin}}_{\xi})\}\;,\qquad\mathbf 1_{(A^{\mathrm{fin}}_{\mathrm{cov}})^c}\leqslant \sum_{b\in\Gamma_C}\mathbf 1\{ b\in E(B^{\mathrm{fin}}_{\xi})\}:
\end{equation*}
a free bond not in $D^{\mathrm{fin}}$ has no endpoint in $\Lambda_1^{\mathrm{fin}}$, so $J^{\mathrm{fin}}\subset E(B^{\mathrm{fin}}_{\xi})$, and each $b\in I_C\cap D^{\mathrm{fin}}\subset\Gamma_C$ contributes $e^{-cd(b,J^{\mathrm{fin}})}\leqslant \sum_{z\in J^{\mathrm{fin}}}e^{-cd(b,z)}$; a contour bond outside $\Omega_1$ or a free input outside $D^{\mathrm{fin}}$ meets $B^{\mathrm{fin}}_{\xi}$ (the argument of Lemma~\ref{lem:c3-7-1} with $I$ replaced by $D^{\mathrm{fin}}$). These two inequalities are combinatorial. The statement that the response of the retreated conditional integral is bounded by $C_{\mathrm{resp},\chi}g_0rD^{\mathrm{fin}}_{\xi}$ is the conditional estimate for the retreated carrier (Section~\ref{sec:c5-4}, with Remark~\ref{rem:c5-3-2} and Lemma~\ref{lem:c5-3-3} as its inputs). For the history average: by Lemma~\ref{lem:c5-3-1}(c)--(d), on every final term
\begin{equation}\label{eq:c5-3-mark-inequality}
    \mathbf 1\{ z\in E(B^{\mathrm{fin}}_{\xi})\}\leqslant \mathbf 1\{ z\in E(\mathfrak B_0)\}+\sum_{e\, : \,x_e\text{ within }102\ell+1\text{ of an endpoint of }z}M_e
\end{equation}
(a bond $z$ meets $(\Lambda_1^{\mathrm{fin}})^c$ either because it meets $\mathfrak B_0=(\Lambda_1^{(0)})^c$ or because one of its endpoints lies in $\Lambda_1^{(0)}\setminus\Lambda_1^{\mathrm{fin}}$, which is within $102\ell$ of $x_e$ for a bond $e$ found bad on the term, $M_e=1$). Multiply by $\rho_{\xi}\mathbf 1_{\mathrm{lab}}$, integrate over $\mathfrak g^I$, sum over labels and histories, apply $\mathcal O$ and divide by $\mathcal Z_0$ (all integrands nonnegative, finitely many labels per history, Tonelli): the first term is label-independent, $\sum_{\mathrm{lab}}\mathbf 1_{\mathrm{lab}}=1$ reduces it to $\mathcal Z_0^{-1}\int\mathcal O[\mathsf w_{\xi}\mathbf 1\{ z\in E(\mathfrak B_0)\}]\,\dd V= p_z^{(0)}\leqslant q$ (Lemma~\ref{lem:c5-3-0}(b)); the second gives $\sum_{e\text{ near }z}p_e^{\vartheta}\leqslant N_{\mathrm{ret}}q_{\vartheta}$ (Lemma~\ref{lem:c5-3-1}(d), Proposition~\ref{prop:c5-2-6}). Hence, on compatible even fine periods exceeding $200\ell$ and for $0< g_0< g_*(\vartheta)$ (Proposition~\ref{prop:c5-2-6}; this threshold also covers Lemma~\ref{lem:c5-3-0}(b)),
\begin{equation}\label{eq:c5-3-coverage-average}
    \mathcal Z_0^{-1}\int\mathcal O\Bigl[\sum_{\mathrm{lab}}\mathsf w^{\mathrm{lab}}_{\xi} D^{\mathrm{fin}}_{\xi}\Bigr]\dd V\leqslant C_LS_cn_C(q+N_{\mathrm{ret}}q_{\vartheta})\;,\qquad\mathcal Z_0^{-1}\int\mathcal O\Bigl[\sum_{\mathrm{lab}}\mathsf w^{\mathrm{lab}}_{\xi}\mathbf 1_{(A^{\mathrm{fin}}_{\mathrm{cov}})^c}\Bigr]\dd V\leqslant C_Ln_C(q+N_{\mathrm{ret}}q_{\vartheta})\;.
\end{equation}
This is the bound \eqref{eq:c3-7-3} with $q$ replaced by $q+N_{\mathrm{ret}}q_{\vartheta}$, which has the same normalised form. The scalar $\kappa_C$ of \eqref{eq:extracted-loop-scalar} does not enter the history average: the ratio $\kappa_C^{-1}\mathcal B_{\xi}[f_{\xi}]$ is formed in Section~\ref{sec:c5-7} from the response bound and this average.

(b) Boundary activities. The bonds that change status are $D\setminus D^{(0)}$ (label-independent, round $0$, Lemma~\ref{lem:c5-3-0}) and, on a final term, $D^{(0)}\setminus D^{\mathrm{fin}}$ (young in the original split, exterior on the term): at most $N_{\mathrm{ret}}$ of the latter per bad bond $e$, each with an endpoint within $102\ell$ of $x_e$ (Lemma~\ref{lem:c5-3-1}(c)--(d)), and $p_e^{\vartheta}\leqslant q_{\vartheta}$. How the connected boundary terms of the retreated conditional integral are localised relative to these bonds, and their bound, are Section~\ref{sec:c5-5}. If that cost on a term is a sum over the bonds $b\in D^{(0)}\setminus D^{\mathrm{fin}}$ of quantities bounded by $Ce^{-cd(b,\Gamma_C)}$ with $C$, $c$ independent of the labels, the fields and $r$, then it is linear in the marks, (c) applies, and its history average is at most $CN_{\mathrm{ret}}\sum_be^{-cd(b,\Gamma_C)}\sup_ep_e^{\vartheta}\leqslant CN_{\mathrm{ret}}S_cC_Ln_Cq_{\vartheta}$, uniform in the volume; without such a localisation the sum over all testable bonds is proportional to the volume, and if the cost were multiplicative in the bad bonds, (c) would not apply. Section~\ref{sec:c5-5} establishes the additive form: in \eqref{eq:c5-5-9} the source-containing connected boundary terms are bounded by sums over $e\in J(\Lambda)$ of $e^{-\mu'\operatorname{dist}(e,C)}$ (and over $e\in J_1(\Lambda)$ of $e^{-\mu_H\operatorname{dist}(e,I_C)/2}$) with prefactors uniform in the exterior data within the term's support, and Section~\ref{sec:c5-7} takes their history average through \eqref{eq:c5-3-mark-inequality} and (c). The bonds $D\setminus D^{(0)}$ carry no marks; they form the label-independent boundary of $\Lambda_1^{(0)}$, belong to $J(\Lambda)$ on every final term, and enter \eqref{eq:c5-5-9} through the sum over $J(\Lambda)$ like the bonds of $J$.

(c) No label entropy for the coverage cost. The collar labels are a partition of unity on each term: $\sum_{B^{(k)}}\prod_{e\in B^{(k)}}\mathbf 1\{ e\text{ bad}\}\prod_{e\notin B^{(k)}}\mathbf 1\{ e\text{ good}\}=1$ pointwise, round by round, so by induction from the terminal rounds (termination: Lemma~\ref{lem:c5-3-1}(a)) $\sum_{\mathrm{lab}}\mathbf 1_{\mathrm{lab}}=1$ pointwise on $\mathfrak g^I\times\mathfrak g^J$, and $\sum_{\mathrm{lab}}\mathsf w^{\mathrm{lab}}_{\xi}=\int_{\mathfrak g^I}\rho_{\xi}\sum_{\mathrm{lab}}\mathbf 1_{\mathrm{lab}}\,\dd a_{\mathrm{old}}=\mathsf w_{\xi}$ (Tonelli; the labels are indicators of finitely many measurable sets). Every term is nonnegative and the cost in (a) is bounded by $\sum_eM_e\times$(a local quantity), linear in the marks; summing the labelled terms back therefore gives $\sum_eM_e\times$(local quantity) under the history average, with no product over labels and no factor depending on the number of tested bonds. This is established for the cost (a) only; for (b) see the conditional sentence of (b) and Section~\ref{sec:c5-5}.

Comparison. $r=\delta_0/g_0= A_1(\log g_0^{-2})^{p_0}$ (CRE p.246, p.247: $\varepsilon_0= g_0A_0(\log g_0^{-2})^{p_0}$, $\delta_0=(A_1/A_0)\varepsilon_0$; $r$ is not used in \eqref{eq:c5-3-1}--\eqref{eq:c5-3-2}). CRE (2.6) (p.255, $n=1$, $m=0$): $g_1\leqslant (1+\beta_0)g_0$ and $g_0\leqslant (1+\beta_0)g_1$, so for small $g_0$, $\tfrac12\log g_0^{-2}\leqslant \log g_1^{-2}\leqslant 2\log g_0^{-2}$. $R_1$ is the smallest power of $L$ with $R_1\geqslant (\log g_1^{-2})^{\mathfrak r}$ (p.246), so $R_1\leqslant L(\log g_1^{-2})^{\mathfrak r}\leqslant L(2\log g_0^{-2})^{\mathfrak r}$ and
\begin{gather}
    \ell= LMR_1\leqslant L^2M(2\log g_0^{-2})^{\mathfrak r}\;,\qquad\text{hence}\qquad\log N_{\mathrm{ret}}\leqslant C+4\mathfrak r\log(2\log g_0^{-2})\label{eq:c5-3-1}\\
    \text{and}\qquad C\ell^4\leqslant C(L^2M)^42^{4\mathfrak r}(\log g_0^{-2})^{4\mathfrak r}\;.\notag
\end{gather}
With $\pi_*\geqslant A_0(\tfrac12\log g_0^{-2})^{p_0}$ and $p_0\geqslant 5\mathfrak r$ (p.246),
\begin{equation}\label{eq:c5-3-2}
    \log(1/q_{\vartheta})\geqslant c_q\vartheta^2A_0^22^{-2p_0}(\log g_0^{-2})^{2p_0}\gg C(L^2M)^42^{4\mathfrak r}(\log g_0^{-2})^{4\mathfrak r}\geqslant C\ell^4\;,
\end{equation}
since $2p_0\geqslant 10\mathfrak r>4\mathfrak r$. Hence for $g_0$ below a threshold depending on $L$, $M$, $M_1$, $M_2$, the group, the source constants (CRE's $A_0$, $A_1$, $\beta_0$, $p_0$, $\mathfrak r$ among them; $M_1$ enters through $g_*(\vartheta)$) and $\vartheta$, $N_{\mathrm{ret}}q_{\vartheta}\leqslant e^{-c_q\vartheta^2\pi_*^2/2}$; and for each $C$ independent of $g_0$ there is a threshold, depending also on $C$, below which $e^{C\ell^4}q_{\vartheta}\leqslant e^{-c_q\vartheta^2\pi_*^2/2}$ ($C$ a constant of $L$, $M$, $M_2$, the group and the source constants; a per-term cost growing with $r$, say like $r^2\propto(\log g_0^{-2})^{2p_0}$, is not covered by this comparison and would have to be paid differently). This is CRE's $p_0\geqslant 5\mathfrak r$ accounting with the present constants.

\subsection{The selected majorant for the translated box}
\label{sec:c5-4}

Fix a history and a labelled term with final region $\Lambda\coloneqq \Lambda_1^{\mathrm{fin}}$ (a filled union of $\ell$-cubes, Lemma~\ref{lem:c5-3-3}(i)), carrier $D\coloneqq D(\Lambda)= D^{\mathrm{fin}}\subset I$, exterior $J(\Lambda)= E_{\Omega,*}\setminus D$ and exterior variable $A^{\Lambda,\Psi}=(\Psi_{\xi}(A_{\mathrm e}),a_{\mathrm{old}}|_{I\setminus D})$ (Lemma~\ref{lem:c5-2-5}). Write $H\coloneqq H_{D(\Lambda)D(\Lambda)}$, $\ell_{\Lambda}\coloneqq H_{D(\Lambda)J(\Lambda)}A^{\Lambda,\Psi}\in\mathfrak g^D$ and, for $t\in[0,1]$,
\begin{gather*}
    P^{\Lambda}_t(a)\coloneqq \sum_{\mathfrak p\text{ near}}\mathfrak p(a+tA^{\Lambda,\Psi})\;,\qquad K^{\Lambda}_{\xi;t}(a)\coloneqq \mathbf 1_{\mathcal D^{\Lambda}}(a)\exp\Bigl[-\tfrac12\langle a,Ha\rangle-t\langle a,\ell_{\Lambda}\rangle+P^{\Lambda}_t(a)\Bigr]\;,\\
    I^{\Lambda}_t\coloneqq \int_{\mathfrak g^D}K^{\Lambda}_{\xi;t}\,\dd a\;,\qquad\dd\nu^{\Lambda}_t\coloneqq K^{\Lambda}_{\xi;t}\,\dd a/I^{\Lambda}_t\;,
\end{gather*}
the retreated split's version of the exterior interpolation \eqref{eq:native-source-boundary-interpolation}: near means near relative to $D(\Lambda)$ (Lemma~\ref{lem:c5-2-5}), $\mathfrak p$ runs over the near primitives of all three kinds, and for $\Lambda=\Lambda_1$ these are the $P_{\xi,t}$, $K_{\xi;t}$, $I_{\xi;t}$, $\nu_{\xi;t}$ of Section~\ref{sec:sourced-effective-interaction} ($\Psi_{\xi}$ is the identity on the bonds read by the near primitives and by $H_{IJ}$, Lemma~\ref{lem:collar-coverage}(i)); at $t=1$, $K^{\Lambda}_{\xi;1}$ is that of Lemma~\ref{lem:c5-2-5}. On the support of the term's exterior cutoffs (the branch indicator $\iota_{\xi}$, the collar indicators and the box factors on $I\setminus D(\Lambda)$, Lemma~\ref{lem:c5-2-5}) the retreated young fibre is nonempty and is the box: by Lemma~\ref{lem:collar-coverage}(iv) the original fibre over the branch data is the box over $I$, and the box factors on $I\setminus D(\Lambda)$ place $(a,a_{\mathrm{old}}|_{I\setminus D(\Lambda)})$ in it for every $a$ with $| a_b|_{\sigma}< r$, $b\in D(\Lambda)$, so the fibre of the retreated split over $A^{\Lambda,\Psi}$ is nonempty and Lemma~\ref{lem:c5-3-3}(iv) gives $\mathbf 1_{\mathcal D^{\Lambda}}(a)=\Theta_D(a)\coloneqq \prod_{b\in D}\chi_b(a)$, $\chi_b(a)\coloneqq \mathbf 1\{| a_b|_{\sigma}< r\}$; off that support the term vanishes, and nothing is claimed for it. (This nonemptiness is used for $I^{\Lambda}_t>0$ in Lemma~\ref{lem:c5-4-0}(c), for $\mathbf 1_{\mathcal D^{\Lambda}}(B+tm)=\Theta^t(B)$ in \eqref{eq:c5-4-2}, and for the identity hard cutoff in Lemma~\ref{lem:c5-4-4}.) The mean shift of the retreated split is $m\coloneqq -H^{-1}\ell_{\Lambda}$ (Section~\ref{sec:c5-1} with $D(\Lambda)$, $J(\Lambda)$), and by Lemma~\ref{lem:c5-1-2} (its proof uses (I1) and (I3) for the principal block $H_{D(\Lambda)D(\Lambda)}$, which hold by principal restriction, and (I2) for the rows $b\in D(\Lambda)\subset I$, Section~\ref{sec:c5-0} and Remark~\ref{rem:c5-3-2}) $m$ depends on $A^{\Lambda,\Psi}|_{J_1(\Lambda)}$ only, $J_1(\Lambda)\subset J_{\mathrm{col}}(\Lambda)$. The final collar is good (Lemma~\ref{lem:c5-3-1}(b)), $| A^{\Lambda,\Psi}(e)|_{\sigma}=| A^{\Lambda}_{\mathrm e}(e)|_{\sigma}<\vartheta r$ for every $e\in J_{\mathrm{col}}(\Lambda)$ ($\Psi_{\xi}$ is the identity there, Lemma~\ref{lem:c5-2-3}(a); on $I\setminus D(\Lambda)$ by \eqref{eq:c-collar-distance}), so \eqref{eq:c5-1-2} of Remark~\ref{rem:c5-1-3} (the choice $c_{\mathrm h}C_H^{\mathrm{near}}\vartheta\leqslant 1/4$), applied at $D(\Lambda)$ by Remark~\ref{rem:c5-3-2}, gives
\begin{equation}\label{eq:c5-4-1}
    \sup_{b\in D(\Lambda)}| m_b|_{\sigma}\leqslant r/4\;,\qquad\text{hence}\qquad\sup_bt| m_b|_{\sigma}\leqslant r/4\quad\text{for every }t\in[0,1].
\end{equation}
Completing the square (Section~\ref{sec:c5-1}): $-\tfrac12\langle a,Ha\rangle-t\langle a,\ell_{\Lambda}\rangle=-\tfrac12\langle a-tm,H(a-tm)\rangle+\tfrac12t^2\langle m,Hm\rangle$.

\begin{lemma}[the final term is a branch for the estimate lemmas]\label{lem:c5-4-0}
(a) $H= H_{D(\Lambda)D(\Lambda)}$ satisfies \eqref{eq:fixed-kernel-inputs}, $\gamma\,\id\leqslant H\leqslant h_+\id$ with $\gamma=\gamma_L/2$, $h_+= m_L$, and the entry range $\rho_L\leqslant R$: this is (I1), (I3) for the principal restriction (Section~\ref{sec:c5-0}). Hence Lemmas~\ref{lem:walk-series}, \ref{lem:conditioning}, \ref{lem:component-blocks} and~\ref{lem:gaussian-payment} hold for $H$ with the constants $n_*$, $R_*$, $C_w$, $\delta_M$, $A_*$, $C_d$, $\alpha_G$ of Section~\ref{sec:sourced-effective-interaction} and Appendix~\ref{app:background-estimates} (these lemmas use only \eqref{eq:fixed-kernel-inputs}, the range and the partitions, as the proof of Lemma~\ref{lem:c2-reference} records). (b) Retained inputs near the carrier. Let $r^{\Lambda}_{\mathrm e}\coloneqq \max\{| A^{\Lambda,\Psi}(b)|_{\sigma}\, : \,b\in J(\Lambda),\ \operatorname{dist}(b,D(\Lambda))\leqslant 3R_{\mathrm c}\}$ (zero if the set is empty), the retreated split's $r_{\mathrm e,\xi}$ (\eqref{eq:convex-native-radius}, Section~\ref{sec:c3-1}). Then $r^{\Lambda}_{\mathrm e}< r$: a bond $b\in J$ within $3R_{\mathrm c}$ of $D(\Lambda)\subset I$ is within $3R_{\mathrm c}$ of $I$, so by Lemma~\ref{lem:collar-coverage}(i),(iii) for the original term it lies in $\mathcal G_{\xi}\setminus S_1^*$, $\Psi_{\xi}$ is the identity at $b$ and $| A_{\mathrm e}(b)|_{\sigma}< r$ on the branch; a bond $b\in I\setminus D(\Lambda)$ carries $a_{\mathrm{old}}(b)$ with $| a_{\mathrm{old}}(b)|_{\sigma}< r$ (its box factor is a factor of the retained cutoff of the term, Lemma~\ref{lem:c5-2-5}) and $\Psi_{\xi}$ is the identity on $I$ (\eqref{eq:c-collar-distance}). So $s^{\Lambda}\coloneqq 1+r+r^{\Lambda}_{\mathrm e}\leqslant 1+2r$. On the collar $J_{\mathrm{col}}(\Lambda)$ (range $\rho_L=4L\leqslant 3R_{\mathrm c}$) the sharper bound $<\vartheta r$ of Lemma~\ref{lem:c5-3-1}(b) holds; it is used only through \eqref{eq:c5-4-1}. (c) Charts and primitive bounds. A row within $2R_{\mathrm c}+2L$ of $D(\Lambda)$ is within $2R_{\mathrm c}+2L$ of $I$, so the applicability paragraph of Appendix~\ref{app:background-estimates} (Lemmas~\ref{lem:c2-plaquette-deviation}, \ref{lem:c2-local-gauge}, \ref{lem:c2-uniform-chart}, which need only that the row's two-block set lies within $3R_{\mathrm c}$ of $I$ and the retained inputs there are below $r$, supplied by (b)) applies to it; hence Lemma~\ref{lem:c2-primitive-bounds} holds for the near primitives of the retreated split under threshold item 3 of Appendix~\ref{app:thresholds}, $4C_{\mathsf C}g_0(1+2r)\leqslant \vartheta_*/2$, with $s_{\xi}$ replaced by $s^{\Lambda}\leqslant 1+2r$: every near primitive $\mathfrak p$ is, at the input $a+tA^{\Lambda,\Psi}$, the restriction to real arguments of a function analytic in $(a,t)$ on a complex neighbourhood of the enlarged box $\{| a_b|_{\sigma}\leqslant 2r\ \forall b\in D(\Lambda)\}\times[0,1]$, with $|\mathfrak p(a+tA^{\Lambda,\Psi})|\leqslant C_pg_0(s^{\Lambda})^3$ for the plaquette and gauge primitives and the smaller bounds $C_pg_0^2(s^{\Lambda})^2$, $C_pg_0s^{\Lambda}$ for the Haar and Jacobian primitives ($g_0\leqslant 1\leqslant s^{\Lambda}$), and the derivative bounds \eqref{eq:convex-primitive-derivatives} with $s^{\Lambda}$; and Lemma~\ref{lem:charts} holds at $D(\Lambda)$, so the exponent of $K^{\Lambda}_{\xi;t}$ is real and finite on the open box and $K^{\Lambda}_{\xi;t}>0$ there, $I^{\Lambda}_t>0$. (d) Hosts. The host $Y_{\mathfrak p}$ of a primitive \eqref{eq:part1-potential-host} and the probe host $Y_b$ of a bond \eqref{eq:c-probe-host} are defined from the lifted stencil and $M$ alone; the localisation family of the retreated split is $\mathbf D_0(\Lambda)\coloneqq \{ Y_{\mathfrak p}\, : \,\mathfrak p\text{ near relative to }D(\Lambda)\}\cup\{ Y_b\, : \,b\in D(\Lambda)\}$, with $E(Y)\coloneqq \{ b\in D(\Lambda)\, : \,b\cap Y\neq\varnothing\}$, and the free-input set of a primitive assigned to $Y$ lies in $Y$ (\eqref{eq:part1-potential-host}), its young part in $E(Y)$. The counts of Lemma~\ref{lem:c2-potential-size}, at most $N_h$ primitives per host and $d_0(Y)\leqslant B_h+4$, are counts over templates and anchor positions, independent of the carrier. Hence, for $Y\in\mathbf D_0(\Lambda)$, $t\in[0,1]$ and $a$ in the local box $\{| a_b|_{\sigma}< r\ \forall b\in E(Y)\}$,
\begin{gather}
    \mathbf V^{\Lambda}_{\xi,t}(Y,a)\coloneqq \sum_{\mathfrak p\text{ near},\,Y_{\mathfrak p}= Y}\mathfrak p(a+tA^{\Lambda,\Psi})\quad\text{satisfies}\notag\\
    |\mathbf V^{\Lambda}_{\xi,t}(Y,a)|\leqslant N_hC_pg_0(s^{\Lambda})^3\leqslant \omega_te^{-\kappa d_0(Y)}\;,\qquad\omega_t\coloneqq C_0g_0(1+2r)^3\;,\label{eq:c5-4-4}
\end{gather}
by (c), $d_0(Y)\leqslant B_h+4$ and $C_0= N_hC_pe^{\kappa(B_h+4)}$ (\eqref{eq:part1-host-multiplicity}); $\mathbf V^{\Lambda}_{\xi,t}(Y,\bigcdot )\coloneqq 0$ off the local box and on hosts carrying no near primitive. For $t=0$ and every $\Lambda$ the input is $a$ alone, with $| a_b|_{\sigma}< r$ on the local box, so (c) gives the sharper bound $|\mathbf V^{\Lambda}_{\xi,0}(Y,a)|\leqslant \omega(g_0)e^{-\kappa d_0(Y)}$ with $\omega(g_0)= C_0g_0(1+r)^3\leqslant \omega_t$ ($s=1+r$ in Lemma~\ref{lem:c2-primitive-bounds}); for $\Lambda=\Lambda_1$ this is the $\mathbf V_{\xi}$ of Section~\ref{sec:sourced-effective-interaction} with its bound $\omega(g_0)$. The factor $(1+2r)^3$ in place of $(1+r)^3$ is the retained input, which enters for $t>0$.
\end{lemma}

\subsubsection{The translated objects}
\label{sec:c5-4-2}
Let $B\coloneqq a-tm$ be the centred variable, $a= B+tm$. Put, for $b\in D(\Lambda)$ and $Y\in\mathbf D_0(\Lambda)$,
\begin{equation*}
    \chi^t_b(B)\coloneqq \chi_b(B+tm)=\mathbf 1\{| B_b+tm_b|_{\sigma}< r\}\;,\qquad\Theta^t(B)\coloneqq \prod_{b\in D(\Lambda)}\chi^t_b(B)\;,\qquad\mathbf W_t(Y,B)\coloneqq \mathbf V^{\Lambda}_{\xi,t}(Y,B+tm);
\end{equation*}
$\chi^t_b$ is a Borel function of $B_b$ alone with values in $\{0,1\}$, $\mathbf W_t(Y,\bigcdot )$ depends on $B$ through $B|_{E(Y)}$ only ($m$ is a fixed vector), and $\mathbf W_t(Y,B)=0$ unless $| B_b+tm_b|_{\sigma}< r$ for all $b\in E(Y)$ (the translated local box), where $|\mathbf W_t(Y,B)|\leqslant \omega_te^{-\kappa d_0(Y)}$ by \eqref{eq:c5-4-4}. Translation invariance of Lebesgue measure gives, for every bounded Borel $F$ on $\mathfrak g^{D(\Lambda)}$,
\begin{equation}\label{eq:c5-4-2}
    \int F(a)K^{\Lambda}_{\xi;t}(a)\,\dd a= e^{\frac12t^2\langle m,Hm\rangle}\int F(B+tm)\,\Theta^t(B)\exp\Bigl[-\tfrac12\langle B,HB\rangle+\sum_{Y\in\mathbf D_0(\Lambda)}\mathbf W_t(Y,B)\Bigr]\dd B\;,
\end{equation}
since $\mathbf 1_{\mathcal D^{\Lambda}}(B+tm)=\Theta^t(B)$ and, where $\Theta^t(B)=1$, $\sum_Y\mathbf W_t(Y,B)=\sum_{\mathfrak p\text{ near}}\mathfrak p(B+tm+tA^{\Lambda,\Psi})= P^{\Lambda}_t(B+tm)$ (each near primitive is assigned to exactly one host). The constant $e^{\frac12t^2\langle m,Hm\rangle}$ is the same in the numerator and the denominator of every conditional expectation,
\begin{equation}\label{eq:c5-4-3}
    \int F\,\dd\nu^{\Lambda}_t=\frac{\int F(B+tm)\,\Theta^te^{-\frac12\langle B,HB\rangle+\sum_Y\mathbf W_t(Y,B)}\,\dd B}{\int\Theta^te^{-\frac12\langle B,HB\rangle+\sum_Y\mathbf W_t(Y,B)}\,\dd B}\;,
\end{equation}
and is never estimated: its logarithm is at most $\tfrac12h_+\sum_bh(m_b)^2\leqslant h_+| D(\Lambda)| r^2/32$ ($h\leqslant |\bigcdot |_{\sigma}$, $t\leqslant 1$), a constant of the order of the carrier volume, which disappears in \eqref{eq:c5-4-3}. Every use of the translated expansion below is through ratios of the form \eqref{eq:c5-4-3} or through quantities normalised by $\Xi^{f,t}_{\Lambda}$ (Corollary~\ref{cor:c5-4-3}).

\begin{lemma}[translated polymer representation]\label{lem:c5-4-1}
With $\mathfrak n_{\Lambda}\coloneqq (2\pi)^{d_G| D(\Lambda)|/2}\det(H)^{-1/2}$,
\begin{gather*}
    \int\Theta^t(B)\exp\Bigl[-\tfrac12\langle B,HB\rangle+\sum_Y\mathbf W_t(Y,B)\Bigr]\dd B=\mathfrak n_{\Lambda}\,\Xi^{f,t}_{\Lambda}\;,\qquad\Xi^{f,t}_{\Lambda}\coloneqq \sum_{\mathcal P\text{ compatible}}\prod_{Z\in\mathcal P}k^{f,t}_{\Lambda}(Z)\;,\\
    k^{f,t}_{\Lambda}(Z)\coloneqq \sum_{(\mathcal D,\mathsf P)\, : \,Z_0\neq\varnothing,\ Z_0'\subseteq Z}\mathcal A^{f,t}_{(\mathcal D,\mathsf P,Z)}[1]\;,
\end{gather*}
where the selections $\mathfrak s=(\mathcal D,\mathsf P,Z)$ are those of Definition~\ref{def:selection} for the carrier $D(\Lambda)$ and the family $\mathbf D_0(\Lambda)$ ($\mathcal D\subseteq\mathbf D_0(\Lambda)$ finite, $\mathsf P\subseteq D(\Lambda)\setminus E(Y_0)$ finite, seed $Z_0$, $E_0= E(Z_0)$, enlargement $Z_0'$, final domain $Z\supseteq Z_0'$ a union of $LM$-cubes), and the translated selected integral of a bounded Borel $F$ on $\mathfrak g^{E_0}$ is \eqref{eq:c-selected-integral} with $\chi_b$ replaced by $\chi^t_b$ and $\mathbf W$ by $\mathbf W_t$:
\begin{multline*}
    \mathcal A^{f,t}_{\mathfrak s}[F]=\prod_{\Delta\in\pi_1(Z)\setminus\pi_1(Z_0')}\int_0^1\dd s_{\Delta}\oint_{|\sigma_{\Delta}|= e^{\kappa_1}}\frac{\dd\sigma_{\Delta}}{2\pi i(\sigma_{\Delta}-s_{\Delta})^2}\prod_{Y\in\mathcal D}\int_0^1\dd t_Y\int\dd\mu_0(X)\int\dd\mu_{C_{E_0}}(B)\\
    \times e^{-\frac12\langle\Gamma_{E_0}(\sigma)X,C_{E_0}\Gamma_{E_0}(\sigma)X\rangle}e^{-\langle B,\Gamma_{E_0}(\sigma)X\rangle}(-1)^{|\mathsf P|}\prod_{b\in E(Y_0)}\chi^t_b(B)\prod_{b\in\mathsf P}(1-\chi^t_b(B))\\
    \times\prod_{Y\in\mathcal D}\mathbf W_t(Y,B)e^{t_Y\mathbf W_t(Y,B)}\,F(B)\;,
\end{multline*}
with $C_{E_0}$, $\Gamma_{E_0}(\sigma)$, $\mu_0$, $\mu_{C_{E_0}}$ built from $H= H_{D(\Lambda)D(\Lambda)}$. Consequently $\Xi^{f,t}_{\Lambda}= e^{-\frac12t^2\langle m,Hm\rangle}I^{\Lambda}_t/\mathfrak n_{\Lambda}>0$. For each selection $\mathfrak s$ and each bounded Borel $F$ on $\mathfrak g^{E_0}$ the translated selected integral $\mathcal A^{f,t}_{\mathfrak s}[F]$ is a well-defined finite complex number. No representation of the numerator of \eqref{eq:c5-4-3} for an insertion is asserted in this lemma: the insertion of a function of the young variables requires the probe-host (forced-seed) expansion of Section~\ref{sec:probe-kernels} for the $t$-family, Section~\ref{sec:c5-5}.
\end{lemma}

\begin{proof}
The proof of Lemma~\ref{lem:polymer} is reproduced with $\chi_b\to\chi^t_b$, $\mathbf W\to\mathbf W_t$, $I\to D(\Lambda)$, $\mathbf D_0\to\mathbf D_0(\Lambda)$. The properties of the cutoff and of the potential family that proof uses are: (i) $\Theta=\prod_b\chi_b$ with each $\chi_b$ a Borel function of $a_b$ alone with values in $[0,1]$: in the expansion $\Theta=\prod_{E(Y_0)}\chi_b\prod_{I\setminus E(Y_0)}(1-(1-\chi_b))$ with the sign $(-1)^{|\mathsf P|}$, and in the factorisation over the components of $Z$, where a bad bond lies in the component containing its incident cubes and a cutoff factor is a function of one bond variable; $\chi^t_b$ has these properties. (ii) $\mathbf W(Y,\bigcdot )$ depends on $a$ through $a|_{E(Y)}$, with $E(Y)\subseteq E_0$ for $Y\in\mathcal D$, and is bounded where $\prod_{b\in E(Y)}\chi_b\neq0$: in the reduction to $a|_{E_0}$ that lets Lemma~\ref{lem:conditioning} produce \eqref{eq:fixed-kernel-density}, and in the Fubini steps (bounded factors times an absolutely integrable density); $\mathbf W_t$ has these properties, with the translated local box in place of the local box, by Section~\ref{sec:c5-4-2}. (iii) Everything else, the walk series (Lemma~\ref{lem:walk-series}), the conditioning identity and the weakened density (Lemma~\ref{lem:conditioning}, \eqref{eq:fixed-kernel-density}), the block structure (Lemma~\ref{lem:component-blocks}), the integrability of the absolute density uniformly on the closed polydisc (Lemma~\ref{lem:gaussian-payment}(ii)), the Cauchy representation of the $\sigma$-derivatives and the vanishing of unseeded components, and the geometry of seeds, enlargements and the coarse cover, involves $H$, the partitions $\pi_0$, $\pi_1$ of the torus and $\sigma$ only, and holds for $H= H_{D(\Lambda)D(\Lambda)}$ by Lemma~\ref{lem:c5-4-0}(a); the host and probe geometry is that of Lemma~\ref{lem:c5-4-0}(d). The expansion of $\exp(\sum_Y\mathbf W_t)$ into Mayer factors $\int_0^1\mathbf W_te^{t_Y\mathbf W_t}\,\dd t_Y$ is the same algebra. Positivity: by \eqref{eq:c5-4-2} with $F=1$, $\mathfrak n_{\Lambda}\Xi^{f,t}_{\Lambda}= e^{-\frac12t^2\langle m,Hm\rangle}I^{\Lambda}_t$, and $I^{\Lambda}_t>0$ by Lemma~\ref{lem:c5-4-0}(c). Well-definedness of $\mathcal A^{f,t}_{\mathfrak s}[F]$: the integrand is a bounded function ($| F|\leqslant \sup| F|$, cutoffs $\leqslant 1$, $|\mathbf W_te^{t_Y\mathbf W_t}|\leqslant \omega_te^{\omega_t}$ by \eqref{eq:c5-4-4}) times the density \eqref{eq:fixed-kernel-density}, whose modulus is integrable uniformly on the closed polydisc by Lemma~\ref{lem:gaussian-payment}(ii) for $H$ (Lemma~\ref{lem:c5-4-0}(a)); the finitely many $s$-, $\sigma$- and $t_Y$-integrals are over compact sets.
\end{proof}

\begin{lemma}[translated selected majorant]\label{lem:c5-4-2}
Let $g_0$ satisfy the thresholds of Appendix~\ref{app:thresholds}, items 1--5. Two of their consequences are used here:
\begin{align}
    &C_0g_0(1+2r)^3\leqslant \omega_*\quad\text{(i.e. $\omega_t\leqslant \omega_*$; threshold item 4(ii) gives $\omega_t\leqslant \omega_*/2$)},\label{eq:c5-4-T}\\
    &\tfrac{9}{16}a_{\mathrm{bad}}\geqslant \max\bigl(160(\kappa+\log(8M^4+8)+1),\,64\log(1/\eta_*)\bigr)\;,\qquad a_{\mathrm{bad}}= c_n\alpha_Gr^2/4\notag\\
    &\qquad\text{(\eqref{eq:native-source-bad-budget} for $\tfrac9{16}a_{\mathrm{bad}}$; threshold item 4(i), which states it for $\tfrac9{128}a_{\mathrm{bad}}$)}.\label{eq:c5-4-Tp}
\end{align}
Both are conditions on $g_0$ alone ($g_0(1+r)^3\to0$ and $a_{\mathrm{bad}}\to\infty$ as $g_0\to0$, $r= A_1(\log g_0^{-2})^{p_0}$), with thresholds depending on $L$, the group, $M$, $\kappa$, $\alpha_G$, $\omega_*$ and the cutoff parameters $A_1$, $p_0$ only (through $r(g_0)$). Then on every labelled term with final region $\Lambda$, for exterior data in the support of the term's collar indicators and exterior cutoffs (where \eqref{eq:c5-4-1} and $r^{\Lambda}_{\mathrm e}< r$ hold, Lemmas~\ref{lem:c5-3-1}(b) and~\ref{lem:c5-3-3}(iii); the term vanishes elsewhere), for every $t\in[0,1]$, every selection $\mathfrak s=(\mathcal D,\mathsf P,Z)$ for $D(\Lambda)$ and every bounded Borel $F$ on $\mathfrak g^{E_0}$,
\begin{equation}\label{eq:c5-4-5}
    |\mathcal A^{f,t}_{\mathfrak s}[F]|\leqslant \sup| F|\cdot e^{-(\kappa_1-1)N_1(Z\setminus Z_0')+N_1(Z)-\frac9{16}a_{\mathrm{bad}}|\mathsf P|}\prod_{Y\in\mathcal D}\omega_*e^{1/2}e^{-\kappa d_0(Y)}\;.
\end{equation}
The right side depends neither on $t$, nor on $m$ or the exterior data $A^{\Lambda,\Psi}$ within that support, nor on the labels; bondwise, the reserve is $(1-t| m_b|_{\sigma}/r)^2a_{\mathrm{bad}}\geqslant \tfrac9{16}a_{\mathrm{bad}}$.
\end{lemma}

\begin{proof}
The proof of Lemma~\ref{lem:c2-selected-majorant} with the substitutions $\chi_b\to\chi^t_b$, $\mathbf W\to\mathbf W_t$ and $I\to D(\Lambda)$ with $\mathbf D_0\to\mathbf D_0(\Lambda)$. The weakening step is unchanged: for each cube $\Delta$ of $Z\setminus Z_0'$ the operation $G\mapsto\int_0^1\dd s\oint G\,\dd\sigma_{\Delta}/(2\pi i(\sigma_{\Delta}-s)^2)$ on the inner integral as a function of $\sigma_{\Delta}$ on the circle $|\sigma_{\Delta}|= e^{\kappa_1}$ is bounded by $e^{-(\kappa_1-1)}$ times the supremum of the inner integral's modulus over the distinguished boundary, by $e^{\kappa_1}(e^{\kappa_1}-1)^{-2}\leqslant e^{-(\kappa_1-1)}$ ($e^{\kappa_1}\geqslant 3$); it needs the continuity of the integrand in $\sigma$ (the walk series, Lemma~\ref{lem:walk-series} for $H$) and the uniform integrability below. On the distinguished boundary, the inner integrand's modulus is bounded factor by factor: $\int_0^1\dd t_Y$ has mass one; $\chi^t_b\leqslant 1$; for $b\in\mathsf P$ and real $B$, by the triangle inequality for the norm $|\bigcdot |_{\sigma}$ and \eqref{eq:c5-4-1},
\begin{equation*}
    1-\chi^t_b(B)=\mathbf 1\{| B_b+tm_b|_{\sigma}\geqslant r\}\leqslant \mathbf 1\{| B_b|_{\sigma}\geqslant r-t| m_b|_{\sigma}\}\leqslant \mathbf 1\{| B_b|_{\sigma}\geqslant (1-t| m_b|_{\sigma}/r)r\}\leqslant \mathbf 1\{| B_b|_{\sigma}\geqslant 3r/4\}\;,
\end{equation*}
and on $\{| X|_{\sigma}\geqslant 3r/4\}$ the group inequality $h(X)^2\geqslant c_n| X|_{\sigma}^2$ of Lemma~\ref{lem:gaussian-payment}(iii) ($c_n=1/N$ serves for both cutoff norms, the conventions of Appendix~\ref{app:background-estimates}) gives $h(X)^2\geqslant \tfrac9{16}c_nr^2$, so $\mathbf 1\{| B_b|_{\sigma}\geqslant 3r/4\}\leqslant \exp[(\gamma_2/2)(h(B_b)^2-\tfrac9{16}c_nr^2)]$ for every $\gamma_2>0$ (the exponent is nonnegative on the set); with $\gamma_2=\alpha_G/2$, $a_{\mathrm{bad}}= c_n\gamma_2r^2/2= c_n\alpha_Gr^2/4$ and $\| B\|^2=\sum_bh(B_b)^2$,
\begin{equation*}
    \prod_{b\in\mathsf P}(1-\chi^t_b(B))\leqslant e^{-\frac9{16}a_{\mathrm{bad}}|\mathsf P|}e^{(\alpha_G/2)\sum_{b\in\mathsf P}h(B_b)^2}\leqslant e^{-\frac9{16}a_{\mathrm{bad}}|\mathsf P|}e^{(\alpha_G/2)\| B\|^2}
\end{equation*}
(this is Lemma~\ref{lem:gaussian-payment}(iii) at the radius $3r/4$ in place of $r$; bondwise the same argument gives the reserve $(1-t| m_b|_{\sigma}/r)^2a_{\mathrm{bad}}$); for $Y\in\mathcal D$, on the support of $\prod_{b\in E(Y)}\chi^t_b$ the point $B+tm$ lies in the local box of $Y$, so $|\mathbf W_t(Y,B)|\leqslant \omega_te^{-\kappa d_0(Y)}\leqslant \omega_*e^{-\kappa d_0(Y)}$ by \eqref{eq:c5-4-4} and \eqref{eq:c5-4-T}, and $|\mathbf W_te^{t_Y\mathbf W_t}|\leqslant \omega_*e^{\omega_*}e^{-\kappa d_0(Y)}\leqslant \omega_*e^{1/2}e^{-\kappa d_0(Y)}$, since $\omega_*\leqslant 1/2$ (the standing choice of Appendix~\ref{app:thresholds}); $| F|\leqslant \sup| F|$; and the integral of the modulus of the density \eqref{eq:fixed-kernel-density} against $e^{(\alpha_G/2)\| B\|^2}$ is at most $e^{N_1(Z)}$ by Lemma~\ref{lem:gaussian-payment}(ii) for $H$ (Lemma~\ref{lem:c5-4-0}(a)), under the conditions on $\delta_M$ and $\alpha_G$ of Appendix~\ref{app:thresholds}. Multiplying the bounds and applying the weakening bound once per cube of $Z\setminus Z_0'$ gives \eqref{eq:c5-4-5}.
\end{proof}

\begin{corollary}[activity bounds and vacuum clusters for the $t$-family]\label{cor:c5-4-3}
Under the hypotheses of Lemma~\ref{lem:c5-4-2}, for every $t\in[0,1]$ and every connected final domain $Z$ for $D(\Lambda)$,
\begin{equation*}
    | k^{f,t}_{\Lambda}(Z)|\leqslant \varepsilon_*e^{-\lambda_*\ell_1(Z)}\;,\qquad\sum_{Z'\not\sim Z}| k^{f,t}_{\Lambda}(Z')| e^{(1+6\theta)N_1(Z')}\leqslant N_1(Z)/8\;,
\end{equation*}
with the constants $\lambda_*$, $\varepsilon_*$, $\theta$ of \eqref{eq:c2-kappa-choice}--\eqref{eq:native-source-vacuum-KP-budget}, and Lemma~\ref{lem:vacuum-clusters} holds for the cluster weights $\Phi^{v,t}_{\Lambda}(\mathfrak c)=(\varphi^T(\mathfrak c)/k!)\prod_{j=1}^kk^{f,t}_{\Lambda}(D_j)$ of a cluster $\mathfrak c=(D_1,\dots,D_k)$ (the symbol $m$ is the shift here): the hitting sum $\sum_{\mathfrak c\text{ hits }\mathcal A}|\Phi^{v,t}_{\Lambda}(\mathfrak c)| e^{2\theta\sum_jN_1(D_j)}\leqslant \tfrac18\sum_{D\in\mathcal A}N_1(D)$ and the identity $(\Xi^{f,t}_{\Lambda})^{[\cup\mathcal A]}/\Xi^{f,t}_{\Lambda}=\exp(-\sum_{\mathfrak c\text{ hits }\mathcal A}\Phi^{v,t}_{\Lambda}(\mathfrak c))$, uniformly in $t$, $m$, the exterior data in the support of the term's collar indicators and exterior cutoffs, and the labels.
\end{corollary}

\begin{proof}
The proof of Lemma~\ref{lem:activity-bounds} uses, about the activities, only the majorant \eqref{eq:fixed-kernel-selected-majorant}, here \eqref{eq:c5-4-5} with $F=1$, with $\omega_*e^{1/2}e^{-\kappa d_0(Y)}\leqslant \eta_*e^{-\kappa\ell_0(Y)}$ ($\ell_0= d_0+5$ and $\eta_*=2e^{1/2+5\kappa}\omega_*$), and the number $a_{\mathrm{bad}}$ through the conditions it draws from \eqref{eq:native-source-bad-budget}: $a_{\mathrm{bad}}\geqslant 160(\kappa+1)>1$ and $64M^4e^{-a_{\mathrm{bad}}/4}\leqslant 1$ (bad-label sum), $a_{\mathrm{bad}}/32\geqslant 5\kappa$ (component reserve), $e^{-a_{\mathrm{bad}}/64}\leqslant \eta_*\leqslant t_0$ (the case $k_j=0$), and $C_ge^{-a_{\mathrm{bad}}/64}\leqslant 162e^{-5(\kappa+1)/2}\leqslant \kappa/40$ (choice of $Y_0$). Each of these holds for every number not smaller than one satisfying \eqref{eq:native-source-bad-budget}; \eqref{eq:c5-4-Tp} says that $\tfrac9{16}a_{\mathrm{bad}}$ satisfies it, so the proof applies verbatim with $a_{\mathrm{bad}}$ replaced by $\tfrac9{16}a_{\mathrm{bad}}$ and $\eta_*$, $t_0$, $\omega_*$, $\kappa$, $\delta$, $\alpha_6$, $C_g$, $M$, $L$ unchanged (the geometric inputs Lemmas~\ref{lem:c2-domain-geometry}, \ref{lem:c2-cover-sum} and~\ref{lem:c2-scale} concern the partitions only; the cover sum is over $\mathbf D_0(\Lambda)\subseteq\mathbf D_0$, every host of the retreated family being a host of the original one). The proof of Lemma~\ref{lem:vacuum-clusters} (Appendix~\ref{app:kp}, Kotecký--Preiss with the pinned-tree bound) uses only the two displayed bounds and the geometry, and nothing else about the branch (proof of Lemma~\ref{lem:c2-reference}).
\end{proof}

\begin{lemma}[the conditional estimate at the retreated carrier]\label{lem:c5-4-4}
Under the thresholds $C_Pg_0(1+2r)\leqslant \gamma/2$ and $r\geqslant 1$ of Appendix~\ref{app:exterior-chessboard} (\eqref{eq:c3-1-2}; Lemma~\ref{lem:c3-3-2} uses $r\geqslant 1$ in $O(g_0s^2/r)= O(g_0s)$; both are conditions on $g_0$), threshold item 3 of Appendix~\ref{app:thresholds} and the hypotheses listed in Lemma~\ref{lem:c5-4-0}, on every labelled term with final region $\Lambda$, for exterior data in the support of the term's collar indicators and exterior cutoffs, and on the covered event $A^{\mathrm{fin}}_{\mathrm{cov}}$ of Section~\ref{sec:c5-3-4}(a) (every contour bond in $\Omega_1$, every free input of the contour in $D(\Lambda)$),
\begin{gather*}
    |\mathcal B^{\Lambda}_{\xi}[f_{\xi}]|\leqslant C_{\mathrm{resp},\chi}g_0rD^{\Lambda}_{\xi}\;,\qquad\mathcal B^{\Lambda}_{\xi}[f]\coloneqq \int f\,\dd\nu^{\Lambda}_1-\int f\,\dd\nu^{\Lambda}_0=\int_0^1\operatorname{Cov}_{\nu^{\Lambda}_t}(f,L^{\Lambda}_t)\,\dd t\;,\\
    D^{\Lambda}_{\xi}\coloneqq \sum_{b\in I_C\cap D(\Lambda)}e^{-(\mu/2)d(b,J(\Lambda))}\;,
\end{gather*}
with the constant $C_{\mathrm{resp},\chi}$ and the rate $\mu= R^{-1}\log(1+\gamma/(4M_*))$, $M_*= m_L+\gamma/2$, of Appendix~\ref{app:exterior-chessboard} (Lemmas~\ref{lem:c3-4-1} and~\ref{lem:exterior-response-bound}), independent of $g_0$, the term, the labels, the volume and $t$; $D^{\Lambda}_{\xi}$ is the $D^{\mathrm{fin}}_{\xi}$ of Section~\ref{sec:c5-3-4}(a), and $L^{\Lambda}_t\coloneqq -\langle\bigcdot ,\ell_{\Lambda}\rangle+\partial_tP^{\Lambda}_t$.
\end{lemma}

\begin{proof}
Sections~\ref{sec:c3-1}--\ref{sec:c3-5} are applied with $(I,J,A_{\mathrm e},H,\ell,P_t,r_{\mathrm e,\xi})$ replaced by $(D(\Lambda),\allowbreak J(\Lambda),\allowbreak A^{\Lambda,\Psi},\allowbreak H_{D(\Lambda)D(\Lambda)},\allowbreak\ell_{\Lambda},\allowbreak P^{\Lambda}_t,\allowbreak r^{\Lambda}_{\mathrm e})$. The inputs they use, each supplied: $\gamma\leqslant H\leqslant m_L$ and the finite range $R$ of $H$ and of the primitive stencils (Lemma~\ref{lem:c5-4-0}(a); Section~\ref{sec:c3-2} takes $R=8L+8\geqslant \rho_L$); the primitive bounds on the enlarged box at the inputs $a+tA^{\Lambda,\Psi}$, uniformly in $t$, with $s^{\Lambda}\leqslant 1+2r$ (Lemma~\ref{lem:c5-4-0}(c)), giving $\| D^2_aP^{\Lambda}_t\|\leqslant C_Pg_0s^{\Lambda}$ and the Schur row/column bounds of Lemma~\ref{lem:c3-3-2}; $r^{\Lambda}_{\mathrm e}< r$ for the score bound $h(\nabla_b\partial_tP^{\Lambda}_t)\leqslant C_Sg_0s^{\Lambda} r^{\Lambda}_{\mathrm e}\mathbf 1_{K^{\Lambda}}(b)$ and $h((\ell_{\Lambda})_b)\leqslant m_Lr^{\Lambda}_{\mathrm e}\mathbf 1_{K^{\Lambda}}(b)$, $K^{\Lambda}\coloneqq \{ b\in D(\Lambda)\, : \,\operatorname{dist}(b,J(\Lambda))\leqslant R\}$ (Lemma~\ref{lem:c5-4-0}(b) and (I2)); the source gradient bound $\sup h(\nabla_bf_{\xi})\leqslant C_{\mathrm{resp},\chi}g_0\mathbf 1_{I_C}(b)$ on the enlarged box, from the rowwise reconstruction and the uniform charts for the rows meeting $D(\Lambda)$ (Lemma~\ref{lem:c5-4-0}(c); the links of the contour depend on $a|_{D(\Lambda)}$ through those rows, the de-younged bonds being constants of the term); the identity $\mathcal B=\int_0^1\operatorname{Cov}\,\dd t$, which is the differentiation of $t\mapsto\int f\,\dd\nu^{\Lambda}_t$ under the integral (the exponent of $K^{\Lambda}_{\xi;t}$ is smooth in $t$ with bounded derivative on the box, Lemma~\ref{lem:c5-4-0}(c)); the extension $T_r$ of Section~\ref{sec:c3-3}, the barrier $\beta_n$ and the Helffer--Sj\"ostrand identity of Section~\ref{sec:c3-4}, which involve only the box $\{| a_b|_{\sigma}< r,\ b\in D(\Lambda)\}$, the Hessian bounds and the Schur bounds; and the shell summation of Lemma~\ref{lem:exterior-response-bound} with $d(c,J(\Lambda))\leqslant R$ on $K^{\Lambda}$. On $A^{\mathrm{fin}}_{\mathrm{cov}}$ every free input of the contour lies in $D(\Lambda)$, so $I_C\subseteq D(\Lambda)$ and $D^{\Lambda}_{\xi}$ sums over all of $I_C$; the intersection in its definition records this. Nothing in Sections~\ref{sec:c3-1}--\ref{sec:c3-5} uses that the carrier is $\Lambda_1^*$.
\end{proof}

\begin{remark}[scope of this subsection]\label{rem:c5-4-5}
The translated expansion gives the vacuum activities $k^{f,t}_{\Lambda}$, their clusters and a majorant for selected integrals with a bounded insertion $F$ on $\mathfrak g^{E_0}$; the hosts are unchanged and the primitives are evaluated at $a= B+tm$ in the original local box, so the chart condition is threshold item 3 at $s^{\Lambda}\leqslant 1+2r$, and the enlarged box $\{| a_b|_{\sigma}\leqslant 2r\}$, not needed for the majorant, is available for analyticity, $B\mapsto B+tm$ mapping $\{| B_b|_{\sigma}\leqslant 7r/4\}$ into it. Nothing here concerns the source-containing connected terms attached to $J(\Lambda)$: $f_{\xi}(B+tm)$ is not a function of $B|_{E_0}$ for a selection unrelated to the contour, so \eqref{eq:c5-4-5} with $F$ the source is not available, and the insertion of the source through the probe hosts of $\mathbf D_0(\Lambda)$, the $t$-centred reference and the marked connected terms are Section~\ref{sec:c5-5}, whose result is \eqref{eq:c5-5-9}: $|\mathcal B^{\Lambda}_{\xi}[f_{\xi}]|$ is bounded by $| S^{\circ}_{\Lambda}| e^{2n_C\bar\varepsilon_t}$ times $C_1g_0r\,\mathfrak S_1(\Lambda)+C_2\bar\varepsilon_tr^{d_G}\mathfrak S_2(\Lambda)$, with $\mathfrak S_1$, $\mathfrak S_2$ the exponentially weighted sums over $J_1(\Lambda)$ and $J(\Lambda)$ defined there; the bookkeeping of the labels in the history sum is Section~\ref{sec:c5-6}. The consequences \eqref{eq:c5-4-T}, \eqref{eq:c5-4-Tp} of threshold item 4 are the form in which it is used wherever the $t$-family is used.
\end{remark}

\subsection{The source-containing connected boundary terms}
\label{sec:c5-5}

\subsubsection{Setting and target}
\label{sec:c5-5-0}
Throughout this subsection fix a history, a labelled term with final region $\Lambda$, carrier $D\coloneqq D(\Lambda)$, exterior $J(\Lambda)$, exterior variable $A^{\Lambda,\Psi}$, shift $m$, and the $t$-family $K^{\Lambda}_{\xi;t}$, $\nu^{\Lambda}_t$, $\mathbf W_t$, $\chi^t_b$ of Section~\ref{sec:c5-4}, with the exterior data in the support of the term's collar indicators and exterior cutoffs (where \eqref{eq:c5-4-1}, $r^{\Lambda}_{\mathrm e}< r$ and Lemma~\ref{lem:c5-4-0} hold; the term vanishes elsewhere, and every bound below is a bound on a factor of the term that is multiplied by those indicators), on the covered event $A^{\mathrm{fin}}_{\mathrm{cov}}$ of Section~\ref{sec:c5-3-4}(a) (every contour bond in $\Omega_1$, every free input of the contour in $D(\Lambda)$; so $I_C\subseteq D(\Lambda)$, and $f_{\xi}(a)= W_{\chi}(C;U_{\xi}(a))$ depends on $a\in\mathfrak g^{D(\Lambda)}$ through $a|_{I_C}$ only, the de-younged bonds $I\setminus D(\Lambda)$ being constants of the term, Lemma~\ref{lem:c5-4-4}). Write
\begin{equation*}
    F^{\Lambda}(t)\coloneqq \int f_{\xi}\,\dd\nu^{\Lambda}_t\;,\qquad\mathcal B^{\Lambda}_{\xi}[f_{\xi}]= F^{\Lambda}(1)-F^{\Lambda}(0)\quad\text{(Lemma~\ref{lem:c5-4-4})},
\end{equation*}
and $n_C$ for the number of traversals of $C$, $I_C$ for the set of free-input bonds of the contour (each within $\ell^{\infty}$ distance $2L+1$ of a contour bond: the forward influence set $\mathcal R_b$ of Section~\ref{sec:probe-kernels} lies within $\ell^{\infty}$ distance $2L+1$ of $b$, applied to the contour bond $r$ with $b\in\mathcal F_r$; hence $\operatorname{dist}(e,I_C)\geqslant \operatorname{dist}(e,C)-(2L+1)$ in the distance convention below), $J_1(\Lambda)\coloneqq \{ e\in J(\Lambda)\, : \,H(b',e)\neq0\text{ for some }b'\in D(\Lambda)\}\subseteq J_{\mathrm{col}}(\Lambda)$ (Lemma~\ref{lem:c5-1-2} at the retreated carrier; Section~\ref{sec:c5-4}: $m$ depends on $A^{\Lambda,\Psi}|_{J_1(\Lambda)}$ only, and $| A^{\Lambda,\Psi}(e)|_{\sigma}=| A^{\Lambda}_{\mathrm e}(e)|_{\sigma}<\vartheta r$ on $J_{\mathrm{col}}(\Lambda)$). The target is a bound on $\mathcal B^{\Lambda}_{\xi}[f_{\xi}]$ relative to the reference gas, with the perimeter factor extracted inside the connected terms before absolute values, of the form (reference)$\times\sum_{e\in J(\Lambda)}e^{-\mu'\operatorname{dist}(e,C)}\times$(polynomial in $r$)$\times$(history weights). The reference used here is $S^{\circ}_{\Lambda}$, the reference group gas at the carrier $D(\Lambda)$ (Lemma~\ref{lem:c5-5-3}); on $A^{\mathrm{fin}}_{\mathrm{cov}}$ it is identified with the regional square scalar at $D(\Lambda)$ by Lemma~\ref{lem:c5-7-1} and with $\kappa_Ce^{E^{\circ}_{\Lambda}}$ by Corollary~\ref{cor:c5-7-2}, under \eqref{eq:c5-7-T}; its comparison with the reference scalar at the carrier $I$ is Lemma~\ref{lem:c5-6-3}, whose ratio factor is not small on every history (Remark~\ref{rem:c5-6-5}(1)).

Inputs beyond Sections~\ref{sec:c5-0}--\ref{sec:c5-4}, used as stated at their stated hypotheses: Lemmas~\ref{lem:probe-measures}, \ref{lem:measure-identity}, Proposition~\ref{prop:centred-expectation} and Section~\ref{sec:joint-extraction} (Definition~\ref{def:groups}, Lemma~\ref{lem:grouped-source}, the reference data, Lemma~\ref{lem:schur}, Proposition~\ref{prop:joint-extraction}); Lemmas~\ref{lem:c4-4-1} and~\ref{lem:c4-7-1} (the dummy comparison \eqref{eq:c4-4-7}, $\sum_{\ell\, : \,Z(\ell)= Z,\,\mathsf J\neq\varnothing}u_*^{q(\ell)}M_{\ell}\leqslant m(Z)\coloneqq \varepsilon_*e^{-\lambda_*\ell_1(Z)}$ with $M_{\ell}=\sup_{|\zeta|\leqslant R_{\Upsilon}}\|\Lambda^{\mathrm f}_{\xi;\ell,\zeta}\|_{\mathrm{TV}}$ the $\zeta$-uniform total-variation majorant of Section~\ref{sec:c4-4}, $u_*=\min(1/4,\omega_*e^{-\kappa s_{\mathrm p}}/(8P_M))$ \eqref{eq:c4-4-6}, proved by the positive host comparison from the selected majorant and the cover/seed sums of Appendix~\ref{app:background-estimates}), Lemma~\ref{lem:c4-6-1} (the closed-contraction factorisation of $R_{\xi}(\mathfrak B)$ and the product inequality $| R_{\xi}(\mathfrak B)|\leqslant \prod_WK_{\chi}\mathfrak d_W$, $K_{\chi}$ the constant of the irreducible character); Lemma~\ref{lem:c4-6-3} and Section~\ref{app:kp} (the reference exclusion ratio $q^{\circ}(\mathfrak B)= S^{\circ[\mathfrak B]}/S^{\circ}=\exp(-\sum_{\Gamma\text{ hits }\mathfrak B}\Phi^{\circ}(\Gamma))$ by Kotecký--Preiss for the reference group gas, and the Kotecký--Preiss theorem in the rooted form of the proof of Lemma~\ref{lem:vacuum-clusters}, the recursion following \eqref{eq:marginal-pinned-bound} with \cite[Proposition~5.3, Theorem~5.4]{FriedliVelenik}: with activities $| a_W| e^{3\eta s(W)/2}$ and node cost $\eta s(W)/8$ it gives the rooted bound $\sum_{\Gamma\ni W}m_W(\Gamma)|\Phi^{\circ}(\Gamma)|\leqslant | a_W| e^{13\eta s(W)/8}$ and the exclusion identity); Lemma~\ref{lem:c4-6-2} (the positive domain/cluster/interval summation giving the footprint and contour pins $\sup_{\Delta}\sum_{W\, : \,\Delta\in F(W)}e^{\nu\sigma(W)}M(W)\leqslant C\varepsilon$ and $\sup_j\sum_{W\, : \,j\in\mathcal H(W)}e^{\nu\sigma(W)}M(W)\leqslant C\varepsilon$ for a positive group majorant $M(W)$ with root weights $\rho\,m(Z)$ and the positive cluster tree weights, $\nu$ the support reserve, $\sigma(W)$ the size, $F(W)$ the footprint; the positive source-containing group summation with the weights $e^{4\eta\mathfrak s(W)}$ under $C\rho LM\leqslant \eta/8$ and the root-incidence bound of Section~\ref{sec:c4-8-2}); Lemma~\ref{lem:c4-6-env} (the connected envelope $\mathfrak E(W)$ of a group has at most $A_E\sigma(W)$ cubes and is connected under maximum-norm adjacency, so a group whose footprint contains the cube $\Delta$ has $\sigma(W)\geqslant \operatorname{dist}(\Delta,C)/(A_ELM)$); the fact that incompatible groups have touching footprints or overlapping intervals is the compatibility of groups of Section~\ref{sec:joint-extraction} (after Definition~\ref{def:groups}: compatible groups have mutually compatible constituent domains and disjoint intervals) with the contact convention of Section~\ref{sec:c4-12-4}; the tree-graph justification following \eqref{eq:marginal-pinned-bound} ($|\varphi^T|\leqslant $ number of spanning trees of the incompatibility graph, used in Appendix~\ref{app:kp}); the far event $\mathsf N_b=\{\operatorname{dist}(\Gamma_C,E(\Lambda_1^c))\leqslant b\}$ of \eqref{eq:c4-2-1} (Lemma~\ref{lem:c4-2-1}), with $b=\lfloor LMR_1/16\rfloor$. The finite-difference bound of Section~\ref{sec:c4-4-1} is not used as an input: Lemma~\ref{lem:c5-5-7z} proves the form needed here from Lemmas~\ref{lem:charts} and~\ref{lem:c2-uniform-chart}.

Geometry. Throughout Sections~\ref{sec:c5-5} and~\ref{sec:c5-6} the square/torus geometry of Section~\ref{sec:extent-weights} is in force, as imposed by Section~\ref{sec:c4-6} and collected with the other torus conditions in Section~\ref{app:thresholds}: $C$ is a coarse square of side $s= nL$ fine bonds, $n_C=4s$, on fine periods $LK_i$ with $K_i\geqslant 3(n+2)$, together with the period condition of Section~\ref{sec:c5-3-1} (compatible even periods exceeding $200\ell$; when $n\geqslant 67MR_1$, $K_i\geqslant 3(n+2)$ already gives $LK_i\geqslant 3L(n+2)>200LMR_1=200\ell$, while for the fixed loop as $g_0\to0$ the condition on the periods is the binding one). The constants $C_{\mathrm{pin}}$, $A_F$, $A_E$ and the hull/extent bounds of Section~\ref{sec:c4-6-2} (Lemmas~\ref{lem:c4-6-env} and~\ref{lem:c4-6-2}) are proved under this geometry, and Lemmas~\ref{lem:c5-5-3}(d), \ref{lem:c5-5-8}, Proposition~\ref{prop:c5-5-9} and Corollary~\ref{cor:c5-5-10} inherit it as a hypothesis. Distance convention: every distance in Sections~\ref{sec:c5-5}--\ref{sec:c5-6} is the $\ell^{\infty}$ distance in fine units; the distance of a bond is measured from its initial site, that of a cube or host from its nearest point, and the distance to $C$ is the distance to the vertex set of $C$ (so $\operatorname{dist}(b,C)=0$ for a contour bond $b$, whose initial site is a vertex, and $\operatorname{dist}(\Delta,C)=0$ for a cube of $\pi_1(\mathcal H)$, the cubes containing vertices of the traversals of $\mathcal H$, Lemma~\ref{lem:c4-6-env}); the counts $4(2s+1)^4$ and the conversions by $LM+1$, $4M+LM$, $2L+1$ below are stated in this convention.

\subsubsection{The decay form of the mean shift}
Lemma~\ref{lem:c5-1-2} bounds $\sup_bh(m_b)$; the dependence on the distance to the exterior is needed.

\begin{lemma}\label{lem:c5-5-1}
For every $b\in D(\Lambda)$,
\begin{equation}\label{eq:c5-5-1}
    h(m_b)\leqslant C_m\vartheta r\sum_{e\in J_1(\Lambda)}e^{-\mu_Hd(b,e)}\;,\qquad C_m\coloneqq \gamma^{-1}m_L\cdot4(2\rho_L+1)^4e^{\mu_H\rho_L}\;,
\end{equation}
with $\mu_H=\gamma/(h_+\rho_L)$ of Lemma~\ref{lem:c5-1-1} and $d$ the $\ell^{\infty}$ distance of initial sites in fine units. Consequently, with $N_{\mathrm{lat}}(\mu)=\sum_{s\geqslant 0}4(2s+1)^4e^{-\mu s}$ of Lemma~\ref{lem:c5-1-2},
\begin{equation}\label{eq:c5-5-2}
    \sum_{b\in D(\Lambda)}h(m_b)e^{-\mu\operatorname{dist}(b,X)}\leqslant C_m\vartheta rN_{\mathrm{lat}}(\mu_H/2)\sum_{e\in J_1(\Lambda)}e^{-\min(\mu_H/2,\mu)\operatorname{dist}(e,X)}
\end{equation}
for every set $X$ of bonds or of sites ($\operatorname{dist}(b,X)$ the $\ell^{\infty}$ distance from the initial site of $b$; Proposition~\ref{prop:c5-5-9} takes $X$ the vertex set of $C$) and $\mu\geqslant 0$; in particular $\sum_{b\in I_C}h(m_b)\leqslant C_m\vartheta rN_{\mathrm{lat}}(\mu_H/2)\sum_{e\in J_1(\Lambda)}e^{-\mu_H\operatorname{dist}(e,I_C)/2}$.
\end{lemma}

\begin{proof}
$m_b=-\sum_{b'\in D(\Lambda)}H^{-1}(b,b')\sum_{e\in J_1(\Lambda)}H(b',e)A^{\Lambda,\Psi}_e$ (Section~\ref{sec:c5-1}, Lemma~\ref{lem:c5-1-2}). By (I1), $H(b',e)\neq0$ forces $d(b',e)\leqslant \rho_L$, hence $d(b,b')\geqslant d(b,e)-\rho_L$; by Lemma~\ref{lem:c5-1-1}, $\| H^{-1}(b,b')\|\leqslant \gamma^{-1}e^{-\mu_Hd(b,b')}\leqslant \gamma^{-1}e^{\mu_H\rho_L}e^{-\mu_Hd(b,e)}$ for such $b'$; the number of $b'$ with $d(b',e)\leqslant \rho_L$ is at most $4(2\rho_L+1)^4$; each entry has $\| H(b',e)\|\leqslant \sum_{b''}\| H(b',b'')\|\leqslant m_L$ by (I2) for the row $b'\in D(\Lambda)$; and $h(A^{\Lambda,\Psi}_e)\leqslant | A^{\Lambda,\Psi}_e|_{\sigma}<\vartheta r$ for $e\in J_1(\Lambda)\subseteq J_{\mathrm{col}}(\Lambda)$ (Lemma~\ref{lem:c5-3-1}(b)). Multiplying gives \eqref{eq:c5-5-1}. For \eqref{eq:c5-5-2}: $e^{-\mu_Hd(b,e)}e^{-\mu\operatorname{dist}(b,X)}\leqslant e^{-\mu_Hd(b,e)/2}e^{-\min(\mu_H/2,\mu)(d(b,e)+\operatorname{dist}(b,X))}\leqslant e^{-\mu_Hd(b,e)/2}e^{-\min(\mu_H/2,\mu)\operatorname{dist}(e,X)}$ by the triangle inequality, and $\sum_{b\in D(\Lambda)}e^{-\mu_Hd(b,e)/2}\leqslant N_{\mathrm{lat}}(\mu_H/2)$ by the lattice count of Lemma~\ref{lem:c5-1-2} (the sum over all bonds dominates the sum over $D(\Lambda)$). The last display is \eqref{eq:c5-5-2} with $X= I_C$ and any $\mu\geqslant \mu_H/2$ (so that $\min(\mu_H/2,\mu)=\mu_H/2$), restricted to $b\in I_C\subseteq D(\Lambda)$, where $\operatorname{dist}(b,I_C)=0$; the terms with $b\notin I_C$ are nonnegative and dropped. Only (I1)--(I3) at the retreated carrier and Lemma~\ref{lem:c5-3-1}(b) are used; no column-sum bound.
\end{proof}

\subsubsection{Thresholds}
\label{sec:c5-5-2}
In addition to the thresholds of Appendix~\ref{app:thresholds}, \eqref{eq:c5-4-T} and \eqref{eq:c5-4-Tp}, assume
\begin{equation}\label{eq:c5-5-T}
    \begin{aligned}
        &\text{(i) }u_t\coloneqq C'_{\chi} g_0(1+2r)\leqslant u_*;\qquad\text{(ii) }\omega_t= C_0g_0(1+2r)^3\leqslant \omega_*/2;\\
        &\text{(iii) }\bar\omega\coloneqq C_0g_0(1+3r)^3\leqslant \min\bigl(1,\,\omega_*e^{-\kappa(B_h+4)}/(e^{9/8}N_{\mathrm{host}})\bigr);\\
        &\text{(iv) }\rho_t\coloneqq u_t/u_*\leqslant \rho_{\mathrm{pin}};\qquad\text{(v) }\bar\varepsilon_t\coloneqq 2K_{\chi} C_{\mathrm{pin}}\rho_t\leqslant 1,
    \end{aligned}
\end{equation}
together with the standing choice $\omega_*\leqslant 1/2$ recorded in Appendix~\ref{app:thresholds} (used in the positive host comparison of Lemma~\ref{lem:c4-7-1}). Here $C'_{\chi}\coloneqq 2N_{\mathrm{tr}}C_{\mathrm{link}}$ is the constant of Lemma~\ref{lem:c5-5-7z} (the finite-difference bound for closed contractions: Cauchy estimates on the row chart with $C_{\mathrm{link}}$, the traversal multiplicity $N_{\mathrm{tr}}$, and the factor 2 for the two $s$-dependent words of the source slot in Lemma~\ref{lem:c5-5-7}(d); it depends on $L$, the group and the character only). $N_{\mathrm{host}}\coloneqq 4^4\cdot20^4$ is the host count of Lemma~\ref{lem:c5-5-6} and Proposition~\ref{prop:c5-5-9} (the number of hosts $Y\in\mathbf D_0(\Lambda)$ with $b\in E(Y)$ for a given bond $b$, and the number within distance $R_{\mathrm{st}}\leqslant 5M$ of a given bond: a host is a box of $M$-cubes of side at most 4 cubes, \eqref{eq:part1-potential-host} and \eqref{eq:c-probe-host}, determined by its lowest cube, one of at most $20^4$ cubes within $9M$ of the bond, and by its side lengths, at most $4^4$ choices). $\rho_{\mathrm{pin}}$ is the root-weight smallness required by the positive summation of Lemma~\ref{lem:c4-6-2} (the condition $C\rho LM\leqslant \eta/8$ there, together with the root-incidence bound of Section~\ref{sec:c4-8-2} required by Lemma~\ref{lem:c4-6-2}, a property of the fixed geometry, and $(3^4A_F+LM+1)C_{\mathrm{pin}}\rho\leqslant \eta/16$ for the reference gas, Lemma~\ref{lem:c5-5-3}(d), with their fixed constants; $3^4$ is the neighbour factor for footprint contact of Section~\ref{sec:c4-12-4} and Lemma~\ref{lem:c4-6-env}(d)), and $C_{\mathrm{pin}}$, $K_{\chi}$ are the constants of Lemma~\ref{lem:c5-5-8}(a) and Lemma~\ref{lem:c5-5-7}(a). The subscript $t$ in $u_t$, $\rho_t$, $\bar\varepsilon_t$ names the $t$-family; none of them depends on $t$. Every condition holds for $g_0$ small, since $g_0(1+r)^3\to0$ ($r= A_1(\log g_0^{-2})^{p_0}$), with thresholds depending on $L$, the group, $M$, $\kappa$, $\omega_*$, $\eta$, $B_h$, $A_1$, $p_0$ and the character only, not on $n$, the periods or the history. Put $\bar s\coloneqq 1+2r+r^{\Lambda}_{\mathrm e}\leqslant 1+3r$.

\subsubsection{Translated probe measures and the centred expectation for the $t$-family}
The probe hosts $Y_b$, $b\in D(\Lambda)$, are in $\mathbf D_0(\Lambda)$ (Lemma~\ref{lem:c5-4-0}(d)); for a union $R$ of $LM$-cubes put $\mathsf L_R\coloneqq \{ b\in D(\Lambda)\, : \,Y_b\subseteq R\}$ (\eqref{eq:probe-eligible-carrier} at the carrier $D(\Lambda)$), and for $\lambda\in\mathfrak g^{D(\Lambda)}_{\mathbb C}$ let $L_Y(\lambda,B)\coloneqq \sum_{b\, : \,Y_b= Y}\langle\lambda_b,B_b\rangle$ (the probe in the centred variable). The family $\mathbf W_t(Y,\bigcdot )+L_Y(\lambda,\bigcdot )$, $Y\in\mathbf D_0(\Lambda)$, has the properties (i)--(iii) of the proof of Lemma~\ref{lem:c5-4-1}: $L_Y(\lambda,\bigcdot )$ depends on $B|_{E(Y)}$ ($Y_b= Y$ gives $b\in E(Y)$), and is bounded on the translated local box of $Y$ ($| B_b|_{\sigma}< r+t| m_b|_{\sigma}\leqslant 5r/4$ there), and nothing else changes. Hence Lemma~\ref{lem:c5-4-1} holds for it, with activities $k^{f,t}_{\Lambda,Z}(\lambda)$ entire in $\lambda$ and $k^{f,t}_{\Lambda,Z}(0)= k^{f,t}_{\Lambda}(Z)$.

\begin{lemma}[probe measures, majorant and dummy comparison for the $t$-family]\label{lem:c5-5-2}
Under the hypotheses of Lemma~\ref{lem:c5-4-2} and \eqref{eq:c5-5-T}, for every $t\in[0,1]$ and every connected final domain $Z$ for $D(\Lambda)$:
\begin{enumerate}[label=(\alph*)]
    \item There is a finite complex Borel measure $\nu^{f,t}_{\Lambda,Z}$ on $\mathfrak g^{\mathsf L_Z}$, supported in the closed translated box $\{| B_b+tm_b|_{\sigma}\leqslant r,\ b\in\mathsf L_Z\}$, with $k^{f,t}_{\Lambda,Z}(\lambda)=\int e^{\sum_{b\in\mathsf L_Z}\langle\lambda_b,B_b\rangle}\,\dd\nu^{f,t}_{\Lambda,Z}(B)$ and $\nu^{f,t}_{\Lambda,Z}(\mathfrak g^{\mathsf L_Z})= k^{f,t}_{\Lambda}(Z)$; explicitly, with the split selections $\ell=(\mathcal D,\mathsf J,\mathsf P,Z)$ of Lemma~\ref{lem:probe-measures} ($\mathsf J\subseteq\mathsf L_Z$ the selected probes, enlarged seed data $Y_0=\bigcup_{Y\in\mathcal D}Y\cup\bigcup_{b\in\mathsf J}Y_b$, $Z_0$, $E_0= E(Z_0)$, $Z_0'$) and $\Lambda^{f,t}_{\Lambda;\ell}$ the complex measure on $\mathfrak g^{E_0}$ defined by the right side of the display of Lemma~\ref{lem:c5-4-1} with these enlarged data (cutoffs $\chi^t_b$ over $E(Y_0)$, the weakened density and the bad set over this seed, Mayer factors for $Y\in\mathcal D$ only, no probe factor), $\nu^{f,t}_{\Lambda,Z}=\sum_{\ell\, : \,Z(\ell)= Z,\,Z_0\neq\varnothing}\int\bigl[\bigotimes_{b\in\mathsf J}(\delta_{B_b}-\delta_0)\otimes\bigotimes_{b\in\mathsf L_Z\setminus\mathsf J}\delta_0\bigr]\dd\Lambda^{f,t}_{\Lambda;\ell}(B)$.
    \item (Probe-enlarged majorant.) For every split selection $\ell$ and every bounded Borel $F$ on $\mathfrak g^{E_0}$,
    \begin{equation}\label{eq:c5-5-3}
        \Bigl|\int F\,\dd\Lambda^{f,t}_{\Lambda;\ell}\Bigr|\leqslant \sup| F|\cdot e^{-(\kappa_1-1)N_1(Z\setminus Z_0')+N_1(Z)-\frac9{16}a_{\mathrm{bad}}|\mathsf P|}\prod_{Y\in\mathcal D}\omega_*e^{1/2}e^{-\kappa d_0(Y)}\;,
    \end{equation}
    uniformly in $t$, $m$, the exterior data within the cutoff support and the labels; the probe hosts $Y_b$, $b\in\mathsf J$, carry no factor.
    \item (Dummy comparison for the $t$-family.) $\sum_{\ell\, : \,Z(\ell)= Z}u_*^{|\mathsf J|}\|\Lambda^{f,t}_{\Lambda;\ell}\|_{\mathrm{TV}}\leqslant m(Z)=\varepsilon_*e^{-\lambda_*\ell_1(Z)}$, the sum over all split selections with final domain $Z$ and nonempty seed $Z_0\neq\varnothing$ ($\mathsf J=\varnothing$ included; the activity sum of Section~\ref{sec:sourced-effective-interaction} excludes $Z_0=\varnothing$, and the $Z_0=\varnothing$ term has zero measure by the Cauchy argument of Lemma~\ref{lem:polymer}). The bound holds with the same constants when $\mathbf W_t$ is replaced by any potentials $\mathbf W'$ with $|\mathbf W'(Y,B)|\leqslant A_Y=(\omega_*/2)e^{-\kappa d_0(Y)}$ on the support of the host's cutoffs, the proof (and that of (b), on which it rests) using $\mathbf W_t$ only through this bound; in particular for $\zeta\mathbf W_0$, uniformly in complex $\zeta$ with $|\zeta|\leqslant R_{\Upsilon}$, wherever $R_{\Upsilon}|\mathbf W_0(Y,B)|\leqslant A_Y$ on that support. The bound is termwise: each $\|\Lambda^{f,t}_{\Lambda;\ell}\|_{\mathrm{TV}}$ is bounded by a majorant depending on the potentials only through $A_Y$ (proof), so the supremum over such $\mathbf W'$, in particular over $|\zeta|\leqslant R_{\Upsilon}$, may be taken inside the sum over $\ell$.
    \item (Centred expectation.) For bounded Borel $g$ on $\mathfrak g^{D(\Lambda)}$, with $\mathfrak Q^t_{\Lambda}(\mathcal A;g)\coloneqq \int g(\mathsf e_{\mathcal A}B)\,\dd\bigl(\bigotimes_{Z\in\mathcal A}\hat\nu^{f,t}_{\Lambda,Z}\bigr)(B)$, $\hat\nu^{f,t}_{\Lambda,Z}\coloneqq \nu^{f,t}_{\Lambda,Z}-k^{f,t}_{\Lambda}(Z)\delta_0$, $q^t_{\Lambda}(\mathcal A)\coloneqq (\Xi^{f,t}_{\Lambda})^{[\cup\mathcal A]}/\Xi^{f,t}_{\Lambda}$,
    \begin{equation*}
        \frac{\int g(B)\,\Theta^te^{-\frac12\langle B,HB\rangle+\sum_Y\mathbf W_t(Y,B)}\,\dd B}{\int\Theta^te^{-\frac12\langle B,HB\rangle+\sum_Y\mathbf W_t(Y,B)}\,\dd B}= g(0)+\sum_{\varnothing\neq\mathcal A\text{ compatible}}q^t_{\Lambda}(\mathcal A)\,\mathfrak Q^t_{\Lambda}(\mathcal A;g)\;,
    \end{equation*}
    and, by \eqref{eq:c5-4-3} with $F= f_{\xi}$,
    \begin{gather}
        F^{\Lambda}(t)= f_{\xi}(tm)+\sum_{\varnothing\neq\mathcal A\text{ compatible}}q^t_{\Lambda}(\mathcal A)\,\mathfrak Q^t_{\Lambda}(\mathcal A;f_{\xi}(\bigcdot +tm))\;,\label{eq:c5-5-4}\\
        \mathfrak Q^t_{\Lambda}(\mathcal A;f_{\xi}(\bigcdot +tm))=\sum_{(\ell_Z)_{Z\in\mathcal A}\, : \,\mathsf J_Z\neq\varnothing}\int\Delta_{\mathsf J}[f_{\xi}(\bigcdot +tm)](B)\,\dd\bigotimes_Z\Lambda^{f,t}_{\Lambda;\ell_Z}(B)\;,\notag
    \end{gather}
    where $\mathsf J=\bigcup_Z\mathsf J_Z$ and $\Delta_{\mathsf J}g(B)\coloneqq \sum_{K\subseteq\mathsf J}(-1)^{|\mathsf J\setminus K|}g(\mathbf 1_KB)$ is the mixed finite difference of \eqref{eq:c4-4-1} (every label outside $K$ set to zero, so the young variable of $f_{\xi}$ in the $K$-term is $\mathbf 1_KB+tm$: the unselected labels sit at $a= tm$).
\end{enumerate}
\end{lemma}

\begin{proof}
(a) is the proof of Lemma~\ref{lem:probe-measures} with $\mathbf V_{\xi}\to\mathbf W_t$, $\chi_b\to\chi^t_b$, $I\to D(\Lambda)$: the split-label expansion at a host, the exact integration \eqref{eq:probe-selected-factor} of the probe Mayer variables, the Laplace transform of $\delta_{B_b}-\delta_0$ and the finiteness of each $\Lambda^{f,t}_{\Lambda;\ell}$ (image of a density whose modulus is integrable by Lemma~\ref{lem:gaussian-payment}(ii) for $H$, times bounded factors: the Cauchy kernels on the circles $|\sigma_{\Delta}|= e^{\kappa_1}$ are bounded by $e^{\kappa_1}(e^{\kappa_1}-1)^{-2}$, and the integrand is jointly Borel in $(X,B,s,\sigma,t_Y)$) use only Lemma~\ref{lem:c5-4-1} and its properties (i)--(iii). The support statement: $\chi^t_b(B)$, $b\in E(Y_0)\supseteq\mathsf J$, forces $| B_b+tm_b|_{\sigma}< r$ on the selected probes, and the unselected labels carry $\delta_0$ ($B_b=0$, and $| tm_b|_{\sigma}\leqslant r/4< r$). At $\lambda=0$ every term with $\mathsf J\neq\varnothing$ has zero mass (each $\delta_{B_b}-\delta_0$ has), so the total mass is $k^{f,t}_{\Lambda,Z}(0)= k^{f,t}_{\Lambda}(Z)$. (b) The proof of Lemma~\ref{lem:c5-4-2} is written for a selection $\mathfrak s=(\mathcal D,\mathsf P,Z)$ and bounds factor by factor; the enlarged seed enters only through $E_0= E(Z_0)$ in the density \eqref{eq:fixed-kernel-density}, whose modulus is paid against $e^{(\alpha_G/2)\| B\|^2}$ by Lemma~\ref{lem:gaussian-payment}(ii) for every seed $Z_0$ and every $Z\supseteq Z_0'$ (the lemma is stated for every seed), through the cutoffs $\chi^t_b$ over $E(Y_0)$, each $\leqslant 1$, and through the bad bonds $\mathsf P\subseteq D(\Lambda)\setminus E(Y_0)$, each paying $e^{-\frac9{16}a_{\mathrm{bad}}}e^{(\alpha_G/2)h(B_b)^2}$ as in Lemma~\ref{lem:c5-4-2}; the Mayer factors are those of $Y\in\mathcal D$ only, each bounded by $\omega_*e^{1/2}e^{-\kappa d_0(Y)}$ on the support of its cutoffs; a probe host contributes no factor; the weakening bound $e^{-(\kappa_1-1)}$ per cube of $Z\setminus Z_0'$ is unchanged. (c) is the positive host comparison of Lemma~\ref{lem:c4-7-1}: at a host $Y$ the sum over nonempty label subsets of the positive weights $v_{\mathrm{ordinary}}= e^{A_Y}-1$, $v_{\mathrm{probe}}= u_*$ is at most $\omega_*e^{1/2}e^{-\kappa d_0(Y)}$, and all split selections with the same active-host set have the same enlarged seed, density, bad set and weakening, so the many-to-one sum is dominated hostwise by the undifferentiated majorant \eqref{eq:c5-5-3} with $F=1$, after which the seed/enlargement/final-domain sums \eqref{eq:native-source-initial-domain-sum}, \eqref{eq:native-source-enlarged-domain-sum} and the $Y'$-sum in the proof of Lemma~\ref{lem:activity-bounds} give $m(Z)$ exactly as in Corollary~\ref{cor:c5-4-3} (whose proof shows that these sums hold with $\frac9{16}a_{\mathrm{bad}}$ in place of $a_{\mathrm{bad}}$ under \eqref{eq:c5-4-Tp}); the comparison needs the ordinary factor's modulus $|\int_0^1\mathbf W_te^{t_Y\mathbf W_t}\,\dd t_Y|=| e^{\mathbf W_t}-1|\leqslant e^{|\mathbf W_t|}-1\leqslant e^{A_Y}-1= v_{\mathrm{ordinary}}$ with $A_Y=(\omega_*/2)e^{-\kappa d_0(Y)}$, which holds because $|\mathbf W_t(Y,B)|\leqslant \omega_te^{-\kappa d_0(Y)}\leqslant A_Y$ on the support of the host's cutoffs by \eqref{eq:c5-4-4} and \eqref{eq:c5-5-T}(ii), and the chain's last step uses $\omega_*\leqslant 1/2$ and $n_Yu_*\leqslant (\omega_*/8)e^{-\kappa d_0(Y)}$ (\eqref{eq:c4-7-2}). The sum includes the split selections with $\mathsf J=\varnothing$ (a host is active through its ordinary label or through probe labels; \eqref{eq:c4-7-3} sums over all nonempty label subsets), so (c) holds with $\sum_{\ell\, : \,Z(\ell)= Z}$ in place of $\sum_{\ell\, : \,Z(\ell)= Z,\,\mathsf J\neq\varnothing}$. (d) is the proof of Lemma~\ref{lem:measure-identity} and Proposition~\ref{prop:centred-expectation} for the translated expansion: both sides of $\mathfrak n_{\Lambda}^{-1}\Theta^te^{-\frac12\langle B,HB\rangle+\sum\mathbf W_t}\,\dd B=\sum_{\mathcal P}(\mathsf e_{\mathcal P})_*\bigotimes_{Z\in\mathcal P}\nu^{f,t}_{\Lambda,Z}$ are finite complex measures with compact support and the same Laplace transform (Lemma~\ref{lem:c5-4-1} for $\mathbf W_t+L(\lambda,\bigcdot )$ and (a)), hence equal; expanding $\nu^{f,t}=\hat\nu^{f,t}+k^{f,t}\delta_0$ and dividing by $\mathfrak n_{\Lambda}\Xi^{f,t}_{\Lambda}$ gives the first display; the second is \eqref{eq:c5-4-3} with $F= f_{\xi}$ (bounded on $\mathfrak g^{D(\Lambda)}$, extended by zero off the enlarged box, which contains the support of every $\nu^{f,t}$: $| B_b+tm_b|_{\sigma}\leqslant r$ and $| tm_b|_{\sigma}\leqslant r/4$ give $| B_b|_{\sigma}\leqslant 5r/4$ and $|(\mathbf 1_KB+tm)_b|_{\sigma}\leqslant 3r/2\leqslant 2r$), and the formula for $\mathfrak Q^t$ in terms of $\Delta_{\mathsf J}$ is the integration of $g$ against the Dirac tensors in (a), zero-mass part (the terms with $\mathsf J_Z=\varnothing$ for some $Z$ are removed by the centring).
\end{proof}

\begin{lemma}[joint extraction for the $t$-family]\label{lem:c5-5-3}
\begin{enumerate}[label=(\alph*)]
    \item For a connected final domain $Z$ for $D(\Lambda)$ the contact set $\mathsf T_Z\coloneqq \{ n\, : \,\mathcal F_{r_n}\cap\mathsf L_Z\neq\varnothing\}$ (\eqref{eq:part1-contact-set}, with the convention $\mathcal F_r=\varnothing$ off $\Omega_1$) is field-independent and $t$-independent; if $\mathcal A$ contains an unmarked $Z$ ($\mathsf T_Z=\varnothing$) then $\mathfrak Q^t_{\Lambda}(\mathcal A;f_{\xi}(\bigcdot +tm))=0$. The groups $W=(\mathcal A(W),(\mathfrak c_1,\dots,\mathfrak c_{k(W)}),\mathcal H(W))$ of Definition~\ref{def:groups} for the carrier $D(\Lambda)$ form one index set for all $t\in[0,1]$ and for the reference data; their weights are $w^t_{\Lambda}(W)\coloneqq ((-1)^{k(W)}/k(W)!)\prod_j\Phi^{v,t}_{\Lambda}(\mathfrak c_j)$ with the cluster weights of Corollary~\ref{cor:c5-4-3}. With $A^t_{\Lambda}(\mathcal S)\coloneqq \prod_{W\in\mathcal S}w^t_{\Lambda}(W)\,\mathfrak Q^t_{\Lambda}(\mathcal A(\mathcal S);f_{\xi}(\bigcdot +tm))$, $A^t_{\Lambda}(\varnothing)\coloneqq f_{\xi}(tm)$, Lemma~\ref{lem:grouped-source} holds: $F^{\Lambda}(t)=\sum_{\mathfrak B\text{ compatible}}A^t_{\Lambda}(\mathfrak B)$, absolutely convergent.
    \item Reference data at the carrier $D(\Lambda)$: the same $\Omega_1$, cubes, hosts $\mathbf D_0(\Lambda)$, constants and cutoffs, with $U^{\circ}_1=\id$, $V^{\circ}_0=\id$, $A^{\circ}=0$, $\Psi^{\circ}=\id$, exterior variable zero on $J(\Lambda)$ (so $m^{\circ}=0$, $\ell^{\circ}=0$, and the reference objects are independent of $t$), $H^{\circ}=(\mathsf H^{\Omega_1}_{\id})_{D(\Lambda)D(\Lambda)}$, $f^{\circ}(a)= W_{\chi}(C;e^{iq_{\id}(a)})$ with links off $\Omega_1$ equal to $\id$. Lemma~\ref{lem:c5-4-0} holds for these data (its (a) by principal restriction of the reference Hessian of Section~\ref{sec:sourced-effective-interaction}, which satisfies \eqref{eq:fixed-kernel-inputs} by Lemma~\ref{lem:c2-reference}; (b) trivially, $r^{\circ}_{\mathrm e}=0$; (c), (d) as stated), and so do Lemma~\ref{lem:c5-4-1}, Lemma~\ref{lem:c5-4-2}, Corollary~\ref{cor:c5-4-3} and Lemma~\ref{lem:c5-5-2} (their proofs use the listed lemmas only). Define $a^{\Lambda}_W\coloneqq w^{\circ}_{\Lambda}(W)\int\chi(P^{\circ}_W(a))\,\dd\bigl(\bigotimes_{Z\in\mathcal A(W)}\hat\nu^{f,\circ}_{\Lambda,Z}\bigr)(a)$ as in \eqref{eq:part1-reference-group}, with $P^{\circ}_W$ the ordered product of reference links along $\mathcal H(W)$ at the young variable $\mathsf e_{\mathcal A(W)}a$. Then Lemma~\ref{lem:schur} holds: $w^{\circ}_{\Lambda}(W)\int\chi(P^{\circ}_W(a)h)\,\dd\bigotimes\hat\nu^{f,\circ}_{\Lambda}= a^{\Lambda}_W\chi(h)$ for $h\in G$, and $S^{\circ}_{\Lambda}\coloneqq \sum_{\mathfrak B\text{ compatible}}\prod_{W\in\mathfrak B}a^{\Lambda}_W=\int f^{\circ} K^{\circ}_{\Lambda}\,\dd a/I^{\circ}_{\Lambda}$ ($K^{\circ}_{\Lambda}$ the kernel $K^{\Lambda}_{\xi;0}$ of Section~\ref{sec:c5-4} at these reference data and $I^{\circ}_{\Lambda}\coloneqq \int K^{\circ}_{\Lambda}\,\dd a$, written out in \eqref{eq:c5-7-1}).
    \item With $R^t_{\Lambda}(\mathfrak B)\coloneqq \sum_{\mathcal S\subseteq\mathfrak B}(-1)^{|\mathfrak B\setminus\mathcal S|}\bigl(\prod_{W\in\mathfrak B\setminus\mathcal S}a^{\Lambda}_W\bigr)A^t_{\Lambda}(\mathcal S)$ and $S^{\circ[\mathfrak B]}_{\Lambda}$ the reference sum over families compatible with every member of $\mathfrak B$,
    \begin{gather}
        F^{\Lambda}(t)= S^{\circ}_{\Lambda} f_{\xi}(tm)+\sum_{\varnothing\neq\mathfrak B\text{ compatible}}S^{\circ[\mathfrak B]}_{\Lambda} R^t_{\Lambda}(\mathfrak B)\;,\qquad\text{hence}\label{eq:c5-5-5}\\
        \mathcal B^{\Lambda}_{\xi}[f_{\xi}]= S^{\circ}_{\Lambda}[f_{\xi}(m)-f_{\xi}(0)]+\sum_{\varnothing\neq\mathfrak B\text{ compatible}}S^{\circ[\mathfrak B]}_{\Lambda}[R^1_{\Lambda}(\mathfrak B)-R^0_{\Lambda}(\mathfrak B)]\;.\notag
    \end{gather}
    \item $| a^{\Lambda}_W|\leqslant M_{\Lambda}(W)$, the positive group majorant of Lemma~\ref{lem:c5-5-7}(d) below, $M_{\Lambda}(W)=(1/k(W)!)\prod_{Z\in\mathcal A(W)}\rho_tm(Z)\prod_j|\Phi|^+_{\Lambda}(\mathfrak c_j)$, defined without reference to the source data. Under \eqref{eq:c5-5-T}(iv) the reference group gas with the activities $a^{\Lambda}_W$ satisfies the Kotecký--Preiss condition, so that $S^{\circ}_{\Lambda}\neq0$, $q^{\circ}_{\Lambda}(\mathfrak B)\coloneqq S^{\circ[\mathfrak B]}_{\Lambda}/S^{\circ}_{\Lambda}=\exp(-\sum_{\Gamma\text{ hits }\mathfrak B}\Phi^{\circ}_{\Lambda}(\Gamma))$ (the cluster expansion of the reference group gas, Lemma~\ref{lem:c4-6-3} and Section~\ref{app:kp}, at the carrier $D(\Lambda)$; a cluster hits $\mathfrak B$ when it contains a group incompatible with a member of $\mathfrak B$), and $| q^{\circ}_{\Lambda}(\mathfrak B)|\leqslant \prod_{W\in\mathfrak B}e^{\theta^{\circ}(W)}$ with $\theta^{\circ}(W)\coloneqq \sum_{\Gamma\text{ hits }W}|\Phi^{\circ}_{\Lambda}(\Gamma)|\leqslant \nu\sigma(W)/16$, where $\sigma(W)$ and $\nu$ are the size and the reserve constant of Lemma~\ref{lem:c5-5-8}; and the rooted bound $\sum_{\Gamma\ni W}|\Phi^{\circ}_{\Lambda}(\Gamma)|\leqslant | a^{\Lambda}_W| e^{13\nu\sigma(W)/8}$ holds. All of (d) holds verbatim for the reference gas at the carrier $I$ (the reference data of Section~\ref{sec:sourced-effective-interaction}, activities $a_W$, cluster coefficients $\Phi^{\circ}_{\xi}$), since its proof uses, at the carrier, only Lemma~\ref{lem:c5-5-7z}(b) at the reference data (Lemma~\ref{lem:charts} at $U=\id$), Lemma~\ref{lem:c5-5-2}(a),(c),(d) at the reference data (Lemmas~\ref{lem:probe-measures}, \ref{lem:measure-identity} and~\ref{lem:c4-7-1} at $I$), the first display of Lemma~\ref{lem:c5-5-8}(a) (the summation of Lemma~\ref{lem:c4-6-2}, whose geometry does not involve the carrier) and the Kotecký--Preiss rooted bound of Section~\ref{app:kp}; each holds at $I$, the case $D(\Lambda)= I$ (no retreat) of these statements.
\end{enumerate}
\end{lemma}

\begin{proof}
(a) A traversal's link depends on the young variable through $a|_{\mathcal F_{r_n}\cap D(\Lambda)}$ (Lemma~\ref{lem:row-locality}; the de-younged bonds are constants); for unmarked $Z$ the function $B\mapsto f_{\xi}(\mathbf 1_KB+tm)$ does not depend on $B|_{\mathsf L_Z}$, so the $K$-sum over the labels of $\mathsf J_Z$ telescopes to zero in \eqref{eq:c5-5-4} (equivalently, $\hat\nu^{f,t}_{\Lambda,Z}$ has zero mass). The groups are built from roots, cluster lists (tuples of domains), contact sets and hulls, none of which involves the fields or $t$; only $w^t_{\Lambda}$ and $\hat\nu^{f,t}_{\Lambda}$ depend on $t$. The regrouping proof of Lemma~\ref{lem:grouped-source} is combinatorial, and its absolute convergence uses $|\mathfrak Q^t_{\Lambda}(\mathcal A;g)|\leqslant \sup| g|\prod_Z\|\hat\nu^{f,t}_{\Lambda,Z}\|$ (finite by Lemma~\ref{lem:c5-5-2}(a)) and the hitting sum $\sum_{\mathfrak c\text{ hits }\mathcal A}|\Phi^{v,t}_{\Lambda}(\mathfrak c)|\leqslant \frac18\sum_{Z\in\mathcal A}N_1(Z)$ of Corollary~\ref{cor:c5-4-3}. (b) The proof of Lemma~\ref{lem:schur} uses the constant gauge $R_u$: $H^{\circ}= R_uH^{\circ} R_u^{-1}$ holds for the principal restriction because $R_u$ acts bondwise; the domains, weights, cutoffs (Ad-invariant norm), the Gaussian density, the probe term and the reconstruction covariance $q_{\id}(R_uB)= R_uq_{\id}(B)$ are as in Section~\ref{sec:sourced-effective-interaction}, and the hosts and partitions are field-independent; Schur's lemma and the Fubini evaluation of the group product are verbatim. (c) is the Boolean inversion and interchange of absolutely convergent sums of Proposition~\ref{prop:joint-extraction}, with the bound $\beta_W$ of its proof formed from $w^t_{\Lambda}$, $w^{\circ}_{\Lambda}$, $\hat\nu^{f,t}_{\Lambda}$, $\hat\nu^{f,\circ}_{\Lambda}$; the $t$-independence of $a^{\Lambda}_W$ and $S^{\circ[\mathfrak B]}_{\Lambda}$ gives the second form. (d) The activity bound: $a^{\Lambda}_W= w^{\circ}_{\Lambda}(W)\int g^{\circ}_W\,\dd\bigotimes_{Z\in\mathcal A(W)}\hat\nu^{f,\circ}_{\Lambda,Z}$ with $g^{\circ}_W(a)\coloneqq \chi(P^{\circ}_W(a))$, a closed Wilson contraction at the reference data; by Lemma~\ref{lem:c5-5-2}(a),(d) for the reference data (part (b)), $\int g^{\circ}_W\,\dd\bigotimes\hat\nu^{\circ}=\sum_{(\ell_Z)_{Z\in\mathcal A(W)}\, : \,\mathsf J_Z\neq\varnothing}\int\Delta_{\mathsf J}g^{\circ}_W\,\dd\bigotimes_Z\Lambda^{f,\circ}_{\Lambda;\ell_Z}$; Lemma~\ref{lem:c5-5-7z}(b) at the reference data gives $\sup|\Delta_{\mathsf J}g^{\circ}_W|\leqslant u_t^{|\mathsf J|}$ (base point $a_0=0$, increments $| B_b|_{\sigma}\leqslant r\leqslant 5r/4$ on the support of $\nu^{f,\circ}$, $m^{\circ}=0$), and Lemma~\ref{lem:c5-5-2}(c) for the reference data with $(u_t/u_*)^{|\mathsf J_Z|}\leqslant \rho_t$ ($\mathsf J_Z\neq\varnothing$, and $\rho_t\leqslant 1$ by \eqref{eq:c5-5-T}(i)) gives $|\int g^{\circ}_W\,\dd\bigotimes\hat\nu^{\circ}|\leqslant \prod_{Z\in\mathcal A(W)}\rho_tm(Z)$; $| w^{\circ}_{\Lambda}(W)|\leqslant (1/k(W)!)\prod_j|\Phi^{v,\circ}_{\Lambda}(\mathfrak c_j)|\leqslant (1/k(W)!)\prod_j|\Phi|^+_{\Lambda}(\mathfrak c_j)$ by Corollary~\ref{cor:c5-4-3} for the reference data and the tree-graph inequality (Lemma~\ref{lem:c5-5-6}(b)). The Kotecký--Preiss condition: by the first display of Lemma~\ref{lem:c5-5-8}(a), which uses only the majorant $M_{\Lambda}$, $\sup_{\Delta}\sum_{W\, : \,\Delta\in F(W)}e^{4\nu\sigma(W)}M_{\Lambda}(W)+\sup_n\sum_{W\, : \,n\in\mathcal H(W)}e^{4\nu\sigma(W)}M_{\Lambda}(W)\leqslant C_{\mathrm{pin}}\rho_t$; two groups are incompatible only if their footprints touch (contact includes adjacency across a face, edge or corner, not only intersection: Section~\ref{sec:c4-12-4}) or their intervals overlap (the compatibility of groups, Section~\ref{sec:joint-extraction} after Definition~\ref{def:groups}; Section~\ref{sec:c4-6-2}); a footprint has at most $A_F\sigma(W)$ cubes (Lemma~\ref{lem:c4-6-env}), so an incompatible $W'$ has a footprint cube in the set $N_1(F(W))$ of cubes touching $F(W)$, of at most $3^4A_F\sigma(W)$ cubes (the neighbour factor $3^4$), or a traversal in $\mathcal H(W)$, of at most $LM\sigma(W)$ traversals ($|\mathcal H(W)|/(LM)\leqslant \sigma(W)$); hence $\sum_{W'\text{ incompatible with }W}| a^{\Lambda}_{W'}| e^{3\nu\sigma(W')/2}e^{\nu\sigma(W')/8}\leqslant \sum_{\Delta\in N_1(F(W))}\sum_{W'\, : \,\Delta\in F(W')}(\dots)+\sum_{n\in\mathcal H(W)}\sum_{W'\, : \,n\in\mathcal H(W')}(\dots)\leqslant (3^4A_F+LM)\sigma(W)C_{\mathrm{pin}}\rho_t\leqslant \nu\sigma(W)/8$ under \eqref{eq:c5-5-T}(iv): this is the Kotecký--Preiss hypothesis for the gas with activities $| a^{\Lambda}_W| e^{3\nu\sigma(W)/2}$ and node cost $\nu\sigma(W)/8$, and the Kotecký--Preiss theorem in the rooted form of Section~\ref{app:kp} (the recursion following \eqref{eq:marginal-pinned-bound} applied to this gas, with \cite[Proposition~5.3, Theorem~5.4]{FriedliVelenik}) gives the absolute convergence of the cluster series, $S^{\circ}_{\Lambda}=\exp(\sum_{\Gamma}\Phi^{\circ}_{\Lambda}(\Gamma))\neq0$, the exclusion identity $S^{\circ[\mathfrak B]}_{\Lambda}=\exp(\sum_{\Gamma\text{ hits no member of }\mathfrak B}\Phi^{\circ}_{\Lambda}(\Gamma))$ (the same with the species incompatible with a member of $\mathfrak B$ excluded), and the rooted bound $\sum_{\Gamma\ni W}m_W(\Gamma)|\Phi^{\circ}_{\Lambda}(\Gamma)|\leqslant | a^{\Lambda}_W| e^{13\nu\sigma(W)/8}$, $m_W(\Gamma)$ the occurrence multiplicity. Hence $\theta^{\circ}(W)\leqslant \sum_{W'= W\text{ or incompatible with }W}\sum_{\Gamma\ni W'}|\Phi^{\circ}_{\Lambda}(\Gamma)|\leqslant \sum_{W'= W\text{ or incompatible with }W}| a^{\Lambda}_{W'}| e^{13\nu\sigma(W')/8}\leqslant (3^4A_F+LM+1)\sigma(W)C_{\mathrm{pin}}\rho_t\leqslant \nu\sigma(W)/16$ under \eqref{eq:c5-5-T}(iv) ($e^{13\nu\sigma/8}\leqslant e^{4\nu\sigma}$, and the term $W'= W$ is bounded by the footprint pin at a cube of $F(W)$).
\end{proof}

\subsubsection{The two terms of \eqref{eq:c5-5-5}}
The first is the shift of the loop at zero fluctuation; the second is the difference of the joint defects, which is where the source-containing connected boundary terms live.

\begin{lemma}[the shift term]\label{lem:c5-5-4}
\begin{equation*}
    | f_{\xi}(m)-f_{\xi}(0)|\leqslant C_{\mathrm{resp},\chi}g_0\sum_{b\in I_C}h(m_b)\leqslant C_{\mathrm{resp},\chi}C_m\vartheta N_{\mathrm{lat}}(\mu_H/2)g_0r\sum_{e\in J_1(\Lambda)}e^{-\mu_H\operatorname{dist}(e,I_C)/2}\;.
\end{equation*}
\end{lemma}

\begin{proof}
$f_{\xi}(m)-f_{\xi}(0)=\int_0^1\sum_{b\in I_C}\langle m_b,\nabla_bf_{\xi}(sm)\rangle\,\dd s$; the points $sm$ lie in the box $| a_b|_{\sigma}\leqslant r/4$ by \eqref{eq:c5-4-1}, inside the enlarged box, where $\sup h(\nabla_bf_{\xi})\leqslant C_{\mathrm{resp},\chi}g_0\mathbf 1_{I_C}(b)$ (\eqref{eq:c3-2-3}, the labelled source-gradient bound of Section~\ref{sec:c3-2}, as used in the proof of Lemma~\ref{lem:c5-4-4}); then the last clause of Lemma~\ref{lem:c5-5-1}.
\end{proof}

\begin{lemma}[conditional-Gaussian shell bound]\label{lem:c5-5-5}
Let $b\in E_0$, $S_b\coloneqq \{ x\in\mathfrak g\, : \,r-| m_b|_{\sigma}\leqslant | x|_{\sigma}< r+| m_b|_{\sigma}\}$, $v\in\mathfrak g^{E_0}$ real, and $0\leqslant \alpha\leqslant \alpha_G$ ($\alpha_G\leqslant \gamma/2$; $h(x)^2=\langle x,x\rangle$ is the squared norm of the invariant inner product). Write $\mu_{C_{E_0}}(\dd B)=\mu_{C'}(\dd B')\,N(\mu(B'),c_b)(\dd B_b)$ for the disintegration of the centred Gaussian of covariance $C_{E_0}=(\iota^*_{E_0}H\iota_{E_0})^{-1}$ (Section~\ref{sec:fixed-first-kernel}, at the carrier $D(\Lambda)$) along the coordinate $b$ ($B'$ the other coordinates, $N(\mu,c)$ the Gaussian on $\mathfrak g$ with mean $\mu$ and covariance $c$), where $c_b=((C_{E_0}^{-1})_{bb})^{-1}=(H_{bb})^{-1}$ and $\mu(B')$ is linear in $B'$. Then for every $B'$,
\begin{gather}
    \int\mathbf 1_{S_b}(B_b)e^{-\langle B_b,v_b\rangle+(\alpha/2)h(B_b)^2}\,N(\mu(B'),c_b)(\dd B_b)\leqslant \varpi_b\int e^{-\langle B_b,v_b\rangle+(\alpha/2)h(B_b)^2}\,N(\mu(B'),c_b)(\dd B_b)\;,\label{eq:c5-5-6}\\
    \varpi_b\coloneqq (h_+/2\pi)^{d_G/2}\operatorname{Leb}(S_b)\leqslant C_Vr^{d_G-1}| m_b|_{\sigma}\;,\notag
\end{gather}
$C_V\coloneqq (h_+/2\pi)^{d_G/2}\,2d_GV_{\sigma}(5/4)^{d_G-1}$, $V_{\sigma}$ the volume of the unit $|\bigcdot |_{\sigma}$-ball of $\mathfrak g$ for the Lebesgue measure of the $h$-inner product. Consequently $\varpi_b\leqslant C_{\varpi} r^{d_G}\sum_{e\in J_1(\Lambda)}e^{-\mu_Hd(b,e)}$, $C_{\varpi}\coloneqq C_Vc_hC_m\vartheta$.
\end{lemma}

\begin{proof}
The conditional law of $B_b$ given $B'$ under a centred Gaussian of covariance $C$ is Gaussian with covariance $((C^{-1})_{bb})^{-1}$ and mean linear in $B'$; here $C^{-1}=\iota^*H\iota$, whose $bb$ block is $H_{bb}$. Completing the square (write $\mu$ for $\mu(B')$): the density of $e^{-\langle x,v_b\rangle+(\alpha/2)\langle x,x\rangle}N(\mu,c_b)(\dd x)$ is a constant times $\exp(-\tfrac12\langle x,(H_{bb}-\alpha\,\id)x\rangle+\langle x,H_{bb}\mu-v_b\rangle)$, and $H_{bb}-\alpha\,\id\geqslant (\gamma-\alpha_G)\id\geqslant (\gamma/2)\id>0$ ($H_{bb}\geqslant \gamma\,\id$ as a principal restriction of $H$, (I3)), so $e^{-\langle x,v_b\rangle+(\alpha/2)\langle x,x\rangle}N(\mu,c_b)(\dd x)= Z_{\alpha} N(\mu_{\alpha},c_{\alpha})(\dd x)$ with $c_{\alpha}\coloneqq (H_{bb}-\alpha\,\id)^{-1}$, $\mu_{\alpha}\coloneqq c_{\alpha}(H_{bb}\mu-v_b)$ and a positive constant $Z_{\alpha}$ (for $\alpha=0$: $Z_0= e^{-\langle\mu,v_b\rangle+\frac12\langle v_b,c_bv_b\rangle}$, $\mu_0=\mu-c_bv_b$); hence the left side of \eqref{eq:c5-5-6} equals $Z_{\alpha} N(\mu_{\alpha},c_{\alpha})(S_b)$ and the right side's integral equals $Z_{\alpha}$; the Gaussian density of $N(\mu_{\alpha},c_{\alpha})$ is at most $(2\pi)^{-d_G/2}\det(c_{\alpha})^{-1/2}=(2\pi)^{-d_G/2}\det(H_{bb}-\alpha\,\id)^{1/2}\leqslant (2\pi)^{-d_G/2}\det(H_{bb})^{1/2}\leqslant (h_+/2\pi)^{d_G/2}$ ($H_{bb}\leqslant h_+\id$ by (I3)), whence $N(\mu_{\alpha},c_{\alpha})(S_b)\leqslant (h_+/2\pi)^{d_G/2}\operatorname{Leb}(S_b)$. By homogeneity of the norm, $\operatorname{Leb}(\{| x|_{\sigma}< s\})= V_{\sigma} s^{d_G}$, so $\operatorname{Leb}(S_b)= V_{\sigma}[(r+\epsilon)^{d_G}-(r-\epsilon)^{d_G}]$ with $\epsilon=| m_b|_{\sigma}\leqslant r/4$, and the mean value theorem gives $\leqslant V_{\sigma}\cdot2\epsilon\cdot d_G(5r/4)^{d_G-1}$. The last bound is $| m_b|_{\sigma}\leqslant c_hh(m_b)$ (Remark~\ref{rem:c5-1-3}) and \eqref{eq:c5-5-1}.
\end{proof}

The point of \eqref{eq:c5-5-6} is that the shell is paid by the Gaussian itself, with the tilt $e^{-\langle B,v\rangle}$ (which in the density \eqref{eq:fixed-kernel-density} is $e^{-\langle B,\operatorname{Re}\Gamma_{E_0}(\sigma)X\rangle}$ after moduli are taken) kept intact: the ratio $\varpi_b$ is of order $r^{d_G-1}| m_b|$ and does not involve the reserve $a_{\mathrm{bad}}$. The alternative sup-over-$B_b$ thin-shell bound would cost $e^{25h_+r^2/32}$ per mark, which $e^{-\frac9{16}a_{\mathrm{bad}}}$ cannot pay ($\alpha_G\leqslant \gamma/2< h_+$), and is not used.

\begin{lemma}[marked differences]\label{lem:c5-5-6}
Define, for a bond $b\in D(\Lambda)$ and a host $Y\in\mathbf D_0(\Lambda)$, the mark weights
\begin{align*}
    \mathfrak m_b&\coloneqq 2\varpi_b\leqslant 2C_{\varpi} r^{d_G}\sum_{e\in J_1(\Lambda)}e^{-\mu_Hd(b,e)}\\
    &\qquad\text{(one $\varpi_b$ for the bond's own translated cutoff,}\\
    &\qquad\text{one for the cutoffs of the hosts containing $b$, paid by \eqref{eq:c5-5-T}(iii))},\\
    \mathfrak m_Y&\coloneqq 2\bar\omega e^{\bar\omega}\cdot e^{\kappa(B_h+4)}/\omega_*\cdot e^{1/8}\\
    &\qquad\text{if some near primitive assigned to $Y$ has an input in $J(\Lambda)$},\\
    \mathfrak m_Y&\coloneqq N_hC_pg_0\bar s^2(1+\bar\omega)e^{\bar\omega}\cdot e^{\kappa(B_h+4)}/\omega_*\cdot e^{1/8}\sum_{b\in E(Y)}h(m_b)\quad\text{otherwise},
\end{align*}
and for a union $Z$ of $LM$-cubes $\mathfrak M_{\Lambda}(Z)\coloneqq \sum_{b\in E(Z)}\mathfrak m_b+\sum_{Y\in\mathbf D_0(\Lambda),\,Y\subseteq Z}\mathfrak m_Y$. Under the hypotheses of Lemma~\ref{lem:c5-5-2}:
\begin{enumerate}[label=(\alph*)]
    \item For every split selection $\ell$ with $Z(\ell)= Z$ ($\mathsf J=\varnothing$ allowed) and every bounded Borel $F$, $|\int F\,\dd(\Lambda^{f,1}_{\Lambda;\ell}-\Lambda^{f,0}_{\Lambda;\ell})|\leqslant \sup| F|\cdot\mathfrak M_{\Lambda}(Z)\cdot e^{-(\kappa_1-1)N_1(Z\setminus Z_0')+N_1(Z)-\frac9{16}a_{\mathrm{bad}}|\mathsf P|}\prod_{Y\in\mathcal D}\omega_*e^{1/2}e^{-\kappa d_0(Y)}$ (the majorant \eqref{eq:c5-5-3} times the sum of the mark weights; the division of a host mark by the host's majorant factor is built into $\mathfrak m_Y$); and $\sum_{\ell\, : \,Z(\ell)= Z}u_*^{|\mathsf J|}\|\Lambda^{f,1}_{\Lambda;\ell}-\Lambda^{f,0}_{\Lambda;\ell}\|_{\mathrm{TV}}\leqslant \mathfrak M_{\Lambda}(Z)m(Z)$.
    \item $| k^{f,1}_{\Lambda}(Z)-k^{f,0}_{\Lambda}(Z)|\leqslant \mathfrak M_{\Lambda}(Z)m(Z)$, and for a cluster $\mathfrak c=(D_1,\dots,D_n)$ with the positive incompatibility-tree majorant $|\Phi|^+_{\Lambda}(\mathfrak c)\coloneqq (1/n!)\sum_{T\in\mathrm{Trees}_n}\prod_{\{ i,j\}\in T}\mathbf 1\{ D_i\not\sim D_j\}\prod_im(D_i)$ (the $M_v(\mathfrak c)$ of Section~\ref{sec:c4-6-2};$|\Phi|^+_{\Lambda}(\mathfrak c)\geqslant \max_t|\Phi^{v,t}_{\Lambda}(\mathfrak c)|$ by the tree-graph inequality $|\varphi^T(\mathfrak c)|\leqslant \#\{$spanning trees of the incompatibility graph of $\mathfrak c\}$, proved in the justification following \eqref{eq:marginal-pinned-bound} and used in Appendix~\ref{app:kp}, and by $| k^{f,t}_{\Lambda}(D_i)|\leqslant m(D_i)$, Corollary~\ref{cor:c5-4-3}), $|\Phi^{v,1}_{\Lambda}(\mathfrak c)-\Phi^{v,0}_{\Lambda}(\mathfrak c)|\leqslant \bigl(\sum_i\mathfrak M_{\Lambda}(D_i)\bigr)|\Phi|^+_{\Lambda}(\mathfrak c)$.
\end{enumerate}
\end{lemma}

\begin{proof}
(a) In the display of Lemma~\ref{lem:c5-4-1} for the enlarged data, the $t$-dependent factors are $\chi^t_b$ ($b\in E(Y_0)$), $1-\chi^t_b$ ($b\in\mathsf P$) and the Mayer factors $M^t_Y\coloneqq \int_0^1\mathbf W_t(Y,B)e^{t_Y\mathbf W_t(Y,B)}\,\dd t_Y$ ($Y\in\mathcal D$), with $\mathbf W_t(Y,B)=\Theta^t_{E(Y)}(B)\mathbf P_t(Y,B)$, $\mathbf P_t(Y,B)\coloneqq \sum_{\mathfrak p\text{ near},\,Y_{\mathfrak p}= Y}\mathfrak p(B+tm+tA^{\Lambda,\Psi})$, $\Theta^t_{E(Y)}\coloneqq \prod_{b\in E(Y)}\chi^t_b$; the density, the Dirac tensor's pushforward (through $F$) and the $s$-, $\sigma$-integrations are $t$-independent. Order these factors (bonds first, then hosts) and telescope: $\prod a_k-\prod b_k=\sum_k(\prod_{i< k}a_i)(a_k-b_k)(\prod_{i> k}b_i)$, $a$ the $t=1$ and $b$ the $t=0$ factors. Each term has exactly one differenced factor, the mark. In every term each bond of $E(Y_0)\cup\mathsf P$ carries a factor $\chi^1_b$, $\chi^0_b$, $1-\chi^t_b$ or a difference of two of these, so on the support $| B_b|_{\sigma}<5r/4$ for $b\in E(Y_0)$ ($| B_b+tm_b|_{\sigma}< r$, $| m_b|_{\sigma}\leqslant r/4$), and every host's primitives are evaluated in the enlarged box (young part of $B+tm+tA^{\Lambda,\Psi}$ of norm $<3r/2\leqslant 2r$), where they are analytic with $|\mathfrak p|\leqslant C_pg_0\bar s^3$ and $\| D\mathfrak p\|\leqslant C_pg_0\bar s^2$ (Lemma~\ref{lem:c5-4-0}(c), Lemma~\ref{lem:c2-primitive-bounds} on the chart containing the enlarged box under threshold item 3), so $|\mathbf P_t(Y,B)|\leqslant N_hC_pg_0\bar s^3\leqslant \bar\omega$ for both $t$. Bond mark at $b$: $|\chi^1_b-\chi^0_b|\leqslant \mathbf 1_{S_b}(B_b)$ and $|(1-\chi^1_b)-(1-\chi^0_b)|\leqslant \mathbf 1_{S_b}(B_b)$, since the two indicators differ only where $| B_b|_{\sigma}$ is within $| m_b|_{\sigma}$ of $r$. Host mark at $Y\in\mathcal D$: in its term the bond factors are at $t=1$, and $E(Y)\subseteq E(Y_0)\setminus\mathsf P$ (Definition~\ref{def:selection}: $\mathsf P\cap E(Y_0)=\varnothing$), so $\Theta^1_{E(Y)}=1$ on the support and
\begin{equation*}
    M^1_Y-M^0_Y=\int_0^1\bigl[\mathbf P_1e^{t_Y\mathbf P_1}-\mathbf P_0e^{t_Y\mathbf P_0}\bigr]\dd t_Y+(1-\Theta^0_{E(Y)})\int_0^1\mathbf P_0e^{t_Y\mathbf P_0}\,\dd t_Y;
\end{equation*}
the smooth part is bounded on the support by $2\bar\omega e^{\bar\omega}$ always, and, when no primitive of $Y$ has an input in $J(\Lambda)$, by $\int_0^1|\partial_t(\mathbf P_te^{t_Y\mathbf P_t})|\,\dd t\leqslant (1+\bar\omega)e^{\bar\omega}\sup_t|\partial_t\mathbf P_t|$ with $\partial_t\mathbf P_t(Y,B)=\sum_{\mathfrak p}D\mathfrak p(B+tm)[m|_{E(Y)}]$ (no exterior input), $|\partial_t\mathbf P_t|\leqslant N_hC_pg_0\bar s^2\sum_{b\in E(Y)}h(m_b)$; the cutoff part satisfies $0\leqslant 1-\Theta^0_{E(Y)}\leqslant \sum_{b\in E(Y)}(1-\chi^0_b)$, and on the support of $\chi^1_b$, $(1-\chi^0_b)\chi^1_b\leqslant \mathbf 1_{S_b}(B_b)$, with $|\int_0^1\mathbf P_0e^{t_Y\mathbf P_0}\,\dd t_Y|\leqslant \bar\omega e^{\bar\omega}$: it is a sum over $b\in E(Y)$ of bond marks $\mathbf 1_{S_b}$ in which the host's factor is replaced by $\bar\omega e^{\bar\omega}$. Now bound each term. Take moduli; every factor other than the mark and the density is bounded by its supremum as in the proof of Lemma~\ref{lem:c5-4-2} (cutoffs $\leqslant 1$; bad bonds by $e^{-\frac9{16}a_{\mathrm{bad}}}e^{(\alpha_G/2)h(B_b)^2}$, from $\mathbf 1\{| B_b|_{\sigma}\geqslant 3r/4\}\leqslant e^{-\frac9{16}a_{\mathrm{bad}}}e^{(\alpha_G/2)h(B_b)^2}$; undifferenced Mayer factors by $\omega_*e^{1/2}e^{-\kappa d_0(Y)}$, valid at either $t$ since $\Theta^t_{E(Y)}$ restricts $B+tm$ to the local box of $Y$, where \eqref{eq:c5-4-4} and \eqref{eq:c5-5-T}(ii) give $|\mathbf W_t|\leqslant (\omega_*/2)e^{-\kappa d_0(Y)}$; $| F|\leqslant \sup| F|$); the modulus of the density is $e^{-\frac12\operatorname{Re}\langle\Gamma(\sigma)X,C\Gamma(\sigma)X\rangle}e^{-\langle B,v\rangle}\,\dd\mu_{C_{E_0}}(B)\,\dd\mu_0(X)$ with $v\coloneqq \operatorname{Re}\Gamma_{E_0}(\sigma)X$ real. At a bond mark $b\in E(Y_0)$ ($b\notin\mathsf P$): integrate over $B_b$ first; by Lemma~\ref{lem:c5-5-5} with $\alpha=0$ the mark $\mathbf 1_{S_b}$ is replaced by $\varpi_b$, leaving $\int e^{-\langle B,v\rangle}e^{(\alpha_G/2)\sum_{b'\in\mathsf P}h(B_{b'})^2}\,\dd\mu_C\leqslant \int e^{-\langle B,v\rangle}e^{(\alpha_G/2)\| B\|^2}\,\dd\mu_C$, which Lemma~\ref{lem:gaussian-payment}(ii) pays with $e^{N_1(Z)}$ together with the $X$-integral. At a bond mark $b\in\mathsf P$: on $S_b$, $| B_b|_{\sigma}\geqslant r-| m_b|_{\sigma}\geqslant 3r/4$, so $\mathbf 1_{S_b}(B_b)\leqslant \mathbf 1_{S_b}(B_b)e^{-\frac9{16}a_{\mathrm{bad}}}e^{(\alpha_G/2)h(B_b)^2}$, the reserve of $b$ is kept, and Lemma~\ref{lem:c5-5-5} with $\alpha=\alpha_G$ (the conditional Gaussian in $B_b$ with the tilt and the factor $e^{(\alpha_G/2)h(B_b)^2}$ on $b$) replaces the mark by $\varpi_b$ with all factors $e^{(\alpha_G/2)h(B_{b'})^2}$, $b'\in\mathsf P$, in place, so the bad-bond payment $e^{-\frac9{16}a_{\mathrm{bad}}|\mathsf P|}$ is complete and the Gaussian payment is as before. Hence a bond-marked term is bounded by $\varpi_b$ times the majorant \eqref{eq:c5-5-3}. At a host mark $Y$, the smooth part is a constant, $2\bar\omega e^{\bar\omega}$ or $N_hC_pg_0\bar s^2(1+\bar\omega)e^{\bar\omega}\sum_{E(Y)}h(m_b)$, in place of the host's majorant factor $\omega_*e^{1/2}e^{-\kappa d_0(Y)}\geqslant \omega_*e^{1/2}e^{-\kappa(B_h+4)}$; the quotient (constant)$/(\omega_*e^{1/2}e^{-\kappa d_0(Y)})$ is at most (constant)$\cdot e^{\kappa(B_h+4)}/\omega_*$ ($d_0(Y)\leqslant B_h+4$, $e^{1/2}\geqslant 1$), which is $\mathfrak m_Y$ without its factor $e^{1/8}$, reserved for the dummy comparison below; the cutoff part gives, for each $b\in E(Y)$, a bond mark $\mathbf 1_{S_b}$ ($b\notin\mathsf P$), paid by $\varpi_b$ as above, with the host's factor $\bar\omega e^{\bar\omega}$ in place of $\omega_*e^{1/2}e^{-\kappa d_0(Y)}$, a quotient $\leqslant \bar\omega e^{\bar\omega}e^{\kappa(B_h+4)}/\omega_*\leqslant e\bar\omega e^{\kappa(B_h+4)}/\omega_*\leqslant e^{-1/8}/N_{\mathrm{host}}\leqslant 1/N_{\mathrm{host}}$ by \eqref{eq:c5-5-T}(iii); the hosts $Y\in\mathcal D$ with $b\in E(Y)$ number at most $N_{\mathrm{host}}$ (Section~\ref{sec:c5-5-2}), so these contribute at most $\varpi_b$ per bond $b$, which with the bond's own mark $\varpi_b$ gives $\mathfrak m_b=2\varpi_b$; the rest of the proof of Lemma~\ref{lem:c5-4-2} (weakening per cube, Gaussian payment) is unchanged. Summing over the marks gives the first assertion with $\mathfrak M_{\Lambda}(Z)$ (every mark bond lies in $E(Y_0)\cup\mathsf P\subseteq E(Z)$, every marked host in $\mathcal D$ lies in $Z_0'\subseteq Z$). For the dummy comparison with a mark, repeat the positive host comparison of Lemma~\ref{lem:c5-5-2}(c) with the marked factor: at the bond's own mark (the differenced factor is $\chi^t_b$ or $1-\chi^t_b$) the hostwise comparison is unchanged (the undifferenced hosts have $v_{\mathrm{ordinary}}= e^{A_Y}-1$ at either $t$, as above) and the mark $\varpi_b$ multiplies the whole; at a bond mark generated by the cutoff part of a host $Y\in\mathcal D$, the host $Y$ is active through its ordinary label with its factor replaced by $\bar\omega e^{\bar\omega}$, so the sum over $Y$'s label subsets containing the ordinary label is $\bar\omega e^{\bar\omega}(1+u_*)^{n_Y}\leqslant \bar\omega e^{\bar\omega}e^{1/8}$, whose quotient by $Y$'s majorant factor $\omega_*e^{1/2}e^{-\kappa d_0(Y)}$ is at most $e^{9/8}\bar\omega e^{\kappa(B_h+4)}/\omega_*\leqslant 1/N_{\mathrm{host}}$ by \eqref{eq:c5-5-T}(iii) (this is what the factor $e^{9/8}$ in (iii) is for), and the at most $N_{\mathrm{host}}$ such hosts again contribute at most $\varpi_b$ per bond, so $\mathfrak m_b=2\varpi_b$ suffices in the dummy comparison as well; at a host mark $Y$ the ordinary weight $v_{\mathrm{ordinary}}$ of $Y$ is replaced by the mark's constant, and the sum over the label subsets $S$ containing the ordinary label (which carries the mark) is $(\text{constant})\cdot\sum_{S'\subseteq\{\text{probe labels of }Y\}}u_*^{| S'|}=(\text{constant})\cdot(1+u_*)^{n_Y}\leqslant (\text{constant})\cdot e^{1/8}$ ($n_Yu_*\leqslant \omega_*/8\leqslant 1/16<1/8$), so the host's contribution is $\leqslant \mathfrak m_Y\cdot\omega_*e^{1/2}e^{-\kappa d_0(Y)}$, the factor $e^{1/8}$ being the one in $\mathfrak m_Y$; the remaining sums give $m(Z)$ as in Lemma~\ref{lem:c5-5-2}(c), and dropping the constraint that the mark lies in the selection only enlarges the sum. (b) $k^{f,t}_{\Lambda}(Z)=\sum_{(\mathcal D,\mathsf P)}\mathcal A^{f,t}_{(\mathcal D,\mathsf P,Z)}[1]$ is the $\mathsf J=\varnothing$, $F=1$ case of (a), summed over the selections, which yields $\mathfrak M_{\Lambda}(Z)m(Z)$ by the same seed/enlargement sums; the cluster bound is the telescoping of the product $\prod_ik^{f,t}_{\Lambda}(D_i)$ with $| k^{f,t}_{\Lambda}(D_i)|\leqslant m(D_i)$ (Corollary~\ref{cor:c5-4-3}) and the tree-graph inequality for $|\varphi^T(\mathfrak c)|$.
\end{proof}

The mark weights are small where it matters: $\mathfrak m_b$ is of order $r^{d_G}$ times the decay from $b$ to $J_1(\Lambda)$; a host carrying an exterior input lies within the stencil range of $J(\Lambda)$ (its near primitive has an input in $J(\Lambda)$ and an input in $D(\Lambda)$), and its mark $\mathfrak m_Y\leqslant C_{\mathrm{bd}}g_0(1+3r)^3$ with $C_{\mathrm{bd}}\coloneqq 2e\cdot e^{\kappa(B_h+4)}e^{1/8}C_0/\omega_*$ is small in the coupling ($g_0(1+r)^3\to0$) though not by a threshold; a host without exterior input has $\mathfrak m_Y\leqslant C'_{\mathrm{bd}}g_0(1+3r)^2\sum_{E(Y)}h(m_b)$ with $C'_{\mathrm{bd}}\coloneqq 2e\cdot N_hC_pe^{\kappa(B_h+4)}e^{1/8}/\omega_*$ ($(1+\bar\omega)e^{\bar\omega}\leqslant 2e$ for $\bar\omega\leqslant 1$), which carries the decay of Lemma~\ref{lem:c5-5-1}.

\begin{lemma}[the finite-difference bound for closed contractions]\label{lem:c5-5-7z}
Let $\pi$ be the irreducible unitary representation with $\chi=\dim(\pi)^{-1}\operatorname{tr}\pi$ ($\pi$ also denotes the Lie-algebra representation, extended complex-linearly to $\mathfrak g_{\mathbb C}=\mathfrak g\oplus i\mathfrak g$, and $\pi$ denotes as well the holomorphic extension to the complexified group, with $\pi(e^{iX})= e^{i\pi(X)}$ for $X\in\mathfrak g_{\mathbb C}$), $h_{\mathbb C}(x+iy)^2\coloneqq h(x)^2+h(y)^2$ the Hermitian extension of $h$ to $\mathfrak g_{\mathbb C}$, $C_{\pi}\coloneqq \sup_{0\neq x\in\mathfrak g}\|\pi(x)\|_{\mathrm{op}}/h(x)$, so that $\|\pi(x+iy)\|_{\mathrm{op}}\leqslant C_{\pi}(h(x)+h(y))\leqslant \sqrt2C_{\pi} h_{\mathbb C}(x+iy)$, $m_*\coloneqq 1+\max_n|\mathcal F_{r_n}|$ (one more than the largest free-input set of a traversal link, $\mathcal F_r$ the full free-input set of the link at $r$, Section~\ref{sec:probe-kernels} and Lemma~\ref{lem:row-locality}; $\mathcal F_r\subseteq\mathcal E_c\cap E_*$ for the row $c$ of $r$, so $m_*$ depends on $L$ only), $N_{\mathrm{tr}}\coloneqq \max_b\#\{\text{traversals }n\text{ of }C\, : \,b\in\mathcal F_{r_n}\}$ ($b$ lies in the two-block set $\mathcal E_c$ of boundedly many rows $c$, and a traversal with input $b$ is the free bond $b$ itself or the designated bond $b_0(c)$ of such a row, each traversed once by the square, so $N_{\mathrm{tr}}$ depends on $L$ only). For a row $c$ within $2R_{\mathrm c}+2L$ of $I$, Lemmas~\ref{lem:charts} and~\ref{lem:row-locality} give the phase $q_U(B)|_{\mathcal E_c}$ as an analytic function $\Phi_c$ of the row phase $A= g_0\mathsf C_UB|_{\mathcal E_c}$ on a complex neighbourhood of the closed ball $\| A\|\leqslant \vartheta_*/2$ (the chart statements of Lemma~\ref{lem:c2-primitive-bounds}: the row solution and the phase are analytic there), with values in $\mathfrak g_{\mathbb C}^{\mathcal E_c}$ and $h_{\mathbb C}(\Phi_c(A)(b))\leqslant Q_*$ for every $b\in\mathcal E_c$ on that ball, $Q_*\coloneqq C^{\mathrm a}_0$ the order-zero bound of Lemma~\ref{lem:c2-uniform-chart}(iii) on the phase, in the Hilbert norm of $\mathfrak g^{\mathcal E_c}_{\mathbb C}$, which dominates $h_{\mathbb C}$ of each component ($Q_*$ depends on $L$ and the group only); put $C_{\Pi}\coloneqq e^{\sqrt2C_{\pi} Q_*}$ (so $\|\pi(e^{i\Phi_c(A)(b)})\|_{\mathrm{op}}=\| e^{i\pi(\Phi_c(A)(b))}\|_{\mathrm{op}}\leqslant e^{\|\pi(\Phi_c(A)(b))\|}\leqslant e^{\sqrt2C_{\pi} h_{\mathbb C}(\Phi_c(A)(b))}\leqslant C_{\Pi}$ on the ball) and $C_{\mathrm{link}}\coloneqq \max(C_{\pi},4m_*C_{\Pi} C_{\mathsf C}/\vartheta_*)$, $C'_{\chi}\coloneqq 2N_{\mathrm{tr}}C_{\mathrm{link}}$; all depend on $L$, the group and the character only. Under threshold item 3 of Appendix~\ref{app:thresholds} ($4C_{\mathsf C}g_0(1+2r)\leqslant \vartheta_*/2$), on $A^{\mathrm{fin}}_{\mathrm{cov}}$ the row $c$ of every traversal $n$ of $C$ has its free-input set $\mathcal F_{r_n}\subseteq D(\Lambda)\subseteq I$ (so the row is within $2L\leqslant 2R_{\mathrm c}+2L$ of $I$: $\mathcal F_{r_n}\subseteq\mathcal E_c$ lies in the two endpoint blocks of the row, of $\ell^{\infty}$ diameter at most $2L$), and the row phase $A= g_0\mathsf C_Ua|_{\mathcal E_c}$ depends on $a|_{\mathcal F_{r_n}}$ only (Lemma~\ref{lem:row-locality}), so the exterior variable $A^{\Lambda,\Psi}$ (supported on $J(\Lambda)$, disjoint from $D(\Lambda)$) does not enter a traversal link at all; for every real $a$ in the enlarged box $\{| a_b|_{\sigma}\leqslant 2r,\ b\in D(\Lambda)\}$ the input of the row lies in the row polydisc $\mathcal U_c=\{| B(b)|<2+4r\ \forall b\in\mathcal E_c\cap E_*\}$ of Lemmas~\ref{lem:charts} and~\ref{lem:c2-primitive-bounds} (the polydisc is in the norm of the enlarged box of Section~\ref{sec:source-phase-primitives}, which is the norm $|\bigcdot |_{\sigma}$ of the enlarged box here, Lemma~\ref{lem:c5-4-0}(c); every input bond has $| a_b|_{\sigma}\leqslant 2r<2+4r$), and the chart statements give $\| A\|\leqslant 2C_{\mathsf C}g_0(1+2r)\leqslant \vartheta_*/4$ for the row phase. At the reference data (Lemma~\ref{lem:c5-5-3}(b)) the exterior variable is zero and the same holds. ``The enlarged box'' below means this.
\begin{enumerate}[label=(\alph*)]
    \item (Link derivatives.) For a traversal link $L_n(a)\coloneqq U_{\xi}(a)(r_n)= e^{iq_U(\Psi_{\xi}(a+A_{\mathrm e}))(r_n)}U_1(r_n)$ (\eqref{eq:c-reconstructed-field}; it depends on $a$ through $a|_{\mathcal F_{r_n}}$ only, Lemma~\ref{lem:row-locality}), every real $a$ in the enlarged box, every $1\leqslant k\leqslant m_*$ (for $k=0$ the factor is unitary, of norm $1$) and every bonds $b_1,\dots,b_k\in\mathcal F_{r_n}$ (repetitions allowed) with $h$-unit directions $e_1,\dots,e_k\in\mathfrak g$: $\|\partial_{e_1,b_1}\cdots\partial_{e_k,b_k}\pi(L_n)(a)\|_{\mathrm{op}}\leqslant (C_{\mathrm{link}}g_0)^k$. Consequently $\|\partial_s\pi(L_n(a+sm))\|\leqslant C_{\mathrm{link}}g_0\sum_{b\in\mathcal F_{r_n}}h(m_b)$ and $\|\pi(L_n(a+m))-\pi(L_n(a))\|\leqslant C_{\mathrm{link}}g_0\sum_{b\in\mathcal F_{r_n}}h(m_b)$ whenever the segment from $a$ to $a+m$ lies in the enlarged box.
    \item (Mixed finite differences, with fixed insertions.) Let $g(a)=\dim(\pi)^{-1}\operatorname{tr}(X_1(a)\cdots X_p(a))$, where each factor $X_i(a)$ is either a traversal link $\pi(L_{n_i}(a))$ of $C$ or its inverse $\pi(L_{n_i}(a))^{-1}=\pi(L_{n_i}(a))^*$ (the orientation-reversed traversals of the Wilson word; at real points its real-direction derivatives are the adjoints of those of $\pi(L_{n_i})$, with the same norms, so (a) applies to it), with every traversal $n$ of $C$ occurring at most once among the link factors (as in the contour word, its inverse and their sub-words), or a fixed matrix $N_i\in\operatorname{End}(V_{\pi})$ not depending on $a$, and put $\| N\|\coloneqq \prod_{i\, : \,X_i\text{ fixed}}\| N_i\|_{\mathrm{op}}$ ($\| N\|=1$ when every fixed factor is unitary; the closed contractions $\chi(L_{n_1}(a)\cdots L_{n_j}(a)h_0)$, $h_0\in G$, with $h_0$ anywhere in the word, and the words with a fixed unitary insertion such as $\pi(P)^{-1}\pi(h)$ for a word $P$ of links at a fixed argument, are of this kind with $\| N\|=1$). Let $\mathsf J$ be a finite set of bonds of $D(\Lambda)$, $a_0$ real with $| a_{0,b}|_{\sigma}\leqslant r/4$ for all $b$ (the base point), and real $B$ with $| B_b|_{\sigma}\leqslant 5r/4$ for $b\in\mathsf J$. Then $\Delta_{\mathsf J}g(B)\coloneqq \sum_{K\subseteq\mathsf J}(-1)^{|\mathsf J\setminus K|}g(\mathbf 1_KB+a_0)$ satisfies
    \begin{gather*}
        |\Delta_{\mathsf J}g(B)|\leqslant (N_{\mathrm{tr}}C_{\mathrm{link}}g_0\cdot5r/4)^{|\mathsf J|}\| N\|\leqslant u_t^{|\mathsf J|}\| N\|\;,\\
        |\Delta_{\mathsf J}[\partial_{e,b}g](B)|\leqslant (N_{\mathrm{tr}}C_{\mathrm{link}}g_0\cdot5r/4)^{|\mathsf J|}N_{\mathrm{tr}}C_{\mathrm{link}}g_0\| N\|\leqslant u_t^{|\mathsf J|}(C'_{\chi}/2)g_0\| N\|
    \end{gather*}
    for every bond $b$ of $D(\Lambda)$ and $h$-unit $e$ ($u_t= C'_{\chi} g_0(1+2r)=2N_{\mathrm{tr}}C_{\mathrm{link}}g_0(1+2r)\geqslant N_{\mathrm{tr}}C_{\mathrm{link}}g_0\cdot5r/4$). Both hold at the reference data ($U=\id$, Lemma~\ref{lem:charts}) with the same constants.
\end{enumerate}
\end{lemma}

\begin{proof}
(a) Write $\pi(L_n(a))=\Pi_n(g_0\mathsf C_Ua|_{\mathcal E_c})\pi(U_1(r_n))$ (on $A^{\mathrm{fin}}_{\mathrm{cov}}$, where no exterior input enters the row, as in the preamble) with $\Pi_n(A)\coloneqq \pi(e^{i\Phi_c(A)(b_0(c))})= e^{i\pi(\Phi_c(A)(b_0(c)))}$ for the designated bond $r_n= b_0(c)$ (the component of the row phase at the designated bond), analytic on the neighbourhood of the closed ball $\| A\|\leqslant \vartheta_*/2$ (composition of the analytic $\Phi_c$ with the entire exponential of the linear map $\pi$ on $\mathfrak g_{\mathbb C}$) with $\|\Pi_n\|\leqslant C_{\Pi}$ there by the bound of the statement; for a free bond $r_n$, $q_U(B)(r_n)= g_0B(r_n)$ (Lemma~\ref{lem:row-locality}), so $\pi(L_n(a))= e^{ig_0\pi(a_{r_n})}\pi(U_1(r_n))$ depends on $a_{r_n}$ alone, and its $k$-fold derivative at a real point is bounded by $(C_{\pi} g_0)^k$ directly (Duhamel's formula: the $k$-th derivative of $e^{ig_0\pi(a)}$ in the $h$-unit directions $e_1,\dots,e_k$ is a sum over the $k!$ orderings of the $k$ insertions $ig_0\pi(e_i)$, each of norm $\leqslant C_{\pi} g_0$, between unitary factors, integrated over a simplex of volume $1/k!$). The directional derivative $\partial_{e,b}$ of $a\mapsto\Pi_n(g_0\mathsf C_Ua)$ is $g_0D\Pi_n(g_0\mathsf C_Ua)[\mathsf C_Ue_b]$ with $\|\mathsf C_Ue_b\|\leqslant C_{\mathsf C}$ (Lemma~\ref{lem:c2-uniform-chart}(i); $e_b$ the vector $e$ in the slot $b$, of Hilbert norm 1), and the $k$-fold mixed derivative is $g_0^kD^k\Pi_n(A)[\mathsf C_Ue_{b_1},\dots,\mathsf C_Ue_{b_k}]$. At a real $a$ in the enlarged box, $\| A\|\leqslant \vartheta_*/4$, so the ball of radius $\vartheta_*/4$ about $A$ lies in the closed ball of radius $\vartheta_*/2$; Cauchy's estimate for the analytic $\Pi_n$ on that ball bounds the homogeneous polynomial $D^k\Pi_n(A)[h,\dots,h]/k!$ by $C_{\Pi}(4\| h\|/\vartheta_*)^k$, and polarisation bounds the symmetric $k$-linear form $D^k\Pi_n(A)/k!$ by $k^k/k!$ times the polynomial's norm; hence $\| D^k\Pi_n(A)[h_1,\dots,h_k]\|\leqslant k^kC_{\Pi}(4/\vartheta_*)^k\prod\| h_i\|\leqslant (m_*C_{\Pi}4/\vartheta_*)^k\prod\| h_i\|$ for $k\leqslant m_*$ ($C_{\Pi}\geqslant 1$), giving $\|\partial_{e_1,b_1}\cdots\partial_{e_k,b_k}\pi(L_n)\|\leqslant (4m_*C_{\Pi} C_{\mathsf C}g_0/\vartheta_*)^k\leqslant (C_{\mathrm{link}}g_0)^k$; $\pi(U_1(r_n))$ is a constant unitary. The derivative and difference bounds are the chain rule and the mean value inequality along the segment, $\partial_s\pi(L_n(a+sm))=\sum_{b\in\mathcal F_{r_n}}h(m_b)\partial_{m_b/h(m_b),b}\pi(L_n)(a+sm)$, each term of norm $\leqslant h(m_b)C_{\mathrm{link}}g_0$ by the case $k=1$. (b) Write $q\coloneqq |\mathsf J|$ and $B^s\coloneqq (s_bB_b)_{b\in\mathsf J}+a_0$ for $s\in[0,1]^{\mathsf J}$. For each $b\in\mathsf J$, the difference in the label $b$ alone is $G(1)-G(0)=\int_0^1\partial_{s_b}G(s_b)\,\dd s_b$ for the function $G$ of $s_b$ at the other labels fixed, and iterating over the labels (the signs $(-1)^{|\mathsf J\setminus K|}$ are exactly the iterated differences) gives $\Delta_{\mathsf J}g(B)=\int_{[0,1]^{\mathsf J}}\partial_{s_{b_1}}\cdots\partial_{s_{b_q}}g(B^s)\,\dd s$ with $\partial_{s_b}g(B^s)=\langle B_b,\nabla_bg\rangle(B^s)= h(B_b)\partial_{B_b/h(B_b),b}g(B^s)$, so $|\Delta_{\mathsf J}g(B)|\leqslant \prod_{b\in\mathsf J}h(B_b)\cdot\sup_s|\partial^{\mathsf J}g(B^s)|$, where $\partial^{\mathsf J}$ is the mixed derivative in the $q$ directions $B_b/h(B_b)$ at the labels $b\in\mathsf J$; every $B^s$ is real and in the enlarged box ($| s_bB_b+a_{0,b}|_{\sigma}\leqslant 5r/4+r/4\leqslant 2r$). Product rule inside the normalised trace: the fixed factors $N_i$ do not depend on $a$, and a link $L_{n_i}$ depends on $a_b$ only if $b\in\mathcal F_{r_{n_i}}$, so $\partial^{\mathsf J}g=\sum_{\text{assignments}}\dim(\pi)^{-1}\operatorname{tr}(X_1^{(\mathsf J_1)}\cdots X_p^{(\mathsf J_p)})$, where $X_i^{(\mathsf J_i)}\coloneqq \partial^{\mathsf J_i}\pi(L_{n_i})$ for a link factor and $X_i^{(\varnothing)}\coloneqq N_i$ for a fixed factor, the sum over the assignments $\mathsf J=\bigsqcup_i\mathsf J_i$ of each label $b$ to one of the link factors with $b\in\mathcal F_{r_{n_i}}$, of which there are at most $N_{\mathrm{tr}}^q$; $|\operatorname{tr}(\prod M_i)|\leqslant \dim(\pi)\prod\| M_i\|_{\mathrm{op}}$, the undifferentiated links are unitary at real points, the fixed factors contribute $\prod\| N_i\|=\| N\|$, and each differentiated link has $|\mathsf J_i|\leqslant |\mathcal F_{r_{n_i}}|\leqslant m_*-1$ directions, so by (a) $\sup|\partial^{\mathsf J}g|\leqslant (N_{\mathrm{tr}}C_{\mathrm{link}}g_0)^q\| N\|$. With $h(B_b)\leqslant | B_b|_{\sigma}\leqslant 5r/4$ ($h\leqslant |\bigcdot |_{\sigma}$, as in the proof of Lemma~\ref{lem:c5-5-1}) this gives $|\Delta_{\mathsf J}g(B)|\leqslant (N_{\mathrm{tr}}C_{\mathrm{link}}g_0\cdot5r/4)^q\| N\|\leqslant (N_{\mathrm{tr}}C_{\mathrm{link}}g_0(1+2r))^q\| N\|\leqslant u_t^q\| N\|$, since $C'_{\chi}=2N_{\mathrm{tr}}C_{\mathrm{link}}\geqslant N_{\mathrm{tr}}C_{\mathrm{link}}$. For $\partial_{e,b}g$ the same argument applies to the mixed derivative $\partial^{\mathsf J}\partial_{e,b}g$ of order $q+1$ (the extra direction at the label $b$, possibly repeating a label of $\mathsf J$; each link then carries at most $m_*$ directions, within (a)): $|\Delta_{\mathsf J}[\partial_{e,b}g](B)|\leqslant (5r/4)^q(N_{\mathrm{tr}}C_{\mathrm{link}}g_0)^{q+1}\| N\|\leqslant u_t^q\cdot N_{\mathrm{tr}}C_{\mathrm{link}}g_0\| N\|= u_t^q(C'_{\chi}/2)g_0\| N\|$. At the reference data Lemma~\ref{lem:charts} holds for $U=\id$ with the same $\vartheta_*$, $C_{\mathsf C}$, $Q_*$, so (a) and (b) hold there verbatim.
\end{proof}

Lemma~\ref{lem:c5-5-7z} replaces the finite-difference bound of Section~\ref{sec:c4-4-1} (whose argument it follows: product rule inside the closed character, Cauchy estimates on the row chart, bounded traversal multiplicity) by a proof from Lemmas~\ref{lem:charts} and~\ref{lem:c2-uniform-chart} at the carrier $D(\Lambda)$; the parity gain of Section~\ref{sec:c4-4-1} is not used.

\begin{lemma}[the defect difference as a sum of marked connected terms]\label{lem:c5-5-7}
Let $\pi$ be the irreducible unitary representation with $\chi=\dim(\pi)^{-1}\operatorname{tr}\pi$, acting on $V_{\pi}$. For a group $W$ and $t\in\{0,1\}$ let $T^t_W(B)$ be the ordered product of the reconstructed links along the consecutive traversals of $\mathcal H(W)$ at the young variable $\mathbf 1_{\mathcal A(W)}B+tm$ (links off $\Omega_1$ equal to their background value), $P^t_W\coloneqq T^t_W(0)$ (young variable $tm$ throughout), and
\begin{equation*}
    \Upsilon^t_W\coloneqq w^t_{\Lambda}(W)\int\pi\bigl(T^t_W(B)(P^t_W)^{-1}\bigr)\,\dd\Bigl(\bigotimes_{Z\in\mathcal A(W)}\hat\nu^{f,t}_{\Lambda,Z}\Bigr)(B)-a^{\Lambda}_W\cdot\id_{V_{\pi}}\in\operatorname{End}(V_{\pi})\;,
\end{equation*}
so that $\dim(\pi)^{-1}\operatorname{tr}(\Upsilon^t_W\pi(h))= c^t_W(h)-a^{\Lambda}_W\chi(h)$ with $c^t_W(h)\coloneqq w^t_{\Lambda}(W)\int\chi(T^t_W(B)(P^t_W)^{-1}h)\,\dd\bigotimes\hat\nu^{f,t}_{\Lambda}$ (\eqref{eq:c4-6-2} at the young variable $tm$).
\begin{enumerate}[label=(\alph*)]
    \item (Burnside.) There is $K_{\chi}$ depending on $\pi$ only with $\|\Upsilon\|_{\mathrm{op}}\leqslant K_{\chi}\sup_{h\in G}|\dim(\pi)^{-1}\operatorname{tr}(\Upsilon\pi(h))|$ for every $\Upsilon\in\operatorname{End}(V_{\pi})$.
    \item For a compatible family $\mathfrak B=\{ W_1,\dots,W_n\}$ with short intervals in cyclic order (every nonempty compatible family is of this kind or a single whole-contour group: compatible groups have disjoint intervals, and a family containing a whole-contour group is a singleton, Section~\ref{sec:joint-extraction} after Definition~\ref{def:groups}), $R^t_{\Lambda}(\mathfrak B)=\dim(\pi)^{-1}\operatorname{tr}(\Upsilon^t_{W_1}U^t_1\Upsilon^t_{W_2}U^t_2\cdots\Upsilon^t_{W_n}U^t_n)$, where $U^t_j\coloneqq \pi(P^t_{W_j}G^t_j)$ with $G^t_j$ the ordered product of the links of the traversals strictly between $\mathcal H(W_j)$ and $\mathcal H(W_{j+1})$ (cyclically; $G^t_j=\id$ if there are none) at the young variable $tm$: the $t$-background word of the interval $\mathcal H(W_j)$ followed by the gap to the next interval (for a whole-contour group, $n=1$, $G^t_1=\id$ and $U^t_1=\pi(P^t_W)$).
    \item Hence
    \begin{gather}
        R^1_{\Lambda}(\mathfrak B)-R^0_{\Lambda}(\mathfrak B)=\sum_{j=1}^n\dim(\pi)^{-1}\operatorname{tr}\bigl(\Upsilon^1_{W_1}U^1_1\cdots(\Upsilon^1_{W_j}-\Upsilon^0_{W_j})U^0_j\cdots\Upsilon^0_{W_n}U^0_n\bigr)\notag\\
        +\sum_j\dim(\pi)^{-1}\operatorname{tr}\bigl(\cdots\Upsilon^1_{W_j}(U^1_j-U^0_j)\Upsilon^0_{W_{j+1}}\cdots\bigr)\;,\notag\\
        | R^1_{\Lambda}(\mathfrak B)-R^0_{\Lambda}(\mathfrak B)|\leqslant K_{\chi}^{|\mathfrak B|}\Bigl[\sum_{W_0\in\mathfrak B}\mathfrak d^{\mathrm{diff}}_{W_0}\prod_{W\neq W_0}\bar{\mathfrak d}_W+C'_{\chi} g_0\Bigl(\sum_{b\in I_C}h(m_b)\Bigr)\prod_{W\in\mathfrak B}\bar{\mathfrak d}_W\Bigr]\;,\label{eq:c5-5-7}
    \end{gather}
    with $\bar{\mathfrak d}_W\coloneqq \max_{t\in\{0,1\}}\sup_h| c^t_W(h)-a^{\Lambda}_W\chi(h)|$ and $\mathfrak d^{\mathrm{diff}}_W\coloneqq \sup_h| c^1_W(h)-c^0_W(h)|$.
    \item (Marks in $\mathfrak d^{\mathrm{diff}}_W$.) With $M^t_{\Lambda}(W)\coloneqq (1/k(W)!)\prod_{Z\in\mathcal A(W)}\rho_tm(Z)\prod_j|\Phi|^+_{\Lambda}(\mathfrak c_j)$ (the positive group majorant with root weight $\rho_tm(Z)$, the same for $t=0$ and $t=1$ since $\rho_t$ does not depend on $t$; write $M_{\Lambda}(W)$ for it; it is the $M(W)$ of Section~\ref{sec:c4-6-2} with $\rho=\rho_t$),
    \begin{gather}
        \mathfrak d^{\mathrm{diff}}_W\leqslant M_{\Lambda}(W)\Bigl[\sum_{Z\in\mathcal A(W)}\mathfrak M_{\Lambda}(Z)+\sum_j\sum_{D\in\mathfrak c_j}\mathfrak M_{\Lambda}(D)+C'_{\chi} g_0\sum_{b\in I_{\mathcal H(W)}}h(m_b)\Bigr]\;,\label{eq:c5-5-8}\\
        I_{\mathcal H(W)}\coloneqq \text{free inputs of the traversals of }\mathcal H(W),\notag
    \end{gather}
    where the last term is the source-shift mark (the root measures contribute their $\prod_Z\rho_tm(Z)$ through the finite-difference factor $u_t^{|\mathsf J|}$ exactly as in $M_{\Lambda}(W)$; the shift adds the factor $C'_{\chi} g_0\sum h(m_b)$).
    \item (Bounds for the undifferenced slots.) $\bar{\mathfrak d}_W\leqslant M_{\Lambda}(W)+| a^{\Lambda}_W|\eqqcolon \bar{\mathfrak b}_W$.
\end{enumerate}
\end{lemma}

\begin{proof}
(a) $\Upsilon\mapsto(h\mapsto\operatorname{tr}(\Upsilon\pi(h)))$ is linear and injective on $\operatorname{End}(V_{\pi})$ ($\pi(G)$ spans $\operatorname{End}(V_{\pi})$ for an irreducible representation, Burnside's theorem), so the sup over $h$ is a norm on the finite-dimensional $\operatorname{End}(V_{\pi})$, equivalent to the operator norm; $K_{\chi}$ is the equivalence constant. (b) is the factorisation of Lemma~\ref{lem:c4-6-1}(b) at the young variable $tm$: a traversal in $\mathcal H(W_j)$ has no contact with a root of another group, so its link uses $W_j$'s root variables (and $tm$); a traversal outside every interval has its link at $tm$; in the expansion of the product, a slot $W_j$ taken from the $\int\pi(T^t_{W_j}(P^t_{W_j})^{-1})$-term contributes $\pi(T^t_{W_j}(B))\pi(G^t_j)$ (the links of $\mathcal H(W_j)$ at the root variables, then the gap at $tm$) and a slot taken from $-a^{\Lambda}_{W_j}\id$ contributes $-a^{\Lambda}_{W_j}\pi(P^t_{W_j}G^t_j)$ (the links of $\mathcal H(W_j)$ and the gap at $tm$), so the normalised trace of the ordered product is $f_{\xi}(\mathbf 1_{\mathcal A(\mathcal S)}B+tm)$ for the subfamily $\mathcal S$ of slots of the first kind, times $\prod_{W\notin\mathcal S}(-a^{\Lambda}_W)$; expanding $\prod(X_W-a_W\id)$ over subsets reproduces the signed sum $R^t_{\Lambda}(\mathfrak B)$ of Lemma~\ref{lem:c5-5-3}(c) by multilinearity of the trace, and the Fubini evaluation of the $\int\pi(\dots)$-slots against the product measure is the one in the proof of Lemma~\ref{lem:schur}. (c) Telescoping of the ordered product in its $2n$ factors; then $|\operatorname{tr}(X_1\cdots X_{2n})|\leqslant \dim(\pi)\prod\| X_i\|_{\mathrm{op}}$, $\| U^t_j\|=1$, $\|\Upsilon^t_W\|\leqslant K_{\chi}\bar{\mathfrak d}_W$, $\|\Upsilon^1_W-\Upsilon^0_W\|\leqslant K_{\chi}\mathfrak d^{\mathrm{diff}}_W$ by (a), and $\| U^1_j-U^0_j\|\leqslant \sum_{\text{traversals in }\mathcal H(W_j)\text{ and in the gap after it}}\|\pi(\text{link at }m)-\pi(\text{link at }0)\|\leqslant C'_{\chi} g_0\sum_{b\in I_C}h(m_b)$ (the links are unitary and the word is their ordered product, so the difference is at most the sum of the link differences; $\|\pi(L_n(m))-\pi(L_n(0))\|\leqslant C_{\mathrm{link}}g_0\sum_{b\in\mathcal F_{r_n}}h(m_b)$ by Lemma~\ref{lem:c5-5-7z}(a), the segment $sm$ lying in the box of radius $r/4$; over $j$ the intervals and gaps partition the traversals of $C$, each traversal's inputs lie in $I_C$, and a bond is an input of at most $N_{\mathrm{tr}}$ traversals, so the total is at most $N_{\mathrm{tr}}C_{\mathrm{link}}g_0\sum_{b\in I_C}h(m_b)\leqslant C'_{\chi} g_0\sum_{b\in I_C}h(m_b)$). (d) $c^1_W-c^0_W$ is telescoped over its $t$-dependent slots: the $k(W)$ cluster weights $\Phi^{v,t}_{\Lambda}(\mathfrak c_j)$, the $|\mathcal A(W)|$ root measures $\hat\nu^{f,t}_{\Lambda,Z}$ (through their split selections $\ell_Z$, Lemma~\ref{lem:c5-5-2}(a)), and the source word $g_t(B)\coloneqq \chi(T^t_W(B)(P^t_W)^{-1}h)$. Cluster slot: $|\Phi^{v,1}-\Phi^{v,0}|(\mathfrak c_j)\leqslant \sum_{D\in\mathfrak c_j}\mathfrak M_{\Lambda}(D)|\Phi|^+(\mathfrak c_j)$ (Lemma~\ref{lem:c5-5-6}(b)), the other clusters at $|\Phi^{v,t}|\leqslant |\Phi|^+$, and the root integral $|\int\Delta_{\mathsf J}g_t\,\dd\bigotimes\Lambda^{f,t}_{\ell_Z}|\leqslant \prod_Zu_t^{|\mathsf J_Z|}\|\Lambda^{f,t}_{\ell_Z}\|_{\mathrm{TV}}$ summed with Lemma~\ref{lem:c5-5-2}(c) and $(u_t/u_*)^{|\mathsf J_Z|}\leqslant \rho_t$ gives $\prod_Z\rho_tm(Z)$. Here $\sup|\Delta_{\mathsf J}g_t|\leqslant u_t^{|\mathsf J|}$ is the first assertion of Lemma~\ref{lem:c5-5-7z}(b) for $g_t(B)=\dim(\pi)^{-1}\operatorname{tr}(\pi(T_W(\mathbf 1_KB+tm))N_t)$ with the fixed unitary factor $N_t\coloneqq \pi(P^t_W)^{-1}\pi(h)$ ($\| N_t\|=1$: $P^t_W$ is the word of $\mathcal H(W)$ at the fixed argument $tm$ and $h$ a fixed group element) at the base point $a_0= tm$ ($| tm_b|_{\sigma}\leqslant r/4$) with the increments $\mathbf 1_KB$, $| B_b|_{\sigma}\leqslant 5r/4$ on the support of $\nu^{f,t}$. Root slot: one root $Z_0$ carries $\Lambda^{f,1}_{\ell}-\Lambda^{f,0}_{\ell}$ and Lemma~\ref{lem:c5-5-6}(a) gives $\sum_{\ell} u_*^{|\mathsf J|}\|\Lambda^1_{\ell}-\Lambda^0_{\ell}\|\leqslant \mathfrak M_{\Lambda}(Z_0)m(Z_0)$, with the same $u_t^{|\mathsf J|}$ from the finite difference, so the root pays $\mathfrak M_{\Lambda}(Z_0)\rho_tm(Z_0)$. Source slot: all measures and weights at one $t$, and the insertion is $\Delta_{\mathsf J}[g_1-g_0](B)=\int_0^1\Delta_{\mathsf J}[\partial_sg_s](B)\,\dd s$ with $g_s(a)\coloneqq \dim(\pi)^{-1}\operatorname{tr}(\pi(T_W(a))\pi(P^s_W)^{-1}\pi(h))$, evaluated in the $K$-term of $\Delta_{\mathsf J}$ at $a=\mathbf 1_KB+sm$, $T_W(a)$ the word of $\mathcal H(W)$ at the young variable $a$ and $P^s_W= T_W(sm)$ (the de-younged bonds are constants, so only the inputs in $I_{\mathcal H(W)}\cap D(\Lambda)$ vary; $\partial_s$ below is the derivative of $s\mapsto g_s(\mathbf 1_KB+sm)$ in each $K$-term). Two things depend on $s$: the argument $\mathbf 1_KB+sm$ of the link word and the fixed factor $N_s\coloneqq \pi(P^s_W)^{-1}\pi(h)$, whose links sit at the argument $sm$ and not at $\mathbf 1_KB+sm$. Hence $\partial_sg_s=\sum_{b\in I_{\mathcal H(W)}\cap D(\Lambda)}h(m_b)\partial_{m_b/h(m_b),b}\tilde g_s+\dim(\pi)^{-1}\operatorname{tr}(\pi(T_W(\mathbf 1_KB+sm))N'_s)$, with $\tilde g_s(a)\coloneqq \dim(\pi)^{-1}\operatorname{tr}(\pi(T_W(a))N_s)$ (fixed unitary insertion, $\| N_s\|=1$) and $N'_s\coloneqq \partial_s[\pi(P^s_W)^{-1}]\pi(h)=-\pi(P^s_W)^{-1}(\partial_s\pi(P^s_W))\pi(P^s_W)^{-1}\pi(h)$, a fixed non-unitary insertion with $\| N'_s\|\leqslant \|\partial_s\pi(P^s_W)\|\leqslant \sum_{n\in\mathcal H(W)}\|\partial_s\pi(L_n(sm))\|\leqslant \sum_{n\in\mathcal H(W)}C_{\mathrm{link}}g_0\sum_{b\in\mathcal F_{r_n}}h(m_b)\leqslant N_{\mathrm{tr}}C_{\mathrm{link}}g_0\sum_{b\in I_{\mathcal H(W)}}h(m_b)$ (product rule for the word of unitaries; Lemma~\ref{lem:c5-5-7z}(a) along the segment $sm$, which lies in the box of radius $r/4$; each bond is an input of at most $N_{\mathrm{tr}}$ traversals). The mixed difference of order $|\mathsf J|$ of the first sum is bounded by the second assertion of Lemma~\ref{lem:c5-5-7z}(b) (base point $sm$, increments $\mathbf 1_KB$, $\| N_s\|=1$): $\leqslant u_t^{|\mathsf J|}(C'_{\chi}/2)g_0\sum_{b\in I_{\mathcal H(W)}}h(m_b)$; that of the second term by the first assertion with the insertion $N'_s$: $\leqslant u_t^{|\mathsf J|}\| N'_s\|\leqslant u_t^{|\mathsf J|}N_{\mathrm{tr}}C_{\mathrm{link}}g_0\sum_{b\in I_{\mathcal H(W)}}h(m_b)= u_t^{|\mathsf J|}(C'_{\chi}/2)g_0\sum_{b\in I_{\mathcal H(W)}}h(m_b)$. Together $\sup|\Delta_{\mathsf J}[\partial_sg_s]|\leqslant u_t^{|\mathsf J|}C'_{\chi} g_0\sum_{b\in I_{\mathcal H(W)}}h(m_b)$: the factor 2 in $C'_{\chi}=2N_{\mathrm{tr}}C_{\mathrm{link}}$ pays the two $s$-dependent words. This gives \eqref{eq:c5-5-8}; the three sums are over the marks of the three kinds. (e) $| c^t_W(h)|\leqslant M^t_{\Lambda}(W)$ by the same root and cluster bounds (no mark), and $| a^{\Lambda}_W\chi(h)|\leqslant | a^{\Lambda}_W|$.
\end{proof}

Nothing in Lemma~\ref{lem:c5-5-7} uses a parity (odd-root) cancellation: the translated cutoffs $\chi^t_b$ are not even in $B$ for $t>0$, so the reversal argument of Section~\ref{sec:c4-4-1} is unavailable and only the first-order gain $\rho_t= u_t/u_*$ per root is taken; this is why $\rho_t= Cg_0(1+2r)$ rather than $Cg_0^2(1+r)^4$ appears. Section~\ref{sec:c5-5R} recovers the second order for the undifferenced slots at $t\in\{0,1\}$ and for the activities $a^{\Lambda}_W$ (Lemma~\ref{lem:c5-5R-3}); the differenced slots and the marked part are unchanged.

\begin{lemma}[pins for the $t$-family]\label{lem:c5-5-8}
Let $F(W)$ be the footprint of $W$ (Section~\ref{sec:c4-6-2}: the $LM$-cubes meeting its roots, cluster constituents, their stencil halos and the stencils of the traversals of $\mathcal H(W)$), $\sigma(W)\coloneqq s(W)= n(W)+|\mathcal H(W)|/(LM)$ its size (\eqref{eq:c4-6-1}; $n(W)$ counts the cubes of the roots and cluster constituents with tuple multiplicity, so $\sigma(W)$ is at least the number of domain occurrences in $W$), $\nu\coloneqq \eta$ the reserve constant of Section~\ref{sec:c4-6-2}, and $C_F\coloneqq A_E$ the envelope constant of Lemma~\ref{lem:c4-6-env} ($|\mathfrak E(W)|\leqslant A_E\sigma(W)$). (a) Under \eqref{eq:c5-5-T}(iv) the positive summation of Lemma~\ref{lem:c4-6-2} applies to $M_{\Lambda}(W)$ (its hypotheses hold: $C\rho_tLM\leqslant \eta/8$ is part of $\rho_t\leqslant \rho_{\mathrm{pin}}$, Section~\ref{sec:c5-5-2}, and the root-incidence bound of Section~\ref{sec:c4-8-2} is a property of the fixed geometry, independent of the carrier; it uses only the root weights $\rho\,m(Z)$ with $\rho=\rho_t\leqslant \rho_{\mathrm{pin}}$, the positive incompatibility-tree weights $|\Phi|^+_{\Lambda}= M_v$, and the field-independent domain/cluster/interval geometry, none of which involves the carrier beyond the cubes): with $C_{\mathrm{pin}}$ the constant of \eqref{eq:c4-6-5},
\begin{equation*}
    \sup_{\Delta}\sum_{W\, : \,\Delta\in F(W)}e^{4\nu\sigma(W)}M_{\Lambda}(W)+\sup_n\sum_{W\, : \,n\in\mathcal H(W)}e^{4\nu\sigma(W)}M_{\Lambda}(W)\leqslant C_{\mathrm{pin}}\rho_t\;.
\end{equation*}
Consequently, with $\bar\varepsilon_t\coloneqq 2K_{\chi} C_{\mathrm{pin}}\rho_t$,
\begin{gather*}
    \sup_{\Delta}\sum_{W\, : \,\Delta\in F(W)}e^{2\nu\sigma(W)}K_{\chi} e^{\theta^{\circ}(W)}\bar{\mathfrak b}_W\leqslant \bar\varepsilon_t\;,\qquad\sup_n\sum_{W\, : \,n\in\mathcal H(W)}e^{2\nu\sigma(W)}K_{\chi} e^{\theta^{\circ}(W)}\bar{\mathfrak b}_W\leqslant \bar\varepsilon_t\;,\\
    \sum_WK_{\chi} e^{\theta^{\circ}(W)}\bar{\mathfrak b}_W\leqslant n_C\bar\varepsilon_t\;,
\end{gather*}
the last because every root of a group is marked (Definition~\ref{def:groups}), so $\mathcal H(W)\neq\varnothing$ and $W$ is counted in $\sum_n$ for some $n$. The same statements hold for the reference gas at the carrier $I$ (the reference data of Section~\ref{sec:sourced-effective-interaction}, activities $a_W$) with the same constants: the proof uses only the lemmas of Section~\ref{sec:sourced-effective-interaction} and Appendix~\ref{app:background-estimates} at the carrier, which hold at $I$. (b) (Decay from the footprint to the contour.) If $\Delta\in F(W)$ then $\sigma(W)\geqslant \operatorname{dist}(\Delta,C)/(C_FLM)$: the envelope $\mathfrak E(W)\supseteq F(W)\cup\pi_1(\mathcal H(W))$ is connected under maximum-norm adjacency of $LM$-cubes and has at most $A_E\sigma(W)$ cubes (Lemma~\ref{lem:c4-6-env}), so its $\ell^{\infty}$ diameter in fine units is at most $A_ELM\sigma(W)$, and it contains $\Delta$ and a cube meeting the contour. Hence for every cube $\Delta$,
\begin{gather*}
    \sum_{W\, : \,\Delta\in F(W)}e^{\nu\sigma(W)}K_{\chi} e^{\theta^{\circ}(W)}\bar{\mathfrak b}_W\leqslant \bar\varepsilon_te^{-\nu\operatorname{dist}(\Delta,C)/(C_FLM)}\;,\\
    \sum_{W\, : \,\Delta\in F(W)}\sigma(W)K_{\chi} e^{\theta^{\circ}(W)}\bar{\mathfrak b}_W\leqslant \nu^{-1}\bar\varepsilon_te^{-\nu\operatorname{dist}(\Delta,C)/(C_FLM)}\;.
\end{gather*}
\end{lemma}

\begin{proof}
(a) The first display is the cited summation with $\rho_t$ for the root weights; the second follows from $\bar{\mathfrak b}_W\leqslant 2M_{\Lambda}(W)$ (Lemma~\ref{lem:c5-5-3}(d) and Lemma~\ref{lem:c5-5-7}(e)) and $e^{\theta^{\circ}(W)}\leqslant e^{\nu\sigma(W)/16}\leqslant e^{2\nu\sigma(W)}$ (Lemma~\ref{lem:c5-5-3}(d)). (b) $e^{2\nu\sigma(W)}\geqslant e^{\nu\sigma(W)}e^{\nu\operatorname{dist}(\Delta,C)/(C_FLM)}$ by the geometric fact, and $\sigma\leqslant \nu^{-1}e^{\nu\sigma}$ ($x\leqslant e^x$).
\end{proof}

\begin{proposition}[the source-containing connected boundary terms]\label{prop:c5-5-9}
In the square/torus geometry of Section~\ref{sec:c5-5-0}, under the hypotheses of Lemmas~\ref{lem:c5-4-0}--\ref{lem:c5-4-4} and Corollary~\ref{cor:c5-4-3}, \eqref{eq:c5-4-T}, \eqref{eq:c5-4-Tp} and \eqref{eq:c5-5-T}, on every labelled term with final region $\Lambda$, for exterior data in the support of the term's collar indicators and exterior cutoffs, on $A^{\mathrm{fin}}_{\mathrm{cov}}$,
\begin{gather}
    |\mathcal B^{\Lambda}_{\xi}[f_{\xi}]|\leqslant | S^{\circ}_{\Lambda}| e^{2n_C\bar\varepsilon_t}\bigl[C_1g_0r\,\mathfrak S_1(\Lambda)+C_2\bar\varepsilon_tr^{d_G}\mathfrak S_2(\Lambda)\bigr]\;,\label{eq:c5-5-9}\\
    \mathfrak S_1(\Lambda)\coloneqq \sum_{e\in J_1(\Lambda)}e^{-\mu_H\operatorname{dist}(e,I_C)/2}\;,\qquad\mathfrak S_2(\Lambda)\coloneqq \sum_{e\in J(\Lambda)}e^{-\mu'\operatorname{dist}(e,C)}\;,\qquad\mu'\coloneqq \min\bigl(\mu_H/4,\,\nu/(2C_FLM)\bigr)\;,\notag
\end{gather}
with $C_1\coloneqq C_m\vartheta N_{\mathrm{lat}}(\mu_H/2)(C_{\mathrm{resp},\chi}+(1+N_{\mathrm{tr}})C'_{\chi})$ and $C_2\coloneqq \nu^{-1}C_{\mathrm{dist}}[4C_{\varpi} N_{\mathrm{lat}}(\mu_H/2)+N_{\mathrm{host}}e^{\nu R_{\mathrm{st}}/(C_FLM)}(64C_{\mathrm{bd}}+16C'_{\mathrm{bd}}C_m\vartheta N_{\mathrm{lat}}(\mu_H/2))]$, $C_{\mathrm{dist}}\coloneqq e^{7\nu/C_F}$, $R_{\mathrm{st}}\coloneqq 4M$ ($\geqslant 4L+3M$, since $M\geqslant 8LM_2\geqslant 8L^2\geqslant 4L$ by \eqref{eq:c-standing-conditions}), constants depending on $L$, the group, $M$, $\kappa$, $\omega_*$, $\nu$, $\vartheta$, the character and the stencil ranges only, in the square/torus geometry of Section~\ref{sec:c5-5-0}; $g_0r\to0$ and $\bar\varepsilon_t=2K_{\chi} C_{\mathrm{pin}}\rho_t\to0$ as $g_0\to0$.
\end{proposition}

\begin{proof}
Divide \eqref{eq:c5-5-5} by $S^{\circ}_{\Lambda}$ (Lemma~\ref{lem:c5-5-3}(d)): $\mathcal B^{\Lambda}/S^{\circ}_{\Lambda}=[f_{\xi}(m)-f_{\xi}(0)]+\sum_{\mathfrak B\neq\varnothing}q^{\circ}_{\Lambda}(\mathfrak B)[R^1_{\Lambda}(\mathfrak B)-R^0_{\Lambda}(\mathfrak B)]$. The first term is Lemma~\ref{lem:c5-5-4}: $\leqslant C_{\mathrm{resp},\chi}C_m\vartheta N_{\mathrm{lat}}(\mu_H/2)g_0r\mathfrak S_1(\Lambda)$. For the second, insert \eqref{eq:c5-5-7} and $| q^{\circ}_{\Lambda}(\mathfrak B)|\leqslant \prod_We^{\theta^{\circ}(W)}$. The background-word part:
\begin{equation*}
    C'_{\chi} g_0\Bigl(\sum_{I_C}h(m_b)\Bigr)\sum_{\mathfrak B}\prod_{W\in\mathfrak B}K_{\chi} e^{\theta^{\circ}(W)}\bar{\mathfrak d}_W\leqslant C'_{\chi} g_0\Bigl(\sum_{I_C}h(m_b)\Bigr)\prod_W(1+K_{\chi} e^{\theta^{\circ}(W)}\bar{\mathfrak b}_W)\leqslant C'_{\chi} g_0\Bigl(\sum_{I_C}h(m_b)\Bigr)e^{n_C\bar\varepsilon_t}
\end{equation*}
(families are sets of distinct groups; Lemma~\ref{lem:c5-5-8}(a)), and $\sum_{I_C}h(m_b)\leqslant C_m\vartheta N_{\mathrm{lat}}(\mu_H/2)r\mathfrak S_1(\Lambda)$ by the last clause of Lemma~\ref{lem:c5-5-1}: together with the first term this is $C_m\vartheta N_{\mathrm{lat}}(\mu_H/2)(C_{\mathrm{resp},\chi}+C'_{\chi})g_0r\mathfrak S_1(\Lambda)e^{n_C\bar\varepsilon_t}$, the first part of the $C_1$ term (the second part is (ii) below). The marked part:
\begin{equation*}
    \sum_{\mathfrak B}\sum_{W_0\in\mathfrak B}\prod_We^{\theta^{\circ}(W)}K_{\chi}\mathfrak d^{\mathrm{diff}}_{W_0}\prod_{W\neq W_0}K_{\chi}\bar{\mathfrak d}_W\leqslant e^{n_C\bar\varepsilon_t}\sum_{W_0}K_{\chi} e^{\theta^{\circ}(W_0)}\mathfrak d^{\mathrm{diff}}_{W_0}
\end{equation*}
(sum first over $\mathfrak B\setminus\{ W_0\}$, as above). Insert \eqref{eq:c5-5-8}. (i) Root and cluster marks. Write $\mathfrak M_{\Lambda}(\Delta)\coloneqq \sum_{b\in E(\Delta)}\mathfrak m_b+\sum_{Y\text{ anchored at }\Delta}\mathfrak m_Y$ for an $LM$-cube $\Delta$ (a host is anchored at the $LM$-cube containing its lowest $M$-cube; a bond meeting two cubes is counted in both), so that $\mathfrak M_{\Lambda}(Z)\leqslant \sum_{\Delta\subseteq Z}\mathfrak M_{\Lambda}(\Delta)$ for every domain $Z$. The occurrences of a cube $\Delta$ among the roots and cluster constituents of $W_0$ number $\operatorname{mult}_{W_0}(\Delta)\leqslant $ (number of domain occurrences in $W_0$) $\leqslant \sigma(W_0)$ (Lemma~\ref{lem:c5-5-8}), so $\sum_{Z\in\mathcal A(W_0)}\mathfrak M_{\Lambda}(Z)+\sum_j\sum_{D\in\mathfrak c_j}\mathfrak M_{\Lambda}(D)\leqslant \sum_{\Delta\in F(W_0)}\operatorname{mult}_{W_0}(\Delta)\mathfrak M_{\Lambda}(\Delta)\leqslant \sigma(W_0)\sum_{\Delta\in F(W_0)}\mathfrak M_{\Lambda}(\Delta)$, and $\sum_{W_0}K_{\chi} e^{\theta^{\circ}(W_0)}M_{\Lambda}(W_0)\sigma(W_0)\sum_{\Delta\in F(W_0)}\mathfrak M_{\Lambda}(\Delta)\leqslant \sum_{\Delta}\mathfrak M_{\Lambda}(\Delta)\sum_{W_0\, : \,\Delta\in F(W_0)}\sigma(W_0)K_{\chi} e^{\theta^{\circ}}\bar{\mathfrak b}_{W_0}\leqslant \nu^{-1}\bar\varepsilon_t\sum_{\Delta}\mathfrak M_{\Lambda}(\Delta)e^{-\nu\operatorname{dist}(\Delta,C)/(C_FLM)}$ by Lemma~\ref{lem:c5-5-8}(b). Convert cube distances to bond and host distances: for $b\in E(\Delta)$, $\operatorname{dist}(\Delta,C)\geqslant \operatorname{dist}(b,C)-LM-1$, and for $Y$ anchored at $\Delta$, $\operatorname{dist}(\Delta,C)\geqslant \operatorname{dist}(Y,C)-4M-LM$ (a host has side at most $4M$), so $e^{-\nu\operatorname{dist}(\Delta,C)/(C_FLM)}\leqslant C_{\mathrm{dist}}e^{-\nu\operatorname{dist}(b,C)/(C_FLM)}$, resp. $\leqslant C_{\mathrm{dist}}e^{-\nu\operatorname{dist}(Y,C)/(C_FLM)}$, with $C_{\mathrm{dist}}\coloneqq e^{7\nu/C_F}$ ($(LM+1)/(LM)\leqslant 2$ and $(4M+LM)/(LM)\leqslant 7/3$, both at most $7$). Now the three kinds of $\mathfrak m$-weights. Bond marks: a bond lies in $E(\Delta)$ for at most two cubes $\Delta$ (no site lies on a cube face), so $\sum_{\Delta}\sum_{b\in E(\Delta)}\mathfrak m_be^{-\nu\operatorname{dist}(\Delta,C)/(C_FLM)}\leqslant 2C_{\mathrm{dist}}\sum_{b\in D(\Lambda)}\mathfrak m_be^{-\nu\operatorname{dist}(b,C)/(C_FLM)}$, and
\begin{multline*}
    \sum_{b\in D(\Lambda)}\mathfrak m_be^{-\nu\operatorname{dist}(b,C)/(C_FLM)}\leqslant 2C_{\varpi} r^{d_G}\sum_b\sum_{e\in J_1}e^{-\mu_Hd(b,e)}e^{-\nu\operatorname{dist}(b,C)/(C_FLM)}\\
    \leqslant 2C_{\varpi} r^{d_G}N_{\mathrm{lat}}(\mu_H/2)\sum_{e\in J_1(\Lambda)}e^{-\mu'\operatorname{dist}(e,C)}
\end{multline*}
(so the bond term of $C_2$ is $4C_{\varpi} N_{\mathrm{lat}}(\mu_H/2)$ after the factor $C_{\mathrm{dist}}$ is taken out) by the splitting in the proof of \eqref{eq:c5-5-2} with $\mu'\leqslant \min(\mu_H/2,\nu/(C_FLM))/2$. Host marks with an exterior input $e\in J(\Lambda)$: the primitive's stencil has $\ell^{\infty}$ diameter at most $4L$ and contains its input $e$, and its host lies in the $M$-cubes meeting the stencil with one further $M$-layer (\eqref{eq:part1-potential-host}), so every point of the host is within $4L+3M\leqslant 4M= R_{\mathrm{st}}$ of $e$ ($M\geqslant 8L^2\geqslant 4L$), $\operatorname{dist}(Y,C)\geqslant \operatorname{dist}(e,C)-R_{\mathrm{st}}$, and there are at most $N_{\mathrm{host}}$ such hosts per $e$ (Section~\ref{sec:c5-5-2}); their total is $\leqslant N_{\mathrm{host}}e^{\nu R_{\mathrm{st}}/(C_FLM)}C_{\mathrm{bd}}g_0(1+3r)^3\sum_{e\in J(\Lambda)}e^{-\mu'\operatorname{dist}(e,C)}$ ($\nu/(C_FLM)\geqslant \mu'$). Host marks without exterior input: $\sum_Y\mathfrak m_Ye^{-\nu\operatorname{dist}(Y,C)/(C_FLM)}\leqslant C'_{\mathrm{bd}}g_0(1+3r)^2\sum_{b\in D(\Lambda)}h(m_b)\cdot\#\{ Y\in\mathbf D_0(\Lambda)\, : \,b\in E(Y)\}\cdot e^{\nu R_{\mathrm{st}}/(C_FLM)}e^{-\nu\operatorname{dist}(b,C)/(C_FLM)}\leqslant C'_{\mathrm{bd}}N_{\mathrm{host}}e^{\nu R_{\mathrm{st}}/(C_FLM)}g_0(1+3r)^2C_m\vartheta rN_{\mathrm{lat}}(\mu_H/2)\sum_{e\in J_1}e^{-\mu'\operatorname{dist}(e,C)}$ by \eqref{eq:c5-5-2} ($\operatorname{dist}(Y,C)\geqslant \operatorname{dist}(b,C)-4M=\operatorname{dist}(b,C)-R_{\mathrm{st}}$ for $b\in E(Y)$, a host having side at most $4M$; at most $N_{\mathrm{host}}$ hosts contain $b$). (ii) Source-shift marks: $\sum_{W_0}K_{\chi} e^{\theta^{\circ}}M_{\Lambda}(W_0)C'_{\chi} g_0\sum_{b\in I_{\mathcal H(W_0)}}h(m_b)\leqslant C'_{\chi} g_0\sum_{b\in I_C}h(m_b)\sum_{n\in n(b)}\sum_{W_0\, : \,n\in\mathcal H(W_0)}K_{\chi} e^{\theta^{\circ}}\bar{\mathfrak b}_{W_0}\leqslant N_{\mathrm{tr}}\bar\varepsilon_tC'_{\chi} g_0\sum_{b\in I_C}h(m_b)$ (the first step is an inequality: a bond $b\in I_{\mathcal H(W_0)}$ is an input of at least one traversal $n\in\mathcal H(W_0)\cap n(b)$, so every term on the left occurs at least once on the right, possibly several times when $b$ is an input of several traversals of $\mathcal H(W_0)$) $\leqslant N_{\mathrm{tr}}C'_{\chi} g_0C_m\vartheta N_{\mathrm{lat}}(\mu_H/2)r\mathfrak S_1(\Lambda)$ ($n(b)$ the set of traversals with input $b$, $| n(b)|\leqslant N_{\mathrm{tr}}$, Lemma~\ref{lem:c5-5-7z}; the contour pin of Lemma~\ref{lem:c5-5-8}(a) for each $n\in n(b)$; $\bar\varepsilon_t\leqslant 1$ by \eqref{eq:c5-5-T}(v); then \eqref{eq:c5-5-2}); this is the second part of the $C_1$ term, $C_1= C_m\vartheta N_{\mathrm{lat}}(\mu_H/2)(C_{\mathrm{resp},\chi}+(1+N_{\mathrm{tr}})C'_{\chi})$. Collecting, with $e^{n_C\bar\varepsilon_t}\leqslant e^{2n_C\bar\varepsilon_t}$, $g_0\leqslant 1$, $g_0(1+3r)^3\leqslant 64r^{d_G}$ and $(1+3r)^2r\leqslant 16r^{d_G}$ for $r\geqslant 1$, $d_G\geqslant 3$, gives \eqref{eq:c5-5-9} with $C_2=\nu^{-1}C_{\mathrm{dist}}[4C_{\varpi} N_{\mathrm{lat}}(\mu_H/2)+N_{\mathrm{host}}e^{\nu R_{\mathrm{st}}/(C_FLM)}(64C_{\mathrm{bd}}+16C'_{\mathrm{bd}}C_m\vartheta N_{\mathrm{lat}}(\mu_H/2))]$.
\end{proof}

\begin{corollary}[the exterior sum on the far event]\label{cor:c5-5-10}
$J(\Lambda)= J\cup(I\setminus D(\Lambda))$ with $J= E_{\Omega,*}\setminus I\subseteq E(\Lambda_1^c)$ (a free bond not in $I$ has both endpoints off $\Lambda_1$, by the bond convention of Section~\ref{sec:c5-0}). On the far event $\mathsf N_b^c$ of \eqref{eq:c4-2-1} (Section~\ref{sec:c4-2}: $\operatorname{dist}(\Gamma_C,E(\Lambda_1^c))> b$, with $b=\lfloor LMR_1/16\rfloor$ as in Lemma~\ref{lem:c4-2-1}) every $e\in J$ has $\operatorname{dist}(e,C)\geqslant b-1$ in the convention of Section~\ref{sec:c5-5-0} (the distance for $N_b$ may be measured between bond sites rather than to the vertex set of $C$, which changes a distance by at most 1), and
\begin{equation*}
    \sum_{e\in J}e^{-\mu'\operatorname{dist}(e,C)}\leqslant \sum_{s\geqslant b-1}\#\{ e\, : \,\operatorname{dist}(e,C)\leqslant s\} e^{-\mu's}\leqslant 4n_C\sum_{s\geqslant b-1}(2s+1)^4e^{-\mu's}\leqslant n_CN_{\mathrm{lat}}(\mu'/2)e^{-\mu'(b-1)/2}
\end{equation*}
(at most $4(2s+1)^4$ bonds within $\ell^{\infty}$ distance $s$ of each of the $n_C$ vertices of $C$). Hence on $A^{\mathrm{fin}}_{\mathrm{cov}}\cap\mathsf N_b^c$,
\begin{gather}
    |\mathcal B^{\Lambda}_{\xi}[f_{\xi}]|\leqslant | S^{\circ}_{\Lambda}| e^{2n_C\bar\varepsilon_t}\bigl(C_1'g_0r+C_2\bar\varepsilon_tr^{d_G}\bigr)\bigl(n_CN_{\mathrm{lat}}(\mu'/2)e^{-\mu'(b-1)/2}+\mathfrak D_{\Lambda}\bigr)\;,\label{eq:c5-5-10}\\
    \mathfrak D_{\Lambda}\coloneqq \sum_{e\in I\setminus D(\Lambda)}e^{-\mu'\operatorname{dist}(e,C)}\leqslant | I\setminus D(\Lambda)|\;,\notag
\end{gather}
using $\mathfrak S_1(\Lambda)\leqslant e^{\mu_H(2L+1)/2}\mathfrak S_2(\Lambda)$ ($\operatorname{dist}(e,I_C)\geqslant \operatorname{dist}(e,C)-(2L+1)$, $\mu_H/2\geqslant \mu'$), with $C_1'\coloneqq C_1e^{\mu_H(2L+1)/2}$. The first summand is volume-independent and, since $b=\lfloor\ell_{\mathrm{mac}}/16\rfloor$ with $\ell_{\mathrm{mac}}= LMR_1$ growing as a fixed positive power of $\log g_1^{-2}$ (the scale integer $R_1$; $g_1\leqslant (1+\beta_0)g_0$ by CRE (2.6), as recalled in Section~\ref{sec:c5-3} before \eqref{eq:c5-3-1}, so $\log g_1^{-2}\geqslant \log g_0^{-2}-2\log(1+\beta_0)$), beats every polynomial in $r$: $r^{d_G}e^{-\mu'(b-1)/2}\leqslant e^{-\mu'(b-1)/4}$ as soon as $b-1\geqslant (4d_G/\mu')\log r$, which holds for $g_0$ small, since $b-1\geqslant \ell_{\mathrm{mac}}/16-2$ grows as a power of $\log g_0^{-2}$ while $\log r= p_0\log\log g_0^{-2}+\log A_1$. The second summand $\mathfrak D_{\Lambda}$ is the term's de-younged bonds weighted by their distance to the contour; it is a quantity of the labelled term, bounded by the number $| I\setminus D(\Lambda)|$ of de-younged bonds, and its average over the labels is the bookkeeping of Section~\ref{sec:c5-6} (Section~\ref{sec:c5-3-4}(a)--(b): a de-younged bond $z$ of $D^{(0)}\setminus D^{\mathrm{fin}}$ is not itself marked; it is attributed to a marked bond $e$ (a collar-test failure $M_e$) with $x_e$ within $102\ell+1$ of an endpoint of $z$, and Section~\ref{sec:c5-3-4}(a) pays the marked bonds attributed to $z$ by $\sum_ep^{\vartheta}_e\leqslant N_{\mathrm{ret}}q_{\vartheta}$, which is the term $N_{\mathrm{ret}}q_{\vartheta}$ of \eqref{eq:c5-3-coverage-average}; the round-0 bonds $D\setminus D^{(0)}$ carry no marks and are paid by $p_z^{(0)}\leqslant q$, Lemma~\ref{lem:c5-3-0}(b)).
\end{corollary}

\begin{remark}[scope of this subsection]\label{rem:c5-5-11}
The marked terms are the source-containing connected boundary terms: each is a closed contraction (Lemma~\ref{lem:c5-5-7}(b)) connected to $J(\Lambda)$ through its mark, a translated cutoff (bond), a potential with an exterior input or with the shift (host), or the shift of the source word, and each mark pays its distance to the contour by the footprint pin (Lemma~\ref{lem:c5-5-8}(b)) or by the decay of $m$ (Lemma~\ref{lem:c5-5-1}); the perimeter factor $S^{\circ}_{\Lambda}$ is extracted before absolute values. What this subsection does not contain: (1) the identification of $S^{\circ}_{\Lambda}$ with $\kappa_Ce^{E^{\circ}_{\Lambda}}$ at the retreated carrier, which is Lemma~\ref{lem:c5-7-1} (the identity with the regional square scalar at $D(\Lambda)$, through Proposition~\ref{prop:regional-square-extraction} with $D= D(\Lambda)$) and Corollary~\ref{cor:c5-7-2} (the form $\kappa_Ce^{E^{\circ}_{\Lambda}}$); the reference data at $D(\Lambda)$ differ from those at $I$ not only in the groups whose roots or clusters meet $I\setminus D(\Lambda)$ but also in $H^{\circ}$ (the principal restriction to $D(\Lambda)$), the eligible labels $\mathsf L_Z$, the host set $\mathbf D_0(\Lambda)$, the seed sets $E_0$ and their covariances $C_{E_0}$, so the activities of the groups near $I\setminus D(\Lambda)$ change as well, and no bound on the ratio $S^{\circ}_{\Lambda}/S^{\circ}_{\xi,C}$ is asserted; (2) the history average of $\mathfrak D_{\Lambda}$, Section~\ref{sec:c5-6}; (3) the complementary events: $(A^{\mathrm{fin}}_{\mathrm{cov}})^c$ is paid by Section~\ref{sec:c5-3-4}(a) in terms of $q+N_{\mathrm{ret}}q_{\vartheta}$, and $\mathsf N_b$ is the far-history exclusion \eqref{eq:c4-2-1} of Lemma~\ref{lem:c4-2-1}, whose probability is paid there; (4) the inputs listed in Section~\ref{sec:c5-5-0} (the pin sums of Lemma~\ref{lem:c4-6-2}, the dummy comparison of Lemma~\ref{lem:c4-4-1}, the closed-contraction factorisation of Lemma~\ref{lem:c4-6-1}, the reference-gas cluster expansion, the envelope and footprint bounds, and the far event) are used as stated, at their stated hypotheses, and are not re-verified here; a failure of any of them would propagate to \eqref{eq:c5-5-9}--\eqref{eq:c5-5-10}; $\bar\varepsilon_t=2K_{\chi} C_{\mathrm{pin}}\rho_t$ is the pin bound of Lemma~\ref{lem:c4-6-2} with $\rho_t$ for the root weights; (5) the labels enter \eqref{eq:c5-5-9}--\eqref{eq:c5-5-10} additively through $\mathfrak D_{\Lambda}$, $\mathfrak S_1(\Lambda)$ and $\mathfrak S_2(\Lambda)$ and multiplicatively through $| S^{\circ}_{\Lambda}|$ ($\Lambda=\Lambda_1^{\mathrm{fin}}$ is label-dependent); the multiplicative dependence is divided out in Section~\ref{sec:c5-6}, which bounds $|\mathcal B^{\Lambda}|/| S^{\circ}_{\Lambda}|$ and uses $| S^{\circ}_{\Lambda}|\leqslant 1$; (6) the square/torus geometry of Section~\ref{sec:c5-5-0} is a standing hypothesis of $C_{\mathrm{pin}}$, $A_F$, $A_E$ and the hull bounds. The factor $e^{2n_C\bar\varepsilon_t}$ in \eqref{eq:c5-5-9} is the cost of summing over the undifferenced groups of a family, $\prod_W(1+K_{\chi} e^{\theta^{\circ}(W)}\bar{\mathfrak b}_W)\leqslant e^{n_C\bar\varepsilon_t}$; at fixed $n_C$ it is $1+O(n_Cg_0r)$, but in the square geometry it is a perimeter rate $2L\bar\varepsilon_t$ per coarse segment ($2n_C\bar\varepsilon_t$ over the $4n$ segments, $n_C=4nL$) with $\bar\varepsilon_t= O(g_0r)$, first order, whereas $\kappa_C= e^{-\delta mP(C)}z_c^4$ has the rate $\delta m$ built from one-mark components of weight $\tau^2$ and the tolerated error $E^{\circ}$ has the rate $C\tau^2$ per segment, $\tau= C[g_0(1+r)^3+e^{-cr^2}]$. The bound \eqref{eq:c5-5-10} is small only while $2n_C\bar\varepsilon_t$ stays below the exponents of its two gains, $\mu'(b-1)/4$ for the far-event summand and, after the label average of $\mathfrak D_{\Lambda}$ (Proposition~\ref{prop:c5-6-4}), $c_q\vartheta^2\pi_*^2/2$; with $\bar\varepsilon_t= O(g_0r)$, $b-1\geqslant \ell/16-2$, $\ell\leqslant L^2M(2\log g_0^{-2})^{\mathfrak r}$ \eqref{eq:c5-3-1} and $\pi_*\geqslant A_0(\tfrac12\log g_0^{-2})^{p_0}$ \eqref{eq:c5-3-2}, this restricts $n_C$ to $O((\log g_0^{-2})^{\mathfrak r}/(g_0r))$ and $O((\log g_0^{-2})^{2p_0}/(g_0r))$ respectively; the statement is about the fixed loop as $g_0\to0$. Section~\ref{sec:c5-5R} brings the rate to second order, $e^{n_C\bar\varepsilon_{\mathrm{ph}}}$ with $\bar\varepsilon_{\mathrm{ph}}= O(g_0^2(1+r)^4)= O(\tau^2)$, the order accepted for $E^{\circ}$; a bound uniform in $n_C$ would need the undifferenced slots resummed against the reference exclusions, which is not done.
\end{remark}

\subsection{Second-order bounds for the undifferenced slots}
\label{sec:c5-5R}

Proposition~\ref{prop:c5-5-9} carries the family-sum factor $e^{n_C\bar\varepsilon_t}$ (written $e^{2n_C\bar\varepsilon_t}$ in \eqref{eq:c5-5-9}), produced by $\prod_W(1+K_{\chi} e^{\theta^{\circ}(W)}\bar{\mathfrak d}_W)\leqslant \exp(\sum_WK_{\chi} e^{\theta^{\circ}(W)}\bar{\mathfrak d}_W)$ with the first-order bound $\bar{\mathfrak d}_W\leqslant M_{\Lambda}(W)+| a^{\Lambda}_W|\leqslant 2M_{\Lambda}(W)$ of Lemma~\ref{lem:c5-5-7}(e) and Lemma~\ref{lem:c5-5-3}(d), whose root weight $\rho_tm(Z)$ has $\rho_t= u_t/u_*= O(g_0r)$. This subsection proves the second-order bound $\bar{\mathfrak d}_W\leqslant 2M^{\mathrm{ph}}_{\Lambda}(W)$, root weight $\rho_{\mathrm{ph}}m(Z)$ with $\rho_{\mathrm{ph}}= O(g_0^2(1+r)^4)$, for the undifferenced slots at $t\in\{0,1\}$ and for the reference activities, root by root inside the selected integrals, by three mechanisms: ($\alpha$) a root with at least two selected labels carries $(u_t/u_*)^{|\mathsf J_Z|}\leqslant \rho_t^2$ from the finite-difference bound as it stands; ($\beta$) a root with a single selected label whose split selection contains an ordinary host ($|\mathsf J_Z|=1$, $\mathcal D_Z\neq\varnothing$; a root with two or more labels is in ($\alpha$) whatever its hosts) carries the factor $2\omega_t/\omega_*$ of its own Mayer factors, through a complex parameter on the selected ordinary primitives, at every $t$ and without parity; ($\gamma$) a single-label root without ordinary host ($\mathcal D_Z=\varnothing$, $|\mathsf J_Z|=1$) has an even measure at $t=0$ and at the reference data, so the term of the source linear in its label integrates to zero and the remainder carries one more factor $u_t$; at $t=1$ its measure is the $t=0$ measure plus marked differences of its translated cutoffs, and each mark is paid by the large-field reserve (a mark at a bond outside $\mathsf P$ costs $e^{-\frac9{16}a_{\mathrm{bad}}}$ against the majorant with the reserve $\frac9{128}a_{\mathrm{bad}}$, a mark at a bad bond is paid by that bond's own factor, and the whole marked difference of the measure costs $2C_Me^{-\frac9{32}a_{\mathrm{bad}}}$, Lemma~\ref{lem:c5-5R-2}(b)), not by $\varpi_b$, so the marks of the undifferenced slots carry neither exterior decay nor label dependence, and the exponent of the family sum stays a constant times $n_C$. The differenced slots ($\mathfrak d^{\mathrm{diff}}_W$, \eqref{eq:c5-5-8}) and the marked part of Proposition~\ref{prop:c5-5-9} are unchanged. The results are Lemmas~\ref{lem:c5-5R-1}--\ref{lem:c5-5R-3} and Proposition~\ref{prop:c5-5R-4}, which restates \eqref{eq:c5-5-9}, \eqref{eq:c5-5-10}, \eqref{eq:c5-6-3}, \eqref{eq:c5-6-4} with $e^{n_C\bar\varepsilon_{\mathrm{ph}}}$ in place of $e^{2n_C\bar\varepsilon_t}$.

Thresholds. In addition to the thresholds of Appendix~\ref{app:thresholds}, \eqref{eq:c5-4-T}, \eqref{eq:c5-4-Tp}, \eqref{eq:c5-5-T} and $\omega_*\leqslant 1/2$ assume
\begin{equation}\label{eq:c5-5R-T}
    \begin{aligned}
        &\text{(i) }\tfrac9{128}a_{\mathrm{bad}}\geqslant \max\bigl(160(\kappa+\log(8M^4+8)+1),\,64\log(1/\eta_*)\bigr)\quad\text{(\eqref{eq:c5-4-Tp} for $a_{\mathrm{bad}}/8$)};\\
        &\text{(ii) }\rho_t+u_t/2+2\omega_t/\omega_*+2C_Me^{-\frac9{32}a_{\mathrm{bad}}}\leqslant 1;\\
        &\text{(iii) }2C_M\leqslant (1+2r)^3,\quad\log g_0^{-2}\geqslant 1,\quad(9c_n\alpha_G/64)A_1^2\log g_0^{-2}\geqslant 1,
    \end{aligned}
\end{equation}
where $C_M\coloneqq 8(4M+1)^4\geqslant \max_b| E(Y_b)|$ (a probe host is a box of $M$-cubes of side at most 4 cubes, \eqref{eq:c-probe-host}, so it has at most $(4M+1)^4$ sites, each an endpoint of eight bonds). Lemmas~\ref{lem:c5-5R-1}--\ref{lem:c5-5R-3} and the bounds of Proposition~\ref{prop:c5-5R-4} use (i)--(ii) only; (iii) is used only for the explicit form of the rate (the display $\rho_{\mathrm{ph}}\leqslant C_{\mathrm{ph}}g_0^2(1+2r)^4$ below and the rate bound of Proposition~\ref{prop:c5-5R-4}). All three are conditions on $g_0$ alone of the kind of \eqref{eq:c5-4-Tp} and \eqref{eq:c5-5-T}: $a_{\mathrm{bad}}= c_n\alpha_Gr^2/4\to\infty$, $r\to\infty$ and $g_0(1+r)^3\to0$ as $g_0\to0$, with thresholds depending on $L$, $M$, the group, the character, $\kappa$, $\omega_*$, $\alpha_G$, $\eta_*$, $A_1$ and $p_0$. Put
\begin{gather*}
    \rho_{\mathrm{ph}}\coloneqq \rho_t\bigl[\rho_t+u_t/2+2\omega_t/\omega_*+2C_Me^{-\frac9{32}a_{\mathrm{bad}}}\bigr]\;,\qquad\bar\varepsilon_{\mathrm{ph}}\coloneqq 2K_{\chi} C_{\mathrm{pin}}\rho_{\mathrm{ph}}\;,\\
    M^{\mathrm{ph}}_{\Lambda}(W)\coloneqq \frac1{k(W)!}\prod_{Z\in\mathcal A(W)}\rho_{\mathrm{ph}}m(Z)\prod_j|\Phi|^+_{\Lambda}(\mathfrak c_j)\;,
\end{gather*}
so that $\rho_{\mathrm{ph}}\leqslant \rho_t\leqslant \rho_{\mathrm{pin}}$ by \eqref{eq:c5-5R-T}(ii) and \eqref{eq:c5-5-T}(iv), $M^{\mathrm{ph}}_{\Lambda}(W)\leqslant M_{\Lambda}(W)$, and, with $\rho_t= C'_{\chi} g_0(1+2r)/u_*$, $u_t= C'_{\chi} g_0(1+2r)$, $\omega_t= C_0g_0(1+2r)^3$ and $a_{\mathrm{bad}}= c_n\alpha_Gr^2/4$,
\begin{gather*}
    \rho_{\mathrm{ph}}\leqslant C_{\mathrm{ph}}g_0^2(1+2r)^4\;,\qquad C_{\mathrm{ph}}\coloneqq (C'_{\chi}/u_*)(C'_{\chi}/u_*+C'_{\chi}/2+2C_0/\omega_*+1)\;,\\
    \text{hence}\qquad\bar\varepsilon_{\mathrm{ph}}\leqslant 32K_{\chi} C_{\mathrm{pin}}C_{\mathrm{ph}}g_0^2(1+r)^4
\end{gather*}
under \eqref{eq:c5-5R-T}(iii): the last bracket term needs $2C_Me^{-9c_n\alpha_Gr^2/128}\leqslant g_0(1+2r)^3$, which holds once $2C_M\leqslant (1+2r)^3$ and $e^{-9c_n\alpha_Gr^2/128}\leqslant g_0$, the latter because $r= A_1(\log g_0^{-2})^{p_0}$ with $p_0\geqslant 1$ and $\log g_0^{-2}\geqslant 1$ gives, with $c_{\mathrm{ph}}\coloneqq 9c_n\alpha_G/128$, $e^{-c_{\mathrm{ph}}r^2}\leqslant e^{-c_{\mathrm{ph}}A_1^2(\log g_0^{-2})^2}= g_0^{2c_{\mathrm{ph}}A_1^2\log g_0^{-2}}\leqslant g_0$ as soon as $2c_{\mathrm{ph}}A_1^2\log g_0^{-2}\geqslant 1$, which is the third condition of \eqref{eq:c5-5R-T}(iii). The first three bracket terms are at most their coefficients times $g_0(1+2r)^3$ without any condition. In $\tau= C[g_0(1+r)^3+e^{-cr^2}]$ (\eqref{eq:marginal-fixed-forest-bound}; its $C$, $c$ are the constants of that bound), $\tau^2\geqslant C^2g_0^2(1+r)^6\geqslant C^2g_0^2(1+r)^4$, so $\bar\varepsilon_{\mathrm{ph}}\leqslant (32K_{\chi} C_{\mathrm{pin}}C_{\mathrm{ph}}/C^2)\tau^2$: the family-sum rate becomes second order, of the order $\tau^2$ of the bound \eqref{eq:regional-square-error} on $E^{\circ}_{D,C_n,T}$.

\subsubsection{Higher derivatives and the two-term Taylor remainder}

\begin{lemma}\label{lem:c5-5R-1}
(a) Lemma~\ref{lem:c5-5-7z}(a) holds for every $k\leqslant 2m_*$ with $2C_{\mathrm{link}}$ in place of $C_{\mathrm{link}}$: $\|\partial_{e_1,b_1}\cdots\partial_{e_k,b_k}\pi(L_n)(a)\|_{\mathrm{op}}\leqslant (2C_{\mathrm{link}}g_0)^k$ for real $a$ in the enlarged box, $b_i\in\mathcal F_{r_n}$ (repetitions allowed), $h$-unit $e_i$. (b) Let $g$ be as in Lemma~\ref{lem:c5-5-7z}(b) (link factors, each traversal at most once, fixed insertions $N$), $a_0$ real with $| a_{0,b}|_{\sigma}\leqslant r/4$, $\mathsf J$ a finite set of bonds of $D(\Lambda)$, and $\varphi(B)\coloneqq g(B+a_0)$ for $B\in\mathfrak g^{\mathsf J}$ with $| B_b|_{\sigma}\leqslant 5r/4$ ($B$ extended by zero off $\mathsf J$). For $b\in\mathsf J$ define, on $C^2$ functions $\Psi$ of $B$ ($D_b$ makes sense on every function, $L_b$ needs one derivative and $R_b$ two, and all three are applied to the analytic $\varphi$ only), $(D_b\Psi)(B)\coloneqq \Psi(B)-\Psi(B|_{B_b\to0})$, $(L_b\Psi)(B)\coloneqq \langle B_b,\nabla_b\Psi(B|_{B_b\to0})\rangle= h(B_b)\partial_{B_b/h(B_b),b}\Psi(B|_{B_b\to0})$ (the right side read as $0$ when $B_b=0$), $(R_b\Psi)(B)\coloneqq \int_0^1(1-s)\partial^2_s[\Psi(B|_{B_b\to sB_b})]\,\dd s$; operators at distinct bonds commute, and $D_b= L_b+R_b$ on $C^2$ functions. Then $\Delta_{\mathsf J}g(B)=(\prod_{b\in\mathsf J}D_b\varphi)(B)$, and for every partition $\mathsf J=\mathsf J'\sqcup\{ b_1,\dots,b_p\}$ ($p\geqslant 0$)
\begin{gather*}
    \Bigl|\Bigl(\prod_{b\in\mathsf J'}D_b\prod_{i=1}^pR_{b_i}\varphi\Bigr)(B)\Bigr|\leqslant 2^{-p}u_t^{|\mathsf J'|+2p}\| N\|\;,\\
    \prod_{b\in\mathsf J}D_b\varphi-\prod_{b\in\mathsf J'}D_b\prod_{i=1}^pR_{b_i}\varphi=\sum_{i=1}^p\prod_{b\in\mathsf J'}D_b\Bigl(\prod_{i'< i}D_{b_{i'}}\Bigr)L_{b_i}\Bigl(\prod_{i'> i}R_{b_{i'}}\Bigr)\varphi\;,
\end{gather*}
where the $i$-th term of the sum is, at fixed values of the coordinates other than $B_{b_i}$, a linear function of $B_{b_i}$: $B_{b_i}\mapsto\langle B_{b_i},c_i(B|_{\mathsf J\setminus b_i})\rangle$. Both (a) and (b) hold at the reference data ($U=\id$, base point $a_0=0$, the link factors $\pi(L_n)$ at the flat phase) with the same constants, as Lemma~\ref{lem:c5-5-7z} does: the proofs use only the charts and the Cauchy estimates of Lemma~\ref{lem:c5-5-7z}, which hold there.
\end{lemma}

\begin{proof}
(a) The proof of Lemma~\ref{lem:c5-5-7z}(a) with $m_*$ replaced by $2m_*$: Cauchy's estimate and polarisation give $\| D^k\Pi_n(A)[h_1,\dots,h_k]\|\leqslant k^kC_{\Pi}(4/\vartheta_*)^k\prod\| h_i\|$ for every $k$, and $k^k\leqslant (2m_*)^k$ for $k\leqslant 2m_*$, so $\| D^k\Pi_n(A)[h_1,\dots,h_k]\|\leqslant (2m_*)^kC_{\Pi}(4/\vartheta_*)^k\prod\| h_i\|$ and $\|\partial_{e_1,b_1}\cdots\partial_{e_k,b_k}\pi(L_n)\|\leqslant (8m_*C_{\Pi} C_{\mathsf C}g_0/\vartheta_*)^k\leqslant (2C_{\mathrm{link}}g_0)^k$; for a free bond the Duhamel bound $(C_{\pi} g_0)^k\leqslant (2C_{\mathrm{link}}g_0)^k$ holds for every $k$. (b) $D_b= L_b+R_b$ is Taylor's formula with integral remainder for $s\mapsto\Psi(B|_{B_b\to sB_b})$ on $[0,1]$. Expanding $\prod_{b\in\mathsf J}D_b$ (value minus value with the $b$-th coordinate zeroed) gives $\sum_{K\subseteq\mathsf J}(-1)^{|\mathsf J\setminus K|}\varphi(\mathbf 1_KB)=\Delta_{\mathsf J}g(B)$. The telescoping identity $\prod_i(L_{b_i}+R_{b_i})-\prod_iR_{b_i}=\sum_i(\prod_{i'< i}(L+R)_{b_{i'}})L_{b_i}(\prod_{i'> i}R_{b_{i'}})$ with $(L+R)_b= D_b$, and the commuting factor $\prod_{\mathsf J'}D_b$, give the second display. $L_{b_i}\Psi$ is linear in $B_{b_i}$ with a coefficient not depending on $B_{b_i}$; the operators at bonds $b\neq b_i$ are linear and act on the other coordinates, so they preserve this form. The bound: $D_b\Psi(B)=\int_0^1\partial_{s_b}\Psi(B|_{B_b\to s_bB_b})\,\dd s_b$ with $\partial_{s_b}= h(B_b)\partial_{B_b/h(B_b),b}$, and $R_{b_i}$ contributes $\int_0^1(1-s_i)\,\dd s_i\cdot h(B_{b_i})^2\partial^2_{e_i,b_i}$, $e_i\coloneqq B_{b_i}/h(B_{b_i})$; iterating, $\prod_{\mathsf J'}D_b\prod_iR_{b_i}\varphi(B)$ is an integral over $[0,1]^{|\mathsf J'|+p}$ with weights $\prod_i(1-s_i)$ of $\prod_{\mathsf J'}h(B_b)\prod_ih(B_{b_i})^2\cdot(\partial^{\mathsf J'}\partial^2_{e_1,b_1}\cdots\partial^2_{e_p,b_p}g)(a_0+B^s)$, every evaluation point real and in the enlarged box ($| s_bB_b+a_{0,b}|_{\sigma}\leqslant 3r/2\leqslant 2r$). The mixed derivative of order $Q\coloneqq |\mathsf J'|+2p$ is bounded as in the proof of Lemma~\ref{lem:c5-5-7z}(b): product rule inside the normalised trace, at most $N_{\mathrm{tr}}^Q$ assignments of the $Q$ directions to link factors whose input set contains the direction's bond, fixed factors contributing $\| N\|$, undifferentiated links unitary, and each differentiated link carrying at most two directions per bond of its input set, hence at most $2|\mathcal F_{r_{n_i}}|\leqslant 2m_*-2$ directions, bounded by (a); so $\sup|\partial^{\dots}g|\leqslant (2N_{\mathrm{tr}}C_{\mathrm{link}}g_0)^Q\| N\|$. With $h(B_b)\leqslant | B_b|_{\sigma}\leqslant 5r/4$, $\int_0^1(1-s)\,\dd s=1/2$ and $(5r/4)\cdot2N_{\mathrm{tr}}C_{\mathrm{link}}g_0\leqslant 2N_{\mathrm{tr}}C_{\mathrm{link}}g_0(1+2r)= u_t$, the bound is $2^{-p}u_t^Q\| N\|$.
\end{proof}

\subsubsection{Evenness, marks paid by the large-field reserve, and the Cauchy parameter}
For a split selection $\ell=(\mathcal D,\mathsf J,\mathsf P,Z)$ for the carrier $D(\Lambda)$ (Lemma~\ref{lem:c5-5-2}(a)) and $x\in\{1/8,1\}$ put
\begin{equation*}
    \mathfrak M^{(x)}(\ell)\coloneqq e^{-(\kappa_1-1)N_1(Z\setminus Z_0')+N_1(Z)-x\frac9{16}a_{\mathrm{bad}}|\mathsf P|}\prod_{Y\in\mathcal D}(e^{A_Y}-1)\;,\qquad A_Y\coloneqq (\omega_*/2)e^{-\kappa d_0(Y)}\;,
\end{equation*}
the right side of \eqref{eq:c5-5-3} with the host factor $e^{A_Y}-1\leqslant \omega_*e^{1/2}e^{-\kappa d_0(Y)}$ (as in the proof of Lemma~\ref{lem:c5-5-2}(c)) and the reserve exponent multiplied by $x$; $\mathfrak M^{(1)}\leqslant \mathfrak M^{(1/8)}$. ($\mathfrak M^{(x)}(\ell)$ is a majorant of a single selection and is not the mark weight $\mathfrak M_{\Lambda}(Z)$ of Section~\ref{sec:c5-5}.)

\begin{lemma}\label{lem:c5-5R-2}
Let $\ell$ be a split selection for $D(\Lambda)$, $\Lambda^{f,t}_{\Lambda;\ell}$ its measure on $\mathfrak g^{E_0}$ ($t\in[0,1]$) and $\Lambda^{f,\circ}_{\Lambda;\ell}$ the reference measure (Lemma~\ref{lem:c5-5-3}(b)); the hypotheses are those of Lemma~\ref{lem:c5-5-2}.
\begin{enumerate}[label=(\alph*)]
    \item (Evenness.) If $\mathcal D=\varnothing$ then $\Lambda^{f,0}_{\Lambda;\ell}$ and $\Lambda^{f,\circ}_{\Lambda;\ell}$ are invariant under $B\mapsto-B$ on $\mathfrak g^{E_0}$: $\int F(-B)\,\dd\Lambda(B)=\int F(B)\,\dd\Lambda(B)$ for bounded Borel $F$. In particular $\int\Psi\,\dd\Lambda=0$ for every bounded Borel $\Psi$ with $\Psi(-B)=-\Psi(B)$, e.g. $\Psi(B)=\langle B_b,c\rangle\mathbf 1\{| B_b|_{\sigma}\leqslant 5r/4\}$ for a bond $b\in E(Y_0)$ and a fixed $c\in\mathfrak g$ (the truncation is even and equals 1 on the support of the measure, where $| B_b|_{\sigma}< r$ by the cutoff $\chi^0_b$; the untruncated linear function is not bounded).
    \item (Marked difference, large-field payment.) If $\mathcal D=\varnothing$ and $|\mathsf J|=1$, then under \eqref{eq:c5-5R-T}(i), for bounded Borel $F$, $|\int F\,\dd(\Lambda^{f,1}_{\Lambda;\ell}-\Lambda^{f,0}_{\Lambda;\ell})|\leqslant \sup| F|\cdot2C_Me^{-\frac9{32}a_{\mathrm{bad}}}\mathfrak M^{(1/8)}(\ell)$.
    \item (Cauchy parameter.) For $\zeta\in\mathbb C$ let $\Lambda^{f,t}_{\Lambda;\ell}(\zeta)$ be defined by the display of Lemma~\ref{lem:c5-4-1} for the enlarged data with each Mayer factor $\mathbf W_t(Y,B)e^{t_Y\mathbf W_t(Y,B)}$, $Y\in\mathcal D$, replaced by $\zeta\mathbf W_t(Y,B)e^{t_Y\zeta\mathbf W_t(Y,B)}$ (every near primitive assigned to a host of $\mathcal D$ multiplied by $\zeta$; density, cutoffs, probes unchanged), so that $\Lambda^{f,t}_{\Lambda;\ell}(1)=\Lambda^{f,t}_{\Lambda;\ell}$; put $R_{\Upsilon}\coloneqq \omega_*/(2\omega_t)\geqslant 1$ (\eqref{eq:c5-5-T}(ii)). For bounded Borel $F$ (complex-valued; the total variation norm below is the supremum over such $F$ with $| F|\leqslant 1$) the function $\zeta\mapsto\int F\,\dd\Lambda^{f,t}_{\Lambda;\ell}(\zeta)$ is entire; it vanishes at $\zeta=0$ if $\mathcal D\neq\varnothing$; and $|\int F\,\dd\Lambda^{f,t}_{\Lambda;\ell}(\zeta)|\leqslant \sup| F|\cdot\mathfrak M^{(1)}(\ell)$ for $|\zeta|\leqslant R_{\Upsilon}$. Consequently $\|\Lambda^{f,t}_{\Lambda;\ell}\|_{\mathrm{TV}}\leqslant \mathfrak M^{(1)}(\ell)$, and if $\mathcal D\neq\varnothing$ then $|\int F\,\dd\Lambda^{f,t}_{\Lambda;\ell}|\leqslant (2\omega_t/\omega_*)\cdot\sup_{|\zeta|\leqslant R_{\Upsilon}}|\int F\,\dd\Lambda^{f,t}_{\Lambda;\ell}(\zeta)|\leqslant (2\omega_t/\omega_*)\sup| F|\mathfrak M^{(1)}(\ell)$. The same holds at the reference data.
    \item (Weakened dummy comparison.) Under \eqref{eq:c5-5R-T}(i), for every connected final domain $Z$ for $D(\Lambda)$ and $x\in\{1/8,1\}$, $\sum_{\ell\, : \,Z(\ell)= Z}u_*^{|\mathsf J|}\mathfrak M^{(x)}(\ell)\leqslant m(Z)$, the sum over all split selections with final domain $Z$ and $Z_0\neq\varnothing$ ($\mathsf J=\varnothing$ included).
\end{enumerate}
\end{lemma}

\begin{proof}
(a) With $\mathcal D=\varnothing$ the display of Lemma~\ref{lem:c5-4-1} for the enlarged data has, inside the $s$- and $\sigma$-integrations (which do not involve $(B,X)$), the integrand $\prod_{b\in E(Y_0)}\chi^0_b(B)\prod_{b\in\mathsf P}(1-\chi^0_b(B))F(B)$ against the density \eqref{eq:fixed-kernel-density} $e^{-\frac12\langle\Gamma_{E_0}(\sigma)X,C_{E_0}\Gamma_{E_0}(\sigma)X\rangle}e^{-\langle B,\Gamma_{E_0}(\sigma)X\rangle}\,\dd\mu_{C_{E_0}}(B)\,\dd\mu_0(X)$. Substitute $(B,X)\to(-B,-X)$: $\mu_0$ is the standard Gaussian measure and $\mu_{C_{E_0}}$ the centred Gaussian measure of covariance $C_{E_0}$ (Section~\ref{sec:fixed-first-kernel}), both invariant; the quadratic form in $X$ and the bilinear form $\langle B,\Gamma_{E_0}(\sigma)X\rangle$ are invariant under the joint sign change; $\chi^0_b(-B)=\mathbf 1\{|-B_b|_{\sigma}< r\}=\chi^0_b(B)$ and likewise $1-\chi^0_b$ ($|\bigcdot |_{\sigma}$ is a norm); there is no Mayer factor. Hence $\int F(-B)\,\dd\Lambda^0=\int F\,\dd\Lambda^0$. At the reference data the display is the same with $H^{\circ}$, $m^{\circ}=0$ and the cutoffs $\chi_b(a)=\mathbf 1\{| a_b|_{\sigma}< r\}$ (Lemma~\ref{lem:c5-5-3}(b)). For odd $\Psi$, $\int\Psi\,\dd\Lambda=\int\Psi(-B)\,\dd\Lambda=-\int\Psi\,\dd\Lambda$. (b) With $\mathcal D=\varnothing$ and $\mathsf J=\{ b_0\}$, the $t$-dependent factors of the display are $\chi^t_b$, $b\in E(Y_0)= E(Y_{b_0})$ (at most $C_M$ bonds), and $1-\chi^t_b$, $b\in\mathsf P$ (proof of Lemma~\ref{lem:c5-5-6}(a)); order them and telescope, $\prod a_k-\prod c_k=\sum_k(\prod_{i< k}a_i)(a_k-c_k)(\prod_{i> k}c_i)$, $a$ the $t=1$ and $c$ the $t=0$ factors: at most $C_M+|\mathsf P|$ terms, each with exactly one differenced factor (the mark at a bond $b$), of modulus $|\chi^1_b-\chi^0_b|\leqslant \mathbf 1_{S_b}(B_b)\leqslant \mathbf 1\{| B_b|_{\sigma}\geqslant r-| m_b|_{\sigma}\}\leqslant \mathbf 1\{| B_b|_{\sigma}\geqslant 3r/4\}\leqslant e^{-\frac9{16}a_{\mathrm{bad}}}e^{(\alpha_G/2)h(B_b)^2}$ (\eqref{eq:c5-4-1} and the proof of Lemma~\ref{lem:c5-4-2}). Bound the modulus of each term factor by factor as in the proof of Lemma~\ref{lem:c5-4-2}: undifferenced cutoffs $\leqslant 1$; undifferenced bad factors $1-\chi^t_{b'}\leqslant e^{-\frac9{16}a_{\mathrm{bad}}}e^{(\alpha_G/2)h(B_{b'})^2}$; the mark as displayed; $| F|\leqslant \sup| F|$; the product of the exponential factors is $e^{-\frac9{16}a_{\mathrm{bad}}n}e^{(\alpha_G/2)\| B\|^2}$ with $n\coloneqq |\mathsf P|+\mathbf 1\{ b\notin\mathsf P\}\geqslant \max(|\mathsf P|,1)$, and the density's modulus against $e^{(\alpha_G/2)\| B\|^2}$ integrates to at most $e^{N_1(Z)}$ (Lemma~\ref{lem:gaussian-payment}(ii) for $H$), the weakening giving $e^{-(\kappa_1-1)N_1(Z\setminus Z_0')}$. Since $\frac9{16}a_{\mathrm{bad}}n\geqslant \frac9{128}a_{\mathrm{bad}}|\mathsf P|+\frac{63}{128}a_{\mathrm{bad}}\cdot\max(|\mathsf P|,1)$, each term is at most $\sup| F|\cdot e^{-\frac{63}{128}a_{\mathrm{bad}}\max(|\mathsf P|,1)}\mathfrak M^{(1/8)}(\ell)$, and the number of terms is at most $C_M+|\mathsf P|\leqslant C_M(1+|\mathsf P|)$. For $|\mathsf P|=0$ this gives $C_Me^{-\frac{63}{128}a_{\mathrm{bad}}}\leqslant C_Me^{-\frac9{32}a_{\mathrm{bad}}}$; for $|\mathsf P|= k\geqslant 1$, $C_M(1+k)e^{-\frac{63}{128}a_{\mathrm{bad}}k}= C_Me^{-\frac9{32}a_{\mathrm{bad}}}\cdot(1+k)e^{-\frac{63k-36}{128}a_{\mathrm{bad}}}\leqslant C_Me^{-\frac9{32}a_{\mathrm{bad}}}\cdot(1+k)e^{-\frac{27}{128}a_{\mathrm{bad}}k}\leqslant 2C_Me^{-\frac9{32}a_{\mathrm{bad}}}$, because $63k-36\geqslant 27k$ for $k\geqslant 1$ and $(1+k)e^{-ck}\leqslant 2e^{-c}\leqslant 2$ for $c\coloneqq \frac{27}{128}a_{\mathrm{bad}}\geqslant 1$ (implied by \eqref{eq:c5-5R-T}(i)). (c) For fixed $(X,B,s,\sigma,(t_Y))$ the integrand is $\prod_{Y\in\mathcal D}\zeta\mathbf W_te^{t_Y\zeta\mathbf W_t}$ times a $\zeta$-independent factor, entire in $\zeta$; $\mathbf W_t(Y,B)=\mathbf V^{\Lambda}_{\xi,t}(Y,B+tm)$ (Section~\ref{sec:c5-4-2}) vanishes off the translated local box of $Y$ and has $|\mathbf W_t(Y,B)|\leqslant \omega_te^{-\kappa d_0(Y)}\leqslant \omega_t$ on it (\eqref{eq:c5-4-4}), so on $|\zeta|\leqslant R$ the integrand's modulus is at most $(R\omega_te^{R\omega_t})^{|\mathcal D|}\sup| F|$ times the density's modulus, which is integrable (Lemma~\ref{lem:gaussian-payment}(ii)); the $s$-, $\sigma$-, $t_Y$-integrals are over compact sets. Dominated convergence gives continuity in $\zeta$ and, with Morera's theorem (or differentiation under the integral sign), analyticity on every disc: the function is entire. At $\zeta=0$ every Mayer factor vanishes and $\mathcal D\neq\varnothing$, so the integral is zero. For $|\zeta|\leqslant R_{\Upsilon}$, integrate $t_Y$ first: $\int_0^1\zeta\mathbf W_te^{t_Y\zeta\mathbf W_t}\,\dd t_Y= e^{\zeta\mathbf W_t}-1$, of modulus $\leqslant e^{|\zeta\mathbf W_t|}-1\leqslant e^{R_{\Upsilon}\omega_te^{-\kappa d_0(Y)}}-1= e^{A_Y}-1$ on the support of $\Theta^t_{E(Y)}$; the other factors are bounded as in the proofs of Lemma~\ref{lem:c5-4-2} and Lemma~\ref{lem:c5-5-2}(b) (cutoffs $\leqslant 1$, bad bonds by $e^{-\frac9{16}a_{\mathrm{bad}}}e^{(\alpha_G/2)h^2}$, Gaussian payment $e^{N_1(Z)}$, weakening), giving $\sup| F|\mathfrak M^{(1)}(\ell)$. The TV bound is the case $\zeta=1$ (for $\mathcal D=\varnothing$ there is no $\zeta$, and the bound is the factor-by-factor bound directly). Schwarz's lemma for $\Phi(\zeta)\coloneqq \int F\,\dd\Lambda^{f,t}_{\Lambda;\ell}(\zeta)$ on $|\zeta|\leqslant R_{\Upsilon}$ with $\Phi(0)=0$ and $|\Phi|\leqslant M_F\coloneqq \sup| F|\mathfrak M^{(1)}(\ell)$ gives $|\Phi(\zeta)|\leqslant M_F|\zeta|/R_{\Upsilon}$, and $\zeta=1$ lies in the disc since $R_{\Upsilon}\geqslant 1$. At the reference data \eqref{eq:c5-4-4} holds by Lemma~\ref{lem:c5-4-0}(d) for these data (Lemma~\ref{lem:c5-5-3}(b)). (d) The proof of Lemma~\ref{lem:c5-5-2}(c): the hostwise comparison bounds the sum over the nonempty label subsets of a host of the weights $e^{A_Y}-1$ (ordinary) and $u_*$ (probe) by $\omega_*e^{1/2}e^{-\kappa d_0(Y)}$, so $\sum_{\ell} u_*^{|\mathsf J|}\mathfrak M^{(x)}(\ell)$ is at most the sum over the ordinary selections $\mathfrak s=(\mathcal D,\mathsf P,Z)$ of the majorant \eqref{eq:c5-4-5} with $x\frac9{16}a_{\mathrm{bad}}$ in place of $\frac9{16}a_{\mathrm{bad}}$; the seed, enlargement and final-domain sums of Appendix~\ref{app:background-estimates} bound this by $m(Z)$ exactly as in Corollary~\ref{cor:c5-4-3}, whose proof applies to every number $\geqslant 1$ in the role of $a_{\mathrm{bad}}$ that satisfies \eqref{eq:native-source-bad-budget}, which $\frac9{128}a_{\mathrm{bad}}$ does by \eqref{eq:c5-5R-T}(i) (and $\frac9{16}a_{\mathrm{bad}}$ by \eqref{eq:c5-4-Tp}).
\end{proof}

\subsubsection{The second-order group bound}

\begin{lemma}\label{lem:c5-5R-3}
Under the hypotheses of Lemma~\ref{lem:c5-5-2}, \eqref{eq:c5-5R-T} and the setting of Section~\ref{sec:c5-5-0} (on $A^{\mathrm{fin}}_{\mathrm{cov}}$), for every group $W$ (Lemma~\ref{lem:c5-5-3}(a)), every $h\in G$ and $t\in\{0,1\}$: $| c^t_W(h)|\leqslant M^{\mathrm{ph}}_{\Lambda}(W)$ and $| a^{\Lambda}_W|\leqslant M^{\mathrm{ph}}_{\Lambda}(W)$. Hence $\bar{\mathfrak d}_W\leqslant 2M^{\mathrm{ph}}_{\Lambda}(W)\eqqcolon \bar{\mathfrak b}^{\mathrm{ph}}_W$ (the bound $\bar{\mathfrak d}_W\leqslant M_{\Lambda}(W)+| a^{\Lambda}_W|$ of Lemma~\ref{lem:c5-5-7}(e) with both terms improved), and Lemma~\ref{lem:c5-5-8} holds with $(M^{\mathrm{ph}}_{\Lambda},\rho_{\mathrm{ph}},\bar\varepsilon_{\mathrm{ph}})$ in place of $(M_{\Lambda},\rho_t,\bar\varepsilon_t)$:
\begin{gather*}
    \sup_{\Delta}\sum_{W\, : \,\Delta\in F(W)}e^{2\nu\sigma(W)}K_{\chi} e^{\theta^{\circ}(W)}\bar{\mathfrak b}^{\mathrm{ph}}_W\leqslant \bar\varepsilon_{\mathrm{ph}}\;,\qquad\sup_n\sum_{W\, : \,n\in\mathcal H(W)}e^{2\nu\sigma(W)}K_{\chi} e^{\theta^{\circ}(W)}\bar{\mathfrak b}^{\mathrm{ph}}_W\leqslant \bar\varepsilon_{\mathrm{ph}}\;,\\
    \sum_WK_{\chi} e^{\theta^{\circ}(W)}\bar{\mathfrak b}^{\mathrm{ph}}_W\leqslant n_C\bar\varepsilon_{\mathrm{ph}}\;,
\end{gather*}
and (b) of Lemma~\ref{lem:c5-5-8} likewise with $\bar\varepsilon_{\mathrm{ph}}$. The differenced bound \eqref{eq:c5-5-8}, its majorant $M_{\Lambda}(W)$ and $\theta^{\circ}(W)\leqslant \nu\sigma(W)/16$ are unchanged.
\end{lemma}

\begin{proof}
Fix $W$, $h$ and $t\in\{0,1\}$. By Lemma~\ref{lem:c5-5-2}(a),(d) and Lemma~\ref{lem:c5-5-3}(a), $c^t_W(h)= w^t_{\Lambda}(W)\sum_{(\ell_Z)_{Z\in\mathcal A(W)}\, : \,\mathsf J_Z\neq\varnothing}\int\Delta_{\mathsf J}g_t\,\dd\bigotimes_Z\Lambda^{f,t}_{\Lambda;\ell_Z}$, with $g_t(a)=\dim(\pi)^{-1}\operatorname{tr}(\pi(T_W(a))N_t)$, $N_t\coloneqq \pi(P^t_W)^{-1}\pi(h)$ fixed unitary, base point $a_0= tm$, $\mathsf J=\bigsqcup_Z\mathsf J_Z$, and $| w^t_{\Lambda}(W)|\leqslant (1/k(W)!)\prod_j|\Phi|^+_{\Lambda}(\mathfrak c_j)$ (Corollary~\ref{cor:c5-4-3} and Lemma~\ref{lem:c5-5-6}(b), as in the proof of Lemma~\ref{lem:c5-5-3}(d)). The integrand $\Delta_{\mathsf J}g_t$ depends on $B$ through the labels $B_{\mathsf J}$ only, and the labels of the root $Z$ are its $\mathsf J_Z\subseteq E_0(\ell_Z)$. Fix a tuple $(\ell_Z)$ and sort its roots into three classes: ($\alpha$) $|\mathsf J_Z|\geqslant 2$; ($\beta$) $|\mathsf J_Z|=1$ and $\mathcal D_Z\neq\varnothing$; ($\gamma$) $|\mathsf J_Z|=1$ and $\mathcal D_Z=\varnothing$, with single label $b_Z$.

Step 1 (splitting the class-($\gamma$) measures; $t=1$ only). For each class-($\gamma$) root write $\Lambda^{f,1}_{\Lambda;\ell_Z}=\Lambda^{f,0}_{\Lambda;\ell_Z}+\delta\Lambda_{\ell_Z}$ and expand the product measure: $\bigotimes_Z\Lambda^{f,1}_{\ell_Z}=\sum_{\mathcal S\subseteq(\gamma)}[\bigotimes_{Z\in\mathcal S}\delta\Lambda_{\ell_Z}]\otimes[\bigotimes_{Z\in(\gamma)\setminus\mathcal S}\Lambda^{f,0}_{\ell_Z}]\otimes[\bigotimes_{Z\notin(\gamma)}\Lambda^{f,1}_{\ell_Z}]$. At $t=0$ only $\mathcal S=\varnothing$ occurs. Fix $\mathcal S$; the roots of $(\gamma)\setminus\mathcal S$ are the parity roots, $Z_1,\dots,Z_p$ with labels $b_1,\dots,b_p$, and $\mathsf J'\coloneqq \mathsf J\setminus\{ b_1,\dots,b_p\}$.

Step 2 (parity). By Lemma~\ref{lem:c5-5R-1}(b) with $\varphi_t\coloneqq g_t(\bigcdot +tm)$ ($| tm_b|_{\sigma}\leqslant r/4$ by \eqref{eq:c5-4-1}; $| B_b|_{\sigma}\leqslant 5r/4$ on the supports of the measures, Lemma~\ref{lem:c5-5-2}(a); every function of the labels below, $\varphi_t$, the $L$- and $R$-terms, $g^{(p)}$, is multiplied by the even truncation $\prod_{b\in\mathsf J}\mathbf 1\{| B_b|_{\sigma}\leqslant 5r/4\}$, which equals 1 on the supports of the measures and makes it bounded Borel on $\mathfrak g^{E_0}$, so that Lemma~\ref{lem:c5-5R-2} applies), $\Delta_{\mathsf J}g_t= g^{(p)}+\sum_{i=1}^p(\text{term linear in }B_{b_i})$, $g^{(p)}\coloneqq \prod_{\mathsf J'}D_b\prod_iR_{b_i}\varphi_t$. The $i$-th linear term depends on the variables of the root $Z_i$ through $B_{b_i}$ alone ($\mathsf J_{Z_i}=\{ b_i\}$) and is odd under $B_{E_0(Z_i)}\to-B_{E_0(Z_i)}$ at fixed values of the other roots' variables; the measures are finite and the integrand is bounded Borel, so by Fubini and Lemma~\ref{lem:c5-5R-2}(a) ($\Lambda^{f,0}_{\ell_{Z_i}}$ is even, $\mathcal D_{Z_i}=\varnothing$) its integral against the product measure of Step 1 vanishes. Hence $\int\Delta_{\mathsf J}g_t\,\dd(\text{product for }\mathcal S)=\int g^{(p)}\,\dd(\text{product for }\mathcal S)$, and $\sup| g^{(p)}|\leqslant 2^{-p}u_t^{|\mathsf J'|+2p}$ ($\| N_t\|=1$).

Step 3 (Cauchy parameter). Let
\begin{equation*}
    \Phi((\zeta_Z)_{Z\in(\beta)})\coloneqq \int g^{(p)}\,\dd\Bigl[\bigotimes_{Z\in(\beta)}\Lambda^{f,t}_{\Lambda;\ell_Z}(\zeta_Z)\otimes(\text{the other factors of the product for }\mathcal S)\Bigr]\;.
\end{equation*}
$\Phi$ is analytic in each $\zeta_Z$ separately on the open polydisc $\{|\zeta_Z|< R_{\Upsilon}\}$: at fixed values of the other roots' variables, the integral over the variables of the root $Z$ is, for every bounded Borel integrand, entire in $\zeta_Z$ with the uniform bound $\sup|\text{integrand}|\cdot\mathfrak M^{(1)}(\ell_Z)$ on $|\zeta_Z|\leqslant R_{\Upsilon}$ (Lemma~\ref{lem:c5-5R-2}(c), a single-root statement), and this uniform bound carries analyticity through the integrals over the other roots (finite measures, Fubini; dominated convergence and Morera's theorem for the outer integral). Analyticity on the open polydisc with a bound on the closed one is all Schwarz's lemma below needs. $\Phi$ vanishes whenever some $\zeta_Z=0$, and on the polydisc $\{|\zeta_Z|\leqslant R_{\Upsilon}\ \forall Z\in(\beta)\}$ satisfies $|\Phi|\leqslant M_{\Phi}\coloneqq \sup| g^{(p)}|\cdot\prod_{Z\in(\beta)}\mathfrak M^{(1)}(\ell_Z)\cdot\prod_{Z\in\mathcal S}\|\delta\Lambda_{\ell_Z}\|_{\mathrm{TV}}\cdot\prod_{Z\in(\alpha)\cup((\gamma)\setminus\mathcal S)}\|\Lambda_{\ell_Z}\|_{\mathrm{TV}}$ (the TV norms of the $t=0$ measures of the parity roots, of the $t=1$ measures of class ($\alpha$)). Schwarz's lemma one variable at a time, from the inside out: write $(\beta)=\{ Z_1,\dots,Z_q\}$; for each fixed $(\zeta_{Z_1},\dots,\zeta_{Z_{q-1}})$ in the polydisc, $\zeta_{Z_q}\mapsto\Phi(\zeta_{Z_1},\dots,\zeta_{Z_{q-1}},\zeta_{Z_q})$ vanishes at 0 and is bounded by $M_{\Phi}$ on $|\zeta_{Z_q}|\leqslant R_{\Upsilon}$, so $|\Phi(\zeta_{Z_1},\dots,\zeta_{Z_{q-1}},1)|\leqslant M_{\Phi}/R_{\Upsilon}$ on the polydisc of the remaining $q-1$ variables; the function $(\zeta_{Z_1},\dots,\zeta_{Z_{q-1}})\mapsto\Phi(\dots,1)$ is again analytic in each variable on the open polydisc, vanishes when any variable is 0 and is bounded by $M_{\Phi}/R_{\Upsilon}$, so the same step in $\zeta_{Z_{q-1}}$ gives $M_{\Phi}/R_{\Upsilon}^2$ at $\zeta_{Z_{q-1}}=1$, and after $q$ steps $|\Phi(1,\dots,1)|\leqslant R_{\Upsilon}^{-q}M_{\Phi}=(2\omega_t/\omega_*)^{|(\beta)|}M_{\Phi}$.

Step 4 (collecting). With $\|\Lambda_{\ell_Z}\|_{\mathrm{TV}}\leqslant \mathfrak M^{(1)}(\ell_Z)\leqslant \mathfrak M^{(1/8)}(\ell_Z)$ (Lemma~\ref{lem:c5-5R-2}(c)) and $\|\delta\Lambda_{\ell_Z}\|_{\mathrm{TV}}\leqslant 2C_Me^{-\frac9{32}a_{\mathrm{bad}}}\mathfrak M^{(1/8)}(\ell_Z)$ (Lemma~\ref{lem:c5-5R-2}(b)), the summand for $((\ell_Z),\mathcal S)$ is at most
\begin{multline*}
    \prod_{Z\in(\alpha)}u_t^{|\mathsf J_Z|}\mathfrak M^{(1/8)}(\ell_Z)\cdot\prod_{Z\in(\beta)}u_t(2\omega_t/\omega_*)\mathfrak M^{(1/8)}(\ell_Z)\\
    \times\prod_{Z\in(\gamma)\setminus\mathcal S}(u_t^2/2)\mathfrak M^{(1/8)}(\ell_Z)\cdot\prod_{Z\in\mathcal S}u_t\cdot2C_Me^{-\frac9{32}a_{\mathrm{bad}}}\mathfrak M^{(1/8)}(\ell_Z)\;.
\end{multline*}
Summing over $\mathcal S\subseteq(\gamma)$, a root of each class contributes
\begin{align*}
    (\gamma):&\quad[u_t^2/2+2C_Mu_te^{-\frac9{32}a_{\mathrm{bad}}}]\mathfrak M^{(1/8)}(\ell_Z)=\rho_t[u_t/2+2C_Me^{-\frac9{32}a_{\mathrm{bad}}}]u_*^{|\mathsf J_Z|}\mathfrak M^{(1/8)}(\ell_Z)\leqslant \rho_{\mathrm{ph}}u_*^{|\mathsf J_Z|}\mathfrak M^{(1/8)}(\ell_Z)\;,\\
    (\beta):&\quad\rho_t(2\omega_t/\omega_*)u_*^{|\mathsf J_Z|}\mathfrak M^{(1/8)}(\ell_Z)\leqslant \rho_{\mathrm{ph}}u_*^{|\mathsf J_Z|}\mathfrak M^{(1/8)}(\ell_Z)\;,\\
    (\alpha):&\quad\rho_t^{|\mathsf J_Z|}u_*^{|\mathsf J_Z|}\mathfrak M^{(1/8)}(\ell_Z)\leqslant \rho_t^2u_*^{|\mathsf J_Z|}\mathfrak M^{(1/8)}(\ell_Z)\leqslant \rho_{\mathrm{ph}}u_*^{|\mathsf J_Z|}\mathfrak M^{(1/8)}(\ell_Z)
\end{align*}
($\rho_t\leqslant 1$ by \eqref{eq:c5-5-T}(i)). Summing over the tuples,
\begin{equation*}
    \sum_{(\ell_Z)}\prod_Z\rho_{\mathrm{ph}}u_*^{|\mathsf J_Z|}\mathfrak M^{(1/8)}(\ell_Z)=\prod_Z\Bigl[\rho_{\mathrm{ph}}\sum_{\ell_Z\, : \,Z(\ell_Z)= Z}u_*^{|\mathsf J_Z|}\mathfrak M^{(1/8)}(\ell_Z)\Bigr]\leqslant \prod_Z\rho_{\mathrm{ph}}m(Z)
\end{equation*}
by Lemma~\ref{lem:c5-5R-2}(d) (the sum over all split selections dominates the sum over those with $\mathsf J_Z\neq\varnothing$). Multiplying by $| w^t_{\Lambda}(W)|$ gives $| c^t_W(h)|\leqslant M^{\mathrm{ph}}_{\Lambda}(W)$.

The reference activity $a^{\Lambda}_W= w^{\circ}_{\Lambda}(W)\int g^{\circ}_W\,\dd\bigotimes_Z\hat\nu^{f,\circ}_{\Lambda,Z}$ (Lemma~\ref{lem:c5-5-3}(b),(d)) is treated by the same four steps at the reference data, where Lemma~\ref{lem:c5-5R-2} holds (Lemma~\ref{lem:c5-5-3}(b)): base point 0, increments $| B_b|_{\sigma}\leqslant r\leqslant 5r/4$, $\mathcal S=\varnothing$ throughout (the reference cutoffs are even, so every class-($\gamma$) root is a parity root), $g^{\circ}_W$ a closed contraction with $\| N\|=1$, and $| w^{\circ}_{\Lambda}(W)|\leqslant (1/k(W)!)\prod_j|\Phi|^+_{\Lambda}(\mathfrak c_j)$.

Lemma~\ref{lem:c5-5-7}(e) improved: $\bar{\mathfrak d}_W=\max_t\sup_h| c^t_W(h)-a^{\Lambda}_W\chi(h)|\leqslant M^{\mathrm{ph}}_{\Lambda}(W)+| a^{\Lambda}_W|\leqslant 2M^{\mathrm{ph}}_{\Lambda}(W)$ by the two bounds just proved. The pins: the first display of Lemma~\ref{lem:c5-5-8}(a) is the positive summation with a root weight $\rho\,m(Z)$, and the inherited summation bound (Lemma~\ref{lem:c4-6-2}) holds in the form ``sum $\leqslant C_{\mathrm{pin}}\rho$ for every $\rho\leqslant \rho_{\mathrm{pin}}$'' (its smallness conditions on $\rho$ are monotone, so every $\rho\leqslant \rho_{\mathrm{pin}}$ satisfies them, and the bound is linear in $\rho$ at each such $\rho$; monotonicity alone would only give $C_{\mathrm{pin}}\rho_t$); applied at $\rho=\rho_{\mathrm{ph}}\leqslant \rho_{\mathrm{pin}}$ it gives $C_{\mathrm{pin}}\rho_{\mathrm{ph}}$. Alternatively, $M^{\mathrm{ph}}_{\Lambda}(W)=(\rho_{\mathrm{ph}}/\rho_t)^{|\mathcal A(W)|}M_{\Lambda}(W)$ with $|\mathcal A(W)|\geqslant 1$ (Definition~\ref{def:groups}, nonempty root set), so the first-order pin at $\rho_t$ gives $C_{\mathrm{pin}}\rho_t(\rho_{\mathrm{ph}}/\rho_t)= C_{\mathrm{pin}}\rho_{\mathrm{ph}}$ as well; this argument is independent of the first. The second display uses $\bar{\mathfrak b}_W\leqslant 2M_{\Lambda}(W)$ and $e^{\theta^{\circ}(W)}\leqslant e^{2\nu\sigma(W)}$, now with $\bar{\mathfrak b}^{\mathrm{ph}}_W=2M^{\mathrm{ph}}_{\Lambda}(W)$ and the unchanged $\theta^{\circ}(W)\leqslant \nu\sigma(W)/16$ (Lemma~\ref{lem:c5-5-3}(d), which uses the first-order $M_{\Lambda}$, $| a^{\Lambda}_W|\leqslant M_{\Lambda}(W)$ and \eqref{eq:c5-5-T}(iv), all still in force); Lemma~\ref{lem:c5-5-8}(b) is the geometric fact applied to the new pin.
\end{proof}

\subsubsection{The repaired bounds}

\begin{proposition}\label{prop:c5-5R-4}
Under the hypotheses of Proposition~\ref{prop:c5-5-9} and \eqref{eq:c5-5R-T}, on every labelled term with final region $\Lambda$, for exterior data in the support of the term's collar indicators and exterior cutoffs, on $A^{\mathrm{fin}}_{\mathrm{cov}}$,
\begin{equation}\label{eq:c5-5-9R}
    |\mathcal B^{\Lambda}_{\xi}[f_{\xi}]|\leqslant | S^{\circ}_{\Lambda}| e^{n_C\bar\varepsilon_{\mathrm{ph}}}\bigl[C_1g_0r\,\mathfrak S_1(\Lambda)+C_2\bar\varepsilon_tr^{d_G}\mathfrak S_2(\Lambda)\bigr]\;,
\end{equation}
with $C_1$, $C_2$, $\mu'$, $\mathfrak S_1$, $\mathfrak S_2$ of Proposition~\ref{prop:c5-5-9}; on $A^{\mathrm{fin}}_{\mathrm{cov}}\cap N_b^c$, \eqref{eq:c5-5-10} holds with $e^{n_C\bar\varepsilon_{\mathrm{ph}}}$ in place of $e^{2n_C\bar\varepsilon_t}$; and, under the hypotheses of Proposition~\ref{prop:c5-6-4} (compatible even fine periods exceeding $200\ell$, $g_0$ below the thresholds of Lemma~\ref{lem:c5-2-4}, of Proposition~\ref{prop:c5-2-6} and of \eqref{eq:c5-2-3}, and for the closing form its further threshold), \eqref{eq:c5-6-3}, \eqref{eq:c5-6-4} and the closing form hold with $e^{n_C\bar\varepsilon_{\mathrm{ph}}}$ in place of $e^{2n_C\bar\varepsilon_t}$. In the square geometry of Section~\ref{sec:c5-5-0} ($n_C=4nL$, $4n$ coarse segments) the family-sum factor has, under \eqref{eq:c5-5R-T}(iii), the perimeter rate $L\bar\varepsilon_{\mathrm{ph}}\leqslant 32LK_{\chi} C_{\mathrm{pin}}C_{\mathrm{ph}}g_0^2(1+r)^4$ per coarse segment, second order in $g_0(1+r)^2$, of the order $\tau^2$ of the bound \eqref{eq:regional-square-error} on $E^{\circ}_{D,C_n,T}$; the first-order constant $\bar\varepsilon_t$ survives only as the prefactor $C_2\bar\varepsilon_t$ of $\mathfrak S_2(\Lambda)$ (through Corollary~\ref{cor:c5-5-10} it multiplies $n_CN_{\mathrm{lat}}(\mu'/2)e^{-\mu'(b-1)/2}+\mathfrak D_{\Lambda}$) and no longer appears in an exponent.
\end{proposition}

\begin{proof}
The proof of Proposition~\ref{prop:c5-5-9} with one change. In \eqref{eq:c5-5-7} the undifferenced slots carry $\bar{\mathfrak d}_W$, bounded by Lemma~\ref{lem:c5-5R-3} instead of Lemma~\ref{lem:c5-5-7}(e); so the two family sums $\prod_W(1+K_{\chi} e^{\theta^{\circ}(W)}\bar{\mathfrak d}_W)$, in the background-word part, and over $\mathfrak B\setminus\{ W_0\}$ in the marked part, are at most $\exp(\sum_WK_{\chi} e^{\theta^{\circ}(W)}\bar{\mathfrak b}^{\mathrm{ph}}_W)\leqslant e^{n_C\bar\varepsilon_{\mathrm{ph}}}$ (the contour pin of Lemma~\ref{lem:c5-5R-3} summed over the $n_C$ traversals, every group having a contact with the contour). The differenced slot $\mathfrak d^{\mathrm{diff}}_{W_0}$ keeps \eqref{eq:c5-5-8} with the first-order $M_{\Lambda}(W_0)$ and the pins of Lemma~\ref{lem:c5-5-8} with $\bar\varepsilon_t$, so (i) and (ii) of that proof and the constants $C_1$, $C_2$ are unchanged, and the closing step $e^{n_C\bar\varepsilon_t}\leqslant e^{2n_C\bar\varepsilon_t}$ is dropped. Corollary~\ref{cor:c5-5-10} and Proposition~\ref{prop:c5-6-4} use \eqref{eq:c5-5-9} only through its right side. The rate: $n_C\bar\varepsilon_{\mathrm{ph}}=4n\cdot L\bar\varepsilon_{\mathrm{ph}}$, and $\bar\varepsilon_{\mathrm{ph}}\leqslant 32K_{\chi} C_{\mathrm{pin}}C_{\mathrm{ph}}g_0^2(1+r)^4$ by the display above under \eqref{eq:c5-5R-T}(iii).
\end{proof}

\begin{remark}[scope of this subsection]\label{rem:c5-5R-5}
The mechanisms (parity of the kernel, the parameter $\zeta$ on the ordinary primitives) are those of Section~\ref{sec:c4-4-1}, re-proved here from the display of Lemma~\ref{lem:c5-4-1}, the factor bounds of Lemma~\ref{lem:c5-4-2} and Schwarz's lemma, and not imported. The differenced slots and the marked part keep the first-order prefactor $\bar\varepsilon_t$ (\eqref{eq:c5-5-8}), harmless because it multiplies the exterior sum and not $n_C$. The family-sum factor remains $e^{n_C\bar\varepsilon_{\mathrm{ph}}}\geqslant 1$ and is not uniform in $n_C$: like $e^{E^{\circ}_{D,C_n,T}}$ (the bound \eqref{eq:regional-square-error} is of order $\tau^2n$), it may grow with the perimeter at a second-order rate; a bound linear in $n_C$ in place of this exponential factor is not proved here. The inherited inputs of Section~\ref{sec:c5-5-0} (the pin sums for every $\rho\leqslant \rho_{\mathrm{pin}}$, Lemma~\ref{lem:c4-6-2}; the host comparison of Section~\ref{sec:c4-7}; the reference-gas expansion; the tree-graph inequality; the square geometry) remain inherited from the sections cited there and are not re-verified in this subsection.\end{remark}

\subsection{The labelled history sum}
\label{sec:c5-6}

The history sum of Section~\ref{sec:sourced-effective-interaction} is the sum over the branch data $\xi=(\Omega_1,R_1,V,V_0,A_{\mathrm e})$ inside the outer operation $\mathcal O$, each term carrying the weight $\mathsf w_{\xi}= e^{c_{\xi}}I_{\xi;1}=\int_{\mathfrak g^I}\rho_{\xi}\,\dd a_{\mathrm{old}}$, with $\rho_{\xi}\coloneqq e^{c_{\xi}}K_{\xi;1}$ the marked density \eqref{eq:c3-7-8} of Lemma~\ref{lem:c3-7-3} (the composed retained cutoff $\chi_{\xi,\mathrm e}$ and $\zeta_{\xi}$ multiply it), and the conditional Wilson expectation $F_{\xi}(1)=\mathcal J^a_{\xi}[f_{\xi}]/\mathcal J^a_{\xi}[1]$ (Lemma~\ref{lem:collar}, \eqref{eq:c-collar-identity}), and the leading term, the defects and the exterior response are read off \eqref{eq:c-one-step-branch-identity}. After Section~\ref{sec:c5-3} every history term is the finite sum of its labelled terms ($\mathrm{lab}=(B^{(0)},B^{(1)},\dots)$, Lemma~\ref{lem:c5-3-1}(a)), with $\sum_{\mathrm{lab}}\mathbf 1_{\mathrm{lab}}=1$ pointwise (Section~\ref{sec:c5-3-4}(c)); a labelled term has the final region $\Lambda\coloneqq \Lambda_1^{\mathrm{fin}}(\mathrm{lab})$, carrier $D(\Lambda)$, exterior $J(\Lambda)= J\cup(I\setminus D(\Lambda))$, and the objects of Lemma~\ref{lem:c5-2-5} and Sections~\ref{sec:c5-4}--\ref{sec:c5-5R}. Write $\mathcal J^{a,\mathrm{lab}}_{\xi}[f]\coloneqq \int_{\mathfrak g^I}\rho_{\xi}\mathbf 1_{\mathrm{lab}}f\,\dd a_{\mathrm{old}}$ for bounded Borel $f$ on $\mathfrak g^I$; then $\mathcal J^a_{\xi}[f]=\sum_{\mathrm{lab}}\mathcal J^{a,\mathrm{lab}}_{\xi}[f]$ (Tonelli for $f\geqslant 0$, then by linearity), and the extra resummation index is $\mathrm{lab}$.

\begin{lemma}[the labelled branch identity]\label{lem:c5-6-1}
On every history and every labelled term, with the convention that $e^{c^{\Lambda}_{\xi}}I^{\Lambda}_1F^{\Lambda}(1)$ and $e^{c^{\Lambda}_{\xi}}I^{\Lambda}_1$ denote the integrals $\int_{\mathfrak g^{D(\Lambda)}}\rho_{\xi}(a,a_{\mathrm{old}}|_{I\setminus D(\Lambda)})f_{\xi}(a)\,\dd a$ and $\int_{\mathfrak g^{D(\Lambda)}}\rho_{\xi}(a,a_{\mathrm{old}}|_{I\setminus D(\Lambda)})\,\dd a$ over the retreated fibre (here $f_{\xi}(a)$ stands for $f_{\xi}(a,a_{\mathrm{old}}|_{I\setminus D(\Lambda)})$; both integrals vanish where the fibre is empty, i.e. where $| a_{\mathrm{old},b}|_{\sigma}\geqslant r$ for some $b\in I\setminus D(\Lambda)$, since $\rho_{\xi}$ carries the box factor $\mathbf 1_{\mathcal D_{\xi}}$, and also wherever the composed retained cutoff or $\zeta_{\xi}$ vanishes, these being factors of $\rho_{\xi}$ which multiply both sides of the identity of Lemma~\ref{lem:c5-2-5}; $c^{\Lambda}_{\xi}$ and $F^{\Lambda}(1)$ themselves are used only where the fibre is nonempty and these cutoffs are positive),
\begin{equation}\label{eq:c5-6-1}
    \mathcal J^{a,\mathrm{lab}}_{\xi}[f_{\xi}]=\int_{\mathfrak g^{I\setminus D(\Lambda)}}\mathbf 1_{\mathrm{lab}}e^{c^{\Lambda}_{\xi}}I^{\Lambda}_1F^{\Lambda}(1)\,\dd a_{\mathrm{old}}|_{I\setminus D(\Lambda)}\;,\qquad\mathsf w^{\mathrm{lab}}_{\xi}=\int_{\mathfrak g^{I\setminus D(\Lambda)}}\mathbf 1_{\mathrm{lab}}e^{c^{\Lambda}_{\xi}}I^{\Lambda}_1\,\dd a_{\mathrm{old}}|_{I\setminus D(\Lambda)}\;,
\end{equation}
and, on $A^{\mathrm{fin}}_{\mathrm{cov}}$ (an event depending on the labels through the region only), on the support of $\mathbf 1_{\mathrm{lab}}$ (which carries the collar condition $| A_{\mathrm e}|_{\sigma}<\vartheta r$ on $J_{\mathrm{col}}(\Lambda)$ required by Section~\ref{sec:c5-4} and Lemma~\ref{lem:c5-4-2}; every use of \eqref{eq:c5-6-2} below is multiplied by $\mathbf 1_{\mathrm{lab}}$), where the fibre is nonempty, under the hypotheses of Lemma~\ref{lem:c5-4-2} and \eqref{eq:c5-5-T} (through Lemma~\ref{lem:c5-5-3}, whose absolute convergence and reference gas rest on the inputs listed in Section~\ref{sec:c5-5-0}),
\begin{equation}\label{eq:c5-6-2}
    F^{\Lambda}(1)= S^{\circ}_{\Lambda} W_{\chi}(C;U_1)+\sum_{\varnothing\neq\mathfrak B\text{ compatible}}S^{\circ[\mathfrak B]}_{\Lambda} R^0_{\Lambda}(\mathfrak B)+\mathcal B^{\Lambda}_{\xi}[f_{\xi}]\;,
\end{equation}
the retreated carrier's version of \eqref{eq:c-one-step-branch-identity}, with $S^{\circ}_{\Lambda}$, $S^{\circ[\mathfrak B]}_{\Lambda}$, $R^0_{\Lambda}(\mathfrak B)$ of Lemma~\ref{lem:c5-5-3} and $\mathcal B^{\Lambda}_{\xi}[f_{\xi}]$ of Lemma~\ref{lem:c5-4-4}, bounded by Propositions~\ref{prop:c5-5-9} and~\ref{prop:c5-5R-4} (the latter under \eqref{eq:c5-5R-T}); here $f_{\xi}(0)= W_{\chi}(C;U_1)$ on $A^{\mathrm{fin}}_{\mathrm{cov}}$ because the free inputs of the contour lie in $D(\Lambda)$, where the free vector vanishes at $a=0$ (proof).
\end{lemma}

\begin{proof}
By Lemma~\ref{lem:c5-2-5}, pointwise $e^{c^{\Lambda}_{\xi}}K^{\Lambda}_{\xi;1}(a)=\mathbf 1_{\mathcal D^{\Lambda}}(a)|\det D\Psi_{\xi}(A_{\mathrm e})|\exp\mathscr E_{\xi}(\Psi_{\xi}(a+A^{\Lambda}_{\mathrm e}))=\rho_{\xi}(a_{\mathrm{old}})$ at $a_{\mathrm{old}}=(a,a_{\mathrm{old}}|_{I\setminus D(\Lambda)})$ ($\rho_{\xi}= e^{c_{\xi}}K_{\xi;1}$ by definition, \eqref{eq:c3-7-8} of Lemma~\ref{lem:c3-7-3}: the right side of \eqref{eq:c3-7-8} and that of the identity of Lemma~\ref{lem:c5-2-5} are the same function of $a_{\mathrm{old}}$ written in the two splits, since $\Psi_{\xi}(a+A^{\Lambda}_{\mathrm e})=\Psi_{\xi}(a_{\mathrm{old}}+A_{\mathrm e})$ in the original variables and the cutoff factors are common), and $f_{\xi}$, $\mathbf 1_{\mathrm{lab}}$ are functions of the same variables ($\mathbf 1_{\mathrm{lab}}$ of $A_{\mathrm e}$ and $a_{\mathrm{old}}|_{I\setminus D(\Lambda)}$ only); integrate over $a\in\mathfrak g^{D(\Lambda)}$ first (Fubini; $\rho_{\xi}\geqslant 0$, $f_{\xi}$ and $\mathbf 1_{\mathrm{lab}}$ bounded): $\int K^{\Lambda}_{\xi;1}f_{\xi}\,\dd a= I^{\Lambda}_1F^{\Lambda}(1)$ by the definition of $\nu^{\Lambda}_1$ in Section~\ref{sec:c5-4} where the fibre is nonempty (both sides vanish otherwise), and with $f=1$ the weight. For \eqref{eq:c5-6-2}: $F^{\Lambda}(1)= F^{\Lambda}(0)+\mathcal B^{\Lambda}_{\xi}[f_{\xi}]$ (Lemma~\ref{lem:c5-4-4}), $F^{\Lambda}(0)= S^{\circ}_{\Lambda} f_{\xi}(0)+\sum_{\mathfrak B}S^{\circ[\mathfrak B]}_{\Lambda} R^0_{\Lambda}(\mathfrak B)$ (Lemma~\ref{lem:c5-5-3}(c) at $t=0$), and $f_{\xi}(0)= W_{\chi}(C;U_1)$ on $A^{\mathrm{fin}}_{\mathrm{cov}}$: at $a=0$ the free vector is ($0$ on $D(\Lambda)$, $A^{\Lambda,\Psi}$ on $J(\Lambda)$), every contour bond is a bond of $\Omega_1$ and its free inputs lie in $D(\Lambda)$, so the argument of Lemma~\ref{lem:covered} with $I$ replaced by $D(\Lambda)$ gives $U_{\xi}(0)= U_1$ on the contour (the de-younged bonds, though nonzero, are not inputs of the contour on $A^{\mathrm{fin}}_{\mathrm{cov}}$).
\end{proof}

\subsubsection{The species and their locality}
The species of the labelled term are the final domains $Z$ for the carrier $D(\Lambda)$ with their activities $k^{f,t}_{\Lambda}(Z)$ and probe measures $\nu^{f,t}_{\Lambda,Z}$ (Lemmas~\ref{lem:c5-4-1}, \ref{lem:c5-5-2}), the groups $W$ with their weights $w^t_{\Lambda}(W)$ and reference activities $a^{\Lambda}_W$ (Lemma~\ref{lem:c5-5-3}), and the marked-polymer objects built from them (the pins of Lemma~\ref{lem:c5-5-8}). They depend on the labels through $D(\Lambda)$ only. Let $\mathcal N_{\mathrm{lab}}$ be the set of $LM$-cubes meeting a bond of $I\setminus D(\Lambda)$ (no site lies on a cube face, so a bond meets exactly the one or two $LM$-cubes containing its endpoints, and $|\mathcal N_{\mathrm{lab}}|\leqslant 2| I\setminus D(\Lambda)|$).

\begin{lemma}[locality of the species]\label{lem:c5-6-2}
Under the hypotheses of Lemma~\ref{lem:c5-5-2} (Lemma~\ref{lem:c5-4-2} and \eqref{eq:c5-5-T}), let $Z$ be a union of $LM$-cubes containing no cube of $\mathcal N_{\mathrm{lab}}$, equivalently: every bond of $I$ with an endpoint in $Z$ lies in $D(\Lambda)$. Then for every selection and split selection with final domain $Z$, the selected integrals at the carrier $D(\Lambda)$ and at the carrier $I$ coincide at $t=0$: $k^{f,0}_{\Lambda}(Z)= k^f_{\xi}(Z)$, $\nu^{f,0}_{\Lambda,Z}=\nu^f_{\xi,Z}$, and likewise for the reference data ($k^{f,\circ}_{\Lambda}(Z)= k^{f,\circ}_{\xi}(Z)$, $\hat\nu^{f,\circ}_{\Lambda,Z}=\hat\nu^{f,\circ}_{\xi,Z}$). Consequently, for a group $W$ none of whose roots and cluster constituents contains a cube of $\mathcal N_{\mathrm{lab}}$ ($W$ is then a group for both carriers: the probe label set $\mathsf L_Z$ of each of its domains is the same at both carriers, since $Y_b\subseteq Z$ gives $b\in E(Y_b)\subseteq E(Z)\subseteq D(\Lambda)$, and its contact sets and hulls are determined by its domains and the contour), $w^0_{\Lambda}(W)= w_{\xi}(W)$, $a^{\Lambda}_W= a_W$, and on $A^{\mathrm{fin}}_{\mathrm{cov}}$ the pairings $\mathfrak Q^0_{\Lambda}(\mathcal A(\mathcal S);f_{\xi})=\mathfrak Q_{\xi}(\mathcal A(\mathcal S);f_{\xi})$ for families $\mathcal S$ of such groups, so $A^0_{\Lambda}(\mathcal S)= A_{\xi}(\mathcal S)$ and $R^0_{\Lambda}(\mathcal S)= R_{\xi}(\mathcal S)$. A species that differs between the two carriers has a root or constituent containing a cube of $\mathcal N_{\mathrm{lab}}$.
\end{lemma}

\begin{proof}
Fix a selection $\mathfrak s=(\mathcal D,\mathsf P,Z)$ (or a split selection) at the carrier $I$ with final domain $Z$; it is a selection at the carrier $D(\Lambda)$ and conversely: the hosts $Y\in\mathcal D$ lie in $Z$ ($Y\subseteq Z_0\subseteq Z_0'\subseteq Z$), their primitives have free-input sets in $Y$, whose young part lies in $E(Y)=\{ b\in\text{carrier}\, : \,b\cap Y\neq\varnothing\}$, the same set for both carriers by hypothesis, so a primitive assigned to $Y$ is near relative to $I$ exactly when it is near relative to $D(\Lambda)$, and $\mathbf D_0(\Lambda)$-hosts inside $Z$ are the $\mathbf D_0$-hosts inside $Z$ (probe hosts $Y_b\subseteq Z$ have $b\in E(Y_b)$, the same set); the bad bonds $\mathsf P$ meet $Z$ (their incident cubes are in $Z_0\subseteq Z$, Definition~\ref{def:selection}), so $\mathsf P\subseteq D(\Lambda)\setminus E(Y_0)$ and $\mathsf P\subseteq I\setminus E(Y_0)$ are the same condition; the seed and $Z_0'$ are formed from $\mathcal D$, $\mathsf P$ alone, and $E_0= E(Z_0)$ (Section~\ref{sec:fixed-first-kernel}) is the same set for both carriers by hypothesis. The integrand of \eqref{eq:c-selected-integral} at $t=0$: the potentials $\mathbf W_0(Y,B)=\sum_{\mathfrak p\, : \,Y_{\mathfrak p}= Y}\mathfrak p(B)$ (exterior input zero in both splits) and the cutoffs are the same functions of $B|_{E_0}$; $C_{E_0}=(\iota^*_{E_0}H\iota_{E_0})^{-1}$ involves the entries of $H$ between bonds of $E_0$, which are in both carriers; $\Gamma_{E_0}(\sigma)=\iota^*_{E_0}H\Pi_{\text{carrier}\setminus E_0}S_{Z_0}(\sigma)$ involves, for a row $e\in E_0$, the entries $H_{eb'}$ with $d(e_-,b'_-)\leqslant R$, so $b'\subset\operatorname{int}D\subset Z$ for the component $D$ of $e$ (proof of Lemma~\ref{lem:component-blocks}), and the walk series $S_{Z_0}(\sigma)_{b'b}$ with $\sigma_{\Delta}=0$ off $Z$, whose nonvanishing walks have $D(\omega)\subset Z$ (every cube meeting $D(\omega)$ is in $Z$, by the weight \eqref{eq:fixed-kernel-walk-weakening}); the blocks $(\mathsf A_H^r)_{b_0b_1}$, $(\mathsf P_H)_{b_jb_{j+1}}$ of a walk are entries of powers of $\id-H/s_H$ restricted to the carrier, and each depends only on the carrier's $H$-coefficients within distance $(n+1)R$ of its first bond, which lie in $D(\omega)\subset Z$ by the interior margin of $D_b$ (Section~\ref{sec:fixed-first-kernel}: $D_b$ contains the bonds and coefficient stencils of every block starting at $b$ with an interior margin); hence every bond entering $S_{Z_0}(\sigma)$ and $\Gamma_{E_0}(\sigma)$ is a bond of $I$ meeting $Z$, in $D(\Lambda)$ by hypothesis, and every coefficient is an entry of $H$ between two such bonds, the same for both carriers (principal restrictions of the one Hessian $H_{\xi}$), so the two carriers give the same $\Gamma_{E_0}(\sigma)$ and the same $C_{E_0}$; the $X$-integral is over $\mathfrak g^{\text{carrier}}$ with the standard Gaussian, but $\Gamma_{E_0}(\sigma)X$ involves $X$ only on bonds of $Z$ (the columns of $\Gamma$ outside $\bigsqcup J_D$ vanish, Lemma~\ref{lem:component-blocks}), so the remaining coordinates integrate to one in both cases. The Dirac tensors of a split selection are the same. Hence $\mathcal A^{f,0}_{\Lambda;\mathfrak s}[F]=\mathcal A^f_{\xi;\mathfrak s}[F]$ for every $F$, giving the activities and measures; the reference case is the same argument with $H^{\circ}$. The group statements follow: $w^0_{\Lambda}(W)$ is the product of the activities of its constituents with the combinatorial factors of its structure (the Ursell coefficients of its cluster constituents and $1/k(W)!$), which depend on the incompatibility graph only, $a^{\Lambda}_W$ an integral of the reference source word (the same function, the reference links depending on the root variables only) against the same measures, and $\mathfrak Q$ involves $f_{\xi}(\mathsf e_{\mathcal A}B)$, which on $A^{\mathrm{fin}}_{\mathrm{cov}}$ depends on $B|_{I_C}$ with $I_C\subseteq D(\Lambda)$ and is the same function for both carriers.
\end{proof}

\begin{lemma}[the reference gas at the retreated carrier]\label{lem:c5-6-3}
Under \eqref{eq:c5-5-T}(iv), $S^{\circ}_{\Lambda}= S^{\circ}_{\xi,C}\cdot\exp(\eta_{\mathrm{lab}})$ with
\begin{multline*}
    |\eta_{\mathrm{lab}}|\leqslant \sum_{\Gamma\text{ containing a group }W\text{ with }F(W)\cap\mathcal N_{\mathrm{lab}}\neq\varnothing}\bigl(|\Phi^{\circ}_{\Lambda}(\Gamma)|+|\Phi^{\circ}_{\xi}(\Gamma)|\bigr)\\
    \leqslant \sum_{\Delta\in\mathcal N_{\mathrm{lab}}}\sum_{\Gamma\text{ hits }\Delta}\bigl(|\Phi^{\circ}_{\Lambda}(\Gamma)|+|\Phi^{\circ}_{\xi}(\Gamma)|\bigr)\leqslant 2| I\setminus D(\Lambda)|\cdot\varepsilon^{\circ}_g\;,
\end{multline*}
where $\varepsilon^{\circ}_g\coloneqq \sup_{\Delta}\sum_{\Gamma\text{ hits }\Delta}(|\Phi^{\circ}_{\Lambda}(\Gamma)|+|\Phi^{\circ}_{\xi}(\Gamma)|)\leqslant 2C_{\mathrm{pin}}\rho_t$: a cluster $\Gamma$ hits the cube $\Delta$ when one of its groups $W$ has $\Delta\in F(W)$; the rooted bound of Lemma~\ref{lem:c5-5-3}(d), $\sum_{\Gamma\ni W}|\Phi^{\circ}(\Gamma)|\leqslant | a_W| e^{13\nu\sigma(W)/8}$, gives $\sum_{\Gamma\text{ hits }\Delta}|\Phi^{\circ}(\Gamma)|\leqslant \sum_{W\, : \,\Delta\in F(W)}| a_W| e^{13\nu\sigma(W)/8}\leqslant \sum_{W\, : \,\Delta\in F(W)}e^{4\nu\sigma(W)}M_{\Lambda}(W)\leqslant C_{\mathrm{pin}}\rho_t$ by $| a_W|\leqslant M_{\Lambda}(W)$ and the first display of Lemma~\ref{lem:c5-5-8}(a), at the carrier $D(\Lambda)$ and, by the closing remarks of Lemma~\ref{lem:c5-5-3}(d) and Lemma~\ref{lem:c5-5-8}(a), at the carrier $I$. Moreover $| S^{\circ}_{\Lambda}|\leqslant 1$.
\end{lemma}

\begin{proof}
Under \eqref{eq:c5-5-T}(iv) the reference group gas satisfies the Kotecký--Preiss condition at both carriers (Lemma~\ref{lem:c5-5-3}(d); the reference group gas at $I$ is that of Section~\ref{sec:sourced-effective-interaction}), so its cluster expansion (Appendix~\ref{app:kp}) gives $\log S^{\circ}=\sum_{\Gamma}\Phi^{\circ}(\Gamma)$, absolutely convergent, at both carriers, and $\log S^{\circ}_{\Lambda}-\log S^{\circ}_{\xi,C}=\sum_{\Gamma}[\Phi^{\circ}_{\Lambda}(\Gamma)-\Phi^{\circ}_{\xi}(\Gamma)]$; a cluster all of whose groups have roots and constituents containing no cube of $\mathcal N_{\mathrm{lab}}$ has $\Phi^{\circ}_{\Lambda}(\Gamma)=\Phi^{\circ}_{\xi}(\Gamma)$ by Lemma~\ref{lem:c5-6-2} (the Ursell coefficient depends on the incompatibility graph only, the group activities coincide), and the groups present for one carrier and not the other contain a cube of $\mathcal N_{\mathrm{lab}}$ as well (a root or constituent met by a de-younged bond). $| S^{\circ}_{\Lambda}|\leqslant 1$: $S^{\circ}_{\Lambda}=\int f^{\circ} K^{\circ}_{\Lambda}\,\dd a/I^{\circ}_{\Lambda}$ with $K^{\circ}_{\Lambda}>0$ on the box (Lemma~\ref{lem:c5-4-0}(c) for the reference data) and $| f^{\circ}|\leqslant 1$ ($\chi$ is the normalised character, $|\chi|\leqslant \chi(\id)=1$).
\end{proof}

\subsubsection{The accounting}
Three quantities acquire the index $\mathrm{lab}$: the history weight ($\mathsf w_{\xi}=\sum_{\mathrm{lab}}\mathsf w^{\mathrm{lab}}_{\xi}$, \eqref{eq:c5-6-1}), the exterior-response/uncovered term $E_{\mathrm{bdry}}$, and the defect term $E_{\mathrm{def}}$. Define
\begin{equation*}
    E^{\mathrm{lab}}_{\mathrm{bdry}}\coloneqq \mathcal Z_0^{-1}\int\mathcal O\Bigl[\sum_{\mathrm{lab}}\int_{\mathfrak g^{I\setminus D(\Lambda)}}\mathbf 1_{\mathrm{lab}}e^{c^{\Lambda}_{\xi}}I^{\Lambda}_1\bigl(\mathcal B^{\Lambda}_{\xi}[f_{\xi}]\mathbf 1_{A^{\mathrm{fin}}_{\mathrm{cov}}}+F^{\Lambda}(1)\mathbf 1_{(A^{\mathrm{fin}}_{\mathrm{cov}})^c}\bigr)\dd a_{\mathrm{old}}\Bigr](V)\,\dd V\;,
\end{equation*}
the labelled version of the $E_{\mathrm{bdry}}$ of Appendix~\ref{app:exterior-chessboard} (the exterior response on covered terms, the whole conditional expectation on uncovered terms, $| F^{\Lambda}(1)|\leqslant 1$; the integrands are read with the convention of Lemma~\ref{lem:c5-6-1}, as $\int_{\mathfrak g^{D(\Lambda)}}\rho_{\xi}(\dots)\,\dd a$, so they are defined and vanish where the fibre is empty), and $E^{\mathrm{lab}}_{\mathrm{def}}$ the same with $\sum_{\mathfrak B}S^{\circ[\mathfrak B]}_{\Lambda} R^0_{\Lambda}(\mathfrak B)\mathbf 1_{A^{\mathrm{fin}}_{\mathrm{cov}}}$ in place of the bracket; the leading term is $\mathcal Z_0^{-1}\int\mathcal O[\sum_{\mathrm{lab}}\int\mathbf 1_{\mathrm{lab}}e^{c^{\Lambda}_{\xi}}I^{\Lambda}_1S^{\circ}_{\Lambda}\mathbf 1_{A^{\mathrm{fin}}_{\mathrm{cov}}}\,\dd a_{\mathrm{old}}\cdot W_{\chi}(C;U_1)]\,\dd V$. By \eqref{eq:c5-6-1}--\eqref{eq:c5-6-2} these three sum to the labelled Wilson numerator, exactly.

\begin{proposition}[the boundary accounting with labels]\label{prop:c5-6-4}
On compatible even fine periods exceeding $200\ell$, under the hypotheses of Proposition~\ref{prop:c5-5-9}, for $g_0$ below the thresholds of Lemma~\ref{lem:c5-2-4}, of Proposition~\ref{prop:c5-2-6} ($g_*(\vartheta)$) and of \eqref{eq:c5-2-3}, with $q$, $q_{\vartheta}$, $N_{\mathrm{ret}}$, $C_L$ of Sections~\ref{sec:c5-2}--\ref{sec:c5-3} $\bar\varepsilon_t$, $C_1$, $C_2$, $\mu'$ of Section~\ref{sec:c5-5}, and $N_{\mathrm{lat}}(\mu')=\sum_{s\geqslant 0}4(2s+1)^4e^{-\mu's}$ of Lemma~\ref{lem:c5-1-2} (the labelled integrands $S^{\circ}_{\Lambda}$, $F^{\Lambda}(t)$, $\mathcal B^{\Lambda}_{\xi}[f_{\xi}]$ are Borel in the branch data: the labels take finitely many values, and on each labelled term these are ratios and differences of the Borel integrals of the measurability paragraph of Section~\ref{sec:sourced-effective-interaction}):
\begin{gather}
    \mathcal Z_0^{-1}\int\mathcal O\Bigl[\sum_{\mathrm{lab}}\int_{\mathfrak g^{I\setminus D(\Lambda)}}\mathbf 1_{\mathrm{lab}}e^{c^{\Lambda}_{\xi}}I^{\Lambda}_1\mathbf 1_{A^{\mathrm{fin}}_{\mathrm{cov}}}\frac{|\mathcal B^{\Lambda}_{\xi}[f_{\xi}]|}{| S^{\circ}_{\Lambda}|}\,\dd a_{\mathrm{old}}\Bigr]\dd V\notag\\
    \leqslant e^{2n_C\bar\varepsilon_t}\bigl(C_1g_0r+C_2\bar\varepsilon_tr^{d_G}\bigr)\bigl(1+e^{\mu_H(2L+1)/2}\bigr)N_{\mathrm{lat}}(\mu')\,n_C(q+N_{\mathrm{ret}}q_{\vartheta})\;,\label{eq:c5-6-3}\\
    \mathcal Z_0^{-1}\int\mathcal O\Bigl[\sum_{\mathrm{lab}}\mathsf w^{\mathrm{lab}}_{\xi}\mathbf 1_{(A^{\mathrm{fin}}_{\mathrm{cov}})^c}\Bigr]\dd V\leqslant C_Ln_C(q+N_{\mathrm{ret}}q_{\vartheta})\quad\text{(\eqref{eq:c5-3-coverage-average})},\notag
\end{gather}
hence, with $| S^{\circ}_{\Lambda}|\leqslant 1$ (Lemma~\ref{lem:c5-6-3}),
\begin{equation}\label{eq:c5-6-4}
    | E^{\mathrm{lab}}_{\mathrm{bdry}}|\leqslant n_C(q+N_{\mathrm{ret}}q_{\vartheta})\Bigl[e^{2n_C\bar\varepsilon_t}\bigl(C_1g_0r+C_2\bar\varepsilon_tr^{d_G}\bigr)\bigl(1+e^{\mu_H(2L+1)/2}\bigr)N_{\mathrm{lat}}(\mu')+C_L\Bigr]\;,
\end{equation}
and for $g_0$ below a further threshold depending on $L$, $M$, $M_2$, the group, the source constants, $\vartheta$ and the character, $r^{d_G}N_{\mathrm{ret}}(q+N_{\mathrm{ret}}q_{\vartheta})\leqslant e^{-c_q\vartheta^2\pi_*^2/2}$ (the factor $N_{\mathrm{ret}}$ is not needed for \eqref{eq:c5-6-4}; it costs only a logarithm, proof, and is kept for the form of the statement), so that, with $g_0\leqslant 1\leqslant r$, $d_G\geqslant 1$ and $\bar\varepsilon_t\leqslant 1$ (\eqref{eq:c5-5-T}(v)), $C_1g_0r+C_2\bar\varepsilon_tr^{d_G}\leqslant (C_1+C_2)r^{d_G}$ and $| E^{\mathrm{lab}}_{\mathrm{bdry}}|\leqslant Cn_Ce^{2n_C\bar\varepsilon_t}e^{-c_q\vartheta^2\pi_*^2/2}$ with $C=(C_1+C_2)(1+e^{\mu_H(2L+1)/2})N_{\mathrm{lat}}(\mu')+C_L$: the form $| E_{\mathrm{bdry}}|\leqslant Cn_C(1+g_0r)e^{-c_*\pi_*^2}$ of Appendix~\ref{app:exterior-chessboard} with the exponent $c_q\vartheta^2/2$ in place of $c_*$ and the family-sum factor $e^{2n_C\bar\varepsilon_t}$ of \eqref{eq:c5-5-9} (it is not a history weight; Proposition~\ref{prop:c5-5R-4} replaces it by $e^{n_C\bar\varepsilon_{\mathrm{ph}}}$ under \eqref{eq:c5-5R-T}). Both perimeter factors grow with $n_C$: the bound is a statement about the fixed loop $C$ as $g_0\to0$, the factor $e^{2n_C\bar\varepsilon_t}\leqslant e^{2n_C}$ being a constant of that loop.
\end{proposition}

\begin{proof}
\eqref{eq:c5-6-3}: on $A^{\mathrm{fin}}_{\mathrm{cov}}$, \eqref{eq:c5-5-9} gives $|\mathcal B^{\Lambda}|/| S^{\circ}_{\Lambda}|\leqslant e^{2n_C\bar\varepsilon_t}[C_1g_0r\mathfrak S_1(\Lambda)+C_2\bar\varepsilon_tr^{d_G}\mathfrak S_2(\Lambda)]$ with $\mathfrak S_1(\Lambda)\leqslant e^{\mu_H(2L+1)/2}\mathfrak S_2(\Lambda)$ (Corollary~\ref{cor:c5-5-10}), and $\mathfrak S_2(\Lambda)=\sum_{e\in J(\Lambda)}e^{-\mu'\operatorname{dist}(e,C)}\leqslant \sum_ee^{-\mu'\operatorname{dist}(e,C)}\mathbf 1\{ e\in E(B^{\mathrm{fin}}_{\xi})\}$, $B^{\mathrm{fin}}_{\xi}\coloneqq (\Lambda_1^{\mathrm{fin}})^c$, because a bond not in $D(\Lambda)$ has no endpoint in $\Lambda_1^{\mathrm{fin}}$ (Section~\ref{sec:c5-3-4}(a)). The prefactor is a constant (independent of the history, the labels and the fields: \eqref{eq:c5-5-9} is uniform in the exterior data within the term's support), and $\mathfrak S_2(\Lambda)$, $S^{\circ}_{\Lambda}$ and $A^{\mathrm{fin}}_{\mathrm{cov}}$ depend on the labels through $\Lambda$ only, so the integral over $a_{\mathrm{old}}|_{I\setminus D(\Lambda)}$ of $\mathbf 1_{\mathrm{lab}}e^{c^{\Lambda}_{\xi}}I^{\Lambda}_1\mathbf 1_{A^{\mathrm{fin}}_{\mathrm{cov}}}|\mathcal B^{\Lambda}|/| S^{\circ}_{\Lambda}|$ is at most (prefactor)$\cdot\mathfrak S_2(\Lambda)\cdot\mathsf w^{\mathrm{lab}}_{\xi}\mathbf 1_{A^{\mathrm{fin}}_{\mathrm{cov}}}$. By \eqref{eq:c5-3-mark-inequality}, $\mathbf 1\{ e\in E(B^{\mathrm{fin}}_{\xi})\}\leqslant \mathbf 1\{ e\in E(\mathfrak B_0)\}+\sum_{e'\, : \,x_{e'}\text{ within }102\ell+1\text{ of an endpoint of }e}M_{e'}$ (Lemma~\ref{lem:c5-3-1}(c)--(d)); multiply by $\mathsf w^{\mathrm{lab}}_{\xi}\geqslant 0$, sum over the labels and histories, apply $\mathcal O$, divide by $\mathcal Z_0$ (Tonelli, finitely many labels per history): the first term gives $p_e^{(0)}\leqslant q$ (Lemma~\ref{lem:c5-3-0}(b)), the second $\sum_{e'}p^{\vartheta}_{e'}\leqslant N_{\mathrm{ret}}q_{\vartheta}$ (Lemma~\ref{lem:c5-3-1}(d), Proposition~\ref{prop:c5-2-6}). Then $\sum_ee^{-\mu'\operatorname{dist}(e,C)}\leqslant n_CN_{\mathrm{lat}}(\mu')$ (at most $4n_C(2s+1)^4$ bonds at $\ell^{\infty}$ distance $\leqslant s$ from the contour's initial sites, as in Corollary~\ref{cor:c5-5-10}). The second display is \eqref{eq:c5-3-coverage-average}. \eqref{eq:c5-6-4} adds the two with $|\mathcal B^{\Lambda}|\leqslant | S^{\circ}_{\Lambda}|\cdot(\dots)$ and $| F^{\Lambda}(1)|\leqslant 1$. The comparison: $\log(r^{d_G}N_{\mathrm{ret}})\leqslant d_G(\log A_1+p_0\log\log g_0^{-2})+C+4\mathfrak r\log(2\log g_0^{-2})$ by \eqref{eq:c5-3-1} and $r= A_1(\log g_0^{-2})^{p_0}$, while $\log(1/q_{\vartheta})\geqslant c_q\vartheta^2\pi_*^2\geqslant c_q\vartheta^2A_0^22^{-2p_0}(\log g_0^{-2})^{2p_0}$ by \eqref{eq:c5-2-3} and \eqref{eq:c5-3-2}, and $q\leqslant q_{\vartheta}$ (proof of Lemma~\ref{lem:c5-2-4}: $\theta_{\vartheta}\leqslant \theta$ for $\vartheta\leqslant 1$, and $z_P$, $z_G$ are decreasing in $\theta$ at fixed $n_W$), so $q+N_{\mathrm{ret}}q_{\vartheta}\leqslant (1+N_{\mathrm{ret}})q_{\vartheta}$ (alternatively, without comparing $q$ and $q_{\vartheta}$: Proposition~\ref{prop:c3-7-5} gives $q\leqslant e^{-c_*\pi_*^2}= e^{-c_q\pi_*^2}\leqslant e^{-c_q\vartheta^2\pi_*^2}$ for $g_0< g_*$ (\eqref{eq:c5-2-3} at $\vartheta=1$), and $q_{\vartheta}\leqslant e^{-c_q\vartheta^2\pi_*^2}$ by \eqref{eq:c5-2-3} itself, so $q+N_{\mathrm{ret}}q_{\vartheta}\leqslant (1+N_{\mathrm{ret}})e^{-c_q\vartheta^2\pi_*^2}$, which serves equally); every logarithm is dominated for small $g_0$, leaving $e^{-c_q\vartheta^2\pi_*^2/2}$. The polynomial prefactor $r^{d_G}$ and the factor $N_{\mathrm{ret}}$ in front are thus paid by the same log comparison ($\log r= O(\log\log g_0^{-2})$, $\log N_{\mathrm{ret}}= O(\log\log g_0^{-2})$ by \eqref{eq:c5-3-1}); they are prefactors, not the cost in the exponent growing with $r$ that the comparison of Section~\ref{sec:c5-3-4} set aside.
\end{proof}

\subsubsection{The labelled defect term}

\begin{proposition}[the labelled defect bound]\label{prop:c5-6-def}
Under the hypotheses of Proposition~\ref{prop:c5-6-4}, \eqref{eq:c5-5-T} and the hypotheses of Lemma~\ref{lem:c5-4-2} (under which Corollary~\ref{cor:c5-4-3} and Lemma~\ref{lem:c5-5-2} are stated), and those of Proposition~\ref{prop:c4-12-3} (the hypotheses stated at the start of Section~\ref{sec:c4-12}, $g_0$ below the thresholds of Lemmas~\ref{lem:c4-2-1}, \ref{lem:c4-6-2} and~\ref{lem:c4-6-3}, $R_{\Upsilon}\geqslant 2$ and $u\leqslant u_*$, $v\leqslant 1$ of Lemma~\ref{lem:c4-4-2}, and the small-coupling conditions of Sections~\ref{sec:c4-10-1} and~\ref{sec:c4-11-5}, the joint radii $\tau_{\mathrm b}$, $\theta_0$ and the source specifications used in Section~\ref{sec:c4-12}), with $\rho$ of \eqref{eq:c4-4-8} satisfying $\rho\leqslant \rho_{\mathrm{pin}}$ (a condition on $g_0$, since $\rho= O(g_0^2(1+r)^4)$), $\varepsilon_1= g_1p_0(g_1)$, $b=\lfloor LMR_1/16\rfloor$, $\pi_*=\min(p_0(g_0),p_0(g_1))$, the event $\mathsf N_b$ of \eqref{eq:c4-2-1} and the constants $C'_L$, $c_*$ of \eqref{eq:c4-2-2} and $C_{\chi}$, $c$ of \eqref{eq:c4-6-8}, write $\mathscr R^{\Lambda}\coloneqq \sum_{\varnothing\neq\mathfrak B\text{ compatible}}S^{\circ[\mathfrak B]}_{\Lambda} R^0_{\Lambda}(\mathfrak B)$ for the labelled defect of \eqref{eq:c5-6-2}.
\begin{enumerate}[label=(\alph*)]
    \item On every labelled term with final region $\Lambda$, for exterior data in the support of the term's collar indicators and exterior cutoffs, $\mathscr R^{\Lambda}= F^{\Lambda}(0)-S^{\circ}_{\Lambda} f_{\xi}(0)$ satisfies $|\mathscr R^{\Lambda}|\leqslant 2$, and on $A^{\mathrm{fin}}_{\mathrm{cov}}$ it depends on the term through $D(\Lambda)$ and on the history through $U_1$, $\Omega_1$, $\Lambda_1$ only (functions of $V$, $V_0$ and the regions), not on $a_{\mathrm{old}}|_{I\setminus D(\Lambda)}$ or $A_{\mathrm e}$; of this, only the independence from $a_{\mathrm{old}}|_{I\setminus D(\Lambda)}$ is used, in (b). On $A^{\mathrm{fin}}_{\mathrm{cov}}\cap\mathsf N_b^c$,
    \begin{equation}\label{eq:c5-6-def-1}
        |\mathscr R^{\Lambda}|\leqslant | S^{\circ}_{\Lambda}|\bigl(e^{C_{\chi} n_C\rho(\varepsilon_1+e^{-cb/(LM)})}-1\bigr)\;.
    \end{equation}
    \item Consequently
    \begin{equation}\label{eq:c5-6-def-2}
        | E^{\mathrm{lab}}_{\mathrm{def}}|\leqslant C'_Ln_Ce^{-(c_*/2)\pi_*^2}+\mathcal Z_0^{-1}\int\mathcal O\Bigl[\sum_{\mathrm{lab}}\mathsf w^{\mathrm{lab}}_{\xi}\mathbf 1_{A^{\mathrm{fin}}_{\mathrm{cov}}\cap\mathsf N_b^c}| S^{\circ}_{\Lambda}|\bigl(e^{C_{\chi} n_C\rho(\varepsilon_1+e^{-cb/(LM)})}-1\bigr)\Bigr]\dd V\leqslant C'_Ln_Ce^{-(c_*/2)\pi_*^2}+e^{C_{\chi} n_C\rho(\varepsilon_1+e^{-cb/(LM)})}-1\;,
    \end{equation}
    the bounds \eqref{eq:c4-12-7}--\eqref{eq:c4-12-8} of Proposition~\ref{prop:c4-12-3} with the same constants; $S^{\circ}_{\Lambda}$ is retained and not identified here.
\end{enumerate}
\end{proposition}

\begin{proof}
(a) By Lemma~\ref{lem:c5-5-3}(c) at exterior parameter $t=0$, $F^{\Lambda}(0)= S^{\circ}_{\Lambda} f_{\xi}(0)+\mathscr R^{\Lambda}$, with $\mathscr R^{\Lambda}$ the defect term of (c), the series over the countable index set of nonempty compatible families (Lemma~\ref{lem:c5-5-3}(a) gives the absolute convergence of the series $\sum_{\mathfrak B}A^t(\mathfrak B)$ and (c) the identity; the absolute summability of the series defining $\mathscr R^{\Lambda}$ on $A^{\mathrm{fin}}_{\mathrm{cov}}\cap\mathsf N_b^c$ is established by the termwise bounds of steps (ii)--(iii) below, the only place the series is used, and on the whole support of the exterior cutoffs it follows from Lemma~\ref{lem:c5-5-7}(a),(b) with Lemma~\ref{lem:c5-5-3}(d) and Lemma~\ref{lem:c5-5-8}(a); elsewhere only the closed form $F^{\Lambda}(0)-S^{\circ}_{\Lambda} f_{\xi}(0)$ and the bound $|\mathscr R^{\Lambda}|\leqslant 2$ enter); the inequality $\sum_{\mathfrak B\neq\varnothing}\prod_{W\in\mathfrak B}x_W\leqslant \exp(\sum_Wx_W)-1$ used below holds for countable sums. At $t=0$ the kernel $K^{\Lambda}_{\xi;0}$ of Section~\ref{sec:c5-4} has no linear term and its primitives are evaluated at $a$ alone (exterior variable zero), so, on $A^{\mathrm{fin}}_{\mathrm{cov}}$, $\nu^{\Lambda}_0$ and $F^{\Lambda}(0)$ depend on $U_1$, $\Omega_1$ and $D(\Lambda)$ only (its box indicator is the product cutoff $\Theta$ over $D(\Lambda)$ on the support of the exterior cutoffs, Lemma~\ref{lem:collar-coverage}(iv) and Lemma~\ref{lem:c5-3-3}(iv)); $S^{\circ}_{\Lambda}$ is built from the reference data at $D(\Lambda)$ (Lemma~\ref{lem:c5-5-3}(b)); and on $A^{\mathrm{fin}}_{\mathrm{cov}}$ every contour link is reconstructed from $a|_{D(\Lambda)}$ at the background $U_1$ (preamble of Lemma~\ref{lem:c5-5-7z}: the free inputs of the contour lie in $D(\Lambda)$, and the de-younged bonds are not inputs), so $f_{\xi}$ and $f_{\xi}(0)= W_{\chi}(C;U_1)$ (Lemma~\ref{lem:c5-6-1}) involve, besides the branch datum $U_1$, neither $a_{\mathrm{old}}|_{I\setminus D(\Lambda)}$ nor $A_{\mathrm e}$. $|\mathscr R^{\Lambda}|\leqslant 2$: $\nu^{\Lambda}_0$ is a probability measure ($K^{\Lambda}_{\xi;0}>0$ on the open box and $I^{\Lambda}_0>0$, Lemma~\ref{lem:c5-4-0}(c)) and $| f_{\xi}|\leqslant 1$, so $| F^{\Lambda}(0)|\leqslant 1$; $| S^{\circ}_{\Lambda}|\leqslant 1$ (Lemma~\ref{lem:c5-6-3}) and $| f_{\xi}(0)|\leqslant 1$.

For \eqref{eq:c5-6-def-1} the three steps of the proof of Lemma~\ref{lem:c4-6-5} are taken at the carrier $D(\Lambda)$. (i) Exclusion ratios: by Lemma~\ref{lem:c5-5-3}(d), under \eqref{eq:c5-5-T}(iv), $S^{\circ}_{\Lambda}\neq0$ and $| S^{\circ[\mathfrak B]}_{\Lambda}/S^{\circ}_{\Lambda}|=| q^{\circ}_{\Lambda}(\mathfrak B)|\leqslant \prod_{W\in\mathfrak B}e^{\theta^{\circ}(W)}$ with $\theta^{\circ}(W)\leqslant \nu\sigma(W)/16\leqslant \eta\mathfrak s(W)/8$ ($\nu=\eta$ and $\sigma(W)=\mathfrak s(W)$, Lemma~\ref{lem:c5-5-8}), the retreated carrier's form of \eqref{eq:c4-6-7}. (ii) Factorisation: by Lemma~\ref{lem:c5-5-7}(b) at exterior parameter zero, $R^0_{\Lambda}(\mathfrak B)=\dim(\pi)^{-1}\operatorname{tr}(\Upsilon^0_{W_1}U^0_1\cdots\Upsilon^0_{W_n}U^0_n)$, and the elementary trace bound $|\dim(\pi)^{-1}\operatorname{tr}(X_1\cdots X_m)|\leqslant \prod_j\| X_j\|$ (as in the proof of Lemma~\ref{lem:c5-5-7}(c)) with $\| U^0_j\|=1$ and $\|\Upsilon^0_W\|\leqslant K_{\chi}\mathfrak d^0_W$ (Lemma~\ref{lem:c5-5-7}(a)), $\mathfrak d^0_W\coloneqq \sup_{h\in G}| c^0_W(h)-a^{\Lambda}_W\chi(h)|$ (the $\mathfrak d_W(U_1)$ of \eqref{eq:c4-6-3} at the carrier $D(\Lambda)$; the superscript $0$ marks the exterior parameter zero at the physical background $U_1$, as in $c^0_W$, whereas the quantities of the coefficient path below carry its parameter $t$), gives $| R^0_{\Lambda}(\mathfrak B)|\leqslant \prod_{W\in\mathfrak B}K_{\chi}\mathfrak d^0_W$; hence, families being sets of distinct groups,
\begin{equation*}
    |\mathscr R^{\Lambda}|\leqslant | S^{\circ}_{\Lambda}|\sum_{\varnothing\neq\mathfrak B}\prod_{W\in\mathfrak B}K_{\chi} e^{\theta^{\circ}(W)}\mathfrak d^0_W\leqslant | S^{\circ}_{\Lambda}|\Bigl[\exp\Bigl(\sum_WK_{\chi} e^{\theta^{\circ}(W)}\mathfrak d^0_W\Bigr)-1\Bigr]\;.
\end{equation*}
(iii) The one-group sum: on $A^{\mathrm{fin}}_{\mathrm{cov}}\cap\mathsf N_b^c$,
\begin{equation}\label{eq:c5-6-def-3}
    \sum_WK_{\chi} e^{\theta^{\circ}(W)}\mathfrak d^0_W\leqslant C_{\chi} n_C\rho\bigl(\varepsilon_1+e^{-cb/(LM)}\bigr)\;,
\end{equation}
Lemma~\ref{lem:c4-6-4} at $D(\Lambda)$. Its proof (Section~\ref{sec:c4-12-2}) rests on \eqref{eq:c4-12-3} for the small interior groups \eqref{eq:c4-8-9}, on \eqref{eq:c4-12-4} for the others and on the weighted pins; the carrier enters these only through the following inputs, each supplied at $D(\Lambda)$.

Hessian along the coefficient path. For a small interior group with comparison rectangle $Q_W$ and path $U_{W,t}= e^{itB_W}$ (Section~\ref{sec:c4-9-1}) put $H_t\coloneqq (\mathsf H^{\Omega_1}_{U_{W,t}})_{D(\Lambda)D(\Lambda)}$. Lemma~\ref{lem:c4-7-3} holds for this compression, a compression of the $I$-compression of Section~\ref{sec:c4-9-1} since $D(\Lambda)\subseteq I$ (Lemma~\ref{lem:c5-2-5}): the two-sided bound passes to every compression by the min--max principle, the flat lower bound $\gamma_L$ included, and the Schur-test bound $c_L\tau_{\mathrm b}$ on the entry differences holds for a principal submatrix, whose row and column sums are at most those of the full matrix; so $\tfrac34\gamma_L\leqslant H_t\leqslant m_L$ and the $R$-locality of the entries (Lemma~\ref{lem:hessian-gauge-locality}) hold for every $t\in[0,1]$. This, with the Schur bound $c_H$ on the row and column norms of the entry derivatives $\mathscr D_{\omega} H_t/\|\omega\|_{\infty}$ (Section~\ref{sec:c4-10-2}, supplied by the finite analytic entries of Section~\ref{sec:c4-10-1}, which the principal submatrix inherits), is all that Sections~\ref{sec:c4-10-2} and~\ref{sec:c4-10-3} use of the background: the differentiated walk series \eqref{eq:c4-10-6}, the derivatives of $C_{E_0}$ and $\Gamma_{E_0}$ (built from $H_t$ with $\Pi_{D(\Lambda)\setminus E_0}$, as in Lemma~\ref{lem:c5-4-1}) and the Gaussian moment payment \eqref{eq:c4-10-14}, whose inputs $H_E-\alpha_G\geqslant \gamma/2$, $\| T_{\alpha}\|\leqslant 1/4$ and $\|\Gamma_1\|\leqslant A_*$ are those of Lemma~\ref{lem:gaussian-payment} for $H_t$, supplied by its proof in Appendix~\ref{app:fixed-kernel-payment} (\eqref{eq:fixed-kernel-reserve-payment}; that proof uses only \eqref{eq:fixed-kernel-inputs}, the two-sided bound with $\gamma=\gamma_L/2\leqslant \tfrac34\gamma_L$, $h_+= m_L$ as in Lemma~\ref{lem:c5-4-0}(a), the range and the partitions), and whose $d\leqslant 8d_G(LM+1)^4N_1(Z)$ is a count over the partitions.

Charts and primitives. Throughout this proof the exterior parameter is zero; the parameter $t\in[0,1]$ of this and the preceding paragraph is that of the coefficient path $U_{W,t}$. Section~\ref{sec:c4-10-1} is stated per stencil type; at the rows meeting $D(\Lambda)$ it applies because the joint chart condition $\| Z\|<2\tau_{\mathrm b}$ holds along the path $U_{W,t}$ at every such row for every $t\in[0,1]$, by the phase test \eqref{eq:c4-10-17}: $\| tB_W\|_{\infty}\leqslant \tau_{\mathrm b}/2$ on $Q_W$ and $B_W=0$ off $Q_W$ (the path runs from the identity to $e^{iB_W}$, not through $U_1$), so that the condition is that of the coefficient path at $I$ (Section~\ref{sec:c4-9-1}); with zero exterior variable ($s=1+r\leqslant s^{\Lambda}$); the hosts are those of $\mathbf D_0(\Lambda)$ (Lemma~\ref{lem:c5-4-0}(d)), defined from the lifted stencils and $M$ alone, with the counts of Lemma~\ref{lem:c2-potential-size}. Hence the parts of Lemmas~\ref{lem:c5-4-0}--\ref{lem:c5-4-2}, Corollary~\ref{cor:c5-4-3} and Lemma~\ref{lem:c5-5-2}(a)--(c) used along the path, namely the selected-integral representation and its finiteness, the activity bound $| k|\leqslant m$ and the hitting-sum bound of Corollary~\ref{cor:c5-4-3}, and the majorants of Lemma~\ref{lem:c5-5-2}(a)--(c), hold at exterior parameter zero at the background $U_{W,t}$ for every $t\in[0,1]$ with the same constants (the normaliser statements, $K^{\Lambda}>0$ and $I^{\Lambda}>0$ of Lemma~\ref{lem:c5-4-0}(c), the identity of Lemma~\ref{lem:c5-4-1} for $\Xi^{f,t}_{\Lambda}$, the ratio identity of Corollary~\ref{cor:c5-4-3} and part (d) of Lemma~\ref{lem:c5-5-2}, which divide by $\Xi^{f,t}_{\Lambda}$ or concern the physical kernel, are not used along the path): their proofs use the background only through Lemma~\ref{lem:c5-4-0}(a) (the Hessian bounds and range, supplied for $H_t$ by the preceding paragraph), (c) (the row charts and primitive bounds, supplied along the path by Section~\ref{sec:c4-10-1} at these rows) and (d) (the hosts). The dummy comparison \eqref{eq:c4-4-7} at $D(\Lambda)$ is Lemma~\ref{lem:c5-5-2}(c) at exterior parameter zero, and it holds with the ordinary potentials $\zeta\mathbf W_0$, $|\zeta|\leqslant R_{\Upsilon}$, with the supremum over $\zeta$ inside the sum over the selections, as \eqref{eq:c4-4-7} requires (each term is bounded by a majorant independent of $\zeta$, Lemma~\ref{lem:c5-5-2}(c)): the proof of Lemma~\ref{lem:c5-5-2}(c), and of the majorant (b) on which it rests, uses the potentials only through the host budget $|\zeta\mathbf W_0(Y,B)|\leqslant A_Y=(\omega_*/2)e^{-\kappa d_0(Y)}$ (as stated in (c)), and at exterior parameter zero $|\mathbf W_0(Y,B)|\leqslant C_0g_0(1+r)^3e^{-\kappa d_0(Y)}$ at $U_1$ for every $\Lambda$ (Lemma~\ref{lem:c5-4-0}(d), whose potentials $\mathbf V^{\Lambda}_{\xi,0}(Y,\bigcdot )$ are the $\mathbf W_0(Y,\bigcdot )$ of Lemma~\ref{lem:c5-4-1}: the input is $a$ alone), while along the path the Taylor-subtracted potentials of Section~\ref{sec:c4-9-1} have the host bound $C^{\partial}_0g_0(1+r)^3e^{-\kappa d_0(Y)}$ (Section~\ref{sec:c4-9-1}); in both cases $R_{\Upsilon} C_0g_0(1+r)^3\leqslant R_{\Upsilon} C^{\partial}_0g_0(1+r)^3=\omega_*/2$, since $R_{\Upsilon}=\omega_*/(2C^{\partial}_0g_0(1+r)^3)$ with $C^{\partial}_0\geqslant C_0$ (\eqref{eq:c4-7-1} and the sentence following it). The positive host comparison of Lemma~\ref{lem:c4-7-1}(a) is a numerical inequality free of the carrier, and Lemma~\ref{lem:c4-7-1}(b),(c) is the same comparison with the differentiated host bound of Section~\ref{sec:c4-10-1} and the moment payment above, with at most $H_MN_1(D)$ active hosts (a count over templates).

Parity and the source. At exterior parameter zero the cutoffs of $K^{\Lambda}_{\xi;0}$ are the $\chi_b$, even (the translated $\chi^t_b$ enter for positive exterior parameter only), the exterior variable is zero, and the density \eqref{eq:fixed-kernel-density} for $H_t$ is even in $(B,X)$ with its determinant and compensator; the selected ordinary potentials carry the parameter $\zeta$ (Section~\ref{sec:c4-4-1}). So Lemmas~\ref{lem:c4-7-2} and~\ref{lem:c4-4-2} hold at $D(\Lambda)$ at exterior parameter zero, at $U_1$ and along the path, $|\mathfrak Q_{\Lambda}(\mathcal A;f)|\leqslant \prod_{D\in\mathcal A}\rho\,m(D)$ at either background with $\rho$ of \eqref{eq:c4-4-8}: the source remainder bounds of Section~\ref{sec:c4-10-4} hold on $A^{\mathrm{fin}}_{\mathrm{cov}}$, where every contour row has its free inputs in $D(\Lambda)$ and the exterior variable does not enter a link (Lemma~\ref{lem:c5-5-7z}), with the row charts of Section~\ref{sec:c4-10-1} along the path (preceding paragraph). Along the path the quantities used are the selected integrals and their majorants, Corollary~\ref{cor:c5-4-3} and Lemma~\ref{lem:c5-5-2}(b),(c) at the background $U_{W,t}$, together with \eqref{eq:c4-7-9} for $m(D)$ and Lemma~\ref{lem:c4-7-1}(c), which hold at $D(\Lambda)$ by the preceding paragraph; the joint ratio $q^t_{\Lambda}(\mathcal A)$ of Lemma~\ref{lem:c5-5-2}(d) and Lemma~\ref{lem:c4-4-4} are not invoked along the path. This gives \eqref{eq:c4-12-1}--\eqref{eq:c4-12-2} at $D(\Lambda)$: $| c^0_W(h)|\leqslant M^{(\rho)}_{\Lambda}(W)$, $| a^{\Lambda}_W|\leqslant M^{(\rho)}_{\Lambda}(W)$ and $|\mathscr D_{\omega} c_W(U_{W,t};h)|\leqslant C_{\chi}\|\omega\|_{\infty}(n(W)+|\mathcal H(W)|)M^{(\rho)}_{\Lambda}(W)$, where $M^{(\rho)}_{\Lambda}(W)$ is the majorant $M_{\Lambda}(W)$ of Lemma~\ref{lem:c5-5-7}(d) with root weight $\rho\,m(Z)$ in place of $\rho_tm(Z)$. The majorant holds for $| c_W(U_{W,t};h)|$ along the path and for $| c^0_W(h)|$ and $| a^{\Lambda}_W|$ at the endpoints: at the physical background $U_1$ by the same argument (the parity argument applies verbatim at $U_1$, the cutoffs being the even $\chi_b$ and the density \eqref{eq:fixed-kernel-density} for $\mathsf H_{D(\Lambda)D(\Lambda)}$ being even, Lemmas~\ref{lem:c5-4-1} and~\ref{lem:c5-5-2}(a), and the source remainders at $U_1$ come from the $U_1$ charts, Lemma~\ref{lem:c5-4-0}(c) and Section~\ref{sec:c4-4-1}), and at the reference data (Lemma~\ref{lem:c5-5-3}(b)); on $A^{\mathrm{fin}}_{\mathrm{cov}}$ the source is the carrier-$I$ source with zero on $I\setminus D(\Lambda)$. The endpoint majorants are what \eqref{eq:c4-12-4}, and the failed-phase-test case of \eqref{eq:c4-12-3}, use for the groups without a path.

Locality, covariance and endpoints. The gauge action $R_g$ preserves each bond subspace and commutes with geometric restrictions, so the principal restriction to $D(\Lambda)$ is conjugated bondwise, $R_g$ acting on each bond subspace separately (Section~\ref{sec:c4-11-2}); the support statements of Section~\ref{sec:c4-10-2} (the seed inverse on $E_0$, the walks inside their domains, the stencils in $Q_W$) concern $H_t$ and the partitions. Hence \eqref{eq:c4-11-5} and the endpoint identities \eqref{eq:c4-11-6} hold at $D(\Lambda)$: $c_W(U_{W,1};h)= c^0_W(h)$, and $c_W(U_{W,0};h)= a^{\Lambda}_W\chi(h)$, the $t=0$ recipe being the reference data of Lemma~\ref{lem:c5-5-3}(b) (the same $D(\Lambda)$, $\Omega_1$, $\mathbf D_0(\Lambda)$ and cutoffs) with its Schur identity; the coefficient path \eqref{eq:c4-9-2} (Lemma~\ref{lem:c4-9-1}) follows by dominated convergence with the majorants above; at $t=1$ the Taylor-subtracted primitives of the path are the near primitives relative to $D(\Lambda)$ of Section~\ref{sec:c4-9-1}, which agree primitive by primitive with the compensated primitives of the retreated split (the locality statement of Section~\ref{sec:c4-9-1}, free of the carrier, with the near set selected relative to $D(\Lambda)$ by Lemma~\ref{lem:c5-2-5}). The small interior groups \eqref{eq:c4-8-9}, the rectangle $Q_W$, its injective lift (Section~\ref{sec:c4-11-4}) and, on $\mathsf N_b^c$, the inclusion of $Q_W$ in $\Lambda_1$ outside $S_1$ (Section~\ref{sec:c4-8-4}) concern the group's domains, the contour and $\Lambda_1$, not the carrier (the groups at $D(\Lambda)$ have the footprints and hulls of their domains, Lemma~\ref{lem:c5-5-8}); the plaquette estimate $\delta_{Q_W}(U_1)\leqslant C\varepsilon_1$ (Section~\ref{sec:c4-11-5}) is a statement about $U_1$ near $I\supseteq D(\Lambda)$. So \eqref{eq:c4-12-3} holds for $\mathfrak d^0_W$ of a small interior group, whether or not it passes the phase test, and \eqref{eq:c4-12-4} for the others.

Pins. The first display of Lemma~\ref{lem:c5-5-8}(a) is the positive summation of Lemma~\ref{lem:c4-6-2} at $D(\Lambda)$ with root weight $\rho_t$; it holds with the bound $C_{\mathrm{pin}}\rho'$ for every root weight $\rho'$ satisfying the smallness conditions of Lemma~\ref{lem:c4-6-2} (which states its bound in the form ``$\leqslant C\rho'$ for every $\rho'$ below the threshold'', the conditions being monotone in the root weight), in particular for $\rho'=\rho$ under the thresholds of Lemma~\ref{lem:c4-6-2} assumed through Proposition~\ref{prop:c4-12-3}. With $e^{\theta^{\circ}(W)}\leqslant e^{\eta\mathfrak s(W)/8}$ the computation \eqref{eq:c4-12-5}--\eqref{eq:c4-12-6} is verbatim: $\sum_WK_{\chi} e^{\theta^{\circ}(W)}\mathfrak d^0_W\leqslant C_{\chi}(\varepsilon_1+e^{-2\eta T_b})\sum_WM^{(\rho)}_{\Lambda}(W)e^{4\eta\mathfrak s(W)}\leqslant C_{\chi} n_C\rho(\varepsilon_1+e^{-cb/(LM)})$, the contour pin summed over the $n_C$ traversals (every group has a nonempty interval). This is \eqref{eq:c5-6-def-3}, and with (i)--(ii) it gives \eqref{eq:c5-6-def-1}.

(b) By (a) and Lemma~\ref{lem:c5-6-1}, on each labelled term the integrand of $E^{\mathrm{lab}}_{\mathrm{def}}$ integrates over $a_{\mathrm{old}}|_{I\setminus D(\Lambda)}$ to $\mathsf w^{\mathrm{lab}}_{\xi}\mathbf 1_{A^{\mathrm{fin}}_{\mathrm{cov}}}\mathscr R^{\Lambda}$ ($A^{\mathrm{fin}}_{\mathrm{cov}}$ depends on the labels through $\Lambda$ only). Split by $\mathsf N_b$, a history event (\eqref{eq:c4-2-1}; $\Lambda_1$ is a function of the history, not of the labels). On $\mathsf N_b$ use $|\mathscr R^{\Lambda}|\leqslant 2$ and $\sum_{\mathrm{lab}}\mathsf w^{\mathrm{lab}}_{\xi}=\mathsf w_{\xi}$ (\eqref{eq:c5-6-1}): the contribution is at most $2\mathcal Z_0^{-1}\int\mathcal O[\mathsf w_{\xi}\mathbf 1_{\mathsf N_b}]\,\dd V=2\mathbb P_{\mathrm h}(\mathsf N_b)$, the probability of $\mathsf N_b$ under the original history law \eqref{eq:c3-6-1}, in which \eqref{eq:c4-2-2} is proved (the labels split the weight $\mathsf w_{\xi}$ without changing the law, and $\mathsf N_b$ is defined by the original $\Lambda_1$), so it is at most $8C_Ln_C(2b+3)^4q_{\mathrm{pt}}\leqslant C'_Ln_Ce^{-(c_*/2)\pi_*^2}$ by \eqref{eq:c4-2-2}. On $\mathsf N_b^c$ use \eqref{eq:c5-6-def-1}; this is the middle expression of \eqref{eq:c5-6-def-2}, and $| S^{\circ}_{\Lambda}|\leqslant 1$, $\sum_{\mathrm{lab}}\mathsf w^{\mathrm{lab}}_{\xi}\mathbf 1_{A^{\mathrm{fin}}_{\mathrm{cov}}\cap\mathsf N_b^c}\leqslant \mathsf w_{\xi}$ and $\mathcal Z_0^{-1}\int\mathcal O[\mathsf w_{\xi}]\,\dd V=1$ (\eqref{eq:c3-6-1} at $X=1$) give the last bound. Not used: the inclusion $D_-\subseteq D(\Lambda)$ (Remark~\ref{rem:c5-7-5}(3); $S^{\circ}_{\Lambda}$ is not identified here) and the locality of Lemma~\ref{lem:c5-6-2} (every object is taken at $D(\Lambda)$ throughout).
\end{proof}

\begin{remark}[scope of this subsection]\label{rem:c5-6-5}
The locality of the species at $t=0$ (Lemma~\ref{lem:c5-6-2}) means that every activity, probe measure, group weight and reference activity of a labelled term at $t=0$, and on $A^{\mathrm{fin}}_{\mathrm{cov}}$ every defect $R^0_{\Lambda}(\mathfrak B)$, coincides with that of Section~\ref{sec:sourced-effective-interaction} unless its domains meet the de-younged bonds, so the extra resummation index touches only the species whose domains contain a cube of $\mathcal N_{\mathrm{lab}}$ (such domains may extend far from $I\setminus D(\Lambda)$; what is local is the marking, not the domain); the $t>0$ species, which enter $\mathcal B^{\Lambda}$ only, depend on $m$ and the exterior data and are bounded, not identified, by Section~\ref{sec:c5-5}. Every term of \eqref{eq:c5-6-3} is a one-point witness probability times a lattice sum, uniform in the volume. Not contained here: (1) Lemma~\ref{lem:c5-6-3} gives $S^{\circ}_{\Lambda}= S^{\circ}_{\xi,C}e^{\eta_{\mathrm{lab}}}$ with $\eta_{\mathrm{lab}}$ marked by the de-younged bonds, but $| I\setminus D(\Lambda)|$ counts the round-0 bonds $D\setminus D^{(0)}$ (Lemma~\ref{lem:c5-3-0}; zero when $\Lambda_1$ is filled in the side-$100\ell$ sense) and up to $N_{\mathrm{ret}}$ bonds per failed collar test, numbers controlled in probability only ($p_z^{(0)}\leqslant q$, $p_e^{\vartheta}\leqslant q_{\vartheta}$ per bond), so $e^{\eta_{\mathrm{lab}}}-1$ is not small on every history, and averaging it over the labels would need a bound on $\eta_{\mathrm{lab}}$ in terms of $| I\setminus D(\Lambda)|$, which is not established, together with exponential moments of $| I\setminus D(\Lambda)|$ under the history measure, which the joint witness bound of Lemma~\ref{lem:c5-7-joint} would supply (Sections~\ref{sec:c5-2}--\ref{sec:c5-3} themselves give one-point probabilities only); this route is not pursued, and the leading coefficient is instead identified at $\Lambda_1^{\mathrm{fin}}$ directly rather than compared to the carrier $I$: Lemma~\ref{lem:c5-7-1} identifies $S^{\circ}_{\Lambda}$ with the regional square scalar at $D(\Lambda)$ and Corollary~\ref{cor:c5-7-2} gives $S^{\circ}_{\Lambda}=\kappa_Ce^{E^{\circ}_{\Lambda}}$ with a label-dependent $E^{\circ}_{\Lambda}$, so that $e^{E^{\circ}_{\Lambda}}$ stays inside $\sum_{\mathrm{lab}}\mathsf w^{\mathrm{lab}}_{\xi} S^{\circ}_{\Lambda}$ as a label-dependent factor (the form of \eqref{eq:retained-leading-source}, Remark~\ref{rem:c5-7-5}(5)), whose label average is bounded in first moment only (Corollary~\ref{cor:c5-7-4}(b)). (2) $E^{\mathrm{lab}}_{\mathrm{def}}$ is bounded in Proposition~\ref{prop:c5-6-def}, the reading of Proposition~\ref{prop:c4-12-3} (\eqref{eq:c4-12-7}, \eqref{eq:c4-12-8}) at the carrier $D(\Lambda)$ ($I\to D(\Lambda)$, $S^{\circ}\to S^{\circ}_{\Lambda}$, $R_{\xi}\to R^0_{\Lambda}$); that reading takes the inputs of Sections~\ref{sec:c4-1}--\ref{sec:c4-12} at their stated status and does not re-verify them (as in (4)), and it does not use the locality of Lemma~\ref{lem:c5-6-2}, whose content is that the species of the two carriers coincide away from $\mathcal N_{\mathrm{lab}}$. (3) The marked-polymer induction of Section~\ref{sec:marginal-perimeter} (the pinned recursion \eqref{eq:marginal-pinned-recursion}) is unchanged in form; its species are the groups at the carrier of the term, and Lemma~\ref{lem:c5-6-2} is what makes the induction's inputs label-independent away from the de-younged bonds; writing the label index into that induction is not done here. (4) The inputs of Section~\ref{sec:c5-5-0} (the pin sums of Lemma~\ref{lem:c4-6-2} for every $\rho\leqslant \rho_{\mathrm{pin}}$, the host comparison of Section~\ref{sec:c4-7}, the reference-gas expansion of Lemma~\ref{lem:c5-5-3}(d), the tree-graph inequality and the square geometry) enter this subsection through Lemmas~\ref{lem:c5-5-3} and~\ref{lem:c5-5-8} and Propositions~\ref{prop:c5-5-9} and~\ref{prop:c5-5R-4}; they are used as stated there and not re-verified here, and a failure of any of them would propagate to \eqref{eq:c5-6-2}--\eqref{eq:c5-6-4} and Lemma~\ref{lem:c5-6-3}.
\end{remark}

\subsection{Identification of the labelled leading coefficient}
\label{sec:c5-7}

The labelled leading coefficient $S^{\circ}_{\Lambda}$ of \eqref{eq:c5-6-2} is identified with $\kappa_Ce^{E^{\circ}_{\Lambda}}$, $\kappa_C\coloneqq e^{-\delta mP(C)}z_c^4$ (the $e^{-\delta m^{\sigma} P(C_n)}(z^{\sigma}_c(g_0,\pi/2))^4$ of \eqref{eq:extracted-loop-scalar} in the cutoff prescription $\sigma$ fixed throughout), on $A^{\mathrm{fin}}_{\mathrm{cov}}$, with a bounded $E^{\circ}_{\Lambda}$, and the bound \eqref{eq:c5-6-3} (in the form of Proposition~\ref{prop:c5-5R-4}) and the leading term of \eqref{eq:c5-6-1}--\eqref{eq:c5-6-2} are restated relative to $\kappa_C$ (Corollaries~\ref{cor:c5-7-4}(a) and~\ref{cor:c5-7-2}(b)). The objects: $S^{\circ}_{\Lambda}=\sum_{\mathfrak B\text{ compatible}}\prod_{W\in\mathfrak B}a^{\Lambda}_W=\int f^{\circ} K^{\circ}_{\Lambda}\,\dd a/I^{\circ}_{\Lambda}$ (Lemma~\ref{lem:c5-5-3}(b); proof of Lemma~\ref{lem:c5-6-3}), the reference data at $D(\Lambda)$ being $H^{\circ}=(\mathsf H^{\Omega_1}_{\id})_{D(\Lambda)D(\Lambda)}$, exterior variable zero on $J(\Lambda)$, $m^{\circ}=\ell^{\circ}=0$, $f^{\circ}(a)= W_{\chi}(C;e^{iq_{\id}(a)})$, so that, by the definition of $K^{\Lambda}_{\xi;t}$ (Section~\ref{sec:c5-4}) at $t=0$ and these data,
\begin{gather}
    K^{\circ}_{\Lambda}(a)=\mathbf 1_{\mathcal D^{\Lambda}}(a)\exp\Bigl[-\tfrac12\langle a,H^{\circ} a\rangle+\sum_{\mathfrak p\text{ near rel. }D(\Lambda)}\mathfrak p(a)\Bigr]\;,\quad a\in\mathfrak g^{D(\Lambda)}\text{ (exterior variable zero)},\label{eq:c5-7-1}\\
    I^{\circ}_{\Lambda}=\int K^{\circ}_{\Lambda}\,\dd a\;,\qquad S^{\circ}_{\Lambda}=\frac{\int f^{\circ} K^{\circ}_{\Lambda}\,\dd a}{I^{\circ}_{\Lambda}}\;,\notag
\end{gather}
the near primitives being those of all three kinds (action, Haar, constraint Jacobian) whose free-input set meets $D(\Lambda)$. Inputs, used at their stated status: (I1) the regional square scalar \eqref{eq:regional-square-scalar}, $S^{\circ}_{D,C_n,T}=\int\Theta^{\sigma}_D(B)e^{P_{D,T}(B)}\chi(\operatorname{Hol}_{U_{g_0,T}(i_{D,T}B)}C_n)\,\dd\mu_{C_{D,T}}(B)/\int\Theta^{\sigma}_D(B)e^{P_{D,T}(B)}\,\dd\mu_{C_{D,T}}(B)$, $C_{D,T}=(i^*_{D,T}H_Ti_{D,T})^{-1}$, defined there for the free bonds $D$ integrated in a covered conditional kernel, i.e. every free-input stencil of the contour lies in $D$ (the free-input stencil of a traversal link $r$ being the coverage stencil $\Lambda_r$ of Section~\ref{sec:remaining-physical-comparison}, which is the coverage of Proposition~\ref{prop:regional-square-extraction}), with the original primitive labels retained including those crossing the boundary of $D$; the proposition's closing sentence states that its estimates use only the covered stencils, the spectral bounds and the small-coupling conditions of the defined flat integral, so that they apply to every covered $D\subseteq D_+$ with the collar, in particular to $D(\Lambda)$ (Lemma~\ref{lem:c5-7-1}); (I2) Proposition~\ref{prop:regional-square-extraction}: under coverage, $K_{\mu}\geqslant 3(n+2)$, the collar conditions \eqref{eq:regional-collar-sizes} ($M\geqslant 8LM_2$, $LMR_1>4R_0+4L^2M_2R_1+8$, $R_0=4(1+R_*+R_H+L)$, enforceable by choosing $M$ with $LM-4L^2M_2>4R_0+8$ before the coupling) and the common small-coupling thresholds of the flat and regional expansions, $S^{\circ}_{D,C_n,T}=\kappa_Ce^{E^{\circ}_{D,C_n,T}}$ with $| E^{\circ}_{D,C_n,T}|\leqslant C\tau^2[\sum_{j=0}^{4n-1}e^{-cd_j}+4ne^{-c\iota_T}+(n+1)e^{-cn}]$ (\eqref{eq:regional-square-error}), $d_j=\max(0,d_T(K_j,D^c)-R_0)$, $\iota_T=\max(0,N_{\min}/4-R_0)$, $K_j$ the union of the free-input stencils of the fine traversals of coarse segment $j$, complements in the torus free carrier, $d_T$ the periodic fine-bond metric ($d_T(b,c)= d_{\infty,T}(x_b,x_c)+\mathbf 1\{\text{directions differ}\}$, Section~\ref{sec:marginal-perimeter}, after \eqref{eq:marginal-covariance-input}), $C$, $c$ independent of $D$, $n$ and the periods; (I3) the identification paragraphs of Section~\ref{sec:regional-identification} between \eqref{eq:regional-torus-bound} and \eqref{eq:regional-kernel-identification}: the comparison holds for every $D_-\subseteq D\subseteq D_+$ ($D_{\pm}$ the free bonds with both/at least one endpoint in $\Lambda_1$), assigning the complementary free variables to the exterior integral, with the following facts at identity data, numbered here for reference: (1) every stencil meeting $D\subseteq D_+$ stays in the regular region and the Hessian on $D$ is $H_{D,T}= i^*_{D,T}H_Ti_{D,T}$ including the diagonal terms of crossing plaquettes; (2) the constraint linearisation gives the flat phase \eqref{eq:flat-phase}, so that the reconstructed links at identity data are $e^{iq_{\id}}$, and the action, Haar and Jacobian terms become $v_T(i_{D,T}B)$ with all primitive labels meeting $D$ retained; (3) the boundary change CRE (1.22) is the identity at identity data, with Jacobian one and zero additional exponent; (4) at zero exterior free field the conditional cutoff is $\Theta^{\sigma}_D$ with $r=\delta_0/g_0$; (5) the exterior quadratic cross term and its compensator vanish at the projection. Each of (1)--(5) is argued there for $D\subseteq D_+$ with the collar \eqref{eq:regional-collar-sizes}: the arguments use that every bond of $D$ has an endpoint in $\Lambda_1$, so that it is tested by a good test cube inside $\Lambda_1$ and its stencils stay in the regular region, and that bad tests and changed-background tests are separated from $D_+$ by the collar; the lower inclusion $D_-\subseteq D$ is not used in them. The primitives of the conditional reference scalar are those of \eqref{eq:regional-periodic-weight}, $P_{D,T}(B)=\sum_{v_T\, : \,\operatorname{supp}v_T\cap D\neq\varnothing}v_T(i_{D,T}B)$, the terms whose free support misses $D$ vanishing at $i_{D,T}B$; (I4) Lemma~\ref{lem:c5-5-3}(b) (the reference data at $D(\Lambda)$), $A^{\mathrm{fin}}_{\mathrm{cov}}$ of Section~\ref{sec:c5-3-4}(a), used in Lemma~\ref{lem:c5-6-1} (every contour bond in $\Omega_1$ with its free-input set $\mathcal F_r\subseteq D(\Lambda)$), Section~\ref{sec:c5-3-4}(a) (a free bond not in $D(\Lambda)$ has no endpoint in $\Lambda_1^{\mathrm{fin}}$, so lies in $E(B^{\mathrm{fin}}_{\xi})$, $B^{\mathrm{fin}}_{\xi}\coloneqq (\Lambda_1^{\mathrm{fin}})^c$), Proposition~\ref{prop:c5-6-4} and Section~\ref{sec:c5-5R}. Thresholds, in addition to items 1--5 of Appendix~\ref{app:thresholds} (with the collar item 4), \eqref{eq:c5-4-T}, \eqref{eq:c5-4-Tp}, \eqref{eq:c5-5-T}, \eqref{eq:c5-5R-T} and the period hypothesis of Section~\ref{sec:c5-3-1}:
\begin{equation}\label{eq:c5-7-T}
    \begin{aligned}
        &\text{(i) the collar conditions \eqref{eq:regional-collar-sizes}: $M\geqslant 8LM_2$ (already in the order of choice of}\\
        &\quad\text{Appendix~\ref{app:thresholds}) and $LMR_1>4R_0+4L^2M_2R_1+8$ with $R_0=4(1+R_*+R_H+L)$}\\
        &\quad\text{(implied by $LM-4L^2M_2>4R_0+8$ since $R_1\geqslant 1$; threshold item 1 of Appendix~\ref{app:thresholds}}\\
        &\quad\text{is the same condition with $R_{\mathrm c}$ in place of $R_0$; both are imposed), with aligned nested partitions and}\\
        &\quad\text{all periods divisible by $LMR_1$ (\eqref{eq:c-standing-conditions}) and the contour a coarse translate}\\
        &\quad\text{of the counterclockwise $(e_1,e_2)$ square of side $nL$ (Section~\ref{sec:c5-5-0});}\\
        &\text{(ii) the common small-coupling thresholds of the flat and regional expansions}\\
        &\quad\text{($g_0C_{\mathsf C}(2+4r)<\vartheta_*/2$, $g_0(2+4r)< a_L$, and the polymer-load inequalities of the proof of}\\
        &\quad\text{Proposition~\ref{prop:regional-square-extraction}), conditions on $g_0$ alone;}\\
        &\text{(iii) $K_{\mu}\geqslant 3(n+2)$ for every coarse period (Section~\ref{sec:c5-5-0}), and, for the $n$-uniform}\\
        &\quad\text{form of Corollary~\ref{cor:c5-7-2}(c) only, $4ne^{-c\iota_T}\leqslant 1$ (a condition on the torus:}\\
        &\quad\text{$N_{\min}\geqslant 4R_0+(4/c)\log(4n)$).}
    \end{aligned}
\end{equation}

\subsubsection{The identification at the retreated carrier}

\begin{lemma}\label{lem:c5-7-1}
Under \eqref{eq:c5-7-T}, on every labelled term with final region $\Lambda$, for exterior data in the support of the term's collar indicators and exterior cutoffs, on $A^{\mathrm{fin}}_{\mathrm{cov}}$: $S^{\circ}_{\Lambda}= S^{\circ}_{D(\Lambda),C,T}$, the regional square scalar (I1) at $D= D(\Lambda)$.
\end{lemma}

\begin{proof}
Both sides are ratios of two integrals over $\mathfrak g^{D(\Lambda)}$ with the same denominator structure; it suffices to identify the four factors of the integrands, numerator and denominator alike, the Gaussian normaliser of $\mu_{C_{D(\Lambda),T}}$ being common to numerator and denominator on the side of (I1) and absent on both on the side of \eqref{eq:c5-7-1}. First, $D(\Lambda)$ qualifies: $D(\Lambda)\subseteq I= D_+(\Lambda_1)$ (the carrier of Section~\ref{sec:sourced-effective-interaction} consists of all free bonds with an endpoint in $\Lambda_1$, the bond convention of Section~\ref{sec:c5-0}; the retreat only removes bonds), so every stencil meeting $D(\Lambda)$ stays in the regular region and every bond of $D(\Lambda)$ has an endpoint in a good cube of $\Lambda_1$, which is all the facts (I3) use (the lower inclusion $D_-\subseteq D$ of the comparison statement is not used in the arguments for the facts (1)--(5), which are made for $D\subseteq D_+$ with the collar, (I3)); the de-younged bonds $I\setminus D(\Lambda)$ are complementary free variables assigned to the exterior integral, which is what the zero exterior variable on $J(\Lambda)$ of Lemma~\ref{lem:c5-5-3}(b) does; and $D(\Lambda)$ is covered on $A^{\mathrm{fin}}_{\mathrm{cov}}$, since every contour bond lies in $\Omega_1$ with $\mathcal F_r\subseteq D(\Lambda)$, and $\mathcal F_r$ contains the coverage stencil $\Lambda_r$ of Section~\ref{sec:remaining-physical-comparison}: $\Lambda_b=\{ b\}=\mathcal F_b$ for a free bond, and for a designated bond $\Lambda_{b_0(c)}=\mathcal S_c\cap E_*$ with $\mathcal S_c$ the positive fine bonds with both endpoints in the two endpoint blocks $\mathcal B_{c_-}\cup\mathcal B_{c_+}$ (Appendix~\ref{sec:perimeter-attempt}), while $\mathcal F_{b_0(c)}=\mathcal E_c\cap E_*$ (proof of Lemma~\ref{lem:covered}) with $\mathcal E_c=\{ r\, : \,\{ r_-,r_+\}\subset\mathcal B_{c_-}\cup\mathcal B_{c_+}\}$ (Section~\ref{sec:regular-hessian-stencils}), so $\mathcal S_c\subseteq\mathcal E_c$ and $\Lambda_r\subseteq\mathcal F_r$; what (I2) needs is $\Lambda_r\subseteq D(\Lambda)$ (Remark~\ref{rem:c5-7-5}(2)). (F1) Precision: $H^{\circ}=(\mathsf H^{\Omega_1}_{\id})_{D(\Lambda)D(\Lambda)}$ is the Hessian at identity of the regular-region action (CRE's constrained quadratic form, $\mathsf H^{\Omega_1}_{U_1}=\mathsf C^*\Delta_1\mathsf C$, Section~\ref{sec:regular-hessian-stencils}, as used in the proof of Lemma~\ref{lem:primitive-sum}) restricted to $D(\Lambda)\times D(\Lambda)$; the proof of Lemma~\ref{lem:c2-reference} states, by Lemma~\ref{lem:c2-regional-entry} at the identity background, that the entries of $H^{\circ}$ with a label in $I$ are $f_{bb'}(0)$, the entries of the flat Hessian of \eqref{eq:flat-covariance} on the branch's torus, so the principal restriction to $D(\Lambda)\subseteq I$ is the flat torus precision on $D(\Lambda)$ including the diagonal contributions of the plaquettes crossing the boundary of $D(\Lambda)$, the same statement as (I3) fact 1 at $D= D(\Lambda)$, which holds because every stencil meeting $D(\Lambda)$ lies in the regular region, where the action at identity is the flat quadratic form; the flat Hessian of Appendix~\ref{app:background-estimates} and $H_T$ are the same operator \eqref{eq:flat-covariance} on the same torus. So $\exp(-\frac12\langle a,H^{\circ} a\rangle)\,\dd a$ and $\dd\mu_{C_{D(\Lambda),T}}(B)$ agree up to the common normaliser, with $a= B$. (F2) Cutoff: at the reference data the cutoff of $K^{\circ}_{\Lambda}$ is the product of the bond tests $\prod_{b\in D(\Lambda)}\mathbf 1\{| a_b|_{\sigma}< r\}$ by the reference-data convention (Lemma~\ref{lem:c5-5-3}(b), the same cutoffs, the reference data not being a branch; Lemma~\ref{lem:c5-4-0}(c) for these data, $K^{\circ}_{\Lambda}>0$ on the open box, proof of Lemma~\ref{lem:c5-6-3}; the $K^{\circ}_{\xi;0}=\mathbf 1_{\mathcal D_I}\exp(\dots)$ of Section~\ref{sec:joint-extraction}, positive on the open box at the carrier $I$ by Lemma~\ref{lem:c2-reference}); on a branch the corresponding statement is Lemma~\ref{lem:c5-3-3}(iv), which is not what is used here. $\Theta^{\sigma}_D$ is the same product of strict bond tests at the same threshold $r=\delta_0/g_0$ in the same norm $|\bigcdot |_{\sigma}$ ((I3) fact (4)). (F3) Primitives: $\sum_{\mathfrak p\text{ near rel. }D(\Lambda)}\mathfrak p(a)$ at zero exterior variable is the sum of the local primitives of Section~\ref{sec:source-phase-primitives} at identity background whose free-input set meets $D(\Lambda)$, evaluated at $a= i_{D(\Lambda),T}B$; $P_{D(\Lambda),T}(i_{D(\Lambda),T}B)$ is the sum of the periodic primitives $v_T$ of \eqref{eq:regional-periodic-weight} (those of \eqref{eq:perimeter-literal-local-terms} on $T$) whose supports meet $D(\Lambda)$ (crossing labels retained), the others vanishing at $i_{D(\Lambda),T}B$. Both sums are labelled by the plaquettes, gauge centres, bonds and rows of the regular region, and they agree once the primitives agree label by label, the gauge terms as the sum over a centre: the primitives of Section~\ref{sec:source-phase-primitives} are, by Lemma~\ref{lem:primitive-sum}, the degree-two Taylor remainders of the action and gauge-fixing terms, the Haar logarithms and the row Jacobian logarithms at the phase $q_{\id}(a)$, and Lemma~\ref{lem:identity-primitives} gives this agreement label by label (energy residuals plaquette by plaquette and, for the gauge terms, centre by centre as the sum over the pairs $(y,x)$; Haar and constraint-Jacobian primitives bond by bond and row by row) for every $B$ with $g_0C_{\mathsf C}\max_b| B(b)|\leqslant \vartheta_*/2$ and $g_0\max_b| B(b)|< a_L$, which the small-coupling conditions of the proof of Proposition~\ref{prop:regional-square-extraction} (part of \eqref{eq:c5-7-T}) guarantee on the integration box; its second consequence is exactly the equality of the two sums at $i_{D(\Lambda),T}B$ (Remark~\ref{rem:c5-7-5}(1)); the far primitives contribute their value at zero exterior variable, a constant common to numerator and denominator of \eqref{eq:c5-7-1}, so whether they are counted is immaterial. (F4) Wilson factor: $f^{\circ}(a)= W_{\chi}(C;e^{iq_{\id}(a)})$ is $\chi$ of the holonomy along $C$ of the links $e^{iq_{\id}(a)_r}$; by (I3) fact 2 the reconstructed links at identity data are $U_{g_0,T}(i_{D(\Lambda),T}B)= e^{iq_{\id}}$ evaluated at the same free field, and on $A^{\mathrm{fin}}_{\mathrm{cov}}$ the link of each contour bond $r$ depends on the free field through $\mathcal F_r\subseteq D(\Lambda)$ only, so the exterior variable (zero) does not enter and $f^{\circ}(a)=\chi(\operatorname{Hol}_{U_{g_0,T}(i_{D(\Lambda),T}B)}C)$. There is no fifth factor: at $t=0$ and $m^{\circ}=\ell^{\circ}=0$ the kernel $K^{\Lambda}_{\xi;0}$ has no linear term and no shift, and facts 3 and 5 of (I3) (boundary change, cross term) concern the identification of the conditional CRE scalar with $S^{\circ}_{D,C,T}$, which is the statement of Section~\ref{sec:regional-identification} itself and is not needed to compare the two explicit integrals \eqref{eq:c5-7-1} and (I1). Numerator and denominator agree factor by factor, so the ratios agree.
\end{proof}

\begin{corollary}\label{cor:c5-7-2}
Under \eqref{eq:c5-7-T} and, for (b), the hypotheses of Proposition~\ref{prop:c5-5R-4} (those of Proposition~\ref{prop:c5-5-9} and \eqref{eq:c5-5R-T}; the collar condition on the exterior data is carried by $\mathbf 1_{\mathrm{lab}}$), on every labelled term, on $A^{\mathrm{fin}}_{\mathrm{cov}}$, with $E^{\circ}_{\Lambda}\coloneqq E^{\circ}_{D(\Lambda),C,T}$ (the error of (I2) at $D= D(\Lambda)$, with the branch of the logarithm of Proposition~\ref{prop:regional-square-extraction}, continued from zero in the traversal-source polynomial \eqref{eq:regional-wilson-polynomial}) and $d_j(\Lambda)\coloneqq \max(0,d_T(K_j,D(\Lambda)^c)-R_0)$:
\begin{enumerate}[label=(\alph*)]
    \item
    \begin{equation}\label{eq:c5-7-2}
        S^{\circ}_{\Lambda}=\kappa_Ce^{E^{\circ}_{\Lambda}}\;,\qquad| E^{\circ}_{\Lambda}|\leqslant C\tau^2\Bigl[\sum_{j=0}^{4n-1}e^{-cd_j(\Lambda)}+4ne^{-c\iota_T}+(n+1)e^{-cn}\Bigr]\leqslant C\tau^2(9n+1)\;,
    \end{equation}
    $E^{\circ}_{\Lambda}$ depends on the labels through $D(\Lambda)$ only (a simple function on the finitely many labels of a history), $\kappa_C\neq0$ (finite perimeter exponent, $z_c$ the exponential of a convergent sum, Section~\ref{sec:remaining-physical-comparison}; also $S^{\circ}_{\Lambda}\neq0$ by (I2)), and $|\kappa_C| e^{-C\tau^2(9n+1)}\leqslant | S^{\circ}_{\Lambda}|\leqslant |\kappa_C| e^{C\tau^2(9n+1)}$ (so $S^{\circ}_{\Lambda}\neq0$ on $A^{\mathrm{fin}}_{\mathrm{cov}}$, consistent with $| S^{\circ}_{\Lambda}|\leqslant 1$ of Lemma~\ref{lem:c5-6-3});
    \item the leading term of \eqref{eq:c5-6-1}--\eqref{eq:c5-6-2} is $\kappa_C\sum_{\mathrm{lab}}e^{E^{\circ}_{\Lambda}}\int_{\mathfrak g^{I\setminus D(\Lambda)}}\mathbf 1_{\mathrm{lab}}e^{c^{\Lambda}_{\xi}}I^{\Lambda}_1\mathbf 1_{A^{\mathrm{fin}}_{\mathrm{cov}}}W_{\chi}(C;U_1)\,\dd a_{\mathrm{old}}$ ($e^{E^{\circ}_{\Lambda}}$ and $\mathbf 1_{A^{\mathrm{fin}}_{\mathrm{cov}}}$ depend on the labels only and leave the $a_{\mathrm{old}}$-integral), and on $A^{\mathrm{fin}}_{\mathrm{cov}}$, by \eqref{eq:c5-5-9R},
    \begin{gather}
        |\kappa_C|^{-1}|\mathcal B^{\Lambda}_{\xi}[f_{\xi}]|\leqslant e^{| E^{\circ}_{\Lambda}|}e^{n_C\bar\varepsilon_{\mathrm{ph}}}\bigl[C_1g_0r\mathfrak S_1(\Lambda)+C_2\bar\varepsilon_tr^{d_G}\mathfrak S_2(\Lambda)\bigr]\label{eq:c5-7-3}\\
        \leqslant e^{C\tau^2(9n+1)}e^{n_C\bar\varepsilon_{\mathrm{ph}}}\bigl[C_1g_0r\mathfrak S_1(\Lambda)+C_2\bar\varepsilon_tr^{d_G}\mathfrak S_2(\Lambda)\bigr]\;;\notag
    \end{gather}
    \item (refined split) $| E^{\circ}_{\Lambda}|\leqslant C\tau^2(2+1/c)+C\tau^2\sum_je^{-cd_j(\Lambda)}$ under $4ne^{-c\iota_T}\leqslant 1$ (\eqref{eq:c5-7-T}(iii); for $n\geqslant 1$ this is $N_{\min}\geqslant 4R_0+(4/c)\log(4n)$), since $(n+1)e^{-cn}\leqslant 1+\sup_{x\geqslant 0}xe^{-cx}=1+1/(ec)\leqslant 1+1/c$: the only part of the error that grows with $n$ at fixed $\tau$ is the boundary sum $\sum_je^{-cd_j(\Lambda)}\leqslant 4n$, each of whose terms is at most 1 and exponentially small in $d_T(K_j,D(\Lambda)^c)-R_0$ (the sum is zero only when $D(\Lambda)^c=\varnothing$, by the convention $e^{-\infty}=0$); it is the quantity averaged in Lemma~\ref{lem:c5-7-3}.
\end{enumerate}
\end{corollary}

\begin{proof}
(a) Lemma~\ref{lem:c5-7-1} and (I2) at $D= D(\Lambda)$, whose hypotheses are \eqref{eq:c5-7-T} and coverage ($A^{\mathrm{fin}}_{\mathrm{cov}}$); the crude bound uses $e^{-cd_j}\leqslant 1$, $e^{-c\iota_T}\leqslant 1$ and $(n+1)e^{-cn}\leqslant n+1$ (as in the constant $B_C$ following \eqref{eq:retained-source-correction}); $D(\Lambda)$ is a function of the labels. (b) The leading term of \eqref{eq:c5-6-2} is $S^{\circ}_{\Lambda} W_{\chi}(C;U_1)$ integrated as in \eqref{eq:c5-6-1}; substitute (a). The bound: \eqref{eq:c5-5-9R} is $|\mathcal B^{\Lambda}|\leqslant | S^{\circ}_{\Lambda}| e^{n_C\bar\varepsilon_{\mathrm{ph}}}[\dots]$ and $| S^{\circ}_{\Lambda}|=|\kappa_C| e^{\operatorname{Re}E^{\circ}_{\Lambda}}\leqslant |\kappa_C| e^{| E^{\circ}_{\Lambda}|}$. (c) The stated elementary inequality.
\end{proof}

\subsubsection{The label average of the boundary distances}

\begin{lemma}\label{lem:c5-7-3}
In the setting of Proposition~\ref{prop:c5-6-4} (square geometry of Section~\ref{sec:c5-5-0}, compatible even fine periods exceeding $200\ell$, the hypotheses of Sections~\ref{sec:c5-2}--\ref{sec:c5-3}), for every $c>0$,
\begin{gather}
    \mathcal Z_0^{-1}\int\mathcal O\Bigl[\sum_{\mathrm{lab}}\mathsf w^{\mathrm{lab}}_{\xi}\mathbf 1_{A^{\mathrm{fin}}_{\mathrm{cov}}}\sum_{j=0}^{4n-1}e^{-cd_j(\Lambda)}\Bigr]\dd V\leqslant e^{c(R_0+\varrho_K)}N_{\mathrm{lat}}(c)\,n_C(q+N_{\mathrm{ret}}q_{\vartheta})\;,\label{eq:c5-7-4}\\
    \varrho_K\coloneqq \max_j\max_{k\in K_j}\min_{r\in C_j}d_T(k,r)\leqslant 2L+1\;,\notag
\end{gather}
with $N_{\mathrm{lat}}(c)=\sum_{s\geqslant 0}4(2s+1)^4e^{-cs}$ of Lemma~\ref{lem:c5-1-2}, $C_j$ the contour bonds of coarse segment $j$, $q$ and $q_{\vartheta}$ of \eqref{eq:c5-2-3} and Lemma~\ref{lem:c5-2-4}, $N_{\mathrm{ret}}$ of Lemma~\ref{lem:c5-3-1}: uniform in the volume, and of the same form as \eqref{eq:c5-6-3}, a lattice sum times one-point witness probabilities.
\end{lemma}

\begin{proof}
Pointwise on a term: $d_j(\Lambda)\geqslant d_T(K_j,D(\Lambda)^c)-R_0$, so $e^{-cd_j(\Lambda)}\leqslant e^{cR_0}e^{-cd_T(K_j,D(\Lambda)^c)}\leqslant e^{cR_0}\sum_{b\notin D(\Lambda)}e^{-cd_T(K_j,b)}$ (the maximum over $b\in D(\Lambda)^c$ of $e^{-cd_T(K_j,b)}$ is at most the sum; if $D(\Lambda)^c=\varnothing$ the left side is zero by the convention $e^{-\infty}=0$). For $k\in K_j$ and $b$, the triangle inequality for $d_T$ gives $d_T(k,b)\geqslant \min_{r\in C_j}d_T(r,b)-\varrho_K$, hence $e^{-cd_T(K_j,b)}\leqslant e^{c\varrho_K}\max_{r\in C_j}e^{-cd_T(r,b)}\leqslant e^{c\varrho_K}\sum_{r\in C_j}e^{-cd_T(r,b)}$, and $\sum_j\sum_{r\in C_j}=\sum_{r\in C}$ because the segments partition the contour bonds. A free bond $b\notin D(\Lambda)$ (the complement is taken in the torus free carrier, \eqref{eq:regional-square-distances}) has no endpoint in $\Lambda_1^{\mathrm{fin}}$ (for $b\in I\setminus D(\Lambda)$ by Section~\ref{sec:c5-3-4}(a); for $b\notin I$ because $I= D_+(\Lambda_1)$ is the set of all free bonds with an endpoint in $\Lambda_1\supseteq\Lambda_1^{\mathrm{fin}}$, the bond convention of Section~\ref{sec:c5-0}; the equality, not only $D(\Lambda)\subseteq I$, is used here), so $\mathbf 1\{ b\notin D(\Lambda)\}\leqslant \mathbf 1\{ b\in E(B^{\mathrm{fin}}_{\xi})\}\leqslant \mathbf 1\{ b\in E(\mathfrak B_0)\}+\sum_{e'\, : \,x_{e'}\text{ within }102\ell+1\text{ of an endpoint of }b}M_{e'}$ (\eqref{eq:c5-3-mark-inequality}, Lemma~\ref{lem:c5-3-1}(c)--(d)). Altogether $\sum_je^{-cd_j(\Lambda)}\leqslant e^{c(R_0+\varrho_K)}\sum_{r\in C}\sum_be^{-cd_T(r,b)}[\mathbf 1\{ b\in E(\mathfrak B_0)\}+\sum_{e'\text{ near }b}M_{e'}]$, every term nonnegative and label-independent except through the marks. Multiply by $\mathsf w^{\mathrm{lab}}_{\xi}\mathbf 1_{A^{\mathrm{fin}}_{\mathrm{cov}}}\leqslant \mathsf w^{\mathrm{lab}}_{\xi}$, sum over the labels and histories, apply $\mathcal O$ and divide by $\mathcal Z_0$ (Tonelli; $\mathsf w^{\mathrm{lab}}_{\xi}\geqslant 0$ by \eqref{eq:c5-6-1} with $\rho_{\xi}\geqslant 0$, and $\mathcal O$ is positive; finitely many labels per history; $\sum_{\mathrm{lab}}\mathsf w^{\mathrm{lab}}_{\xi}=\mathsf w_{\xi}$, Section~\ref{sec:c5-3-4}(c)): for each fixed $b$ the first term gives $p^{(0)}_b\leqslant q$ (Lemma~\ref{lem:c5-3-0}(b)) and the second $\sum_{e'\text{ near }b}p^{\vartheta}_{e'}\leqslant N_{\mathrm{ret}}q_{\vartheta}$ (Lemma~\ref{lem:c5-3-1}(d), Proposition~\ref{prop:c5-2-6}), exactly as in the proof of Proposition~\ref{prop:c5-6-4}. Finally $\sum_be^{-cd_T(r,b)}\leqslant \sum_{s\geqslant 0}4(2s+1)^4e^{-cs}= N_{\mathrm{lat}}(c)$ for each contour bond $r$ ($d_T(r,b)\geqslant \| x_r-x_b\|_{\infty}$, and at most $4(2s+1)^4$ bonds have initial site within $\ell^{\infty}$ distance $s$ of $x_r$), and $\sum_{r\in C}1= n_C$. The bound on $\varrho_K$: for a free contour bond $r$, $\Lambda_r=\{ r\}$ and $d_T(r,r)=0$; for a designated $r= b_0(c)$, the stencil bonds and $r$ itself have both endpoints in the two endpoint blocks of $c$ ($\Lambda_{b_0(c)}=\mathcal S_c\cap E_*$ with $\mathcal S_c$ the bonds with both endpoints in the two endpoint blocks of $c$, Section~\ref{sec:remaining-physical-comparison}; $r= b_0(c)$ is a bond of the row $c$), two adjacent blocks of side $L$, any two sites of which are within $\ell^{\infty}$ distance $2L-1$ of each other, so $d_T(k,r)\leqslant 2L-1+1=2L\leqslant 2L+1$ ($d_T(b,c)= d_{\infty,T}(x_b,x_c)+\mathbf 1\{\text{directions differ}\}$, (I2)).
\end{proof}

\subsubsection{Joint witness bounds and the exponential moment of the boundary sum}

\begin{lemma}[joint witness bound]\label{lem:c5-7-joint}
In the setting of Proposition~\ref{prop:c5-2-6} and Lemma~\ref{lem:c5-3-0} (compatible even fine periods, $0< g_0< g_*(\vartheta)$), call a mark either $M_e$ for a bond $e$ (Section~\ref{sec:c5-2-3}) or $a_b\coloneqq \mathbf 1\{ b\in E(\mathfrak B_0)\}$ for a bond $b$ (Lemma~\ref{lem:c5-3-0}), located at the initial site of its bond. Put $\sigma_K\coloneqq 2K\ell+4$ and $C_K\coloneqq 4(2\sigma_K+1)^4=4(4K\ell+9)^4$, an upper bound on the number of bonds whose initial site lies within $\ell^{\infty}$ distance $\sigma_K$ of a given site. For every finite set $\mathsf S$ of marks,
\begin{equation}\label{eq:c5-7-joint}
    \mathcal Z_0^{-1}\int\mathcal O\Bigl[\sum_{\mathrm{lab}}\mathsf w^{\mathrm{lab}}_{\xi}\prod_{m\in\mathsf S}m\Bigr]\dd V\leqslant q_{\vartheta}^{\,|\mathsf S'|}\leqslant q_{\vartheta}^{\,|\mathsf S|/(2C_K)}\;,
\end{equation}
where the first inequality holds for every subset $\mathsf S'$ of $\mathsf S$ whose marks are at pairwise $\ell^{\infty}$ distance $>\sigma_K$, and the second for a maximal such subset (which has $|\mathsf S'|\geqslant |\mathsf S|/(2C_K)$: every mark of $\mathsf S$ lies within $\sigma_K$ of a mark of $\mathsf S'$, at most $C_K$ bonds lie within $\sigma_K$ of a given site, and a bond carries at most two marks; the count $C_K$ is an upper bound on every torus, so this holds also on periods at most $2\sigma_K+1$, where a separated subset has at most one mark and the second inequality carries only the rate $q_{\vartheta}^{1/(2C_K)}$). The same holds, without labels and with $q$ in place of $q_{\vartheta}$, for the marks $\mathbf 1\{ z\in E(B_{\xi})\}$ of Section~\ref{sec:c3-7} ($B_{\xi}=\Lambda_1^c$, \eqref{eq:c3-7-1}), for $g_0$ below the thresholds of Lemmas~\ref{lem:c3-7-2} and~\ref{lem:c3-7-4}, on compatible even fine periods (Section~\ref{sec:c3-7}): for every finite set $\mathsf S$ of bonds and every subset $\mathsf S'$ of $\mathsf S$ at pairwise $\ell^{\infty}$ distance $>\sigma_K$,
\begin{equation}\label{eq:c5-7-joint-I}
    \mathbb P_{\mathrm h}\bigl(\mathsf S\subseteq E(B_{\xi})\bigr)=\mathcal Z_0^{-1}\int\mathcal O\Bigl[\mathsf w_{\xi}\prod_{z\in\mathsf S}\mathbf 1\{ z\in E(B_{\xi})\}\Bigr]\dd V\leqslant q^{\,|\mathsf S'|}\;,\qquad\text{and}\qquad\mathbb P_{\mathrm h}\bigl(\mathsf S\subseteq E(B_{\xi})\bigr)\leqslant q^{\,|\mathsf S|/C_K}\;.
\end{equation}
The second inequality of \eqref{eq:c5-7-joint} uses $q_{\vartheta}\leqslant 1$: $q_{\vartheta}\leqslant e^{-c_q\vartheta^2\pi_*^2}$ for $g_0< g_*(\vartheta)$ by \eqref{eq:c5-2-3}. The second bound of \eqref{eq:c5-7-joint-I} follows from the first when $q\leqslant 1$, with $\mathsf S'$ a maximal separated subset (one mark per bond, so $|\mathsf S'|\geqslant |\mathsf S|/C_K$, and $q^{|\mathsf S'|}\leqslant q^{|\mathsf S|/C_K}$ for $q\leqslant 1$), in particular for $g_0< g_*$, the threshold of Proposition~\ref{prop:c3-7-5}, where $q\leqslant e^{-c_*\pi_*^2}$ with $c_*= c_q$ (this is \eqref{eq:c5-2-3} at $\vartheta=1$, where $\theta_1=\theta$ and $q_1= q$). When $q>1$ it holds trivially, the left side being a probability, at most $1\leqslant q^{|\mathsf S|/C_K}$ (with equality only for $\mathsf S=\varnothing$); the two bounds of \eqref{eq:c5-7-joint-I} are then not a chain ($q^{|\mathsf S'|}> q^{|\mathsf S|/C_K}$ when $|\mathsf S'|>|\mathsf S|/C_K$), which is why they are stated separately. So both bounds of \eqref{eq:c5-7-joint-I} hold for every $g_0$ below the thresholds of Lemmas~\ref{lem:c3-7-2} and~\ref{lem:c3-7-4}, the second being informative when $q\leqslant 1$.
\end{lemma}

\begin{proof}
Fix a history and a labelled term. $\prod_{m\in\mathsf S}m$ is a nonnegative function of $U$, of the regions and of the labels ($a_b$ depends on the regions only, $M_e$ on the labels and on $U$), and $\prod_{\mathsf S}m\leqslant \prod_{\mathsf S'}m$, the marks taking the values $0$ and $1$. The reversal of Lemma~\ref{lem:c3-7-3}, applied in the proof of Proposition~\ref{prop:c5-2-6} to $\mathsf w^{\mathrm{lab}}_{\xi} M_e$ and in Lemma~\ref{lem:c5-3-0}(b) to $\mathsf w_{\xi} a_b$, is a change of variables for integrals of (density $\times$ nonnegative function of $U$, the regions and the labels), so it applies to $\mathsf w^{\mathrm{lab}}_{\xi}\prod_{\mathsf S'}m$. After the reversal, on the support of the term's cutoffs, each $M_e$ with $e\in\mathsf S'$ is dominated by $\mathbf 1_{W_e^{\vartheta}}$ (Lemma~\ref{lem:c5-2-4}: a plaquette or gauge-variable witness at threshold $\theta_{\vartheta}$ within $K\ell$ of $e$) and each $a_b$ by $\mathbf 1_{W_b}\leqslant \mathbf 1_{W_b^{\vartheta}}$ (Lemma~\ref{lem:c5-3-0}(b): a witness at threshold $\theta\geqslant \theta_{\vartheta}$ within $K\ell$ of $b$; the event at the larger threshold is contained in the event at the smaller one). The events $W^{\vartheta}$ are label-free, so, as in the proof of Proposition~\ref{prop:c5-2-6}, summing the nonnegative labelled terms and histories (their partitions sum to one) gives $\mathcal Z_0^{-1}\int\mathcal O[\sum_{\mathrm{lab}}\mathsf w^{\mathrm{lab}}_{\xi}\prod_{\mathsf S}m]\,\dd V\leqslant (\mu_0\otimes q_g^{\otimes\mathcal I}\otimes\mathrm{Haar}^{\otimes T})(\bigcap_{m\in\mathsf S'}W^{\vartheta}_m)$. Here $W^{\vartheta}_m$ is taken as the union, over the plaquette labels $p$ and the gauge-variable labels $w$ all of whose sites (the four corners of $p$; the centre of $w$ and the sites of its variable) lie within $\ell^{\infty}$ distance $K\ell+1$ of the initial site $x_m$ of the mark, of the single-witness events $\{| U(\partial p)-\id|_{\sigma}\geqslant \theta_{\vartheta}\}$ and $\{| w-\id|_{\sigma}\geqslant \theta_{\vartheta}\}$. The dominations above hold with this union: the proofs of Lemmas~\ref{lem:c3-7-2} and~\ref{lem:c5-2-4} place every site of the witnessing label within a distance of the bond that is a fixed multiple of $\ell$ far below $K\ell$ (the witness lies in the $6\ell$-neighbourhood $N(\square')$ of a cube of side at most $\ell/8$, its sites within $7\ell$ of the initial vertex of the bond and the gauge-label centres within $O(L)\leqslant \ell$ more, Lemma~\ref{lem:c5-2-4}; the neighbourhood of Lemma~\ref{lem:c3-7-2} is larger by a fixed multiple of $\ell$, still below $K\ell$ with $K=256$), so the witnessing label is among those of the union. The number of labels in the union is at most $7(2K\ell+3)^4\leqslant n_W$: at most $(2K\ell+3)^4$ sites lie within $K\ell+1$ of $x_m$, a plaquette is fixed by its base site and one of six plane orientations, and a gauge label by the noncentral site of its variable, which lies within $K\ell+1$ of $x_m$ and determines the centre (the centre of its block); a centre alone carries $L^4-1$ labels and would not fix one. So the intersection is contained in a union of at most $n_W^{|\mathsf S'|}$ intersections of single-witness events, one label per mark of $\mathsf S'$. Fix such a choice. Two marks of $\mathsf S'$ are at distance $>\sigma_K=2(K\ell+1)+2$, so their labels have all their sites at distance $>2$ from each other: the chosen gauge variables are distinct, and the chosen plaquettes have distinct base cubes $C_x$ (the closed unit cube whose minimal corner $x$ is the plaquette's base site). Under the product measure the probability of the intersection factorises into the $\mu_0$-probability of the plaquette events and the product of the $q_g$-probabilities of the gauge-variable events, each of the latter at most $z_G(\theta_{\vartheta})$ \eqref{eq:c3-7-13}. For the plaquette events apply the chessboard estimate \eqref{eq:c3-7-chessboard} with $f_x$, for each chosen base cube $C_x$, the indicator $\mathbf 1\{| U(\partial p_0)-\id|_{\sigma}\geqslant \theta_{\vartheta}\}$ of the plaquette $p_0$ of $C_0$ whose transport $f_x^{[x]}$ is the indicator of the chosen plaquette of $C_x$, and $f_x=1$ for the other cubes: every plaquette of $C_x$ is the transport of a plaquette of $C_0$ of the same plane orientation, because a reflection in direction $\rho$ carries the transverse coordinate $t\in\{0,1\}$ of a face to $1-t$, so the face of $C_0$ with transverse coordinates $(t_{\rho})$ is carried to the face of $C_x$ with transverse coordinates $x_{\rho}+t_{\rho}$ or $x_{\rho}+1-t_{\rho}$ according to the parity of $x_{\rho}$, and the indicator does not depend on the orientation or the base point of the plaquette. Each root $(\int\prod_yf_x^{[y]}\,\dd\mu_0)^{1/V_T}$ of a disseminated expectation is at most $z_P(\theta_{\vartheta})$ by the computation in the proof of Lemma~\ref{lem:c3-7-4}, which applies to every plaquette of $C_0$ (its dissemination contains $V_T/4$ distinct plaquettes of its orientation, each with $s_G\geqslant b_G\theta_{\vartheta}^2$ on the event, and \eqref{eq:c3-7-11} bounds $\mathcal Z_0$ below). Hence the fixed choice has probability at most $\max(z_P(\theta_{\vartheta}),z_G(\theta_{\vartheta}))^{|\mathsf S'|}$, the union bound over the choices gives $(n_W\max(z_P(\theta_{\vartheta}),z_G(\theta_{\vartheta})))^{|\mathsf S'|}= q_{\vartheta}^{|\mathsf S'|}$ \eqref{eq:c5-2-2}, and $q_{\vartheta}\leqslant 1$ with $|\mathsf S'|\geqslant |\mathsf S|/(2C_K)$ gives the second inequality. As in Lemma~\ref{lem:c3-7-3}, the product measure is a representation of the integral before the history conditions; no independence of the marks under the history measure is asserted, the factorisation is that of the physical measure after the reversal, and the correlation of distinct plaquettes under $\mu_0$ is paid by the chessboard estimate. For \eqref{eq:c5-7-joint-I} the same argument runs without labels: the reversal of Lemma~\ref{lem:c3-7-3} with the product of the marks $\mathbf 1\{ z\in E(B_{\xi})\}$, $z\in\mathsf S'$, inserted (each a function of the regions and the deterministic filling), each mark dominated by $\mathbf 1_{W_z}$ (Lemma~\ref{lem:c3-7-2}: a witness at threshold $\theta$ within $K\ell$ of $z$; \eqref{eq:c3-7-9}), the intersection $\bigcap_{z\in\mathsf S'}W_z$, with $W_z$ now the union of the single-witness events at threshold $\theta$ over the labels all of whose sites lie within $K\ell+1$ of the initial site of $z$ (it contains the witnessing label of Lemma~\ref{lem:c3-7-2} by the geometry above, so the domination holds with it), a union of at most $n_W^{|\mathsf S'|}$ intersections of single-witness events at threshold $\theta$ with the same separation, and the chessboard and product bounds as above at threshold $\theta$, giving $(n_W\max(z_P,z_G))^{|\mathsf S'|}= q^{|\mathsf S'|}$.
\end{proof}

\begin{lemma}[exponential moment of the boundary sum]\label{lem:c5-7-expmoment}
In the setting of Lemma~\ref{lem:c5-7-3}, for $c>0$ put $R_0'\coloneqq \lceil R_0\rceil+2L+1$ (an integer, at least $R_0+\varrho_K$ since $\varrho_K\leqslant 2L+1$ (Lemma~\ref{lem:c5-7-3}), and independent of the contour and of $n$; the lattice distances $d_T(r,b)$ exceeding $R_0'$ are then at least $R_0'+1$) and, for a bond $b$ and a bond $e$,
\begin{equation}\label{eq:c5-7-boundary-marks}
    \varphi_c(b)\coloneqq \sum_{r\in C}e^{-c(d_T(r,b)-R_0')_+}\leqslant \Phi_c\coloneqq 4(2R_0'+1)(1+1/c)+8/c^2\;,\qquad\Phi_c(e)\coloneqq \sum_{b\, : \,x_e\text{ within }102\ell+1\text{ of an endpoint of }b}\varphi_c(b)\leqslant N_{\mathrm{ret}}\Phi_c\;,
\end{equation}
and $N'_{\mathrm{lat}}(c)\coloneqq \sum_{s\geqslant 0}4(2s+1)^4e^{-c(s-R_0')_+}<\infty$. Then, on the support of the indicators $\mathbf 1_{\mathrm{lab}}$ of every labelled term with final region $\Lambda$ (there every bond found bad on the term has $M_e=1$; every use below carries the factor $\mathsf w^{\mathrm{lab}}_{\xi}$, which contains $\mathbf 1_{\mathrm{lab}}$),
\begin{equation}\label{eq:c5-7-pointwise-marks}
    \sum_{j=0}^{4n-1}e^{-cd_j(\Lambda)}\leqslant \sum_b\varphi_c(b)\,a_b+\sum_e\Phi_c(e)\,M_e\;,
\end{equation}
and for every $\lambda>0$
\begin{equation}\label{eq:c5-7-expmoment}
    \mathcal Z_0^{-1}\int\mathcal O\Bigl[\sum_{\mathrm{lab}}\mathsf w^{\mathrm{lab}}_{\xi}\mathbf 1_{A^{\mathrm{fin}}_{\mathrm{cov}}}\exp\Bigl(\lambda\sum_{j=0}^{4n-1}e^{-cd_j(\Lambda)}\Bigr)\Bigr]\dd V\leqslant \exp\bigl[n_C\,\lambda(1+N_{\mathrm{ret}})N'_{\mathrm{lat}}(c)\,e^{\lambda N_{\mathrm{ret}}\Phi_c}\,q_{\vartheta}^{1/(2C_K)}\bigr]\;,
\end{equation}
uniform in the volume; the bound \eqref{eq:c5-7-expmoment} holds with $\mathbf 1_{A^{\mathrm{fin}}_{\mathrm{cov}}}$ omitted, the proof using the indicator only through $\mathbf 1\leqslant 1$. At the original carrier ($D= I= D_+(\Lambda_1)$, the $d_j$ of \eqref{eq:regional-square-distances} at $D= I$, which are those of Proposition~\ref{prop:c4c-5-1}; no labels; compatible even fine periods, Section~\ref{sec:c3-7}; otherwise in the setting of Lemma~\ref{lem:c5-7-3}: the square $C$ with its $n_C$ contour bonds, the distances $d_T$, the constant $R_0$ behind $R_0'$ and the bound $\varrho_K\leqslant 2L+1$ used in the proof, and $\varphi_c$, $\Phi_c$, $N'_{\mathrm{lat}}(c)$ as defined above), pointwise on every branch and for every $\lambda>0$, and again with the second bound holding when $\mathbf 1_{A_{\mathrm{cov}}}$ is omitted,
\begin{equation}\label{eq:c5-7-expmoment-I}
    \sum_{j=0}^{4n-1}e^{-cd_j}\leqslant \sum_b\varphi_c(b)\mathbf 1\{ b\in E(B_{\xi})\}\;,\qquad\mathcal Z_0^{-1}\int\mathcal O\Bigl[\mathsf w_{\xi}\mathbf 1_{A_{\mathrm{cov}}}\exp\Bigl(\lambda\sum_{j=0}^{4n-1}e^{-cd_j}\Bigr)\Bigr]\dd V\leqslant \exp\bigl[n_C\,\lambda N'_{\mathrm{lat}}(c)\,e^{\lambda\Phi_c}\,q^{1/C_K}\bigr]\;,
\end{equation}
for $g_0$ below the thresholds of Lemmas~\ref{lem:c3-7-2} and~\ref{lem:c3-7-4}.
\end{lemma}

\begin{proof}
The bounds in \eqref{eq:c5-7-boundary-marks}: the initial sites of the contour bonds lie on the four lines of the square, each meeting an $\ell^{\infty}$ ball of radius $s$ in at most $2s+1$ sites, so at most $4(2s+1)$ contour bonds have $d_T(r,b)\leqslant s$ ($d_T(r,b)\geqslant \| x_r-x_b\|_{\infty}$); hence $\varphi_c(b)\leqslant 4(2R_0'+1)+\sum_{u\geqslant 1}4(2(R_0'+u)+1)e^{-cu}\leqslant 4(2R_0'+1)(1+1/c)+8/c^2$, using $\sum_{u\geqslant 1}e^{-cu}\leqslant 1/c$ and $\sum_{u\geqslant 1}ue^{-cu}\leqslant 1/c^2$. The bonds $b$ with an endpoint within $102\ell+1$ of $x_e$ have their initial site within $102\ell+2$ of $x_e$ and number at most $4(204\ell+5)^4\leqslant 4(207\ell)^4= N_{\mathrm{ret}}$ for $\ell\geqslant 2$ (an elementary count; $N_{\mathrm{ret}}$ is the constant of Lemma~\ref{lem:c5-3-1}(d)). \eqref{eq:c5-7-pointwise-marks}: $d_j(\Lambda)=(d_T(K_j,D(\Lambda)^c)-R_0)_+$ with $d_T(K_j,D(\Lambda)^c)=\min_{b\notin D(\Lambda)}d_T(K_j,b)$, and $x\mapsto e^{-c(x-R_0)_+}$ is nonincreasing, so $e^{-cd_j(\Lambda)}\leqslant \sum_{b\notin D(\Lambda)}e^{-c(d_T(K_j,b)-R_0)_+}$ (zero if $D(\Lambda)^c=\varnothing$); for $k\in K_j$, $d_T(k,b)\geqslant \min_{r\in C_j}d_T(r,b)-\varrho_K$ (proof of Lemma~\ref{lem:c5-7-3}), so $e^{-c(d_T(K_j,b)-R_0)_+}\leqslant e^{-c(d_T(C_j,b)-R_0')_+}\leqslant \sum_{r\in C_j}e^{-c(d_T(r,b)-R_0')_+}$, and $\sum_j\sum_{r\in C_j}=\sum_{r\in C}$. Thus $\sum_je^{-cd_j(\Lambda)}\leqslant \sum_b\varphi_c(b)\mathbf 1\{ b\notin D(\Lambda)\}$, the sum over the free bonds of the torus free carrier, which is enlarged below to the sum over all bonds $b$ of the torus (the marks $a_b$ and $M_e$ are defined for every bond, on every history and labelled term, Section~\ref{sec:c5-2-3} and Lemma~\ref{lem:c5-3-0}, so the index set of the marks is the same for every history, as the exchange of the history sum with the expansion below requires; the added terms are nonnegative, and the counts below run over all bonds), and $\mathbf 1\{ b\notin D(\Lambda)\}\leqslant a_b+\sum_{e\, : \,x_e\text{ within }102\ell+1\text{ of an endpoint of }b}M_e$ on every labelled term (proof of Lemma~\ref{lem:c5-7-3}, from \eqref{eq:c5-3-mark-inequality}); exchanging the two sums gives \eqref{eq:c5-7-pointwise-marks}. \eqref{eq:c5-7-expmoment}: by \eqref{eq:c5-7-pointwise-marks} and $\mathbf 1_{A^{\mathrm{fin}}_{\mathrm{cov}}}\leqslant 1$ the integrand is at most $\mathsf w^{\mathrm{lab}}_{\xi}\prod_b(1+(e^{\lambda\varphi_c(b)}-1)a_b)\prod_e(1+(e^{\lambda\Phi_c(e)}-1)M_e)$, since $e^{\lambda wm}=1+(e^{\lambda w}-1)m$ for $m\in\{0,1\}$; the products are finite. Expand them into the sum over the sets $\mathsf S$ of marks of $\prod_{m\in\mathsf S}(e^{\lambda w_m}-1)\prod_{m\in\mathsf S}m$, $w_m$ the weight $\varphi_c(b)$ or $\Phi_c(e)$ of the mark; every term is nonnegative, so the history average (positive; finitely many labels per history; Tonelli) is at most $\sum_{\mathsf S}\prod_{m\in\mathsf S}(e^{\lambda w_m}-1)\,q_{\vartheta}^{|\mathsf S|/(2C_K)}$ by Lemma~\ref{lem:c5-7-joint}, which equals $\prod_m(1+(e^{\lambda w_m}-1)\hat q)$ with $\hat q\coloneqq q_{\vartheta}^{1/(2C_K)}$, and $\prod_m(1+(e^{\lambda w_m}-1)\hat q)\leqslant \exp(\hat q\sum_m(e^{\lambda w_m}-1))$. Finally $e^{x}-1\leqslant xe^{x}$ for $x\geqslant 0$, $w_m\leqslant N_{\mathrm{ret}}\Phi_c$ for every mark, and $\sum_mw_m=\sum_b\varphi_c(b)+\sum_e\Phi_c(e)\leqslant (1+N_{\mathrm{ret}})\sum_b\varphi_c(b)\leqslant (1+N_{\mathrm{ret}})\,n_CN'_{\mathrm{lat}}(c)$: $\sum_e\Phi_c(e)=\sum_b\varphi_c(b)\,\#\{ e\, : \,x_e\text{ within }102\ell+1\text{ of an endpoint of }b\}\leqslant N_{\mathrm{ret}}\sum_b\varphi_c(b)$ (Lemma~\ref{lem:c5-3-1}(d)), and $\sum_b\varphi_c(b)=\sum_{r\in C}\sum_be^{-c(d_T(r,b)-R_0')_+}\leqslant n_CN'_{\mathrm{lat}}(c)$, since at most $4(2s+1)^4$ bonds have initial site within $\ell^{\infty}$ distance $s$ of $x_r$ (as in the proof of Lemma~\ref{lem:c5-7-3}). \eqref{eq:c5-7-expmoment-I}: the first steps are the same with $D(\Lambda)$ replaced by $I$; a free bond $b\notin I$ has no endpoint in $\Lambda_1$ ($I$ is the set of all free bonds with an endpoint in $\Lambda_1$, the bond convention of Section~\ref{sec:c5-0}), so both its endpoints lie in $B_{\xi}=\Lambda_1^c$ and $\mathbf 1\{ b\notin I\}\leqslant \mathbf 1\{ b\in E(B_{\xi})\}$, which gives the pointwise bound; the expansion is the same with the marks $\mathbf 1\{ b\in E(B_{\xi})\}$, one per bond with weight $\varphi_c(b)\leqslant \Phi_c$ and $\sum_b\varphi_c(b)\leqslant n_CN'_{\mathrm{lat}}(c)$, the history average $\mathcal Z_0^{-1}\int\mathcal O[\mathsf w_{\xi}\,\bigcdot \,]\,\dd V$ in place of the labelled one, and \eqref{eq:c5-7-joint-I} in place of \eqref{eq:c5-7-joint}.
\end{proof}

\begin{corollary}[bounds relative to $\kappa_C$]\label{cor:c5-7-4}
Under \eqref{eq:c5-7-T} and the hypotheses of Proposition~\ref{prop:c5-5R-4} and Lemma~\ref{lem:c5-7-3}:
\begin{enumerate}[label=(\alph*)]
    \item (crude form) $\mathcal Z_0^{-1}\int\mathcal O[\sum_{\mathrm{lab}}\int\mathbf 1_{\mathrm{lab}}e^{c^{\Lambda}_{\xi}}I^{\Lambda}_1\mathbf 1_{A^{\mathrm{fin}}_{\mathrm{cov}}}|\kappa_C|^{-1}|\mathcal B^{\Lambda}_{\xi}[f_{\xi}]|\,\dd a_{\mathrm{old}}]\,\dd V\leqslant e^{C\tau^2(9n+1)}\times$(the right side of \eqref{eq:c5-6-3} with $e^{n_C\bar\varepsilon_{\mathrm{ph}}}$) $= e^{C\tau^2(9n+1)}e^{n_C\bar\varepsilon_{\mathrm{ph}}}(C_1g_0r+C_2\bar\varepsilon_tr^{d_G})(1+e^{\mu_H(2L+1)/2})N_{\mathrm{lat}}(\mu')n_C(q+N_{\mathrm{ret}}q_{\vartheta})$, by \eqref{eq:c5-7-3} and the statement \eqref{eq:c5-6-3} of Proposition~\ref{prop:c5-6-4} (monotonicity of its integrand) with $|\mathcal B^{\Lambda}|/| S^{\circ}_{\Lambda}|$ replaced by $|\kappa_C|^{-1}|\mathcal B^{\Lambda}|\leqslant e^{C\tau^2(9n+1)}|\mathcal B^{\Lambda}|/| S^{\circ}_{\Lambda}|$; by the comparison of Proposition~\ref{prop:c5-6-4} ($r^{d_G}N_{\mathrm{ret}}(q+N_{\mathrm{ret}}q_{\vartheta})\leqslant e^{-c_q\vartheta^2\pi_*^2/2}$ for $g_0$ below a further threshold; $N_{\mathrm{ret}}$ and $r^{d_G}$ themselves diverge as $g_0\to0$) this is at most $C'n_Ce^{n_C\bar\varepsilon_{\mathrm{ph}}}e^{C\tau^2(9n+1)}e^{-c_q\vartheta^2\pi_*^2/2}$ with the constant $C'$ of the closing form of Proposition~\ref{prop:c5-6-4}, and $e^{C\tau^2(9n+1)}$ is the factor carried in $B_C=|\kappa_C|\exp(C\tau^2(9n+1))$ following \eqref{eq:retained-source-correction}.
    \item (first moment of the leading error) $\mathcal Z_0^{-1}\int\mathcal O[\sum_{\mathrm{lab}}\mathsf w^{\mathrm{lab}}_{\xi}\mathbf 1_{A^{\mathrm{fin}}_{\mathrm{cov}}}| E^{\circ}_{\Lambda}|]\,\dd V\leqslant C\tau^2(2+1/c)+C\tau^2e^{c(R_0+\varrho_K)}N_{\mathrm{lat}}(c)n_C(q+N_{\mathrm{ret}}q_{\vartheta})$ under \eqref{eq:c5-7-T}(iii), by Corollary~\ref{cor:c5-7-2}(c), Lemma~\ref{lem:c5-7-3} and the normalisation $\mathcal Z_0^{-1}\int\mathcal O[\sum_{\mathrm{lab}}\mathsf w^{\mathrm{lab}}_{\xi}\mathbf 1_{A^{\mathrm{fin}}_{\mathrm{cov}}}]\,\dd V\leqslant \mathcal Z_0^{-1}\int\mathcal O[\mathsf w_{\xi}]\,\dd V=1$ ($\sum_{\mathrm{lab}}\mathsf w^{\mathrm{lab}}_{\xi}=\mathsf w_{\xi}$, Section~\ref{sec:c5-3-4}(c); $\mathcal Z_0=\int\mathcal O[\mathcal J^{\mathrm a}_{\xi}[1]]\,\dd V=\int\mathcal O[\mathsf w_{\xi}]\,\dd V$ by the measurability paragraph of Section~\ref{sec:branch-data} and \eqref{eq:c-collar-identity}; $\mathcal O$ is positive and $\mathsf w^{\mathrm{lab}}_{\xi}\geqslant 0$, as in the proof of Proposition~\ref{prop:c5-6-4}): the labelled leading error has, in label average, a part $O(\tau^2)$ uniform in $n$ and a part linear in $n_C$ paid by the one-point witness probabilities (small for $g_0$ below the threshold of the comparison of Proposition~\ref{prop:c5-6-4}, at fixed $n_C$); the first moment is not uniform in $n$. The exponential moment is (c).
    \item (exponential moment of the leading error) Under \eqref{eq:c5-7-T}(iii), with the constants $C$, $c$ of \eqref{eq:c5-7-2}, $N'_{\mathrm{lat}}(c)$, $\Phi_c$, $C_K$ of Lemmas~\ref{lem:c5-7-joint} and~\ref{lem:c5-7-expmoment}, and
    \begin{equation}\label{eq:c5-7-delta}
        \delta_q\coloneqq C\tau^2(1+N_{\mathrm{ret}})N'_{\mathrm{lat}}(c)\,e^{C\tau^2N_{\mathrm{ret}}\Phi_c}\,q_{\vartheta}^{1/(2C_K)}\;,
    \end{equation}
    \begin{equation}\label{eq:c5-7-expbound}
        \mathcal Z_0^{-1}\int\mathcal O\Bigl[\sum_{\mathrm{lab}}\mathsf w^{\mathrm{lab}}_{\xi}\mathbf 1_{A^{\mathrm{fin}}_{\mathrm{cov}}}e^{| E^{\circ}_{\Lambda}|}\Bigr]\dd V\leqslant \exp\bigl[C\tau^2(2+1/c)+n_C\delta_q\bigr]\;,
    \end{equation}
    by Corollary~\ref{cor:c5-7-2}(c) ($| E^{\circ}_{\Lambda}|\leqslant C\tau^2(2+1/c)+C\tau^2\sum_je^{-cd_j(\Lambda)}$) and Lemma~\ref{lem:c5-7-expmoment} at $\lambda= C\tau^2$. Consequently the leading term of \eqref{eq:c5-6-1}--\eqref{eq:c5-6-2} in the form of Corollary~\ref{cor:c5-7-2}(b) satisfies
    \begin{equation}\label{eq:c5-7-leading-modulus}
        |\kappa_C|^{-1}\Bigl|\mathcal Z_0^{-1}\int\mathcal O\Bigl[\kappa_C\sum_{\mathrm{lab}}e^{E^{\circ}_{\Lambda}}\int\mathbf 1_{\mathrm{lab}}e^{c^{\Lambda}_{\xi}}I^{\Lambda}_1\mathbf 1_{A^{\mathrm{fin}}_{\mathrm{cov}}}W_{\chi}(C;U_1)\,\dd a_{\mathrm{old}}\Bigr]\dd V\Bigr|\leqslant \exp\bigl[C\tau^2(2+1/c)+n_C\delta_q\bigr]
    \end{equation}
    (the left side is the modulus of the leading term of Corollary~\ref{cor:c5-7-2}(b) divided by $|\kappa_C|$; $| W_{\chi}|\leqslant 1$, $| e^{E^{\circ}_{\Lambda}}|\leqslant e^{| E^{\circ}_{\Lambda}|}$, $E^{\circ}_{\Lambda}$ and $\mathbf 1_{A^{\mathrm{fin}}_{\mathrm{cov}}}$ depend on the labels only, and the $a_{\mathrm{old}}$-integral of $\mathbf 1_{\mathrm{lab}}e^{c^{\Lambda}_{\xi}}I^{\Lambda}_1$ is $\mathsf w^{\mathrm{lab}}_{\xi}$, Section~\ref{sec:c5-2-3}), and $\mathcal Z_0^{-1}\int\mathcal O[\sum_{\mathrm{lab}}\mathsf w^{\mathrm{lab}}_{\xi}\mathbf 1_{A^{\mathrm{fin}}_{\mathrm{cov}}}| e^{E^{\circ}_{\Lambda}}-1|]\,\dd V\leqslant \exp[C\tau^2(2+1/c)+n_C\delta_q]-1$: pointwise on every labelled term, $\mathbf 1_{A^{\mathrm{fin}}_{\mathrm{cov}}}| e^{E^{\circ}_{\Lambda}}-1|\leqslant \mathbf 1_{A^{\mathrm{fin}}_{\mathrm{cov}}}(e^{| E^{\circ}_{\Lambda}|}-1)\leqslant e^{C\tau^2(2+1/c)}\exp\bigl(C\tau^2\sum_je^{-cd_j(\Lambda)}\bigr)-1$ by Corollary~\ref{cor:c5-7-2}(c) and $\mathbf 1\leqslant 1$ (the last member is nonnegative), and the label average of the last member is at most $\exp[C\tau^2(2+1/c)+n_C\delta_q]-1$ by \eqref{eq:c5-7-expmoment} with the indicator omitted and the exact normalisation $\mathcal Z_0^{-1}\int\mathcal O[\sum_{\mathrm{lab}}\mathsf w^{\mathrm{lab}}_{\xi}]\,\dd V=1$ of (b). (The indicator is not subtracted: $\mathcal Z_0^{-1}\int\mathcal O[\sum_{\mathrm{lab}}\mathsf w^{\mathrm{lab}}_{\xi}\mathbf 1_{A^{\mathrm{fin}}_{\mathrm{cov}}}]\,\dd V$ may be less than one.) Comparison: since $\tau<1$ (fixed after \eqref{eq:marginal-fixed-forest-bound}), $\log\delta_q\leqslant \log(\tau^{-2}\delta_q)\leqslant \log C+\log(1+N_{\mathrm{ret}})+\log N'_{\mathrm{lat}}(c)+C\tau^2N_{\mathrm{ret}}\Phi_c-c_q\vartheta^2\pi_*^2/(2C_K)$, where $N'_{\mathrm{lat}}(c)$ and $\Phi_c$ are constants of $c$ and $R_0'=\lceil R_0\rceil+2L+1\leqslant R_0+2L+2$ ($R_0$ the fixed range constant of \eqref{eq:regional-collar-sizes}), $\log N_{\mathrm{ret}}= O(\log\log g_0^{-2})$ \eqref{eq:c5-3-1}, $\tau^2N_{\mathrm{ret}}\to0$ as $g_0\to0$ ($\tau= C[g_0(1+r)^3+e^{-cr^2}]$, $r= A_1(\log g_0^{-2})^{p_0}$), and $C_K=4(4K\ell+9)^4\leqslant 4(4K+9)^4\ell^4\leqslant 4(4K+9)^4(L^2M)^42^{4\mathfrak r}(\log g_0^{-2})^{4\mathfrak r}$ while $\pi_*^2\geqslant A_0^22^{-2p_0}(\log g_0^{-2})^{2p_0}$ with $2p_0-4\mathfrak r\geqslant 6\mathfrak r\geqslant 12$ (\eqref{eq:c5-3-1}--\eqref{eq:c5-3-2}); so for $g_0$ below a further threshold depending on $L$, $M$, $M_1$, $M_2$, the group, the character (through the constant $C$ of \eqref{eq:c5-7-2}), the source constants, $\vartheta$ and $c$,
    \begin{equation}\label{eq:c5-7-delta-small}
        \delta_q\leqslant \exp\bigl[-c_q\vartheta^2\pi_*^2/(4C_K)\bigr]\leqslant \exp\bigl[-c'(\log g_0^{-2})^{2p_0-4\mathfrak r}\bigr]\;,
    \end{equation}
    smaller than every power of $g_0$, with $c'>0$ depending on $L$, $M$, $M_1$, $M_2$, the group, the character, the source constants and $\vartheta$. The $n$-dependence of the label average of $e^{| E^{\circ}_{\Lambda}|}$ is thus the perimeter-rate factor $e^{n_C\delta_q}$, of the same kind as the factor $n_C(q+N_{\mathrm{ret}}q_{\vartheta})$ of (a) and of \eqref{eq:c5-6-4}; a bound independent of $n$ is not claimed, and the perimeter rate $\delta_q$ is governed by the $2C_K$-th root $q_{\vartheta}^{1/(2C_K)}$ of the one-point rate $q_{\vartheta}$ of bad collar bonds and round-$0$ bonds along the contour, the separation costing the factor $2C_K$, of order $(\log g_0^{-2})^{4\mathfrak r}$, in the exponent. At the original carrier $D= I$, under the conventions of Section~\ref{sec:c4c-1}, on compatible even fine periods (Section~\ref{sec:c3-7}) and for $g_0< g_*$ (Proposition~\ref{prop:c3-7-5}), the same holds for the identity \eqref{eq:c4c-5-3} of Corollary~\ref{cor:c4c-5-2}, with $E^{\circ}_{\xi}= E^{\circ}_{I,C_n,T}$ of Proposition~\ref{prop:c4c-5-1} (whose bound has the form of \eqref{eq:c5-7-2}, with constants $C$, $c$ fixed before the coupling, so Corollary~\ref{cor:c5-7-2}(c) applies verbatim under $4ne^{-c\iota_T}\leqslant 1$), the history expectation $\mathbb E_{\mathrm h}$ of that corollary (weight $\mathsf w_{\xi}$, physical $\mathcal Z_0$) and \eqref{eq:c5-7-expmoment-I} at $\lambda= C\tau^2$:
    \begin{equation}\label{eq:c5-7-expbound-I}
        \mathbb E_{\mathrm h}\bigl[\mathbf 1_{A_{\mathrm{cov}}}e^{| E^{\circ}_{\xi}|}\bigr]\leqslant \exp\bigl[C\tau^2(2+1/c)+n_C\delta^I_q\bigr]\;,\qquad\delta^I_q\coloneqq C\tau^2N'_{\mathrm{lat}}(c)\,e^{C\tau^2\Phi_c}\,q^{1/C_K}\leqslant \exp\bigl[-c_q\pi_*^2/(2C_K)\bigr]\;,
    \end{equation}
    hence the leading term $\kappa_{C_n}\mathbb E_{\mathrm h}[\mathbf 1_{A_{\mathrm{cov}}}e^{E^{\circ}_{\xi}}W_{\chi}(C_n;U_1)]$ of \eqref{eq:c4c-5-3} has modulus at most $|\kappa_{C_n}|\exp[C\tau^2(2+1/c)+n_C\delta^I_q]$, that is $|\mathbb E_{\mathrm h}[\mathbf 1_{A_{\mathrm{cov}}}e^{E^{\circ}_{\xi}}W_{\chi}(C_n;U_1)]|\leqslant \exp[C\tau^2(2+1/c)+n_C\delta^I_q]$, and $\mathbb E_{\mathrm h}[\mathbf 1_{A_{\mathrm{cov}}}| e^{E^{\circ}_{\xi}}-1|]\leqslant \exp[C\tau^2(2+1/c)+n_C\delta^I_q]-1$ by the pointwise bound $\mathbf 1_{A_{\mathrm{cov}}}| e^{E^{\circ}_{\xi}}-1|\leqslant e^{C\tau^2(2+1/c)}\exp(C\tau^2\sum_je^{-cd_j})-1$ (Corollary~\ref{cor:c5-7-2}(c) for $E^{\circ}_{\xi}$ and $\mathbf 1\leqslant 1$), \eqref{eq:c5-7-expmoment-I} with the indicator omitted and $\mathbb E_{\mathrm h}[1]=1$ ($\mathbb E_{\mathrm h}$ is a probability expectation, $\mathcal Z_0=\int\mathcal O[\mathsf w_{\xi}]\,\dd V$); the last inequality of \eqref{eq:c5-7-expbound-I} holds for $g_0$ below a further threshold, depending on $L$, $M$, $M_1$, $M_2$, the group, the character, the source constants and $c$, by the comparison above with $q\leqslant e^{-c_*\pi_*^2}= e^{-c_q\pi_*^2}$ for $g_0< g_*$ (Proposition~\ref{prop:c3-7-5}, $c_*= c_q$; this is \eqref{eq:c5-2-3} at $\vartheta=1$, where $\theta_1=\theta$ and $q_1= q$, as in the proof of Proposition~\ref{prop:c5-6-4}); $g_*$ is the threshold named in (H2) of the main theorem.
\end{enumerate}
\end{corollary}

\begin{remark}[scope of this subsection]\label{rem:c5-7-5}
The identity $S^{\circ}_{\Lambda}= S^{\circ}_{D(\Lambda),C,T}$ on $A^{\mathrm{fin}}_{\mathrm{cov}}$ (Lemma~\ref{lem:c5-7-1}) is a factorwise comparison of the two explicit integrals (precision, cutoff, primitives, Wilson factor) using the facts (I3) at $D= D(\Lambda)\subseteq D_+(\Lambda_1)$; the identification deferred in the closing remark of Section~\ref{sec:sourced-effective-interaction} (after Lemma~\ref{lem:covered}) is, at the retreated carrier, reduced to these four factor identities. (1) Proposition~\ref{prop:regional-square-extraction} and the facts of the identification paragraph are used at their stated status; in particular (F3) is the label-by-label agreement of Lemma~\ref{lem:identity-primitives} between the primitives of Section~\ref{sec:source-phase-primitives} at the identity background (the exponent of Section~\ref{sec:sourced-effective-interaction}, CRE (1.14)--(1.16), Lemma~\ref{lem:primitive-sum}) and the terms of \eqref{eq:perimeter-literal-local-terms} (energy residuals plaquette by plaquette, gauge terms centre by centre as sums over the pairs, Haar and constraint-Jacobian primitives bond by bond and row by row), the same lemma that the second paragraph of the proof of Proposition~\ref{prop:regional-square-extraction} invokes for its own kernel. (2) The inclusion $\Lambda_r\subseteq\mathcal F_r$ (the coverage stencil of Section~\ref{sec:remaining-physical-comparison} inside the free-input set of the traversal link), which is what makes $A^{\mathrm{fin}}_{\mathrm{cov}}$ imply the coverage of (I2), follows, for a free bond $r$, from $\Lambda_r=\{ r\}=\mathcal F_r$ and, for a designated bond, from $\mathcal F_{b_0(c)}=\mathcal E_c\cap E_*$ \eqref{eq:regular-stencil-extension-support} and $\Lambda_{b_0(c)}=\mathcal S_c\cap E_*$, i.e. from $\mathcal S_c\subseteq\mathcal E_c$, which holds by the two definitions (Appendix~\ref{sec:perimeter-attempt}: $\mathcal S_c$ consists of the positive fine bonds with both endpoints in $\mathcal B_{c_-}\cup\mathcal B_{c_+}$; Section~\ref{sec:regular-hessian-stencils}: $\mathcal E_c$ consists of all bonds with both endpoints there), as recorded in the proof of Lemma~\ref{lem:c5-7-1}. (3) The comparison statement of Section~\ref{sec:regional-identification} is for $D_-\subseteq D\subseteq D_+$; $D(\Lambda)\subseteq D_+$ but $D_-\subseteq D(\Lambda)$ fails whenever a bond with both endpoints in $\Lambda_1$ is de-younged; Lemma~\ref{lem:c5-7-1} uses only $D(\Lambda)\subseteq D_+$ with the collar, which is what the arguments for the facts (1)--(5) summarised in (I3) use. (4) The label average of $e^{| E^{\circ}_{\Lambda}|}$ is bounded by $\exp[C\tau^2(2+1/c)+n_C\delta_q]$ with $\delta_q$ smaller than every power of $g_0$ (Corollary~\ref{cor:c5-7-4}(c), from the joint witness bound of Lemma~\ref{lem:c5-7-joint}, which is the chessboard estimate \eqref{eq:c3-7-chessboard} applied to several plaquettes at once after the reversal of Lemma~\ref{lem:c3-7-3}); its $n$-dependence is a perimeter rate, and a bound independent of $n$ is not claimed. (5) $E^{\mathrm{lab}}_{\mathrm{def}}$ is bounded in Proposition~\ref{prop:c5-6-def} with $S^{\circ}_{\Lambda}$ retained ($| S^{\circ}_{\Lambda}|\leqslant 1$) and is untouched by this subsection; \eqref{eq:retained-leading-source} already keeps $e^{E^{\circ}_{\xi}}$ inside the history expectation as a label-dependent factor, which is the form Corollary~\ref{cor:c5-7-2}(b) gives at the retreated carrier. (6) \eqref{eq:c5-7-T}(i) adds the $R_0$-collar condition to the $R_{\mathrm c}$-condition of threshold item 1 of Appendix~\ref{app:thresholds} (listed there among the conditions depending on the coupling through $R_1$); both are implied by conditions on $M$ alone, $LM-4L^2M_2>4R_0+8$ and $LM-4L^2M_2>4R_{\mathrm c}+8$, since $R_1\geqslant 1$, and are compatible (choose $M$ after both; $R_0=4(1+R_*+R_H+L)$ is fixed before $M$ in the order of Appendix~\ref{app:thresholds}, $R_*$ and $R_H$ being the ranges of Section~\ref{sec:marginal-perimeter}); whether $R_0\leqslant R_{\mathrm c}$ or conversely is not needed. (7) Further inputs used as stated: the reference-data cutoff convention behind (F2); $I= D_+(\Lambda_1)$ in Lemma~\ref{lem:c5-7-3}; the geometry of $\Lambda_{b_0(c)}$ behind $\varrho_K$; the square orientation and $M\geqslant 8LM_2$ of \eqref{eq:c5-7-T}(i); the normalisation $\mathcal Z_0=\int\mathcal O[\mathsf w_{\xi}]\,\dd V$ in Corollary~\ref{cor:c5-7-4}(b); \eqref{eq:regional-periodic-weight}, \eqref{eq:flat-phase}, \eqref{eq:perimeter-literal-local-terms}, \eqref{eq:marginal-fixed-forest-bound} ($\tau$), Lemma~\ref{lem:two-block} ($\mathcal E_c$), \eqref{eq:regular-stencil-extension-support}, and Lemmas~\ref{lem:c5-1-2} and~\ref{lem:c5-4-0}(c).
\end{remark}

\section{Two Disjoint Squares}\label{sec:two-loops}

This section extends the extraction of Section~\ref{sec:joint-extraction} and Appendices~\ref{app:signed-defects} and~\ref{app:regional-scalar} from one square to the product of the Wilson functions of two vertex-disjoint squares in planes parallel to that of $e_1,e_2$. Nothing in the single-square sections is changed; the objects of this section are defined for the pair and compared with the single-square objects of each square. Section~\ref{sec:two-loops-retreated} records the corresponding statements at the retreated carrier of Section~\ref{sec:one-step-insertion}; Theorem~\ref{thm:two-loop-theorem} does not use them.

\subsection{Setting}\label{sec:two-loops-setting}
Let $C_1,C_2$ be coarse translates of the counterclockwise squares with vertices $0,n_iLe_1,n_iL(e_1+e_2),n_iLe_2$, $i=1,2$, in fine-lattice units (so both lie in planes parallel to that of $e_1,e_2$); let $V(C_i)$ be the set of fine sites of $C_i$, the endpoints of its bonds, and assume $V(C_1)\cap V(C_2)=\varnothing$. Let $\chi_1,\chi_2$ be normalised characters of irreducible unitary representations $\pi_1,\pi_2$ of $G$, $\chi_i=\dim(\pi_i)^{-1}\operatorname{tr}\pi_i$. The reversed orientation of $C_i$ is the case $\bar\chi_i$, since $\chi_i(U^{-1})=\overline{\chi_i(U)}$. Put $C= C_1\sqcup C_2$ and
\begin{equation}\label{eq:two-loop-observable}
    W_{\chi_1,\chi_2}(C;U)=\chi_1\bigl(\Hol_UC_1\bigr)\,\chi_2\bigl(\Hol_UC_2\bigr)\;.
\end{equation}
The traversal set is $\mathcal I_C=\mathcal I_{C_1}\sqcup\mathcal I_{C_2}$, the traversals of each $C_i$ numbered cyclically from the base point used for the single-square objects of $C_i$; $n_{C_i}=4n_iL$ and $n_C= n_{C_1}+n_{C_2}$. Each square satisfies the square/torus geometry of Section~\ref{sec:extent-weights} with its own side, $K_{\mu}\geqslant 3(n_i+2)$ for every coarse period and $i=1,2$; no condition relating the positions of the two squares is imposed beyond $V(C_1)\cap V(C_2)=\varnothing$. For $a\in\mathcal I_{C_i}$ let $\delta_a$ be the periodic $\ell^{\infty}$ fine distance from the initial site of the traversal $a$ to $V(C_{3-i})$, and put
\begin{equation}\label{eq:two-loop-cross-sum}
    \Delta_{\times}=\sum_{a\in\mathcal I_C}e^{-c_{\times}\delta_a/(LM)}\leqslant n_Ce^{-c_{\times} d/(LM)}\;,\qquad d=\min_{a\in\mathcal I_C}\delta_a\;,
\end{equation}
with the constant $c_{\times}>0$ of Lemma~\ref{lem:two-loop-cross}, fixed before the coupling.

A superscript or subscript $i$ marks the single-square objects of the square $C_i$ with the character $\chi_i$: the contact sets $\mathsf T^i_D=\mathsf T_D\cap\mathcal I_{C_i}$, the groups of Definition~\ref{def:groups} for the contour $C_i$ alone, their reference scalars $a^{(i)}_W$ of \eqref{eq:part1-reference-group}, the reference gas $S^{\circ}_{\xi,C_i}$ of \eqref{eq:part1-reference-gas}, and the single-square lemmas of Appendix~\ref{app:signed-defects} with $C= C_i$. All of them are formed on the same branch $\xi$ with the same carrier, cubes, hosts, measures and cluster occurrences: the domains, probe measures $\widehat\nu^{\mathrm f}_{\xi,D}$, vacuum clusters and reference data of Sections~\ref{sec:fixed-first-kernel}--\ref{sec:joint-extraction} do not involve the observable, which enters only through the contact sets, $f_{\xi}$, $f^{\circ}_{\xi}$ and the reference scalars $a_W$. The coupling is below the thresholds of the single-square results for $\chi_1$ and for $\chi_2$ and below those of this section.

Covered and far are taken for both squares: $\xi\in A_{\mathrm{cov}}$ means that $\xi$ is covered (Lemma~\ref{lem:covered}) for the contour $C= C_1\sqcup C_2$, which, the conditions being imposed bond by bond, is the intersection of the covered events of $C_1$ and $C_2$; a history is far when $\operatorname{dist}(\Gamma_C,E(\Lambda_1^c))> b$ with $\Gamma_C=\Gamma_{C_1}\cup\Gamma_{C_2}$, that is far for both squares.

\subsection{Groups for two squares}\label{sec:two-loops-groups}
For the pair, $f_{\xi}(a)= W_{\chi_1,\chi_2}(C;U_{\xi}(a))$, and the contact set $\mathsf T_D$ of \eqref{eq:part1-contact-set} is taken over $n\in\mathcal I_C$. The argument preceding Definition~\ref{def:groups} applies verbatim up to the formation of the components: unmarked domains contribute zero, $q_{\xi}(\mathcal A)$ is expanded by \eqref{eq:c-cluster-expansion}, and the graph on the marked roots and the cluster occurrences of a term has components $\gamma$, each with a marked root and the contact set $\mathsf T_{\gamma}\subseteq\mathcal I_C$, the union of the contact sets of its roots. For $i=1,2$ let $H^i_{\gamma}=\varnothing$ if $\mathsf T_{\gamma}\cap\mathcal I_{C_i}=\varnothing$; otherwise $H^i_{\gamma}$ is the hull of $\mathsf T_{\gamma}\cap\mathcal I_{C_i}$ in the cyclic set $\mathcal I_{C_i}$, defined as in Section~\ref{sec:joint-extraction}: the shortest cyclic interval containing it when that interval has fewer than $n_{C_i}/2$ positions, and $\mathcal I_{C_i}$ otherwise. Put $H_{\gamma}= H^1_{\gamma}\sqcup H^2_{\gamma}$. Two components overlap when $H^i_{\gamma}\cap H^i_{\gamma'}\neq\varnothing$ for some $i$; the overlap classes are the connected components of the overlap graph on the components of the term. A set $K$ of components is long on $C_i$ when $\bigl|\bigcup_{\gamma\in K}H^i_{\gamma}\bigr|\geqslant n_{C_i}/2$.

\begin{definition}[groups for two squares]\label{def:two-loop-groups}
A partition $\mathfrak P$ of the components of a term is admissible if (A1) every overlap class lies in one block, and (A2) a block that is long on $C_i$ contains every component $\gamma$ with $\mathsf T_{\gamma}\cap\mathcal I_{C_i}\neq\varnothing$. The groups of the term are the blocks of its finest admissible partition (Lemma~\ref{lem:two-loop-grouping}(a)). A group $W$ consists of its root family $\mathcal A(W)$, its ordered list of cluster occurrences, and its interval set $\mathcal H(W)=\mathcal H_1(W)\sqcup\mathcal H_2(W)$, where $\mathcal H_i(W)=\mathcal I_{C_i}$ if $W$ is long on $C_i$, read from the base point of $C_i$, and $\mathcal H_i(W)=\bigcup_{\gamma\in W}H^i_{\gamma}$ otherwise, a finite union of cyclic intervals (a cyclic interval or $\mathcal I_{C_i}$ for a $C_i$-group, by Lemma~\ref{lem:two-loop-grouping}(d)); as in Definition~\ref{def:groups}, these data arise from a term, and conversely the term with root family $\mathcal A(W)$ and cluster list of $W$ has $W$ as its only group, with the same interval set (Lemma~\ref{lem:two-loop-grouping}(a)). Its weight $w_{\xi}(W)$ is \eqref{eq:part1-group-weight}. $W$ is a $C_i$-group if $\mathcal H_{3-i}(W)=\varnothing$ and a cross group if $\mathcal H_1(W)\neq\varnothing\neq\mathcal H_2(W)$. Every group has $\mathcal H(W)\neq\varnothing$ (each component has a marked root), so it is exactly one of a $C_1$-group, a $C_2$-group and a cross group. Two groups are compatible when all their constituent domains are mutually compatible and $\mathcal H_i(W)\cap\mathcal H_i(W')=\varnothing$ for $i=1,2$. A compatible group family is a finite set of pairwise compatible groups; the singleton clause of Definition~\ref{def:groups} is not imposed (for families of $C_i$-groups it is implied, Lemma~\ref{lem:two-loop-grouping}(c)). For a compatible group family, $\mathcal A(\mathcal S)$ and $A_{\xi}(\mathcal S)$ are \eqref{eq:part1-closed-family-pairing}.
\end{definition}

\begin{lemma}[two-square grouping]\label{lem:two-loop-grouping}
(a) Every term has a finest admissible partition; each of its blocks that is long on neither square is an overlap class. If no component of the term has a contact on $C_{3-i}$, it is the grouping of Section~\ref{sec:joint-extraction} for the square $C_i$. For a group $W$ of a term, the term $T_W$ with root family $\mathcal A(W)$ and the cluster list of $W$ has $W$ as its only group, with the same interval set. (b) The groups of a term are pairwise compatible, and Lemma~\ref{lem:grouped-source} holds for the pair: $F_{\xi}(0)=\sum_{\mathfrak B\text{ compatible}}A_{\xi}(\mathfrak B)$, absolutely convergent. (c) A $C_i$-group of the pair is a group of the single square $C_i$ (Definition~\ref{def:groups} with $C= C_i$) none of whose roots has a contact on $C_{3-i}$; conversely every such single-square group is a $C_i$-group of the pair, with the same interval, weight and footprint (the footprint of a group of the pair is defined in Section~\ref{sec:two-loops-pins}). Two $C_i$-groups are compatible for the pair if and only if they are compatible for the square $C_i$; a $C_1$-group and a $C_2$-group are compatible if and only if their constituent domains are mutually compatible. (d) If $W$ is a $C_i$-group that is not long on $C_i$, then $\mathcal H_i(W)$ is a cyclic interval of fewer than $n_{C_i}/2$ traversals.
\end{lemma}

\begin{proof}
(a) The one-block partition is admissible. If $\mathfrak P$ and $\mathfrak Q$ are admissible, so is their common refinement, whose blocks are the nonempty sets $P\cap Q$: (A1) holds because every overlap class lies in one block of each; for (A2), if $P\cap Q$ is long on $C_i$, then so are $P\supseteq P\cap Q$ and $Q\supseteq P\cap Q$ (the union of hulls is monotone), so both contain every component with a contact on $C_i$, and so does $P\cap Q$. There are finitely many components, so the common refinement of all admissible partitions is admissible and is the finest. If a block $P$ of the finest partition is long on neither square and contains two overlap classes, replace $P$ by its overlap classes: (A1) still holds, the new blocks are long on neither square (they are subsets of $P$), and the other blocks are unchanged, so the new partition is admissible and strictly finer, a contradiction. If no component has a contact on $C_{3-i}$, condition (A2) on $C_{3-i}$ is void; if no overlap class is long the overlap classes form an admissible partition, hence the finest, and these are the classes of Section~\ref{sec:joint-extraction} for $C_i$ with the intervals $\mathcal H(K)$ there, which equal the unions of the members' hulls; if some overlap class is long, every admissible partition has a long block containing it, which by (A2) contains every component (each has a contact on $C_i$), so the finest partition is the single class of all components, the whole-contour group of Section~\ref{sec:joint-extraction}. The term of a group: let $\mathfrak P$ be the finest admissible partition of a term $T$ and $W\in\mathfrak P$. Every cluster occurrence of $W$ hits a root of $T$, which is joined to it by an edge and so lies in its component, hence in $W$; so $T_W$ is a term of \eqref{eq:c-cluster-expansion}. Edges are defined pairwise, so the graph of $T_W$ is the subgraph of that of $T$ induced on $W$, and $W$ being a union of components of $T$, the components of $T_W$ are those of $T$ in $W$, with the same contact sets, hulls and long-ness; by (A1) $W$ is a union of overlap classes of $T$, which are the overlap classes of $T_W$. Let $\mathfrak P'$ be an admissible partition of $T_W$ and $\mathfrak P''=\mathfrak P'\cup(\mathfrak P\setminus\{ W\})$. (A1) holds for $\mathfrak P''$, and (A2) for its blocks outside $W$; if $Q\in\mathfrak P'$ is long on $C_i$, then $W\supseteq Q$ is long on $C_i$, so $W$ contains every component of $T$ with a contact on $C_i$ by (A2) for $\mathfrak P$, and all of them lie in $Q$ by (A2) for $\mathfrak P'$. So $\mathfrak P''$ is admissible for $T$, $\mathfrak P$ refines it, and $\mathfrak P'=\{ W\}$. Long-ness and hulls being the same in $T$ and $T_W$, the interval set is the same.

(b) Distinct groups of a term have disjoint intervals: if neither is long on $C_i$, their $C_i$-parts are unions of hulls of components in different blocks, hence of different overlap classes by (A1), and such hulls are disjoint; if one is long on $C_i$, it contains every component with a contact on $C_i$, so the other has $\mathcal H_i=\varnothing$. Domains in different groups are compatible as in the proof of Lemma~\ref{lem:grouped-source} (no edge of the term's graph joins different components). Every component's hull satisfies $H^i_{\gamma}\subseteq\mathcal H_i(W)$ for its group $W$: trivially if $W$ is long on $C_i$, and by definition otherwise. Reversibility: let $\mathfrak B=\{ W_1,\ldots,W_m\}$ be a compatible family, each $W_j$ being the only group of its own term by (a), and let $\mathfrak P$ be the finest admissible partition of the union term, whose components are those of the $W_j$ (compatibility excludes edges between groups). $\mathfrak B$ is admissible for the union term: components of different $W_j$ have hulls in disjoint intervals, so they do not overlap and (A1) holds; if $W_j$ is long on $C_i$, then $\mathcal H_i(W_j)=\mathcal I_{C_i}$, so every other $W_k$ has $\mathcal H_i(W_k)=\varnothing$, no contact on $C_i$, and $W_j$ contains every component with a contact on $C_i$. Hence $\mathfrak P$ refines $\mathfrak B$. The restriction of $\mathfrak P$ to the components of $W_j$ is admissible for the term of $W_j$ alone: its overlap classes are those of the union term inside $W_j$, and a block of $\mathfrak P$ inside $W_j$ that is long on $C_i$ contains every component of the union term with a contact on $C_i$, in particular those of $W_j$. Since $\{ W_j\}$ is the finest admissible partition of that term, $\mathfrak P$ restricted to $W_j$ is the single block $W_j$, so $\mathfrak P=\mathfrak B$, with the same intervals. The rest of the proof of Lemma~\ref{lem:grouped-source} (interleaving count, multiplicities of identical occurrences, absolute convergence) does not involve the intervals and applies verbatim.

(c) Let $W$ be a $C_i$-group of the pair. Its roots are marked and $\mathcal H_{3-i}(W)=\varnothing$, so every root $D$ has $\mathsf T_D=\mathsf T^i_D\neq\varnothing$, and the components, their hulls on $C_i$, the overlap graph and the long-ness on $C_i$ of the term $T_W$ of $W$ are the same for the pair and for the square $C_i$. So $T_W$ is a term of \eqref{eq:c-cluster-expansion} for the square $C_i$ (its roots are marked for $C_i$, and its clusters hit $\mathcal A(W)$ by (a)), no component of which has a contact on $C_{3-i}$; by (a) its finest admissible partition $\{ W\}$ is its grouping of Section~\ref{sec:joint-extraction} for $C_i$, with $\mathcal H(W)=\mathcal H_i(W)$. Conversely, if $W$ is a group of $C_i$ whose roots have no contact on $C_{3-i}$, its term (Definition~\ref{def:groups}) is a term of the pair, with $\mathsf T_D=\mathsf T^i_D\neq\varnothing$ for its roots, and by (a) its pair grouping is its grouping for $C_i$, the single group $W$, with $\mathcal H_i(W)=\mathcal H(W)$ and $\mathcal H_{3-i}(W)=\varnothing$. The weight \eqref{eq:part1-group-weight} and the footprint of Section~\ref{sec:c4-6} depend on the roots, the cluster list and the interval only. Compatibility for the pair of two $C_i$-groups is compatibility of domains and disjointness of the $C_i$-intervals, which is Definition~\ref{def:groups} for $C_i$ (its singleton clause for a whole-contour group being implied by the disjointness, since every group has a nonempty interval); for a $C_1$-group and a $C_2$-group the interval condition is void. (d) $W$ is not long on $C_{3-i}$, since $\mathcal H_{3-i}(W)=\varnothing$, so by (a) it is an overlap class; its members have hulls on $C_i$ only, so they are chained by intersections of these hulls, and $\mathcal H_i(W)$ is a union of chained cyclic intervals with fewer than $n_{C_i}/2$ positions in total, which is a cyclic interval by the argument of Section~\ref{sec:joint-extraction}.
\end{proof}

\subsection{Reference scalars and the joint extraction}\label{sec:two-loops-extraction}
For a $C_i$-group $W$ let $P^{\circ}_W(a)$ be the ordered product of the reference links along the consecutive traversals of $\mathcal H_i(W)$, in the cyclic order of $C_i$ starting from the first traversal of the interval (from the base point of $C_i$ if $W$ is long on $C_i$), as in Section~\ref{sec:joint-extraction}, and define $a_W$ by \eqref{eq:part1-reference-group} with $\chi_i$ in place of $\chi$; $\mathcal H_i(W)$ is a cyclic interval or $\mathcal I_{C_i}$ by Lemma~\ref{lem:two-loop-grouping}(d), and $a_W= a^{(i)}_W$ by Lemma~\ref{lem:two-loop-grouping}(c), since the reference links and the measures $\widehat\nu^{\mathrm f,\circ}_{\xi,D}$ do not involve the observable (Lemma~\ref{lem:probe-measures}). For a cross group put $a_W\coloneqq 0$. With these numbers define, for the pair,
\begin{equation}\label{eq:two-loop-reference-gas}
    S^{\circ}_{\xi,C}=\sum_{\mathfrak B\text{ compatible}}\prod_{W\in\mathfrak B}a_W\;,
\end{equation}
in which a family containing a cross group contributes zero, the restricted sums $S^{\circ,[\mathfrak B]}_{\xi,C}$ over the families compatible with every member of $\mathfrak B$, and the joint defect $R_{\xi}(\mathfrak B)$ by \eqref{eq:part1-joint-defect}.

\begin{lemma}[Schur identity and joint extraction for two squares]\label{lem:two-loop-schur}
(a) For a $C_i$-group $W$ and $h\in G$, \eqref{eq:part1-closed-schur} holds with $\chi_i$ in place of $\chi$. (b) Proposition~\ref{prop:joint-extraction} holds for the pair: \eqref{eq:part1-scalar-extraction} and \eqref{eq:c-one-step-branch-identity} hold with $S^{\circ}_{\xi,C}$ of \eqref{eq:two-loop-reference-gas}, the sums converging absolutely. (c) At the reference data, $R^{\circ}_{\xi}(\mathfrak B)=0$ for every nonempty compatible family $\mathfrak B$ without cross groups.
\end{lemma}

\begin{proof}
(a) The proof of Lemma~\ref{lem:schur} for the square $C_i$: it uses only the invariance of the reference measures under the constant gauge rotations $R_u$, the covariance $P^{\circ}_W(R_ua)= uP^{\circ}_W(a)u^{-1}$ and Schur's lemma for $\pi_i$. (b) The proof of Proposition~\ref{prop:joint-extraction} is the Boolean inversion of \eqref{eq:part1-joint-defect} and an absolutely convergent rearrangement; it uses only $| a_W|\leqslant \beta_W$, which holds for $a_W=0$, the finiteness of $\sum_W\beta_W$ (a group is determined by its root family, cluster list and interval, Lemma~\ref{lem:vacuum-clusters} and Lemma~\ref{lem:c2-reference}) and Lemma~\ref{lem:two-loop-grouping}(b) in place of Lemma~\ref{lem:grouped-source}. (c) Let $\mathcal S$ be a compatible family of $C_1$- and $C_2$-groups. The contacts on $C_i$ of the roots of a $C_i$-group $W$ lie in $\mathcal H_i(W)$ (Lemma~\ref{lem:two-loop-grouping}(b)), and the roots of a $C_{3-i}$-group have no contact on $C_i$. So, as in the proof of Lemma~\ref{lem:schur}, a traversal outside every interval of $\mathcal S$ has the reference link $\id$, and a traversal in $\mathcal H_i(W)$ has at $\mathsf e_{\mathcal A(\mathcal S)}a$ its link at $\mathsf e_{\mathcal A(W)}a$, since it lies in no contact set of a root of another group; the links of $\mathcal H_i(W)$ multiply to $P^{\circ}_W(a)$. Number the $C_i$-groups of $\mathcal S$ in the cyclic order of $C_i$; by cyclicity of the trace the reference Wilson function at $\mathsf e_{\mathcal A(\mathcal S)}a$ is $\prod_{i=1,2}\chi_i\bigl(\prod_{W\in\mathcal S,\ C_i\text{-group}}P^{\circ}_W(a)\bigr)$, the product over the $C_i$-groups in that order (an empty product being $\id$), and a family containing a group long on $C_i$ has no other $C_i$-group. Integrate one group at a time by Fubini, with (a) for the group integrated, the product of the other $P^{\circ}$ of the same square as $h$, and the character of the other square as a constant factor: $A^{\circ}_{\xi}(\mathcal S)=\prod_{W\in\mathcal S}a_W$. Hence $R^{\circ}_{\xi}(\mathfrak B)=\prod_{W\in\mathfrak B}a_W\sum_{\mathcal S\subseteq\mathfrak B}(-1)^{|\mathfrak B\setminus\mathcal S|}=0$ for $\mathfrak B\neq\varnothing$.
\end{proof}

For a family containing a cross group, $R^{\circ}_{\xi}(\mathfrak B)$ need not vanish (for instance $R^{\circ}_{\xi}(\{ W\})= A^{\circ}_{\xi}(\{ W\})$ for a cross group $W$); the defects of such families are bounded directly (Lemmas~\ref{lem:two-loop-defect} and~\ref{lem:two-loop-cross}). In place of the second equality of \eqref{eq:part1-reference-gas}, Lemma~\ref{lem:two-loop-grouping}(b) at the reference data (its proof uses only the field-independent geometry, as in Section~\ref{sec:joint-extraction}) and the Fubini step of (c) give
\[
    \frac{\int f^{\circ}_{\xi} K^{\circ}_{\xi;0}\,\dd a}{I^{\circ}_{\xi;0}}= S^{\circ}_{\xi,C}+\sum_{\mathfrak B\text{ compatible, with a cross group}}A^{\circ}_{\xi}(\mathfrak B)\;.
\]
The choice $a_W\coloneqq 0$ is made for every cross group, and no reduction of the reference integrals of cross families to products of numbers is used. The pair scalar is identified with the single-square scalars by Lemma~\ref{lem:two-loop-cross}(b).

\subsection{The defects of a family with cross groups}\label{sec:two-loops-defects}
For a $C_i$-group $W$ the closed contraction $c_W(U_1;h)$ of \eqref{eq:c4-6-2}, its defect $d_W$ and $\mathfrak d_W(U_1)$ of \eqref{eq:c4-6-3} are formed with $\chi_i$ and the interval $\mathcal H_i(W)$; by Lemma~\ref{lem:two-loop-grouping}(c) they are the single-square quantities of $C_i$. For every group, cross groups included, $M(W)$ is the positive majorant of Section~\ref{sec:c4-6-2}: the product of the root weights $\rho\,m(D)$, the cluster weights $M_v(\mathfrak c_j)$ and $1/k(W)!$, with $\rho$ of Lemma~\ref{lem:two-loop-source}; for $C_i$-groups this only enlarges the single-square majorant, so every single-square bound by $M(W)$ still holds. Put $K_{\chi}=\max(K_{\chi_1},K_{\chi_2})$ with the constants of Lemma~\ref{lem:c4-6-1}(a).

\begin{lemma}[sources with two closed words]\label{lem:two-loop-source}
Let $f(a)=\prod_{i=1,2}\dim(\pi_i)^{-1}\operatorname{tr}Y_i(a)$, where $Y_i(a)$ is an ordered product of the matrices $\pi_i(U_{\xi}(a)(r))^{\pm1}$ of a set of reconstructed contour links of $C_i$, each traversal used at most once, and of fixed unitary matrices of $V_{\pi_i}$. Then \eqref{eq:c4-4-3}, \eqref{eq:c4-10-8}--\eqref{eq:c4-10-9}, the undifferentiated bound \eqref{eq:c4-7-6} and the last statement of Lemma~\ref{lem:c4-7-2}, and Lemma~\ref{lem:c4-4-2} hold for $f$, with $C_{\chi}$ enlarged once to a constant $C^{\mathrm{src}}_{\chi_1\chi_2}$ of $(\chi_1,\chi_2)$, so that $u= C^{\mathrm{src}}_{\chi_1\chi_2}g_0(1+r)$, and $g_0$ so small that $u\leqslant u_*$ and $v\leqslant 1$; in particular $|\mathfrak Q_{\xi}(\mathcal A;f)|\leqslant \prod_{D\in\mathcal A}\rho\,m(D)$ for every nonempty compatible root family $\mathcal A$, with $\rho$ of \eqref{eq:c4-4-8} formed with this $u$.
\end{lemma}

\begin{proof}
$f$ is the normalised trace of $Y_1(a)\otimes Y_2(a)$ on $V_{\pi_1}\otimes V_{\pi_2}$, and $| f|\leqslant 1$; since $\operatorname{tr}(A\otimes B)=\operatorname{tr}A\operatorname{tr}B$, the normalised trace on $V_{\pi_1}\otimes V_{\pi_2}$ satisfies $|\dim^{-1}\operatorname{tr}(M_1\cdots M_k)|\leqslant \prod_j\| M_j\|_{\mathrm{op}}$. The proof of \eqref{eq:c4-10-8} uses only this inequality, the unitarity of the untouched factors, the row locality and the fixed derivative constants of the rows. For the row locality: $V(C_1)\cap V(C_2)=\varnothing$, so the bond sets of $C_1$ and $C_2$ are disjoint, each bond carries at most one traversal of $C_1\cup C_2$, and the number of reconstructed links a free label affects is the geometric count of bonds $r$ with the label in $\mathcal F_r$, at most $K$ as for one square, each traversal being used at most once in the two words together; the row input count, which fixes the derivative order of \eqref{eq:c4-10-4}, is unchanged. Then the product rule distributes the derivatives over the links of both words, with at most $K^q$ row assignments, and every untouched factor, a reconstructed link or a fixed matrix, is unitary in $V_{\pi_1}\otimes V_{\pi_2}$. The remainder bounds \eqref{eq:c4-10-9} follow by the same Taylor argument. The proofs of \eqref{eq:c4-7-6}, of the last statement of Lemma~\ref{lem:c4-7-2} and of Lemma~\ref{lem:c4-4-2} use, besides these bounds, only properties of the selected measures (the reversal of a root with empty ordinary selection, the Cauchy difference in $\zeta_j$, the majorants $M_{\ell}$ and \eqref{eq:c4-4-7}) and the vanishing of $\Delta_{\mathsf J}f$ when a selected label does not influence the words; none of them uses the form of $f$ otherwise.
\end{proof}

\begin{lemma}[defect factorisation for two squares]\label{lem:two-loop-defect}
On a branch covered for $C_1\sqcup C_2$ (Section~\ref{sec:two-loops-setting}), for every compatible family $\mathfrak B$ with cross groups $\mathfrak B_{\times}\subseteq\mathfrak B$,
\begin{equation}\label{eq:two-loop-defect}
    | R_{\xi}(\mathfrak B)|\leqslant \prod_{W\in\mathfrak B\setminus\mathfrak B_{\times}}K_{\chi}\mathfrak d_W(U_1)\prod_{W\in\mathfrak B_{\times}}M(W)\;.
\end{equation}
\end{lemma}

\begin{proof}
As in the proof of Lemma~\ref{lem:c4-6-1}(b), a traversal of $\mathcal H(W)$, $W\in\mathfrak B$, has no contact with a root of another group of $\mathfrak B$ (the intervals are disjoint and contain the contacts), so its reconstructed link depends only on the root variables of $W$, and a traversal outside every interval has its background link $U_1$ (Lemma~\ref{lem:covered}, Section~\ref{sec:c4-8-1}). For a $C_i$-group $W\in\mathfrak B$ let $X_W= w_{\xi}(W)\int\pi_i(T_WP_W^{-1})\,\dd\bigotimes_{D\in\mathcal A(W)}\widehat\nu^{\mathrm f}_{\xi,D}$ and $\Upsilon_W= X_W-a_W\id$, with $T_W$, $P_W$ of Section~\ref{sec:c4-6-1}. Since $a_W=0$ for cross groups, the term of $\mathcal S$ in \eqref{eq:part1-joint-defect} vanishes unless $\mathcal S\supseteq\mathfrak B_{\times}$. Fix the root variables $a_{\times}$ of the cross groups. For $i=1,2$ let $Y_i$ be the ordered product, in the cyclic order of $C_i$ starting at the first traversal of the interval of a $C_i$-group of $\mathfrak B$ if there is one and at the base point of $C_i$ otherwise (the trace is cyclic; as in the proof of Lemma~\ref{lem:c4-6-1}(b)), of the following factors: for a $C_i$-group $W\in\mathfrak B$, the matrix $\Upsilon_W\pi_i(P_W)$ at the position of $\mathcal H_i(W)$; for a traversal in $\mathcal H_i(W)$ of a cross group $W$, its reconstructed link at $a_{\times}$; for every other traversal, its background link. Expanding each $\Upsilon_W= X_W-a_W\id$ and using Fubini over the product measures of the $C_1$- and $C_2$-groups, exactly as in the proof of Lemma~\ref{lem:c4-6-1}(b), the term of $\mathcal S\supseteq\mathfrak B_{\times}$ is reproduced, because $f_{\xi}$ at $\mathsf e_{\mathcal A(\mathcal S)}a$ is the product of the two normalised traces and each $X_W$ enters exactly one of them; hence
\begin{equation*}
    R_{\xi}(\mathfrak B)=\Bigl(\prod_{W\in\mathfrak B_{\times}}w_{\xi}(W)\Bigr)\int\prod_{i=1,2}\frac{\operatorname{tr}Y_i(a_{\times})}{\dim\pi_i}\,\dd\bigotimes_{W\in\mathfrak B_{\times}}\bigotimes_{D\in\mathcal A(W)}\widehat\nu^{\mathrm f}_{\xi,D}(a_{\times})\;.
\end{equation*}
The integrand is multilinear in the matrices $\Upsilon_W$, $W\in\mathfrak B\setminus\mathfrak B_{\times}$. Write $\Upsilon_W=\|\Upsilon_W\|_{\mathrm{op}}T_W'$ with $\| T'_W\|_{\mathrm{op}}\leqslant 1$, and each contraction $T'= V| T'|$ (polar decomposition, $V$ unitary) as $T'=\frac12(Ve^{iA}+Ve^{-iA})$ with $A=\arccos| T'|$, an average of two unitary matrices. The integrand is then $\prod_W\|\Upsilon_W\|_{\mathrm{op}}$ times an average of $2^{|\mathfrak B\setminus\mathfrak B_{\times}|}$ functions of the form of Lemma~\ref{lem:two-loop-source}, depending on $a_{\times}$ only through the links of the cross groups' traversals. If $\mathfrak B_{\times}\neq\varnothing$, that lemma applied to the compatible root family $\mathcal A(\mathfrak B_{\times})$ bounds each such integral by $\prod_{D\in\mathcal A(\mathfrak B_{\times})}\rho\,m(D)$ in modulus; if $\mathfrak B_{\times}=\varnothing$, there is no integral, and each of the unitary words has normalised traces of modulus at most one (for $\mathfrak B=\varnothing$ the bound is $| f_{\xi}(0)|\leqslant 1$). Further $| w_{\xi}(W)|\leqslant \prod_jM_v(\mathfrak c_j)/k(W)!$: by \eqref{eq:part1-vacuum-cluster} the cluster weight is $\varphi^T/m!$ times the product of the activities, $|\varphi^T|$ is at most the number of spanning trees of the incompatibility graph (proof of Lemma~\ref{lem:vacuum-clusters}, Section~\ref{app:kp}), and the activities are bounded by $m(D)$ by \eqref{eq:fixed-kernel-raw-activity-payment} of Lemma~\ref{lem:activity-bounds}. Together with $\|\Upsilon_W\|_{\mathrm{op}}\leqslant K_{\chi_i}\mathfrak d_W(U_1)\leqslant K_{\chi}\mathfrak d_W(U_1)$ for a $C_i$-group (Lemma~\ref{lem:c4-6-1}(a) for the square $C_i$) this is \eqref{eq:two-loop-defect}; the bound holds with $K_{\chi_i}$ on the $C_i$-groups.
\end{proof}

\subsection{Pins for the groups of two squares}\label{sec:two-loops-pins}
For a group $W$ of the pair, $n(W)$, $\mathfrak s(W)= n(W)+|\mathcal H(W)|/(LM)$, the footprint $F(W)$, $X(W)$, $\pi_1(\mathcal H(W))$ and the envelope $\mathfrak E(W)$ are defined as in Section~\ref{sec:c4-6} and Lemma~\ref{lem:c4-6-env}, with $\mathcal H(W)=\mathcal H_1(W)\sqcup\mathcal H_2(W)$ and the traversals of both squares; for components, $t_{\gamma}$, $n_{\gamma}$, $X_{\gamma}$, $F(\gamma)$ and $\mathfrak s(\gamma)= n_{\gamma}+| H_{\gamma}|/(LM)$ are those of Section~\ref{sec:c4-8-2} with $H_{\gamma}= H^1_{\gamma}\sqcup H^2_{\gamma}$, and $\sigma_K=\sum_{\gamma\in K}\mathfrak s(\gamma)$. Put $A'_E=17(2R_{\mathrm e}+1)^4$.

\begin{lemma}[pins for two squares]\label{lem:two-loop-pins}
Under the hypotheses of Lemma~\ref{lem:c4-6-2}, with its constant $C$ enlarged by a fixed factor:
(a) Lemma~\ref{lem:c4-6-env} holds for every group $W$ of the pair with $16n(W)+|\mathcal H(W)|/(LM)$ and $17\mathfrak s(W)$ in place of the bounds $8+|\mathcal H(W)|/(LM)$ and $9\mathfrak s(W)$ of its part (a), $A'_E$ in place of $A_F$ and $A_E$, and $\operatorname{dist}(\Delta,C)$ the distance to $V(C_1)\cup V(C_2)$; if $W$ is a cross group, $\pi_1(\mathcal H(W))\subseteq\mathfrak E(W)$ contains a cube containing a site of $V(C_1)$ and a cube containing a site of $V(C_2)$.
(b) Lemma~\ref{lem:c4-8-2} holds for the components of the pair, with $| H^i_{\gamma}|\leqslant A_{\mathrm h}LMn_{\gamma}$ for $i=1,2$, the hull pin \eqref{eq:c4-8-5} for $a\in\mathcal I_{C_i}$ taken with the contour distance $d_{C_i}$ of $C_i$, and the sum over $a'\in\mathcal I_C$ in \eqref{eq:c4-8-6}.
(c) Lemma~\ref{lem:c4-8-3}(a) holds for the overlap classes of the pair. A class long on $C_i$ satisfies $\sigma_K\geqslant n_{C_i}/(2LM)$ and $\sum_{K\text{ long on }C_i}t_Ke^{34\eta\sigma_K}\leqslant C\rho n_{C_i}e^{-15\eta n_{C_i}/(LM)}$. The groups long on some square have total weighted mass $\sum_{W\text{ long}}M(W)e^{4\eta\mathfrak s(W)}\leqslant 2C\rho\bigl(n_{C_1}e^{-10\eta n_{C_1}/(LM)}+n_{C_2}e^{-10\eta n_{C_2}/(LM)}\bigr)\leqslant C'\rho$ with $C'=2CLM/(5e\eta)$, uniformly in $n_1,n_2$.
(d) The pins \eqref{eq:c4-6-5} hold for the groups of the pair, with $\mathcal I_C$ the traversals of both squares; for every group $W$ of the pair, cross groups included, $\sum_{W'\not\sim W}| a_{W'}| e^{3\eta\mathfrak s(W')}\leqslant \frac18\eta\mathfrak s(W)$, the sum over the groups of the pair; and Lemma~\ref{lem:c4-6-3} and Lemma~\ref{lem:c4-9-2} hold for the reference gas \eqref{eq:two-loop-reference-gas}, with $C\rho n_C$ in \eqref{eq:c4-9-4}.
\end{lemma}

\begin{proof}
(a) Count. On each square, $\mathcal H_i(W)$ is empty, or $\mathcal I_{C_i}$, or, if $W$ is not long on $C_i$, a nonempty union of member hulls with fewer than $n_{C_i}/2$ positions in total; in the second case each member hull is a cyclic interval (none is the whole of $\mathcal I_{C_i}$), each maximal cyclic interval of the union has fewer than $n_{C_i}/2$ positions because the union does, and there are at most $| W|$ of them, where $| W|$, the number of components of $W$, is at most $n(W)=\sum_{\gamma} n_{\gamma}$ since $n_{\gamma}\geqslant 1$. As in the proof of Lemma~\ref{lem:c4-6-env}(a), a cyclic interval of fewer than $n_{C_i}/2$ traversals meets each side in one run and lies in at most $8+t/(LM)$ cubes for $t$ traversals, and so does $\mathcal I_{C_i}$; summing over at most $2n(W)$ intervals, $|\pi_1(\mathcal H(W))|\leqslant 16n(W)+|\mathcal H(W)|/(LM)$, so $| X(W)\cup\pi_1(\mathcal H(W))|\leqslant 17\mathfrak s(W)$ and $| F(W)|\leqslant |\mathfrak E(W)|\leqslant A'_E\mathfrak s(W)$. Connectivity. For a component $\gamma$, $X_{\gamma}$ is connected (proof of Lemma~\ref{lem:c4-6-env}(b)); $\pi_1(H^i_{\gamma})$ is connected, $H^i_{\gamma}$ being a cyclic interval or $\mathcal I_{C_i}$, and it contains the cube of the initial vertex of a contact of $\gamma$ on $C_i$, which is at index distance at most one from a cube of a root of $\gamma$; so $X_{\gamma}\cup\pi_1(H_{\gamma})$ is connected. Two overlapping components share a traversal, hence the cube of its initial vertex. So $X(W)\cup\pi_1(\mathcal H(W))$ is connected if $W$ is an overlap class, which by Lemma~\ref{lem:two-loop-grouping}(a) covers every group long on neither square. Let $W$ be long on $C_i$. Each overlap class in $W$ has a component with a contact on $C_i$, or $W$ is also long on $C_{3-i}$ and the class has a contact on $C_{3-i}$: otherwise splitting this class off $W$ gives an admissible partition (the remainder is nonempty, since it contains the components of $W$ with a contact on $C_i$; it keeps every component with a contact on a square on which $W$ is long and keeps the same union of hulls there, the class is long on no square, being a subset of $W$ without contacts on those squares and with the union of its hulls on the other square inside that of $W$), contradicting the finest choice. If $W$ is long on $C_i$ only, every class in $W$ meets the connected set $\pi_1(\mathcal I_{C_i})$, so the union is connected. If $W$ is long on both squares and no class in $W$ has contacts on both, then splitting $W$ into its classes with contacts on $C_1$ and those with contacts on $C_2$ is admissible (each part contains every component with a contact on the square on which it is long, and has no contact on the other square), a contradiction; so some class of $W$ meets both $\pi_1(\mathcal I_{C_1})$ and $\pi_1(\mathcal I_{C_2})$, and the union is connected. Hence $\mathfrak E(W)$ is connected. The statements (c) and (d) of Lemma~\ref{lem:c4-6-env} then follow by its proof. For a cross group both $\mathcal H_1(W)$ and $\mathcal H_2(W)$ are nonempty, and $\pi_1(\mathcal H(W))$ contains the cubes of their vertices.

(b) The proof of Lemma~\ref{lem:c4-8-2}. Part (a) there does not involve hulls. In (b) the contacts of $\gamma$ on $C_i$ lie within $2L$ of $X_{\gamma}$, and the square geometry of $C_i$ (Section~\ref{sec:extent-weights}) gives $| H^i_{\gamma}|\leqslant A_{\mathrm h}LMn_{\gamma}$ and $d_{C_i}(a,a')\leqslant A_{\mathrm h}LMn_{\gamma}$ for $a\in H^i_{\gamma}$ and a contact $a'$ of $\gamma$ on $C_i$; the contour sum is over $\mathcal I_{C_i}$. In \eqref{eq:c4-8-6} the sum over the traversals of each square is bounded as there, and there are two squares. For part (c) of Lemma~\ref{lem:c4-8-2}: for the pair $\mathfrak s(\gamma)\leqslant (1+2A_{\mathrm h})n_{\gamma}\leqslant 2A_{\mathrm s}n_{\gamma}$, so $100\eta\mathfrak s(\gamma)\leqslant 200\eta A_{\mathrm s}n_{\gamma}<\varpi n_{\gamma}/40$ by $1000\eta A_{\mathrm s}<\varpi/8$, and the pins with the weight $e^{100\eta\mathfrak s(\gamma)}$ are dominated by those of (b) with the weight $e^{\varpi n_{\gamma}/4}$, with a margin $e^{\varpi n_{\gamma}/8}$; $\eta$ is not changed. In the footprint pin, at most $CLM$ traversals of each square have stencils meeting a given cube, so at most $2CLM$ in all, and the footprint pin holds with the constant doubled.

(c) The proof of Lemma~\ref{lem:c4-8-3}(a) applies with $H_{\gamma}= H^1_{\gamma}\sqcup H^2_{\gamma}$: the components whose hulls meet $H_{\gamma}$ have weighted mass at most $C\rho| H_{\gamma}|$ by the hull pins of (b) at the traversals of $H_{\gamma}$, and an overlap class long on neither square has the interval $\bigcup_{\gamma} H_{\gamma}$, of size at most $\sum_{\gamma}| H_{\gamma}|\leqslant LM\sigma_K$. A class long on $C_i$ has $\sigma_K\geqslant \sum_{\gamma}| H^i_{\gamma}|/(LM)\geqslant n_{C_i}/(2LM)$; anchoring it at a traversal of $\bigcup_{\gamma} H^i_{\gamma}$ ($n_{C_i}$ choices) and using the class pin with the weight $e^{64\eta\sigma_K}$ gives the stated sum, the exponent $30\eta\sigma_K$ being at least $15\eta n_{C_i}/(LM)$. Long groups: if no overlap class of a term is long, the overlap classes are admissible and no group is long. Let $W$ be long; it contains an overlap class $K_0$ long on some square $C_{i_0}$: otherwise replacing $W$ by its overlap classes keeps (A1), the new blocks are long on no square, and the other blocks are unchanged, so the new partition is admissible, and it is strictly finer unless $W$ is a single class, which is then not long; the first contradicts the finest choice, the second the assumption that $W$ is long. $W$ is long on $C_{i_0}$, and by the splitting argument of (a) it is the union of the overlap classes with a contact on a square on which it is long. As in Lemma~\ref{lem:c4-8-3}(c), the sum of $M(W)$ over the groups with a given set of components is at most $\prod_{\gamma} t_{\gamma}$, and $\mathfrak s(W)\leqslant \sum_{\gamma\in W}\mathfrak s(\gamma)+\sum_{i:\,W\text{ long on }C_i}n_{C_i}/(LM)$. If $W$ is long on $C_{i_0}$ only, the classes other than $K_0$ are anchored at traversals of $C_{i_0}$ and sum to at most $e^{C\rho n_{C_{i_0}}}$ with their weights $e^{64\eta\sigma}$, the interval weight is $e^{4\eta n_{C_{i_0}}/(LM)}$, and $K_0$ contributes $C\rho n_{C_{i_0}}e^{-15\eta n_{C_{i_0}}/(LM)}$; with $C\rho LM\leqslant \eta$ the product is at most $C\rho n_{C_{i_0}}e^{-10\eta n_{C_{i_0}}/(LM)}$. If $W$ is long on both squares, let $j\neq i_0$; since $\sum_{\gamma\in W}| H^j_{\gamma}|\geqslant n_{C_j}/2$ and $| H^j_{\gamma}|\leqslant LM\mathfrak s(\gamma)$, $1\leqslant e^{-15\eta n_{C_j}/(LM)}\prod_{\gamma\in W}e^{30\eta\mathfrak s(\gamma)}$; inserting this factor, the classes other than $K_0$ are anchored at traversals of $C_1$ or $C_2$ and sum to at most $e^{C\rho n_C}$ with the weights $e^{64\eta\sigma}$, the interval weight is $e^{4\eta n_C/(LM)}$, and $K_0$ contributes $C\rho n_{C_{i_0}}e^{-15\eta n_{C_{i_0}}/(LM)}$ with the weight $e^{34\eta\sigma_{K_0}}$; the product is at most $C\rho n_{C_{i_0}}e^{-10\eta n_C/(LM)}\leqslant C\rho n_{C_{i_0}}e^{-10\eta n_{C_{i_0}}/(LM)}$. The groups long on $C_i$ only thus contribute at most $C\rho n_{C_i}e^{-10\eta n_{C_i}/(LM)}$, and the groups long on both squares at most that for each $i_0= i$; summing over the four pairs (set of long squares, $i_0$) gives the bound, and $xe^{-10\eta x/(LM)}\leqslant LM/(10e\eta)$ for $x= n_{C_1},n_{C_2}$ gives $C'=2CLM/(5e\eta)$.

(d) The pins over groups long on neither square are the class pins of (c), as in Lemma~\ref{lem:c4-8-3}: such a group is an overlap class $K$ (Lemma~\ref{lem:two-loop-grouping}(a)), the sum of $M(W)$ over the groups with class $K$ is at most $t_K$, $\mathfrak s(W)\leqslant \sigma_K$ and $F(W)\subseteq\bigcup_{\gamma\in K}F(\gamma)$; the long groups are bounded by their total mass (c), so \eqref{eq:c4-6-5} holds with the constant $C+C'$. For the second statement, a group $W'\not\sim W$ has a footprint cube in $\mathfrak N_1(F(W))$ or a traversal in $\mathcal H(W)$ by (a) and Lemma~\ref{lem:c4-6-env}(d), so the pins give at most $C\rho(3^4| F(W)|+|\mathcal H(W)|)\leqslant C\rho(3^4A'_E+LM)\mathfrak s(W)$, which is at most $\eta\mathfrak s(W)/8$ under $C\rho LM\leqslant \eta/8$ with the constant $C$ of (c) replaced by $C(1+2LM/(5e\eta))(1+3^4A'_E/(LM))$; this factor, together with the doubling in (b), is the enlargement of the lemma, fixed with $\eta$ before the coupling (and $e^{3\eta\mathfrak s}\leqslant e^{4\eta\mathfrak s}$, $| a_{W'}|\leqslant M(W')$). The proofs of Lemmas~\ref{lem:c4-6-3} and~\ref{lem:c4-9-2} then apply verbatim; the species of the reference gas are the groups with $a_W\neq0$, all of them $C_1$- or $C_2$-groups, and the exclusion ratio \eqref{eq:c4-6-7} holds for every compatible family $\mathfrak B$ of groups of the pair, cross groups included, because the criterion was proved for every $W$.
\end{proof}

\subsection{The cross terms}\label{sec:two-loops-cross}
Covered and far histories are those of Section~\ref{sec:two-loops-setting}. A single-square group $W$ of $C_i$ touches $C_{3-i}$ if some root $D$ of $W$ has $\mathsf T_D\cap\mathcal I_{C_{3-i}}\neq\varnothing$, with the contact set $\mathsf T_D$ of the pair. Put $c_{\times}=2\eta/A'_E$.

\begin{lemma}[the cross terms]\label{lem:two-loop-cross}
Under the hypotheses of Lemma~\ref{lem:two-loop-pins} and, for (c), those of Lemma~\ref{lem:c4-6-4} for each square, with a constant $C_{\times}$ fixed before the coupling and independent of $n_1$, $n_2$ and the positions of the squares, and with $C_{\chi}=\max_iC_{\chi_i}$ and $c=\min_ic_i$ formed from the constants $C_{\chi_i}$, $c_i$ of \eqref{eq:c4-6-8} for the square $C_i$:
(a) On every branch $\xi$, for every cross group $W$ of the pair and every $a\in\mathcal H(W)$, $\mathfrak s(W)>\delta_a/(A'_ELM)$; and $\sum_{W\text{ cross}}M(W)e^{\mathfrak a_W/8}\leqslant C_{\times}\rho\Delta_{\times}$.
(b) Let $\mathscr L$ and $\mathscr L_i$ be the connected-cluster logarithms of the reference gas \eqref{eq:two-loop-reference-gas} of the pair and of the reference gas of the single square $C_i$ (Lemma~\ref{lem:c4-9-2}), and $E^{\times}_{\xi}=\mathscr L-\mathscr L_1-\mathscr L_2$. Then, on every branch $\xi$,
\begin{equation}\label{eq:two-loop-factorisation}
    S^{\circ}_{\xi,C}= S^{\circ}_{\xi,C_1}S^{\circ}_{\xi,C_2}e^{E^{\times}_{\xi}}\;,\qquad| E^{\times}_{\xi}|\leqslant C_{\times}\rho\Delta_{\times}\;;
\end{equation}
in particular $| S^{\circ}_{\xi,C}|\leqslant e^{C_{\times}\rho\Delta_{\times}}$ on every branch.
(c) On covered far histories, with $\mathscr R_{\xi}$ the family sum \eqref{eq:c4-1-1} of the pair,
\begin{equation}\label{eq:two-loop-family-sum}
    |\mathscr R_{\xi}|\leqslant | S^{\circ}_{\xi,C}|\Bigl[\exp\Bigl(C_{\chi} n_C\rho\bigl(\varepsilon_1+e^{-cb/(LM)}\bigr)+C_{\times}\rho\Delta_{\times}\Bigr)-1\Bigr]\;.
\end{equation}
\end{lemma}

\begin{proof}
Distance from an envelope. Let $\mathfrak X$ be a finite set of cubes, connected under adjacency, with $|\mathfrak X|\leqslant k$. As in the proof of Lemma~\ref{lem:c4-6-env}(c), any two sites in cubes of $\mathfrak X$ are at periodic fine $\ell^{\infty}$ distance less than $kLM$. Hence if $\mathfrak X$ contains the cube of the initial vertex of a traversal $a\in\mathcal I_{C_i}$ and a cube containing a site of $V(C_{3-i})$, then $\delta_a< kLM$.

(a) For a cross group $W$, $\mathfrak E(W)$ is connected, has at most $A'_E\mathfrak s(W)$ cubes and contains the cubes of the initial vertices of the traversals of $\mathcal H(W)$ and a cube containing a site of $V(C_1)$ and one containing a site of $V(C_2)$ (Lemma~\ref{lem:two-loop-pins}(a)); so $\delta_a< A'_ELM\mathfrak s(W)$ for every $a\in\mathcal H(W)$. Every group has a nonempty interval, so, pinning a cross group at a traversal of its interval and using $\mathfrak a_W/8\leqslant \eta\mathfrak s(W)$, $e^{-3\eta\mathfrak s(W)}\leqslant e^{-3\eta\delta_a/(A'_ELM)}\leqslant e^{-c_{\times}\delta_a/(LM)}$ and the pin \eqref{eq:c4-6-5} of the pair (Lemma~\ref{lem:two-loop-pins}(d)),
\[
    \sum_{W\text{ cross}}M(W)e^{\mathfrak a_W/8}\leqslant \sum_{a\in\mathcal I_C}e^{-c_{\times}\delta_a/(LM)}\sum_{W\,|\,a\in\mathcal H(W)}M(W)e^{4\eta\mathfrak s(W)}\leqslant C\rho\Delta_{\times}\;.
\]

The pins of Lemma~\ref{lem:two-loop-pins} and \eqref{eq:c4-6-5}, the criterion \eqref{eq:two-loop-weighted-kp} below and Lemma~\ref{lem:c4-9-2} use only $| a_W|\leqslant M(W)$ (Lemma~\ref{lem:c4-4-2} at the reference data, Lemma~\ref{lem:c2-reference}), the smallness of $\rho$ and the field-independent geometry; coverage and farness enter only Lemmas~\ref{lem:two-loop-defect} and~\ref{lem:c4-6-4}, so (a) and (b) hold on every branch.

(b) The three gases. The species of the pair gas with nonzero activity are its $C_1$- and $C_2$-groups; the species of the gas of $C_i$ are the single-square groups of $C_i$, with the activities $a^{(i)}_W$. By Lemma~\ref{lem:two-loop-grouping}(c) and the equality $a_W= a^{(i)}_W$ of Section~\ref{sec:two-loops-extraction}, the $C_i$-groups of the pair are the single-square groups of $C_i$ that do not touch $C_{3-i}$, with the same activities, and the compatibility relation of the pair restricted to them is that of the square $C_i$. Each gas satisfies the Kotecký--Preiss criterion in the following weighted form: for every species $W$ of that gas,
\begin{equation}\label{eq:two-loop-weighted-kp}
    \sum_{W'\not\sim W}| a_{W'}| e^{3\eta\mathfrak s(W')}\leqslant \frac18\eta\mathfrak s(W)\;,
\end{equation}
for the pair by Lemma~\ref{lem:two-loop-pins}(d), and for the square $C_i$ by the proof of \eqref{eq:c4-6-6}, which bounds $e^{\mathfrak a_{W'}}$ by $e^{4\eta\mathfrak s(W')}$ and so bounds $e^{3\eta\mathfrak s(W')}$ in the same way. This is the criterion of Section~\ref{app:kp} for the weights $\widetilde z_W=| a_W| e^{2\eta\mathfrak s(W)}$ and sizes $\eta\mathfrak s(W)$, so the rooted-tree argument there gives, for every species $W$ of each gas,
\begin{equation}\label{eq:two-loop-pinned-cluster}
    \sum_{\mathfrak c}m_W(\mathfrak c)|\Psi(\mathfrak c)| e^{2\eta\sum_{W'\in\mathfrak c}\mathfrak s(W')}\leqslant | a_W| e^{3\eta\mathfrak s(W)}\;,
\end{equation}
the sum over the cluster tuples $\mathfrak c$ of that gas, $\Psi(\mathfrak c)=\frac{\varphi^T(\mathfrak c)}{m!}\prod_{W'\in\mathfrak c}a_{W'}$, $m_W(\mathfrak c)$ the multiplicity of $W$ in $\mathfrak c$. The logarithms $\mathscr L$, $\mathscr L_i$ are the absolutely convergent series $\sum_{\mathfrak c}\Psi(\mathfrak c)$ of the three gases, $S^{\circ}_{\xi,C}= e^{\mathscr L}$ and $S^{\circ}_{\xi,C_i}= e^{\mathscr L_i}$ (Lemma~\ref{lem:c4-9-2}, for the pair by Lemma~\ref{lem:two-loop-pins}(d)), which gives the identity in \eqref{eq:two-loop-factorisation}.

Cancellation. A cluster tuple of the pair gas all of whose members are $C_i$-groups is a cluster tuple of the gas of $C_i$ with the same $\Psi$: the Ursell coefficient depends only on the incompatibility graph of the tuple, which is the same for both gases, and the activities agree. Conversely a tuple of the gas of $C_i$ none of whose members touches $C_{3-i}$ is such a tuple of the pair gas. Hence $E^{\times}_{\xi}=\sum_{\mathfrak c\in\mathcal X_0}\Psi(\mathfrak c)-\sum_{i=1,2}\sum_{\mathfrak c\in\mathcal X_i}\Psi(\mathfrak c)$, where $\mathcal X_0$ is the set of tuples of the pair gas containing a $C_1$-group and a $C_2$-group, and $\mathcal X_i$ the set of tuples of the gas of $C_i$ containing a group that touches $C_{3-i}$; the rearrangement is justified by absolute convergence.

Extent of the bridging tuples. Let $\mathfrak c\in\mathcal X_0\cup\mathcal X_1\cup\mathcal X_2$, let $W\in\mathfrak c$ and $a\in\mathcal H(W)$, $a\in\mathcal I_{C_i}$. The union $\mathfrak X$ of the envelopes of the members of $\mathfrak c$ is connected: each envelope is connected (Lemma~\ref{lem:two-loop-pins}(a) for the pair, Lemma~\ref{lem:c4-6-env}(b) for a single square), the incompatibility graph of a tuple with $\varphi^T(\mathfrak c)\neq0$ is connected, and two incompatible members have footprints at index distance at most one or a common traversal in their intervals (Lemma~\ref{lem:c4-6-env}(d), for the pair by Lemma~\ref{lem:two-loop-pins}(a)), the footprints and the cubes of the traversal vertices lying in the envelopes. $\mathfrak X$ contains the cube of the initial vertex of $a$, and a cube containing a site of $V(C_{3-i})$: for $\mathfrak c\in\mathcal X_0$, the cube of a traversal vertex of the interval of a member that is a $C_{3-i}$-group; for $\mathfrak c\in\mathcal X_j$ (then $i= j$), the cube of the initial vertex of a contact $n\in\mathsf T_D\cap\mathcal I_{C_{3-i}}$ of a root $D$ of a member $W_0$ touching $C_{3-i}$, which is within $2L$ of $D$ (proof of Lemma~\ref{lem:c4-6-env}: the stencil bond of the contact has its endpoints in $D$), hence at index distance at most one from a cube of $X(W_0)$ and in $\mathfrak E(W_0)$, since $R_{\mathrm e}\geqslant 1$. With $|\mathfrak X|\leqslant A'_E\sum_{W'\in\mathfrak c}\mathfrak s(W')$ ($A_E\leqslant A'_E$), the distance bound gives $\sum_{W'\in\mathfrak c}\mathfrak s(W')>\delta_a/(A'_ELM)$.

Bound. Every member of a tuple has a nonempty interval. Pinning each bridging tuple at a member $W$ and a traversal $a\in\mathcal H(W)$, by \eqref{eq:two-loop-pinned-cluster} and the extent bound,
\[
    \sum_{\mathfrak c\in\mathcal X_0}|\Psi(\mathfrak c)|\leqslant \sum_{a\in\mathcal I_C}e^{-2\eta\delta_a/(A'_ELM)}\sum_{W\,|\,a\in\mathcal H(W)}\sum_{\mathfrak c}m_W(\mathfrak c)|\Psi(\mathfrak c)| e^{2\eta\sum_{W'\in\mathfrak c}\mathfrak s(W')}\leqslant \sum_{a\in\mathcal I_C}e^{-c_{\times}\delta_a/(LM)}\sum_{W\,|\,a\in\mathcal H(W)}M(W)e^{3\eta\mathfrak s(W)}\;,
\]
the inner sums over the species of the pair gas, and the pin \eqref{eq:c4-6-5} of the pair bounds the last sum by $C\rho\Delta_{\times}$. The same computation in the gas of $C_i$, with the single-square pin \eqref{eq:c4-6-5} for $C_i$ (with $| a^{(i)}_W|\leqslant M_i(W)$ and $\rho_i\leqslant \rho$, since \eqref{eq:c4-4-8} is increasing in $u$) and the traversals of $\mathcal I_{C_i}$, bounds $\sum_{\mathfrak c\in\mathcal X_i}|\Psi(\mathfrak c)|$ by $C\rho\sum_{a\in\mathcal I_{C_i}}e^{-c_{\times}\delta_a/(LM)}$. Summing, $| E^{\times}_{\xi}|\leqslant 2C\rho\Delta_{\times}$. The series are those of the countable index sets; their absolute convergence is Lemma~\ref{lem:c4-9-2}, and no truncation is needed. Finally $| S^{\circ}_{\xi,C_i}|\leqslant 1$ on every branch, by the second equality of \eqref{eq:part1-reference-gas} and $K^{\circ}_{\xi;0}>0$ (Section~\ref{sec:c4-9-4}: a normalised integral of a normalised character against a positive kernel), so $| S^{\circ}_{\xi,C}|\leqslant e^{| E^{\times}_{\xi}|}\leqslant e^{C_{\times}\rho\Delta_{\times}}$.

(c) By Lemma~\ref{lem:two-loop-schur}(b) the family sum of the pair is \eqref{eq:c4-1-1}, $\mathscr R_{\xi}=\sum_{\mathfrak B\neq\varnothing}S^{\circ,[\mathfrak B]}_{\xi,C}R_{\xi}(\mathfrak B)$ over the compatible families of groups of the pair. By the ratio bound \eqref{eq:c4-6-7} for the pair (Lemma~\ref{lem:two-loop-pins}(d), valid for families containing cross groups) and Lemma~\ref{lem:two-loop-defect},
\[
    \bigl| S^{\circ,[\mathfrak B]}_{\xi,C}R_{\xi}(\mathfrak B)\bigr|\leqslant | S^{\circ}_{\xi,C}|\prod_{i=1,2}\ \prod_{W\in\mathfrak B,\ C_i\text{-group}}K_{\chi_i}\mathfrak d_W(U_1)e^{\mathfrak a_W/8}\prod_{W\in\mathfrak B,\ \text{cross}}M(W)e^{\mathfrak a_W/8}\;,
\]
with $K_{\chi_i}$ on the $C_i$-groups (the last sentence of the proof of Lemma~\ref{lem:two-loop-defect}). The $C_i$-groups of the pair are single-square groups of $C_i$ with the same $\mathfrak d_W$ and $\mathfrak a_W$, and the pair's $\rho$ dominates that of the square $C_i$, as in (b), so Lemma~\ref{lem:c4-6-4} for $C_i$ (on histories covered and far for both squares) bounds the sum of the first factors over the $C_i$-groups by $C_{\chi_i}n_{C_i}\rho(\varepsilon_1+e^{-c_ib/(LM)})$; the sum over all non-cross groups is at most $C_{\chi} n_C\rho(\varepsilon_1+e^{-cb/(LM)})$, and (a) bounds the sum of the second factors. Dropping the compatibility constraints, the sum over nonempty families is at most the exponential of the one-group sum minus one, as in the proof of Lemma~\ref{lem:c4-6-5}. The constant $C_{\times}$ is the largest of those above.
\end{proof}

\begin{corollary}[the scalar of the pair]\label{cor:two-loop-scalar}
Under the hypotheses of Lemma~\ref{lem:two-loop-cross} and of Proposition~\ref{prop:c4c-5-1} for each square, on $A_{\mathrm{cov}}$,
\begin{equation}\label{eq:two-loop-scalar}
    S^{\circ}_{\xi,C}=\kappa_{C_1}\kappa_{C_2}\,e^{E^{\circ}_{\xi,1}+E^{\circ}_{\xi,2}+E^{\times}_{\xi}}\;,
\end{equation}
with $E^{\circ}_{\xi,i}$ the error \eqref{eq:c4c-5-2} of the square $C_i$, bounded by \eqref{eq:c4c-5-1} with $n_i$, $J_{C_i}$ and the carrier distances $d_j$ of the segments of $C_i$, and $| E^{\times}_{\xi}|\leqslant C_{\times}\rho\Delta_{\times}$.
\end{corollary}

\begin{proof}
On $A_{\mathrm{cov}}$ the branch is covered for each square, so Proposition~\ref{prop:c4c-5-1} gives $S^{\circ}_{\xi,C_i}=\kappa_{C_i}e^{E^{\circ}_{\xi,i}}$; insert this in \eqref{eq:two-loop-factorisation}.
\end{proof}

\subsection{The retreated carrier}\label{sec:two-loops-retreated}
On a labelled term with final region $\Lambda$ (Section~\ref{sec:one-step-insertion}) the carrier is $D(\Lambda)$ of Section~\ref{sec:c5-2-3}, a function of $\Lambda$ and $I$. The labels, hence $\Lambda=\Lambda_1^{\mathrm{fin}}$, come from the collar tests of Section~\ref{sec:c5-3-1}, which test the fluctuation variables of collar bonds only, and the labelled weight (Lemma~\ref{lem:c5-2-5}) contains no observable; so the labelled terms, their carriers and their weights do not depend on the contour, and both squares share them. The groups of the pair at $D(\Lambda)$ are those of Definition~\ref{def:two-loop-groups} formed from the domains for $D(\Lambda)$ and the contact sets $\mathsf T_Z$ of Lemma~\ref{lem:c5-5-3}(a), taken over $n\in\mathcal I_C$; $a^{\Lambda}_W$ is the reference scalar of Lemma~\ref{lem:c5-5-3}(b) with $\chi_i$ for a $C_i$-group and $a^{\Lambda}_W\coloneqq 0$ for a cross group, and $S^{\circ}_{\Lambda,C}\coloneqq \sum_{\mathfrak B\text{ compatible}}\prod_{W\in\mathfrak B}a^{\Lambda}_W$. $S^{\circ}_{\Lambda,C_i}$ is the labelled scalar of the square $C_i$, and $\mathscr R^{\Lambda}=\sum_{\varnothing\neq\mathfrak B\text{ compatible}}S^{\circ[\mathfrak B]}_{\Lambda,C}R^0_{\Lambda}(\mathfrak B)$ is the labelled defect \eqref{eq:c5-6-2} of the pair. $A^{\mathrm{fin}}_{\mathrm{cov}}$ (Section~\ref{sec:c5-3-4}) and $\mathsf N_b$ \eqref{eq:c4-2-1} are taken for $C= C_1\sqcup C_2$; the first is defined bond by bond, so it is the intersection of the single-square events, and the second, defined through $\Gamma_C=\Gamma_{C_1}\cup\Gamma_{C_2}$, is their union, so $\mathsf N_b^c$ is the intersection of the single-square far events.

\begin{lemma}[the pair at the retreated carrier]\label{lem:two-loop-retreated}
Under the hypotheses of Proposition~\ref{prop:c5-6-def} and of \eqref{eq:c5-5-T} for each square, and for $g_0$ below the thresholds of Lemmas~\ref{lem:two-loop-source} and~\ref{lem:two-loop-pins}, with their constants, on every labelled term, for exterior data in the support of the term's collar indicators and exterior cutoffs:
(a) $S^{\circ}_{\Lambda,C}= S^{\circ}_{\Lambda,C_1}S^{\circ}_{\Lambda,C_2}e^{E^{\times}_{\Lambda}}$ with $| E^{\times}_{\Lambda}|\leqslant C_{\times}\rho\Delta_{\times}$, and $| S^{\circ}_{\Lambda,C}|\leqslant e^{C_{\times}\rho\Delta_{\times}}$;
(b) on $A^{\mathrm{fin}}_{\mathrm{cov}}\cap\mathsf N_b^c$, $\mathscr R^{\Lambda}$ satisfies \eqref{eq:two-loop-family-sum} with $S^{\circ}_{\Lambda,C}$ in place of $S^{\circ}_{\xi,C}$;
(c) under \eqref{eq:c5-7-T} for each square, on $A^{\mathrm{fin}}_{\mathrm{cov}}$, $S^{\circ}_{\Lambda,C}=\kappa_{C_1}\kappa_{C_2}e^{E^{\circ}_{\Lambda,1}+E^{\circ}_{\Lambda,2}+E^{\times}_{\Lambda}}$, with $E^{\circ}_{\Lambda,i}$ the error of Corollary~\ref{cor:c5-7-2} for the square $C_i$.
\end{lemma}

\begin{proof}
(i) Lemma~\ref{lem:two-loop-grouping} and Lemma~\ref{lem:two-loop-pins}(a) do not involve the carrier: they use the contact sets, which are field- and $t$-independent (Lemma~\ref{lem:c5-5-3}(a)), and the footprints and hulls of the domains (Lemma~\ref{lem:c5-5-8}). (ii) Lemma~\ref{lem:two-loop-schur} at $D(\Lambda)$: the Schur identity at the reference data of Lemma~\ref{lem:c5-5-3}(b) for each $\pi_i$, as in the proof of Lemma~\ref{lem:two-loop-schur}(a), and the Boolean inversion with $a^{\Lambda}_W\coloneqq 0$ for cross groups, Lemma~\ref{lem:c5-5-3}(c) with Lemma~\ref{lem:two-loop-grouping}(b) in place of Lemma~\ref{lem:grouped-source}. (iii) Lemma~\ref{lem:two-loop-source} at $D(\Lambda)$, at exterior parameter zero, on $A^{\mathrm{fin}}_{\mathrm{cov}}$: its proof uses $| f|\leqslant 1$, the normalised-trace inequality, the row locality, the unitarity of the untouched factors and properties of the selected measures; the last are supplied at $D(\Lambda)$ by the paragraphs ``Charts and primitives'' and ``Parity and the source'' of the proof of Proposition~\ref{prop:c5-6-def} (even cutoffs $\chi_b$ and density \eqref{eq:fixed-kernel-density} at exterior parameter zero, the dummy comparison \eqref{eq:c4-4-7} as Lemma~\ref{lem:c5-5-2}(c), the majorants, and the source remainder bounds of Section~\ref{sec:c4-10-4} on $A^{\mathrm{fin}}_{\mathrm{cov}}$, where every contour row has its free inputs in $D(\Lambda)$ and no exterior variable enters a link, Lemma~\ref{lem:c5-5-7z}); those paragraphs use the source only through the properties listed in the proof of Lemma~\ref{lem:two-loop-source}. So $|\mathfrak Q^0_{\Lambda}(\mathcal A;f)|\leqslant \prod_{Z\in\mathcal A}\rho\,m(Z)$ for the two-word $f$. (iv) Lemma~\ref{lem:two-loop-defect} at $D(\Lambda)$: on $A^{\mathrm{fin}}_{\mathrm{cov}}$ the contour links are reconstructed from the young variable on $D(\Lambda)$ at the background $U_1$ (proof of Proposition~\ref{prop:c5-6-def}(a), Lemma~\ref{lem:c5-6-1}), so a traversal outside the intervals has the link $U_1$ and one in $\mathcal H(W)$ depends on the roots of $W$ only; the weights obey $| w^0_{\Lambda}(W)|\leqslant \prod_j|\Phi|^+_{\Lambda}(\mathfrak c_j)/k(W)!$ (Corollary~\ref{cor:c5-4-3}, Lemma~\ref{lem:c5-5-6}); with (iii) the proof gives \eqref{eq:two-loop-defect} with $\mathfrak d^0_W$ and $M^{(\rho)}_{\Lambda}(W)$. (v) Lemma~\ref{lem:two-loop-pins}(b)--(d) at $D(\Lambda)$: these are the summations of Lemmas~\ref{lem:c4-8-2} and~\ref{lem:c4-8-3} with the anchoring of the pair, and they use only the root weights $\rho\,m(Z)$, the tree weights $|\Phi|^+_{\Lambda}= M_v$ and the field-independent domain, cluster and interval geometry, which is the reason Lemma~\ref{lem:c5-5-8}(a) gives for one square; so they hold for $M^{(\rho)}_{\Lambda}$ with the constants of Lemma~\ref{lem:two-loop-pins}, and, since $\sigma(W)=\mathfrak s(W)$, $\nu=\eta$ and $C_F= A_E$ there, $c_{\times}$ is unchanged. In particular the weighted criterion \eqref{eq:two-loop-weighted-kp} holds at $D(\Lambda)$ for the pair, from $| a^{\Lambda}_W|\leqslant M^{(\rho)}_{\Lambda}(W)$ and the pins with $3^4A'_E+LM$ in place of $3^4A_F+LM$, and for each square as in the proof of Lemma~\ref{lem:two-loop-cross}(b). (vi) With (i)--(v), the proof of Lemma~\ref{lem:two-loop-cross}(a),(b) gives (a), with $| S^{\circ}_{\Lambda,C_i}|\leqslant 1$ (Lemma~\ref{lem:c5-6-3}) in place of $| S^{\circ}_{\xi,C_i}|\leqslant 1$; for (b), the proof of Lemma~\ref{lem:two-loop-cross}(c) uses \eqref{eq:c5-6-def-3} for each square (Lemma~\ref{lem:c4-6-4} at $D(\Lambda)$ in the form with $e^{\mathfrak a_W/8}$, proved in the paragraph ``Pins'' of the proof of Proposition~\ref{prop:c5-6-def}). (c) is (a) with Corollary~\ref{cor:c5-7-2}(a) for each square, $S^{\circ}_{\Lambda,C_i}=\kappa_{C_i}e^{E^{\circ}_{\Lambda,i}}$ on $A^{\mathrm{fin}}_{\mathrm{cov}}$.
\end{proof}

Lemma~\ref{lem:two-loop-retreated} is not used in the proof of Theorem~\ref{thm:two-loop-theorem}, which works at the carrier $I$ throughout (Section~\ref{sec:two-loops-proof}), as Theorem~\ref{thm:main-theorem} does. For the pair it gives, in (c), the identification that Remark~\ref{rem:main-regime}(2) records for one square (Corollary~\ref{cor:c5-7-2}(a)); in (a), the factorisation and the bound of $S^{\circ}_{\Lambda,C}$; and in (b), the far labelled family sum, the pair inputs of \eqref{eq:c5-6-def-2}. The labelled defect bound \eqref{eq:c5-6-def-2} and Corollary~\ref{cor:c5-7-4} are not restated for the pair.

\subsection{Proof of Theorem~\ref{thm:two-loop-theorem}}\label{sec:two-loops-proof}
The single-contour statements used by Theorem~\ref{thm:main-theorem} outside the extraction (Lemmas~\ref{lem:c3-2-1}, \ref{lem:c3-3-2}, \ref{lem:exterior-response-bound}, Corollary~\ref{cor:c3-6-1}, Lemmas~\ref{lem:c3-7-1}, \ref{lem:c3-7-4}, Proposition~\ref{prop:c3-7-5}, Lemma~\ref{lem:c4-2-1}, \eqref{eq:c4-8-1} and Lemma~\ref{lem:c5-7-expmoment}) involve the contour through its observable $f_{\xi}$, the set $I_C$ of its free inputs \eqref{eq:c3-2-1}, the set $\Gamma_C$ of its contour bonds and free-input templates (Section~\ref{sec:c3-7}) and its segments, and, in Lemma~\ref{lem:c4-2-1} and \eqref{eq:c4-8-1}, through the bound $| S^{\circ}_{\xi,C}|\leqslant 1$ of \eqref{eq:part1-reference-gas}, which for the pair is replaced by $| S^{\circ}_{\xi,C}|\leqslant e^{C_{\times}\rho\Delta_{\times}}$ of (b) of Lemma~\ref{lem:two-loop-inputs}. For the pair, $f_{\xi}$ is \eqref{eq:two-loop-observable} at the reconstructed field, $I_C= I_{C_1}\cup I_{C_2}$, $\Gamma_C=\Gamma_{C_1}\cup\Gamma_{C_2}$, so that $|\Gamma_C|\leqslant C_Ln_{C_1}+C_Ln_{C_2}= C_Ln_C$, and $J_C= J_{C_1}\sqcup J_{C_2}$ is the set of the $4n_1+4n_2$ coarse segments of the two squares, $d_j$ being the distance \eqref{eq:c4c-1-1} of the square of the segment $j$ at the common carrier $I$. The constants $C$, $c$ are those of Theorem~\ref{thm:two-loop-theorem}.

\begin{lemma}[single-contour inputs for the pair]\label{lem:two-loop-inputs}
Under the hypotheses of Theorem~\ref{thm:two-loop-theorem}:
(a) \eqref{eq:c3-2-3} holds for the pair with $C_{\mathrm{resp},\chi_1}+C_{\mathrm{resp},\chi_2}$ in place of $C_{\mathrm{resp},\chi}$, and so do Lemma~\ref{lem:c3-3-2}, Lemma~\ref{lem:exterior-response-bound} and Corollary~\ref{cor:c3-6-1}; \eqref{eq:c3-7-14} holds for the pair with $C_{\mathrm{resp},\chi_1\chi_2}$.
(b) On $A_{\mathrm{cov}}$, $f_{\xi}(0)= W_{\chi_1,\chi_2}(C;U_1)$ and $| S^{\circ}_{\xi,C}|\leqslant e^{C_{\times}\rho\Delta_{\times}}$, and \eqref{eq:c4-8-1} holds for the pair, every term being integrable against the history weight.
(c) On $A_{\mathrm{cov}}$, $|\mathscr R_{\xi}|\leqslant 1+e^{C_{\times}\rho\Delta_{\times}}$; and $\mathbb P_{\mathrm h}(\mathsf N_b)\leqslant 4C_Ln_C(2b+3)^4q_{\mathrm{pt}}$ and $| E^{\mathrm{near}}_{\mathrm{def}}|\leqslant e^{C_{\times}\rho\Delta_{\times}}C'_Ln_Ce^{-(c_*/2)\pi_*^2}$, with $\mathsf N_b$ and $E^{\mathrm{near}}_{\mathrm{def}}$ of Section~\ref{sec:c4-2} formed with $\Gamma_C$ and $\mathscr R_{\xi}$ of the pair.
(d) Put $\varphi_c(b)\coloneqq \sum_{r\in C_1\sqcup C_2}e^{-c(d_T(r,b)-R_0')_+}$, the sum of the functions $\varphi_c$ of \eqref{eq:c5-7-boundary-marks} of the two squares. Pointwise on every branch, $\sum_{j\in J_C}e^{-cd_j}\leqslant \sum_b\varphi_c(b)\mathbf 1\{ b\in E(B_{\xi})\}$, with $\varphi_c(b)\leqslant 2\Phi_c$ and $\sum_b\varphi_c(b)\leqslant n_CN'_{\mathrm{lat}}(c)$; and for every $\lambda>0$, also with $\mathbf 1_{A_{\mathrm{cov}}}$ omitted,
\[
    \mathcal Z_0^{-1}\int\mathcal O\Bigl[\mathsf w_{\xi}\mathbf 1_{A_{\mathrm{cov}}}\exp\Bigl(\lambda\sum_{j\in J_C}e^{-cd_j}\Bigr)\Bigr]\dd V\leqslant \exp\bigl[n_C\lambda N'_{\mathrm{lat}}(c)\,e^{2\lambda\Phi_c}\,q^{1/C_K}\bigr]\;.
\]
(e) For $g_0$ below a further threshold depending on $L$, $M$, $M_1$, $M_2$, the group, the characters, the source constants and $c$,
\begin{equation}\label{eq:two-loop-moment}
    \delta^{I,2}_q\coloneqq C\tau^2N'_{\mathrm{lat}}(c)\,e^{2C\tau^2\Phi_c}\,q^{1/C_K}\leqslant \exp\bigl[-c_q\pi_*^2/(2C_K)\bigr]\;.
\end{equation}
\end{lemma}

\begin{proof}
(a) On the enlarged box the real reconstruction gives group-valued links (proof of Lemma~\ref{lem:c3-3-2}), so $|\chi_i(Y_i)|\leqslant 1$ for the ordered link product $Y_i$ of $C_i$ and $| f_{\xi}|\leqslant 1$. By the product rule, $\nabla_bf_{\xi}=(\nabla_b\chi_1(Y_1))\chi_2(Y_2)+\chi_1(Y_1)\nabla_b\chi_2(Y_2)$, and on a branch covered for the pair, which is covered for each square (Section~\ref{sec:two-loops-setting}), the proof of \eqref{eq:c3-2-3} for the square $C_i$ gives $h(\nabla_b\chi_i(Y_i))\leqslant C_{\mathrm{resp},\chi_i}g_0\mathbf 1_{I_{C_i}}(b)$, so $h(\nabla_bf_{\xi})\leqslant |\chi_2(Y_2)|\,h(\nabla_b\chi_1(Y_1))+|\chi_1(Y_1)|\,h(\nabla_b\chi_2(Y_2))\leqslant (C_{\mathrm{resp},\chi_1}\mathbf 1_{I_{C_1}}(b)+C_{\mathrm{resp},\chi_2}\mathbf 1_{I_{C_2}}(b))g_0\leqslant (C_{\mathrm{resp},\chi_1}+C_{\mathrm{resp},\chi_2})g_0\mathbf 1_{I_C}(b)$, which is \eqref{eq:c3-2-3} for the pair. Lemma~\ref{lem:c3-3-2} uses the observable through $\widetilde f= f_{\xi}\circ T$, $|\widetilde f|\leqslant 1$ and \eqref{eq:c3-2-3} composed with $DT$; Lemma~\ref{lem:exterior-response-bound} uses it through Lemma~\ref{lem:c3-3-2} and the support $I_C$; Corollary~\ref{cor:c3-6-1} is that lemma averaged. The enlargement in the proof of Lemma~\ref{lem:exterior-response-bound} multiplies the constant of \eqref{eq:c3-2-3} by a factor independent of the character, so the enlarged pair constant is the sum of the enlarged constants of $\chi_1$ and $\chi_2$. The proof of Proposition~\ref{prop:c3-7-5} combines Corollary~\ref{cor:c3-6-1} with Lemmas~\ref{lem:c3-7-1} and~\ref{lem:c3-7-4}, whose bounds are sums of one-point marks over the members of $\Gamma_C$; with $|\Gamma_C|\leqslant C_Ln_C$ it gives \eqref{eq:c3-7-14} with $C_L\max((C_{\mathrm{resp},\chi_1}+C_{\mathrm{resp},\chi_2})S_c,1)= C_{\mathrm{resp},\chi_1\chi_2}$.

(b) Lemma~\ref{lem:covered} is proved bond by bond, so on $A_{\mathrm{cov}}$, $U_{\xi}(0)(r)= U_1(r)$ for every bond $r$ of $C_1$ and of $C_2$, and $f_{\xi}(0)= W_{\chi_1,\chi_2}(C;U_1)$. By Lemma~\ref{lem:two-loop-cross}(b), $| S^{\circ}_{\xi,C}|\leqslant | S^{\circ}_{\xi,C_1}|| S^{\circ}_{\xi,C_2}| e^{| E^{\times}_{\xi}|}\leqslant e^{C_{\times}\rho\Delta_{\times}}$, where $| S^{\circ}_{\xi,C_i}|\leqslant 1$ because it is a normalised integral of a normalised character against the positive reference kernel (proof of Proposition~\ref{prop:c4-12-3}). The derivation of \eqref{eq:c4-8-1} in Section~\ref{sec:c4-8-1} multiplies \eqref{eq:c-one-step-branch-identity}, which holds for the pair by Lemma~\ref{lem:two-loop-schur}(b), by $\mathbf 1_{A_{\mathrm{cov}}}$ and writes $\mathcal J^{\mathrm a}_{\xi}[f_{\xi}]=\mathsf w_{\xi} F_{\xi}(1)$ on $A^c_{\mathrm{cov}}$; \eqref{eq:native-source-physical-ratio} holds with $W_{\chi_1,\chi_2}(C)$ in place of $W_{\chi}(C)$, $f_{\xi}$ being jointly Borel with $| f_{\xi}|\leqslant 1$ (the measurability paragraph of Section~\ref{sec:branch-data}). Its integrability uses $| F_{\xi}(s)|\leqslant 1$ ($F_{\xi}(s)=\int f_{\xi}\,\dd\nu_{\xi;s}$ with $| f_{\xi}|\leqslant 1$), $|\mathcal B_{\xi}[f_{\xi}]|\leqslant (C_{\mathrm{resp},\chi_1}+C_{\mathrm{resp},\chi_2})g_0r| I_C|$ on $A_{\mathrm{cov}}$ by (a), and a bound on $S^{\circ}_{\xi,C}f_{\xi}(0)$, which is now $| S^{\circ}_{\xi,C}f_{\xi}(0)|\leqslant e^{C_{\times}\rho\Delta_{\times}}$ in place of $3$.

(c) On $A_{\mathrm{cov}}$, $\mathscr R_{\xi}= F_{\xi}(0)-S^{\circ}_{\xi,C}f_{\xi}(0)$ by \eqref{eq:part1-scalar-extraction} for the pair (Lemma~\ref{lem:two-loop-schur}(b)), so $|\mathscr R_{\xi}|\leqslant 1+e^{C_{\times}\rho\Delta_{\times}}\leqslant 2e^{C_{\times}\rho\Delta_{\times}}$ by (b). The union bound of the proof of Lemma~\ref{lem:c4-2-1}, over the at most $C_Ln_C$ members of $\Gamma_C$, gives the bound on $\mathbb P_{\mathrm h}(\mathsf N_b)$. Positivity of $\mathbb E_{\mathrm h}$ then gives $| E^{\mathrm{near}}_{\mathrm{def}}|\leqslant 2e^{C_{\times}\rho\Delta_{\times}}\cdot4C_Ln_C(2b+3)^4q_{\mathrm{pt}}$, and the last step of that proof gives $8C_L(2b+3)^4q_{\mathrm{pt}}\leqslant C'_Le^{-(c_*/2)\pi_*^2}$.

(d) The first inequality of \eqref{eq:c5-7-expmoment-I} is proved segment by segment, the term of a segment being bounded through the contour bonds of its own square (proof of Lemma~\ref{lem:c5-7-expmoment}). The two squares' inequalities at the common carrier add up to the pointwise bound. Each square's $\varphi_c$ is at most $\Phi_c$ and has sum at most $n_{C_i}N'_{\mathrm{lat}}(c)$, by the counts of that proof. The rest of that proof applies with the weights $w_b=\varphi_c(b)\leqslant 2\Phi_c$: the expansion over the marks $\mathbf 1\{ b\in E(B_{\xi})\}$, one per bond, and the joint bound \eqref{eq:c5-7-joint-I}, which holds for every finite set of bonds. They bound the average by $\exp(\hat q\sum_b(e^{\lambda w_b}-1))$ with $\hat q= q^{1/C_K}$. Since $e^x-1\leqslant xe^x$ for $x\geqslant 0$, this is at most $\exp(\hat q\lambda e^{2\lambda\Phi_c}\sum_bw_b)$, and the indicator enters only through $\mathbf 1\leqslant 1$.

(e) Since $\tau<1$ and $q\leqslant e^{-c_*\pi_*^2}= e^{-c_q\pi_*^2}$ (Proposition~\ref{prop:c3-7-5}), $\log\delta^{I,2}_q\leqslant \log(CN'_{\mathrm{lat}}(c))+2C\Phi_c-c_q\pi_*^2/C_K$. Moreover $\pi_*^2/C_K\geqslant c'(\log g_0^{-2})^{2p_0-4\mathfrak r}$ with $2p_0-4\mathfrak r\geqslant 12$ (the arithmetic preceding \eqref{eq:c5-7-delta-small}, Remark~\ref{rem:main-regime}(1)). So $c_q\pi_*^2/(2C_K)\geqslant \log(CN'_{\mathrm{lat}}(c))+2C\Phi_c$ for small $g_0$, which is \eqref{eq:two-loop-moment}.
\end{proof}

\begin{proof}[Proof of Theorem~\ref{thm:two-loop-theorem}]
(i) By Lemma~\ref{lem:two-loop-inputs}(b), $\mathbb E_{\varepsilon}[W_{\chi_1,\chi_2}(C)]=\mathbb E_{\mathrm h}[\mathbf 1_{A_{\mathrm{cov}}}S^{\circ}_{\xi,C}W_{\chi_1,\chi_2}(C;U_1)]+E_{\mathrm{def}}+E_{\mathrm{bdry}}$, with all terms integrable. On $A_{\mathrm{cov}}$ substitute \eqref{eq:two-loop-scalar} (Corollary~\ref{cor:two-loop-scalar}). Its hypotheses are those of Lemma~\ref{lem:two-loop-cross} and, for each square, those of Proposition~\ref{prop:c4c-5-1}, which (H2) and (H3) for that square imply as in the proof of Theorem~\ref{thm:main-theorem}. After the substitution the first integrand is still exactly $\mathbf 1_{A_{\mathrm{cov}}}S^{\circ}_{\xi,C}f_{\xi}(0)$. The statements on $\kappa_{C_i}$ are Theorem~\ref{thm:main-theorem}(i) for $(C_i,\chi_i)$.

(ii) The pair's $A_{\mathrm{cov}}$ is contained in the covered event of each square, on which Proposition~\ref{prop:c4c-5-1} bounds $E^{\circ}_{\xi,i}$ with the constants $C_i$, $c_i$ of $\chi_i$; these bounds are at most the stated ones, since $C_i\leqslant C$ and $e^{-c_ix}\leqslant e^{-cx}$ for $x\geqslant 0$. The elementary bounds give $C\tau^2(9n_i+1)$, as in Theorem~\ref{thm:main-theorem}. The bound on $E^{\times}_{\xi}$ is Lemma~\ref{lem:two-loop-cross}(b) with \eqref{eq:two-loop-cross-sum}. Split $E_{\mathrm{def}}= E^{\mathrm{near}}_{\mathrm{def}}+E^{\mathrm{far}}_{\mathrm{def}}$ as in Section~\ref{sec:c4-2}, both terms being integrable by Lemma~\ref{lem:two-loop-inputs}(c). The near term is bounded in Lemma~\ref{lem:two-loop-inputs}(c). The histories of $A_{\mathrm{cov}}\cap\mathsf N_b^c$ are covered and far for the pair. On them \eqref{eq:two-loop-family-sum} and Lemma~\ref{lem:two-loop-inputs}(b) give $|\mathscr R_{\xi}|\leqslant e^{C_{\times}\rho\Delta_{\times}}(e^{X+C_{\times}\rho\Delta_{\times}}-1)$ with $X= C_{\chi} n_C\rho(\varepsilon_1+e^{-cb/(LM)})$. So $| E^{\mathrm{far}}_{\mathrm{def}}|\leqslant e^{C_{\times}\rho\Delta_{\times}}(e^{X+C_{\times}\rho\Delta_{\times}}-1)$, by the positivity of $\mathbb E_{\mathrm h}$ and $\mathbb E_{\mathrm h}[1]=1$. The bound on $E_{\mathrm{bdry}}$ is \eqref{eq:c3-7-14} for the pair (Lemma~\ref{lem:two-loop-inputs}(a)).

(iii) The first bound follows from the positivity of $\mathbb E_{\mathrm h}$, $\mathbb E_{\mathrm h}[1]=1$, $| W_{\chi_1,\chi_2}|\leqslant 1$ and $| e^z|\leqslant e^{| z|}$ with (ii). Under the condition of (H3) for (iii) for the square $C_i$, $4n_ie^{-c_i\iota_T}\leqslant 1$ (\eqref{eq:c5-7-T}(iii)). So Corollary~\ref{cor:c5-7-2}(c), applied to $E^{\circ}_{\xi,i}$ as in the proof of Theorem~\ref{thm:main-theorem}, gives $| E^{\circ}_{\xi,i}|\leqslant C_i\tau^2(2+1/c_i)+C_i\tau^2\sum_{j\in J_{C_i}}e^{-c_id_j}\leqslant C\tau^2(2+1/c)+C\tau^2\sum_{j\in J_{C_i}}e^{-cd_j}$ on $A_{\mathrm{cov}}$. Hence $\mathbf 1_{A_{\mathrm{cov}}}e^{| E^{\circ}_{\xi,1}|+| E^{\circ}_{\xi,2}|}\leqslant e^{2C\tau^2(2+1/c)}\mathbf 1_{A_{\mathrm{cov}}}\exp(C\tau^2\sum_{j\in J_C}e^{-cd_j})$, and Lemma~\ref{lem:two-loop-inputs}(d) at $\lambda= C\tau^2$, with (e), gives \eqref{eq:thm-two-moment}. The modulus of the leading term is at most $|\kappa_{C_1}\kappa_{C_2}| e^{C_{\times}\rho\Delta_{\times}}\mathbb E_{\mathrm h}[\mathbf 1_{A_{\mathrm{cov}}}e^{| E^{\circ}_{\xi,1}|+| E^{\circ}_{\xi,2}|}]$, by $| W_{\chi_1,\chi_2}|\leqslant 1$ and (ii). Pointwise on $A_{\mathrm{cov}}$, with $z= E^{\circ}_{\xi,1}+E^{\circ}_{\xi,2}+E^{\times}_{\xi}$, $| e^z-1|\leqslant e^{| z|}-1\leqslant e^{2C\tau^2(2+1/c)+C_{\times}\rho\Delta_{\times}}\exp(C\tau^2\sum_{j\in J_C}e^{-cd_j})-1$, a nonnegative function. Its $\mathbb E_{\mathrm h}$-average without the indicator is at most $\exp[2C\tau^2(2+1/c)+C_{\times}\rho\Delta_{\times}+n_C\delta^{I,2}_q]-1$, by Lemma~\ref{lem:two-loop-inputs}(d) and $\mathbb E_{\mathrm h}[1]=1$.
\end{proof}

\section{Towards the next step}\label{sec:next-step}
The statement of this section is a conjecture; nothing in the other sections uses it. It concerns loops whose side is a power of the block size, $n= L^k$ with an integer $k\geqslant 0$, after $k+1$ steps (the letter $K$ is RG1's exponent in (H3) and the index of $\sigma_K$, $C_K$ in Lemma~\ref{lem:c5-7-joint}). On a torus of (H3) the lattice $L^{k+1}\mathbb Z^4$ and its plaquettes are defined ($K_{\mu}=2L^{m-1}\geqslant 3(n+2)$ forces $m\geqslant k+2$, so $L^{k+1}$ divides every period $N_{\mu}=2L^m$), and a translation of $C_n$ does not change $\mathbb E_{\varepsilon}[W_{\chi}(C_n)]$ (the Wilson measure $\rho_0(U)\,\dd U$ on the periodic torus is invariant under lattice translations, its action being a sum over all plaquettes and the Haar measure a product measure; Section~\ref{sec:extent-weights} uses the remaining lattice symmetries), so we take the corners of $C_n$ in $L^{k+1}\mathbb Z^4$; then $C_n$, of side $L^{k+1}$ fine units, is a contour of the lattice of spacing $L^j\varepsilon$ for every $j\leqslant k+1$ and the boundary of a single plaquette of the lattice of spacing $L^{k+1}\varepsilon$. The scales $\tau$ and $\pi_*$ are those of the notation paragraph of Section~\ref{sec:main-result}, at the bare coupling $g_0$.

\begin{conjecture}[the loop along the scales]\label{conj:iteration}
Let $G$, $\chi$ be as in (H1) and fix the cutoff prescription $\sigma$ and the constants as in (H2). There are $C,c>0$, depending on $L$, the group, the character, $\sigma$ and the constants fixed in (H2) only, thresholds $g^*_j>0$ ($j\geqslant 1$) and, for $0< g_0< g^*_j$, real numbers $\delta m_j$ and complex numbers $\zeta_j$, depending on $g_0$, $\varepsilon$, $j$, $L$, the group, the character, $\sigma$ and the fixed prescriptions listed in Theorem~\ref{thm:main-theorem}(i) only, with $z_j\coloneqq e^{\zeta_j}$ and
\begin{equation}\label{eq:conj-iteration-coefficients}
    \delta m_1=\delta m^{\sigma}(g_0,\varepsilon)\;,\qquad\zeta_1=\frac14\sum_{(v,w)\in\mathscr C_{\square}}\zeta^{\sigma}_{v,w}(g_0)\;,\qquad|\delta m_j|\,L^j\varepsilon\leqslant C\tau^2\;,\qquad|\zeta_j|\leqslant C\tau^2\;,
\end{equation}
so that $z_1= z^{\sigma}_{\mathrm c}(g_0,\pi/2)$ \eqref{eq:corner-square-factor}, such that for every integer $k\geqslant 0$, with $n= L^k$, there is $g_{\mathrm{it}}(k)\leqslant \min_{j\leqslant k+1}g^*_j$ such that for $0< g_0< g_{\mathrm{it}}(k)$, on every torus of (H3) whose exponent $m$ exceeds a bound depending on $k$ and $g_0$,
\begin{equation}\label{eq:conj-iteration}
    \begin{gathered}
        \mathbb E_{\varepsilon}[W_{\chi}(C_n)]=\kappa^{(k+1)}_{C_n}(1+R)+R'\;,\qquad
        \kappa^{(k+1)}_{C_n}=\exp\Bigl(-P(C_n)\sum_{j=1}^{k+1}\delta m_j\Bigr)\prod_{j=1}^{k+1}z_j^4\;,\\
        | R|\leqslant C\bigl[(k+2)\tau^2+n_Ce^{-c\pi_*^2}\bigr]\;,\qquad| R'|\leqslant Cn_Ce^{-c\pi_*^2}\;.
    \end{gathered}
\end{equation}
\end{conjecture}

\begin{remark}[content of the conjecture and what a proof needs]\label{rem:next-step}
(1) The $j$-th coefficients do not depend on the size of the loop: the perimeter part is a sum over the scales of perimeter coefficients, each of size $\tau^2$ in units of its own lattice spacing, so $| P(C_n)\sum_j\delta m_j|\leqslant 4C\tau^2n(1+L^{-1}+L^{-2}+\cdots)$, while the corner factors contribute one factor $z_j^4$ per scale, $k+1=\log_L(nL)$ in all. For $k=0$ the statement is Theorem~\ref{thm:main-theorem}, with the relative form of the far-history bound of Remark~\ref{rem:main-regime}(1), which places the second term of the bound on $E_{\mathrm{def}}$ in $\kappa_{C_1}R$, together with the claim that the leading expectation $\mathbb E_{\mathrm h}[\mathbf 1_{A_{\mathrm{cov}}}e^{E^{\circ}_{\xi}}W_{\chi}(C_1;U_1)]$ of the single-plaquette loop is $1+O(\tau^2+n_Ce^{-c\pi_*^2})$; the first step gives the bounds \eqref{eq:thm-main-bounds}, the upper bounds of (iii) and, by Corollaries~\ref{cor:main-factorisation} and~\ref{cor:main-leading-size} at $n=1$, $1+O(\tau^2+L^2\varepsilon_1+(L^2+LC_L)q+n_C\delta^{\sharp}_q)$, in which the area term $L^2\varepsilon_1$ of the bound, of order $g_0(\log g_0^{-2})^{p_0}$ by CRE (2.6) as quoted in Section~\ref{app:thresholds}, is not $O(\tau^2)$; whether the leading expectation itself is $1+O(\tau^2)$ is not decided here. (2) For $j\geqslant 2$ the coefficients $\delta m_j$, $z_j$ are not defined in this paper; they are to be the perimeter coefficient and the corner factor of the $j$-th step, computed from the density after $j-1$ steps as \eqref{eq:marginal-square-mass} and \eqref{eq:corner-square-factor} compute them from the flat reference integrals $N^{\sigma}_{C,T}$, $D^{\sigma}_T$ of \eqref{eq:marginal-literal-ratio}, in which case the conjecture says that the coefficients of every step have the size that Sections~\ref{sec:marginal-perimeter} and~\ref{sec:marginal-corner} prove for the first. At the stated precision the corner factors are not separately determined: $| z_j^4-1|=| e^{4\zeta_j}-1|\leqslant e^{4C\tau^2}-1$, so $|\prod_{j\leqslant k+1}z_j^4-1|\leqslant e^{4C(k+1)\tau^2}-1\leqslant 8C(k+1)\tau^2$ once $4C(k+1)\tau^2\leqslant 1$, and replacing every $z_j$ by $1$ replaces $1+R$ by $\prod_jz_j^4(1+R)$, with $|\prod_jz_j^4(1+R)-1|\leqslant 8C(k+1)\tau^2(1+| R|)+| R|$, which a larger $C$ in the bound on $R$ absorbs below a threshold depending on $k$; the $z_j$ record where the term $(k+2)\tau^2$ of $R$ comes from. With the threshold $g_{\mathrm{it}}(k)$ depending on $k$, $R'$ is not needed either: $|\kappa^{(k+1)}_{C_n}|\geqslant e^{-8C\tau^2(n+k+1)}$ ($L\geqslant 2$), so once $8C\tau^2(n+k+1)\leqslant 1$ the bound on $R'$ is at most $eCn_Ce^{-c\pi_*^2}|\kappa^{(k+1)}_{C_n}|$ and $R'$ can be moved into $\kappa^{(k+1)}_{C_n}R$ at the cost of replacing $C$ by $(1+e)C$; it is substantive in the form with a threshold independent of $k$. Since $\delta m_1=\delta m^{\sigma}$ is given, what the statement determines beyond the first step and beyond $O((k+2)\tau^2)$ is $P(C_n)\sum_{j\geqslant 2}\delta m_j$, of size at most $4C\tau^2(n-1)/(L-1)$, modulo $O((k+2)\tau^2)$. For $k\leqslant 1$ this is at most $4C\tau^2$, so nothing beyond the first step is determined; in general the scales $j$ with $L^{k+1-j}\leqslant k+2$ contribute at most $4C\tau^2(k+2)L/(L-1)$, also of order $(k+2)\tau^2$, and the content lies in the scales $j\geqslant 2$ with $L^{k+1-j}> k+2$. (3) A proof by $k+1$ applications of the first step needs the following ((c) only for a threshold independent of $k$), none of which is proved here beyond the first step. (a) A local form of the leading term: on the covered event, $W_{\chi}(C_n;U_1)e^{E^{\partial}_{\xi}}$ as a function of the coarse field $V$ near the loop, a coarse Wilson loop times the exponential of a sum over the coarse bonds of the loop of terms depending on $V$ within a fixed distance of the bond, analytic on the small-field chart, gauge covariant and bounded by $C\tau^2$, with the history labels absorbed into the large-field expansion of $\rho_1$; this is what turns the output of the step, the pair of the complex measure $\kappa_{C_n}^{-1}N_1$ and the measure $\rho_1$ (Remark~\ref{rem:main-regime}(3)), into an input of the same type. (b) General-step analogues of (B1)--(B4), which Section 2 of \cite{BalabanCRE} does not state as such: it formulates the inductive assumption with the operations corresponding to $\mathbf T_1$ left undefined, only their general properties being formulated (p.254), their inductive definition being deferred to its next section and forthcoming papers (pp.254--255), and with modifications connected with new terms obtained by the $\mathbf R$-operations (p.254), among them the admitted case $\Lambda_j=\Omega_j$ (p.255); the sets $S_j\subset\Omega_j\cap\Lambda_j^c$ and the functions $\chi(\Omega_j\cap\Lambda_j^c,S_j)$, $\zeta(\Omega_j^c)$ are the general-step forms of the first-step objects of CRE (1.29) (pp.254--255); the geometry and the scales are CRE (2.1)--(2.9), with the couplings of the first $k+1$ scales compared through the first inequality $g_n\leqslant (1+g_n^2\beta'(n-m))^{1/2}g_m$ of CRE (2.6) (p.255), $g_{k+1}^2\beta'(k+1)$ being small for $g_0< g_{\mathrm{it}}(k)$ (Appendix~\ref{app:thresholds} quotes (2.6) for the first step). (c) The defect: the far-history part of \eqref{eq:c4-12-7} is relative to $|\kappa_{C_n}|$ but of order $x= C_{\chi} n_C\rho(\varepsilon_1+e^{-cb/(LM)})$ when $x\leqslant 1$, polynomially small per contour bond (Remark~\ref{rem:main-regime}(1)); for a threshold independent of $k$ it must be shown to be a perimeter-local correction, absorbed in $\delta m_2$ and $z_2$ at the next step. Below a threshold depending on $k$ it is at most $C\tau^2$ relative to $|\kappa_{C_n}|$ ($\rho\leqslant C\tau^2$, $\varepsilon_1\to0$ and $e^{-cb/(LM)}\to0$ at fixed $n_C$) and goes into $R$, as at $k=0$. (d) The absolute terms $E_{\mathrm{bdry}}$ and the near-history part of $E_{\mathrm{def}}$ at the $j$-th step, of order $4L^{k+2-j}e^{-c\pi_{*,j}^2}$ for the loop of $4L^{k+2-j}$ bonds of the lattice of spacing $L^{j-1}\varepsilon$, with $\pi_{*,j}=\min(p_0(g_{j-1}),p_0(g_j))$ the scale of that step, $\pi_{*,j}^2\geqslant \pi_*^2/2$ once $g_i^2\beta'i\leqslant 1$ for $i\in\{ j-1,j\}$ and $(1-\log2/\log g_0^{-2})^{2p_0}\geqslant 1/2$ (the first inequality of CRE (2.6) with $m=0$, $n= i\geqslant 1$ gives $g_i^2\leqslant (1+g_i^2\beta'i)g_0^2\leqslant 2g_0^2$, hence $\log g_i^{-2}\geqslant \log g_0^{-2}-\log2$ and $p_0(g_i)^2\geqslant p_0(g_0)^2/2\geqslant \pi_*^2/2$; for $i=0$ this is $p_0(g_0)\geqslant \pi_*$), and entering multiplied by the leading coefficients of the steps $1,\ldots,j-1$, whose product must be bounded; so that their sum over the steps is at most $Cn_Ce^{-c\pi_*^2}$; their relative size is discussed in Remark~\ref{rem:main-regime}(1). (e) At every step the remainder of Corollary~\ref{cor:main-factorisation}, which at the first step is at most $e^{n_C\delta^{\sharp}_q}-1$, and the torus and corner part $E^{\circ}_{T,n}$, at most $C\tau^2(2+1/c)$, one source of the term $(k+2)\tau^2$. (f) The endpoint: the leading expectation of the single-plaquette loop at the last step is $1+O(\tau^2)$; at the first step Corollary~\ref{cor:main-leading-size} gives $1+O(\tau^2+L^2\varepsilon_1+(L^2+LC_L)q+n_C\delta^{\sharp}_q)$ only, and its area term $L^2\varepsilon_1$, which comes from the deterministic bound on the plaquettes of the background, must be replaced by a bound of order $\tau^2$. (g) The terms of the $\mathbf R$-operations at every step, by which the general representation modifies the first-step representation CRE (1.28)--(1.30) (p.254).
\end{remark}

\appendix
\section{The Flat Covariance}\label{app:flat-covariance}

The bulk coefficient uses the flat fluctuation operator on the infinite lattice with the same periodic block grid and constraint extension. We first obtain a uniform lower bound on compatible finite tori, then construct that operator and compare its inverse with the finite-torus inverses. This is a quadratic first-step calculation at the flat background.

The spatial dimension is four throughout this comparison. For a site field $\lambda$, set $(D\lambda)(x,x')=\lambda(x')-\lambda(x)$. For an oriented plaquette $p$, set $(\operatorname{curl}a)(p)= a(\partial p)$, where $a(\gamma)$ is the signed sum along an oriented path $\gamma$. All squared norms in this comparison sum the Lie-algebra inner product over positive bonds, plaquettes or sites of the indicated carrier. Define the mean path sum by
\begin{equation}\label{eq:covariance-mean-path}
    (\Phi a)(y,x)=\frac{1}{|\mathcal G(y,x)|}
        \sum_{\gamma\in\mathcal G(y,x)}a(\gamma)\;.
\end{equation}
RG1 (0.8) and (2.5) identify $G^{(2)}(a)=\|\Phi a\|^2$. Differentiating RG1 (0.4), with $x'= x+Le_{\mu}$ for $c=(y,y+Le_{\mu})$, gives
\begin{equation}\label{eq:covariance-two-averages}
    \begin{aligned}
        (Q_1a)(c)&= L^{-5}\sum_{x\in\mathcal B_y}a([x,x'])\;,\\
        (Qa)(c)&=(Q_1a)(c)
            +L^{-5}\sum_{x\in\mathcal B_y}
                ((\Phi a)(y,x)-(\Phi a)(y+Le_{\mu},x'))\;.
    \end{aligned}
\end{equation}
The operator $Q_1$ is the straight-path average in PRT2 (2.125) \cite{BalabanPRT2}; it is not identified with $Q$. Telescoping these paths gives $Q(D\lambda)(c)=(\lambda(c_+)-\lambda(c_-))/L$.

For each block choose its lower corner $o_y= y-\frac{L-1}{2}\sum_{\mu=1}^4 e_{\mu}$. Let $\kappa_{y,x}$ traverse the coordinate displacements from $o_y$ to $x$ in the order $1,\ldots,4$. The union of these paths is a spanning tree. Define
\begin{equation}\label{eq:covariance-tree-gauge}
    \lambda(x)= a(\kappa_{y,x})-a(\kappa_{y,y}),
    \qquad b= a-D\lambda,
    \qquad E(y,x)=(\Phi a)(y,x)-\lambda(x)\;.
\end{equation}
The block partition makes $\lambda$ a periodic site field, and $\lambda(y)=0$. Every tree bond of $b$ vanishes. Moreover $Qb= Qa$ and $\operatorname{curl}b=\operatorname{curl}a$.

The Lie-algebra element $E(y,x)$ is a mean of the signed sums of $a$ around the closed walks obtained by adjoining a mean path from $y$ to $x$ to the reverse tree path. Each walk lies inside its block and has length at most $12(L-1)$. To fill it, commute adjacent steps in different directions and cancel opposite steps in each direction. Every commutation adds one unit plaquette. It remains inside the block, because exchanging two different directions preserves the intermediate coordinate ranges. Sorting the directions takes at most the square of the walk length commutations. Thus a filling has integer plaquette coefficients with total absolute sum at most $J_L=144(L-1)^2$. Cauchy--Schwarz and averaging give
\begin{equation}\label{eq:covariance-filling-bound}
    \| E\|^2\leqslant K_L\|\operatorname{curl}a\|^2,
    \qquad K_L= L^4J_L^2\;.
\end{equation}
There are $L^4$ site tests per block and the interior plaquettes of different blocks are disjoint, so this constant is independent of their number. Also $\|\lambda\|^2\leqslant 2\|\Phi a\|^2+2\| E\|^2$ and $\| D\lambda\|^2\leqslant 16\|\lambda\|^2$.

For $Qa=0$, equation \eqref{eq:covariance-two-averages} applied to $b$ has $\Phi b= E$. Jensen's inequality and the eight incidences of a block with coarse bonds imply
\begin{equation}\label{eq:covariance-constraint-error}
    \| Q_1b\|^2\leqslant 16L^{-6}\| E\|^2\;.
\end{equation}
For the block-tree gauge, the following periodic bound holds with $A_L=384L^5$:
\begin{equation}\label{eq:covariance-hard-bound}
    \| b\|^2\leqslant A_L
        (\|\operatorname{curl}b\|^2+L^2\| Q_1b\|^2)\;.
\end{equation}
For the proof, the coordinate permutation $R(x_1,x_2,x_3,x_4)=(x_4,x_3,x_2,x_1)$ reverses the block-tree order. Its pushforward on positive bonds is $(R_*b)(Rx,Rx+e_{5-\mu})= b(x,x+e_{\mu})$. It preserves the bond and curl norms, and \eqref{eq:covariance-two-averages} gives $(Q_1R_*b)(Rc)=(Q_1b)(c)$. Thus it suffices to prove the same bound for the reversed order $4,\ldots,1$, after translating a block corner to the origin. The inner product remains the fixed $\tr_N$ inner product throughout.

\paragraph{Direct face estimates in the positive-bond norm.}
The fine lattice has $N_{\mu}= LK_{\mu}$ sites in direction $\mu$, with odd $L\geqslant 3$ and $K_{\mu}\geqslant 3$. The preceding axis reversal permutes the periods to $(N_4,N_3,N_2,N_1)$ and reduces the proof to the tree order $4,\ldots,1$. Its inverse restores the original order at the end. A local calculation uses the translated block $B_0=\{0,\ldots,L-1\}^4$. All norms and sums below count positive bond or plaquette labels once, with the fixed $\mathfrak g$ inner product.

For a block $B$, let $\mathcal E_B^+$ be its positive fine bonds with both endpoints in $B$, and let $\mathcal P_B^+$ be its positive plaquettes with all vertices in $B$. For a boundary face $\Sigma\subset B$, let $\mathcal E_{\Sigma}^+$ be its positive fine bonds with both endpoints in $\Sigma$. Define
\begin{equation}\label{eq:direct-block-face-energies}
    I_B=\sum_{r\in \mathcal E_B^+}\| b(r)\|_{\mathfrak g}^2,
    \qquad
    T_{\Sigma}=\sum_{r\in \mathcal E_{\Sigma}^+}\| b(r)\|_{\mathfrak g}^2,
    \qquad
    F_B=\sum_{p\in \mathcal P_B^+}
        \|(\operatorname{curl}b)(p)\|_{\mathfrak g}^2\;.
\end{equation}
These are scalar quadratic forms in the existing tree-gauged field. For the flat comparison, a linear gauge change $a\mapsto a+D\omega$ preserving its coarse average has constant values $\omega(y)$ on the connected coarse lattice. Equation \eqref{eq:covariance-tree-gauge} then leaves $b$ unchanged. Constant colour conjugation sends $b$ to $\operatorname{Ad}_g b$ and preserves every inner product below. Only the flat reference covariance is estimated here.

\paragraph{The block-face trace bound.}
Write $b_{\nu}(x)= b(x,x+e_{\nu})$ for an internal positive bond of $B_0$, so $x_{\nu}\leqslant L-2$. The source tree condition makes $b_{\nu}(x)=0$ whenever $x_1=\cdots= x_{\nu-1}=0$; in particular $b_1=0$. For $k<\nu$ and integers $0\leqslant t< x_k$, define
\begin{equation}\label{eq:direct-telescoping-point}
    u^{(k)}(x,t)=(0,\ldots,0,t,x_{k+1},\ldots,x_4)\;.
\end{equation}
Here the $t$ is the $k$th component. Let $F_{k\nu}(u)$ be the curl on the positive $(k,\nu)$ plaquette based at $u$. On the plane $u_1=\cdots= u_{k-1}=0$, both $k$-direction sides of that plaquette vanish by the tree condition. Successive telescoping therefore gives the identity used before PRT2 (2.123) \cite{BalabanPRT2}:
\begin{equation}\label{eq:direct-tree-telescoping}
    b_{\nu}(x)=
        \sum_{k=1}^{\nu-1}\sum_{t=0}^{x_k-1}
            F_{k\nu}(u^{(k)}(x,t))\;.
\end{equation}
Every displayed plaquette lies inside $B_0$. The number of terms is at most $3(L-1)$.

Fix a face $x_{\mu}=\alpha$, where $\alpha\in \{0,L-1\}$, and a tangential component $\nu\ne \mu$. For a fixed $(k,\nu)$ plaquette in \eqref{eq:direct-tree-telescoping}, its multiplicity in the face sum is bounded by
\begin{equation}\label{eq:direct-face-multiplicities}
    \begin{array}{c| c}
        \mu< k & L^{k-2}\,L\leqslant L^2\\
        \mu= k & L^{k-1}\leqslant L^2\\
        \mu> k,\ \mu\ne \nu & L^{k-1}\,L\leqslant L^2
    \end{array}\;.
\end{equation}
In the first row, one of the $k-1$ otherwise free earlier components is fixed by the face. In the second row, $x_k$ is fixed, so a given $t$ occurs at most once. In the third row, the two distinct indices $\mu,\nu$ exceed $k$, hence $k\leqslant 2$. In all rows, later components are fixed by the plaquette label. Without the face restriction the multiplicity is at most $L^{k-1}L\leqslant L^3$. For different $\nu$, the plaquette pairs $(k,\nu)$ are disjoint. Cauchy--Schwarz in \eqref{eq:direct-tree-telescoping}, followed by these counts, proves
\begin{equation}\label{eq:direct-block-face-bound}
    I_B\leqslant 3(L-1)L^3F_B\leqslant 4L^4F_B,
    \qquad
    T_{\Sigma}\leqslant 3(L-1)L^2F_B\leqslant 4L^3F_B\;.
\end{equation}
The first estimate is the positive-label version of PRT2 (2.123). The second is the required face restriction of its telescoping calculation.

\paragraph{The interface Poincare estimate.}
Let $F=\{0,\ldots,L-1\}^3$ and let $\mathcal E_F^+$ be its positive nearest-neighbour edges. The Neumann graph Laplacian $\mathcal L_F \colon \mathfrak g^F \to \mathfrak g^F$ acts componentwise. For $s\in F$, it is given by
\begin{equation}\label{eq:direct-face-laplacian}
    (\mathcal L_Fv)(s)=
        \sum_{\substack{s'\in F\\ | s'-s|_1=1}}
            (v(s)-v(s'))\;.
\end{equation}
It has quadratic form $\sum_{(s,s')\in \mathcal E_F^+}\| v(s')-v(s)\|^2$. For $m\in \{0,\ldots,L-1\}^3$, its real scalar eigenbasis and eigenvalues are
\begin{equation}\label{eq:direct-face-spectrum}
    v_m(s)=\prod_{i=1}^3
        \cos\left(\frac{\pi m_i(s_i+1/2)}{L}\right),
    \qquad
    \lambda_m=4\sum_{i=1}^3
        \sin^2\left(\frac{\pi m_i}{2L}\right)\;.
\end{equation}
The one-dimensional cosine identity verifies the eigenvalue at interior and endpoint vertices; finite cosine orthogonality gives a basis. Only $m=0$ is constant. Since $\sin t\geqslant 2t/\pi$ for $0\leqslant t\leqslant \pi/2$, every nonzero eigenvalue is at least $4/L^2$. Apply this to each orthonormal colour component. For $\bar v= L^{-3}\sum_{s\in F}v(s)$, it follows that
\begin{equation}\label{eq:direct-face-poincare}
    \sum_{s\in F}\| v(s)-\bar v\|^2
        \leqslant \frac{L^2}{4}
            \sum_{(s,s')\in \mathcal E_F^+}
                \| v(s')-v(s)\|^2\;.
\end{equation}

\paragraph{Crossing bonds and the exact average.}
For a positive coarse bond $c$ in direction $\mu$, write its two endpoint blocks as $B_-= o+\{0,\ldots,L-1\}^4$ and $B_+= B_-+Le_{\mu}$. The transverse label set is $F_c=\{0,\ldots,L-1\}^{\{\nu\ne \mu\}}$, identified with $F$ by increasing transverse index. In that label set, $e_{\nu}$ denotes its transverse coordinate unit vector. For $s\in F_c$, put
\begin{equation}\label{eq:direct-interface-data}
    x_c(s)= o+(L-1)e_{\mu}+\sum_{\nu\ne \mu}s_{\nu} e_{\nu},
    \qquad
    t_c(s)= b(x_c(s),x_c(s)+e_{\mu}),
    \qquad
    J_c=\sum_{s\in F_c}\| t_c(s)\|^2\;.
\end{equation}
Let $\Sigma_-$ and $\Sigma_+$ be the facing boundary faces in $B_-$ and $B_+$. For $\nu\ne \mu$ and $s,s+e_{\nu}\in F_c$, the plaquette based at $x_c(s)$ has curl
\begin{equation}\label{eq:direct-interface-curl}
    t_c(s)-t_c(s+e_{\nu})
        +b_{\nu}(x_c(s)+e_{\mu})-b_{\nu}(x_c(s))\;.
\end{equation}
For $\nu<\mu$ its sign is opposite to the positive plaquette convention; its squared norm is unchanged. Let $F_c^{\mathrm{cr}}$ be the sum of these squared curls, each transverse edge once. The inequality $\| u+v+w\|^2\leqslant 3(\| u\|^2+\| v\|^2+\| w\|^2)$ gives
\begin{equation}\label{eq:direct-interface-gradient}
    \sum_{\nu\ne \mu}
        \sum_{\substack{s\in F_c\\s+e_{\nu}\in F_c}}
            \| t_c(s+e_{\nu})-t_c(s)\|^2
        \leqslant 3(F_c^{\mathrm{cr}}+T_{\Sigma_-}+T_{\Sigma_+})\;.
\end{equation}

Define $I_{B,\mu}$ as the part of $I_B$ in direction $\mu$, and set $\bar t_c= L^{-3}\sum_{s\in F_c}t_c(s)$. The straight-path formula \eqref{eq:covariance-two-averages} has the exact decomposition
\begin{equation}\label{eq:direct-interface-average}
    \begin{aligned}
        (Q_1b)(c)&= L^{-1}\bar t_c+q_c\;,\\
        q_c&= L^{-5}\sum_{s\in F_c}
                \sum_{i=0}^{L-2}(i+1)
                    b_{\mu}\left(o+ie_{\mu}+
                        \sum_{\nu\ne \mu}s_{\nu} e_{\nu}\right)\\
            &\quad+L^{-5}\sum_{s\in F_c}\sum_{j=0}^{L-2}(L-1-j)
                    b_{\mu}\left(o+(L+j)e_{\mu}+
                        \sum_{\nu\ne \mu}s_{\nu} e_{\nu}\right)
            \;.
    \end{aligned}
\end{equation}
Indeed, an internal lower-block bond at position $i$ belongs to $i+1$ starting paths, and an internal upper-block bond at position $j$ belongs to $L-1-j$ paths. Each crossing bond belongs to $L$ paths. This is the positive-label content of PRT2 (2.125). There are at most $2L^4$ internal bonds in this sum, and every coefficient is at most $L^{-4}$. Cauchy--Schwarz therefore yields
\begin{equation}\label{eq:direct-interface-mean}
    \| q_c\|^2\leqslant 2L^{-4}(I_{B_-,\mu}+I_{B_+,\mu}),
    \qquad
    L^3\|\bar t_c\|^2
        \leqslant 2L^5\|(Q_1b)(c)\|^2
            +4L(I_{B_-,\mu}+I_{B_+,\mu})\;.
\end{equation}
Combining the exact mean/variance decomposition of $J_c$ with \eqref{eq:direct-face-poincare}, \eqref{eq:direct-interface-gradient} and \eqref{eq:direct-interface-mean} gives
\begin{equation}\label{eq:direct-interface-bound}
    J_c\leqslant \frac{3L^2}{4}
        (F_c^{\mathrm{cr}}+T_{\Sigma_-}+T_{\Sigma_+})
        +2L^5\|(Q_1b)(c)\|^2
        +4L(I_{B_-,\mu}+I_{B_+,\mu})\;.
\end{equation}

\paragraph{The periodic sum and its payments.}
All sums over $c$ run over the positive coarse bonds $E_1$, and $\mu(c)$ is the direction of $c$. Let $\mathscr I=\sum_B I_B$, $\mathscr F=\|\operatorname{curl}b\|^2$ and $\mathscr Q=\| Q_1b\|^2$. Every positive fine bond is either internal to exactly one block or crosses exactly one positive coarse interface, so $\| b\|^2=\mathscr I+\sum_cJ_c$. Every plaquette counted by $F_c^{\mathrm{cr}}$ crosses only the $\mu$-interface: its transverse edge remains inside the face. It is therefore counted for at most one $c$. Plaquettes crossing two coarse interfaces are not needed. Each block has eight incident faces, and an internal $\mu$-bond enters the last term of \eqref{eq:direct-interface-bound} only for the incoming and outgoing $\mu$-interfaces. Consequently,
\begin{equation}\label{eq:direct-interface-incidences}
    \begin{aligned}
        \mathscr I&\leqslant 4L^4\mathscr F\;,\\
        \sum_cF_c^{\mathrm{cr}}&\leqslant \mathscr F\;,\\
        \sum_c(T_{\Sigma_-}+T_{\Sigma_+})&\leqslant 32L^3\mathscr F\;,\\
        \sum_{c \in E_1}
            (I_{B_-,\mu(c)}+I_{B_+,\mu(c)})&=2\mathscr I\;.
    \end{aligned}
\end{equation}
Thus summing \eqref{eq:direct-interface-bound} gives the explicit payment
\begin{equation}\label{eq:direct-hard-bound-payment}
    \begin{aligned}
        \| b\|^2
        &\leqslant (1+8L)\mathscr I
            +\frac{3L^2}{4}\mathscr F
            +24L^5\mathscr F+2L^5\mathscr Q\\
        &\leqslant \left(4L^4+56L^5+\frac{3L^2}{4}\right)\mathscr F
            +2L^5\mathscr Q\\
        &\leqslant 384L^5(\mathscr F+L^2\mathscr Q)\;.
    \end{aligned}
\end{equation}
The last inequality uses $L\geqslant 3$; its coefficients are bounded by $61L^5$ and $2L^5$, respectively. This proves \eqref{eq:covariance-hard-bound}.

Combining the proved bound \eqref{eq:covariance-hard-bound} with the constraint-error bound gives $\| b\|^2\leqslant A_L(1+16L^{-4}K_L)\|\operatorname{curl}a\|^2$. Since $a= b+D\lambda$, the bound $\| a\|^2\leqslant 2\| b\|^2+2\| D\lambda\|^2$ proves, for $Qa=0$,
\begin{equation}\label{eq:flat-coercivity}
    \| a\|_2^2\leqslant C_L
        \bigl(\|\operatorname{curl}a\|_2^2+G^{(2)}(a)\bigr)\;.
\end{equation}
An explicit admissible constant is
\begin{equation}\label{eq:covariance-coercivity-constant}
    C_L=\max(64,\,2A_L(1+16L^{-4}K_L)+64K_L)\;.
\end{equation}
At the flat background, differentiating the original action and exponential gauge functional in \eqref{eq:action-primitive} gives $\langle a,\Delta_1a\rangle=\|\operatorname{curl}a\|^2+\|\Phi a\|^2$. Substituting $a=\mathsf C B$ in \eqref{eq:flat-coercivity} uses $Q\mathsf C B=0$ and $\|\mathsf C B\|\geqslant \| B\|$, because the free bonds are unchanged. Consequently the free Hessian satisfies
\begin{equation}\label{eq:covariance-uniform-gap}
    \mathsf H\geqslant \gamma_L I,
    \qquad \gamma_L= C_L^{-1}\;.
\end{equation}
\paragraph{Construction on the infinite lattice.}
Let $E_{\infty}=\mathbb Z^4\times\{1,2,3,4\}$ denote positive fine bonds, with $(x,\mu)$ representing $(x,x+e_{\mu})$. The coarse centres are $L\mathbb Z^4$, with the centred blocks and path families specified above. Remove the designated bonds to obtain $E_{*,\infty}$. The square-summable field spaces use the same real Lie-algebra inner product as the finite-torus calculation. Initially on finitely supported fields, define $\Phi_{\infty}$ and $Q_{\infty}$ by \eqref{eq:covariance-mean-path}--\eqref{eq:covariance-two-averages}, and $\mathsf C_{\infty}$ by \eqref{eq:constraint-extension}. The designated coefficient is $Q_{0,\infty}= L^{-4}I$, so the inverse in that definition is fibrewise. Put
\begin{equation}\label{eq:bulk-operator-definition}
    \Delta_{\infty}=\operatorname{curl}^*\operatorname{curl}
        +\Phi_{\infty}^*\Phi_{\infty},
    \qquad
    \mathsf H_{\infty}=\mathsf C_{\infty}^*\Delta_{\infty}\mathsf C_{\infty}\;.
\end{equation}
The coefficient $Q_{0,\infty}$ can be checked directly in \eqref{eq:covariance-two-averages}. The within-block paths entering $\Phi$ contain no designated bond. For $b_0(c)$, precisely $L$ of the translated straight paths in the sum for $Q_1(c)$ traverse that bond: their transverse coordinates equal those of the coarse centres, and their longitudinal initial positions range over the $L$ sites of the block. No translated path for another coarse bond contains it. Hence its coefficient is $L^{-5}L= L^{-4}$ and all off-diagonal coefficients vanish.

The action normalisation also agrees with the finite-torus Hessian. At the identity, the first plaquette phase is $\operatorname{curl}a$ and the first mean-transport phase is $\Phi a$. The mixed second derivative of $1-\operatorname{Re}\tr_N e^{i(sX+tY)}$ at zero is $\tr_N(XY)$, the prescribed real inner product. The terms from second derivatives of the phases multiply the zero first derivative of this class function. Summing over the plaquettes and mean transports gives $\Delta_1=\operatorname{curl}^*\operatorname{curl}+\Phi^*\Phi$ with the same normalisation on every torus.

Each map has finite range, bounded coefficients and finitely many entry types modulo $L\mathbb Z^4$. Its row and column absolute sums are bounded, so it extends uniquely to a bounded operator on its square-summable carrier. In particular $\mathsf H_{\infty}$ is bounded and self-adjoint on all of $\ell^2(E_{*,\infty};\mathfrak g)$.

For every finitely supported $B$, the support of $\mathsf C_{\infty} B$ and all action terms meeting that latter support embed without wraparound in a sufficiently large compatible torus. Their quadratic form and norm are unchanged. Equation~\eqref{eq:covariance-uniform-gap} therefore applies to this same field. Density of finitely supported fields gives
\begin{equation}\label{eq:bulk-uniform-gap}
    \gamma_L I\leqslant \mathsf H_{\infty}\leqslant h_L I,
    \qquad
    h_L=\max_{e\in E_{*,\infty}/L\mathbb Z^4}
        \sum_{f\in E_{*,\infty}}\|\mathsf H_{\infty}(e,f)\|\;.
\end{equation}
The maximum uses representatives of the finitely many free-bond types. Symmetry gives the same column bound. Thus $h_L$ is finite and independent of the torus and coupling, with $h_L\geqslant \gamma_L>0$. Set $b_L= h_L+\gamma_L$ and $\vartheta_L=1-\gamma_L/b_L$. Then $0<\vartheta_L<1$ and $\| I-\mathsf H_{\infty}/b_L\|\leqslant \vartheta_L$. The inverse is
\begin{equation}\label{eq:bulk-neumann-inverse}
    \mathsf H_{\infty}^{-1}
        =\frac{1}{b_L}\sum_{j=0}^{\infty}
            (I-\mathsf H_{\infty}/b_L)^j\;.
\end{equation}
The series converges in operator norm. The fine-field covariance is the bounded positive operator $\mathsf K_{\infty}=\mathsf C_{\infty}\mathsf H_{\infty}^{-1}\mathsf C_{\infty}^*$. At the flat background every colour component has the same scalar operator; denote its scalar bond kernel by $K_{\infty}(b,b')$.

\paragraph{Decay with finite constants.}
Let $\ell(b)\in\mathbb Z^4$ be the coarse index of the block containing the initial point of $b$, and put $d(b,b')=\|\ell(b)-\ell(b')\|_{\infty}$. Let $R_H$ be the maximum of $1$ and the distances $d(e,f)$ for nonzero entries of $\mathsf H_{\infty}$. Let $R_C$ be the maximum of $d(b,e)$ for nonzero entries of $\mathsf C_{\infty}$. These maxima can be taken over the finitely many entry types modulo coarse translation. Define
\begin{equation}\label{eq:bulk-decay-constants}
    c_L^{\mathsf C}=\max_{b\in E_{\infty}/L\mathbb Z^4}
        \sum_{e\in E_{*,\infty}}\|\mathsf C_{\infty}(b,e)\|,
    \qquad
    \mu_L=\frac{-\log\vartheta_L}{R_H},
    \qquad
    A_L^{\mathrm{ker}}=\frac{(c_L^{\mathsf C})^2}{\gamma_L}e^{2\mu_LR_C}\;.
\end{equation}
The $j$th power in \eqref{eq:bulk-neumann-inverse} has zero $(e,f)$ entry if $jR_H< d(e,f)$, and its norm is at most $\vartheta_L^j$. Summing the remaining terms bounds the inverse entry by $\gamma_L^{-1}e^{-\mu_Ld(e,f)}$. Using the two finite rows of $\mathsf C_{\infty}$ and $d(e,f)\geqslant d(b,b')-2R_C$ yields
\begin{equation}\label{eq:bulk-kernel-decay}
    | K_{\infty}(b,b')|
        \leqslant A_L^{\mathrm{ker}}e^{-\mu_Ld(b,b')}\;.
\end{equation}
This constructs the inverse and decay in the exponential-gauge free carrier.

\paragraph{Comparison with the periodic covariance.}
Take a compatible torus with fine periods $N_{\alpha}= LK_{\alpha}$ and $K_{\alpha}\geqslant 3$. For $z\in\mathbb Z^4$, write $Nz=(N_1z_1,\ldots,N_4z_4)$. The local torus definitions of $Q,\Phi,\mathsf C,\Delta_1$ are the periodisations of their infinite-lattice definitions. Periodisation preserves their finite-range products. The absolutely convergent periodisation of the free inverse therefore solves the finite-torus inverse equation: applying $\mathsf H_T$ gives the periodised Kronecker delta. Uniqueness of the inverse and multiplication by the two constraint-extension maps give
\begin{equation}\label{eq:bulk-periodisation}
    K_T(b,b')=\sum_{z\in\mathbb Z^4}K_{\infty}(b,b'+Nz)\;.
\end{equation}
Bonds on the right have fixed lifts and unchanged orientations. Let $K_{\min}=\min_{\alpha} K_{\alpha}$ and define
\begin{equation}\label{eq:bulk-image-constant}
    B_L^{\mathrm{img}}
        =\sum_{z\in\mathbb Z^4\setminus\{0\}}
            e^{-\mu_L(\| z\|_{\infty}-1)}<\infty\;.
\end{equation}
The shell of maximum norm $j$ has $(2j+1)^4-(2j-1)^4$ points. Since $d(b,b'+Nz)\geqslant K_{\min}\| z\|_{\infty}-d(b,b')$, the image sum satisfies
\begin{equation}\label{eq:bulk-image-error}
    | K_T(b,b')-K_{\infty}(b,b')|
        \leqslant A_L^{\mathrm{ker}}B_L^{\mathrm{img}}
            e^{-\mu_L(K_{\min}-d(b,b'))}\;.
\end{equation}
Thus every fixed finite collection of covariance entries converges as all coarse periods tend to infinity. The limiting connected-cluster construction uses the infinite-lattice covariance, not a periodically repeated finite-volume covariance.

\section{Background Estimates at the First Step}
\label{app:background-estimates}

This appendix proves the estimates stated in Section~\ref{sec:sourced-effective-interaction}: Lemmas~\ref{lem:regular-bounds}, \ref{lem:charts}, \ref{lem:gaussian-payment}, \ref{lem:activity-bounds} and~\ref{lem:vacuum-clusters}. It fixes every constant named there that Section~\ref{sec:fixed-first-kernel} does not fix itself, and places all of them in the one order of choice stated in Section~\ref{app:thresholds}: $L$ and the group with the group constants $c_{\mathrm h}=\sqrt N$, $c_{\mathrm n}=1/N$, $\varrho_G=\pi/\sqrt N$, $c$, $r_G$, $c_G$ below; the constants of the lemmas of this appendix in the order of their proofs; $M_1$, $M_2$ and $M$ with the constants chosen after them, including the collar constants of item 4 there; the coupling $g_0$, through the thresholds collected there; and last the torus, whose periods satisfy the conditions of \eqref{eq:c-standing-conditions} (multiples of $LMR_1$, at least $8M$; the scale $R_1$ depends on $g_0$, so the period conditions are imposed after $g_0$). Group constants: $c$ is a constant with $\| e^{iX}-\id\|_{\mathrm{op}}\leqslant c\,h(X)$ for all $X\in\mathfrak g$ ($c=\sqrt N$ serves, as $\| e^{iX}-\id\|_{\mathrm{op}}\leqslant \| X\|_{\mathrm{op}}\leqslant \sqrt Nh(X)$). For every $N\times N$ matrix $Y$ put $h(Y)=(N^{-1}\tr Y^*Y)^{1/2}$, which extends the normalised Hilbert--Schmidt norm of $\mathfrak g$ to all matrices; then $\| Y\|_{\mathrm{op}}\leqslant (\tr Y^*Y)^{1/2}=\sqrt N\,h(Y)$, so $\| Y\|_{\mathrm{op}}\leqslant c_{\mathrm h}| Y|$ for every matrix $Y$ and either cutoff norm $|\bigcdot |$ (the operator norm or $h$), not only for Hermitian $Y$; this is used below for differences of unitaries. The exponential map restricts to an analytic diffeomorphism from the ball $\{ X\in\mathfrak g\, | \,h(X)<\varrho_G\}$ onto an open neighbourhood $N_G$ of the identity in $G$ (the conventions of Section~\ref{sec:sourced-effective-interaction}), with inverse $\log_G$; since $G$ is compact and $\| g-\id\|_{\mathrm{op}}$ vanishes on $G$ only at the identity, there is $r_G\in(0,1]$ with $\{ g\in G\, | \,\| g-\id\|_{\mathrm{op}}\leqslant r_G\}\subset N_G$. The constant $c_G$ with $h(\log_Gg)\leqslant c_G\| g-\id\|_{\mathrm{op}}$ on that set exists: $\log_G$ is analytic, so Lipschitz with some constant $c'$ on $\{\| g-\id\|_{\mathrm{op}}\leqslant r'\}$ for some $r'\in(0,r_G]$, and on $r'\leqslant \| g-\id\|_{\mathrm{op}}\leqslant r_G$ one has $h(\log_Gg)<\varrho_G\leqslant (\varrho_G/r')\| g-\id\|_{\mathrm{op}}$; take $c_G=\max(c',\varrho_G/r')$. The radius $r_M$ of the local group mean is fixed in Lemma~\ref{lem:c2-uniform-chart}(iv). The phase norm of Lemmas~\ref{lem:two-block} and~\ref{lem:hessian-gauge-locality} is $\| A\|_{\infty}=\max_bh(A_b)$, and operator norms of entries and their derivatives are taken with respect to the Hilbert norms induced by $h$. Throughout, a branch is a branch in the sense of Section~\ref{sec:branch-data} (an admissible pair with $(A_{\mathrm e},a)\in\mathcal D_{\xi}$ for some $a$); the reference data of Section~\ref{sec:joint-extraction} are treated in Section~\ref{app:reference-check}.

\subsection{Cubes, localisation domains and tree lengths}
\label{app:domain-geometry}

The partitions $\pi_0$ (side $M$) and $\pi_1$ (side $LM$) of Section~\ref{sec:fixed-first-kernel} have cubes $Q_z= Mz+[-M/2,M/2]^4$ and $\Delta_z= LMz+[-LM/2,LM/2]^4$, reduced modulo the periods, with $M_1$ and $M$ odd. Cube faces therefore lie at half-integer fine coordinates: every lattice site lies in the interior of exactly one cube of each partition, a fine bond meets at most two cubes of each partition, and a fine bond lies in the interior of a union of cubes exactly when both its endpoints lie in cubes of that union. A cube has $M^4$ sites and meets at most $8M^4$ positive fine bonds. Two cubes of the same partition are wall adjacent when they share a three-dimensional face; a localisation domain of $\pi_i$ is a finite nonempty union of cubes of $\pi_i$ that is connected under wall adjacency (RG1 p.257). A localisation domain $Y$ of $\pi_i$ has $N_i(Y)$ cubes and, following RG1 p.257, tree length
\begin{equation}\label{eq:c2-tree-length}
    d_i(Y)=\inf\{\,|\tau|\, | \,\tau\subset Y\text{ a connected finite union of segments meeting every cube of }Y\,\},
\end{equation}
the Euclidean length $|\tau|$ being measured in units of the cube side. Put $\ell_i(Y)= d_i(Y)+5$. Two domains are compatible when they share neither a cube nor a wall (RG2 (2.11)); a domain is incompatible with itself.

\begin{lemma}[geometry of localisation domains]\label{lem:c2-domain-geometry}
Let $Y,Y_1,\ldots,Y_m$ be localisation domains of one partition and $C_{\mathrm g}=162$.
\begin{enumerate}
\item $N_i(Y)\leqslant C_{\mathrm g}\ell_i(Y)$ and $\ell_i(Y)\leqslant N_i(Y)+4$.
\item If $\bigcup_jY_j$ is a localisation domain, then $\ell_i(\bigcup_jY_j)\leqslant \sum_j\ell_i(Y_j)$.
\item If $Z$ is a localisation domain containing $Y_1,\ldots,Y_k$ ($k\geqslant 0$, pairwise disjoint) and $Z\setminus\bigcup_jY_j$ consists of $n$ cubes, then $\ell_i(Z)\leqslant \sum_j\ell_i(Y_j)+5n$.
\item The number of localisation domains with $n$ cubes containing a given cube is at most $64^{n-1}$, and for $s> C_{\mathrm g}\log64$
\begin{equation}\label{eq:native-source-animal-sum}
    \sum_{Y\ni Q}e^{-s\ell_i(Y)}\leqslant B(s)=\frac{e^{-s/C_{\mathrm g}}}{1-64e^{-s/C_{\mathrm g}}}
\end{equation}
for every cube $Q$, on every torus. For $s\geqslant C_{\mathrm g}(\log64+\log2)$, $B(s)\leqslant 2e^{-s/C_{\mathrm g}}$.
\end{enumerate}
\end{lemma}

\begin{proof}
(i) Let $\tau$ realise \eqref{eq:c2-tree-length} up to $\varepsilon\in(0,1/2)$. A closed traversal of $\tau$ (each segment twice) has length at most $2(d_i(Y)+\varepsilon)$; place sample points along it at arclength spacing one, at most $2d_i(Y)+2$ of them. Every cube of $Y$ contains a point of $\tau$, which is within arclength $1/2$, hence Euclidean and $\ell^{\infty}$ distance $1/2$, of a sample; the centre of that cube is then within $\ell^{\infty}$ distance one of the sample. Cube centres form a unit lattice, so at most $3^4=81$ cubes are assigned to each sample, and $N_i(Y)\leqslant 81(2d_i(Y)+2)\leqslant 162\,\ell_i(Y)$. For the second inequality join the centres of wall-adjacent cubes along a spanning tree of the wall-adjacency graph; each of the $N_i(Y)-1$ edges has length one and lies in $Y$, so $d_i(Y)\leqslant N_i(Y)-1$.

(ii) Take trees $\tau_j\subset Y_j$ with $|\tau_j|\leqslant d_i(Y_j)+\varepsilon$. The graph on $\{ Y_j\}$ with an edge when $Y_j\cap Y_k\ne\varnothing$ is connected: a wall-adjacent chain of cubes of the union passes from $Y_j$ to $Y_k$ through consecutive cubes that share a wall, hence through domains that intersect. For each edge of a spanning tree pick $x\in Y_j\cap Y_k$, a cube $Q_j\ni x$ of $Y_j$ and a point $p_j\in\tau_j\cap Q_j$, and add the segment $[p_j,x]\subset Q_j$ (length at most the diameter $2$ of a unit cube in $\mathbb R^4$), and likewise for $k$. The result is a connected union of segments inside $\bigcup Y_j$ meeting every cube, of length at most $\sum_j(d_i(Y_j)+\varepsilon)+4(m-1)$. Letting $\varepsilon\to0$, $d_i(\bigcup Y_j)+5\leqslant \sum_jd_i(Y_j)+4m+1\leqslant \sum_j(d_i(Y_j)+5)$ for $m\geqslant 1$.

(iii) Apply (ii) to the family consisting of $Y_1,\ldots,Y_k$ and the $n$ single cubes, each of which has $\ell_i=5$.

(iv) A depth-first traversal of a spanning tree of the wall-adjacency graph of a domain with $n$ cubes, started at the given cube, is a word of length $2(n-1)$ over the eight wall directions, and determines the domain; this gives $64^{n-1}$. By (i), $e^{-s\ell_i(Y)}\leqslant e^{-sN_i(Y)/C_{\mathrm g}}$, so the rooted sum is at most $\sum_{n\geqslant 1}64^{n-1}e^{-sn/C_{\mathrm g}}= B(s)$. The last claim is $64e^{-s/C_{\mathrm g}}\leqslant 1/2$.
\end{proof}

\begin{lemma}[scale comparison]\label{lem:c2-scale}
Let $Z$ be a localisation domain of $\pi_0$, $\widetilde Z$ the union of the $M$-cubes meeting $Z$, and $Z'$ the union of the $LM$-cubes meeting $\operatorname{int}\widetilde Z$. This is the enlargement rule of Section~\ref{sec:fixed-first-kernel}, which takes the $LM$-cubes meeting $\operatorname{int}Q^{\sim}$ for some cube $Q$ of $Z$: indeed $\widetilde Z=\bigcup_{Q\subseteq Z}Q^{\sim}$, an $LM$-cube meeting $\operatorname{int}Q^{\sim}$ meets $\operatorname{int}\widetilde Z$, and conversely an $LM$-cube $\Delta$ meeting $\operatorname{int}\widetilde Z$ contains a nonempty open subset of $\widetilde Z$, a finite union of the closed sets $Q^{\sim}$, one of which therefore contains a nonempty open subset of $\Delta$ (a finite union of closed sets with empty interiors has empty interior), so $\Delta$ meets $\operatorname{int}Q^{\sim}$ for that $Q$. The same argument shows, for a union $Z_0$ of $M$-cubes with wall components $Z_{0,j}$, that $Z_0'=\bigcup_jZ_{0,j}'$. Then $Z\subset\operatorname{int}\widetilde Z$, $\widetilde Z\subset Z'$, $Z'$ is a localisation domain of $\pi_1$, every point of $\widetilde Z$ is within $\ell^{\infty}$ distance $2M$ of every tree realising $d_0(Z)$ up to any $\varepsilon>0$, and
\begin{equation}\label{eq:c2-scale-inequality}
    \ell_1(Z')\leqslant a^{\mathrm{sc}}_L\ell_0(Z)+b^{\mathrm{sc}}_L,\qquad a^{\mathrm{sc}}_L=\frac{2049}{L},\quad b^{\mathrm{sc}}_L=5+\frac{1024}{L}\;.
\end{equation}
\end{lemma}

\begin{proof}
Every $M$-cube meeting a point of $Z$ belongs to $\widetilde Z$, so $Z\subset\operatorname{int}\widetilde Z$. An $M$-cube $\widetilde Q$ of $\widetilde Z$ lies in a single $LM$-cube, which meets $\operatorname{int}\widetilde Q\subset\operatorname{int}\widetilde Z$; hence $\widetilde Z\subset Z'$. If $\widetilde Q$ meets $Q\subset Z$ only at a lower-dimensional face, the cubes between them sharing walls also meet $Q$, so $\widetilde Z$ is wall connected, $\operatorname{int}\widetilde Z$ is connected, and the $LM$-cubes meeting a connected open set are wall connected (a path avoiding lower-dimensional faces passes between consecutive cubes through a common wall). Let $\tau\subset Z$ meet every $M$-cube of $Z$ with $|\tau|\leqslant M(d_0(Z)+\varepsilon)$ in fine units. For $x\in\widetilde Z$ pick $\widetilde Q\ni x$, a cube $Q\subset Z$ meeting $\widetilde Q$, $y\in\widetilde Q\cap Q$ and $p\in\tau\cap Q$; then $| x-p|_{\infty}\leqslant | x-y|_{\infty}+| y-p|_{\infty}\leqslant 2M$.

For \eqref{eq:c2-scale-inequality}, place samples along a closed traversal of $\tau$ at arclength spacing $LM$; there are at most $2(d_0(Z)+\varepsilon)/L+1$ of them, and every point of $\tau$ is within Euclidean distance $LM/2$ of a sample. For $\Delta\subset Z'$ pick $x_{\Delta}\in\Delta\cap\operatorname{int}\widetilde Z$; by the preceding paragraph it is within $\ell^{\infty}$ distance $2M+LM/2\leqslant 7LM/6$ of a sample ($L\geqslant 3$), so the centre of $\Delta$ is within $\ell^{\infty}$ distance $5LM/3<2LM$ of that sample. The centres of $LM$-cubes form a lattice of spacing $LM$, and an open $\ell^{\infty}$ ball of radius $2LM$ contains at most $4^4=256$ of its points. Hence $N_1(Z')\leqslant 256(2(d_0(Z)+\varepsilon)/L+1)$. Now build a connected union of segments in $Z'$ meeting every $LM$-cube of $Z'$: start from $\tau\subset Z\subset Z'$ and, for each $\Delta\subset Z'$ not met by $\tau$, add the segments $[x_{\Delta},y]\subset\widetilde Q\subset Z'$ and $[y,p]\subset Q\subset Z'$ chosen above, each of length at most $2M$. Its length is at most $M(d_0(Z)+\varepsilon)+4MN_1(Z')$, so
\begin{equation*}
    d_1(Z')\leqslant \frac{d_0(Z)+\varepsilon}{L}+\frac{4}{L}\cdot256\Bigl(\frac{2(d_0(Z)+\varepsilon)}{L}+1\Bigr)\leqslant \frac{2049}{L}(d_0(Z)+\varepsilon)+\frac{1024}{L}\;.
\end{equation*}
Let $\varepsilon\to0$, add $5$ to both sides and use $d_0(Z)\leqslant \ell_0(Z)$.
\end{proof}

\subsection{The regular-background bounds}
\label{app:regular-covariance}

\begin{lemma}[entry range and stencils]\label{lem:c2-entry-range}
For free bonds $b,b'$ let $z(b)$ be the coarse label of the block containing $b_-$ and $f_{bb'}(U)$ the entry function \eqref{eq:c-entry-function} with its entry stencil $\mathcal S_{bb'}$ (the union of the background sets of the terms of \eqref{eq:regular-stencil-quadratic-decomposition} whose free-input set contains both $b$ and $b'$). If $| z(b)-z(b')|_1>2$ then $\mathcal T_{bb'}=\varnothing$ and $f_{bb'}=0$. The stencil of a nonzero entry consists of bonds $e$ with $| z(e)-z(b)|_1\leqslant 2$; in fine units all bonds and sites of $\mathcal S_{bb'}$, and both endpoints of $b$ and $b'$, lie within $\ell^{\infty}$ distance $3L$ of $b_-$, so $\rho_L=4L$ is an admissible entry range in Lemma~\ref{lem:hessian-gauge-locality}. Each row has at most $164L^4$ nonzero entries.
\end{lemma}

\begin{proof}
By \eqref{eq:regular-stencil-quadratic-decomposition} and its support sets, a term containing the free input $b$ is $\mathfrak a_p$ with $b\in\mathcal F_r$ for a boundary bond $r$ of $p$, $\mathfrak g_y$ with $b\in\mathcal E_y$, or $\mathfrak t_c$ with $b\in\mathcal E_c\cap E_*$. The free inputs of a bond $r$ have block labels among the endpoint blocks of $r$ (Lemma~\ref{lem:two-block}): $z(r)$ for a free bond, $\{ c_-,c_+\}$ for $r= b_0(c)$. The four vertices of a fine plaquette in directions $\mu,\nu$ have block labels differing only in the $\mu$ and $\nu$ coordinates, by at most one each, so these labels, and the endpoint blocks of its boundary bonds, lie in a set of $\ell^1$ diameter at most two. The free inputs of $\mathfrak g_y$ lie in one block, those of $\mathfrak t_c$ in the two blocks of $c$. Hence two free inputs of one term have $| z(b)-z(b')|_1\leqslant 2$, and the background bonds of that term (plaquette bonds, the sets $\mathcal E_c$, and for $\mathfrak t_c$ the plaquettes meeting $b_0(c)$, whose vertices lie in the two blocks of $c$) have block labels within $\ell^1$ distance two of $z(b)$. A site of a block with label within $\ell^1$ distance two of $z(b)$ is within $\ell^{\infty}$ distance $2L+(L-1)<3L$ of $b_-$; the endpoints of $b$ and $b'$ are such sites or within one of them. The $\ell^1$ ball of radius two in $\mathbb Z^4$ has $41$ points and a block has at most $4L^4$ positive bonds with initial site in it.
\end{proof}

\begin{lemma}[plaquette deviation on the collar]\label{lem:c2-plaquette-deviation}
Let $B_3$, $\delta_{\mathrm c}$ and the $O(1)$ be the constants of the comparison following CRE (1.19), as recorded in (B3)(iii): for every $L^2M_2R_1$-cube $\square'\subseteq\Omega_1^{\sim-1}\setminus R_1$ (CRE's $\square'\subset R_1^c$, for which $\square'^{\sim4}\subseteq\Omega_1$, p.250), $| U_1(e)U_{1,\square'}(e)^{-1}-1|< O(1)B_3e^{-\delta_{\mathrm c}LM_2R_1}\varepsilon_1$ on the bonds $e$ of $\square'^{\sim2}$ (p.250). CRE states no dependence for these constants: it reuses the symbols $B_3$ and $\delta$ at every step (p.268, the regularity bound $|\partial U_{k+1}-1|<2B_3\varepsilon_{k+1}(L^{-1}\eta)^2$ and the comparison with $B_3\exp(-\delta LM_2R_{k+1})44d^2B_3\varepsilon_{k+1}$) and derives the comparison (3.19) from these facts and the bounds (106)--(108), (159)--(163) of [12] (p.268). This paper uses one property of them only, the reading recorded in (B3)(iii): $B_3$, $\delta_{\mathrm c}$ and the $O(1)$ do not depend on the coupling $g_0$ or on the branch, so that at fixed $L$ and $M_2$ the factor $e^{-\delta_{\mathrm c}LM_2R_1}$ tends to zero as $g_0\to0$. CRE's $\delta$ is written $\delta_{\mathrm c}$ here to distinguish it from the Mayer constant $\delta$ of Section~\ref{app:mayer-constants}. Reduce the coupling threshold until $c_{\mathrm h}\,O(1)B_3e^{-\delta_{\mathrm c}LM_2R_1}\leqslant 1/4$, which holds for small $g_0$ because the constants do not depend on $g_0$ and $R_1\geqslant (\log g_1^{-2})^{\mathfrak r}\to\infty$ as $g_0\to0$ (CRE (2.5) and (2.6), p.255: $g_1\leqslant (1+\beta_0)g_0$ and $g_0\leqslant (1+\beta_0)g_1$, so $g_1\to0$ with $g_0$), and let $R_{\mathrm c}$ satisfy Lemma~\ref{lem:collar-coverage}. Here $\varepsilon_1= g_1p_0(g_1)$ is CRE's constant of the same name (p.246, CRE (2.4)) and $g_1$ the first-step coupling. On every branch, every fine plaquette $p$ having a bond within distance $3R_{\mathrm c}$ of $I$ satisfies
\begin{equation}\label{eq:c2-plaquette-deviation}
    \| U_1(\partial p)-\id\|_{\mathrm{op}}\leqslant 2c_{\mathrm h}\varepsilon_1\;.
\end{equation}
\end{lemma}

\begin{proof}
Let $b$ be such a bond and $\square'$ the cube of the partition of CRE (1.18) into cubes of side $L^2M_2R_1$ that contains $b_-$ (a partition, so the cube is unique). By Lemma~\ref{lem:collar-coverage}(i) (stated for $d(b,I)\leqslant 3R_{\mathrm c}$; the standing condition $LMR_1>4R_{\mathrm c}+4L^2M_2R_1+8$ gives $3R_{\mathrm c}+1+L^2M_2R_1< LMR_1-1$), $b_-\in\Omega_1^{\sim-1}$ and $\square'$ is a good cube of CRE (1.18), $\square'\subset R_1^c\subset\Omega_1^{\sim-1}$; since $M\geqslant 8LM_2$, the layer of $LMR_1$-cubes around $\Omega_1^{\sim-1}$ inside $\Omega_1$ contains $\square'^{\sim4}$, as CRE states after (1.19). The outer weight \eqref{eq:native-source-explicit-outer-operation} contains the background test $\chi_1(\Omega_1)$ of CRE (2.17) at $k=1$, a function of $V$, and it is nonzero on a branch: the pair is admissible, so $D^{(11)}(V,V_0)\ne\varnothing$ (Section~\ref{sec:branch-data}), and $\chi_1(\Omega_1)\ne0$ at every point of $D^{(11)}$. That test is written on the lattice $T_{\eta}$, $\eta= L^{-1}$, whose unit is the block side, so its cubes of side $LM_2R_1$ have fine side $L^2M_2R_1$ and the plaquette threshold $\varepsilon_1\eta^2$ is $\varepsilon_1L^{-2}$; its partition is ``compatible with all other partitions'' (CRE p.257; CRE does not define the term, pp.245 and 257, and it is read here as: the cubes of the coarser of two compatible partitions are unions of cubes of the finer), in particular with the partition of (1.18) into cubes of the same side, and two compatible partitions into cubes of equal side coincide. Hence on the branch $\sup_{p'\subset\square'^{\sim}}| U_{1,\square'}(\partial p')-1|<\varepsilon_1L^{-2}$ for the local minimum $U_{1,\square'}$ of CRE (2.16); a gauge transformation conjugates the plaquette holonomy and $|\bigcdot |$ is $\operatorname{Ad}$-invariant, so the bound holds for the axial-gauge representative of $U_{1,\square'}$ (p.247) to which the comparison (B3)(iii) refers. The plaquette $p$ contains $b$, hence lies in $\square'^{\sim}$. The comparison (B3)(iii), stated for the cubes $\square'\subset R_1^c$, for which $\square'^{\sim4}\subset\Omega_1$ and $U_1$, $U_{1,\square'}$ are determined by the same averages on $\square'^{\sim3}$, applies to this $\square'$ and gives $| U_1(e)U_{1,\square'}(e)^{-1}-1|< O(1)B_3e^{-\delta_{\mathrm c}LM_2R_1}\varepsilon_1$ on the bonds $e$ of $\square'^{\sim2}$, hence $\| U_1(e)-U_{1,\square'}(e)\|_{\mathrm{op}}=\| U_1(e)U_{1,\square'}(e)^{-1}-\id\|_{\mathrm{op}}\leqslant c_{\mathrm h}| U_1(e)U_{1,\square'}(e)^{-1}-\id|\leqslant \varepsilon_1/4$ there (unitary invariance of the operator norm, and the norm comparison of the preamble, valid for every matrix). Writing the plaquette holonomy as an ordered product of four unitaries and telescoping, $\| U_1(\partial p)-U_{1,\square'}(\partial p)\|_{\mathrm{op}}\leqslant \sum_{e\subset\partial p}\| U_1(e)-U_{1,\square'}(e)\|_{\mathrm{op}}\leqslant \varepsilon_1$, and $c_{\mathrm h}\varepsilon_1L^{-2}+\varepsilon_1\leqslant 2c_{\mathrm h}\varepsilon_1$.
\end{proof}

\begin{lemma}[local gauge on a rectangle]\label{lem:c2-local-gauge}
Let $\mathfrak P$ be a product of fine integer intervals of side at most $s_{\mathfrak P}$ (a lifted rectangle, injective modulo the periods), and let $U$ be a configuration whose plaquettes inside $\mathfrak P$ satisfy $\| U(\partial p)-\id\|_{\mathrm{op}}\leqslant 2c_{\mathrm h}\varepsilon_1$ with $6c_{\mathrm h}s_{\mathfrak P}\varepsilon_1\leqslant r_G$. There is a site gauge $u$ on $\mathfrak P$ such that every bond $b\subset\mathfrak P$ has $U^u(b)= e^{iA_b}$ with $A_b\in\mathfrak g$ and
\begin{equation}\label{eq:c2-local-gauge-bound}
    h(A_b)\leqslant \widehat\mu_{\mathfrak P}=6c_Gc_{\mathrm h}s_{\mathfrak P}\varepsilon_1\;.
\end{equation}
\end{lemma}

\begin{proof}
Fix the corner $x_0$ of $\mathfrak P$ with minimal coordinates. For $x\in\mathfrak P$ let $P(x)$ be the lattice path from $x_0$ to $x$ that runs first in direction $e_1$, then $e_2,e_3,e_4$; it stays in $\mathfrak P$. These paths form a spanning tree; define $u(x_0)=\id$ and $u(x')= u(x)U(x,x')$ along tree bonds oriented away from the root. Then $U^u(x,x')= u(x)U(x,x')u(x')^{-1}=\id$ on tree bonds, and for any bond $(x,x')\subset\mathfrak P$, $U^u(x,x')= U(P(x))U(x,x')U(P(x'))^{-1}$ is the holonomy of the closed loop $P(x)\cdot(x,x')\cdot P(x')^{-1}$ based at $x_0$. Let $(x,x')$ have direction $e_{\mu}$, so $x'= x+e_{\mu}$. The two tree paths share their runs in directions $\nu<\mu$ and their first $x_{\mu}-x_{0,\mu}$ steps in direction $\mu$; after conjugation by the holonomy of that common initial segment, which preserves $\|\bigcdot -\id\|_{\mathrm{op}}$, the loop becomes $B\cdot(x,x')\cdot B'^{-1}\cdot(y+e_{\mu},y)$, where $B=(b_1,\ldots,b_k)$ is the staircase from $y=(x_1,\ldots,x_{\mu},x_{0,\mu+1},\ldots,x_{0,4})$ to $x$ in directions $\nu>\mu$, $B'$ its translate by $e_{\mu}$, and $k\leqslant 3s_{\mathfrak P}$. Writing $\mathcal L_j$ for the loop built from $(b_1,\ldots,b_j)$, one checks $\mathcal L_j=(b_1\cdots b_{j-1})\,\partial p_j\,(b_1\cdots b_{j-1})^{-1}\cdot\mathcal L_{j-1}$, where $\partial p_j= b_je_{\mu} b_j'^{-1}e_{\mu}^{-1}$ is the boundary of the plaquette spanned by $b_j$ and $e_{\mu}$, which lies in $\mathfrak P$. Hence the holonomy of the loop is an ordered product of $k$ conjugates of plaquette holonomies, and telescoping gives $\| U^u(x,x')-\id\|_{\mathrm{op}}\leqslant 2kc_{\mathrm h}\varepsilon_1\leqslant 6c_{\mathrm h}s_{\mathfrak P}\varepsilon_1\leqslant r_G$. The element $g= U^u(x,x')\in G$ therefore lies in $N_G$, and $A_b=\log_Gg\in\mathfrak g$ is the unique $X$ in the ball $h(X)<\varrho_G$ with $e^{iX}= g$ (the group chart of the conventions), with $h(A_b)\leqslant c_G\| g-\id\|_{\mathrm{op}}\leqslant 6c_Gc_{\mathrm h}s_{\mathfrak P}\varepsilon_1$.
\end{proof}

The proof uses $\varepsilon_1$ only through the plaquette hypothesis and the condition $6c_{\mathrm h}s_{\mathfrak P}\varepsilon_1\leqslant r_G$; so the lemma holds with any $\varepsilon>0$ in place of $\varepsilon_1$, with the bound $6c_Gc_{\mathrm h}s_{\mathfrak P}\varepsilon$ in \eqref{eq:c2-local-gauge-bound}. Lemma~\ref{lem:c2-fibre-smallness} uses it with $\varepsilon=\varepsilon_0/2$.

Here are the fixed geometric choices. Reduce $\tau_L$ of Lemma~\ref{lem:two-block} so that
\begin{equation}\label{eq:c2-tau-final}
    \tau_L\leqslant \min\Bigl(1,\ 2c_Gr_G,\ \frac{r_G}{2(L-1)c},\ \frac{r_M}{4(L-1)c\,c_G}\Bigr)\;,
\end{equation}
with $c$, $r_G$, $c_G$ the group constants of the preamble (depending on the group only) and $r_M$ the radius of Lemma~\ref{lem:c2-uniform-chart}(iv) (depending on $L$ and the group); this is the final value of $\tau_L$, used in every statement. Put
\begin{equation}\label{eq:regular-covariance-radius-replay}
    c_{1,L}=320000\,m_L,\qquad
    R_{\mathrm p}=\max\Bigl(2,\Bigl\lceil2\sqrt{c_{1,L}/\gamma_L}\Bigr\rceil\Bigr),\qquad
    s_{\mathfrak P}= L(2R_{\mathrm p}+5)+1,\qquad
    C_{L,R_{\mathrm c}}=6c_Gc_{\mathrm h}s_{\mathfrak P}\;,
\end{equation}
and choose the collar range $R_{\mathrm c}\geqslant \max(2R_{\mathrm{loc}}+1,\ s_{\mathfrak P}+2)$. Require every coarse period to exceed $2R_{\mathrm p}+8$, so that the coarse $\ell^{\infty}$ ball of radius $R_{\mathrm p}+3$ about any coarse label lifts injectively, and impose on $\varepsilon_1$
\begin{equation}\label{eq:convex-background-thresholds}
    \widehat\mu= C_{L,R_{\mathrm c}}\varepsilon_1<\frac{\tau_L}{2},\qquad
    \max(1,c_L)\widehat\mu<\frac{\gamma_L}{4}\;.
\end{equation}
Since $\tau_L\leqslant 2c_Gr_G$, the first condition gives $6c_{\mathrm h}s_{\mathfrak P}\varepsilon_1=\widehat\mu/c_G< r_G$, the hypothesis of Lemma~\ref{lem:c2-local-gauge}.

\begin{lemma}[regional entries]\label{lem:c2-regional-entry}
Let $\Omega$ be a region and $U$ a background on which the regional form $\mathsf H^{\Omega}_U$ is defined (Section~\ref{sec:regular-hessian-stencils}: $\mathsf M(U)$ defined and $\alpha_c(U)$ invertible on the rows of $\Omega^{[1]}$, the local group means $\mathcal U_U(y,x)$ defined for $y\in\Omega^{(1)}$). For all $(b,b')\in\Omega^*\times\Omega^*$,
\begin{equation}\label{eq:convex-regional-entry}
    (\mathsf H^{\Omega}_U)_{bb'}= f_{bb'}(U|_{\mathcal S_{bb'}})\;,
\end{equation}
the entry function \eqref{eq:c-entry-function} evaluated on the background restricted to the entry stencil $\mathcal S_{bb'}$ of Lemma~\ref{lem:c2-entry-range}, and $f_{bb'}=0$ when $| z(b)-z(b')|_1>2$. In particular, on every branch, $(H_{\xi})_{bb'}= f_{bb'}(U_1|_{\mathcal S_{bb'}})$ for all $b,b'\in\Omega_1^*$, and for $b,b'\in I$ the stencil $\mathcal S_{bb'}$ lies inside $\Omega_1$ (within $3L< R_{\mathrm c}$ of $I$, Lemma~\ref{lem:collar-coverage}(i)).
\end{lemma}

\begin{proof}
This is the compression identity \eqref{eq:c-regional-compression} of Section~\ref{sec:regular-hessian-stencils}, whose proof is recalled: the regional form is the sum of the terms of \eqref{eq:regular-stencil-quadratic-decomposition} with index in $\Omega$ (plaquettes meeting $\Omega$, centres in $\Omega$, rows in $\Omega^{[1]}$), its entry at $(b,b')$ is the sum of the polarised entries of those of these terms whose free-input set contains $b$ and $b'$, and every term whose free-input set contains a bond $b$ of $\Omega^*$ has its index in $\Omega$: a plaquette term $\mathfrak a_p$ with $b\in\mathcal F_r$, $r\subset\partial p$, has $r= b$ or $r= b_0(c)$ with $b\in\mathcal E_c$, and in either case $p$ contains a bond with an endpoint in the block of the endpoint of $b$ that lies in $\Omega$ ($b$ has such an endpoint, being a bond of $\Omega$; the bonds of $\mathcal E_c$ have both endpoints in the two blocks of $c$, and $b_0(c)$ has one endpoint in each), and that block lies in $\Omega$, a union of blocks; a gauge term $\mathfrak g_y$ with $b\in\mathcal E_y$ has $y\in\Omega$; a row term $\mathfrak t_c$ with $b\in\mathcal E_c\cap E_*$ has an endpoint block of $c$ containing an endpoint of $b$, so $c\in\Omega^{[1]}$. Hence the set of terms contributing to the regional entry is exactly $\mathcal T_{bb'}$, and the entry is $f_{bb'}(U|_{\mathcal S_{bb'}})$, a finite formula in the background on the stencil in which the rowwise elimination \eqref{eq:regular-stencil-elimination} uses only the rows of the terms in $\mathcal T_{bb'}$, which are rows of $\Omega^{[1]}$, where $\alpha_c(U)$ is invertible: the rows of $\mathfrak g_y$ and $\mathfrak t_c$ are rows with a centre in $\Omega$ by the above, and a row $c'$ of $\mathfrak a_p$ has $b_0(c')\subset\partial p$, whose endpoints are the two sites of the segment between the centres of $c'$ next to the common face of its blocks; the other two vertices of $p$ are their translates by one unit in a direction perpendicular to $c'$, which stay in the same two blocks as $L\geqslant 3$; so all four vertices of $p$ lie in the two blocks of $c'$, one of which meets $\Omega$ since $p$ does, hence lies in $\Omega$, and $c'\in\Omega^{[1]}$. The vanishing for $| z(b)-z(b')|_1>2$ is Lemma~\ref{lem:c2-entry-range}. For the last claim, by Lemma~\ref{lem:c2-entry-range} every site of $\mathcal S_{bb'}$ is within $3L$ of $b_-\in I$, hence within $3R_{\mathrm c}$ of $I$ and in $\Omega_1^{\sim-1}$ by Lemma~\ref{lem:collar-coverage}(i).
\end{proof}

\begin{proof}[Proof of Lemma~\ref{lem:regular-bounds}]
Extend $B\in\mathfrak g^I$ by zero. By Lemma~\ref{lem:c2-regional-entry}, $\langle B,H_{\xi} B\rangle=\sum_{b,b'\in I}\langle B_b,f_{bb'}(U_1|_{\mathcal S_{bb'}})B_{b'}\rangle$ with the local entries, whose stencils lie inside $\Omega_1$; the vanishing $H_{\xi,II}(b,b')=0$ for $d(b,b')>\rho_L$ is Lemma~\ref{lem:c2-entry-range}. The quantities $m_L$, $c_L$ of \eqref{eq:regular-covariance-entry-constants} bound the row sums $\sum_{b'}$ of $\| f_{bb'}(A)\|$ and $\| Df_{bb'}(A)\|$ over $\| A\|_{\infty}\leqslant \tau_L$; since $f_{b'b}(A)= f_{bb'}(A)^*$ (the form is symmetric) and hence $Df_{b'b}(A)[X]=(Df_{bb'}(A)[X])^*$, the column sums obey the same bounds.

\emph{Partition of unity.} Let $\mathcal Z_T$ be the coarse grid and $\psi(u)=\prod_{\mu=1}^4(1-| u_{\mu}|/R_{\mathrm p})_+$ on the lifted coarse ball of radius $R_{\mathrm p}$, extended by zero and periodised (injectively, by the period condition). For $w\in\mathcal Z_T$ put
\begin{equation}\label{eq:regular-covariance-partition}
    \chi_w(b)=\frac{\psi(z(b)-w)}{S^{1/2}},\qquad S=\sum_{u}\psi(u)^2=\Bigl(\sum_{k\in\mathbb Z}(1-| k|/R_{\mathrm p})_+^2\Bigr)^4\geqslant \Bigl(\frac{R_{\mathrm p}}{4}\Bigr)^4\;,
\end{equation}
the last inequality because at least $R_{\mathrm p}$ integers $k$ satisfy $| k|\leqslant R_{\mathrm p}/2$ and each contributes at least $1/4$. Then $\sum_w\chi_w(b)^2=1$ for every $b$. A one-step difference $|\psi(u)-\psi(u+e_{\nu})|$ is at most $1/R_{\mathrm p}$, since only one factor changes and the others are at most one, and it vanishes unless $u$ or $u+e_{\nu}$ lies in the support, a set of at most $(2R_{\mathrm p}+2)^4\leqslant (5R_{\mathrm p})^4$ points. Telescoping along a coarse path of $n=| z(b)-z(b')|_1$ steps and using $(\sum_{i=1}^na_i)^2\leqslant n\sum a_i^2$,
\begin{equation}\label{eq:regular-covariance-cutoff-error}
    \sum_w(\chi_w(b)-\chi_w(b'))^2\leqslant \frac{1}{S}\,n\sum_{i=1}^n\sum_w(\text{one-step difference})^2\leqslant \frac{256}{R_{\mathrm p}^4}\,n^2\,\frac{(5R_{\mathrm p})^4}{R_{\mathrm p}^2}=\frac{160000}{R_{\mathrm p}^2}\,| z(b)-z(b')|_1^2\;.
\end{equation}
Since $\sum_w\chi_w(b)\chi_w(b')-1=-\frac12\sum_w(\chi_w(b)-\chi_w(b'))^2$, Lemma~\ref{lem:c2-entry-range} ($n\leqslant 2$ on nonzero entries) and the Schur bound $\sum_{b,b'}\| f_{bb'}\| h(B_b)h(B_{b'})\leqslant m_L\| B\|^2$, valid because row and column sums of $\| f_{bb'}\|$ are at most $m_L$ (shown below), give
\begin{equation}\label{eq:regular-covariance-localisation-error}
    \Bigl|\sum_w\langle\chi_wB,H_{\xi}\chi_wB\rangle-\langle B,H_{\xi} B\rangle\Bigr|\leqslant \frac12\cdot\frac{160000\cdot4}{R_{\mathrm p}^2}\,m_L\| B\|^2=\frac{c_{1,L}}{R_{\mathrm p}^2}\| B\|^2\leqslant \frac{\gamma_L}{4}\| B\|^2\;.
\end{equation}

\emph{Local comparison.} Fix $w$ with $\chi_wB\ne0$; since $B$ is supported on $I$, some $b\in I$ has $\chi_w(b)\neq0$, so the patch $\{ b\, | \,| z(b)-w|_{\infty}\leqslant R_{\mathrm p}\}$ meets $I$. Let $\mathfrak P_w$ be the lifted rectangle of fine sites in blocks with $| z-w|_{\infty}\leqslant R_{\mathrm p}+2$, of side at most $L(2R_{\mathrm p}+5)+1= s_{\mathfrak P}$; it lies within fine distance $s_{\mathfrak P}+1\leqslant 2R_{\mathrm c}-1$ of $I$, so Lemma~\ref{lem:c2-plaquette-deviation} applies to all its plaquettes and Lemma~\ref{lem:c2-local-gauge} gives a gauge $u_w$ on $\mathfrak P_w$ with phases $h(A_e)\leqslant \widehat\mu$. Extend $u_w$ by the identity. By Lemma~\ref{lem:c2-entry-range} the stencils of all entries with both labels in the patch lie in $\mathfrak P_w$. With $T_w= R_{u_w,*}$ the orthogonal action of Lemma~\ref{lem:hessian-gauge-locality},
\begin{equation*}
    \langle\chi_wB,H_{\xi}\chi_wB\rangle=\sum_{b,b'}\langle(T_w\chi_wB)_b,f_{bb'}(A)(T_w\chi_wB)_{b'}\rangle
    =\langle T_w\chi_wB,\mathsf H_{\id}T_w\chi_wB\rangle+\sum_{b,b'}\langle(T_w\chi_wB)_b,(f_{bb'}(A)-f_{bb'}(0))(T_w\chi_wB)_{b'}\rangle\;,
\end{equation*}
where $f_{bb'}(A)$ is the entry at the gauge-transformed background $U_1^{u_w}= e^{iA}$ on the stencil; the first equality is the entrywise gauge covariance \eqref{eq:c-entry-gauge-covariance}, $f_{bb'}(U^u)=\operatorname{Ad}_{u(b_-)}f_{bb'}(U)\operatorname{Ad}_{u(b'_-)}^{-1}$, together with $(T_w\chi_wB)_b=\operatorname{Ad}_{u_w(b_-)}\chi_w(b)B_b$ and the invariance of the inner product.
The map $A\mapsto f_{bb'}(A)$ is analytic on $\| A\|_{\infty}\leqslant \tau_L$ (Lemma~\ref{lem:hessian-gauge-locality}) and $\| A\|_{\infty}\leqslant \widehat\mu<\tau_L/2$, so $\| f_{bb'}(A)-f_{bb'}(0)\|\leqslant \sup_{t\in[0,1]}\| Df_{bb'}(tA)\|\,\widehat\mu$, and the row and column sums of these bounds are at most $c_L\widehat\mu$ by \eqref{eq:regular-covariance-entry-constants}; the Schur test gives
\begin{equation}\label{eq:regular-covariance-local-error}
    \Bigl|\langle\chi_wB,H_{\xi}\chi_wB\rangle-\langle T_w\chi_wB,\mathsf H_{\id}T_w\chi_wB\rangle\Bigr|\leqslant c_L\widehat\mu\,\|\chi_wB\|^2\;.
\end{equation}
The operator $\mathsf H_{\id}$ is the flat Hessian $\mathsf H$ of \eqref{eq:flat-covariance}: at the identity $J_{\id}=0$, every relative mean $\mathcal U_{\id}(y,x)$ is the identity, so $\mathscr G_{\id}$ is the gauge functional of \eqref{eq:action-primitive}, and $\mathsf C_{\id}$ is \eqref{eq:constraint-extension}. By \eqref{eq:covariance-uniform-gap} on the branch's torus, $\langle T_w\chi_wB,\mathsf H_{\id}T_w\chi_wB\rangle\geqslant \gamma_L\|\chi_wB\|^2$. Summing over $w$ and using $\sum_w\|\chi_wB\|^2=\| B\|^2$ and \eqref{eq:convex-background-thresholds},
\begin{equation}\label{eq:regular-covariance-lower-bound}
    \langle B,H_{\xi} B\rangle\geqslant \Bigl(\gamma_L-\frac{c_{1,L}}{R_{\mathrm p}^2}-c_L\widehat\mu\Bigr)\| B\|^2\geqslant \frac{\gamma_L}{2}\| B\|^2\;.
\end{equation}

\emph{Upper bound.} For a row $b\in I$ all entries $f_{bb'}$, $b'$ any free bond, have stencils in $\mathfrak P_{z(b)}$ (the patch of $w= z(b)$ meets $I$ at $b$; a nonzero entry has its stencil within $\ell^{\infty}$ distance $3L$ of $b_-$ by Lemma~\ref{lem:c2-entry-range}, hence in blocks $z$ with $| z-z(b)|_{\infty}\leqslant 3\leqslant R_{\mathrm p}+2$); in the gauge $u_{z(b)}$ of that patch their phases are at most $\widehat\mu\leqslant \tau_L$, and $\| f_{bb'}\|$ is unchanged by the fibre rotations of \eqref{eq:c-entry-gauge-covariance}, which are orthogonal. Hence, by Lemma~\ref{lem:c2-regional-entry} on $\Omega_1^*\times\Omega_1^*$ and \eqref{eq:regular-covariance-entry-constants}, $\sum_{b'\in\Omega_1^*}\|(H_{\xi})_{bb'}\|=\sum_{b'\in\Omega_1^*}\| f_{bb'}(U_1|_{\mathcal S_{bb'}})\|\leqslant m_L$, which is the row-sum statement of the lemma; in particular $\sum_{b'\in I}\| f_{bb'}(U_1)\|\leqslant m_L$, and by the symmetry noted at the start the column sums of $H_{\xi,II}$ obey the same bound. The Schur test gives $\| H_{\xi,II}\|\leqslant m_L$. Both bounds pass to principal restrictions: the lower bound by compression, the upper by the Schur test on fewer entries. This proves \eqref{eq:c-regular-two-sided} and the range statement.
\end{proof}

\subsection{Uniform charts and primitive bounds}
\label{app:charts}

Primitives are evaluated on the input $A= g_0\mathsf C_UB$ with the free vector $B= a+tA_{\mathrm e}$, $t\in[0,1]$ (the vacuum potentials of Section~\ref{sec:selected-integrals} have $t=0$). For a coarse row $c$ the relevant input is $A|_{\mathcal E_c}$, with the Hilbert norm on $\mathfrak g^{\mathcal E_c}$ (written $\| A|_{\mathcal E_c}\|$ or $\| A\|$). Lemma~\ref{lem:charts} asks, for every row $c$ within $2R_{\mathrm c}+2L$ of $I$ (every bond of $\mathcal E_c$ within $2R_{\mathrm c}+2L$ of $I$), for the existence and uniqueness of the small solution ($h(d_c)<\rho_{\mathrm d}=2K_*\vartheta_*^2$) of the row equation for every real phase with $\| A|_{\mathcal E_c}\|\leqslant \vartheta_*/2$, given by a function $\varphi_c$ analytic on the complex ball of radius $\vartheta_*$ and real on real phases, the positivity of the determinant, and the analyticity in $B|_{\mathcal E_c\cap E_*}$ on the row polydisc $\mathcal U_c=\{ B\, | \,| B(b)|<2+4r\ \forall b\in\mathcal E_c\cap E_*\}$ (which $g_0\mathsf C_U$ maps into that phase ball under threshold item 3) of the row solution, the phase, the two logarithms and the plaquette terms of the plaquettes with a bond in $\mathcal E_c$; for every centre $y$ within $2R_{\mathrm c}+2L$ of $I$, for the existence of the relative means and the analyticity of the gauge term on the block polydisc $\mathcal U_y=\{| B(b)|<2+4r\ \forall b\in\mathcal E_y\}$; in each case as compositions of phase functions with derivatives up to order five bounded on the closed phase ball $\| A\|\leqslant \vartheta_*/2$ by constants depending on $L$ and the group only, with the primitive sizes \eqref{eq:c2-primitive-sizes}; Lemmas~\ref{lem:c2-uniform-chart} and~\ref{lem:c2-primitive-bounds} supply this.

\begin{lemma}[uniform implicit chart]\label{lem:c2-uniform-chart}
There are constants $\varrho_*,\vartheta_*>0$, $H_*,K_*\geqslant 1$, $C_{\mathsf C}\geqslant \max(1,(4L^3(L-1))^{1/2})$, $C_{\mathsf D}$, $r_M,r'_M>0$ with $r'_M\leqslant \varrho_G$, and derivative constants $C^{\mathsf D}_k$ for $0\leqslant k\leqslant 6$ and $C^{\mathrm a}_k$ for $0\leqslant k\leqslant 5$, depending on $L$ and the group only, with $2\vartheta_*<\varrho_*$, $2\vartheta_*\leqslant \varrho_G$, $\vartheta_*\leqslant \tau_L/2$ and $4K_*H_*\vartheta_*<1$, such that for every background $U$ admitting a site gauge $u$ in which the phases $\phi$, $U^u= e^{i\phi}$, satisfy $\|\phi\|_{\infty}\leqslant \tau_L/2$ on $\mathcal E_c$, the following statements hold. The operator and derivative bounds in (i)--(iii) transfer under the orthogonal fibre gauge action. In (iv), the near-identity assertions for open paths and the mean itself are made in the chosen gauge $u$; only the relative-mean and gauge-functional conclusions are transferred to $U$. All formulas depending on $\vartheta_*$, including $\rho_{\mathrm d}$ and $C^{\mathsf D}_k$, use its final reduced value. Norms of phases on $\mathcal E_c$ are the Hilbert norms of $\mathfrak g^{\mathcal E_c}$ and of its complexification (so $\| d\|= h(d)$ for a single $d\in\mathfrak g$).
\begin{enumerate}
\item $\| h_U(c)\|\leqslant H_*$; $\|\mathsf C_UB\|_{\mathcal E_c}\leqslant C_{\mathsf C}\max_{b\in\mathcal E_c\cap E_*}| B_b|$ and $\|\mathsf C_UB\|_{\mathcal E_c}\leqslant C_{\mathsf C}\| B|_{\mathcal E_c\cap E_*}\|$ (Hilbert norms); and $\mathcal N_{U,c}$ is analytic on the open complex ball $\| A\|<2\varrho_*$ of $\mathfrak g^{\mathcal E_c}_{\mathbb C}$ with $\mathcal N_{U,c}(0)=0$, $D\mathcal N_{U,c}(0)=0$ and $\| D^2\mathcal N_{U,c}\|\leqslant K_*$ on the closed ball $\| A\|\leqslant \varrho_*$; $\mathcal N_{U,c}$ is real on real $A$.
\item On the open complex ball $\| A\|<\vartheta_*$ there is an analytic function $\varphi_c=\widetilde D_{U,c}$, real on real $A$, with $\varphi_c(0)=0$, $D\varphi_c(0)=0$, such that for every $A$ in that ball $\varphi_c(A)$ is the unique solution $d$ of the row equation \eqref{eq:native-source-local-implicit} with $\| d\|\leqslant 2K_*\vartheta_*^2=\rho_{\mathrm d}$ (in particular, for real $A$ the unique small solution in the sense of Section~\ref{sec:source-phase-primitives}); it agrees with the implicit-function germ of Lemma~\ref{lem:row-locality} (constructed there from Lemma~\ref{lem:two-block}, $\mathcal N_{U,c}(0)=0$, $D\mathcal N_{U,c}(0)=0$ and the analyticity of $\mathcal N_{U,c}$ near zero, without the present lemma), and satisfies
\begin{equation}\label{eq:convex-uniform-implicit-chart}
    \|\widetilde D_{U,c}(A)\|\leqslant 2K_*\| A\|^2,\qquad
    \| D_A\widetilde D_{U,c}(A)\|\leqslant C_{\mathsf D}\| A\|\;,\qquad C_{\mathsf D}=4K_*\;.
\end{equation}
\item On the closed real ball $\| A\|\leqslant \vartheta_*/2$ the determinant $\det_{\mathfrak g}(\id-h_U(c)\partial_{A(b_0(c))}\widetilde D_{U,c}(A))$ is positive. On a complex neighbourhood of the closed complex ball $\| A\|\leqslant \vartheta_*/2$ the functions $\widetilde D_{U,c}$, the phase $q_U(A)|_{\mathcal E_c}= A-h_U\widetilde D_U(A)$, the normalised Haar logarithm $\log\bigl(\sigma(q_U(A)(b))/\sigma(0)\bigr)$, $b\in\mathcal E_c$, and the logarithm of that determinant are analytic (the two logarithms being the real logarithms on real $A$ and their analytic continuations on the complex neighbourhood), and on the closed ball their derivatives of order $k\leqslant 5$ are bounded in operator norm by $C^{\mathrm a}_k$, uniformly in the background; for $\widetilde D_{U,c}$ this holds for $k\leqslant 6$ with $C^{\mathsf D}_k= e^kk!\,2K_*\vartheta_*^2(2/\vartheta_*)^k$. The two logarithms vanish at $A=0$, and the first derivative of the Haar logarithm vanishes there ($\sigma(0)>0$ is the value at the identity of the Haar density $\sigma$, which \eqref{eq:native-source-haar-primitive} subtracts); no statement is made about the first derivative of the determinant logarithm at $A=0$.
\item For a block centre $y$ whose block $\mathcal B_y$ lies in the gauge patch of $u$ (phases $\|\phi\|_{\infty}\leqslant \tau_L/2$ on $\mathcal E_y$), put $q^u= R_uq$, where $(R_uq)(b)=\operatorname{Ad}_{u(b_-)}q(b)$, so $\| q^u|_{\mathcal E_y}\|=\| q|_{\mathcal E_y}\|$. For every real $q$ with $\| q|_{\mathcal E_y}\|\leqslant \vartheta_*$, the path transports $(e^{iq^u}U^u)(\gamma)$, $\gamma\in\mathcal G(y,x)$, $x\in\mathcal B_y$, are within $r_G$ of $\id$ in operator norm, their logarithms $\log_G$ have $h< r_M$, and the local group mean $\mathcal U_{e^{iq^u}U^u}(y,x)$ of \eqref{eq:action-primitive} (RG1 (0.10)--(0.11)), the unique solution near the diagonal value, exists, of the form $e^{iX}$ with $h(X)< r'_M\leqslant \varrho_G$; the relative mean $\mathcal U_{e^{iq}U}(y,x)\mathcal U_U(y,x)^{-1}$ lies within $r_G$ of $\id$ in operator norm (the domain of $\log_G$), and it and the gauge term $\mathscr G_{y,U}(q)$ are restrictions to real $q$ of functions analytic on the open complex ball $\| q|_{\mathcal E_y}\|<2\vartheta_*$, with derivatives up to order five bounded by $C^{\mathrm a}_k$ on the closed complex ball of radius $\vartheta_*$, uniformly in the background, and $\mathscr G_{y,U}(0)=0$, $D\mathscr G_{y,U}(0)=0$.
\end{enumerate}
\end{lemma}

\begin{proof}
There are finitely many row types modulo $L\mathbb Z^4$. For each type the maps $U\mapsto h_U(c)$, $\mathsf C_U$ and $(U,A)\mapsto\mathcal N_{U,c}(A)$ are compositions of exponentials, relative block averages, the fibre inverse $\alpha_c(U)^{-1}$ and the chart logarithm, all analytic for phases $\|\phi\|_{\infty}\leqslant \tau_L$ and $A$ near zero (Lemma~\ref{lem:two-block} gives the singular-value bound $L^{-4}/2$ on $\alpha_c$), and real for real arguments. Under the simultaneous gauge transformation of Lemma~\ref{lem:hessian-gauge-locality}, $Q_U,h_U,\mathsf C_U$ and $\mathcal N_U$ intertwine with the orthogonal fibre action $R_{u,*}$ (proof of that lemma), which preserves the Hilbert norms and the maximum norm $\max_b| B_b|$ ($|\bigcdot |$ is $\operatorname{Ad}$-invariant); so the bounds in (i)--(iii) may be proved in the gauge $u$ and transferred to $U$. Part (iv) distinguishes open path transports in that gauge from the relative means and the gauge functional. The set of phase configurations with $\|\phi\|_{\infty}\leqslant \tau_L/2$ on $\mathcal E_c$ is compact, so a common complex radius $2\varrho_*$ of analyticity and bounds $H_*,K_*,C_{\mathsf C}$ exist (the two bounds on $\mathsf C_U$ are two operator norms of the same linear map, each bounded on the compact set; $C_{\mathsf C}$ is the larger, increased if necessary to $\max(1,(4L^3(L-1))^{1/2})$, which only weakens the inequalities); $\mathcal N_{U,c}(0)=0$ because the zero fluctuation reproduces the background averages, and $D\mathcal N_{U,c}(0)=0$ as in the proof of Lemma~\ref{lem:row-locality}.

(ii) Fix $A$ in the complex ball with $\| A\|\leqslant \vartheta\leqslant \vartheta_*$ and consider $\Phi(d)=\mathcal N_{U,c}(A-h_Ud)$ on the closed complex ball $\| d\|\leqslant 2K_*\vartheta^2$. Then $\| A-h_Ud\|\leqslant \vartheta+2H_*K_*\vartheta^2\leqslant 2\vartheta\leqslant \varrho_*$; Taylor's formula with $\mathcal N(0)=0$, $D\mathcal N(0)=0$ gives $\|\Phi(d)\|\leqslant \frac12K_*(2\vartheta)^2=2K_*\vartheta^2$, and $\| D\Phi(d)\|\leqslant \| D\mathcal N(A-h_Ud)\| H_*\leqslant K_*\cdot2\vartheta H_*<\frac12$. So $\Phi$ is a contraction of the closed ball into itself, with a unique fixed point; taking $\vartheta=\vartheta_*$ gives the uniqueness in the ball $\| d\|\leqslant \rho_{\mathrm d}$ for every $A$ with $\| A\|<\vartheta_*$, and taking $\vartheta=\| A\|$ gives the first bound of \eqref{eq:convex-uniform-implicit-chart}; the fixed point $\varphi_c(A)$ is analytic in $A$ on the open ball as the uniform limit of the analytic iterates $\Phi^n(0)$ (each iterate is analytic in $A$, and the convergence is geometric, uniformly in $A$), and it is real for real $A$ because the iterates are. Differentiating the fixed-point equation, $D_A\widetilde D=(\id+D\mathcal N(x)h_U)^{-1}D\mathcal N(x)$ with $x= A-h_U\widetilde D$, $\| D\mathcal N(x)\|\leqslant K_*\| x\|\leqslant 2K_*\| A\|$ and $\|(\id+D\mathcal N(x)h_U)^{-1}\|\leqslant 2$; this is the second bound, and $D\varphi_c(0)=0$. The germ of Lemma~\ref{lem:row-locality} is the solution near zero, and uniqueness in the ball identifies the two.

(iii) Cauchy's estimate on the ball of radius $\vartheta_*$, where $\|\widetilde D_{U,c}\|\leqslant 2K_*\vartheta_*^2$, bounds the $k$-th derivative at a point of the closed ball of radius $\vartheta_*/2$ in a fixed direction $v$, $\| v\|\leqslant 1$, by $k!\,2K_*\vartheta_*^2(\vartheta_*/2)^{-k}$ for every $k$, and the polarisation identity for symmetric $k$-linear maps bounds the operator norm of $D^k\widetilde D_{U,c}$ by $e^k$ times the diagonal bound; this gives $C^{\mathsf D}_k$ for $0\leqslant k\leqslant 6$. The constants $C^{\mathrm a}_k$ bounding the other functions are chosen separately after (iv). The phase $q_U(A)|_{\mathcal E_c}= A-h_U\widetilde D_U(A)$ is a composition of $\widetilde D_{U,c}$ with the fixed linear map $h_U(c)$ of norm at most $H_*$, so it is analytic on the ball of radius $\vartheta_*$ and its derivatives are bounded by the chain rule. The Haar density: $\sigma(q)\,\dd q$ is the pullback of Haar measure under $q\mapsto e^{iq}$ on the chart ball $h(q)<\varrho_G$, so $\sigma(q)=\sigma(0)|\det D(q\mapsto e^{iq})(q)|$, the determinant of the differential taken with respect to left-invariant frames; the determinant is analytic in $q$ and nonvanishing on the chart ball (the chart has invertible differential, conventions of Section~\ref{sec:sourced-effective-interaction}), so it has constant sign on the connected ball, positive since it equals one at $q=0$; thus $\sigma$ is analytic and positive there, and even, $\sigma(-q)=\sigma(q)$, because Haar measure is invariant under $g\mapsto g^{-1}$ and $e^{-iq}=(e^{iq})^{-1}$. Hence $\log(\sigma/\sigma(0))$, the real logarithm, extends analytically to a complex neighbourhood of the real chart ball, vanishes at $0$ with vanishing first derivative (evenness), and its composition with $q_U(A)(b)$ (a phase with $h(q_U(A)(b))\leqslant \| A\|+H_*2K_*\| A\|^2\leqslant \vartheta_*$ on the ball) is analytic on a complex neighbourhood of the closed ball $\| A\|\leqslant \vartheta_*/2$ after reducing $\vartheta_*$ so that the compositions stay in the domain of analyticity, with derivatives bounded on the closed ball by the chain rule, and its first derivative at $A=0$ is $D\log\sigma(0)\circ Dq_U(0)=0$. The determinant argument $\id-h_U(c)\partial_{A(b_0(c))}\widetilde D_{U,c}(A)$ equals $\id$ at $A=0$ since $D\widetilde D(0)=0$, and is within $H_*C_{\mathsf D}\vartheta_*/2<1/2$ of $\id$ in operator norm on the ball; for real $A$ it is a real endomorphism of $\mathfrak g$ within $1/2$ of $\id$, whose eigenvalues have positive real part, so its determinant is positive, and on a complex neighbourhood of the closed ball $\| A\|\leqslant \vartheta_*/2$ the determinant stays in the disc of radius $1/2$ about $1$ (after reducing $\vartheta_*$ once more if needed, by continuity and compactness), where $\log$ is analytic; so the real logarithm of the determinant extends analytically to that neighbourhood and vanishes at $A=0$, where the determinant is one. Its derivatives of order $k\leqslant 5$ are bounded on the closed ball by the chain rule, since the argument contains $\partial_{A(b_0(c))}\widetilde D_{U,c}$ and so involves the derivatives of $\widetilde D_{U,c}$ up to order six, bounded above. Its first derivative at $A=0$, $-\operatorname{tr}_{\mathfrak g}\bigl(h_U(c)\partial_{A(b_0(c))}D\widetilde D_{U,c}(0)[\bigcdot ]\bigr)$, involves the second derivative $D^2\widetilde D_{U,c}(0)= D^2\mathcal N_{U,c}(0)$, which does not vanish in general (RG1 p.267: the expansion of $\widetilde D$ begins with quadratic terms), and nothing is claimed about it. The bounds are uniform in the background by the compactness of the set of phase configurations. All reductions of $\vartheta_*$ are by constants depending on $L$ and the group only, over the finitely many row types.

(iv) Set $q^u= R_uq$; the fibre action is orthogonal on real phases and unitary on their complexifications, so it preserves the phase balls and operator norms of multilinear derivatives. In the gauge $u$ the links of $\mathcal B_y$ are $e^{i\phi}$ with $\|\phi\|_{\infty}\leqslant \tau_L/2$, and $h(q_b)\leqslant \| q|_{\mathcal E_y}\|\leqslant \vartheta_*\leqslant \tau_L/2$ for every $b\in\mathcal E_y$. A path of $\mathcal G(y,x)$ has at most $2(L-1)$ bonds (four coordinate runs of length at most $(L-1)/2$), and for unitaries $\| U_1\cdots U_n-\id\|_{\mathrm{op}}\leqslant \sum_j\| U_j-\id\|_{\mathrm{op}}$, so every path transport $g_{\gamma}=(e^{iq^u}U^u)(\gamma)$ lies within $2(L-1)c(\tau_L/2+\vartheta_*)\leqslant 2(L-1)c\,\tau_L\leqslant r_G$ of $\id$ in operator norm, by \eqref{eq:c2-tau-final}; hence $A_{\gamma}=\log_Gg_{\gamma}$ is defined with $h(A_{\gamma})\leqslant c_G\cdot2(L-1)c\,\tau_L\leqslant r_M/2$, again by \eqref{eq:c2-tau-final}. The local group mean of a family $(g_{\gamma})_{\gamma\in\mathcal G(y,x)}$ is the solution $\mathcal U$ near the diagonal value of $\sum_{\gamma}\log_G(g_{\gamma}\mathcal U^{-1})=0$; write $g_{\gamma}= e^{iA_{\gamma}}$, $\mathcal U= e^{iX}$ and $F(X,(A_{\gamma}))=\sum_{\gamma}\log_G(e^{iA_{\gamma}}e^{-iX})$, an analytic map on a neighbourhood of $(0,0)$ in $\mathfrak g\times\mathfrak g^{\mathcal G(y,x)}$ with $F(0,0)=0$ and $D_XF(0,0)=-n\,\id$, $n=|\mathcal G(y,x)|$ (as $D\log_G(\id)=\id$). Choose the product neighbourhood in the complexified Lie-algebra variables so that every $e^{iA_{\gamma}}e^{-iX}$ lies in the analytic logarithm domain there. The analytic implicit function theorem on that neighbourhood gives radii $r_M,r'_M>0$, with $r'_M\leqslant \varrho_G$ (any smaller solution radius is admissible together with a correspondingly smaller $r_M$), and an analytic map $X= X((A_{\gamma}))$ on $\max_{\gamma} h(A_{\gamma})< r_M$, with $h(X)< r'_M$, such that $X((A_{\gamma}))$ is the unique solution of $F=0$ with $h(X)< r'_M$; there are finitely many values of $n$ (one per $x\in\mathcal B_y$ modulo translations), so $r_M$, $r'_M$ and the derivative bounds of $X$ are constants of the group and $L$. This solution is identified with the fixed mean used in Stage 1 as follows. RG1 (0.8) places that specified analytic mean near the identity for input families near the identity, and (0.10) says it solves this equation. Reduce $r_M$, before the final choice of $\tau_L$, so that all families $\{ e^{iA_{\gamma}}\}$ with $\max_{\gamma} h(A_{\gamma})< r_M$ lie in the small-diameter domain of the specified mean, that this mean has logarithm of norm below $r'_M$, and that all relative logarithms use the common chart. Use the operator-norm diameter $\max_{\gamma,\gamma'}\| g_{\gamma}-g_{\gamma'}\|_{\mathrm{op}}$; it is unchanged by the two-sided multiplication $g_{\gamma}\mapsto u(y)g_{\gamma} u(x)^{-1}$. Thus the original-gauge families lie in the same small-diameter domain. Their relative products with the mean transform by conjugation at $y$, which preserves the $\operatorname{Ad}$-invariant logarithm chart, so the original-gauge relative logarithms are defined as well. Local uniqueness from the implicit function theorem then identifies it with $e^{iX}$. Its left-right covariance is the stated property RG1 (0.6), not a consequence of uniqueness on an unspecified chart domain. Thus for any admissible gauge $u$ its value in the original gauge is $u(y)^{-1}\mathcal U_{U^u}(y,x)u(x)$; the value is independent of the choice of admissible gauge because both expressions are values of the same specified mean. Since $h(A_{\gamma})\leqslant r_M/2< r_M$ for the transport families above, the mean $\mathcal U= e^{iX((A_{\gamma}))}$ exists, lies in the chart domain, and is the unique solution near the diagonal value in the sense of \eqref{eq:action-primitive}. Write $A_{\gamma}(q^u)=\log_G((e^{iq^u}U^u)(\gamma))$ and $X(q^u)= X((A_{\gamma}(q^u))_{\gamma})$. Define $M_x(q)=\mathcal U_{e^{iq}U}(y,x)\mathcal U_U(y,x)^{-1}$ and $M_x^u(q^u)=\mathcal U_{e^{iq^u}U^u}(y,x)\mathcal U_{U^u}(y,x)^{-1}$. Gauge covariance gives $M_x^u(q^u)= u(y)M_x(q)u(y)^{-1}= e^{iX(q^u)}e^{-iX(0)}$, with inverse $M_x^u(q^u)^{-1}= e^{iX(0)}e^{-iX(q^u)}$; both are compositions of the analytic map $X$ with $\log_G$, the exponentials of the phases and the fixed links, so they extend analytically to a complex ball $\| q|_{\mathcal E_y}\|<2\vartheta_*$ (after reducing $\vartheta_*$ by a constant of the group and $L$). For real $q$, $M_x(q)$ is unitary and $\| M_x(q)-\id\|_{\mathrm{op}}=\| e^{iX(q^u)}-e^{iX(0)}\|_{\mathrm{op}}\leqslant C\| q|_{\mathcal E_y}\|$ with a constant $C$ of $L$ and the group (the derivatives of $X$ on the ball $\max_{\gamma} h(A_{\gamma})\leqslant r_M/2$, of $\log_G$ on the $r_G$-ball and of the products of exponentials are bounded), so after reducing $\vartheta_*$ once more $M_x(q)$ lies within $r_G$ of $\id$ for $\| q|_{\mathcal E_y}\|\leqslant \vartheta_*$. The gauge term is $\mathscr G_{y,U}(q)=\sum_x(1-\operatorname{Re}\tr_NM_x(q))$; for real $q$, $\operatorname{Re}\tr_NM_x(q)=\frac12(\tr_NM_x(q)+\tr_NM_x(q)^{-1})$ because $M_x(q)$ is unitary, and the right side is analytic in $q$ on the complex ball: this is the analytic continuation meant. The derivatives up to order five of $M_x$ and $\mathscr G_{y,U}$ are bounded on the closed ball of radius $\vartheta_*$ by constants $C^{\mathrm a}_k$ uniform in the background, because they are analytic jointly in $q$ and in the phases $\phi$ of the links of $\mathcal B_y$ in the gauge $u$, and the set $\|\phi\|_{\infty}\leqslant \tau_L/2$ is compact. $\mathscr G_{y,U}(0)=0$ because $M_x(0)=\id$, and $D\mathscr G_{y,U}(0)=0$ because each summand is nonnegative on real $q$ and vanishes at $q=0$. The gauge $u$ multiplies every transport from $y$ to $x$ on the left by $u(y)$ and on the right by $u(x)^{-1}$, and the mean equation is covariant under this ($\mathcal U\mapsto u(y)\mathcal Uu(x)^{-1}$, so $M_x\mapsto u(y)M_xu(y)^{-1}$, which preserves $\| M_x-\id\|_{\mathrm{op}}$ and $\tr_NM_x$), so the relative-mean and gauge-functional bounds hold for $U$ itself. The near-identity path and unnormalised-mean assertions remain in the gauge $u$. For each $0\leqslant k\leqslant 5$, choose $C^{\mathrm a}_k$ to be at least one and the maximum of $C^{\mathsf D}_k$ and the uniform derivative bounds for the phase, the two logarithms, the relative means and the gauge functional obtained in (iii)--(iv). Their finiteness and uniformity follow from the compact phase sets, the finitely many templates and the stated analytic neighbourhoods. No explicit value of $C^{\mathsf D}_k$ is asserted for these other functions.
\end{proof}

Let $c$ be a row within $2R_{\mathrm c}+2L$ of $I$ (a near row); by Section~\ref{sec:source-phase-primitives} every row of a near primitive (a primitive whose free-input set meets $I$) is within $2R_{\mathrm c}$ of $I$, since its bonds are within $4L+3\leqslant 2R_{\mathrm c}$ of a free-input bond in $I$, and every row of a plaquette with a bond in $\mathcal E_c$ for such a row $c$ is within $2R_{\mathrm c}+2L$ of $I$ (its bonds are within $2L-1+1$ of a bond of $\mathcal E_c$). The two-block set $\mathcal E_c$ lies inside a lifted rectangle of side $2L+1\leqslant s_{\mathfrak P}$ all of whose bonds are within $2R_{\mathrm c}+2L+2L+1\leqslant 3R_{\mathrm c}$ of $I$ (as $R_{\mathrm c}\geqslant 2R_{\mathrm{loc}}+1=4L+5$); Lemmas~\ref{lem:c2-plaquette-deviation} and~\ref{lem:c2-local-gauge} give a gauge on it with phases at most $\widehat\mu<\tau_L/2$ (the bound of Lemma~\ref{lem:c2-local-gauge} is monotone in the side), and Lemma~\ref{lem:collar-coverage}(iii) (range $3R_{\mathrm c}$) bounds the retained inputs $A_{\mathrm e}(b)$, $b\in\mathcal E_c\cap J$, by $r_{\mathrm e,\xi}$. Likewise a block $\mathcal B_y$ with centre $y$ within $2R_{\mathrm c}+2L$ of $I$ is a lifted rectangle of side $L\leqslant s_{\mathfrak P}$ whose bonds are within $3R_{\mathrm c}$ of $I$, so Lemmas~\ref{lem:c2-plaquette-deviation} and~\ref{lem:c2-local-gauge} apply to it directly. For real $B\in\mathcal U_y$, the block input satisfies $\| g_0B|_{\mathcal E_y}\|\leqslant 2C_{\mathsf C}g_0(1+2r)\leqslant \vartheta_*/4$ by threshold item 3 and $C_{\mathsf C}\geqslant |\mathcal E_y|^{1/2}$. Hence the real relative-mean conclusions of (iv) apply on the caller's entire block polydisc; its analytic continuation is supplied by the same lemma. So Lemma~\ref{lem:c2-uniform-chart} applies to every row and every centre within $2R_{\mathrm c}+2L$ of $I$, in particular to every row of every near primitive and of every plaquette with a bond in the two-block set of such a row. Put $s_{\xi}=1+r+r_{\mathrm e,\xi}\leqslant 1+2r$. The next lemma shows that Lemma~\ref{lem:c2-uniform-chart}(i)--(iii) apply at every row of $\Omega_1^{[1]}$, near or far, and that the arguments of the whole fibre $D^{(11)}(V,V_0)$ lie in its balls.

\begin{lemma}[fibre smallness]\label{lem:c2-fibre-smallness}
Let $n_{\mathcal E}$ be the largest number of bonds of a two-block set $\mathcal E_c$ of \eqref{eq:regular-stencil-bonds} (a constant of $L$), let $O(L^2)$ be CRE's constant of p.247 in (B1)(ii), and let $g_0$ satisfy threshold items 5 and 6 of Section~\ref{app:thresholds}; item 6 is
\begin{equation}\label{eq:c2-fibre-thresholds}
    3c_Gc_{\mathrm h}s_{\mathfrak P}\varepsilon_0\leqslant \tfrac12\tau_L,\qquad
    n_{\mathcal E}^{1/2}c_Gc_{\mathrm h}\,2\,O(L^2)\,\varepsilon_0\leqslant \tfrac14\vartheta_*\;.
\end{equation}
The first condition gives $3c_{\mathrm h}s_{\mathfrak P}\varepsilon_0\leqslant \tau_L/(2c_G)\leqslant r_G$, since $\tau_L\leqslant 2c_Gr_G$ by \eqref{eq:c2-tau-final}. Let $(\Omega_1,V,V_0)$ be an admissible pair, with the background $U_1= U_1(V,V_0)$ of (B1) on the bonds of $\Omega_1$ and $V_0$ on the other bonds. Phases $A\in\mathfrak g^{\Omega_1}$ are extended by zero, and $\varphi(A)\coloneqq (\varphi_c(A|_{\mathcal E_c}))_{c\in\Omega_1^{[1]}}$ wherever every component is defined.
\begin{enumerate}
\item For every row $c\in\Omega_1^{[1]}$, $U_1$ has a site gauge on $\mathfrak P_c=\mathcal B_{c_-}\cup\mathcal B_{c_+}$, the union of the two blocks of $c$, with phases $\|\phi\|_{\infty}\leqslant \tau_L/2$ on $\mathcal E_c$. Hence Lemma~\ref{lem:c2-uniform-chart}(i)--(iii) hold at $c$ for $U= U_1$: the function $\varphi_c=\widetilde D_{U_1,c}$ of Lemma~\ref{lem:c2-uniform-chart}(ii) is analytic on the ball $\| A|_{\mathcal E_c}\|<\vartheta_*$, real on real phases, with $\varphi_c(0)=0$, $D\varphi_c(0)=0$, and for every $A$ in that ball $\varphi_c(A|_{\mathcal E_c})$ is the unique solution $d_c$ of the row equation of $c$ with $h(d_c)\leqslant \rho_{\mathrm d}$.
\item Let $U\in D^{(11)}(V,V_0)$, $A'=\frac1i\log(UU_1^{-1})$ on the bonds of $\Omega_1$ (defined by (B1)), $d=\mathcal N_{U_1}(A')$ and $A= A'+h_{U_1}d$, so that $A= g_0\mathsf C_{U_1}B$ with $B= g_0^{-1}A'|_{\Omega_1^*}$ (the decomposition step of the converse argument of Section~\ref{sec:branch-data} after \eqref{eq:c-free-variable-domain}, up to the definition of $B$, which uses (B1), the constraint, $LQ_{U_1}h_{U_1}=\id$ and the invertibility of $\alpha_c(U_1)$ in the elimination maps \eqref{eq:regular-stencil-elimination}, given by (i) and Lemma~\ref{lem:two-block} as in the Remark under (B2), and neither the cutoff (1.25) nor (iii)). Then for every row $c\in\Omega_1^{[1]}$
\begin{equation*}
    \| A'|_{\mathcal E_c}\|\leqslant \tfrac14\vartheta_*,\qquad
    h(d_c)\leqslant \tfrac1{32}K_*\vartheta_*^2<\rho_{\mathrm d},\qquad
    \| A|_{\mathcal E_c}\|<\tfrac12\vartheta_*,\qquad
    d_c=\varphi_c(A|_{\mathcal E_c})\;.
\end{equation*}
\item For every $B\in\mathfrak g^{\Omega_1^*}$, every family $d$ of chart class at $A= g_0\mathsf C_{U_1}B$ (Section~\ref{sec:branch-data}) with $U_d(B)\in D^{(11)}(V,V_0)$ is $d=\varphi(A)$.
\item Let $\mathcal N_{\Theta}=\{ A\in\mathfrak g^{\Omega_1}:\| A|_{\mathcal E_c}\|<\vartheta_*/2\text{ for every }c\in\Omega_1^{[1]}\}$; it is open, and by (ii) it contains the argument $A$ of every $U\in D^{(11)}(V,V_0)$. On $\mathcal N_{\Theta}$ the family $\varphi$ is analytic, real on real phases, and solves the row equations \eqref{eq:native-source-local-implicit}; the phase $\Theta(A)= A-h_{U_1}\varphi(A)$ has $\|\Theta(A)|_{\mathcal E_c}\|<\vartheta_*$ at every row and is chart-valued, $h(\Theta(A)(b))<\vartheta_*<\varrho_G$ on every bond $b$ of $\Omega_1$; for every row, $\| h_{U_1}(c)\partial_{A(b_0(c))}\varphi_c(A|_{\mathcal E_c})\|\leqslant 4H_*K_*\| A|_{\mathcal E_c}\|<1$; and
\begin{equation}\label{eq:c2-fibre-jacobian}
    \det\Bigl(\id-h_{U_1}\frac{\delta}{\delta A}\varphi(A)\Bigr)=\prod_{c\in\Omega_1^{[1]}}\det\nolimits_{\mathfrak g}\bigl(\id-h_{U_1}(c)\partial_{A(b_0(c))}\varphi_c(A|_{\mathcal E_c})\bigr)>0\;.
\end{equation}
\end{enumerate}
\end{lemma}

\begin{proof}
(i) The set $\mathfrak P_c$ is a lifted rectangle of side at most $2L+1\leqslant s_{\mathfrak P}$, injective modulo the periods (every fine period is at least $8M>2L+1$), and it contains every bond of $\mathcal E_c$, both endpoints of such a bond lying in $\mathcal B_{c_-}\cup\mathcal B_{c_+}$ by \eqref{eq:regular-stencil-bonds}. The two blocks of $c$ lie in $\Omega_1^{\sim}$ ((B1)(iv)) and $U_1\in D^{(11)}(V,V_0)$ ((B1)(iii)), so (B1)(i) gives $| U_1(\partial p)-1|<\varepsilon_0$, hence $\| U_1(\partial p)-\id\|_{\mathrm{op}}\leqslant c_{\mathrm h}\varepsilon_0=2c_{\mathrm h}(\varepsilon_0/2)$ by the norm comparison of the preamble, for every plaquette $p\subset\mathfrak P_c$. By the consequence $3c_{\mathrm h}s_{\mathfrak P}\varepsilon_0\leqslant r_G$ of \eqref{eq:c2-fibre-thresholds} noted after it, $6c_{\mathrm h}s_{\mathfrak P}(\varepsilon_0/2)\leqslant r_G$, so Lemma~\ref{lem:c2-local-gauge} with $\varepsilon_0/2$ in place of $\varepsilon_1$ (the remark after its proof) gives a site gauge on $\mathfrak P_c$ with phases $h(A_b)\leqslant 3c_Gc_{\mathrm h}s_{\mathfrak P}\varepsilon_0\leqslant \tau_L/2$, by the first condition. This is the hypothesis of Lemma~\ref{lem:c2-uniform-chart} at $c$. The phases of $\mathfrak g^{\Omega_1}$, extended by zero, form the subspace of $\mathfrak g^{\mathcal E_c}$ of the phases vanishing on the bonds of $\mathcal E_c$ off $\Omega_1$; for $A$ in it and $d\in\mathfrak g$ the phase $A-h_{U_1}(c)d$ lies in it, since $b_0(c)$ is a bond of $\Omega_1$, and Lemma~\ref{lem:c2-uniform-chart}(i)--(iii) hold on it by restriction.

(ii) By (B1)(ii)--(iii), $| U'(b)-1|<2\,O(L^2)\varepsilon_0$ for $U'= UU_1^{-1}$ on every bond $b$ of $\Omega_1$: every such bond, the designated bonds included, is a bond of $\square'^{\sim2}$ for the $L^2M_2R_1$-cube $\square'\subset\Omega_1$ that contains an endpoint of $b$ in $\Omega_1$ ($\Omega_1$ is a union of $LMR_1$-cubes, p.248, each a union of $L^2M_2R_1$-cubes: CRE's partitions are compatible, p.245, read in (H4) as nested, and all sides are powers of $L$, the paragraph following \eqref{eq:c-standing-conditions}). By item 5, $2\,O(L^2)\varepsilon_0< r_G/\sqrt N$, so $\| U'(b)-\id\|_{\mathrm{op}}\leqslant c_{\mathrm h}| U'(b)-1|< r_G$ and $h(A'(b))\leqslant c_G\| U'(b)-\id\|_{\mathrm{op}}\leqslant c_Gc_{\mathrm h}\,2\,O(L^2)\varepsilon_0$ (preamble); $\mathcal E_c$ has at most $n_{\mathcal E}$ bonds, and $A'$ vanishes on those off $\Omega_1$, so the second condition of \eqref{eq:c2-fibre-thresholds} gives $\| A'|_{\mathcal E_c}\|\leqslant \vartheta_*/4<\varrho_*$. The row component $d_c=\mathcal N_{U_1,c}(A')$ depends on $A'|_{\mathcal E_c}$ only (Lemma~\ref{lem:two-block}), and Taylor's formula with Lemma~\ref{lem:c2-uniform-chart}(i), which applies by (i), gives $h(d_c)\leqslant \frac12K_*(\vartheta_*/4)^2=\frac1{32}K_*\vartheta_*^2<2K_*\vartheta_*^2=\rho_{\mathrm d}$. The only designated bond in $\mathcal E_c$ is $b_0(c)$ (Lemma~\ref{lem:two-block}) and $h_{U_1}$ has its outputs at the designated bonds, so $A|_{\mathcal E_c}= A'|_{\mathcal E_c}+h_{U_1}(c)d_c$, and $\| A|_{\mathcal E_c}\|\leqslant \vartheta_*/4+H_*K_*\vartheta_*^2/32<\vartheta_*/2$ by $4K_*H_*\vartheta_*<1$. Finally $A'|_{\mathcal E_c}= A|_{\mathcal E_c}-h_{U_1}(c)d_c$, so $d_c=\mathcal N_{U_1,c}(A|_{\mathcal E_c}-h_{U_1}(c)d_c)$ solves the row equation of $c$ at $A$, with $h(d_c)\leqslant \rho_{\mathrm d}$ and $\| A|_{\mathcal E_c}\|<\vartheta_*$; by the uniqueness in (i), $d_c=\varphi_c(A|_{\mathcal E_c})$.

(iii) Since $d$ is of chart class, $q_d(B)= A-h_{U_1}d$ is the chart logarithm of $U_d(B)U_1^{-1}$ on the bonds of $\Omega_1$, so it is the phase $A'$ of (ii) for $U= U_d(B)$. The row equations give $d_c=\mathcal N_{U_1,c}(A-h_{U_1}d)=\mathcal N_{U_1,c}(A')$, that is $d=\mathcal N_{U_1}(A')$, and $A= A'+h_{U_1}d$. So $A'$, $d$ and $A$ are the quantities of (ii) for $U= U_d(B)$, and (ii) gives $d_c=\varphi_c(A|_{\mathcal E_c})$ for every row.

(iv) $\mathcal N_{\Theta}$ is a finite intersection of preimages of open balls under restriction maps, hence open, and it contains the arguments by (ii). By (i) each component $\varphi_c(A|_{\mathcal E_c})$ is analytic and real on $\mathcal N_{\Theta}$ and solves the row equation of $c$ at $A$; the row equations decouple ($(A-h_{U_1}\varphi(A))|_{\mathcal E_c}= A|_{\mathcal E_c}-h_{U_1}(c)\varphi_c(A|_{\mathcal E_c})$, Lemma~\ref{lem:two-block}), so $\varphi$ solves the system \eqref{eq:native-source-local-implicit}. By \eqref{eq:convex-uniform-implicit-chart} and Lemma~\ref{lem:c2-uniform-chart}(i), $\|\Theta(A)|_{\mathcal E_c}\|\leqslant \| A|_{\mathcal E_c}\|+2H_*K_*\| A|_{\mathcal E_c}\|^2<\vartheta_*/2+H_*K_*\vartheta_*^2/2<\vartheta_*$. Every bond $b$ of $\Omega_1$ lies in $\mathcal E_c$ for a row $c\in\Omega_1^{[1]}$: $b$ has an endpoint in $\Omega_1$, whose block $\mathcal B_y$ lies in $\Omega_1$ ($\Omega_1$ is a union of blocks); if both endpoints of $b$ lie in $\mathcal B_y$, any coarse bond $c$ with $c_-= y$ serves, and otherwise the coarse bond joining $y$ and the centre of the block of the other endpoint does; in both cases a block centre of $c$ lies in $\Omega_1$, so $c\in\Omega_1^{[1]}$, and $b\in\mathcal E_c$ by \eqref{eq:regular-stencil-bonds}. Hence $h(\Theta(A)(b))\leqslant \|\Theta(A)|_{\mathcal E_c}\|<\vartheta_*\leqslant \varrho_G/2$. By \eqref{eq:convex-uniform-implicit-chart} and Lemma~\ref{lem:c2-uniform-chart}(i), $\| h_{U_1}(c)\partial_{A(b_0(c))}\varphi_c\|\leqslant H_*C_{\mathsf D}\| A|_{\mathcal E_c}\|=4H_*K_*\| A|_{\mathcal E_c}\|<2H_*K_*\vartheta_*<1$. The matrix $\frac{\delta}{\delta A}\varphi$ has rows indexed by $\Omega_1^{[1]}$ and columns by the bonds of $\Omega_1$, $h_{U_1}$ maps row $c$ to the bond $b_0(c)$, and $\partial_{A(b_0(c'))}\varphi_c=0$ for $c'\neq c$, since $\varphi_c$ depends on $A|_{\mathcal E_c}$, which contains no designated bond other than $b_0(c)$ (Lemma~\ref{lem:two-block}). So, in the ordering (free bonds, designated bonds), $\id-h_{U_1}\frac{\delta}{\delta A}\varphi$ is block triangular, with identity block on the free bonds and designated block $\operatorname{diag}_c(\id-h_{U_1}(c)\partial_{A(b_0(c))}\varphi_c)$; this gives the product in \eqref{eq:c2-fibre-jacobian}. Each factor is positive by Lemma~\ref{lem:c2-uniform-chart}(iii), which applies by (i), since $\| A|_{\mathcal E_c}\|<\vartheta_*/2$.
\end{proof}

\begin{lemma}[primitive bounds]\label{lem:c2-primitive-bounds}
Let $g_0$ satisfy threshold item 3 of Section~\ref{app:thresholds}, $4C_{\mathsf C}g_0(1+2r)\leqslant \vartheta_*/2$ (hence $4C_{\mathsf C}g_0s_{\xi}\leqslant \vartheta_*/2$, and $g_0s_{\xi}\leqslant \vartheta_*/(8C_{\mathsf C})\leqslant 1$, since $C_{\mathsf C}\geqslant 1$ and $\vartheta_*\leqslant \tau_L/2\leqslant 1/2$ by \eqref{eq:c2-tau-final}). Then on the enlarged box $\{| a_b|\leqslant 2r\ \forall b\in I\}$ and for $t\in[0,1]$, every primitive $p_{\alpha}$ with free-input set $S_{\alpha}$ meeting $I$ is, at the input $a+tA_{\mathrm e}$, the restriction to real arguments of a function analytic in $(a,t)$ on a complex neighbourhood of the enlarged box times $[0,1]$, and with a constant $C_{\mathrm p}$ depending on $L$ and the group only,
\begin{equation}\label{eq:c2-primitive-sizes}
    | p_{\xi,\alpha}(a+tA_{\mathrm e})|\leqslant C_{\mathrm p}g_0s_{\xi}^3,\qquad
    | p^{\mathrm H}_{\xi,b}(a+tA_{\mathrm e})|\leqslant C_{\mathrm p}g_0^2s_{\xi}^2,\qquad
    | p^{\mathrm J}_{\xi,c}(a+tA_{\mathrm e})|\leqslant C_{\mathrm p}g_0s_{\xi}\;,
\end{equation}
and for each of the three kinds, in the Hilbert norms on $\mathfrak g^{S_{\alpha}}$,
\begin{equation}\label{eq:convex-primitive-derivatives}
    \| D_a^2p_{\alpha}\|\leqslant C_{\mathrm p}g_0s_{\xi},\qquad
    \| D_ap_{\alpha}\|\leqslant C_{\mathrm p}g_0s_{\xi}^2,\qquad
    | p_{\alpha}(a+tA_{\mathrm e})-p_{\alpha}(tA_{\mathrm e})|\leqslant C_{\mathrm p}g_0s_{\xi}^3\;,
\end{equation}
and $\| D_a^3p_{\alpha}\|\leqslant C_{\mathrm p}g_0$, $\| D_a^4p_{\alpha}\|\leqslant C_{\mathrm p}g_0^2$; the same four derivative bounds hold for the derivatives $D_B^kp_{\alpha}$, $k=1,\dots,4$, in the full input $B= a+tA_{\mathrm e}$ on $\mathfrak g^{S_{\alpha}}$, retained inputs included (the input of the primitive is linear in $B$ with the same norm bound, see the proof). Moreover, for every row $c$ within $2R_{\mathrm c}+2L$ of $I$ and every real phase $A$ with $\| A|_{\mathcal E_c}\|\leqslant \vartheta_*/2$, the row equation at $A$ has exactly one small solution $\widetilde D_{U,c}(A)=\varphi_c(A|_{\mathcal E_c})$ ($h(d_c)<\rho_{\mathrm d}=2K_*\vartheta_*^2$), with $\varphi_c$ analytic on the complex ball of radius $\vartheta_*$, real on real phases, $\varphi_c(0)=0$, $D\varphi_c(0)=0$; $g_0\mathsf C_U\mathcal U_c$ lies in that phase ball; the determinant $\det_{\mathfrak g}(\id-h_U(c)\partial_{A(b_0(c))}\widetilde D_{U,c}(A))$ is positive there; and the row solution, the phase $q_U|_{\mathcal E_c}$, the normalised Haar logarithm and the real determinant logarithm are analytic in $B|_{\mathcal E_c\cap E_*}$ on $\mathcal U_c$, as are, when $c$ is within $2R_{\mathrm c}$ of $I$, the plaquette terms $s_{p,U}(q_U(\bigcdot ))$ of the plaquettes $p$ with a bond in $\mathcal E_c$ (on the intersection of the polydiscs of the rows of $p$, all of which are rows within $2R_{\mathrm c}+2L$ of $I$); for every centre $y$ within $2R_{\mathrm c}+2L$ of $I$ and every $B\in\mathcal U_y$ the relative means exist, lie within $r_G$ of $\id$ in operator norm, and $\mathscr G_{y,U}(g_0B|_{\mathcal E_y})$ is analytic in $B|_{\mathcal E_y}$; all are compositions, with the linear map $B\mapsto g_0\mathsf C_UB$ (of norm at most $C_{\mathsf C}$ for the Hilbert norms, by Lemma~\ref{lem:c2-uniform-chart}(i)), of functions of the phase analytic on a complex neighbourhood of the closed ball $\| A\|\leqslant \vartheta_*/2$ (on the complex ball of radius $\vartheta_*$ for the row solution and the phase) with derivatives up to order five bounded on that closed ball by constants depending on $L$ and the group only, so their derivatives in $B$ up to order five are bounded on the polydiscs by such constants. Where $(A_{\mathrm e},a)\in\mathcal D_{\xi}$ these objects are CRE's under the reading of (B4)(i), by (B4)(i) and the uniqueness of the small solution: for every row within $2R_{\mathrm c}+2L$ of $I$ the phase $g_0\mathsf C_U\Psi_{\xi}(a+A_{\mathrm e})$ lies in the ball $\| A|_{\mathcal E_c}\|\leqslant \vartheta_*/4$, because $a$ is in the open box and $\Psi_{\xi}(a+A_{\mathrm e})= a+A_{\mathrm e}$ with $| A_{\mathrm e}(b)|< r$ on the two-block set (Lemma~\ref{lem:collar-coverage}(i),(iii),(iv)), so $\Psi_{\xi}(a+A_{\mathrm e})\in\mathcal U_c$ and the bound on $g_0\mathsf C_U\mathcal U_c$ applies. This is Lemma~\ref{lem:charts}.
\end{lemma}

\begin{proof}
\emph{The input of a primitive.} On the enlarged box, $| a_b+tA_{\mathrm e}(b)|\leqslant 2r+r_{\mathrm e,\xi}\leqslant 2s_{\xi}$ for every bond $b$ within $3R_{\mathrm c}$ of $I$ (Lemma~\ref{lem:collar-coverage}(iii) bounds the retained inputs by $r_{\mathrm e,\xi}$), in particular for the free bonds of the two-block sets of the near rows and for the free bonds of the plaquettes and blocks of the near primitives. The rows of a near primitive are within $2R_{\mathrm c}$ of $I$ (applicability paragraph). The \emph{input} $\mathfrak u$ of a primitive is the vector $g_0\mathsf C_UB$ restricted to the following bond set $S$: for a plaquette primitive $E_{\xi,p}$, $p_{\xi,p}$ the union of the sets $\mathcal E_c$ over the rows $c$ of $p$ (at most four) together with the free bonds of $\partial p$ not in those sets; for the Haar primitive $p^{\mathrm H}_{\xi,b}$ the set $\mathcal E_c$ when $b= b_0(c)$ and $\{ b\}$ when $b$ is free; for the Jacobian primitive $p^{\mathrm J}_{\xi,c}$ the set $\mathcal E_c$; for a gauge primitive $E_{\xi,y}$, $p_{\xi,y}$ the set $\mathcal E_y$ (on which $\mathsf C_U$ is the identity, all its bonds being free). The primitive is a function of the phase $q_U(B)$ on its bonds, and that phase is a function of the input: on $\mathcal E_c$ it is $\mathfrak u|_{\mathcal E_c}-h_U\widetilde D_{U,c}(\mathfrak u|_{\mathcal E_c})$ (Lemma~\ref{lem:c2-uniform-chart}(ii)), and on a free bond outside the row sets it is $\mathfrak u_b= g_0B(b)$. With $B= a+tA_{\mathrm e}$, by Lemma~\ref{lem:c2-uniform-chart}(i) each row component satisfies $\|\mathfrak u|_{\mathcal E_c}\|\leqslant C_{\mathsf C}g_0\cdot2s_{\xi}\leqslant \vartheta_*/4$ and each free-bond component outside the row sets $h(\mathfrak u_b)\leqslant |\mathfrak u_b|\leqslant 2g_0s_{\xi}\leqslant \vartheta_*/4$ (threshold item 3 and $C_{\mathsf C}\geqslant 1$), so for the plaquette, Haar and Jacobian primitives (at most four rows and four further free bonds)
\begin{equation*}
    \|\mathfrak u\|\leqslant 16C_{\mathsf C}g_0s_{\xi},\qquad\| D_a\mathfrak u\|\leqslant 5C_{\mathsf C}g_0,\qquad\|\partial_t\mathfrak u\|\leqslant 16C_{\mathsf C}g_0s_{\xi}\;,
\end{equation*}
the second in the operator norm from the Hilbert norm of $\mathfrak g^{S_{\alpha}}$ (the Hilbert bound of Lemma~\ref{lem:c2-uniform-chart}(i) on each of at most four rows, plus the identity on the free bonds); for the gauge primitive, whose input $g_0B|_{\mathcal E_y}$ has $|\mathcal E_y|=4L^3(L-1)$ free bonds, $\|\mathfrak u\|\leqslant g_0|\mathcal E_y|^{1/2}\cdot2s_{\xi}\leqslant 2C_{\mathsf C}g_0s_{\xi}\leqslant \vartheta_*/4$ (as $C_{\mathsf C}\geqslant (4L^3(L-1))^{1/2}$ and $h\leqslant |\bigcdot |$), $\| D_a\mathfrak u\|\leqslant g_0$ and $\|\partial_t\mathfrak u\|\leqslant 2C_{\mathsf C}g_0s_{\xi}$, so the three displayed bounds hold for all four kinds. The input is linear in $B$, hence affine in $(a,t)$ and linear in $a$ at fixed $t$; $D_a\mathfrak u$ is the restriction of $D_B\mathfrak u= g_0\mathsf C_U|_{S}$ to the young directions, and the same count gives $\| D_B\mathfrak u\|\leqslant 5C_{\mathsf C}g_0$ (respectively $\leqslant g_0$ for the gauge primitive), so every bound below on a derivative in $a$ holds verbatim for the derivative of the same order in $B$ on $\mathfrak g^{S_{\alpha}}$, retained inputs included. Each primitive is $\varphi(\mathfrak u)$ for a function $\varphi$ of the input analytic on a complex neighbourhood of the closed set $\mathcal P=\{\mathfrak u\in\mathfrak g^S_{\mathbb C}:\|\mathfrak u|_{\mathcal E_c}\|\leqslant \vartheta_*/2\text{ for each row }c\text{ of the primitive},\ h(\mathfrak u_b)\leqslant \vartheta_*/2\text{ for each free bond }b\text{ of }S\text{ outside the row sets}\}$ (for the gauge primitive, the closed ball $\|\mathfrak u\|\leqslant \vartheta_*$ of $\mathfrak g^{\mathcal E_y}_{\mathbb C}$), with derivatives up to order five bounded on $\mathcal P$ by constants of $L$ and the group: Lemma~\ref{lem:c2-uniform-chart}(iii) for the row solution, the phase and the two logarithms, each a function of one row component (the Haar logarithm at a free bond outside the row sets is the function of (iii) at a free bond of a two-block set, where $q_U(A)(b)= A(b)$, a function of the single component $\mathfrak u_b$); for the plaquette term, with $W(q)=(e^{iq}U)(\partial p)$, $s_{p,U}(q)=1-\operatorname{Re}\tr_NW(q)$ equals $1-\frac12(\tr_NW(q)+\tr_NW(q)^{-1})$ for real $q$, $W(q)$ being unitary, and $W(q)$ and $W(q)^{-1}$ are ordered products of the fixed links and of $e^{\pm iq(b)}$, entire functions of the phases $q(b)$, which are analytic functions of one row component (with $h(q_U(A)(b))\leqslant \| A\|+H_*2K_*\| A\|^2\leqslant \vartheta_*$ on the ball) or are free-bond components of $\mathfrak u$; this is the analytic function meant by $s_{p,U}(q_U(\bigcdot ))$, with derivatives bounded on $\mathcal P$ by the chain rule; for the gauge term, Lemma~\ref{lem:c2-uniform-chart}(iv). The segment from $0$ to $\mathfrak u$ lies in $\mathcal P$, each component of $\mathfrak u$ being within $\vartheta_*/4$ of zero. Hence every primitive at the input $a+tA_{\mathrm e}$ is the restriction to real $(a,t)$ of a function analytic in $(a,t)$ on a complex neighbourhood of the enlarged box times $[0,1]$ (the composition of $\varphi$ with the affine map), and the bounds below are uniform there.

\emph{Energy residual.} $E_{\xi,\alpha}(a+tA_{\mathrm e})= g_0^{-2}\varphi(\mathfrak u)$ with $\varphi(\mathfrak u)=-[s_{\alpha,U}(q_U(\mathfrak u))-s_{\alpha,U}(0)]$ for a plaquette and $\varphi(\mathfrak u)=-\mathscr G_{y,U}(\mathfrak u)$ for a centre. Since $\mathfrak u$ is linear in $B$, the Taylor polynomial of degree two in $B$ at $B=0$ subtracted in \eqref{eq:native-source-energy-residual} is that of $\varphi$ at $0$ composed with the linear map, so $p_{\xi,\alpha}= g_0^{-2}[\varphi(\mathfrak u)-D\varphi(0)\mathfrak u-\frac12D^2\varphi(0)[\mathfrak u,\mathfrak u]]$, the second-order Taylor residual at the input, bounded by $g_0^{-2}\frac16\| D^3\varphi\|\|\mathfrak u\|^3\leqslant Cg_0s_{\xi}^3$ (the segment from $0$ to $\mathfrak u$ lies in $\mathcal P$). Its first derivative in $a$ is $g_0^{-2}(D\varphi(\mathfrak u)-D\varphi(0)-D^2\varphi(0)\mathfrak u)D_a\mathfrak u$, bounded by $g_0^{-2}\frac12\| D^3\varphi\|\|\mathfrak u\|^2\cdot5C_{\mathsf C}g_0\leqslant Cg_0s_{\xi}^2$; its second derivative is $g_0^{-2}(D^2\varphi(\mathfrak u)-D^2\varphi(0))[D_a\mathfrak u,D_a\mathfrak u]$, bounded by $g_0^{-2}\| D^3\varphi\|\|\mathfrak u\|(5C_{\mathsf C}g_0)^2\leqslant Cg_0s_{\xi}$; its third and fourth derivatives are $g_0^{-2}D^3\varphi(\mathfrak u)[D_a\mathfrak u]^3$ and $g_0^{-2}D^4\varphi(\mathfrak u)[D_a\mathfrak u]^4$, bounded by $Cg_0$ and $Cg_0^2$. The Taylor polynomial subtracted is that of the germ (Section~\ref{sec:source-phase-primitives}), which is the Taylor polynomial of $\varphi$ at $0$ by Lemma~\ref{lem:c2-uniform-chart}(ii).

\emph{Haar logarithm.} $\psi(\mathfrak u)=\log\sigma(q_U(\mathfrak u)(b))-\log\sigma(0)$, the primitive \eqref{eq:native-source-haar-primitive}, has $\psi(0)=0$ and $D\psi(0)=0$ (Lemma~\ref{lem:c2-uniform-chart}(iii): $\sigma$ is even and $Dq_U(0)$ is linear). Hence $|\psi(\mathfrak u)|\leqslant \frac12\| D^2\psi\|\|\mathfrak u\|^2\leqslant Cg_0^2s_{\xi}^2$, $\| D_a\psi\|\leqslant \| D^2\psi\|\|\mathfrak u\|\cdot5C_{\mathsf C}g_0\leqslant Cg_0^2s_{\xi}\leqslant Cg_0s_{\xi}^2$, $\| D_a^2\psi\|\leqslant \| D^2\psi\|(5C_{\mathsf C}g_0)^2\leqslant Cg_0s_{\xi}$, and the third and fourth derivatives are bounded by $Cg_0^3\leqslant Cg_0$ and $Cg_0^4\leqslant Cg_0^2$.

\emph{Constraint logarithm.} $\lambda(\mathfrak u)=\log|\det_{\mathfrak g}(\id-h_U(c)\partial_{A(b_0(c))}\widetilde D_{U,c}(\mathfrak u))|=\log\det_{\mathfrak g}(\cdots)$ (the determinant is positive, Lemma~\ref{lem:c2-uniform-chart}(iii)) has $\lambda(0)=0$ and derivatives up to order five bounded on $\mathcal P$; its first derivative at $0$ need not vanish. Hence $|\lambda(\mathfrak u)|\leqslant \| D\lambda\|\|\mathfrak u\|\leqslant Cg_0s_{\xi}$, $\| D_a\lambda\|\leqslant \| D\lambda\|\cdot5C_{\mathsf C}g_0\leqslant Cg_0\leqslant Cg_0s_{\xi}^2$, $\| D_a^2\lambda\|\leqslant \| D^2\lambda\|(5C_{\mathsf C}g_0)^2\leqslant Cg_0^2\leqslant Cg_0s_{\xi}$ (as $g_0\leqslant 1$ and $s_{\xi}\geqslant 1$), and the third and fourth derivatives are bounded by $Cg_0^3\leqslant Cg_0$ and $Cg_0^4\leqslant Cg_0^2$. Let $C'_{\mathrm p}$ be the maximum over the finitely many templates of the constants $C$ above and put $C_{\mathrm p}=2C'_{\mathrm p}$. Then \eqref{eq:c2-primitive-sizes}, the first two bounds of \eqref{eq:convex-primitive-derivatives} and the third- and fourth-order bounds hold with room to spare, and the last bound of \eqref{eq:convex-primitive-derivatives} is the difference of two values each bounded by $C'_{\mathrm p}g_0s_{\xi}^3$ (for all three kinds, since $s_{\xi}\geqslant 1$ and $g_0\leqslant 1$), hence by $2C'_{\mathrm p}g_0s_{\xi}^3= C_{\mathrm p}g_0s_{\xi}^3$.

\emph{The chart statements.} Let $c$ be a row within $2R_{\mathrm c}+2L$ of $I$; by the applicability paragraph Lemma~\ref{lem:c2-uniform-chart} applies to it, and (ii)--(iii) give, for every real $A$ with $\| A|_{\mathcal E_c}\|\leqslant \vartheta_*/2$, the unique small solution $\varphi_c(A|_{\mathcal E_c})$ with $\varphi_c$ analytic on the complex ball of radius $\vartheta_*$ and real on real phases, $\varphi_c(0)=0$, $D\varphi_c(0)=0$, the positivity of the determinant, and the analyticity of the row solution, the phase and the two logarithms on a complex neighbourhood of the closed ball of radius $\vartheta_*/2$, with derivatives up to order five bounded on that ball. For $B\in\mathcal U_c$, $\max_{b\in\mathcal E_c\cap E_*}| B(b)|<2+4r=2(1+2r)$ and $\| g_0\mathsf C_UB\|_{\mathcal E_c}\leqslant 2C_{\mathsf C}g_0(1+2r)\leqslant \vartheta_*/4$ by threshold item 3, so $g_0\mathsf C_U\mathcal U_c$ lies in the phase ball; $B\mapsto g_0\mathsf C_UB$ is linear with norm at most $g_0C_{\mathsf C}\leqslant C_{\mathsf C}$ for the Hilbert norms (Lemma~\ref{lem:c2-uniform-chart}(i)), so the derivatives in $B$ are bounded by the same constants times $C_{\mathsf C}^{\,k}$. For $c$ within $2R_{\mathrm c}$ of $I$, the plaquette term of a plaquette $p$ with a bond in $\mathcal E_c$ is analytic in $B$ on the intersection of the polydiscs of the rows of $p$, as above, and these rows are within $2R_{\mathrm c}+2L$ of $I$ (applicability paragraph), so the chart applies to each of them. For a centre $y$ within $2R_{\mathrm c}+2L$ of $I$, Lemma~\ref{lem:c2-local-gauge} applied to the block $\mathcal B_y$ (applicability paragraph) gives a gauge with phases at most $\widehat\mu<\tau_L/2$ on $\mathcal E_y$; every bond of $\mathcal E_y$ is free, so $q_U(B)= g_0B$ there and $\| g_0B|_{\mathcal E_y}\|\leqslant g_0|\mathcal E_y|^{1/2}\max_b| B(b)|\leqslant g_0(4L^3(L-1))^{1/2}\cdot2(1+2r)\leqslant 2C_{\mathsf C}g_0(1+2r)\leqslant \vartheta_*/4$ under threshold item 3 (using $h\leqslant |\bigcdot |$, $|\mathcal E_y|=4L^3(L-1)$ and $C_{\mathsf C}\geqslant (4L^3(L-1))^{1/2}$); Lemma~\ref{lem:c2-uniform-chart}(iv) then gives the existence of the relative means within $r_G$ of $\id$ and the analyticity and derivative bounds of $\mathscr G_{y,U}(g_0B|_{\mathcal E_y})$ on $\mathcal U_y$. The values at $q=0$ of the plaquette and gauge terms are subtracted in the primitives. Finally, where $(A_{\mathrm e},a)\in\mathcal D_{\xi}$ and $c$ is within $2R_{\mathrm c}+2L$ of $I$, the phase $A= g_0\mathsf C_U\Psi_{\xi}(a+A_{\mathrm e})$ has $\| A|_{\mathcal E_c}\|\leqslant \vartheta_*/4$ (statement of the lemma), $\varphi_c(A|_{\mathcal E_c})$ is the row solution $\widetilde D_c(A)$ of (B4)(i), which is CRE's under the reading of (B4)(i), by Lemma~\ref{lem:c2-fibre-smallness}(iii), and the phase, the Haar density, the Jacobian and the plaquette and gauge terms of CRE (1.14)--(1.15) are the same functions of it.
\end{proof}

\begin{lemma}[host multiplicity and potential size]\label{lem:c2-potential-size}
Let $\mathfrak T$ be the set of primitive templates (kind, label and free-input shape, including the shapes truncated by $\Omega_1$) modulo $L\mathbb Z^4$, anchor each template at one of its input bonds, and let $D_{\mathrm h}$ be one plus the maximum $\ell^{\infty}$ diameter of a lifted host \eqref{eq:part1-potential-host} and $B_{\mathrm h}$ the maximum number of $M$-cubes in one, both over $\mathfrak T$. With
\begin{equation}\label{eq:part1-host-multiplicity}
    N_{\mathrm h}=|\mathfrak T|(2\lceil D_{\mathrm h}\rceil+3)^4,\qquad
    C_0= N_{\mathrm h}C_{\mathrm p}e^{\kappa(B_{\mathrm h}+4)},\qquad
    \omega(g_0)= C_0g_0(1+r)^3\;,
\end{equation}
at most $N_{\mathrm h}$ primitives are assigned to any host $Y$, $d_0(Y)\leqslant B_{\mathrm h}+4$, and on the enlarged box, for every host,
\begin{equation}\label{eq:native-source-interaction-size}
    \sup|\mathbf V_{\xi}(Y,a)|\leqslant \omega(g_0)e^{-\kappa d_0(Y)}\;.
\end{equation}
\end{lemma}

\begin{proof}
Fix a vertex of $Y$. Every primitive assigned to $Y$ has its stencil, hence its anchor, in the lifted host, within $\ell^{\infty}$ distance $D_{\mathrm h}$ of that vertex; there are at most $(2\lceil D_{\mathrm h}\rceil+1)^4$ such anchor positions, and a template with an anchor position determines the primitive. Projection to the torus only identifies positions. By Lemma~\ref{lem:c2-domain-geometry}(i), $d_0(Y)\leqslant N_0(Y)-1\leqslant B_{\mathrm h}$. At $t=0$, $s_{\xi}$ may be replaced by $1+r$ in Lemma~\ref{lem:c2-primitive-bounds}, since no retained input enters; each of the three kinds is then bounded by $C_{\mathrm p}g_0(1+r)^3$, and $|\mathbf V_{\xi}(Y,a)|\leqslant N_{\mathrm h}C_{\mathrm p}g_0(1+r)^3\leqslant N_{\mathrm h}C_{\mathrm p}g_0(1+r)^3e^{\kappa(B_{\mathrm h}+4)}e^{-\kappa d_0(Y)}$.
\end{proof}

The definition gives $C_0\geqslant C_{\mathrm p}$, since $N_{\mathrm h}\geqslant 1$ and the exponential factor is at least one. The constant $C_0$ depends on $L$, the group, $M$ and $\kappa$ (through $B_{\mathrm h}$, $D_{\mathrm h}$), not on the volume, the branch or $g_0$. The hypothesis of Lemma~\ref{lem:ibp} holds for $f= f_{\xi}$ on the enlarged box, which is a neighbourhood of the closed cutoff box: the exponent $\mathsf Q_{\xi,t}$ is a quadratic polynomial plus $P_{\xi,t}$, a finite sum of near primitives, each analytic there by Lemma~\ref{lem:c2-primitive-bounds}; and $f_{\xi}(a)= W_{\chi}(C;U_{\xi}(a))$ is a character of a product of the reconstructed links \eqref{eq:c-reconstructed-field} over the contour bonds, each of which depends on $a$ through the rows of a free-input set meeting $I$ (Lemma~\ref{lem:collar}(i) puts $\Psi_{\xi}(a+A_{\mathrm e})= a+A_{\mathrm e}$ on those rows, so the phase is \eqref{eq:native-collar-source-phase}), where $q_U$ is analytic in $a$ by Lemma~\ref{lem:c2-uniform-chart}(ii)--(iii); contour bonds whose rows do not meet $I$ do not depend on $a$. Both factors are therefore analytic on the enlarged box, uniformly in $t\in[0,1]$.

\subsection{Proof of the Gaussian payment}
\label{app:fixed-kernel-payment}

\begin{proof}[Proof of Lemma~\ref{lem:gaussian-payment}]
\emph{Tail of the walk series.} Throughout, $E_0$, $C_{E_0}$, $\Gamma_{E_0}(\sigma)$ are as in \eqref{eq:fixed-kernel-C-Gamma} and the sentence preceding it (their block decomposition is \eqref{eq:fixed-kernel-exact-blocks}), and $S_{Z_0}(1)=\mathsf G^{1/2}$. Let $\mathsf N$ be the set of bonds $b'\in I$ with $d_{\infty,T}(b'_-,Z_0)\leqslant R+1$. Since $\iota_{E_0}^*H$ has nonzero entries $H(b,b')$ only for $d(b_-,b'_-)\leqslant R$ and $b\in E_0$ meets $Z_0$, $\Gamma_{E_0}(\sigma)-\Gamma_{E_0}(1)=\iota_{E_0}^*H\Pi_{I\setminus E_0}\Pi_{\mathsf N}(S_{Z_0}(\sigma)-\mathsf G^{1/2})$. A walk $\omega=(r;b_0,\ldots,b_{m+1})$ with $b_0\in\mathsf N$ has $| w_{\omega}(\sigma)-1|\ne0$ only if $D(\omega)$ meets a cube $\Delta\not\subseteq Z_0'$. The open $M$-neighbourhood of $Z_0$ is $\bigcup_{Q\subseteq Z_0}\operatorname{int}Q^{\sim}$, and by definition every $LM$-cube meeting it lies in $Z_0'$ (Section~\ref{sec:fixed-first-kernel}). Hence $D(\omega)$ contains a point at distance at least $M$ from $Z_0$, and at least $M-R-1$ from $b_{0,-}$. Each block of $\mathsf A_H^r$ ($r< n_*$) or $\mathsf P_H$ vanishes beyond distance $R_*$, and $D_b$ lies within $\rho_{\mathrm w}$ of $b_-$, so $D(\omega)$ lies within $(m+1)R_*+\rho_{\mathrm w}$ of $b_{0,-}$. By \eqref{eq:fixed-kernel-collar-size}, $M/4\geqslant \rho_{\mathrm w}+R+R_*+2$, so $(m+1)R_*\geqslant M-R-1-\rho_{\mathrm w}\geqslant \frac34M+R_*+1$ and $m> M/(2R_*)\geqslant N_M$. The absolute block row sum of the terms with $m\geqslant N_M$ is, by the bounds in the proof of Lemma~\ref{lem:walk-series} and $| w_{\omega}-1|\leqslant 2a_*^{m+1}$,
\begin{equation}\label{eq:fixed-kernel-rooted-tail}
    s_H^{-1/2}\sum_{r< n_*}\sum_{m\geqslant N_M}a_{n_*m+r}\,2a_*^{m+1}V_*\beta^m
    \leqslant 2s_H^{-1/2}\frac{n_*V_*a_*}{1-a_*\beta}4^{-N_M}=\delta_M\;,
\end{equation}
using $a_*\beta\leqslant 1/4$. The absolute block column sums of these terms obey the same bound, because every walk in the tail has $m\geqslant N_M$ whatever its first bond. Schur's test gives $\|\Pi_{\mathsf N}(S_{Z_0}(\sigma)-\mathsf G^{1/2})\|\leqslant \delta_M$. With $\| H\|\leqslant h_+$ this is the first bound of (i); with $\| C_{E_0}\|\leqslant \gamma^{-1}$, since $\iota_{E_0}^*H\iota_{E_0}\geqslant \gamma$, the second follows.

\emph{Reserve identity.} Write $\Gamma=\Gamma_{E_0}(\sigma)=\Gamma_1+i\Gamma_2$ with real $\Gamma_1,\Gamma_2$; $\Gamma_{E_0}(1)$ is real, so $\Gamma_2=\operatorname{Im}(\Gamma_{E_0}(\sigma)-\Gamma_{E_0}(1))$. The quadratic form in \eqref{eq:fixed-kernel-density} is bilinear, so for real $X$, $\operatorname{Re}\langle\Gamma X,C_{E_0}\Gamma X\rangle=\langle\Gamma_1X,C_{E_0}\Gamma_1X\rangle-\langle\Gamma_2X,C_{E_0}\Gamma_2X\rangle$ and $| e^{-\langle B,\Gamma X\rangle}|= e^{-\langle B,\Gamma_1X\rangle}$ for real $B$. For $0\leqslant \alpha_G<\gamma$ put $C_{E_0,\alpha_G}=(\iota_{E_0}^*H\iota_{E_0}-\alpha_G)^{-1}$. The Gaussian integral over $B$ with density $(2\pi)^{-d_G| E_0|/2}\det(C_{E_0})^{-1/2}e^{-\frac12\langle B,(C_{E_0}^{-1}-\alpha_G)B\rangle-\langle B,\Gamma_1X\rangle}$ equals $\det_{E_0}(\id-\alpha_GC_{E_0})^{-1/2}e^{\frac12\langle\Gamma_1X,C_{E_0,\alpha_G}\Gamma_1X\rangle}$. Collecting the $X$-exponent,
\begin{equation}\label{eq:fixed-kernel-actual-reserve-form}
    \int\Bigl| e^{-\frac12\langle\Gamma X,C_{E_0}\Gamma X\rangle}e^{-\langle B,\Gamma X\rangle}\Bigr| e^{\frac{\alpha_G}{2}\| B\|^2}\dd\mu_{C_{E_0}}(B)
    =\det_{E_0}(\id-\alpha_GC_{E_0})^{-1/2}e^{\frac12\langle X,T_{\alpha_G}X\rangle},\quad
    T_{\alpha_G}=\Gamma_1^{\mathsf T}(C_{E_0,\alpha_G}-C_{E_0})\Gamma_1+\Gamma_2^{\mathsf T}C_{E_0}\Gamma_2\;,
\end{equation}
a real symmetric positive semidefinite matrix on $\mathfrak g^I$ ($C_{E_0,\alpha_G}\geqslant C_{E_0}$ for $\alpha_G\geqslant 0$). By Lemma~\ref{lem:component-blocks} the columns of $\Gamma$ outside $J_Z$ vanish, so $T_{\alpha_G}$ is supported on $J_Z$, and $\int e^{\frac12\langle X,T_{\alpha_G}X\rangle}\dd\mu_0(X)=\det_{J_Z}(\id-T_{\alpha_G})^{-1/2}$ when $\| T_{\alpha_G}\|<1$. Integrating \eqref{eq:fixed-kernel-actual-reserve-form} against $\mu_0$ therefore gives the integral of (ii) in closed form, $\det_{E_0}(\id-\alpha_GC_{E_0})^{-1/2}\det_{J_Z}(\id-T_{\alpha_G})^{-1/2}$, once $\| T_{\alpha_G}\|<1$ is shown.

\emph{Norm of $T_{\alpha_G}$.} The resolvent identity gives $C_{E_0,\alpha_G}-C_{E_0}=\alpha_GC_{E_0,\alpha_G}C_{E_0}$, of norm at most $\alpha_G/(\gamma(\gamma-\alpha_G))$. By (i) and $\delta_M\leqslant 1$, $\|\Gamma_1\|\leqslant \|\Gamma_{E_0}(1)\|+h_+\delta_M\leqslant h_+\gamma^{-1/2}+h_+= A_*$, since $\|\Gamma_{E_0}(1)\|\leqslant \| H\|\|\mathsf G^{1/2}\|\leqslant h_+\gamma^{-1/2}$ ($\mathsf G= H^{-1}$ and $H\geqslant \gamma$); and $\|\Gamma_2\|\leqslant h_+\delta_M$. Hence, for $\alpha_G\leqslant \min(\gamma/2,\gamma^2/(16A_*^2))$,
\begin{equation}\label{eq:fixed-kernel-reserve-payment}
    \| T_{\alpha_G}\|\leqslant \frac{\alpha_GA_*^2}{\gamma(\gamma-\alpha_G)}+\frac{h_+^2\delta_M^2}{\gamma}\leqslant \frac{2\alpha_GA_*^2}{\gamma^2}+\frac18\leqslant \frac14\;.
\end{equation}

\emph{Determinant bounds.} The eigenvalues of $\alpha_GC_{E_0}$ lie in $[0,1/2]$ and $-\frac12\log(1-t)\leqslant t$ for $t\leqslant 1/2$, so $\log\det_{E_0}(\id-\alpha_GC_{E_0})^{-1/2}\leqslant \alpha_G\operatorname{tr}C_{E_0}\leqslant (\alpha_G/\gamma)d_G| E_0|$. The eigenvalues of $T_{\alpha_G}$ lie in $[0,1/4]$ and $-\frac12\log(1-t)\leqslant \frac23t$ there, so $\log\det_{J_Z}(\id-T_{\alpha_G})^{-1/2}\leqslant \frac23\operatorname{tr}T_{\alpha_G}\leqslant \frac23d_G| J_Z|\| T_{\alpha_G}\|$. With $E_0\subseteq J_Z$ (Lemma~\ref{lem:component-blocks}),
\begin{equation}\label{eq:fixed-kernel-trace-payment}
    \log\bigl[\det_{E_0}(\id-\alpha_GC_{E_0})^{-1/2}\det_{J_Z}(\id-T_{\alpha_G})^{-1/2}\bigr]
    \leqslant d_G| J_Z|\Bigl(\frac{\alpha_G}{\gamma}+\frac{4\alpha_GA_*^2}{3\gamma^2}+\frac{2h_+^2\delta_M^2}{3\gamma}\Bigr)
    \leqslant C_{\mathrm d}(\alpha_G+\delta_M^2)| J_Z|,\qquad
    C_{\mathrm d}= d_G\max\Bigl(\frac1\gamma+\frac{4A_*^2}{3\gamma^2},\frac{2h_+^2}{3\gamma}\Bigr)\;.
\end{equation}
A bond of $J_Z$ has its initial site in a cube of $Z$, so $| J_Z|\leqslant 4(LM+1)^4N_1(Z)$. Since $n_*,R_*,C_{\mathrm w}$ are fixed before $M$ and $\delta_M=2C_{\mathrm w}4^{-\lfloor M/(2R_*)\rfloor}$ decays exponentially in $M$, there is $M$ with $4C_{\mathrm d}\delta_M^2(LM+1)^4\leqslant 1/2$; then $\alpha_G\leqslant (8C_{\mathrm d}(LM+1)^4)^{-1}$ gives the bound $e^{N_1(Z)}$. Both $C_{E_0}$ and $\Gamma_{E_0}$ are block diagonal over the components $D$ of $Z$ by \eqref{eq:fixed-kernel-exact-blocks}, hence so are $T_{\alpha_G}$ and both determinants, which proves the factorisation.

\emph{Bad-bond reserve.} On $\{| B_b|\geqslant r\}$, $h(B_b)^2\geqslant c_{\mathrm n}| B_b|^2\geqslant c_{\mathrm n}r^2$, so $\mathbf 1_{\{| B_b|\geqslant r\}}\leqslant e^{\frac{\gamma_2}{2}(h(B_b)^2-c_{\mathrm n}r^2)}$; the product over $\mathsf P$ is (iii).
\end{proof}

\subsection{The Mayer constants and the activity bounds}
\label{app:mayer-constants}

Fix, in this order and depending on $L$ and the group only: $C_{\mathrm g}=162$ and $\delta=1/40$; a decay rate $\kappa\geqslant 3$ with
\begin{equation}\label{eq:c2-kappa-choice}
    \delta\kappa\geqslant 2C_{\mathrm g}\log128,\qquad
    \lambda'=(1-5\delta)\frac{L\kappa}{2049},\qquad
    \lambda_*=\lambda'-2C_{\mathrm g}\geqslant 2C_{\mathrm g}\log128\;;
\end{equation}
then
\begin{equation}\label{eq:c2-mayer-constants}
    \alpha_6=\min\Bigl(\frac18,\frac{\delta\kappa}{2C_{\mathrm g}B(\delta\kappa/2)}\Bigr),\qquad
    \kappa_1-1\geqslant 5\lambda'+1,\qquad
    \theta=\frac{\lambda_*}{16C_{\mathrm g}},\qquad
    C_{\mathrm s}=\frac{2049L^3C_{\mathrm g}B(\delta\kappa)e^{L\kappa/2}}{\delta\kappa}\;,
\end{equation}
and a number $\omega_*>0$ with
\begin{equation}\label{eq:native-source-vacuum-KP-budget}
    \eta_*=2e^{1/2+5\kappa}\omega_*\leqslant \frac{\alpha_6}{4},\qquad
    \varepsilon_*=\frac{C_{\mathrm s}\eta_*}{\alpha_6}\leqslant \min\Bigl(\frac14,\frac{1}{72B(\lambda_*/2)}\Bigr)\;;
\end{equation}
these conditions involve only $\kappa$, $\alpha_6$, $C_{\mathrm s}$ and $\lambda_*$, so $\omega_*$, $\eta_*$ and $\varepsilon_*= C_{\mathrm s}t_0$ are fixed before $M_1$, $M$ and the Gaussian reserve $\alpha_G$, as Lemma~\ref{lem:activity-bounds} states. Put $t_0=\eta_*/\alpha_6\leqslant 1/4$ and, once $\alpha_G$ of Lemma~\ref{lem:gaussian-payment} is fixed after $M$, $\gamma_2=\alpha_G/2$, so $a_{\mathrm{bad}}= c_{\mathrm n}\alpha_Gr^2/4$. The coupling is then taken so small that $\omega(g_0)\leqslant \omega_*$, the chart condition of Lemma~\ref{lem:c2-primitive-bounds} holds, and
\begin{equation}\label{eq:native-source-bad-budget}
    a_{\mathrm{bad}}\geqslant \max\Bigl(160\bigl(\kappa+\log(8M^4+8)+1\bigr),\ 64\log\frac1{\eta_*}\Bigr)\;.
\end{equation}
These are compatible because $r= A_1(\log g_0^{-2})^{p_0}$ (the radius of the literal cutoff \eqref{eq:literal-cutoff}) tends to infinity while $g_0(1+r)^3\to0$. Note $B(\delta\kappa/2)\leqslant 2e^{-\delta\kappa/(2C_{\mathrm g})}$, $B(\delta\kappa)\leqslant 2e^{-\delta\kappa/C_{\mathrm g}}$ and $B(\lambda_*/2)\leqslant 2e^{-\lambda_*/(2C_{\mathrm g})}$ by Lemma~\ref{lem:c2-domain-geometry}(iv) and \eqref{eq:c2-kappa-choice}.

\begin{lemma}[cover sum]\label{lem:c2-cover-sum}
For every nonempty union $Y_0$ of $M$-cubes,
\begin{equation}\label{eq:native-source-cover-sum}
    \sum_{\mathcal D\subseteq\mathbf D_0\, | \,\bigcup_{Y\in\mathcal D}Y= Y_0}\ \prod_{Y\in\mathcal D}\alpha_6e^{-\delta\kappa\ell_0(Y)}\leqslant 1\;.
\end{equation}
\end{lemma}

\begin{proof}
The members of $\mathbf D_0$ are localisation domains of $\pi_0$ and $\mathcal D$ is a set of distinct domains. For a cover, $\sum_{Y\in\mathcal D}\ell_0(Y)\geqslant \sum_YN_0(Y)/C_{\mathrm g}\geqslant N_0(Y_0)/C_{\mathrm g}$ by Lemma~\ref{lem:c2-domain-geometry}(i), so $\prod_Ye^{-\delta\kappa\ell_0(Y)/2}\leqslant e^{-\delta\kappa N_0(Y_0)/(2C_{\mathrm g})}$. Dropping the cover constraint, the remaining sum over sets of domains contained in $Y_0$ is at most $\prod_{Y\subseteq Y_0}(1+\alpha_6e^{-\delta\kappa\ell_0(Y)/2})\leqslant \exp\bigl(\alpha_6\sum_{Q\subseteq Y_0}\sum_{Y\ni Q}e^{-\delta\kappa\ell_0(Y)/2}\bigr)\leqslant \exp(\alpha_6B(\delta\kappa/2)N_0(Y_0))$ by \eqref{eq:native-source-animal-sum}. The product of the two bounds is at most one by the choice of $\alpha_6$.
\end{proof}

\begin{lemma}[selected majorant]\label{lem:c2-selected-majorant}
On every branch, with $\mathbf W=\mathbf V_{\xi}$, $\omega(g_0)\leqslant \omega_*$ and the chart condition, every selected integral of Definition~\ref{def:selection} with $F=1$ satisfies
\begin{equation}\label{eq:fixed-kernel-selected-majorant}
    |\mathcal A^{\mathrm f}_{(\mathcal D,\mathsf P,Z)}[1]|
    \leqslant e^{-(\kappa_1-1)N_1(Z\setminus Z_0')+N_1(Z)-a_{\mathrm{bad}}|\mathsf P|}\prod_{Y\in\mathcal D}\omega_*e^{1/2}e^{-\kappa d_0(Y)}\;.
\end{equation}
\end{lemma}

\begin{proof}
The weakening in \eqref{eq:c-selected-integral} is, for each cube $\Delta$ of $Z\setminus Z_0'$, the operation $G\mapsto\int_0^1\dd s\oint_{|\sigma_{\Delta}|= e^{\kappa_1}}\frac{G\,\dd\sigma_{\Delta}}{2\pi i(\sigma_{\Delta}-s)^2}$ applied to the inner integral (the $t_Y$-, $X$- and $B$-integrals) as a function of $\sigma_{\Delta}$ on the circle; it needs only the measurability of the inner integral in $\sigma$ on the distinguished boundary (the integrand is continuous in $\sigma$, since $\Gamma_{E_0}(\sigma)$ is the walk series of Lemma~\ref{lem:walk-series}, and Fubini applies to the iterated integral because the integral of the modulus of the inner integrand is bounded uniformly on the distinguished boundary, by the bounds below) and the bound $\sup_{s\in[0,1],|\sigma|= e^{\kappa_1}}|\sigma-s|^{-2}\cdot e^{\kappa_1}\leqslant e^{\kappa_1}(e^{\kappa_1}-1)^{-2}\leqslant e^{-(\kappa_1-1)}$, the last step because $e^{\kappa_1}\geqslant 3$. Applying it for each cube of $Z\setminus Z_0'$ bounds \eqref{eq:c-selected-integral} by $e^{-(\kappa_1-1)N_1(Z\setminus Z_0')}$ times the supremum over the distinguished boundary $\{|\sigma_{\Delta}|= e^{\kappa_1}\ \forall\Delta\in\pi_1(Z)\setminus\pi_1(Z_0')\}$ (with $\sigma_{\Delta}=0$ off $Z$, $\sigma_{\Delta}=1$ on $Z_0'$) of the integral of the inner integrand's modulus, which lies in the closed polydisc of Lemma~\ref{lem:gaussian-payment}. There, the Mayer integrals $\int_0^1\dd t_Y$ have mass one; $\chi_b\leqslant 1$; by Lemma~\ref{lem:gaussian-payment}(iii) with $\gamma_2\leqslant \alpha_G$, $\prod_{b\in\mathsf P}(1-\chi_b)\leqslant e^{-a_{\mathrm{bad}}|\mathsf P|}e^{\frac{\alpha_G}{2}\| B\|^2}$; on the box $|\mathbf V_{\xi}(Y,B)e^{t_Y\mathbf V_{\xi}(Y,B)}|\leqslant \omega(g_0)e^{-\kappa d_0(Y)}e^{\omega(g_0)}\leqslant \omega_*e^{1/2}e^{-\kappa d_0(Y)}$ by Lemma~\ref{lem:c2-potential-size} and $\omega_*\leqslant 1/2$ (which follows from $\eta_*\leqslant 1/32$); and the integral of the absolute density against $e^{\alpha_G\| B\|^2/2}$ is at most $e^{N_1(Z)}$ by Lemma~\ref{lem:gaussian-payment}(ii).
\end{proof}

\begin{proof}[Proof of Lemma~\ref{lem:activity-bounds}]
Fix a connected final domain $Z$. Each selection $(\mathcal D,\mathsf P)$ in \eqref{eq:c-vacuum-activity} determines $Y_0=\bigcup\mathcal D$, the seed $Z_0$ with wall components $Z_{0,1},\ldots,Z_{0,k'}$, the components $Y_i$ of $Y_0$ (each inside one $Z_{0,j}$), $n_j= N_0(Z_{0,j}\setminus Y_0)$ and $n=\sum_jn_j$. By Lemma~\ref{lem:c2-selected-majorant} and $\omega_*e^{1/2}e^{-\kappa d_0(Y)}\leqslant \eta_*e^{-\kappa\ell_0(Y)}$, the sum over selections is organised as $\sum_{Y'\subseteq Z}\sum_{Z_0\, | \,Z_0'= Y'}\sum_{Y_0\subseteq Z_0}\sum_{\mathcal D\, | \,\bigcup\mathcal D= Y_0}\sum_{\mathsf P}$ over unions of cubes $Y'$ of $\pi_1$ and $Y_0$ of $\pi_0$, and each inner sum is bounded in turn.

\emph{Potential labels.} Split $\eta_*e^{-\kappa\ell_0(Y)}= t_0e^{-(1-\delta)\kappa\ell_0(Y)}\cdot\alpha_6e^{-\delta\kappa\ell_0(Y)}$. Each $Y\in\mathcal D$ lies in one component $Y_i$, and $\sum_{Y\subseteq Y_i}\ell_0(Y)\geqslant \ell_0(Y_i)$ by Lemma~\ref{lem:c2-domain-geometry}(ii), so with $t_0\leqslant 1$ and at least one label per component, $\prod_{Y\in\mathcal D}t_0e^{-(1-\delta)\kappa\ell_0(Y)}\leqslant \prod_it_0e^{-(1-\delta)\kappa\ell_0(Y_i)}$. Lemma~\ref{lem:c2-cover-sum} sums the second factors to at most one.

\emph{Bad-bond labels.} Given $Y_0$ and $Z_0$, every $b\in\mathsf P$ has $b\cap Y_0=\varnothing$ and both endpoints in cubes of $Z_0$, hence in $Z_0\setminus Y_0$; there are at most $4M^4n$ such bonds. Every cube of $Z_0\setminus Y_0$ meets a bond of $\mathsf P$ and a bond meets at most two cubes, so $|\mathsf P|\geqslant n/2$. Hence
\begin{equation}\label{eq:native-source-bad-label-sum}
    \sum_{\mathsf P}e^{-a_{\mathrm{bad}}|\mathsf P|/2}\leqslant e^{-a_{\mathrm{bad}}n/8}\sum_{\mathsf P}e^{-a_{\mathrm{bad}}|\mathsf P|/4}\leqslant \exp\Bigl[-\frac{a_{\mathrm{bad}}}{8}n+4M^4e^{-a_{\mathrm{bad}}/4}n\Bigr]\leqslant e^{-a_{\mathrm{bad}}n/16}\;,
\end{equation}
since $4M^4e^{-a_{\mathrm{bad}}/4}\leqslant a_{\mathrm{bad}}/16$: by \eqref{eq:native-source-bad-budget}, $a_{\mathrm{bad}}\geqslant 160(\kappa+1)>1$ and $64M^4e^{-a_{\mathrm{bad}}/4}\leqslant 1$ (as $a_{\mathrm{bad}}/4\geqslant 40\log(8M^4+8)\geqslant \log(64M^4)$). The unused half $e^{-a_{\mathrm{bad}}|\mathsf P|/2}\leqslant e^{-a_{\mathrm{bad}}n/4}$ leaves in total $e^{-5a_{\mathrm{bad}}n/16}=\prod_je^{-5a_{\mathrm{bad}}n_j/16}$, of which $e^{-a_{\mathrm{bad}}n_j/32}$, $e^{-a_{\mathrm{bad}}n_j/64}$ and $e^{-a_{\mathrm{bad}}n_j/64}$ are used below.

\emph{Component reserve.} For a component $Z_{0,j}$ containing $k_j\geqslant 1$ components $Y_i$ of $Y_0$, Lemma~\ref{lem:c2-domain-geometry}(iii) gives $\ell_0(Z_{0,j})\leqslant \sum_{i}\ell_0(Y_i)+5n_j$, so with $a_{\mathrm{bad}}/32\geqslant 5\kappa$,
\begin{equation}\label{eq:native-source-component-reserve}
    t_0^{k_j}e^{-(1-\delta)\kappa\sum_i\ell_0(Y_i)}e^{-a_{\mathrm{bad}}n_j/32}\leqslant t_0e^{-(1-\delta)\kappa\ell_0(Z_{0,j})}\;.
\end{equation}
If $k_j=0$ then $n_j\geqslant 1$, $\ell_0(Z_{0,j})\leqslant 5n_j$, and $e^{-a_{\mathrm{bad}}n_j/32}e^{-a_{\mathrm{bad}}n_j/64}\leqslant e^{-5\kappa n_j}\eta_*\leqslant t_0e^{-(1-\delta)\kappa\ell_0(Z_{0,j})}$, using $e^{-a_{\mathrm{bad}}/64}\leqslant \eta_*\leqslant t_0$.

\emph{Choice of $Y_0$.} With the reserve $e^{-a_{\mathrm{bad}}n/64}$, $\sum_{Y_0\subseteq Z_0}e^{-a_{\mathrm{bad}}N_0(Z_0\setminus Y_0)/64}=(1+e^{-a_{\mathrm{bad}}/64})^{N_0(Z_0)}\leqslant \exp(e^{-a_{\mathrm{bad}}/64}C_{\mathrm g}\sum_j\ell_0(Z_{0,j}))\leqslant \prod_je^{\delta\kappa\ell_0(Z_{0,j})}$, since $C_{\mathrm g}e^{-a_{\mathrm{bad}}/64}\leqslant 162e^{-5(\kappa+1)/2}\leqslant \kappa/40$ for $\kappa\geqslant 3$. Altogether, writing $\mathfrak M(Z_0,Z)$ for the sum of \eqref{eq:fixed-kernel-selected-majorant} over the selections with seed $Z_0$,
\begin{equation}\label{eq:native-source-initial-domain-sum}
    \mathfrak M(Z_0,Z)\leqslant e^{-(\kappa_1-1)N_1(Z\setminus Z_0')+N_1(Z)}\prod_{j}t_0e^{-(1-2\delta)\kappa\ell_0(Z_{0,j})}\;.
\end{equation}

\emph{Scale change.} By Lemma~\ref{lem:c2-scale}, the enlargement $Z_0'$ of Section~\ref{sec:fixed-first-kernel} (the $LM$-cubes meeting $\operatorname{int}Q^{\sim}$ for some cube $Q$ of $Z_0$) equals $\bigcup_jZ_{0,j}'$, where $Z_{0,j}'$ is the enlargement of the component $Z_{0,j}$ in the sense of that lemma, so \eqref{eq:c2-scale-inequality} applies to each $Z_{0,j}$. Each $Z_{0,j}'$ is connected, so every component $Y'_i$ of $Y'= Z_0'$ is the union of the $Z_{0,j}'$ it contains. Split $e^{-(1-2\delta)\kappa\ell_0(Z_{0,j})}= e^{-\delta\kappa\ell_0(Z_{0,j})}e^{-(1-3\delta)\kappa\ell_0(Z_{0,j})}$ and apply \eqref{eq:c2-scale-inequality} in the form $\ell_0(Z_{0,j})\geqslant \frac{L}{2049}\ell_1(Z_{0,j}')-\frac{5L+1024}{2049}$, with $(5L+1024)/2049\leqslant L/2$ for $L\geqslant 3$:
\begin{equation*}
    \prod_{j\, | \,Z_{0,j}'\subseteq Y'_i}e^{-(1-3\delta)\kappa\ell_0(Z_{0,j})}\leqslant e^{L\kappa/2\cdot k'_i}\prod_je^{-(1-3\delta)\frac{L\kappa}{2049}\ell_1(Z_{0,j}')}\leqslant e^{L\kappa k'_i/2}e^{-(1-3\delta)\frac{L\kappa}{2049}\ell_1(Y'_i)}\;,
\end{equation*}
by Lemma~\ref{lem:c2-domain-geometry}(ii), where $k'_i$ counts the fine components in $Y'_i$. The remaining factors $t_0e^{L\kappa/2}e^{-\delta\kappa\ell_0(Z_{0,j})}$ pay for the sum over the set of fine components inside $Y'_i$: each contains an $M$-cube of $Y'_i$, of which there are $L^4N_1(Y'_i)$, so by \eqref{eq:native-source-animal-sum} the sum over single components is $x_i\leqslant L^4N_1(Y'_i)t_0e^{L\kappa/2}B(\delta\kappa)\leqslant \varepsilon_*\frac{\delta L\kappa}{2049}\ell_1(Y'_i)$, using $N_1\leqslant C_{\mathrm g}\ell_1$ and the definitions of $C_{\mathrm s}$ and $\varepsilon_*= C_{\mathrm s}t_0$, and the sum over nonempty sets is at most $e^{x_i}-1\leqslant x_ie^{x_i}$. Write $\Lambda_i=\frac{L\kappa}{2049}\ell_1(Y'_i)$ and $y=\delta\Lambda_i$, so $x_i\leqslant \varepsilon_*y\leqslant y$. Then $x_ie^{x_i}e^{-(1-3\delta)\Lambda_i}\leqslant \varepsilon_*ye^{y}e^{-(1-3\delta)\Lambda_i}=\varepsilon_*ye^{-(1-4\delta)\Lambda_i}=\varepsilon_*(ye^{-y})e^{-(1-5\delta)\Lambda_i}\leqslant \varepsilon_*e^{-\lambda'\ell_1(Y'_i)}$, using $ye^{-y}\leqslant 1$, so
\begin{equation}\label{eq:native-source-enlarged-domain-sum}
    \sum_{Z_0\, | \,Z_0'= Y'}\mathfrak M(Z_0,Z)\leqslant e^{-(\kappa_1-1)N_1(Z\setminus Y')+N_1(Z)}\prod_{Y'_i\in\operatorname{cc}(Y')}\varepsilon_*e^{-\lambda'\ell_1(Y'_i)}\;.
\end{equation}

\emph{Choice of $Y'$.} For nonempty $Y'\subseteq Z$ with $n'= N_1(Z\setminus Y')$ cubes outside, Lemma~\ref{lem:c2-domain-geometry}(iii) gives $\ell_1(Z)\leqslant \sum_i\ell_1(Y'_i)+5n'$, so $\prod_ie^{-\lambda'\ell_1(Y'_i)}\leqslant e^{-\lambda'\ell_1(Z)}e^{5\lambda'n'}$, and $\varepsilon_*^{\#\operatorname{cc}(Y')}\leqslant \varepsilon_*$. With $\kappa_1-1\geqslant 5\lambda'+1$,
\begin{equation*}
    | k^{\mathrm f}_{\xi}(Z)|\leqslant \varepsilon_*e^{N_1(Z)}e^{-\lambda'\ell_1(Z)}\sum_{\varnothing\ne Y'\subseteq Z}e^{-N_1(Z\setminus Y')}
    \leqslant \varepsilon_*e^{(1+e^{-1})N_1(Z)}e^{-\lambda'\ell_1(Z)}\leqslant \varepsilon_*e^{-(\lambda'-2C_{\mathrm g})\ell_1(Z)}\;,
\end{equation*}
which is the first inequality of \eqref{eq:fixed-kernel-raw-activity-payment}.

\emph{Criterion.} Let $Z$ be any union of $LM$-cubes, connected or not (the lemma states the case of a connected final domain; the proof of Lemma~\ref{lem:vacuum-clusters} uses the general case). A domain $Z'$ incompatible with $Z$ contains one of the at most $9N_1(Z)$ cubes of $Z$ or wall adjacent to $Z$. Since $(1+6\theta)N_1(Z')\leqslant (1+6\theta)C_{\mathrm g}\ell_1(Z')\leqslant \frac{\lambda_*}{2}\ell_1(Z')$ by $C_{\mathrm g}\leqslant \lambda_*/8$ and $6\theta C_{\mathrm g}=3\lambda_*/8$, \eqref{eq:native-source-animal-sum} gives
\begin{equation}\label{eq:native-source-weighted-vacuum-criterion}
    \sum_{Z'\not\sim Z}| k^{\mathrm f}_{\xi}(Z')| e^{(1+6\theta)N_1(Z')}\leqslant 9N_1(Z)\varepsilon_*B(\lambda_*/2)\leqslant \frac18N_1(Z)\;.
\end{equation}
The reference data are treated in Section~\ref{app:reference-check}.
\end{proof}

\subsection{Proof of the vacuum-cluster bounds}
\label{app:kp}

\begin{proof}[Proof of Lemma~\ref{lem:vacuum-clusters}]
Species are the connected final domains $D$, with weights $\widetilde z_D=| k^{\mathrm f}_{\xi}(D)| e^{2\theta N_1(D)}$ and sizes $a(D)= N_1(D)/8$; two species are incompatible when they share a cube or a wall, and each with itself. By \eqref{eq:native-source-weighted-vacuum-criterion} and $2\theta+\frac18\leqslant 1+6\theta$,
\begin{equation*}
    \sum_{D'\not\sim D}\widetilde z_{D'}e^{a(D')}\leqslant \sum_{D'\not\sim D}| k^{\mathrm f}_{\xi}(D')| e^{(1+6\theta)N_1(D')}\leqslant a(D)\;.
\end{equation*}
Tuples of species are written $\mathfrak c=(D_1,\ldots,D_m)$ in this proof ($\gamma$ is the spectral bound $\gamma_L/2$). The argument of Section~\ref{sec:marginal-perimeter} following \eqref{eq:marginal-pinned-bound} applies verbatim to these species, weights and sizes: the absolute Ursell coefficient $|\varphi^T(\mathfrak c)|$ is bounded by the number of spanning trees of the incompatibility graph of $\mathfrak c$, the rooted-tree generating sum with root $D$ is the monotone limit of the recursion \eqref{eq:marginal-pinned-recursion} with $\widetilde z$ in place of $|\widetilde z_p|$, the displayed criterion gives $T^{(k)}_D\leqslant \widetilde z_De^{a(D)}$ inductively, and the factor $1/m!$ in \eqref{eq:part1-vacuum-cluster} together with the multiplicity $m_D(\mathfrak c)$ of $D$ in the tuple reproduces the rooted sum. Hence, for every connected final domain $D$,
\begin{equation}\label{eq:c2-pinned-vacuum}
    \sum_{\mathfrak c}m_D(\mathfrak c)|\Phi^{\mathrm v}_{\xi}(\mathfrak c)| e^{2\theta\sum_jN_1(D_j)}\leqslant | k^{\mathrm f}_{\xi}(D)| e^{2\theta N_1(D)+N_1(D)/8}\;.
\end{equation}
Here $\mathcal A$ is any finite set of unions of $LM$-cubes, connected or not; the criterion \eqref{eq:native-source-weighted-vacuum-criterion} was proved for every such union. A tuple hitting $\mathcal A$ contains a member $D'$ incompatible with some $D\in\mathcal A$, so
\begin{equation*}
    \sum_{\mathfrak c\text{ hits }\mathcal A}|\Phi^{\mathrm v}_{\xi}(\mathfrak c)| e^{2\theta\sum_jN_1(D_j)}
    \leqslant \sum_{D\in\mathcal A}\sum_{D'\not\sim D}\sum_{\mathfrak c}m_{D'}(\mathfrak c)|\Phi^{\mathrm v}_{\xi}(\mathfrak c)| e^{2\theta\sum_jN_1(D_j)}
    \leqslant \sum_{D\in\mathcal A}\sum_{D'\not\sim D}| k^{\mathrm f}_{\xi}(D')| e^{(2\theta+\frac18)N_1(D')}\leqslant \frac18\sum_{D\in\mathcal A}N_1(D)\;.
\end{equation*}
For \eqref{eq:part1-joint-vacuum-exponential}, at fixed volume there are finitely many species. Let $\mathcal S_+$ be the set of species with $k^{\mathrm f}_{\xi}(D)\ne0$; species with $k^{\mathrm f}_{\xi}(D)=0$ contribute to neither side (a tuple containing one has $\Phi^{\mathrm v}_{\xi}=0$, and the gas does not see them), so they are discarded. For complex activities $z=(z_D)_{D\in\mathcal S_+}$, the hard-core gas $\Xi(z)=\sum_{\mathcal P\text{ compatible}}\prod_{D\in\mathcal P}z_D$ is a polynomial with $\Xi(0)=1$: the polymer partition function (5.3) of \cite[Chapter~5]{FriedliVelenik} for the polymers $\mathcal S_+$ with the hard-core interaction $\zeta(D,D')=0$ on incompatible and $\zeta(D,D')=1$ on compatible pairs (which satisfies (5.1)--(5.2) there). Let $S(z)=\sum_{\mathfrak c}\frac{\varphi^T(\mathfrak c)}{m!}\prod_jz_{D_j}$, the sum over all ordered tuples $\mathfrak c=(D_1,\ldots,D_m)$, $m\geqslant 1$, of members of $\mathcal S_+$; this is the cluster series of that gas, since the Ursell function of \cite[(5.6)]{FriedliVelenik} is $\frac1{m!}\sum_G\prod_{\{ i,j\}\in G}(\zeta(D_i,D_j)-1)=\frac{\varphi^T(\mathfrak c)}{m!}$, the sum over the connected graphs $G$ on $\{1,\ldots,m\}$, a factor $\zeta-1$ being $-1$ on an incompatible pair and $0$ otherwise. \cite[Proposition~5.3]{FriedliVelenik} is the identity $\Xi(z)=\exp(S(z))$ between formal power series in the variables $z_D$ (the book states it as a formal identity; its proof, (5.4)--(5.8) there, rearranges finitely many terms for each monomial). By \eqref{eq:c2-pinned-vacuum} with $z_D$ in place of $k^{\mathrm f}_{\xi}(D)$ (the bound is monotone in the weights and only uses $| z_D|\leqslant | k^{\mathrm f}_{\xi}(D)|$ through the criterion; equivalently \cite[Theorem~5.4]{FriedliVelenik}, whose hypothesis (5.10) is the first display of this proof and whose conclusion (5.9) is this absolute convergence), $S$ converges absolutely at $| z_D|=| k^{\mathrm f}_{\xi}(D)|$, hence absolutely and uniformly on the closed polydisc $\overline P=\{| z_D|\leqslant | k^{\mathrm f}_{\xi}(D)|\ \forall D\in\mathcal S_+\}$ (Weierstrass M-test), so $S$ is continuous on $\overline P$, analytic on the open polydisc $P$, and its Taylor series at $0$ is the formal series $S$ grouped by monomials. The analytic functions $\Xi$ and $\exp\circ S$ on the connected open set $P\ni0$ have the same Taylor series at $0$, by the formal identity, so they coincide on $P$ (identity theorem) and, by continuity, on $\overline P$; in particular at $z_D= k^{\mathrm f}_{\xi}(D)$, where $\Xi(k^{\mathrm f}_{\xi})=\Xi^{\mathrm f}_{\xi}$ and $S(k^{\mathrm f}_{\xi})=\sum_{\mathfrak c}\Phi^{\mathrm v}_{\xi}(\mathfrak c)$ by \eqref{eq:part1-vacuum-cluster}. The restricted gas $(\Xi^{\mathrm f}_{\xi})^{[\cup\mathcal A]}$ is the same gas with $z_D=0$ for the $D$ incompatible with $\bigcup\mathcal A$, which are exactly the $D$ incompatible with some member of $\mathcal A$; its expansion is the sum over the tuples none of whose members hits $\mathcal A$. The quotient is therefore $\exp(-\sum_{\mathfrak c\text{ hits }\mathcal A}\Phi^{\mathrm v}_{\xi}(\mathfrak c))$.
\end{proof}

\subsection{The reference data}
\label{app:reference-check}

The reference data of Section~\ref{sec:joint-extraction} consist of the identity background $U_1^{\circ}=\id$ on every bond (the constrained minimiser for identity averages, with zero action and zero gauge functional), $V_0^{\circ}=\id$, zero retained input $A^{\circ}_{\mathrm e}=0$ with $\Psi^{\circ}_{\xi}=\id$, and the same $I$, $\Omega_1$, partitions, hosts, cutoff radius $r$ and constants as the underlying branch. They are not a branch, so the conclusions of the lemmas are checked directly. The reference Hessian is the regional Hessian \eqref{eq:c-regional-hessian} at the identity background on the branch's $\Omega_1$, $H^{\circ}_{\xi}=\mathsf H^{\Omega_1}_{\id}$, and $H^{\circ}$ denotes its restriction to $I$. Since $I$, $\Omega_1$ and $R_{\mathrm c}$ are those of the branch, the geometric facts used in Lemma~\ref{lem:c2-regional-entry} (for $b\in I$ every site of a stencil $\mathcal S_{bb'}$ lies within $3L$ of $b_-$, hence within $3R_{\mathrm c}$ of $I$ and in $\Omega_1^{\sim-1}$, by Lemmas~\ref{lem:c2-entry-range} and~\ref{lem:collar-coverage}(i)) hold for the reference data as well; Lemma~\ref{lem:c2-regional-entry} is stated for every background on which the regional form is defined, so it applies to the identity background without change.

\begin{lemma}[reference data]\label{lem:c2-reference}
The reference precision $H^{\circ}$ satisfies \eqref{eq:c-regular-two-sided}, as do its principal restrictions, and has the range and gauge law of Lemma~\ref{lem:hessian-gauge-locality}; Lemma~\ref{lem:charts} holds for the reference data (at $U=\id$, $A_{\mathrm e}=0$); Lemma~\ref{lem:gaussian-payment} holds for $H^{\circ}$ with the same $\gamma,h_+,R,M,\alpha_G$; and \eqref{eq:fixed-kernel-raw-activity-payment} and Lemma~\ref{lem:vacuum-clusters} hold for the reference activities $k^{\mathrm f,\circ}_{\xi}$ with the same constants. The reference row solutions, phases and primitives are real, and the reference normalisations $K^{\circ}_{\xi;0}$, $I^{\circ}_{\xi;0}$ of Section~\ref{sec:joint-extraction} are positive.
\end{lemma}

\begin{proof}
By Lemma~\ref{lem:c2-regional-entry} with the identity background, the entries of $H^{\circ}$ with a label in $I$ are $f_{bb'}(0)$, the entries of the flat Hessian $\mathsf H$ of \eqref{eq:flat-covariance} on the branch's torus. The compression of $\mathsf H\geqslant \gamma_LI$ \eqref{eq:covariance-uniform-gap} to $I$, and further to any subset of $I$, gives the lower bound $\gamma_L\geqslant \gamma_L/2$ for $H^{\circ}$ and for every principal restriction, and the Schur bound with row sums $\sum_{b'}\| f_{bb'}(0)\|\leqslant m_L$ gives the upper bound for $H^{\circ}$ and every principal restriction (a principal restriction has smaller row sums); the range is $\rho_L$ by Lemma~\ref{lem:c2-entry-range}. The gauge law of Lemma~\ref{lem:hessian-gauge-locality} is the entrywise covariance \eqref{eq:c-entry-gauge-covariance}, proved there for every background, hence for $\id$: $f_{bb'}(\id^u)=\mathrm{Ad}_{u(b_-)}f_{bb'}(\id)\mathrm{Ad}^{-1}_{u(b'_-)}$. Lemma~\ref{lem:gaussian-payment} and Lemmas~\ref{lem:walk-series}, \ref{lem:conditioning}, \ref{lem:component-blocks} use only \eqref{eq:fixed-kernel-inputs}, the range and the partitions, so they hold for $H^{\circ}$. Lemma~\ref{lem:c2-uniform-chart} applies to the identity background (its hypothesis, a site gauge with phases at most $\tau_L/2$ on the two-block sets and blocks concerned, is satisfied by $u=\id$, $\phi=0$; the geometric facts of the applicability paragraph depend only on $I$, $\Omega_1$, $R_{\mathrm c}$), so Lemma~\ref{lem:charts} holds at $U=\id$, and Lemma~\ref{lem:c2-primitive-bounds} holds with $r_{\mathrm e}^{\circ}=0$ and $s^{\circ}=1+r\leqslant s_{\xi}$; the hosts and the host multiplicity are those of Lemma~\ref{lem:c2-potential-size}, so \eqref{eq:native-source-interaction-size} holds for $\mathbf V^{\circ}_{\xi}$. The proofs of Lemmas~\ref{lem:c2-selected-majorant}, \ref{lem:activity-bounds} and~\ref{lem:vacuum-clusters} use nothing else about the branch. Reality: the row solutions are the limits of the contraction iterates of Lemma~\ref{lem:c2-uniform-chart}(ii), which are real for real phases, so for real inputs the phases, the Haar logarithms, the determinant logarithms (positive determinants) and the plaquette and gauge terms are real, and so are the primitives and $\mathbf V^{\circ}_{\xi}$. Positivity: $K^{\circ}_{\xi;0}=\mathbf 1_{\mathcal D_I}\exp(-\frac12\langle a,H^{\circ} a\rangle+P^{\circ}_{\xi,0}(a))$, the function \eqref{eq:native-source-boundary-interpolation} for the reference data at $t=0$, is a real exponential on the open box, the reference primitives being real, hence positive there; and $I^{\circ}_{\xi;0}=\int K^{\circ}_{\xi;0}\,\dd a$ is positive and finite ($H^{\circ}\geqslant \gamma$ and the exponent is bounded on the box by Lemma~\ref{lem:charts}).
\end{proof}

\subsection{The thresholds}
\label{app:thresholds}

Before the coupling, in this order: $L>11$ odd (RG1 p.251, the range in which Ba{\l}aban's results are stated); the group constants $c_{\mathrm h}=\sqrt N$, $c_{\mathrm n}=1/N$, $\varrho_G=\pi/\sqrt N$, $c$, $r_G$, $c_G$, and $r_M,r'_M$ (Lemma~\ref{lem:c2-uniform-chart}(iv); these two depend on $L$ and the group); $\tau_L$ by \eqref{eq:c2-tau-final}, with $m_L$, $c_L$ (Lemma~\ref{lem:hessian-gauge-locality}) and $\gamma_L$ (Appendix~\ref{app:flat-covariance}); $R_{\mathrm{loc}}=2L+2$, $R_{\mathrm p}$, $s_{\mathfrak P}$, $R_{\mathrm c}\geqslant \max(2R_{\mathrm{loc}}+1,s_{\mathfrak P}+2,4L+5)$ and $C_{L,R_{\mathrm c}}$; the two-block size $n_{\mathcal E}$ of Lemma~\ref{lem:c2-fibre-smallness} (a constant of $L$); the ranges $R_*$ and $R_H$ of Section~\ref{sec:marginal-perimeter} (the perimeter range bounding the primitive-label diameters and the source-to-stencil distances, distinct from the walk range $R_*= n_*R$ below, and the range of the flat Hessian, whose entries couple the bonds of a common plaquette or block; both are constants of $L$), hence $R_0=4(1+R_*+R_H+L)$; $\vartheta_*\leqslant \tau_L/2$, $\rho_{\mathrm d}=2K_*\vartheta_*^2$, $C_{\mathsf C}$, $C_{\mathsf D}$, $C^{\mathrm a}_k$ ($0\leqslant k\leqslant 5$) and $C^{\mathsf D}_k$ ($0\leqslant k\leqslant 6$) of Lemma~\ref{lem:c2-uniform-chart} and $C_{\mathrm p}$ of Lemma~\ref{lem:c2-primitive-bounds}; $\gamma=\gamma_L/2$, $h_+= m_L$, $R$; $C_{\mathrm g}$, $\delta$, $\kappa$, $\lambda'$, $\lambda_*$, $\alpha_6$, $\kappa_1$, $\theta$, $C_{\mathrm s}$ by \eqref{eq:c2-kappa-choice}--\eqref{eq:c2-mayer-constants}, then $\omega_*$, $\eta_*$, $\varepsilon_*$, $t_0$ by \eqref{eq:native-source-vacuum-KP-budget}, with $\omega_*\leqslant 1/2$ (the standing choice of the collar estimates with a Wilson insertion, Section~\ref{sec:one-step-insertion}) (these involve only $\kappa$, $\alpha_6$, $C_{\mathrm s}$, $\lambda_*$, so they are fixed before $M_1$ and $M$, as Lemma~\ref{lem:activity-bounds} states); $n_*$, $R_*$, $C_{\mathrm w}$ (Section~\ref{sec:fixed-first-kernel}), $A_*$, $C_{\mathrm d}$ (Lemma~\ref{lem:gaussian-payment}); $M_1$ as CRE's $M_1$ (p.245), $M_2$ and $M$ odd powers of $L$ with $M_1< M_2< M$, $M/M_1$ integral, $M\geqslant 8LM_2$, CRE's requirement that $M$ be large (p.245), $\rho_{\mathrm w}= R_*+R+1+2M_1$, \eqref{eq:fixed-kernel-collar-size}, $\delta_M\leqslant 1$, $h_+^2\delta_M^2/\gamma\leqslant 1/8$ and $4C_{\mathrm d}\delta_M^2(LM+1)^4\leqslant 1/2$; $\alpha_G$ with $0<\alpha_G\leqslant \min(\gamma/2,\gamma^2/(16A_*^2),(8C_{\mathrm d}(LM+1)^4)^{-1})$, $\gamma_2=\alpha_G/2$; then $N_{\mathrm h}$ and $C_0$ of \eqref{eq:part1-host-multiplicity}, whose hosts depend on the chosen $M$; then the constants of the collar estimates with a Wilson insertion (Section~\ref{sec:one-step-insertion}, where they are introduced), in this order: the finite-difference constant $C'_{\chi}=2N_{\mathrm{tr}}C_{\mathrm{link}}$ and the character constant $K_{\chi}$, the dummy weight $u_*$ of the host comparison, the support reserve $\eta$ of the positive group summation, its pin constant $C_{\mathrm{pin}}$ and then the pin radius $\rho_{\mathrm{pin}}$ (defined through $C_{\mathrm{pin}}$ and $\eta$), the host count $N_{\mathrm{host}}=4^4\cdot20^4$ and the probe-host bond count $C_M=8(4M+1)^4$, all fixed before the coupling, each depending on $L$, $M$, $\kappa$, $\omega_*$, the group, the character and the constants earlier in this order only. The coupling $g_0$ is then taken below a threshold, and the torus is chosen last: every fine period a multiple of $LMR_1$ and at least $8M$ (\eqref{eq:c-standing-conditions}; the partitions of Section~\ref{sec:fixed-first-kernel} need a multiple of $LM$, the hosts of Section~\ref{sec:source-phase-primitives} need $8M$, and CRE's partitions into $LMR_1$-cubes need the multiple of $LMR_1$; $R_1= R_1(g_0)$, so the period condition is imposed after the coupling) and every coarse period exceeding $\max(2,2R_{\mathrm p}+8)$. The statements about a square $C_n$ of side $nL$ (Proposition~\ref{prop:regional-square-extraction}, Appendices~\ref{app:signed-defects} and~\ref{app:regional-scalar}, Sections~\ref{sec:c5-3}--\ref{sec:c5-6}) assume in addition the square/torus geometry of Section~\ref{sec:extent-weights}: every coarse period $K_{\mu}\geqslant 3(n+2)$, so that the square and its fixed neighbourhood have unique lifts and the loop is contractible; and Appendix~\ref{app:exterior-chessboard} (Section~\ref{sec:c3-7}) and Sections~\ref{sec:c5-3}--\ref{sec:c5-6} assume compatible even fine periods, exceeding $200\ell$ with $\ell= LMR_1$ where stated (Section~\ref{sec:c5-3-1}). These are lower bounds on the periods in terms of the fixed loop and the coupling; for a given loop and coupling they hold on every sufficiently large compatible torus, and nothing is let tend to infinity. In the threshold conditions, $\varepsilon_1= g_1p_0(g_1)$ with $g_1= g_1(g_0)$ by CRE (1.27); $R_1$ is CRE's, the smallest power of $L$ with $R_1\geqslant (\log g_1^{-2})^{\mathfrak r}$ (p.246, with CRE's exponent written $\mathfrak r$; $\mathfrak r\geqslant 2$ at p.245, $p_0\geqslant 5\mathfrak r$ at p.246, and CRE (2.5) for the general step); and the cutoff radius is $r= A_1(\log g_0^{-2})^{p_0}$, the radius of the literal cutoff \eqref{eq:literal-cutoff}; it is $\delta_0/g_0$ for CRE's $\delta_k= g_kA_1/A_0\,p_0(g_k)$ as printed at p.265 (of the equality chain printed for $\delta_0$ at p.247 only the first member differs from this, its last member $(A_1/A_0)\varepsilon_0$ agrees; the reading $r=\delta_0/g_0$ identifies $\Theta_{\xi}$ with the test (1.18) and is used throughout, while below only $r\to\infty$ and $g_0(1+r)^3\to0$ enter, which hold for either reading, except item 4(vi), which is stated for the explicit form $r= A_1(\log g_0^{-2})^{p_0}$). The threshold conditions are:
\begin{enumerate}
\item $LMR_1>4R_{\mathrm c}+4L^2M_2R_1+8$ (Lemma~\ref{lem:collar-coverage}), $LM_2R_1\geqslant M$ (so that the $LM$-cubes nest in the $L^2M_2R_1$-cubes, \eqref{eq:c-standing-conditions}) and $c_{\mathrm h}O(1)B_3e^{-\delta_{\mathrm c}LM_2R_1}\leqslant 1/4$ (Lemma~\ref{lem:c2-plaquette-deviation}); both depend on the coupling through $R_1$;
\item \eqref{eq:convex-background-thresholds}, the hypotheses of Lemma~\ref{lem:regular-bounds}, which \eqref{eq:c-standing-conditions} lists as well;
\item $4C_{\mathsf C}g_0(1+2r)\leqslant \vartheta_*/2$ (Lemma~\ref{lem:c2-primitive-bounds});
\item the collar thresholds, with $a_{\mathrm{bad}}= c_{\mathrm n}\alpha_Gr^2/4$, $u_t\coloneqq C'_{\chi} g_0(1+2r)$, $\rho_t\coloneqq u_t/u_*$ and $\omega_t\coloneqq C_0g_0(1+2r)^3$: (i) \eqref{eq:native-source-bad-budget} for $(9/128)a_{\mathrm{bad}}$ in the role of $a_{\mathrm{bad}}$, that is $(9/128)a_{\mathrm{bad}}\geqslant \max(160(\kappa+\log(8M^4+8)+1),64\log(1/\eta_*))$, which implies the budget for $a_{\mathrm{bad}}$ and for $(9/16)a_{\mathrm{bad}}$; (ii) $\rho_t+u_t/2+2\omega_t/\omega_*+2C_Me^{-(9/32)a_{\mathrm{bad}}}\leqslant 1$, which implies $u_t\leqslant u_*$, $\omega_t\leqslant \omega_*/2$ and the former condition $C_0g_0(1+r)^3\leqslant \omega_*$; (iii) $C_0g_0(1+3r)^3\leqslant \min(1,\omega_*e^{-\kappa(B_{\mathrm h}+4)}/(e^{9/8}N_{\mathrm{host}}))$; (iv) $\rho_t\leqslant \rho_{\mathrm{pin}}$; (v) $2K_{\chi} C_{\mathrm{pin}}\rho_t\leqslant 1$; (vi) $2C_M\leqslant (1+2r)^3$, $\log g_0^{-2}\geqslant 1$ and $(9c_{\mathrm n}\alpha_G/64)A_1^2\log g_0^{-2}\geqslant 1$, under which the second-order family-sum rate of the collar estimates takes its explicit form (these are the conditions (4.T), (4.T$'$), (5.T)(i)--(v) and (R.T)(i)--(iii) of the collar estimates, Section~\ref{sec:one-step-insertion}, in reduced form);
\item CRE's own conditions, taken over as they stand: $\varepsilon_0$ small (Theorem 1 of \cite{BalabanVAR}, CRE's [15], behind (B1); and, the undisplayed constants $O(L^2)$ of p.247 and $O(1)$ of \cite[p.~25]{BalabanAVG} being independent of the coupling by the reading of (B1), so small that $2\,O(L^2)\varepsilon_0$ and $O(1)L^2\varepsilon_0$ lie below $\varrho_G$, $r_G/\sqrt N$ and the small-diameter radius of the group mean, the domain statement of (B1)); $MR_1$ large (p.251, the bound on $\mathbf H_{1,Ax}$ at the tolerance $\delta_0/40$ of \eqref{eq:boundary-bijection-inputs}, behind the change of variables (1.22), Lemma~\ref{lem:boundary-bijection}; the $O(1)$, $B_3$ and $\delta$ of that bound are taken to be independent of the coupling and of the branch, the reading of (B3)(iii) extended to this bound and named in (H4), without which the condition is not shown satisfiable as $g_0\to0$); $A_0/A_1$ and $M_2R_1$ large (the two qualitative conditions of (B3), pp.248 and 250; the second requires $O(1)B_3e^{-\delta_{\mathrm c}LM_2R_1}\varepsilon_1$ to be much smaller than $\delta_0$, and $\varepsilon_1/\delta_0\leqslant (1+\beta_0)A_0/A_1$ stays bounded by CRE (2.8) and $\delta_0=(A_1/A_0)\varepsilon_0$, so $R_1\to\infty$ at fixed $L$, $M_2$ suffices, the constants being independent of the coupling by (B3)(iii)). The conditions on $MR_1$ and $M_2R_1$ depend on the coupling through $R_1$ and are met for small $g_0$ at fixed $M$, $M_2$, since $R_1\to\infty$; the condition on $A_0/A_1$ is met by the choice before the coupling; CRE's conditions on $M_1$, $M_2$, $M$ themselves are not listed by CRE (p.245: $M_1$ the cube size of the previous papers, $M_2$ not too much larger, $M$ much larger than $M_2$, ``dictated by many conditions involving this constant'') and are taken over as the hypothesis that CRE's construction is valid for the values chosen above, which follow CRE's prescription.
\item the fibre conditions \eqref{eq:c2-fibre-thresholds} of Lemma~\ref{lem:c2-fibre-smallness}: $3c_Gc_{\mathrm h}s_{\mathfrak P}\varepsilon_0\leqslant \tau_L/2$ (which gives $3c_{\mathrm h}s_{\mathfrak P}\varepsilon_0\leqslant r_G$, by $\tau_L\leqslant 2c_Gr_G$) and $n_{\mathcal E}^{1/2}c_Gc_{\mathrm h}\,2\,O(L^2)\varepsilon_0\leqslant \vartheta_*/4$, with $n_{\mathcal E}$ the largest number of bonds of a two-block set and $O(L^2)$ the constant of p.247 of item 5.
\end{enumerate}
All hold for small $g_0$; for item 4(i)--(v) because $a_{\mathrm{bad}}\to\infty$ and $g_0(1+r)^3\to0$ (so every term of (ii) tends to zero, the fourth because $a_{\mathrm{bad}}\to\infty$), for item 4(vi) because $r\to\infty$ and $\log g_0^{-2}\to\infty$, and for item 6 because $\varepsilon_0= g_0A_0(\log g_0^{-2})^{p_0}\to0$ while its other constants are fixed before the coupling, $O(L^2)$ being independent of the coupling and of the branch by the reading of (B1). Under (vi), with $c=9c_{\mathrm n}\alpha_G/128$ and $p_0\geqslant 1$, $e^{-cr^2}\leqslant e^{-cA_1^2(\log g_0^{-2})^2}= g_0^{2cA_1^2\log g_0^{-2}}\leqslant g_0$, so the fourth term of (ii), $2C_Me^{-(9/32)a_{\mathrm{bad}}}=2C_Me^{-cr^2}$, is at most $g_0(1+2r)^3$: this is the estimate behind the explicit second-order rate of the collar estimates, not a condition needed for the satisfiability of (ii). The first-step coupling is defined by CRE (1.27), $g_0^{-2}= g_1^{-2}+\beta_1(g_0)$, and CRE (2.6) (p.255, with $n=1$, $m=0$) states $g_1\leqslant (1+\beta_0)g_0$ and $g_0\leqslant (1+\beta_0)g_1$, where $\beta_0>0$ is a constant that "can be chosen arbitrarily small, if $g$ is sufficiently small" (p.255); CRE derives (2.6) from the renormalisation group equations (0.20) of [I] and the properties of the $\beta$-functions (p.255), and it is taken over as CRE's statement, not derived here. Hence $g_1\to0$ with $g_0$, so $R_1\to\infty$; and CRE (2.8) (p.256, with $n=1$, $m=0$) gives $\varepsilon_1\leqslant (1+\beta_0)\varepsilon_0$, so $\varepsilon_1= g_1A_0(\log g_1^{-2})^{p_0}\to0$ with $\varepsilon_0= g_0A_0(\log g_0^{-2})^{p_0}$; while $r\to\infty$ and $g_0(1+r)^3\to0$. (CRE's $p_0$ denotes both the function $p_0(g)= A_0(\log g^{-2})^{p_0}$ and the integer exponent in it; the notation is CRE's and is kept.) The threshold depends on $L$, the group (through $c_{\mathrm h}$, $c_{\mathrm n}$, $\varrho_G$, $c$, $r_G$, $c_G$, $r_M$, $d_G$), the choices $\tau_L$, $\vartheta_*$, $R_{\mathrm c}$, $\kappa$, $\alpha_G$, $\omega_*$ made above, $M_1$, $M_2$, $M$, the collar constants $C'_{\chi}$, $K_{\chi}$, $u_*$, $\rho_{\mathrm{pin}}$, $C_{\mathrm{pin}}$, $\eta$, $N_{\mathrm{host}}$, $C_M$, and, through $\varepsilon_1$, $R_1$ and $r$, on CRE's constants $A_0$, $A_1$, $p_0$, $\mathfrak r$, the map $g_0\mapsto g_1$ of CRE (1.27), the constants $B_3$, $\delta_{\mathrm c}$ and the $O(1)$ of the comparison after (1.19) (p.268; CRE states no dependence for them, Lemma~\ref{lem:c2-plaquette-deviation} uses only their independence from the coupling, the reading (B3)(iii), and the threshold is stated with whatever dependence on $L$, the group and $M_2$ they have), the threshold $\varepsilon_0$ of Theorem 1 of \cite{BalabanVAR} (CRE's [15]) behind (B1), the constant $O(L^2)$ of p.247 (items 5 and 6), the constants of the bound on $H_{1,\mathrm{Ax}}$ of p.251 behind (B4)(ii), and the tolerances ``much smaller than $\delta_0$'' of p.250 behind (B3)(ii).

Further conditions are imposed where the appendices and Sections~\ref{sec:c5-2}--\ref{sec:c5-7} state them; each is a condition that $g_0$ lie below a threshold depending on $L$, the group, the character, $M_1$, $M_2$, $M$, the constants chosen above and CRE's constants (and on the parameter $\vartheta$ where indicated), holding for all small $g_0$, or a lower bound on the periods in terms of the loop and the coupling, holding on every sufficiently large compatible torus. They are: $\delta_0\leqslant \delta_G$ (Lemma~\ref{lem:boundary-bijection}); \eqref{eq:c3-1-2} ($r\geqslant 1$ and $C_Pg_0(1+2r)\leqslant \gamma/2$, Section~\ref{sec:c3-1}); the threshold $g_*$ of Proposition~\ref{prop:c3-7-5}, which contains the threshold of Proposition~\ref{prop:c3-9-1}, the coupling conditions of CRE (2.6) and the smallness giving $q\leqslant e^{-c_*\pi_*^2}$ (Section~\ref{sec:c3-7}); the thresholds of Lemma~\ref{lem:c4-2-1} (its last inequality), the smallness conditions of Lemma~\ref{lem:c4-6-2} on the root weight and the smallness of Lemma~\ref{lem:c4-6-3} (Appendix~\ref{app:signed-defects}, Proposition~\ref{prop:c4-12-3}); the further threshold of \eqref{eq:c5-7-expbound-I} (Corollary~\ref{cor:c5-7-4}(c)); the thresholds of Lemmas~\ref{lem:c3-7-seed}, \ref{lem:c3-7-2} and~\ref{lem:c3-7-4} (Section~\ref{sec:c3-7}; every condition there is an upper bound on the witness threshold $\theta$ of \eqref{eq:c3-7-4} or a lower bound on $R_1$); the thresholds of Lemmas~\ref{lem:c5-2-3}(b) and~\ref{lem:c5-2-4}, the threshold $g_*(\vartheta)$ of Proposition~\ref{prop:c5-2-6} and the smallness required for \eqref{eq:c5-2-3} (Section~\ref{sec:c5-2}); \eqref{eq:c5-4-T}, \eqref{eq:c5-4-Tp}, \eqref{eq:c5-5-T} and \eqref{eq:c5-5R-T}, which are the collar conditions of item 4 in the forms used there; the period hypothesis of Section~\ref{sec:c5-3-1} (compatible even fine periods exceeding $200\ell$, $\ell= LMR_1$); and \eqref{eq:c5-7-T} (the collar condition with $R_0=4(1+R_*+R_H+L)$ in place of $R_{\mathrm c}$, the small-coupling conditions $g_0C_{\mathsf C}(2+4r)<\vartheta_*/2$ and $g_0(2+4r)< a_L$ with the polymer-load inequalities of the proof of Proposition~\ref{prop:regional-square-extraction}, and the torus conditions of its item (iii)); the coupling conditions of the flat and regional expansions where they are imposed ($\tau<1$, the original chart condition, the smallness of $K_1^{s_*}\tau$ and the straight- and bent-path thresholds of Sections~\ref{sec:marginal-perimeter} and~\ref{sec:marginal-corner}, and the threshold of Proposition~\ref{prop:regional-coefficient}); $R_{\Upsilon}\geqslant 2$, $u\leqslant u_*$, $v\leqslant 1$ and the small-coupling conditions of Sections~\ref{sec:c4-10-1} and~\ref{sec:c4-11-5} (Proposition~\ref{prop:c4-12-3}); the further thresholds of the closing forms of Proposition~\ref{prop:c5-6-4} and of Corollary~\ref{cor:c5-7-4}(a), (c); and the further threshold of Corollary~\ref{cor:main-factorisation} ($C\tau^2L^4C_K\Phi_c\leqslant 1$ and the bound on $\delta^{\sharp}_q$), each stated and discussed at the place cited. The list of conditions on the parameters is the enumeration above together with these.

\section{The Local Interaction at the Identity}
\label{sec:perimeter-attempt}

\paragraph{Locality of the first-step interaction.}
The trace in the bare action is $\tr_N= N^{-1}\operatorname{Tr}$ in its fixed unitary embedding, so $\tr_N(\id)=1$. Choose a radius $a_{\mathrm{ph}}$ for the input $A$ of $\widetilde D$, measured in the maximum fine-bond Hilbert--Schmidt norm, on which the local implicit solutions, principal logarithms and composed action and Haar primitives are analytic, with nonzero Jacobian determinants and Haar densities. The radius $a_L$ is a free-input radius: take $a_L\leqslant a_{\mathrm{ph}}/\max(1,\|\mathsf C\|_{\infty\to\infty})$, where the operator norm maps the maximum free-bond Hilbert--Schmidt norm to the maximum fine-bond Hilbert--Schmidt norm. This norm is finite and depends on $L$ only, by the fixed row stencils. The Lie-algebra norm is the normalised Hilbert--Schmidt norm, and $| A|_{\mathrm{HS}}\leqslant | A|_{\mathrm{op}}$. Thus $c_{\mathrm{cmp}}=1$ is a common norm-comparison constant for both cutoff prescriptions (it is not the group constant $c_G$ of Appendix~\ref{app:background-estimates}), and $g_0r(g_0)< a_L$ suffices. The cutoff sets themselves are unchanged. The componentwise implicit solutions assemble to a small solution in this maximum norm. Uniqueness of the analytic germ in RG1 identifies it with $\widetilde D$ on the smaller common chart.
The following calculation concerns exactly the identity-background, all-regular integral below. It does not transfer that integral to a nonflat or inherited branch. Let $\mathcal S_c$ contain the positive fine bonds with both endpoints in $\mathcal B_{c_-}\cup\mathcal B_{c_+}$, and put $\Lambda_c=\mathcal S_c\cap E_*$. At the present torus sizes, the only designated bond in $\mathcal S_c$ is $b_0(c)$. Indeed, an internal path in either block crosses no block interface, and the translated length-$L$ path in RG1 (0.4) crosses only the interface between these two blocks. The reference coarse path crosses the same interface. Thus $\mathcal Q_c$ and $\mathcal N_c$ depend only on the fine bonds in $\mathcal S_c$. In the defining equation for $\widetilde D_c(A)$, the only unknown correction entering those bonds is $\widetilde D_c(A)$ itself. The analytic implicit-function theorem therefore gives a function of $A|_{\mathcal S_c}$ alone, with its radius and constants independent of the number of blocks. Here $h_c= L^3I$ at the identity, since $Q_0= L^{-4}I$.

For a free fine bond $b$, set $\Lambda_b=\{ b\}$. For a designated bond $b= b_0(c)$, set $\Lambda_b=\Lambda_c$. Both $(\mathsf CB)(b)$ and $q_{g_0}(B)(b)$ depend only on $B|_{\Lambda_b}$. The index $\alpha$ runs over the plaquettes and the pairs $(y,x)$ in \eqref{eq:action-primitive}. Define $f_{\alpha}(q)$ to be the corresponding summand of $\mathscr A(e^{iq})$, as a real analytic scalar function near zero in $\mathfrak g^E$; use its unique holomorphic extension for complex arguments. Let $\Lambda_{\alpha}$ be the union of the sets $\Lambda_b$ over the bonds in that summand's plaquette or path family. Each $\Lambda_{\alpha}$ lies in at most four adjacent blocks. Each free bond belongs to a bounded number of the sets $\Lambda_{\alpha},\Lambda_b,\Lambda_c$, with a bound depending on $L$, not on the torus.

Write $\psi(A)= A-h\widetilde D(A)$. Define the following local scalar functions by their real analytic branches at zero:
\begin{equation}\label{eq:perimeter-literal-local-terms}
    \begin{aligned}
        v_{\alpha}(B)&=-\frac{1}{g_0^2}
            \left(f_{\alpha}(\psi(g_0\mathsf CB))
                -\frac{g_0^2}{2}
                    D^2f_{\alpha}(0)[\mathsf CB,\mathsf CB]\right)\;,\\
        j_b(B)&=\log j_H(q_{g_0}(B)(b))\;,\\
        d_c(B)&=\log\det_{\mathfrak g}
            \left(I-L^3\partial_{A(b_0(c))}
                \widetilde D_c(g_0\mathsf CB)\right)\;.
    \end{aligned}
\end{equation}
Their supports are respectively $\Lambda_{\alpha},\Lambda_b,\Lambda_c$. The full fine Jacobian is block triangular: its nonidentity rows are designated rows, and its designated-to-designated block is diagonal in $c$. Its determinant is consequently the product of the determinants in $d_c$. Each real determinant is positive in the chart connected to zero, so this branch agrees with the absolute determinant in \eqref{eq:literal-residual}. Summing the action Hessians gives $\mathsf H$, and hence
\begin{equation}\label{eq:perimeter-literal-local-sum}
    \mathcal P_{g_0}(B)
        =\sum_{\alpha} v_{\alpha}(B)+\sum_{b\in E}j_b(B)
            +\sum_{c\in E_1}d_c(B)\;.
\end{equation}


There is no linear term in either $j_b$ or $d_c$. Haar density in exponential variables has zero derivative at zero. For $d_c$, the quadratic part of $\widetilde D_c$ is the quadratic part of $\mathcal N_c$. Expanding RG1 (0.4) to degree two gives a linear combination of Lie brackets of the bond variables. Differentiating with respect to the designated variable gives a linear combination of adjoint maps. Their traces vanish for a compact Lie algebra. Multiplication by the scalar $L^3$ does not change this conclusion. Each $f_{\alpha}$ and its first derivative vanish at zero, and $\psi(A)= A+O(A^2)$, so the subtracted action term begins at degree three.

Use the positive Lie-algebra scalar-product norm for the following estimate. There are constants $a_L,C_L^{\mathrm{loc}}>0$, depending on $L,G$ and the fixed averaging prescription, such that, on the complex polydisc $g_0\max_b\| B(b)\|< a_L$, each term $v$ of \eqref{eq:perimeter-literal-local-sum}, with its indicated support $\Lambda$, satisfies
\begin{equation}\label{eq:perimeter-literal-local-bound}
    | v(B)|\leqslant C_L^{\mathrm{loc}}
        \left(g_0\max_{b\in\Lambda}\| B(b)\|+g_0^2\right)
        \sum_{b\in\Lambda}\| B(b)\|^2\;.
\end{equation}
To obtain the estimate, take a fixed smaller complex chart for the finitely many local templates. Their convergent Taylor series and the vanishing orders just proved give a cubic remainder divided by $g_0^2$ for $v_{\alpha}$ and quadratic remainders for $j_b,d_c$. Equivalence of the finite-dimensional local maximum and Euclidean norms gives the displayed inequality. The bounded overlap then yields a torus-independent quadratic majorant on the cutoff. If $c_{\mathrm{cmp}}$ is the fixed norm-equivalence constant from the CRE bond norm to this norm (the constant of the preceding paragraph), then
\begin{equation}\label{eq:perimeter-literal-total-reserve}
    \sum_{\alpha}| v_{\alpha}(B)|+\sum_b| j_b(B)|+\sum_c| d_c(B)|
        \leqslant C_L^{\mathrm{ov}}
            (c_{\mathrm{cmp}}g_0r(g_0)+g_0^2)\| B\|^2\;.
\end{equation}
The constant $C_L^{\mathrm{ov}}$ includes the bounded overlap and is independent of the torus. These estimates concern the flat reference integral and supply the bounded local potentials for its Gaussian-marginal expansion.

The local scalars are evaluated in the identity-background frame and are invariant under simultaneous constant colour conjugation. This follows from equivariance of RG1's logarithm, averages and implicit correction, cyclicity of the action traces, and conjugation invariance of Haar density and determinants. Their extension to the gauge orbit of this background is unambiguous: if $U^u_0(b)= u(b_-)u(b_+)^{-1}$ and $A^u(b)= u(b_-)A(b)u(b_-)^{-1}$, define the transported scalar at $(U^u_0,A^u)$ to have its identity-frame value at $A$. Two choices of $u$ representing the same background differ by a constant right factor on the connected torus, and give the same scalar by the stated constant-conjugation invariance. This specifies gauge-invariant scalar contractions on this orbit only; it is not an extension to arbitrary nonflat backgrounds. No open Wilson segment is assigned a separate activity or propagated by this decomposition.

\section{The Exterior Response and its History Average}\label{app:exterior-chessboard}
This appendix bounds the exterior response $\mathcal B_{\xi}[f_{\xi}]$ of \eqref{eq:c-exterior-response} on every branch (Lemma~\ref{lem:exterior-response-bound}, Section~\ref{sec:c3-5}), and pays its history average, together with the uncovered branches of Lemma~\ref{lem:covered}, by one-point witness probabilities on compatible even fine periods (Proposition~\ref{prop:c3-7-5}, Section~\ref{sec:c3-7}). Section~\ref{sec:c3-9} proves the deterministic comparison of the two minima used by the witness argument. The inputs are the Hessian, Schur, chart and primitive bounds of Appendix~\ref{app:background-estimates} and the collar statements of Section~\ref{sec:boundary-law-collar}; Remark~\ref{rem:c3-8-1} records what is and is not controlled. Section~\ref{sec:one-step-insertion} applies Sections~\ref{sec:c3-1}--\ref{sec:c3-5} at the retreated carrier (Lemma~\ref{lem:c5-4-4}) and replaces the history bound of Proposition~\ref{prop:c3-7-5} by its labelled version (Proposition~\ref{prop:c5-6-4}).

\subsection{Objects and the conditional law}\label{sec:c3-1}
Fix a branch with its data $\xi$ (Section~\ref{sec:branch-data}): the young carrier $I$, the retained free carrier $J$, the exterior field $A_{\mathrm e}\in\mathfrak g^J$, extended by zero on every other free bond, the reconstructed field $U_{\xi}(a)$ of \eqref{eq:c-reconstructed-field}, and the cutoff norm $|\bigcdot |_{\sigma}$ (the bond norm $|\bigcdot |$ of \eqref{eq:literal-cutoff}, operator or normalised Hilbert--Schmidt, fixed once). The inner product $\langle\bigcdot ,\bigcdot \rangle$, the norm $h$ and the product Hilbert norms are those of Section~\ref{sec:sourced-effective-interaction}; $h\leqslant |\bigcdot |_{\sigma}$ for both permitted norms. The cutoff radius is
\begin{equation}\label{eq:c3-1-1}
    r=\delta_0/g_0=(A_1/A_0)\,p_0(g_0)= A_1(\log g_0^{-2})^{p_0}\;,
\end{equation}
(the first-step cutoff normalisation \eqref{eq:cre-cutoff-normalisation}, in which $p_0$ denotes both the function $p_0(g)= A_0(\log g^{-2})^{p_0}$ and its exponent, as in CRE), so $r\to\infty$ and $g_0r=\delta_0\to0$ as $g_0\to0$. Every derivative, overlap, range and Hessian constant below is independent of $g_0$ once $L$, the group, the character and the source constants are fixed; the two conditions
\begin{equation}\label{eq:c3-1-2}
    r\geqslant 1\qquad\text{and}\qquad C_Pg_0(1+2r)\leqslant \gamma/2\;,
\end{equation}
with $C_P$ the constant of Lemma~\ref{lem:c3-2-1} as enlarged in Lemma~\ref{lem:c3-3-2} (fixed there independently of $g_0$; \eqref{eq:c3-1-2} is imposed with the enlarged constant) and $\gamma=\gamma_L/2$, therefore hold for all $g_0$ below a threshold (the exponent $p_0$ is positive, $p_0\geqslant 5\mathfrak r$ by Section~\ref{app:thresholds}, so $r\to\infty$ while $g_0r\to0$), and they are assumed throughout this appendix, together with the standing conditions \eqref{eq:c-standing-conditions} and the coupling thresholds of Section~\ref{app:thresholds} under which Lemmas~\ref{lem:regular-bounds}, \ref{lem:collar-coverage} and~\ref{lem:c2-primitive-bounds} hold (in particular threshold item 3 there, $4C_{\mathsf C}g_0(1+2r)\leqslant \vartheta_*/2$). If $I=\varnothing$ then $F_{\xi}(1)= F_{\xi}(0)$ (\eqref{eq:convex-response-operands}) and the response \eqref{eq:c-exterior-response} is zero; below $I\neq\varnothing$.

Write
\begin{equation}\label{eq:c3-1-3}
    H= H_{\xi,II},\qquad \ell_{\xi}= H_{\xi,IJ}A_{\mathrm e},\qquad j_{\xi}(a)=\langle\ell_{\xi},a\rangle,\qquad s_{\xi}=1+r+r_{\mathrm e,\xi}\leqslant 1+2r\;,
\end{equation}
with $r_{\mathrm e,\xi}< r$ the retained radius \eqref{eq:convex-native-radius}, the maximum of $| A_{\mathrm e}(b)|_{\sigma}$ over the retained bonds within $3R_{\mathrm c}$ of $I$ (zero if there is none). No bound is imposed on farther retained fields; every retained input entering $H_{\xi,IJ}$ or a near primitive lies in this collar (Lemma~\ref{lem:collar}(i)). The collar bound holds on every branch, independently of contour coverage: by Lemma~\ref{lem:collar-coverage}(i) every retained bond within $3R_{\mathrm c}$ of $I$ lies in $\mathcal G_{\xi}\setminus S_1^*$ ($\mathcal G_{\xi}$ the union of the good test carriers, \eqref{eq:c-good-region}), its factor in $\chi_{\xi,\mathrm e}$ is $\mathbf 1\{|\Psi_{\xi}(A_{\mathrm e})(b)|< r\}$ (\eqref{eq:retained-cutoff-composition}) and $\Psi_{\xi}$, the identity off $S_1^*$ (\eqref{eq:native-source-exact-boundary-map}), is the identity there, so a nonzero retained weight forces $| A_{\mathrm e}(b)|< r$; this is Lemma~\ref{lem:collar-coverage}(iii), not a condition on the covered branches of Lemma~\ref{lem:covered}. On the open box $\mathcal D_I=\{ a\in\mathfrak g^I\, | \,| a_b|_{\sigma}< r\ \forall b\in I\}$ of \eqref{eq:native-source-boundary-interpolation} the interpolating law of that display, which at $t=1$ is the conditional law of the branch (Lemma~\ref{lem:collar}(ii)), is, exactly, for $t\in[0,1]$,
\begin{equation}\label{eq:c3-1-4}
    \dd\nu_{\xi;t}(a)= I_{\xi;t}^{-1}e^{\mathsf Q_{\xi,t}(a)}\,\dd a\;,\qquad
    \mathsf Q_{\xi,t}(a)=-\tfrac12\langle a,Ha\rangle-t\langle\ell_{\xi},a\rangle+P_{\xi,t}(a)\;,
\end{equation}
zero off $\mathcal D_I$ (the indicator in \eqref{eq:native-source-boundary-interpolation}), with $\mathsf Q_{\xi,t}$ as in \eqref{eq:native-collar-flux}; $I_{\xi;t}$ is the normaliser of \eqref{eq:native-source-boundary-interpolation}, not a flat-reference normaliser, and $\dd a$ is the product Lebesgue measure of the fixed inner product. Throughout $f_{\xi}(a)= W_{\chi}(C;U_{\xi}(a))$ (\eqref{eq:convex-response-operands}) with the exterior data fixed at their branch values, defined where $U_{\xi}(a)$ is (Lemma~\ref{lem:charts}), which includes $\mathcal D_I$, and extended by zero elsewhere (immaterial below, since $\nu_{\xi;t}$ vanishes off $\mathcal D_I$); it is bounded by one, and
\begin{equation}\label{eq:c3-1-5}
    L_{\xi,t}=-j_{\xi}+\partial_tP_{\xi,t}\;,\qquad
    \mathcal B_{\xi}[f]=\int_0^1\operatorname{Cov}_{\nu_{\xi;t}}(f,L_{\xi,t})\,\dd t
\end{equation}
(Lemma~\ref{lem:boundary-response}); $L_{\xi,t}=\partial_t\mathsf Q_{\xi,t}$ is called the score below. Derivatives and operator norms refer to the product Hilbert structure, complexified for complex scalar gradients; the fixed matrix-norm equivalence factors of Appendix~\ref{app:background-estimates} are absorbed into the constants before $g_0$ is chosen. The inputs are Lemmas~\ref{lem:regular-bounds} and~\ref{lem:c2-primitive-bounds} under their stated hypotheses (in force by the preceding paragraph): $\gamma\,\id\leqslant H\leqslant m_L\,\id$ with $\gamma=\gamma_L/2$; finite interaction range; and bounded local first and second derivatives of the near primitives on the enlarged box $| a_b|_{\sigma}\leqslant 2r$ at the inputs $a+tA_{\mathrm e}$, $t\in[0,1]$. They are used at the background $U_1$ of the branch; no evaluation at the identity background is made.

\subsection{Local derivative bounds}\label{sec:c3-2}
Let $S_{\alpha}$ be the full free-input stencil, retained inputs included, of a near primitive $\mathfrak p_{\alpha}$ (the $p_{\alpha}$ of Lemma~\ref{lem:c2-primitive-bounds}), let $N_{\mathrm{in}}$ bound its size and $N_{\mathrm{ov}}$ the number of primitive stencils meeting any free bond; both are fixed by $L$ and the group. The distance $d(b,b')$ of two bonds is the periodic $\ell^{\infty}$ distance of their initial sites, plus one for differing directions (the sum of two pseudometrics; a metric, because a positive bond is determined by its initial site and its direction); $d(b,S)=\inf_{b'\in S}d(b,b')$ for a bond set $S$, with $d(b,\varnothing)=+\infty$, and this $d$ is the only bond distance used in this appendix. Let $R=8L+8$, a range dominating in this metric every near-primitive stencil diameter (at most $4L+1$, Lemma~\ref{lem:collar-coverage}, $R_{\mathrm{loc}}=2L+2$) and the Hessian range (at most $\rho_L+1=4L+1$ in this metric, $\rho_L=4L$ being the $\ell^{\infty}$ range of initial sites in (I1) of Section~\ref{sec:c5-0}: Lemmas~\ref{lem:c2-entry-range} and~\ref{lem:hessian-gauge-locality}). Put
\begin{equation}\label{eq:c3-2-1}
    \mathsf K_{\xi}=\{ b\in I\, : \,d(b,J)\leqslant R\}\;,\qquad
    I_C=\{ b\in I\, : \,b\in\mathcal F_r\text{ for some traversal }r\text{ of }C\}\;,
\end{equation}
with $\mathcal F_r$ the free-input stencil \eqref{eq:regular-stencil-extension-support} of the contour bond $r$ (a traversal bond, as in Section~\ref{sec:source-phase-primitives}; the letter $r$ denotes a traversal bond only inside \eqref{eq:c3-2-1} and in the stencils $\mathcal F_r$, and the cutoff radius $r$ of \eqref{eq:c3-1-1} everywhere else); the source depends on the young field only through $I_C$ (Lemma~\ref{lem:row-locality}).

\begin{lemma}[local bounds]\label{lem:c3-2-1}
There are constants $C_P$, $C_S$ (depending on $L$ and the group only) and $C_{\mathrm{resp},\chi}$ (depending on $L$, the group, the character and the maximal multiplicity with which $C$ traverses a link only, the latter $1$ for the square $C_n$) such that on the enlarged box $\{| a_b|_{\sigma}\leqslant 2r\ \forall b\in I\}$, uniformly in $t\in[0,1]$ and in the branch (the fourth bound on every covered branch, Lemma~\ref{lem:covered}),
\begin{gather}
    \| D^2_aP_{\xi,t}\|\leqslant C_Pg_0s_{\xi}\;,\qquad
    h(\nabla_b\partial_tP_{\xi,t})\leqslant C_Sg_0s_{\xi}\,r_{\mathrm e,\xi}\mathbf 1_{\mathsf K_{\xi}}(b)\;,\qquad
    h((\ell_{\xi})_b)\leqslant m_Lr_{\mathrm e,\xi}\mathbf 1_{\mathsf K_{\xi}}(b)\;,\label{eq:c3-2-2}\\
    \sup_a h(\nabla_bf_{\xi}(a))\leqslant C_{\mathrm{resp},\chi}g_0\mathbf 1_{I_C}(b)\;.\label{eq:c3-2-3}
\end{gather}
Under \eqref{eq:c3-1-2} the function $-\mathsf Q_{\xi,t}$ is strongly convex on $\mathcal D_I$ with Hessian at least $\rho=\gamma/2$.
\end{lemma}

\begin{proof}
The first bound of \eqref{eq:c3-2-2} is the second-derivative bound of Lemma~\ref{lem:c2-primitive-bounds} on the enlarged box with the finite overlap $N_{\mathrm{ov}}$. For the second, $\partial_tP_{\xi,t}(a)=\sum_{\alpha} D\mathfrak p_{\alpha}(a+tA_{\mathrm e})[A_{\mathrm e}]$ (by \eqref{eq:native-source-boundary-interpolation}, $P_{\xi,t}$ has no other $t$-dependence), so each summand of $\nabla_b\partial_tP_{\xi,t}$ is $D^2\mathfrak p_{\alpha}$, the full second derivative on $\mathfrak g^{S_{\alpha}}$ (retained inputs included) bounded by Lemma~\ref{lem:c2-primitive-bounds}, applied to a young direction at $b$ and to $A_{\mathrm e}|_{S_{\alpha}}$; it vanishes unless $S_{\alpha}$ contains a retained input, which lies in the collar (stencil diameter at most $4L+1<3R_{\mathrm c}$), so $h(A_{\mathrm e}(b'))\leqslant | A_{\mathrm e}(b')|_{\sigma}\leqslant r_{\mathrm e,\xi}$ for each such input ($h\leqslant |\bigcdot |_{\sigma}$ for both permitted norms, Section~\ref{sec:c3-1}) and $h(A_{\mathrm e}|_{S_{\alpha}})\leqslant \sqrt{N_{\mathrm{in}}}\,r_{\mathrm e,\xi}$, and summing the at most $N_{\mathrm{ov}}$ stencils meeting $b$ gives the bound; the support statement follows from the range $R$ of the stencils. The third bound is the row-sum bound $\sum_{b'\in\Omega_1^*}\|(H_{\xi})_{bb'}\|\leqslant m_L$ for the rows $b\in I$, the input (I2) of Section~\ref{sec:c5-0} (the row-sum statement of Lemma~\ref{lem:regular-bounds}), applied to $A_{\mathrm e}$ on the collar: $(H_{\xi,IJ})_{bb'}=0$ unless the initial sites are within $\ell^{\infty}$ distance $\rho_L$ ((I1) of Section~\ref{sec:c5-0}), so a nonzero entry has $d(b',I)\leqslant \rho_L+1=4L+1<3R_{\mathrm c}$ ($R_{\mathrm c}\geqslant 2R_{\mathrm{loc}}+1=4L+5$) and $h(A_{\mathrm e}(b'))\leqslant | A_{\mathrm e}(b')|_{\sigma}\leqslant r_{\mathrm e,\xi}$, as for the second bound, and $\rho_L+1\leqslant R$ gives the support statement. For \eqref{eq:c3-2-3}, on a covered branch every contour bond $r$ lies in $\Omega_1$ with $\mathcal F_r\subseteq I$ (Lemma~\ref{lem:covered}), so each contour link is reconstructed by a row within $R_{\mathrm{loc}}=2L+2\leqslant 2R_{\mathrm c}+2L$ of $I$, a row of the charts of Lemma~\ref{lem:c2-primitive-bounds}; differentiate the normalised character of the ordered link product: by the rowwise reconstruction (Lemma~\ref{lem:row-locality}) and the uniform charts (Lemma~\ref{lem:c2-uniform-chart}) on the enlarged box with the exterior fixed, whose inputs lie in the phase ball of the charts by Lemma~\ref{lem:c2-primitive-bounds}, each affected link has a phase derivative bounded by $Cg_0$, a fixed input affects a bounded number of links, each link is traversed by $C$ at most as often as its maximal link multiplicity (once for the square $C_n$, a simple closed contour; the multiplicity enters $C_{\mathrm{resp},\chi}$ as a factor), the derivative of the character is bounded in the invariant tangent norm on the compact group, and the other links remain unitary inside the closed contraction; no open Wilson segment is estimated separately. The convexity statement is $\gamma\,\id\leqslant H$ (Lemma~\ref{lem:regular-bounds}) together with $\| D^2P_{\xi,t}\|\leqslant C_Pg_0s_{\xi}\leqslant C_Pg_0(1+2r)\leqslant \gamma/2$, using $s_{\xi}\leqslant 1+2r$ (Lemma~\ref{lem:collar-coverage}(iii)) and \eqref{eq:c3-1-2}.
\end{proof}

\subsection{An extension within the proved chart}\label{sec:c3-3}
The enlarged-box input matters: a radial extension bounded only in $h$ can leave the operator-norm box, and a radius proportional to $3s_{\xi}$ can exceed the enlarged box. Let $B_{\sigma}=\{ x\in\mathfrak g\, : \,| x|_{\sigma}\leqslant 1\}$, let $\Pi_r$ be the Hilbert metric projection onto the closed convex set $\tfrac32rB_{\sigma}$, let $\phi$ be a fixed nonnegative smooth $h$-radial mollifier on $\mathfrak g$ supported in the unit $h$-ball with integral one (in this appendix $\phi$ is only this mollifier, not the phase of Lemma~\ref{lem:c2-uniform-chart}), $\phi_s(y)= s^{-d_G}\phi(y/s)$ for $s>0$ with $d_G=\dim\mathfrak g$, and put $\epsilon=1/(4\sqrt N)$. Define
\begin{equation}\label{eq:c3-3-1}
    T_r(x)=\int_{\mathfrak g}\Pi_r(x-y)\,\phi_{\epsilon r}(y)\,\dd y\qquad(x\in\mathfrak g)\;.
\end{equation}

\begin{lemma}[the extension map]\label{lem:c3-3-1}
$T_r$ is smooth and $\operatorname{Ad}$-equivariant, and
\begin{equation}\label{eq:c3-3-2}
    T_r(x)= x\ \text{ if }| x|_{\sigma}\leqslant r,\qquad | T_r(x)|_{\sigma}\leqslant \tfrac32r,\qquad \| DT_r\|\leqslant 1,\qquad \| D^2T_r\|\leqslant C_{\sigma}/r\;,
\end{equation}
with $C_{\sigma}=4\sqrt N\,\|\nabla\phi\|_{L^1}$, which depends on $N$ and the fixed mollifier only, not on $r$ or $g_0$.
\end{lemma}

\begin{proof}
Both norms are $\operatorname{Ad}$-invariant, so $B_{\sigma}$ is $\operatorname{Ad}$-invariant and $\operatorname{Ad}$ is $h$-orthogonal; hence metric projection and the $h$-radial mollifier commute with the adjoint action, which gives the equivariance. If $h(y)\leqslant \epsilon r$ then $| y|_{\sigma}\leqslant \sqrt N\,h(y)\leqslant r/4$ (Section~\ref{sec:c5-0}: $|\bigcdot |_{\sigma}\leqslant \sqrt N\,h$), so for $| x|_{\sigma}\leqslant r$ one has $| x-y|_{\sigma}\leqslant \tfrac54r<\tfrac32r$ and $\Pi_r(x-y)= x-y$ (indeed for all $| x|_{\sigma}<\tfrac54r$); the convolution then equals $x$ because the mollifier has integral one and, being $h$-radial, zero first moment. The range statement is the convexity of $\tfrac32rB_{\sigma}$: $T_r(x)$ is a probability average of its points. Metric projection is $1$-Lipschitz for $h$ and convolution preserves this, giving $\| DT_r\|\leqslant 1$ in the $h$-operator norm; moving one derivative onto the mollifier of scale $\epsilon r$ bounds $\| D^2T_r\|$ by $\| D\Pi_r\|_{\infty}\|\nabla\phi_{\epsilon r}\|_{L^1}\leqslant \|\nabla\phi\|_{L^1}/(\epsilon r)= C_{\sigma}/r$.
\end{proof}

Apply $T_r$ bondwise to the young inputs only and write $T$ for the product map on $\mathfrak g^I$. Define
\begin{equation}\label{eq:c3-3-3}
    \widetilde P_t(a)=\sum_{\alpha\ \mathrm{near}}\bigl[\mathfrak p_{\alpha}(Ta+tA_{\mathrm e})-\mathfrak p_{\alpha}(tA_{\mathrm e})\bigr],\qquad
    \widetilde f(a)= f_{\xi}(Ta),\qquad
    \widetilde L_t(a)=-\langle\ell_{\xi},a\rangle+\partial_t\widetilde P_t(a)\;.
\end{equation}

\begin{lemma}[the extended objects]\label{lem:c3-3-2}
$\widetilde P_t$, $\widetilde f$ and $\widetilde L_t$ are defined and smooth on all of $\mathfrak g^I$, $|\widetilde f|\leqslant 1$, and on $\mathcal D_I$: $\widetilde f= f_{\xi}$, $\widetilde P_t= P_{\xi,t}-P_{\xi,t}(0)$ (the scalar $P_{\xi,t}(0)$ is removed only in the normalised conditional law; the unnormalised integral keeps it) and $\widetilde L_t-L_{\xi,t}=-\partial_tP_{\xi,t}(0)$, a constant, so that $\operatorname{Cov}_{\nu_{\xi;t}}(\widetilde f,\widetilde L_t)=\operatorname{Cov}_{\nu_{\xi;t}}(f_{\xi},L_{\xi,t})$. After enlarging $C_P$ (before $g_0$ is chosen; the enlarged constant depends on $L$, the group, $C_{\sigma}$, $C_{\mathrm p}$, $N_{\mathrm{ov}}$ and $N_{\mathrm{in}}$, all fixed before the coupling, and \eqref{eq:c3-1-2} is imposed with it), uniformly in $t\in[0,1]$ and $a\in\mathfrak g^I$: $\| D^2\widetilde P_t\|\leqslant C_Pg_0s_{\xi}$, the block Schur row and column sums of $D^2\widetilde P_t$ are at most $C_Pg_0s_{\xi}$ with range $R$,
\begin{equation}\label{eq:c3-3-4}
    h(\nabla_b\partial_t\widetilde P_t)\leqslant C_Sg_0s_{\xi}\,r_{\mathrm e,\xi}\mathbf 1_{\mathsf K_{\xi}}(b)\;,\qquad
    h(\nabla_b\widetilde f)\leqslant C_{\mathrm{resp},\chi}g_0\mathbf 1_{I_C}(b)\;.
\end{equation}
\end{lemma}

\begin{proof}
The image $Ta$ has young cutoff norm at most $\tfrac32r<2r$, so it lies in the enlarged box of Lemma~\ref{lem:c2-primitive-bounds} (threshold item 3 of Section~\ref{app:thresholds}, in force by Section~\ref{sec:c3-1}), where every near primitive and every reconstruction row within $2R_{\mathrm c}+2L$ of $I$ is analytic, the exterior fields being fixed; the real reconstruction gives group-valued links, whence $|\widetilde f|\leqslant 1$. On $\mathcal D_I$, $T$ is the identity (Lemma~\ref{lem:c3-3-1}), which gives the three identities; a constant shift of one operand leaves the covariance unchanged. The Hessian of $\widetilde P_t$ is $D^2\mathfrak p_{\alpha}(DT\bigcdot ,DT\bigcdot )+D\mathfrak p_{\alpha}\,D^2T$ summed over the near primitives: the first term is $O(g_0s_{\xi})$ in operator norm by Lemma~\ref{lem:c3-2-1} and $\| DT\|\leqslant 1$, and per stencil by the second-derivative bound of Lemma~\ref{lem:c2-primitive-bounds} (the per-stencil form is what the Schur sums below use), the second is $O(g_0s_{\xi}^2/r)= O(g_0s_{\xi})$ by the first-derivative bound of Lemma~\ref{lem:c2-primitive-bounds}, $\| D^2T\|\leqslant C_{\sigma}/r$, $r\geqslant 1$ and $s_{\xi}\leqslant 1+2r\leqslant 3r$. Each input belongs to at most $N_{\mathrm{ov}}$ stencils and each stencil has at most $N_{\mathrm{in}}$ blocks, so summing at most $N_{\mathrm{ov}}N_{\mathrm{in}}$ blocks with this chain-rule bound gives the operator and Schur bounds; the supports and the range are those of the original stencils. The mixed derivative $\nabla_b\partial_t\widetilde P_t$ contains no second derivative of $T$, since $T$ does not depend on $t$ and is bondwise: $\nabla_b\partial_t\widetilde P_t(a)= DT_r(a_b)^*(\nabla_b\partial_tP_{\xi,t})(Ta)$, so \eqref{eq:c3-2-2} of Lemma~\ref{lem:c3-2-1} at $Ta$ and $\| DT\|\leqslant 1$ give the first bound of \eqref{eq:c3-3-4}; the second is \eqref{eq:c3-2-3} composed with $DT$.
\end{proof}

\subsection{Covariance decay at the hard boundary}\label{sec:c3-4}
For a positive integer $n$ define the barrier and the barrier laws
\begin{gather}
    \beta_n(a)= n\sum_{b\in I}\bigl[\operatorname{dist}_h(\,\bigcdot \,,rB_{\sigma})^2*\phi_{n^{-2}}\bigr](a_b)\;,\label{eq:c3-4-1}\\
    V_{n,t}(a)=\tfrac12\langle a,Ha\rangle+t\langle\ell_{\xi},a\rangle-\widetilde P_t(a)+\beta_n(a)\;,\qquad
    \dd\nu_{n,t}= Z_{n,t}^{-1}e^{-V_{n,t}(a)}\,\dd a\quad\text{on }\mathfrak g^I\;.\notag
\end{gather}
The squared distance to a convex set is convex and convolution with a nonnegative mollifier preserves convexity, so the barrier Hessian is positive semidefinite; it is diagonal in the bond index because $\beta_n$ is a sum of functions of the single bond variables. For fixed $n$ the potential is smooth with bounded Hessian (that of the barrier is bounded by $2n$: the gradient $2(x-\Pi_r(x))$ of the squared distance is $2$-Lipschitz, and mollification is an average of translates) and Gaussian confinement, its Hessian is at least $\rho I$ with $\rho=\gamma/2$ as in Section~\ref{sec:c3-2} ($\gamma\,\id\leqslant H$ by Lemma~\ref{lem:regular-bounds}, and $\| D^2\widetilde P_t\|\leqslant C_Pg_0s_{\xi}\leqslant C_Pg_0(1+2r)\leqslant \gamma/2$ by Lemma~\ref{lem:c3-3-2} and \eqref{eq:c3-1-2}--\eqref{eq:c3-1-3}), and its off-diagonal block row and column sums are at most
\begin{equation}\label{eq:c3-4-2}
    M_*= m_L+\gamma/2\;,
\end{equation}
with range at most $R$ in the bond distance $d$ of Section~\ref{sec:c3-2}, independently of $n$, $t$, the volume and the branch (the barrier is diagonal; the row sums of $H$ are at most $m_L$ by (I2) of Section~\ref{sec:c5-0}, its column sums equal its row sums since $H_{bb'}= H_{b'b}^{\mathsf T}$ have equal norms, and its range is $\rho_L=4L\leqslant R$ by (I1) there; the Schur sums of $D^2\widetilde P_t$ are at most $C_Pg_0s_{\xi}\leqslant \gamma/2$ with range $R$ by Lemma~\ref{lem:c3-3-2}).

\begin{lemma}[Helffer--Sj\"ostrand identity and decay]\label{lem:c3-4-1}
Let $\mathscr A_{n,t}=\mathscr L_{n,t}\otimes I+D^2V_{n,t}$ with $\mathscr L_{n,t}=-\Delta+\nabla V_{n,t}\cdot\nabla$, on the Hilbert direct sum over $b\in I$ of $L^2(\nu_{n,t};\mathfrak g)$. Then $\mathscr A_{n,t}\geqslant \rho I$ ($\mathscr L_{n,t}=\nabla^*\nabla\geqslant 0$ on $L^2(\nu_{n,t})$ and $D^2V_{n,t}\geqslant \rho I$), and for real $u,v$ in the weighted space $H^1(\nu_{n,t})$
\begin{equation}\label{eq:c3-4-3}
    \operatorname{Cov}_{\nu_{n,t}}(u,v)=\langle\nabla u,\mathscr A_{n,t}^{-1}\nabla v\rangle_{L^2(\nu_{n,t})}\;.
\end{equation}
With $\mu= R^{-1}\log(1+\gamma/(4M_*))$ and $\Pi_b$ the projection onto the $b$-component,
\begin{equation}\label{eq:c3-4-4}
    \|\Pi_b\mathscr A_{n,t}^{-1}\Pi_c\|\leqslant \frac4\gamma e^{-\mu d(b,c)}\;,\qquad
    |\operatorname{Cov}_{\nu_{n,t}}(u,v)|\leqslant \frac4\gamma\sum_{b,c\in I}e^{-\mu d(b,c)}\|\nabla_bu\|_{L^2(\nu_{n,t})}\|\nabla_cv\|_{L^2(\nu_{n,t})}\;.
\end{equation}
For complex $u$ and real $v$ (the case used: $u=\widetilde f$, the Wilson character, and $v=\widetilde L_t$), applying \eqref{eq:c3-4-4} to the real and imaginary parts of $u$ gives the same bound with the gradient norm of $u$ read as the sum of the gradient norms of its two parts; each is at most the gradient norm of $u$, a factor at most two that Section~\ref{sec:c3-5} absorbs into $C_{\mathrm{resp},\chi}$.
\end{lemma}

\begin{proof}
The identity \eqref{eq:c3-4-3} is \cite[Proposition 4.5]{Johnsen} (Theorem 1.3 there at $k=0$) for the real-valued $C^2$ potential $V_{n,t}+\log Z_{n,t}$ (normalised; the additive constant leaves the derivatives unchanged; $V_{n,t}$ is real because its quadratic, linear and barrier terms are real and $\widetilde P_t$ is real, the primitives being real on real inputs with a positive determinant, Lemma~\ref{lem:c2-uniform-chart}(iii) and Lemma~\ref{lem:c2-primitive-bounds}). Its conditions are a real $C^2$ normalised potential and the closed range of the operator $A_0$ there; the latter follows from $A_1\geqslant \rho I$, which is $\mathscr A_{n,t}\geqslant \rho I$ here, by \cite[Theorem 1.3]{Johnsen}, as explained there after (4.39). The theorem gives the projection onto the orthogonal complement of the constants as $d_0^*(A_1|_{\overline{\operatorname{ran}}\,d_0})^{-1}d_0$, and at fixed $n$ the bounded Hessian identifies $A_1$ with the displayed operator $\mathscr A_{n,t}$: $A_1$ is the Friedrichs extension of \cite[Lemma 4.2]{Johnsen} with the action (4.22) there, $\mathscr L_{n,t}\otimes I+D^2V_{n,t}$ on gradient fields, and since $D^2V_{n,t}$ is a bounded multiplication operator at fixed $n$ its form domain is that of $\mathscr L_{n,t}\otimes I$ and $A_1$ is the operator sum on the domain of $\mathscr L_{n,t}\otimes I$, a domain the bounded field-independent weights $W_c$ below preserve. The operands are the source $\widetilde f$, bounded with bounded gradient (Lemma~\ref{lem:c3-3-2}), and the score $\widetilde L_t=-\langle\ell_{\xi},\bigcdot \rangle+\partial_t\widetilde P_t$, which grows at most linearly and has bounded gradient by \eqref{eq:c3-3-4}; both lie in $H^1(\nu_{n,t})$. For a fixed $c\in I$ let $W_c$ multiply the component $b$ by $e^{\mu d(b,c)}$; it is constant in field space, hence commutes with the scalar generator $\mathscr L_{n,t}\otimes I$ and with every diagonal Hessian block, the barrier included, and the off-diagonal Schur bounds give $\| W_c\mathscr A_{n,t}W_c^{-1}-\mathscr A_{n,t}\|\leqslant M_*(e^{\mu R}-1)\leqslant \rho/2$ by the choice of $\mu$. $W_c$ and its inverse are bounded at finite volume, so a resolvent Neumann series gives $\| W_c\mathscr A_{n,t}^{-1}W_c^{-1}\|\leqslant (\rho-\rho/2)^{-1}=4/\gamma$; compressing by $\Pi_b$ on the left and $\Pi_c$ on the right, where $W_c$ acts as $e^{\mu d(b,c)}$ and as $e^{\mu d(c,c)}=1$, gives the first bound of \eqref{eq:c3-4-4}; the second follows from \eqref{eq:c3-4-3} and the Cauchy--Schwarz inequality. The large diagonal barrier never enters the decay constant.
\end{proof}

\begin{lemma}[the barrier limit]\label{lem:c3-4-2}
For $u,v$ bounded or of linear growth with bounded gradients, $\operatorname{Cov}_{\nu_{n,t}}(u,v)\to\operatorname{Cov}_{\nu_{\xi;t}}(u,v)$ as $n\to\infty$. Consequently the second bound of \eqref{eq:c3-4-4} holds for $\nu_{\xi;t}$ on the hard box $\mathcal D_I$, with the gradient norms taken as global suprema, $\|\nabla_bu\|_{\infty}\|\nabla_cv\|_{\infty}$ in place of the $L^2$ norms.
\end{lemma}

\begin{proof}
At finite volume $\beta_n\to0$ on the closed cutoff body $\{| a_b|_{\sigma}\leqslant r\ \forall b\}$ and $\beta_n\to\infty$ outside it, the boundary of the body has Lebesgue measure zero, and $\beta_n\geqslant 0$ together with the Gaussian confinement supplies an integrable majorant for $u$, $v$ and $uv$ (bounded or of linear growth); dominated convergence gives the convergence of the unnormalised integrals, and of the normaliser to $\int_{\mathcal D_I}e^{-V^{\circ}_t}\,\dd a>0$ (the box has nonempty interior), where $V^{\circ}_t= V_{n,t}-\beta_n$ is the potential without the barrier, independent of $n$, hence of the covariances, to those of the law with density $e^{-V^{\circ}_t}$ on $\mathcal D_I$, which is $\nu_{\xi;t}$ by Lemma~\ref{lem:c3-3-2} ($\widetilde P_t= P_{\xi,t}-P_{\xi,t}(0)$ on $\mathcal D_I$, and \eqref{eq:c3-1-4}). On the right of \eqref{eq:c3-4-4} the $L^2(\nu_{n,t})$ norms depend on $n$, but they are bounded by the global suprema $\|\nabla_bu\|_{\infty}$, $\|\nabla_cv\|_{\infty}$, which do not, and $\mu$ and $4/\gamma$ are independent of $n$, so the bound passes to the limit. No translated cutoff and no new exterior law is introduced.
\end{proof}

\subsection{The response}\label{sec:c3-5}
Put
\begin{equation}\label{eq:c3-5-1}
    D_{\xi}=\sum_{b\in I_C}e^{-\mu d(b,J)/2}\;,
\end{equation}
with $\mu$ of Lemma~\ref{lem:c3-4-1} and the bond distance $d$ of Section~\ref{sec:c3-2}; $D_{\xi}=0$ when $J=\varnothing$ (by $d(b,\varnothing)=+\infty$; then also $A_{\mathrm e}=0$ and the score vanishes).

\begin{lemma}[the exterior response]\label{lem:exterior-response-bound}
Under \eqref{eq:c3-1-2}, on every covered branch (Lemma~\ref{lem:covered}; the bound is used on covered branches only, Section~\ref{sec:c3-6} and Lemma~\ref{lem:c5-4-4}) and for every $t\in[0,1]$,
\begin{equation}\label{eq:c3-5-2}
    |\operatorname{Cov}_{\nu_{\xi;t}}(f_{\xi},L_{\xi,t})|\leqslant C_{\mathrm{resp},\chi}\,g_0r\sum_{b\in I_C}e^{-\mu d(b,J)/2}\;,\qquad
    |\mathcal B_{\xi}[f_{\xi}]|\leqslant C_{\mathrm{resp},\chi}\,g_0r\,D_{\xi}\;,
\end{equation}
after enlarging the constant $C_{\mathrm{resp},\chi}$ of Lemma~\ref{lem:c3-2-1} (the enlarged constant keeps the name); it depends on $L$, $\gamma$, $m_L$, $M_*$, $R$, $N_{\mathrm{in}}$, $N_{\mathrm{ov}}$, $C_P$, $C_S$, $C_{\sigma}$, the character, the group and the fixed source parameters, and is independent of $g_0$, the branch, the volume, $t$ and the barrier index.
\end{lemma}

\begin{proof}
By Lemma~\ref{lem:c3-3-2}, $\operatorname{Cov}_{\nu_{\xi;t}}(f_{\xi},L_{\xi,t})=\operatorname{Cov}_{\nu_{\xi;t}}(\widetilde f,\widetilde L_t)$, and Lemma~\ref{lem:c3-4-2} bounds it by $(8/\gamma)\sum_{b,c\in I}e^{-\mu d(b,c)}\|\nabla_b\widetilde f\|_{\infty}\|\nabla_c\widetilde L_t\|_{\infty}$: the bound \eqref{eq:c3-4-4} is applied to the real and imaginary parts of the complex $\widetilde f$ and the real $\widetilde L_t$ (Lemma~\ref{lem:c3-4-1}), the sum of the two gradient norms being at most $2\|\nabla_b\widetilde f\|_{\infty}$; this is the factor two that Lemma~\ref{lem:c3-4-1} leaves to this section. By \eqref{eq:c3-3-4} and \eqref{eq:c3-2-2}, $\|\nabla_c\widetilde L_t\|_{\infty}\leqslant (m_L+C_Sg_0s_{\xi})r_{\mathrm e,\xi}\mathbf 1_{\mathsf K_{\xi}}(c)\leqslant (m_L+C_S\gamma/(2C_P))\,r\,\mathbf 1_{\mathsf K_{\xi}}(c)$ (as $r_{\mathrm e,\xi}< r$ and $C_Sg_0s_{\xi}\leqslant C_Sg_0(1+2r)\leqslant C_S\gamma/(2C_P)$ by \eqref{eq:c3-1-2}; the coefficient is a constant fixed before $g_0$) and $\|\nabla_b\widetilde f\|_{\infty}\leqslant C_{\mathrm{resp},\chi}g_0\mathbf 1_{I_C}(b)$ with the constant of Lemma~\ref{lem:c3-2-1}. For the shell step, $d(c,J)\leqslant R$ on $\mathsf K_{\xi}$, so $d(b,c)\geqslant d(b,J)-R$ and
\begin{equation*}
    e^{-\mu d(b,c)}= e^{-\mu d(b,c)/2}e^{-\mu d(b,c)/2}\leqslant e^{\mu R/2}e^{-\mu d(b,J)/2}\,e^{-\mu d(b,c)/2}\;;
\end{equation*}
the sum of $e^{-\mu d(b,c)/2}$ over all periodic bonds $c$ is bounded by the infinite-lattice shell sum $S_{\mu/2}=\sup_b\sum_ce^{-\mu d(b,c)/2}$ (the nearest-lift map from torus bonds to lattice bonds is injective and the terms are nonnegative), uniformly in the periods. Collecting the constants, $(8/\gamma)(m_L+C_S\gamma/(2C_P))e^{\mu R/2}S_{\mu/2}$ times the constant of Lemma~\ref{lem:c3-2-1}, which depends on $\mu$ through $e^{\mu R/2}$ and $S_{\mu/2}$ and hence on $M_*$, gives the first bound of \eqref{eq:c3-5-2} with the enlarged $C_{\mathrm{resp},\chi}$; integrating over $t\in[0,1]$ in \eqref{eq:c3-1-5} gives the second.
\end{proof}

\begin{remark}[the surface terms are retained]\label{rem:c3-5-1}
In the notation \eqref{eq:native-collar-mean-direction}--\eqref{eq:native-collar-flux}, Lemma~\ref{lem:ibp} (its hypothesis holds: $f_{\xi}$ is $C^1$ on the enlarged box $\{| a_b|_{\sigma}\leqslant 2r\}$, a neighbourhood of $\overline{\mathcal D_I}$, by the charts of Lemma~\ref{lem:c2-primitive-bounds} (Lemma~\ref{lem:charts}) and Lemma~\ref{lem:c2-uniform-chart}, as in the proof of Lemma~\ref{lem:c3-2-1}) writes the integrand of the response as
\begin{equation}\label{eq:c3-5-3}
    -\mathbb E_{\nu_{\xi;t}}[D_{m_{\xi}}f_{\xi}]+\operatorname{Cov}_{\nu_{\xi;t}}(f_{\xi},\partial_tP_{\xi,t}-D_{m_{\xi}}P_{\xi,t})+\mathcal T_{\xi,t}[f_{\xi}]-\mathbb E_{\nu_{\xi;t}}[f_{\xi}]\,\mathcal T_{\xi,t}[1]\;,
\end{equation}
with the hard-cutoff flux $\mathcal T_{\xi,t}[u]= I_{\xi;t}^{-1}\int_{\partial\mathcal D_I}u\,e^{\mathsf Q_{\xi,t}}\langle m_{\xi},n\rangle\,\dd S$. Lemma~\ref{lem:exterior-response-bound} bounds this whole expression: neither flux is set to zero or estimated separately. The extension $T_r$ of Lemma~\ref{lem:c3-3-1} is the identity on the neighbourhood $\{| a_b|_{\sigma}<\tfrac54r\}$ of $\overline{\mathcal D_I}$ and the barrier limit of Lemma~\ref{lem:c3-4-2} returns the original measure, so the boundary traces and the normalised covariance in \eqref{eq:c3-5-3} are those of Section~\ref{sec:boundary-law-collar}.
\end{remark}

\subsection{The physical substitution}\label{sec:c3-6}
Let $\mathsf w_{\xi}= e^{c_{\xi}}I_{\xi;1}=\mathcal J^{\mathrm a}_{\xi}[1]$ be the branch weight \eqref{eq:c-collar-identity} and $\mathcal O$ the outer operation of \eqref{eq:native-source-explicit-outer-operation}, so that $\mathcal Z_0=\int\mathcal O[\mathsf w_{\xi}](V)\,\dd V$ (\eqref{eq:native-source-physical-ratio}). Define the history expectation and probability
\begin{equation}\label{eq:c3-6-1}
    \mathbb E_{\mathrm h}[X]=\mathcal Z_0^{-1}\int\mathcal O[\mathsf w_{\xi} X](V)\,\dd V\;,\qquad
    \mathbb P_{\mathrm h}(A)=\mathbb E_{\mathrm h}[\mathbf 1_A]\;,
\end{equation}
for nonnegative or bounded Borel functions $X$ of the branch data; $\mathcal O$ has a nonnegative kernel (its Haar and constraint factors are measures and its cutoffs and exponentials are nonnegative, \eqref{eq:native-source-explicit-outer-operation}) and $0<\mathcal Z_0<\infty$ (the measurability paragraph of Section~\ref{sec:branch-data}), so $\mathbb E_{\mathrm h}$ is positive, $|\mathbb E_{\mathrm h}[X]|\leqslant \sup| X|$ and $\mathbb E_{\mathrm h}[1]=1$. Let $A_{\mathrm{cov}}$ be the set of covered branches of Lemma~\ref{lem:covered}: every contour bond is a bond of $\Omega_1$ and $\mathcal F_r\subseteq I$ for every contour bond $r$. The event $A_{\mathrm{cov}}$ is determined by $\Omega_1$, $R_1$ and the fixed contour templates ($I=\Lambda_1^*$, the free bonds with an endpoint in $\Lambda_1$, Section~\ref{sec:branch-data}, with $\Lambda_1$ derived from $\Omega_1$ and $R_1$ by CRE (1.20) and its filling; the stencils $\mathcal F_r$ are fixed); it and $D_{\xi}$ are measurable functions of the outer branch data, constant in the young integration, and $0\leqslant D_{\xi}\leqslant | I_C|$, a bound by the number of free inputs of the contour. The response bound is applied on covered branches only. The combined exterior-response and uncovered contribution to $\mathbb E_{\varepsilon}[W_{\chi}(C)]$ of \eqref{eq:native-source-physical-ratio} is
\begin{equation}\label{eq:c3-6-2}
    E_{\mathrm{bdry}}=\mathbb E_{\mathrm h}\bigl[\mathbf 1_{A_{\mathrm{cov}}}\mathcal B_{\xi}[f_{\xi}]+\mathbf 1_{A_{\mathrm{cov}}^c}F_{\xi}(1)\bigr]\;,
\end{equation}
the exterior response on the covered branches and the whole conditional expectation on the uncovered ones (\eqref{eq:c-exterior-response}, \eqref{eq:c-one-step-branch-identity}); the integrand is a bounded Borel function of the branch data ($F_{\xi}(t)=\int f_{\xi} K_{\xi;t}\,\dd a/I_{\xi;t}$ is Borel in the branch data for each $t$, by the composition argument of Section~\ref{sec:branch-data} applied to the interpolated kernel of \eqref{eq:convex-response-operands} and Fubini's theorem, and $|\mathcal B_{\xi}[f_{\xi}]|\leqslant C_{\mathrm{resp},\chi}g_0rD_{\xi}$ is bounded), so \eqref{eq:c3-6-1} applies to it. No bound on $\mathbb E_{\mathrm h}[\mathbf 1_{A_{\mathrm{cov}}}D_{\xi}]$ or on $\mathbb P_{\mathrm h}(A_{\mathrm{cov}}^c)$ is used in this subsection; both are bounded in Section~\ref{sec:c3-7} (Proposition~\ref{prop:c3-7-5}).

\begin{corollary}\label{cor:c3-6-1}
Under \eqref{eq:c3-1-2},
\begin{gather}
    \Bigl|\mathcal Z_0^{-1}\int\mathcal O\bigl[\mathsf w_{\xi}\mathbf 1_{A_{\mathrm{cov}}}\mathcal B_{\xi}[f_{\xi}]\bigr](V)\,\dd V\Bigr|\leqslant C_{\mathrm{resp},\chi}\,g_0r\,\mathbb E_{\mathrm h}[\mathbf 1_{A_{\mathrm{cov}}}D_{\xi}]\;,\label{eq:c3-6-3}\\
    | E_{\mathrm{bdry}}|\leqslant C_{\mathrm{resp},\chi}\,g_0r\,\mathbb E_{\mathrm h}[\mathbf 1_{A_{\mathrm{cov}}}D_{\xi}]+\mathbb P_{\mathrm h}(A_{\mathrm{cov}}^c)\;.\notag
\end{gather}
\end{corollary}

\begin{proof}
$\mathsf w_{\xi}\geqslant 0$ and $\mathcal O$ is positive, so Lemma~\ref{lem:exterior-response-bound} gives the first bound pointwise on the branches; the triangle inequality, \eqref{eq:c3-6-3}, positivity and $| F_{\xi}(1)|\leqslant 1$ give the second. No conditional or flat-reference normaliser replaces $\mathcal Z_0$, and $c_{\xi}$ stays inside $\mathsf w_{\xi}$.
\end{proof}

\subsection{The history average by one-point witnesses}\label{sec:c3-7}
This section assumes compatible even fine periods: every fine period is a multiple of $\ell= LMR_1$ and at least $8M$ (\eqref{eq:c-standing-conditions}) and even. This is a restriction relative to the compatible periods of Section~\ref{sec:sourced-effective-interaction}; it includes the torus with periods $2L^m$ when the partitions fit. Lemma~\ref{lem:exterior-response-bound} needs no parity restriction, and no odd-period history bound is asserted.

Write $\ell= LMR_1$ for the macro side ($R_1$ the scale integer), $B_{\xi}=\Lambda_1^c$ for the filled bad region, $E(X)$ for the set of positive fine bonds with at least one endpoint in a set $X$ of sites, and $n_C$ for the number of contour traversals. Let $\Gamma_C$ be the deterministic set of the contour bonds together with the free-input templates of their reconstruction rows, projected periodically from the torus templates before any history is selected (so $\mathcal F_r\subseteq\Gamma_C$ for every contour bond $r$, with $\mathcal F_r$ the free-input set of \eqref{eq:regular-stencil-extension-support}, the set appearing in the covered event of Lemma~\ref{lem:covered}); there is a constant $C_L$, depending on $L$ only (the templates are fixed lattice stencils), with $|\Gamma_C|\leqslant C_Ln_C$, and $I_C\subseteq\Gamma_C$ on every branch (the young designated variables are eliminated, so their free-input stencils, not the designated bonds, are the members of $\Gamma_C$). For a positive bond $z$ define the marked probability
\begin{equation}\label{eq:c3-7-1}
    p_z=\mathcal Z_0^{-1}\int\mathcal O\bigl[\mathsf w_{\xi}\mathbf 1\{ z\in E(B_{\xi})\}\bigr](V)\,\dd V=\mathbb P_{\mathrm h}(z\in E(B_{\xi}))\;;
\end{equation}
the mark depends on the regions and on the deterministic filling, not on the young variable.

\begin{lemma}[reduction to one-point marks]\label{lem:c3-7-1}
Let $c=\mu/2$ and $S_c=\sup_b\sum_ze^{-cd(b,z)}$, the periodic bond shell sum, bounded by its infinite-lattice counterpart uniformly in the periods. Pointwise on every branch
\begin{equation}\label{eq:c3-7-2}
    D_{\xi}\leqslant \sum_{b\in\Gamma_C}\sum_ze^{-cd(b,z)}\mathbf 1\{ z\in E(B_{\xi})\}\;,\qquad
    \mathbf 1_{A_{\mathrm{cov}}^c}\leqslant \sum_{b\in\Gamma_C}\mathbf 1\{ b\in E(B_{\xi})\}\;.
\end{equation}
Hence, if $p_z\leqslant q$ for every positive bond $z$,
\begin{equation}\label{eq:c3-7-3}
    \mathbb E_{\mathrm h}[D_{\xi}]\leqslant C_LS_cn_Cq\;,\qquad \mathbb P_{\mathrm h}(A_{\mathrm{cov}}^c)\leqslant C_Ln_Cq\;.
\end{equation}
\end{lemma}

\begin{proof}
$J\subseteq E(B_{\xi})$: a free bond not in $I$ has no endpoint in $\Lambda_1$. So each $b\in I_C\subseteq\Gamma_C$ contributes $e^{-cd(b,J)}\leqslant \sum_{z\in J}e^{-cd(b,z)}\leqslant \sum_ze^{-cd(b,z)}\mathbf 1\{ z\in E(B_{\xi})\}$. If coverage fails, either a contour bond is not a bond of $\Omega_1$, and then it meets $B_{\xi}$ since $\Lambda_1\subseteq\Omega_1$, or a free input of the contour is not in $I$ (some $b\in\mathcal F_r\setminus I$ for a contour bond $r$, and $\mathcal F_r\subseteq\Gamma_C$), and then it is retained in $J$ or lies outside $\Omega_1$ and meets $B_{\xi}$; in both cases some $b\in\Gamma_C$ lies in $E(B_{\xi})$. Taking $\mathbb E_{\mathrm h}$ of \eqref{eq:c3-7-2} (Tonelli; all integrands nonnegative) and using $p_z\leqslant q$ and $|\Gamma_C|\leqslant C_Ln_C$ gives \eqref{eq:c3-7-3}.
\end{proof}

\subsubsection{Cutoff failures give local witnesses}\label{sec:c3-7-1}
The witness argument operates on the original unresummed cutoffs CRE (1.1), (1.4), (1.8), (1.18), before the boundary substitution; it is proved in Lemmas~\ref{lem:c3-7-block}, \ref{lem:c3-7-seed} and~\ref{lem:c3-7-2} below. Its conclusion: a bad eligible test forces a fine plaquette with $| U(\partial p)-\id|_{\sigma}\geqslant \theta$ in its neighbourhood, or an active mean-transport variable with $| w-\id|_{\sigma}\geqslant \theta$ there (the $P_0$ and $P_1$ alternatives need no gauge-variable witness), with the threshold
\begin{equation}\label{eq:c3-7-4}
    \theta=\min\Bigl(\varepsilon_0,\ \frac{\varepsilon_1}{4B_3N^{3/2}L^2},\ \frac{\delta_0}{2C_Q}\Bigr)\;,\qquad
    C_Q=(9K_L+2)(1+2B_3N^{3/2}L^2)\;,
\end{equation}
where $K_L$ is the plaquette-filling constant of the loops of RG1 (0.4), (0.11) and of the bond loops of one block or of two adjacent blocks, defined in the paragraph ``Constants of the block comparison'' below. The norm is the cutoff norm $|\bigcdot |_{\sigma}$, and the mean and comparison constants are chosen for it. The seed is used on the original earlier-cutoff and averaging-constraint support only, not for arbitrary frozen exterior data. Its common local width $w$ is $LM_2R_1$ or $L^2M_2R_1$, with one convention throughout, and its macro width is $\ell$; these are the conventions of the branch hypotheses of Section~\ref{sec:branch-data}. For a cube $Q$ and $k\geqslant 1$, $Q_w^{[k]}$ denotes $Q$ enlarged by $k$ layers of width $w$ (CRE's $\square'^{\sim k}$ when $w= L^2M_2R_1$), so that $\operatorname{side}(Q_w^{[k]})=\operatorname{side}(Q)+2kw$ and $Q_w^{[2]}$ is at $\ell^{\infty}$ distance $2w$ from the complement of $Q_w^{[4]}$; no mixed-width version is asserted. An eligible cube is a cube $\square'$ of CRE's partition into $L^2M_2R_1$-cubes carrying the test (1.18), which CRE performs for $\square'\subset\Omega_1^{\sim-1}$ (p.250), so every eligible cube satisfies the hypothesis $Q\subseteq\Omega_1^{\sim-1}$ of Proposition~\ref{prop:c3-9-1}.

\emph{Constants of the block comparison.} For a block centre $y$ and a site $x$ of its block let $\mathbf G(y,x)$ be RG1's finite family of shortest paths from $y$ to $x$ (RG1 p.252, the definition preceding (0.4), used in (0.11)): for each permutation $\pi$ of the axes the path moves first along $e_{\pi(1)}$ by the full coordinate difference, then along $e_{\pi(2)}$, and so on; each coordinate is monotone along the path, so the path stays in the block $B(y)$. Let $\Gamma_x\in\mathbf G(y,x)$ be a fixed representative ($\Gamma_y$ the empty path, with transport $\id$). Let $K_L$ bound the number of plaquettes needed to fill each closed loop of the following finite family, each filling chosen inside the union of the block or blocks containing the loop (a lattice box, in which every closed lattice loop bounds a plaquette filling): the difference of two paths of one $\mathbf G(y,x)$; the loop $\Gamma_x\cup b\cup(-\Gamma_{x'})$ of a block-internal bond $b=(x,x')$; the loop $\Gamma_x\cup b\cup(-\Gamma_{x'})\cup(-c)$ of an interface bond $b=(x,x')$ of the coarse bond $c$, $x\in B(c_-)$, $x'\in B(c_+)$, $\Gamma_x\in\mathbf G(c_-,x)$, $\Gamma_{x'}\in\mathbf G(c_+,x')$; and the loops $\Gamma\cup[x,x']\cup(-\Gamma')\cup(-c)$ of RG1 (0.4), for a site $x\in B(c_-)$, with $[x,x']$ the parallel transport of the coarse bond $c$ to $x$ (RG1 p.252, the sentence preceding the definition of the averages: the straight contour of $L$ fine bonds in the direction of $c$ from $x$ to $x'= x+(c_+-c_-)\in B(c_+)$) and paths $\Gamma\in\mathbf G(c_-,x)$, $\Gamma'\in\mathbf G(c_+,x')$. Each loop lies in one block or in the union $B(c_-)\cup B(c_+)$ of two adjacent blocks (the contour $[x,x']$ runs inside this union, a box of sides $2L,L,L,L$), so $K_L$ depends on $L$ only. The transports of a unitary field $F$ along two paths with the same endpoints differ in $|\bigcdot |_{\sigma}$ by at most the sum of $| F(\partial p)-\id|_{\sigma}$ over the plaquettes filling their loop (unitary invariance of the norm and telescoping), hence by at most $K_L\eta$ when $| F(\partial p)-\id|_{\sigma}\leqslant \eta$ for every plaquette $p$ with its four vertices in $B(c_-)\cup B(c_+)$ (``the plaquettes of the two blocks''; this includes the plaquettes straddling the interface). For the group mean $M$ of RG1 (0.5)--(0.10) let $\mathcal N_L$ be the finite set of cardinalities $|\mathbf G(y,x)|$ and, for $n\in\mathcal N_L$, $\mathfrak m_n(A_1,\ldots,A_n)=(1/i)\log M(e^{iA_1},\ldots,e^{iA_n})$; it is analytic near zero with $\mathfrak m_n(0)=0$ and derivative $A\mapsto n^{-1}\sum_jA_j$ at zero (RG1 (0.8)). Let $r_{\mathrm M}>0$ be a radius such that every $\mathfrak m_n$, $n\in\mathcal N_L$, is analytic on the closed maximum-norm ball of $n$-tuples of radius $r_{\mathrm M}$ (the minimum over the finite set $\mathcal N_L$ of such radii), and $C_{\mathrm M}$ half the maximum over $n\in\mathcal N_L$ of the second-derivative norms on these balls; Taylor's theorem gives $|\mathfrak m_n(A)|_{\sigma}\leqslant \max_j| A_j|_{\sigma}+C_{\mathrm M}(\max_j| A_j|_{\sigma})^2$ for $\max_j| A_j|_{\sigma}\leqslant r_{\mathrm M}$. On the principal unitary chart $|(1/i)\log W|_{\sigma}\leqslant | W-\id|_{\sigma}+C_{\log}| W-\id|_{\sigma}^2$ with a fixed $C_{\log}$: eigenvalue by eigenvalue, $2\arcsin(t/2)\leqslant t+Ct^2$ for $0\leqslant t\leqslant 2$, $t_j=| e^{i\varphi_j}-1|$, which gives the operator-norm case directly and the normalised Hilbert--Schmidt case with a constant depending on $N$ ($\max_jt_j\leqslant \sqrt N| W-\id|_{\sigma}$). Put $C'_{\mathrm M}= C_{\log}+4C_{\mathrm M}$. If $| W_jW_k^{-1}-\id|_{\sigma}\leqslant \rho$ for all $j$, with $\rho\leqslant 1/\max(1,C_{\log})$ and $2\rho\leqslant r_{\mathrm M}$, then $| M(W_1,\ldots,W_n)-W_k|_{\sigma}\leqslant \rho+C'_{\mathrm M}\rho^2$: translate by $W_k^{-1}$ (RG1 (0.6)), so that the inputs $A_j=(1/i)\log(W_jW_k^{-1})$ have norm at most $\rho+C_{\log}\rho^2\leqslant 2\rho$; the phase $H=\mathfrak m_n(A)$ is Hermitian (RG1 (0.9), the $W_j$ being unitary), so $| e^{iH}-\id|_{\sigma}\leqslant | H|_{\sigma}\leqslant \rho+C_{\log}\rho^2+4C_{\mathrm M}\rho^2$. Fix $\eta_0>0$ with $C'_{\mathrm M}K_L\eta_0\leqslant 1$, $2K_L\eta_0\leqslant r_{\mathrm M}$, $K_L\eta_0\leqslant 1$, $C_{\log}K_L\eta_0\leqslant 1$ (so $\rho= K_L\eta_0\leqslant 1/\max(1,C_{\log})$) and $2K_L\eta_0\leqslant r_G/\sqrt N$, so that every $W\in G$ with $| W-\id|_{\sigma}\leqslant 2K_L\eta_0$ lies in the chart domain: $\| W-\id\|_{\mathrm{op}}\leqslant \sqrt N| W-\id|_{\sigma}\leqslant r_G$, and $r_G$ is the radius of Appendix~\ref{app:background-estimates} with this property (the eigenvalue bound alone does not place $W$ in the chart domain of a general compact $G$, the preamble) (every logarithm above is taken at such a $W$: the loops have $| W-\id|_{\sigma}\leqslant K_L\eta_0$, the translated means $| W_jW_k^{-1}-\id|_{\sigma}\leqslant \rho\leqslant K_L\eta_0$); for $\eta\leqslant \eta_0$ and $\rho= K_L\eta$ this gives $\rho+C'_{\mathrm M}\rho^2\leqslant 2K_L\eta$. The constants $K_L$, $r_{\mathrm M}$, $C_{\mathrm M}$, $C_{\log}$, $\eta_0$ depend on $L$, the group and the norm only.

\begin{lemma}[block comparison]\label{lem:c3-7-block}
Let $F$ be a unitary field on the bonds of two blocks $B(c_-)$, $B(c_+)$ joined by the coarse bond $c$ (or of one block), with $| F(\partial p)-\id|_{\sigma}\leqslant \eta\leqslant \eta_0$ for every plaquette with its four vertices in $B(c_-)\cup B(c_+)$ and $| F(y,x)-\id|_{\sigma}\leqslant \eta$ for its averaged contour variables $F(y,x)= M(\{ F(\Gamma)\}_{\Gamma\in\mathbf G(y,x)})$ (RG1 (0.11)), $y\in\{ c_-,c_+\}$. Then, with $C_{\mathrm G}=9K_L+2$ and $M(F,c)$ the block average of RG1 (0.4),
\begin{equation}\label{eq:c3-7-block}
    | F(b)-\id|_{\sigma}\leqslant C_{\mathrm G}\eta\ \text{ for a block-internal bond }b,\qquad
    | F(b)-M(F,c)|_{\sigma}\leqslant C_{\mathrm G}\eta\ \text{ for an interface bond }b\in B(c)\;.
\end{equation}
\end{lemma}

\begin{proof}
Two paths of $\mathbf G(y,x)$ have transports within $K_L\eta$ of each other, so the group-mean estimate with $\rho= K_L\eta$ puts each $F(\Gamma)$, $\Gamma\in\mathbf G(y,x)$, within $2K_L\eta$ of $F(y,x)$, hence within $(2K_L+1)\eta$ of the identity. For an internal bond $b=(x,x')$ the loop $\Gamma_x\cup b\cup(-\Gamma_{x'})$ gives $| F(\Gamma_x)F(b)F(\Gamma_{x'})^{-1}-\id|_{\sigma}\leqslant K_L\eta$, and replacing $F(\Gamma_x)$, $F(\Gamma_{x'})$ by the identity costs $2(2K_L+1)\eta$, so $| F(b)-\id|_{\sigma}\leqslant (5K_L+2)\eta$. For an interface bond $b=(x,x')$ the loop $\Gamma_x\cup b\cup(-\Gamma_{x'})\cup(-c)$ gives in the same way $| F(b)-F(c)|_{\sigma}\leqslant (5K_L+2)\eta$, $F(c)$ the transport along the straight centre-to-centre path. By RG1 (0.4), $M(F,c)= e^{iH}F(c)$ with $H$ the average, with weights summing to one (the weight $L^{-4}$ of each site $x\in B(c_-)$ and the path weights of (0.4)), of the logarithms $(1/i)\log F(\Gamma\cup[x,x']\cup(-\Gamma')\cup(-c))$ over the sites $x\in B(c_-)$ and the path pairs $\Gamma\in\mathbf G(c_-,x)$, $\Gamma'\in\mathbf G(c_+,x')$, $[x,x']$ the transported coarse bond; each of these loops belongs to the family defining $K_L$, so each logarithm has norm at most $K_L\eta+C_{\log}(K_L\eta)^2\leqslant 2K_L\eta$ and is Hermitian, so $| M(F,c)-F(c)|_{\sigma}\leqslant | H|_{\sigma}\leqslant 2K_L\eta\leqslant 4K_L\eta$. Together $| F(b)-M(F,c)|_{\sigma}\leqslant (9K_L+2)\eta$.
\end{proof}

\begin{lemma}[witnesses for bad tests]\label{lem:c3-7-seed}
Let $k_N=\sqrt N$ (so $|\bigcdot |_{\sigma}$ and the norms of AVG and VAR are within factors $k_N$ of each other: each of these is the operator norm or the Hilbert--Schmidt norm in one of its two normalisations, and $h\leqslant |\bigcdot |_{\sigma}\leqslant \|\bigcdot \|_{\mathrm{op}}\leqslant \sqrt N\,h$ with the unnormalised Hilbert--Schmidt norm equal to $\sqrt N\,h$, the conventions of Section~\ref{sec:sourced-effective-interaction}, so any two of the three norms are within a factor $\sqrt N$ of each other), let $\alpha=2k_N^2L^2\theta$ with $\theta$ of \eqref{eq:c3-7-4}, and let $g_0$ be small enough that $LM_2R_1>4LM_1+4L+8$, $\theta\leqslant \eta_0$, $2B_3k_N^3L^2\theta\leqslant \eta_0$, $k_N\theta\leqslant c_2'$ and $C_0L^2k_N\theta\leqslant 1$ (the constants of AVG Proposition~1, (51)) and $\alpha\leqslant a_1$ (the constant of VAR Theorem~1; then $B_3\alpha\leqslant B_3a_1\leqslant a_0$); every condition is an upper bound on $\theta$ or a lower bound on $R_1$, and all hold for small $g_0$ since $\theta\leqslant \varepsilon_0\to0$ and $R_1\to\infty$ (CRE (2.4)--(2.8), the definition of $R_1$). For a cube $Q$ let $N(Q)$ be the set of fine sites within distance $6\ell$ of $Q$. On the support of the averaging constraint $\delta(\bar UV^{-1})$ of CRE (1.6), (1.11), for the gauge-fixed field $\widetilde U$ of that integrand:
(a) a $P_0$ test is good whenever every plaquette of its defining enlargement has $| U(\partial p)-\id|_{\sigma}<\theta$;
(b) if every plaquette $p\subseteq N(Q)$ has $| U(\partial p)-\id|_{\sigma}<\theta$, where $Q$ is any cube of CRE's partition into $L^2M_2R_1$-cubes (in particular any cube carrying the test (1.4); the position of $Q$ is not used), the $P_1$ test of $Q$ is good, and $| U_{1,Q}(\partial p)-\id|_{\sigma}\leqslant 2B_3k_N^3\theta\leqslant \varepsilon_1L^{-2}/2$ on $Q_w^{[1]}$ and on $Q_w^{[2]}$, its incident bonds and their block stencils (for a bond $b$, its block stencil is the block or the pair of adjacent blocks containing its endpoints, with the plaquettes having their four vertices in their union and the averaged contour variables of their centres; these are the sets to which Lemma~\ref{lem:c3-7-block} is applied in (c));
(c) if moreover every active averaged contour variable $w_{yx}=\widetilde U(y,x)$ with $x\in N(Q)$ has $| w_{yx}-\id|_{\sigma}<\theta$, then, for a $Q_1$- or $R_1$-eligible cube $Q$,
\begin{equation}\label{eq:c3-7-seed}
    \sup_{b\in(Q_w^{[2]})^*}|\widetilde U(b)U_{1,Q}(b)^{-1}-\id|_{\sigma}\leqslant C_Q\theta\leqslant \frac{\delta_0}{2}<2\delta_0\;,
\end{equation}
so the $Q_1$ test of $Q$ is good. The $R_1$ test is treated in the proof of Lemma~\ref{lem:c3-7-2}. The hypotheses enter only through the stated upper bounds, so the lemma holds verbatim with $\theta$ replaced by any $\theta'\leqslant \theta$ and $C_Q\theta$ by $C_Q\theta'$ in \eqref{eq:c3-7-seed}.
\end{lemma}

\begin{proof}
(a) is $\theta\leqslant \varepsilon_0$. Collar payments: $w\leqslant a= L^2M_2R_1$ and $M\geqslant 8LM_2$ give $4w+2L+4\leqslant 4a+2L+4<\ell$, and $w\geqslant LM_2R_1>4LM_1+4L+8$ gives $2w-2LM_1>2L+4$. The $\sim$ operations on $P_0'$, $P_1'$ and $B(P_1^1)^c$ are by $LMR_1$-cubes (CRE pp.246 and 248), so the first inequality places $Q_w^{[4]}$ and the blocks meeting it inside $B(P_1^1)$ for a $Q_1$-eligible cube $Q\subseteq B(P_1^1)^{\sim-1}$ (CRE (1.8)) and, by CRE (1.10) ($\Omega_1\subseteq B(P_1^1)^{\sim-3}$), for an $R_1$-eligible cube $Q\subseteq\Omega_1^{\sim-1}$ as well; the second places $Q_w^{[2]}$, its incident bonds and their block stencils at distance more than $2L+4$ inside $Q_w^{[4]}$ eroded by $2LM_1$. All these sets lie in $N(Q)$, since they are within $4w+2L+4\leqslant 4a+2L+4<\ell<6\ell$ of $Q$. The placement inside $B(P_1^1)$ is used in (c) only; (b) uses the hypothesis on $N(Q)$ alone.

(b) The block average is gauge covariant, $M(U^v,c)= v(c_-)M(U,c)v(c_+)^{-1}$ for a gauge transformation $v$: in RG1 (0.4) every loop is based at $c_-$, so its holonomy is conjugated by $v(c_-)$, the chart logarithm is conjugation covariant, and $U(c)$ is a transport from $c_-$ to $c_+$. So the coarse plaquette deviations $|\bar U(\partial p')-\id|_{\sigma}$ of $\bar U= M(U,\bigcdot )$ coincide for $U$ and its gauge transform $\widetilde U$, as do the fine plaquette deviations, and the plaquette hypothesis on $N(Q)$ may be read for either field. AVG Proposition~1, (51), in the averaging norm (at most $k_N|\bigcdot |_{\sigma}$) with $\alpha_0= k_N\theta\leqslant c_2'$ gives $|\bar U(\partial p')-\id|\leqslant L^2k_N\theta+C_0(L^2k_N\theta)^2\leqslant 2L^2k_N\theta$ for every coarse plaquette $p'$ whose blocks lie in $N(Q)$; a further factor $k_N$ converts this to the VAR norm, so on the support of $\delta(\bar UV^{-1})$ the coarse data $V$ on $Q_w^{[4]}$ satisfy the VAR data bound $\alpha=2k_N^2L^2\theta$. The determining data of $U_{1,Q}$ are those of CRE (1.2): the minimal determining set $\mathbf B_1(Q_w^{[4]})$ with the data $M^{\bigcdot }(Q_1^{s*}V)$, no exterior data. The extension (1.3) has identity block-internal links and interface links $V(c)$; its fine plaquettes are the identity (two internal and two interface links across one coarse bond) or coarse plaquette holonomies (four interface links at a corner of four blocks), and its block averages are $V$ (the loops of RG1 (0.4) are the identity on it). So the ordinary and mixed determining plaquettes of VAR (7) obey the bound $\alpha\leqslant a_1$, and VAR Theorem~1 gives a minimal critical orbit in the class (8) with bound $B_3\alpha$. CRE defines $U_{1,Q}$ (p.256, (2.12)) as an element of the unique minimal orbit of the local action over the regular configurations $U$ with $M_{\mathbf B}(U)= V$, and refers for ``regular'' and for the variational theory to VAR; we read ``regular'' as membership in the class VAR (6), the space (2) of VAR with a radius $\varepsilon$ satisfying $B_3\alpha\leqslant \varepsilon\leqslant a_0$, the setting of VAR Theorem~1 (such $\varepsilon$ exist, as $B_3\alpha\leqslant B_3a_1\leqslant a_0$). Theorem~1 states that the minimal orbit of the class (8) is the unique critical orbit of the class (6); a minimal configuration is critical, so the minimal orbit over the regular configurations is this orbit, $U_{1,Q}$ satisfies the bounds (8) with VAR's parameter (written $\varepsilon_1$ in VAR; not CRE's $\varepsilon_1= g_1p_0(g_1)$) equal to $\alpha$, and the axial gauge condition (CRE p.247) fixes the representative. CRE's construction of $\mathbf B_1$ (p.256--257, after (2.12); $Q_w^{[4]}$ is a union of $w$-cubes, $w\in\{ LM_2R_1,L^2M_2R_1\}$ a multiple of $LM_1$, hence a union of $LM_1$-cubes and of $M_1$-cubes, $M_1< M_2$ being powers of $L$) gives that the coarse determining interior contains $(Q_w^{[4]})^{\sim-2}$ with the erosion in $M_1$-cubes of $T_{L^{-1}}$, of side $LM_1$ in fine units, so the loss is at most $2LM_1$, and by the collar payment $Q_w^{[2]}$, its incident bonds and their block stencils lie inside it. There VAR (8) bounds the fine plaquettes of the minimum by $B_3\alpha L^{-2}$ in the VAR norm, that is by $k_NB_3\alpha L^{-2}=2B_3k_N^3\theta\leqslant \varepsilon_1L^{-2}/2$ in $|\bigcdot |_{\sigma}$ by \eqref{eq:c3-7-4}; on $Q_w^{[1]}$ this is the good $P_1$ test of CRE (1.4).

(c) The blocks meeting $(Q_w^{[2]})^*$ are active by the collar payments. Apply Lemma~\ref{lem:c3-7-block} to $\widetilde U$ with $\eta=\theta$ (plaquettes and averaged contour variables below $\theta$) and to $U_{1,Q}$ with $\eta=2B_3k_N^3L^2\theta\leqslant \eta_0$ (its plaquette deviations are at most $2B_3k_N^3\theta$ by (b), and its averaged contour variables are the identity in the axial gauge). On the support of the averaging constraint $M(\widetilde U,c)= V(c)$, and $M(U_{1,Q},c)= V(c)$ on the coarse bonds of the determining interior (the constraint of the local problem, data (1.3)). For an internal bond compare both fields with the identity, for an interface bond both with $V(c)$: $|\widetilde U(b)U_{1,Q}(b)^{-1}-\id|_{\sigma}=|\widetilde U(b)-U_{1,Q}(b)|_{\sigma}\leqslant C_{\mathrm G}\theta(1+2B_3k_N^3L^2)= C_Q\theta\leqslant \delta_0/2$ by \eqref{eq:c3-7-4}, the good $Q_1$ test of CRE (1.8). The final sentence of the statement holds because every step used $\theta$ only through the hypotheses and through the premise on the witnesses.
\end{proof}

\begin{lemma}[a bad test forces a witness]\label{lem:c3-7-2}
For $g_0$ below a threshold depending on $L$, the group, $M_1$, $M_2$, $M$ and the source constants, on the original cutoff support: every bad $P_0$, $P_1$, $Q_1$ or $R_1$ test forces a plaquette or gauge-variable witness at threshold $\theta$ within distance $K\ell$ of every bond $z\in E(B_{\xi})$ whose membership it causes, with the fixed $K=256$; and the number of plaquette and gauge-variable labels in such a neighbourhood is at most
\begin{equation}\label{eq:c3-7-5}
    n_W=7(2K\ell+7)^4\;.
\end{equation}
\end{lemma}

\begin{proof}
The $P_0$, $P_1$ and $Q_1$ alternatives are Lemma~\ref{lem:c3-7-seed}(a)--(c). For the $R_1$ test, Lemma~\ref{lem:c3-7-seed}(c) is combined with the comparison of the two minima on the twice enlarged carrier of every eligible cube $Q$, uniformly in its label, proved in Proposition~\ref{prop:c3-9-1} below:
\begin{equation}\label{eq:c3-7-6}
    \sup_{b\in(Q_w^{[2]})^*}| U_1(b)U_{1,Q}(b)^{-1}-\id|_{\sigma}\leqslant C_{L,G}\varepsilon_1e^{-c_{L,G}LM_2R_1}<\delta_0/10\;.
\end{equation}
Of the characteristic factors of CRE (1.11) it uses only their nonvanishing at $U_1$ itself, the membership $U_1\in D^{(11)}(V,V_0)$ of (B1) (the reading of CRE p.248 recorded there and in (H4)); it does not assume that $U_1$ lies in the support of the fluctuation cutoffs (1.18) or (1.25) at zero fluctuation. On the native principal chart a bad $R_1$ label of the cube $Q$ forces $| g_0A(b)|_{\sigma}\geqslant \delta_0$ at a bond $b\in(Q_w^{[2]})^*$, the carrier of the test in CRE (1.18); the eigenvalue inequality $| e^{iX}-\id|_{\sigma}\geqslant (2/\pi)| X|_{\sigma}$, valid when every eigenvalue of $X$ lies in $[-\pi,\pi]$ (as on the chart) in both permitted norms, and \eqref{eq:c3-7-6} give, for the gauge-fixed field $\widetilde U= e^{ig_0A}U_1$ at the free bond $b$ (CRE (1.13): $U'= UU_1^{-1}= e^{iA'}$ with $A'= g_0A$ the fluctuation field; the variable tested in (1.18) is the free variable $A(b)$ of (1.15), and the fluctuation at a free bond is exactly $e^{ig_0A(b)}$, because the substitution (1.14), $A'= B-h_{U_1}\widetilde D(B)$ with $B= g_0\mathsf C_{U_1}A$, changes the phase only at the designated bonds: $h_{U_1}$ is supported on them and $\mathsf C_{U_1}A$ agrees with $A$ on the free bonds, \eqref{eq:constraint-extension} at the background, Section~\ref{sec:regular-hessian-stencils}),
\begin{equation}\label{eq:c3-7-7}
    |\widetilde U(b)U_{1,Q}(b)^{-1}-\id|_{\sigma}\geqslant (2/\pi-1/10)\delta_0>\delta_0/2\;.
\end{equation}
If every witness in $N(Q)$ were below $\theta$, \eqref{eq:c3-7-seed} (Lemma~\ref{lem:c3-7-seed}(c) for the $R_1$-eligible cube $Q$; the earlier cutoffs and $\Omega_1$ are fixed term by term in the sum (1.11), which resums $P_0$, $P_1$, $Q_1$ at fixed $\Omega_1$, and $\Omega_1\subseteq B(P_1^1)^{\sim-3}$ holds for each term; the comparison \eqref{eq:c3-7-6} holds whatever the $R_1$ label of $Q$) would give $|\widetilde U(b)U_{1,Q}(b)^{-1}-\id|_{\sigma}\leqslant C_Q\theta\leqslant \delta_0/2$, a contradiction; so the label is excluded. CRE (1.1), (1.4), (1.8), (1.18) are sharp characteristic partitions at their printed thresholds, so these inequalities make the corresponding bad factor vanish; off-chart configurations are not added to (1.18), which is introduced after the logarithmic reconstruction (1.13)--(1.15). For the distances: before filling, CRE (1.10), (1.20) give the bad region as $Q_1'^{\sim4}\cup(B(P_1^1)^c)^{\sim5}\cup R_1'^{\sim1}$, with $B(P_1^1)^c\subseteq P_0'^{\sim2}\cup P_1'^{\sim1}$ (the definition of $P_1^1$ following CRE (1.4), as printed, gives $B(P_1^1)^c= P_0'^{\sim2}$; the reading consistent with the active-region construction of CRE (3.2)--(3.3) gives $P_0'^{\sim2}\cup P_1'^{\sim}$; both are contained in $P_0'^{\sim2}\cup P_1'^{\sim1}$), so a point of this region is within $10\ell$ of a bad test. The modification after (1.20) fills only components contained in cubes of side less than $100\ell$, so every added point is within $100\ell$ of an original unfilled bad point; this also holds if small merged components are filled again, since by induction every current component contains an original bad point and each fill is performed inside a cube of side less than $100\ell$ containing that component (every point of the rectangular hull of such a component is within $100\ell$ of a chosen original bad point of it, and overlapping hulls or later merges preserve the presence of an original point), so there is no cumulative distance loss. CRE (2.9) gives, for CRE's scale $R_0$ (the range of CRE's expansions; not this paper's fixed range constant $R_0$ of \eqref{eq:regional-collar-sizes}), $R_0\leqslant (L+1)R_1$, hence $R_0\leqslant LR_1$ as both are powers of $L$ (the conventions of Section~\ref{sec:sourced-effective-interaction}). Including each test's witness neighbourhood and incident bonds, every $z\in E(B_{\xi})$ forces a witness within $K\ell$ of $z$ with $K=256$: a point of the filled bad region is within $100\ell$ of an original bad point, which is within $10\ell$ of a bad test cube $Q$ of side at most $a\leqslant \ell/8$, whose witness lies in $N(Q)$, within $6\ell$ of $Q$; an incident bond adds one; in all $100\ell+10\ell+a+6\ell+1\leqslant 117\ell<256\ell$. The count \eqref{eq:c3-7-5} of the plaquette and gauge-variable labels in that neighbourhood is an overcount on a small periodic torus, which is harmless.
\end{proof}

\subsubsection{The marked probability retains the physical normaliser}\label{sec:c3-7-2}

\begin{lemma}[the marked density and its reversal]\label{lem:c3-7-3}
Let $A^{\Psi}_{\mathrm e}=\Psi_{\xi}(A_{\mathrm e})$ and let $\mathscr E_{\xi}(B)$ be the pre-boundary-change exponent of CRE (1.15), its outer prefactor excluded (\eqref{eq:native-collar-full-exponent}). With $c_{\xi}$ of \eqref{eq:native-collar-boundary-constant}, pointwise on $\mathfrak g^I$,
\begin{equation}\label{eq:c3-7-8}
    e^{c_{\xi}}K_{\xi;1}(a)=\mathbf 1_{\mathcal D_I}(a)\,|\det D\Psi_{\xi}(A_{\mathrm e})|\,\exp\mathscr E_{\xi}\bigl(\Psi_{\xi}(a+A_{\mathrm e})\bigr)\;,
\end{equation}
and the composed retained cutoff $\chi_{\xi,\mathrm e}$ and $\zeta_{\xi}$ multiply both sides; integrating over $a$,
\begin{equation}\label{eq:c3-7-8-weight}
    \mathsf w_{\xi}= e^{c_{\xi}}I_{\xi;1}=\int e^{c_{\xi}}K_{\xi;1}(a)\,\dd a\;,
\end{equation}
the branch constant $c_{\xi}$ and every far primitive present (\eqref{eq:c-collar-identity}). The reversal of the changes of variables carried out in the proof below is a change of variables for the integral of the density times any nonnegative measurable function of $U$, of the regions and of the labels of a term, performed term by term before the characteristic partitions are summed; it is stated here for the mark $\mathbf 1\{ z\in E(B_{\xi})\}$ and is used in that general form in Proposition~\ref{prop:c5-2-6} and Lemma~\ref{lem:c5-7-joint}. Together with Lemma~\ref{lem:c3-7-2}, on compatible even fine periods, for every positive bond $z$,
\begin{equation}\label{eq:c3-7-9}
    p_z\leqslant \bigl(\mu_0\otimes q_g^{\otimes\mathcal I}\otimes\mathrm{Haar}^{\otimes T}\bigr)(W_z)\;,\qquad
    \mu_0(\dd U)=\mathcal Z_0^{-1}e^{-\mathsf E}e^{-S(U)/g_0^2}\,\dd U\;,
\end{equation}
where $W_z$ is the event that some plaquette or some auxiliary gauge variable in the deterministic $K\ell$-neighbourhood of $z$ is a witness at threshold $\theta$, $\mathcal I$ is the set of pairs $(y,x)$ of a block centre $y$ and a noncentral site $x$ of its block, $T$ the set of centres and unused sites, and
\begin{equation}\label{eq:c3-7-10}
    q_g(\dd w)= z_g^{-1}\mathbf 1\{| w-\id|_{\sigma}<\varepsilon_0\} e^{-s_G(w)/g_0^2}\,\dd w\;,\qquad s_G(w)=1-\operatorname{Re}\tr_Nw\;,
\end{equation}
with $z_g$ the integral of the displayed density, CRE's gauge normaliser.
\end{lemma}

\begin{proof}
$\Psi_{\xi}$ is the identity in every near stencil and its support $S_1^*$ is disjoint from $I$, so $\Psi_{\xi}(a+A_{\mathrm e})= a+A^{\Psi}_{\mathrm e}$ (\eqref{eq:c-bondwise-shift}). Expanding the quadratic form of $\mathscr E_{\xi}$ into its $II$, $IJ$, $JJ$ blocks, leaving every far primitive at $A^{\Psi}_{\mathrm e}$ and retaining the boundary determinant gives \eqref{eq:c3-7-8}, which is the computation of Lemma~\ref{lem:collar}(ii); the retained cutoff and $\zeta_{\xi}$ are factors of the integrand outside the $a$-integral. Together with the outer prefactor $\kappa_{\xi}$ (\eqref{eq:native-source-outer-prefactor}) the right side recombines into the transformed CRE (1.15) density, boundary Jacobian included. Now reverse the changes of variables of Section~\ref{sec:branch-data} with the mark $\mathbf 1\{ z\in E(B_{\xi})\}$ inserted. The reversal is performed in each term of the history sum separately, where the regions, the filling and the labels are fixed, and each change of variables is an identity between integrals of a density times a nonnegative measurable integrand (the term's partitions and cutoffs multiply the integrand and are carried along), so it holds for the mark, and equally for any nonnegative measurable function of $U$, the regions and the labels in its place; the partitions are summed only after the domination by a label-free event below. For the present mark this is also seen directly: it commutes with the earlier resummation because the final regions and the filling are fixed in each term. Undo $\Psi_{\xi}$, the free-variable constraint elimination, the scaling, the nonlinear constraint linearisation of (1.14), the logarithmic chart of (1.13) with its Haar density, and the translation about $U_1$, in this order; the retained cutoff stays composed with $\Psi_{\xi}$ and $\zeta$ remains in that composition, boundary-fluctuation dependence included. The marked integral is then CRE (1.6) with its four original partitions and the site-gauge integration, no determinant or normalisation discarded; the unmarked chain is $\mathcal Z_0= e^{-\mathsf E}\int e^{-S(U)/g_0^2}\,\dd U$. For each active block centre $y$ and noncentral $x$ let $w_{yx}$ be its mean transport after the gauge change; RG1 (0.6) and CRE (1.5) give the Haar substitution $v(x)= w_{yx}^{-1}h_yU(y,x)$, $v(y)= h_y$, with Jacobian one, and the law of $w$ is \eqref{eq:c3-7-10}. Add independent $q_g$ variables of integral one for the inactive indices and independent Haar variables for the centres and unused sites. The earlier $P_0$, $P_1$ choices depend on $U$ and its coarse gauge orbit, not on these auxiliary variables. By Lemma~\ref{lem:c3-7-2} the mark is dominated by $\mathbf 1_{W_z}$, and $W_z$ does not depend on the later $Q_1$, $R_1$ labels; only after this enlargement are the characteristic partitions summed to one. This gives \eqref{eq:c3-7-9}. The product measure is a representation of the integral before the history conditions, not an independence statement about them; plaquette witness norms and the conjugation-invariant gauge-variable norms are gauge-invariant scalars under the simultaneous gauge transformations, and no open Wilson contraction is estimated.
\end{proof}

\subsubsection{One-point tails and summation}\label{sec:c3-7-3}

\begin{lemma}[single-witness tails]\label{lem:c3-7-4}
Let $d_G=\dim G$, $b_G=1/(2N)$, so that $s_G(w)\geqslant b_G| w-\id|_{\sigma}^2$ ($s_G(w)=1-\operatorname{Re}\tr_Nw$ with the normalised trace $\tr_N= N^{-1}\operatorname{Tr}$, \eqref{eq:c3-7-10}; the action is $S(U)=\sum_ps_G(U(\partial p))$ over the unoriented plaquettes, \eqref{eq:bare-action}, so that $\mathcal Z_0=\int\rho_0\,\dd U= e^{-\mathsf E}\int e^{-S/g_0^2}\,\dd U$ by \eqref{eq:one-step-physical-partition} with $\mathsf E= E_{\varepsilon}$ the vacuum scalar of \eqref{eq:bare-action}; $\dd U$, $\dd w$ are normalised Haar measures, and $\mu_0$ and $q_g$ are probability measures by their normalisations $\mathcal Z_0$ and $z_g$ \eqref{eq:c3-7-10}). There are constants $c_H$, $g_H>0$ depending on the group only such that the Haar measure of the operator ball of radius $g$ about $\id$ is at least $c_Hg^{d_G}$ for $0< g\leqslant g_H$: in the chart, $\dd w=\sigma(A)\,\dd A$ with the continuous Haar density $\sigma$, $\sigma(0)=\sigma_0>0$ (the steps listed after \eqref{eq:native-source-physical-ratio}); choose $g_H<\varrho_G$ with $\sigma\geqslant \sigma_0/2$ on $\{\| A\|_{\mathrm{op}}\leqslant g_H\}$, a subset of the chart ball as $h\leqslant \|\bigcdot \|_{\mathrm{op}}$; for $g\leqslant g_H$ the map $A\mapsto e^{iA}$ sends $\{\| A\|_{\mathrm{op}}< g\}$ injectively into the operator ball of radius $g$ ($\| e^{iA}-\id\|_{\mathrm{op}}\leqslant \| A\|_{\mathrm{op}}$), so that ball has Haar measure at least $(\sigma_0/2)\operatorname{Leb}\{\| A\|_{\mathrm{op}}<1\}\,g^{d_G}\eqqcolon c_Hg^{d_G}$. Let $g_0\leqslant g_H$ be small enough that the operator ball of radius $g_0$ lies in $\{| w-\id|_{\sigma}<\varepsilon_0\}$; this holds for small $g_0$, since the two norms are comparable with fixed constants and $\varepsilon_0/g_0= A_0(\log g_0^{-2})^{p_0}\to\infty$, $A_0>0$ (\eqref{eq:cre-cutoff-normalisation}, CRE p.246: $\varepsilon_0= g_0A_0(\log g_0^{-2})^{p_0}$). With $V_T$ the number of fine sites,
\begin{equation}\label{eq:c3-7-11}
    \mathcal Z_0\geqslant e^{-\mathsf E}(c_Hg_0^{d_G})^{4V_T}e^{-48V_T}\;.
\end{equation}
On compatible even fine periods, for every threshold $\theta'\in(0,\varepsilon_0]$, every plaquette $p$ and every auxiliary gauge variable $w$,
\begin{gather}
    \mu_0\bigl(| U(\partial p)-\id|_{\sigma}\geqslant \theta'\bigr)\leqslant z_P(\theta')\coloneqq \exp\Bigl[-\frac{b_G\theta'^2}{4g_0^2}+4d_G\log g_0^{-1}+4\log c_H^{-1}+48\Bigr]\;,\label{eq:c3-7-12}\\
    q_g\bigl(| w-\id|_{\sigma}\geqslant \theta'\bigr)\leqslant z_G(\theta')\coloneqq e^{1/2}(c_Hg_0^{d_G})^{-1}e^{-b_G\theta'^2/g_0^2}\;,\label{eq:c3-7-13}
\end{gather}
and therefore, with $z_P\coloneqq z_P(\theta)$ and $z_G\coloneqq z_G(\theta)$ at the threshold $\theta\leqslant \varepsilon_0$ of \eqref{eq:c3-7-4} and for $g_0$ below the threshold of Lemma~\ref{lem:c3-7-2} (which \eqref{eq:c3-7-9} inherits through Lemma~\ref{lem:c3-7-3}), $p_z\leqslant q\coloneqq n_W\max(z_P,z_G)$ for every positive bond $z$.
\end{lemma}

\begin{proof}
On the event that every fine link lies in the operator ball of radius $g_0$ every plaquette has $\| U(\partial p)-\id\|_{\mathrm{op}}\leqslant 4g_0$ (telescoping over its four links) and contributes $s_G(U(\partial p))=\frac1{2N}\sum_j| e^{i\varphi_j}-1|^2\leqslant \frac12\| U(\partial p)-\id\|_{\mathrm{op}}^2\leqslant 8g_0^2$ to $S$ ($e^{i\varphi_j}$ the eigenvalues of $U(\partial p)$), and there are at most $6V_T$ plaquettes and $4V_T$ links, which gives \eqref{eq:c3-7-11}. Fix $\theta'\in(0,\varepsilon_0]$. For \eqref{eq:c3-7-12} use reflection positivity and the chessboard estimate of $\mu_0$ on the even fine periods $N_{\mu}$. For a coordinate reflection $\vartheta$ through the planes $x_{\mu}=0$, $x_{\mu}= N_{\mu}/2$ put $(\vartheta U)(x,x')= U(\vartheta x,\vartheta x')$ (inverse values on reversed bonds; $s_G(W)$ and $| W-\id|_{\sigma}$ are invariant under $W\mapsto W^{-1}$ and under conjugation, so the plaquette terms of $S$ and the plaquette indicators below do not depend on the orientation or the base point of a plaquette) and $\Theta F= F\circ\vartheta$. Every plaquette lies in one of the two closed half-tori; write $U_{\partial}$ for the links in the two planes, $U_+$ for the remaining links of one half, and $S_+$ for the sum of $s_G(U(\partial p))$ over the plaquettes of that closed half, the plaquettes lying in the planes (functions of $U_{\partial}$) counted with weight $1/2$, so that $S= S_++\Theta S_+$. For a bounded $F$ of the links $(U_+,U_{\partial})$, since $\Theta$ fixes $U_{\partial}$ and carries the Haar measure of $U_+$ to that of the other half,
\begin{equation}\label{eq:c3-7-rp}
    \int\overline F\,\Theta F\,\dd\mu_0=\frac{e^{-\mathsf E}}{\mathcal Z_0}\int\Bigl|\int Fe^{-S_+/g_0^2}\,\dd U_+\Bigr|^2\,\dd U_{\partial}\geqslant 0\;.
\end{equation}
Let $C_x= x+[0,1]^4$, $x$ a fine site, be the closed unit cubes and, for a bounded function $f$ of the links of $C_0$, let $f^{[x]}$ be its transport to $C_x$ by successive reflections (the four coordinate reflections commute and the periods are even, so $f^{[x]}$ depends neither on the order of the directions nor on the winding of the path of reflections; and the reflection through any site plane $x_{\mu}= m$ carries the transport to the slab at position $k$ in direction $\mu$ to the transport to position $2m-1-k$, since every composition of integer-plane reflections carrying position $0$ to position $k$ has orientation $(-1)^k$ in that direction and so coincides with the defining transport). For nonnegative bounded functions $f_x$ of the links of $C_0$, indexed by the fine sites $x$,
\begin{equation}\label{eq:c3-7-chessboard}
    \int\prod_{x}f_x^{[x]}\,\dd\mu_0\leqslant \prod_{x}\Bigl(\int\prod_{y}f_x^{[y]}\,\dd\mu_0\Bigr)^{1/V_T}\;,
\end{equation}
the chessboard estimate. It is proved here, by the argument of \cite[Theorem~10.11, p.~445]{FriedliVelenik} with \eqref{eq:c3-7-rp} as the reflection positivity; that theorem is stated for a spin system on the sites with reflections through planes between sites, and is not applied as stated to the link variables with the site reflections used here. In one direction, with $N_{\mu}=2h$: by translation invariance of $\mu_0$, \eqref{eq:c3-7-rp} holds for the reflection through any pair of antipodal site planes $x_{\mu}= k$, $x_{\mu}= k+h$; let $\Phi(f_0,\ldots,f_{2h-1})$ be the $\mu_0$-expectation of the reflected product of functions of the $2h$ closed unit slabs (slab $j$ carrying $f_j^{[j]}$, the transport of $f_j$ to slab $j$); adjacent slabs share only plane links; the two planes bounding $h$ consecutive slabs are fixed by the reflection through them and the plane links interior to these slabs belong to $U_+$, so the product over any $h$ consecutive slabs is a function of $(U_+,U_{\partial})$ for the reflection through the planes bounding it, and \eqref{eq:c3-7-rp} with the Cauchy--Schwarz inequality for the sesquilinear form $(F,F')\mapsto\int\overline F\,\Theta F'\dd\mu_0$ bounds $\Phi^2$ by the product of the two values of $\Phi$ obtained by reflecting either half of the tuple onto the other. Add a positive constant to each function and divide it by the $2h$-th root of its fully disseminated expectation (the normalised functions $\widetilde f_j$ have $\Phi(\widetilde f_j,\ldots,\widetilde f_j)=1$); among the finitely many tuples with entries from $\{\widetilde f_0,\ldots,\widetilde f_{2h-1}\}$ take a maximiser of $\Phi$ with a longest constant run (runs counted cyclically), of length $r$; both tuples obtained by reflecting either half onto the other are maximisers as well (the maximum is positive, every function being bounded below by its added constant, and the Cauchy--Schwarz bound, the square of the maximum being at most the product of the two reflected values, forces both to equal the maximum). If $r\leqslant h$, reflecting the half that ends with the run onto the other half gives a maximiser with a constant run of length $2r> r$, contradicting the choice of $r$; so $r> h$, a half lying inside the run is reflected, and the resulting maximiser is constant: the maximum is $\Phi(\widetilde f_j,\ldots,\widetilde f_j)=1$ for some $j$, which gives the one-dimensional inequality for the normalised functions; letting the added constants decrease to zero (dominated convergence) gives it for the $f_j$. Applying it successively in the four directions (the reflections commute, and the slab functions produced after disseminating in the earlier directions are again nonnegative and bounded) gives the exponent $1/V_T$. Every plaquette is a translate of a plaquette of $C_0$ and $\mu_0$ is translation invariant, so it suffices to apply \eqref{eq:c3-7-chessboard} with $f_0=\mathbf 1\{| U(\partial p)-\id|_{\sigma}\geqslant \theta'\}$ for a plaquette $p$ of $C_0$ and $f_x=1$ for $x\neq0$: the dissemination $\{ f_0^{[y]}\}_{y}$ contains $V_T/4$ distinct plaquettes of the plane orientation $(\mu\nu)$ of $p$ (the two transverse coordinates of a copy range over one parity class; repeated copies of an indicator are harmless), each with $s_G\geqslant b_G\theta'^2$ on the event, and $S\geqslant 0$ termwise, so the disseminated expectation is at most $e^{-\mathsf E}e^{-b_G\theta'^2V_T/(4g_0^2)}/\mathcal Z_0$; dividing by \eqref{eq:c3-7-11} and taking the $1/V_T$-th root gives $z_P(\theta')$, a normalised probability estimate. For \eqref{eq:c3-7-13} restrict one gauge variable to its operator ball of radius $g_0$, where its action is at most $g_0^2/2$ and its cutoff factor is one (the ball lies in $\{| w-\id|_{\sigma}<\varepsilon_0\}$ by the hypothesis on $g_0$), which bounds $z_g$ from below by $c_Hg_0^{d_G}e^{-1/2}$; on the event $| w-\id|_{\sigma}\geqslant \theta'$ the action is at least $b_G\theta'^2$, which gives $z_G(\theta')$. Neither step uses anything about $\theta'$ beyond $0<\theta'\leqslant \varepsilon_0$ ($q_g$ is carried by the ball $| w-\id|_{\sigma}<\varepsilon_0$, so for larger thresholds the event in \eqref{eq:c3-7-13} is empty). The union bound over the at most $n_W$ witness labels of Lemma~\ref{lem:c3-7-2}, without any independence between witnesses, and \eqref{eq:c3-7-9} at the threshold $\theta$ give $p_z\leqslant q$.
\end{proof}

\begin{proposition}[the exterior-response and uncovered term]\label{prop:c3-7-5}
Let $\pi_*=\min(p_0(g_0),p_0(g_1))$. There are $c_{\theta}$, $c_*$, $g_*>0$, fixed after $L$, the group, the character, $M_1$, $M_2$, $M$ and the source constants, such that $\theta\geqslant c_{\theta} g_0\pi_*$ and $q\leqslant e^{-c_*\pi_*^2}$ for $0< g_0< g_*$, and on compatible even fine periods, for $0< g_0< g_*$ and under \eqref{eq:c3-1-2},
\begin{equation}\label{eq:c3-7-14}
    | E_{\mathrm{bdry}}|\leqslant C_{\mathrm{resp},\chi}\,n_C(1+g_0r)\,e^{-c_*\pi_*^2}\;,
\end{equation}
after replacing $C_{\mathrm{resp},\chi}$ by $C_L\max(C_{\mathrm{resp},\chi}S_c,1)$; $g_*$ is at most the thresholds of Lemma~\ref{lem:c3-7-2}, the $g_H$ and the smallness condition of Lemma~\ref{lem:c3-7-4}, the threshold of Proposition~\ref{prop:c3-9-1} and the coupling below which the CRE (2.6) comparison used in the proof holds. The constants are uniform in the compatible even periods, the contour location and the first-step histories; the factor $n_C$ is retained.
\end{proposition}

\begin{proof}
CRE (2.6)--(2.9) and $\delta_0=(A_1/A_0)\varepsilon_0$ give $\theta\geqslant c_{\theta} g_0\pi_*$ for small coupling: each of the three terms of \eqref{eq:c3-7-4} is at least a fixed constant times $g_0p_0(g_0)$ or $g_1p_0(g_1)$, and $g_1/g_0$ is bounded above and below by CRE (2.6). Since $\ell= LMR_1$ grows as a fixed power of $\log g_1^{-2}$ (\eqref{eq:c5-3-1}), $\log n_W= O(\log\log g_0^{-2})$; with CRE's exponent $p_0\geqslant 5\mathfrak r\geqslant 10$, $\pi_*^2\geqslant A_0^2(\tfrac12\log g_0^{-2})^{2p_0}$ dominates $4d_G\log g_0^{-1}$ and $\log n_W$ for small $g_0$, so $q= n_W\max(z_P,z_G)\leqslant e^{-c_*\pi_*^2}$ with $c_*= b_Gc_{\theta}^2/8$, for $g_0< g_*$. Then Corollary~\ref{cor:c3-6-1}, Lemma~\ref{lem:c3-7-1} and Lemma~\ref{lem:c3-7-4} give $| E_{\mathrm{bdry}}|\leqslant C_{\mathrm{resp},\chi}g_0r\,C_LS_cn_Cq+C_Ln_Cq\leqslant C_L\max(C_{\mathrm{resp},\chi}S_c,1)\,n_C(1+g_0r)\,q$, and $q\leqslant e^{-c_*\pi_*^2}$ gives \eqref{eq:c3-7-14}.
\end{proof}

\begin{remark}[scope]\label{rem:c3-8-1}
Lemma~\ref{lem:exterior-response-bound} is proved from the Hessian, Schur, enlarged-chart and primitive derivative bounds of Appendix~\ref{app:background-estimates} and the collar cutoff of Section~\ref{sec:boundary-law-collar}, with both hard-cutoff surface terms inside the bounded covariance; its physical caller is restricted to $A_{\mathrm{cov}}$. Proposition~\ref{prop:c3-7-5} holds on compatible even-period tori, from the single-witness tails of Lemma~\ref{lem:c3-7-4} alone, without a component probability theorem, and uses the reconstruction and measure identities of Section~\ref{sec:sourced-effective-interaction}; the marked density \eqref{eq:c3-7-8} is displayed without equating any older two-parameter normaliser with $I_{\xi;1}$, and the filling proximity of Lemma~\ref{lem:c3-7-2} holds for a simultaneous or an iterated filling of small components. The $R_1$-label step of Lemma~\ref{lem:c3-7-2} rests on Proposition~\ref{prop:c3-9-1}, which uses local VAR data, with the local hypotheses of VAR's existence and uniqueness theorem verified where this paper's inputs allow and otherwise taken over as CRE's application of VAR to its determining sets (p.256), uses VAR's derivative-decay results and CRE's identification (3.17)--(3.19) as stated, and accepts Ba{\l}aban's gauge-conversion results at their stated regular-chart hypotheses through RG1's extension statements; of the characteristic factors of CRE (1.11) it uses their nonvanishing at the minimum $U_1$ itself, the membership $U_1\in D^{(11)}(V,V_0)$ of (B1) (the reading of CRE p.248 recorded in (B1)(iii) and (H4)), and nothing about the fluctuation cutoffs (1.18), (1.25) at zero fluctuation. This appendix controls the exterior-response and uncovered term only: it does not control the joint defects $R_{\xi}(\mathfrak B)$ of \eqref{eq:part1-joint-defect} and does not identify the regional scalar $S^{\circ}_{\xi,C}$. Section~\ref{sec:one-step-insertion} replaces Proposition~\ref{prop:c3-7-5} on the labelled route only (at the original carrier the proposition stands and bounds $E_{\mathrm{bdry}}$ of \eqref{eq:c3-6-2} in Theorem~\ref{thm:main-theorem}(ii)): its history weight is the labelled weight of the second collar test with the retreat of $\Lambda_1$, its boundary bound is Proposition~\ref{prop:c5-6-4} with the label accounting \eqref{eq:c5-6-3}--\eqref{eq:c5-6-4}, in which the one-point cost $q$ becomes $q+N_{\mathrm{ret}}q_{\vartheta}$, the constant $C_{\mathrm{resp},\chi}$ enters its $C_1$, a prefactor $r^{d_G}$ and a family-sum factor appear, and the leading coefficient is identified at the retreated carrier $D(\Lambda)$ (Section~\ref{sec:c5-7}); Sections~\ref{sec:c3-1}--\ref{sec:c3-6}, the tails of Section~\ref{sec:c3-7} and Proposition~\ref{prop:c3-9-1} are used there as stated, with the covered event $A_{\mathrm{cov}}$ replaced by the labelled event $A^{\mathrm{fin}}_{\mathrm{cov}}$ (Section~\ref{sec:c5-3-4}(a)).
\end{remark}

\subsection{The minimum comparison from local data}\label{sec:c3-9}
This section proves the deterministic comparison \eqref{eq:c3-7-6} of the two minima from local data. Of the branch hypothesis (B1) of Section~\ref{sec:branch-data} it uses the membership $U_1\in D^{(11)}(V,V_0)$ of the minimum in CRE's domain of integration in (1.11) (CRE p.248, by Theorem 1 of VAR) with the smallness statement (i) derived from it in (B1), the averaging constraint and the source axial mean gauge, which place $U_1$ in VAR's regularity class near $\Omega_1$ and so identify it with the regular critical orbit constructed below; it does not use, and does not claim, that $U_1$ lies in the support of the fluctuation cutoffs (1.18) or (1.25) at zero fluctuation, which is false on $R_1$ (the large-fluctuation tests there exclude the configurations near $U_1$, Section~\ref{sec:branch-data}). The sources are VAR \cite{BalabanVAR}, CRE \cite{BalabanCRE}, RG1 \cite{BalabanRG1} and the averaging paper AVG \cite{BalabanAVG}; Ba{\l}aban's existence and uniqueness result (VAR Theorem 1) is used at local hypotheses verified below where this paper's inputs allow, and otherwise as CRE's application of VAR to its determining sets (p.256, where CRE states that all the results of [15] are used for its averaging); his derivative-decay results (VAR Proposition 9 with (172) and (190)), the gauge transformation of VAR Proposition 2 and CRE's identification (3.17)--(3.19) are used as stated, and his gauge-conversion results are accepted at their stated regular-chart hypotheses through RG1's extension statements; no claim to have re-proved the averaging paper is made.

\begin{proposition}[comparison of the global and local minima]\label{prop:c3-9-1}
There are constants $C_{L,G}$, $c_{L,G}>0$, depending on $L$, $G$, $M_1$, $M_2$ and the source parameters and fixed before $g_0$, independent of the fluctuation, its $R_1$ label, the exterior history and the torus, such that, on compatible even fine periods (every fine period is then at least $2\ell$, which the lifts below need), for every admissible $(\Omega_1,V,V_0)$ and every eligible cube $Q\subseteq\Omega_1^{\sim-1}$, on its twice enlarged free-bond carrier,
\begin{equation}\label{eq:c3-9-1}
    \sup_{b\in(Q_w^{[2]})^*}| U_1(b)U_{1,Q}(b)^{-1}-\id|_{\sigma}\leqslant C_{L,G}\varepsilon_1e^{-c_{L,G}LM_2R_1}<\delta_0/10\;,
\end{equation}
for $g_0$ below a threshold depending on $L$, $G$, $M_1$, $M_2$ and the source constants (the conditions $\alpha\leqslant a_1$, $B_3\alpha\leqslant a_0$, $CL^2\varepsilon_0\leqslant 1$, $\ell>\max(2L+4,RM_1)$, $LM_2R_1\geqslant 2LM_1+C_L$ and the chart reductions of the proof; this threshold is part of $g_*$ of Proposition~\ref{prop:c3-7-5}), the last inequality by Section~\ref{sec:c3-9-5}.
\end{proposition}

The fine periods are multiples of $\ell$ (compatibility, Section~\ref{sec:sourced-effective-interaction}) and, by the hypothesis of the proposition, even (the restriction of Section~\ref{sec:c3-7}), and $\ell= LMR_1$ is odd ($L$ is odd and $M$, $R_1$ are powers of $L$, \eqref{eq:c-standing-conditions}), so every compatible fine period is at least $2\ell$; since $w\leqslant L^2M_2R_1$ and $M\geqslant 8LM_2$ (\eqref{eq:c-standing-conditions}), $\operatorname{side}(Q_w^{[4]})\leqslant 9L^2M_2R_1\leqslant 9\ell/8<2\ell$, so this determining box and its fixed-range stencils have lifted copies before the next periodic image for small coupling. If $\Omega_1$ is the full torus there is no exterior collar to freeze and the same argument uses the periodic VAR problem. The quantity in \eqref{eq:c3-9-1} concerns a simultaneously gauge-covariant pair of links: their relative product transforms by conjugation at the initial endpoint, so its norm is gauge invariant; it is not an estimate on an open Wilson source. The proof occupies the rest of this section.

\subsubsection{Source inputs and the local-domain verification}\label{sec:c3-9-1}
$U_1$ is the critical minimum of CRE (1.12), the unique critical point in the domain of integration in (1.11), supplied by VAR Theorem 1 with the prescribed $V$, $V_0$ and the source axial mean gauge (CRE p.248, recorded in (B1) as $U_1\in D^{(11)}(V,V_0)$). This membership gives, by (B1)(i), $| U_1(\partial p)-\id|<\varepsilon_0$ for every plaquette $p\subset\Omega_1^{\sim}$ and the axial condition $| U_1(y,x)-\id|<\varepsilon_0$ at the block centres of $\Omega_1^{\sim(1)}$, so that $U_1$ lies in VAR's regularity class (6) with radius $C_L\varepsilon_0$ on VAR's outer domain $\Omega_0$ below, by the same conjugation and counting argument as for $\widehat U$ in the local-domain verification (VAR counts as plaquettes of a region those with a corner in it, so $\Omega_0$ is taken to be the enlarged domain $\Omega^+$ shrunk by $\ell/2$, a union of $M_1$-blocks every plaquette with a corner in which lies in $\Omega_1^{\sim}$; the remaining separation $\ell/2$ from $\Omega_1$ keeps VAR's spacing condition (1)); the localisation below verifies the hypotheses of VAR Theorem 1 for the localised problem near $\Omega_1$ and does not involve a distant large-field exterior. The plaquette variables of $U_1$ on $\Omega_1^{\sim}$ are thus within $\varepsilon_0\leqslant (1+\beta_0)^2\varepsilon_1$ of the identity (CRE (2.8)). The local minimum $U_{1,Q}$ is CRE (1.2) with determining set $\mathfrak B_Q=\mathbf B_1(Q_w^{[4]})$ and reference data $D^0_Q=\dot M(Q_1^{s*}V)|_{\mathfrak B_Q}$; it has the corresponding VAR mixed-scale regularity, with the scale-zero boundary bound described below, and the same coarse averages $V$ on the interior determining region (the interior coarse labels of $Q_w^{[4]}$ lie in $\Omega_1$, where the constraint $\mathsf M(U_1)= V$ holds: $Q\subseteq\Omega_1^{\sim-1}$ is at $\ell^{\infty}$ distance at least $\ell$ from $\Omega_1^c$, and $4w\leqslant 4L^2M_2R_1\leqslant \ell/2$ by $M\geqslant 8LM_2$, \eqref{eq:c-standing-conditions}; CRE states this containment at p.250). These are the hypotheses of CRE's derivation (3.17)--(3.19), whose only field is the minimum, not the integrated fluctuation; at $k=0$ the fine spacing is $\eta=1$ and $V^{(0)}= U_1$, $V^{(0)}_Q= U_{1,Q}$, and no preceding effective interaction is needed.

The local-domain verification. Work before the $P_0$, $P_1$, $Q_1$ resummation, at any supported history contributing to the fixed $\Omega_1$, $V$, $V_0$, with fine field $\widehat U$. Take the enlarged fine domain $\Omega^+=\Omega_1^{\sim1}$ (the macro enlargement of CRE (1.10)) and fix $\widehat U= V_0$ on the collar $\Omega^+\setminus\Omega_1$. Outside $\Omega^+$ retain the original field but omit its action terms from the variational problem: once $\ell>2L+4$, every plaquette or block-gauge term containing an unfrozen variable lies in $\Omega^+$, so the omitted terms are constants, and the constrained first variation is exactly that of CRE (1.12), with no change of boundary condition or integration law (the gauge functional $G$ of (1.12) has zero first variation at $U_1$, each of its nonnegative summands vanishing there since $U_1(y,x)=\id$ in the source axial mean gauge, CRE p.248, so the first variations of $A+G$ and of $A$ agree on the constraint tangent space at $U_1$). CRE (1.10) and the following sentence separate $\Omega_1$ from the earlier bad regions by at least two macro layers, so every plaquette and averaging stencil in $\Omega^+$ lies in the $P_0$-good region, where at this unresummed stage $|\widehat U(\partial p)-\id|_{\sigma}<\varepsilon_0$. If the canonical outer operation has already omitted these factors, they equal one on the support of the retained original-field factors ((B3)(i), the form in which CRE's removal statement before (1.11), p.248, is recorded), so restoring them changes no integral; since the history summands are nonnegative, a point of nonzero resummed density has at least one supported unresummed history. This restoration concerns $\widehat U$, not the minimum $U_1$, and controls the prescribed collar $V_0$; the distant exterior is never required to be regular. Conjugation and the bounded number of incident plaquettes give $| D^*_{\widehat U}\partial\widehat U|\leqslant C_{\mathrm{inc}}\varepsilon_0$ on the collar and the variable region ($C_{\mathrm{inc}}$ the incidence count, not the Gaussian constant $C_{\mathrm d}$), and multiplying by the largest of the fixed scale weights $L^2$, $L^3$ puts $\widehat U$ in VAR's class (6) with radius $C_L\varepsilon_0$. Its fixed data are exactly $V_0$ on the fine collar and $V$ on the interior coarse labels, since the original averaging constraint is retained.

Use the positive-bond carriers $E(X)$. The fixed fine collar carrier is $E(\Omega^+)\setminus E(\Omega_1)$, so no crossing bond that is a variable of CRE (1.12) is fixed. The hierarchy for VAR has outer fine domain $\Omega^+$, inner block domain $\Omega_1$ and determining set this bond-set difference at scale zero together with $\Omega_1$ at scale one, with the mixed-bond convention of VAR (7); its spacing condition (1) holds because the separating macro layer has width $\ell$, exceeding the fixed $RM_1$ of VAR for small coupling. For a mixed determining plaquette $p'$, replace every coarse bond of $p'$ outside the scale-one determining set by the block average of the scale-zero data, as prescribed in VAR (7): the original constraint makes the resulting loop an ordinary fine plaquette of $\widehat U$ or a coarse plaquette of its block average, and AVG Proposition 1, (51), bounds the latter, in AVG's operator norm, by $L^2\alpha_0+C(L^2\alpha_0)^2\leqslant 2L^2\alpha_0$ with $\alpha_0=\sqrt N\varepsilon_0$ for $g_0$ small ($CL^2\alpha_0\leqslant 1$; the conversion $\|\bigcdot \|_{\mathrm{op}}\leqslant \sqrt Nh$ of the conventions covers the case of the normalised Hilbert--Schmidt cutoff norm), which is also CRE's explicit bound before (1.2); the same $O_L(\varepsilon_0)$ bound follows from the elementary two-block comparison of Lemma~\ref{lem:c3-7-block}. VAR's constraint (3) and AVG's Proposition 1 are stated for Ba{\l}aban's earlier averaging; their use with the averaging RG1 (0.4) of CRE (1.12) is CRE's own transfer (p.256, all the results of [15] used for its averaging) and RG1's extension statement after (0.12), accepted as such. So the entire determining set of the localised problem has VAR data regularity at most $C_L\varepsilon_0$ ($C_L$ including $\sqrt N$), and no regularity of the rest of $\Omega_1^c$ is involved.

VAR Theorem 1 therefore constructs the regular critical orbit for these data, at radius $B_3C_L\varepsilon_0$, unique in the regularity class (6) of radius $\varepsilon_0^{\mathrm V}\coloneqq \max(C_L\varepsilon_0,B_3C_L\varepsilon_0)\leqslant a_0$ for small $g_0$, a class containing $U_1$, with both the plaquette and the divergence bounds (2). The minimum $U_1$ of CRE (1.12) lies on this orbit: it is critical for the localised problem, since the localisation changes the action by constants only, for every admissible variation, and it lies in VAR's regularity class on $\Omega^+$ by its membership $U_1\in D^{(11)}(V,V_0)$ (the first paragraph of this subsection); the uniqueness in the class identifies the two. Restricted to $Q_w^{[4]}$ the same bound gives membership in the local uniqueness class, with fixed scale and norm conversions (that CRE's determining set $\mathbf B_1(Q_w^{[4]})$ of (1.2), whose fine layer has width between $LM_1$ and $2LM_1$, p.256, satisfies VAR's spacing condition (1), is CRE's application of VAR to its own determining sets, p.256, taken over as such); since $\varepsilon_0\leqslant (1+\beta_0)^2\varepsilon_1$ by CRE (2.8) (p.256, with $n=1$, $m=0$), the radius is $C_{L,G}\varepsilon_1$, below VAR's $a_0$ for small coupling. The fine boundary layer of $U_{1,Q}$ has VAR's scale-zero bound, not an $\varepsilon_1L^{-2}$ bound: the plaquettes of its straight-extension data are trivial or equal to coarse plaquette holonomies of $V$, bounded by $2L^2\varepsilon_0$, so its mixed-scale data are regular with the same parameter $C_{L,G}\varepsilon_1$. What is used of the characteristic factors of CRE (1.11) at $U_1$ is their nonvanishing at $U_1$ itself, CRE's statement $U_1\in D^{(11)}(V,V_0)$ of p.248 ((B1)); nothing is assumed about the fluctuation cutoffs (1.18), (1.25) at zero fluctuation, and no branch at $A=0$ is introduced.

\subsubsection{The boundary data and their size}\label{sec:c3-9-2}
Let $D_Q=\dot M(U_1)|_{\mathfrak B_Q}$ and define on the determining set
\begin{equation}\label{eq:c3-9-2}
    B_Q(e)=\tfrac1i\log\bigl(D_Q(e)D^0_Q(e)^{-1}\bigr)\;.
\end{equation}
Both data are defined by the source averaging operation on this regular configuration; the determining interior carries scale-one coarse bonds and its boundary layer the finer determining variables, as defined after CRE (2.12); the mixed determining plaquettes and the local averages are those of CRE (1.2), (2.12)--(2.16). For the straight extension of CRE (1.3) every intrablock contour is the identity, every centre-to-centre path crossing a specified interface has transport $V(c)$, and every relative closed holonomy of RG1 (0.4) is the identity, so the logarithmic correction vanishes and $\mathsf M(Q_1^{s*}V)(c)= V(c)$; on the interior both data equal $V$, hence $B_Q=0$ there, and $\operatorname{supp}B_Q$ lies in the boundary layer of width $2LM_1$ of $Q_w^{[4]}$, the support statement preceding CRE (3.18). The principal logarithm is legitimate because the reasoning after CRE (3.17) bounds the argument of the logarithm by $C_{L,G}\varepsilon_1$ (CRE's printed coefficient $44d^2B_3$ belongs to its own regularity coefficient; the enlarged local coefficient is kept inside $C_{L,G}$). At one step this can also be read from the two-block comparison: in the axial mean gauge a regular minimum is within $C_{L,G}\varepsilon_1$ of the identity on internal bonds and of $V(c)$ on interface bonds, which are exactly the values of $Q_1^{s*}V$; the bound follows from the regularity of the minimum and the equality of the coarse averages, not from a fluctuation cutoff. Thus
\begin{equation}\label{eq:c3-9-3}
    \| B_Q\|_{\infty}\leqslant C_{L,G}\varepsilon_1\;.
\end{equation}
To meet VAR (172) with its absolute constant $C_1$, take its small parameter $\alpha= K_{\alpha}\varepsilon_0$ with $K_{\alpha}$ dominating all determining-data, uniqueness-class and $B_Q$ bounds (for the last, $\varepsilon_1\leqslant (1+\beta_0)\varepsilon_0$ by CRE (2.8)); then $| e^{itB_Q}-\id|\leqslant | tB_Q|\leqslant \alpha$ and $C_1=1$ suffices. VAR Proposition 9 and (190) are applied with this $\alpha$, not with CRE's identically named $\varepsilon_1$; for small $g_0$, $\alpha\leqslant a_1$ and $B_3\alpha\leqslant a_0$, and $\varepsilon_0\leqslant (1+\beta_0)^2\varepsilon_1$ (CRE (2.8)) converts the resulting $C\alpha$ bounds into $C'\varepsilon_1$; this uses VAR's printed absolute-$C_1$ statement without extending it. No bound is imposed on the exterior integration variables away from this local regular domain.

\subsubsection{Integrating the local derivative estimate}\label{sec:c3-9-3}
Let $\mathscr H_Q(B)$ be the relative Landau-gauge logarithm of the local minimal-configuration map of VAR Proposition 9 on $\mathfrak B_Q$ with reference $U_{1,Q}$, normalised by $\mathscr H_Q(0)=0$; this is the map of CRE (3.17), with its $L^{-1}$ scale factor at $k=0$. CRE (3.17) identifies the map at $B_Q$ with the restriction of $U_1$ up to the source gauge restoration: both lie on the unique regular critical orbit with determining data $D_Q$, the former satisfying the constraint up to the gauge transformation of VAR Proposition 2. For the criticality, a local admissible variation fixes the determining fine layer and the averages on the interior coarse rows; extended by zero outside $Q_w^{[4]}$ it preserves every global average and exterior value in CRE (1.12), because the fixed layer is wider than the plaquette and averaging range. The first variation of the global action vanishes and its local part is the local action variation; the gauge functional of the source axial mean gauge, in which the variational problem CRE (1.12) is posed (as recalled in the local-domain verification of Section~\ref{sec:c3-9-1}), has zero first variation at an axial mean configuration, where each of its nonnegative summands attains zero. So the restriction of $U_1$ is critical for the local action, it lies in the uniqueness class of Section~\ref{sec:c3-9-1} by its regularity there (from $U_1\in D^{(11)}(V,V_0)$), and the uniqueness identifies its orbit with the local analytic solution. The identity concerns the minimum $U_1$ and its regularity; no membership of the integrated fine field in a cutoff event is involved.

VAR (190), Proposition 9, gives the functional derivative with the mixed-scale weights; at one step the scale labels are $0$ and $1$, so the powers of $L$ in those weights are fixed constants, and a determining bond $e$ of the outer layer and an interior bond $b$ satisfy the bound $C_{L,G}e^{-c_{L,G}d(b,e)}$, the unit-density measure at the finer level and the scale-one determining measure contributing fixed $L$-dependent factors. Integrating along $tB_Q$, $t\in[0,1]$, and summing the labels,
\begin{equation}\label{eq:c3-9-4}
    |\mathscr H_Q(B_Q;b)|_{\sigma}\leqslant C_{L,G}\varepsilon_1\sum_{e\in\operatorname{supp}B_Q}e^{-c_{L,G}d(b,e)}\leqslant C'_{L,G}\varepsilon_1e^{-c'_{L,G}d(b,\operatorname{supp}B_Q)}\;,
\end{equation}
where for the second inequality each exponential is split into two halves, $e^{-c_{L,G}d(b,e)}\leqslant e^{-c_{L,G}d(b,\operatorname{supp}B_Q)/2}e^{-c_{L,G}d(b,e)/2}$, and the remaining sum over $e$ is bounded by the infinite-lattice sum $\sum_ee^{-c_{L,G}d(b,e)/2}$ over the bonds of $\mathbb Z^4$ (four per site) through injective nearest lifts, a constant depending on $c_{L,G}$ only; so $c'_{L,G}= c_{L,G}/2$ and no volume-dependent constant enters. For a bond $b\in(Q_w^{[2]})^*$ (its initial vertex is within one unit of $Q_w^{[2]}$, a difference absorbed in $C_L$),
\begin{equation}\label{eq:c3-9-5}
    d(b,\operatorname{supp}B_Q)\geqslant 2w-2LM_1-C_L\geqslant LM_2R_1
\end{equation}
once $w\geqslant LM_2R_1$ (the collar-width convention of Section~\ref{sec:c3-7}) and $LM_2R_1\geqslant 2LM_1+C_L$, which holds for small $g_0$ since $R_1\to\infty$ (Appendix~\ref{app:thresholds}); the first inequality is the distance from $Q_w^{[2]}$ to the complement of $Q_w^{[4]}$ less the width $2LM_1$ of the boundary layer carrying $\operatorname{supp}B_Q$ and the bond offsets. This covers either candidate common width $w$ and never uses a good/bad $R_1$ label.

\subsubsection{Restoring the source gauge}\label{sec:c3-9-4}
The transformation of VAR Proposition 2 is not centre-fixing: its condition is $(\overline{R_0u})^j=\id$ on $\Lambda_j$, in particular $u=\id$ on the fine determining layer (VAR p.281), and at one step its interior condition $\overline{R_0u}(y)=\id$ does not imply $u(y)=\id$. Keep the restoring transformation $u_Q$ of CRE (3.17), defined by this source normalisation: the configuration $e^{iL^{-1}\mathscr H_Q(B_Q)}U_{1,Q}$ of Proposition 9 is in the relative gauge of Proposition 2 and satisfies the constraint with data $D_Q$ up to that gauge, so it lies on the orbit of the restriction of $U_1$, and $u_Q$ is the transformation carrying it to $U_1$, whose values at the block centres need not be the identity and are retained. At the first step, with CRE's $\eta=1$ and its convention (1.17), on the required bonds
\begin{equation}\label{eq:c3-9-6}
    U_1(b)= u_Q(b_-)^{-1}e^{iL^{-1}\mathscr H_Q(B_Q;b)}U_{1,Q}(b)\,u_Q(b_+)\;.
\end{equation}
The local inputs of the transferred source estimates are the uniqueness and regularity class of Section~\ref{sec:c3-9-1}, the small supported $B_Q$ of Section~\ref{sec:c3-9-2} and the response bound \eqref{eq:c3-9-4} on both endpoint blocks and the crossing bond. Fix an analytic mean neighbourhood radius $\rho_{\mathrm{av}}>0$ for the finite source contour families, with derivative bounds on its closed half-radius neighbourhood, and choose $g_0$ after it so that the reference path-family diameter and its perturbation are both less than $\rho_{\mathrm{av}}/4$; this is the regular-chart reduction. The block-local gauge inequality \eqref{eq:c3-9-7} below is the transferred source result; the following path description explains its one-step support and its constants, and does not assert AVG's single-contour formula unchanged for the modified mean. In AVG the equations determining $u_Q$ are (77)--(87), with its expression in the relative fields given in (106)--(108): at one step the product in (106) has a single $j=0$ factor, a reference-unitary conjugation of the inverse block average followed by a local relative path transport, and the centre factor (87) is retained. Unitary conjugations preserve the norm, each path has length at most $(2d+1)L$, so the telescoping product identity bounds the relative deviation by $C_L$ times the largest relative-link deviation in the block, and the analytic averaging bound of (161)--(163) bounds the block-average deviation by the same quantity after reducing the chart. The transferred result reads
\begin{equation}\label{eq:c3-9-7}
    \sup_{x\in\mathcal B_y}| u_Q(x)-\id|_{\sigma}\leqslant C_{L,G}\sup_{b\in E(\mathcal B_y)}|\mathscr H_Q(B_Q;b)|_{\sigma}\;,
\end{equation}
the factor $L^{-1}$ of CRE (3.17) absorbed into the constant. The restoring map is relative to the reference pair: under a simultaneous gauge change $h$ of field and reference its representative is $u^h_Q(x)= h(x)u_Q(x)h(x)^{-1}$ (conjugate the relative action, VAR (16)--(17), and its reference-dependent gauge condition), so $| u_Q(x)-\id|_{\sigma}$ is a conjugation-invariant scalar of the relative construction, used only inside this gauge restoration; the centre contribution $u_Q(y)$ is retained. AVG uses its older averaging definition; the application here rests on Ba{\l}aban's extension statements, RG1 after (0.10) (RG1's considerations and results, and those of the earlier papers, hold for every group mean with the properties (0.5)--(0.9)) and RG1 after (0.12) (the average (0.12) has properties similar to those of (0.4), all results of RG1's [12] hold for it, with proofs mostly unchanged and otherwise with minor modifications, and both definitions are equally good for RG1's purposes; the coverage of (0.4) rests on this last remark). These are blanket transfer statements, none of which displays a restoring bound, so \eqref{eq:c3-9-7} is an accepted transfer of AVG's results to the RG1 averaging and mean gauge, used as that source result and not by asserting that the older individual-contour identities stay unchanged; CRE (3.17)--(3.19) uses the transferred results in the same way for the later scheme, as does RG1 for its gauge conversions ((3.26)--(3.27) with AVG (77), and (3.39) with AVG (97), retaining the centre-block gauge factors; RG1's reference [12] is AVG). No change of the physical averaging operation is made.

For the bond $b$ of \eqref{eq:c3-9-1} apply \eqref{eq:c3-9-4} on both endpoint blocks, on $b$ itself and on the paths entering the gauge factors, which enlarges the input set by at most $2L+2$, already included in the $C_L$ of \eqref{eq:c3-9-5}. The relative exponential and both endpoint factors of \eqref{eq:c3-9-6} then give $| U_1(b)U_{1,Q}(b)^{-1}-\id|_{\sigma}\leqslant C_{L,G}\varepsilon_1e^{-c_{L,G}LM_2R_1}$; both endpoint gauge factors of CRE (3.18) are retained, and the quantity is the simultaneous-gauge-invariant norm of the relative pair of links, not a comparison modulo an unspecified gauge transformation.

\subsubsection{The caller and the denominator}\label{sec:c3-9-5}
CRE (2.6), (2.8) bound $\varepsilon_1/\delta_0$ by fixed source constants and $R_1$ grows as a power of $\log g_1^{-2}$, so the bound is $<\delta_0/10$ for small $g_0$, which completes the proof of Proposition~\ref{prop:c3-9-1}. In Lemma~\ref{lem:c3-7-2} the same norm $|\bigcdot |_{\sigma}$ is used for the CRE (1.18) test, the comparison and the no-witness bound (the norm convention fixed before that lemma), and the eigenvalue inequality $| e^{iX}-\id|_{\sigma}\geqslant (2/\pi)| X|_{\sigma}$, valid when every eigenvalue of $X$ lies in $[-\pi,\pi]$, as on the native chart (the eigenvalues of a chart element lie in $(-\pi,\pi)$, the conventions of Section~\ref{sec:sourced-effective-interaction}), holds in both permitted norms without a norm conversion; the replay ($| g_0A(b)|_{\sigma}\geqslant \delta_0$ on a bad $R_1$ test, \eqref{eq:c3-7-7}, and the contradiction with the no-witness two-block bound $C_Q\theta\leqslant \delta_0/2$) is carried out in the proof of that lemma. The physical integral is never altered: for each original history the comparison is used inside its cutoff-support implication, the numerator of \eqref{eq:c3-7-1} remains $\int\mathcal O[\mathsf w_{\xi}\mathbf 1\{ z\in E(B_{\xi})\}]\,\dd V$ with $\mathsf w_{\xi}= e^{c_{\xi}}I_{\xi;1}$ and its denominator $\mathcal Z_0$; no $\zeta$ factor is removed, evaluated at $U_1$ or divided out, and the covariance estimate of Section~\ref{sec:c3-5} is unchanged. Lemmas~\ref{lem:c3-7-3} and~\ref{lem:c3-7-4} and Proposition~\ref{prop:c3-7-5} then apply with their existing hypotheses and even-period scope.

\section{The Signed Integrated Wilson Defects}\label{app:signed-defects}
This appendix bounds the defect term of the branch identity \eqref{eq:c-one-step-branch-identity} at the carrier $I$ of Section~\ref{sec:sourced-effective-interaction}, under the original history law: the joint defects $R_{\xi}(\mathfrak B)$ of \eqref{eq:part1-joint-defect} are summed with their restricted reference sums before absolute values are taken, each closed group is compared with its flat value by one background derivative along a local coefficient path, and the comparison is paid by the Gaussian reserve, the positive host comparison and the positive group summation of Sections~\ref{sec:c4-4}, \ref{sec:c4-7} and~\ref{sec:c4-8}. The result is Proposition~\ref{prop:c4-12-3}: on compatible even fine periods and for the square geometry of Section~\ref{sec:extent-weights},
\begin{equation*}
    | E_{\mathrm{def}}|\leqslant C'_Ln_Ce^{-(c_*/2)\pi_*^2}+e^{C_{\chi} n_C\rho(\varepsilon_1+e^{-cb/(LM)})}-1\;,
\end{equation*}
with $\rho= O(g_0^2(1+r)^4)$, $\varepsilon_1= g_1p_0(g_1)$, $b=\lfloor LMR_1/16\rfloor$, $\pi_*=\min(p_0(g_0),p_0(g_1))$, $C'_L=8C_L$ with the contour constant $C_L$ of Section~\ref{sec:c3-7} and the tail exponent $c_*$ of Proposition~\ref{prop:c3-7-5}, and $C_{\chi}$, $c$ fixed before the coupling. The regional scalar $S^{\circ}_{\xi,C}$ is retained inside the expectation and is not identified with the perimeter/corner coefficient here; the exponential is not replaced by its linear term unless its argument is bounded. Section~\ref{sec:one-step-insertion} keeps the history law and splits its weight $\mathsf w_{\xi}$ into labelled terms ($\sum_{\mathrm{lab}}\mathsf w^{\mathrm{lab}}_{\xi}=\mathsf w_{\xi}$, \eqref{eq:c5-6-1}) on the branches where a second collar test fails, retreating $\Lambda_1$ and the carrier $I$ to a labelled carrier $D(\Lambda)$; the reading of this appendix at $D(\Lambda)$ with the labelled weight $\mathbf 1_{\mathrm{lab}}e^{c^{\Lambda}_{\xi}}I^{\Lambda}_1$ (with $I\to D(\Lambda)$, $S^{\circ}_{\xi,C}\to S^{\circ}_{\Lambda}$, $R_{\xi}\to R^0_{\Lambda}$), which that section needs for its labelled defect $E^{\mathrm{lab}}_{\mathrm{def}}$, is Proposition~\ref{prop:c5-6-def}: it lists the inputs of Sections~\ref{sec:c4-2}--\ref{sec:c4-12} and where each is supplied at $D(\Lambda)$. Sections~\ref{sec:c4-2}--\ref{sec:c4-12} refer to the original law. The appendix shares the names $\varpi$, $\eta$, $C_{\mathrm{pin}}$, $K_{\chi}$ and $M(W)$ with Section~\ref{sec:c5-5}, which uses the same positive summation at $D(\Lambda)$; the objects here are those at $I$.

\subsection{Operand and retained normalisation}\label{sec:c4-1}
Use exactly the groups $W$ of Definition~\ref{def:groups}, the compatible families $\mathfrak B$, the closed pairings $A_{\xi}(\mathcal S)$ of \eqref{eq:part1-closed-family-pairing}, the scalar references $a_W$ of \eqref{eq:part1-reference-group}, the restricted reference sums $S^{\circ,[\mathfrak B]}_{\xi,C}$ and the joint defects $R_{\xi}(\mathfrak B)$ of Proposition~\ref{prop:joint-extraction}. Every source term, the Wilson contraction $f_{\xi}$ and each group's closed contraction of Definition~\ref{def:groups}, is a gauge-invariant closed contraction (Section~\ref{sec:c4-11-2}). Put
\begin{equation}\label{eq:c4-1-1}
    \mathscr R_{\xi}=\sum_{\substack{\mathfrak B\ne\varnothing\\\mathfrak B\ \mathrm{compatible}}}S^{\circ,[\mathfrak B]}_{\xi,C}R_{\xi}(\mathfrak B),\qquad
    E_{\mathrm{def}}=\frac{1}{\mathcal Z_0}\int\mathcal O\bigl[\mathsf w_{\xi}\mathbf 1_{A_{\mathrm{cov}}}\mathscr R_{\xi}\bigr](V)\,\dd V\;.
\end{equation}
The weight is the branch weight $\mathsf w_{\xi}= e^{c_{\xi}}I_{\xi;1}$ of \eqref{eq:c-collar-identity}, the outer operation $\mathcal O$ of \eqref{eq:native-source-explicit-outer-operation} contains the original cutoffs, $\zeta_{\xi}$ with its boundary-fluctuation dependence, the gauge and Haar factors, the Jacobians, the vacuum prefactor and the deferred integrations, $\mathcal Z_0=\int\mathcal O[\mathsf w_{\xi}]\,\dd V$, and $\mathbb E_{\mathrm h}$, $\mathbb P_{\mathrm h}$ are the history expectation and probability \eqref{eq:c3-6-1} for this positive weight; $A_{\mathrm{cov}}$ is the covered event of Lemma~\ref{lem:covered} (every contour bond is a bond of $\Omega_1$ and every free input of the contour lies in $I$). The weight is positive: $K_{\xi;1}$ is the exponential of a real function times indicators (Section~\ref{sec:boundary-law-collar}) and $c_{\xi}$ is real, the logarithm of the positive $a$-independent factor collected in the proof of \eqref{eq:c-collar-identity}. At fixed finite volume the signed sum \eqref{eq:c4-1-1} converges absolutely: the family sum \eqref{eq:part1-grouped-source} converges absolutely, and the proof of Proposition~\ref{prop:joint-extraction} obtains \eqref{eq:c4-1-1} from it by Boolean inversion over finitely many root families, each with an absolutely convergent cluster-list sum; no uniform absolute summability is inferred from that fact.

\subsection{Histories near the bad region}\label{sec:c4-2}
Let $B_{\xi}=\Lambda_1^c$ be the filled bad region and $E(B_{\xi})$ the set of positive fine bonds with at least one endpoint in $B_{\xi}$. Let $\Gamma_C$ be the deterministic set of contour bonds and free-input templates of Section~\ref{sec:c3-7}, with $|\Gamma_C|\leqslant C_Ln_C$, $n_C$ the number of contour traversals; distances are those of the periodic fine-bond metric $d_T$ of \eqref{eq:marginal-covariance-input}, $d_T(b,c)= d_{\infty,T}(x,y)+\mathbf 1\{ i\ne j\}$ for $b=(x,i)$, $c=(y,j)$. For an integer $b\geqslant 0$ define
\begin{equation}\label{eq:c4-2-1}
    \mathsf N_b=\{\xi\, | \,\operatorname{dist}(\Gamma_C,E(B_{\xi}))\leqslant b\}\;,
\end{equation}
with distance to the empty set infinite. This is a measurable history event ($\Lambda_1$ is a function of the history labels); it changes neither the covered event nor the law. Write $E^{\mathrm{near}}_{\mathrm{def}}=\mathbb E_{\mathrm h}[\mathbf 1_{A_{\mathrm{cov}}}\mathbf 1_{\mathsf N_b}\mathscr R_{\xi}]$ and $E^{\mathrm{far}}_{\mathrm{def}}=\mathbb E_{\mathrm h}[\mathbf 1_{A_{\mathrm{cov}}}\mathbf 1_{\mathsf N_b^c}\mathscr R_{\xi}]$, so that $E_{\mathrm{def}}= E^{\mathrm{near}}_{\mathrm{def}}+E^{\mathrm{far}}_{\mathrm{def}}$ (both integrable, by the bound $|\mathscr R_{\xi}|\leqslant 2$ of Lemma~\ref{lem:c4-2-1}).

\begin{lemma}[near histories]\label{lem:c4-2-1}
Let $b\geqslant 0$ be an integer and let $q_{\mathrm{pt}}$ be the one-point bound of Lemma~\ref{lem:c3-7-4} ($\mathbb P_{\mathrm h}(z\in E(B_{\xi}))= p_z\leqslant q_{\mathrm{pt}}$ for every bond $z$, on compatible even fine periods and for $g_0\leqslant g_H$), with $q_{\mathrm{pt}}\leqslant e^{-c_*\pi_*^2}$ for $0< g_0< g_*$ by Proposition~\ref{prop:c3-7-5}. Then $|\mathscr R_{\xi}|\leqslant 2$ pointwise on every branch, and
\begin{equation}\label{eq:c4-2-2}
    \mathbb P_{\mathrm h}(\mathsf N_b)\leqslant 4C_Ln_C(2b+3)^4q_{\mathrm{pt}},\qquad
    | E^{\mathrm{near}}_{\mathrm{def}}|\leqslant 8C_Ln_C(2b+3)^4q_{\mathrm{pt}}\leqslant C'_Ln_Ce^{-(c_*/2)\pi_*^2}\;,
\end{equation}
the last inequality for $b=\lfloor LMR_1/16\rfloor$, $C'_L\coloneqq 8C_L$ and $g_0$ below a threshold depending on $L$, $M_1$, $M_2$, $M$, the group, the character and the source constants.
\end{lemma}

\begin{proof}
By \eqref{eq:part1-scalar-extraction}, $\mathscr R_{\xi}= F_{\xi}(0)-S^{\circ}_{\xi,C}f_{\xi}(0)$. Both $F_{\xi}(0)=\int f_{\xi}\,\dd\nu_{\xi;0}$ (\eqref{eq:convex-response-operands}) and $S^{\circ}_{\xi,C}$ (\eqref{eq:part1-reference-gas}) are normalised integrals of a character of modulus at most one against the real positive kernels $K_{\xi;0}$ (the exponential of a real function times indicators, Section~\ref{sec:boundary-law-collar}) and $K^{\circ}_{\xi;0}$ (Lemma~\ref{lem:c2-reference}), and $| f_{\xi}(0)|\leqslant 1$; hence $|\mathscr R_{\xi}|\leqslant 2$. This bound is used only after the joint defects have been summed; it is not a bound on any $R_{\xi}(\mathfrak B)$, group or restricted reference factor. The members of $\Gamma_C$ are bonds (contour bonds and the bonds of the free-input templates), and there are at most $4(2b+3)^4$ positive fine-bond labels within $d_T$-distance $b$ of a fixed fine bond (an overcount on short periods), so a union bound over $\Gamma_C$ and the one-point estimate $p_z\leqslant q_{\mathrm{pt}}$ of Lemma~\ref{lem:c3-7-4} ($p_z$ of \eqref{eq:c3-7-1}) give the first bound; positivity of $\mathsf w_{\xi}$ and $\mathcal O$, used only after the signed sum has been formed, and $|\mathscr R_{\xi}|\leqslant 2$ give the second. For the last: $b=\lfloor LMR_1/16\rfloor$ exceeds every fixed source stencil for small $g_0$, and $\log(2b+3)= O(\log\log g_0^{-2})$ since $\ell= LMR_1\leqslant L^2M(2\log g_0^{-2})^{\mathfrak r}$ (\eqref{eq:c5-3-1}), while $\pi_*=\min(p_0(g_0),p_0(g_1))\geqslant A_0(\tfrac12\log g_0^{-2})^{p_0}$ with $p_0>0$, by $p_0(g)= A_0(\log g^{-2})^{p_0}$ (\eqref{eq:c3-1-1}) and $\log g_1^{-2}\geqslant \frac12\log g_0^{-2}$ (as at \eqref{eq:c5-3-1}); so $4\log(2b+3)\leqslant (c_*/2)\pi_*^2$, that is $8C_L(2b+3)^4e^{-c_*\pi_*^2}\leqslant C'_Le^{-(c_*/2)\pi_*^2}$ with $C'_L=8C_L$, for small $g_0$ at fixed source constants. The constants are uniform in the compatible even-period tori and the histories; no lower bound on a physical denominator is substituted in this proof (the normaliser bound \eqref{eq:c3-7-11} enters only through the cited one-point estimate), no conditional scalar is cancelled across histories, and no pointwise smallness of a Wilson-loop random variable is claimed.
\end{proof}

\subsection{The interior operand}\label{sec:c4-3}
The remaining operand is $E^{\mathrm{far}}_{\mathrm{def}}$, with the full signed group-family sum retained under its original weight. A group's random-field support is not its background support: its Wilson contraction carries the background transports along its contour interval $\mathcal H(W)$ as well. The flat comparisons of Sections~\ref{sec:c4-4} and~\ref{sec:c4-6} compare each closed group with its flat value along the group's own interval transport closed back to its base (Remark~\ref{rem:c4-5-1}, Lemma~\ref{lem:c4-9-1}), so these transports, in their order, are part of the compared word: no group, in particular none meeting separated sides of the square, is flattened on its random support alone while the intervening transports are changed. The analytic step is to insert the background derivative of the closed source-containing selected integrals and of their joint vacuum ratio before taking absolute values or summing group families (Section~\ref{sec:c4-4}), and to sum the groups afterwards (Sections~\ref{sec:c4-6}, \ref{sec:c4-8} and~\ref{sec:c4-9}). The derivative bounds used are those of Sections~\ref{sec:c4-7} and~\ref{sec:c4-10}.

\subsection{The joint signed selected-integral estimate}\label{sec:c4-4}
Fix a compatible root family $\mathcal A=\{ D_1,\ldots,D_m\}$. For each $D_j$ use the selected measures $\Lambda^{\mathrm f}_{\xi;\ell_j}$ of \eqref{eq:probe-selection-measure} with $\ell_j=(\mathcal D_j,\mathsf J_j,\mathsf P_j,Z_j)$, $Z_j= D_j$ and $\mathsf J_j\ne\varnothing$; put $q_j=|\mathsf J_j|$ and $q=\sum_jq_j$. The centred measure $\widehat\nu^{\mathrm f}_{\xi,D}=\nu^{\mathrm f}_{\xi,D}-k^{\mathrm f}_{\xi}(D)\delta_{0,\mathsf L_D}$ of \eqref{eq:fixed-probe-zero-mass} is exactly the sum of the nonempty-probe selections of \eqref{eq:probe-selection-measure}: the empty-probe selections give the mass times $\delta_0$ and disappear under centring. This is the measure identity of Lemma~\ref{lem:probe-measures}; no new expansion is introduced. Define the finite-difference operator on a bounded function $f$ of the young variable by
\begin{equation}\label{eq:c4-4-1}
    \Delta_{\mathsf J}f(B)=\sum_{K\subseteq\mathsf J}(-1)^{|\mathsf J\setminus K|}f(\mathbf 1_KB)\;,
\end{equation}
every unselected young label being zero; for several roots compose these operators on their disjoint label sets, writing $\Delta_{\boldsymbol{\mathsf J}}f$. The integrand is a sum of closed Wilson contractions with the ordering and background transports retained; the source may carry a fixed complementary closed holonomy $h\in G$, transforming by conjugation at the base, and all bounds are uniform in $h$. With \eqref{eq:fixed-probe-active-pairing},
\begin{equation}\label{eq:c4-4-2}
    \mathfrak Q_{\xi}(\mathcal A;f)=\sum_{(\ell_j)_{j\leqslant m}\, | \,Z_j= D_j,\ \mathsf J_j\ne\varnothing}\int\Delta_{\boldsymbol{\mathsf J}}f(B)\,\dd\bigotimes_{j=1}^m\Lambda^{\mathrm f}_{\xi;\ell_j}(B_j)\;,
\end{equation}
the integral of each Dirac tensor $\bigotimes_{b\in\mathsf J}(\delta_{B_b}-\delta_0)$ against $f$ being the finite difference \eqref{eq:c4-4-1}.

\subsubsection{The powers supplied by the signed source}\label{sec:c4-4-1}
Set $u= C_{\chi} g_0(1+r)$, increasing the fixed row/character constant $C_{\chi}$ once (Section~\ref{sec:c4-10-4}). On the real protected box the row reconstruction depends on $g_0B$ and has analytic derivative bounds on a larger fixed phase chart (Lemma~\ref{lem:charts}, Lemma~\ref{lem:c2-uniform-chart}). Each selected free label affects a bounded number of reconstructed contour links, each link has a bounded input set, and for the square the traversal multiplicity is bounded; applying the product rule inside the closed character therefore gives $C_{\chi}^q$ derivative assignments, not $n_C^q$ or $q!$. Derivative orders on any one row are bounded by twice its fixed input count plus one, including the remainders below (Section~\ref{sec:c4-10-4}); the higher row derivatives are obtained from the bounded analytic phase map on its larger fixed ball by Cauchy estimates (Lemma~\ref{lem:c2-uniform-chart}(i)--(ii): $\widetilde D_{U,c}$ is analytic on the ball of radius $\vartheta_*$ with $\|\widetilde D_{U,c}(A)\|\leqslant 2K_*\| A\|^2$ and $\| h_U(c)\|\leqslant H_*$, so Cauchy's estimates on the ball of radius $\vartheta_*/2$ bound every order; the finite list of derivative constants of item (iii) is not relied on beyond order five), and the maximum order used is fixed by $L$. All untouched transports are unitary matrices inside the same closed contraction. For the closed character $f$ of the source (Lemma~\ref{lem:c4-4-2} below) the mixed integral remainder formula \eqref{eq:c4-10-7} gives
\begin{equation}\label{eq:c4-4-3}
    \sup|\Delta_{\boldsymbol{\mathsf J}}f|\leqslant u^q\;,
\end{equation}
the supremum over the real protected box, which contains the support of every selected measure in the selected variables (the measure \eqref{eq:probe-selection-measure} carries the cutoffs $\chi_b$ over $E(Y_0)$ in the variables it integrates; its bad cutoffs $1-\chi_b$ act on the labels of $\mathsf P$, which are integrated out, not on the selected probes $\mathsf J$).

For a root with odd $q_j$ take the homogeneous term of degree $q_j$ in its selected variables $B_{D_j}$, the other root variables fixed; it is odd under $B_{D_j}\to-B_{D_j}$, and the difference from it carries an additional factor $u$ (\eqref{eq:c4-10-9}; the choice of the extra increment is bounded by $q_j\leqslant 2^{q_j}$, absorbed in $C_{\chi}$). The same holds simultaneously for any set of odd roots, the product-rule allocation count remaining exponential in $q$. Every term is a closed contraction; no open interval is estimated separately. Each selected Gaussian kernel \eqref{eq:fixed-kernel-density} with its cutoffs is invariant under $(B_{D_j},X_{D_j})\to(-B_{D_j},-X_{D_j})$ once its ordinary Mayer factors are removed: the determinant, the quadratic cross term, the compensator, the real cutoffs $\chi_b$, $1-\chi_b$ and the spatial Cauchy integrations are even. The tensor product in \eqref{eq:c4-4-2} gives separate integration variables and a separate reversal for each root. Thus the leading odd part integrates to zero in the signed selected integral, also when several roots are present.

Let $\omega(g_0)= C_0^{\partial}g_0(1+r)^3$ with the host derivative constant $C_0^{\partial}\geqslant C_0$ of Section~\ref{sec:c4-7-1}, let $\omega_*$ be the fixed host budget of \eqref{eq:native-source-vacuum-KP-budget}, and put $R_{\Upsilon}=\omega_*/(2\omega(g_0))$; choose $g_0$ so that $R_{\Upsilon}\geqslant 2$. In root $j$ multiply the selected ordinary potentials $\mathbf V_{\xi}(Y)$, $Y\in\mathcal D_j$, by an independent complex parameter $\zeta_j$, $|\zeta_j|\leqslant R_{\Upsilon}$, leaving the Gaussian kernel, cutoffs and probes unchanged. If $\mathcal D_j\ne\varnothing$ the selected integral vanishes at $\zeta_j=0$ (each Mayer factor $\int_0^1\zeta\mathbf V e^{t\zeta\mathbf V}\dd t= e^{\zeta\mathbf V}-1$ does), so the Cauchy difference between $\zeta_j=0$ and $\zeta_j=1$ of the analytic function of $\zeta_j$ costs at most $2/(R_{\Upsilon}-1)$ times its supremum on the disc; if $\mathcal D_j=\varnothing$ the density is the even Gaussian/cutoff kernel and the reversal applies. Independent parameters on disjoint roots fit the same local potential budget simultaneously (Section~\ref{sec:c4-7-1}). Therefore an odd $q_j$ costs
\begin{equation}\label{eq:c4-4-4}
    v= u+\frac{2}{R_{\Upsilon}-1}\;.
\end{equation}
No separate cancellation of a group coefficient is assumed; the cancellations occur inside the original signed selected integrals. Both $u$ and $v$ tend to zero as $g_0\to0$, since $g_0(1+r)^3\to0$ by \eqref{eq:c3-1-1}. ($\omega$ as a background direction, with $\|\omega\|_{\infty}$ and $\mathscr D_{\omega}$, and $\omega(g_0)$, $\omega_*$ as host bounds are distinguished by their arguments.)

\subsubsection{Numerical majorants paid with the same Gaussian reserve}\label{sec:c4-4-2}
For a selected root $\ell$ define
\begin{equation}\label{eq:c4-4-5}
    M_{\ell}=\sup_{|\zeta|\leqslant R_{\Upsilon}}\|\Lambda^{\mathrm f}_{\xi;\ell,\zeta}\|_{\mathrm{TV}}\;,
\end{equation}
where $\Lambda^{\mathrm f}_{\xi;\ell,\zeta}$ is the measure \eqref{eq:probe-selection-measure} with each selected ordinary potential replaced by $\zeta$ times that potential; the cutoffs, Gaussian density and normalisers, compensator, integration carriers, weakening contours and seed data are unchanged. The total variation is a numerical majorant used after the signed source cancellation; it is not a bound on the centred Dirac tensors.

\begin{lemma}[dummy comparison]\label{lem:c4-4-1}
There are at most $P_M=4(4M+3)^4$ probe labels with a given host (a label $b$ lies in the interior of its host $Y_b$, a rectangle of at most four $M$-cubes per side, \eqref{eq:c-probe-host}), and each probe host has $N_0(Y_b)\leqslant 4^4\leqslant s_{\mathrm p}=5^4$ cubes (\eqref{eq:c-probe-host}). Let $R_{\Upsilon}\geqslant 2$ and let $g_0$ be below the thresholds of Lemmas~\ref{lem:activity-bounds} and~\ref{lem:c2-selected-majorant}. Choose once, after $M$ and before $g_0$,
\begin{equation}\label{eq:c4-4-6}
    u_*=\min\Bigl(\frac14,\ \frac{\omega_*e^{-\kappa s_{\mathrm p}}}{8P_M}\Bigr)\;.
\end{equation}
Then for every connected final domain $D$, with $m(D)=\varepsilon_*e^{-\lambda_*\ell_1(D)}$ the majorant of Lemma~\ref{lem:activity-bounds},
\begin{equation}\label{eq:c4-4-7}
    \sum_{\ell\, | \,Z(\ell)= D,\ \mathsf J(\ell)\ne\varnothing}u_*^{q(\ell)}M_{\ell}\leqslant m(D)\;.
\end{equation}
\end{lemma}

\begin{proof}
The number $u_*$ is a fixed comparison weight, not a measure. For every probe label $b$ insert, in the numerical comparison only, the dummy potential $p^+_b= u_*$ on its protected host box and zero outside; its selected Mayer factor $e^{u_*}-1$ is at least $u_*$. The summed dummy potential at a host is at most $P_Mu_*\leqslant \omega_*e^{-\kappa s_{\mathrm p}}/8\leqslant \omega_*e^{-\kappa d_0(Y)}/8$, the input $B_Y$ of \eqref{eq:c4-7-2} (for a probe host $d_0(Y)\leqslant N_0(Y)-1< s_{\mathrm p}$: its cubes form a rectangular array, and a path through the centres of wall-adjacent cubes visiting each once meets every cube and has length $N_0(Y)-1$ in \eqref{eq:c2-tree-length}), while the complex ordinary potentials use at most half the budget $\omega_*e^{-\kappa d_0(Y)}$ on $|\zeta|\leqslant R_{\Upsilon}$ (Lemma~\ref{lem:c2-potential-size} and $R_{\Upsilon}\omega(g_0)=\omega_*/2$). Split the labels with a common host exactly as in the proof of Lemma~\ref{lem:probe-measures}; the seeds, conditioning carrier and Gaussian kernel depend on the active hosts only and are unchanged. Sum all split selections with the same active-host set first: Lemma~\ref{lem:c4-7-1} proves that their many-to-one positive sum is dominated hostwise by the undifferentiated majorant \eqref{eq:fixed-kernel-selected-majorant}. That majorant pays its determinant and compensator with the Gaussian reserve $\alpha_G$ (Lemma~\ref{lem:gaussian-payment}(ii)), and the cover and seed sums \eqref{eq:native-source-initial-domain-sum}, \eqref{eq:native-source-enlarged-domain-sum} with the final $Y'$-sum in the proof of Lemma~\ref{lem:activity-bounds} (Appendix~\ref{app:background-estimates}) sum the seeds and final domains to $m(D)$. All parameters in $m(D)$ are the existing positive-majorant parameters, and \eqref{eq:c4-4-7} holds with the exact $\lambda_*$. (The polynomial factors inserted in Sections~\ref{sec:c4-7} and~\ref{sec:c4-8} are paid there by a fixed decrease of the remaining decay of $m(D)$.) Nothing is changed in the physical potential or in a normalisation of $E_{\mathrm{def}}$.
\end{proof}

\begin{lemma}[joint selected-integral bound]\label{lem:c4-4-2}
Let $f$ be a normalised character of an ordered closed reconstructed word, optionally closed with its inverse background interval and a fixed complementary $h$ (Section~\ref{sec:c4-7-2}). For $g_0$ small enough that $u\leqslant u_*$ and $v\leqslant 1$, put
\begin{equation}\label{eq:c4-4-8}
    \rho=\Bigl(\frac{u}{u_*}\Bigr)^2+\frac{u}{u_*}\,v\;.
\end{equation}
Then for every compatible root family $\mathcal A$ and every choice of selections $(\ell_j)$ with $\mathsf J_j\ne\varnothing$,
\begin{equation}\label{eq:c4-4-9}
    \Bigl|\int\Delta_{\boldsymbol{\mathsf J}}f(B)\,\dd\bigotimes_{j=1}^m\Lambda^{\mathrm f}_{\xi;\ell_j}(B_j)\Bigr|\leqslant \prod_{j=1}^mu^{q_j}v^{q_j\bmod2}M_{\ell_j},\qquad
    |\mathfrak Q_{\xi}(\mathcal A;f)|\leqslant \prod_{D\in\mathcal A}\rho\,m(D)\;.
\end{equation}
\end{lemma}

\begin{proof}
The first bound combines the reversal, the Cauchy differences and \eqref{eq:c4-4-3}--\eqref{eq:c4-4-4} as written out in Section~\ref{sec:c4-7-2}: the roots with odd $q_j$ and empty ordinary selection lose their leading odd part to the even kernel and keep the remainder $u^{q_j+1}$, the roots with odd $q_j$ and nonempty ordinary selection pay $u^{q_j}\cdot2/(R_{\Upsilon}-1)$ by the Cauchy difference in $\zeta_j$ (the supremum on the disc being $u^{q_j}M_{\ell_j}$ by \eqref{eq:c4-4-5}), the roots with even $q_j$ pay $u^{q_j}M_{\ell_j}$. The bound retains the full joint source through each cancellation. For every nonempty probe set $(u/u_*)^{q_j}v^{q_j\bmod2}\leqslant \rho$ ($q_j=1$ gives $(u/u_*)v$, $q_j\geqslant 2$ gives $(u/u_*)^2$). Summing over the selections of every root with \eqref{eq:c4-4-2} and \eqref{eq:c4-4-7} gives the second bound.
\end{proof}

The gain is $\rho= O(g_0^2(1+r)^4)$ with constants fixed before the coupling; it is obtained after the signed integration, not by replacing a probe tensor by its total variation $2^{|\mathsf J|}$.

\subsubsection{Background derivative and the joint vacuum ratio}\label{sec:c4-4-3}
For the flat subtraction take a real background direction $\omega$ along one of the globally small analytic extensions of Section~\ref{sec:c4-9-1}; the geometry, carriers, cutoff norm, means, Gaussian reserve and walk recipe are unchanged. Let $L_f$ be the number of occurrences of contour background links in the closed word whose local coefficient arguments vary; when a group's relative interval is compared this includes the interval and its inverse closing transport. Complementary closed holonomies are held fixed, and no outer history variable or physical weight is differentiated. Write $\mathscr D_{\omega}$ for this derivative.

\begin{lemma}[one background derivative]\label{lem:c4-4-3}
On the extensions of Section~\ref{sec:c4-9-1}, with $\rho$, $m(D)$ of Lemma~\ref{lem:c4-4-2} and a constant $C$ fixed before the coupling,
\begin{equation}\label{eq:c4-4-10}
    |\mathscr D_{\omega}\mathfrak Q_{\xi}(\mathcal A;f)|\leqslant C\|\omega\|_{\infty}\Bigl(L_f+\sum_{D\in\mathcal A}N_1(D)\Bigr)\prod_{D\in\mathcal A}\rho\,m(D)\;.
\end{equation}
\end{lemma}

\begin{proof}
The derivative of the selected Gaussian factors (density, normaliser, compensator) inserts a scalar polynomial of degree at most two in $(B,X)$, including the determinant trace; its moment is paid by the joint coercive Gaussian bound of Section~\ref{sec:c4-10-3} at cost $CN_1(D)\|\omega\|_{\infty}$. Differentiated walk coefficients sum in the convergent differentiated series of Section~\ref{sec:c4-10-2}. The derivative of a host Mayer $t$-integrand is $\zeta(\mathscr D_{\omega}\mathbf V)e^{t\zeta\mathbf V}(1+t\zeta\mathbf V)$, whose $t$-integral $\zeta(\mathscr D_{\omega}\mathbf V)e^{\zeta\mathbf V}$ vanishes at $\zeta=0$; its size is bounded by the differentiated host comparison of Lemma~\ref{lem:c4-7-1}. A Gaussian derivative is even under each component reversal and the leading source derivative has the parity of the undifferentiated term (\eqref{eq:c4-10-10}), so the signed powers of Section~\ref{sec:c4-4-1} survive one derivative, with the additional factor $C\|\omega\|_{\infty}(L_f+\sum_DN_1(D))$: the source derivative contributes $L_f$ plus the number of selected labels, the latter at most $4(LM+1)^4\sum_DN_1(D)$ (Section~\ref{sec:c4-7-1}), and the host, Gaussian and walk derivatives contribute one factor $N_1(D)$ each by Lemma~\ref{lem:c4-7-1}. No division by a potential, activity or character is used. The factor $N_1(D)$ is kept explicitly; it is paid by $m(D)$ in Section~\ref{sec:c4-7}.
\end{proof}

The vacuum ratio is kept as the joint ratio of the whole root family, not a product of one-root ratios: by \eqref{eq:fixed-probe-active-pairing} and \eqref{eq:part1-joint-vacuum-exponential},
\begin{equation}\label{eq:c4-4-11}
    q_{\xi}(\mathcal A)=\exp\Bigl[-\sum_{\mathfrak c\text{ hits }\mathcal A}\Phi^{\mathrm v}_{\xi}(\mathfrak c)\Bigr]\;,
\end{equation}
a cluster hitting several roots being counted once, and its derivative is the derivative of this exponential.

\begin{lemma}[the joint ratio and its derivative]\label{lem:c4-4-4}
With the constant $C_M$ of Lemma~\ref{lem:c4-7-1},
\begin{equation}\label{eq:c4-4-12}
    | q_{\xi}(\mathcal A)|\leqslant e^{\sum_{D\in\mathcal A}N_1(D)/8},\qquad
    |\mathscr D_{\omega} q_{\xi}(\mathcal A)|\leqslant \frac{C_M}{2}\|\omega\|_{\infty}\sum_{D\in\mathcal A}N_1(D)\,e^{\sum_{D\in\mathcal A}N_1(D)/8}\;,
\end{equation}
and consequently
\begin{equation}\label{eq:c4-4-13}
    \bigl|\mathscr D_{\omega}[q_{\xi}(\mathcal A)\mathfrak Q_{\xi}(\mathcal A;f)]\bigr|\leqslant C\|\omega\|_{\infty}\Bigl(L_f+\sum_{D\in\mathcal A}N_1(D)\Bigr)\prod_{D\in\mathcal A}\bigl[\rho\,m(D)e^{N_1(D)/8}\bigr]\;.
\end{equation}
\end{lemma}

\begin{proof}
The first bound is Lemma~\ref{lem:vacuum-clusters}. The second is Lemma~\ref{lem:c4-7-4}, which differentiates one vacuum activity in its connected-cluster tree majorant, keeping every occurrence and symmetry factor and marking the single derivative in the fixed positive cluster majorant, without shrinking a previously chosen parameter or dividing by an activity. The third follows from the product rule with Lemmas~\ref{lem:c4-4-2} and~\ref{lem:c4-4-3} and the two bounds of \eqref{eq:c4-4-12}, with the constant $C$ of \eqref{eq:c4-4-10} increased by $C_M/2$.
\end{proof}

The actual joint ratio is differentiated before estimating. The normaliser $I_{\xi;0}$ encoded by this ratio is not cancelled against the physical history factor $\mathsf w_{\xi}= e^{c_{\xi}}I_{\xi;1}$.

\subsection{The closed group subtraction and its signed caller}\label{sec:c4-6}
\begin{remark}[what the subtraction must pay]\label{rem:c4-5-1}
Section~\ref{sec:c4-4} supplies the derivative estimate when the flat subtraction of each closed group is made inside $\mathscr R_{\xi}$. The group-family summation (this section and Sections~\ref{sec:c4-8}, \ref{sec:c4-9}) must pay the overlap of contour intervals and the spatial decorations, and must control the restricted reference sums $S^{\circ,[\mathfrak B]}_{\xi,C}$ with their regional flat kernel; no product of one-component ratios and no uncontrolled exponential in the loop length is admissible. The group contraction for this caller uses the interval transport relative to its background interval, closed back to its base, with an arbitrary complementary closed holonomy $h$; its flat value is $a_W\chi(h)$ by Lemma~\ref{lem:schur}. The product bound for such closed contractions is the finite-dimensional norm equivalence of Lemma~\ref{lem:c4-6-1}; no independent open matrix norm is an inductive carrier, and the multi-group factorisation is checked against \eqref{eq:part1-closed-family-pairing}, not presumed from spatial independence of the physical law.
\end{remark}

\emph{Geometry.} The square/torus geometry of Section~\ref{sec:extent-weights} is in force: $C$ is a coarse square of side $s= nL$ fine bonds on fine periods $LK_i$ with $K_i\geqslant 3(n+2)$, and $n_C=4s$. For a group $W$, $\mathcal H(W)$ is its interval (Definition~\ref{def:groups}: a cyclic interval of fewer than $n_C/2$ contour traversals, or the whole contour read from the base point) and $|\mathcal H(W)|$ the number of traversals in it. No loop separation hypothesis is added. Contour hull length is paid by connected spatial extent as in \eqref{eq:marginal-hull-extent}: contacts of diameter $< s/2$ occupy one side or two adjacent sides and fit in an arc of length at most $4\operatorname{diam}+1$; otherwise $n_C\leqslant 8\operatorname{diam}$. Periodic distances between square contacts equal their lifted distances.

For a group $W$ let $n(W)$ be the sum of $N_1$ over its roots and over every constituent domain of every occurrence in its cluster list, with multiplicities, and define
\begin{equation}\label{eq:c4-6-1}
    \mathfrak s(W)= n(W)+\frac{|\mathcal H(W)|}{LM}\;.
\end{equation}
The footprint $F(W)$ is the set of $LM$-cubes meeting a constituent domain, its fixed coefficient stencil halo, or the contour/free-input stencil of a traversal of $\mathcal H(W)$ (``meeting'' in the site sense: a cube containing a site or a bond endpoint of the set; the ranges are made explicit in Lemma~\ref{lem:c4-6-env}; the comparison box $Q_W$ of Section~\ref{sec:c4-8-4} carries its own fixed enlargement, absorbed in $C_{\mathrm{box}}$, covering the Hessian stencils and reconstruction rows consulted by the coefficient, and is not part of the footprint); the source geometry gives $| F(W)|\leqslant C\mathfrak s(W)$ (Lemma~\ref{lem:c4-6-env}(a)). For a component of the roots/cluster-occurrence graph the union of its domains is wall connected, so its diameter is at most $CLM$ times its cube count, and its contact hull is bounded by the square estimate above; overlapping hulls have a connected interval union, which pays the hull of a group by the sum of its component counts. If all components are merged because some hull is long (Section~\ref{sec:joint-extraction}), the sum of the component sizes in that class is at least $cn_C/(LM)$. Hence $|\mathcal H(W)|/(LM)\leqslant Cn(W)$. Coefficients whose local rectangle fails to lift are kept in the large-domain class of Section~\ref{sec:c4-8-4}.

\subsubsection{The closed relative interval functional}\label{sec:c4-6-1}
For $W$ use its root variables only, every other young variable zero, in the reconstructed links along $\mathcal H(W)$. Let $T_W(a;U_1)$ be the ordered product of these links (the interval transport) and $P_W(U_1)$ the same product at zero young variable, both with their source orientations and with the same two endpoint vertices. The relative product $T_W(a;U_1)P_W(U_1)^{-1}$ is based and closed at the initial vertex. Define the scalar closed contraction
\begin{equation}\label{eq:c4-6-2}
    c_W(U_1;h)= w_{\xi}(W)\int\chi\bigl(T_W(a;U_1)P_W(U_1)^{-1}h\bigr)\,\dd\Bigl(\bigotimes_{D\in\mathcal A(W)}\widehat\nu^{\mathrm f}_{\xi,D}\Bigr)(a)\;,
\end{equation}
$h\in G$ a complementary closed holonomy at that base, transforming by simultaneous conjugation; the product $T_WP_W^{-1}$ and $h$ transform at one base, so the displayed scalar is gauge invariant and no open interval functional is estimated. Put
\begin{equation}\label{eq:c4-6-3}
    d_W(U_1;h)= c_W(U_1;h)-a_W\chi(h),\qquad\mathfrak d_W(U_1)=\sup_{h\in G}| d_W(U_1;h)|\;.
\end{equation}
At the reference data, where the same formula is taken with the reference weights and measures, $c_W(\id;h)= a_W\chi(h)$ by Lemma~\ref{lem:schur} (there the interval transport at the root variables is the $P^{\circ}_W(a)$ of that lemma and the background transport is the identity). The weights and measures in $c_W$ are those of the joint extraction, including every cluster occurrence of the group; they are not a newly normalised conditional law.

\begin{lemma}[factorisation and product inequality]\label{lem:c4-6-1}
Let $\pi$ be the irreducible unitary representation with $\chi=\dim(\pi)^{-1}\operatorname{tr}\pi$. (a) There is a constant $K_{\chi}$ of the character such that $\|\Upsilon\|_{\mathrm{op}}\leqslant K_{\chi}\sup_{h\in G}|\dim(\pi)^{-1}\operatorname{tr}(\Upsilon\pi(h))|$ for every $\Upsilon\in\operatorname{End}(V_{\pi})$. (b) For a compatible family $\mathfrak B=\{ W_1,\ldots,W_m\}$ on a covered branch, $R_{\xi}(\mathfrak B)$ is the normalised trace of the ordered product of the matrices $\Upsilon_{W_j}\pi(G_j)$, where $\Upsilon_W= w_{\xi}(W)\int\pi(T_WP_W^{-1})\,\dd\bigotimes\widehat\nu^{\mathrm f}_{\xi}-a_W\id$ and $G_j$ is the background transport of the gap after $\mathcal H(W_j)$ composed with $P_{W_j}$. (c) Consequently
\begin{equation}\label{eq:c4-6-4}
    | R_{\xi}(\mathfrak B)|\leqslant \prod_{W\in\mathfrak B}K_{\chi}\mathfrak d_W(U_1)\;.
\end{equation}
\end{lemma}

\begin{proof}
(a) $\Upsilon\mapsto(h\mapsto\operatorname{tr}(\Upsilon\pi(h)))$ is linear and injective on $\operatorname{End}(V_{\pi})$, since $\pi(G)$ spans $\operatorname{End}(V_{\pi})$ for an irreducible representation (Burnside); so the supremum over $h$ is a norm on the finite-dimensional space $\operatorname{End}(V_{\pi})$, equivalent to the operator norm, and $K_{\chi}$ is the equivalence constant. (b) In a compatible family a traversal of $\mathcal H(W)$ has no contact with a root of a different group (that contact hull would meet $\mathcal H(W)$), so its reconstructed link uses the root variables of $W$ only (Lemma~\ref{lem:row-locality}); a traversal outside every interval has its background link. Number the groups in cyclic order and start the ordered link product at the first traversal of $\mathcal H(W_1)$ (cyclicity of the trace). Expanding the ordered product $\prod_j(X_j-a_{W_j}\id)\pi(G_j)$ with $X_j= w_{\xi}(W_j)\int\pi(T_{W_j}P_{W_j}^{-1})\dd\bigotimes\widehat\nu^{\mathrm f}$ over subsets $\mathcal S$ of slots, a slot of the first kind contributes the links of $\mathcal H(W_j)$ at the root variables followed by the gap transport, a slot of the second kind contributes $-a_{W_j}$ times the links of $\mathcal H(W_j)$ at zero young variable followed by the gap; by multilinearity of the trace and Fubini over the product measure (as in the proof of Lemma~\ref{lem:schur}) the normalised trace of the term of $\mathcal S$ is $(\prod_{W\notin\mathcal S}(-a_W))\,(\prod_{W\in\mathcal S}w_{\xi}(W))\,\mathfrak Q_{\xi}(\mathcal A(\mathcal S);f_{\xi})=(\prod_{W\notin\mathcal S}(-a_W))A_{\xi}(\mathcal S)$ by \eqref{eq:part1-closed-family-pairing}, since on a covered branch the links outside the intervals are the background links at $U_1$ (Section~\ref{sec:c4-8-1}; Lemma~\ref{lem:covered} gives $f_{\xi}(0)= W_{\chi}(C;U_1)$). Summing over $\mathcal S$ reproduces \eqref{eq:part1-joint-defect}. This is an algebraic factorisation of the existing polymer integrals, not independence of physical field regions. (c) $|\operatorname{tr}(Y_1\cdots Y_{2m})|\leqslant \dim(\pi)\prod\| Y_i\|_{\mathrm{op}}$, the gap transports are unitary, and $\|\Upsilon_W\|_{\mathrm{op}}\leqslant K_{\chi}\mathfrak d_W(U_1)$ by (a), since $\dim(\pi)^{-1}\operatorname{tr}(\Upsilon_W\pi(h))= d_W(U_1;h)$. Matrices are used to prove the finite-dimensional multilinear identity; no separately estimated endpoint-covariant matrix or open Wilson carrier is introduced.
\end{proof}

\subsubsection{Positive graph sums for the reference factors}\label{sec:c4-6-2}
This is the summation with the estimates of Sections~\ref{sec:c4-4} and~\ref{sec:c4-7}; it concerns the source-containing groups, not the vacuum gas alone. A cluster occurrence $\mathfrak c$ has cube cost $n(\mathfrak c)=\sum_{D\in\mathfrak c}N_1(D)$; its positive weight $M_v(\mathfrak c)$ is the tree majorant of its Ursell coefficient: for the ordered tuple defining $\mathfrak c$, divide by the tuple factorial and sum over the incompatibility spanning trees, multiplying the majorants $m(D)$ of the activities (the positive tree sum of the proof of Lemma~\ref{lem:vacuum-clusters}, Section~\ref{app:kp}). Source roots have the weights $\rho\,m(D)$ of Lemma~\ref{lem:c4-4-2}. For each $W$ the product of its root weights, its cluster weights $M_v(\mathfrak c_j)$ and $1/k(W)!$ is denoted $M(W)$; it bounds $| a_W|$, the pairings $| w_{\xi}(W)\mathfrak Q_{\xi}(\mathcal A(W);f)|$ for the words of Lemma~\ref{lem:c4-4-2}, and the undifferentiated closed source functional $| c_W(U_1;h)|$, the physical and reference data having the same numerical majorant (Lemma~\ref{lem:c4-4-2} applies at the reference data too, whose hypotheses Lemma~\ref{lem:c2-reference} verifies).

\begin{lemma}[footprints and envelopes]\label{lem:c4-6-env}
The cubes are those of the partition $\pi_1$ (side $LM$, Section~\ref{app:domain-geometry}): $\Delta_z= LMz+[-LM/2,LM/2]^4$ reduced modulo the fine periods, which are multiples of $LM$ by \eqref{eq:c-standing-conditions}; every site lies in exactly one cube, and the index distance of two cubes is the periodic maximum-norm distance of their indices $z$. For a set $S$ of sites, $\pi_1(S)$ is the set of cubes containing a site of $S$; throughout, a cube meets a set of sites or bonds (as in the footprint definition following \eqref{eq:c4-6-1}) when it contains a site or a bond endpoint of the set, the site reading; and $\pi_1(\mathcal H(W))$ abbreviates $\pi_1$ of the vertex set of the traversals of $\mathcal H(W)$. For a set $X$ of cubes, $\mathfrak N_R(X)$ is the set of cubes at index distance at most $R$ from a cube of $X$, so $|\mathfrak N_R(X)|\leqslant (2R+1)^4| X|$. The three ingredients of the footprint $F(W)$ have the following fixed ranges, in fine $\ell^{\infty}$ units. The free-input stencil $\mathcal F_r$ of a bond $r$ \eqref{eq:regular-stencil-extension-support} is $\{ r\}$ or $\mathcal E_c\cap E_*$, bonds with both endpoints in two adjacent blocks of side $L$ \eqref{eq:regular-stencil-bonds}, or empty (for a contour bond that is not a bond of $\Omega_1$, the convention of Section~\ref{sec:joint-extraction}), so every endpoint of a bond of $\mathcal F_r\cup\{ r\}$ is within $2L$ of the initial vertex of $r$; this is the contour/free-input stencil of a traversal. The coefficient stencil halo of a domain $D$ is the union of the hosts \eqref{eq:part1-potential-host}, \eqref{eq:c-probe-host} whose base lies within $R_{\mathrm h}$ of $D$, where $R_{\mathrm h}$ is the maximal base-to-host offset of Section~\ref{sec:c4-7-1}, fixed by the host constructions: the maximal $\ell^{\infty}$ distance from the base of a host to a point of the host, so that each host lies within $R_{\mathrm h}$ of its base (the active hosts of a selection with final domain $D$ are among these hosts, their bases lying in the $R_{\mathrm h}$-neighbourhood of $D$); the halo therefore lies within $2R_{\mathrm h}$ of $D$. A root contact $n\in\mathsf T_D$ \eqref{eq:part1-contact-set} has a bond $b\in\mathcal F_{r_n}\cap\mathsf L_D$, whose endpoints lie in $D$ ($Y_b\subseteq D$ by \eqref{eq:probe-eligible-carrier}, and $Y_b$ contains $b$ by its definition \eqref{eq:c-probe-host}, a rectangle containing the neighbourhood of the lift of $b$), so the initial vertex of $r_n$ is within $2L$ of $D$. Put $R_{\mathrm e}=\lceil2\max(R_{\mathrm h},L)/(LM)\rceil$, a constant of $L$ and $M$ fixed before the coupling, $A_F= A_E=9(2R_{\mathrm e}+1)^4$, $X(W)=\bigcup_D\pi_1(D)$ over the constituent domains $D$ of $W$ (roots and cluster constituents; a set, so $| X(W)|\leqslant n(W)$, which counts with multiplicity) and $\mathfrak E(W)=\mathfrak N_{R_{\mathrm e}}(X(W)\cup\pi_1(\mathcal H(W)))$, the envelope. Then: (a) $|\pi_1(\mathcal H(W))|\leqslant 8+|\mathcal H(W)|/(LM)$, $| X(W)\cup\pi_1(\mathcal H(W))|\leqslant 9\mathfrak s(W)$, $F(W)\subseteq \mathfrak E(W)$ and $| F(W)|\leqslant |\mathfrak E(W)|\leqslant A_E\mathfrak s(W)$; (b) $\mathfrak E(W)$ is connected under maximum-norm adjacency of cubes (two distinct cubes are adjacent when their index distance is one); (c) $\mathcal H(W)\neq\varnothing$, and if $\Delta\in F(W)$ then $\operatorname{dist}(\Delta,C)\leqslant A_ELM\,\mathfrak s(W)$, that is $\mathfrak s(W)\geqslant \operatorname{dist}(\Delta,C)/(A_ELM)$, where $\operatorname{dist}(\Delta,C)$ is the $\ell^{\infty}$ distance in fine units from the nearest site of $\Delta$ to the vertex set of $C$ (the proof bounds the distance from every site of $\Delta$); (d) two incompatible groups have footprints in contact (sharing a cube, or adjacent across a face, edge or corner, that is at index distance at most one) or overlapping intervals, so an incompatible $W'$ has a footprint cube in $\mathfrak N_1(F(W))$, a set of at most $3^4A_F\mathfrak s(W)$ cubes, or a traversal in $\mathcal H(W)$.
\end{lemma}

\begin{proof}
Index conversion: if sites $x\in\Delta_z$ and $y\in\Delta_{z'}$ are at fine distance at most $\rho$, then, for lifts of $x$, $y$ realising that distance and the indices $z$, $z'$ of their lifted cubes, $| LM(z-z')|_{\infty}\leqslant \rho+2\cdot\frac{LM-1}2<\rho+LM$, so the index distance of $\Delta_z$ and $\Delta_{z'}$ is at most $\lceil\rho/(LM)\rceil$; in particular the cube of a site within $2L$ of a cube $\Delta$ is at index distance at most $\lceil2L/(LM)\rceil=1$ from $\Delta$ ($M\geqslant 8$), and a cube meeting the halo of $D$ is at index distance at most $\lceil2R_{\mathrm h}/(LM)\rceil\leqslant R_{\mathrm e}$ from a cube of $D$. Every root $D$ of $W$ is marked, that is $\mathsf T_D\neq\varnothing$ (\eqref{eq:part1-contact-set} and the clause following it; the construction preceding Definition~\ref{def:groups}: every component has a marked root, and the groups are formed from the components), and $\mathsf T_D\subseteq\mathcal H(W)$, because $\mathcal H(W)$ is the hull of the union of the contacts of the group's components, which include the contact sets of its roots (the same construction, from ``Every component has a marked root'' to the definition of $\mathcal H(K)$). In particular $\mathcal H(W)\neq\varnothing$. (a) $\mathcal H(W)$ is a cyclic interval of fewer than $n_C/2$ traversals of the square $C$, or the whole contour (Definition~\ref{def:groups}), so the vertices of its traversals on each side form one run of consecutive collinear sites (the whole side, for the whole contour), and $C$ does not wrap around the torus ($LK_i\geqslant 3L(n+2)> nL$). A run of $t$ traversals along a side has $t+1$ consecutive collinear vertices, which lie in at most $\lfloor t/(LM)\rfloor+2$ cubes, since the cubes cut the side into runs of $LM$ consecutive sites; summing over the four sides, $|\pi_1(\mathcal H(W))|\leqslant 8+|\mathcal H(W)|/(LM)$. Hence $| X(W)\cup\pi_1(\mathcal H(W))|\leqslant n(W)+8+|\mathcal H(W)|/(LM)=\mathfrak s(W)+8\leqslant 9\mathfrak s(W)$ by \eqref{eq:c4-6-1} and $\mathfrak s(W)\geqslant n(W)\geqslant 1$ (a group has a nonempty root family, Definition~\ref{def:groups}, and each root is a connected final domain, a nonempty union of cubes, Section~\ref{sec:fixed-first-kernel}). A cube of $F(W)$ (its definition follows \eqref{eq:c4-6-1}) is a cube of a constituent domain, in $X(W)$; or meets the halo of a constituent domain, at index distance at most $R_{\mathrm e}$ from $X(W)$; or contains an endpoint of a bond of the stencil of a traversal of $\mathcal H(W)$, within $2L$ of the initial vertex of that traversal, hence at index distance at most $1\leqslant R_{\mathrm e}$ from $\pi_1(\mathcal H(W))$. So $F(W)\subseteq \mathfrak E(W)$ and $|\mathfrak E(W)|\leqslant (2R_{\mathrm e}+1)^4\cdot9\mathfrak s(W)= A_E\mathfrak s(W)$. (b) The set $X(W)\cup\pi_1(\mathcal H(W))$ is connected under adjacency: the two vertices of a traversal are at fine distance one and consecutive traversals of $\mathcal H(W)$ share a vertex, so the vertices lie in the same or in adjacent cubes, and $\pi_1(\mathcal H(W))$ is connected; each constituent domain is wall connected (Section~\ref{sec:fixed-first-kernel}), so its cubes are connected; within a cluster occurrence the constituent domains are linked by incompatibilities (the incompatibility graph of a cluster occurrence is connected: the Ursell coefficient of Lemma~\ref{lem:vacuum-clusters} is a sum over the connected spanning subgraphs of that graph, so a tuple with a disconnected graph has $\varphi^T=0$ and is dropped from \eqref{eq:c-cluster-expansion}, as stated after that display), a cluster occurrence of $W$ hits a root of $W$ (Definition~\ref{def:groups}: some constituent is incompatible with a root), and incompatible domains share a cube or a wall (Section~\ref{sec:selected-integrals}), that is contain cubes at index distance at most one; finally every root $D$ has a contact $n\in\mathsf T_D\subseteq\mathcal H(W)$ (shown above), the initial vertex of $r_n$ being within $2L$ of $D$, so a cube of $\pi_1(\mathcal H(W))$ is at index distance at most one from a cube of $D$. Thus every constituent domain is joined to $\pi_1(\mathcal H(W))$ through adjacent cubes, and $\mathfrak E(W)$, the union of the connected boxes $\mathfrak N_{R_{\mathrm e}}(\{\Delta\})$ over the cubes $\Delta$ of a connected set, each containing its centre $\Delta$, is connected. (c) $\mathcal H(W)\neq\varnothing$ was shown above. A connected set of $k$ cubes contains a path of at most $k$ cubes between any two of them, so the cubes of any two of its sites are at index distance at most $k-1$ and the sites are at fine distance less than $kLM$; with $k\leqslant A_E\mathfrak s(W)$, the $\ell^{\infty}$ diameter of $\mathfrak E(W)$ is at most $A_ELM\,\mathfrak s(W)$. $\mathfrak E(W)$ contains $\Delta\in F(W)$ and the nonempty set $\pi_1(\mathcal H(W))$, whose cubes contain vertices of $C$. (d) Two groups are compatible when all their constituent domains are mutually compatible and their intervals are disjoint, and a family containing a whole-contour group is a singleton (Definition~\ref{def:groups}). So if $W$ and $W'$ are incompatible, either one of them is a whole-contour group, and then the intervals overlap, as $\mathcal H(W)$ and $\mathcal H(W')$ are nonempty by (c); or the intervals overlap, so $W'$ has a traversal in $\mathcal H(W)$; or some constituent domains $D$ of $W$ and $D'$ of $W'$ share a cube or a wall, so a cube of $D'\subseteq F(W')$ lies at index distance at most one from a cube of $D\subseteq F(W)$, that is in $\mathfrak N_1(F(W))$, and $|\mathfrak N_1(F(W))|\leqslant 3^4| F(W)|\leqslant 3^4A_F\mathfrak s(W)$.
\end{proof}

\begin{lemma}[positive group sums]\label{lem:c4-6-2}
Choose $\eta>0$ after the fixed geometry and forest decay constants, as in Section~\ref{sec:c4-8-2} ($1000\eta A_{\mathrm s}<\varpi/8$ there). Under the smallness conditions $C\rho LM\leqslant \eta/8$ and the root-incidence bound of Section~\ref{sec:c4-8-2}, which hold for $g_0$ small with constants fixed before the coupling and independent of the loop length,
\begin{equation}\label{eq:c4-6-5}
    \sup_{n\in\mathcal I_C}\sum_{W\, | \,n\in\mathcal H(W)}M(W)e^{4\eta\mathfrak s(W)}+\sup_{\Delta\in\pi_1}\sum_{W\, | \,\Delta\in F(W)}M(W)e^{4\eta\mathfrak s(W)}\leqslant C\rho\;,
\end{equation}
with a constant $C$ fixed before the coupling (the constants of Lemma~\ref{lem:c4-8-2}(b) and of Section~\ref{sec:c4-8-3}), where $\mathcal I_C$ is the set of contour traversals; the bound holds in the form ``$\leqslant C\rho$ for every $\rho$ below the threshold'', the conditions being monotone in $\rho$. For footprint contact, wall adjacency and corner contact are included (Section~\ref{sec:c4-12-4}).
\end{lemma}

\begin{proof}
This is the combinatorial payment of Sections~\ref{sec:c4-8-2} and~\ref{sec:c4-8-3}; its structure is as follows. First form the connected components of the roots/cluster-occurrence graph before hull merging. The factorials are the exponential-generating factorials of these objects: roots form an unordered distinct family and cluster occurrences carry $1/k!$. Dropping the compatibility constraints is an upper bound, and a connected graph indicator is at most its sum of incompatibility spanning trees. For a root node $D$ the incident vacuum nodes have weighted mass at most $N_1(D)/8$ by the hit bound of Lemma~\ref{lem:vacuum-clusters} (its rooted form is \eqref{eq:c2-pinned-vacuum}), in the exponentially weighted form \eqref{eq:c4-7-9} with weight $e^{n(\mathfrak c)/4}$ that Lemma~\ref{lem:c4-8-2} uses (Section~\ref{sec:c4-12-4}); for a vacuum node $\mathfrak c$ the same estimate summed over its constituents gives $n(\mathfrak c)/8$; the added source-root mass is at most $C\rho$ times the same cost by the animal sum \eqref{eq:native-source-animal-sum} and Lemma~\ref{lem:c4-4-2}, hence at most that cost$/8$ for small $g_0$. The two-species positive tree recursion is thus bounded with node cost one quarter of the cube cost: if $t_v$ is a node weight and $a_v$ its cost, the iterates $T^{(k+1)}_v= t_v\exp(\sum_{v'\sim v}T^{(k)}_{v'})$ satisfy $T^{(k)}_v\leqslant t_ve^{a_v}$, because $\sum_{v'\sim v}t_{v'}e^{a_{v'}}\leqslant a_v$, and their monotone limit sums the rooted trees. This pays the vacuum decorations and leaves a factor $\rho$ for a component with a marked root. The same recursion pays the polynomial loss when a component is pinned at a point of its hull rather than at a contact: the hull length is bounded by the spatial extent, so a root lies within $CLM$ times the cube cost of that point, and the root locations and the polynomial in the cube count are absorbed by the unused exponential decay; a spatial footprint pin is handled by a tree path from the marked root to the pinned cube, and the square has bounded local traversal density, so the sum over source contacts weighted by that distance is uniform in $n_C$. Then merge the components whose hull intervals overlap, summing the connected-overlap indicators by interval spanning trees with the same recursion: a hull $\mathcal H$ meets components of total weighted mass at most $C\rho|\mathcal H|$, its node cost contains $\eta|\mathcal H|/(LM)$, and $C\rho LM\leqslant \eta/8$ pays the branching; the weighted union size is at most the sum of the component sizes, and the exponential-generating factorials agree with the interleavings of the cluster lists exactly as in Lemma~\ref{lem:grouped-source}. The exceptional merge (a long hull merges every remaining component into the whole-contour group) is paid in Lemma~\ref{lem:c4-8-3}: a long class costs at least $cn_C/(LM)$ in the sum of its component sizes, the arbitrary remaining components cost at most $e^{C\rho n_C}$, and at small $g_0$ this is absorbed into a fraction of the long-class decay, not left as a loop-length factor. Lemmas~\ref{lem:c4-8-2} and~\ref{lem:c4-8-3} carry out these steps with the explicit weights.
\end{proof}

\begin{lemma}[the reference group gas]\label{lem:c4-6-3}
Put $\mathfrak a_W=\eta\mathfrak s(W)$. Under Lemma~\ref{lem:c4-6-2}, for every group $W$,
\begin{equation}\label{eq:c4-6-6}
    \sum_{W'\not\sim W}| a_{W'}| e^{\mathfrak a_{W'}}\leqslant C\rho\bigl(| F(W)|+|\mathcal H(W)|\bigr)\leqslant \frac18\mathfrak a_W
\end{equation}
after choosing $g_0$ small, every constant fixed before the coupling and independently of the loop length. Hence the reference group gas \eqref{eq:part1-reference-gas} has a convergent connected-cluster logarithm, $S^{\circ}_{\xi,C}\ne0$, and for every compatible family $\mathfrak B$
\begin{equation}\label{eq:c4-6-7}
    \Bigl|\frac{S^{\circ,[\mathfrak B]}_{\xi,C}}{S^{\circ}_{\xi,C}}\Bigr|\leqslant \exp\Bigl(\frac18\sum_{W\in\mathfrak B}\mathfrak a_W\Bigr)\;.
\end{equation}
\end{lemma}

\begin{proof}
Incompatibility of $W$ and $W'$ implies contact of their footprints or overlap of their intervals, so the two pin sums \eqref{eq:c4-6-5} (with $| a_{W'}|\leqslant M(W')$ and $e^{\mathfrak a_{W'}}\leqslant e^{4\eta\mathfrak s(W')}$) give the first inequality of \eqref{eq:c4-6-6}, and $| F(W)|+|\mathcal H(W)|\leqslant C_1LM\mathfrak s(W)$ with $C\rho LM\leqslant \eta/8$, the constant $C$ there taken to include $C_1$, gives the second. This is the Kotecký--Preiss criterion for the species $W$ with activities $a_W$ and costs $\mathfrak a_W$; the pinned-tree argument of Section~\ref{app:kp} gives the convergent logarithm and the exclusion identity, whose countable-species form is Lemma~\ref{lem:c4-9-2}. The restricted ratio is the single joint exclusion ratio of that gas, and the pinned cluster estimate counts each cluster hitting $\mathfrak B$ once, giving \eqref{eq:c4-6-7}. There is no product-of-ratios assertion; this is the step that prevents a factor $e^{C\rho n_C}$ from replacing the flat scalar in $E_{\mathrm{def}}$.
\end{proof}

\subsubsection{Subtraction and its curvature payment}\label{sec:c4-6-3}
For a group $W$ its local comparison box $Q_W$ contains every constituent domain's coefficient halo and every background transport along $\mathcal H(W)$, including the inverse interval used to close the source; its lifted side is at most $CLMn(W)$ (Section~\ref{sec:c4-8-4}). For the good interior boxes of Section~\ref{sec:c4-8-4}, contained in $\Lambda_1$ and disjoint from $S_1$, root a tree gauge at the base of the complementary holonomy. The ordered plaquette filling of Lemma~\ref{lem:c4-10-5} gives a phase bound $C\operatorname{diam}(Q_W)^2\delta_{Q_W}(U_1)$, where $\delta_Q(U)$ is the maximum plaquette deviation $| U(\partial p)-\id|_{\sigma}$ over the plaquettes of $Q$. The source is a closed contraction throughout. The derivative bound of Section~\ref{sec:c4-8-4} (\eqref{eq:c4-8-10} with Lemmas~\ref{lem:c4-7-4} and~\ref{lem:c4-4-3}), applied to this coefficient with its cluster list and the whole relative interval retained, costs the factor $n(W)+|\mathcal H(W)|$ and the factor $\rho$ for each marked root; integrating the real phase path from the flat projection to $U_1$ (Lemma~\ref{lem:c4-9-1}) gives $\mathfrak d_W\leqslant C\delta_{Q_W}\operatorname{diam}(Q_W)^2$ times that polynomial times $M(W)$. If the phase bound exceeds the fixed interpolation chart, the two undifferentiated endpoint bounds give the same inequality with a larger constant (Section~\ref{sec:c4-12-2}). No $g_0$-dependent restriction on the group size is added.

The regional carrier $I$ is retained in the reference coefficient. The comparison is used only for groups whose box, with every coefficient and source halo, lies in $\Lambda_1$ and misses $S_1$; boxes crossing that integration boundary are kept in the large-domain class and are not interpolated, which leaves the target unchanged because their full contribution is paid by the tail below. On an interior box the centred kernel depends on the branch fields only through $U_1$ there: the retained field and the boundary change are absent from its source rows and primitive arguments, and $V_0$ enters only through that background (Section~\ref{sec:c4-8-1}). On far histories $\operatorname{dist}(\Gamma_C,E(\Lambda_1^c))> b$, and Section~\ref{sec:c4-8-4} uses the good comparison only when $n(W)< b/(4C_{\mathrm{box}}LM)$, which places every required stencil in an interior lifted box; all other groups, which have $n(W)\geqslant b/(4C_{\mathrm{box}}LM)$, keep both endpoint coefficients and the undifferentiated estimate, no assertion about their background continuation being made, and the reserved exponential pays this tail.

\begin{lemma}[the group pin sum]\label{lem:c4-6-4}
On covered far histories, with $\varepsilon_1= g_1p_0(g_1)$, $b=\lfloor LMR_1/16\rfloor$ and constants $C_{\chi}$, $c$ fixed before the coupling,
\begin{equation}\label{eq:c4-6-8}
    \sum_WK_{\chi}\mathfrak d_W(U_1)e^{\mathfrak a_W/8}\leqslant C_{\chi} n_C\rho\bigl(\varepsilon_1+e^{-cb/(LM)}\bigr)\;.
\end{equation}
\end{lemma}

The proof is Section~\ref{sec:c4-12-2}, from the local estimate of Section~\ref{sec:c4-8-4} and the pins of Lemma~\ref{lem:c4-6-2}. The polynomial factors $n(W)^k$, the squared box size and the interval length are paid by the reserved domain decay, not placed in a loop-dependent constant.

\subsubsection{The signed integrated bound}\label{sec:c4-6-4}
\begin{lemma}[the family sum]\label{lem:c4-6-5}
On covered far histories,
\begin{equation}\label{eq:c4-6-9}
    |\mathscr R_{\xi}|\leqslant | S^{\circ}_{\xi,C}|\Bigl[\exp\Bigl(C_{\chi} n_C\rho\bigl(\varepsilon_1+e^{-cb/(LM)}\bigr)\Bigr)-1\Bigr]\;.
\end{equation}
\end{lemma}

\begin{proof}
Use the sum \eqref{eq:c4-1-1}: absolute values are taken only after the flat subtraction has been made in each $R_{\xi}(\mathfrak B)$ and the joint reference ratio \eqref{eq:c4-6-7} has been retained. By \eqref{eq:c4-6-4} and \eqref{eq:c4-6-7}, $| S^{\circ,[\mathfrak B]}_{\xi,C}R_{\xi}(\mathfrak B)|\leqslant | S^{\circ}_{\xi,C}|\prod_{W\in\mathfrak B}K_{\chi}\mathfrak d_W(U_1)e^{\mathfrak a_W/8}$, and the sum over nonempty compatible families is bounded by the exponential series of the one-group sum \eqref{eq:c4-6-8} minus its zeroth term. The restricted reference sums, not a new vacuum normalisation, supply the factor $| S^{\circ}_{\xi,C}|$; at no stage is $\mathsf w_{\xi}$ cancelled across histories. The countable-group limit is Lemma~\ref{lem:c4-9-2}.
\end{proof}

Together with Lemma~\ref{lem:c4-2-1} this gives the first form of the defect bound, \eqref{eq:c4-12-7} below, with $\rho$ of \eqref{eq:c4-4-8} and all physical history weights and $\mathcal Z_0$ retained. It leaves the regional scalar $S^{\circ}_{\xi,C}$ inside the expectation; the identification of that scalar with the perimeter/corner coefficient is Appendix~\ref{app:regional-scalar} at the carrier $I$, and, for the labelled scalar $S^{\circ}_{\Lambda}$ at the retreated carrier, Section~\ref{sec:c5-7} (the defect bound at $D(\Lambda)$ is Proposition~\ref{prop:c5-6-def}); neither is part of this appendix. The reconstructed fine background loop is not replaced by the coarse loop.

\subsection{The joint kernel payment}\label{sec:c4-7}
This section supplies the quantitative steps used in Section~\ref{sec:c4-4}; it adds no assumption on $E_{\mathrm{def}}$.

\subsubsection{Host budget and fixed parameter order}\label{sec:c4-7-1}
$C_{\mathrm p}$ bounds one primitive (Lemma~\ref{lem:c2-primitive-bounds}) and $C_0= N_{\mathrm h}C_{\mathrm p}e^{\kappa(B_{\mathrm h}+4)}$ bounds their sum at a host (Lemma~\ref{lem:c2-potential-size}). Let $C^{\partial}_{\mathrm p}$ be the maximum of the finite local coefficient-derivative bounds of Section~\ref{sec:c4-10-1}: the energy primitive is $g_0$ times the third-order Taylor remainder in the protected variables, and its background derivative has the same factor $g_0(1+r)^3$ because its local analytic coefficients and their background derivatives are bounded on the larger fixed row chart; the Haar primitive and its background derivative start at order $g_0^2(1+r)^2$; the constraint-Jacobian primitive and its derivative are bounded by $Cg_0(1+r)$. Set
\begin{equation}\label{eq:c4-7-1}
    C^{\partial}_0= N_{\mathrm h}\max(C_{\mathrm p},C^{\partial}_{\mathrm p})e^{\kappa(B_{\mathrm h}+4)},\qquad
    \omega(g_0)= C^{\partial}_0g_0(1+r)^3,\qquad R_{\Upsilon}=\frac{\omega_*}{2\omega(g_0)}\;,
\end{equation}
$C^{\partial}_0$ chosen after $M$ and before $g_0$, as $C_0$ is ($\omega(g_0)$ and $R_{\Upsilon}$ then depend on $g_0$ through $g_0(1+r)^3$); neither $\omega_*$ nor $\varepsilon_*$ nor any earlier choice is shrunk after $M$. Summing the same finite host assignments proves $\sup|\mathscr D_{\omega}\mathbf V_{\xi}(Y)|\leqslant \omega(g_0)e^{-\kappa d_0(Y)}\|\omega\|_{\infty}$ on the background neighbourhood of Section~\ref{sec:c4-10-1} (phases within $\tau_{\mathrm b}$ of the branch background, Lemma~\ref{lem:c4-7-3}), where the primitive bounds hold with the same constants, so that $\omega(g_0)\geqslant C_0g_0(1+r)^3$ also bounds $\sup|\mathbf V_{\xi}(Y)| e^{\kappa d_0(Y)}$ there (\eqref{eq:native-source-interaction-size}); in Sections~\ref{sec:c4-4} and~\ref{sec:c4-6} this $\omega(g_0)$ replaces the $\omega(g_0)$ of \eqref{eq:part1-host-multiplicity}, which it dominates. The coupling threshold is $\omega(g_0)\leqslant \omega_*/4$, that is $R_{\Upsilon}\geqslant 2$ (Lemma~\ref{lem:c4-4-1}). Let $N_{\mathrm{type}}$ be the number of tagged ordinary primitive and probe types based in one fine translation cell of the fixed $L$-periodic template and $R_{\mathrm h}$ the maximal base-to-host offset, both fixed by the host constructions \eqref{eq:part1-potential-host}, \eqref{eq:c-probe-host}. Distinct active hosts of a selection with final domain $D$ have a labelled base in the fixed $R_{\mathrm h}$-neighbourhood of $D$, so their number is at most $H_MN_1(D)$ with $H_M= N_{\mathrm{type}}(LM+2R_{\mathrm h}+3)^4$, and the selected probe count is at most the number of fine bonds in $D$, $q\leqslant 4(LM+1)^4\sum_DN_1(D)$. The quantities $\omega(g_0)$ and $R_{\Upsilon}$ of \eqref{eq:c4-7-1} are used throughout Sections~\ref{sec:c4-4} and~\ref{sec:c4-6}; only the coupling threshold changes.

\begin{lemma}[positive host comparison]\label{lem:c4-7-1}
For a host $Y$ let $n_Y\leqslant P_M$ be the number of probe labels with host $Y$ and put
\begin{equation}\label{eq:c4-7-2}
    A_Y=\frac{\omega_*}{2}e^{-\kappa d_0(Y)},\qquad B_Y= n_Yu_*\leqslant \frac{\omega_*}{8}e^{-\kappa d_0(Y)}
\end{equation}
(the inequality by $n_Y\leqslant P_M$, $u_*\leqslant \omega_*e^{-\kappa s_{\mathrm p}}/(8P_M)$ from \eqref{eq:c4-4-6}, and $d_0(Y)\leqslant N_0(Y)-1< s_{\mathrm p}$ for a probe host, \eqref{eq:c2-tree-length} and Lemma~\ref{lem:c4-4-1}). (a) With $v_{\mathrm{ordinary}}= e^{A_Y}-1$ and $v_{\mathrm{probe}}= u_*$, the sum over the nonempty label subsets at $Y$ satisfies
\begin{equation}\label{eq:c4-7-3}
    \sum_{\varnothing\ne S\subseteq\{\mathrm{ordinary},\mathrm{probes}\}}\prod_{\ell'\in S}v_{\ell'}\leqslant e^{A_Y}(1+u_*)^{n_Y}-1\leqslant e^{A_Y+B_Y}-1\leqslant (A_Y+B_Y)e^{A_Y+B_Y}\leqslant \omega_*e^{1/2}e^{-\kappa d_0(Y)}\;.
\end{equation}
(b) For a differentiated host, the summed weight of the label subsets at $Y$ is bounded by
\begin{equation}\label{eq:c4-7-4}
    R_{\Upsilon}\sup|\mathscr D_{\omega}\mathbf V_{\xi}(Y)|\,e^{A_Y}(1+u_*)^{n_Y}\leqslant A_Ye^{A_Y+B_Y}\|\omega\|_{\infty}\leqslant \omega_*e^{1/2}e^{-\kappa d_0(Y)}\|\omega\|_{\infty}\;.
\end{equation}
(c) There is a constant $C_M$, fixed after $M$ and before $g_0$, such that for every connected final domain $D$
\begin{equation}\label{eq:c4-7-5}
    \sum_{\ell\, | \,Z(\ell)= D}u_*^{q(\ell)}\sup_{|\zeta|\leqslant R_{\Upsilon}}\|\mathscr D_{\omega}\Lambda^{\mathrm f}_{\xi;\ell,\zeta}\|_{\mathrm{TV}}\leqslant C_MN_1(D)\|\omega\|_{\infty}\,\varepsilon_*e^{-\lambda_*\ell_1(D)}\;,
\end{equation}
the derivative including every Gaussian normaliser and the compensator; and $|\mathscr D_{\omega} k^{\mathrm f}_{\xi}(D)|\leqslant C_MN_1(D)\|\omega\|_{\infty} m(D)$ for the vacuum activities, without shrinking $\varepsilon_*$.
\end{lemma}

\begin{proof}
The bound $n_Y\leqslant P_M$ holds because a probe label $b$ lies in its host, whose fine side is at most $4M$, and there are four positive directions per site. After taking absolute values inside the real Mayer $t$-integral the ordinary factor is at most $e^{A_Y}-1$ on $|\zeta|\leqslant R_{\Upsilon}$ (as $R_{\Upsilon}\omega(g_0)=\omega_*/2$). The last inequality of \eqref{eq:c4-7-3} uses $\omega_*\leqslant 1/2$ (the standing choice of Appendix~\ref{app:thresholds}) and $A_Y+B_Y\leqslant \frac58\omega_*e^{-\kappa d_0(Y)}$. This is a positive numerical comparison, not the Mayer factor of a complex combined potential. All split selections with the same active-host set have the same $Y_0$, seed, conditioning carrier and Gaussian kernel and, at fixed bad-bond choice $\mathsf P$ and final domain $Z$, the same final-domain weakening (proof of Lemma~\ref{lem:probe-measures}), so their many-to-one sum is dominated at each host by \eqref{eq:c4-7-3} before the cover and seed sums of Appendix~\ref{app:background-estimates} are applied; this is the domination used in Lemma~\ref{lem:c4-4-1}, the final domain sum being the chain \eqref{eq:native-source-initial-domain-sum}, \eqref{eq:native-source-enlarged-domain-sum} and the final $Y'$-sum in the proof of Lemma~\ref{lem:activity-bounds}. (b) is the same comparison with the differentiated host bound of Section~\ref{sec:c4-7-1}. (c) There are at most $H_MN_1(D)$ active hosts; the Gaussian and walk derivatives have the same factor after their polynomial moment is paid (Sections~\ref{sec:c4-10-2}, \ref{sec:c4-10-3}): the uniform Gaussian coercivity gives polynomial moments at most $C_MN_1(D)$ times the same exponential volume majorant $e^{N_1(D)}$, so no fresh exponential reserve is spent. Executing the differentiated sum with the positive host comparison (b) gives \eqref{eq:c4-7-5}, and the same sum over the selections with $\mathsf J=\varnothing$ gives the vacuum-activity bound.
\end{proof}

\subsubsection{The multi-root cancellation}\label{sec:c4-7-2}
The input $f$ is a single normalised character of the ordered closed reconstructed word, optionally closed with its inverse background interval and a fixed complementary $h$; for a linear combination the bound is multiplied by the sum of the absolute coefficients, and an arbitrary unbounded sum is not an input. On covered branches its contour input labels lie in $I$, so the retained $A_{\mathrm e}$ and the boundary map do not enter the source rows (Lemma~\ref{lem:covered}); the centred kernel $K_{\xi;0}$ likewise uses zero retained input, and the reference group construction uses the same row and mean conventions.

\begin{lemma}[commuting projections]\label{lem:c4-7-2}
Let $O_0$ be the set of roots with odd $q_j$ and empty ordinary selection, $O_1$ the set of roots with odd $q_j$ and nonempty ordinary selection. Then
\begin{equation}\label{eq:c4-7-6}
    \Bigl|\int\Delta_{\boldsymbol{\mathsf J}}f\,\dd\bigotimes_j\Lambda^{\mathrm f}_{\xi;\ell_j}\Bigr|\leqslant u^q\prod_{j\in O_0\cup O_1}\Bigl(u+\frac{2}{R_{\Upsilon}-1}\Bigr)\prod_jM_{\ell_j}\;,
\end{equation}
with the exponential-in-$q$ constants absorbed in $C_{\chi}$; after one background derivative the same holds with the factor $M_{\ell_j}$ of the differentiated root replaced by $\sup_{|\zeta|\leqslant R_{\Upsilon}}\|\mathscr D_{\omega}\Lambda^{\mathrm f}_{\xi;\ell_j,\zeta}\|_{\mathrm{TV}}$, which \eqref{eq:c4-7-5} sums, and with the additional factor of Lemma~\ref{lem:c4-4-3}. Moreover $\Delta_{\boldsymbol{\mathsf J}}f=0$ if some selected label does not influence the closed source word, so every surviving root has a contact with $\Gamma_C$.
\end{lemma}

\begin{proof}
For $j\in O_0$ let $P_j$ take the homogeneous part of degree $q_j$ of $\Delta_{\boldsymbol{\mathsf J}}f$ in the selected variables of root $j$, the other root variables fixed; these projections commute. Expand $\prod_{j\in O_0}(P_j+(1-P_j))$. A term containing $P_j$ is odd under the reversal of exactly root $j$ and integrates to zero against that root's even source-free density; only $\prod_{j\in O_0}(1-P_j)$ remains, with one extra $u$ per such root by the mixed Taylor remainder \eqref{eq:c4-10-9} and the fixed row-overlap count. For $j\in O_1$ the independent $\zeta_j$-integral vanishes at $\zeta_j=0$, so the multivariable Cauchy difference supplies $\prod_{j\in O_1}2/(R_{\Upsilon}-1)$ times the supremum $\prod M_{\ell_j}$ of \eqref{eq:c4-4-5}. Together this is \eqref{eq:c4-7-6}, and for even $q_j$ only $u^{q_j}M_{\ell_j}$ is paid. After one background derivative the Gaussian derivative is even, the leading source derivative has the same parity in the selected variables (a background derivative hits a row coefficient or a background occurrence of the closed word, not the selected root variables, Section~\ref{sec:c4-10-4}; the resulting bound is \eqref{eq:c4-10-10}) and a differentiated Mayer factor still vanishes at $\zeta_j=0$, so the argument repeats. The last statement is the pairwise cancellation of $K$ with $K\cup\{ b\}$ in the $K$-sum \eqref{eq:c4-4-1} for a label $b$ of a root that the word does not see; it is the pin used in the $\mathcal A$- and $W$-sums, and no sum over unmarked roots in the volume is claimed to be uniform.
\end{proof}

\subsubsection{The common Hessian bounds on the extension}\label{sec:c4-7-3}
\begin{lemma}[Hessian bounds along the coefficient path]\label{lem:c4-7-3}
Choose $\tau_{\mathrm b}>0$ once, after $\tau_L$, $c_L$, $\gamma_L$ and before $g_0$, with $\tau_{\mathrm b}\leqslant \tau_L/4$, $c_L\tau_{\mathrm b}\leqslant \gamma_L/4$ and the chart conditions of Section~\ref{sec:c4-10-1} (the matrix-logarithm neighbourhood and the entry-chart norm, each a smallness condition on $\tau_{\mathrm b}$ in terms of constants fixed before $g_0$); this one $\tau_{\mathrm b}$ is used throughout Sections~\ref{sec:c4-9}--\ref{sec:c4-12}. For every $t\in[0,1]$ and every globally small extension $U= e^{itB}U^{\circ}$ with $\| B\|_{\infty}\leqslant \tau_{\mathrm b}$ of the flat background on a comparison box, every compression of the regional Hessian to the same $I$ satisfies
\begin{equation}\label{eq:c4-7-7}
    \frac34\gamma_L\,\id\leqslant H_{e^{itB}}\leqslant m_L\,\id\;.
\end{equation}
\end{lemma}

\begin{proof}
On physical branches Lemma~\ref{lem:regular-bounds} gives $\gamma_L/2\leqslant H_{\xi,II}\leqslant m_L$ with the force--constraint and soft-gauge terms. The local entry formula $f_{bb'}(A)$ of Lemma~\ref{lem:c2-regional-entry} has row sums at most $m_L$ on the phase box $\tau_L$ by the definition of $m_L$ in \eqref{eq:regular-covariance-entry-constants}, column sums equal to its row sums since $f_{bb'}= f_{b'b}^{\mathsf T}$ have equal norms, entries vanishing beyond the range of that lemma, and Lipschitz bound $c_L$ in the phase; so the extension has operator difference from $H^{\circ}=\mathsf H^{\Omega_1}_{\id}$ at most $c_L\tau_{\mathrm b}\leqslant \gamma_L/4$ by the Schur test, uniformly in $t$. The flat lower bound $\gamma_L$ (\eqref{eq:covariance-uniform-gap}) passes to every compression to $I$ by the compression argument in the proof of Lemma~\ref{lem:c2-reference}, and with the difference bound this gives \eqref{eq:c4-7-7}. The recipe $\gamma=\gamma_L/2$, $h_+= m_L$ of Section~\ref{sec:fixed-first-kernel} therefore works at the physical endpoint and throughout the extension; no new expansion instantiation or coefficientwise identification of different recipes is used, and nothing is asserted about the global closeness of the physical background to the identity.
\end{proof}

\subsubsection{Differentiating the joint vacuum ratio}\label{sec:c4-7-4}
$\Phi^{\mathrm v}_{\xi}(\mathfrak c)$ is the Ursell expression $\varphi^T(\mathfrak c)\prod_jk^{\mathrm f}_{\xi}(D_j)/m!$ of \eqref{eq:part1-vacuum-cluster} with the fixed-kernel vacuum activities; that formula is used directly.

\begin{lemma}[derivative of the cluster sum]\label{lem:c4-7-4}
Let $M_v(\mathfrak c)$ be the positive incompatibility-tree majorant of $\Phi^{\mathrm v}_{\xi}(\mathfrak c)$ built from $m(D)$ (Section~\ref{sec:c4-6-2}) and $n(\mathfrak c)=\sum_{D\in\mathfrak c}N_1(D)$. Then
\begin{equation}\label{eq:c4-7-8}
    |\mathscr D_{\omega}\Phi^{\mathrm v}_{\xi}(\mathfrak c)|\leqslant C_M\|\omega\|_{\infty}\,n(\mathfrak c)M_v(\mathfrak c)\;,
\end{equation}
and for every finite set $\mathcal A$ of connected final domains
\begin{equation}\label{eq:c4-7-9}
    \sum_{\mathfrak c\text{ hits }\mathcal A}M_v(\mathfrak c)e^{n(\mathfrak c)/2}\leqslant \frac18\sum_{D\in\mathcal A}N_1(D),\qquad
    \sum_{\mathfrak c\text{ hits }\mathcal A}|\mathscr D_{\omega}\Phi^{\mathrm v}_{\xi}(\mathfrak c)|\leqslant \frac{C_M}{2}\|\omega\|_{\infty}\sum_{D\in\mathcal A}N_1(D)\;.
\end{equation}
\end{lemma}

\begin{proof}
\eqref{eq:c4-7-8} is the product rule over the constituents with the activity derivative bound of Lemma~\ref{lem:c4-7-1}(c): a sum over the possible differentiated vertices, not a division by an activity. The numerical criterion of Lemma~\ref{lem:activity-bounds}, \eqref{eq:native-source-weighted-vacuum-criterion}, is stated for $| k^{\mathrm f}_{\xi}|$; its proof uses only $| k^{\mathrm f}_{\xi}(D)|\leqslant m(D)=\varepsilon_*e^{-\lambda_*\ell_1(D)}$, so it gives $\sum_{D'\not\sim D}m(D')e^{(1+6\theta)N_1(D')}\leqslant N_1(D)/8$ as well. Apply the rooted-tree proof of Lemma~\ref{lem:vacuum-clusters} (Appendix~\ref{app:kp}) to the activity $m(D)e^{N_1(D)/2}$ with node cost $N_1(D)/8$, the combined exponent $5/8$ being below $1+6\theta$ (Section~\ref{sec:c4-12-4} records this $n/2$ form beside the $n/4$ form). This gives the first bound of \eqref{eq:c4-7-9}, the family version following by summing the single-root bound over $D\in\mathcal A$ (clusters hitting several roots are overcounted in this majorant only, not in the exact ratio). Since $n\leqslant 4e^{n/4}\leqslant 4e^{n/2}$, inserting \eqref{eq:c4-7-8} gives the second bound. Inserting the second bound into the derivative of the exponential \eqref{eq:c4-4-11} proves \eqref{eq:c4-4-12}: $C_M$ is an outside prefactor for a single derivative and is not put into every cluster weight, so no $M$-dependent shrinking of $\omega_*$ or $\varepsilon_*$ is needed after those constants have been fixed. The same $m(D)$ controls the undifferentiated activity along the whole extension ($| k^{\mathrm f}_{\xi}(D)|\leqslant m(D)$ there by Lemma~\ref{lem:c4-7-3} and the host bounds of Section~\ref{sec:c4-10}) and its one-derivative polynomial majorant, the polynomial being paid by the displayed exponential moment.
\end{proof}

\subsection{Quantitative group and caller checks}\label{sec:c4-8}
\subsubsection{Carrier, source and physical weights}\label{sec:c4-8-1}
All group comparisons and flat-subtraction identities are made on $A_{\mathrm{cov}}$. There the source rows have their free inputs in $I$, so Lemma~\ref{lem:covered} gives $U_{\xi}(0)(r)= U_1(r)$ for every contour bond $r$, and the free vector $B= a+\Psi_{\xi}(A_{\mathrm e})$ equals $a$ on $I$ ($\Psi_{\xi}(A_{\mathrm e})$ vanishes on $I$, proof of that lemma), so neither $A_{\mathrm e}$ nor $\Psi_{\xi}$ enters the fluctuating source rows; Lemma~\ref{lem:collar} and \eqref{eq:c-collar-identity} give $K_{\xi;1}$, and $K_{\xi;0}$ is the $s=0$ member of the interpolation \eqref{eq:c4-11-1}, with $H= H_{\xi,II}$, $P_{\xi,0}(a)=\sum_{\mathfrak p\ \mathrm{near}}\mathfrak p(a)$ and the chosen $r$-box, every primitive argument having zero retained input. The boundary Jacobian, the far primitives and the retained quadratic term are in $c_{\xi}$, not in the centred coefficients, and every fixed-kernel operation, probe measure and vacuum activity is built from those centred objects. Hence the local background data of a coefficient are $U_1$, the unchanged carrier geometry and the fixed coupling/cutoff/mean prescription; $V_0$ is not replaced in the physical integral, and its influence on these local coefficients is through $U_1$.

The local-data comparison of Section~\ref{sec:c4-11} applies on a box containing all source and coefficient stencils and lying in $\Lambda_1$ outside $S_1$. The fixed-kernel support statements (Section~\ref{sec:c4-10-2}) place every surviving walk coefficient, seed inverse and primitive stencil there, and the regional $I$ and $\Omega_1$ are kept on both sides of the interpolation. At the flat endpoint these local formulas are the reference formulas for $a_W$, since their arguments contain no retained field, boundary map or additional $V_0$: an equality of the selected coefficient, not of physical densities. The force--constraint and relative soft-gauge terms are those of the regional Hessian \eqref{eq:c-regional-hessian} and are included throughout.

Under a simultaneous site gauge each Hessian entry, restricted inverse and fixed walk coefficient is conjugated by the free-field adjoint action (Lemma~\ref{lem:hessian-gauge-locality}, Lemma~\ref{lem:walk-series}), the protected cutoffs are invariant, and the row solution is covariant; orthogonal changes of the integration variables leave the closed selected integrals unchanged, including their spatial Cauchy integrals, and their Laplace transforms give the covariance of each probe measure; the vacuum activities and Ursell products are scalar invariants; finally $T_WP_W^{-1}$ and $h$ are simultaneously conjugated at the interval base. This is the gauge law of the group closed contraction factor by factor (Section~\ref{sec:c4-11-2}); no straight-side assumption enters it.

$\mathcal O$ is the positive outer operation \eqref{eq:native-source-explicit-outer-operation}: the Haar and constraint factors are measures and the displayed cutoffs and exponentials are nonnegative, and $0<\mathcal Z_0<\infty$ (finite by the measurability paragraph of Section~\ref{sec:branch-data}; $\mathcal Z_0=\int\mathcal O[\mathcal J^{\mathrm a}_{\xi}[1]]\,\dd V=\int\mathcal O[\mathsf w_{\xi}]\,\dd V$ is the normaliser of \eqref{eq:native-source-physical-ratio}, the integral of a positive function, since $\nu_{\xi;1}$ is a probability measure and so $\mathcal J^{\mathrm a}_{\xi}[1]=\mathsf w_{\xi}$ by \eqref{eq:c-collar-identity}), so $\mathbb E_{\mathrm h}$ is a positive normalised functional. Multiply \eqref{eq:c-one-step-branch-identity} by $\mathbf 1_{A_{\mathrm{cov}}}$ and write $\mathcal J^{\mathrm a}_{\xi}[f_{\xi}]=\mathsf w_{\xi} F_{\xi}(1)$ on $A^c_{\mathrm{cov}}$; every term is integrable against the finite positive measure $\mathcal O[\mathsf w_{\xi}\,\bigcdot \,]\,\dd V$, because $| S^{\circ}_{\xi,C}f_{\xi}(0)|=| F_{\xi}(0)-\mathscr R_{\xi}|\leqslant 3$ by \eqref{eq:part1-scalar-extraction}, $| F_{\xi}(0)|\leqslant 1$ and Lemma~\ref{lem:c4-2-1}, $| F_{\xi}(1)|\leqslant 1$, and $|\mathcal B_{\xi}[f_{\xi}]|\leqslant C_{\mathrm{resp},\chi}g_0r| I_C|$ on $A_{\mathrm{cov}}$ (Section~\ref{sec:c3-6}). With $E_{\mathrm{bdry}}$ of \eqref{eq:c3-6-2}, the full identity is
\begin{equation}\label{eq:c4-8-1}
    \mathbb E_{\varepsilon}[W_{\chi}(C)]=\mathbb E_{\mathrm h}\bigl[\mathbf 1_{A_{\mathrm{cov}}}S^{\circ}_{\xi,C}f_{\xi}(0)\bigr]+E_{\mathrm{def}}+E_{\mathrm{bdry}}\;,
\end{equation}
by \eqref{eq:c-one-step-branch-identity} and \eqref{eq:native-source-physical-ratio}. The far uncovered event is empty: if coverage fails, either a contour bond $r$ is not a bond of $\Omega_1\supseteq\Lambda_1$, so that $r$ has an endpoint in $\Lambda_1^c= B_{\xi}$, or a bond of a stencil $\mathcal F_r\not\subseteq I$ is a free bond with no endpoint in $\Lambda_1$ ($I=\Lambda_1^*$, Section~\ref{sec:branch-data}); in both cases an element of $\Gamma_C$ (a contour bond or a bond of a free-input stencil) lies in $E(B_{\xi})$, so $\operatorname{dist}(\Gamma_C,E(B_{\xi}))=0$ and $A_{\mathrm{cov}}^c\subseteq\mathsf N_0\subseteq\mathsf N_b$ (\eqref{eq:c4-2-1}; the events $\mathsf N_b$ increase with $b$). The exterior response is not dropped: it is the separately controlled $E_{\mathrm{bdry}}$ of \eqref{eq:c3-6-2} and Appendix~\ref{app:exterior-chessboard}, not a part of $E_{\mathrm{def}}$.

\subsubsection{Two quantitative pins}\label{sec:c4-8-2}
A component $\gamma$ consists of a nonempty unordered root family $\mathcal A(\gamma)$ and an ordered list $(\mathfrak c_1,\ldots,\mathfrak c_k)$ of cluster occurrences with connected incidence graph as in Section~\ref{sec:joint-extraction}; every root is a marked domain, so $\gamma$ has at least one marked root and $\mathsf T_{\gamma}\neq\varnothing$ (only families of marked domains contribute, Section~\ref{sec:joint-extraction}). Here $n_{\gamma}$ counts every constituent cube with its occurrence multiplicity, $\mathsf T_{\gamma}$ is the union of the contact sets of its roots, $H_{\gamma}$ the hull of $\mathsf T_{\gamma}$ (a cyclic interval of fewer than $n_C/2$ traversals, or the whole contour, Section~\ref{sec:joint-extraction}), $X_{\gamma}$ the union of its domains (roots and cluster constituents), and $F(\gamma)$ its footprint: the $LM$-cubes meeting a domain of $X_{\gamma}$, its fixed coefficient stencil halo, or the contour/free-input stencil of a traversal of $H_{\gamma}$, as $F(W)$ in Section~\ref{sec:c4-6}, with the contact conventions of Section~\ref{sec:c4-12-4} (wall adjacency and corner contact) used in \eqref{eq:c4-6-5}. A class of components (defined below) is long when the union of the hulls of its members has at least $n_C/2$ positions, in which case all components merge into the whole-contour group (Section~\ref{sec:joint-extraction}). Define
\begin{equation}\label{eq:c4-8-2}
    t_{\gamma}=\frac{\rho^{|\mathcal A(\gamma)|}}{k!}\prod_{D\in\mathcal A(\gamma)}m(D)\prod_{j=1}^kM_v(\mathfrak c_j)\;.
\end{equation}
Summing unordered distinct roots is overcounted by ordered root tuples divided by their factorial. A class $K$ is an unordered family of components with connected hull-overlap graph, with weight $t_K=\prod_{\gamma\in K}t_{\gamma}$; for the positive recursion the sum over classes is overcounted by ordered tuples divided by the tuple factorial, dropping the disjointness restrictions, which only enlarges the sum and fixes every multiplicity convention. Fix
\begin{equation}\label{eq:c4-8-3}
    0<\varpi\leqslant \min\Bigl(\frac1{16},\frac{\lambda_*}{64C_{\mathrm g}}\Bigr)\;,
\end{equation}
with $C_{\mathrm g}$ of Lemma~\ref{lem:c2-domain-geometry}.

\begin{lemma}[component pins]\label{lem:c4-8-2}
There are constants $C_{\mathrm{pin}}$, $A_{\mathrm h}$, $A_{\mathrm x}$ fixed before the coupling such that, for $\rho$ below a threshold depending on them: (a) for every root $D_0$,
\begin{equation}\label{eq:c4-8-4}
    \sum_{\gamma\, | \,D_0\in\mathcal A(\gamma)}t_{\gamma} e^{\varpi n_{\gamma}}\leqslant \rho\,m(D_0)e^{(\varpi+1/4)N_1(D_0)},\qquad
    \sup_{a\in\mathcal I_C}\sum_{\gamma\, | \,a\in\mathsf T_{\gamma}}t_{\gamma} e^{\varpi n_{\gamma}}\leqslant C_{\mathrm{pin}}\rho\;;
\end{equation}
(b) $| H_{\gamma}|\leqslant A_{\mathrm h}LMn_{\gamma}$, and the distance from a source contact of $\gamma$ to any cube of $X_{\gamma}$, halos included, is at most $A_{\mathrm x}LMn_{\gamma}$; with $\theta_{\mathrm h}=\varpi/(4A_{\mathrm h}LM)$, $\theta_{\mathrm x}=\varpi/(4A_{\mathrm x}LM)$ and $d_C$ the contour distance,
\begin{gather}
    \sum_{\gamma\, | \,a\in H_{\gamma}}t_{\gamma} e^{\varpi n_{\gamma}/4}\leqslant \sum_{a'\in\mathcal I_C}e^{-\theta_{\mathrm h}d_C(a,a')}\sum_{\gamma\, | \,a'\in\mathsf T_{\gamma}}t_{\gamma} e^{\varpi n_{\gamma}/2}\leqslant C_{\mathrm{pin}}\rho\Bigl(1+\frac{2}{e^{\theta_{\mathrm h}}-1}\Bigr)\;,\label{eq:c4-8-5}\\
    \sum_{\gamma\, | \,\Delta\subset X_{\gamma}}t_{\gamma} e^{\varpi n_{\gamma}/4}\leqslant C_{\mathrm{pin}}\rho\sum_{a'\in\mathcal I_C}e^{-\theta_{\mathrm x}d_T(a',\Delta)}\leqslant C\rho\;,\label{eq:c4-8-6}
\end{gather}
with $d_T(a',\Delta)$ the periodic fine distance (Section~\ref{sec:c4-2}) from the bond of the traversal $a'$ to the cube $\Delta$, and the footprint pin $\sum_{\gamma\, | \,\Delta\in F(\gamma)}t_{\gamma} e^{\varpi n_{\gamma}/4}\leqslant C\rho$ with the constant $C$ enlarged by the stencil factor $CLM$; (c) with $A_{\mathrm s}=1+A_{\mathrm h}$ and $\eta>0$ chosen after these constants with $1000\eta A_{\mathrm s}<\varpi/8$, the three pins of (b) hold with the weight $e^{100\eta\mathfrak s(\gamma)}$, $\mathfrak s(\gamma)= n_{\gamma}+| H_{\gamma}|/(LM)$, leaving a fixed exponential margin.
\end{lemma}

\begin{proof}
(a) The sum is a two-species positive tree recursion, the recursion of the proof of Lemma~\ref{lem:vacuum-clusters} (Appendix~\ref{app:kp}) with two kinds of nodes: the roots $D$, with weight $\rho\,m(D)e^{\varpi N_1(D)}$ and cost $N_1(D)/4$, and the cluster occurrences $\mathfrak c$, with weight $M_v(\mathfrak c)e^{\varpi n(\mathfrak c)}$ and cost $n(\mathfrak c)/4$; a connected incidence graph is bounded by the sum of its spanning trees, the factorials are the exponential-generating factorials of the unordered roots and the ordered occurrence list, and the iterates $T^{(k+1)}_v= t_v\exp(\sum_{v'\sim v}T^{(k)}_{v'})$ satisfy $T^{(k)}_v\leqslant t_ve^{a_v}$ once $\sum_{v'\sim v}t_{v'}e^{a_{v'}}\leqslant a_v$ for every node $v$ (weights $t_v$, costs $a_v$). The criterion holds. The occurrences hitting a node have weighted mass at most one half of its cost, by \eqref{eq:c4-7-9} applied to the node's constituents, since $e^{\varpi n}e^{n/4}\leqslant e^{n/2}$ ($\varpi+1/4<1/2$). A root hitting a node is incompatible with one of its constituents, so it contains one of the at most $9N_1$ cubes meeting a constituent cube or sharing a wall with it ($N_1$ the node's cube count); for a fixed cube $Q$,
\begin{equation*}
    \sum_{D\ni Q}m(D)e^{(\varpi+1/4)N_1(D)}\leqslant \varepsilon_*\sum_{D\ni Q}e^{-(\lambda_*-(\varpi+1/4)C_{\mathrm g})\ell_1(D)}\leqslant \varepsilon_*B(\lambda_*/2)\eqqcolon C_{\mathrm a}\;,
\end{equation*}
by $N_1(D)\leqslant C_{\mathrm g}\ell_1(D)$ (Lemma~\ref{lem:c2-domain-geometry}(1); the final domains are localisation domains) and the animal bound \eqref{eq:native-source-animal-sum}, applicable because $\lambda_*-(\varpi+1/4)C_{\mathrm g}\geqslant \lambda_*-\tfrac5{16}C_{\mathrm g}\geqslant \lambda_*/2> C_{\mathrm g}\log64$ by $\varpi\leqslant 1/16$ and $\lambda_*\geqslant 2C_{\mathrm g}\log128$ (\eqref{eq:c2-kappa-choice}). So the roots hitting a node have weighted mass at most $9\rho C_{\mathrm a}N_1\leqslant N_1/32$, one eighth of the cost, once $9\rho C_{\mathrm a}\leqslant 1/32$, a condition on $\rho$ alone, that is on $g_0$. The rooted sum at $D_0$ is therefore at most $\rho\,m(D_0)e^{\varpi N_1(D_0)}e^{N_1(D_0)/4}$, the first bound of \eqref{eq:c4-8-4}. A root $D_0$ with source contact $a$ has a bond of $\mathsf L_{D_0}$ in the stencil $\mathcal F_{r_a}$ (\eqref{eq:part1-contact-set}), so it contains one of the at most $C_{\mathrm{st}}$ cubes meeting that stencil ($C_{\mathrm{st}}$ a constant of $L$ and $M$); summing the first bound over these roots by the display gives the second bound with $C_{\mathrm{pin}}= C_{\mathrm{st}}C_{\mathrm a}$. (b) The two geometric bounds follow from the square geometry and the wall connectivity of a component's domains (the geometry paragraph of Section~\ref{sec:c4-6}). If a traversal $a$ lies in $H_{\gamma}$ and $a'$ is a contact of $\gamma$, then $d_C(a,a')\leqslant A_{\mathrm h}LMn_{\gamma}$, so $e^{\varpi n_{\gamma}/4}\leqslant e^{\varpi n_{\gamma}/2}e^{-\theta_{\mathrm h}d_C(a,a')}$, and the contour sum $\sum_{a'}e^{-\theta_{\mathrm h}d_C(a,a')}\leqslant 1+2/(e^{\theta_{\mathrm h}}-1)$ is uniform in $n_C$. For a cube $\Delta\subset X_{\gamma}$ the analogous inequality holds with $\theta_{\mathrm x}$ and the torus distance; for each of the four straight sides the sum over $a'$ is at most $1+2/(e^{\theta_{\mathrm x}}-1)$ up to the fixed cube/halo offset, bounded by the full one-dimensional periodic lattice sum independently of $n$ and the periods. This retains $\rho$ even when $\Delta$ lies in a cluster domain. A cube $\Delta\in F(\gamma)$ either meets a domain of $X_{\gamma}$ or its stencil halo (a fixed offset, absorbed as above) or meets the stencil of a traversal of $H_{\gamma}$, and at most $CLM$ traversals have stencils meeting a given cube; so the footprint pin is at most \eqref{eq:c4-8-6} plus $CLM$ times the hull pin \eqref{eq:c4-8-5}, which is the stated bound with the enlarged constant. (c) $100\eta\mathfrak s(\gamma)\leqslant 100\eta A_{\mathrm s}n_{\gamma}<\varpi n_{\gamma}/8$, so the enhanced weights are dominated by the $e^{\varpi n_{\gamma}/4}$ of (b) with a margin $e^{\varpi n_{\gamma}/8}$. All constants are fixed before $g_0$; the only condition on $\rho$ in this lemma is $9\rho C_{\mathrm a}\leqslant 1/32$ from (a), and Lemma~\ref{lem:c4-8-3} adds $C\rho LM\leqslant \eta/8$ with the constant $C$ of the pins (b).
\end{proof}

\subsubsection{Short classes and the exceptional whole-contour merge}\label{sec:c4-8-3}
\begin{lemma}[class pins]\label{lem:c4-8-3}
Let $C$ be the constant of the pins (b)--(c) of Lemma~\ref{lem:c4-8-2} and assume $C\rho LM\leqslant \eta/8$. Put $\sigma_K=\sum_{\gamma\in K}\mathfrak s(\gamma)$ for a class $K$. Then: (a) for every overlap-connected class $K$, before the promotion of a long class to the whole contour, the hull and footprint pin sums $\sup_{a\in\mathcal I_C}\sum_{K\, | \,a\in\bigcup_{\gamma\in K}H_{\gamma}}t_Ke^{64\eta\sigma_K}$ and $\sup_{\Delta}\sum_{K\, | \,\Delta\in\bigcup_{\gamma\in K}F(\gamma)}t_Ke^{64\eta\sigma_K}$ are at most $C\rho$; (b) a long class satisfies
\begin{equation}\label{eq:c4-8-7}
    \sigma_K\geqslant \frac{1}{LM}\sum_{\gamma\in K}| H_{\gamma}|\geqslant \frac{n_C}{2LM},\qquad
    \sum_{K\ \mathrm{long}}t_Ke^{4\eta\sigma_K}\leqslant C\rho n_Ce^{-30\eta n_C/(LM)}\;;
\end{equation}
(c) the whole-contour groups created by the exceptional merge have total weighted mass $\sum_{W\ \mathrm{whole}}M(W)e^{4\eta\mathfrak s(W)}$ at most
\begin{equation}\label{eq:c4-8-8}
    C\rho n_Ce^{-30\eta n_C/(LM)}e^{4\eta n_C/(LM)}e^{C\rho n_C}\leqslant C\rho n_Ce^{-25\eta n_C/(LM)}\leqslant C'\rho\;,
\end{equation}
uniformly in $n_C$. A group $W$ with a short interval determines a class $K$ (Section~\ref{sec:joint-extraction}), and the groups with the same class $K$ are interleavings of the cluster lists of its components: their number is at most $k(W)!/\prod_{\gamma\in K}k_{\gamma}!$ with $k(W)=\sum_{\gamma\in K}k_{\gamma}$ (the interleaving count of Lemma~\ref{lem:grouped-source}), and each has $M(W)=\prod_{\gamma\in K}\bigl(\rho^{|\mathcal A(\gamma)|}\prod_{D\in\mathcal A(\gamma)}m(D)\prod_jM_v(\mathfrak c_j)\bigr)/k(W)!$ ($M(W)$ is the product of the root weights $\rho\,m(D)$, the cluster weights $M_v(\mathfrak c_j)$ and $1/k(W)!$, Section~\ref{sec:c4-6-2}), so the sum of $M(W)$ over the groups with class $K$ is at most $\prod_{\gamma\in K}\bigl(\rho^{|\mathcal A(\gamma)|}\prod_{D}m(D)\prod_jM_v(\mathfrak c_j)/k_{\gamma}!\bigr)= t_K$ (the per-group inequality $M(W)\leqslant t_K$ alone would not account for this many-to-one correspondence); moreover every such $W$ has $\mathfrak s(W)\leqslant \sigma_K$ ($n(W)=\sum_{\gamma\in K}n_{\gamma}$, and $\mathcal H(W)$, the union of the overlapping hulls $H_{\gamma}$, has at most $\sum_{\gamma\in K}| H_{\gamma}|$ traversals); its footprint $F(W)$ lies in $\bigcup_{\gamma\in K}F(\gamma)$. Hence the two sums of \eqref{eq:c4-6-5} over short groups are bounded by the pin sums of (a) ($e^{4\eta\sigma_K}\leqslant e^{64\eta\sigma_K}$), and the sums over whole-contour groups by (c): together, (a) and (c) prove the two weighted group pins \eqref{eq:c4-6-5}.
\end{lemma}

\begin{proof}
(a) Run the positive overlap-tree recursion on the decorated components with the weights enhanced by $e^{64\eta\mathfrak s(\gamma)}$. The pin of Lemma~\ref{lem:c4-8-2}(c) controls the overlap sum of a hull: the components whose hulls meet $H_{\gamma}$ have weighted mass at most $C\rho| H_{\gamma}|\leqslant \eta\mathfrak s(\gamma)/8$ by the hypothesis $C\rho LM\leqslant \eta/8$, so the rooted recursion over components (node weight $t_{\gamma} e^{64\eta\mathfrak s(\gamma)}$, node cost $\eta\mathfrak s(\gamma)/8$; the pins of Lemma~\ref{lem:c4-8-2}(c) carry the weight $e^{100\eta\mathfrak s(\gamma)}\geqslant e^{64\eta\mathfrak s(\gamma)}e^{\eta\mathfrak s(\gamma)/8}$, which covers the node cost) gives the hull and footprint pins for every class before promotion: the sum over the classes containing a component pinned by the hull pin at $a$, or by the footprint pin at $\Delta$, is the rooted sum with that component as root. No property of the union interval enters, so the class pins hold for long classes as well; for a short class the interval $\mathcal H(K)$ is the union of its hull intervals (Section~\ref{sec:joint-extraction}), so its size is bounded by the additive class cost $\sigma_K$. (b) The first inequality is the definition of $\mathfrak s$; the second is the definition of a long class (the union of its hulls has at least $n_C/2$ positions, Section~\ref{sec:joint-extraction}). Anchor the class at a traversal of its union interval; there are at most $n_C$ anchors, and the class pin (a) with the weight $e^{64\eta\sigma_K}$ gives the second bound of \eqref{eq:c4-8-7}, the remaining exponent $60\eta\sigma_K$ being at least $30\eta n_C/(LM)$. (c) Once a long class occurs, every other component is merged into the whole-contour group; the other components with their $4\eta$ weights sum to at most $e^{C\rho n_C}$ (each has a contact on the contour and is pinned there by Lemma~\ref{lem:c4-8-2}(c), so the sum over the families of further components is at most $\prod_{a\in\mathcal I_C}\exp(C\rho)$), and the whole-contour interval weight costs $e^{4\eta n_C/(LM)}$. The middle inequality of \eqref{eq:c4-8-8} uses $C\rho LM\leqslant \eta$, implied by the hypothesis, and the final maximum is finite uniformly in $n_C$ because $LM$ and $\eta$ are fixed. If a spatial pin lies in an additional component, root that component by its footprint pin at a further cost $C\rho$, the anchored long-class factor and remaining exponential being paid by \eqref{eq:c4-8-8}; if it lies in the long class, use the footprint pin in place of the hull pin; pins on the whole contour are covered by the total whole-group bound. This proves the exceptional-merge pin without assigning its cost to a single original component. Polynomial moments are allowed by using a smaller part of the same fixed enhanced-weight margin.
\end{proof}

\subsubsection{Interior boxes, endpoints and the differentiated group weight}\label{sec:c4-8-4}
Let $C_{\mathrm{box}}$ be a fixed geometric constant for the group footprint and the row/contour halos such that a liftable enclosing comparison rectangle has diameter at most $C_{\mathrm{box}}LMn(W)$, increased so that a group with
\begin{equation}\label{eq:c4-8-9}
    n(W)<\frac{b}{4C_{\mathrm{box}}LM}
\end{equation}
has its entire rectangle and stencils within $b/2$ of a source contact and an injective lift (Section~\ref{sec:c4-11-4}). On a far history this box lies in $\Lambda_1$, hence in $\Omega_1$, and outside $S_1=(\Omega_1\setminus\Omega_1^{\sim-1})\cup R_1$; its plaquettes satisfy $\delta_{Q_W}(U_1)\leqslant C\varepsilon_1$ by Lemma~\ref{lem:c2-plaquette-deviation} and the distance check of Section~\ref{sec:c4-11-5}, and, for a group passing the phase test of Lemma~\ref{lem:c4-9-1} (the groups failing it are treated in Section~\ref{sec:c4-12-2}), the Hessian bounds \eqref{eq:c4-7-7} of Lemma~\ref{lem:c4-7-3} hold on the globally small extension used for its comparison. All other groups keep their undifferentiated difference, $| c_W|+| a_W\chi|\leqslant 2M(W)$ (Section~\ref{sec:c4-6-2}); their cube count $n(W)$ is at least $b/(4C_{\mathrm{box}}LM)$, and if a rectangle fails to lift the same lower bound holds, since $N_{\min}=\min_iN_i\geqslant 2LMR_1>32b$ (Section~\ref{sec:c4-9-2}) and a wrapping connected network has length at least a fixed fraction of $N_{\min}$. No interpolation is made across the integration boundary.

For a good box use the coefficient locality and covariance of Sections~\ref{sec:c4-8-1} and~\ref{sec:c4-11}: root the tree gauge at the initial interval base, keep $h$ fixed, and compare $c_W(U_1;h)$ with $c_W(\id;h)= a_W\chi(h)$ by the small-phase extension and the fundamental theorem of calculus (Lemma~\ref{lem:c4-9-1}, for a group passing the phase test). It is the same regional reference coefficient, with $I$, all selection indices and the local geometry retained; no bulk comparison replaces it by a different coefficient. The group weight is differentiated by the product rule, each occurrence retained:
\begin{equation}\label{eq:c4-8-10}
    \mathscr D_{\omega} w_{\xi}(W)=\frac{(-1)^{k(W)}}{k(W)!}\sum_{j=1}^{k(W)}\mathscr D_{\omega}\Phi^{\mathrm v}_{\xi}(\mathfrak c_j)\prod_{l\ne j}\Phi^{\mathrm v}_{\xi}(\mathfrak c_l)\;,
\end{equation}
bounded by Lemma~\ref{lem:c4-7-4} by $C_M\|\omega\|_{\infty}\sum_jn(\mathfrak c_j)$ times the positive product majorant, without dividing by any $\Phi$. With the joint source derivative of Lemma~\ref{lem:c4-4-3}, whose factor $L_f+\sum_DN_1(D)$ is at most $C(|\mathcal H(W)|+n(W))$ because the number $L_f$ of varying background occurrences in the relative word is at most $2|\mathcal H(W)|$ times the fixed reconstruction-row overlap constant (Section~\ref{sec:c4-9-1}, the sentence defining $L_f$) and $\sum_DN_1(D)\leqslant n(W)$, this gives $|\mathscr D_{\omega} c_W|\leqslant C\|\omega\|_{\infty}(n(W)+|\mathcal H(W)|)M(W)$ uniformly in $h$ (\eqref{eq:c4-12-2}); the curvature phase bound contributes $C\operatorname{diam}(Q_W)^2\delta_{Q_W}$, and with $|\mathcal H(W)|\leqslant CLMn(W)$ all these losses are a fixed constant times $\varepsilon_1$ times a fixed polynomial in $\mathfrak s(W)$, absorbed by the enhanced pin ($x^ke^{-\eta x}\leqslant k!\eta^{-k}$). For a large group the factor $e^{-\eta b/(8C_{\mathrm{box}}LM)}$ is extracted from its reserved weight $e^{\eta\mathfrak s(W)/2}$. Combining the two classes and summing over a contour pin gives \eqref{eq:c4-6-8}; the computation is Section~\ref{sec:c4-12-2}. The same matrices and the same cutoff norm are used at the physical and flat endpoints, and every source estimated is the closed functional $c_W$: no partial interval is replaced by a scalar before its closed subtraction is formed.

\subsection{The coefficient extension and the countable-gas limit}\label{sec:c4-9}
\subsubsection{Background scope and the coefficient path}\label{sec:c4-9-1}
Sections~\ref{sec:c4-4} and~\ref{sec:c4-7} give the joint source and one-derivative bounds at the physical centred background and on every globally small real extension described here, for each closed character $\chi(T_WP_W^{-1}h)$ uniformly in $h$; they do not assert a path from the entire physical background to the identity in one global small chart.

Fix a good interior group of Section~\ref{sec:c4-8-4} and its enclosing fine rectangle $Q_W$, which contains all coefficient, reconstruction, averaging and interval/closing stencils. By the support statements of Section~\ref{sec:c4-10-2}, for the fixed selection indices and the fixed spatial Cauchy coefficient defining the group only the background links on $Q_W$ are used: the restricted seed inverse uses entries on $E_0$, every walk coefficient has its domains in the final component, and the potential and source rows have their finite stencils. The regional $I$, $\Omega_1$ and every auxiliary parameter are held fixed. Gauge $U_1$ on $Q_W$ by its rooted tree gauge $g_W$ with $g_W(\mathrm{base})=\id$ (Lemma~\ref{lem:c4-10-5}). If $C_{\mathrm{tg}}\operatorname{diam}(Q_W)^2\delta_{Q_W}\leqslant \tau_{\mathrm b}/2$ (the phase test; since $\operatorname{diam}(Q_W)$ is at least the largest side $R$ of $Q_W$, it implies the hypothesis $C_{\mathrm{tg}}R^2\delta_{Q_W}\leqslant \tau_{\mathrm b}/2$ of Lemma~\ref{lem:c4-10-5}), write the principal phases $B_W(b)=\frac1i\log U_1^{g_W}(b)$ on the rectangle and extend $B_W=0$ on every other bond, and define the coefficient argument and its direction by
\begin{equation}\label{eq:c4-9-1}
    U_{W,t}(b)= e^{itB_W(b)},\qquad\omega_W(b)= B_W(b),\qquad\|\omega_W\|_{\infty}=\max_bh(B_W(b)),\qquad0\leqslant t\leqslant 1\;;
\end{equation}
the norm is the invariant norm $h$, the comparison with the cutoff norm being absorbed in $C$. Set $L_f$ equal to the number of occurrences of background links in the relative word $T_WP_W^{-1}$ on which this argument varies, $L_f\leqslant 2|\mathcal H(W)|$ times the fixed reconstruction-row overlap constant. The complementary $h$ is held fixed at the root base.

At each $t$ form $H_t=\iota_I^*\mathsf H^{\Omega_1}_{U_{W,t}}\iota_I$, the row reconstruction $q_{U_{W,t}}$ and the near primitives at $U_{W,t}$, with zero retained input and the same protected $r$-box. For a generic noncritical $U_{W,t}$ the energy primitive is defined by subtracting its constant, linear and quadratic Taylor terms as in \eqref{eq:native-source-energy-residual}, the Hessian keeping the force--constraint term; the Haar and constraint-Jacobian primitives use their analytic row formulas. These definitions specify an analytic coefficient extension. This auxiliary kernel is not identified with the physical CRE integral at $t\in(0,1)$, and the original fixed-$V$ constraint is not imposed along the path; at the physical endpoint each compensated primitive agrees with its Taylor-subtracted formula by locality, criticality being used in the primitive-sum identity \eqref{eq:native-source-primitive-sum} for the physical exponent and not to assert that $U_{W,1}$ is a global minimum; at the flat endpoint the formulas agree with the reference formulas. Lemma~\ref{lem:c4-7-3} gives $\gamma_L/2\leqslant H_t\leqslant m_L$; the designated inverses, reconstruction and means use the fixed smaller row charts since $\| B_W\|_{\infty}\leqslant \tau_{\mathrm b}/2$ and $\tau_{\mathrm b}\leqslant \tau_L/4$; the host norms and their derivatives are bounded by $C^{\partial}_0g_0(1+r)^3$ on these charts; so the Gaussian payment, the positive forest budget and the source estimates hold uniformly in $t$ with the same recipe. The force term is one of the local entries and is not omitted off the minimum.

\begin{lemma}[the coefficient path]\label{lem:c4-9-1}
For a small interior group (\eqref{eq:c4-8-9}, a good box of Section~\ref{sec:c4-8-4}) passing the phase test of Section~\ref{sec:c4-9-1}, and for each $h\in G$ in the closed character $\chi(T_WP_W^{-1}h)$,
\begin{equation}\label{eq:c4-9-2}
    d_W(U_1;h)=\int_0^1\mathscr D_{\omega_W}c_W(U_{W,t};h)\,\dd t\;,
\end{equation}
and $|\mathscr D_{\omega_W}c_W(U_{W,t};h)|\leqslant C\|\omega_W\|_{\infty}(L_f+n(W))M(W)$ uniformly in $t$ and $h$.
\end{lemma}

\begin{proof}
At $t=1$ the group coefficient agrees with the original one by coefficient locality on $Q_W$ and simultaneous gauge covariance (Section~\ref{sec:c4-11-3}): the values outside $Q_W$ can be replaced in the argument because no term of the spatial Cauchy coefficient consults them, and the root gauge $g_W(\mathrm{base})=\id$ leaves $h$ unchanged. At $t=0$ the coefficient is exactly the regional reference one, with the same $I$ and the same group index $W$, so $c_W(U_{W,0};h)= a_W\chi(h)$. The fundamental theorem of calculus applies: by the dominated-convergence hypotheses stated in the last paragraph of Section~\ref{sec:c4-11-3} and the termwise differentiation statement of Lemma~\ref{lem:c4-9-2}, $t\mapsto c_W(U_{W,t};h)$ is continuously differentiable on $[0,1]$ with derivative $\mathscr D_{\omega_W}c_W(U_{W,t};h)$. The derivative of $c_W$ differentiates $\prod_j\Phi^{\mathrm v}(\mathfrak c_j)$, giving the occurrence sum \eqref{eq:c4-8-10} with majorant $C\|\omega\|_{\infty}(\sum_jn(\mathfrak c_j))M(W)$, and the root integral, giving $C\|\omega\|_{\infty}(L_f+\sum_{D}N_1(D))M(W)$ by Lemmas~\ref{lem:c4-4-3} and~\ref{lem:c4-7-1}; the sum is the displayed bound. This is a cluster-list bound for the same group weight, not an identification of a one-root vacuum ratio with it: the joint-ratio estimate of Lemma~\ref{lem:c4-7-4} controls the original family normalisation, the occurrence estimate controls its regrouping. If the phase test fails, the two endpoint estimates give the curvature bound of Section~\ref{sec:c4-12-2}; no continuation is claimed in that regime. The path is the connection between the local physical coefficient and the globally small extension; it does not assume global near-identity physical fields.
\end{proof}

\subsubsection{Parameter and period conventions}\label{sec:c4-9-2}
All results inherit the standing conditions \eqref{eq:c-standing-conditions} and the compatible even fine periods of Section~\ref{sec:c3-7}: every fine period $N_i$ is an even integer multiple of $\ell= LMR_1$, and $\ell$ is odd ($L$ odd, $M$ and $R_1$ powers of $L$). If $N_i= q_i\ell$ then $q_i$ is even and positive, so
\begin{equation}\label{eq:c4-9-3}
    N_{\min}=\min_iN_i\geqslant 2LMR_1>32\lfloor LMR_1/16\rfloor=32b\;,
\end{equation}
the strict middle inequality because $\ell$ is odd. This is independent of the square-fit condition $K_i\geqslant 3(n+2)$; both are retained, with the source-row and coarse-grid conditions of Section~\ref{sec:sourced-effective-interaction}. Thus the good comparison box and every fixed stencil have a unique lift under the small-group criterion \eqref{eq:c4-8-9}, and the periodic failure lies in the large-group class; no condition on the fluctuating field or a denominator is added. $\Gamma_C$ consists of the individual positive contour bonds and all their fixed input templates, $C_L$ including the number and size of the templates, and the radius-$b$ ball of Lemma~\ref{lem:c4-2-1} is around each individual bond, so the count $4(2b+3)^4$ needs no further $n_C$-dependent factor. The input of Appendix~\ref{app:exterior-chessboard} is the one-point probability $p_z\leqslant q\leqslant e^{-c_*\pi_*^2}$ of Lemma~\ref{lem:c3-7-4} and Proposition~\ref{prop:c3-7-5} (this $q$ is the $q_{\mathrm{pt}}$ of Lemma~\ref{lem:c4-2-1}) with $p_0(g)= A_0(\log g^{-2})^{p_0}$, $p_0\geqslant 5\mathfrak r\geqslant 10$. $E_{\mathrm{bdry}}$ of \eqref{eq:c3-6-2} is bounded separately by Proposition~\ref{prop:c3-7-5}; $E_{\mathrm{def}}$ contains only the covered centred-kernel subtraction, so neither the exterior response nor the uncovered contribution is dropped from \eqref{eq:c4-8-1}.

\subsubsection{Countably many group species and differentiation}\label{sec:c4-9-3}
\begin{lemma}[countable species]\label{lem:c4-9-2}
At a fixed finite torus let $\mathcal W_N=\{ W\, | \,n(W)\leqslant N\}$, a finite set, with partition sum $S_N$ and connected-cluster logarithm $\mathscr L_N$, $S_N= e^{\mathscr L_N}$ (the cluster logarithm is written $\mathscr L$, the block scale being $L$). Under the hypotheses of Lemma~\ref{lem:c4-6-3} ($C\rho LM\leqslant \eta/8$ and $g_0$ small),
\begin{equation}\label{eq:c4-9-4}
    \sum_W| a_W| e^{\mathfrak a_W/8}\leqslant C\rho n_C<\infty,\qquad \mathscr L_N\to\mathscr L,\qquad S_N\to S^{\circ}_{\xi,C}= e^{\mathscr L}\ne0\;,
\end{equation}
the restricted identity $S^{\circ,[\mathfrak B]}_{\xi,C}= e^{\mathscr L^{[\mathfrak B]}}$ ($\mathscr L^{[\mathfrak B]}$ the connected-cluster logarithm of the species compatible with $\mathfrak B$) and the ratio bound \eqref{eq:c4-6-7} hold for the countable index set, and the background derivatives of Lemmas~\ref{lem:c4-4-3}, \ref{lem:c4-7-1} and~\ref{lem:c4-7-4} may be taken termwise and integrated in $t$ before the countable limit.
\end{lemma}

\begin{proof}
$\mathcal W_N$ is finite: the domain set is finite, and $n(W)$ bounds the number of roots, occurrences and members of every occurrence. Its gas satisfies the same Kotecký--Preiss inequality, deleting species only decreasing the positive sums. Every group has a nonempty contour interval, $| a_W|\leqslant M(W)$ (Section~\ref{sec:c4-6-2}) and $e^{\mathfrak a_W/8}= e^{\eta\mathfrak s(W)/8}\leqslant e^{4\eta\mathfrak s(W)}$, so the pin \eqref{eq:c4-6-5} summed over the $n_C$ traversals gives the first bound of \eqref{eq:c4-9-4}; the rooted-tree estimate bounds the total absolute connected-cluster series by this finite sum up to a fixed constant; each finite cluster tuple is present for all large $N$, so absolute domination gives $\mathscr L_N\to\mathscr L$; independently the full gas sum over finite compatible families is absolutely bounded by $\exp(\sum_W| a_W|)$, and dominated convergence gives $S_N\to S^{\circ}_{\xi,C}$. The same argument after excluding every species incompatible with a fixed compatible family $\mathfrak B$ gives the restricted identity; the logarithmic difference is the connected-cluster sum of tuples hitting $\mathfrak B$, counted once, with absolute sum at most $\sum_{W\in\mathfrak B}\mathfrak a_W/8$ by the pinned estimate. For the derivatives, Lemmas~\ref{lem:c4-7-1} and~\ref{lem:c4-7-4} give summable positive majorants independent of $t$ on the extension; on each finite truncation the product and integral derivatives are justified by the uniform Gaussian coercivity (Section~\ref{sec:c4-10-3}) and the fixed protected box, and the walk derivative series converges on its unchanged contours (Section~\ref{sec:c4-10-2}); the cluster majorant $n(\mathfrak c)M_v(\mathfrak c)$ is paid by its exponential moment and the group majorant has only the polynomial in $\mathfrak s(W)$, paid by the reserved weight. Uniform absolute convergence permits termwise differentiation and $t$-integration before the countable limit; no conditionally convergent series is differentiated.
\end{proof}

\subsubsection{Integrability at the physical caller}\label{sec:c4-9-4}
The near-history contribution is bounded after forming $\mathscr R_{\xi}$, without exchanging its sum with the outer integral (Lemma~\ref{lem:c4-2-1}). On covered far histories the proof of Lemma~\ref{lem:c4-6-5} bounds $\sum_{\mathfrak B\ne\varnothing}| S^{\circ,[\mathfrak B]}_{\xi,C}R_{\xi}(\mathfrak B)|$ by $| S^{\circ}_{\xi,C}|$ times the exponential of the one-group sum \eqref{eq:c4-6-9} minus one; since $| S^{\circ}_{\xi,C}|\leqslant 1$ ($S^{\circ}_{\xi,C}$ is the normalised integral \eqref{eq:part1-reference-gas} of $f^{\circ}_{\xi}$, a normalised character of modulus at most one, against the positive reference kernel $K^{\circ}_{\xi;0}$) and the exponential is uniform in the branch fields and retained variables, it is integrable against the normalised positive history law at each fixed loop and cutoff, which justifies the far-history sum/integral exchanges. The physical $\mathcal Z_0$ and $\mathsf w_{\xi}= e^{c_{\xi}}I_{\xi;1}$ remain those of Section~\ref{sec:c4-1}.

\subsection{Analytic inputs for the coefficient extension}\label{sec:c4-10}
\subsubsection{Joint background/phase analyticity and primitive derivatives}\label{sec:c4-10-1}
Fix one of the finitely many row/primitive stencil types in its translated fine-lattice frame and write the background on the stencil as $U_Z(b)= e^{iZ_b}$; background directions use the maximum $h$-norm on this finite label set. Choose $\tau_{\mathrm b}$ and a phase radius $\theta_0>0$ so that $\| Z\|<2\tau_{\mathrm b}$ and $\| A\|<2\theta_0$ keep every finite path product and relative product entering the bond average and the mean of RG1 (0.4), (0.11) (\eqref{eq:nonlinear-average}) inside the matrix-logarithm neighbourhood $\| W-\id\|_{\mathrm{op}}<1/2$; this is possible uniformly over the finite types, every word having bounded length and $e^{iA}e^{iZ}$ being analytic in both variables. The radii $\tau_{\mathrm b}$, $\theta_0$ used below are those of the final joint domain, after the reductions of this section and of Section~\ref{sec:c4-11-5}; in particular $2\tau_{\mathrm b}\leqslant \tau_L/2$ (Lemma~\ref{lem:c4-7-3}) and $2\theta_0\leqslant \vartheta_*/2$ (Section~\ref{sec:c4-11-5}), so that Lemma~\ref{lem:c2-uniform-chart}(ii)--(iv) applies at each real $Z$ with $\| Z\|<2\tau_{\mathrm b}$ and each $A$ with $\| A\|<2\theta_0$.

For the mean of RG1 (0.11) write its defining equation $\sum_j\log(W_jM^{-1})=0$. At $W_j=\id$, $M=\id$ its derivative in the Lie-algebra parameter of $M$ is $-n\,\id$, so the parametric analytic implicit function theorem gives a jointly analytic local mean in the background and fluctuating phases; it agrees with the mean of Lemma~\ref{lem:c2-uniform-chart}(iv), both taking the value $\id$ at zero and solving that equation in the same smaller neighbourhood. The joint radii are shrunk so that the mean stays in that neighbourhood; no global uniqueness is inferred from a logarithm chart.

The bond constraint on a row $c$ is the jointly analytic function
\begin{equation}\label{eq:c4-10-1}
    \mathcal Q_Z(A)=\frac1i\log\bigl(\mathsf M(e^{iA}U_Z,c)\mathsf M(U_Z,c)^{-1}\bigr),\qquad LQ_Z= D_A\mathcal Q_Z(0),\qquad N_Z(A)=\mathcal Q_Z(A)-LQ_ZA\;,
\end{equation}
$\mathsf M$ the fixed bond average of RG1 (0.4), not the finite-group mean, with $N_Z(0)= D_AN_Z(0)=0$ for every $Z$. The coefficient of the designated phase in $LQ_Z$ is uniformly invertible on a smaller background polydisc, by the identity-background designated inverse (Lemma~\ref{lem:two-block}) and continuity, so the designated inverse $h_Z$ with $LQ_Zh_Z=\id$ and $\mathsf C_Z=\iota_{\mathrm{free}}-h_ZLQ_Z\iota_{\mathrm{free}}$ are analytic in $Z$ with uniformly bounded first background derivatives; these formulas fix their maps and normalisations, which are those of Section~\ref{sec:regular-hessian-stencils}, without a new free-variable convention. The row solution is defined by
\begin{equation}\label{eq:c4-10-2}
    \widetilde D_Z(A)= N_Z\bigl(A-h_Z\widetilde D_Z(A)\bigr)\;.
\end{equation}
On smaller joint balls the derivative of the right side in $\widetilde D$ has norm at most $1/2$ uniformly in $Z$, since $D_AN_Z(A)= O(\| A\|)$; uniform analytic contraction (the parametric implicit function theorem) gives joint analyticity and $\|\widetilde D_Z(A)\|\leqslant C\| A\|^2$ (Lemma~\ref{lem:c2-uniform-chart}(ii) at each real $Z$). Cauchy estimates on the larger joint balls, followed by the order-two zero in $A$, give for $l=0,1$ and each fixed $k$ needed below
\begin{equation}\label{eq:c4-10-3}
    \| D_Z^lD_A^k\widetilde D_Z(A)\|\leqslant C_{l,k}\| A\|^{\max(2-k,0)}\;.
\end{equation}
All constants are maxima over finitely many types, independent of the volume, the history and $g_0$; the dimensions of the types depend on $L$ and $G$ only. The protected free input satisfies $g_0\|\mathsf C_Za\|<\theta_0/2$ after reducing the coupling, and the phase $q_Z(a)= g_0\mathsf C_Za-h_Z\widetilde D_Z(g_0\mathsf C_Za)$ satisfies
\begin{equation}\label{eq:c4-10-4}
    \| D_Zq_Z(a)\|\leqslant Cg_0(1+r),\qquad\| D_Z^lD_a^kq_Z(a)\|\leqslant C_{l,k}g_0^k
\end{equation}
for $l=0,1$ and $1\leqslant k\leqslant 2K+1$, $K$ the fixed row input count of Section~\ref{sec:c4-10-4}. These are joint background/young-field derivative bounds; they are not the assertion that the finite list of phase-only derivative constants of Lemma~\ref{lem:c2-uniform-chart} contains them.

Analytic traces of $X$ and $X^{-1}$ extend $\operatorname{Re}\operatorname{tr}X$ from unitary fields, where it equals $\operatorname{tr}(X+X^{-1})/2$; the energy and relative gauge expressions are therefore jointly analytic finite word functions. For the local energy function $E_Z(s)$ in the dimensionless free phase $s$ (the action and gauge-fixing functional of \eqref{eq:action-primitive} at the background $U_Z$), the compensated primitive has the Taylor-remainder form
\begin{equation}\label{eq:c4-10-5}
    p^{\mathrm E}_Z(a)=-\frac{g_0}{2}\int_0^1(1-t)^2D_s^3E_Z(tg_0a)[a,a,a]\,\dd t\;.
\end{equation}
Its first background derivative has the same factor $g_0$, since $D_ZD_s^3E_Z$ is uniformly bounded on the smaller joint domain; hence $| p^{\mathrm E}|+\| D_Zp^{\mathrm E}\|\leqslant Cg_0(1+r)^3$, the relative means and the soft-gauge term being included in $E_Z$. The Haar function $\log(\sigma(q)/\sigma(0))$ is analytic in $q$ with an order-two zero, so composed with $q_Z$ it and its $Z$-derivative are at most $Cg_0^2(1+r)^2$. The constraint-Jacobian logarithm is $\log\det(\id-h_ZD_A\widetilde D_Z)$ on a matrix ball of radius $<1/2$; its value and $Z$-derivative are at most $Cg_0(1+r)$, as $D_A\widetilde D_Z= O(g_0(1+r))$ and $D_ZD_A\widetilde D_Z$ has the same order, and the real determinant is positive on that ball, so the analytic logarithm restricts to the real logarithm of Lemma~\ref{lem:c2-uniform-chart}(iii). Summing the finite primitive types assigned to one host gives the host derivative bound with $C^{\partial}_0= N_{\mathrm h}\max(C_{\mathrm p},C^{\partial}_{\mathrm p})e^{\kappa(B_{\mathrm h}+4)}$ of Section~\ref{sec:c4-7-1}; this constant and the finite source derivative constants are chosen after $M$ and $\tau_{\mathrm b}$, before $g_0$, and no earlier Mayer, walk or reserve choice depends on them.

\subsubsection{Differentiated fixed walks and exact locality}\label{sec:c4-10-2}
Use the series \eqref{eq:c-weighted-walk-series} with its parameters $s_H=2h_+$, $\mathsf A_H=\id-H/s_H$, $\mathsf P_H=\mathsf A_H^{n_*}$, $q_H=1-\gamma/s_H$, $\beta= V_*q_H^{n_*}$ ($V_*= V_{n_*}$), $a_*= e^{81\kappa_1}$ and $n_*$, fixed from the common input bounds and not re-evaluated at the varied background. Of the background these constructions use only the input bounds \eqref{eq:fixed-kernel-inputs}, $\gamma\leqslant H\leqslant h_+$, and the $R$-locality of the entries; both hold along $U_{W,t}$, the two-sided bound by Lemma~\ref{lem:c4-7-3} and the locality by the entry statement of Lemma~\ref{lem:hessian-gauge-locality}, whose radius $R$ does not depend on the background. Include every geometrically allowed tuple of bond indices, including coefficients vanishing at a particular background, so that differentiation does not change the index set. The weights $w_{\omega}(\sigma)$ and domains $D(\omega)$ depend only on the labels and the fixed geometry. Let $c_H$ bound the Schur row and column norms of $\mathscr D_{\omega} H/\|\omega\|_{\infty}$, supplied by the finite analytic entries of Section~\ref{sec:c4-10-1}. Then $\mathscr D_{\omega}\mathsf A_H=-\mathscr D_{\omega} H/s_H$ and $\mathscr D_{\omega}\mathsf A_H^j=\sum_{l=0}^{j-1}\mathsf A_H^l(\mathscr D_{\omega}\mathsf A_H)\mathsf A_H^{j-1-l}$, a block sum of which is bounded by $V_jjq_H^{j-1}c_H\|\omega\|_{\infty}/s_H$. Put $d_*= V_*n_*q_H^{n_*-1}c_H/s_H$ and let $d_{\mathrm{ini}}$ be the maximum of the corresponding bounds for $j< n_*$. Differentiating the grouped walk products gives the absolute Schur bound
\begin{equation}\label{eq:c4-10-6}
    \|\mathscr D_{\omega} S_{Z_0}(\sigma)\|_{\mathrm{Schur}}\leqslant C^{\partial}_{\mathrm w}\|\omega\|_{\infty},\qquad
    C^{\partial}_{\mathrm w}= s_H^{-1/2}n_*\Bigl(\frac{a_*d_{\mathrm{ini}}}{1-a_*\beta}+\frac{a_*^2V_*d_*}{(1-a_*\beta)^2}\Bigr)\;,
\end{equation}
the denominators positive because $a_*\beta\leqslant 1/4$ (\eqref{eq:fixed-kernel-grouping}); the same bound holds on the spatial Cauchy circles. The derivative of $C_{E_0}$ is $-C_{E_0}(\mathscr D_{\omega} H_{E_0})C_{E_0}$, with $H_{E_0}=\iota_{E_0}^*H\iota_{E_0}$, and that of $\Gamma_{E_0}$ is the product rule on $\iota_{E_0}^*H\Pi_{I\setminus E_0}S_{Z_0}$ (\eqref{eq:fixed-kernel-C-Gamma}); their operator norms are fixed constants times $\|\omega\|_{\infty}$.

For a selected final domain $Z$ the spatial variables outside $Z$ are zero and the protected seed enlargement $Z_0'$ lies in $Z$; a walk with a coefficient domain meeting a cube outside $Z$ has a zero spatial factor, which remains zero after background differentiation. Every surviving walk, including every hidden fine step in its grouped powers, has its coefficient stencil inside its declared domains by the construction of $D_b$ (Section~\ref{sec:fixed-first-kernel}). The restricted inverse $C_{E_0}$ consults only entries on the seed $E_0$, included in the comparison rectangle; it is not a global inverse restricted afterwards, and the normalised Gaussian determinant is $\det C_{E_0}$, not a determinant of $H$ on all of $I$. All potential and source stencils and every constituent vacuum domain of $c_W$ lie in $Q_W$. These facts prove the locality used at $t=1$ in Lemma~\ref{lem:c4-9-1}, with the fixed regional $I$: no surviving selected coefficient consults the background outside $Q_W$, and no field-dependent global spectral quantity appears among the recipe parameters. The common full Gaussian prefactor $\mathfrak n_{\xi}$ (Section~\ref{sec:fixed-first-kernel}) has cancelled inside the exact conditional ratio defining $F_{\xi}(0)$ (\eqref{eq:fixed-probe-centred-source}); it is not a missing local activity factor, and that conditional cancellation does not cancel $\mathsf w_{\xi}= e^{c_{\xi}}I_{\xi;1}$ or the physical $\mathcal Z_0$.

\subsubsection{The complex-Gaussian moment payment}\label{sec:c4-10-3}
On a selected component write $E$ for its seed carrier $E_D$, $J_Z$ for its auxiliary carrier $J_D$ (the block of $\bigsqcup_DJ_D$ belonging to this component, Lemma~\ref{lem:component-blocks}; in this subsection $J_Z$ denotes the carrier of one component), $H_E=\iota_E^*H\iota_E$, $C_E= H_E^{-1}$ and $\Gamma=\Gamma_1+i\Gamma_2$ for the block of $\Gamma_{E_0}(\sigma)$, the pairings complex bilinear. Before cutoffs and Mayer factors the joint density with respect to real Lebesgue measure $\dd B\,\dd X$ is the two centred Gaussian normalisations times
\begin{equation}\label{eq:c4-10-11}
    \exp\Bigl[-\frac12\langle B,H_EB\rangle-\langle B,\Gamma X\rangle-\frac12\langle\Gamma X,C_E\Gamma X\rangle-\frac12\| X\|^2\Bigr]\;,
\end{equation}
call it $G_{\sigma}$. The columns of $\Gamma$ outside $J_Z$, and those of its derivative, vanish identically by Lemma~\ref{lem:component-blocks}, so the auxiliary standard Gaussians integrate to one and no moment cost proportional to $| I|$ arises. With $\alpha=\alpha_G$ and $C_{\alpha}=(H_E-\alpha\id)^{-1}$, the absolute density times $e^{\alpha\| B\|^2/2}$ is a positive real Gaussian with the original determinant prefactor; in the variables $y= B+C_{\alpha}\Gamma_1X$ its exponent is
\begin{equation}\label{eq:c4-10-12}
    -\frac12\langle y,(H_E-\alpha)y\rangle-\frac12\langle X,(\id-T_{\alpha})X\rangle\;,\qquad T_{\alpha}=\Gamma_1^{\mathsf T}(C_{\alpha}-C_E)\Gamma_1+\Gamma_2^{\mathsf T}C_E\Gamma_2\;,
\end{equation}
the reserve matrix of Appendix~\ref{app:fixed-kernel-payment}. With $H_E-\alpha\geqslant \gamma/2$ (from $H\geqslant \gamma$, \eqref{eq:fixed-kernel-inputs}, and $\alpha_G\leqslant \gamma/2$, Lemma~\ref{lem:gaussian-payment}(ii)), $\| T_{\alpha}\|\leqslant 1/4$ (\eqref{eq:fixed-kernel-reserve-payment}) and $\|\Gamma_1\|\leqslant A_*$ (the norm paragraph of Appendix~\ref{app:fixed-kernel-payment}) and
\begin{equation}\label{eq:c4-10-13}
    c_{\mathrm G}=\frac{\min(\gamma/2,3/4)}{2(1+2A_*/\gamma)^2}>0\;,
\end{equation}
the real quadratic form is at least $c_{\mathrm G}(\| B\|^2+\| X\|^2)$. Its total unnormalised mass is the determinant product bounded by $e^{N_1(Z)}$ in Lemma~\ref{lem:gaussian-payment}(ii), and its normalised covariance is at most $c_{\mathrm G}^{-1}\id$. With $d= d_G(| E|+| J_Z|)$,
\begin{equation}\label{eq:c4-10-14}
    \int(1+\| B\|^2+\| X\|^2)e^{\alpha_G\| B\|^2/2}| G_{\sigma}(B,X)|\,\dd B\,\dd X\leqslant (1+d/c_{\mathrm G})e^{N_1(Z)}\;,
\end{equation}
and $d\leqslant 8d_G(LM+1)^4N_1(Z)$, so the loss is linear in the cube count $N_1(Z)$ with a fixed $M$-dependent constant; it uses the positive $3/4$ margin in $\id-T_{\alpha}$, not a new exponential reserve. The logarithmic derivative of the Gaussian normalisation contributes $\operatorname{Tr}(C_E\mathscr D_{\omega} H_E)/2$; the exponent derivative contributes $-\langle B,(\mathscr D_{\omega} H_E)B\rangle/2$, $-\langle B,(\mathscr D_{\omega}\Gamma)X\rangle$ and $-\mathscr D_{\omega}\langle\Gamma X,C_E\Gamma X\rangle/2$ with $\mathscr D_{\omega} C_E=-C_E(\mathscr D_{\omega} H_E)C_E$; each is bounded by $C\|\omega\|_{\infty}(d+\| B\|^2+\| X\|^2)$ by Section~\ref{sec:c4-10-2}. Hence \eqref{eq:c4-10-14} proves the differentiated Gaussian payment $C_MN_1(Z)\|\omega\|_{\infty} e^{N_1(Z)}$ on the same Cauchy circles and carriers, all normalisers and the negative compensator $-\frac12\langle\Gamma X,C_E\Gamma X\rangle$ differentiated. The protected cutoff domain is fixed throughout; no cutoff translation or surface term is suppressed in this derivative.

\subsubsection{Source remainder bounds}\label{sec:c4-10-4}
Let $\mathsf J=(b_1,\ldots,b_q)$ be distinct selected labels and let $a'$ denote all other selected root variables. The mixed fundamental theorem of calculus gives the exact identity
\begin{equation}\label{eq:c4-10-7}
    \Delta_{\mathsf J}f(B,a')=\int_{[0,1]^q}D_{b_1}\cdots D_{b_q}f\Bigl(a'+\sum_{j=1}^qs_jB_{b_j}\Bigr)[B_{b_1},\ldots,B_{b_q}]\,\dd\mathbf s\;,
\end{equation}
the sum in the argument meaning the canonical embedding into the free carrier. Let $K$ be a fixed bound for the number of source rows affected by one free label and for the number of inputs of one row. Only the source $f_W(a;h)=\chi(T_W(a)P_W^{-1}h)$, or the original full closed Wilson character, is used (a normalised character, so $| f|\leqslant 1$); untouched factors are unitary in its representation. The derivative product rule has at most $K^q$ row assignments, and the derivatives within a row have order at most $2K+1$, whose constants \eqref{eq:c4-10-4} are fixed. Consequently, with a fixed $C_{\chi}$ enlarged once,
\begin{equation}\label{eq:c4-10-8}
    \| D_{b_1}\cdots D_{b_q}f\|\leqslant (C_{\chi} g_0)^q,\qquad\sup|\Delta_{\mathsf J}f|\leqslant (C_{\chi} g_0r)^q\;.
\end{equation}
For root $j$ let $P_j$ be its homogeneous degree-$q_j$ projection after taking $\Delta_{\mathsf J_j}$. Subtracting the integrand of \eqref{eq:c4-10-7} at zero root input and integrating its first derivative along the segment $t(\mathbf sB)$ gives $q_j$ terms, each with one additional derivative and one additional selected $B$ factor. Thus, for any set $O$ of roots,
\begin{equation}\label{eq:c4-10-9}
    \sup\Bigl|\prod_{j\in O}(1-P_j)\Delta_{\boldsymbol{\mathsf J}}f\Bigr|\leqslant u^{q+| O|},\qquad u= C_{\chi} g_0(1+r)\;,
\end{equation}
the choices of the extra increment, at most $\prod_{j\in O}q_j\leqslant 2^q$, being absorbed in $C_{\chi}$. The row-order bound $2K+1$ covers these remainders simultaneously because the roots have disjoint selected labels and a row has at most $K$ inputs. A single background derivative hits either a row coefficient or one of the $L_f$ background occurrences of the closed word; \eqref{eq:c4-10-4} and the product rule give
\begin{equation}\label{eq:c4-10-10}
    \sup\Bigl|\mathscr D_{\omega}\prod_{j\in O}(1-P_j)\Delta_{\boldsymbol{\mathsf J}}f\Bigr|\leqslant C_{\chi}\|\omega\|_{\infty}(L_f+q)\,u^{q+| O|}\;.
\end{equation}
These are the inequalities used in the commuting-projection cancellation of Lemma~\ref{lem:c4-7-2}. Roots with an empty ordinary selection retain only these remainders under their even Gaussian integration; roots with an ordinary selection use their independent Cauchy vanishing. This recovers the product $u^{q_j}v^{q_j\bmod2}$, and its one-background-derivative version, before any total variation majorant is summed; every term remains a closed contraction with the same fixed $h$ and the original word order.

\subsubsection{Tree gauge on a comparison rectangle}\label{sec:c4-10-5}
\begin{lemma}[tree gauge]\label{lem:c4-10-5}
Let $Q$ be an injectively lifted fine rectangle of largest side $R\geqslant 1$ with a root $o\in Q$, let $U$ be a unitary field, and let $\delta_Q(U)=\max_{p\subset Q}| U(\partial p)-\id|_{\sigma}$. For $x\in Q$ let $P_x$ run from $o$ to $x$ in coordinate order, each coordinate run oriented by its displacement; these paths form a rooted tree. Put $g(o)=\id$, $g(x)= U(P_x)$, so that $U^g(b)= U(P_x)U(b)U(P_{x'})^{-1}$ for $b=(x,x')$ is the based closed transport at $o$ of the fundamental loop of $b$. Then
\begin{equation}\label{eq:c4-10-15}
    U^g(b)=\prod_{j=1}^{m_b}G_jU(\partial p_j)^{\epsilon_j}G_j^{-1},\qquad m_b\leqslant 3R,\qquad\epsilon_j\in\{1,-1\}\;,
\end{equation}
with every plaquette $p_j$ in $Q$, and
\begin{equation}\label{eq:c4-10-16}
    | U^g(b)-\id|_{\sigma}\leqslant 3R\,\delta_Q(U)\;.
\end{equation}
There is a constant $C_{\mathrm{tg}}$, depending only on the fixed group chart and $\tau_{\mathrm b}$, such that $C_{\mathrm{tg}}R^2\delta_Q(U)\leqslant \tau_{\mathrm b}/2$ implies that every $U^g(b)$, $b\subset Q$, lies in the group chart, with principal phase $B_b$ satisfying
\begin{equation}\label{eq:c4-10-17}
    \max_{b\subset Q}h(B_b)\leqslant C_{\mathrm{tg}}R^2\delta_Q(U)\leqslant \tau_{\mathrm b}/2\;.
\end{equation}
\end{lemma}

\begin{proof}
For a bond of direction $\mu$, move its last $\mu$-step past the later coordinate runs of $P_x$. Each interchange with a perpendicular fine edge replaces the ordered word by one plaquette transport conjugated by the preceding word; all plaquettes stay in $Q$, and the number of interchanges is at most $\sum_{\nu>\mu}| x_{\nu}-o_{\nu}|\leqslant 3R$. Negative coordinate runs reverse a plaquette orientation without altering this count. Induction on the interchanges gives the exact nonabelian word \eqref{eq:c4-10-15}, the order being that of the step interchanges and $G_j$ the preceding transport to the plaquette, so every factor is a conjugated closed plaquette; no commutative plaquette replacement is used. Unitary invariance of $|\bigcdot |_{\sigma}$ and telescoping give \eqref{eq:c4-10-16}, with a constant independent of $R$. Norm conversion to the operator norm and the fixed local group logarithm give $h(\log U^g(b))\leqslant 3C_G\sqrt N\,R\,\delta_Q$ whenever that distance is inside the group chart; since $R\geqslant 1$, requiring $C_{\mathrm{tg}}\geqslant 3C_G\sqrt N$ and $C_{\mathrm{tg}}\geqslant 6\sqrt N\tau_{\mathrm b}/r_G$ makes $C_{\mathrm{tg}}R^2\delta_Q\leqslant \tau_{\mathrm b}/2$ sufficient for the chart condition and for \eqref{eq:c4-10-17}. The constants do not grow with $b$ or the scale $R_1$. If the sufficient condition fails, the two-endpoint coefficient bound \eqref{eq:c4-12-2} is used instead of the interpolation (Section~\ref{sec:c4-12-2}), so the lemma covers growing good boxes without imposing their smallness as a hypothesis. The root $g(o)=\id$ leaves the complementary $h$ fixed in every closed contraction.
\end{proof}

\subsection{Identification of the centred coefficient and its local extension}\label{sec:c4-11}
This section concerns $c_W$ of Section~\ref{sec:c4-6-1}, with $W$ fixed on a covered branch; its first caller is the interpolation identity of Lemma~\ref{lem:c4-9-1}. The conclusions retain the hypotheses of Section~\ref{sec:sourced-effective-interaction} (its Section~\ref{sec:joint-extraction} included) and Appendix~\ref{app:background-estimates}, with their fixed charts, carriers, cutoffs and parameters, and do not identify the regional scalar $S^{\circ}_{\xi,C}$ with the perimeter/corner scalar.

\subsubsection{Which exterior dependence is removed}\label{sec:c4-11-1}
Write $s$ for the exterior-response parameter and $t$ for the background-path parameter. Section~\ref{sec:sourced-effective-interaction} defines
\begin{equation}\label{eq:c4-11-1}
    K_{\xi;s}(a)=\Theta_{\xi}(a)\exp\Bigl(-\frac12\langle a,H_{\xi,II}a\rangle-s\langle a,H_{\xi,IJ}A_{\mathrm e}\rangle+\sum_{\mathfrak p\ \mathrm{near}}\mathfrak p(a+sA_{\mathrm e})\Bigr)\;,
\end{equation}
with $U_1$ held fixed in the $s$-interpolation. The collar identity and the boundary-response identity give, with $F_{\xi}(s)$ defined with the same Wilson function $f_{\xi}$,
\begin{equation}\label{eq:c4-11-2}
    \mathcal J^{\mathrm a}_{\xi}[1]= e^{c_{\xi}}I_{\xi;1},\qquad F_{\xi}(1)= F_{\xi}(0)+\mathcal B_{\xi}[f_{\xi}]\;.
\end{equation}
On a covered branch every contour row has its entire free-input set in $I$. Since $\Psi_{\xi}$ is the identity on $I$ and $A_{\mathrm e}$ vanishes there, row locality implies, for every $a$ in the integration box, that every contour link in $f_{\xi}(a)$ is its reconstruction from $a$ alone at background $U_1$: this is the argument of the proof of Lemma~\ref{lem:covered} with $a$ retained instead of set to zero, using Lemma~\ref{lem:row-locality}, and it uses no estimate. Thus both the source and $K_{\xi;0}$ have zero retained free input, while the dependence of $U_1$ itself on the exterior determining data remains. Every selected measure in $c_W$ is constructed from $K_{\xi;0}$: its potentials are the near primitives at $a$, its Hessian is $H_{\xi,II}$ and its cutoff is $\Theta_{\xi}$. The factors involving $\Psi_{\xi}(A_{\mathrm e})$, the far primitives and the boundary-map Jacobian are in $c_{\xi}$, and the interaction between $a$ and $A_{\mathrm e}$ is accounted for by $\mathcal B_{\xi}$; they are not factors of $c_W$. Consequently the physical comparison keeps
\begin{equation}\label{eq:c4-11-3}
    \frac1{\mathcal Z_0}\int\mathcal O\bigl[e^{c_{\xi}}I_{\xi;1}F_{\xi}(0)\bigr]\dd V+\frac1{\mathcal Z_0}\int\mathcal O\bigl[e^{c_{\xi}}I_{\xi;1}\mathcal B_{\xi}[f_{\xi}]\bigr]\dd V\;,
\end{equation}
and does not replace $I_{\xi;1}$ by $I_{\xi;0}$. This is the role of the boundary separation in CRE (1.26)--(1.30), not a claim that the exterior input of the original conditional integral vanishes.

\subsubsection{Covariance of the selected coefficient}\label{sec:c4-11-2}
For a site gauge $g$ let $R_ga(b)=\operatorname{Ad}_{g(b_-)}a(b)$ on free bonds; $R_g$ is orthogonal for $h$, preserves each free bond subspace and commutes with geometric restrictions. The intertwinings in the proof of Lemma~\ref{lem:hessian-gauge-locality}, $Q_{U^g}R_g= R^{(1)}_gQ_U$, $\mathcal N_{U^g}\circ R_g= R^{(1)}_g\mathcal N_U$, $h_{U^g}= R_gh_U(R^{(1)}_g)^{-1}$ and $\mathsf C_{U^g}= R_g\mathsf C_UR_{g,*}^{-1}$ with $(R^{(1)}_gv)(c)=\operatorname{Ad}_{g(c_-)}v(c)$ (the row covariance of Section~\ref{sec:regular-hessian-stencils}, used in the same way in the proof of Lemma~\ref{lem:schur}), give $q_{U^g}(R_ga)= R_gq_U(a)$, branch included: $R_gq_U(a)$ solves the transformed row equation and lies in the same invariant small ball, and the uniqueness of the small row solution there (Lemma~\ref{lem:row-locality}; the uniqueness ball of Lemma~\ref{lem:c2-uniform-chart}(ii)) identifies it with the selected solution. The joint fixed point of Section~\ref{sec:c4-10-1} is identified with this solution by choosing its phase radius so that its $O(\| A\|^2)$ bound lies inside the uniqueness ball of Lemma~\ref{lem:c2-uniform-chart} (Section~\ref{sec:c4-11-5}); the inverse and mean branches are those fixed in the row construction.

The action of each plaquette is invariant and each relative soft-gauge mean is conjugated at its base, so the constant, linear and quadratic Taylor coefficients of the energy primitives transform with their arguments, the force--constraint term included; subtracting them preserves the covariance of the compensated primitives even at noncritical backgrounds. The Haar density is conjugation invariant. The constraint-Jacobian endomorphism transforms by conjugation between the same free and designated spaces, so its determinant and its selected logarithm are unchanged, the logarithm branch being fixed by continuation from zero within the nonvanishing chart. Hence each primitive and each host sum $\mathbf V_{\xi}(Y)$ is invariant under the simultaneous transformation of background and free variables.

The Hessian, each restricted seed inverse and every fixed walk transform by $R_g$ on their indices, by the entry identities of Lemma~\ref{lem:hessian-gauge-locality} and the geometric walk formula; all geometric choices, Cauchy variables and Mayer parameters are unchanged. In the selected integral change variables $B\mapsto R_gB$, $X\mapsto R_gX$: Lebesgue and standard Gaussian measures, the two allowed cutoff norms (\eqref{eq:marginal-cutoff-invariance}), the conditional Gaussian determinant and the cross and compensator terms are preserved, so each selected measure is pushed forward by $R_g$, and in the probe formula $\delta_B-\delta_0$ is pushed to $\delta_{R_gB}-\delta_0$. Summing selections,
\begin{equation}\label{eq:c4-11-4}
    \widehat\nu^{\mathrm f}_{U^g,D}=(R_g)_*\widehat\nu^{\mathrm f}_{U,D}\;.
\end{equation}
The vacuum activities, their Ursell products and $w_{\xi}(W)$ are scalar invariants, the incompatibility graph and every numerical combinatorial factor being field independent. Both $T_WP_W^{-1}$ and the complementary $h$ transform at the initial interval vertex $o$, so
\begin{equation}\label{eq:c4-11-5}
    c_W(U^g;g(o)hg(o)^{-1})= c_W(U;h)\;.
\end{equation}
The rooted comparison gauge has $g(o)=\id$. Open transports occur only in this equality inside the closed character; no norm of an open transport is used.

\subsubsection{Domain and endpoint}\label{sec:c4-11-3}
Section~\ref{sec:c4-8-4} defines the small interior groups by \eqref{eq:c4-8-9}. The support of such a group includes every root domain, each vacuum-cluster domain occurrence, their coefficient halos and the whole contour interval $\mathcal H(W)$ with its closing transport; the hull-extent bound is needed here, since a box enclosing only the random inputs would not suffice. Under \eqref{eq:c4-8-9} the enclosing rectangle has diameter at most $C_{\mathrm{box}}LMn(W)$, lies within $b/2$ of a source contact and is injectively liftable under the period convention of Section~\ref{sec:c4-9-2}; on the far-history event the box and its halos lie in the regular region. All other groups, including winding or boundary-reaching groups, stay in the large-domain tail of Section~\ref{sec:c4-8-4}. This is the small-local-domain versus large/distant-domain separation of RG1 (3.5)--(3.9) and of CRE's boundary separation; no interpolation is asserted for a group in the tail.

On a small interior group first use its rooted tree gauge. If $C_{\mathrm{tg}}\operatorname{diam}(Q_W)^2\delta_{Q_W}(U_1)\leqslant \tau_{\mathrm b}/2$, define $B_W$ and $U_{W,t}$ as in \eqref{eq:c4-9-1}. With $\tau_{\mathrm b}\leqslant \tau_L/4$ and $c_L\tau_{\mathrm b}\leqslant \gamma_L/4$ (Lemma~\ref{lem:c4-7-3}) the entry Lipschitz estimate and the flat compression gap supply the two-sided Hessian bounds throughout the path. That the path need not consist of constrained minima is allowed only because its compensated primitive formulas are defined on the joint chart; physical criticality identifies those formulas at the endpoint, not at intermediate $t$.

For each selected coefficient Section~\ref{sec:c4-10-2} fixes all walk parameters and restricts every surviving walk and coefficient stencil to its final domain; the seed inverse is the inverse on the seed, not the restriction of a global inverse. Thus, on the supports included in $Q_W$, the $t=1$ recipe equals the gauge-transformed physical centred recipe, its values on bonds outside $Q_W$ cannot enter the coefficient, and \eqref{eq:c4-11-5} returns the coefficient to the original gauge. At $t=0$ the recipe is the regional reference recipe with exactly the same $I$, $\Omega_1$, $W$, cutoffs and constants. This yields the endpoint identities
\begin{equation}\label{eq:c4-11-6}
    c_W(U_{W,1};h)= c_W(U_1;h),\qquad c_W(U_{W,0};h)= a_W\chi(h)\;,
\end{equation}
equalities of the selected coefficient, not of the outer history density. If the tree-gauge phase condition fails, the two-endpoint majorant of Section~\ref{sec:c4-6-3} is used; if the box condition fails, the whole group is in the tail. Neither case is assigned an endpoint identity it does not satisfy.

For differentiation under the selected integrals, the norm of $X$ in Section~\ref{sec:c4-10-3} is the norm of $X$ restricted to $J_Z$; all other auxiliary columns integrate out with mass one and have zero derivative. The common strictly positive $c_{\mathrm G}$ of \eqref{eq:c4-10-13} supplies a background-independent Gaussian majorant for that finite-dimensional integral and its first derivative; the selected sums use the summable majorants of Lemma~\ref{lem:c4-7-1} and Section~\ref{sec:c4-7-4}, and the vacuum occurrence sums those of Lemma~\ref{lem:c4-7-4} and Section~\ref{sec:c4-8-4}. These are the dominated-convergence hypotheses used in the proof of Lemma~\ref{lem:c4-9-1}; the group derivative includes the derivative of $w_{\xi}(W)$. RG1 (3.21), (4.35)--(4.36) estimate changes of reconstruction and domain rather than declaring them zero; no such change is made in the identities above, the regional carrier and each local recipe being fixed. A failure of the stated finite-support or covariance identity would invalidate the endpoint step and would require a retained comparison error, which the tree-gauge estimate alone could not absorb.

\subsubsection{Endpoint conventions and the lifting claim}\label{sec:c4-11-4}
The endpoint statement is restricted to covered far histories and the small interior groups of Section~\ref{sec:c4-8-4}. Extend the rooted site gauge from $Q_W$ to the whole torus by the identity on sites outside $Q_W$. Every background bond consulted by the coefficient has both endpoints inside $Q_W$: in defining $Q_W$ the enclosing rectangle is enlarged to contain both endpoints of each Hessian stencil, row set, block-mean path, plaquette and interval or closing-path bond, an enlargement that is part of the fixed halo included in $C_{\mathrm{box}}$. Boundary bonds of the site-gauge extension are not consulted by the coefficient.

Only the finite formula $w_{\xi}(W)=(-1)^{k(W)}/k(W)!$ times the product of the specified vacuum cluster occurrences is varied along $U_{W,t}$; neither the global partition function nor the global exclusion ratio $q_{\xi}(\mathcal A)$ is an operand of the interpolation. Lemma~\ref{lem:c4-7-4} supplies the sum of the occurrence derivative majorants and \eqref{eq:c4-8-10} applies the finite product rule; the exclusion ratio was expanded at the original background before the group construction. Thus no identity or positivity statement for a global partition function at $U_{W,t}$ is needed to identify either endpoint. The regional scalar $a_W$ is the one defined in Section~\ref{sec:joint-extraction} from $w^{\circ}_{\xi}(W)$ and the same root measures at the identity background (\eqref{eq:part1-reference-group}); it is not an infinite-volume coefficient, and its closed Schur identity \eqref{eq:part1-closed-schur} is the $t=0$ identity used here.

For the lifting claim, represent each constituent connected domain by a tree of its constituent coarse cubes and adjoining faces and root the union at a source contact. Each root or cluster-domain tree has at most its number of cubes minus one edges; incompatibility joins have bounded range, and the finite reconstruction halos add a fixed range at their contacts. The contour-hull extent bound bounds the additional contour edges joining the components by $CLMn(W)$; the network contains the whole interval $\mathcal H(W)$ and its closing path, not only their contacts. The resulting connected bond network has a spanning traversal of length at most $CLMn(W)$ with $C$ fixed. If this is less than $N_{\min}/4$, lift the traversal from its root: every lifted vertex lies within that distance of the root, and two lifts of the same vertex could differ only by a nonzero period vector, of length at least $N_{\min}$, while the two paths giving them have total length less than $N_{\min}/2$, a contradiction. Every point of the enclosing coordinate rectangle and of its stipulated halos is within $\operatorname{diam}(Q_W)\leqslant C_{\mathrm{box}}LMn(W)$ of the root, and two lifts of such a point would differ by a nonzero period vector, of length at least $N_{\min}$, while being joined through the root by two segments of total length at most $2C_{\mathrm{box}}LMn(W)< b< N_{\min}$ (\eqref{eq:c4-8-9}, \eqref{eq:c4-9-3}); so the rectangle and its halos project injectively. $C_{\mathrm{box}}$ is chosen to include these constants; with \eqref{eq:c4-8-9} and \eqref{eq:c4-9-3} the small-group criterion implies this sufficient lifting condition, and failure to lift cannot occur in the small-group class. No curvature estimate for a winding loop is used. In applying the entry estimates, $\tau_{\mathrm b}$ is taken small enough for both the $h$-maximum and the stated entry-chart norm, using $h\leqslant |\bigcdot |_{\sigma}\leqslant \sqrt Nh$, before $g_0$ is chosen. Geometrically allowed zero walks re-index the same series and do not change its value or support domains.

\subsubsection{Source and radius specifications}\label{sec:c4-11-5}
The imported mean covariance is RG1 (0.6), $\mathsf M(\{ uU_jv\})= u\mathsf M(\{ U_j\})v$ for the specified mean on its small-diameter domain, which RG1 (0.10) asserts for the selected mean; it is used for unitary site gauges with no smallness restriction on the gauge values, the relative distances and the real chart domains being preserved. The cutoff invariance is \eqref{eq:marginal-cutoff-invariance}.

The curvature input of Section~\ref{sec:c4-8-4} is Lemma~\ref{lem:c2-plaquette-deviation}. For a plaquette of $Q_W$ choose one of its vertices $x$; $x\in\Lambda_1$ since $Q_W\subset\Lambda_1$. Within the $L$-block containing $x$ choose a neighbouring vertex still in that block ($L\geqslant 3$). That internal fine bond is free, since designated bonds join distinct blocks, and belongs to $I$ because it has the endpoint $x\in\Lambda_1$; so the plaquette has a bond within two fine-bond metric units of $I$, within the $3R_{\mathrm c}$ the lemma requires. This yields the bound $\delta_{Q_W}(U_1)\leqslant C\varepsilon_1$, with the fixed conversion between the lemma's operator norm and the norm $|\bigcdot |_{\sigma}$ of Lemma~\ref{lem:c4-10-5}.

For the branch identification of Section~\ref{sec:c4-11-2} let $C_{\mathrm{D,joint}}$ be the constant in $\|\widetilde D_Z(A)\|\leqslant C_{\mathrm{D,joint}}\| A\|^2$ on the joint chart of Section~\ref{sec:c4-10-1}. After choosing that chart reduce $\theta_0$ so that $2\theta_0\leqslant \vartheta_*/2$ and $4C_{\mathrm{D,joint}}\theta_0^2<\rho_{\mathrm d}$; the norms of $Z$ and $A$ are those of Lemma~\ref{lem:c2-uniform-chart}, in which $\vartheta_*$, $\rho_{\mathrm d}$ and $\tau_L$ are measured, and the background radius satisfies $2\tau_{\mathrm b}\leqslant \tau_L/2$ (Lemma~\ref{lem:c4-7-3}), so every real background with $\| Z\|<2\tau_{\mathrm b}$ lies in the chart of that lemma. At real backgrounds the joint solution for $\| A\|<2\theta_0$ then lies in the uniqueness ball of Lemma~\ref{lem:c2-uniform-chart} and is its selected solution. The coupling threshold is chosen after these radii so that the phase inputs $g_0\mathsf C_UB$ of the enlarged free box remain in $\| A\|<\theta_0$; no expansion parameter is enlarged after this choice.

\subsection{Quantitative assembly at the signed physical caller}\label{sec:c4-12}
This section proves the first-step estimate for $E_{\mathrm{def}}$ of \eqref{eq:c4-1-1}. It uses the hypotheses of Section~\ref{sec:sourced-effective-interaction} (its Section~\ref{sec:joint-extraction} included) and Appendix~\ref{app:background-estimates}, the compatible even-period and square-fit conventions of Section~\ref{sec:c4-9-2}, and the same irreducible normalised character. Constants may depend on the fixed group, the character, $L$, $M$ and the fixed expansion and chart choices, but not on the coupling, the square length, the compatible torus size or the history; the coupling is chosen last. The regional scalar $S^{\circ}_{\xi,C}$ is retained; no identification with a perimeter/corner scalar is asserted.

\subsubsection{The joint derivative supplied to the group coefficient}\label{sec:c4-12-1}
The quantitative derivative is obtained from Lemmas~\ref{lem:c4-7-1} and~\ref{lem:c4-7-2} and Section~\ref{sec:c4-10}. For each selected root the source is $\Delta_{\mathsf J}f$, a closed character difference; \eqref{eq:c4-10-9} bounds its mixed Taylor remainder by $u^q$ and supplies an extra $u$ for each odd leading term removed by reversal. The root's Gaussian derivative is even in its joint $(B,X)$ variables, including its determinant and compensator derivatives, and a derivative of a selected ordinary potential carries a factor $\zeta$ and vanishes at $\zeta=0$; so the independent-root parity/Cauchy argument of Lemma~\ref{lem:c4-7-2} applies after one background derivative. Lemma~\ref{lem:c4-7-1}(b),(c) (\eqref{eq:c4-7-4}, \eqref{eq:c4-7-5}) bounds a differentiated potential without dividing by that potential; \eqref{eq:c4-10-6} sums the derivatives of all geometrically indexed walks on the same spatial contours; \eqref{eq:c4-10-14} pays each Gaussian derivative by $C_MN_1(D)$ times the same $e^{N_1(D)}$ determinant majorant, with $X$ restricted to $J_Z$, all other auxiliary variables integrating to one. Host and Gaussian derivatives therefore insert one factor $N_1(D)$, not one new constant per selected label. The source derivative inserts the number $L_f$ of varying contour/closing links plus the number of selected labels, which Section~\ref{sec:c4-7-1} bounds by $C_M\sum_DN_1(D)$. The resulting joint estimate is
\begin{equation}\label{eq:c4-12-1}
    |\mathscr D_{\omega}\mathfrak Q_{\xi}(\mathcal A;f)|\leqslant C_{\chi}\|\omega\|_{\infty}\Bigl(L_f+\sum_{D\in\mathcal A}N_1(D)\Bigr)\prod_{D\in\mathcal A}\rho\,m(D)\;,
\end{equation}
with $\rho$ of \eqref{eq:c4-4-8} and $m(D)=\varepsilon_*e^{-\lambda_*\ell_1(D)}$ of Lemma~\ref{lem:c4-4-1}; the pairing $\mathfrak Q_{\xi}(\mathcal A;f)$ and its source are evaluated at the background $U_{W,t}$ of the fixed coefficient path of Section~\ref{sec:c4-11} in place of $U_1$ (the subscript $\xi$ is kept for brevity), and the estimate differentiates neither an outer history weight nor a global exclusion ratio. For a vacuum occurrence $\mathfrak c$ the selected-integral derivative gives \eqref{eq:c4-7-8}; differentiating the finite product $w_{\xi}(W)$ by \eqref{eq:c4-8-10}, its occurrences $\Phi^{\mathrm v}(\mathfrak c_j)$ likewise evaluated at $U_{W,t}$, and adding the root-integral derivative yields, uniformly for $h\in G$,
\begin{equation}\label{eq:c4-12-2}
    | c_W(U;h)|\leqslant M(W),\qquad|\mathscr D_{\omega} c_W(U;h)|\leqslant C_{\chi}\|\omega\|_{\infty}\bigl(n(W)+|\mathcal H(W)|\bigr)M(W)\;.
\end{equation}
The undifferentiated estimate holds at the original physical background and at the regional flat reference. $M(W)$ contains $\rho$ for every marked root and the original factorials; polynomial factors have not been summed or hidden in $M(W)$, and the factor from the relative interval's inverse word, at most a fixed multiple of $|\mathcal H(W)|$, is included in $C_{\chi}$.

\subsubsection{Payment of the local and tail losses}\label{sec:c4-12-2}
\begin{proof}[Proof of Lemma~\ref{lem:c4-6-4}]
Write $s=\mathfrak s(W)\geqslant 1$, so that $\mathfrak a_W=\eta s$. The fixed hull and box bounds give $n(W)+|\mathcal H(W)|\leqslant CLMs$ and $\operatorname{diam}(Q_W)\leqslant C_{\mathrm{box}}LMn(W)\leqslant C_{\mathrm{box}}LMs$ (Section~\ref{sec:c4-8-4}). The endpoint identities \eqref{eq:c4-11-6}, the path integral \eqref{eq:c4-9-2} with the derivative bound \eqref{eq:c4-12-2} along it, and the plaquette estimate $\delta_{Q_W}(U_1)\leqslant C\varepsilon_1$ of Section~\ref{sec:c4-11-5} with the phase bound \eqref{eq:c4-10-17} therefore give, for a small interior group passing the phase test,
\begin{equation}\label{eq:c4-12-3}
    \mathfrak d_W(U_1)\leqslant C_{\chi}\varepsilon_1s^3M(W)\;.
\end{equation}
If the phase test of such a group fails, then $C_{\mathrm{tg}}\operatorname{diam}(Q_W)^2\delta_{Q_W}>\tau_{\mathrm b}/2$; with $\delta_{Q_W}\leqslant C\varepsilon_1$ and $\operatorname{diam}(Q_W)\leqslant CLMs$ this gives $1\leqslant C\varepsilon_1s^2$ with constants fixed before the coupling, so the two undifferentiated endpoint bounds \eqref{eq:c4-12-2} give \eqref{eq:c4-12-3} again, as $s\geqslant 1$. No curvature comparison is used for a group outside the small interior class; for such a group
\begin{equation}\label{eq:c4-12-4}
    \mathfrak d_W(U_1)\leqslant 2M(W),\qquad s\geqslant T_b\coloneqq \frac{b}{4C_{\mathrm{box}}LM}\;.
\end{equation}
For the local polynomial payment let
\begin{equation}\label{eq:c4-12-5}
    C_{\eta,3}=\sup_{s\geqslant 1}s^3e^{-(31\eta/8)s}<\infty\;,
\end{equation}
so that $s^3e^{\eta s/8}\leqslant C_{\eta,3}e^{4\eta s}$; for the tail, $s\geqslant T_b$ gives $e^{\eta s/8}\leqslant e^{-2\eta T_b}e^{4\eta s}$. Hence the two weighted pin estimates of Lemmas~\ref{lem:c4-8-2} and~\ref{lem:c4-8-3} imply
\begin{equation}\label{eq:c4-12-6}
    \sum_WK_{\chi}\mathfrak d_W(U_1)e^{\mathfrak a_W/8}\leqslant C_{\chi}\bigl(\varepsilon_1+e^{-2\eta T_b}\bigr)\sum_WM(W)e^{4\eta\mathfrak s(W)}\leqslant C_{\chi} n_C\rho\bigl(\varepsilon_1+e^{-cb/(LM)}\bigr),\qquad c=\frac{\eta}{2C_{\mathrm{box}}}\;.
\end{equation}
The last summation uses that every $W$ has a nonempty contour interval: summing the contour-pin estimate \eqref{eq:c4-6-5} over the $n_C$ traversals counts every group at least once. The cubic moment and the tail use the existing $4\eta$ weight, not an additional size-dependent hypothesis or a loss in the scalar $S^{\circ}_{\xi,C}$. The exceptional whole-contour merge is the one paid in Lemma~\ref{lem:c4-8-3}, where the long class pays both the arbitrary extra components and the whole-contour weight. This is \eqref{eq:c4-6-8}.
\end{proof}

\subsubsection{The signed family sum and the physical expectation}\label{sec:c4-12-3}
\begin{proposition}[the signed defect bound]\label{prop:c4-12-3}
Under the hypotheses stated at the start of this section, for $g_0$ below the thresholds of Lemmas~\ref{lem:c4-2-1}, \ref{lem:c4-6-2} and~\ref{lem:c4-6-3}, with $R_{\Upsilon}\geqslant 2$ (Section~\ref{sec:c4-7-1}, Lemma~\ref{lem:c4-4-1}), $u\leqslant u_*$ and $v\leqslant 1$ (Lemma~\ref{lem:c4-4-2}) and the small-coupling conditions of Sections~\ref{sec:c4-10-1} and~\ref{sec:c4-11-5}, with $\rho$ of \eqref{eq:c4-4-8}, $\rho= O(g_0^2(1+r)^4)$, $\varepsilon_1= g_1p_0(g_1)$, $b=\lfloor LMR_1/16\rfloor$, $\pi_*=\min(p_0(g_0),p_0(g_1))$ and the constants $C'_L$, $c_*$ of \eqref{eq:c4-2-2} and $C_{\chi}$, $c$ of \eqref{eq:c4-6-8},
\begin{equation}\label{eq:c4-12-7}
    | E_{\mathrm{def}}|\leqslant C'_Ln_Ce^{-(c_*/2)\pi_*^2}+\frac1{\mathcal Z_0}\int\mathcal O\Bigl[\mathsf w_{\xi}\mathbf 1_{A_{\mathrm{cov}}\cap\mathsf N_b^c}| S^{\circ}_{\xi,C}|\Bigl(e^{C_{\chi} n_C\rho(\varepsilon_1+e^{-cb/(LM)})}-1\Bigr)\Bigr](V)\,\dd V\;,
\end{equation}
and consequently
\begin{equation}\label{eq:c4-12-8}
    | E_{\mathrm{def}}|\leqslant C'_Ln_Ce^{-(c_*/2)\pi_*^2}+e^{C_{\chi} n_C\rho(\varepsilon_1+e^{-cb/(LM)})}-1\;.
\end{equation}
\end{proposition}

\begin{proof}
Lemma~\ref{lem:c4-6-3} supplies the single joint ratio \eqref{eq:c4-6-7}. Combined with the closed-contraction product bound \eqref{eq:c4-6-4}, formed after each $d_W$, the nonempty compatible-family sum is bounded by the exponential series of the summable one-group weight \eqref{eq:c4-6-8}, which on covered far histories is \eqref{eq:c4-6-9} (Lemma~\ref{lem:c4-6-5}); the countable-group limit is the dominated finite-species limit of Lemma~\ref{lem:c4-9-2}, with the same positive majorant and the nonzero reference partition function supplied by its convergent logarithm, so no quotient is formed at a zero partition function. Combining the near-history estimate \eqref{eq:c4-2-2} with the positivity of the original history measure gives \eqref{eq:c4-12-7}; all history weights, deferred integrations and the physical normalisation are those of \eqref{eq:c4-1-1}, and the exact source on the left was summed before absolute values were taken. The regional reference scalar $S^{\circ}_{\xi,C}=\int f^{\circ}_{\xi} K^{\circ}_{\xi;0}\,\dd a/I^{\circ}_{\xi;0}$ (\eqref{eq:part1-reference-gas}) is a normalised integral of $f^{\circ}_{\xi}$, a normalised character of modulus at most one, against the positive reference kernel $K^{\circ}_{\xi;0}$, so $| S^{\circ}_{\xi,C}|\leqslant 1$, and the same physical normalisation yields \eqref{eq:c4-12-8}.
\end{proof}

The exponential in \eqref{eq:c4-12-8} is not replaced by its linear term unless its argument is bounded; in particular no assertion of UV summability for a fixed physical loop is made. The conclusion has exactly the hypotheses and quantitative scope stated at the beginning of this section.

\subsubsection{Fixed reserve conventions}\label{sec:c4-12-4}
The exponential vacuum estimate used in Lemma~\ref{lem:c4-8-2} holds with $n(\mathfrak c)/2$ as well as with $n(\mathfrak c)/4$: the rooted-tree proof of Lemma~\ref{lem:activity-bounds} applied to the numerical activity $m(D)e^{N_1(D)/2}$ with node cost $N_1(D)/8$ has combined exponent $5/8$, below the $1+6\theta$ of the activity payment \eqref{eq:fixed-kernel-raw-activity-payment}, which gives the first bound of \eqref{eq:c4-7-9}. No reserve of Appendix~\ref{app:background-estimates} is re-spent. Adjacency allowance: in the footprint-contact sums of Lemma~\ref{lem:c4-6-2} (\eqref{eq:c4-6-5}) and Lemma~\ref{lem:c4-6-3}, footprint contact includes wall adjacency of two $LM$-cubes (a shared face, edge or corner), not only intersection of the two cube sets; counting every cube touching a given one at a face, an edge or a corner gives at most $3^4-1$ neighbours, so the spatial pin is summed over the neighbouring cubes as well, which inserts a fixed factor in the constant $C$ multiplying $\rho$. The smallness condition $C\rho LM\leqslant \eta/8$ of Lemma~\ref{lem:c4-6-2} is imposed after incorporating this factor; the allowance does not alter $\rho$, the reserve exponent or the loop-length dependence. The tail rate is $c=\eta/(2C_{\mathrm{box}})$ as in \eqref{eq:c4-12-6}; earlier weaker tail displays remain consequences of it with a reduced positive $c$.

\section{The Regional Scalar of the Covered Branch}\label{app:regional-scalar}
This appendix identifies the scalar $S^{\circ}_{\xi,C}$ of the branch identity \eqref{eq:c-one-step-branch-identity}, defined in Section~\ref{sec:joint-extraction} as the reference-gas sum \eqref{eq:part1-reference-gas}, with the regional square scalar \eqref{eq:regional-square-scalar} at the carrier $I$, and thereby with the perimeter/corner coefficient $\kappa_C$ of \eqref{eq:extracted-loop-scalar} up to the error of Proposition~\ref{prop:regional-square-extraction}. It is the statement used in Section~\ref{sec:remaining-physical-comparison} at \eqref{eq:retained-leading-source}; its labelled counterpart at the retreated carrier $D(\Lambda)$ is Lemma~\ref{lem:c5-7-1} with Corollary~\ref{cor:c5-7-2}. The quantitative comparison of the regional square scalar with the bulk square is the proof of Proposition~\ref{prop:regional-square-extraction}, in particular \eqref{eq:regional-square-boundary-comparison} and \eqref{eq:regional-square-periodic-comparison}; nothing is added to it here.

\subsection{Conventions}\label{sec:c4c-1}
Fix one of the two cutoff prescriptions $\sigma$, the bare coupling $g_0$, the irreducible normalised character $\chi$, the averaging and mean prescription of RG1 and the designated-bond convention, as in the rest of the paper, with the standing conditions \eqref{eq:c-standing-conditions} and the collar conditions \eqref{eq:regional-collar-sizes}. The fine periods are $N_{\mu}= LK_{\mu}$, and $C= C_n$ is a coarse translate of the counterclockwise square with vertices $0$, $nLe_1$, $nL(e_1+e_2)$, $nLe_2$, with $n$ coarse segments of length $L$ on each side and $K_{\mu}\geqslant 3(n+2)$ for every coarse period; its physical perimeter is $P(C_n)=4nL\varepsilon$. Let $A_{\mathrm{cov}}$ be the set of covered branches of Lemma~\ref{lem:covered} and let $\xi\in A_{\mathrm{cov}}$. The young carrier is $D= I=\Lambda_1^*$ (Section~\ref{sec:sourced-effective-interaction}), the set of free bonds with an endpoint in $\Lambda_1$, which is the carrier $D_+$ of Section~\ref{sec:regional-identification} and contains its $D_-$; $D^c$ is its complement in the torus free carrier $E_{*,T}$; distances are those of the periodic fine-bond metric $d_T$ of \eqref{eq:marginal-covariance-input}. For a coarse segment $j$ of the cyclic index set $J_C=\{0,\ldots,4n-1\}$ let $K_j$ be the union of the fixed free-input stencils of its fine traversals, and with the fixed range $R_0=4(1+R_*+R_H+L)$ of Section~\ref{sec:regional-identification} (the paragraph following \eqref{eq:regional-collar-sizes}; the $R_0$ of \eqref{eq:regional-square-distances}) put
\begin{equation}\label{eq:c4c-1-1}
    d_j=\max\bigl(0,d_T(K_j,I^c)-R_0\bigr),\qquad\iota_T=\max\bigl(0,N_{\min}/4-R_0\bigr),\qquad N_{\min}=\min_{\mu} N_{\mu}\;,
\end{equation}
the distance to the empty set being infinite with $e^{-\infty}=0$; there is no distance to a side endpoint in $d_j$. The constants below depend on the character, the group, $L$ and the forest parameters only, not on $n$, the compatible periods or the shape of $I$; $g_0$ is chosen after them, below the common flat and regional thresholds. $\tau$ is the marked-forest parameter of \eqref{eq:marginal-fixed-forest-bound}; enlarging a fixed majorant constant multiplies $\tau$ by a fixed factor, absorbed in the constants and the coupling threshold.

\subsection{Identification with the regional square integral}\label{sec:c4c-2}
\begin{lemma}[the reference-gas scalar at the carrier $I$]\label{lem:c4c-2-1}
On $A_{\mathrm{cov}}$, under the conventions of Section~\ref{sec:c4c-1} and the thresholds displayed in the proof of Proposition~\ref{prop:regional-square-extraction} (also \eqref{eq:c5-7-T}(ii)),
\begin{equation}\label{eq:c4c-2-1}
    S^{\circ}_{\xi,C_n}= S^{\circ}_{I,C_n,T}=\frac{\displaystyle\int\Theta^{\sigma}_I(B)e^{P_{I,T}(B)}\chi\bigl(\Hol_{U_{g_0,T}(i_{I,T}B)}C_n\bigr)\,\dd\mu_{C_{I,T}}(B)}{\displaystyle\int\Theta^{\sigma}_I(B)e^{P_{I,T}(B)}\,\dd\mu_{C_{I,T}}(B)}\;,
\end{equation}
the right side being \eqref{eq:regional-square-scalar} at $D= I$, with every original primitive label retained, including the labels crossing the boundary of $I$, and with the exterior link value the identity.
\end{lemma}

\begin{proof}
By \eqref{eq:part1-reference-gas}, $S^{\circ}_{\xi,C_n}=\int f^{\circ}_{\xi} K^{\circ}_{\xi;0}\,\dd a/I^{\circ}_{\xi;0}$, the normalised expectation of the full reference Wilson loop under the centred kernel at the identity background, whose precision is $H^{\circ}_{\xi}=\iota_I^*\mathsf H^{\Omega_1}_{\id}\iota_I$, whose potentials are the near primitives evaluated at zero retained input and whose cutoff is the bond cutoff on $I$. Both sides of \eqref{eq:c4c-2-1} are ratios of two integrals over $\mathfrak g^I$, so it suffices to identify the factors of the integrands. Precision: by Lemma~\ref{lem:c2-regional-entry} at the identity background the entries of $H^{\circ}_{\xi}$ are $f_{bb'}(0)$, the entries of the flat Hessian \eqref{eq:flat-covariance} on the branch's torus (as in the proof of Lemma~\ref{lem:c2-reference}), and the collar puts every precision and primitive stencil touching $I$ inside the regular region (the identification paragraph of Section~\ref{sec:regional-identification} preceding \eqref{eq:regional-kernel-identification}), so $H^{\circ}_{\xi}= i_{I,T}^*H_Ti_{I,T}$, the diagonal contributions of the plaquettes crossing the boundary of $I$ included; the Gaussian normaliser of $\mu_{C_{I,T}}$ is common to numerator and denominator and cancels within the ratio, not against the outer physical weight. This is the full-square form of \eqref{eq:regional-kernel-identification}: the kernel identification uses only that the source is a closed contraction with the given input stencils, not the straightness of the segment. Reconstruction: the row solutions $\widetilde D_{U,c}$ of Section~\ref{sec:source-phase-primitives} (Lemma~\ref{lem:charts}) at the identity background and those of Section~\ref{sec:regional-identification} agree as analytic germs at zero by their defining equation and uniqueness; the Hilbert-ball domain of Lemma~\ref{lem:charts} (reference version) and the maximum-norm domain of the flat expansion have a convex common phase domain containing the closed cutoff box at the stated thresholds, and analytic continuation on that domain identifies the two solutions throughout the box. Primitives: Section~\ref{sec:source-phase-primitives} groups the gauge primitives by their centre $y$ while \eqref{eq:regional-periodic-weight} indexes them by the pairs $(y,x)$; the sums coincide, the degree-two Taylor residual being linear in the local term, and every primitive whose free support misses $I$ vanishes at $i_{I,T}B$ while every primitive meeting $I$ is retained on both sides. Cutoff: at zero exterior free field the conditional cutoff is $\Theta^{\sigma}_I$ at the threshold $r=\delta_0/g_0$ in the norm $|\bigcdot |_{\sigma}$, as in Section~\ref{sec:regional-identification}. Wilson factor: on a covered branch every contour traversal has its free-input set in $I$, so the proof of Lemma~\ref{lem:covered} applies verbatim with $a$ retained instead of set to zero and gives that each contour link of $f^{\circ}_{\xi}(a)$ is its reconstruction from $a$ alone, that is, $U_{g_0,T}(i_{I,T}B)$ at $B= a$; the boundary map is the identity by the reference-data convention $\Psi^{\circ}=\id$ of Section~\ref{sec:joint-extraction}. Numerator and denominator agree factor by factor, so the ratios agree. The equality is one of integral ratios; no identification of a separately chosen logarithm of the reference group gas of Section~\ref{sec:joint-extraction} with the connected source logarithm of the proof of Proposition~\ref{prop:regional-square-extraction} is used.
\end{proof}

\subsection{The scalar of the branch identity}\label{sec:c4c-5}
\begin{proposition}[the regional scalar of the covered branch]\label{prop:c4c-5-1}
Under the conventions of Section~\ref{sec:c4c-1} and the hypotheses of Proposition~\ref{prop:regional-square-extraction}, on $A_{\mathrm{cov}}$,
\begin{gather}
    S^{\circ}_{\xi,C_n}= e^{-\delta m^{\sigma}(g_0,\varepsilon)P(C_n)}\bigl(z^{\sigma}_{\mathrm c}(g_0,\pi/2)\bigr)^4e^{E^{\circ}_{\xi}}=\kappa_{C_n}e^{E^{\circ}_{\xi}}\;,\label{eq:c4c-5-1}\\
    | E^{\circ}_{\xi}|\leqslant C\tau^2\Bigl[\sum_{j\in J_C}e^{-cd_j}+4ne^{-c\iota_T}+(n+1)e^{-cn}\Bigr]\;,\notag
\end{gather}
with $\kappa_{C_n}\coloneqq e^{-\delta m^{\sigma}(g_0,\varepsilon)P(C_n)}(z^{\sigma}_{\mathrm c}(g_0,\pi/2))^4$ the coefficient of \eqref{eq:extracted-loop-scalar}, built from the coefficients of \eqref{eq:marginal-square-mass} and \eqref{eq:corner-square-factor}, and
\begin{equation}\label{eq:c4c-5-2}
    E^{\circ}_{\xi}= E^{\circ}_{I,C_n,T}=(F_{I,C,T}-F_{C,T})+(F_{C,T}-F_{C,\infty}(n))+E^{\sigma}_n\;,
\end{equation}
the error of Proposition~\ref{prop:regional-square-extraction} at $D= I$: as in its proof, $F_{I,C,T}=\log S^{\circ}_{I,C_n,T}$ on the connected source branch, $F_{C,T}=\log(N^{\sigma}_{C_n,T}/D^{\sigma}_T)$, $F_{C,\infty}(n)= F_{\square,\infty}(n)$ is the bulk logarithm \eqref{eq:corner-fixed-square-limit}, and $E^{\sigma}_n\coloneqq F_{\square,\infty}(n)-n\sum_{v\in\mathscr V_{\square}}\kappa_v-\sum_{(v,w)\in\mathscr C_{\square}}\zeta_{v,w}$ is the corner-extraction error, $| E^{\sigma}_n|\leqslant C\tau^2(n+1)e^{-cn}$ by \eqref{eq:corner-square-extraction-bound}, with $n\sum_v\kappa_v=-\delta m^{\sigma} P(C_n)$ (\eqref{eq:marginal-square-mass}) and $\sum_{(v,w)}\zeta_{v,w}=4\log z^{\sigma}_{\mathrm c}$ (\eqref{eq:corner-square-factor}), so that $F_{I,C,T}=\log\kappa_{C_n}+E^{\circ}_{\xi}$. The scalar and its nonzero logarithm are those of the regional conditional reference integral; no physical history weight or large-field indicator is evaluated at flat fields, and the bound does not assume the loop to be far from the integration boundary: $d_j=0$ retains a boundary contribution of size $C\tau^2$.
\end{proposition}

\begin{proof}
Lemma~\ref{lem:c4c-2-1} and Proposition~\ref{prop:regional-square-extraction} at $D= I$, which is covered on $A_{\mathrm{cov}}$ and is the carrier $D_+$ itself (Section~\ref{sec:c4c-1}), so $D_-\subseteq D\subseteq D_+$ holds with the collar; the three terms of \eqref{eq:c4c-5-2} are \eqref{eq:regional-square-boundary-comparison}, \eqref{eq:regional-square-periodic-comparison} and the corner extraction \eqref{eq:corner-square-extraction-bound}, the first two retaining the changes of kernel and of summation range as errors before the translation-invariant coefficients are identified, as in RG1 (4.35)--(4.37), while the exterior operation and its boundary terms are kept as they stand in this paper (CRE's (1.26)--(1.30), which perform CRE's own separation, are not used, except (1.27) for $g_1$; Section~\ref{sec:remaining-physical-comparison}).
\end{proof}

\begin{corollary}[the physical identity with the scalar substituted]\label{cor:c4c-5-2}
Under the hypotheses of Proposition~\ref{prop:c4c-5-1}, with $E_{\mathrm{def}}$ of \eqref{eq:c4-1-1} and $E_{\mathrm{bdry}}$ of \eqref{eq:c3-6-2},
\begin{equation}\label{eq:c4c-5-3}
    \mathbb E_{\varepsilon}[W_{\chi}(C_n)]=\kappa_{C_n}\,\mathbb E_{\mathrm h}\bigl[\mathbf 1_{A_{\mathrm{cov}}}e^{E^{\circ}_{\xi}}W_{\chi}(C_n;U_1)\bigr]+E_{\mathrm{def}}+E_{\mathrm{bdry}}\;,
\end{equation}
where $\mathbb E_{\mathrm h}$ uses the original outer operation, the weight $\mathsf w_{\xi}= e^{c_{\xi}}I_{\xi;1}$ and the physical $\mathcal Z_0$. Every term is integrable against the normalised history weight, and the reconstructed fine loop $W_{\chi}(C_n;U_1)$ is retained.
\end{corollary}

\begin{proof}
By \eqref{eq:c4-8-1}, $\mathbb E_{\varepsilon}[W_{\chi}(C_n)]=\mathbb E_{\mathrm h}[\mathbf 1_{A_{\mathrm{cov}}}S^{\circ}_{\xi,C_n}f_{\xi}(0)]+E_{\mathrm{def}}+E_{\mathrm{bdry}}$, with $f_{\xi}(0)= W_{\chi}(C_n;U_1)$ on $A_{\mathrm{cov}}$ (Lemma~\ref{lem:covered}); substitute \eqref{eq:c4c-5-1}. $E^{\circ}_{\xi}$ is a function of the finite geometry at fixed torus and reference parameters, hence of the history labels through $I$ alone, and is independent of the field values at fixed labels, so the first term is measurable; its bound has at most $4n$ boundary terms, $| E^{\circ}_{\xi}|\leqslant C\tau^2(9n+1)$, so the factor $e^{E^{\circ}_{\xi}}$ is bounded and the term integrable. For the other terms, $| S^{\circ}_{\xi,C_n}|\leqslant 1$ because it is the normalised integral of a character of modulus at most one against the positive reference kernel $K^{\circ}_{\xi;0}$ (Lemma~\ref{lem:c2-reference}); $| F_{\xi}(s)|\leqslant 1$ since $F_{\xi}(s)=\int f_{\xi}\,\dd\nu_{\xi;s}$ (\eqref{eq:convex-response-operands}) with $\nu_{\xi;s}$ a probability measure, and $| f_{\xi}(0)|\leqslant 1$; hence $|\mathcal B_{\xi}[f_{\xi}]|=| F_{\xi}(1)-F_{\xi}(0)|\leqslant 2$ (\eqref{eq:c-exterior-response}) and the summed signed defect $| F_{\xi}(0)-S^{\circ}_{\xi,C_n}f_{\xi}(0)|\leqslant 2$, and both are integrable against the unchanged probability history weight. After the substitution the first integrand is still exactly $S^{\circ}_{\xi,C_n}f_{\xi}(0)$, so no independent bound on $e^{E^{\circ}_{\xi}}$ is needed to justify the identity. No new denominator estimate or cancellation is invoked.
\end{proof}

\begin{remark}[scope]\label{rem:c4c-8-1}
The callers of \eqref{eq:c4c-5-1} at $D= I$ are the unlabelled branch identity at \eqref{eq:retained-leading-source} and Corollary~\ref{cor:c4c-5-2}, through which (i) and the first bound of (ii) of the main theorem use it; Appendix~\ref{app:signed-defects} carries $S^{\circ}_{\xi,C}$ inside the expectation and does not use \eqref{eq:c4c-5-1}. Section~\ref{sec:one-step-insertion} carries out the same identification at the retreated carrier $D(\Lambda)\subseteq I$ of its labelled terms (Lemma~\ref{lem:c5-7-1}: factorwise agreement of precision, cutoff, primitives and Wilson factor; Corollary~\ref{cor:c5-7-2}: $S^{\circ}_{\Lambda}=\kappa_Ce^{E^{\circ}_{\Lambda}}$ with the error at $D(\Lambda)$; Lemma~\ref{lem:c5-7-3}: the label average of the boundary sum paid by one-point witness probabilities). The ingredients are the same in both cases: the primitives of \eqref{eq:regional-weight}, the regional stencil identity and the collar. The identity \eqref{eq:c4c-5-3} compares the physical expectation with the reconstructed fine loop at the background $U_1$; a comparison with the bare coarse loop, or convergence over all scales, is not asserted.
\end{remark}

\begin{thebibliography}{99}
    \bibitem{KoteckyPreiss}
    R.~Koteck\'y and D.~Preiss,
    \emph{Cluster Expansion for Abstract Polymer Models},
    Commun. Math. Phys. \textbf{103} (1986), 491--498.
    \bibitem{AbdesselamRivasseau}
    A.~Abdesselam and V.~Rivasseau,
    \emph{Trees, Forests and Jungles: A Botanical Garden for Cluster Expansions},
    in \emph{Constructive Physics}, Lecture Notes in Physics \textbf{446} (1995),
    arXiv:hep-th/9409094.
    \bibitem{BalabanCRE}
    T.~Ba{\l}aban,
    \emph{Convergent Renormalization Expansions for Lattice Gauge Theories},
    Commun. Math. Phys. \textbf{119} (1988), 243--285.

    \bibitem{BalabanRG1}
    T.~Ba{\l}aban,
    \emph{Renormalization Group Approach to Lattice Gauge Field Theories. I. Generation of Effective Actions in a Small Field Approximation and a Coupling Constant Renormalization in Four Dimensions},
    Commun. Math. Phys. \textbf{109} (1987), 249--301.

    \bibitem{BalabanPRT2}
    T.~Ba{\l}aban,
    \emph{Propagators and Renormalization Transformations for Lattice Gauge Theories. II},
    Commun. Math. Phys. \textbf{96} (1984), 223--250.

    \bibitem{BalabanVAR}
    T.~Ba{\l}aban,
    \emph{The Variational Problem and Background Fields in Renormalization Group Method for Lattice Gauge Theories},
    Commun. Math. Phys. \textbf{102} (1985), 277--309.

    \bibitem{BalabanAVG}
    T.~Ba{\l}aban,
    \emph{Averaging Operations for Lattice Gauge Theories},
    Commun. Math. Phys. \textbf{98} (1985), 17--51.

    \bibitem{BalabanREG}
    T.~Ba{\l}aban,
    \emph{Spaces of Regular Gauge Field Configurations on a Lattice and Gauge Fixing Conditions},
    Commun. Math. Phys. \textbf{99} (1985), 75--102.

    \bibitem{Johnsen}
    J.~Johnsen,
    \emph{On the Spectral Properties of Witten-Laplacians, their Range Projections and Brascamp--Lieb's Inequality},
    Integral Equations Operator Theory \textbf{36} (2000), 288--324; arXiv:1703.10011v2 (the numbering of Proposition 4.5, Theorem 1.3 and (4.39) used in Appendix~\ref{app:exterior-chessboard} is that of v2).

    \bibitem{FriedliVelenik}
    S.~Friedli and Y.~Velenik,
    \emph{Statistical Mechanics of Lattice Systems: A Concrete Mathematical Introduction},
    Cambridge University Press, Cambridge, 2017.

    \bibitem{Kechris}
    A.~S.~Kechris,
    \emph{Classical Descriptive Set Theory},
    Graduate Texts in Mathematics \textbf{156}, Springer, New York, 1995.

    \end{thebibliography}
\end{document}
